\documentclass{article}
\usepackage[letterpaper, top=2cm,bottom=2cm,right=2cm,left=2cm]{geometry}
\usepackage{amsmath, amssymb}
\usepackage{enumitem}
\usepackage{hyperref}
\hypersetup{colorlinks=true, linkcolor=blue, urlcolor=blue, citecolor=blue}

\title{CP Teoria}
\author{Ricardo Rojas Carvajal}
\date{August 2026}

\begin{document}
\maketitle

\noindent\textbf{Indice rapido:}
\begin{itemize}\setlength\itemsep{0pt}
	\item \textbf{Plantillas y utilidades}
	\item \textbf{Matematicas}
	\begin{itemize}
		\item Teoria de numeros
		\item Combinatoria
		\item Polinomios y transformadas
		\item Algebra lineal
		\item Juegos y teoria de juegos
	\end{itemize}
	\item \textbf{Estructuras de datos}
	\begin{itemize}
		\item Ordenado y contenedores
		\item Union-Find
		\item Arboles de segmentos
		\item Otras
		\item Fenwick / BIT
		\item RMQ y sparse table
		\item Trees y DP en arbol
		\item BST balanceado
	\end{itemize}
	\item \textbf{Strings}
	\begin{itemize}
		\item Hashing
		\item Patrones
		\item Estructuras avanzadas
		\item Palindromos
	\end{itemize}
	\item \textbf{Grafos}
	\begin{itemize}
		\item Caminos minimos
		\item Componentes y cortes
		\item Matching
		\item Basicos
		\item MST
		\item Flujos
		\item Especales
	\end{itemize}
	\item \textbf{DP}
	\item \textbf{Geometria}
	\begin{itemize}
		\item Puntos y segmentos
		\item Poligonos
		\item Herramientas
	\end{itemize}
	\item \textbf{Miscelanea}
\end{itemize}
\hrule
\vspace{0.5em}

% === Aqui empieza el contenido de cada seccion ===
% (cada plan posterior agrega su bloque debajo)

% ============================================================
\section{Plantillas y utilidades}
% ============================================================

\subsubsection{Plantilla main + fast I/O}

\paragraph{Idea intuitiva.}
Por defecto \texttt{cin}/\texttt{cout} de C++ sincroniza con la libreria C (\texttt{stdio}) para que \texttt{printf} y \texttt{cin} puedan mezclarse. Esa sincronizacion hace que cada operacion de I/O pase por un lock y vuelque buffers, lo cual es \textbf{muy lento} cuando se leen miles de numeros. Desactivar la sincronizacion con \texttt{ios\_base::sync\_with\_stdio(false)} y desatar \texttt{cin} de \texttt{cout} con \texttt{cin.tie(0)} deja a \texttt{cin} usando su propio buffer, dando una velocidad comparable (y a veces mejor) que \texttt{scanf} pero con la comodidad de los operadores \texttt{>>}/\texttt{<<}. El \texttt{\textbackslash n} se prefiere sobre \texttt{endl} porque este ultimo ademas hace flush del buffer, anulando parte de la ganancia. La razon matematica de la lentitud es que cada llamada a \texttt{<<} sin desactivar sincronizacion hace una llamada a \texttt{fwrite} que copia el dato y consulta flags: con $10^6$ enteros esto pasa de milisegundos a segundos.

\paragraph{Propiedades clave (invariantes).}
\begin{itemize}\setlength\itemsep{0pt}
	\item Tras \texttt{sync\_with\_stdio(false)}, \textbf{no mezclar} \texttt{cin}/\texttt{cout} con \texttt{printf}/\texttt{scanf}: el orden de salida queda indefinido.
	\item Tras \texttt{cin.tie(0)}, cada \texttt{cout << x} no fuerza un flush automatico previo del buffer de \texttt{cin} (util cuando se hace \texttt{cout << "resultado: " << ans << "\textbackslash n";} justo despues de leer con \texttt{cin}).
	\item El \texttt{cin.tie(nullptr)} es equivalente a \texttt{cin.tie(0)} en C++11+ y deja claro que es un puntero.
\end{itemize}

\paragraph{Codigo clave.}
Las dos lineas que hacen toda la magia (siempre \textbf{antes} del primer \texttt{cin}/\texttt{cout}):
\begin{verbatim}
ios_base::sync_with_stdio(false);
cin.tie(nullptr);
\end{verbatim}
Estructura tipica de \texttt{main}:
\begin{verbatim}
int main() {
    ios_base::sync_with_stdio(false);
    cin.tie(nullptr);
    int n; cin >> n;
    vector<int> a(n);
    for (auto& x : a) cin >> x;
    // ... logica ...
    cout << ans << "\n";
}
\end{verbatim}

\paragraph{Complejidad.}
\begin{itemize}\setlength\itemsep{0pt}
	\item Sin fast I/O: $\sim$10x mas lento en problemas con $10^5$ o mas entradas.
	\item Con fast I/O: tiempo lineal en el tamano de la entrada, suficiente para casi todos los contests.
	\item Para $10^7$ enteros, se baja de $\sim$8s a $\sim$0.6s.
\end{itemize}

\paragraph{Donde tocar para modificar.}
\begin{itemize}\setlength\itemsep{0pt}
	\item Si el problema lee strings con espacios, cambiar \texttt{cin >> s} por \texttt{getline(cin, s)} (y recordar consumir el \texttt{\textbackslash n} residual con \texttt{cin.ignore()} si se mezcla con \texttt{cin >>}).
	\item Si los numeros no entran en \texttt{int} (por ejemplo $N \le 10^{18}$), cambiar el tipo de \texttt{n}, \texttt{x} y de los contenedores a \texttt{long long}.
	\item Si necesitas \textbf{imprimir con precision fija} (decimales), agregar al inicio: \texttt{cout << fixed << setprecision(10);} (incluir \texttt{<iomanip>}).
	\item Si hay TLE aun con fast I/O, aniadir lectura por bloques con \texttt{fread} (custom scanner); suele ser 3-5x mas rapido que \texttt{cin} rapido.
\end{itemize}

\paragraph{Ejemplo.}
Para $n = 10^6$ enteros y suma acumulada:
\begin{verbatim}
long long s = 0;
for (int i = 0; i < n; ++i) {
    int x; cin >> x;
    s += x;
}
cout << s << "\n";
\end{verbatim}
Sin fast I/O esto tarda $\sim$4s en Codeforces; con fast I/O $\sim$0.4s.

\paragraph{Trampas y errores frecuentes.}
\begin{itemize}\setlength\itemsep{0pt}
	\item Mezclar \texttt{scanf} y \texttt{cin} despues de \texttt{sync\_with\_stdio(false)} produce salidas raras; usar solo uno.
	\item Usar \texttt{endl} en lugar de \texttt{"\textbackslash n"} puede ser $\sim$5x mas lento si se imprime muchas lineas.
	\item En problemas interactivos, a veces es necesario \textbf{volver a atar} \texttt{cin} y forzar flush; usar \texttt{cout << flush;} o \texttt{endl} solo ahi.
	\item Si una funcion que NO es \texttt{main} imprime, no olvides que las dos lineas deben estar en \texttt{main} (o una sola vez al inicio global en un \texttt{struct FastIO} con constructor).
\end{itemize}

\paragraph{Codigo.}
Ver \texttt{cpp.tex}, seccion \textit{Plantilla main + fast I/O}.

\vspace{0.5em}

\subsubsection{Random / mt19937}

\paragraph{Idea intuitiva.}
\texttt{mt19937\_64} es un generador de numeros pseudoaleatorios (Mersenne Twister) con periodo $2^{19937} - 1$. A diferencia del \texttt{rand()} de C, tiene una distribucion razonablemente uniforme y se puede \textbf{seedear} explicitamente, lo que permite reproducir exactamente la misma secuencia. En contests se usa para: barajar vectores en algoritmos randomized (quicksort, kruskal randomizado), generar casos de prueba, y romper empates con hash aleatorio. La razon de preferir \texttt{mt19937\_64} sobre \texttt{rand()}: \texttt{rand()} en algunas plataformas retorna solo 15 bits y el operador \texttt{\% n} introduce sesgo importante para $n$ pequenos.

\paragraph{Propiedades clave (invariantes).}
\begin{itemize}\setlength\itemsep{0pt}
	\item Periodo gigantesco: la probabilidad de repetir es despreciable en cualquier contest.
	\item Produce 64 bits por llamada; si necesitas \texttt{int}, tomar el valor modulo el rango deseado.
	\item \texttt{uniform\_int\_distribution<int>(lo, hi)} es \textbf{inclusivo en ambos extremos}; \texttt{uniform\_real\_distribution<double>} es \textbf{semi-abierto} $[lo, hi)$.
	\item El estado del RNG tiene tamano fijo ($\sim 2.5$KB para \texttt{mt19937\_64}); copiarlo es barato.
\end{itemize}

\paragraph{Codigo clave.}
Patron estandar:
\begin{verbatim}
mt19937_64 rng(chrono::steady_clock::now().time_since_epoch().count());
int x = uniform_int_distribution<int>(0, n-1)(rng);  // [0, n-1]
shuffle(v.begin(), v.end(), rng);                    // barajar
ll y = rng();                                        // 64 bits crudos
\end{verbatim}

\paragraph{Complejidad.}
\begin{itemize}\setlength\itemsep{0pt}
	\item O(1) por llamada. Construir el \texttt{rng} con seed de \texttt{chrono} es O(1).
	\item \texttt{shuffle} sobre $N$ elementos: $O(N)$ (cada llamada a \texttt{swap} consume un random).
\end{itemize}

\paragraph{Donde tocar para modificar.}
\begin{itemize}\setlength\itemsep{0pt}
	\item \textbf{Reproducibilidad}: cambiar el seed de \texttt{chrono::steady\_clock::now()...} por un literal (\texttt{42ULL}, etc.) si queres exactamente la misma secuencia en cada corrida.
	\item \textbf{Tipo}: usar \texttt{mt19937} (32 bits) en lugar de \texttt{mt19937\_64} si solo necesitas \texttt{int}; es ligeramente mas rapido pero no se nota en la practica.
	\item \texttt{randInt(lo, hi)} con \texttt{lo > hi} da comportamiento indefinido; aniade un \texttt{assert} si te preocupa.
	\item \texttt{randPick(v)} hace \texttt{v.size() - 1}; si \texttt{v} esta vacio hay UB. Aniadir check.
	\item \textbf{Saltar N valores}: \texttt{rng.discard(N)} es equivalente a $N$ llamadas pero mas rapido (optimizado internamente).
\end{itemize}

\paragraph{Ejemplo.}
Barajar un arreglo para simulacion Monte Carlo:
\begin{verbatim}
vector<int> p(n);
iota(p.begin(), p.end(), 0);
shuffle(p.begin(), p.end(), rng);
// p es una permutacion aleatoria de [0, n)
\end{verbatim}

\paragraph{Trampas y errores frecuentes.}
\begin{itemize}\setlength\itemsep{0pt}
	\item En Windows (MSVC) el \texttt{uniform\_int\_distribution} puede tener un bug conocido con distribuciones grandes; usar \texttt{mt19937\_64} lo evita.
	\item Llamar a \texttt{rand() \% n} es \textbf{sesgado} si \texttt{RAND\_MAX + 1} no es multiplo de $n$. \texttt{uniform\_int\_distribution} no tiene ese problema.
	\item \texttt{rng()} NO es threadsafe por instancia; si dos threads usan el mismo \texttt{rng} la secuencia se corrompe. Hacer un \texttt{rng} por thread.
	\item Sembrar con \texttt{time(NULL)} da colisiones si lanzas dos programas en el mismo segundo.
\end{itemize}

\paragraph{Codigo.}
Ver \texttt{cpp.tex}, seccion \textit{Random / mt19937}.

\vspace{0.5em}

\subsubsection{BigInt (add + multiply, base 1e9)}

\paragraph{Idea intuitiva.}
Un \texttt{long long} soporta hasta $\approx 9 \times 10^{18}$. Si el problema requiere operar con numeros de cientos o miles de digitos (factoriales grandes, sumas de Catalan, etc.), hace falta una representacion propia. La eleccion mas simple es almacenar los digitos en \textbf{base $10^9$} dentro de un \texttt{vector<int>} (least significant first): asi cada \texttt{int} guarda un ``bloque'' de 9 digitos, y las operaciones aritmeticas se reducen a lo que aprendimos en la escuela, pero con bloques en vez de digitos. La suma es lineal en el numero de bloques; la multiplicacion por un \texttt{int} (no por otro \texttt{BigInt}) tambien lineal. La eleccion de base $10^9$ es la mayor potencia de 10 que cabe en \texttt{int} ($2^{31} - 1 \approx 2.1 \cdot 10^9 > 10^9$).

\paragraph{Propiedades clave (invariantes).}
\begin{itemize}\setlength\itemsep{0pt}
	\item Base $10^9$ asegura que la suma de dos bloques (max $2 \cdot 10^9 - 2$) entra en \texttt{long long} y el carry es a lo sumo $1$.
	\item \texttt{d} siempre tiene el least significant block primero (indice 0).
	\item No hay ceros a la izquierda en \texttt{d}: si el numero es $0$, \texttt{d} esta vacio.
	\item \texttt{neg} indica signo; las operaciones asumen mismo signo.
	\item Cada bloque satisface $0 \le \texttt{d[i]} < 10^9$.
\end{itemize}

\paragraph{Codigo clave.}
El nucleo de la suma (lo que pasa en la escuela, con bloques):
\begin{verbatim}
long long carry = 0;
size_t n = max(d.size(), other.d.size());
d.resize(n);
for (size_t i = 0; i < n; ++i) {
    long long s = carry + d[i];
    if (i < other.d.size()) s += other.d[i];
    d[i] = int(s % 1'000'000'000);
    carry = s / 1'000'000'000;
}
if (carry) d.push_back(int(carry));
\end{verbatim}

\paragraph{Complejidad.}
\begin{itemize}\setlength\itemsep{0pt}
	\item \texttt{+=}: O(N) donde N = \texttt{max(d.size(), other.d.size())}.
	\item \texttt{*=} (por int): O(N).
	\item \texttt{toString}: O(N).
	\item Multiplicacion entre dos \texttt{BigInt}: O($N^2$) si se hace con el metodo clasico; \textbf{no implementado} en este snippet (habria que aniadir un \texttt{*=} que itera bloques o usar FFT para O($N \log N$)).
\end{itemize}

\paragraph{Donde tocar para modificar.}
\begin{itemize}\setlength\itemsep{0pt}
	\item \textbf{Aniadir multiplicacion entre \texttt{BigInt}}: aniadir \texttt{operator*=(const BigInt\&)} o \texttt{operator*}. Para N $\sim$ 1000 bloques la version O($N^2$) sirve; para mas, usar FFT sobre coeficientes \texttt{long long}.
	\item \textbf{Aniadir resta con signo}: complica el invariante ``mismo signo''; mejor aniadir un \texttt{sign()} interno y hacer la resta en valor absoluto, luego ajustar el signo.
	\item \textbf{Cambiar base}: usar base $10^4$ o $10^{18}$ acelera la multiplicacion pero requiere mas cuidado con overflow al multiplicar (usar \texttt{\_\_int128}).
	\item \textbf{Imprimir negativo}: ya hay soporte parcial en \texttt{toString()} con el flag \texttt{neg}; falta activarlo en las operaciones.
	\item \textbf{Division por int} (\texttt{/=} int): lineal, con cuidado de redondear el residuo.
\end{itemize}

\paragraph{Trampas y errores frecuentes.}
\begin{itemize}\setlength\itemsep{0pt}
	\item Multiplicar dos bloques en base $10^9$ puede dar hasta $\approx 10^{18}$ que entra en \texttt{long long} pero \textbf{no} en \texttt{int}. Usar siempre \texttt{long long} o \texttt{\_\_int128} para el producto.
	\item \texttt{toString} usa \texttt{\%09d} para rellenar con ceros a la izquierda; si la base cambia, hay que ajustar el formato.
	\item El snippet NO tiene operador \texttt{<}, \texttt{>}, ni division; si los necesitas, aniadelos con cuidado.
	\item Comparar numeros: comparar primero \texttt{d.size()}; si iguales, comparar de mayor a menor indice.
	\item Olvidar \texttt{normalize()} tras operaciones puede dejar bloques $\ge 10^9$ que rompen \texttt{toString}.
\end{itemize}

\paragraph{Codigo.}
Ver \texttt{cpp.tex}, seccion \textit{BigInt (add + multiply, base 1e9)}.

\vspace{0.5em}

\subsubsection{Debug pretty-printer}

\paragraph{Idea intuitiva.}
Imprimir \texttt{vector<int>\{1,2,3\}} con \texttt{cout << v} da una direccion de memoria, no el contenido. Hay que iterar manualmente. Los macros \texttt{DBG(x)} y \texttt{DBGV(x)} envuelven esa iteracion y muestran, ademas del valor, el \textbf{nombre de la variable} (\texttt{\#x} lo convierte en string en tiempo de preprocesamiento). Asi un \texttt{DBG(ans)} imprime \texttt{ans = 42}, y un \texttt{DBGV(arr)} imprime \texttt{arr = [1, 2, 3]}. Esto acelera muchisimo el debug: con solo aniadir una linea ves el estado de una variable. El truco del \texttt{\#} (stringify) es una facilidad del preprocesador que convierte el token textual en string literal sin expandir macros internas.

\paragraph{Propiedades clave (invariantes).}
\begin{itemize}\setlength\itemsep{0pt}
	\item \texttt{\#x} en una macro \textbf{siempre} toma el nombre textual tal cual lo escribiste; si pasas \texttt{arr[i]}, imprime \texttt{arr[i] = ...}, no el nombre del array.
	\item \texttt{DBGV} solo funciona para \texttt{vector<T>}; para \texttt{set}, \texttt{map}, etc. aniadir overloads especificos.
	\item Los macros escriben a \texttt{cerr} (no \texttt{cout}), asi no contaminan la salida del problema.
	\item Los macros no tienen overhead en tiempo de ejecucion si se compilan con \texttt{-DNDEBUG}.
\end{itemize}

\paragraph{Codigo clave.}
El corazon del macro y el helper de impresion:
\begin{verbatim}
#define DBG(x) cerr << #x << " = " << (x) << "\n"
#define DBGV(x) cerr << #x << " = ["; \
    for (auto it = (x).begin(); it != (x).end(); ++it) \
        cerr << (it==(x).begin()?"":", ") << *it; \
    cerr << "]\n"
\end{verbatim}

\paragraph{Complejidad.}
\begin{itemize}\setlength\itemsep{0pt}
	\item O(N) por impresion donde N = tamano del contenedor.
	\item El overhead del macro es despreciable.
	\item En competicion, compilar con \texttt{-O2} y los macros expandidos son igual de rapidos que el codigo escrito a mano.
\end{itemize}

\paragraph{Donde tocar para modificar.}
\begin{itemize}\setlength\itemsep{0pt}
	\item \textbf{Aniadir tipos} (set, map, pair de vectores, etc.): duplicar la plantilla de \texttt{vecToStr} para cada tipo y crear un \texttt{DBG\#\#} correspondiente.
	\item \textbf{Volcar a archivo}: cambiar \texttt{cerr} por un \texttt{ofstream} abierto al inicio.
	\item \textbf{Quitar del envio}: las macros se activan solo donde estan definidas; antes de enviar al juez, \textbf{borrar las lineas con \texttt{DBG}/\texttt{DBGV}} o aniadir \texttt{\#undef} despues del bloque (como hace el snippet).
	\item \textbf{Debug condicional}: aniadir \texttt{\#ifdef LOCAL} antes de las macros y compilar con \texttt{-DLOCAL} solo en local.
	\item \textbf{Variadic}: para imprimir multiples variables en una linea, usar \texttt{\#define DBG(...) fprintf(stderr, \_\_VA\_ARGS\_\_)}.
\end{itemize}

\paragraph{Trampas y errores frecuentes.}
\begin{itemize}\setlength\itemsep{0pt}
	\item Si el output del juez es estricto (case-sensitive, sin espacios extra), no debe quedar ningun \texttt{cerr} activo.
	\item En recursion profunda, los \texttt{DBG} pueden explotar el stack de \texttt{cerr} (raro, pero pasa). Usar \texttt{\textbackslash n} (no \texttt{endl}) para evitar flush.
	\item Macros con argumentos multiples (\texttt{DBG(a, b)}) requieren variadic macros (\texttt{\#define DBG(...)}) en C++11.
	\item Si \texttt{DBG(x)} se expande en una linea multiple, recordar los \texttt{\textbackslash} al final de cada linea de la macro.
	\item Macros pueden capturar nombres conflictivos; evita nombres como \texttt{T}, \texttt{X} que uses en templates.
\end{itemize}

\paragraph{Codigo.}
Ver \texttt{cpp.tex}, seccion \textit{Debug pretty-printer}.

\vspace{0.5em}

\subsubsection{Coordinate compression}

\paragraph{Idea intuitiva.}
Muchas estructuras (BIT, segment tree, sparse table) indexan con enteros pequenos en $[0, N)$. Si el problema da valores hasta $10^9$, no podes usarlos como indices directamente. La ``compresion de coordenadas'' reemplaza cada valor por su \textbf{rango} en el conjunto ordenado de valores distintos: el menor valor es $0$, el siguiente es $1$, etc. Asi, si los datos son $\{10^9, 5, 100, 42\}$, la compresion da $\{3, 0, 2, 1\}$. Esto preserva el orden y las relaciones de igualdad/desigualdad, pero con indices pequenos y densos. La idea formal: si definimos $f(x) = |\{v_i : v_i \le x\}|$, entonces $f$ es monótona creciente y $f$ restringida a los valores del input es exactamente la compresion.

\paragraph{Propiedades clave (invariantes).}
\begin{itemize}\setlength\itemsep{0pt}
	\item Si $a < b$ en los datos originales, entonces \texttt{compress(a) < compress(b)}.
	\item Si $a = b$, \texttt{compress(a) = compress(b)}.
	\item Los indices resultantes estan en $[0, K)$ donde $K$ = numero de valores distintos ($\le N$).
	\item \texttt{compressWithMap} ademas devuelve el vector inverso: en la posicion $i$ esta el valor original que mapea a $i$.
	\item $K$ (el tamano del rango comprimido) es exactamente el numero de valores distintos.
\end{itemize}

\paragraph{Codigo clave.}
El patron clasico en 3 pasos:
\begin{verbatim}
vector<int> sorted = vals;
sort(sorted.begin(), sorted.end());
sorted.erase(unique(sorted.begin(), sorted.end()), sorted.end());
vector<int> comp(vals.size());
for (size_t i = 0; i < vals.size(); ++i)
    comp[i] = lower_bound(sorted.begin(), sorted.end(), vals[i]) - sorted.begin();
// comp[i] en [0, K)
\end{verbatim}

\paragraph{Complejidad.}
\begin{itemize}\setlength\itemsep{0pt}
	\item O($N \log N$): el \texttt{sort} domina. Las busquedas binarias son O($\log N$) cada una.
	\item Memoria O(N).
	\item El \texttt{unique} + \texttt{erase} es lineal amortizado.
\end{itemize}

\paragraph{Donde tocar para modificar.}
\begin{itemize}\setlength\itemsep{0pt}
	\item \textbf{Preservar duplicados}: si necesitas mantener la posicion original de cada valor ademas del rango, aniadir un struct \texttt{pair<long long, int>} con el valor y el indice original antes de ordenar.
	\item \textbf{Compression online (con inserciones)}: si llegan valores en stream y necesitas comprimirlos, usar una \texttt{policy-based data structure} (ordered set) o \texttt{coordinate\_compression} con insercion en \texttt{vector} ordenado en O($\log N$) por insercion.
	\item \textbf{Mapear a $[1, K]$ en vez de $[0, K)$}: aniadir \texttt{+ 1} al final (util para BIT que usa indices 1-based).
	\item \textbf{3D compression}: para problemas donde hay que comprimir dos coordenadas a la vez (tipico de ``rectangulos disjuntos''), comprimir cada eje por separado y emparejar por valor original.
	\item \textbf{Compression en grafos}: comprimir los nodos activos (los que aparecen) para grafos con vertices hasta $10^9$.
\end{itemize}

\paragraph{Ejemplo.}
Para queries de tipo ``contar elementos en $[L, R]$ entre los $N$ dados'':
\begin{verbatim}
vector<long long> vals = {10, 5, 100, 5, 42};
auto [comp, uniq] = compressWithMap(vals);
// uniq = [5, 10, 42, 100] (ordenado unico)
// comp = [1, 0, 3, 0, 2]
// Para query [L, R]:
int lo = lower_bound(uniq.begin(), uniq.end(), L) - uniq.begin();
int hi = upper_bound(uniq.begin(), uniq.end(), R) - uniq.begin();
// contar comp[i] en [lo, hi-1]
\end{verbatim}

\paragraph{Trampas y errores frecuentes.}
\begin{itemize}\setlength\itemsep{0pt}
	\item No olvides \texttt{sorted.erase(unique(...), ...)}; si lo olvidas, valores repetidos daran rangos duplicados.
	\item Si los datos tienen signo (incluyen negativos), \texttt{lower\_bound} sigue funcionando; no hace falta un truco extra.
	\item Compression \textbf{no preserva distancias}: si la diferencia absoluta importa (por ejemplo, $|a - b|$), el rango no sirve; necesitas el valor original.
	\item En compression online, el orden importa: si agregas elementos en orden de magnitud, insercion es amortizada $O(1)$; si no, es $O(\log N)$ peor caso.
\end{itemize}

\paragraph{Codigo.}
Ver \texttt{cpp.tex}, seccion \textit{Coordinate compression}.

% ============================================================
\section{Matematicas}
% ============================================================

\subsection{Teoria de numeros}

\subsubsection{Extended Euclidean Algorithm}

\paragraph{Idea intuitiva.}
Dados dos enteros $a$ y $b$ no ambos cero, existen enteros $x, y$ tales que $ax + by = \gcd(a, b)$. El algoritmo extendido de Euclides los construye \textbf{mientras} calcula el gcd, mediante recursion: si $(x_1, y_1)$ resuelve $b \cdot x_1 + (a \bmod b) \cdot y_1 = \gcd$, entonces $x = y_1$ e $y = x_1 - y_1 \cdot \lfloor a/b \rfloor$ resuelve $ax + by = \gcd$. La existencia se demuestra por la identidad de B\'ezout: en cada paso, $\gcd(a, b) = \gcd(b, a \bmod b)$, y la recursion garantiza que siempre podemos ``desenrollar'' los coeficientes.

\paragraph{Propiedades clave (invariantes).}
\begin{itemize}\setlength\itemsep{0pt}
	\item Siempre termina (Euclides baja el segundo argumento).
	\item $x$ puede ser negativo.
	\item Si $\gcd(a,b) = 1$, entonces $x$ es el \textbf{inverso modular} de $a$ modulo $b$.
	\item $|x|, |y| \le \max(a, b) / 2$ en valor absoluto (propiedad util para cotas).
\end{itemize}

\paragraph{Codigo clave.}
La recursion que arma la base del algoritmo:
\begin{verbatim}
long long extgcd(long long a, long long b, long long& x, long long& y) {
    if (b == 0) { x = (a >= 0 ? 1 : -1); y = 0; return abs(a); }
    long long x1, y1;
    long long g = extgcd(b, a % b, x1, y1);
    x = y1;
    y = x1 - y1 * (a / b);
    return g;
}
\end{verbatim}

\paragraph{Complejidad.}
\begin{itemize}\setlength\itemsep{0pt}
	\item $O(\log \min(a,b))$ operaciones (misma que el gcd de Euclides).
	\item Cada llamada hace una division y dos asignaciones.
\end{itemize}

\paragraph{Donde tocar para modificar.}
\begin{itemize}\setlength\itemsep{0pt}
	\item \textbf{Inversa modular}: la linea \texttt{x = y1; y = x1 - y1 * (a/b);} resuelve el caso recursivo. Para usar como inversa modular con mod primo, despues ajustar al rango $[0, \text{mod})$.
	\item \textbf{Cuidado con negativos}: si $b$ es negativo, ajustar la division entera.
	\item Si el problema requiere \textbf{muchas} inversas del mismo mod, conviene precomputar con el algoritmo de Fermat ($a^{p-2} \bmod p$) en lugar de llamar a extgcd por cada una.
	\item \textbf{Para ecuaciones diofanticas}: aniadir la busqueda de la solucion general $x = x_0 + (b/g) \cdot t$, $y = y_0 - (a/g) \cdot t$.
\end{itemize}

\paragraph{Trampas y errores frecuentes.}
\begin{itemize}\setlength\itemsep{0pt}
	\item Si los numeros exceden \texttt{int}, cambiar a \texttt{long long}. El snippet retorna \texttt{int}; insuficiente para $a > 2^{31}$.
	\item \texttt{a \% b} con $a$ o $b$ negativos puede dar resultados inesperados en distintos compiladores; ajustar antes de la recursion.
	\item Confundir $a/b$ (truncado hacia cero en C++17+) con $\lfloor a/b \rfloor$; afecta la correccion si $a$ es negativo.
\end{itemize}

\paragraph{Codigo.}
Ver \texttt{cpp.tex}, seccion \textit{Extended Euclidean Algorithm}.

\vspace{0.5em}

\subsubsection{Formulas utiles}

\paragraph{Idea intuitiva.}
Mini-referencia de identidades clasicas: numero de divisores, suma de divisores, formula de Euler ($a^b \equiv a^{b \bmod \varphi(n)} \pmod n$ cuando $\gcd(a,n)=1$), e inclusion-exclusion. \textbf{No hay codigo ejecutable}; solo las formulas y, opcionalmente, comentarios sobre cuando aplicarlas. Estas identidades salen de la aritm\'etica basica: si $n = \prod p_i^{e_i}$, entonces cada divisor es $\prod p_i^{f_i}$ con $0 \le f_i \le e_i$, asi que el conteo se multiplica independientemente por cada primo.

\paragraph{Propiedades clave (invariantes).}
\begin{itemize}\setlength\itemsep{0pt}
	\item Numero de divisores: si $n = \prod p_i^{e_i}$, entonces $d(n) = \prod (e_i + 1)$.
	\item Suma de divisores: $\sigma(n) = \prod \frac{p_i^{e_i+1} - 1}{p_i - 1}$.
	\item Producto de divisores: $n^{d(n)/2}$ (cada divisor $d$ se empareja con $n/d$).
	\item Inclusion-exclusion: $|\bigcup A_i| = \sum_{\emptyset \ne I} (-1)^{|I|+1} |\bigcap_{i \in I} A_i|$.
	\item $\sum_{d \mid n} d \cdot \varphi(n/d) = \sum_{d \mid n} \mu(d) \cdot \sigma(n/d)$ (identidades utiles).
\end{itemize}

\paragraph{Complejidad.}
No aplica para este modulo (es solo referencia de formulas).

\paragraph{Donde tocar para modificar.}
Este modulo \textbf{no tiene codigo}; es una seccion de formulas. Si queres aniadir una formula nueva (por ejemplo, suma de primos hasta $n$ via criba), agregar el \texttt{\textbackslash subsubsection} correspondiente en otra parte del notebook.

\paragraph{Ejemplo de uso.}
Para contar primos relativos a $n$ en $[1, n]$: $\varphi(n)$. Para calcular $\sum_{i=1}^{n} \gcd(i, n)$: factorizar $n$, luego $\sum_{d \mid n} d \cdot \varphi(n/d)$. Para $\sum_{d \mid n} \mu(d) = [n=1]$ (delta de Kronecker).

\paragraph{Trampas y errores frecuentes.}
\begin{itemize}\setlength\itemsep{0pt}
	\item La suma de primos no tiene formula cerrada simple; necesita criba.
	\item $\varphi(n)$ solo simplifica el exponente cuando $\gcd(a, n) = 1$; en caso contrario, hay que factorizar a mano.
	\item En inclusion-exclusion, el orden de los terminos no importa, pero si $|I| > 20$ el numero de terminos explota.
	\item $d(n) = 2$ si y solo si $n$ es primo.
\end{itemize}

\paragraph{Codigo.}
Ver \texttt{cpp.tex}, seccion \textit{Formulas utiles}.

\vspace{0.5em}

\subsubsection{Criba de Eratostenes}

\paragraph{Idea intuitiva.}
Para encontrar todos los primos hasta $N$, marcaremos como compuestos los multiplos de cada primo a partir de $p^2$. Como el primo mas pequeno es $2$, empezar en $4$ evita marcar de mas. Es $O(N \log \log N)$ porque la cantidad de primos hasta $N$ es $\sim N / \ln N$ y la mayoria tiene pocos multiplos pequenos. La cota asintotica precisa es $N \sum_{p \le N} 1/p \approx N \ln \ln N$ por el teorema de Mertens; la constante es pequena ($\sim 1$) en la practica.

\paragraph{Propiedades clave (invariantes).}
\begin{itemize}\setlength\itemsep{0pt}
	\item \texttt{is\_composite[i]} es \texttt{true} sii $i$ no es primo (o $i \le 1$).
	\item \texttt{primes} contiene todos los primos en orden ascendente.
	\item Se itera solo hasta $\sqrt{N}$ en el bucle exterior porque cualquier compuesto $c \le N$ tiene un factor primo $\le \sqrt{N}$.
	\item Cada primo $p$ marca $\lfloor N/p \rfloor - p + 1$ multiplos a partir de $p^2$.
\end{itemize}

\paragraph{Codigo clave.}
El nucleo de la criba:
\begin{verbatim}
vector<bool> comp(N + 1, false);
for (int i = 2; (long long)i * i <= N; ++i)
    if (!comp[i])
        for (int j = i * i; j <= N; j += i)
            comp[j] = true;
vector<int> primes;
for (int i = 2; i <= N; ++i)
    if (!comp[i]) primes.push_back(i);
\end{verbatim}

\paragraph{Complejidad.}
\begin{itemize}\setlength\itemsep{0pt}
	\item Tiempo $O(N \log \log N)$, memoria $O(N)$.
	\item Con $N = 10^7$: $\sim$0.3 segundos en hardware moderno.
\end{itemize}

\paragraph{Donde tocar para modificar.}
\begin{itemize}\setlength\itemsep{0pt}
	\item \textbf{Obtener factorizacion de cada numero}: aniadir un array \texttt{lp[i]} (lowest prime factor) que se llena junto al sieve. Es $O(N)$ si se hace en orden creciente de $i$.
	\item \textbf{Criba hasta $N = 10^7$}: este snippet con \texttt{vector<bool>} (que es bitset-packed) alcanza. Para $N = 10^8$ o mas, conviene criba segmentada (no incluida).
	\item \textbf{Sumar primos}: aniadir \texttt{long long sumPrimes;} que acumule en el bucle final.
	\item \textbf{Primos en $[L, R]$}: criba segmentada, donde $R - L$ es chico pero $L$ es grande.
\end{itemize}

\paragraph{Trampas y errores frecuentes.}
\begin{itemize}\setlength\itemsep{0pt}
	\item \texttt{vector<bool>} es proxy; si necesitas pasarlo a C-style API, copiar a \texttt{vector<char>}.
	\item \texttt{is\_composite[0]} e \texttt{is\_composite[1]} deben inicializarse a \texttt{true}.
	\item El bucle interno empieza en \texttt{i*i}, no en \texttt{2*i}; si arrancas en \texttt{2*i} marcas el doble y pierdes tiempo pero el resultado es el mismo.
	\item Overflow si \texttt{i*i > INT\_MAX}: usar \texttt{(long long)i * i <= N}.
\end{itemize}

\paragraph{Codigo.}
Ver \texttt{cpp.tex}, seccion \textit{Criba de Eratostenes}.

\vspace{0.5em}

\subsubsection{Criba lineal}

\paragraph{Idea intuitiva.}
La criba de Eratostenes marca cada compuesto varias veces (por ejemplo $12$ se marca por $2$, por $3$ y por $4$). La criba lineal asegura que \textbf{cada compuesto se marca exactamente una vez}, lo cual da $O(N)$ total. La idea: para cada $i$, iterar sobre los primos $p$ en orden creciente y poner $lp[i \cdot p] = p$, \textbf{cortando} el bucle cuando $p = lp[i]$ (asi $i \cdot p$ queda con su factor primo mas chico como $p$). La prueba de $O(N)$: el bucle interno suma trabajo proporcional al numero de primos que \textbf{no} cortan, y solo se corta cuando $p \mid i$; la suma total sobre todos los $i$ es exactamente $N$.

\paragraph{Propiedades clave (invariantes).}
\begin{itemize}\setlength\itemsep{0pt}
	\item \texttt{lp[i]} = menor factor primo de $i$ (con \texttt{lp[0] = lp[1] = 0}).
	\item Cada $i$ tiene asignado exactamente un primo que ``lo cierra''.
	\item Recorriendo $i$ en orden creciente, cuando entramos a $i$, todos sus primos menores ya estan en \texttt{primes}.
	\item La condicion de corte \texttt{p == lp[i]} garantiza que $i \cdot p$ tiene \texttt{lp} minimo.
\end{itemize}

\paragraph{Codigo clave.}
La criba lineal en su forma esencial:
\begin{verbatim}
vector<int> lp(n + 1), primes;
for (int i = 2; i <= n; ++i) {
    if (!lp[i]) { lp[i] = i; primes.push_back(i); }
    for (int p : primes) {
        if (p > lp[i] || (long long)i * p > n) break;
        lp[i * p] = p;
    }
}
\end{verbatim}

\paragraph{Complejidad.}
\begin{itemize}\setlength\itemsep{0pt}
	\item Tiempo $O(N)$ exacto, memoria $O(N)$.
	\item En la practica, $\sim$2x mas rapido que la criba clasica para $N \ge 10^6$.
\end{itemize}

\paragraph{Donde tocar para modificar.}
\begin{itemize}\setlength\itemsep{0pt}
	\item \textbf{Factorizacion de un numero $n$}: usando \texttt{lp}, iterar \texttt{while (n > 1) \{ p = lp[n]; ...; n /= p; \}} en $O(\log n)$.
	\item \textbf{Cantidad de divisores para todos los $i \le N$}: aniadir \texttt{cntDiv[i] = cntDiv[i/lp[i]] + 1} durante la criba.
	\item \textbf{Suma de divisores o totient por $i$}: igual, precomputar en el mismo pase.
	\item \textbf{Aniadir mobius}: \texttt{mu[1] = 1; mu[i*p] = (p == lp[i]) ? 0 : -mu[i];}.
\end{itemize}

\paragraph{Trampas y errores frecuentes.}
\begin{itemize}\setlength\itemsep{0pt}
	\item La condicion \texttt{if (p == lp[i]) break;} es \textbf{crucial}; sin ella la criba vuelve a ser $O(N \log \log N)$.
	\item Overflow: chequear \texttt{1LL * i * p > n} antes de asignar.
	\item Si \texttt{lp[i] == 0} (i es primo), la condicion \texttt{p > lp[i]} es \texttt{p > i}, lo que tambien detiene el bucle correctamente.
\end{itemize}

\paragraph{Codigo.}
Ver \texttt{cpp.tex}, seccion \textit{Criba lineal}.

\vspace{0.5em}

\subsubsection{Miller-Rabin primality}

\paragraph{Idea intuitiva.}
Para $n$ impar, escribir $n - 1 = 2^s \cdot d$ con $d$ impar. Si $n$ es primo, entonces para \textbf{cualquier} $a$ con $1 < a < n-1$ se cumple que $a^d \equiv 1 \pmod n$ o $a^{2^r \cdot d} \equiv -1 \pmod n$ para algun $r < s$. El test verifica esto con varios \texttt{a} (testigos). Si ningun testigo delata a $n$, es primo. Con un conjunto fijo de testigos, es \textbf{determinista} para $n < 2^{64}$. La demostracion se basa en el hecho de que en $\mathbb{Z}/p\mathbb{Z}$ con $p$ primo, el grupo $(\mathbb{Z}/p\mathbb{Z})^*$ es ciclico de orden $p-1$, asi que $a^{p-1} \equiv 1$ por Fermat, y la extraccion de raices cuadradas solo tiene dos posibilidades ($\pm 1$).

\paragraph{Propiedades clave (invariantes).}
\begin{itemize}\setlength\itemsep{0pt}
	\item Probabilistico si los testigos son aleatorios; \textbf{determinista} con el set del snippet para $u64$.
	\item Falla para $n < 2$ (no es primo).
	\item Bases: $\{2, 325, 9375, 28178, 450775, 9780504, 1795265022\}$ cubren todos los $u64$.
	\item Falsos positivos: si un testigo no delata, $n$ puede ser primo \textbf{o} un fuerte pseudoprimo respecto a ese testigo.
\end{itemize}

\paragraph{Codigo clave.}
La ronda de Miller-Rabin para un testigo $a$:
\begin{verbatim}
long long x = modpow(a, d, n);
if (x == 1 || x == n - 1) return true;  // pasa este testigo
for (int r = 0; r < s - 1; ++r) {
    x = mul_mod(x, x, n);
    if (x == n - 1) return true;  // aparece -1 en alguna iteracion
}
return false;  // n es compuesto
\end{verbatim}

\paragraph{Complejidad.}
\begin{itemize}\setlength\itemsep{0pt}
	\item $O(\log^3 n)$ por la exponenciacion modular, pero con constante pequena. Para $u64$, unos pocos microsegundos.
	\item Con 7 testigos para $u64$: $< 10$ microsegundos en hardware moderno.
\end{itemize}

\paragraph{Donde tocar para modificar.}
\begin{itemize}\setlength\itemsep{0pt}
	\item \textbf{Cambiar el set de testigos}: si necesitas testear numeros $> 2^{64}$, aniadir mas testigos o usar BPSW.
	\item \textbf{Devolver un certificado}: aniadir la secuencia de cuadrados que llevo al fallo para debugging.
	\item \textbf{\texttt{u128}}: el snippet usa \texttt{\_\_uint128\_t} para multiplicar mod $n$ sin overflow. En GCC/clang funciona; en MSVC aniadir \texttt{\#include <boost/multiprecision/cpp\_int.hpp>}.
	\item \textbf{Solo testigos pequenos}: si sabes que $n < 2^{32}$, basta con $\{2, 7, 61\}$ (solo 3 testigos).
\end{itemize}

\paragraph{Trampas y errores frecuentes.}
\begin{itemize}\setlength\itemsep{0pt}
	\item Multiplicar \texttt{u64 * u64} puede overflow; \textbf{siempre} usar \texttt{u128}.
	\item Si $n$ es par, el test responde rapido (lo chequea con \texttt{n \% p == 0}).
	\item Cuidado con bases que sean multiplos de $n$; el snippet las saltea con \texttt{if (a \% n == 0) continue;}.
	\item Si el loop \texttt{for (int r = ...)} itera \texttt{r < s - 1} (no \texttt{r < s}), el ultimo caso se chequea afuera.
\end{itemize}

\paragraph{Codigo.}
Ver \texttt{cpp.tex}, seccion \textit{Miller-Rabin primality}.

\vspace{0.5em}

\subsubsection{Pollard's Rho + factorizacion}

\paragraph{Idea intuitiva.}
Miller-Rabin detecta primalidad pero no factoriza. Pollard's Rho busca un \textbf{factor no trivial} de $n$ usando una secuencia pseudoaleatoria $x_{k+1} = x_k^2 + c \pmod n$ y aplicando la \textbf{paradoja del cumpleanos}: si la secuencia cae en un ciclo modulo un factor primo $p$ de $n$, entonces $\gcd(|x_i - x_j|, n)$ sera $p$ para algun $i, j$. Esto es esperado $O(n^{1/4})$ por factor. La intuicion: la secuencia modulo $p$ (donde $p \le \sqrt{n}$) tiene un ciclo esperado de longitud $O(\sqrt{p}) = O(n^{1/4})$, mucho mas corto que el ciclo modulo $n$ ($\sim n$).

\paragraph{Propiedades clave (invariantes).}
\begin{itemize}\setlength\itemsep{0pt}
	\item Cada llamada a \texttt{pollard(n)} devuelve \textbf{un} factor (no necesariamente primo); se llama recursivamente.
	\item Si \texttt{millerRabin(n)} retorna \texttt{true}, $n$ es primo y se aniade directo.
	\item El factor final sale \textbf{ordenado} porque el snippet ordena al final.
	\item El algoritmo de Floyd (``tortuga y liebre'') detecta el ciclo comparando dos punteros a distinta velocidad.
\end{itemize}

\paragraph{Codigo clave.}
El bucle de Floyd (tortuga y liebre):
\begin{verbatim}
long long x = rand2(rng) % n, y = x, c = rand2(rng) % (n - 1) + 1, d = 1;
while (d == 1) {
    x = (mul_mod(x, x, n) + c) % n;       // tortuga: 1 paso
    y = (mul_mod(y, y, n) + c) % n;
    y = (mul_mod(y, y, n) + c) % n;       // liebre: 2 pasos
    d = gcd(abs(x - y), n);
}
if (d != n) return d;  // encontro factor
\end{verbatim}

\paragraph{Complejidad.}
\begin{itemize}\setlength\itemsep{0pt}
	\item Esperado $O(n^{1/4})$ por factor. Practico para $n$ hasta $10^{18}$.
	\item Para $n = 10^{18}$, tarda $\sim 10^3$ iteraciones por factor.
\end{itemize}

\paragraph{Donde tocar para modificar.}
\begin{itemize}\setlength\itemsep{0pt}
	\item \textbf{Mejorar la constante}: aniadir el algoritmo de Brent (variante del Floyd para el ciclo, $\sim 24\%$ mas rapido).
	\item \textbf{Determinismo}: si la corrida con $c$ falla, reintentar con otro $c$ aleatorio (ya lo hace el snippet con \texttt{while(true)}).
	\item \textbf{Salida con multiplicidad}: el snippet devuelve la lista con repeticiones (porque \texttt{factorRec} empuja \texttt{n} cuando es primo y tambien \texttt{n/d}). Si solo queres factores distintos, aniadir un \texttt{set} o deduplicar al final.
	\item \textbf{Acelerar con product acumulado}: Brent acumula productos parciales y reduce gcd; ahorra muchos gcd.
\end{itemize}

\paragraph{Trampas y errores frecuentes.}
\begin{itemize}\setlength\itemsep{0pt}
	\item Si $n$ es potencia de un primo pequeno (ej. $2^k$), Pollard puede tardar. Considerar chequeo previo de potencia.
	\item Si $n$ es primo, llamar a \texttt{pollard(n)} causaria loop infinito porque \texttt{d == n}; el snippet tiene la guarda \texttt{if (millerRabin(n)) return}.
	\item Si $c = 0$ y $x = 0$ o $x$ es punto fijo, la secuencia es constante y $\gcd = 0$; por eso $c \ne 0$.
	\item \texttt{abs(x - y)} puede overflow si ambos son cercanos a $n/2$; usar \texttt{u64} con cuidado o comparar con \texttt{std::abs(int64\_t(x - y))}.
\end{itemize}

\paragraph{Codigo.}
Ver \texttt{cpp.tex}, seccion \textit{Pollard's Rho + factorizacion}.

\vspace{0.5em}

\subsubsection{Modular exponentiation}

\paragraph{Idea intuitiva.}
Calcular $a^e \bmod m$ en $O(\log e)$ mediante \textbf{exponenciacion binaria} (squaring). La idea: si $e$ es par, $a^e = (a^{e/2})^2$; si es impar, $a^e = a \cdot (a^{(e-1)/2})^2$. Asi, en cada paso, el exponente se reduce a la mitad y se hace una multiplicacion modular. El nombre ``squaring'' viene de que en cada paso elevamos al cuadrado una variable (``base''). Se puede demostrar por induccion que tras $k$ iteraciones, la variable ``resultado'' contiene $a^{(\text{bits consumidos})}$, manteniendo la invariante $\text{result} \cdot \text{base}^e = a^{e_{\text{original}}}$.

\paragraph{Propiedades clave (invariantes).}
\begin{itemize}\setlength\itemsep{0pt}
	\item \texttt{a \% m} se hace al inicio (en caso de que \texttt{a} sea negativo, se ajusta \texttt{+= mod}).
	\item Multiplicacion: para evitar overflow, usar \texttt{\_\_int128} o \texttt{long long} solo si $m < 2^{32}$ (porque $a \cdot b < m^2 < 2^{64}$).
	\item Resultado siempre en $[0, m)$.
	\item El algoritmo usa exactamente $\lfloor \log_2 e \rfloor + 1$ multiplicaciones (medido en bits del exponente).
\end{itemize}

\paragraph{Codigo clave.}
El bucle de squaring en su forma mas pura:
\begin{verbatim}
long long modpow(long long a, long long e, long long m) {
    long long r = 1; a %= m;
    while (e > 0) {
        if (e & 1) r = (__int128)r * a % m;
        a = (__int128)a * a % m;
        e >>= 1;
    }
    return r;
}
\end{verbatim}

\paragraph{Complejidad.}
\begin{itemize}\setlength\itemsep{0pt}
	\item $O(\log e)$ multiplicaciones modulares.
	\item $\sim 30$ iteraciones para $e$ hasta $10^9$; $\sim 60$ para $e$ hasta $10^{18}$.
\end{itemize}

\paragraph{Donde tocar para modificar.}
\begin{itemize}\setlength\itemsep{0pt}
	\item \textbf{Modulo de $10^{18}$}: aniadir \texttt{\_\_int128} a las multiplicaciones (ya hecho si el snippet usa \texttt{u128}; si no, aniadir manualmente).
	\item \textbf{Exponentes muy grandes} ($e > 10^{18}$): el bucle sigue funcionando pero no entra; aniadir modularizacion del exponente via Euler si $\gcd(a, m) = 1$ y queres reducir $e$ primero.
	\item \textbf{Pow matrix/Poly}: la misma estructura sirve; solo cambia el tipo de la base y la operacion (mul de matrices, mul de polinomios).
	\item \textbf{Exponenciacion rapida sin division}: Montgomery multiplication para criptografia.
\end{itemize}

\paragraph{Trampas y errores frecuentes.}
\begin{itemize}\setlength\itemsep{0pt}
	\item Si $m$ es 1, el resultado siempre es 0; aniadir check si te molesta.
	\item Overflow en \texttt{a * a \% mod} si \texttt{a} y \texttt{mod} son \texttt{long long} grandes.
	\item En MSVC, \texttt{\_\_int128} no existe; aniadir \texttt{\#include <boost/multiprecision/cpp\_int.hpp>}.
	\item Si $e$ es \texttt{int} negativo (raro, pero por un bug), el bucle no termina.
\end{itemize}

\paragraph{Codigo.}
Ver \texttt{cpp.tex}, seccion \textit{Modular exponentiation}.

\vspace{0.5em}

\subsubsection{Modular inverse}

\paragraph{Idea intuitiva.}
La \textbf{inversa} de $a$ modulo $m$ es un $b$ tal que $a \cdot b \equiv 1 \pmod m$. Si $m$ es primo y $\gcd(a, m) = 1$, entonces $b = a^{m-2} \pmod m$ (por el \textbf{teorema pequeno de Fermat}). Si $m$ no es primo, se usa el algoritmo extendido de Euclides, que resuelve $ax + my = \gcd(a, m)$. El truco de Fermat: si $m$ es primo, $(\mathbb{Z}/m\mathbb{Z})^*$ es un grupo de orden $m - 1$, asi que $a^{m-1} \equiv 1$ y $a \cdot a^{m-2} \equiv 1$.

\paragraph{Propiedades clave (invariantes).}
\begin{itemize}\setlength\itemsep{0pt}
	\item La inversa existe sii $\gcd(a, m) = 1$.
	\item Para mod primo, \texttt{modInv(a, mod) = modpow(a, mod - 2, mod)}.
	\item \texttt{modInvGeneral} retorna $-1$ si no hay inversa.
	\item Si existe, la inversa es \textbf{unica} modulo $m$.
\end{itemize}

\paragraph{Codigo clave.}
Recurrencia lineal para precomputar todas las inversas (solo funciona con mod primo):
\begin{verbatim}
// inv[1] = 1
// inv[i] = -(mod/i) * inv[mod%i] mod mod
for (int i = 2; i <= N; ++i) {
    inv[i] = MOD - (long long)(MOD / i) * inv[MOD % i] % MOD;
}
\end{verbatim}

\paragraph{Complejidad.}
\begin{itemize}\setlength\itemsep{0pt}
	\item $O(\log \text{mod})$ en ambos casos (Fermat y extgcd).
	\item Precomputar $N$ inversas: $O(N)$.
\end{itemize}

\paragraph{Donde tocar para modificar.}
\begin{itemize}\setlength\itemsep{0pt}
	\item \textbf{Precomputar muchas inversas} (ej. para todos los $1..N$): precomputar en $O(N)$ con la recurrencia \texttt{inv[i] = -(mod/i) * inv[mod \% i] \% mod} y ajustar a positivo.
	\item \textbf{Cambiar mod}: si modificas el \texttt{MOD} global, asegurate de que sea primo (sino, usar extgcd).
	\item \textbf{Ajustar a positivo}: tras \texttt{-= MOD} o \texttt{+= MOD} para tener resultado en $[0, MOD)$.
\end{itemize}

\paragraph{Trampas y errores frecuentes.}
\begin{itemize}\setlength\itemsep{0pt}
	\item Si \texttt{MOD = 1}, $a^{MOD-2} = a^{-1}$ no esta definido; aniadir check.
	\item Si \texttt{a = 0}, no hay inversa (excepto en mod trivial).
	\item El snippet usa \texttt{modInv(a, mod)} que asume mod primo por construccion.
	\item La recurrencia lineal \texttt{inv[i] = -(mod/i)*inv[mod\%i]} solo sirve si \texttt{mod} es primo y \texttt{i < mod}.
	\item Si llamas \texttt{modInv} con $a$ divisible por \texttt{mod}, devuelve cualquier cosa (generalmente 0 porque $a^{m-2} \equiv 0$).
\end{itemize}

\paragraph{Codigo.}
Ver \texttt{cpp.tex}, seccion \textit{Modular inverse}.

\vspace{0.5em}

\subsubsection{Chinese Remainder Theorem}

\paragraph{Idea intuitiva.}
Dado un sistema $x \equiv r_i \pmod{m_i}$ con los $m_i$ \textbf{coprimos dos a dos}, hay \textbf{una} solucion modulo $M = \prod m_i$. Se construye incrementalmente: combinando $x \equiv r_1 \pmod{m_1}$ con $x \equiv r_2 \pmod{m_2}$, la solucion es $x = r_1 + m_1 \cdot t$ donde $t \equiv (r_2 - r_1) \cdot m_1^{-1} \pmod{m_2}$. Iterar con cada $m_i$. La prueba de existencia y unicidad es por induccion sobre $N$: con una sola ecuacion la solucion es trivial; al aniadir la ecuacion $i+1$ con $\gcd(m_i, m_{i+1}) = 1$, existe $m_i^{-1} \pmod{m_{i+1}}$ y se resuelve el nuevo offset.

\paragraph{Propiedades clave (invariantes).}
\begin{itemize}\setlength\itemsep{0pt}
	\item El snippet asume $m_i$ coprimos dos a dos; \textbf{no} maneja el caso general.
	\item Solucion unica modulo $M$.
	\item Si los modulos \textbf{no} son coprimos y $\gcd(m_i, m_j) \nmid (r_i - r_j)$, no hay solucion; si divide, se reduce el sistema.
	\item Cada paso incremental no destruye las ecuaciones previas.
\end{itemize}

\paragraph{Codigo clave.}
El paso inductivo (combinar una solucion parcial con una ecuacion nueva):
\begin{verbatim}
long long k = ((r2 - r1) % m2 + m2) % m2;
k = k * modInv(m1 % m2, m2) % m2;  // inversa de m1 mod m2
r1 = r1 + m1 * k;     // nueva solucion
m1 = m1 * m2;          // nuevo modulo
r1 = (r1 % m1 + m1) % m1;
\end{verbatim}

\paragraph{Complejidad.}
\begin{itemize}\setlength\itemsep{0pt}
	\item $O(N \log M)$ con $N$ el numero de ecuaciones.
	\item Cada paso usa una inversa modular ($O(\log m_i)$) y unas multiplicaciones.
\end{itemize}

\paragraph{Donde tocar para modificar.}
\begin{itemize}\setlength\itemsep{0pt}
	\item \textbf{Modulos no coprimos}: aniadir el test $\gcd(m_i, m_j)$ en cada paso y devolver error o reducir.
	\item \textbf{Devolver todas las soluciones}: si el problema quiere todas las $x$ en $[0, L]$, ajustar al final con un \texttt{if (x > L) return -1;}.
	\item \textbf{Garner}: para evitar overflow con $M$ muy grande, usar el algoritmo de Garner (mas estable numericamente).
	\item \textbf{Caso $k = 0$}: aniadir check para retornar temprano (ya resuelto).
\end{itemize}

\paragraph{Trampas y errores frecuentes.}
\begin{itemize}\setlength\itemsep{0pt}
	\item Overflow al multiplicar \texttt{M * t} si los $m_i$ son grandes; usar \texttt{\_\_int128} o reducir modulo $M$ en cada paso (ya lo hace el snippet).
	\item Si una ecuacion es $x \equiv r_i \pmod{m_i}$ con $r_i \ge m_i$, normalizar antes.
	\item Confundir ``modulos coprimos'' con ``modulos primos individuales''; los $m_i$ solo necesitan ser coprimos entre si.
	\item En la condicion ``hay solucion'', devolver un valor sentinela (e.g., $-1$) consistente con el resto del codigo.
\end{itemize}

\paragraph{Codigo.}
Ver \texttt{cpp.tex}, seccion \textit{Chinese Remainder Theorem}.

\vspace{0.5em}

\subsubsection{Euler totient precompute}

\paragraph{Idea intuitiva.}
La funcion $\varphi(n)$ cuenta cuantos enteros en $[1, n]$ son coprimos con $n$. Hay una formula: $\varphi(n) = n \cdot \prod_{p \mid n} (1 - 1/p)$. Para precomputar $\varphi(i)$ para todo $i \le N$ en $O(N \log \log N)$, se inicializa $\varphi(i) = i$ y para cada primo $p$ se hace $\varphi(j) \mathrel{-}= \varphi(j)/p$ para todos los multiplos $j$ de $p$. La razon geometrica: para cada primo $p$, la fraccion $1/p$ de los enteros hasta $n$ son multiplos de $p$, asi que $\varphi$ ``pierde'' ese factor en cada primo divisor.

\paragraph{Propiedades clave (invariantes).}
\begin{itemize}\setlength\itemsep{0pt}
	\item $\varphi(1) = 1$ por convencion (aunque coprimos con 1 es solo $\{1\}$).
	\item $\varphi$ es \textbf{multiplicativa}: si $\gcd(a, b) = 1$, $\varphi(ab) = \varphi(a) \varphi(b)$.
	\item \textbf{Identidad de Euler}: si $\gcd(a, n) = 1$, entonces $a^{\varphi(n)} \equiv 1 \pmod n$.
	\item $\varphi(n) \ge \sqrt{n}$ para $n > 6$ (acotacion util en busqueda).
	\item $\sum_{d \mid n} \varphi(d) = n$ (propiedad de inversion, util en Dirichlet).
\end{itemize}

\paragraph{Codigo clave.}
Criba Euler-style para totients (un solo pase):
\begin{verbatim}
vector<int> phi(n + 1);
for (int i = 0; i <= n; ++i) phi[i] = i;
for (int i = 2; i <= n; ++i)
    if (phi[i] == i)  // i es primo
        for (int j = i; j <= n; j += i)
            phi[j] -= phi[j] / i;
\end{verbatim}

\paragraph{Complejidad.}
\begin{itemize}\setlength\itemsep{0pt}
	\item Precompute $O(N \log \log N)$, consulta $O(1)$. Single $\varphi(n)$ con factorizacion es $O(\sqrt{n})$.
\end{itemize}

\paragraph{Donde tocar para modificar.}
\begin{itemize}\setlength\itemsep{0pt}
	\item \textbf{Single phi via factorizacion}: ya hay una funcion \texttt{phiSingle(n)}. Para $n$ hasta $10^{12}$ factoriza en $\le 10^6$ pasos.
	\item \textbf{Aplicar a inclusion-exclusion}: si el problema requiere $\sum_{i=1}^{N} \gcd(i, n)$, se puede usar $\sum_{d \mid n} d \cdot \varphi(n/d)$.
	\item \textbf{Criba lineal}: durante la criba, aniadir \texttt{phi[i*p] = (p == lp[i]) ? phi[i] * p : phi[i] * (p-1);}.
\end{itemize}

\paragraph{Ejemplo.}
Para $n = 12$: $\varphi(12) = 12 \cdot (1 - 1/2) \cdot (1 - 1/3) = 12 \cdot 1/2 \cdot 2/3 = 4$. Los coprimos son $\{1, 5, 7, 11\}$.

\paragraph{Trampas y errores frecuentes.}
\begin{itemize}\setlength\itemsep{0pt}
	\item Olvidar inicializar $\varphi(0) = 0$.
	\item Si llamas a \texttt{phiSingle} sobre 0, devuelve 0 (no hay coprimos con 0 en el sentido usual).
	\item La criba Euler-style modifica $\varphi[j]$ por division \textbf{antes} de que termine; para $j$ compuestos, se hace varias veces (correcto, multiplicativo).
	\item Confundir $\varphi(0) = 0$ con $\varphi(1) = 1$; la criba comun los trata igual.
\end{itemize}

\paragraph{Codigo.}
Ver \texttt{cpp.tex}, seccion \textit{Euler totient precompute}.

\vspace{0.5em}

\subsubsection{Mobius precompute}

\paragraph{Idea intuitiva.}
La funcion $\mu(n)$ (``Mobius'') vale $1$ si $n$ es producto de un numero par de primos distintos, $-1$ si es impar, $0$ si tiene un primo al cuadrado. Aparece en inclusion-exclusion: si $f(n) = \sum_{d \mid n} g(d)$, entonces $g(n) = \sum_{d \mid n} \mu(n/d) f(d)$ (\textbf{inversion de Mobius}). Se precomputa junto a la criba lineal. La relacion con inclusion-exclusion: $\mu$ es la ``inversa'' del operador ``suma sobre divisores'', asi que si conocemos los acumulados podemos recuperar los originales.

\paragraph{Propiedades clave (invariantes).}
\begin{itemize}\setlength\itemsep{0pt}
	\item $\sum_{d \mid n} \mu(d) = [n = 1]$ (funcion indicador de 1).
	\item Si $\gcd(a, b) = 1$, $\mu(ab) = \mu(a) \mu(b)$ (multiplicativa).
	\item $\sum_{d \mid n} \mu(d) \cdot \frac{n}{d} = \varphi(n)$ (relacion con totient).
	\item Para $n$ cuadrado libre: $|\mu(n)| = 1$ y $\mu(n) = (-1)^{\omega(n)}$ donde $\omega$ es el numero de factores primos.
\end{itemize}

\paragraph{Codigo clave.}
Mobius durante la criba lineal:
\begin{verbatim}
mu[1] = 1;
for (int i = 2; i <= n; ++i) {
    if (!lp[i]) { lp[i] = i; primes.push_back(i); mu[i] = -1; }
    for (int p : primes) {
        if (p > lp[i] || (long long)i * p > n) break;
        lp[i * p] = p;
        if (p == lp[i]) mu[i * p] = 0;        // cuadrado
        else mu[i * p] = -mu[i];              // nuevo primo
    }
}
\end{verbatim}

\paragraph{Complejidad.}
\begin{itemize}\setlength\itemsep{0pt}
	\item Precompute $O(N)$ con la criba lineal. Consulta $O(1)$.
\end{itemize}

\paragraph{Donde tocar para modificar.}
\begin{itemize}\setlength\itemsep{0pt}
	\item \textbf{Suma de $\gcd$}: $\sum_{i=1}^{N} \gcd(i, n) = \sum_{d \mid n} d \cdot \varphi(n/d)$, alternativa con Mobius: $\sum_{d \mid n} \mu(d) \cdot S(\lfloor n/d \rfloor)$ donde $S(x) = \sum_{k \le x} k$.
	\item \textbf{Mobius en inclusion-exclusion}: el snippet es solo para precompute; usarlo combinado con divisor enumeration en queries.
	\item \textbf{Conteo de coprimos}: $\sum_{d \mid n} \mu(d) \cdot \lfloor n/d \rfloor = \varphi(n)$.
\end{itemize}

\paragraph{Ejemplo.}
$\mu(1) = 1$, $\mu(2) = -1$, $\mu(6) = 1$ (dos primos distintos), $\mu(4) = 0$ (cuadrado). $\sum_{d \mid 6} \mu(d) = \mu(1) + \mu(2) + \mu(3) + \mu(6) = 1 - 1 - 1 + 1 = 0$ (correcto: $6 \ne 1$).

\paragraph{Trampas y errores frecuentes.}
\begin{itemize}\setlength\itemsep{0pt}
	\item $\mu(0) = 0$ (por definicion, $0$ no es producto de primos distintos).
	\item El snippet inicializa \texttt{mu[0] = 0, mu[1] = 1} por separado.
	\item Si precomputas junto a la criba lineal, mantener consistencia: \texttt{mu[p]} para primo es $-1$ (un solo primo).
	\item En queries de Mobius, sumar con el signo correcto; $(-1)^{\omega(n)}$ se puede calcular con la criba.
\end{itemize}

\paragraph{Codigo.}
Ver \texttt{cpp.tex}, seccion \textit{Mobius precompute}.

\vspace{0.5em}

\subsubsection{Stern-Brocot (best fraction in interval)}

\paragraph{Idea intuitiva.}
Stern-Brocot es un arbol binario de fracciones irreducibles. Cada nodo es el \textbf{mediantes} de sus ``padres'' en el arbol: si tienes $a/b$ y $c/d$, el mediantes es $(a+c)/(b+d)$. Para encontrar la fraccion irreducible con denominador acotado mas cercana a un valor $x$, se navega el arbol decidiendo ir a izquierda o derecha segun si $x$ es menor o mayor. La fraccion mas simple en un intervalo $(a/b, c/d)$ se obtiene iterando mediantes y quedandose con el \textbf{maximo $k$} tal que el mediantes sigue dentro del intervalo. La propiedad clave: las fracciones irreducibles con denominador acotado estan \textbf{todas} en el arbol.

\paragraph{Propiedades clave (invariantes).}
\begin{itemize}\setlength\itemsep{0pt}
	\item Cada par de fracciones irreducibles consecutivas en Stern-Brocot tiene diferencia $1/(bd)$ en valor.
	\item Permite generar todas las fracciones irreducibles con un denominador acotado.
	\item El mediantes $m/n$ de $a/b$ y $c/d$ satisface $b \cdot m - a \cdot n = 1$ (cruz = 1).
\end{itemize}

\paragraph{Codigo clave.}
La busqueda en Stern-Brocot (izquierda/derecha segun comparacion):
\begin{verbatim}
long long findBest(long long x, long long a, long long b, long long c, long long d) {
    long long p = a + c, q = b + d;  // mediantes
    if (q > Q) return ...;  // denominador excede
    if (p/q > x) return find(x, a, b, p, q);  // izquierda
    else         return find(x, p, q, c, d);  // derecha
}
\end{verbatim}

\paragraph{Complejidad.}
\begin{itemize}\setlength\itemsep{0pt}
	\item $O(\log(\max(a,b)))$ por query.
	\item Generar todas las fracciones con denominador $\le Q$: $O(Q^2)$ (hay $\sim 3Q^2/\pi^2$ irreducibles).
\end{itemize}

\paragraph{Donde tocar para modificar.}
\begin{itemize}\setlength\itemsep{0pt}
	\item \textbf{Acotar denominador}: aniadir un parametro \texttt{Q} (denominador maximo) y ajustar \texttt{k = min(k, Q/...)}.
	\item \textbf{Devolver la fraccion irreducible (no reducida)}: aniadir division por $\gcd(p, q)$ al final.
	\item \textbf{Buscar fraccion mas cercana}: aniadir comparacion entre $\text{actual}$ y $\text{candidato}$ al estilo del BS.
\end{itemize}

\paragraph{Trampas y errores frecuentes.}
\begin{itemize}\setlength\itemsep{0pt}
	\item El snippet tiene una aproximacion cruda; un mejor implementation debe iterar \textbf{exactamente} el maximo $k$ permitido.
	\item Cuidado con overflow en \texttt{(p - a) / b} si los numeros son grandes.
	\item El mediantes es irreducible sii $\gcd(a+c, b+d) = 1$; en el arbol de Stern-Brocot se preserva.
	\item La raiz del arbol es $1/1$; los hijos son $1/2$ y $2/1$.
\end{itemize}

\paragraph{Codigo.}
Ver \texttt{cpp.tex}, seccion \textit{Stern-Brocot (best fraction in interval)}.

\subsection{Combinatoria}

\subsubsection{Catalan (formula)}

\paragraph{Idea intuitiva.}
Los numeros de Catalan cuentan una gran variedad de estructuras (arboles binarios, parentesis balanceados, triangulaciones de poligono). La formula es $C_n = \frac{1}{n+1}\binom{2n}{n}$. La version recursiva es $C_n = \sum_{i=0}^{n-1} C_i C_{n-1-i}$, que da una DP en $O(n^2)$. La conexion con $\binom{2n}{n}$ viene del \textbf{lattice path counting}: $C_n$ es el numero de caminos de $(0,0)$ a $(n,n)$ que no cruzan la diagonal, y los caminos sin esa restriccion son $\binom{2n}{n}$ (elegir cuando subir).

\paragraph{Propiedades clave (invariantes).}
\begin{itemize}\setlength\itemsep{0pt}
	\item $C_0 = C_1 = 1$.
	\item $C_n$ crece rapidamente; $C_{19}$ ya pasa el millon.
	\item La formula cerrada evita la DP pero requiere multiplicacion que \textbf{siempre} es exacta (la division por $i+2$ no deja residuo).
	\item Funcion generadora: $C(x) = \sum_{n \ge 0} C_n x^n$ satisface $C(x) = 1 + x C(x)^2$.
\end{itemize}

\paragraph{Codigo clave.}
Formula iterativa (sin division intermedia):
\begin{verbatim}
long long catalanLL(int n) {
    long long c = 1;
    for (int i = 0; i < n; ++i) {
        c = c * 2 * (2*i + 1) / (i + 2);  // division siempre exacta
    }
    return c;
}
\end{verbatim}

\paragraph{Complejidad.}
\begin{itemize}\setlength\itemsep{0pt}
	\item DP: $O(N^2)$ tiempo, $O(N)$ memoria.
	\item Formula cerrada (iterativa): $O(N)$ tiempo, $O(1)$ memoria.
	\item Para un solo $C_n$: $O(n)$ (no $O(n^2)$ con DP).
\end{itemize}

\paragraph{Donde tocar para modificar.}
\begin{itemize}\setlength\itemsep{0pt}
	\item \textbf{Catalan mod p}: aniadir \texttt{MOD} y hacer \texttt{\% MOD} en cada multiplicacion. La division \texttt{/(i+2)} no se puede hacer mod $p$ directamente; conviene usar el inverso modular: \texttt{ans = ans * 2 * (2*i+1) \% MOD * inv[i+2] \% MOD;}.
	\item \textbf{Catalan con doble precision}: \texttt{catalanLL} retorna \texttt{long long}; para $n \ge 35$ desborda.
	\item \textbf{Catalan via binomial}: usar el snippet de factoriales: $C_n = \binom{2n}{n} - \binom{2n}{n+1}$ (numericamente estable mod primo).
\end{itemize}

\paragraph{Trampas y errores frecuentes.}
\begin{itemize}\setlength\itemsep{0pt}
	\item La division entera \texttt{(i+2)} en \texttt{catalanLL} es exacta solo porque matematicamente el cociente es entero. Si cambia el orden de operaciones (por ejemplo, modular primero), la division falla.
	\item Para Catalan mod $p$ NO primo, no se puede usar la recurrencia con division; usar DP pura.
	\item $C_{35}$ ya excede $2^{63}$; usar BigInt si lo necesitas.
\end{itemize}

\paragraph{Codigo.}
Ver \texttt{cpp.tex}, seccion \textit{Catalan (formula)}.

\vspace{0.5em}

\subsubsection{Factoriales + nCr mod p}

\paragraph{Idea intuitiva.}
Si $p$ es primo, $\binom{n}{k} \bmod p$ se puede calcular en $O(1)$ con factoriales precomputados: $\binom{n}{k} = \frac{n!}{k! (n-k)!} = n! \cdot (k!)^{-1} \cdot ((n-k)!)^{-1} \pmod p$. Las inversas se precomputan una vez con $fact^{-1}[n] = fact[n]^{p-2}$ y se propaga hacia abajo. La optimizacion clave es propagar $invFact[i-1] = invFact[i] \cdot i$ hacia abajo en vez de calcular $i^{p-2}$ para cada $i$.

\paragraph{Propiedades clave (invariantes).}
\begin{itemize}\setlength\itemsep{0pt}
	\item Precompute en $O(N)$, query en $O(1)$.
	\item Asume $p$ primo; si $p$ no es primo, la inversa modular no siempre existe.
	\item $\binom{n}{k} = 0$ si $k > n$.
	\item Simetria: $\binom{n}{k} = \binom{n}{n-k}$ (util para reducir $k$).
	\item $\sum_k \binom{n}{k} = 2^n$.
\end{itemize}

\paragraph{Codigo clave.}
Construir factoriales + inversas en $O(N)$:
\begin{verbatim}
fact[0] = 1;
for (int i = 1; i <= N; ++i) fact[i] = (long long)fact[i-1] * i % MOD;
invFact[N] = modpow(fact[N], MOD - 2, MOD);
for (int i = N - 1; i >= 0; --i) invFact[i] = (long long)invFact[i+1] * (i+1) % MOD;

long long nCr(int n, int k) {
    if (k < 0 || k > n) return 0;
    return fact[n] * invFact[k] % MOD * invFact[n-k] % MOD;
}
\end{verbatim}

\paragraph{Complejidad.}
\begin{itemize}\setlength\itemsep{0pt}
	\item Precompute $O(N)$, query $O(1)$.
	\item Memoria $O(N)$.
\end{itemize}

\paragraph{Donde tocar para modificar.}
\begin{itemize}\setlength\itemsep{0pt}
	\item \textbf{$n$ grande con $p$ primo pequenio} ($\le 10^6$): usar Lucas theorem.
	\item \textbf{Cambiar mod}: aniadir nuevo \texttt{MOD} y regenerar \texttt{fact}/\texttt{invFact}.
	\item \textbf{Sumar $\binom{n}{k}$ sobre $k$}: usar DP o suma geometrica.
	\item \textbf{Permutaciones}: $P(n, k) = \binom{n}{k} \cdot k! = fact[n] \cdot invFact[n-k]$.
	\item \textbf{Recurrence}: $\binom{n}{k} = \binom{n-1}{k-1} + \binom{n-1}{k}$ (util en DP).
\end{itemize}

\paragraph{Trampas y errores frecuentes.}
\begin{itemize}\setlength\itemsep{0pt}
	\item \texttt{MOD = 1} no es primo, falla en \texttt{modpow(fact[n], MOD - 2)}.
	\item \texttt{buildFact(0)} deja \texttt{fact = \{1\}}; \texttt{nCr(0, 0) = 1} OK.
	\item Si $N > p$, el factorial crece y eventualmente $\binom{n}{k} \bmod p$ puede no ser el esperado sin Lucas.
	\item Para $k > N$ pero $k \le n$, el resultado es $0$; chequealo en la funcion.
\end{itemize}

\paragraph{Codigo.}
Ver \texttt{cpp.tex}, seccion \textit{Factoriales + nCr mod p}.

\vspace{0.5em}

\subsubsection{Lucas theorem}

\paragraph{Idea intuitiva.}
Para $p$ primo y $n, k$ grandes, $\binom{n}{k} \bmod p$ se calcula descomponiendo $n$ y $k$ en base $p$: $n = n_0 + n_1 p + n_2 p^2 + \dots$, $k = k_0 + k_1 p + \dots$. Entonces $\binom{n}{k} \equiv \prod_i \binom{n_i}{k_i} \pmod p$ (\textbf{Lucas}). Cada $\binom{n_i}{k_i}$ se calcula con el snippet anterior. La demostracion usa que $(1+x)^n \equiv \prod_i (1+x^{p^i})^{n_i}$ en $\mathbb{F}_p[x]$ (por el ``freshman's dream'' $a^p \equiv a \pmod p$).

\paragraph{Propiedades clave (invariantes).}
\begin{itemize}\setlength\itemsep{0pt}
	\item Solo funciona con $p$ primo.
	\item Si algun $k_i > n_i$, el coeficiente es 0.
	\item Reduce el problema a $\log_p(n)$ coeficientes chiquitos.
	\item Cada $\binom{n_i}{k_i}$ se computa con factoriales precomputados hasta $p-1$.
\end{itemize}

\paragraph{Codigo clave.}
Lucas iterativo (descomponiendo en base $p$):
\begin{verbatim}
long long lucas(long long n, long long k, long long p) {
    long long res = 1;
    while (n > 0 || k > 0) {
        int ni = n % p, ki = k % p;
        if (ki > ni) return 0;
        res = res * nCr_small(ni, ki) % p;  // usa fact hasta p-1
        n /= p; k /= p;
    }
    return res;
}
\end{verbatim}

\paragraph{Complejidad.}
\begin{itemize}\setlength\itemsep{0pt}
	\item $O(\log_p n)$ multiplicaciones, cada una requiere un $\binom{n_i}{k_i}$ que es $O(1)$ con factoriales precomputados.
	\item Para $p = 10^6$ y $n = 10^{18}$: $\sim 12$ iteraciones.
\end{itemize}

\paragraph{Donde tocar para modificar.}
\begin{itemize}\setlength\itemsep{0pt}
	\item \textbf{$p$ no primo}: requiere extender Lucas a modulos compuestos (no trivial, basado en factorizacion).
	\item \textbf{Devolver la suma}: si necesitas $\sum_k \binom{n}{k} \bmod p$, es $2^n \bmod p$ via Fermat (cuando $p$ es primo).
	\item \textbf{Multiples queries}: precomputar todos los $\binom{n_i}{k_i}$ en una tabla y consultar.
\end{itemize}

\paragraph{Trampas y errores frecuentes.}
\begin{itemize}\setlength\itemsep{0pt}
	\item \texttt{buildFact(MOD)} solo precomputa hasta \texttt{MOD}, no hasta \texttt{n}. Pero como cada $\binom{n_i}{k_i}$ tiene $n_i < p$, eso es suficiente.
	\item Si $p$ es muy chico ($\le 100$), los factoriales repiten valores (\texttt{fact[p] = 0}); eso no rompe Lucas pero aniadir logica.
	\item Lucas \textbf{no} funciona para modulos no primos sin extension (algoritmo basado en factorizacion).
\end{itemize}

\paragraph{Codigo.}
Ver \texttt{cpp.tex}, seccion \textit{Lucas theorem}.

\vspace{0.5em}

\subsubsection{Stirling S2 (particiones)}

\paragraph{Idea intuitiva.}
$S(n, k)$ = numero de formas de particionar $n$ elementos en $k$ subconjuntos no vacios. La recurrencia es $S(n, k) = S(n-1, k-1) + k \cdot S(n-1, k)$ (o pones el elemento $n$-esimo solo en un nuevo subconjunto, o lo metes en uno de los $k$ existentes). La recurrencia sale de un caso especial de la recurrencia general de inclusion-exclusion sobre funciones sobreyectivas.

\paragraph{Propiedades clave (invariantes).}
\begin{itemize}\setlength\itemsep{0pt}
	\item $S(n, 0) = [n = 0]$.
	\item $S(n, n) = 1$, $S(n, 1) = 1$.
	\item $\sum_k S(n, k) = B_n$ (numero de Bell).
	\item $S(n, k) = \frac{1}{k!} \sum_{i=0}^{k} (-1)^{k-i} \binom{k}{i} i^n$ (formula con inclusion-exclusion).
	\item $\sum_{k=0}^{n} S(n, k) \cdot k! = $ numero de funciones sobreyectivas de $[n]$ a $[k]$.
\end{itemize}

\paragraph{Codigo clave.}
DP bidimensional clasica:
\begin{verbatim}
S2[0][0] = 1;
for (int n = 1; n <= N; ++n)
    for (int k = 1; k <= n; ++k)
        S2[n][k] = (S2[n-1][k-1] + k * S2[n-1][k]) % MOD;
\end{verbatim}

\paragraph{Complejidad.}
\begin{itemize}\setlength\itemsep{0pt}
	\item $O(N^2)$ tiempo y memoria.
	\item Para $N = 5000$: $\sim 25$MB si se usa \texttt{int}; con \texttt{long long} el doble.
\end{itemize}

\paragraph{Donde tocar para modificar.}
\begin{itemize}\setlength\itemsep{0pt}
	\item \textbf{Stirling de primera especie}: recurrencia distinta; aniadir nueva tabla.
	\item \textbf{Memory optimization}: solo guardar la fila anterior (DP en espacio $O(N)$).
	\item \textbf{Output sin mod}: si el conteo es chico, retornar \texttt{long long} y omitir \texttt{\% MOD}.
	\item \textbf{Bell numbers}: $\sum_k S(n, k)$ por fila.
\end{itemize}

\paragraph{Trampas y errores frecuentes.}
\begin{itemize}\setlength\itemsep{0pt}
	\item Olvidar inicializar \texttt{S2[0][0] = 1} y el resto en 0.
	\item Overflow sin modulo.
	\item $S(n, k)$ crece muy rapido; sin modulo desborda para $n \ge 20$.
	\item El termino $k \cdot S(n-1, k)$ es donde se acumula la mayoria del costo; en modulo $p$, esto es $O(1)$ por iteracion.
\end{itemize}

\paragraph{Codigo.}
Ver \texttt{cpp.tex}, seccion \textit{Stirling S2 (particiones)}.

\vspace{0.5em}

\subsubsection{Inclusion-exclusion over subsets}

\paragraph{Idea intuitiva.}
Para contar elementos que cumplen \textbf{alguna} de $n$ propiedades, la formula clasica de inclusion-exclusion es: $|\bigcup A_i| = \sum_{\emptyset \ne I} (-1)^{|I|+1} |\bigcap_{i \in I} A_i|$. Esto se computa iterando sobre todos los $2^n - 1$ subconjuntos no vacios y aplicando el signo correspondiente. La demostracion usa que la formula del PIE cancela exactamente las sobreposiciones multiples, dejando solo los elementos que cumplen al menos una propiedad.

\paragraph{Propiedades clave (invariantes).}
\begin{itemize}\setlength\itemsep{0pt}
	\item El usuario pasa una funcion \texttt{f(mask)} que cuenta cuantos elementos cumplen \textbf{exactamente} las propiedades en \texttt{mask}.
	\item El signo es $(-1)^{|I|+1}$, equivalente a $+1$ si $|I|$ impar, $-1$ si par.
	\item Para $n \le 20$ es factible ($2^{20} \approx 10^6$); para mas, usar otra tecnica (Mobius, formula cerrada).
	\item Variantes: ``todos cumplen'' = $(-1)^{|I|}$ (signo invertido).
\end{itemize}

\paragraph{Codigo clave.}
Inclusion-exclusion clasica sobre todos los subconjuntos:
\begin{verbatim}
long long ans = 0;
for (int mask = 1; mask < (1 << n); ++mask) {
    long long cnt = f(mask);  // |interseccion|
    if (__builtin_popcount(mask) % 2 == 1) ans += cnt;
    else ans -= cnt;
}
return ans;
\end{verbatim}

\paragraph{Complejidad.}
\begin{itemize}\setlength\itemsep{0pt}
	\item $O(2^n \cdot \text{costo de } f)$.
	\item Memoria $O(1)$ (no se almacena subconjunto).
\end{itemize}

\paragraph{Donde tocar para modificar.}
\begin{itemize}\setlength\itemsep{0pt}
	\item \textbf{Subset de subconjuntos} (elementos con un subconjunto de propiedades): usar bitmask para mapear cada elemento a un mask, luego contar ocurrencias por mask y sumar con signo.
	\item \textbf{Aniadir la interseccion}: en lugar de $|\bigcap A_i|$ exacta, aproximar o usar otra estructura.
	\item \textbf{Cambiar el signo}: si tu problema tiene inclusion-exclusion ``inversa'', invertir el patron $\pm$.
	\item \textbf{Version con Mobius}: si las propiedades son divisibilidad, $\mu(d)$ reemplaza el patron $\pm$.
\end{itemize}

\paragraph{Ejemplo.}
Contar elementos de $[1, N]$ no divisibles por ninguno de $\{2, 3, 5\}$:
\[ N - \lfloor N/2 \rfloor - \lfloor N/3 \rfloor - \lfloor N/5 \rfloor + \lfloor N/6 \rfloor + \lfloor N/10 \rfloor + \lfloor N/15 \rfloor - \lfloor N/30 \rfloor. \]

\paragraph{Trampas y errores frecuentes.}
\begin{itemize}\setlength\itemsep{0pt}
	\item El snippet dado es solo el \textbf{esqueleto} (\texttt{f} no esta implementado); hay que adaptarlo al problema.
	\item Para $n > 25$ la iteracion sobre $2^n$ se vuelve intratable; buscar otra estructura.
	\item El signo se calcula con el bit menos significativo (LSB) si la mascara lo tiene, o con \texttt{popcount} modulo 2.
	\item Si $|I|$ es par pero muy grande, considerar $2^{|I|}$ ``pesado''; aniadir memoizacion en \texttt{f}.
\end{itemize}

\paragraph{Codigo.}
Ver \texttt{cpp.tex}, seccion \textit{Inclusion-exclusion over subsets}.

\subsection{Polinomios y transformadas}

\subsubsection{FFT (iterativo, complex)}

\paragraph{Idea intuitiva.}
La \textbf{Transformada Rapida de Fourier} calcula el DFT de un vector en $O(N \log N)$. La inversa tambien. La convolucion de dos polinomios $A(x) \cdot B(x)$ se obtiene como $\text{IFFT}(\text{FFT}(A) \cdot \text{FFT}(B))$ componente a componente. La version \textbf{iterativa} reordena los elementos por bit-reversal antes y aplica ``mariposas'' (combina pares en $O(1)$). La intuicion matematica: cualquier secuencia $a_0, a_1, \dots, a_{n-1}$ se puede escribir como $\sum_k c_k \omega^{jk}$ donde $\omega = e^{2\pi i / n}$ y $c_k$ son los coeficientes de Fourier; el FFT calcula $c_k$ rapidamente usando el ``divide y venceras'' sobre la paridad de los indices.

\paragraph{Propiedades clave (invariantes).}
\begin{itemize}\setlength\itemsep{0pt}
	\item Trabaja con \texttt{complex<double>}; hay errores de redondeo que se acumulan.
	\item Tamano $N$ debe ser potencia de 2 (se hace \texttt{n <<= 1} hasta alcanzar \texttt{a.size() + b.size()}).
	\item Resultado se redondea con \texttt{llround}.
	\item Parseval: $\sum |a_i|^2 = \frac{1}{N} \sum |A_k|^2$ (util para verificar).
\end{itemize}

\paragraph{Codigo clave.}
El corazon del FFT iterativo (mariposa):
\begin{verbatim}
for (int len = 2; len <= n; len <<= 1) {
    double ang = 2 * PI / len * invert;  // invert = -1 para FFT, +1 para IFFT
    complex<double> wlen(cos(ang), sin(ang));
    for (int i = 0; i < n; i += len) {
        complex<double> w(1);
        for (int j = 0; j < len/2; ++j) {
            complex<double> u = a[i+j], v = a[i+j+len/2] * w;
            a[i+j] = u + v;
            a[i+j+len/2] = u - v;
            w *= wlen;
        }
    }
}
if (invert) for (auto& x : a) x /= n;
\end{verbatim}

\paragraph{Complejidad.}
\begin{itemize}\setlength\itemsep{0pt}
	\item $O(N \log N)$ con $N$ la primera potencia de 2 que cubre el resultado.
	\item Constante $\sim 5$ para el iterative, $\sim 1$ para recursive (pero recursive tiene overhead).
\end{itemize}

\paragraph{Donde tocar para modificar.}
\begin{itemize}\setlength\itemsep{0pt}
	\item \textbf{Modulo arbitrario}: aniadir tres FFTs y ``FFT con tres modulos'' para reconstruir el resultado con CRT (mejor precision).
	\item \textbf{Tamano mas pequeno}: si sabes que el resultado cabe en $K < a.size() + b.size() - 1$, podes usar $N = 2^{\lceil \log_2 K \rceil}$ para ahorrar.
	\item \textbf{Cambiar precision}: aniadir \texttt{long double} y \texttt{acosl(-1)} para mayor precision (marginal).
	\item \textbf{Bluestein}: para tamano no potencia de 2, bluestein permite NTT/FFT exacto.
\end{itemize}

\paragraph{Trampas y errores frecuentes.}
\begin{itemize}\setlength\itemsep{0pt}
	\item Coeficientes \textbf{negativos}: \texttt{llround(-0.5)} da 0, no -1; aniadir \texttt{floor(x + 0.5)} si hay negativos.
	\item Overflow al convertir a \texttt{long long} si los coeficientes son grandes ($\ge 2^{52}$).
	\item La precision de \texttt{double} llega a $\sim 10^{-15}$; con $N = 10^5$ el error acumulado puede ser $\sim 10^{-3}$, problematico para enteros $> 10^{15}$.
	\item Si los coeficientes son \texttt{int} ($\le 10^9$), FFT funciona; para mas, usar NTT mod primes.
\end{itemize}

\paragraph{Codigo.}
Ver \texttt{cpp.tex}, seccion \textit{FFT (iterativo, complex)}.

\vspace{0.5em}

\subsubsection{NTT (mod 998244353)}

\paragraph{Idea intuitiva.}
La \textbf{Number Theoretic Transform} es la version entera de FFT. En vez de raices de unidad complejas, usa raices modulo un primo $p$ con $p - 1$ divisible por una potencia grande de 2. El primo $998244353 = 119 \cdot 2^{23} + 1$ sirve para NTT hasta $2^{23} = 8 \cdot 10^6$. La existencia de tales primos se basa en buscar $p$ con $p - 1$ altamente divisible por 2, lo cual requiere que $p$ sea grande.

\paragraph{Propiedades clave (invariantes).}
\begin{itemize}\setlength\itemsep{0pt}
	\item \textbf{Exacto} modulo $p$ (sin redondeo).
	\item Raiz primitiva $g = 3$.
	\item Misma estructura que FFT (bit-reversal + mariposas).
	\item $\omega_n = g^{(p-1)/n}$ es una raiz $n$-esima de la unidad en $\mathbb{F}_p$.
\end{itemize}

\paragraph{Codigo clave.}
NTT iterativo (la mariposa es analoga a FFT, con multiplicacion modular):
\begin{verbatim}
for (int len = 2; len <= n; len <<= 1) {
    long long wlen = modpow(ROOT, (MOD - 1) / len, MOD);
    if (invert) wlen = modpow(wlen, MOD - 2, MOD);
    for (int i = 0; i < n; i += len) {
        long long w = 1;
        for (int j = 0; j < len / 2; ++j) {
            long long u = a[i+j], v = a[i+j+len/2] * w % MOD;
            a[i+j] = (u + v) % MOD;
            a[i+j+len/2] = (u - v + MOD) % MOD;
            w = w * wlen % MOD;
        }
    }
}
if (invert) {
    long long n_inv = modpow(n, MOD - 2, MOD);
    for (auto& x : a) x = x * n_inv % MOD;
}
\end{verbatim}

\paragraph{Complejidad.}
\begin{itemize}\setlength\itemsep{0pt}
	\item $O(N \log N)$.
	\item Para 998244353 y $N \le 2^{23}$: un par de milisegundos.
\end{itemize}

\paragraph{Donde tocar para modificar.}
\begin{itemize}\setlength\itemsep{0pt}
	\item \textbf{Otro primo}: $985661441 = 235 \cdot 2^{22} + 1$, raiz $3$; $754974721 = 45 \cdot 2^{24} + 1$, raiz $11$. Ajustar \texttt{MOD} y \texttt{PRIMITIVE\_ROOT}.
	\item \textbf{Multiplicacion de varias cosas}: NTT es lineal; aniadir cache si llamas con los mismos vectores.
	\item \textbf{Poly inv (Newton)}: ver el modulo de polynomial utilities.
	\item \textbf{Tres NTTs + CRT}: si necesitas modulo $10^9 + 7$ con coeficientes grandes, hacer NTT con tres primos y reconstruir.
\end{itemize}

\paragraph{Trampas y errores frecuentes.}
\begin{itemize}\setlength\itemsep{0pt}
	\item El ``exponente magico'' \texttt{(MOD - 1) / len} requiere que \texttt{MOD - 1} sea divisible por \texttt{len}; para 998244353, \texttt{len} debe ser potencia de 2 hasta $2^{23}$.
	\item Si \texttt{len} no divide \texttt{MOD - 1}, el algoritmo da basura.
	\item Confundir la raiz primitiva $g$ con la raiz $n$-esima $g^{(p-1)/n}$; esta ultima depende del nivel.
	\item En el inverse, olvidar el factor $1/n$ da basura escalada.
\end{itemize}

\paragraph{Codigo.}
Ver \texttt{cpp.tex}, seccion \textit{NTT (mod 998244353)}.

\vspace{0.5em}

\subsubsection{Polynomial utilities (mod p)}

\paragraph{Idea intuitiva.}
Capa de utilidades sobre NTT/FFT: suma, resta, \textbf{inversa} (Newton iterativo), derivada, integral. La inversa usa que si $B \cdot B^{-1} \equiv 1 \pmod{x^n}$, entonces aproximar a $2n$ y doblar la precision en cada paso. La recurrencia de Newton: si $r_k$ es inversa mod $x^{2^k}$, entonces $r_{k+1} \equiv r_k \cdot (2 - B \cdot r_k) \pmod{x^{2^{k+1}}}$. La intuicion: $r_{k+1}$ satisface $B \cdot r_{k+1} \equiv 1 + (B \cdot r_k - 1)^2 \pmod{x^{2^{k+1}}}$ y el error cuadratico decae mas rapido.

\paragraph{Propiedades clave (invariantes).}
\begin{itemize}\setlength\itemsep{0pt}
	\item \texttt{polyInv(b, n)}: asume \texttt{b[0]} invertible (no cero).
	\item Usa NTT internamente; misma complejidad $O(N \log N)$.
	\item \texttt{polyInteg} necesita el inverso de \texttt{i+1} mod $p$; se precomputa en una \texttt{static vector} que crece a demanda.
	\item Cada iteracion de Newton dobla la precision.
\end{itemize}

\paragraph{Codigo clave.}
Newton iterativo para inversa polinomial:
\begin{verbatim}
vector<int> polyInv(const vector<int>& b, int n) {
    vector<int> r(1, modpow(b[0], MOD - 2, MOD));
    int cur = 1;
    while (cur < n) {
        cur <<= 1;
        vector<int> b_cut(b.begin(), b.begin() + min((int)b.size(), cur));
        vector<int> nr = r * r % MOD * b_cut % MOD;  // 2*r - r^2*b
        // nr = 2*r - r^2 * b mod x^cur
        nr.resize(cur);
        r = (2LL * r[0] * ...);  // simplificacion algebraica
        r = nr;
    }
    r.resize(n);
    return r;
}
\end{verbatim}

\paragraph{Complejidad.}
\begin{itemize}\setlength\itemsep{0pt}
	\item Operaciones polinomiales $O(N \log N)$.
	\item Newton converge en $O(\log n)$ pasos; cada paso $O(N \log N)$.
\end{itemize}

\paragraph{Donde tocar para modificar.}
\begin{itemize}\setlength\itemsep{0pt}
	\item \textbf{Division}: aniadir \texttt{polyDiv(a, b)} usando la inversa de $b$ truncada.
	\item \textbf{Multi-point evaluation}: $O(N \log^2 N)$ con divide y venceras.
	\item \textbf{Composition}: muy dificil en general; no recomendado en contests.
	\item \textbf{Logaritmo polinomial}: $\log(f) = \int f'/f$, requiere inversa de $f$.
\end{itemize}

\paragraph{Trampas y errores frecuentes.}
\begin{itemize}\setlength\itemsep{0pt}
	\item \texttt{polyInv} asume \texttt{b[0]} no es 0 modulo $p$; sino, no existe la inversa.
	\item Olvidar \texttt{r.resize(cur)} deja residuos de la iteracion anterior.
	\item En Newton, $r$ se modifica en cada paso; no reusar la misma variable.
	\item Los vectores de NTT deben tener tamano potencia de 2; aniadir padding si hace falta.
\end{itemize}

\paragraph{Codigo.}
Ver \texttt{cpp.tex}, seccion \textit{Polynomial utilities (mod p)}.

\subsection{Algebra lineal}

\subsubsection{Gauss-Jordan (mod p)}

\paragraph{Idea intuitiva.}
Reduccion gaussiana para resolver sistemas $Ax = b$ y calcular determinante y rango. En la version \textbf{Gauss-Jordan}, se elimina \textbf{tambien} hacia arriba (no solo triangular inferior). El algoritmo opera sobre la matriz aumentada $[A | b]$ modulo $p$. La demostracion: la operacion ``aniadir multiplo de una fila a otra'' preserva el espacio de soluciones, asi que tras aplicar las operaciones de eliminacion, la matriz reducida tiene las mismas soluciones que la original.

\paragraph{Propiedades clave (invariantes).}
\begin{itemize}\setlength\itemsep{0pt}
	\item Rango \texttt{r}: numero de filas no cero despues de la reduccion.
	\item \textbf{3 casos} para $Ax = b$: $\sin$ solucion ($r < n$ y ultima columna no cero en filas nulas), solucion unica ($r = n$), infinitas ($r < n$ y filas nulas tienen $b = 0$).
	\item Determinante = producto de la diagonal (con signo si se swappean filas).
	\item Las variables libres quedan con valor $0$ por defecto.
\end{itemize}

\paragraph{Codigo clave.}
Eliminacion Gauss-Jordan (con pivoteo parcial):
\begin{verbatim}
for (int col = 0, row = 0; col < n && row < n; ++col) {
    int sel = row;
    while (sel < n && fabs(a[sel][col]) < eps) ++sel;
    if (sel == n) continue;
    swap(a[sel], a[row]); swap(b[sel], b[row]);
    // normalizar fila row
    long long inv = modpow(a[row][col], MOD - 2, MOD);
    for (int j = col; j < n; ++j) a[row][j] = a[row][j] * inv % MOD;
    b[row] = b[row] * inv % MOD;
    // eliminar otras filas
    for (int i = 0; i < n; ++i) if (i != row && a[i][col])
        for (int j = col; j < n; ++j)
            a[i][j] = (a[i][j] - a[row][j] * a[i][col]) % MOD;
    b[i] = (b[i] - b[row] * a[i][col]) % MOD;
    ++row;
}
\end{verbatim}

\paragraph{Complejidad.}
\begin{itemize}\setlength\itemsep{0pt}
	\item $O(N^3)$ por cubica, $O(N^2)$ para triangularizar.
	\item Para $N \le 500$: $\sim 10^8$ operaciones, factible.
\end{itemize}

\paragraph{Donde tocar para modificar.}
\begin{itemize}\setlength\itemsep{0pt}
	\item \textbf{Sin modulo}: aniadir \texttt{double} y \texttt{abs} para determinantes reales.
	\item \textbf{Encontrar la solucion particular con libres en 0}: el snippet ya lo hace; si quieres una \textbf{otra} particular, aniadir un parametro \texttt{freeVars} con valores iniciales.
	\item \textbf{Rango sobre matriz sin aumentar}: pasar solo \texttt{A}, no \texttt{aug}.
	\item \textbf{Matriz de Vandermonde}: estructura especial con inversa explicita.
\end{itemize}

\paragraph{Trampas y errores frecuentes.}
\begin{itemize}\setlength\itemsep{0pt}
	\item Si el pivote es 0 pero hay un swap posible, \textbf{no olvidar} el swap.
	\item Division por cero: solo se divide por \texttt{inv(a[row][col])}, que existe porque \texttt{a[row][col]} es no cero y $p$ es primo (por eso \texttt{modpow(a, p-2, p)}).
	\item El signo del determinante se invierte por cada swap; aniadir multiplicacion por $-1$.
	\item Sin pivoteo, matrices mal condicionadas pueden dar errores numericos grandes (con doubles).
\end{itemize}

\paragraph{Codigo.}
Ver \texttt{cpp.tex}, seccion \textit{Gauss-Jordan (mod p)}.

\vspace{0.5em}

\subsubsection{Matrix exponentiation}

\paragraph{Idea intuitiva.}
Para calcular $M^k$ en $O(N^3 \log k)$, se usa la misma idea que \texttt{modpow} pero la multiplicacion es de matrices. Asi, $M^k$ se descompone en bits y se multiplican las matrices correspondientes. La idea subyacente: $M^k$ actua sobre vectores; si tenemos una recurrencia lineal $\mathbf{v}_{n+1} = A \mathbf{v}_n$, entonces $\mathbf{v}_k = A^k \mathbf{v}_0$, asi que calcular $A^k$ nos da la respuesta en $O(\log k)$.

\paragraph{Propiedades clave (invariantes).}
\begin{itemize}\setlength\itemsep{0pt}
	\item Multiplicacion asociativa: $(M^a)(M^b) = M^{a+b}$.
	\item Identidad es la matriz diagonal con 1s.
	\item Si $M$ es la matriz de transicion de un sistema lineal (recurrencia, automata), $M^k$ da el estado tras $k$ pasos.
	\item Si $M$ es diagonalizable, $M^k = P D^k P^{-1}$; util para closed-form.
\end{itemize}

\paragraph{Codigo clave.}
El squaring matricial:
\begin{verbatim}
Mat operator*(const Mat& o) const {
    Mat r(n, vector<long long>(n, 0));
    for (int i = 0; i < n; ++i)
        for (int k = 0; k < n; ++k) if (a[i][k])
            for (int j = 0; j < n; ++j)
                r[i][j] = (r[i][j] + a[i][k] * o.a[k][j]) % MOD;
    return r;
}
Mat mpow(Mat M, long long e) {
    Mat R = identity(M.n);
    while (e > 0) {
        if (e & 1) R = R * M;
        M = M * M;
        e >>= 1;
    }
    return R;
}
\end{verbatim}

\paragraph{Complejidad.}
\begin{itemize}\setlength\itemsep{0pt}
	\item $O(N^3 \log k)$.
	\item Para $N = 50$, $k = 10^{18}$: $\sim 6 \cdot 10^6$ operaciones.
\end{itemize}

\paragraph{Donde tocar para modificar.}
\begin{itemize}\setlength\itemsep{0pt}
	\item \textbf{Aplicar a recurrencia}: para $F_n = a F_{n-1} + b F_{n-2}$, construir $M = \begin{pmatrix} a & b \\ 1 & 0 \end{pmatrix}$, $M^n$ da $F_n$ en \texttt{M[0][0] * F\_1 + M[0][1] * F\_0}.
	\item \textbf{Multiplicacion sin la optimizacion \texttt{if (a[i][k])}}: para matrices densas, la guarda no ayuda; quitarla.
	\item \textbf{Matrices pequenas} ($N \le 3$): hardcodear la multiplicacion puede ser mas rapido.
	\item \textbf{Camino count en grafo}: $A^k$ da el numero de caminos de largo $k$ entre pares de nodos.
\end{itemize}

\paragraph{Trampas y errores frecuentes.}
\begin{itemize}\setlength\itemsep{0pt}
	\item Multiplicacion \texttt{a[i][k] * b[k][j]} puede overflow; el snippet usa \texttt{\% MOD} al final, pero el producto intermedio puede pasar $2^{63}$ si los valores son grandes y \texttt{MOD} es chico. Mejor hacer el modulo en cada paso.
	\item Para grafos con $N = 10^3$ y aristas, $A^k$ no entra en memoria si no se usa sparse matrix.
	\item Construir mal la matriz de transicion es la fuente mas comun de bugs; verificar con casos pequenos.
\end{itemize}

\paragraph{Codigo.}
Ver \texttt{cpp.tex}, seccion \textit{Matrix exponentiation}.

\subsection{Juegos y teoria de juegos}

\subsubsection{Sprague-Grundy template}

\paragraph{Idea intuitiva.}
Para juegos \textbf{imparciales} (ambos jugadores tienen los mismos movimientos), el teorema de Sprague-Grundy dice que cada posicion tiene un numero (el \textbf{Grundy} o \textbf{mex}) y la posicion es \textbf{ganadora} para el jugador actual sii el XOR de los Grundy de los componentes es no cero. El \textbf{mex} (``minimum excludant'') es el menor entero no negativo que no esta en el conjunto de Grundy de las posiciones alcanzables. La demostracion: para cada posicion, hay un movimiento a cada Grundy menor; la posicion perdedora es exactamente Grundy = 0.

\paragraph{Propiedades clave (invariantes).}
\begin{itemize}\setlength\itemsep{0pt}
	\item Posiciones terminales (sin movimientos) tienen Grundy = 0.
	\item La composicion XOR de Grundys da el Grundy del juego disjunto.
	\item \textbf{Caso base}: si \texttt{state} no tiene transiciones, Grundy = 0.
	\item Si Grundy $\ne 0$, hay un movimiento a una posicion con Grundy $0$ (estrategia ganadora).
	\item Si Grundy $= 0$, todos los movimientos van a Grundy $\ne 0$.
\end{itemize}

\paragraph{Codigo clave.}
Memoizacion clasica:
\begin{verbatim}
unordered_map<long long, long long> memo;
long long grundy(long long state) {
    if (memo.count(state)) return memo[state];
    set<long long> reachable;
    // populate 'reachable' with Grundy of all moves
    long long g = 0;
    while (reachable.count(g)) ++g;
    return memo[state] = g;
}
// Ganador: grundy(start) != 0
\end{verbatim}

\paragraph{Complejidad.}
\begin{itemize}\setlength\itemsep{0pt}
	\item Depende del juego: $O(\text{estados} \cdot \text{transiciones})$ con memoization.
	\item Para Nim con $N$ heaps: $O(N)$.
\end{itemize}

\paragraph{Donde tocar para modificar.}
\begin{itemize}\setlength\itemsep{0pt}
	\item \textbf{Adaptar al juego concreto}: reemplazar el comentario \texttt{// populate 'reachable' ...} con la logica del juego (ej. para Nim: aniadir todos los heap sizes reducidos).
	\item \textbf{Memoization con map}: usar \texttt{unordered\_map<long long, long long>} si el estado es grande.
	\item \textbf{Computar Grundy de varios estados independientes}: para Nim, es el XOR de los heap sizes; para Kayles, hay tabla precomputada.
	\item \textbf{Encontrar movimiento ganador}: aniadir busqueda del Grundy $0$ accesible.
\end{itemize}

\paragraph{Trampas y errores frecuentes.}
\begin{itemize}\setlength\itemsep{0pt}
	\item Olvidar memoizar: el algoritmo sin cache es exponencial.
	\item Considerar \texttt{unordered\_map} vs \texttt{map} segun el tamano del espacio de estados.
	\item Si dos juegos tienen el mismo Grundy, no necesariamente son ``iguales'' (mex lo captura).
	\item \texttt{mex} sobre un set vacio es 0 (caso base).
\end{itemize}

\paragraph{Codigo.}
Ver \texttt{cpp.tex}, seccion \textit{Sprague-Grundy template}.

\vspace{0.5em}

\subsubsection{Small game tables (Kayles SG values)}

\paragraph{Idea intuitiva.}
Cuando el juego tiene \textbf{pocas} posiciones conocidas, conviene \textbf{hardcodear} la tabla de Grundys para responder en $O(1)$ por query. Kayles es un juego clasico donde tiras un bolo y tumbas uno o dos adyacentes; sus Grundys para $n \le 71$ son bien conocidos y estan precomputados. La razon de precomputar: el Grundy de Kayles satisface una recurrencia compleja y no se puede calcular con formula cerrada simple.

\paragraph{Propiedades clave (invariantes).}
\begin{itemize}\setlength\itemsep{0pt}
	\item Solo util para Kayles; otros juegos necesitan su propia tabla.
	\item La tabla se obtuvo experimentalmente (o por recurrencia) y se almacena literal en el codigo.
	\item Los Grundys tienen \textbf{periodo} tras cierto $n$ (e.g., Kayles tiene periodo $12$ tras $n = 72$).
\end{itemize}

\paragraph{Codigo clave.}
Si el problema tiene $n$ en $[1, N]$ con $N$ conocido, precomputar la tabla:
\begin{verbatim}
vector<int> grundy(N + 1, 0);
for (int n = 1; n <= N; ++n) {
    set<int> s;
    for (int k = 0; k < n; ++k) {
        // un bolo en k: kayles se parte en [0, k-1] y [k+1, n-1]
        s.insert(grundy[k] ^ grundy[n-1-k]);
        if (k+1 < n) s.insert(grundy[k] ^ grundy[n-2-k]);
    }
    int g = 0; while (s.count(g)) ++g;
    grundy[n] = g;
}
\end{verbatim}

\paragraph{Complejidad.}
\begin{itemize}\setlength\itemsep{0pt}
	\item $O(1)$ por query (con tabla hardcodeada).
	\item Precompute $O(N^2)$.
\end{itemize}

\paragraph{Donde tocar para modificar.}
\begin{itemize}\setlength\itemsep{0pt}
	\item \textbf{Otro juego}: aniadir otra tabla literal (Nim, Dawson's Kayles, Grundy's Game, etc.).
	\item \textbf{Tabla generada por codigo}: aniadir un precompute al inicio que llene la tabla con el template Sprague-Grundy.
	\item \textbf{Usar periodicidad}: si el Grundy tiene periodo $P$ a partir de $N_0$, aniadir \texttt{int g = (n >= N0) ? tabla[(n - N0) \% P] : tabla[n];}.
\end{itemize}

\paragraph{Trampas y errores frecuentes.}
\begin{itemize}\setlength\itemsep{0pt}
	\item Si el problema da $n$ fuera de rango, aniadir bounds check antes de indexar.
	\item La tabla es especifica de la variante de Kayles (original, Dawson's, etc.); verificar cual es.
	\item En algunos juegos (Wythoff), la estrategia no es Grundy sino un par de moves explicitos.
\end{itemize}

\paragraph{Codigo.}
Ver \texttt{cpp.tex}, seccion \textit{Small game tables (Kayles SG values)}.

% ============================================================
\section{Estructuras de datos}
% ============================================================

\subsection{Ordenado y contenedores}

\subsubsection{Ordered Set (pbds)}

\paragraph{Idea intuitiva.}
La \textbf{policy-based data structure} de GCC es un arbol rojo-negro (rb\_tree) con dos operaciones extra: \texttt{find\_by\_order(k)} devuelve un iterador al $k$-esimo elemento (0-indexed) en $O(\log n)$, y \texttt{order\_of\_key(x)} cuenta cuantos elementos son \textbf{menores} que $x$ en $O(\log n)$. Equivale a tener un \texttt{set} con ``rango'' y ``rango por posicion'' gratis. La estructura subyacente: cada nodo lleva un campo \texttt{prioridad} (aleatoria o fija) y el arbol mantiene la propiedad de heap por prioridad; eso da las estadisticas de orden en $O(\log n)$.

\paragraph{Propiedades clave (invariantes).}
\begin{itemize}\setlength\itemsep{0pt}
	\item Solo en GCC/clang (header \texttt{ext/pb\_ds/...}); \textbf{no} en MSVC.
	\item \texttt{less\_equal<int>} permite multiset; ojo con \texttt{erase} que ahi usa \texttt{upper\_bound}.
	\item Tamano maximo $\sim 10^7$ elementos (limite de la libreria).
	\item \texttt{find\_by\_order(k)} con $k \ge$ tamano devuelve \texttt{s.end()}.
\end{itemize}

\paragraph{Codigo clave.}
Declaracion y uso tipico:
\begin{verbatim}
#include <ext/pb_ds/assoc_container.hpp>
#include <ext/pb_ds/tree_policy.hpp>
using namespace __gnu_pbds;
tree<int, null_type, less<int>, rb_tree_tag, tree_order_statistics_node_update> s;
s.insert(5);
cout << s.order_of_key(3) << "\n";      // 0
cout << *s.find_by_order(0) << "\n";    // 5
\end{verbatim}

\paragraph{Complejidad.}
\begin{itemize}\setlength\itemsep{0pt}
	\item $O(\log n)$ por operacion.
	\item Memoria: $\sim 5x$ un set normal (cada nodo guarda prioridad).
\end{itemize}

\paragraph{Donde tocar para modificar.}
\begin{itemize}\setlength\itemsep{0pt}
	\item \textbf{Multiset}: cambiar \texttt{less<int>} por \texttt{less\_equal<int>} y borrar con \texttt{s.erase(s.upper\_bound(x))}.
	\item \textbf{Estadisticas de orden}: \texttt{find\_by\_order(k)} da el $k$-esimo, util para mantener el $k$-esimo menor dinamicamente.
	\item \textbf{Acceso por rango}: aniadir un wrapper que combine \texttt{lower\_bound} + \texttt{order\_of\_key}.
	\item \textbf{Tipo custom}: aniadir hash si quieres una unordered\_map con order statistics (no funciona directo).
\end{itemize}

\paragraph{Trampas y errores frecuentes.}
\begin{itemize}\setlength\itemsep{0pt}
	\item \texttt{find\_by\_order(n)} (con $n$ fuera de rango) es UB.
	\item No olvidar \texttt{\#include <ext/pb\_ds/assoc\_container.hpp>} y \texttt{<ext/pb\_ds/trie\_policy.hpp>}.
	\item El namespace \texttt{\_\_gnu\_pbds} es feo; ``using'' acelera pero contamina el scope global.
	\item Algunos jueces online (Codeforces, AtCoder) lo soportan; otros no.
\end{itemize}

\paragraph{Codigo.}
Ver \texttt{cpp.tex}, seccion \textit{Ordered Set (pbds)}.

\vspace{0.5em}

\subsubsection{Basis (XOR basis)}

\paragraph{Idea intuitiva.}
Una \textbf{base de XOR} mantiene un conjunto de mascaras (numeros) tal que cualquier XOR de un subconjunto del input se puede expresar como XOR de un subconjunto de la base. La idea es similar a la base de un espacio vectorial sobre $\mathbb{F}_2$. Construir: insertar cada mascara $m$ recorriendo los bits de mayor a menor; si el bit $i$ de $m$ esta ``libre'' en la base, se aniade; sino, se hace $m \oplus basis[i]$ y se continua. La estructura formal: el conjunto de XORs posibles es un subespacio vectorial sobre $\mathbb{F}_2$, y la base lo describe con $\le \log_2(\text{max value})$ elementos.

\paragraph{Propiedades clave (invariantes).}
\begin{itemize}\setlength\itemsep{0pt}
	\item Tamano maximo de la base: $\le 64$ (bits) o $\le \log_2(\text{max value})$.
	\item Base \textbf{minimal}: la construccion garantiza que cada vector tiene un unico bit alto distinto.
	\item El subespacio generado es invariante.
	\item Numero de elementos en la base = dimension del subespacio generado.
\end{itemize}

\paragraph{Codigo clave.}
Algoritmo de insercion (Gaussian elimination sobre $\mathbb{F}_2$):
\begin{verbatim}
void insert(long long x) {
    for (int i = MAXB; i >= 0; --i) {
        if (!(x & (1LL << i))) continue;
        if (!basis[i]) { basis[i] = x; return; }
        x ^= basis[i];
    }
    // x es combinacion lineal de los basis actuales
}
\end{verbatim}

\paragraph{Complejidad.}
\begin{itemize}\setlength\itemsep{0pt}
	\item Insertar una mascara: $O(\text{bits})$. Construir la base completa: $O(N \cdot \text{bits})$.
	\item Query (maximo XOR): $O(\text{bits})$ si la base esta ordenada.
\end{itemize}

\paragraph{Donde tocar para modificar.}
\begin{itemize}\setlength\itemsep{0pt}
	\item \textbf{Obtener el maximo XOR de un subconjunto}: ordenar la base por bit alto descendente y aplicar XOR si aumenta.
	\item \textbf{Obtener el minimo XOR positivo}: el minimo distinto de 0 es el minimo \texttt{basis[i]} no cero.
	\item \textbf{Cuantos subconjuntos dan XOR $k$}: con la base, reducir $k$ y contar grados de libertad.
	\item \textbf{XOR basis en 32 bits}: cambiar \texttt{int d = 64} a \texttt{int d = 32}.
	\item \textbf{K-esimo XOR}: ordenar la base y contar las combinaciones $\le k$ (aplicacion clasica).
\end{itemize}

\paragraph{Trampas y errores frecuentes.}
\begin{itemize}\setlength\itemsep{0pt}
	\item Si los numeros son \texttt{long long}, usar \texttt{1LL << i}.
	\item Si se quiere la base \textbf{reducible} (que no necesariamente tenga vectores con unico bit alto), usar otra variante.
	\item Insertar en orden descendente de bits asegura la propiedad ``unico bit alto''.
	\item Si dos mascaras tienen el mismo bit alto, una se descarta; la informacion se preserva en el XOR.
\end{itemize}

\paragraph{Codigo.}
Ver \texttt{cpp.tex}, seccion \textit{Basis (XOR basis)}.

\vspace{0.5em}

\subsubsection{Trie}

\paragraph{Idea intuitiva.}
Un \textbf{trie} (o prefix tree) almacena strings donde cada nodo tiene hijos por cada caracter posible (tipicamente 26 letras). Permite insercion, busqueda exacta y busqueda por prefijo en $O(|word|)$. La estructura formal: cada nodo representa un prefijo; las aristas son los caracteres; el flag \texttt{isEnd} indica si ese prefijo completo es una palabra valida. La ventaja sobre un set de strings: encontrar todos los strings con un prefijo dado en tiempo proporcional al prefijo + output.

\paragraph{Propiedades clave (invariantes).}
\begin{itemize}\setlength\itemsep{0pt}
	\item Cada nodo tiene \texttt{links[26]} (punteros a hijos) y un flag \texttt{isEnd}.
	\item Memoria: $O(\text{total chars})$; cada nodo cuesta $\sim 26$ punteros (puede comprimirse con \texttt{map<char, Node*>} si el alfabeto es disperso).
	\item Las busquedas descienden siguiendo caracteres; si algun enlace es nulo, falla.
	\item El camino de la raiz a un nodo es exactamente el prefijo que el nodo representa.
\end{itemize}

\paragraph{Codigo clave.}
Insercion y busqueda:
\begin{verbatim}
void insert(const string& s) {
    Node* v = root;
    for (char c : s) {
        int k = c - 'a';
        if (!v->links[k]) v->links[k] = new Node();
        v = v->links[k];
    }
    v->isEnd = true;
}
bool search(const string& s) {
    Node* v = root;
    for (char c : s) {
        int k = c - 'a';
        if (!v->links[k]) return false;
        v = v->links[k];
    }
    return v->isEnd;
}
\end{verbatim}

\paragraph{Complejidad.}
\begin{itemize}\setlength\itemsep{0pt}
	\item Insercion y busqueda: $O(|s|)$ donde $|s|$ es el largo del string.
	\item Memoria total: $O(\sum |s_i|)$ nodos para $N$ strings.
\end{itemize}

\paragraph{Donde tocar para modificar.}
\begin{itemize}\setlength\itemsep{0pt}
	\item \textbf{Alfabeto mas grande}: cambiar \texttt{links[26]} por \texttt{links[256]} y eliminar \texttt{ch - 'a'}.
	\item \textbf{Trie de bits}: para problemas con mascaras, usar \texttt{links[2]} y trabajar con bits en vez de caracteres.
	\item \textbf{Trie persistente}: aniadir \texttt{vector<Node*> versions;} y clonar nodos en cada insercion (como el persistent segtree).
	\item \textbf{Delete}: implementar recursion que borre nodos sin hijos; cuidado con nodos compartidos por varias palabras.
	\item \textbf{Prefijos con peso}: aniadir \texttt{sumWeight} en cada nodo para queries de suma por prefijo.
\end{itemize}

\paragraph{Trampas y errores frecuentes.}
\begin{itemize}\setlength\itemsep{0pt}
	\item En este snippet los nodos se crean con \texttt{new} y nunca se liberan; en contests esta bien, pero en produccion aniadir destructor.
	\item El alfabeto tiene que estar acotado; sino, el trie explota en memoria.
	\item Si los strings son de tamano $\le 10^6$, el trie tendra nodos por cada caracter, no solo por palabra.
	\item \texttt{search(s)} vs \texttt{startsWith(s)}: el primero requiere \texttt{isEnd}, el segundo no.
\end{itemize}

\paragraph{Codigo.}
Ver \texttt{cpp.tex}, seccion \textit{Trie}.

\subsection{Union-Find}

\subsubsection{DSU classico}

\paragraph{Idea intuitiva.}
El \textbf{Disjoint Set Union} (Union-Find) mantiene una particion de elementos en conjuntos disjuntos. Las dos operaciones basicas son \texttt{find(x)} (devuelve el representante del conjunto de $x$) y \texttt{union(x, y)} (fusiona los conjuntos). La optimizacion \textbf{path compression} aplana el arbol durante \texttt{find}, y \textbf{union by size/rank} mantiene el arbol balanceado. Combinadas, dan costo \textbf{amortizado} $O(\alpha(n))$ por operacion, casi $O(1)$. La cota de Ackermann aparece porque las alturas de los arboles son \textbf{extremadamente pequenas} (del orden de $\alpha(n)$ que para $n < 10^{600}$ es $\le 5$).

\paragraph{Propiedades clave (invariantes).}
\begin{itemize}\setlength\itemsep{0pt}
	\item \texttt{parent[i]} apunta al padre en el arbol; la raiz es \texttt{parent[i] == i}.
	\item \texttt{size[i]} solo es valido en la raiz.
	\item \textbf{Representante unico}: cada conjunto tiene exactamente una raiz.
	\item Union by size garantiza que la altura del arbol es $O(\log n)$ peor caso (antes de path compression).
\end{itemize}

\paragraph{Codigo clave.}
DSU con union by size y path compression:
\begin{verbatim}
int find(int x) { return parent[x] == x ? x : parent[x] = find(parent[x]); }
void unite(int a, int b) {
    a = find(a); b = find(b);
    if (a == b) return;
    if (size[a] < size[b]) swap(a, b);
    parent[b] = a;
    size[a] += size[b];
}
bool connected(int a, int b) { return find(a) == find(b); }
\end{verbatim}

\paragraph{Complejidad.}
\begin{itemize}\setlength\itemsep{0pt}
	\item $O(\alpha(n))$ amortizado, donde $\alpha$ es la inversa de Ackermann.
	\item En la practica: $\sim 4$ operaciones por llamada para $n \le 10^6$.
\end{itemize}

\paragraph{Donde tocar para modificar.}
\begin{itemize}\setlength\itemsep{0pt}
	\item \textbf{Aniadir informacion al conjunto} (ej. tamano del conjunto ya esta; suma de los elementos; etc.): aniadir campos al struct y mantenerlos en \texttt{union}.
	\item \textbf{DSU para borrar}: borrar el ultimo elemento es trivial, borrar en medio requiere reasignar o rehacer.
	\item \textbf{Weighted DSU}: para resolver ecuaciones tipo $x \equiv a \pmod n$ con restricciones, mantener el ``peso'' relativo al padre.
	\item \textbf{DSU on next}: aniadir \texttt{next[i]} para saltar elementos borrados (proximo no en el mismo set).
\end{itemize}

\paragraph{Trampas y errores frecuentes.}
\begin{itemize}\setlength\itemsep{0pt}
	\item \texttt{find} recursivo puede overflow del stack para $n$ muy grande; cambiar a iterativo si pasa.
	\item Olvidar inicializar \texttt{size[i] = 1}.
	\item Si \texttt{size[a] == size[b]}, aniadir tie-break (e.g., por id) para determinismo.
	\item Confundir union por size con union por rank; en size, se compara el tamano del set, no la altura.
\end{itemize}

\paragraph{Codigo.}
Ver \texttt{cpp.tex}, seccion \textit{DSU classico}.

\vspace{0.5em}

\subsection{Arboles de segmentos}

\subsubsection{Segment Tree Add + Sum}

\paragraph{Idea intuitiva.}
Un \textbf{segment tree} es un arbol binario donde cada nodo cubre un rango del arreglo. La raiz cubre todo $[0, n)$; cada nodo interno $[l, r]$ tiene dos hijos $[l, m]$ y $[m+1, r]$. Para mantener la \textbf{suma} del rango y soportar \textbf{add en rango} eficientemente, se usa \textbf{lazy propagation}: cuando un update cubre totalmente un nodo, se aniade el delta a un campo \texttt{lazy} y se actualiza el nodo. La propagacion real se hace ``bajando'' solo cuando es necesario. La intuicion geometrica: el arbol es una particion jerarquica del arreglo, y cada nivel dobla la granularidad.

\paragraph{Propiedades clave (invariantes).}
\begin{itemize}\setlength\itemsep{0pt}
	\item \texttt{st[node].val}: suma del rango del nodo.
	\item \texttt{st[node].lazy}: pendiente de propagar a los hijos.
	\item La \texttt{push} aplica el lazy al nodo actual y propaga a los hijos.
	\item El arreglo tiene tamano $\le 4n$ (la cota $4n$ es pesimista; $2n$ es suficiente para tamano potencia de 2).
\end{itemize}

\paragraph{Codigo clave.}
El corazon del lazy propagation:
\begin{verbatim}
void push(int node, int l, int r) {
    if (st[node].lazy == 0 || l == r) return;
    int m = (l + r) / 2;
    st[node*2].val += st[node].lazy * (m - l + 1);
    st[node*2].lazy += st[node].lazy;
    st[node*2+1].val += st[node].lazy * (r - m);
    st[node*2+1].lazy += st[node].lazy;
    st[node].lazy = 0;
}
\end{verbatim}

\paragraph{Complejidad.}
\begin{itemize}\setlength\itemsep{0pt}
	\item Build $O(n)$, update y query $O(\log n)$ amortizado.
	\item Memoria $O(4n)$.
\end{itemize}

\paragraph{Donde tocar para modificar.}
\begin{itemize}\setlength\itemsep{0pt}
	\item \textbf{Cambiar suma por max/min/xor}: solo cambiar la linea \texttt{st[node].val = st[node*2].val + st[node*2+1].val;} por la operacion deseada (con identidad apropriada).
	\item \textbf{Cambiar add range por set range}: cambiar el \texttt{push} y la operacion en el rango.
	\item \textbf{Aniadir operacion ``maximo en rango''}: aniadir nuevo metodo \texttt{queryMax(l, r)}.
	\item \textbf{Tamano no potencia de 2}: el snippet fuerza $n$ a potencia de 2 con \texttt{while (\_\_builtin\_popcount(n) != 1) n++;}; si no quieres padding, aniadir la logica de ``nodo vacio''.
	\item \textbf{Segment tree con templates}: parametrizar la operacion y la identidad para reutilizar.
\end{itemize}

\paragraph{Trampas y errores frecuentes.}
\begin{itemize}\setlength\itemsep{0pt}
	\item En \texttt{push}, la condicion \texttt{if(l != r)} es crucial: no propagar a hijos si el nodo es hoja (evita UB por \texttt{st[node*2]} fuera de rango).
	\item Olvidar \texttt{push} antes de acceder a los hijos da valores incorrectos.
	\item Si combinas dos operaciones (e.g., add + set), el orden de aplicacion importa; documentar.
	\item El \texttt{lazy} inicial debe ser la identidad de la operacion (0 para add, INF para set).
\end{itemize}

\paragraph{Codigo.}
Ver \texttt{cpp.tex}, seccion \textit{Segment Tree Add + Sum}.

\vspace{0.5em}

\subsubsection{Segment Tree Lazy Assign + Sum}

\paragraph{Idea intuitiva.}
Misma estructura que el anterior pero la operacion es \textbf{assign} (setear todos los elementos del rango a un valor), no add. Esto \textbf{invalida} cualquier lazy add pendiente: si haces \texttt{assign} sobre un rango donde ya hay \texttt{add} pendiente, el \texttt{add} se descarta. La razon: \texttt{assign} define un nuevo ``valor base'', asi que cualquier acumulacion previa pierde sentido.

\paragraph{Propiedades clave (invariantes).}
\begin{itemize}\setlength\itemsep{0pt}
	\item \texttt{lazy == LLONG\_MAX} indica ``sin asignacion pendiente'' (valor centinela).
	\item Al asignar, se reemplaza el valor del nodo por \texttt{(r-l+1) * val} y se borran lazies anteriores.
	\item Solo se aplica assign en el nodo si no se esta propagando otra operacion.
\end{itemize}

\paragraph{Codigo clave.}
El push de assign (anula lazies en hijos):
\begin{verbatim}
void push(int node, int l, int r) {
    if (st[node].lazy == LLONG_MAX || l == r) return;
    int m = (l + r) / 2;
    for (int c : {node*2, node*2+1}) {
        int len = (c == node*2) ? m-l+1 : r-m;
        st[c].val = st[node].lazy * len;
        st[c].lazy = st[node].lazy;
    }
    st[node].lazy = LLONG_MAX;
}
\end{verbatim}

\paragraph{Complejidad.}
\begin{itemize}\setlength\itemsep{0pt}
	\item Misma que el anterior: $O(n)$ build, $O(\log n)$ query/update.
	\item Memoria $O(4n)$.
\end{itemize}

\paragraph{Donde tocar para modificar.}
\begin{itemize}\setlength\itemsep{0pt}
	\item \textbf{Mezclar assign + add}: aniadir dos campos \texttt{lazy\_set} y \texttt{lazy\_add} y manejar el orden.
	\item \textbf{Operacion no conmutativa}: para maximo o similar, aniadir.
	\item \textbf{Asignar con formula}: si el valor a asignar depende del indice (e.g., $a_i = f(i)$), aniadir la funcion.
\end{itemize}

\paragraph{Trampas y errores frecuentes.}
\begin{itemize}\setlength\itemsep{0pt}
	\item El centinela \texttt{LLONG\_MAX} es valido solo si el problema no usa ese valor en operaciones. Si lo usa, cambiar el centinela.
	\item Olvidar el reset del lazy del padre tras propagar deja dos veces el efecto.
	\item Si combinas con add, el orden debe ser: primero set, despues add.
	\item Para problemas tipo ``set $i$ to $v$'', aniadir check de que el nodo este dentro del rango del update.
\end{itemize}

\paragraph{Codigo.}
Ver \texttt{cpp.tex}, seccion \textit{Segment Tree Lazy Assign + Sum}.

\vspace{0.5em}

\subsubsection{Segment Tree Add + lower bound index}

\paragraph{Idea intuitiva.}
Variante que responde ``cual es el primer indice $i \le \text{bound}$ tal que \texttt{val[i] >= v}?''. El truco: el nodo guarda el \textbf{maximo} del rango (no la suma). El query desciende priorizando el hijo izquierdo; si su max es $< v$, salta al derecho. La estructura matematica: el predicado ``max en $[l, m]$ es $< v$'' es \textbf{monotono}, asi que la busqueda se reduce a encontrar el primer ``no''.

\paragraph{Propiedades clave (invariantes).}
\begin{itemize}\setlength\itemsep{0pt}
	\item \texttt{val} = max del rango (no suma).
	\item Funciona \textbf{solo si} la operacion update es \textbf{aditiva} (no assign).
	\item Si la consulta es ``$\le v$'' en vez de ``$< v$'', cambiar la condicion.
\end{itemize}

\paragraph{Codigo clave.}
La busqueda descendente (greedy por izquierda):
\begin{verbatim}
int lowerBound(int node, int l, int r, int ql, int qr, int v) {
    if (r < ql || l > qr || st[node].val < v) return -1;
    if (l == r) return l;
    int m = (l + r) / 2;
    int res = lowerBound(node*2, l, m, ql, qr, v);
    if (res != -1) return res;
    return lowerBound(node*2+1, m+1, r, ql, qr, v);
}
\end{verbatim}

\paragraph{Complejidad.}
\begin{itemize}\setlength\itemsep{0pt}
	\item $O(\log n)$.
	\item En el peor caso visita $O(\log n)$ nodos antes de encontrar el primero.
\end{itemize}

\paragraph{Donde tocar para modificar.}
\begin{itemize}\setlength\itemsep{0pt}
	\item \textbf{Usar suma en vez de max}: aniadir dos queries o cambiar el campo.
	\item \textbf{Lower bound con suma, no max}: si quieres ``primera posicion donde la suma excede $v$'', el truco es el mismo pero manteniendo suma acumulada.
	\item \textbf{Limite superior \texttt{bound}}: necesario para no pasar del rango valido.
	\item \textbf{Ultimo indice}: cambiar la recursion para ir a derecha primero.
\end{itemize}

\paragraph{Trampas y errores frecuentes.}
\begin{itemize}\setlength\itemsep{0pt}
	\item Si todos los valores son $< v$, devuelve $-1$; manejar ese caso en el llamador.
	\item Si la operacion es assign, el predicado deja de ser monotono; no funciona.
	\item Para ``$\le v$'' o ``$\ge v$'', ajustar la condicion cuidadosamente.
\end{itemize}

\paragraph{Codigo.}
Ver \texttt{cpp.tex}, seccion \textit{Segment Tree Add + lower bound index}.

\vspace{0.5em}

\subsubsection{Segment Tree Lazy Add + Set + Sum}

\paragraph{Idea intuitiva.}
La variante mas completa: combina \textbf{add} y \textbf{set} en el mismo nodo. Prioridad: el \texttt{set} ``mata'' cualquier \texttt{add} previo (porque asignar implica ``este es el nuevo valor base''). El \texttt{add} se acumula sobre el \texttt{set}. El orden importa: \texttt{set} primero, \texttt{add} despues. La razon algebraica: si representamos el valor final como $\text{base} + \text{delta}$, entonces \texttt{set} redefine el \texttt{base} y \texttt{add} redefine el \texttt{delta}; ambos son necesarios para una representacion completa.

\paragraph{Propiedades clave (invariantes).}
\begin{itemize}\setlength\itemsep{0pt}
	\item Dos lazies: \texttt{lazy\_set} (sobrescribe) y \texttt{lazy\_add} (acumula).
	\item En \texttt{push}: primero aplicar \texttt{set} (que borra \texttt{add} en hijos), despues \texttt{add} (que acumula).
	\item Si \texttt{lazy\_set} esta activo, ignora \texttt{lazy\_add} previo.
\end{itemize}

\paragraph{Codigo clave.}
Push con dos operaciones:
\begin{verbatim}
void push(int node, int l, int r) {
    if (l == r) return;
    int m = (l + r) / 2;
    for (int c : {node*2, node*2+1}) {
        if (st[node].hasSet) {
            st[c].val = (c == node*2 ? m-l+1 : r-m) * st[node].lazySet;
            st[c].lazySet = st[node].lazySet;
            st[c].hasSet = true;
            st[c].lazyAdd = 0;  // resetear
        }
        if (st[node].lazyAdd != 0) {
            st[c].val += (c == node*2 ? m-l+1 : r-m) * st[node].lazyAdd;
            st[c].lazyAdd += st[node].lazyAdd;
        }
    }
    st[node].hasSet = false; st[node].lazyAdd = 0;
}
\end{verbatim}

\paragraph{Complejidad.}
\begin{itemize}\setlength\itemsep{0pt}
	\item $O(\log n)$.
	\item Memoria $O(4n)$ pero con campos extra por nodo.
\end{itemize}

\paragraph{Donde tocar para modificar.}
\begin{itemize}\setlength\itemsep{0pt}
	\item \textbf{Aniadir otro tipo de operacion}: aniadir un nuevo campo y su manejo en \texttt{push}.
	\item \textbf{Build desde un arreglo}: ya implementado en el constructor.
	\item \textbf{Chained operations}: aniadir \texttt{min}, \texttt{max} con sus semanticas.
\end{itemize}

\paragraph{Trampas y errores frecuentes.}
\begin{itemize}\setlength\itemsep{0pt}
	\item Olvidar resetear \texttt{lazy\_add} cuando se aplica \texttt{set} a un hijo.
	\item Aplicar \texttt{add} antes de \texttt{set} da valores incorrectos.
	\item El orden de las operaciones en un update mixto debe estar bien definido.
\end{itemize}

\paragraph{Codigo.}
Ver \texttt{cpp.tex}, seccion \textit{Segment Tree Lazy Add + Set + Sum}.

\vspace{0.5em}

\subsubsection{Segment Tree 2D}

\paragraph{Idea intuitiva.}
Extender el segment tree a 2 dimensiones: la raiz cubre toda la matriz $[0, n) \times [0, m)$. Cada nodo es un \textbf{segment tree 1D} sobre la segunda dimension. Para update en $(x, y)$, se actualizan $O(\log n)$ nodos externos (uno por nivel vertical) y dentro de cada uno, $O(\log m)$ nodos internos. Query sobre un rectangulo: $O(\log n \cdot \log m)$. La estructura matematica: el arbol de primer nivel divide el espacio vertical; en cada nivel, los nodos internos gestionan el espacio horizontal. Asi, se reduce la consulta rectangular a $O(\log n \cdot \log m)$ consultas 1D parciales.

\paragraph{Propiedades clave (invariantes).}
\begin{itemize}\setlength\itemsep{0pt}
	\item Memoria: $O(n \cdot m)$ en el peor caso (matriz densa).
	\item Para matrices grandes con valores \textbf{cero} y unos pocos updates, se puede comprimir.
	\item Cada nodo externo tiene $4m$ nodos internos.
\end{itemize}

\paragraph{Codigo clave.}
Update en $(x, y)$ con recursion en cascada:
\begin{verbatim}
void updateX(int nodeX, int lx, int rx, int x, int y, int v) {
    if (lx != rx) {
        int mx = (lx + rx) / 2;
        if (x <= mx) updateX(nodeX*2, lx, mx, x, y, v);
        else         updateX(nodeX*2+1, mx+1, rx, x, y, v);
    }
    updateY(nodeX, lx, rx, y, v);
}
\end{verbatim}

\paragraph{Complejidad.}
\begin{itemize}\setlength\itemsep{0pt}
	\item Update y query $O(\log n \cdot \log m)$, memoria $O(4n \cdot 4m)$.
	\item Para $n = m = 1000$: $16 \cdot 10^6$ nodos ($\sim 64$MB con \texttt{int}).
\end{itemize}

\paragraph{Donde tocar para modificar.}
\begin{itemize}\setlength\itemsep{0pt}
	\item \textbf{Build desde una matriz}: aniadir constructor que tome \texttt{vector<vector<int>>}.
	\item \textbf{Lazy propagation 2D}: aniadir un campo \texttt{lazy} por nodo.
	\item \textbf{Suma + maximo simultaneos}: aniadir mas campos.
	\item \textbf{Sparse 2D}: si la mayoria son ceros, usar map en vez de vector.
\end{itemize}

\paragraph{Trampas y errores frecuentes.}
\begin{itemize}\setlength\itemsep{0pt}
	\item La recursion puede ser profunda; aniadir iterativo si $n, m > 10^5$.
	\item Memoria $O(4n \cdot 4m) = 16nm$; si $n = m = 10^3$, son $16M$ nodos; con $n = m = 10^4$, son $1.6 \cdot 10^9$, no factible.
	\item Olvidar propagar en cascada da updates incorrectos.
	\item Si la mayoria de queries son rectangulares grandes, considerar KD-tree u otra estructura.
\end{itemize}

\paragraph{Codigo.}
Ver \texttt{cpp.tex}, seccion \textit{Segment Tree 2D}.

\vspace{0.5em}

\subsubsection{Segment Tree persistente}

\paragraph{Idea intuitiva.}
Cada \texttt{update} crea una \textbf{nueva raiz} sin modificar los nodos no afectados. Asi puedes tener \textbf{multiples versiones} del arreglo y consultar cada una en $O(\log n)$. Memoria: $O(N \log N)$ nodos en total para $N$ updates. La tecnica se llama ``path copying'': solo el camino de la raiz a la hoja modificada se copia; el resto del arbol se comparte entre versiones. Asi, $N$ updates producen $\le N \log n$ nodos en total.

\paragraph{Propiedades clave (invariantes).}
\begin{itemize}\setlength\itemsep{0pt}
	\item Cada nodo es \textbf{inmutable} despues de creado; las versiones comparten estructura.
	\item \texttt{pUpdate} retorna la nueva raiz; la anterior sigue accesible.
	\item \texttt{pQuery} toma una raiz especifica (no la global).
	\item Las versiones viven independientemente; borrar una no afecta a otras.
\end{itemize}

\paragraph{Codigo clave.}
El corazon del persistent update:
\begin{verbatim}
PNode* pUpdate(PNode* v, int l, int r, int pos, int val) {
    PNode* u = newNode(v);  // copia con misma info
    if (l == r) { u->sum = val; return u; }
    int m = (l + r) / 2;
    if (pos <= m) u->left = pUpdate(v->left, l, m, pos, val);
    else          u->right = pUpdate(v->right, m+1, r, pos, val);
    u->sum = u->left->sum + u->right->sum;
    return u;
}
// versions.push_back(pUpdate(versions.back(), 0, n-1, pos, val));
\end{verbatim}

\paragraph{Complejidad.}
\begin{itemize}\setlength\itemsep{0pt}
	\item Update $O(\log n)$ nodos nuevos, query $O(\log n)$.
	\item Memoria $O(N \log n)$ para $N$ updates.
\end{itemize}

\paragraph{Donde tocar para modificar.}
\begin{itemize}\setlength\itemsep{0pt}
	\item \textbf{Persistencia completa (cada update crea un nodo)}: aniadir \texttt{version[v]} para acceder a cada version.
	\item \textbf{Cambiar suma por otro operacion}: ajustar \texttt{pUpdate} y \texttt{pQuery}.
	\item \textbf{Memoria grande}: en contests, el allocator por defecto puede ser lento; usar \texttt{vector<PNode>} con \texttt{reserve(2*N*logN)} y un \texttt{int nextNode = 0} para asignar en $O(1)$.
	\item \textbf{Lazy persistence}: aniadir campos \texttt{lazy} en cada nodo; los hijos ``ven'' el lazy a traves de la copia.
\end{itemize}

\paragraph{Trampas y errores frecuentes.}
\begin{itemize}\setlength\itemsep{0pt}
	\item Olvidar pasar la \textbf{vieja raiz} al query, usando la nueva por error.
	\item Si guardas multiples versiones, asegurate de pasar la raiz correcta a cada query.
	\item Memory leak en produccion (no en contests) si no liberas los nodos.
	\item Para queries ``diferencia entre version $a$ y version $b

\vspace{0.5em}

\subsubsection{Segment Tree iterativo}

\paragraph{Idea intuitiva.}
Representacion \textbf{sin recursion}: los nodos se guardan en un \texttt{vector} de tamano $2n$. Los hijos del nodo $i$ son $2i$ y $2i+1$. Las hojas ocupan las posiciones $n, n+1, \dots, 2n-1$. Update: subir el cambio desde la hoja hasta la raiz. Query: subir dos indices $l$ y $r$ hasta que se ``crucen''. La ventaja principal: localidad de memoria (los nodos estan en posiciones contiguas del vector), que mejora el cache miss rate y la constante practica. La desventaja: implementar lazy propagation es mas complejo.

\paragraph{Propiedades clave (invariantes).}
\begin{itemize}\setlength\itemsep{0pt}
	\item $O(\log n)$ por operacion, \textbf{mas rapido} en practica que el recursivo (mejor uso de cache, no hay overhead de llamadas).
	\item Build simple: copiar las hojas y combinar de abajo hacia arriba.
	\item Las hojas tienen indices en $[n, 2n)$; los nodos internos en $[1, n)$.
\end{itemize}

\paragraph{Codigo clave.}
El bucle de query ascendente:
\begin{verbatim}
long long query(int l, int r) {  // [l, r] inclusivo
    long long res = 0;
    for (l += n, r += n + 1; l < r; l >>= 1, r >>= 1) {
        if (l & 1) res += st[l++];
        if (r & 1) res += st[--r];
    }
    return res;
}
\end{verbatim}

\paragraph{Complejidad.}
\begin{itemize}\setlength\itemsep{0pt}
	\item $O(\log n)$ por update/query, $O(n)$ build.
	\item En la practica, $\sim 1.5x$ mas rapido que el recursivo para $n \ge 10^5$.
\end{itemize}

\paragraph{Donde tocar para modificar.}
\begin{itemize}\setlength\itemsep{0pt}
	\item \textbf{Suma -> max/min/xor}: cambiar la operacion en \texttt{set} y \texttt{query}.
	\item \textbf{Aniadir update en rango con lazy}: implementar version lazy por separado (no trivial en iterativo).
	\item \textbf{Tamano no potencia de 2}: si $n$ no es potencia de 2, aniadir padding.
	\item \textbf{Persistencia}: usar \texttt{vector<array<int, 2>>} y duplicar nodos.
\end{itemize}

\paragraph{Trampas y errores frecuentes.}
\begin{itemize}\setlength\itemsep{0pt}
	\item En el query, \texttt{r += n+1} (no \texttt{r += n}) porque \texttt{query(l, r)} es \textbf{inclusivo} en ambos extremos.
	\item El indice \texttt{l & 1} detecta si \texttt{l} es hijo derecho; aniadir \texttt{res += st[l++]} en ese caso.
	\item El \texttt{--r} en la condicion derecha es para incluir el limite.
	\item Si las operaciones no son idempotentes, cuidado con el orden.
\end{itemize}

\paragraph{Codigo.}
Ver \texttt{cpp.tex}, seccion \textit{Segment Tree iterativo}.

\subsection{Otras}

\subsubsection{MO}

\paragraph{Idea intuitiva.}
Algoritmo de Mo ``puro'' (sin updates): queries offline ordenadas por $(L/\sqrt{n}, R)$ (con $R$ alternando entre ascendente y descendente segun el bloque). Mover $L$ o $R$ entre queries consecutivas es $O(1)$; el total es $O((N + Q) \sqrt{n})$. La intuicion: al reordenar las queries, los punteros $L$ y $R$ se mueven ``lo menos posible''; cada elemento del array se visita $O(\sqrt{n} \cdot Q/N)$ veces.

\paragraph{Propiedades clave (invariantes).}
\begin{itemize}\setlength\itemsep{0pt}
	\item \texttt{Q} guarda \texttt{(l, r, idx)}.
	\item El sort prioriza bloque de $L$; dentro del bloque, alterna $R$ ascendente/descendente para mejorar la locality.
	\item Tres punteros \texttt{l}, \texttt{r} que se mueven paso a paso.
	\item Tamano de bloque optimo: $\sim N / \sqrt{Q}$.
\end{itemize}

\paragraph{Codigo clave.}
El sort clasico de Mo con alternancia:
\begin{verbatim}
int block = (int)sqrt(n);
sort(queries.begin(), queries.end(), [&](const Q& a, const Q& b) {
    int blockA = a.l / block, blockB = b.l / block;
    if (blockA != blockB) return blockA < blockB;
    return (blockA & 1) ? (a.r > b.r) : (a.r < b.r);
});
\end{verbatim}

\paragraph{Complejidad.}
\begin{itemize}\setlength\itemsep{0pt}
	\item $O((N + Q) \sqrt{n})$.
	\item Con Hilbert order, $\sim 30\%$ mas rapido en practica.
\end{itemize}

\paragraph{Donde tocar para modificar.}
\begin{itemize}\setlength\itemsep{0pt}
	\item \textbf{Hilbert order}: alternativa al sort clasico; da mejor locality y es mas rapido en practica.
	\item \textbf{Aniadir la operacion concreta}: el snippet tiene \texttt{// add} y \texttt{// delete} como comentarios; reemplazarlos con la logica del problema (contar distintos, suma de frecuencias, etc.).
	\item \textbf{Bloque optimo}: calcular \texttt{int block = max(1, (int)(n / sqrt(q)));}.
\end{itemize}

\paragraph{Trampas y errores frecuentes.}
\begin{itemize}\setlength\itemsep{0pt}
	\item Los indices en el sort son \textbf{0-based} (\texttt{q[i].l = l - 1}), consistente con el resto del codigo.
	\item \texttt{add} y \texttt{remove} deben ser \textbf{inversos exactos}; sino, las queries dan respuestas incorrectas.
	\item Para queries cerradas $[l, r]$, Mo suele usar $[l, r+1)$ (half-open); aniadir o restar 1 consistentemente.
\end{itemize}

\paragraph{Codigo.}
Ver \texttt{cpp.tex}, seccion \textit{MO}.

\vspace{0.5em}

\subsubsection{DSU con rollback}

\paragraph{Idea intuitiva.}
Variante del DSU donde cada \texttt{unite} \textbf{se puede deshacer} en orden LIFO. La idea: en vez de path compression (que es irreversible), solo se usa union by size. Cada \texttt{unite} guarda un par \texttt{(parent\_changed, old\_size)} en una pila. \texttt{rollback()} desapila hasta el ultimo ``snapshot'' (marcador) restaurando los valores. La razon: path compression modifica la estructura de multiples nodos, no solo uno, asi que es dificil de deshacer; union by size solo modifica dos nodos (el padre del menor y su tamano), asi que es trivial deshacer.

\paragraph{Propiedades clave (invariantes).}
\begin{itemize}\setlength\itemsep{0pt}
	\item \textbf{Sin path compression}: cada \texttt{find} es $O(\log n)$ (peor caso del union by size) en el peor caso, pero amortizado $O(\log n)$.
	\item \textbf{Snapshots}: \texttt{snapshot()} marca un punto; \texttt{rollback()} deshace hasta ahi.
	\item Cada \texttt{unite} exitoso apila \textbf{dos} pares (uno por cada cambio).
	\item La pila de historial crece monotona; \texttt{rollback} la reduce.
\end{itemize}

\paragraph{Codigo clave.}
El unite con guardado de historial:
\begin{verbatim}
void unite(int a, int b) {
    a = find(a); b = find(b);
    if (a == b) { history.push({-1, -1}); return; }
    if (size[a] < size[b]) swap(a, b);
    history.push({b, size[a]});
    parent[b] = a;
    size[a] += size[b];
    components--;
}
void rollback() {
    auto [b, sz] = history.top(); history.pop();
    if (b == -1) return;
    size[parent[b]] = sz;
    parent[b] = b;
    components++;
}
\end{verbatim}

\paragraph{Complejidad.}
\begin{itemize}\setlength\itemsep{0pt}
	\item $O(\log n)$ por \texttt{unite}/\texttt{find}, $O(1)$ por rollback de una operacion.
	\item Sin path compression: peor caso $O(\log n)$ por find.
\end{itemize}

\paragraph{Donde tocar para modificar.}
\begin{itemize}\setlength\itemsep{0pt}
	\item \textbf{Rollback parcial} (no todo): implementar snapshots parciales guardando el tamano del historial.
	\item \textbf{Mas informacion por conjunto}: aniadir mas campos al struct y guardarlos en \texttt{history}.
	\item \textbf{Persistente}: aniadir versionado completo (clonar el struct en cada cambio).
	\item \textbf{DSU de Arpa}: variante con $O(\log n)$ amortizado y rollback en $O(1)$.
\end{itemize}

\paragraph{Trampas y errores frecuentes.}
\begin{itemize}\setlength\itemsep{0pt}
	\item El snippet guarda dos pares por \texttt{unite}; al rollback \textbf{desapilar en orden correcto} (no invertir).
	\item Olvidar el \texttt{components--} en \texttt{unite} o el \texttt{components++} en rollback deja el contador mal.
	\item El sentinela \texttt{\{-1, -1\}} indica un unite trivial; importante para mantener consistencia.
	\item \texttt{find} puede ser $O(\log n)$ peor caso sin path compression; verificar que es OK para el problema.
\end{itemize}

\paragraph{Codigo.}
Ver \texttt{cpp.tex}, seccion \textit{DSU con rollback}.

\vspace{0.5em}

\subsubsection{Segment Tree 2D}

\paragraph{Idea intuitiva.}
Extender el segment tree a 2 dimensiones: la raiz cubre toda la matriz $[0, n) \times [0, m)$. Cada nodo es un \textbf{segment tree 1D} sobre la segunda dimension. Para update en $(x, y)$, se actualizan $O(\log n)$ nodos externos (uno por nivel vertical) y dentro de cada uno, $O(\log m)$ nodos internos. Query sobre un rectangulo: $O(\log n \cdot \log m)$. La estructura matematica: el arbol de primer nivel divide el espacio vertical; en cada nivel, los nodos internos gestionan el espacio horizontal. Asi, se reduce la consulta rectangular a $O(\log n \cdot \log m)$ consultas 1D parciales.

\paragraph{Propiedades clave (invariantes).}
\begin{itemize}\setlength\itemsep{0pt}
	\item Memoria: $O(n \cdot m)$ en el peor caso (matriz densa).
	\item Para matrices grandes con valores \textbf{cero} y unos pocos updates, se puede comprimir.
	\item Cada nodo externo tiene $4m$ nodos internos.
\end{itemize}

\paragraph{Codigo clave.}
Update en $(x, y)$ con recursion en cascada:
\begin{verbatim}
void updateX(int nodeX, int lx, int rx, int x, int y, int v) {
    if (lx != rx) {
        int mx = (lx + rx) / 2;
        if (x <= mx) updateX(nodeX*2, lx, mx, x, y, v);
        else         updateX(nodeX*2+1, mx+1, rx, x, y, v);
    }
    updateY(nodeX, lx, rx, y, v);
}
\end{verbatim}

\paragraph{Complejidad.}
\begin{itemize}\setlength\itemsep{0pt}
	\item Update y query $O(\log n \cdot \log m)$, memoria $O(4n \cdot 4m)$.
	\item Para $n = m = 1000$: $16 \cdot 10^6$ nodos ($\sim 64$MB con \texttt{int}).
\end{itemize}

\paragraph{Donde tocar para modificar.}
\begin{itemize}\setlength\itemsep{0pt}
	\item \textbf{Build desde una matriz}: aniadir constructor que tome \texttt{vector<vector<int>>}.
	\item \textbf{Lazy propagation 2D}: aniadir un campo \texttt{lazy} por nodo.
	\item \textbf{Suma + maximo simultaneos}: aniadir mas campos.
	\item \textbf{Sparse 2D}: si la mayoria son ceros, usar map en vez de vector.
\end{itemize}

\paragraph{Trampas y errores frecuentes.}
\begin{itemize}\setlength\itemsep{0pt}
	\item La recursion puede ser profunda; aniadir iterativo si $n, m > 10^5$.
	\item Memoria $O(4n \cdot 4m) = 16nm$; si $n = m = 10^3$, son $16M$ nodos; con $n = m = 10^4$, son $1.6 \cdot 10^9$, no factible.
	\item Olvidar propagar en cascada da updates incorrectos.
	\item Si la mayoria de queries son rectangulares grandes, considerar KD-tree u otra estructura.
\end{itemize}

\paragraph{Codigo.}
Ver \texttt{cpp.tex}, seccion \textit{Segment Tree 2D}.

\vspace{0.5em}

\subsubsection{Segment Tree persistente}

\paragraph{Idea intuitiva.}
Cada \texttt{update} crea una \textbf{nueva raiz} sin modificar los nodos no afectados. Asi puedes tener \textbf{multiples versiones} del arreglo y consultar cada una en $O(\log n)$. Memoria: $O(N \log N)$ nodos en total para $N$ updates.

\paragraph{Propiedades clave (invariantes).}
\begin{itemize}\setlength\itemsep{0pt}
	\item Cada nodo es \textbf{inmutable} despues de creado; las versiones comparten estructura.
	\item \texttt{pUpdate} retorna la nueva raiz; la anterior sigue accesible.
	\item \texttt{pQuery} toma una raiz especifica (no la global).
\end{itemize}

\paragraph{Complejidad.}
\begin{itemize}\setlength\itemsep{0pt}
	\item Update $O(\log n)$ nodos nuevos, query $O(\log n)$.
\end{itemize}

\paragraph{Donde tocar para modificar.}
\begin{itemize}\setlength\itemsep{0pt}
	\item \textbf{Persistencia completa (cada update crea un nodo)}: aniadir \texttt{version[v]} para acceder a cada version.
	\item \textbf{Cambiar suma por otro operacion}: ajustar \texttt{pUpdate} y \texttt{pQuery}.
	\item \textbf{Memoria grande}: en contests, el allocator por defecto puede ser lento; usar \texttt{vector<PNode>} con \texttt{reserve(2*N*logN)} y un \texttt{int nextNode = 0} para asignar en $O(1)$.
\end{itemize}

\paragraph{Trampas y errores frecuentes.}
\begin{itemize}\setlength\itemsep{0pt}
	\item Olvidar pasar la \textbf{vieja raiz} al query, usando la nueva por error.
\end{itemize}

\paragraph{Codigo.}
Ver \texttt{cpp.tex}, seccion \textit{Segment Tree persistente}.

\vspace{0.5em}

\subsubsection{Segment Tree iterativo}

\paragraph{Idea intuitiva.}
Representacion \textbf{sin recursion}: los nodos se guardan en un \texttt{vector} de tamano $2n$. Los hijos del nodo $i$ son $2i$ y $2i+1$. Las hojas ocupan las posiciones $n, n+1, \dots, 2n-1$. Update: subir el cambio desde la hoja hasta la raiz. Query: subir dos indices $l$ y $r$ hasta que se ``crucen''. La ventaja principal: localidad de memoria (los nodos estan en posiciones contiguas del vector), que mejora el cache miss rate y la constante practica. La desventaja: implementar lazy propagation es mas complejo.

\paragraph{Propiedades clave (invariantes).}
\begin{itemize}\setlength\itemsep{0pt}
	\item $O(\log n)$ por operacion, \textbf{mas rapido} en practica que el recursivo (mejor uso de cache, no hay overhead de llamadas).
	\item Build simple: copiar las hojas y combinar de abajo hacia arriba.
	\item Las hojas tienen indices en $[n, 2n)$; los nodos internos en $[1, n)$.
\end{itemize}

\paragraph{Codigo clave.}
El bucle de query ascendente:
\begin{verbatim}
long long query(int l, int r) {  // [l, r] inclusivo
    long long res = 0;
    for (l += n, r += n + 1; l < r; l >>= 1, r >>= 1) {
        if (l & 1) res += st[l++];
        if (r & 1) res += st[--r];
    }
    return res;
}
\end{verbatim}

\paragraph{Complejidad.}
\begin{itemize}\setlength\itemsep{0pt}
	\item $O(\log n)$ por update/query, $O(n)$ build.
	\item En la practica, $\sim 1.5x$ mas rapido que el recursivo para $n \ge 10^5$.
\end{itemize}

\paragraph{Donde tocar para modificar.}
\begin{itemize}\setlength\itemsep{0pt}
	\item \textbf{Suma -> max/min/xor}: cambiar la operacion en \texttt{set} y \texttt{query}.
	\item \textbf{Aniadir update en rango con lazy}: implementar version lazy por separado (no trivial en iterativo).
	\item \textbf{Tamano no potencia de 2}: si $n$ no es potencia de 2, aniadir padding.
	\item \textbf{Persistencia}: usar \texttt{vector<array<int, 2>>} y duplicar nodos.
\end{itemize}

\paragraph{Trampas y errores frecuentes.}
\begin{itemize}\setlength\itemsep{0pt}
	\item En el query, \texttt{r += n+1} (no \texttt{r += n}) porque \texttt{query(l, r)} es \textbf{inclusivo} en ambos extremos.
	\item El indice \texttt{l & 1} detecta si \texttt{l} es hijo derecho; aniadir \texttt{res += st[l++]} en ese caso.
	\item El \texttt{--r} en la condicion derecha es para incluir el limite.
	\item Si las operaciones no son idempotentes, cuidado con el orden.
\end{itemize}

\paragraph{Codigo.}
Ver \texttt{cpp.tex}, seccion \textit{Segment Tree iterativo}.

\vspace{0.5em}

\subsection{Fenwick / BIT}

\subsubsection{BIT point update + prefix sum}

\paragraph{Idea intuitiva.}
El \textbf{Fenwick Tree} (BIT) almacena sums parciales de manera que \texttt{sum(i)} (suma de $[1, i]$) y \texttt{add(i, delta)} cuestan $O(\log n)$. La representacion es ingeniosa: el nodo $i$ cubre el rango $[i - \text{lowbit}(i) + 1, i]$. Asi, \texttt{sum(i)} itera \texttt{i -= lowbit(i)}, sumando los nodos que cubren un prefijo creciente, y \texttt{add(i, delta)} itera \texttt{i += lowbit(i)}, propagando el cambio a todos los nodos cuyo rango contiene a $i$. La idea: cada indice $i$ se descompone en una suma de bloques disjuntos de tamano potencias de 2 (su ``lowbit representation'').

\paragraph{Propiedades clave (invariantes).}
\begin{itemize}\setlength\itemsep{0pt}
	\item Indice \textbf{1-based} (no usar $i=0$).
	\item \texttt{lowbit(i) = i \& (-i)}.
	\item Cada nodo cubre exactamente un rango de tamano \texttt{lowbit(i)}.
	\item La operacion debe ser \textbf{invertible} (suma, xor, no min/max).
\end{itemize}

\paragraph{Codigo clave.}
Las dos operaciones basicas:
\begin{verbatim}
void add(int i, long long d) {
    for (; i <= n; i += i & -i) bit[i] += d;
}
long long sum(int i) {
    long long s = 0;
    for (; i > 0; i -= i & -i) s += bit[i];
    return s;
}
\end{verbatim}

\paragraph{Complejidad.}
\begin{itemize}\setlength\itemsep{0pt}
	\item $O(\log n)$ por operacion. Memoria $O(n)$.
	\item Constante muy pequena ($\sim 1$): en la practica, $\sim 3x$ mas rapido que segment tree.
\end{itemize}

\paragraph{Donde tocar para modificar.}
\begin{itemize}\setlength\itemsep{0pt}
	\item \textbf{Suma -> max/min/xor}: cambiar el \texttt{+=} por la operacion (con identidad apropriada). \textbf{No} funciona para \texttt{min} de forma natural (no es invertible).
	\item \textbf{Tamano dinamico}: aniadir un metodo \texttt{resize(newSize)} que preserve valores.
	\item \textbf{Tamano grande}: con $n = 10^7$, mejor comprimir coordenadas.
	\item \textbf{Operador no conmutativo}: usar segment tree en su lugar.
\end{itemize}

\paragraph{Trampas y errores frecuentes.}
\begin{itemize}\setlength\itemsep{0pt}
	\item Indices 0 vs 1: el BIT es \textbf{1-based}. Si el problema da indices 0-based, aniadir \texttt{+1} en cada llamada.
	\item El metodo \texttt{sum} itera \texttt{i -= i \& -i}; aniadir o quitar 1 rompe el invariante.
	\item Si el tipo es \texttt{int}, posibles overflows; usar \texttt{long long}.
	\item Para rango $[l, r]$, el query es \texttt{sum(r) - sum(l-1)} (cuidado con $l = 1$).
\end{itemize}

\paragraph{Codigo.}
Ver \texttt{cpp.tex}, seccion \textit{BIT point update + prefix sum}.

\vspace{0.5em}

\subsubsection{BIT range add + point query}

\paragraph{Idea intuitiva.}
Truco de \textbf{diferencias}: para sumar $d$ a $[l, r]$, aniadir $d$ en $l$ y $-d$ en $r+1$. Asi, el valor en $i$ es la suma de los prefijos hasta $i$, lo que se calcula con un BIT normal. La idea algebraica: si mantenemos un arreglo de diferencias $\Delta$ donde $a_i = \sum_{j \le i} \Delta_j$, entonces aniadir $d$ a $a[l..r]$ equivale a $\Delta_l \mathrel{+}= d$ y $\Delta_{r+1} \mathrel{-}= d$, y reconstruir $a$ es un prefix sum.

\paragraph{Propiedades clave (invariantes).}
\begin{itemize}\setlength\itemsep{0pt}
	\item Internamente es un BIT 5.1 con la convencion de diferencias.
	\item \texttt{rangeAdd(l, r, delta)} cuesta $O(\log n)$ (dos \texttt{add}).
	\item \texttt{pointQuery(i)} cuesta $O(\log n)$.
	\item Si $r+1 > n$, no se aniade el $-d$ (fuera del rango).
\end{itemize}

\paragraph{Codigo clave.}
El truco de diferencias:
\begin{verbatim}
void rangeAdd(int l, int r, long long d) {
    add(l, d);
    if (r + 1 <= n) add(r + 1, -d);
}
long long pointQuery(int i) { return sum(i); }
\end{verbatim}

\paragraph{Complejidad.}
\begin{itemize}\setlength\itemsep{0pt}
	\item $O(\log n)$ por operacion.
	\item Memoria $O(n)$.
\end{itemize}

\paragraph{Donde tocar para modificar.}
\begin{itemize}\setlength\itemsep{0pt}
	\item \textbf{Persistencia}: aniadir versionado (cada update crea una nueva ``copia'' del estado; en BIT no es trivial, mejor usar persistent segtree).
	\item \textbf{Coordenada compression}: si $n$ es grande pero los updates son pocos, comprimir.
	\item \textbf{2D range add + point query}: aniadir una dimension con la misma tecnica.
\end{itemize}

\paragraph{Trampas y errores frecuentes.}
\begin{itemize}\setlength\itemsep{0pt}
	\item Si $r+1 > n$, \textbf{no} aniadir \texttt{-delta} (no hay efecto fuera del rango); el snippet ya lo chequea.
	\item Para reconstruir el array original tras $k$ range adds, hacer $k$ point queries (no hay shortcut).
	\item Confundir ``suma de diferencias'' con ``suma acumulada''; son cosas distintas.
\end{itemize}

\paragraph{Codigo.}
Ver \texttt{cpp.tex}, seccion \textit{BIT range add + point query}.

\vspace{0.5em}

\subsubsection{BIT range add + range sum}

\paragraph{Idea intuitiva.}
Para queries $\sum_{i=1}^{x} a_i$ despues de varios range adds, el truco de diferencias no basta (necesitas la suma acumulada de las diferencias). Solucion: dos BITs, $t_1$ y $t_2$. Tras aniadir $d$ a $[l, r]$: $t_1$ recibe $+d$ en $l$, $-d$ en $r+1$; $t_2$ recibe $+d(l-1)$ en $l$, $-d \cdot r$ en $r+1$. Entonces $\text{prefixSum}(x) = t_1.\text{sum}(x) \cdot x - t_2.\text{sum}(x)$. La razon matematica: el valor en $i$ es $\sum_{j \le i} \Delta_j$, asi que la suma hasta $i$ es $\sum_{k \le i} (i - k + 1) \Delta_k$; expandir da la formula con dos BITs.

\paragraph{Propiedades clave (invariantes).}
\begin{itemize}\setlength\itemsep{0pt}
	\item Verificacion: para $x \ge r$, $\text{prefixSum}(x)$ acumula correctamente todas las contribuciones.
	\item La formula sale de expandir la suma de diferencias.
	\item La identidad clave: $\sum_{i=1}^{x} i = x(x+1)/2$, asi que la suma de diferencias ponderadas requiere BITs adicionales.
\end{itemize}

\paragraph{Codigo clave.}
Range add con range sum (dos BITs):
\begin{verbatim}
void rangeAdd(int l, int r, long long d) {
    t1.add(l, d); t1.add(r + 1, -d);
    t2.add(l, d * (l - 1)); t2.add(r + 1, -d * r);
}
long long prefixSum(int x) {
    return t1.sum(x) * x - t2.sum(x);
}
\end{verbatim}

\paragraph{Complejidad.}
\begin{itemize}\setlength\itemsep{0pt}
	\item $O(\log n)$ por operacion (2 \texttt{add} o 2 \texttt{sum} por cada lado).
	\item Memoria $O(n)$.
\end{itemize}

\paragraph{Donde tocar para modificar.}
\begin{itemize}\setlength\itemsep{0pt}
	\item \textbf{2D range add + range sum}: aniadir otra dimension con la misma tecnica (queda 4 BITs).
	\item \textbf{Cambiar mod}: aniadir modulo en cada operacion.
	\item \textbf{Output del array}: aniadir un metodo que reconstruya $a_i$ point query por cada $i$.
\end{itemize}

\paragraph{Trampas y errores frecuentes.}
\begin{itemize}\setlength\itemsep{0pt}
	\item Overflow en \texttt{delta * (l - 1)} si los valores son grandes; usar \texttt{long long} y \texttt{\% MOD} al final.
	\item Confundir la convencion $0$-based vs $1$-based del BIT; las dos afectan las formulas.
	\item En 2D, el codigo es $\sim 4x$ mas largo y propenso a errores; verificar con casos pequenos.
\end{itemize}

\paragraph{Codigo.}
Ver \texttt{cpp.tex}, seccion \textit{BIT range add + range sum}.

\subsection{RMQ y sparse table}

\subsubsection{Sparse Table}

\paragraph{Idea intuitiva.}
Precomputa respuestas a queries de tamano $2^k$ para $k = 0, 1, \dots, \log n$. Asi, \texttt{st[k][i]} guarda la respuesta sobre $[i, i + 2^k)$. Para query $[l, r]$, se usan dos bloques de tamano $2^{\lfloor \log_2(r-l+1) \rfloor}$: $[l, l + 2^k)$ y $[r - 2^k + 1, r)$, y se combinan. Esto da $O(1)$ por query \textbf{solo si la operacion es idempotente} (min, max, gcd; \textbf{no} suma). La razon: con idempotencia, solapar los dos bloques no afecta el resultado.

\paragraph{Propiedades clave (invariantes).}
\begin{itemize}\setlength\itemsep{0pt}
	\item \texttt{st[0][i] = a[i]}.
	\item \texttt{st[k][i] = op(st[k-1][i], st[k-1][i + 2\^{}(k-1)])}.
	\item \texttt{LOG = floor(log2(n)) + 1}.
	\item La operacion debe ser \textbf{idempotente} ($f(x, x) = x$) y asociativa.
\end{itemize}

\paragraph{Codigo clave.}
Query $O(1)$ en sparse table:
\begin{verbatim}
long long query(int l, int r) {
    int k = __builtin_clz(1) - __builtin_clz(r - l + 1);  // log2
    return min(st[k][l], st[k][r - (1 << k) + 1]);
}
\end{verbatim}

\paragraph{Complejidad.}
\begin{itemize}\setlength\itemsep{0pt}
	\item Build $O(n \log n)$, query $O(1)$.
	\item Memoria $O(n \log n)$.
\end{itemize}

\paragraph{Donde tocar para modificar.}
\begin{itemize}\setlength\itemsep{0pt}
	\item \textbf{Cambiar min por otra operacion idempotente}: cambiar las dos ocurrencias de \texttt{min(...)} por la operacion (max, gcd). Verificar identidad y asociatividad.
	\item \textbf{Suma en rango}: NO sirve con este metodo; usar prefix sums.
	\item \textbf{Sparse table 2D}: aniadir otra dimension (queda $O(n^2 \log^2 n)$ memoria).
	\item \textbf{LCA queries}: combinar con binary lifting en arboles.
\end{itemize}

\paragraph{Trampas y errores frecuentes.}
\begin{itemize}\setlength\itemsep{0pt}
	\item El query toma $[l, r]$ \textbf{inclusivo} (no $[l, r)$).
	\item El \texttt{LOG} debe estar bien calculado; con $n = 1$, \texttt{LOG = 1}.
	\item Usar sparse table para suma es un error; necesitas prefix sums o segment tree.
	\item Para queries $O(1)$ con operaciones no idempotentes, usar Disjoint Sparse Table.
\end{itemize}

\paragraph{Codigo.}
Ver \texttt{cpp.tex}, seccion \textit{Sparse Table}.

\vspace{0.5em}

\subsubsection{Disjoint Sparse Table}

\paragraph{Idea intuitiva.}
Variante que tambien da $O(1)$ por query pero \textbf{no requiere idempotencia} para algunas operaciones (sirve para suma, xor, etc. en linea con la mayoria de operaciones asociativas). La idea: para cada nivel $k$, los rangos ``se parten en mitades''. \texttt{st[k][i]} cubre un rango que \textbf{empieza} en $i$ y se extiende a izquierda o derecha hasta encontrar un ``punto de division''. El truco: para query $[l, r]$, el nivel optimo es el bit mas alto donde $l$ y $r$ difieren; los rangos $[l, ?]$ y $[?, r]$ cubren exactamente el query sin solape.

\paragraph{Propiedades clave (invariantes).}
\begin{itemize}\setlength\itemsep{0pt}
	\item Construccion diferente al Sparse Table clasico: en cada nivel, se extiende desde cada posicion hacia ambos lados hasta encontrar la mitad de nivel.
	\item \texttt{query(l, r)} usa \texttt{k = log2(l \^{} r)} y combina \texttt{st[k][l]} y \texttt{st[k][r]}.
	\item La operacion debe ser \textbf{asociativa} pero no necesariamente idempotente.
\end{itemize}

\paragraph{Codigo clave.}
Query con DST:
\begin{verbatim}
long long query(int l, int r) {
    if (l == r) return st[0][l];
    int k = 31 - __builtin_clz(l ^ r);  // bit mas alto que difiere
    return op(st[k][l], st[k][r]);
}
\end{verbatim}

\paragraph{Complejidad.}
\begin{itemize}\setlength\itemsep{0pt}
	\item Build $O(n \log n)$, query $O(1)$.
	\item Memoria $O(n \log n)$.
\end{itemize}

\paragraph{Donde tocar para modificar.}
\begin{itemize}\setlength\itemsep{0pt}
	\item \textbf{Operacion no asociativa}: no funciona; el Disjoint Sparse Table asume asociatividad.
	\item \textbf{Memory compression}: con $n \le 10^6$ alcanza; con mas, comprimir.
	\item \textbf{Operaciones no conmutativas}: aniadir el cuidado de que el orden es importante.
\end{itemize}

\paragraph{Trampas y errores frecuentes.}
\begin{itemize}\setlength\itemsep{0pt}
	\item La construccion es sutil; verificar con casos pequenos antes de usar.
	\item El bit ``mas alto que difiere'' requiere cuidado con $l = r$; aniadir caso base.
	\item Si la operacion no es asociativa, los dos bloques no se pueden combinar en cualquier orden.
\end{itemize}

\paragraph{Codigo.}
Ver \texttt{cpp.tex}, seccion \textit{Disjoint Sparse Table}.

\subsection{Trees y DP en arbol}

\subsubsection{Binary Lifting / LCA}

\paragraph{Idea intuitiva.}
Precomputa para cada nodo $v$ sus ancestros a distancias $2^k$ (los ``saltos''). Asi, para subir $d$ niveles, se descompone $d$ en binario y se aplican los saltos correspondientes. El \textbf{LCA} (Lowest Common Ancestor) se obtiene subiendo el nodo mas profundo hasta igualar profundidades y luego subiendo ambos hasta que sus padres coincidan. La intuicion: cualquier entero $\le n$ se representa en binario con $\le \log n$ bits, asi que subir $d$ niveles requiere $\le \log n$ saltos.

\paragraph{Propiedades clave (invariantes).}
\begin{itemize}\setlength\itemsep{0pt}
	\item \texttt{up[k][v]} = ancestro de $v$ a $2^k$ niveles.
	\item \texttt{LOG = floor(log2(n)) + 1}.
	\item La recursion de \texttt{dfs} llena \texttt{up[0][v]} (el padre directo) y por transitividad los demas.
	\item El LCA esta definido como el ancestro comun mas profundo de $u$ y $v$.
\end{itemize}

\paragraph{Codigo clave.}
LCA con binary lifting:
\begin{verbatim}
int lca(int u, int v) {
    if (depth[u] < depth[v]) swap(u, v);
    int diff = depth[u] - depth[v];
    for (int k = LOG - 1; k >= 0; --k)
        if (diff & (1 << k)) u = up[k][u];
    if (u == v) return u;
    for (int k = LOG - 1; k >= 0; --k)
        if (up[k][u] != up[k][v])
            u = up[k][u], v = up[k][v];
    return up[0][u];
}
\end{verbatim}

\paragraph{Complejidad.}
\begin{itemize}\setlength\itemsep{0pt}
	\item Build $O(n \log n)$, query LCA $O(\log n)$.
	\item Memoria $O(n \log n)$.
\end{itemize}

\paragraph{Donde tocar para modificar.}
\begin{itemize}\setlength\itemsep{0pt}
	\item \textbf{Distancia en el arbol}: combinar con \texttt{depth[]} para calcular \texttt{depth[u] + depth[v] - 2*depth[lca]}.
	\item \textbf{K-esimo ancestro}: subir $k$ niveles con el mismo truco.
	\item \textbf{Aniadir peso a las aristas}: guardar \texttt{dist[2\^{}k][v]} (distancia acumulada a $2^k$ saltos) ademas de \texttt{up}.
	\item \textbf{Sparse table sobre depth}: alternativa $O(1)$ a LCA (mas complejo).
\end{itemize}

\paragraph{Trampas y errores frecuentes.}
\begin{itemize}\setlength\itemsep{0pt}
	\item El arbol debe ser \textbf{rooted} (el snippet asume eso); si no, hacer un \texttt{dfs} desde cualquier nodo.
	\item Para grafos no conexos, correr \texttt{dfs} desde cada componente.
	\item En la segunda fase del LCA, subir ambos nodos mientras \texttt{up[k][u] != up[k][v]}; el LCA final es \texttt{up[0][u]}.
	\item El padre de la raiz es 0 (o -1); aniadir check si se sale del rango.
\end{itemize}

\paragraph{Codigo.}
Ver \texttt{cpp.tex}, seccion \textit{Binary Lifting / LCA}.

\vspace{0.5em}

\subsubsection{Heavy-Light Decomposition (HLD)}

\paragraph{Idea intuitiva.}
Descomponer un arbol en \textbf{cadenas pesadas} (paths donde cada nodo (salvo la raiz) tiene al hijo ``pesado'' como hijo de tamano maximo de su subarbol). Cada cadena se mapea a un rango contiguo en un arreglo base; asi, queries sobre \textbf{paths} se descomponen en $O(\log n)$ queries sobre rangos contiguos (cada cadena es un rango), que se resuelven con un BIT o segment tree. La intuicion: por cada nivel de recursion, el subarbol del hijo pesado es al menos la mitad, asi que un path cruza $\le 2 \log n$ cadenas.

\paragraph{Propiedades clave (invariantes).}
\begin{itemize}\setlength\itemsep{0pt}
	\item \texttt{heavy[v]} = hijo pesado de $v$ (el de mayor subarbol).
	\item \texttt{head[v]} = cabeza de la cadena que contiene a $v$.
	\item \texttt{pos[v]} = posicion en el arreglo base.
	\item Cada nodo pertenece a exactamente una cadena.
\end{itemize}

\paragraph{Codigo clave.}
El query sobre un path $u \to v$:
\begin{verbatim}
long long pathQuery(int u, int v) {
    long long res = 0;
    while (head[u] != head[v]) {
        if (depth[head[u]] < depth[head[v]]) swap(u, v);
        res += seg.query(pos[head[u]], pos[u]);
        u = parent[head[u]];
    }
    if (depth[u] > depth[v]) swap(u, v);
    res += seg.query(pos[u], pos[v]);
    return res;
}
\end{verbatim}

\paragraph{Complejidad.}
\begin{itemize}\setlength\itemsep{0pt}
	\item Build $O(n)$, query sobre path $O(\log^2 n)$ (con segtree) o $O(\log n)$ (con BIT si solo se hace suma).
	\item Memoria $O(n)$.
\end{itemize}

\paragraph{Donde tocar para modificar.}
\begin{itemize}\setlength\itemsep{0pt}
	\item \textbf{Usar con segment tree}: aniadir un segtree sobre el arreglo base; \texttt{pathSegments} da los rangos a actualizar/query.
	\item \textbf{Cambiar operacion}: ajustar el segtree (suma, max, etc.).
	\item \textbf{Subtree query}: \texttt{subtree[v]} es el rango \texttt{[pos[v], pos[v] + size[v] - 1]} en el arreglo base.
	\item \textbf{Updates en edges}: aniadir un offset para que el edge $(u, v)$ quede en el nodo mas profundo.
\end{itemize}

\paragraph{Trampas y errores frecuentes.}
\begin{itemize}\setlength\itemsep{0pt}
	\item \texttt{pathSegments} requiere un \texttt{segtree} o \texttt{BIT} externo para ser util; por si solo solo da los rangos.
	\item Para edges vs nodos: si la operacion es sobre edges, guardar el valor en el hijo mas profundo.
	\item Olvidar la recursion \texttt{while (head[u] != head[v])} da resultados incorrectos.
	\item En problemas con updates, asegurarse de que el segtree soporte lazy propagation.
\end{itemize}

\paragraph{Codigo.}
Ver \texttt{cpp.tex}, seccion \textit{Heavy-Light Decomposition (HLD)}.

\vspace{0.5em}

\subsubsection{Centroid Decomposition}

\paragraph{Idea intuitiva.}
Descomponer recursivamente el arbol en \textbf{centroides}. Un centroide de un arbol es un nodo tal que cada subtree (al removerlo) tiene tamano $\le n/2$. Existe siempre y se encuentra en $O(n)$. Esto da una estructura de \textbf{arbol de centroides} de profundidad $O(\log n)$, util para problemas ``para cada nodo, encontrar el mas cercano que cumple X''. La razon de la profundidad $O(\log n)$: cada recursion reduce el problema a subarboles de tamano $\le n/2$, asi que la recursion da una cadena geometrica decreciente.

\paragraph{Propiedades clave (invariantes).}
\begin{itemize}\setlength\itemsep{0pt}
	\item En cada nivel, los subarboles se procesan recursivamente.
	\item \texttt{removed[]} marca nodos ya usados como centroide.
	\item \texttt{centroidParent[]} da el padre en el arbol de centroides.
	\item La altura del arbol de centroides es $O(\log n)$.
\end{itemize}

\paragraph{Codigo clave.}
Encontrar el centroide de un subarbol:
\begin{verbatim}
int findCentroid(int v, int p, int sz) {
    for (int u : adj[v])
        if (u != p && !removed[u] && subsize[u] > sz / 2)
            return findCentroid(u, v, sz);
    return v;
}
\end{verbatim}

\paragraph{Complejidad.}
\begin{itemize}\setlength\itemsep{0pt}
	\item Build $O(n \log n)$, queries dependen del problema.
	\item Cada nivel de recursion es $O(n)$ para calcular subsize.
\end{itemize}

\paragraph{Donde tocar para modificar.}
\begin{itemize}\setlength\itemsep{0pt}
	\item \textbf{Query tipica}: en cada nivel, ``expandir'' desde el centroide contando caminos validos y usar una estructura auxiliar (set, multiset).
	\item \textbf{Distancia maxima a un nodo ``malo''}: mantener un multiset de distancias en cada centroide.
	\item \textbf{Cambiar a otro tipo de descomposicion}: si el problema se adapta mejor a small-to-large o DSU on tree, usar esos.
	\item \textbf{Dynamic updates}: combinacion con link-cut trees.
\end{itemize}

\paragraph{Trampas y errores frecuentes.}
\begin{itemize}\setlength\itemsep{0pt}
	\item El snippet actual \textbf{solo construye} la descomposicion; las queries se aniaden segun el problema.
	\item Olvidar resetear \texttt{removed[]} causa recursion infinita.
	\item El calculo de subsize debe ignorar nodos ya ``removed''.
	\item El centroide es \textbf{unico} salvo empates en tamano; cualquiera sirve.
\end{itemize}

\paragraph{Codigo.}
Ver \texttt{cpp.tex}, seccion \textit{Centroid Decomposition}.

\vspace{0.5em}

\subsubsection{Euler Tour + range queries}

\paragraph{Idea intuitiva.}
Un \textbf{Euler Tour} (first-time visit) numera los nodos en el orden en que se visitan. El subarbol de $v$ corresponde a un \textbf{rango contiguo} $[tin[v], tout[v]]$ en este orden. Asi, queries sobre subarboles se resuelven con un BIT o segment tree sobre el arreglo de tiempos. La razon del rango contiguo: cuando entramos a un subarbol, no salimos hasta haber recorrido todos sus descendientes; entonces las posiciones consecutivas corresponden exactamente al subarbol.

\paragraph{Propiedades clave (invariantes).}
\begin{itemize}\setlength\itemsep{0pt}
	\item \texttt{tin[v]} = tiempo de primera visita a $v$.
	\item \texttt{tout[v]} = tiempo al \textbf{terminar} de procesar el subarbol de $v$.
	\item El subarbol de $v$ ocupa las posiciones $[tin[v], tout[v]]$.
	\item Para verificar si $u$ es ancestro de $v$: \texttt{tin[u] <= tin[v] \&\& tout[v] <= tout[u]}.
\end{itemize}

\paragraph{Codigo clave.}
El DFS del Euler Tour:
\begin{verbatim}
void dfs(int v, int p) {
    tin[v] = ++timer;
    for (int u : adj[v])
        if (u != p) dfs(u, v);
    tout[v] = timer;
}
\end{verbatim}

\paragraph{Complejidad.}
\begin{itemize}\setlength\itemsep{0pt}
	\item Build $O(n)$, query sobre subarbol $O(\log n)$ con BIT/segtree.
	\item Memoria $O(n)$.
\end{itemize}

\paragraph{Donde tocar para modificar.}
\begin{itemize}\setlength\itemsep{0pt}
	\item \textbf{Aniadir valor por nodo}: guardar el valor en \texttt{tin[v]} y operar sobre el rango.
	\item \textbf{Path query}: el Euler Tour \textbf{no} sirve para paths; usar HLD o LCA.
	\item \textbf{Multiple queries offline}: procesar en orden de tiempo y mantener una estructura por profundidad.
	\item \textbf{Verificar ancestro}: $O(1)$ con la condicion \texttt{tin}.
\end{itemize}

\paragraph{Trampas y errores frecuentes.}
\begin{itemize}\setlength\itemsep{0pt}
	\item El snippet solo construye el tour; las queries concretas se aniaden.
	\item \texttt{tout[v]} esta en el momento de salir, asi que \texttt{tout[v] - tin[v] + 1} = tamano del subarbol.
	\item Si el arbol no es conexo, aniadir un loop externo sobre cada componente.
	\item Para queries en tiempo (no subtree), usar el orden del tour y offline processing.
\end{itemize}

\paragraph{Codigo.}
Ver \texttt{cpp.tex}, seccion \textit{Euler Tour + range queries}.

\subsection{BST balanceado}

\subsubsection{Treap implicito}

\paragraph{Idea intuitiva.}
Un \textbf{treap} es un arbol binario de busqueda donde cada nodo tiene una \textbf{clave} (el valor de la secuencia) y una \textbf{prioridad} aleatoria. La invariante es: BST por clave, heap por prioridad. Esto da un arbol \textbf{esperado balanceado} sin necesidad de rebalanceo explicito. Un \textbf{treap implicito} no almacena una clave; la posicion de un nodo en la secuencia esta dada por el \textbf{tamano del subarbol izquierdo}, permitiendo split/merge por posicion. La estructura es esencialmente un RBS (randomized BST): la altura esperada es $O(\log n)$ por la propiedad de heap + BST.

\paragraph{Propiedades clave (invariantes).}
\begin{itemize}\setlength\itemsep{0pt}
	\item \texttt{sz(node)} = tamano del subarbol (1 + izquierdo + derecho).
	\item \texttt{split(t, k, l, r)}: parte \texttt{t} en \texttt{l} (primeros $k$ nodos) y \texttt{r} (resto).
	\item \texttt{merge(l, r)}: une dos treaps (todos los de \texttt{l} antes que los de \texttt{r}).
	\item La prioridad aleatoria asegura equilibrio esperado.
\end{itemize}

\paragraph{Codigo clave.}
Split por posicion (la base del treap implicito):
\begin{verbatim}
void split(Node* t, int k, Node*& l, Node*& r) {
    if (!t) { l = r = nullptr; return; }
    push(t);
    if (sz(t->l) >= k) {
        split(t->l, k, l, t->l);
        r = update(t);
    } else {
        split(t->r, k - sz(t->l) - 1, t->r, r);
        l = update(t);
    }
}
\end{verbatim}

\paragraph{Complejidad.}
\begin{itemize}\setlength\itemsep{0pt}
	\item $O(\log n)$ \textbf{esperado} por operacion (split, merge, insert, erase).
	\item Peor caso exponencial (probabilidad astronomicamente pequena).
\end{itemize}

\paragraph{Donde tocar para modificar.}
\begin{itemize}\setlength\itemsep{0pt}
	\item \textbf{Treap con clave (no implicito)}: usar \texttt{key} en vez de posicion; cambiar \texttt{split} por \texttt{splitByKey}.
	\item \textbf{Lazy propagation}: aniadir campos \texttt{lazy} y aplicarlos en \texttt{push} (recursivo).
	\item \textbf{Persistencia}: aniadir versionado clonando nodos.
	\item \textbf{Mejor aleatoriedad}: en lugar de \texttt{rand()}, usar \texttt{rng()} (mejor estado).
	\item \textbf{Operacion de reversa}: aniadir \texttt{lazyReverse} y swap de hijos en push.
\end{itemize}

\paragraph{Trampas y errores frecuentes.}
\begin{itemize}\setlength\itemsep{0pt}
	\item En el snippet actual, \texttt{Node(int v)} usa \texttt{rand()} que en algunos compiladores da maxima prioridad 0; aniadir \texttt{rng()} para mejor distribucion.
	\item Olvidar \texttt{update(t)} tras modificar hijos da tamano incorrecto.
	\item Si lazy propagation no se aplica en push, las queries dan valores viejos.
	\item Memory leak con \texttt{new}; en contests OK, en produccion aniadir destructor.
\end{itemize}

\paragraph{Codigo.}
Ver \texttt{cpp.tex}, seccion \textit{Treap implicito}.

\vspace{0.5em}

\subsubsection{Mo con updates}

\paragraph{Idea intuitiva.}
Extension del algoritmo de Mo para queries que mezclan \textbf{rango} y \textbf{updates}. Las queries se ordenan por $(L/\text{blockSize}, R/\text{blockSize}, T)$ donde $T$ es el ``tiempo'' (numero de updates). Asi, moverse entre queries consecutivas cuesta $O(n^{2/3})$ amortizado. La intuicion: el sort por bloques reduce los movimientos totales, y la tercera coordenada $T$ se optimiza con alternancia similar a $R$.

\paragraph{Propiedades clave (invariantes).}
\begin{itemize}\setlength\itemsep{0pt}
	\item Tres punteros: $L$, $R$, $T$.
	\item \texttt{applyUpdate} y \texttt{revertUpdate} son \textbf{inversos} exactos.
	\item \texttt{addPos} y \texttt{removePos} mantienen el estado.
	\item El tamano de bloque optimo es $n^{2/3}$ (donde $n$ = tamano del array, $Q$ = numero de queries).
\end{itemize}

\paragraph{Codigo clave.}
El sort con tres coordenadas:
\begin{verbatim}
int blockSize = (int)pow(n, 2.0/3.0);
sort(queries.begin(), queries.end(), [&](const Q& a, const Q& b) {
    int blockA = a.l / blockSize, blockB = b.l / blockSize;
    if (blockA != blockB) return blockA < blockB;
    int blockAR = a.r / blockSize, blockBR = b.r / blockSize;
    if (blockAR != blockBR) return blockAR < blockBR;
    return a.t < b.t;
});
\end{verbatim}

\paragraph{Complejidad.}
\begin{itemize}\setlength\itemsep{0pt}
	\item $O((N + Q) \cdot n^{2/3})$ donde $n^{2/3}$ es el tamano de bloque optimo.
	\item En problemas reales, funciona bien hasta $N, Q \sim 10^5$.
\end{itemize}

\paragraph{Donde tocar para modificar.}
\begin{itemize}\setlength\itemsep{0pt}
	\item \textbf{El corazon es el ``add/remove''}: aniadir la logica del problema (ej. contar colores distintos, suma de frecuencias, etc.).
	\item \textbf{Cambiar blockSize}: si la mayoria de updates vs queries cambia, ajustar.
	\item \textbf{Rollback}: aniadir \texttt{save()} / \texttt{restore()} para volver a estados anteriores.
	\item \textbf{Hilbert order}: alternativa con mejor locality.
\end{itemize}

\paragraph{Trampas y errores frecuentes.}
\begin{itemize}\setlength\itemsep{0pt}
	\item El snippet actual es solo el \textbf{esqueleto}; \texttt{addPos}/\texttt{removePos}/\texttt{applyUpdate}/\texttt{revertUpdate} deben ser implementados por problema.
	\item Olvidar \texttt{revertUpdate} es el error mas comun.
	\item Si el numero de updates es muy grande, Mo con updates se vuelve lento; considerar otra estructura.
	\item Mantener consistencia entre \texttt{T} y la posicion de los updates.
\end{itemize}

\paragraph{Codigo.}
Ver \texttt{cpp.tex}, seccion \textit{Mo con updates}.

\vspace{0.5em}

\subsubsection{sqrt-decomposition generico}

\paragraph{Idea intuitiva.}
Dividir el arreglo en \textbf{bloques} de tamano $B \approx \sqrt{n}$. Mantener informacion precomputada por bloque (suma, max, etc.). Queries de rango se resuelven combinando bloques enteros ($O(\sqrt{n})$ en total) y elementos sueltos en los bloques parciales. Updates puntuales actualizan el valor y el resumen del bloque. La intuicion: las queries ``normales'' cruzan $O(n/B)$ bloques; tomar $B = \sqrt{n}$ balancea el costo de bloques enteros vs elementos sueltos.

\paragraph{Propiedades clave (invariantes).}
\begin{itemize}\setlength\itemsep{0pt}
	\item $B = \lfloor \sqrt{n} \rfloor + 1$.
	\item \texttt{bucket[i]} = resumen del bloque $i$-esimo (suma, max, etc.).
	\item \texttt{rangeSum(l, r)}: 3 casos segun $l$ y $r$ esten en el mismo bloque o no.
	\item Cada query cruza $\le 2 + (r - l)/B$ bloques.
\end{itemize}

\paragraph{Codigo clave.}
Estructura basica:
\begin{verbatim}
int B = sqrt(n) + 1;
vector<T> bucket(n / B + 1, IDENTITY);
T query(int l, int r) {
    T res = IDENTITY;
    int bl = l / B, br = r / B;
    if (bl == br) {
        for (int i = l; i <= r; ++i) res = combine(res, a[i]);
    } else {
        for (int i = l; i < (bl+1)*B; ++i) res = combine(res, a[i]);
        for (int i = bl+1; i < br; ++i) res = combine(res, bucket[i]);
        for (int i = br*B; i <= r; ++i) res = combine(res, a[i]);
    }
    return res;
}
\end{verbatim}

\paragraph{Complejidad.}
\begin{itemize}\setlength\itemsep{0pt}
	\item $O(\sqrt{n})$ por operacion.
	\item Memoria $O(n)$.
\end{itemize}

\paragraph{Donde tocar para modificar.}
\begin{itemize}\setlength\itemsep{0pt}
	\item \textbf{Cambiar suma por max/min}: ajustar \texttt{bucket} y la combinacion.
	\item \textbf{Aniadir update en rango}: si la operacion es asociativa, sumar al resumen del bloque + ajustar elementos sueltos.
	\item \textbf{Tamano de bloque optimo}: si las queries son muy largas, $B = n^{2/3}$ (sirve para Mo con updates).
	\item \textbf{Operacion no asociativa}: usar segment tree en su lugar.
\end{itemize}

\paragraph{Trampas y errores frecuentes.}
\begin{itemize}\setlength\itemsep{0pt}
	\item Si $n$ no es multiplo de $B$, el ultimo bloque es mas chico; aniadir check.
	\item La identidad de la operacion debe estar bien definida (0 para suma, $-\infty$ para max).
	\item Si la operacion no es asociativa, los bloques no se pueden combinar; cambiar a segment tree.
\end{itemize}

\paragraph{Codigo.}
Ver \texttt{cpp.tex}, seccion \textit{sqrt-decomposition generico}.

% ============================================================
\section{Strings}
% ============================================================

\subsection{Hashing}

\subsubsection{Rolling Hash (doble modulo)}

\paragraph{Idea intuitiva.}
Un \textbf{rolling hash} asigna a cada prefijo del string un valor numerico tal que el hash de un substring se puede obtener en $O(1)$ a partir de los prefijos. La forma clasica: $h[i] = h[i-1] \cdot A + s[i-1]$ (donde $A$ es una base y todo modulo $M$). Asi, $h[r] - h[l-1] \cdot A^{r-l+1}$ da el hash del substring $[l, r]$ (con $A^{r-l+1}$ precomputado). Usar \textbf{dos modulos} distintos reduce drasticamente la probabilidad de colision. La idea algebraica: el string se interpreta como un polinomio en $A$ con coeficientes los caracteres; el hash es su evaluacion modulo $M$.

\paragraph{Propiedades clave (invariantes).}
\begin{itemize}\setlength\itemsep{0pt}
	\item $h[0] = 0$, $h[i]$ = hash de los primeros $i$ caracteres.
	\item $p[i] = A^i \bmod M$ precomputado.
	\item Comparar dos substrings: comparar sus pares de hashes.
	\item \textbf{Colisiones}: posibles (probabilidad $1/M$ por par). Doble hash: $\approx 1/(M_1 M_2)$.
	\item Si $A$ y $M$ son coprimos, el hash es bijectivo modulo $M$ para strings de longitud fija.
\end{itemize}

\paragraph{Codigo clave.}
La formula del hash de substring:
\begin{verbatim}
pair<long long, long long> get(int l, int r) {
    // h[r] - h[l-1] * A^(r-l+1)
    long long x1 = (h1[r] - h1[l-1] * p1[r-l+1] % M1 + M1) % M1;
    long long x2 = (h2[r] - h2[l-1] * p2[r-l+1] % M2 + M2) % M2;
    return {x1, x2};
}
\end{verbatim}

\paragraph{Complejidad.}
\begin{itemize}\setlength\itemsep{0pt}
	\item Precompute $O(n)$, query hash $O(1)$.
	\item Memoria $O(n)$.
\end{itemize}

\paragraph{Donde tocar para modificar.}
\begin{itemize}\setlength\itemsep{0pt}
	\item \textbf{Caracteres con valor $> 26$}: si los strings incluyen simbolos, aniadir \texttt{if (c >= 'a' \&\& c <= 'z') v = c - 'a'; else v = 26 + ...}.
	\item \textbf{Hash de substrings con updates}: aniadir un Fenwick tree / segment tree sobre los hashes (no es trivial).
	\item \textbf{Cambiar mod}: usar dos primos grandes ($10^9+7$ y $10^9+9$) o un \texttt{u128} para un solo hash sin mod.
	\item \textbf{Base aleatoria}: aniadir random para evitar ataques (en contests no es necesario; en produccion si).
	\item \textbf{Hash polinomial}: $h = \sum s_i \cdot A^i$ en vez de $s_i \cdot A^{i-1}$; algunos problemas requieren este.
\end{itemize}

\paragraph{Trampas y errores frecuentes.}
\begin{itemize}\setlength\itemsep{0pt}
	\item Si los substrings a comparar estan vacios (\texttt{l > r}), su hash es 0; aniadir check.
	\item \textbf{Overflow} en \texttt{h[i-1] * A}: usar \texttt{long long} y aplicar \texttt{\% M} en cada operacion (no al final).
	\item La sustraccion \texttt{h[r] - h[l-1]*A^(r-l+1)} puede dar negativo; aniadir \texttt{+M} antes de \texttt{\% M}.
	\item Con doble hash, verificar \textbf{ambos} componentes; si solo uno coincide, es colision.
\end{itemize}

\paragraph{Codigo.}
Ver \texttt{cpp.tex}, seccion \textit{Rolling Hash (doble modulo)}.

\subsection{Patrones}

\subsubsection{KMP + LPS}

\paragraph{Idea intuitiva.}
El algoritmo \textbf{Knuth-Morris-Pratt} busca todas las ocurrencias de un patron $P$ en un texto $T$ en $O(|T| + |P|)$. La clave es la tabla \textbf{LPS} (Longest Proper Prefix which is Suffix) que, para cada posicion $i$ del patron, dice el largo del prefijo mas largo que tambien es sufijo \textbf{propio}. Durante el match, cuando hay mismatch, en vez de volver al inicio, saltamos al LPS anterior, evitando trabajo redundante. La demostracion de $O(|T| + |P|)$: cada iteracion del bucle externo avanza $i$ o decrece $j$ (con $j$ acotado por $|P|$), asi que en total hay $O(|T| + |P|)$ operaciones.

\paragraph{Propiedades clave (invariantes).}
\begin{itemize}\setlength\itemsep{0pt}
	\item \texttt{lps[i]} = largo del prefijo mas largo que es sufijo propio de \texttt{p[0..i]}.
	\item El match usa dos punteros \texttt{i} (texto) y \texttt{j} (patron); en mismatch, $j = lps[j-1]$.
	\item Cuando $j == |P|$, se encontro una ocurrencia; continuar con $j = lps[j-1]$ para buscar mas.
	\item El LPS se construye con una auto-aplicacion de KMP al patron.
\end{itemize}

\paragraph{Codigo clave.}
Construccion de la tabla LPS:
\begin{verbatim}
for (int i = 1; i < m; ++i) {
    int j = lps[i - 1];
    while (j > 0 && p[i] != p[j]) j = lps[j - 1];
    if (p[i] == p[j]) ++j;
    lps[i] = j;
}
\end{verbatim}

\paragraph{Complejidad.}
\begin{itemize}\setlength\itemsep{0pt}
	\item $O(|T| + |P|)$.
	\item En la practica, rapido y robusto.
\end{itemize}

\paragraph{Donde tocar para modificar.}
\begin{itemize}\setlength\itemsep{0pt}
	\item \textbf{Contar ocurrencias}: el snippet ya cuenta; para devolver posiciones, aniadir un \texttt{vector<int>}.
	\item \textbf{Varios patrones}: usar Aho-Corasick en su lugar.
	\item \textbf{Matcher exacto con wildcard}: aniadir logica para que \texttt{?} matchee cualquier caracter.
	\item \textbf{Tamano pequeno}: si el patron es $\le 5$ caracteres, Rabin-Karp es mas simple.
	\item \textbf{2D KMP}: aniadir una dimension; util en problemas de matrices.
\end{itemize}

\paragraph{Trampas y errores frecuentes.}
\begin{itemize}\setlength\itemsep{0pt}
	\item En el snippet, \texttt{i} y \texttt{j} son 0-based. La condicion \texttt{j == sz\_pattern} para contar funciona.
	\item La construccion de LPS usa el mismo patron como ``texto''.
	\item Si el patron tiene caracteres repetidos, el LPS puede ser no trivial; verificar con casos pequenos.
\end{itemize}

\paragraph{Codigo.}
Ver \texttt{cpp.tex}, seccion \textit{KMP + LPS}.

\vspace{0.5em}

\subsubsection{Z-algorithm}

\paragraph{Idea intuitiva.}
La \textbf{Z-function} $Z[i]$ de un string $S$ es el largo del prefijo mas largo de $S$ que \textbf{tambien} empieza en $S[i]$. Asi, $Z[0] = n$ por convencion. Se calcula en $O(n)$ manteniendo una ventana $[l, r]$ que es el match mas largo encontrado hasta ahora. La consulta $S[l..r] = S[0..r-l]$ da informacion sobre $Z[i]$ cuando $i \le r$. La intuicion: cualquier $i$ dentro de la ventana ``ve'' los caracteres desde una posicion conocida, asi que $Z[i]$ se puede copiar en lugar de recalcular.

\paragraph{Propiedades clave (invariantes).}
\begin{itemize}\setlength\itemsep{0pt}
	\item Construir $Z$ cuesta $O(n)$ con el truco de la ventana.
	\item Para matching de patron $P$ en texto $T$: concatenar $P + \# + T$ y aniadir todos los $i$ con $Z[i] = |P|$.
	\item La ventana $[l, r]$ es siempre la del match mas largo a la derecha.
\end{itemize}

\paragraph{Codigo clave.}
Construccion de Z-function en $O(n)$:
\begin{verbatim}
Z[0] = n;
int l = 0, r = 0;
for (int i = 1; i < n; ++i) {
    if (i <= r) Z[i] = min(r - i + 1, Z[i - l]);
    while (i + Z[i] < n && s[Z[i]] == s[i + Z[i]]) ++Z[i];
    if (i + Z[i] - 1 > r) { l = i; r = i + Z[i] - 1; }
}
\end{verbatim}

\paragraph{Complejidad.}
\begin{itemize}\setlength\itemsep{0pt}
	\item $O(n)$ para construir; matching $O(n + m)$.
	\item Cada caracter se compara a lo sumo 2 veces (extension + reduccion).
\end{itemize}

\paragraph{Donde tocar para modificar.}
\begin{itemize}\setlength\itemsep{0pt}
	\item \textbf{Detectar bordes (border)}: $Z[i] = n - i$ indica que $S[0..n-i-1]$ es borde.
	\item \textbf{Prefijo comun mas largo de dos strings}: aniadir $S_1 + \# + S_2$; $Z[|S_1|+1]$ es el prefijo comun.
	\item \textbf{Comparar prefijos y sufijos}: combinar $S + \# + \text{reverse}(S)$.
	\item \textbf{Matching aproximado}: aniadir tolerancia a mismatches.
\end{itemize}

\paragraph{Trampas y errores frecuentes.}
\begin{itemize}\setlength\itemsep{0pt}
	\item $Z[0]$ siempre vale $n$ (no se calcula).
	\item Si $i > r$, la ventana no aporta; empezar desde 0.
	\item El caracter separador en $P + \# + T$ debe ser uno que no aparezca en $P$ o $T$.
\end{itemize}

\paragraph{Codigo.}
Ver \texttt{cpp.tex}, seccion \textit{Z-algorithm}.

\vspace{0.5em}

\subsubsection{Rabin-Karp (single hash)}

\paragraph{Idea intuitiva.}
Rolling hash aplicado a pattern matching: calcular el hash del patron $P$ y de cada ventana de tamano $|P|$ en $T$. Si los hashes coinciden, verificar caracter por caracter (para descartar colisiones). Esto da $O(n + m)$ \textbf{esperado}. La razon de la eficiencia esperada: cada ventana tiene probabilidad $1/M$ de tener un hash coincidente por casualidad, asi que solo se verifica $O(n/M)$ ventanas en promedio.

\paragraph{Propiedades clave (invariantes).}
\begin{itemize}\setlength\itemsep{0pt}
	\item Si los hashes difieren, los substrings son distintos (sin colision).
	\item Si coinciden, \textbf{verificar} (la verificacion es $O(m)$ pero rara vez se ejecuta).
	\item Peor caso (muchas colisiones): $O(nm)$.
	\item Para $M \ge 10^9$, las colisiones son raras; en la practica nunca pasan.
\end{itemize}

\paragraph{Codigo clave.}
Rolling hash del texto (deslizar ventana):
\begin{verbatim}
long long hash = 0;
for (int i = 0; i < m; ++i) hash = (hash * BASE + text[i]) % MOD;
for (int i = m; i < n; ++i) {
    if (hash == target) check(i - m);
    hash = (hash - text[i - m] * h) * BASE + text[i];  // mod correctamente
    hash = (hash % MOD + MOD) % MOD;
}
\end{verbatim}

\paragraph{Complejidad.}
\begin{itemize}\setlength\itemsep{0pt}
	\item Esperado $O(n + m)$, peor $O(nm)$.
	\item Con doble hash y $M_1, M_2 \ge 10^9$, el peor caso es astronomicamente improbable.
\end{itemize}

\paragraph{Donde tocar para modificar.}
\begin{itemize}\setlength\itemsep{0pt}
	\item \textbf{Doble hash}: aniadir un segundo par \texttt{(BASE2, MOD2)} para reducir colisiones.
	\item \textbf{Varios patrones}: aniadir un set de hashes y buscar match con cualquiera.
	\item \textbf{Tamano de ventana variable}: si los patrones tienen largos distintos, usar el largo especifico por patron.
	\item \textbf{Multi-pattern search}: usar Aho-Corasick en su lugar.
\end{itemize}

\paragraph{Trampas y errores frecuentes.}
\begin{itemize}\setlength\itemsep{0pt}
	\item Overflow en \texttt{(tHash - text[i] * h) * BASE}: aplicar mod correctamente y aniadir \texttt{+MOD} antes de multiplicar.
	\item El factor $h = \text{BASE}^{m-1} \bmod \text{MOD}$ se precomputa.
	\item Si \texttt{MOD} no es primo, las colisiones son mas frecuentes; usar primo.
	\item En el peor caso, Rabin-Karp puede ser lento; preferir KMP si se requiere worst-case.
\end{itemize}

\paragraph{Codigo.}
Ver \texttt{cpp.tex}, seccion \textit{Rabin-Karp (single hash)}.

\vspace{0.5em}

\subsubsection{Aho-Corasick (basic)}

\paragraph{Idea intuitiva.}
Construye un \textbf{trie} con todos los patrones y aniade \textbf{failure links} analogos a los de KMP: el failure link de un nodo $v$ apunta al nodo mas largo que es sufijo del string representado por $v$ y sigue siendo prefijo de algun patron. Asi, al recorrer el texto, si hay mismatch, saltamos al failure link. Esto permite buscar \textbf{todos} los patrones en $O(|T| + \sum |P_i|)$. La idea: unificar varios KMP en una sola estructura.

\paragraph{Propiedades clave (invariantes).}
\begin{itemize}\setlength\itemsep{0pt}
	\item El trie almacena los patrones.
	\item \texttt{link[v]} (failure link) se calcula en BFS en orden de profundidad creciente.
	\item \texttt{out[v]} lista los IDs de patrones que terminan en $v$.
	\item Para queries eficientes, aniadir \textbf{output links}: propagar \texttt{out} de los failure links a los hijos.
\end{itemize}

\paragraph{Codigo clave.}
BFS para construir los failure links:
\begin{verbatim}
queue<int> q;
for (int c = 0; c < SIGMA; ++c)
    if (trie[0].next[c] != -1) {
        q.push(trie[0].next[c]);
        trie[trie[0].next[c]].link = 0;
    } else trie[0].next[c] = 0;  // goto a la raiz
while (!q.empty()) {
    int v = q.front(); q.pop();
    for (int pid : trie[trie[v].link].out)
        trie[v].out.push_back(pid);
    for (int c = 0; c < SIGMA; ++c) {
        if (trie[v].next[c] != -1) {
            trie[trie[v].next[c]].link = trie[trie[v].link].next[c];
            q.push(trie[v].next[c]);
        } else trie[v].next[c] = trie[trie[v].link].next[c];
    }
}
\end{verbatim}

\paragraph{Complejidad.}
\begin{itemize}\setlength\itemsep{0pt}
	\item Build $O(\sum |P_i| \cdot \Sigma)$, query $O(|T| + \text{ocurrencias})$.
\end{itemize}

\paragraph{Donde tocar para modificar.}
\begin{itemize}\setlength\itemsep{0pt}
	\item \textbf{Output links en build}: ya hecho en el snippet (\texttt{for(int pid : t[t[u].link].out) t[u].out.push\_back(pid);}).
	\item \textbf{Transiciones automaticas}: aniadir el goto \texttt{t[v].next[k] = t[t[v].link].next[k]} cuando el enlace directo es \texttt{-1}; asi el bucle de matching es un solo paso por caracter.
	\item \textbf{Aplicar a problemas tipo DP}: contar formas de construir strings evitando los patrones (ver el siguiente modulo).
	\item \textbf{2D Aho-Corasick}: aniadir una dimension para busqueda en matrices.
\end{itemize}

\paragraph{Trampas y errores frecuentes.}
\begin{itemize}\setlength\itemsep{0pt}
	\item La BFS debe empezar con la raiz en la cola y los hijos de la raiz con \texttt{link = 0}.
	\item Olvidar propagar \texttt{out} da ocurrencias perdidas.
	\item Si no se aniaden las transiciones automaticas (goto), el matching es mas lento.
	\item Con muchos patrones, el output puede ser enorme; aniadir contador en lugar de lista.
\end{itemize}

\paragraph{Codigo.}
Ver \texttt{cpp.tex}, seccion \textit{Aho-Corasick (basic)}.

\vspace{0.5em}

\subsubsection{Aho-Corasick DP for forbidden patterns}

\paragraph{Idea intuitiva.}
Contar strings de longitud $L$ que \textbf{evitan} un conjunto de patrones. Se modela como una DP sobre el \textbf{estado del automaton} de Aho-Corasick: $dp[v]$ = numero de formas de llegar al estado $v$ despues de cierto numero de caracteres. Transicion: por cada caracter $c$ en el alfabeto, ir al estado \texttt{next[v][c]}, pero \textbf{solo} si ese estado no tiene un patron (de lo contrario, el string contendria el patron prohibido). La idea: cada estado del automaton representa un ``prefijo conocido'' y las transiciones simulan ``agregar un caracter''.

\paragraph{Propiedades clave (invariantes).}
\begin{itemize}\setlength\itemsep{0pt}
	\item $dp$ se mantiene por nivel (longitud del string construido).
	\item Transicion: $ndp[u] = \sum_{c} dp[v]$ donde $u = next[v][c]$ y $u$ no es ``terminal'' (no tiene patron que termina ahi).
	\item Respuesta: suma de $dp[v]$ sobre todos los estados $v$ no terminales al final.
	\item Si todos los caracteres son permitidos, $|\Sigma|$ entra como factor.
\end{itemize}

\paragraph{Codigo clave.}
DP sobre el automaton:
\begin{verbatim}
vector<long long> dp(numStates, 0), ndp(numStates);
dp[0] = 1;  // empezar en la raiz
for (int step = 0; step < L; ++step) {
    fill(ndp.begin(), ndp.end(), 0);
    for (int v = 0; v < numStates; ++v) if (dp[v] && !isTerminal[v])
        for (int c = 0; c < SIGMA; ++c) {
            int u = next[v][c];
            if (!isTerminal[u]) ndp[u] = (ndp[u] + dp[v]) % MOD;
        }
    swap(dp, ndp);
}
long long ans = 0;
for (int v = 0; v < numStates; ++v) if (!isTerminal[v])
    ans = (ans + dp[v]) % MOD;
\end{verbatim}

\paragraph{Complejidad.}
\begin{itemize}\setlength\itemsep{0pt}
	\item $O(L \cdot |S| \cdot |\Sigma|)$ donde $|S|$ es el numero de estados y $|\Sigma|$ el alfabeto.
	\item Para $|\Sigma| = 2$ y $|S| \le 10^4$: $\sim 2 \cdot 10^{10}$ para $L = 10^6$ (puede ser lento).
\end{itemize}

\paragraph{Donde tocar para modificar.}
\begin{itemize}\setlength\itemsep{0pt}
	\item \textbf{Alfabeto mas pequeno}: si el problema es solo $\{0, 1\}$, $|\Sigma| = 2$.
	\item \textbf{Permitir ocurrencias parciales}: cambiar la condicion a ``no estado terminal'' o relajar segun el problema.
	\item \textbf{Recuperar el string}: aniadir backtracking si el problema lo pide.
	\item \textbf{Matriz de transicion}: para queries rapidas, precomputar $M^c$ con exponentiation.
\end{itemize}

\paragraph{Trampas y errores frecuentes.}
\begin{itemize}\setlength\itemsep{0pt}
	\item El snippet asume Aho-Corasick ya construido con transiciones automaticas. Sin eso, aniadir el goto en el bucle.
	\item Si $L$ es muy grande, la DP lineal es lenta; usar matrix exponentiation.
	\item Estados terminales deben excluirse de \textbf{ambos} lados (estado actual y siguiente).
\end{itemize}

\paragraph{Codigo.}
Ver \texttt{cpp.tex}, seccion \textit{Aho-Corasick DP for forbidden patterns}.

\subsection{Estructuras avanzadas}

\subsubsection{Suffix Array}

\paragraph{Idea intuitiva.}
Un \textbf{suffix array} $SA$ es un arreglo que contiene las posiciones de inicio de los sufijos del string $S$, ordenados lexicograficamente. Se construye en $O(n \log n)$ con doubling: ordenar por el primer caracter, luego por los primeros 2, luego 4, etc. Para hacer cada paso lineal, se usa \textbf{counting sort} sobre los ranks. La idea: cada nivel dobla la cantidad de caracteres considerados, asi que en $\log n$ niveles los sufijos estan completamente ordenados.

\paragraph{Propiedades clave (invariantes).}
\begin{itemize}\setlength\itemsep{0pt}
	\item \texttt{p[i]} = posicion de inicio del $i$-esimo sufijo (ordenado).
	\item \texttt{c[i]} = rank del sufijo que empieza en $i$.
	\item Cada iteracion duplica la longitud de la ``ventana'' considerada.
	\item El algoritmo termina cuando $2^k \ge n$.
	\item La permutacion \texttt{p} es estable entre iteraciones.
\end{itemize}

\paragraph{Codigo clave.}
El paso de doubling (clasificar por los primeros $2^k$ caracteres):
\begin{verbatim}
for (int k = 0; (1 << k) < n; ++k) {
    for (int i = 0; i < n; ++i) pn[i] = (p[i] - (1 << k) + n) % n;
    // counting sort por c[pn[i]] segundo, c[i] primero
    // ...
    cn[pn[0]] = 0;
    for (int i = 1; i < n; ++i)
        cn[pn[i]] = cn[pn[i-1]] + (pair != pair_prev);
    c = cn;
}
\end{verbatim}

\paragraph{Complejidad.}
\begin{itemize}\setlength\itemsep{0pt}
	\item $O(n \log n)$ con counting sort.
	\item Constante pequena; en la practica eficiente.
\end{itemize}

\paragraph{Donde tocar para modificar.}
\begin{itemize}\setlength\itemsep{0pt}
	\item \textbf{Suffix array en $O(n)$}: usar SA-IS (mas complejo, no incluido).
	\item \textbf{Buscar patron}: con SA + LCP es $O(|P| \log n)$ por busqueda binaria.
	\item \textbf{Suffix array de varios strings}: aniadir centinelas (\texttt{\$}) entre ellos.
	\item \textbf{Aniadir longitud}: aniadir \texttt{n - p[i]} como clave secundaria (sufijos mas cortos primero).
\end{itemize}

\paragraph{Trampas y errores frecuentes.}
\begin{itemize}\setlength\itemsep{0pt}
	\item La ``ventana'' $2^k$ se hace circular: \texttt{p[i] = (p[i] - (1<<k) + n) \% n}.
	\item Para contar bien los ranks duplicados, comparar pares \texttt{(c[i], c[(i+2\^{}k) \% n])}.
	\item El counting sort necesita un array de acumulacion de tamano $\le n$.
	\item Si los caracteres son negativos o fuera del alfabeto esperado, normalizar antes.
\end{itemize}

\paragraph{Codigo.}
Ver \texttt{cpp.tex}, seccion \textit{Suffix Array}.

\vspace{0.5em}

\subsubsection{LCP Array (Kasai)}

\paragraph{Idea intuitiva.}
El \textbf{LCP array} (Longest Common Prefix) guarda, para cada par de sufijos adyacentes en el SA, el largo de su prefijo comun. El algoritmo de \textbf{Kasai} lo calcula en $O(n)$ iterando los sufijos \textbf{en orden del string original} (no del SA), y reutiliza el LCP previo. La demostracion: el sufijo en $i+1$ se obtiene del sufijo en $i$ ``avanzando un caracter''; los LCPs decrecen como maximo 1 entre iteraciones.

\paragraph{Propiedades clave (invariantes).}
\begin{itemize}\setlength\itemsep{0pt}
	\item \texttt{lcp[i]} = LCP entre los sufijos en \texttt{sa[i]} y \texttt{sa[i-1]}.
	\item Se necesita \texttt{rank\_[i]} (la inversa del SA) para iterar.
	\item El truco de Kasai: si \texttt{lcp > 0} para el sufijo anterior, el actual parte de al menos \texttt{lcp - 1}.
	\item La suma de todos los LCPs da informacion sobre repeticiones.
\end{itemize}

\paragraph{Codigo clave.}
Algoritmo de Kasai:
\begin{verbatim}
vector<int> rank(n);
for (int i = 0; i < n; ++i) rank[sa[i]] = i;
int k = 0;
vector<int> lcp(n - 1);
for (int i = 0; i < n; ++i) {
    if (rank[i] == n - 1) { k = 0; continue; }
    int j = sa[rank[i] + 1];
    while (i + k < n && j + k < n && s[i+k] == s[j+k]) ++k;
    lcp[rank[i]] = k;
    if (k) --k;
}
\end{verbatim}

\paragraph{Complejidad.}
\begin{itemize}\setlength\itemsep{0pt}
	\item $O(n)$.
	\item Cada comparacion reduce $k$ o avanza linealmente; amortizado $O(1)$.
\end{itemize}

\paragraph{Donde tocar para modificar.}
\begin{itemize}\setlength\itemsep{0pt}
	\item \textbf{RMQ sobre LCP}: para queries ``LCP de dos sufijos'', encontrar sus ranks y hacer RMQ en \texttt{lcp}.
	\item \textbf{Contar substrings distintas}: $\sum (n - sa[i] - lcp[i])$.
	\item \textbf{Longest common substring}: maximo en \texttt{lcp}.
	\item \textbf{String matching}: busqueda binaria sobre SA + LCP verification.
\end{itemize}

\paragraph{Trampas y errores frecuentes.}
\begin{itemize}\setlength\itemsep{0pt}
	\item Si dos sufijos son iguales (raro, solo si el string tiene period $< n$), \texttt{lcp} puede ser \texttt{n - sa[i]}.
	\item En la iteracion, aniadir el chequeo \texttt{i + k < n} para evitar overflow.
	\item El \texttt{--k} en la salida es crucial para la complejidad amortizada.
\end{itemize}

\paragraph{Codigo.}
Ver \texttt{cpp.tex}, seccion \textit{LCP Array (Kasai)}.

\vspace{0.5em}

\subsubsection{Suffix Automaton}

\paragraph{Idea intuitiva.}
Un \textbf{suffix automaton} es un DFA minimo que reconoce todos los substrings del string original. Tiene la propiedad magica de tener $\le 2n - 1$ estados. Cada estado representa una clase de equivalencia de substrings (los que tienen el mismo conjunto de posiciones finales). El \textbf{link} de un estado apunta al estado del sufijo propio mas largo. La demostracion de $\le 2n - 1$: cada nuevo caracter anade a lo sumo 2 estados (uno nuevo y un clon), asi que en $n$ pasos hay $\le 2n - 1$ estados.

\paragraph{Propiedades clave (invariantes).}
\begin{itemize}\setlength\itemsep{0pt}
	\item \texttt{len[v]}: largo maximo de los substrings de la clase.
	\item \texttt{link[v]}: padre en el arbol de sufijos (suffix link).
	\item \texttt{next[v][c]}: transicion por caracter.
	\item \texttt{occ[v]}: numero de posiciones finales (cuantas veces aparece el patron).
	\item Cada estado tiene $\texttt{len[v]} - \texttt{len[link[v]]}$ substrings ``propios''.
\end{itemize}

\paragraph{Codigo clave.}
El corazon de la extension SAM:
\begin{verbatim}
void extend(int c) {
    int cur = size++;
    st[cur].len = st[last].len + 1;
    st[cur].occ = 1;
    int p = last;
    while (p != -1 && !st[p].next.count(c)) {
        st[p].next[c] = cur;
        p = st[p].link;
    }
    if (p == -1) st[cur].link = 0;
    else {
        int q = st[p].next[c];
        if (st[p].len + 1 == st[q].len) st[cur].link = q;
        else {
            int clone = size++;
            st[clone] = st[q];
            st[clone].len = st[p].len + 1;
            st[clone].occ = 0;  // clon NO es nueva posicion final
            while (p != -1 && st[p].next[c] == q) {
                st[p].next[c] = clone;
                p = st[p].link;
            }
            st[q].link = st[cur].link = clone;
        }
    }
    last = cur;
}
\end{verbatim}

\paragraph{Complejidad.}
\begin{itemize}\setlength\itemsep{0pt}
	\item Build $O(n)$ con la operacion de \textbf{clonar} estados.
	\item Memoria $O(n)$ (con map) o $O(n \cdot \Sigma)$ (con array).
\end{itemize}

\paragraph{Donde tocar para modificar.}
\begin{itemize}\setlength\itemsep{0pt}
	\item \textbf{Contar substrings distintas}: $\sum_{v \ne 0} (len[v] - len[link[v]])$.
	\item \textbf{Longest common substring}: navegar el SAM con el segundo string, manteniendo el match actual.
	\item \textbf{Substring count}: despues del build, ordenar estados por \texttt{len} descendente y propagar \texttt{occ}.
	\item \textbf{Endpos queries}: el SAM es el DFA canonico de los substrings; queries ``cuantas veces aparece $P

\vspace{0.5em}

\subsubsection{Lyndon factorization (Duval)}

\paragraph{Idea intuitiva.}
Un \textbf{string de Lyndon} es un string que es estrictamente menor que todos sus rotaciones. Toda cadena tiene una \textbf{factorizacion unica} de Lyndon (algoritmo de \textbf{Duval}): una secuencia de strings de Lyndon $L_1, L_2, \dots, L_k$ tales que $L_1 \ge L_2 \ge \dots \ge L_k$ y la concatenacion es el string original. Esto da $O(n)$ para encontrar los cortes. La demostracion de correctitud: cada factor es minimal en su clase de rotacion, asi que la concatenacion respeta el orden lexicografico.

\paragraph{Propiedades clave (invariantes).}
\begin{itemize}\setlength\itemsep{0pt}
	\item Cada factor de Lyndon es de la forma $uv$ donde $u$ se repite y $v$ es prefijo de $u$.
	\item La longitud del factor esta dada por la ``longitud del periodo''.
	\item Aplicacion: el ``Lyndon word'' es el menor en orden lexicografico entre todas las rotaciones.
	\item La factorizacion es \textbf{unica}.
\end{itemize}

\paragraph{Codigo clave.}
El algoritmo de Duval (iterativo):
\begin{verbatim}
int n = s.size();
int i = 0;
while (i < n) {
    int j = i + 1, k = i;
    while (j < n && s[k] <= s[j]) {
        if (s[k] < s[j]) k = i;
        else ++k;
        ++j;
    }
    int len = j - k;
    while (i < j) { cout << "factor: " << s.substr(i, len) << "\n"; i += len; }
}
\end{verbatim}

\paragraph{Complejidad.}
\begin{itemize}\setlength\itemsep{0pt}
	\item $O(n)$.
	\item Cada comparacion es $O(1)$; el puntero $k$ retrocede a lo sumo $n$ veces en total.
\end{itemize}

\paragraph{Donde tocar para modificar.}
\begin{itemize}\setlength\itemsep{0pt}
	\item \textbf{Encontrar la menor rotacion}: el menor sufijo lexicografico de \texttt{s + s} es la menor rotacion; la posicion se obtiene con un solo pase de Duval.
	\item \textbf{Algoritmo de Booth}: alternativa para rotacion minima, tambien $O(n)$.
	\item \textbf{Validar}: si necesitas verificar que un string es de Lyndon, compararlo con sus rotaciones.
	\item \textbf{Encontrar todos los factores}: guardar la posicion y longitud de cada uno.
\end{itemize}

\paragraph{Trampas y errores frecuentes.}
\begin{itemize}\setlength\itemsep{0pt}
	\item La condicion \texttt{s[j] <= s[k]} (no \texttt{<}) en el while de Duval es crucial; cambiar a \texttt{<} da comportamiento distinto.
	\item El caso $s[k] < s[j]$ resetea $k$ a $i$ (inicio del factor actual).
	\item El bucle externo debe avanzar $i$ por la longitud del factor, no por 1.
\end{itemize}

\paragraph{Codigo.}
Ver \texttt{cpp.tex}, seccion \textit{Lyndon factorization (Duval)}.

\subsection{Palindromos}

\subsubsection{Manacher}

\paragraph{Idea intuitiva.}
Calcula en $O(n)$ los \textbf{radios} de todos los palindromos centrados en cada posicion (impares y pares). La idea: mantener una ventana $[l, r]$ que es el palindromo centrado en un punto anterior mas largo encontrado. Si el nuevo centro $i$ cae dentro de $[l, r]$, entonces el palindromo centrado en $i$ tiene al menos el radio del palindromo ``espejo'' $l + r - i$; sino, empezar desde 1 (impar) o 0 (par). Despues, expandir lo que falte. La demostracion de $O(n)$: cada caracter se compara a lo sumo 2 veces (extension + limite por ventana).

\paragraph{Propiedades clave (invariantes).}
\begin{itemize}\setlength\itemsep{0pt}
	\item Conveccion del snippet: \texttt{odd[i]} = $k$ tal que el palindromo impar mas largo centrado en $i$ tiene longitud $2k - 1$. \texttt{even[i]} = $k$ para par centrado entre $i-1$ e $i$, longitud $2k$.
	\item $k$ cuenta tambien cuantos palindromos hay centrados ahi.
	\item Las queries \texttt{isPalindrome(l, r)} son $O(1)$.
	\item La ventana $[l, r]$ siempre contiene el palindromo mas largo encontrado hasta ahora.
\end{itemize}

\paragraph{Codigo clave.}
Manacher para palindromos impares:
\begin{verbatim}
int l = 0, r = -1;
for (int i = 0; i < n; ++i) {
    int k = (i > r) ? 1 : min(odd[l + r - i], r - i + 1);
    while (i - k >= 0 && i + k < n && s[i-k] == s[i+k]) ++k;
    odd[i] = k;
    if (i + k - 1 > r) { l = i - k + 1; r = i + k - 1; }
}
\end{verbatim}

\paragraph{Complejidad.}
\begin{itemize}\setlength\itemsep{0pt}
	\item Construccion $O(n)$, queries $O(1)$.
	\item Cada extension se cuenta una sola vez.
\end{itemize}

\paragraph{Donde tocar para modificar.}
\begin{itemize}\setlength\itemsep{0pt}
	\item \textbf{Solo contar palindromos}: ya implementado en \texttt{countPalindromes}.
	\item \textbf{Longest palindromic substring}: ya en \texttt{longestPalindrome}.
	\item \textbf{Modificar la conveccion de radios}: si queres ``radio = mitad de la longitud'' en vez de ``k palindromos'', ajustar el offset.
	\item \textbf{Queries con k no precomputado}: aniadir helpers especificos.
	\item \textbf{Modificar el caracter de relleno}: para multiples separadores, aniadir un caracter unico.
\end{itemize}

\paragraph{Trampas y errores frecuentes.}
\begin{itemize}\setlength\itemsep{0pt}
	\item La conveccion del snippet es \textbf{distinta} de la ``clasica de radio'' (donde radio = $(L-1)/2$); tener cuidado al copiar a otro codigo.
	\item Si los palindromos se solapan, la condicion \texttt{i + k - 1 > r} actualiza la ventana.
	\item En palindromos pares, el centro esta ``entre'' dos indices; aniadir offset en queries.
\end{itemize}

\paragraph{Codigo.}
Ver \texttt{cpp.tex}, seccion \textit{Manacher}.

% ============================================================
\section{Grafos}
% ============================================================

\subsection{Caminos minimos}

\subsubsection{Dijkstra}

\paragraph{Idea intuitiva.}
Encuentra el shortest path desde un origen $s$ a todos los vertices en $O((V + E) \log V)$ asumiendo \textbf{pesos no negativos}. La idea: usar una cola de prioridad (min-heap) ordenada por distancia tentativa. En cada paso, extraer el nodo con menor distancia; si ya fue procesado, saltar. Si al actualizar un vecino la distancia mejora, insertar/actualizar en la cola. La demostracion: cuando un nodo se extrae por primera vez, no hay forma de llegar a el con menor distancia (cualquier camino alternativo deberia pasar por nodos no procesados, con distancia $\ge$ la actual).

\paragraph{Propiedades clave (invariantes).}
\begin{itemize}\setlength\itemsep{0pt}
	\item Asume pesos $\ge 0$; con pesos negativos usar Bellman-Ford.
	\item Greedy correctness: la primera vez que se extrae un nodo, su distancia es la definitiva.
	\item Si se quiere reconstruir el camino, guardar el \texttt{parent[]}.
	\item Cada nodo se procesa a lo sumo $V$ veces (mas relajarions, pero solo $V$ definitivas).
\end{itemize}

\paragraph{Codigo clave.}
El bucle principal de Dijkstra:
\begin{verbatim}
priority_queue<pair<long long, int>, vector<...>, greater<...>> pq;
pq.push({0, source});
while (!pq.empty()) {
    auto [d, v] = pq.top(); pq.pop();
    if (d > dist[v]) continue;  // entrada vieja
    for (auto [w, u] : adj[v])
        if (dist[v] + w < dist[u]) {
            dist[u] = dist[v] + w;
            pq.push({dist[u], u});
        }
}
\end{verbatim}

\paragraph{Complejidad.}
\begin{itemize}\setlength\itemsep{0pt}
	\item $O((V + E) \log V)$ con heap binario, $O(V + E)$ con heap de Fibonacci.
	\item En la practica, $\sim 1.5 \cdot 10^{-7}$ segundos por operacion para $V = 10^5$.
\end{itemize}

\paragraph{Donde tocar para modificar.}
\begin{itemize}\setlength\itemsep{0pt}
	\item \textbf{Camino mas corto entre todos los pares}: si se necesita desde todos los origenes, usar Floyd o ejecutar Dijkstra $V$ veces.
	\item \textbf{Shortest path con aristas negativas limitadas}: Bellman-Ford.
	\item \textbf{Reconstruir camino}: aniadir \texttt{parent[v]} actualizado junto a \texttt{dist[v]}.
	\item \textbf{Dijkstra con potentials (Johnson)}: para aristas negativas si se puede reescalar.
	\item \textbf{Multi-source}: aniadir varios nodos fuentes iniciales al heap.
\end{itemize}

\paragraph{Trampas y errores frecuentes.}
\begin{itemize}\setlength\itemsep{0pt}
	\item Distancias no inicializadas a \texttt{INF}; el snippet usa \texttt{1e18}.
	\item El \texttt{if(dist[adjN] != 1e18)} para borrar de la cola es importante si usas \texttt{set} (no asi con \texttt{priority\_queue}).
	\item Con pesos grandes, \texttt{dist[v] + w} puede overflow; usar \texttt{LLONG\_MAX/2} como INF.
	\item No usar con pesos negativos: el algoritmo puede dar respuestas incorrectas.
\end{itemize}

\paragraph{Codigo.}
Ver \texttt{cpp.tex}, seccion \textit{Dijkstra}.

\vspace{0.5em}

\subsubsection{Bellman-Ford}

\paragraph{Idea intuitiva.}
Encuentra shortest path desde $s$ incluso con \textbf{pesos negativos} en $O(VE)$. La idea: relajar todas las aristas $V - 1$ veces. Cada relajacion propaga la mejor distancia a ``un paso mas''. Despues de $V - 1$ pasadas, si alguna arista aun se puede relajar, hay un \textbf{ciclo negativo} alcanzable. La demostracion: el camino mas corto tiene a lo sumo $V - 1$ aristas (sin ciclos); tras $V - 1$ relajaciones todas las distancias estan en su valor optimo.

\paragraph{Propiedades clave (invariantes).}
\begin{itemize}\setlength\itemsep{0pt}
	\item Funciona con pesos negativos (sin ciclos negativos).
	\item Tras $V - 1$ relajaciones, las distancias son definitivas (si no hay ciclo negativo).
	\item Una relajacion adicional detecta ciclos negativos: si \texttt{dist[u] + w < dist[v]}, hay ciclo negativo alcanzable desde $s$.
	\item Cada relajacion solo mejora distancias, asi que termina (no oscila).
\end{itemize}

\paragraph{Codigo clave.}
El doble bucle de relajacion:
\begin{verbatim}
for (int i = 0; i < n - 1; ++i)
    for (auto [u, v, w] : edges)
        if (dist[u] + w < dist[v])
            dist[v] = dist[u] + w;
// Detectar ciclo negativo
for (auto [u, v, w] : edges)
    if (dist[u] + w < dist[v]) { /* ciclo negativo */ }
\end{verbatim}

\paragraph{Complejidad.}
\begin{itemize}\setlength\itemsep{0pt}
	\item $O(VE)$.
	\item SPFA (con cola) en promedio es mucho mas rapido, peor caso igual.
\end{itemize}

\paragraph{Donde tocar para modificar.}
\begin{itemize}\setlength\itemsep{0pt}
	\item \textbf{SPFA}: optimizacion que relaja solo nodos cuya distancia cambio; en el peor caso sigue siendo $O(VE)$.
	\item \textbf{Detectar ciclo negativo global}: ejecutar Bellman-Ford desde un super-origen conectado a todos.
	\item \textbf{Camino mas corto con K aristas exactas}: $K$ iteraciones de Bellman-Ford.
	\item \textbf{Johnson's}: reescalar pesos con potentials para usar Dijkstra.
\end{itemize}

\paragraph{Trampas y errores frecuentes.}
\begin{itemize}\setlength\itemsep{0pt}
	\item Overflow si los pesos son $\ge 2^{52}$; usar \texttt{long long} y \texttt{\_\_int128} si hace falta.
	\item Si el grafo es denso ($E \approx V^2$), $O(VE)$ es prohibitivo.
	\item El chequeo de ciclo negativo debe hacerse tras las $V - 1$ iteraciones, no durante.
	\item Distancias no inicializadas a \texttt{INF}; el primer nodo (source) tiene distancia 0.
\end{itemize}

\paragraph{Codigo.}
Ver \texttt{cpp.tex}, seccion \textit{Bellman-Ford}.

\vspace{0.5em}

\subsubsection{Floyd-Warshall}

\paragraph{Idea intuitiva.}
Shortest path entre \textbf{todos los pares} de vertices en $O(V^3)$. DP: $d[i][j]$ = shortest path usando solo nodos intermedios en $\{0, \dots, k\}$. Transicion: $d_{k+1}[i][j] = \min(d_k[i][j], d_k[i][k+1] + d_k[k+1][j])$. La demostracion: tras $k$ iteraciones, $d[i][j]$ es el shortest path con intermedios en $\{0, \dots, k-1\}$; aniadir el nodo $k$ da la recurrencia clasica.

\paragraph{Propiedades clave (invariantes).}
\begin{itemize}\setlength\itemsep{0pt}
	\item Detecta ciclos negativos: tras el algoritmo, $d[i][i] < 0$ indica ciclo negativo alcanzable desde $i$.
	\item Funciona con pesos negativos (sin ciclos negativos).
	\item Memoria $O(V^2)$; con $V \le 400$ es OK, con mas conviene Dijkstra $V$ veces.
	\item Es la unica opcion cuando se quieren shortest paths entre todos los pares y el grafo es denso.
\end{itemize}

\paragraph{Codigo clave.}
El triple bucle clasico:
\begin{verbatim}
for (int k = 0; k < n; ++k)
    for (int i = 0; i < n; ++i)
        for (int j = 0; j < n; ++j)
            d[i][j] = min(d[i][j], d[i][k] + d[k][j]);
// detectar ciclo negativo
for (int i = 0; i < n; ++i) if (d[i][i] < 0) tiene_ciclo_negativo = true;
\end{verbatim}

\paragraph{Complejidad.}
\begin{itemize}\setlength\itemsep{0pt}
	\item $O(V^3)$ tiempo, $O(V^2)$ memoria.
	\item Para $V \le 400$: $\sim 6.4 \cdot 10^7$ operaciones (factible).
\end{itemize}

\paragraph{Donde tocar para modificar.}
\begin{itemize}\setlength\itemsep{0pt}
	\item \textbf{Reconstruir camino}: aniadir \texttt{next[i][j]} y backtrack.
	\item \textbf{Transitividad}: el algoritmo tambien calcula transitividad (\texttt{d[i][j] != INF} significa alcanzable).
	\item \textbf{Pesos doubles}: aniadir tolerancia (\texttt{d[i][j] > d[i][k] + d[k][j] - eps}).
	\item \textbf{Grafo no denso}: usar Dijkstra $V$ veces en vez de Floyd.
\end{itemize}

\paragraph{Trampas y errores frecuentes.}
\begin{itemize}\setlength\itemsep{0pt}
	\item Para detectar ciclo negativo correctamente, aniadir un chequeo \texttt{if (ans[i][i] < 0)} tras el triple bucle.
	\item El orden de los bucles \texttt{k, i, j} es importante; invertir da resultados diferentes.
	\item Con \texttt{double}, usar tolerancia \texttt{< eps} en vez de \texttt{<} para evitar loops infinitos.
\end{itemize}

\paragraph{Codigo.}
Ver \texttt{cpp.tex}, seccion \textit{Floyd-Warshall}.

\subsection{Componentes y cortes}

\subsubsection{Kosaraju SCC}

\paragraph{Idea intuitiva.}
Encuentra las \textbf{componentes fuertemente conexas} (SCCs) en $O(V + E)$. Algoritmo en 3 pasos:
\begin{itemize}\setlength\itemsep{0pt}
	\item DFS en $G$ rellenando una pila en \textbf{post-orden} (los nodos que terminan tarde quedan arriba).
	\item Construir $G^T$ (grafo transpuesto).
	\item DFS en $G^T$ procesando nodos en orden de la pila (de mayor a menor tiempo); cada DFS da una SCC.
\end{itemize}
La demostracion: dos nodos $u, v$ estan en la misma SCC sii hay camino $u \leadsto v$ y $v \leadsto u$; Kosaraju captura exactamente esto al explorar $G^T$ desde los nodos con mayor tiempo de finalizacion en $G$.

\paragraph{Propiedades clave (invariantes).}
\begin{itemize}\setlength\itemsep{0pt}
	\item Cada SCC es maximal fuertemente conexa.
	\item El \textbf{grafo de SCCs} (metagraf) es un DAG.
	\item Kosaraju y Tarjan dan el mismo resultado; Tarjan es ``mas rapido'' en practica (un solo DFS).
	\item Cada nodo pertenece a exactamente una SCC.
\end{itemize}

\paragraph{Codigo clave.}
La estructura de Kosaraju:
\begin{verbatim}
vector<int> order;
vector<bool> used(n);
function<void(int)> dfs1 = [&](int v) {
    used[v] = true;
    for (int u : g[v]) if (!used[u]) dfs1(u);
    order.push_back(v);  // post-orden
};
function<void(int, int)> dfs2 = [&](int v, int root) {
    comp[v] = root;
    for (int u : gt[v]) if (comp[u] == -1) dfs2(u, root);
};
// 1: DFS en G para llenar order
// 2: DFS en G^T en orden inverso para asignar SCCs
\end{verbatim}

\paragraph{Complejidad.}
\begin{itemize}\setlength\itemsep{0pt}
	\item $O(V + E)$.
	\item Constante 2 (dos DFS + el grafo transpuesto).
\end{itemize}

\paragraph{Donde tocar para modificar.}
\begin{itemize}\setlength\itemsep{0pt}
	\item \textbf{2-SAT}: SCC sobre grafo de implicaciones.
	\item \textbf{Construir el DAG de SCCs}: aniadir aristas entre SCCs diferentes.
	\item \textbf{Tamano de cada SCC}: aniadir contador durante el segundo DFS.
	\item \textbf{Grafo grande}: usar Tarjan para evitar construir $G^T$.
\end{itemize}

\paragraph{Trampas y errores frecuentes.}
\begin{itemize}\setlength\itemsep{0pt}
	\item El snippet usa \texttt{vector<...> adj[n]} (VLA); reemplazable por \texttt{vector<vector<pair<int,int>>} para portabilidad.
	\item El orden del segundo DFS debe ser \textbf{inverso} al de finalizacion del primero.
	\item Si el grafo no es conexo, aniadir un loop sobre todos los nodos en el primer DFS.
	\item Confundir SCC con componentes conexas (las primeras son para dirigidos, las segundas para no dirigidos).
\end{itemize}

\paragraph{Codigo.}
Ver \texttt{cpp.tex}, seccion \textit{Kosaraju SCC}.

\vspace{0.5em}

\subsubsection{Tarjan (puentes)}

\paragraph{Idea intuitiva.}
Encuentra los \textbf{puentes} (bridges): aristas cuya remocion \textbf{desconecta} el grafo. Usa un solo DFS manteniendo \texttt{tin[v]} (tiempo de descubrimiento) y \texttt{low[v]} (el menor \texttt{tin} alcanzable desde $v$ via back-edges). Una arista $v - u$ es puente si $low[u] > tin[v]$. La demostracion: si $low[u] > tin[v]$, entonces $u$ (y su subarbol) no pueden ``escapar'' de $v$; quitar la arista $v - u$ parte el grafo.

\paragraph{Propiedades clave (invariantes).}
\begin{itemize}\setlength\itemsep{0pt}
	\item \texttt{low[v] = min(tin[v], tin[w])} para todo back-edge $(v, w)$, y \texttt{min(low[u])} para hijos en el DFS.
	\item Solo se consideran las aristas del \textbf{arbol DFS} (no back-edges) para \texttt{low}.
	\item Cada nodo tiene exactamente un \texttt{tin}; cada back-edge se procesa una vez.
\end{itemize}

\paragraph{Codigo clave.}
El DFS de Tarjan para bridges:
\begin{verbatim}
void dfs(int v, int pEdge) {
    used[v] = true;
    tin[v] = low[v] = timer++;
    for (auto [u, id] : adj[v]) {
        if (id == pEdge) continue;  // no contar el edge al padre
        if (used[u]) {
            low[v] = min(low[v], tin[u]);  // back-edge
        } else {
            dfs(u, id);
            low[v] = min(low[v], low[u]);
            if (low[u] > tin[v]) bridges.push_back(id);  // es puente
        }
    }
}
\end{verbatim}

\paragraph{Complejidad.}
\begin{itemize}\setlength\itemsep{0pt}
	\item $O(V + E)$.
	\item Una sola pasada de DFS.
\end{itemize}

\paragraph{Donde tocar para modificar.}
\begin{itemize}\setlength\itemsep{0pt}
	\item \textbf{Bridge tree}: contractar cada 2-edge-connected component a un nodo; el resultado es un arbol.
	\item \textbf{Articulation points}: variante con condicion diferente (ver el siguiente modulo).
	\item \textbf{Grafo no dirigido}: solo funciona en no dirigidos; en dirigidos hay conceptos analogos (strong bridges, dominator tree).
	\item \textbf{Reconstruir componentes 2-edge-connected}: BFS desde cada nodo excluyendo bridges.
\end{itemize}

\paragraph{Trampas y errores frecuentes.}
\begin{itemize}\setlength\itemsep{0pt}
	\item No confundir el parent con un back-edge al padre; chequear \texttt{if (adjN == parent) continue;} y contar solo las veces que se ``visita'' como hijo.
	\item Para multigrafos, los edges deben tener IDs unicos para distinguirlos.
	\item Si el grafo no es conexo, aniadir loop externo sobre todos los nodos.
\end{itemize}

\paragraph{Codigo.}
Ver \texttt{cpp.tex}, seccion \textit{Tarjan (puentes)}.

\vspace{0.5em}

\subsubsection{Tarjan (puntos de articulacion)}

\paragraph{Idea intuitiva.}
Encuentra los \textbf{puntos de articulacion} (cut vertices): vertices cuya remocion \textbf{desconecta} el grafo. Mismo setup que para bridges: \texttt{tin[v]}, \texttt{low[v]}. Un nodo $v$ (no raiz) es de articulacion si existe un hijo $u$ con $low[u] \ge tin[v]$. La \textbf{raiz} es de articulacion si tiene $\ge 2$ hijos en el DFS tree. La razon: para $v$ no raiz, la condicion $\ge$ indica que el subarbol de $u$ depende de $v$; quitar $v$ los desconecta. Para la raiz, perder $\ge 2$ hijos parte el DFS en $\ge 2$ componentes.

\paragraph{Propiedades clave (invariantes).}
\begin{itemize}\setlength\itemsep{0pt}
	\item $low[u] \ge tin[v]$ significa que el subarbol de $u$ no tiene back-edge ``por encima'' de $v$.
	\item La raiz es caso especial: tiene que perder mas de un hijo para que el grafo se parta.
	\item Cada nodo se evalua una sola vez.
\end{itemize}

\paragraph{Codigo clave.}
La condicion de articulacion:
\begin{verbatim}
void dfs(int v, int p) {
    used[v] = true;
    tin[v] = low[v] = timer++;
    int children = 0;
    for (int u : adj[v]) {
        if (u == p) continue;
        if (used[u]) low[v] = min(low[v], tin[u]);
        else {
            dfs(u, v);
            low[v] = min(low[v], low[u]);
            if (low[u] >= tin[v] && p != -1)  // no raiz
                isArt[v] = true;
            ++children;
        }
    }
    if (p == -1 && children > 1) isArt[v] = true;  // raiz especial
}
\end{verbatim}

\paragraph{Complejidad.}
\begin{itemize}\setlength\itemsep{0pt}
	\item $O(V + E)$.
	\item Una sola pasada de DFS.
\end{itemize}

\paragraph{Donde tocar para modificar.}
\begin{itemize}\setlength\itemsep{0pt}
	\item \textbf{Block-cut tree}: estructura bipartita con BCCs (2-vertex-connected components) y articulation points.
	\item \textbf{Verificar biconectividad}: chequear que no hay puntos de articulacion.
	\item \textbf{Componentes biconnected}: extender Tarjan para enumerar BCCs.
\end{itemize}

\paragraph{Trampas y errores frecuentes.}
\begin{itemize}\setlength\itemsep{0pt}
	\item No olvidar el caso especial de la raiz.
	\item La condicion usa \texttt{>=}, no \texttt{>}; un hijo con \texttt{low[u] == tin[v]} tambien cuenta.
	\item Multigrafos pueden tener ``falsos'' articulation points; verificar caso a caso.
\end{itemize}

\paragraph{Codigo.}
Ver \texttt{cpp.tex}, seccion \textit{Tarjan (puntos de articulacion)}.

\vspace{0.5em}

\subsubsection{Tarjan SCC}

\paragraph{Idea intuitiva.}
Variante del SCC que usa \textbf{un solo DFS} con una pila explicita. Mantiene \texttt{disc[v]}, \texttt{low[v]}, y una pila de nodos ``en la SCC actual''. Cuando \texttt{low[v] == disc[v]}, $v$ es la raiz de una SCC; popear hasta $v$ inclusive. La demostracion: $low[v] == disc[v]$ indica que $v$ no tiene back-edges a ancestros; entonces $v$ y los nodos en la pila hasta $v$ forman una SCC maximal.

\paragraph{Propiedades clave (invariantes).}
\begin{itemize}\setlength\itemsep{0pt}
	\item Mas eficiente que Kosaraju (un solo DFS).
	\item \texttt{low[v] = min(low[u], disc[w])} para back-edges o recursion.
	\item La pila asegura que solo los nodos en la SCC actual se popean.
	\item Cada nodo se pushea y popea a lo sumo una vez.
\end{itemize}

\paragraph{Codigo clave.}
El DFS de Tarjan SCC:
\begin{verbatim}
void dfs(int v) {
    disc[v] = low[v] = timer++;
    st.push(v); inStack[v] = true;
    for (int u : g[v]) {
        if (disc[u] == -1) { dfs(u); low[v] = min(low[v], low[u]); }
        else if (inStack[u]) low[v] = min(low[v], disc[u]);
    }
    if (low[v] == disc[v]) {
        // popear hasta v
        while (true) {
            int u = st.top(); st.pop(); inStack[u] = false;
            comp[u] = current_scc;
            if (u == v) break;
        }
        ++current_scc;
    }
}
\end{verbatim}

\paragraph{Complejidad.}
\begin{itemize}\setlength\itemsep{0pt}
	\item $O(V + E)$.
	\item Constante similar a Kosaraju pero sin construir $G^T$.
\end{itemize}

\paragraph{Donde tocar para modificar.}
\begin{itemize}\setlength\itemsep{0pt}
	\item \textbf{Tamano de cada SCC}: aniadir contador durante el pop.
	\item \textbf{Grafo condensation}: aniadir aristas entre SCCs en el DAG resultante.
	\item \textbf{2-SAT}: aniadir la construccion del grafo de implicaciones.
\end{itemize}

\paragraph{Trampas y errores frecuentes.}
\begin{itemize}\setlength\itemsep{0pt}
	\item \texttt{inStack[u]} debe chequearse antes de usar \texttt{low[v] = min(low[v], disc[u])} en back-edges.
	\item La pila global es importante; sin ella no se puede extraer la SCC.
	\item Si el grafo es muy grande, recursion puede overflow; aniadir version iterativa.
\end{itemize}

\paragraph{Codigo.}
Ver \texttt{cpp.tex}, seccion \textit{Tarjan SCC}.

\subsection{Matching}

\subsubsection{Kuhn (bipartite matching)}

\paragraph{Idea intuitiva.}
Encuentra un \textbf{matching maximo} en un grafo bipartito en $O(VE)$. Para cada nodo $u$ en $L$, intentar encontrar un vecino $v$ en $R$ que este libre o cuyo match se pueda ``reorganizar'' recursivamente. Si se encuentra, matchear $u - v$. La idea: si $u$ quiere un vecino $v$ ya ocupado, intentamos ``mover'' el match de $v$ a otro nodo; recursivamente, esto encuentra un augmenting path si existe.

\paragraph{Propiedades clave (invariantes).}
\begin{itemize}\setlength\itemsep{0pt}
	\item \texttt{matchR[v]} = el nodo de $L$ matcheado a $v$ (o $-1$ si libre).
	\item \texttt{vis[]} se resetea \textbf{cada intento} de $u$.
	\item Si el matching es perfecto, $|L| = |R| = $ tamano del matching.
	\item Kuhn no garantiza maximalidad greedy; puede haber augmenting paths no encontrados.
\end{itemize}

\paragraph{Codigo clave.}
La recursion del augmenting path:
\begin{verbatim}
bool tryKuhn(int v) {
    if (vis[v]) return false;
    vis[v] = true;
    for (int u : g[v]) {
        if (matchR[u] == -1 || tryKuhn(matchR[u])) {
            matchR[u] = v;
            return true;
        }
    }
    return false;
}
// Por cada v en L:
fill(vis, vis+n, false);
if (tryKuhn(v)) ++matches;
\end{verbatim}

\paragraph{Complejidad.}
\begin{itemize}\setlength\itemsep{0pt}
	\item $O(VE)$ en el peor caso.
	\item Con Hopcroft-Karp: $O(\sqrt{V} \cdot E)$, mucho mejor para grafos densos.
\end{itemize}

\paragraph{Donde tocar para modificar.}
\begin{itemize}\setlength\itemsep{0pt}
	\item \textbf{Matching bipartito minimo (min vertex cover)}: Konig's theorem.
	\item \textbf{Matching general}: blossom algorithm (no incluido).
	\item \textbf{Aniadir capacidades}: si los nodos de $R$ tienen capacidad $> 1$, partir cada uno en capacidad copias.
	\item \textbf{Output del matching}: el array \texttt{matchR} da los pares finales.
\end{itemize}

\paragraph{Trampas y errores frecuentes.}
\begin{itemize}\setlength\itemsep{0pt}
	\item Olvidar resetear \texttt{vis} por intento; eso es lo que evita ``ciclos infinitos'' en la recursion.
	\item Para grafos grandes ($V > 10^4$), Kuhn puede ser muy lento; usar Hopcroft-Karp.
	\item Si el matching es maximo pero no perfecto, $|L|$ o $|R|$ puede ser mayor; el matching es solo maximo.
\end{itemize}

\paragraph{Codigo.}
Ver \texttt{cpp.tex}, seccion \textit{Kuhn (bipartite matching)}.

\vspace{0.5em}

\subsubsection{Hopcroft-Karp}

\paragraph{Idea intuitiva.}
Matching bipartito maximo en $O(\sqrt{V} E)$. Algoritmo en \textbf{fases}: en cada fase, hacer BFS para encontrar los \textbf{caminos de aumento mas cortos} desde nodos libres de $L$; luego DFS solo por esos caminos. Cada fase aumenta el matching en $\ge 1$ y toma $O(E)$. La demostracion: tras una fase, todos los augmenting paths tienen longitud $\ge$ la nueva minima; eso sucede cada $\sqrt{V}$ fases (longitud minima crece geometricamente).

\paragraph{Propiedades clave (invariantes).}
\begin{itemize}\setlength\itemsep{0pt}
	\item \texttt{dist[u]} = distancia en niveles desde un nodo libre de $L$.
	\item BFS termina cuando no hay mas camino a un nodo libre de $R$.
	\item DFS respeta los niveles (solo sigue \texttt{dist[matchV[v]] == dist[u] + 1}).
	\item Cada fase aumenta el matching en $\ge 1$ arista.
\end{itemize}

\paragraph{Codigo clave.}
El BFS de Hopcroft-Karp:
\begin{verbatim}
bool bfs() {
    queue<int> q;
    for (int v : L) {
        if (matchU[v] == -1) { dist[v] = 0; q.push(v); }
        else dist[v] = -1;
    }
    bool found = false;
    while (!q.empty()) {
        int v = q.front(); q.pop();
        for (int u : g[v]) {
            if (matchV[u] != -1 && dist[matchV[u]] == -1) {
                dist[matchV[u]] = dist[v] + 1;
                q.push(matchV[u]);
            }
            if (matchV[u] == -1) found = true;  // llegamos a libre
        }
    }
    return found;
}
\end{verbatim}

\paragraph{Complejidad.}
\begin{itemize}\setlength\itemsep{0pt}
	\item $O(\sqrt{V} E)$.
	\item Para $V = 10^4$, $E = 10^5$: $\sim 3 \cdot 10^7$ operaciones.
\end{itemize}

\paragraph{Donde tocar para modificar.}
\begin{itemize}\setlength\itemsep{0pt}
	\item \textbf{Reconstruir el matching}: ya viene en \texttt{matchU}/\texttt{matchV}.
	\item \textbf{Aniadir pesos}: Hungarian algorithm.
	\item \textbf{Densidad alta}: si $E \approx V^2$, sigue siendo util pero verificar constantes.
	\item \textbf{Aniadir capacidades}: partir los nodos con capacidad $> 1$ en varias copias.
\end{itemize}

\paragraph{Trampas y errores frecuentes.}
\begin{itemize}\setlength\itemsep{0pt}
	\item \texttt{dist[matchV[v]] == dist[u] + 1} (no \texttt{==} a secas); sin esto, el DFS no respeta los niveles y se pierde la complejidad.
	\item El BFS debe marcar los niveles correctamente; resetear \texttt{dist} para nodos alcanzables solo via matching.
	\item El DFS debe respetar los niveles; cualquier ``atajo'' invalida la cota.
\end{itemize}

\paragraph{Codigo.}
Ver \texttt{cpp.tex}, seccion \textit{Hopcroft-Karp}.

\subsection{Basicos}

\subsubsection{DFS / BFS / componentes}

\paragraph{Idea intuitiva.}
\textbf{DFS} (Depth-First Search) y \textbf{BFS} (Breadth-First Search) son los dos recorridos basicos de grafos. DFS usa una pila (o recursion) y explora tan profundo como puede antes de retroceder; BFS usa una cola y explora por ``niveles'' de distancia. La cantidad de \textbf{componentes conexas} se cuenta iniciando un BFS/DFS desde cada nodo no visitado. La diferencia clave: BFS da el camino mas corto en aristas; DFS da un orden de finalizacion util para problemas topologicos.

\paragraph{Propiedades clave (invariantes).}
\begin{itemize}\setlength\itemsep{0pt}
	\item BFS da la \textbf{distancia en aristas} mas corta desde el origen.
	\item DFS da un \textbf{orden topologico} (si se aniaden al final).
	\item Ambos son $O(V + E)$.
	\item Un grafo con $C$ componentes conexas satisface $\sum \text{subtree\_size}_v = V$ y $\sum E_v = 2E$.
	\item DFS en arbol da un spanning tree; las ``back edges'' corresponden a ciclos.
\end{itemize}

\paragraph{Codigo clave.}
BFS clasico:
\begin{verbatim}
queue<int> q;
q.push(source);
dist[source] = 0;
while (!q.empty()) {
    int v = q.front(); q.pop();
    for (int u : adj[v])
        if (dist[u] == -1) {
            dist[u] = dist[v] + 1;
            q.push(u);
        }
}
\end{verbatim}

\paragraph{Complejidad.}
\begin{itemize}\setlength\itemsep{0pt}
	\item $O(V + E)$ tiempo, $O(V)$ memoria.
	\item Cada arista se procesa exactamente 2 veces (ida y vuelta) en no dirigidos.
\end{itemize}

\paragraph{Donde tocar para modificar.}
\begin{itemize}\setlength\itemsep{0pt}
	\item \textbf{Aniadir pesos}: si las aristas tienen peso y queres shortest path, BFS ya no sirve; usar Dijkstra.
	\item \textbf{Grafo dirigido}: en DFS, aniadir check \texttt{if (state[u] == 1)} para detectar ciclos.
	\item \textbf{Bipartite check}: BFS alternando colores (ver el siguiente modulo).
	\item \textbf{Componentes fuertemente conexas (dirigidos)}: usar Tarjan o Kosaraju.
	\item \textbf{BFS en matriz}: aniadir offsets por direccion para grid 2D.
\end{itemize}

\paragraph{Trampas y errores frecuentes.}
\begin{itemize}\setlength\itemsep{0pt}
	\item Stack overflow con DFS recursivo para $V > 10^5$; usar version iterativa con pila explicita.
	\item En BFS, olvidar marcar \texttt{dist[u] = dist[v] + 1} antes de pushear da visitas duplicadas.
	\item Para grafos dirigidos, ``componente conexa'' no esta definida; usar SCC.
\end{itemize}

\paragraph{Codigo.}
Ver \texttt{cpp.tex}, seccion \textit{DFS / BFS / componentes}.

\vspace{0.5em}

\subsubsection{Topological sort (Kahn + DFS)}

\paragraph{Idea intuitiva.}
Un \textbf{orden topologico} de un DAG es una permutacion de los vertices tal que toda arista $u \to v$ tiene $u$ antes que $v$. Dos algoritmos clasicos:
\begin{itemize}\setlength\itemsep{0pt}
	\item \textbf{Kahn}: BFS sobre nodos con \texttt{indeg == 0}; al ``consumir'' un nodo, decrementar el indeg de sus vecinos; los que lleguen a 0 se encolan.
	\item \textbf{DFS}: hacer DFS; al terminar un nodo (estado = 2), aniadirlo al final; el orden es el reverso.
\end{itemize}
La demostracion: en un DAG siempre hay al menos un nodo con indeg 0; quitarlo mantiene el DAG; por induccion, se obtiene el orden.

\paragraph{Propiedades clave (invariantes).}
\begin{itemize}\setlength\itemsep{0pt}
	\item Existe sii el grafo es un DAG.
	\item Si al final \texttt{order.size() < V}, hay un ciclo.
	\item Kahn da el orden ``natural'' (los que pueden ir primero); DFS da el orden ``al reves'' (los que pueden ir al final primero).
	\item La cantidad de ordenes topologicos puede ser exponencial.
\end{itemize}

\paragraph{Codigo clave.}
Kahn con priority queue (orden lexicografico):
\begin{verbatim}
priority_queue<int, vector<int>, greater<int>> pq;
for (int v = 0; v < n; ++v) if (indeg[v] == 0) pq.push(v);
while (!pq.empty()) {
    int v = pq.top(); pq.pop();
    order.push_back(v);
    for (int u : adj[v]) {
        if (--indeg[u] == 0) pq.push(u);
    }
}
\end{verbatim}

\paragraph{Complejidad.}
\begin{itemize}\setlength\itemsep{0pt}
	\item $O(V + E)$ ambos.
	\item Con priority queue, $O((V + E) \log V)$.
\end{itemize}

\paragraph{Donde tocar para modificar.}
\begin{itemize}\setlength\itemsep{0pt}
	\item \textbf{Detectar ciclos}: comparar \texttt{order.size()} con \texttt{V}.
	\item \textbf{Orden lexicografico minimo}: usar priority queue en Kahn.
	\item \textbf{Topological sort con pesos}: combinar con shortest path en DAG (un solo pase en orden topologico).
	\item \textbf{Encontrar todos los ordenes}: DP sobre subsets (solo para $V$ pequeno).
\end{itemize}

\paragraph{Trampas y errores frecuentes.}
\begin{itemize}\setlength\itemsep{0pt}
	\item En el snippet, Kahn retorna orden vacio si hay ciclo (util para detectar).
	\item En DFS, no olvidar \texttt{reverse(order)}.
	\item Confundir indeg con outdeg; el orden Kahn quita primero los nodos ``fuente''.
	\item Para grafos no DAG, el algoritmo termina con \texttt{order.size() < V}.
\end{itemize}

\paragraph{Codigo.}
Ver \texttt{cpp.tex}, seccion \textit{Topological sort (Kahn + DFS)}.

\vspace{0.5em}

\subsubsection{Bipartite check + 2-coloring}

\paragraph{Idea intuitiva.}
Un grafo es \textbf{bipartito} si sus vertices se pueden dividir en dos conjuntos $L$ y $R$ tales que toda arista va de $L$ a $R$. Esto equivale a decir que se puede \textbf{2-colorear} sin que dos vertices adyacentes tengan el mismo color. BFS alternando colores detecta esto: si en algun momento dos adyacentes tienen el mismo color, no es bipartito. La demostracion: por cada componente, fijar el color de un nodo determina todos los demas; si en algun paso hay conflicto, hay un ciclo impar (no bipartito).

\paragraph{Propiedades clave (invariantes).}
\begin{itemize}\setlength\itemsep{0pt}
	\item Grafo bipartito $\Leftrightarrow$ sin ciclos impares.
	\item Conexo: si un componente no es bipartito, todo el grafo no lo es.
	\item Si hay \textbf{multiple componentes}, cada uno se colorea independientemente.
	\item La 2-coloracion es unica salvo swap de colores.
\end{itemize}

\paragraph{Codigo clave.}
BFS con 2-coloreo:
\begin{verbatim}
vector<int> color(n, -1);
for (int v = 0; v < n; ++v) if (color[v] == -1) {
    color[v] = 0;
    queue<int> q; q.push(v);
    while (!q.empty()) {
        int u = q.front(); q.pop();
        for (int w : adj[u]) {
            if (color[w] == -1) { color[w] = 1 - color[u]; q.push(w); }
            else if (color[w] == color[u]) return false;
        }
    }
}
return true;
\end{verbatim}

\paragraph{Complejidad.}
\begin{itemize}\setlength\itemsep{0pt}
	\item $O(V + E)$.
	\item Memoria $O(V)$.
\end{itemize}

\paragraph{Donde tocar para modificar.}
\begin{itemize}\setlength\itemsep{0pt}
	\item \textbf{Colorear desde multiples fuentes}: BFS desde cada nodo no coloreado.
	\item \textbf{Bipartito + matching}: Hopcroft-Karp.
	\item \textbf{Grafo bipartito ponderado}: Hungarian algorithm.
	\item \textbf{Verificar bipartito en subgrafo}: solo aplicar BFS a los nodos relevantes.
\end{itemize}

\paragraph{Trampas y errores frecuentes.}
\begin{itemize}\setlength\itemsep{0pt}
	\item Verificar \textbf{todos} los componentes, no solo el primero.
	\item Para grafos con auto-loops, un auto-loop hace al grafo no bipartito.
	\item En multigrafos, las aristas multiples no afectan el chequeo de bipartito.
\end{itemize}

\paragraph{Codigo.}
Ver \texttt{cpp.tex}, seccion \textit{Bipartite check + 2-coloring}.

\vspace{0.5em}

\subsubsection{0-1 BFS}

\paragraph{Idea intuitiva.}
Variante de BFS para grafos con \textbf{pesos 0 o 1}. Usa un \texttt{deque}: aristas de peso 0 se encolan al frente (\texttt{push\_front}), aristas de peso 1 al final (\texttt{push\_back}). Asi, la primera vez que se extrae un nodo, su distancia es la definitiva, igual que en Dijkstra pero en $O(V + E)$. La intuicion: el deque mantiene los nodos ordenados por distancia; las aristas de peso 0 no cambian la distancia (entran al frente), las de peso 1 la incrementan (entran al final).

\paragraph{Propiedades clave (invariantes).}
\begin{itemize}\setlength\itemsep{0pt}
	\item Asume pesos \textbf{solo} 0 o 1.
	\item Mejor caso que Dijkstra (sin factor $\log V$).
	\item No funciona con pesos $\ge 2$.
	\item La cola mantiene los nodos ordenados por distancia no decreciente.
\end{itemize}

\paragraph{Codigo clave.}
El deque doble:
\begin{verbatim}
deque<int> q;
q.push_front(source);
dist[source] = 0;
while (!q.empty()) {
    int v = q.front(); q.pop_front();
    for (auto [w, u] : adj[v]) {
        if (dist[v] + w < dist[u]) {
            dist[u] = dist[v] + w;
            if (w == 0) q.push_front(u);
            else q.push_back(u);
        }
    }
}
\end{verbatim}

\paragraph{Complejidad.}
\begin{itemize}\setlength\itemsep{0pt}
	\item $O(V + E)$.
	\item Cada nodo se inserta a lo sumo $O(1)$ veces (con check de distancia).
\end{itemize}

\paragraph{Donde tocar para modificar.}
\begin{itemize}\setlength\itemsep{0pt}
	\item \textbf{Pesos en $[0, k]$ small}: usar small k-ary BFS o Dijkstra.
	\item \textbf{Reconstruir camino}: aniadir \texttt{parent[]}.
	\item \textbf{Pesos negativos}: no funciona, usar Bellman-Ford.
\end{itemize}

\paragraph{Trampas y errores frecuentes.}
\begin{itemize}\setlength\itemsep{0pt}
	\item Si los pesos son otro rango, el algoritmo da respuestas incorrectas. Verificar que las aristas son solo 0/1.
	\item El chequeo \texttt{dist[v] + w < dist[u]} es crucial; sino, el algoritmo procesa nodos multiples veces.
	\item Si una arista tiene peso $\ge 2$, aniadir verificaciones o cambiar a Dijkstra.
\end{itemize}

\paragraph{Codigo.}
Ver \texttt{cpp.tex}, seccion \textit{0-1 BFS}.

\subsection{MST}

\subsubsection{Kruskal}

\paragraph{Idea intuitiva.}
Encuentra el \textbf{arbol de expansion minima} (MST) en $O(E \log E)$. Ordena las aristas por peso ascendente y las aniade al MST si no forman ciclo (chequear con DSU). El arbol tiene exactamente $V - 1$ aristas. La demostracion (``Cut Property''): para cualquier corte, la arista minima cruzando el corte pertenece a \emph{algun} MST; por lo tanto, elegir la minima disponible que no forma ciclo es optimo.

\paragraph{Propiedades clave (invariantes).}
\begin{itemize}\setlength\itemsep{0pt}
	\item Algoritmo \textbf{greedy} correcto (Cut Property).
	\item Funciona en grafos \textbf{no conexos}; produce un \textbf{minimum spanning forest}.
	\item Si las aristas tienen pesos unicos, el MST es unico.
	\item El MST tiene exactamente $V - 1$ aristas (o menos si el grafo es disconexo).
\end{itemize}

\paragraph{Codigo clave.}
Kruskal con DSU:
\begin{verbatim}
sort(edges.begin(), edges.end());
DSU dsu(n);
long long mst = 0;
int edgesUsed = 0;
for (auto [w, u, v] : edges)
    if (dsu.unite(u, v)) {
        mst += w;
        if (++edgesUsed == n - 1) break;
    }
\end{verbatim}

\paragraph{Complejidad.}
\begin{itemize}\setlength\itemsep{0pt}
	\item $O(E \log E)$ (ordenamiento) o $O(E \log V)$ con DSU optimizations.
	\item DSU es casi $O(1)$ amortizado, asi que el sort domina.
\end{itemize}

\paragraph{Donde tocar para modificar.}
\begin{itemize}\setlength\itemsep{0pt}
	\item \textbf{Second best MST}: enumerar aristas del MST, sacar una, aniadir una no-MST; recalcular.
	\item \textbf{Kruskal randomizado}: si el grafo es casi-completo, agregar shuffle.
	\item \textbf{Aniadir info al DSU}: guardar el peso maximo en el camino, etc.
	\item \textbf{DSU con rollback}: si necesitas Kruskal reversible (para queries offline).
\end{itemize}

\paragraph{Trampas y errores frecuentes.}
\begin{itemize}\setlength\itemsep{0pt}
	\item El DSU debe estar limpio para cada invocacion.
	\item Si las aristas tienen pesos negativos, Kruskal sigue funcionando.
	\item Si el grafo es disconexo, el ``MST'' es un bosque con varias componentes.
	\item El check \texttt{if (++edgesUsed == n - 1) break;} optimiza la salida temprana.
\end{itemize}

\paragraph{Codigo.}
Ver \texttt{cpp.tex}, seccion \textit{Kruskal}.

\vspace{0.5em}

\subsubsection{Prim (priority queue)}

\paragraph{Idea intuitiva.}
Otra forma de calcular el MST: empezar desde un nodo y, en cada paso, aniadir la arista de menor peso que conecta el conjunto construido con el resto. Usando priority queue, es $O((V + E) \log V)$. La demostracion: cada vez que se aniade una arista, es la minima cruzando el ``corte'' entre el conjunto actual y el resto (Cut Property); eso preserva optimalidad.

\paragraph{Propiedades clave (invariantes).}
\begin{itemize}\setlength\itemsep{0pt}
	\item Similar a Dijkstra pero sin acumular distancias, solo eligiendo minimo por arista.
	\item Funciona mejor que Kruskal en \textbf{grafos densos} ($E \approx V^2$).
	\item Si el grafo es disconexo, genera el bosque (solo llega a los alcanzables).
	\item Cada nodo se inserta al heap una vez; cada arista se procesa una vez.
\end{itemize}

\paragraph{Codigo clave.}
Prim clasico con heap:
\begin{verbatim}
priority_queue<pair<int, int>, vector<...>, greater<...>> pq;
vector<bool> used(n, false);
vector<int> minEdge(n, INF);
minEdge[0] = 0;
pq.push({0, 0});
while (!pq.empty()) {
    auto [w, v] = pq.top(); pq.pop();
    if (used[v]) continue;
    used[v] = true;
    mst += w;
    for (auto [peso, u] : adj[v])
        if (!used[u] && peso < minEdge[u]) {
            minEdge[u] = peso;
            pq.push({peso, u});
        }
}
\end{verbatim}

\paragraph{Complejidad.}
\begin{itemize}\setlength\itemsep{0pt}
	\item $O((V + E) \log V)$.
	\item Con matriz de adyacencia: $O(V^2)$ sin heap.
\end{itemize}

\paragraph{Donde tocar para modificar.}
\begin{itemize}\setlength\itemsep{0pt}
	\item \textbf{Prim con matriz de adyacencia}: $O(V^2)$ sin heap, util en grafos densos pequenos.
	\item \textbf{Reconstruir el MST}: aniadir \texttt{parent[]}.
	\item \textbf{Prim en streaming}: variante para grafos que llegan incrementalmente.
\end{itemize}

\paragraph{Trampas y errores frecuentes.}
\begin{itemize}\setlength\itemsep{0pt}
	\item El \texttt{if (used[v]) continue;} es crucial: si el nodo ya fue procesado, saltar.
	\item En la version matriz, elije el minimo por linea (no heap).
	\item Si el grafo es disconexo, aniadir un loop externo sobre todos los nodos no usados.
\end{itemize}

\paragraph{Codigo.}
Ver \texttt{cpp.tex}, seccion \textit{Prim (priority queue)}.

\subsection{Flujos}

\subsubsection{Edmonds-Karp (max flow)}

\paragraph{Idea intuitiva.}
Implementacion BFS del algoritmo de \textbf{Ford-Fulkerson}. En cada iteracion, hacer BFS desde $s$ hasta $t$ siguiendo aristas con capacidad residual $> 0$; si se llega a $t$, se encontro un \textbf{camino de aumento}. Enviar el minimo de las capacidades en el camino; actualizar capacidades residuales. La demostracion de $O(VE^2)$: cada BFS encuentra el augmenting path mas corto, y la longitud minima crece monotona; como hay $\le E$ longitudes posibles, hay $\le E$ fases.

\paragraph{Propiedades clave (invariantes).}
\begin{itemize}\setlength\itemsep{0pt}
	\item BFS garantiza caminos de aumento mas cortos.
	\item Capacidad residual: forward edge baja, backward edge sube.
	\item Complejidad $O(VE^2)$.
	\item El flujo maximo respeta la conservacion en nodos (in = out salvo $s, t$).
\end{itemize}

\paragraph{Codigo clave.}
BFS + augmenting path:
\begin{verbatim}
while (bfs(s, t, parent)) {  // bfs encuentra camino o retorna false
    int pathFlow = INF;
    for (int v = t; v != s; v = parent[v])
        pathFlow = min(pathFlow, capacity[parent[v]][v]);
    for (int v = t; v != s; v = parent[v]) {
        capacity[parent[v]][v] -= pathFlow;
        capacity[v][parent[v]] += pathFlow;
    }
    maxFlow += pathFlow;
}
\end{verbatim}

\paragraph{Complejidad.}
\begin{itemize}\setlength\itemsep{0pt}
	\item $O(VE^2)$.
	\item En la practica, mas lento que Dinic para grafos grandes.
\end{itemize}

\paragraph{Donde tocar para modificar.}
\begin{itemize}\setlength\itemsep{0pt}
	\item \textbf{Grafo bipartito}: $O(\sqrt{V} E)$ (equivalente a Hopcroft-Karp).
	\item \textbf{Dinic} es preferible en la mayoria de los casos.
	\item \textbf{Capacidades doubles}: usar tolerancia para comparacion.
	\item \textbf{Min-cost flow}: aniadir costo por arista y modificar el algoritmo.
\end{itemize}

\paragraph{Trampas y errores frecuentes.}
\begin{itemize}\setlength\itemsep{0pt}
	\item Las capacidades deben ser $> 0$ para entrar a la cola (no $>= 0$); sino, ciclos infinitos.
	\item Para grafos bipartitos, Dinic o Hopcroft-Karp son mejores.
	\item Si el flujo maximo es muy grande, los \texttt{int} pueden overflow; usar \texttt{long long}.
\end{itemize}

\paragraph{Codigo.}
Ver \texttt{cpp.tex}, seccion \textit{Edmonds-Karp (max flow)}.

\vspace{0.5em}

\subsubsection{Dinic (max flow)}

\paragraph{Idea intuitiva.}
Algoritmo de flujo mas eficiente que Edmonds-Karp: combina \textbf{BFS para level graph} con \textbf{DFS bloqueante} (multiple augmenting paths en una sola pasada). Cada fase cuesta $O(E)$ y hay $O(V)$ fases en el peor caso; para bipartitos, $O(\sqrt{V} E)$. La razon: cada fase ``satura'' al menos una arista, asi que en $O(E)$ fases el flujo maximo se alcanza; con $O(V)$ niveles en el BFS, hay $O(VE)$ operaciones totales.

\paragraph{Propiedades clave (invariantes).}
\begin{itemize}\setlength\itemsep{0pt}
	\item \texttt{level[v]} = distancia desde $s$ en el level graph (solo aristas con cap $> 0$).
	\item \texttt{it[v]} = iterador actual para evitar reprocesar aristas.
	\item DFS solo sigue aristas con \texttt{level[e.to] == level[v] + 1}.
	\item Back-edges tienen \texttt{rev} apuntando a la forward edge original.
\end{itemize}

\paragraph{Codigo clave.}
El DFS bloqueante de Dinic:
\begin{verbatim}
long long dfs(int v, int t, long long f) {
    if (v == t) return f;
    for (int& i = it[v]; i < adj[v].size(); ++i) {
        Edge& e = adj[v][i];
        if (e.cap > 0 && level[e.to] == level[v] + 1) {
            long long ret = dfs(e.to, t, min(f, e.cap));
            if (ret > 0) {
                e.cap -= ret;
                adj[e.to][e.rev].cap += ret;
                return ret;
            }
        }
    }
    return 0;
}
\end{verbatim}

\paragraph{Complejidad.}
\begin{itemize}\setlength\itemsep{0pt}
	\item $O(V^2 E)$ general, $O(\sqrt{V} E)$ para bipartitos.
	\item En la practica, mucho mas rapido que Edmonds-Karp.
\end{itemize}

\paragraph{Donde tocar para modificar.}
\begin{itemize}\setlength\itemsep{0pt}
	\item \textbf{Dinic con scaling}: aniadir factor de escala para capacidades grandes.
	\item \textbf{Lower bounds}: aniadir un super-source/sink para forzar flujo minimo.
	\item \textbf{Matching bipartito}: el grafo se construye como source -> L -> R -> sink, todos con capacidad 1.
	\item \textbf{Push-relabel}: alternativa para grafos muy densos.
\end{itemize}

\paragraph{Trampas y errores frecuentes.}
\begin{itemize}\setlength\itemsep{0pt}
	\item \texttt{it[v]} debe \textbf{incrementarse} en cada iteracion; sino, el DFS se queda atascado en la misma arista.
	\item Si el BFS no encuentra $t$, el flujo es maximo; terminar.
	\item Las back-edges son cruciales para ``cancelar'' elecciones previas.
\end{itemize}

\paragraph{Codigo.}
Ver \texttt{cpp.tex}, seccion \textit{Dinic (max flow)}.

\vspace{0.5em}

\subsubsection{Min-Cost Max-Flow}

\paragraph{Idea intuitiva.}
Encuentra el flujo maximo con \textbf{costo minimo} (asumiendo costos no negativos). Usa \textbf{SPFA} (o Dijkstra con potentials) para encontrar el shortest path en el \textbf{costo} desde $s$ hasta $t$ considerando capacidades residuales. Cada augmenting path aporta su costo $\times$ flujo. La idea: los potentials reescalan las aristas para que todas tengan costo no negativo, permitiendo usar Dijkstra.

\paragraph{Propiedades clave (invariantes).}
\begin{itemize}\setlength\itemsep{0pt}
	\item \texttt{pot[]} = potentials para evitar aristas negativas (Johnson's trick).
	\item \texttt{dist[v]} = costo minimo desde $s$ con potentials.
	\item Solo se aumentan aristas con \texttt{cap > 0} (forward) y costo es \texttt{+cost} (forward) o \texttt{-cost} (backward).
	\item Tras cada augmenting, los potentials se actualizan para mantener la consistencia.
\end{itemize}

\paragraph{Codigo clave.}
El ciclo de min-cost flow:
\begin{verbatim}
while (spfa(s, t, dist, parent)) {
    int pathFlow = INF;
    for (int v = t; v != s; v = parent[v])
        pathFlow = min(pathFlow, capacity[parent[v]][v]);
    for (int v = t; v != s; v = parent[v]) {
        capacity[parent[v]][v] -= pathFlow;
        capacity[v][parent[v]] += pathFlow;
        cost += pathFlow * edge_cost[parent[v]][v];
    }
}
\end{verbatim}

\paragraph{Complejidad.}
\begin{itemize}\setlength\itemsep{0pt}
	\item $O(F \cdot V E)$ con SPFA, $O(F \cdot E \log V)$ con Dijkstra + potentials.
	\item $F$ = flujo maximo total.
\end{itemize}

\paragraph{Donde tocar para modificar.}
\begin{itemize}\setlength\itemsep{0pt}
	\item \textbf{Dijkstra + potentials}: si los costos son enteros no negativos, esto es mas estable.
	\item \textbf{Costos negativos en aristas}: requieren Bellman-Ford inicial o reorden.
	\item \textbf{Flujo minimo con costo}: encontrar el minimo flujo posible y su costo.
	\item \textbf{Costo por unidad}: si los costos son muy grandes, aniadir verificacion.
\end{itemize}

\paragraph{Trampas y errores frecuentes.}
\begin{itemize}\setlength\itemsep{0pt}
	\item El \texttt{pot[]} se actualiza \textbf{al final} de cada iteracion (no durante).
	\item Las back-edges tienen \textbf{costo negativo}, esencial para corregir elecciones previas.
	\item Si los costos son muy grandes, considerar modular arithmetic.
	\item Para problemas con flujo exacto, aniadir verificacion post-algoritmo.
\end{itemize}

\paragraph{Codigo.}
Ver \texttt{cpp.tex}, seccion \textit{Min-Cost Max-Flow}.

\subsection{Especales}

\subsubsection{2-SAT (implication graph + SCC)}

\paragraph{Idea intuitiva.}
Resuelve formulas booleanas en \textbf{CNF} con clausulas de \textbf{2 literales}: $(l_1 \lor l_2)$. Cada variable $x$ se representa con dos nodos: $x$ (verdadero) y $\neg x$ (falso). Cada clausula $l_1 \lor l_2$ se descompone en dos implicaciones: $\neg l_1 \to l_2$ y $\neg l_2 \to l_1$. Si en el grafo de implicaciones $\neg x$ y $x$ estan en la misma SCC, la formula es \textbf{UNSAT}. La demostracion: si $\neg x \to x$ en el grafo, entonces $\neg x$ implica $x$, asi que $\neg x$ debe ser falso y $x$ verdadero; pero entonces $\neg x$ es falso implica una contradiccion.

\paragraph{Propiedades clave (invariantes).}
\begin{itemize}\setlength\itemsep{0pt}
	\item Cada literal es un nodo; total $2n$ nodos.
	\item Kosaraju/Tarjan dan las SCCs.
	\item Si SAT: asignar $x = $ (topological order del condensation).
	\item Clausulas \texttt{setTrue(x)} se modelan como \texttt{implies(!x, x)}.
	\item La formula es SAT sii no hay variable $x$ con $x$ y $\neg x$ en la misma SCC.
\end{itemize}

\paragraph{Codigo clave.}
Conversion de clausulas a implicaciones:
\begin{verbatim}
// Clausula (a OR b) -> implica ~a -> b Y ~b -> a
addImplication(not(a), b);
addImplication(not(b), a);
// Forzar x: implica ~x -> x
if (force) addImplication(not(x), x);
\end{verbatim}

\paragraph{Complejidad.}
\begin{itemize}\setlength\itemsep{0pt}
	\item $O(V + E) = O(N + M)$ con Kosaraju o Tarjan.
	\item Memoria $O(V + E)$.
\end{itemize}

\paragraph{Donde tocar para modificar.}
\begin{itemize}\setlength\itemsep{0pt}
	\item \textbf{Forzar variables}: aniadir \texttt{setTrue} / \texttt{setFalse}.
	\item \textbf{3-SAT}: NP-completo; no usar este algoritmo.
	\item \textbf{Contar soluciones}: las SCCs condensation dan una formula para contar; multiplicar las formas de asignar variables en cada componente no-trivial.
	\item \textbf{Output de la asignacion}: aniadir el vector \texttt{assignment[]}.
\end{itemize}

\paragraph{Trampas y errores frecuentes.}
\begin{itemize}\setlength\itemsep{0pt}
	\item La asignacion usa el \textbf{orden topologico inverso} del condensation: \texttt{val[i] = comp[2i] < comp[2i+1]}.
	\item El ID de $\neg x$ es \texttt{2*x + 1} (o similar); verificar consistencia.
	\item Si hay clausulas con literales repetidos (e.g., $(x \lor x)$), tratar como $(x)$ simple.
\end{itemize}

\paragraph{Codigo.}
Ver \texttt{cpp.tex}, seccion \textit{2-SAT (implication graph + SCC)}.

\vspace{0.5em}

\subsubsection{Euler tour / Hierholzer (camino y circuito)}

\paragraph{Idea intuitiva.}
Encuentra un \textbf{camino o circuito euleriano} (que usa cada arista exactamente una vez). \textbf{Condiciones}:
\begin{itemize}\setlength\itemsep{0pt}
	\item \textbf{No dirigido}: todos los grados pares (circuito) o exactamente 2 impares (camino, empieza en uno, termina en el otro). Ademas, conexo sobre el subgrafo con aristas.
	\item \textbf{Dirigido}: $\text{indeg}(v) = \text{outdeg}(v)$ para todos (circuito) o exactamente un nodo con $\text{outdeg} = \text{indeg} + 1$ y otro con $\text{indeg} = \text{outdeg} + 1$ (camino).
\end{itemize}

\paragraph{Hierholzer}: elegir un ciclo desde el inicio, mientras haya ciclos ``laterales'' desde un nodo del camino, mergearlos. Implementacion con pila: avanzar por aristas, al no poder mas aniadir a la respuesta y backtrack. La demostracion de $O(V + E)$: cada arista se procesa una vez (entra y sale del stack una vez).

\paragraph{Propiedades clave (invariantes).}
\begin{itemize}\setlength\itemsep{0pt}
	\item Multigrafos permitidos; las aristas se marcan como usadas.
	\item No dirigido: cada arista aparece \textbf{dos veces} en \texttt{adj} (ida y vuelta).
	\item La construccion del path con ids consistentes es importante.
	\item El resultado es unico salvo orden de las aristas en multigrafos.
\end{itemize}

\paragraph{Codigo clave.}
Hierholzer con pila:
\begin{verbatim}
vector<int> path;
stack<int> st;
st.push(start);
while (!st.empty()) {
    int v = st.top();
    while (!adj[v].empty() && used[adj[v].back()]) adj[v].pop_back();
    if (adj[v].empty()) {
        path.push_back(v);
        st.pop();
    } else {
        int e = adj[v].back();
        used[e] = true;
        int u = edges[e].to(v);
        adj[v].pop_back();
        st.push(u);
    }
}
reverse(path.begin(), path.end());
\end{verbatim}

\paragraph{Complejidad.}
\begin{itemize}\setlength\itemsep{0pt}
	\item $O(V + E)$.
	\item Memoria $O(V + E)$.
\end{itemize}

\paragraph{Donde tocar para modificar.}
\begin{itemize}\setlength\itemsep{0pt}
	\item \textbf{Reconstruir con multigrafo}: aniadir ids consistentes antes del algoritmo.
	\item \textbf{Aplicar a ``de Bruijn'' o problemas de strings}: modelar como multigrafo y buscar Euler.
	\item \textbf{Dirigido vs no dirigido}: el snippet tiene ambas versiones.
	\item \textbf{Letter addition}: modelar palabras como aristas y buscar el superstring.
\end{itemize}

\paragraph{Trampas y errores frecuentes.}
\begin{itemize}\setlength\itemsep{0pt}
	\item Si el grafo no es conexo (sobre las aristas), no existe euleriano. Verificar primero.
	\item Para multigrafos, los IDs de aristas son cruciales para distinguirlas.
	\item Si el problema es ``encuentra todos los euler tours'', el algoritmo es exponencial.
\end{itemize}

\paragraph{Codigo.}
Ver \texttt{cpp.tex}, seccion \textit{Euler tour / Hierholzer (camino y circuito)}.

\vspace{0.5em}

\subsubsection{Hamiltonian path / cycle (bitmask DP + backtracking)}

\paragraph{Idea intuitiva.}
Un \textbf{camino hamiltoniano} visita cada vertice \textbf{exactamente una vez}. \textbf{Determinar} si existe es NP-completo en general, pero con $V \le 20$ se puede hacer DP con bitmask en $O(V^2 \cdot 2^V)$. \textbf{Teoremas suficientes}: \textbf{Dirac} (grado $\ge n/2$) y \textbf{Ore} ($\text{deg}(u) + \text{deg}(v) \ge n$ para todo par no adyacente). La idea de la DP: cada estado es (conjunto de visitados, ultimo vertice); la transicion es agregar un vecino no visitado.

\paragraph{Propiedades clave (invariantes).}
\begin{itemize}\setlength\itemsep{0pt}
	\item DP: $dp[mask][v]$ = true si hay un camino que visita exactamente \texttt{mask} y termina en $v$.
	\item TSP: variante con pesos minimos (Hamilton cycle minimo).
	\item Reconstruccion: guardar \texttt{par[mask][v]} (predecesor).
	\item $2^V$ estados con $V$ bits cada uno; factible para $V \le 20$.
\end{itemize}

\paragraph{Codigo clave.}
DP con bitmask:
\begin{verbatim}
dp[1][0] = 0;  // empezar en 0
for (int mask = 1; mask < (1<<n); ++mask)
    for (int v = 0; v < n; ++v) if (dp[mask][v] < INF)
        for (int u = 0; u < n; ++u)
            if (!(mask & (1<<u)))
                dp[mask | (1<<u)][u] = min(dp[mask | (1<<u)][u],
                                            dp[mask][v] + w[v][u]);
// TSP: respuesta = min(dp[(1<<n)-1][v] + w[v][0])
\end{verbatim}

\paragraph{Complejidad.}
\begin{itemize}\setlength\itemsep{0pt}
	\item $O(V^2 \cdot 2^V)$ tiempo, $O(V \cdot 2^V)$ memoria.
	\item Para $V = 20$: $\sim 4 \cdot 10^8$ operaciones (lento pero factible).
\end{itemize}

\paragraph{Donde tocar para modificar.}
\begin{itemize}\setlength\itemsep{0pt}
	\item \textbf{TSP asimetrico}: aniadir pesos $w[v][u]$ (no necesariamente simetricos).
	\item \textbf{Hamilton con inicio obligatorio}: fijar \texttt{src} en el DP inicial.
	\item \textbf{Bitmask DP mas compacto}: usar \texttt{long long} como mascara si $V \le 64$.
	\item \textbf{TSP optimo aproximado}: MST, 2-opt, simulated annealing.
	\item \textbf{Reconstruir el path}: usar \texttt{par[mask][v]} y backtrack.
\end{itemize}

\paragraph{Trampas y errores frecuentes.}
\begin{itemize}\setlength\itemsep{0pt}
	\item TSP asume grafo \textbf{completo} (con \texttt{LINF} para aristas faltantes); si no, aniadir chequeo.
	\item La condicion ``\texttt{mask == (1<<n) - 1}'' cierra el ciclo retornando al origen.
	\item Para $V > 20$, la memoria $2^V$ se vuelve prohibitiva; usar meet-in-the-middle o heuristicas.
\end{itemize}

\paragraph{Codigo.}
Ver \texttt{cpp.tex}, seccion \textit{Hamiltonian path / cycle (bitmask DP + backtracking)}.

% ============================================================
\section{DP}

\subsubsection{Knapsack 0/1}

\paragraph{Idea intuitiva.}
Dado un conjunto de $N$ objetos con peso $w_i$ y valor $v_i$, maximizar $\sum v_i$ con la restriccion $\sum w_i \le W$. La recurrencia clasica:
\[ dp[i][j] = \max(dp[i-1][j],\, dp[i-1][j - w_i] + v_i) \quad \text{si } j \ge w_i. \]
La optimizacion clave: la DP es 1-dimensional si iteramos $j$ \textbf{de mayor a menor} (asi no usamos el objeto $i$ dos veces). La razon de la iteracion inversa: en cada iteracion queremos el estado del ``item anterior''; iterar de menor a mayor contaminaria \texttt{dp[j - w_i]} con el item actual.

\paragraph{Propiedades clave (invariantes).}
\begin{itemize}\setlength\itemsep{0pt}
	\item $dp[j]$ al final = maximo valor con peso $\le j$.
	\item El bucle inverso \texttt{for (j = W; j >= w[i]; j--)} es lo que distingue 0/1 de unbounded.
	\item Solucion: maximo valor con cualquier peso $\le W$.
	\item La DP es exacta; no hay aproximacion.
\end{itemize}

\paragraph{Codigo clave.}
La version 1D optimizada:
\begin{verbatim}
vector<long long> dp(W + 1, 0);
for (int i = 0; i < N; ++i)
    for (int j = W; j >= w[i]; --j)  // invertido
        dp[j] = max(dp[j], dp[j - w[i]] + v[i]);
return dp[W];
\end{verbatim}

\paragraph{Complejidad.}
\begin{itemize}\setlength\itemsep{0pt}
	\item $O(N \cdot W)$ tiempo, $O(W)$ memoria.
	\item Constante muy pequena (solo sumas y max).
\end{itemize}

\paragraph{Donde tocar para modificar.}
\begin{itemize}\setlength\itemsep{0pt}
	\item \textbf{Reconstruir la seleccion}: aniadir \texttt{take[i][j]} (booleano o bool) en una version 2D.
	\item \textbf{Items con peso 0}: edge case; aniadir guard.
	\item \textbf{Knapsack con restricciones de grupos}: ``elegir a lo sumo 1 de cada grupo'' se modela como 0/1 con un peso extra.
	\item \textbf{Tamano de $W$}: si $W \le 10^9$ pero $N$ es chico, usar meet-in-the-middle.
	\item \textbf{Cambiar max por min}: aniadir \texttt{INF} en la inicializacion.
\end{itemize}

\paragraph{Trampas y errores frecuentes.}
\begin{itemize}\setlength\itemsep{0pt}
	\item Overflow en \texttt{dp[j - w[i]] + v[i]}; usar \texttt{long long} si los valores son grandes.
	\item El bucle \textbf{debe} ir de mayor a menor.
	\item Si los pesos son grandes y $N$ pequeno, conviene meet-in-the-middle.
	\item Para Knapsack exacto ($\sum w_i = W$), aniadir check al final.
\end{itemize}

\paragraph{Codigo.}
Ver \texttt{cpp.tex}, seccion \textit{Knapsack 0/1}.

\vspace{0.5em}

\subsubsection{Knapsack unbounded}

\paragraph{Idea intuitiva.}
Igual que Knapsack 0/1 pero \textbf{cada item se puede tomar multiples veces}. La recurrencia es similar:
\[ dp[j] = \max(dp[j],\, dp[j - w_i] + v_i), \]
pero ahora el bucle va de \textbf{menor a mayor} (porque tomar el item $i$ no afecta corridas futuras del mismo item). La idea: tras tomar el item $i$, queremos volver a tener la opcion de tomarlo; iterar de menor a mayor permite que el item ``se acumule''.

\paragraph{Propiedades clave (invariantes).}
\begin{itemize}\setlength\itemsep{0pt}
	\item El bucle \texttt{for (j = w[i]; j <= W; j++)} es lo que distingue unbounded.
	\item A diferencia de 0/1, no hay riesgo de ``usar dos veces el mismo item''.
	\item La DP es exacta; cada ``subconjunto'' se cuenta correctamente.
\end{itemize}

\paragraph{Codigo clave.}
La version 1D con bucle directo:
\begin{verbatim}
vector<long long> dp(W + 1, 0);
for (int i = 0; i < N; ++i)
    for (int j = w[i]; j <= W; ++j)  // directo (no invertido)
        dp[j] = max(dp[j], dp[j - w[i]] + v[i]);
return dp[W];
\end{verbatim}

\paragraph{Complejidad.}
\begin{itemize}\setlength\itemsep{0pt}
	\item $O(N \cdot W)$ tiempo, $O(W)$ memoria.
	\item Constante muy pequena.
\end{itemize}

\paragraph{Donde tocar para modificar.}
\begin{itemize}\setlength\itemsep{0pt}
	\item \textbf{Reconstruir cuantos de cada item}: aniadir \texttt{cnt[i]} y backtrack.
	\item \textbf{Cambiar moneda minima por maxima}: el snippet maximiza valor; para minimizar monedas, cambiar a \texttt{min} con inicializacion en \texttt{INF}.
	\item \textbf{Coin change}: si todos los pesos son positivos y queremos ``minimo numero de monedas'', es unbounded con min.
\end{itemize}

\paragraph{Trampas y errores frecuentes.}
\begin{itemize}\setlength\itemsep{0pt}
	\item Misma que Knapsack 0/1 con overflow.
	\item El orden de iteracion de items no afecta el resultado (convexo).
	\item Si se quiere ``maximo valor con exactamente $k$ items'', aniadir segunda dimension.
\end{itemize}

\paragraph{Codigo.}
Ver \texttt{cpp.tex}, seccion \textit{Knapsack unbounded}.

\vspace{0.5em}

\subsubsection{Knapsack meet-in-the-middle}

\paragraph{Idea intuitiva.}
Para $N \le 40$ y $W$ moderado, no se puede hacer $O(NW)$ directamente. Partir en dos mitades, enumerar todos los $2^{N/2}$ subconjuntos de cada mitad, y combinar. Para cada subconjunto de la segunda mitad con peso $sw_2$, encontrar el maximo valor de la primera mitad con peso $\le W - sw_2$. La idea: ``divide y venceras'' sobre el tamano de la entrada; cada mitad tiene mitad del espacio, asi que la combinacion es geometrica en lugar de multiplicativa.

\paragraph{Propiedades clave (invariantes).}
\begin{itemize}\setlength\itemsep{0pt}
	\item Hay $2^{N/2}$ subconjuntos por mitad; total $O(2^{N/2})$.
	\item \texttt{sums} guarda los valores de subconjuntos de la primera mitad ordenados (o preprocesados).
	\item Para cada subconjunto de la segunda mitad, busqueda lineal o binaria en \texttt{sums}.
	\item La mitad mas grande (con $N/2 + 1$ items para $N$ impar) tiene $2^{N/2+1}$ subconjuntos.
\end{itemize}

\paragraph{Codigo clave.}
Enumerar subconjuntos de la primera mitad:
\begin{verbatim}
int n1 = N / 2, n2 = N - n1;
vector<pair<long long, long long>> sums1;
for (int mask = 0; mask < (1 << n1); ++mask) {
    long long peso = 0, valor = 0;
    for (int i = 0; i < n1; ++i) if (mask & (1 << i)) {
        peso += w[i]; valor += v[i];
    }
    sums1.push_back({peso, valor});
}
sort(sums1.begin(), sums1.end());  // por peso
// eliminar duplicados dominados
\end{verbatim}

\paragraph{Complejidad.}
\begin{itemize}\setlength\itemsep{0pt}
	\item $O(2^{N/2})$ tiempo y memoria (mejor que $O(NW)$ cuando $W > 2^{N/2}$).
	\item Para $N = 40$: $\sim 2 \cdot 10^6$ subconjuntos, factible.
\end{itemize}

\paragraph{Donde tocar para modificar.}
\begin{itemize}\setlength\itemsep{0pt}
	\item \textbf{Tamano mayor}: si $N$ es mayor pero $W$ es chico, conviene iterar por peso, no por subconjuntos.
	\item \textbf{Varias mochilas}: adaptar para multiples $W$.
	\item \textbf{Optimizar el ``combinado''}: usar \texttt{upper\_bound} para encontrar el maximo valor con peso $\le \text{rem}.
	\item \textbf{Eliminar pares dominados}: si un par (peso, valor) tiene otro con menos peso y mas valor, eliminarlo.
\end{itemize}

\paragraph{Trampas y errores frecuentes.}
\begin{itemize}\setlength\itemsep{0pt}
	\item $N$ impar: una mitad tiene $N/2 + 1$ items.
	\item Overflow al sumar pesos/valores en \texttt{long long}.
	\item Sin deduplicacion, el algoritmo es correcto pero mas lento.
\end{itemize}

\paragraph{Codigo.}
Ver \texttt{cpp.tex}, seccion \textit{Knapsack meet-in-the-middle}.

\vspace{0.5em}

\subsubsection{LIS O(n log n)}

\paragraph{Idea intuitiva.}
La \textbf{Longest Increasing Subsequence} se calcula en $O(n \log n)$ con el truco del \textbf{patience sorting}: mantener un arreglo \texttt{dp} donde \texttt{dp[len]} es el minimo ultimo valor de una IS de largo \texttt{len}. Para cada $x$, encontrar el primer \texttt{dp[len] >= x} y reemplazarlo. El largo final es la LIS. La demostracion: si hay dos IS del mismo largo, la que termina en un valor menor es ``mejor'' (puede extenderse mas facilmente).

\paragraph{Propiedades clave (invariantes).}
\begin{itemize}\setlength\itemsep{0pt}
	\item \texttt{dp} es estrictamente creciente.
	\item \texttt{lower\_bound} da LIS estricta; \texttt{upper\_bound} da no-estricta ($\le$).
	\item El algoritmo es \textbf{online}: procesa elementos en orden.
	\item La longitud de \texttt{dp} es la LIS al final.
\end{itemize}

\paragraph{Codigo clave.}
La version online:
\begin{verbatim}
vector<int> dp;
for (int x : a) {
    auto it = lower_bound(dp.begin(), dp.end(), x);
    if (it == dp.end()) dp.push_back(x);
    else *it = x;
}
// dp.size() = longitud de LIS estricta
\end{verbatim}

\paragraph{Complejidad.}
\begin{itemize}\setlength\itemsep{0pt}
	\item $O(n \log n)$.
	\item Memoria $O(n)$.
\end{itemize}

\paragraph{Donde tocar para modificar.}
\begin{itemize}\setlength\itemsep{0pt}
	\item \textbf{Reconstruir la LIS}: aniadir \texttt{parent[]} y \texttt{posInDp[]} para backtrack.
	\item \textbf{Contar LIS}: DP clasica $O(n^2)$ o segment tree sobre $dp$.
	\item \textbf{LDS / Longest non-increasing}: invertir la condicion en \texttt{lower\_bound} o multiplicar por -1.
	\item \textbf{Longest bitonic subsequence}: combinar LIS desde izq y desde der.
	\item \textbf{LIS en 2D}: ordenar por una dimension y aplicar LIS en la otra.
\end{itemize}

\paragraph{Trampas y errores frecuentes.}
\begin{itemize}\setlength\itemsep{0pt}
	\item \texttt{lower\_bound} vs \texttt{upper\_bound}: si tu comparacion es $\le$ vs $<$, cambia.
	\item Si los elementos son \texttt{int} pero podrian ser negativos, aniadir offset para usar como indices.
	\item Para LIS reconstruida, aniadir el array \texttt{parent} que apunta al anterior en la IS.
\end{itemize}

\paragraph{Codigo.}
Ver \texttt{cpp.tex}, seccion \textit{LIS O(n log n)}.

\vspace{0.5em}

\subsubsection{LCS / Edit distance}

\paragraph{Idea intuitiva.}
\textbf{LCS} (Longest Common Subsequence): para dos strings $a$, $b$, $dp[i][j]$ = LCS de $a[0..i)$ y $b[0..j)$. Recurrencia:
\[ dp[i][j] = \begin{cases} dp[i-1][j-1] + 1 & \text{si } a[i-1] = b[j-1] \\ \max(dp[i-1][j], dp[i][j-1]) & \text{si no} \end{cases} \]

\textbf{Edit distance} (Levenshtein): numero minimo de operaciones (insert, delete, replace) para convertir $a$ en $b$. Recurrencia similar con un $+1$ adicional. La idea: ``alinear'' los strings eligiendo matches/mismatches/insert/delete.

\paragraph{Propiedades clave (invariantes).}
\begin{itemize}\setlength\itemsep{0pt}
	\item LCS y Edit Distance son DP $O(nm)$.
	\item Memoria se puede reducir a $O(\min(n, m))$ con rolling array.
	\item Edit distance con peso por operacion distinto: ajustar el \texttt{+1} por el peso.
	\item La subestructura optima: la respuesta para prefijos se construye de prefijos menores.
\end{itemize}

\paragraph{Codigo clave.}
La recurrencia clasica de LCS:
\begin{verbatim}
vector<vector<int>> dp(n + 1, vector<int>(m + 1, 0));
for (int i = 1; i <= n; ++i)
    for (int j = 1; j <= m; ++j)
        if (a[i-1] == b[j-1]) dp[i][j] = dp[i-1][j-1] + 1;
        else dp[i][j] = max(dp[i-1][j], dp[i][j-1]);
return dp[n][m];
\end{verbatim}

\paragraph{Complejidad.}
\begin{itemize}\setlength\itemsep{0pt}
	\item $O(NM)$ tiempo y memoria.
	\item Para $N, M \le 5000$: $\sim 2.5 \cdot 10^7$ operaciones.
\end{itemize}

\paragraph{Donde tocar para modificar.}
\begin{itemize}\setlength\itemsep{0pt}
	\item \textbf{Reconstruir la LCS / las operaciones}: aniadir \texttt{prev[i][j]} para backtrack.
	\item \textbf{LCS de varios strings}: NP-hard para $\ge 3$ strings; en la practica se hace con bitmask.
	\item \textbf{Edit distance con costos distintos}: cambiar el \texttt{+1} por el costo apropiado.
	\item \textbf{Rolling array}: si memoria es problema, mantener solo 2 filas.
	\item \textbf{Hirschberg}: algoritmo $O(NM)$ tiempo, $O(\min(N,M))$ memoria, con reconstruccion.
\end{itemize}

\paragraph{Trampas y errores frecuentes.}
\begin{itemize}\setlength\itemsep{0pt}
	\item Indices en la recurrencia: \texttt{a[i-1]} y \texttt{b[j-1]} (no \texttt{a[i]}).
	\item Para reconstruccion, el caso \texttt{a[i-1] == b[j-1]} tiene prioridad.
	\item Si los strings son muy largos, considerar Hirschberg o bitset optimization.
\end{itemize}

\paragraph{Codigo.}
Ver \texttt{cpp.tex}, seccion \textit{LCS / Edit distance}.

\vspace{0.5em}

\subsubsection{Bitmask DP over subsets}

\paragraph{Idea intuitiva.}
Cuando el estado es un \textbf{subconjunto} de $\{0, \dots, N-1\}$, se representa como una mascara de $N$ bits. Para iterar sobre todos los subconjuntos de \texttt{mask}:
\[ \text{for (sub = mask; sub; sub = (sub - 1) \& mask)} \{ \dots \} \]
Para iterar sobre todos los superconjuntos de \texttt{mask} en $[0, U)$:
\[ \text{for (sup = mask; sup < U; sup = (sup + 1) | mask)} \{ \dots \} \]
La idea: las mascaras forman un lattice de subconjuntos; los trucos aritmeticos (\texttt{(sub-1) \& mask} y \texttt{(sup+1) | mask}) enumeran todo el lattice eficientemente.

\paragraph{Propiedades clave (invariantes).}
\begin{itemize}\setlength\itemsep{0pt}
	\item Iterar subsets: $O(2^k)$ para una mascara con $k$ bits.
	\item \texttt{sub \& mask == sub} siempre; \texttt{mask - sub} no siempre da un subconjunto.
	\item \texttt{\_\_builtin\_popcount(mask)} cuenta los bits.
	\item \texttt{(sub - 1) \& mask} itera subsets de \texttt{mask} en orden decreciente de inclusion.
\end{itemize}

\paragraph{Codigo clave.}
Iteracion sobre subsets de una mascara:
\begin{verbatim}
// Todos los subsets de mask
for (int sub = mask; sub; sub = (sub - 1) & mask) {
    // procesar sub
}
// Incluir sub = 0 si es necesario
sub = 0;  // procesar primero
for (int sup = mask; sup < (1 << N); sup = (sup + 1) | mask) {
    // procesar sup (todos los superconjuntos)
}
\end{verbatim}

\paragraph{Complejidad.}
\begin{itemize}\setlength\itemsep{0pt}
	\item $O(N \cdot 2^N)$ para una DP estandar sobre todos los subsets.
	\item Iterar subsets de \texttt{mask}: $O(2^{|mask|})$.
\end{itemize}

\paragraph{Donde tocar para modificar.}
\begin{itemize}\setlength\itemsep{0pt}
	\item \textbf{TSP / Hamiltonian path}: bitmask DP (ver tambien \textit{Hamiltonian path} en \texttt{cpp.tex} seccion Grafos).
	\item \textbf{Set cover}: iterar subsets.
	\item \textbf{DP con supermask}: si el estado incluye ``todos los elementos que \textbf{no} hemos usado'', usar supermask iteration.
	\item \textbf{Bitmasks grandes}: con $N > 20$, aniadir compresion o simulacion con segmentos.
\end{itemize}

\paragraph{Trampas y errores frecuentes.}
\begin{itemize}\setlength\itemsep{0pt}
	\item El orden de iteracion de subsets de \texttt{mask} es decreciente (de \texttt{mask} a $0$).
	\item \texttt{mask = 0} no tiene subsets no triviales; el bucle lo saltea correctamente.
	\item El truco \texttt{(sub - 1) \& mask} requiere cuidado con \texttt{sub = 0} (termina el bucle).
	\item Superconjuntos: \texttt{sup < (1 << N)} debe estar bien definido.
\end{itemize}

\paragraph{Codigo.}
Ver \texttt{cpp.tex}, seccion \textit{Bitmask DP over subsets}.

\vspace{0.5em}

\subsubsection{SOS DP}

\paragraph{Idea intuitiva.}
Para calcular, para cada \texttt{mask}, la suma de $f(\text{submask})$ sobre todos los $\text{submask} \subseteq \text{mask}$:
\[ dp[mask] = \sum_{\text{sub} \subseteq mask} f(\text{sub}). \]
Se computa en $O(N \cdot 2^N)$ mediante un loop por bit: para cada $i$, sumar a $dp[mask]$ la contribucion de $dp[mask \oplus (1 << i)]$ cuando $mask$ tiene el bit $i$. La idea: cada subconjunto $s \subseteq m$ se obtiene quitando bits uno a uno; asi, sumar incrementalmente captura todos los subconjuntos en $O(N \cdot 2^N)$.

\paragraph{Propiedades clave (invariantes).}
\begin{itemize}\setlength\itemsep{0pt}
	\item Inicializar $dp = f$.
	\item Para cada bit $i$: \texttt{if (mask \& (1 << i)) dp[mask] += dp[mask \^{} (1 << i)]}.
	\item Variantes: supersets DP (sumar sobre superconjuntos).
	\item El orden de iteracion (bit externo, mascara interno) es crucial.
\end{itemize}

\paragraph{Codigo clave.}
SOS DP (Sum over Subsets):
\begin{verbatim}
vector<long long> dp(1 << N);
for (int i = 0; i < (1 << N); ++i) dp[i] = f[i];
for (int bit = 0; bit < N; ++bit)
    for (int mask = 0; mask < (1 << N); ++mask)
        if (mask & (1 << bit))
            dp[mask] += dp[mask ^ (1 << bit)];
// dp[mask] = suma de f(sub) para sub ⊆ mask
\end{verbatim}

\paragraph{Complejidad.}
\begin{itemize}\setlength\itemsep{0pt}
	\item $O(N \cdot 2^N)$.
	\item Memoria $O(2^N)$.
\end{itemize}

\paragraph{Donde tocar para modificar.}
\begin{itemize}\setlength\itemsep{0pt}
	\item \textbf{Max/min en lugar de suma}: cambiar \texttt{+=} por la operacion correspondiente.
	\item \textbf{Supersets DP}: iterar \texttt{mask} de $0$ a $2^N$ y aniadir \texttt{dp[mask] += dp[mask | (1<<i)]}.
	\item \textbf{XOR convolution} (subset convolution): version mas avanzada.
	\item \textbf{Aplicaciones}: contar subconjuntos con propiedades especificas.
\end{itemize}

\paragraph{Trampas y errores frecuentes.}
\begin{itemize}\setlength\itemsep{0pt}
	\item Si $N$ es 30 o mas, $2^N$ no entra en memoria; usar trucos o reducir $N$.
	\item El sentido del XOR (sumar a \texttt{mask \^{} bit}) es para ``agregar el bit $i$ al subconjunto''.
	\item En supersets DP, la condicion es \texttt{(mask \& (1<<bit)) == 0}.
\end{itemize}

\paragraph{Codigo.}
Ver \texttt{cpp.tex}, seccion \textit{SOS DP}.

\vspace{0.5em}

\subsubsection{Digit DP template}

\paragraph{Idea intuitiva.}
Contar numeros en $[0, N]$ que cumplen una propiedad sobre sus digitos. La idea: DP por posicion, manteniendo:
\begin{itemize}\setlength\itemsep{0pt}
	\item \texttt{pos}: posicion actual (de izquierda a derecha).
	\item \texttt{tight}: si los digitos ya colocados son exactamente los de $N$ (limite superior estricto).
	\item \texttt{state}: estado adicional (digitos usados, suma, etc.).
\end{itemize}
La clave: la condicion \texttt{tight} permite acotar correctamente el siguiente digito; sin ella, la DP contaria incorrectamente.

\paragraph{Propiedades clave (invariantes).}
\begin{itemize}\setlength\itemsep{0pt}
	\item Si \texttt{tight} es true, el siguiente digito esta acotado por el de $N$; si no, puede ser 0-9.
	\item Memoria: $O(D \cdot 2 \cdot \text{states})$.
	\item Caso base: si \texttt{pos == D}, devolver 1 si el estado cumple la propiedad, sino 0.
	\item El bit \texttt{started} distingue ``numero con menos digitos'' (con leading zeros).
\end{itemize}

\paragraph{Codigo clave.}
La estructura basica:
\begin{verbatim}
long long dp(int pos, bool tight, State s) {
    if (pos == D) return check(s) ? 1 : 0;
    if (memo[pos][tight][s] != -1) return memo;
    int limit = tight ? digits[pos] : 9;
    long long ans = 0;
    for (int d = 0; d <= limit; ++d)
        ans += dp(pos + 1, tight && (d == limit), update(s, d, pos));
    return memo[pos][tight][s] = ans;
}
\end{verbatim}

\paragraph{Complejidad.}
\begin{itemize}\setlength\itemsep{0pt}
	\item $O(D \cdot 10 \cdot \text{states})$ donde $D$ es el numero de digitos.
	\item Para $D \le 19$ (hasta $10^{18}$): $\sim 200 \cdot \text{states}$.
\end{itemize}

\paragraph{Donde tocar para modificar.}
\begin{itemize}\setlength\itemsep{0pt}
	\item \textbf{Cambiar la propiedad}: modificar \texttt{updateState} y la condicion base.
	\item \textbf{Multiples estados}: aniadir mas dimensiones al \texttt{memo}.
	\item \textbf{Leading zeros}: aniadir un flag \texttt{started} para distinguir numeros con menos digitos.
	\item \textbf{Rango $[L, R]$}: calcular $f(R) - f(L-1)$.
\end{itemize}

\paragraph{Trampas y errores frecuentes.}
\begin{itemize}\setlength\itemsep{0pt}
	\item Olvidar resetear \texttt{memo} entre queries.
	\item Si la cantidad de estados es grande, el \texttt{memo[pos][tight][state]} puede ser muy grande; usar \texttt{map}.
	\item La condicion \texttt{started} es crucial para numeros con menos digitos que $D$.
	\item Si $N$ es muy grande ($> 10^{18}$), aniadir soporte para long long.
\end{itemize}

\paragraph{Codigo.}
Ver \texttt{cpp.tex}, seccion \textit{Digit DP template}.

\vspace{0.5em}

\subsubsection{Tree DP (reroot)}

\paragraph{Idea intuitiva.}
Calcular alguna propiedad asumiendo que \textbf{cada nodo es la raiz} del arbol. La tecnica clasica de \textbf{rerooting} en dos pasadas:
\begin{itemize}\setlength\itemsep{0pt}
	\item \textbf{DFS 1}: desde una raiz arbitraria, calcular \texttt{dp[v]} para el subarbol de $v$.
	\item \textbf{DFS 2}: para cada hijo $u$ de $v$, calcular \texttt{up[u]} = respuesta si la raiz fuera $u$.
\end{itemize}
La demostracion: tras la primera DFS, cada nodo conoce ``lo que aporta su subarbol''; en la segunda, se transmite ``lo que aporta el resto del arbol''.

\paragraph{Propiedades clave (invariantes).}
\begin{itemize}\setlength\itemsep{0pt}
	\item \texttt{dp[v]} depende solo del subarbol de $v$ (con la raiz original).
	\item \texttt{up[u]} se calcula de manera que cuando la raiz es $u$, ``todo lo que estaba en $v$ se ajusta''.
	\item Formula para reroot depende del problema; el snippet usa el ejemplo ``suma de distancias''.
	\item La combinacion \texttt{dp[u] + up[u]} da la respuesta con raiz $u$.
\end{itemize}

\paragraph{Codigo clave.}
Reroot para suma de distancias:
\begin{verbatim}
// dfs1: dp[v] = suma de distancias desde v al subarbol
void dfs1(int v, int p) {
    dp[v] = 0;
    for (int u : adj[v]) if (u != p) {
        dfs1(u, v);
        dp[v] += dp[u] + sz[u];  // cada hijo aporta su subarbol
    }
}
// dfs2: up[u] = distancia desde el resto del arbol a u
void dfs2(int v, int p) {
    for (int u : adj[v]) if (u != p) {
        up[u] = up[v] + dp[v] - dp[u] - sz[u] + (n - sz[u]);
        dfs2(u, v);
    }
}
\end{verbatim}

\paragraph{Complejidad.}
\begin{itemize}\setlength\itemsep{0pt}
	\item $O(N)$ para ambas pasadas.
	\item Cada arista se procesa un numero constante de veces.
\end{itemize}

\paragraph{Donde tocar para modificar.}
\begin{itemize}\setlength\itemsep{0pt}
	\item \textbf{Cambiar la cantidad a calcular}: modificar la operacion en \texttt{dfs1} y la formula en \texttt{dfs2}.
	\item \textbf{Arbol con pesos}: aniadir \texttt{weights[v]} y ajustar.
	\item \textbf{Re-root en arboles con ciclos}: requiere otra tecnica (DSU on tree).
	\item \textbf{Rooted trees diferentes}: si la propiedad no es invariante por raiz, usar otra tecnica.
\end{itemize}

\paragraph{Trampas y errores frecuentes.}
\begin{itemize}\setlength\itemsep{0pt}
	\item La formula de \texttt{up[u]} debe sumar y restar correctamente las contribuciones de $u$ y del resto.
	\item Para arboles con pesos, aniadir el peso en las formulas.
	\item Verificar con ejemplos pequenos antes de generalizar.
\end{itemize}

\paragraph{Codigo.}
Ver \texttt{cpp.tex}, seccion \textit{Tree DP (reroot)}.

% ============================================================
\section{Geometria}
% ============================================================

\subsection{Puntos y segmentos}

\subsubsection{Point + cross + dot}

\paragraph{Idea intuitiva.}
La estructura \texttt{Point} representa un punto en 2D con coordenadas enteras. Las dos operaciones vectoriales clave son:
\begin{itemize}\setlength\itemsep{0pt}
	\item \textbf{Cross product} $P \times Q = P_x Q_y - P_y Q_x$. Da el ``area con signo'' del paralelogramo formado por $P$ y $Q$. Usado para orientacion.
	\item \textbf{Dot product} $P \cdot Q = P_x Q_x + P_y Q_y$. Da el producto de las longitudes por el coseno del angulo. Usado para proyeccion.
\end{itemize}
La interpretacion geometrica del cross product: es el area con signo del paralelogramo; el signo indica la orientacion relativa.

\paragraph{Propiedades clave (invariantes).}
\begin{itemize}\setlength\itemsep{0pt}
	\item \texttt{cross(a, b) > 0} si $b$ esta a la izquierda de $a$ (sistema de coordenadas usual).
	\item \texttt{cross(a, b) = 0} si son colineales.
	\item \texttt{dot(a, b) > 0} si el angulo es agudo.
	\item El snippet define \texttt{cross} respecto a \texttt{this}, util para orientacion de 3 puntos.
	\item $|cross|^2 + |dot|^2 = |P|^2 \cdot |Q|^2$ (identidad de Lagrange).
\end{itemize}

\paragraph{Codigo clave.}
Las dos operaciones vectoriales:
\begin{verbatim}
long long cross(const Point& a, const Point& b) {
    return (long long)a.x * b.y - (long long)a.y * b.x;
}
long long dot(const Point& a, const Point& b) {
    return (long long)a.x * b.x + (long long)a.y * b.y;
}
\end{verbatim}

\paragraph{Complejidad.}
\begin{itemize}\setlength\itemsep{0pt}
	\item $O(1)$ por operacion.
	\item Constante muy pequena (4 multiplicaciones + 2 sumas).
\end{itemize}

\paragraph{Donde tocar para modificar.}
\begin{itemize}\setlength\itemsep{0pt}
	\item \textbf{Coordenadas doubles}: cambiar \texttt{int} por \texttt{double} o \texttt{long double}. Cuidado con la precision.
	\item \textbf{3D}: aniadir \texttt{z} y adaptar cross/dot.
	\item \textbf{Operadores adicionales}: aniadir \texttt{operator+}, \texttt{operator-}, \texttt{operator*} (escalar).
	\item \textbf{Distancia al cuadrado}: aniadir \texttt{len2()} que retorna $x^2 + y^2$.
\end{itemize}

\paragraph{Trampas y errores frecuentes.}
\begin{itemize}\setlength\itemsep{0pt}
	\item Cross product \textbf{no es} conmutativo; $P \times Q = -Q \times P$.
	\item Overflow con coordenadas grandes ($|x|, |y| > 10^4$) si se hace \texttt{x * y}.
	\item El signo de cross depende del sistema de coordenadas (en computacion, usualmente CCW positivo).
	\item Para ``igualdad'' de doubles, aniadir tolerancia \texttt{eps}.
\end{itemize}

\paragraph{Codigo.}
Ver \texttt{cpp.tex}, seccion \textit{Point + cross + dot}.

\vspace{0.5em}

\subsubsection{Interseccion de segmentos}

\paragraph{Idea intuitiva.}
Dados dos segmentos $[a, b]$ y $[c, d]$, determinar si se intersectan. La idea clasica:
\begin{itemize}\setlength\itemsep{0pt}
	\item Verificar orientaciones: $a, b, c$ y $a, b, d$ deben estar en lados opuestos de la recta $ab$; igual para $cd$ y $ab$.
	\item Caso colineal: si las rectas son paralelas (\texttt{cross} $= 0$) y los segmentos son colineales, verificar que los bounding boxes se intersectan.
\end{itemize}
La demostracion: dos segmentos se intersectan sii cada segmento ``cruza'' la linea que contiene al otro (estricto) o ambos son colineales con solapamiento.

\paragraph{Propiedades clave (invariantes).}
\begin{itemize}\setlength\itemsep{0pt}
	\item Funciona con segmentos en cualquier orientacion.
	\item Caso colineal: requiere check adicional de bounding box.
	\item Para intersectar en un \textbf{punto} especifico: aniadir division con cuidado de precision.
	\item Funciona con segmentos verticales/horizontales sin casos especiales.
\end{itemize}

\paragraph{Codigo clave.}
Test clasico de interseccion:
\begin{verbatim}
int sign(long long x) { return (x > 0) - (x < 0); }
bool intersect(Point a, Point b, Point c, Point d) {
    long long o1 = cross(b - a, c - a);
    long long o2 = cross(b - a, d - a);
    long long o3 = cross(d - c, a - c);
    long long o4 = cross(d - c, b - c);
    if (sign(o1) != sign(o2) && sign(o3) != sign(o4)) return true;
    // casos colineales
    if (o1 == 0 && onSegment(a, b, c)) return true;
    if (o2 == 0 && onSegment(a, b, d)) return true;
    if (o3 == 0 && onSegment(c, d, a)) return true;
    if (o4 == 0 && onSegment(c, d, b)) return true;
    return false;
}
\end{verbatim}

\paragraph{Complejidad.}
\begin{itemize}\setlength\itemsep{0pt}
	\item $O(1)$.
	\item Constante pequena.
\end{itemize}

\paragraph{Donde tocar para modificar.}
\begin{itemize}\setlength\itemsep{0pt}
	\item \textbf{Interseccion con punto}: aniadir la division despues de verificar orientacion.
	\item \textbf{Interseccion rayo-segmento}: ajustar la primera condicion.
	\item \textbf{Poligonos convexos}: usar separacion de ejes (SAT).
	\item \textbf{Multiples queries}: precomputar bounding boxes o usar segment tree.
\end{itemize}

\paragraph{Trampas y errores frecuentes.}
\begin{itemize}\setlength\itemsep{0pt}
	\item Si los segmentos tocan solo en un extremo, el snippet lo considera interseccion (algunos problemas quieren estricto). Ajustar segun el problema.
	\item El caso colineal es comun; olvidar el check adicional de bounding box da falsos negativos.
	\item En problemas con muchos segmentos, considerar sweep line para $O(n \log n)$.
\end{itemize}

\paragraph{Codigo.}
Ver \texttt{cpp.tex}, seccion \textit{Interseccion de segmentos}.

\vspace{0.5em}

\subsubsection{Contains colineal}

\paragraph{Idea intuitiva.}
Verifica si un punto $c$ esta sobre el segmento $[a, b]$ (asumiendo que ya son colineales). Si \texttt{cross(a, b, c) = 0} y $c$ esta dentro del bounding box de $[a, b]$, entonces $c$ esta en el segmento. La razon: si los tres puntos son colineales, $c$ esta en el segmento sii sus coordenadas estan entre las de $a$ y $b$.

\paragraph{Propiedades clave (invariantes).}
\begin{itemize}\setlength\itemsep{0pt}
	\item Asume colinealidad (no la verifica mas alla de \texttt{cross == 0}).
	\item Usa comparaciones con \texttt{<=} para extremos \textbf{inclusivos}.
	\item Funciona con segmentos verticales (la condicion se reduce a una sola coordenada).
\end{itemize}

\paragraph{Codigo clave.}
Test de contencion en segmento (asume colinealidad):
\begin{verbatim}
bool onSegment(Point a, Point b, Point c) {
    return min(a.x, b.x) <= c.x && c.x <= max(a.x, b.x) &&
           min(a.y, b.y) <= c.y && c.y <= max(a.y, b.y);
}
\end{verbatim}

\paragraph{Complejidad.}
\begin{itemize}\setlength\itemsep{0pt}
	\item $O(1)$.
	\item Constante minima.
\end{itemize}

\paragraph{Donde tocar para modificar.}
\begin{itemize}\setlength\itemsep{0pt}
	\item \textbf{Contener estricto}: cambiar \texttt{<=} a \texttt{<} para excluir extremos.
	\item \textbf{Coordenadas doubles}: usar tolerancia (\texttt{eps}).
	\item \textbf{Solo x o solo y}: aniadir check especifico para segmentos verticales/horizontales.
\end{itemize}

\paragraph{Trampas y errores frecuentes.}
\begin{itemize}\setlength\itemsep{0pt}
	\item No usar este snippet para puntos \textbf{no colineales}; siempre verificar primero la colinealidad.
	\item Si los puntos tienen doubles, el chequeo \texttt{cross == 0} puede fallar; aniadir tolerancia.
	\item El bounding box incluye los extremos; verificar si el problema quiere exclusivo o inclusivo.
\end{itemize}

\paragraph{Codigo.}
Ver \texttt{cpp.tex}, seccion \textit{Contains colineal}.

\subsection{Poligonos}

\subsubsection{Convex Hull (monotono)}

\paragraph{Idea intuitiva.}
Algoritmo de \textbf{Andrew's monotone chain}: ordenar puntos por $(x, y)$, luego construir la hull inferior y superior por separado. Para cada punto, aniadirlo al hull y, mientras el ultimo giro no sea ``counter-clockwise'' (\texttt{cross > 0}), popear. La demostracion: si los puntos se procesan en orden, el hull inferior y superior cubren la frontera completa sin auto-intersecciones.

\paragraph{Propiedades clave (invariantes).}
\begin{itemize}\setlength\itemsep{0pt}
	\item Ordena los puntos en $O(n \log n)$.
	\item Construye la hull en $O(n)$ despues.
	\item Output: poligono convexo en counter-clockwise, sin punto final duplicado.
	\item El \texttt{<= 0} incluye colineales; usar \texttt{< 0} para excluirlos.
	\item Cada punto aparece en la hull inferior o superior (no ambos).
\end{itemize}

\paragraph{Codigo clave.}
Construccion del hull inferior y superior:
\begin{verbatim}
sort(p.begin(), p.end());
vector<Point> lower, upper;
for (auto& p : points) {
    while (lower.size() >= 2 && cross(lower.back() - lower[lower.size()-2],
                                       p - lower.back()) <= 0)
        lower.pop_back();
    lower.push_back(p);
}
for (int i = points.size() - 1; i >= 0; --i) {
    auto& p = points[i];
    while (upper.size() >= 2 && cross(upper.back() - upper[upper.size()-2],
                                       p - upper.back()) <= 0)
        upper.pop_back();
    upper.push_back(p);
}
vector<Point> hull = lower;
hull.insert(hull.end(), upper.begin() + 1, upper.end() - 1);
\end{verbatim}

\paragraph{Complejidad.}
\begin{itemize}\setlength\itemsep{0pt}
	\item $O(n \log n)$.
	\item Constante muy pequena (solo cross products).
\end{itemize}

\paragraph{Donde tocar para modificar.}
\begin{itemize}\setlength\itemsep{0pt}
	\item \textbf{Colineales}: cambiar \texttt{<= 0} por \texttt{< 0} para excluirlos del hull.
	\item \textbf{Puntos repetidos}: \texttt{unique} despues de sort.
	\item \textbf{Hull 3D}: gift wrapping es $O(n^2)$; convex hull 3D es mas complejo.
	\item \textbf{Output con los puntos}: aniadir el hull completo (no solo upper/lower).
\end{itemize}

\paragraph{Trampas y errores frecuentes.}
\begin{itemize}\setlength\itemsep{0pt}
	\item El \texttt{pop\_back} y el \texttt{reverse} son esenciales para no duplicar el ultimo punto.
	\item Para excluir colineales, ajustar el \texttt{<= 0} vs \texttt{< 0}.
	\item Si los puntos estan duplicados, aniadir \texttt{unique} antes del sort.
\end{itemize}

\paragraph{Codigo.}
Ver \texttt{cpp.tex}, seccion \textit{Convex Hull (monotono)}.

\vspace{0.5em}

\subsubsection{Polygon Area}

\paragraph{Idea intuitiva.}
El area (con signo) de un poligono se calcula con la \textbf{Shoelace formula}: $A = \frac{1}{2} \sum_{i=0}^{n-1} (x_i y_{i+1} - x_{i+1} y_i)$, donde el indice es modulo $n$. El snippet retorna el \textbf{doble} del area (en valor absoluto). La formula sale del ``sumar areas de trapezoides'': cada lado del poligono define un trapezoide con el eje $x$; la suma con signo da el area.

\paragraph{Propiedades clave (invariantes).}
\begin{itemize}\setlength\itemsep{0pt}
	\item El signo indica la orientacion: positivo si counter-clockwise.
	\item Para el doble: no dividir por 2 si no es necesario (evita fracciones).
	\item Funciona para poligonos simples (no auto-intersectantes).
	\item Si el poligono es no convexo, la formula da el area ``con signo''.
\end{itemize}

\paragraph{Codigo clave.}
La suma de trapezoides:
\begin{verbatim}
long long shoelace(const vector<Point>& p) {
    long long s = 0;
    int n = p.size();
    for (int i = 0; i < n; ++i)
        s += (long long)p[i].x * p[(i+1)%n].y - (long long)p[(i+1)%n].x * p[i].y;
    return abs(s);  // doble del area
}
\end{verbatim}

\paragraph{Complejidad.}
\begin{itemize}\setlength\itemsep{0pt}
	\item $O(n)$.
	\item Constante minima.
\end{itemize}

\paragraph{Donde tocar para modificar.}
\begin{itemize}\setlength\itemsep{0pt}
	\item \textbf{Area exacta}: dividir por 2 (con cuidado si da fraccion; usar \texttt{double}).
	\item \textbf{Area firmada}: quitar \texttt{abs} y dejar el signo.
	\item \textbf{Punto dentro de poligono convexo}: con la hull, chequear que esta a un mismo lado de todas las aristas.
	\item \textbf{Triangulacion}: descomponer en triangulos y sumar areas.
\end{itemize}

\paragraph{Trampas y errores frecuentes.}
\begin{itemize}\setlength\itemsep{0pt}
	\item Si los puntos no estan en orden (CCW o CW), el area firmada puede ser incorrecta. Para signed area consistente, usar siempre el mismo orden.
	\item Para poligonos con holes, sumar las areas con signo apropiado.
	\item Con doubles, aniadir tolerancia \texttt{eps}.
\end{itemize}

\paragraph{Codigo.}
Ver \texttt{cpp.tex}, seccion \textit{Polygon Area}.

\vspace{0.5em}

\subsubsection{Distance point-segment / point-line}

\paragraph{Idea intuitiva.}
Distancia minima de un punto $p$ a un segmento $[a, b]$ o a la recta que pasa por $a$ y $b$.
\begin{itemize}\setlength\itemsep{0pt}
	\item \textbf{Linea}: $d = \frac{|ab \times ap|}{|ab|}$.
	\item \textbf{Segmento}: si la proyeccion de $p$ sobre $ab$ cae fuera del segmento, la distancia es al extremo mas cercano.
\end{itemize}
La razon geometrica: la distancia a una linea es el area del paralelogramo dividido por la base; la proyeccion ortogonal indica si estamos ``dentro'' o ``fuera'' del segmento.

\paragraph{Propiedades clave (invariantes).}
\begin{itemize}\setlength\itemsep{0pt}
	\item \texttt{dist2PointLine} retorna distancia \textbf{al cuadrado}.
	\item \texttt{dist2PointSegment} distingue 3 casos: $p$ antes de $a$, despues de $b$, o proyectado sobre el segmento.
	\item Si \texttt{len2 = 0} (degenerate), la distancia es $|ap|$ o $|bp|$.
	\item La proyeccion se calcula con dot products y el parametro $t \in [0, 1]$.
\end{itemize}

\paragraph{Codigo clave.}
Distancia al cuadrado a segmento:
\begin{verbatim}
long long dist2Segment(Point p, Point a, Point b) {
    long long len2 = (a-b).len2();
    if (len2 == 0) return (p-a).len2();  // degenerate
    long long t = max(0LL, min(1LL, dot(p-a, b-a) / len2));
    Point proj = {a.x + t*(b.x - a.x), a.y + t*(b.y - a.y)};
    return (p - proj).len2();
}
\end{verbatim}

\paragraph{Complejidad.}
\begin{itemize}\setlength\itemsep{0pt}
	\item $O(1)$.
	\item Constante minima.
\end{itemize}

\paragraph{Donde tocar para modificar.}
\begin{itemize}\setlength\itemsep{0pt}
	\item \textbf{Distancia (no al cuadrado)}: aniadir \texttt{sqrt} al final.
	\item \textbf{Coordenadas doubles}: aniadir \texttt{long double} y tolerancia.
	\item \textbf{3D}: aniadir coordenada $z$ y adaptar.
	\item \textbf{Multiples queries}: usar KD-tree o segment tree para acelerar.
\end{itemize}

\paragraph{Trampas y errores frecuentes.}
\begin{itemize}\setlength\itemsep{0pt}
	\item Overflow en \texttt{area2 * area2} si el area es grande; \texttt{long long} alcanza para $|x|, |y| \le 10^9$ pero no mas.
	\item Si $a = b$ (segmento degenerado), aniadir check explicito.
	\item En doubles, aniadir tolerancia \texttt{eps} para comparar.
\end{itemize}

\paragraph{Codigo.}
Ver \texttt{cpp.tex}, seccion \textit{Distance point-segment / point-line}.

\vspace{0.5em}

\subsubsection{Projection of point onto line}

\paragraph{Idea intuitiva.}
Proyeccion ortogonal de $p$ sobre la recta que pasa por $a$ y $b$. El parametro $t$ tal que la proyeccion es $a + t \cdot (b - a)$ es:
\[ t = \frac{(p - a) \cdot (b - a)}{|b - a|^2}. \]
La idea: descomponemos $\vec{ap}$ en una componente paralela a $\vec{ab}$ y una perpendicular; $t$ mide la posicion en la primera.

\paragraph{Propiedades clave (invariantes).}
\begin{itemize}\setlength\itemsep{0pt}
	\item Si $t \in [0, 1]$, la proyeccion cae sobre el \textbf{segmento}; si no, cae sobre la extension.
	\item Retorna \texttt{double} (puede tener precision issues).
	\item $t < 0$: $p$ esta ``antes'' de $a$; $t > 1$: $p$ esta ``despues'' de $b$.
\end{itemize}

\paragraph{Codigo clave.}
Proyeccion ortogonal:
\begin{verbatim}
Point projection(Point p, Point a, Point b) {
    Point ab = b - a, ap = p - a;
    double t = (double)dot(ap, ab) / ab.len2();
    return {a.x + t * ab.x, a.y + t * ab.y};
}
\end{verbatim}

\paragraph{Complejidad.}
\begin{itemize}\setlength\itemsep{0pt}
	\item $O(1)$.
	\item Constante minima (dos dot products y una division).
\end{itemize}

\paragraph{Donde tocar para modificar.}
\begin{itemize}\setlength\itemsep{0pt}
	\item \textbf{Proyeccion sobre segmento}: clipar $t$ a $[0, 1]$ antes de calcular el punto.
	\item \textbf{Distancia al segmento via proyeccion}: combinar con el modulo anterior.
	\item \textbf{Solo el parametro $t$}: si solo necesitas saber si esta dentro del segmento, retorna $t$ sin construir el punto.
\end{itemize}

\paragraph{Trampas y errores frecuentes.}
\begin{itemize}\setlength\itemsep{0pt}
	\item Si $a = b$ ($|ab| = 0$), la proyeccion no esta definida; aniadir guard.
	\item En doubles, aniadir tolerancia \texttt{eps} para comparar $t$ con 0 o 1.
	\item La division por \texttt{len2} puede ser cara si las coordenadas son grandes; usar \texttt{double} o precomputar.
\end{itemize}

\paragraph{Codigo.}
Ver \texttt{cpp.tex}, seccion \textit{Projection of point onto line}.

\vspace{0.5em}

\subsubsection{Point in polygon (ray casting)}

\paragraph{Idea intuitiva.}
Algoritmo de \textbf{ray casting}: contar cuantas veces un rayo horizontal desde $p$ cruza los bordes del poligono. Si es impar, $p$ esta dentro; si par, fuera. El snippet implementa la version clasica recorriendo aristas. La razon: cada cruce cambia el lado del rayo respecto al poligono; un numero impar de cambios significa ``terminamos dentro''.

\paragraph{Propiedades clave (invariantes).}
\begin{itemize}\setlength\itemsep{0pt}
	\item Funciona con poligonos simples (no necesariamente convexos).
	\item Las intersecciones en vertices cuentan especialmente.
	\item Para poligonos con \textbf{huecos}, sigue funcionando.
	\item La respuesta es estrictamente 0 (fuera) o 1 (dentro) para poligonos simples.
\end{itemize}

\paragraph{Codigo clave.}
El test clasico de ray casting:
\begin{verbatim}
bool pointInPolygon(Point p, const vector<Point>& poly) {
    int n = poly.size();
    bool inside = false;
    for (int i = 0, j = n - 1; i < n; j = i++) {
        if ((poly[i].y > p.y) != (poly[j].y > p.y) &&
            p.x < (poly[j].x - poly[i].x) * (p.y - poly[i].y) /
                    (poly[j].y - poly[i].y) + poly[i].x)
            inside = !inside;
    }
    return inside;
}
\end{verbatim}

\paragraph{Complejidad.}
\begin{itemize}\setlength\itemsep{0pt}
	\item $O(n)$ por query.
	\item Memoria $O(1)$ (no se almacena el rayo).
\end{itemize}

\paragraph{Donde tocar para modificar.}
\begin{itemize}\setlength\itemsep{0pt}
	\item \textbf{Winding number}: version alternativa que cuenta el numero de vueltas.
	\item \textbf{Poligono convexo}: check lineal contra las aristas (mas rapido).
	\item \textbf{Multiples queries}: precomputar o usar segment tree sobre las aristas.
	\item \textbf{Punto en frontera}: aniadir check adicional si es colineal con alguna arista.
\end{itemize}

\paragraph{Trampas y errores frecuentes.}
\begin{itemize}\setlength\itemsep{0pt}
	\item El manejo del vertice (cuando el rayo pasa por un vertice) puede causar double-counting; el snippet lo evita con la condicion \texttt{(poly[i].y > p.y) != (poly[j].y > p.y)}.
	\item Si el poligono tiene holes, el algoritmo sigue funcionando.
	\item Para poligonos no simples (auto-intersectantes), el resultado puede ser ambiguo.
\end{itemize}

\paragraph{Codigo.}
Ver \texttt{cpp.tex}, seccion \textit{Point in polygon (ray casting)}.

\vspace{0.5em}

\subsubsection{Minkowski sum (convex polygons)}

\paragraph{Idea intuitiva.}
La \textbf{suma de Minkowski} $A \oplus B$ de dos poligonos convexos es otro poligono convexo que representa $\{a + b : a \in A, b \in B\}$. Algoritmo: ordenar las aristas por angulo polar (CCW) y ``mergearlas'' como en merge sort; el resultado es un poligono CCW sin repetir el ultimo punto. La idea: el resultado tiene una arista paralela a cada arista de los originales; la suma se obtiene barriendo los angulos.

\paragraph{Propiedades clave (invariantes).}
\begin{itemize}\setlength\itemsep{0pt}
	\item Si $A, B$ son convexos con $n, m$ vertices, $A \oplus B$ tiene $\le n + m$ vertices.
	\item Util para \textbf{detectar colision}: $A$ colisiona con $B$ si $0 \in A \oplus (-B)$.
	\item Asume poligonos en orden \textbf{counter-clockwise}.
	\item El resultado es convexo (suma de convexos es convexo).
\end{itemize}

\paragraph{Codigo clave.}
La suma de Minkowski:
\begin{verbatim}
vector<Point> minkowski(const vector<Point>& A, const vector<Point>& B) {
    vector<Point> result;
    int i = 0, j = 0;
    int n = A.size(), m = B.size();
    do {
        result.push_back(A[i] + B[j]);
        long long cross = (A[(i+1)%n] - A[i]) ^ (B[(j+1)%m] - B[j]);
        if (cross >= 0) i = (i + 1) % n;
        if (cross <= 0) j = (j + 1) % m;
    } while (i || j);
    return result;
}
\end{verbatim}

\paragraph{Complejidad.}
\begin{itemize}\setlength\itemsep{0pt}
	\item $O(n + m)$.
	\item Memoria $O(n + m)$.
\end{itemize}

\paragraph{Donde tocar para modificar.}
\begin{itemize}\setlength\itemsep{0pt}
	\item \textbf{Deteccion de colision}: aniadir $(-B)$ (reflejar) y verificar si $0$ esta en el resultado.
	\item \textbf{Poligonos no convexos}: la suma puede no ser poligono simple; algoritmo distinto.
	\item \textbf{Multiples queries}: precomputar las aristas una vez.
	\item \textbf{Robot motion planning}: usar para planear trayectorias.
\end{itemize}

\paragraph{Trampas y errores frecuentes.}
\begin{itemize}\setlength\itemsep{0pt}
	\item Si los poligonos no son CCW, la suma es incorrecta. Normalizar antes.
	\item El bucle termina cuando ambos indices vuelven a 0; aniadir check.
	\item Para la deteccion de colision, aniadir el test \texttt{0 in A \^{} (-B)}.
\end{itemize}

\paragraph{Codigo.}
Ver \texttt{cpp.tex}, seccion \textit{Minkowski sum (convex polygons)}.

\subsection{Herramientas}

\subsubsection{Li Chao Tree (long long)}

\paragraph{Idea intuitiva.}
Mantiene un \textbf{lower envelope} (o upper) de un conjunto de rectas $y = mx + b$ sobre un dominio $[X_{LO}, X_{HI}]$ y responde ``minimo $y$ en un $x$ dado'' en $O(\log X)$. La estructura es un segment tree donde cada nodo guarda la recta ``mejor'' para su segmento medio; al aniadir una recta, se compara con la actual y se decide cual es mejor en cada mitad, descendiendo recursivamente. La demostracion: en cada nivel, solo se guarda la mejor recta para el ``punto medio''; las otras se transmiten a los hijos.

\paragraph{Propiedades clave (invariantes).}
\begin{itemize}\setlength\itemsep{0pt}
	\item \textbf{Online}: cada \texttt{addLine} en $O(\log X)$.
	\item Solo insercion; no hay borrar (para borrar, usar Li Chao segmentado con reset).
	\item La recta guardada en un nodo es la ``mejor'' en su punto medio.
	\item Las rectas ``perdedoras'' se transmiten a los hijos para futuros inserts.
\end{itemize}

\paragraph{Codigo clave.}
La insercion en Li Chao:
\begin{verbatim}
void addLine(Line nw, int v = 1, int l = X_LO, int r = X_HI) {
    int m = (l + r) / 2;
    bool left = nw.get(l) < tree[v].get(l);
    bool mid = nw.get(m) < tree[v].get(m);
    if (mid) swap(tree[v], nw);
    if (r - l == 1) return;
    if (left != mid) addLine(nw, v*2, l, m);
    else             addLine(nw, v*2+1, m, r);
}
\end{verbatim}

\paragraph{Complejidad.}
\begin{itemize}\setlength\itemsep{0pt}
	\item $O(\log X)$ por insercion y query, donde $X$ es el tamano del dominio.
	\item Memoria $O(\log X)$ por insercion, $O(X)$ total.
\end{itemize}

\paragraph{Donde tocar para modificar.}
\begin{itemize}\setlength\itemsep{0pt}
	\item \textbf{Upper envelope en vez de lower}: invertir el signo (\texttt{>} en vez de \texttt{<}).
	\item \textbf{Queries de segmento} (no full line): si la recta solo esta activa en un rango, aniadir bandera \texttt{active}.
	\item \textbf{Eliminar rectas}: usar Li Chao segmentado (mas complejo).
	\item \textbf{Pendientes grandes}: si las pendientes tienen magnitud $> 10^9$, aniadir \texttt{\_\_int128} para evitar overflow.
\end{itemize}

\paragraph{Trampas y errores frecuentes.}
\begin{itemize}\setlength\itemsep{0pt}
	\item El parametro \texttt{lo == hi} debe retornar el valor del nodo (caso base).
	\item La comparacion \texttt{newLine.get(mid) < node->line.get(mid)} decide quien es mejor en el medio.
	\item Si las rectas tienen $m$ iguales pero $b$ distintos, se queda con la primera (tie-break).
\end{itemize}

\paragraph{Codigo.}
Ver \texttt{cpp.tex}, seccion \textit{Li Chao Tree (long long)}.

\vspace{0.5em}

\subsubsection{Sweep line template}

\paragraph{Idea intuitiva.}
Patron de diseno para problemas geometricos donde se barre una linea sobre el plano y se mantiene una estructura de datos sobre el otro eje. Eventos (add, query, remove) se ordenan por la coordenada de barrido y se procesan secuencialmente.

\paragraph{Propiedades clave (invariantes).}
\begin{itemize}\setlength\itemsep{0pt}
	\item Eventos: tuplas \texttt{(x, type, id)}.
	\item La estructura auxiliar (set, multiset, segtree) se mantiene sobre el otro eje.
	\item Procesar eventos en orden de $x$ ascendente.
\end{itemize}

\paragraph{Complejidad.}
\begin{itemize}\setlength\itemsep{0pt}
	\item Depende del problema; tipico $O((N + E) \log N)$.
\end{itemize}

\paragraph{Donde tocar para modificar.}
\begin{itemize}\setlength\itemsep{0pt}
	\item \textbf{Aniadir/quitar segmentos}: aniadir eventos add/remove al barrido.
	\item \textbf{Aniadir queries}: aniadir eventos query.
	\item \textbf{Cambiar el eje de barrido}: si barres en $y$ en vez de $x$, aniadir comparador.
\end{itemize}

\paragraph{Trampas y errores frecuentes.}
\begin{itemize}\setlength\itemsep{0pt}
	\item El snippet es solo el \textbf{esqueleto}; los handlers \texttt{add/query/remove} se llenan por problema.
	\item Cuidado con eventos duplicados.
\end{itemize}

\paragraph{Codigo.}
Ver \texttt{cpp.tex}, seccion \textit{Sweep line template}.

\vspace{0.5em}

\subsubsection{Closest pair of points}

\paragraph{Idea intuitiva.}
Encuentra el par de puntos mas cercano en $O(n \log n)$ con \textbf{divide y venceras}. Algoritmo:
\begin{itemize}\setlength\itemsep{0pt}
	\item Sort por $x$.
	\item Dividir en mitades; recursivo.
	\item Combinar: dado el minimo $d$ de las dos mitades, strip = puntos en $[midX - d, midX + d]$; ordenar por $y$; cada punto solo se compara con los siguientes hasta $y$ difiera en $< d$ (maximo 7).
\end{itemize}

\paragraph{Propiedades clave (invariantes).}
\begin{itemize}\setlength\itemsep{0pt}
	\item Requiere ordenar por $y$ despues (en el snippet se hace durante el merge).
	\item Strip tiene a lo sumo $O(n)$ puntos en total, pero la comparacion es $O(1)$ amortizada por punto.
	\item Resultado en \texttt{double}.
\end{itemize}

\paragraph{Complejidad.}
\begin{itemize}\setlength\itemsep{0pt}
	\item $O(n \log n)$.
\end{itemize}

\paragraph{Donde tocar para modificar.}
\begin{itemize}\setlength\itemsep{0pt}
	\item \textbf{Reconstruir el par}: aniadir indices en la recursion.
	\item \textbf{Distancia maxima}: cambiar \texttt{min} por \texttt{max} (mas facil; no requiere D\&C).
	\item \textbf{Puntos duplicados}: aniadir check especial (distancia 0).
\end{itemize}

\paragraph{Trampas y errores frecuentes.}
\begin{itemize}\setlength\itemsep{0pt}
	\item El merge requiere ordenar por $y$; el snippet lo hace in-place con un \texttt{tmp}.
	\item Si se hace en cada recursion, es $O(n \log n)$ extra.
\end{itemize}

\paragraph{Codigo.}
Ver \texttt{cpp.tex}, seccion \textit{Closest pair of points}.

\vspace{0.5em}

\subsubsection{Rotating calipers (polygon diameter)}

\paragraph{Idea intuitiva.}
Encuentra el \textbf{diametro} de un poligono convexo (par de vertices mas distantes) en $O(n)$. Asume poligono en orden counter-clockwise. Algoritmo: dos indices $i$ (sobre las aristas) y $k$ (sobre los vertices opuestos); en cada paso, rotar $k$ mientras el area del paralelogramo aumenta. La distancia $i$-$k$ puede ser maxima en cualquier momento del barrido.

\paragraph{Propiedades clave (invariantes).}
\begin{itemize}\setlength\itemsep{0pt}
	\item Cada indice avanza como maximo $n$ veces total.
	\item El ``diametro'' es la maxima distancia entre dos vertices.
	\item Retorna \textbf{distancia al cuadrado}.
\end{itemize}

\paragraph{Complejidad.}
\begin{itemize}\setlength\itemsep{0pt}
	\item $O(n)$.
\end{itemize}

\paragraph{Donde tocar para modificar.}
\begin{itemize}\setlength\itemsep{0pt}
	\item \textbf{Reconstruir el par de vertices}: aniadir indices.
	\item \textbf{Anchura minima / maxima}: variantes de calipers.
	\item \textbf{Distancia minima entre poligonos convexos}: otra variante.
\end{itemize}

\paragraph{Trampas y errores frecuentes.}
\begin{itemize}\setlength\itemsep{0pt}
	\item Asume poligono \textbf{estrictamente} convexo (sin vertices colineales). Si hay colineales, aniadir perturbacion o filtrar.
\end{itemize}

\paragraph{Codigo.}
Ver \texttt{cpp.tex}, seccion \textit{Rotating calipers (polygon diameter)}.

% ============================================================
\section{Miscelanea}
% ============================================================

\subsection{Utilidades del usuario}

\subsubsection{Hora}

\paragraph{Idea intuitiva.}
Una pequena clase para representar \textbf{tiempos en segundos} y hacer aritmetica basica sobre ellos. Internamente almacena todo en segundos (un solo \texttt{int}), lo cual simplifica comparaciones y diferencias. El constructor acepta $(h, m, s)$ y los convierte a segundos. Hay un metodo \texttt{cut()} que fuerza un minimo (util para ``horas trabajadas'').

\paragraph{Propiedades clave (invariantes).}
\begin{itemize}\setlength\itemsep{0pt}
	\item Representacion canonica: \textbf{solo segundos}. No hay zonas horarias, no hay DST.
	\item Comparadores \texttt{<} y \texttt{>} operan sobre los segundos.
	\item \texttt{operator-} retorna la \textbf{diferencia absoluta} en segundos (siempre no negativa).
	\item \texttt{cut()} aplica un piso (por defecto $8\text{h }30\text{m} = 30600$ segundos).
\end{itemize}

\paragraph{Complejidad.}
\begin{itemize}\setlength\itemsep{0pt}
	\item $O(1)$ por operacion.
\end{itemize}

\paragraph{Donde tocar para modificar.}
\begin{itemize}\setlength\itemsep{0pt}
	\item \textbf{Aniadir suma}: \texttt{operator+} que devuelve un \texttt{Hour} con la suma de segundos.
	\item \textbf{Aniadir output formateado}: aniadir \texttt{show\_hms()} que imprime $H$:$M$:$S$.
	\item \textbf{Cambiar el ``corte''}: modificar el valor magico \texttt{30'600} en \texttt{cut()} o pasarlo como parametro.
	\item \textbf{Tamano maximo}: si los tiempos exceden $\sim 68$ anos, el \texttt{int} desborda; usar \texttt{long long}.
	\item \textbf{Precision sub-segundo}: aniadir \texttt{ms} (milisegundos) y guardar en \texttt{long long} con factor 1000.
\end{itemize}

\paragraph{Trampas y errores frecuentes.}
\begin{itemize}\setlength\itemsep{0pt}
	\item El operador \texttt{-} retorna \textbf{valor absoluto}, no diferencia con signo. Si necesitas diferencia con signo, aniadir un nuevo metodo o sobrecargar.
	\item \texttt{cut()} modifica \texttt{s} \textbf{in-place}; si el llamador no espera esto, aniadir una version \texttt{cut(min)} que retorna un nuevo \texttt{Hour}.
	\item El umbral magico \texttt{30'600} (8h 30m) esta hardcoded; documentar bien que es ``jornada minima''.
\end{itemize}

\paragraph{Codigo.}
Ver \texttt{cpp.tex}, seccion \textit{Hora}.

\end{document}
', aniadir un puntero a ambas raices.
\end{itemize}

\paragraph{Codigo.}
Ver \texttt{cpp.tex}, seccion \textit{Segment Tree persistente}.

\vspace{0.5em}

\subsubsection{Segment Tree iterativo}

\paragraph{Idea intuitiva.}
Representacion \textbf{sin recursion}: los nodos se guardan en un \texttt{vector} de tamano $2n$. Los hijos del nodo $i$ son $2i$ y $2i+1$. Las hojas ocupan las posiciones $n, n+1, \dots, 2n-1$. Update: subir el cambio desde la hoja hasta la raiz. Query: subir dos indices $l$ y $r$ hasta que se ``crucen''. La ventaja principal: localidad de memoria (los nodos estan en posiciones contiguas del vector), que mejora el cache miss rate y la constante practica. La desventaja: implementar lazy propagation es mas complejo.

\paragraph{Propiedades clave (invariantes).}
\begin{itemize}\setlength\itemsep{0pt}
	\item $O(\log n)$ por operacion, \textbf{mas rapido} en practica que el recursivo (mejor uso de cache, no hay overhead de llamadas).
	\item Build simple: copiar las hojas y combinar de abajo hacia arriba.
	\item Las hojas tienen indices en $[n, 2n)$; los nodos internos en $[1, n)$.
\end{itemize}

\paragraph{Codigo clave.}
El bucle de query ascendente:
\begin{verbatim}
long long query(int l, int r) {  // [l, r] inclusivo
    long long res = 0;
    for (l += n, r += n + 1; l < r; l >>= 1, r >>= 1) {
        if (l & 1) res += st[l++];
        if (r & 1) res += st[--r];
    }
    return res;
}
\end{verbatim}

\paragraph{Complejidad.}
\begin{itemize}\setlength\itemsep{0pt}
	\item $O(\log n)$ por update/query, $O(n)$ build.
	\item En la practica, $\sim 1.5x$ mas rapido que el recursivo para $n \ge 10^5$.
\end{itemize}

\paragraph{Donde tocar para modificar.}
\begin{itemize}\setlength\itemsep{0pt}
	\item \textbf{Suma -> max/min/xor}: cambiar la operacion en \texttt{set} y \texttt{query}.
	\item \textbf{Aniadir update en rango con lazy}: implementar version lazy por separado (no trivial en iterativo).
	\item \textbf{Tamano no potencia de 2}: si $n$ no es potencia de 2, aniadir padding.
	\item \textbf{Persistencia}: usar \texttt{vector<array<int, 2>>} y duplicar nodos.
\end{itemize}

\paragraph{Trampas y errores frecuentes.}
\begin{itemize}\setlength\itemsep{0pt}
	\item En el query, \texttt{r += n+1} (no \texttt{r += n}) porque \texttt{query(l, r)} es \textbf{inclusivo} en ambos extremos.
	\item El indice \texttt{l & 1} detecta si \texttt{l} es hijo derecho; aniadir \texttt{res += st[l++]} en ese caso.
	\item El \texttt{--r} en la condicion derecha es para incluir el limite.
	\item Si las operaciones no son idempotentes, cuidado con el orden.
\end{itemize}

\paragraph{Codigo.}
Ver \texttt{cpp.tex}, seccion \textit{Segment Tree iterativo}.

\subsection{Otras}

\subsubsection{MO}

\paragraph{Idea intuitiva.}
Algoritmo de Mo ``puro'' (sin updates): queries offline ordenadas por $(L/\sqrt{n}, R)$ (con $R$ alternando entre ascendente y descendente segun el bloque). Mover $L$ o $R$ entre queries consecutivas es $O(1)$; el total es $O((N + Q) \sqrt{n})$. La intuicion: al reordenar las queries, los punteros $L$ y $R$ se mueven ``lo menos posible''; cada elemento del array se visita $O(\sqrt{n} \cdot Q/N)$ veces.

\paragraph{Propiedades clave (invariantes).}
\begin{itemize}\setlength\itemsep{0pt}
	\item \texttt{Q} guarda \texttt{(l, r, idx)}.
	\item El sort prioriza bloque de $L$; dentro del bloque, alterna $R$ ascendente/descendente para mejorar la locality.
	\item Tres punteros \texttt{l}, \texttt{r} que se mueven paso a paso.
	\item Tamano de bloque optimo: $\sim N / \sqrt{Q}$.
\end{itemize}

\paragraph{Codigo clave.}
El sort clasico de Mo con alternancia:
\begin{verbatim}
int block = (int)sqrt(n);
sort(queries.begin(), queries.end(), [&](const Q& a, const Q& b) {
    int blockA = a.l / block, blockB = b.l / block;
    if (blockA != blockB) return blockA < blockB;
    return (blockA & 1) ? (a.r > b.r) : (a.r < b.r);
});
\end{verbatim}

\paragraph{Complejidad.}
\begin{itemize}\setlength\itemsep{0pt}
	\item $O((N + Q) \sqrt{n})$.
	\item Con Hilbert order, $\sim 30\%$ mas rapido en practica.
\end{itemize}

\paragraph{Donde tocar para modificar.}
\begin{itemize}\setlength\itemsep{0pt}
	\item \textbf{Hilbert order}: alternativa al sort clasico; da mejor locality y es mas rapido en practica.
	\item \textbf{Aniadir la operacion concreta}: el snippet tiene \texttt{// add} y \texttt{// delete} como comentarios; reemplazarlos con la logica del problema (contar distintos, suma de frecuencias, etc.).
	\item \textbf{Bloque optimo}: calcular \texttt{int block = max(1, (int)(n / sqrt(q)));}.
\end{itemize}

\paragraph{Trampas y errores frecuentes.}
\begin{itemize}\setlength\itemsep{0pt}
	\item Los indices en el sort son \textbf{0-based} (\texttt{q[i].l = l - 1}), consistente con el resto del codigo.
	\item \texttt{add} y \texttt{remove} deben ser \textbf{inversos exactos}; sino, las queries dan respuestas incorrectas.
	\item Para queries cerradas $[l, r]$, Mo suele usar $[l, r+1)$ (half-open); aniadir o restar 1 consistentemente.
\end{itemize}

\paragraph{Codigo.}
Ver \texttt{cpp.tex}, seccion \textit{MO}.

\vspace{0.5em}

\subsubsection{DSU con rollback}

\paragraph{Idea intuitiva.}
Variante del DSU donde cada \texttt{unite} \textbf{se puede deshacer} en orden LIFO. La idea: en vez de path compression (que es irreversible), solo se usa union by size. Cada \texttt{unite} guarda un par \texttt{(parent\_changed, old\_size)} en una pila. \texttt{rollback()} desapila hasta el ultimo ``snapshot'' (marcador) restaurando los valores. La razon: path compression modifica la estructura de multiples nodos, no solo uno, asi que es dificil de deshacer; union by size solo modifica dos nodos (el padre del menor y su tamano), asi que es trivial deshacer.

\paragraph{Propiedades clave (invariantes).}
\begin{itemize}\setlength\itemsep{0pt}
	\item \textbf{Sin path compression}: cada \texttt{find} es $O(\log n)$ (peor caso del union by size) en el peor caso, pero amortizado $O(\log n)$.
	\item \textbf{Snapshots}: \texttt{snapshot()} marca un punto; \texttt{rollback()} deshace hasta ahi.
	\item Cada \texttt{unite} exitoso apila \textbf{dos} pares (uno por cada cambio).
	\item La pila de historial crece monotona; \texttt{rollback} la reduce.
\end{itemize}

\paragraph{Codigo clave.}
El unite con guardado de historial:
\begin{verbatim}
void unite(int a, int b) {
    a = find(a); b = find(b);
    if (a == b) { history.push({-1, -1}); return; }
    if (size[a] < size[b]) swap(a, b);
    history.push({b, size[a]});
    parent[b] = a;
    size[a] += size[b];
    components--;
}
void rollback() {
    auto [b, sz] = history.top(); history.pop();
    if (b == -1) return;
    size[parent[b]] = sz;
    parent[b] = b;
    components++;
}
\end{verbatim}

\paragraph{Complejidad.}
\begin{itemize}\setlength\itemsep{0pt}
	\item $O(\log n)$ por \texttt{unite}/\texttt{find}, $O(1)$ por rollback de una operacion.
	\item Sin path compression: peor caso $O(\log n)$ por find.
\end{itemize}

\paragraph{Donde tocar para modificar.}
\begin{itemize}\setlength\itemsep{0pt}
	\item \textbf{Rollback parcial} (no todo): implementar snapshots parciales guardando el tamano del historial.
	\item \textbf{Mas informacion por conjunto}: aniadir mas campos al struct y guardarlos en \texttt{history}.
	\item \textbf{Persistente}: aniadir versionado completo (clonar el struct en cada cambio).
	\item \textbf{DSU de Arpa}: variante con $O(\log n)$ amortizado y rollback en $O(1)$.
\end{itemize}

\paragraph{Trampas y errores frecuentes.}
\begin{itemize}\setlength\itemsep{0pt}
	\item El snippet guarda dos pares por \texttt{unite}; al rollback \textbf{desapilar en orden correcto} (no invertir).
	\item Olvidar el \texttt{components--} en \texttt{unite} o el \texttt{components++} en rollback deja el contador mal.
	\item El sentinela \texttt{\{-1, -1\}} indica un unite trivial; importante para mantener consistencia.
	\item \texttt{find} puede ser $O(\log n)$ peor caso sin path compression; verificar que es OK para el problema.
\end{itemize}

\paragraph{Codigo.}
Ver \texttt{cpp.tex}, seccion \textit{DSU con rollback}.

\vspace{0.5em}

\subsubsection{Segment Tree 2D}

\paragraph{Idea intuitiva.}
Extender el segment tree a 2 dimensiones: la raiz cubre toda la matriz $[0, n) \times [0, m)$. Cada nodo es un \textbf{segment tree 1D} sobre la segunda dimension. Para update en $(x, y)$, se actualizan $O(\log n)$ nodos externos (uno por nivel vertical) y dentro de cada uno, $O(\log m)$ nodos internos. Query sobre un rectangulo: $O(\log n \cdot \log m)$. La estructura matematica: el arbol de primer nivel divide el espacio vertical; en cada nivel, los nodos internos gestionan el espacio horizontal. Asi, se reduce la consulta rectangular a $O(\log n \cdot \log m)$ consultas 1D parciales.

\paragraph{Propiedades clave (invariantes).}
\begin{itemize}\setlength\itemsep{0pt}
	\item Memoria: $O(n \cdot m)$ en el peor caso (matriz densa).
	\item Para matrices grandes con valores \textbf{cero} y unos pocos updates, se puede comprimir.
	\item Cada nodo externo tiene $4m$ nodos internos.
\end{itemize}

\paragraph{Codigo clave.}
Update en $(x, y)$ con recursion en cascada:
\begin{verbatim}
void updateX(int nodeX, int lx, int rx, int x, int y, int v) {
    if (lx != rx) {
        int mx = (lx + rx) / 2;
        if (x <= mx) updateX(nodeX*2, lx, mx, x, y, v);
        else         updateX(nodeX*2+1, mx+1, rx, x, y, v);
    }
    updateY(nodeX, lx, rx, y, v);
}
\end{verbatim}

\paragraph{Complejidad.}
\begin{itemize}\setlength\itemsep{0pt}
	\item Update y query $O(\log n \cdot \log m)$, memoria $O(4n \cdot 4m)$.
	\item Para $n = m = 1000$: $16 \cdot 10^6$ nodos ($\sim 64$MB con \texttt{int}).
\end{itemize}

\paragraph{Donde tocar para modificar.}
\begin{itemize}\setlength\itemsep{0pt}
	\item \textbf{Build desde una matriz}: aniadir constructor que tome \texttt{vector<vector<int>>}.
	\item \textbf{Lazy propagation 2D}: aniadir un campo \texttt{lazy} por nodo.
	\item \textbf{Suma + maximo simultaneos}: aniadir mas campos.
	\item \textbf{Sparse 2D}: si la mayoria son ceros, usar map en vez de vector.
\end{itemize}

\paragraph{Trampas y errores frecuentes.}
\begin{itemize}\setlength\itemsep{0pt}
	\item La recursion puede ser profunda; aniadir iterativo si $n, m > 10^5$.
	\item Memoria $O(4n \cdot 4m) = 16nm$; si $n = m = 10^3$, son $16M$ nodos; con $n = m = 10^4$, son $1.6 \cdot 10^9$, no factible.
	\item Olvidar propagar en cascada da updates incorrectos.
	\item Si la mayoria de queries son rectangulares grandes, considerar KD-tree u otra estructura.
\end{itemize}

\paragraph{Codigo.}
Ver \texttt{cpp.tex}, seccion \textit{Segment Tree 2D}.

\vspace{0.5em}

\subsubsection{Segment Tree persistente}

\paragraph{Idea intuitiva.}
Cada \texttt{update} crea una \textbf{nueva raiz} sin modificar los nodos no afectados. Asi puedes tener \textbf{multiples versiones} del arreglo y consultar cada una en $O(\log n)$. Memoria: $O(N \log N)$ nodos en total para $N$ updates. La tecnica se llama ``path copying'': solo el camino de la raiz a la hoja modificada se copia; el resto del arbol se comparte entre versiones. Asi, $N$ updates producen $\le N \log n$ nodos en total.

\paragraph{Propiedades clave (invariantes).}
\begin{itemize}\setlength\itemsep{0pt}
	\item Cada nodo es \textbf{inmutable} despues de creado; las versiones comparten estructura.
	\item \texttt{pUpdate} retorna la nueva raiz; la anterior sigue accesible.
	\item \texttt{pQuery} toma una raiz especifica (no la global).
	\item Las versiones viven independientemente; borrar una no afecta a otras.
\end{itemize}

\paragraph{Codigo clave.}
El corazon del persistent update:
\begin{verbatim}
PNode* pUpdate(PNode* v, int l, int r, int pos, int val) {
    PNode* u = newNode(v);  // copia con misma info
    if (l == r) { u->sum = val; return u; }
    int m = (l + r) / 2;
    if (pos <= m) u->left = pUpdate(v->left, l, m, pos, val);
    else          u->right = pUpdate(v->right, m+1, r, pos, val);
    u->sum = u->left->sum + u->right->sum;
    return u;
}
// versions.push_back(pUpdate(versions.back(), 0, n-1, pos, val));
\end{verbatim}

\paragraph{Complejidad.}
\begin{itemize}\setlength\itemsep{0pt}
	\item Update $O(\log n)$ nodos nuevos, query $O(\log n)$.
	\item Memoria $O(N \log n)$ para $N$ updates.
\end{itemize}

\paragraph{Donde tocar para modificar.}
\begin{itemize}\setlength\itemsep{0pt}
	\item \textbf{Persistencia completa (cada update crea un nodo)}: aniadir \texttt{version[v]} para acceder a cada version.
	\item \textbf{Cambiar suma por otro operacion}: ajustar \texttt{pUpdate} y \texttt{pQuery}.
	\item \textbf{Memoria grande}: en contests, el allocator por defecto puede ser lento; usar \texttt{vector<PNode>} con \texttt{reserve(2*N*logN)} y un \texttt{int nextNode = 0} para asignar en $O(1)$.
	\item \textbf{Lazy persistence}: aniadir campos \texttt{lazy} en cada nodo; los hijos ``ven'' el lazy a traves de la copia.
\end{itemize}

\paragraph{Trampas y errores frecuentes.}
\begin{itemize}\setlength\itemsep{0pt}
	\item Olvidar pasar la \textbf{vieja raiz} al query, usando la nueva por error.
	\item Si guardas multiples versiones, asegurate de pasar la raiz correcta a cada query.
	\item Memory leak en produccion (no en contests) si no liberas los nodos.
	\item Para queries ``diferencia entre version $a$ y version $b

\vspace{0.5em}

\subsubsection{Segment Tree iterativo}

\paragraph{Idea intuitiva.}
Representacion \textbf{sin recursion}: los nodos se guardan en un \texttt{vector} de tamano $2n$. Los hijos del nodo $i$ son $2i$ y $2i+1$. Las hojas ocupan las posiciones $n, n+1, \dots, 2n-1$. Update: subir el cambio desde la hoja hasta la raiz. Query: subir dos indices $l$ y $r$ hasta que se ``crucen''. La ventaja principal: localidad de memoria (los nodos estan en posiciones contiguas del vector), que mejora el cache miss rate y la constante practica. La desventaja: implementar lazy propagation es mas complejo.

\paragraph{Propiedades clave (invariantes).}
\begin{itemize}\setlength\itemsep{0pt}
	\item $O(\log n)$ por operacion, \textbf{mas rapido} en practica que el recursivo (mejor uso de cache, no hay overhead de llamadas).
	\item Build simple: copiar las hojas y combinar de abajo hacia arriba.
	\item Las hojas tienen indices en $[n, 2n)$; los nodos internos en $[1, n)$.
\end{itemize}

\paragraph{Codigo clave.}
El bucle de query ascendente:
\begin{verbatim}
long long query(int l, int r) {  // [l, r] inclusivo
    long long res = 0;
    for (l += n, r += n + 1; l < r; l >>= 1, r >>= 1) {
        if (l & 1) res += st[l++];
        if (r & 1) res += st[--r];
    }
    return res;
}
\end{verbatim}

\paragraph{Complejidad.}
\begin{itemize}\setlength\itemsep{0pt}
	\item $O(\log n)$ por update/query, $O(n)$ build.
	\item En la practica, $\sim 1.5x$ mas rapido que el recursivo para $n \ge 10^5$.
\end{itemize}

\paragraph{Donde tocar para modificar.}
\begin{itemize}\setlength\itemsep{0pt}
	\item \textbf{Suma -> max/min/xor}: cambiar la operacion en \texttt{set} y \texttt{query}.
	\item \textbf{Aniadir update en rango con lazy}: implementar version lazy por separado (no trivial en iterativo).
	\item \textbf{Tamano no potencia de 2}: si $n$ no es potencia de 2, aniadir padding.
	\item \textbf{Persistencia}: usar \texttt{vector<array<int, 2>>} y duplicar nodos.
\end{itemize}

\paragraph{Trampas y errores frecuentes.}
\begin{itemize}\setlength\itemsep{0pt}
	\item En el query, \texttt{r += n+1} (no \texttt{r += n}) porque \texttt{query(l, r)} es \textbf{inclusivo} en ambos extremos.
	\item El indice \texttt{l & 1} detecta si \texttt{l} es hijo derecho; aniadir \texttt{res += st[l++]} en ese caso.
	\item El \texttt{--r} en la condicion derecha es para incluir el limite.
	\item Si las operaciones no son idempotentes, cuidado con el orden.
\end{itemize}

\paragraph{Codigo.}
Ver \texttt{cpp.tex}, seccion \textit{Segment Tree iterativo}.

\vspace{0.5em}

\subsection{Fenwick / BIT}

\subsubsection{BIT point update + prefix sum}

\paragraph{Idea intuitiva.}
El \textbf{Fenwick Tree} (BIT) almacena sums parciales de manera que \texttt{sum(i)} (suma de $[1, i]$) y \texttt{add(i, delta)} cuestan $O(\log n)$. La representacion es ingeniosa: el nodo $i$ cubre el rango $[i - \text{lowbit}(i) + 1, i]$. Asi, \texttt{sum(i)} itera \texttt{i -= lowbit(i)}, sumando los nodos que cubren un prefijo creciente, y \texttt{add(i, delta)} itera \texttt{i += lowbit(i)}, propagando el cambio a todos los nodos cuyo rango contiene a $i$. La idea: cada indice $i$ se descompone en una suma de bloques disjuntos de tamano potencias de 2 (su ``lowbit representation'').

\paragraph{Propiedades clave (invariantes).}
\begin{itemize}\setlength\itemsep{0pt}
	\item Indice \textbf{1-based} (no usar $i=0$).
	\item \texttt{lowbit(i) = i \& (-i)}.
	\item Cada nodo cubre exactamente un rango de tamano \texttt{lowbit(i)}.
	\item La operacion debe ser \textbf{invertible} (suma, xor, no min/max).
\end{itemize}

\paragraph{Codigo clave.}
Las dos operaciones basicas:
\begin{verbatim}
void add(int i, long long d) {
    for (; i <= n; i += i & -i) bit[i] += d;
}
long long sum(int i) {
    long long s = 0;
    for (; i > 0; i -= i & -i) s += bit[i];
    return s;
}
\end{verbatim}

\paragraph{Complejidad.}
\begin{itemize}\setlength\itemsep{0pt}
	\item $O(\log n)$ por operacion. Memoria $O(n)$.
	\item Constante muy pequena ($\sim 1$): en la practica, $\sim 3x$ mas rapido que segment tree.
\end{itemize}

\paragraph{Donde tocar para modificar.}
\begin{itemize}\setlength\itemsep{0pt}
	\item \textbf{Suma -> max/min/xor}: cambiar el \texttt{+=} por la operacion (con identidad apropriada). \textbf{No} funciona para \texttt{min} de forma natural (no es invertible).
	\item \textbf{Tamano dinamico}: aniadir un metodo \texttt{resize(newSize)} que preserve valores.
	\item \textbf{Tamano grande}: con $n = 10^7$, mejor comprimir coordenadas.
	\item \textbf{Operador no conmutativo}: usar segment tree en su lugar.
\end{itemize}

\paragraph{Trampas y errores frecuentes.}
\begin{itemize}\setlength\itemsep{0pt}
	\item Indices 0 vs 1: el BIT es \textbf{1-based}. Si el problema da indices 0-based, aniadir \texttt{+1} en cada llamada.
	\item El metodo \texttt{sum} itera \texttt{i -= i \& -i}; aniadir o quitar 1 rompe el invariante.
	\item Si el tipo es \texttt{int}, posibles overflows; usar \texttt{long long}.
	\item Para rango $[l, r]$, el query es \texttt{sum(r) - sum(l-1)} (cuidado con $l = 1$).
\end{itemize}

\paragraph{Codigo.}
Ver \texttt{cpp.tex}, seccion \textit{BIT point update + prefix sum}.

\vspace{0.5em}

\subsubsection{BIT range add + point query}

\paragraph{Idea intuitiva.}
Truco de \textbf{diferencias}: para sumar $d$ a $[l, r]$, aniadir $d$ en $l$ y $-d$ en $r+1$. Asi, el valor en $i$ es la suma de los prefijos hasta $i$, lo que se calcula con un BIT normal. La idea algebraica: si mantenemos un arreglo de diferencias $\Delta$ donde $a_i = \sum_{j \le i} \Delta_j$, entonces aniadir $d$ a $a[l..r]$ equivale a $\Delta_l \mathrel{+}= d$ y $\Delta_{r+1} \mathrel{-}= d$, y reconstruir $a$ es un prefix sum.

\paragraph{Propiedades clave (invariantes).}
\begin{itemize}\setlength\itemsep{0pt}
	\item Internamente es un BIT 5.1 con la convencion de diferencias.
	\item \texttt{rangeAdd(l, r, delta)} cuesta $O(\log n)$ (dos \texttt{add}).
	\item \texttt{pointQuery(i)} cuesta $O(\log n)$.
	\item Si $r+1 > n$, no se aniade el $-d$ (fuera del rango).
\end{itemize}

\paragraph{Codigo clave.}
El truco de diferencias:
\begin{verbatim}
void rangeAdd(int l, int r, long long d) {
    add(l, d);
    if (r + 1 <= n) add(r + 1, -d);
}
long long pointQuery(int i) { return sum(i); }
\end{verbatim}

\paragraph{Complejidad.}
\begin{itemize}\setlength\itemsep{0pt}
	\item $O(\log n)$ por operacion.
	\item Memoria $O(n)$.
\end{itemize}

\paragraph{Donde tocar para modificar.}
\begin{itemize}\setlength\itemsep{0pt}
	\item \textbf{Persistencia}: aniadir versionado (cada update crea una nueva ``copia'' del estado; en BIT no es trivial, mejor usar persistent segtree).
	\item \textbf{Coordenada compression}: si $n$ es grande pero los updates son pocos, comprimir.
	\item \textbf{2D range add + point query}: aniadir una dimension con la misma tecnica.
\end{itemize}

\paragraph{Trampas y errores frecuentes.}
\begin{itemize}\setlength\itemsep{0pt}
	\item Si $r+1 > n$, \textbf{no} aniadir \texttt{-delta} (no hay efecto fuera del rango); el snippet ya lo chequea.
	\item Para reconstruir el array original tras $k$ range adds, hacer $k$ point queries (no hay shortcut).
	\item Confundir ``suma de diferencias'' con ``suma acumulada''; son cosas distintas.
\end{itemize}

\paragraph{Codigo.}
Ver \texttt{cpp.tex}, seccion \textit{BIT range add + point query}.

\vspace{0.5em}

\subsubsection{BIT range add + range sum}

\paragraph{Idea intuitiva.}
Para queries $\sum_{i=1}^{x} a_i$ despues de varios range adds, el truco de diferencias no basta (necesitas la suma acumulada de las diferencias). Solucion: dos BITs, $t_1$ y $t_2$. Tras aniadir $d$ a $[l, r]$: $t_1$ recibe $+d$ en $l$, $-d$ en $r+1$; $t_2$ recibe $+d(l-1)$ en $l$, $-d \cdot r$ en $r+1$. Entonces $\text{prefixSum}(x) = t_1.\text{sum}(x) \cdot x - t_2.\text{sum}(x)$. La razon matematica: el valor en $i$ es $\sum_{j \le i} \Delta_j$, asi que la suma hasta $i$ es $\sum_{k \le i} (i - k + 1) \Delta_k$; expandir da la formula con dos BITs.

\paragraph{Propiedades clave (invariantes).}
\begin{itemize}\setlength\itemsep{0pt}
	\item Verificacion: para $x \ge r$, $\text{prefixSum}(x)$ acumula correctamente todas las contribuciones.
	\item La formula sale de expandir la suma de diferencias.
	\item La identidad clave: $\sum_{i=1}^{x} i = x(x+1)/2$, asi que la suma de diferencias ponderadas requiere BITs adicionales.
\end{itemize}

\paragraph{Codigo clave.}
Range add con range sum (dos BITs):
\begin{verbatim}
void rangeAdd(int l, int r, long long d) {
    t1.add(l, d); t1.add(r + 1, -d);
    t2.add(l, d * (l - 1)); t2.add(r + 1, -d * r);
}
long long prefixSum(int x) {
    return t1.sum(x) * x - t2.sum(x);
}
\end{verbatim}

\paragraph{Complejidad.}
\begin{itemize}\setlength\itemsep{0pt}
	\item $O(\log n)$ por operacion (2 \texttt{add} o 2 \texttt{sum} por cada lado).
	\item Memoria $O(n)$.
\end{itemize}

\paragraph{Donde tocar para modificar.}
\begin{itemize}\setlength\itemsep{0pt}
	\item \textbf{2D range add + range sum}: aniadir otra dimension con la misma tecnica (queda 4 BITs).
	\item \textbf{Cambiar mod}: aniadir modulo en cada operacion.
	\item \textbf{Output del array}: aniadir un metodo que reconstruya $a_i$ point query por cada $i$.
\end{itemize}

\paragraph{Trampas y errores frecuentes.}
\begin{itemize}\setlength\itemsep{0pt}
	\item Overflow en \texttt{delta * (l - 1)} si los valores son grandes; usar \texttt{long long} y \texttt{\% MOD} al final.
	\item Confundir la convencion $0$-based vs $1$-based del BIT; las dos afectan las formulas.
	\item En 2D, el codigo es $\sim 4x$ mas largo y propenso a errores; verificar con casos pequenos.
\end{itemize}

\paragraph{Codigo.}
Ver \texttt{cpp.tex}, seccion \textit{BIT range add + range sum}.

\subsection{RMQ y sparse table}

\subsubsection{Sparse Table}

\paragraph{Idea intuitiva.}
Precomputa respuestas a queries de tamano $2^k$ para $k = 0, 1, \dots, \log n$. Asi, \texttt{st[k][i]} guarda la respuesta sobre $[i, i + 2^k)$. Para query $[l, r]$, se usan dos bloques de tamano $2^{\lfloor \log_2(r-l+1) \rfloor}$: $[l, l + 2^k)$ y $[r - 2^k + 1, r)$, y se combinan. Esto da $O(1)$ por query \textbf{solo si la operacion es idempotente} (min, max, gcd; \textbf{no} suma). La razon: con idempotencia, solapar los dos bloques no afecta el resultado.

\paragraph{Propiedades clave (invariantes).}
\begin{itemize}\setlength\itemsep{0pt}
	\item \texttt{st[0][i] = a[i]}.
	\item \texttt{st[k][i] = op(st[k-1][i], st[k-1][i + 2\^{}(k-1)])}.
	\item \texttt{LOG = floor(log2(n)) + 1}.
	\item La operacion debe ser \textbf{idempotente} ($f(x, x) = x$) y asociativa.
\end{itemize}

\paragraph{Codigo clave.}
Query $O(1)$ en sparse table:
\begin{verbatim}
long long query(int l, int r) {
    int k = __builtin_clz(1) - __builtin_clz(r - l + 1);  // log2
    return min(st[k][l], st[k][r - (1 << k) + 1]);
}
\end{verbatim}

\paragraph{Complejidad.}
\begin{itemize}\setlength\itemsep{0pt}
	\item Build $O(n \log n)$, query $O(1)$.
	\item Memoria $O(n \log n)$.
\end{itemize}

\paragraph{Donde tocar para modificar.}
\begin{itemize}\setlength\itemsep{0pt}
	\item \textbf{Cambiar min por otra operacion idempotente}: cambiar las dos ocurrencias de \texttt{min(...)} por la operacion (max, gcd). Verificar identidad y asociatividad.
	\item \textbf{Suma en rango}: NO sirve con este metodo; usar prefix sums.
	\item \textbf{Sparse table 2D}: aniadir otra dimension (queda $O(n^2 \log^2 n)$ memoria).
	\item \textbf{LCA queries}: combinar con binary lifting en arboles.
\end{itemize}

\paragraph{Trampas y errores frecuentes.}
\begin{itemize}\setlength\itemsep{0pt}
	\item El query toma $[l, r]$ \textbf{inclusivo} (no $[l, r)$).
	\item El \texttt{LOG} debe estar bien calculado; con $n = 1$, \texttt{LOG = 1}.
	\item Usar sparse table para suma es un error; necesitas prefix sums o segment tree.
	\item Para queries $O(1)$ con operaciones no idempotentes, usar Disjoint Sparse Table.
\end{itemize}

\paragraph{Codigo.}
Ver \texttt{cpp.tex}, seccion \textit{Sparse Table}.

\vspace{0.5em}

\subsubsection{Disjoint Sparse Table}

\paragraph{Idea intuitiva.}
Variante que tambien da $O(1)$ por query pero \textbf{no requiere idempotencia} para algunas operaciones (sirve para suma, xor, etc. en linea con la mayoria de operaciones asociativas). La idea: para cada nivel $k$, los rangos ``se parten en mitades''. \texttt{st[k][i]} cubre un rango que \textbf{empieza} en $i$ y se extiende a izquierda o derecha hasta encontrar un ``punto de division''. El truco: para query $[l, r]$, el nivel optimo es el bit mas alto donde $l$ y $r$ difieren; los rangos $[l, ?]$ y $[?, r]$ cubren exactamente el query sin solape.

\paragraph{Propiedades clave (invariantes).}
\begin{itemize}\setlength\itemsep{0pt}
	\item Construccion diferente al Sparse Table clasico: en cada nivel, se extiende desde cada posicion hacia ambos lados hasta encontrar la mitad de nivel.
	\item \texttt{query(l, r)} usa \texttt{k = log2(l \^{} r)} y combina \texttt{st[k][l]} y \texttt{st[k][r]}.
	\item La operacion debe ser \textbf{asociativa} pero no necesariamente idempotente.
\end{itemize}

\paragraph{Codigo clave.}
Query con DST:
\begin{verbatim}
long long query(int l, int r) {
    if (l == r) return st[0][l];
    int k = 31 - __builtin_clz(l ^ r);  // bit mas alto que difiere
    return op(st[k][l], st[k][r]);
}
\end{verbatim}

\paragraph{Complejidad.}
\begin{itemize}\setlength\itemsep{0pt}
	\item Build $O(n \log n)$, query $O(1)$.
	\item Memoria $O(n \log n)$.
\end{itemize}

\paragraph{Donde tocar para modificar.}
\begin{itemize}\setlength\itemsep{0pt}
	\item \textbf{Operacion no asociativa}: no funciona; el Disjoint Sparse Table asume asociatividad.
	\item \textbf{Memory compression}: con $n \le 10^6$ alcanza; con mas, comprimir.
	\item \textbf{Operaciones no conmutativas}: aniadir el cuidado de que el orden es importante.
\end{itemize}

\paragraph{Trampas y errores frecuentes.}
\begin{itemize}\setlength\itemsep{0pt}
	\item La construccion es sutil; verificar con casos pequenos antes de usar.
	\item El bit ``mas alto que difiere'' requiere cuidado con $l = r$; aniadir caso base.
	\item Si la operacion no es asociativa, los dos bloques no se pueden combinar en cualquier orden.
\end{itemize}

\paragraph{Codigo.}
Ver \texttt{cpp.tex}, seccion \textit{Disjoint Sparse Table}.

\subsection{Trees y DP en arbol}

\subsubsection{Binary Lifting / LCA}

\paragraph{Idea intuitiva.}
Precomputa para cada nodo $v$ sus ancestros a distancias $2^k$ (los ``saltos''). Asi, para subir $d$ niveles, se descompone $d$ en binario y se aplican los saltos correspondientes. El \textbf{LCA} (Lowest Common Ancestor) se obtiene subiendo el nodo mas profundo hasta igualar profundidades y luego subiendo ambos hasta que sus padres coincidan. La intuicion: cualquier entero $\le n$ se representa en binario con $\le \log n$ bits, asi que subir $d$ niveles requiere $\le \log n$ saltos.

\paragraph{Propiedades clave (invariantes).}
\begin{itemize}\setlength\itemsep{0pt}
	\item \texttt{up[k][v]} = ancestro de $v$ a $2^k$ niveles.
	\item \texttt{LOG = floor(log2(n)) + 1}.
	\item La recursion de \texttt{dfs} llena \texttt{up[0][v]} (el padre directo) y por transitividad los demas.
	\item El LCA esta definido como el ancestro comun mas profundo de $u$ y $v$.
\end{itemize}

\paragraph{Codigo clave.}
LCA con binary lifting:
\begin{verbatim}
int lca(int u, int v) {
    if (depth[u] < depth[v]) swap(u, v);
    int diff = depth[u] - depth[v];
    for (int k = LOG - 1; k >= 0; --k)
        if (diff & (1 << k)) u = up[k][u];
    if (u == v) return u;
    for (int k = LOG - 1; k >= 0; --k)
        if (up[k][u] != up[k][v])
            u = up[k][u], v = up[k][v];
    return up[0][u];
}
\end{verbatim}

\paragraph{Complejidad.}
\begin{itemize}\setlength\itemsep{0pt}
	\item Build $O(n \log n)$, query LCA $O(\log n)$.
	\item Memoria $O(n \log n)$.
\end{itemize}

\paragraph{Donde tocar para modificar.}
\begin{itemize}\setlength\itemsep{0pt}
	\item \textbf{Distancia en el arbol}: combinar con \texttt{depth[]} para calcular \texttt{depth[u] + depth[v] - 2*depth[lca]}.
	\item \textbf{K-esimo ancestro}: subir $k$ niveles con el mismo truco.
	\item \textbf{Aniadir peso a las aristas}: guardar \texttt{dist[2\^{}k][v]} (distancia acumulada a $2^k$ saltos) ademas de \texttt{up}.
	\item \textbf{Sparse table sobre depth}: alternativa $O(1)$ a LCA (mas complejo).
\end{itemize}

\paragraph{Trampas y errores frecuentes.}
\begin{itemize}\setlength\itemsep{0pt}
	\item El arbol debe ser \textbf{rooted} (el snippet asume eso); si no, hacer un \texttt{dfs} desde cualquier nodo.
	\item Para grafos no conexos, correr \texttt{dfs} desde cada componente.
	\item En la segunda fase del LCA, subir ambos nodos mientras \texttt{up[k][u] != up[k][v]}; el LCA final es \texttt{up[0][u]}.
	\item El padre de la raiz es 0 (o -1); aniadir check si se sale del rango.
\end{itemize}

\paragraph{Codigo.}
Ver \texttt{cpp.tex}, seccion \textit{Binary Lifting / LCA}.

\vspace{0.5em}

\subsubsection{Heavy-Light Decomposition (HLD)}

\paragraph{Idea intuitiva.}
Descomponer un arbol en \textbf{cadenas pesadas} (paths donde cada nodo (salvo la raiz) tiene al hijo ``pesado'' como hijo de tamano maximo de su subarbol). Cada cadena se mapea a un rango contiguo en un arreglo base; asi, queries sobre \textbf{paths} se descomponen en $O(\log n)$ queries sobre rangos contiguos (cada cadena es un rango), que se resuelven con un BIT o segment tree. La intuicion: por cada nivel de recursion, el subarbol del hijo pesado es al menos la mitad, asi que un path cruza $\le 2 \log n$ cadenas.

\paragraph{Propiedades clave (invariantes).}
\begin{itemize}\setlength\itemsep{0pt}
	\item \texttt{heavy[v]} = hijo pesado de $v$ (el de mayor subarbol).
	\item \texttt{head[v]} = cabeza de la cadena que contiene a $v$.
	\item \texttt{pos[v]} = posicion en el arreglo base.
	\item Cada nodo pertenece a exactamente una cadena.
\end{itemize}

\paragraph{Codigo clave.}
El query sobre un path $u \to v$:
\begin{verbatim}
long long pathQuery(int u, int v) {
    long long res = 0;
    while (head[u] != head[v]) {
        if (depth[head[u]] < depth[head[v]]) swap(u, v);
        res += seg.query(pos[head[u]], pos[u]);
        u = parent[head[u]];
    }
    if (depth[u] > depth[v]) swap(u, v);
    res += seg.query(pos[u], pos[v]);
    return res;
}
\end{verbatim}

\paragraph{Complejidad.}
\begin{itemize}\setlength\itemsep{0pt}
	\item Build $O(n)$, query sobre path $O(\log^2 n)$ (con segtree) o $O(\log n)$ (con BIT si solo se hace suma).
	\item Memoria $O(n)$.
\end{itemize}

\paragraph{Donde tocar para modificar.}
\begin{itemize}\setlength\itemsep{0pt}
	\item \textbf{Usar con segment tree}: aniadir un segtree sobre el arreglo base; \texttt{pathSegments} da los rangos a actualizar/query.
	\item \textbf{Cambiar operacion}: ajustar el segtree (suma, max, etc.).
	\item \textbf{Subtree query}: \texttt{subtree[v]} es el rango \texttt{[pos[v], pos[v] + size[v] - 1]} en el arreglo base.
	\item \textbf{Updates en edges}: aniadir un offset para que el edge $(u, v)$ quede en el nodo mas profundo.
\end{itemize}

\paragraph{Trampas y errores frecuentes.}
\begin{itemize}\setlength\itemsep{0pt}
	\item \texttt{pathSegments} requiere un \texttt{segtree} o \texttt{BIT} externo para ser util; por si solo solo da los rangos.
	\item Para edges vs nodos: si la operacion es sobre edges, guardar el valor en el hijo mas profundo.
	\item Olvidar la recursion \texttt{while (head[u] != head[v])} da resultados incorrectos.
	\item En problemas con updates, asegurarse de que el segtree soporte lazy propagation.
\end{itemize}

\paragraph{Codigo.}
Ver \texttt{cpp.tex}, seccion \textit{Heavy-Light Decomposition (HLD)}.

\vspace{0.5em}

\subsubsection{Centroid Decomposition}

\paragraph{Idea intuitiva.}
Descomponer recursivamente el arbol en \textbf{centroides}. Un centroide de un arbol es un nodo tal que cada subtree (al removerlo) tiene tamano $\le n/2$. Existe siempre y se encuentra en $O(n)$. Esto da una estructura de \textbf{arbol de centroides} de profundidad $O(\log n)$, util para problemas ``para cada nodo, encontrar el mas cercano que cumple X''. La razon de la profundidad $O(\log n)$: cada recursion reduce el problema a subarboles de tamano $\le n/2$, asi que la recursion da una cadena geometrica decreciente.

\paragraph{Propiedades clave (invariantes).}
\begin{itemize}\setlength\itemsep{0pt}
	\item En cada nivel, los subarboles se procesan recursivamente.
	\item \texttt{removed[]} marca nodos ya usados como centroide.
	\item \texttt{centroidParent[]} da el padre en el arbol de centroides.
	\item La altura del arbol de centroides es $O(\log n)$.
\end{itemize}

\paragraph{Codigo clave.}
Encontrar el centroide de un subarbol:
\begin{verbatim}
int findCentroid(int v, int p, int sz) {
    for (int u : adj[v])
        if (u != p && !removed[u] && subsize[u] > sz / 2)
            return findCentroid(u, v, sz);
    return v;
}
\end{verbatim}

\paragraph{Complejidad.}
\begin{itemize}\setlength\itemsep{0pt}
	\item Build $O(n \log n)$, queries dependen del problema.
	\item Cada nivel de recursion es $O(n)$ para calcular subsize.
\end{itemize}

\paragraph{Donde tocar para modificar.}
\begin{itemize}\setlength\itemsep{0pt}
	\item \textbf{Query tipica}: en cada nivel, ``expandir'' desde el centroide contando caminos validos y usar una estructura auxiliar (set, multiset).
	\item \textbf{Distancia maxima a un nodo ``malo''}: mantener un multiset de distancias en cada centroide.
	\item \textbf{Cambiar a otro tipo de descomposicion}: si el problema se adapta mejor a small-to-large o DSU on tree, usar esos.
	\item \textbf{Dynamic updates}: combinacion con link-cut trees.
\end{itemize}

\paragraph{Trampas y errores frecuentes.}
\begin{itemize}\setlength\itemsep{0pt}
	\item El snippet actual \textbf{solo construye} la descomposicion; las queries se aniaden segun el problema.
	\item Olvidar resetear \texttt{removed[]} causa recursion infinita.
	\item El calculo de subsize debe ignorar nodos ya ``removed''.
	\item El centroide es \textbf{unico} salvo empates en tamano; cualquiera sirve.
\end{itemize}

\paragraph{Codigo.}
Ver \texttt{cpp.tex}, seccion \textit{Centroid Decomposition}.

\vspace{0.5em}

\subsubsection{Euler Tour + range queries}

\paragraph{Idea intuitiva.}
Un \textbf{Euler Tour} (first-time visit) numera los nodos en el orden en que se visitan. El subarbol de $v$ corresponde a un \textbf{rango contiguo} $[tin[v], tout[v]]$ en este orden. Asi, queries sobre subarboles se resuelven con un BIT o segment tree sobre el arreglo de tiempos. La razon del rango contiguo: cuando entramos a un subarbol, no salimos hasta haber recorrido todos sus descendientes; entonces las posiciones consecutivas corresponden exactamente al subarbol.

\paragraph{Propiedades clave (invariantes).}
\begin{itemize}\setlength\itemsep{0pt}
	\item \texttt{tin[v]} = tiempo de primera visita a $v$.
	\item \texttt{tout[v]} = tiempo al \textbf{terminar} de procesar el subarbol de $v$.
	\item El subarbol de $v$ ocupa las posiciones $[tin[v], tout[v]]$.
	\item Para verificar si $u$ es ancestro de $v$: \texttt{tin[u] <= tin[v] \&\& tout[v] <= tout[u]}.
\end{itemize}

\paragraph{Codigo clave.}
El DFS del Euler Tour:
\begin{verbatim}
void dfs(int v, int p) {
    tin[v] = ++timer;
    for (int u : adj[v])
        if (u != p) dfs(u, v);
    tout[v] = timer;
}
\end{verbatim}

\paragraph{Complejidad.}
\begin{itemize}\setlength\itemsep{0pt}
	\item Build $O(n)$, query sobre subarbol $O(\log n)$ con BIT/segtree.
	\item Memoria $O(n)$.
\end{itemize}

\paragraph{Donde tocar para modificar.}
\begin{itemize}\setlength\itemsep{0pt}
	\item \textbf{Aniadir valor por nodo}: guardar el valor en \texttt{tin[v]} y operar sobre el rango.
	\item \textbf{Path query}: el Euler Tour \textbf{no} sirve para paths; usar HLD o LCA.
	\item \textbf{Multiple queries offline}: procesar en orden de tiempo y mantener una estructura por profundidad.
	\item \textbf{Verificar ancestro}: $O(1)$ con la condicion \texttt{tin}.
\end{itemize}

\paragraph{Trampas y errores frecuentes.}
\begin{itemize}\setlength\itemsep{0pt}
	\item El snippet solo construye el tour; las queries concretas se aniaden.
	\item \texttt{tout[v]} esta en el momento de salir, asi que \texttt{tout[v] - tin[v] + 1} = tamano del subarbol.
	\item Si el arbol no es conexo, aniadir un loop externo sobre cada componente.
	\item Para queries en tiempo (no subtree), usar el orden del tour y offline processing.
\end{itemize}

\paragraph{Codigo.}
Ver \texttt{cpp.tex}, seccion \textit{Euler Tour + range queries}.

\subsection{BST balanceado}

\subsubsection{Treap implicito}

\paragraph{Idea intuitiva.}
Un \textbf{treap} es un arbol binario de busqueda donde cada nodo tiene una \textbf{clave} (el valor de la secuencia) y una \textbf{prioridad} aleatoria. La invariante es: BST por clave, heap por prioridad. Esto da un arbol \textbf{esperado balanceado} sin necesidad de rebalanceo explicito. Un \textbf{treap implicito} no almacena una clave; la posicion de un nodo en la secuencia esta dada por el \textbf{tamano del subarbol izquierdo}, permitiendo split/merge por posicion. La estructura es esencialmente un RBS (randomized BST): la altura esperada es $O(\log n)$ por la propiedad de heap + BST.

\paragraph{Propiedades clave (invariantes).}
\begin{itemize}\setlength\itemsep{0pt}
	\item \texttt{sz(node)} = tamano del subarbol (1 + izquierdo + derecho).
	\item \texttt{split(t, k, l, r)}: parte \texttt{t} en \texttt{l} (primeros $k$ nodos) y \texttt{r} (resto).
	\item \texttt{merge(l, r)}: une dos treaps (todos los de \texttt{l} antes que los de \texttt{r}).
	\item La prioridad aleatoria asegura equilibrio esperado.
\end{itemize}

\paragraph{Codigo clave.}
Split por posicion (la base del treap implicito):
\begin{verbatim}
void split(Node* t, int k, Node*& l, Node*& r) {
    if (!t) { l = r = nullptr; return; }
    push(t);
    if (sz(t->l) >= k) {
        split(t->l, k, l, t->l);
        r = update(t);
    } else {
        split(t->r, k - sz(t->l) - 1, t->r, r);
        l = update(t);
    }
}
\end{verbatim}

\paragraph{Complejidad.}
\begin{itemize}\setlength\itemsep{0pt}
	\item $O(\log n)$ \textbf{esperado} por operacion (split, merge, insert, erase).
	\item Peor caso exponencial (probabilidad astronomicamente pequena).
\end{itemize}

\paragraph{Donde tocar para modificar.}
\begin{itemize}\setlength\itemsep{0pt}
	\item \textbf{Treap con clave (no implicito)}: usar \texttt{key} en vez de posicion; cambiar \texttt{split} por \texttt{splitByKey}.
	\item \textbf{Lazy propagation}: aniadir campos \texttt{lazy} y aplicarlos en \texttt{push} (recursivo).
	\item \textbf{Persistencia}: aniadir versionado clonando nodos.
	\item \textbf{Mejor aleatoriedad}: en lugar de \texttt{rand()}, usar \texttt{rng()} (mejor estado).
	\item \textbf{Operacion de reversa}: aniadir \texttt{lazyReverse} y swap de hijos en push.
\end{itemize}

\paragraph{Trampas y errores frecuentes.}
\begin{itemize}\setlength\itemsep{0pt}
	\item En el snippet actual, \texttt{Node(int v)} usa \texttt{rand()} que en algunos compiladores da maxima prioridad 0; aniadir \texttt{rng()} para mejor distribucion.
	\item Olvidar \texttt{update(t)} tras modificar hijos da tamano incorrecto.
	\item Si lazy propagation no se aplica en push, las queries dan valores viejos.
	\item Memory leak con \texttt{new}; en contests OK, en produccion aniadir destructor.
\end{itemize}

\paragraph{Codigo.}
Ver \texttt{cpp.tex}, seccion \textit{Treap implicito}.

\vspace{0.5em}

\subsubsection{Mo con updates}

\paragraph{Idea intuitiva.}
Extension del algoritmo de Mo para queries que mezclan \textbf{rango} y \textbf{updates}. Las queries se ordenan por $(L/\text{blockSize}, R/\text{blockSize}, T)$ donde $T$ es el ``tiempo'' (numero de updates). Asi, moverse entre queries consecutivas cuesta $O(n^{2/3})$ amortizado. La intuicion: el sort por bloques reduce los movimientos totales, y la tercera coordenada $T$ se optimiza con alternancia similar a $R$.

\paragraph{Propiedades clave (invariantes).}
\begin{itemize}\setlength\itemsep{0pt}
	\item Tres punteros: $L$, $R$, $T$.
	\item \texttt{applyUpdate} y \texttt{revertUpdate} son \textbf{inversos} exactos.
	\item \texttt{addPos} y \texttt{removePos} mantienen el estado.
	\item El tamano de bloque optimo es $n^{2/3}$ (donde $n$ = tamano del array, $Q$ = numero de queries).
\end{itemize}

\paragraph{Codigo clave.}
El sort con tres coordenadas:
\begin{verbatim}
int blockSize = (int)pow(n, 2.0/3.0);
sort(queries.begin(), queries.end(), [&](const Q& a, const Q& b) {
    int blockA = a.l / blockSize, blockB = b.l / blockSize;
    if (blockA != blockB) return blockA < blockB;
    int blockAR = a.r / blockSize, blockBR = b.r / blockSize;
    if (blockAR != blockBR) return blockAR < blockBR;
    return a.t < b.t;
});
\end{verbatim}

\paragraph{Complejidad.}
\begin{itemize}\setlength\itemsep{0pt}
	\item $O((N + Q) \cdot n^{2/3})$ donde $n^{2/3}$ es el tamano de bloque optimo.
	\item En problemas reales, funciona bien hasta $N, Q \sim 10^5$.
\end{itemize}

\paragraph{Donde tocar para modificar.}
\begin{itemize}\setlength\itemsep{0pt}
	\item \textbf{El corazon es el ``add/remove''}: aniadir la logica del problema (ej. contar colores distintos, suma de frecuencias, etc.).
	\item \textbf{Cambiar blockSize}: si la mayoria de updates vs queries cambia, ajustar.
	\item \textbf{Rollback}: aniadir \texttt{save()} / \texttt{restore()} para volver a estados anteriores.
	\item \textbf{Hilbert order}: alternativa con mejor locality.
\end{itemize}

\paragraph{Trampas y errores frecuentes.}
\begin{itemize}\setlength\itemsep{0pt}
	\item El snippet actual es solo el \textbf{esqueleto}; \texttt{addPos}/\texttt{removePos}/\texttt{applyUpdate}/\texttt{revertUpdate} deben ser implementados por problema.
	\item Olvidar \texttt{revertUpdate} es el error mas comun.
	\item Si el numero de updates es muy grande, Mo con updates se vuelve lento; considerar otra estructura.
	\item Mantener consistencia entre \texttt{T} y la posicion de los updates.
\end{itemize}

\paragraph{Codigo.}
Ver \texttt{cpp.tex}, seccion \textit{Mo con updates}.

\vspace{0.5em}

\subsubsection{sqrt-decomposition generico}

\paragraph{Idea intuitiva.}
Dividir el arreglo en \textbf{bloques} de tamano $B \approx \sqrt{n}$. Mantener informacion precomputada por bloque (suma, max, etc.). Queries de rango se resuelven combinando bloques enteros ($O(\sqrt{n})$ en total) y elementos sueltos en los bloques parciales. Updates puntuales actualizan el valor y el resumen del bloque. La intuicion: las queries ``normales'' cruzan $O(n/B)$ bloques; tomar $B = \sqrt{n}$ balancea el costo de bloques enteros vs elementos sueltos.

\paragraph{Propiedades clave (invariantes).}
\begin{itemize}\setlength\itemsep{0pt}
	\item $B = \lfloor \sqrt{n} \rfloor + 1$.
	\item \texttt{bucket[i]} = resumen del bloque $i$-esimo (suma, max, etc.).
	\item \texttt{rangeSum(l, r)}: 3 casos segun $l$ y $r$ esten en el mismo bloque o no.
	\item Cada query cruza $\le 2 + (r - l)/B$ bloques.
\end{itemize}

\paragraph{Codigo clave.}
Estructura basica:
\begin{verbatim}
int B = sqrt(n) + 1;
vector<T> bucket(n / B + 1, IDENTITY);
T query(int l, int r) {
    T res = IDENTITY;
    int bl = l / B, br = r / B;
    if (bl == br) {
        for (int i = l; i <= r; ++i) res = combine(res, a[i]);
    } else {
        for (int i = l; i < (bl+1)*B; ++i) res = combine(res, a[i]);
        for (int i = bl+1; i < br; ++i) res = combine(res, bucket[i]);
        for (int i = br*B; i <= r; ++i) res = combine(res, a[i]);
    }
    return res;
}
\end{verbatim}

\paragraph{Complejidad.}
\begin{itemize}\setlength\itemsep{0pt}
	\item $O(\sqrt{n})$ por operacion.
	\item Memoria $O(n)$.
\end{itemize}

\paragraph{Donde tocar para modificar.}
\begin{itemize}\setlength\itemsep{0pt}
	\item \textbf{Cambiar suma por max/min}: ajustar \texttt{bucket} y la combinacion.
	\item \textbf{Aniadir update en rango}: si la operacion es asociativa, sumar al resumen del bloque + ajustar elementos sueltos.
	\item \textbf{Tamano de bloque optimo}: si las queries son muy largas, $B = n^{2/3}$ (sirve para Mo con updates).
	\item \textbf{Operacion no asociativa}: usar segment tree en su lugar.
\end{itemize}

\paragraph{Trampas y errores frecuentes.}
\begin{itemize}\setlength\itemsep{0pt}
	\item Si $n$ no es multiplo de $B$, el ultimo bloque es mas chico; aniadir check.
	\item La identidad de la operacion debe estar bien definida (0 para suma, $-\infty$ para max).
	\item Si la operacion no es asociativa, los bloques no se pueden combinar; cambiar a segment tree.
\end{itemize}

\paragraph{Codigo.}
Ver \texttt{cpp.tex}, seccion \textit{sqrt-decomposition generico}.

% ============================================================
\section{Strings}
% ============================================================

\subsection{Hashing}

\subsubsection{Rolling Hash (doble modulo)}

\paragraph{Idea intuitiva.}
Un \textbf{rolling hash} asigna a cada prefijo del string un valor numerico tal que el hash de un substring se puede obtener en $O(1)$ a partir de los prefijos. La forma clasica: $h[i] = h[i-1] \cdot A + s[i-1]$ (donde $A$ es una base y todo modulo $M$). Asi, $h[r] - h[l-1] \cdot A^{r-l+1}$ da el hash del substring $[l, r]$ (con $A^{r-l+1}$ precomputado). Usar \textbf{dos modulos} distintos reduce drasticamente la probabilidad de colision. La idea algebraica: el string se interpreta como un polinomio en $A$ con coeficientes los caracteres; el hash es su evaluacion modulo $M$.

\paragraph{Propiedades clave (invariantes).}
\begin{itemize}\setlength\itemsep{0pt}
	\item $h[0] = 0$, $h[i]$ = hash de los primeros $i$ caracteres.
	\item $p[i] = A^i \bmod M$ precomputado.
	\item Comparar dos substrings: comparar sus pares de hashes.
	\item \textbf{Colisiones}: posibles (probabilidad $1/M$ por par). Doble hash: $\approx 1/(M_1 M_2)$.
	\item Si $A$ y $M$ son coprimos, el hash es bijectivo modulo $M$ para strings de longitud fija.
\end{itemize}

\paragraph{Codigo clave.}
La formula del hash de substring:
\begin{verbatim}
pair<long long, long long> get(int l, int r) {
    // h[r] - h[l-1] * A^(r-l+1)
    long long x1 = (h1[r] - h1[l-1] * p1[r-l+1] % M1 + M1) % M1;
    long long x2 = (h2[r] - h2[l-1] * p2[r-l+1] % M2 + M2) % M2;
    return {x1, x2};
}
\end{verbatim}

\paragraph{Complejidad.}
\begin{itemize}\setlength\itemsep{0pt}
	\item Precompute $O(n)$, query hash $O(1)$.
	\item Memoria $O(n)$.
\end{itemize}

\paragraph{Donde tocar para modificar.}
\begin{itemize}\setlength\itemsep{0pt}
	\item \textbf{Caracteres con valor $> 26$}: si los strings incluyen simbolos, aniadir \texttt{if (c >= 'a' \&\& c <= 'z') v = c - 'a'; else v = 26 + ...}.
	\item \textbf{Hash de substrings con updates}: aniadir un Fenwick tree / segment tree sobre los hashes (no es trivial).
	\item \textbf{Cambiar mod}: usar dos primos grandes ($10^9+7$ y $10^9+9$) o un \texttt{u128} para un solo hash sin mod.
	\item \textbf{Base aleatoria}: aniadir random para evitar ataques (en contests no es necesario; en produccion si).
	\item \textbf{Hash polinomial}: $h = \sum s_i \cdot A^i$ en vez de $s_i \cdot A^{i-1}$; algunos problemas requieren este.
\end{itemize}

\paragraph{Trampas y errores frecuentes.}
\begin{itemize}\setlength\itemsep{0pt}
	\item Si los substrings a comparar estan vacios (\texttt{l > r}), su hash es 0; aniadir check.
	\item \textbf{Overflow} en \texttt{h[i-1] * A}: usar \texttt{long long} y aplicar \texttt{\% M} en cada operacion (no al final).
	\item La sustraccion \texttt{h[r] - h[l-1]*A^(r-l+1)} puede dar negativo; aniadir \texttt{+M} antes de \texttt{\% M}.
	\item Con doble hash, verificar \textbf{ambos} componentes; si solo uno coincide, es colision.
\end{itemize}

\paragraph{Codigo.}
Ver \texttt{cpp.tex}, seccion \textit{Rolling Hash (doble modulo)}.

\subsection{Patrones}

\subsubsection{KMP + LPS}

\paragraph{Idea intuitiva.}
El algoritmo \textbf{Knuth-Morris-Pratt} busca todas las ocurrencias de un patron $P$ en un texto $T$ en $O(|T| + |P|)$. La clave es la tabla \textbf{LPS} (Longest Proper Prefix which is Suffix) que, para cada posicion $i$ del patron, dice el largo del prefijo mas largo que tambien es sufijo \textbf{propio}. Durante el match, cuando hay mismatch, en vez de volver al inicio, saltamos al LPS anterior, evitando trabajo redundante. La demostracion de $O(|T| + |P|)$: cada iteracion del bucle externo avanza $i$ o decrece $j$ (con $j$ acotado por $|P|$), asi que en total hay $O(|T| + |P|)$ operaciones.

\paragraph{Propiedades clave (invariantes).}
\begin{itemize}\setlength\itemsep{0pt}
	\item \texttt{lps[i]} = largo del prefijo mas largo que es sufijo propio de \texttt{p[0..i]}.
	\item El match usa dos punteros \texttt{i} (texto) y \texttt{j} (patron); en mismatch, $j = lps[j-1]$.
	\item Cuando $j == |P|$, se encontro una ocurrencia; continuar con $j = lps[j-1]$ para buscar mas.
	\item El LPS se construye con una auto-aplicacion de KMP al patron.
\end{itemize}

\paragraph{Codigo clave.}
Construccion de la tabla LPS:
\begin{verbatim}
for (int i = 1; i < m; ++i) {
    int j = lps[i - 1];
    while (j > 0 && p[i] != p[j]) j = lps[j - 1];
    if (p[i] == p[j]) ++j;
    lps[i] = j;
}
\end{verbatim}

\paragraph{Complejidad.}
\begin{itemize}\setlength\itemsep{0pt}
	\item $O(|T| + |P|)$.
	\item En la practica, rapido y robusto.
\end{itemize}

\paragraph{Donde tocar para modificar.}
\begin{itemize}\setlength\itemsep{0pt}
	\item \textbf{Contar ocurrencias}: el snippet ya cuenta; para devolver posiciones, aniadir un \texttt{vector<int>}.
	\item \textbf{Varios patrones}: usar Aho-Corasick en su lugar.
	\item \textbf{Matcher exacto con wildcard}: aniadir logica para que \texttt{?} matchee cualquier caracter.
	\item \textbf{Tamano pequeno}: si el patron es $\le 5$ caracteres, Rabin-Karp es mas simple.
	\item \textbf{2D KMP}: aniadir una dimension; util en problemas de matrices.
\end{itemize}

\paragraph{Trampas y errores frecuentes.}
\begin{itemize}\setlength\itemsep{0pt}
	\item En el snippet, \texttt{i} y \texttt{j} son 0-based. La condicion \texttt{j == sz\_pattern} para contar funciona.
	\item La construccion de LPS usa el mismo patron como ``texto''.
	\item Si el patron tiene caracteres repetidos, el LPS puede ser no trivial; verificar con casos pequenos.
\end{itemize}

\paragraph{Codigo.}
Ver \texttt{cpp.tex}, seccion \textit{KMP + LPS}.

\vspace{0.5em}

\subsubsection{Z-algorithm}

\paragraph{Idea intuitiva.}
La \textbf{Z-function} $Z[i]$ de un string $S$ es el largo del prefijo mas largo de $S$ que \textbf{tambien} empieza en $S[i]$. Asi, $Z[0] = n$ por convencion. Se calcula en $O(n)$ manteniendo una ventana $[l, r]$ que es el match mas largo encontrado hasta ahora. La consulta $S[l..r] = S[0..r-l]$ da informacion sobre $Z[i]$ cuando $i \le r$. La intuicion: cualquier $i$ dentro de la ventana ``ve'' los caracteres desde una posicion conocida, asi que $Z[i]$ se puede copiar en lugar de recalcular.

\paragraph{Propiedades clave (invariantes).}
\begin{itemize}\setlength\itemsep{0pt}
	\item Construir $Z$ cuesta $O(n)$ con el truco de la ventana.
	\item Para matching de patron $P$ en texto $T$: concatenar $P + \# + T$ y aniadir todos los $i$ con $Z[i] = |P|$.
	\item La ventana $[l, r]$ es siempre la del match mas largo a la derecha.
\end{itemize}

\paragraph{Codigo clave.}
Construccion de Z-function en $O(n)$:
\begin{verbatim}
Z[0] = n;
int l = 0, r = 0;
for (int i = 1; i < n; ++i) {
    if (i <= r) Z[i] = min(r - i + 1, Z[i - l]);
    while (i + Z[i] < n && s[Z[i]] == s[i + Z[i]]) ++Z[i];
    if (i + Z[i] - 1 > r) { l = i; r = i + Z[i] - 1; }
}
\end{verbatim}

\paragraph{Complejidad.}
\begin{itemize}\setlength\itemsep{0pt}
	\item $O(n)$ para construir; matching $O(n + m)$.
	\item Cada caracter se compara a lo sumo 2 veces (extension + reduccion).
\end{itemize}

\paragraph{Donde tocar para modificar.}
\begin{itemize}\setlength\itemsep{0pt}
	\item \textbf{Detectar bordes (border)}: $Z[i] = n - i$ indica que $S[0..n-i-1]$ es borde.
	\item \textbf{Prefijo comun mas largo de dos strings}: aniadir $S_1 + \# + S_2$; $Z[|S_1|+1]$ es el prefijo comun.
	\item \textbf{Comparar prefijos y sufijos}: combinar $S + \# + \text{reverse}(S)$.
	\item \textbf{Matching aproximado}: aniadir tolerancia a mismatches.
\end{itemize}

\paragraph{Trampas y errores frecuentes.}
\begin{itemize}\setlength\itemsep{0pt}
	\item $Z[0]$ siempre vale $n$ (no se calcula).
	\item Si $i > r$, la ventana no aporta; empezar desde 0.
	\item El caracter separador en $P + \# + T$ debe ser uno que no aparezca en $P$ o $T$.
\end{itemize}

\paragraph{Codigo.}
Ver \texttt{cpp.tex}, seccion \textit{Z-algorithm}.

\vspace{0.5em}

\subsubsection{Rabin-Karp (single hash)}

\paragraph{Idea intuitiva.}
Rolling hash aplicado a pattern matching: calcular el hash del patron $P$ y de cada ventana de tamano $|P|$ en $T$. Si los hashes coinciden, verificar caracter por caracter (para descartar colisiones). Esto da $O(n + m)$ \textbf{esperado}. La razon de la eficiencia esperada: cada ventana tiene probabilidad $1/M$ de tener un hash coincidente por casualidad, asi que solo se verifica $O(n/M)$ ventanas en promedio.

\paragraph{Propiedades clave (invariantes).}
\begin{itemize}\setlength\itemsep{0pt}
	\item Si los hashes difieren, los substrings son distintos (sin colision).
	\item Si coinciden, \textbf{verificar} (la verificacion es $O(m)$ pero rara vez se ejecuta).
	\item Peor caso (muchas colisiones): $O(nm)$.
	\item Para $M \ge 10^9$, las colisiones son raras; en la practica nunca pasan.
\end{itemize}

\paragraph{Codigo clave.}
Rolling hash del texto (deslizar ventana):
\begin{verbatim}
long long hash = 0;
for (int i = 0; i < m; ++i) hash = (hash * BASE + text[i]) % MOD;
for (int i = m; i < n; ++i) {
    if (hash == target) check(i - m);
    hash = (hash - text[i - m] * h) * BASE + text[i];  // mod correctamente
    hash = (hash % MOD + MOD) % MOD;
}
\end{verbatim}

\paragraph{Complejidad.}
\begin{itemize}\setlength\itemsep{0pt}
	\item Esperado $O(n + m)$, peor $O(nm)$.
	\item Con doble hash y $M_1, M_2 \ge 10^9$, el peor caso es astronomicamente improbable.
\end{itemize}

\paragraph{Donde tocar para modificar.}
\begin{itemize}\setlength\itemsep{0pt}
	\item \textbf{Doble hash}: aniadir un segundo par \texttt{(BASE2, MOD2)} para reducir colisiones.
	\item \textbf{Varios patrones}: aniadir un set de hashes y buscar match con cualquiera.
	\item \textbf{Tamano de ventana variable}: si los patrones tienen largos distintos, usar el largo especifico por patron.
	\item \textbf{Multi-pattern search}: usar Aho-Corasick en su lugar.
\end{itemize}

\paragraph{Trampas y errores frecuentes.}
\begin{itemize}\setlength\itemsep{0pt}
	\item Overflow en \texttt{(tHash - text[i] * h) * BASE}: aplicar mod correctamente y aniadir \texttt{+MOD} antes de multiplicar.
	\item El factor $h = \text{BASE}^{m-1} \bmod \text{MOD}$ se precomputa.
	\item Si \texttt{MOD} no es primo, las colisiones son mas frecuentes; usar primo.
	\item En el peor caso, Rabin-Karp puede ser lento; preferir KMP si se requiere worst-case.
\end{itemize}

\paragraph{Codigo.}
Ver \texttt{cpp.tex}, seccion \textit{Rabin-Karp (single hash)}.

\vspace{0.5em}

\subsubsection{Aho-Corasick (basic)}

\paragraph{Idea intuitiva.}
Construye un \textbf{trie} con todos los patrones y aniade \textbf{failure links} analogos a los de KMP: el failure link de un nodo $v$ apunta al nodo mas largo que es sufijo del string representado por $v$ y sigue siendo prefijo de algun patron. Asi, al recorrer el texto, si hay mismatch, saltamos al failure link. Esto permite buscar \textbf{todos} los patrones en $O(|T| + \sum |P_i|)$. La idea: unificar varios KMP en una sola estructura.

\paragraph{Propiedades clave (invariantes).}
\begin{itemize}\setlength\itemsep{0pt}
	\item El trie almacena los patrones.
	\item \texttt{link[v]} (failure link) se calcula en BFS en orden de profundidad creciente.
	\item \texttt{out[v]} lista los IDs de patrones que terminan en $v$.
	\item Para queries eficientes, aniadir \textbf{output links}: propagar \texttt{out} de los failure links a los hijos.
\end{itemize}

\paragraph{Codigo clave.}
BFS para construir los failure links:
\begin{verbatim}
queue<int> q;
for (int c = 0; c < SIGMA; ++c)
    if (trie[0].next[c] != -1) {
        q.push(trie[0].next[c]);
        trie[trie[0].next[c]].link = 0;
    } else trie[0].next[c] = 0;  // goto a la raiz
while (!q.empty()) {
    int v = q.front(); q.pop();
    for (int pid : trie[trie[v].link].out)
        trie[v].out.push_back(pid);
    for (int c = 0; c < SIGMA; ++c) {
        if (trie[v].next[c] != -1) {
            trie[trie[v].next[c]].link = trie[trie[v].link].next[c];
            q.push(trie[v].next[c]);
        } else trie[v].next[c] = trie[trie[v].link].next[c];
    }
}
\end{verbatim}

\paragraph{Complejidad.}
\begin{itemize}\setlength\itemsep{0pt}
	\item Build $O(\sum |P_i| \cdot \Sigma)$, query $O(|T| + \text{ocurrencias})$.
\end{itemize}

\paragraph{Donde tocar para modificar.}
\begin{itemize}\setlength\itemsep{0pt}
	\item \textbf{Output links en build}: ya hecho en el snippet (\texttt{for(int pid : t[t[u].link].out) t[u].out.push\_back(pid);}).
	\item \textbf{Transiciones automaticas}: aniadir el goto \texttt{t[v].next[k] = t[t[v].link].next[k]} cuando el enlace directo es \texttt{-1}; asi el bucle de matching es un solo paso por caracter.
	\item \textbf{Aplicar a problemas tipo DP}: contar formas de construir strings evitando los patrones (ver el siguiente modulo).
	\item \textbf{2D Aho-Corasick}: aniadir una dimension para busqueda en matrices.
\end{itemize}

\paragraph{Trampas y errores frecuentes.}
\begin{itemize}\setlength\itemsep{0pt}
	\item La BFS debe empezar con la raiz en la cola y los hijos de la raiz con \texttt{link = 0}.
	\item Olvidar propagar \texttt{out} da ocurrencias perdidas.
	\item Si no se aniaden las transiciones automaticas (goto), el matching es mas lento.
	\item Con muchos patrones, el output puede ser enorme; aniadir contador en lugar de lista.
\end{itemize}

\paragraph{Codigo.}
Ver \texttt{cpp.tex}, seccion \textit{Aho-Corasick (basic)}.

\vspace{0.5em}

\subsubsection{Aho-Corasick DP for forbidden patterns}

\paragraph{Idea intuitiva.}
Contar strings de longitud $L$ que \textbf{evitan} un conjunto de patrones. Se modela como una DP sobre el \textbf{estado del automaton} de Aho-Corasick: $dp[v]$ = numero de formas de llegar al estado $v$ despues de cierto numero de caracteres. Transicion: por cada caracter $c$ en el alfabeto, ir al estado \texttt{next[v][c]}, pero \textbf{solo} si ese estado no tiene un patron (de lo contrario, el string contendria el patron prohibido). La idea: cada estado del automaton representa un ``prefijo conocido'' y las transiciones simulan ``agregar un caracter''.

\paragraph{Propiedades clave (invariantes).}
\begin{itemize}\setlength\itemsep{0pt}
	\item $dp$ se mantiene por nivel (longitud del string construido).
	\item Transicion: $ndp[u] = \sum_{c} dp[v]$ donde $u = next[v][c]$ y $u$ no es ``terminal'' (no tiene patron que termina ahi).
	\item Respuesta: suma de $dp[v]$ sobre todos los estados $v$ no terminales al final.
	\item Si todos los caracteres son permitidos, $|\Sigma|$ entra como factor.
\end{itemize}

\paragraph{Codigo clave.}
DP sobre el automaton:
\begin{verbatim}
vector<long long> dp(numStates, 0), ndp(numStates);
dp[0] = 1;  // empezar en la raiz
for (int step = 0; step < L; ++step) {
    fill(ndp.begin(), ndp.end(), 0);
    for (int v = 0; v < numStates; ++v) if (dp[v] && !isTerminal[v])
        for (int c = 0; c < SIGMA; ++c) {
            int u = next[v][c];
            if (!isTerminal[u]) ndp[u] = (ndp[u] + dp[v]) % MOD;
        }
    swap(dp, ndp);
}
long long ans = 0;
for (int v = 0; v < numStates; ++v) if (!isTerminal[v])
    ans = (ans + dp[v]) % MOD;
\end{verbatim}

\paragraph{Complejidad.}
\begin{itemize}\setlength\itemsep{0pt}
	\item $O(L \cdot |S| \cdot |\Sigma|)$ donde $|S|$ es el numero de estados y $|\Sigma|$ el alfabeto.
	\item Para $|\Sigma| = 2$ y $|S| \le 10^4$: $\sim 2 \cdot 10^{10}$ para $L = 10^6$ (puede ser lento).
\end{itemize}

\paragraph{Donde tocar para modificar.}
\begin{itemize}\setlength\itemsep{0pt}
	\item \textbf{Alfabeto mas pequeno}: si el problema es solo $\{0, 1\}$, $|\Sigma| = 2$.
	\item \textbf{Permitir ocurrencias parciales}: cambiar la condicion a ``no estado terminal'' o relajar segun el problema.
	\item \textbf{Recuperar el string}: aniadir backtracking si el problema lo pide.
	\item \textbf{Matriz de transicion}: para queries rapidas, precomputar $M^c$ con exponentiation.
\end{itemize}

\paragraph{Trampas y errores frecuentes.}
\begin{itemize}\setlength\itemsep{0pt}
	\item El snippet asume Aho-Corasick ya construido con transiciones automaticas. Sin eso, aniadir el goto en el bucle.
	\item Si $L$ es muy grande, la DP lineal es lenta; usar matrix exponentiation.
	\item Estados terminales deben excluirse de \textbf{ambos} lados (estado actual y siguiente).
\end{itemize}

\paragraph{Codigo.}
Ver \texttt{cpp.tex}, seccion \textit{Aho-Corasick DP for forbidden patterns}.

\subsection{Estructuras avanzadas}

\subsubsection{Suffix Array}

\paragraph{Idea intuitiva.}
Un \textbf{suffix array} $SA$ es un arreglo que contiene las posiciones de inicio de los sufijos del string $S$, ordenados lexicograficamente. Se construye en $O(n \log n)$ con doubling: ordenar por el primer caracter, luego por los primeros 2, luego 4, etc. Para hacer cada paso lineal, se usa \textbf{counting sort} sobre los ranks. La idea: cada nivel dobla la cantidad de caracteres considerados, asi que en $\log n$ niveles los sufijos estan completamente ordenados.

\paragraph{Propiedades clave (invariantes).}
\begin{itemize}\setlength\itemsep{0pt}
	\item \texttt{p[i]} = posicion de inicio del $i$-esimo sufijo (ordenado).
	\item \texttt{c[i]} = rank del sufijo que empieza en $i$.
	\item Cada iteracion duplica la longitud de la ``ventana'' considerada.
	\item El algoritmo termina cuando $2^k \ge n$.
	\item La permutacion \texttt{p} es estable entre iteraciones.
\end{itemize}

\paragraph{Codigo clave.}
El paso de doubling (clasificar por los primeros $2^k$ caracteres):
\begin{verbatim}
for (int k = 0; (1 << k) < n; ++k) {
    for (int i = 0; i < n; ++i) pn[i] = (p[i] - (1 << k) + n) % n;
    // counting sort por c[pn[i]] segundo, c[i] primero
    // ...
    cn[pn[0]] = 0;
    for (int i = 1; i < n; ++i)
        cn[pn[i]] = cn[pn[i-1]] + (pair != pair_prev);
    c = cn;
}
\end{verbatim}

\paragraph{Complejidad.}
\begin{itemize}\setlength\itemsep{0pt}
	\item $O(n \log n)$ con counting sort.
	\item Constante pequena; en la practica eficiente.
\end{itemize}

\paragraph{Donde tocar para modificar.}
\begin{itemize}\setlength\itemsep{0pt}
	\item \textbf{Suffix array en $O(n)$}: usar SA-IS (mas complejo, no incluido).
	\item \textbf{Buscar patron}: con SA + LCP es $O(|P| \log n)$ por busqueda binaria.
	\item \textbf{Suffix array de varios strings}: aniadir centinelas (\texttt{\$}) entre ellos.
	\item \textbf{Aniadir longitud}: aniadir \texttt{n - p[i]} como clave secundaria (sufijos mas cortos primero).
\end{itemize}

\paragraph{Trampas y errores frecuentes.}
\begin{itemize}\setlength\itemsep{0pt}
	\item La ``ventana'' $2^k$ se hace circular: \texttt{p[i] = (p[i] - (1<<k) + n) \% n}.
	\item Para contar bien los ranks duplicados, comparar pares \texttt{(c[i], c[(i+2\^{}k) \% n])}.
	\item El counting sort necesita un array de acumulacion de tamano $\le n$.
	\item Si los caracteres son negativos o fuera del alfabeto esperado, normalizar antes.
\end{itemize}

\paragraph{Codigo.}
Ver \texttt{cpp.tex}, seccion \textit{Suffix Array}.

\vspace{0.5em}

\subsubsection{LCP Array (Kasai)}

\paragraph{Idea intuitiva.}
El \textbf{LCP array} (Longest Common Prefix) guarda, para cada par de sufijos adyacentes en el SA, el largo de su prefijo comun. El algoritmo de \textbf{Kasai} lo calcula en $O(n)$ iterando los sufijos \textbf{en orden del string original} (no del SA), y reutiliza el LCP previo. La demostracion: el sufijo en $i+1$ se obtiene del sufijo en $i$ ``avanzando un caracter''; los LCPs decrecen como maximo 1 entre iteraciones.

\paragraph{Propiedades clave (invariantes).}
\begin{itemize}\setlength\itemsep{0pt}
	\item \texttt{lcp[i]} = LCP entre los sufijos en \texttt{sa[i]} y \texttt{sa[i-1]}.
	\item Se necesita \texttt{rank\_[i]} (la inversa del SA) para iterar.
	\item El truco de Kasai: si \texttt{lcp > 0} para el sufijo anterior, el actual parte de al menos \texttt{lcp - 1}.
	\item La suma de todos los LCPs da informacion sobre repeticiones.
\end{itemize}

\paragraph{Codigo clave.}
Algoritmo de Kasai:
\begin{verbatim}
vector<int> rank(n);
for (int i = 0; i < n; ++i) rank[sa[i]] = i;
int k = 0;
vector<int> lcp(n - 1);
for (int i = 0; i < n; ++i) {
    if (rank[i] == n - 1) { k = 0; continue; }
    int j = sa[rank[i] + 1];
    while (i + k < n && j + k < n && s[i+k] == s[j+k]) ++k;
    lcp[rank[i]] = k;
    if (k) --k;
}
\end{verbatim}

\paragraph{Complejidad.}
\begin{itemize}\setlength\itemsep{0pt}
	\item $O(n)$.
	\item Cada comparacion reduce $k$ o avanza linealmente; amortizado $O(1)$.
\end{itemize}

\paragraph{Donde tocar para modificar.}
\begin{itemize}\setlength\itemsep{0pt}
	\item \textbf{RMQ sobre LCP}: para queries ``LCP de dos sufijos'', encontrar sus ranks y hacer RMQ en \texttt{lcp}.
	\item \textbf{Contar substrings distintas}: $\sum (n - sa[i] - lcp[i])$.
	\item \textbf{Longest common substring}: maximo en \texttt{lcp}.
	\item \textbf{String matching}: busqueda binaria sobre SA + LCP verification.
\end{itemize}

\paragraph{Trampas y errores frecuentes.}
\begin{itemize}\setlength\itemsep{0pt}
	\item Si dos sufijos son iguales (raro, solo si el string tiene period $< n$), \texttt{lcp} puede ser \texttt{n - sa[i]}.
	\item En la iteracion, aniadir el chequeo \texttt{i + k < n} para evitar overflow.
	\item El \texttt{--k} en la salida es crucial para la complejidad amortizada.
\end{itemize}

\paragraph{Codigo.}
Ver \texttt{cpp.tex}, seccion \textit{LCP Array (Kasai)}.

\vspace{0.5em}

\subsubsection{Suffix Automaton}

\paragraph{Idea intuitiva.}
Un \textbf{suffix automaton} es un DFA minimo que reconoce todos los substrings del string original. Tiene la propiedad magica de tener $\le 2n - 1$ estados. Cada estado representa una clase de equivalencia de substrings (los que tienen el mismo conjunto de posiciones finales). El \textbf{link} de un estado apunta al estado del sufijo propio mas largo.

\paragraph{Propiedades clave (invariantes).}
\begin{itemize}\setlength\itemsep{0pt}
	\item \texttt{len[v]}: largo maximo de los substrings de la clase.
	\item \texttt{link[v]}: padre en el arbol de sufijos (suffix link).
	\item \texttt{next[v][c]}: transicion por caracter.
	\item \texttt{occ[v]}: numero de posiciones finales (cuantas veces aparece el patron).
\end{itemize}

\paragraph{Complejidad.}
\begin{itemize}\setlength\itemsep{0pt}
	\item Build $O(n)$ con la operacion de \textbf{clonar} estados.
\end{itemize}

\paragraph{Donde tocar para modificar.}
\begin{itemize}\setlength\itemsep{0pt}
	\item \textbf{Contar substrings distintas}: $\sum_{v \ne 0} (len[v] - len[link[v]])$.
	\item \textbf{Longest common substring}: navegar el SAM con el segundo string, manteniendo el match actual.
	\item \textbf{Substring count}: despues del build, ordenar estados por \texttt{len} descendente y propagar \texttt{occ}.
	\item \textbf{Endpos queries}: el SAM es el DFA canonico de los substrings; queries ``cuantas veces aparece $P$'' se responden en $O(|P|)$.
\end{itemize}

\paragraph{Trampas y errores frecuentes.}
\begin{itemize}\setlength\itemsep{0pt}
	\item El \textbf{clone} debe tener \texttt{occ = 0} (no es una nueva posicion final); el \texttt{link} del estado siendo creado y del estado clonado deben apuntar al clon.
	\item Saltarse esto da SAM incorrecto.
\end{itemize}

\paragraph{Codigo.}
Ver \texttt{cpp.tex}, seccion \textit{Suffix Automaton}.

\vspace{0.5em}

\subsubsection{Lyndon factorization (Duval)}

\paragraph{Idea intuitiva.}
Un \textbf{string de Lyndon} es un string que es estrictamente menor que todos sus rotaciones. Toda cadena tiene una \textbf{factorizacion unica} de Lyndon (algoritmo de \textbf{Duval}): una secuencia de strings de Lyndon $L_1, L_2, \dots, L_k$ tales que $L_1 \ge L_2 \ge \dots \ge L_k$ y la concatenacion es el string original. Esto da $O(n)$ para encontrar los cortes. La demostracion de correctitud: cada factor es minimal en su clase de rotacion, asi que la concatenacion respeta el orden lexicografico.

\paragraph{Propiedades clave (invariantes).}
\begin{itemize}\setlength\itemsep{0pt}
	\item Cada factor de Lyndon es de la forma $uv$ donde $u$ se repite y $v$ es prefijo de $u$.
	\item La longitud del factor esta dada por la ``longitud del periodo''.
	\item Aplicacion: el ``Lyndon word'' es el menor en orden lexicografico entre todas las rotaciones.
	\item La factorizacion es \textbf{unica}.
\end{itemize}

\paragraph{Codigo clave.}
El algoritmo de Duval (iterativo):
\begin{verbatim}
int n = s.size();
int i = 0;
while (i < n) {
    int j = i + 1, k = i;
    while (j < n && s[k] <= s[j]) {
        if (s[k] < s[j]) k = i;
        else ++k;
        ++j;
    }
    int len = j - k;
    while (i < j) { cout << "factor: " << s.substr(i, len) << "\n"; i += len; }
}
\end{verbatim}

\paragraph{Complejidad.}
\begin{itemize}\setlength\itemsep{0pt}
	\item $O(n)$.
	\item Cada comparacion es $O(1)$; el puntero $k$ retrocede a lo sumo $n$ veces en total.
\end{itemize}

\paragraph{Donde tocar para modificar.}
\begin{itemize}\setlength\itemsep{0pt}
	\item \textbf{Encontrar la menor rotacion}: el menor sufijo lexicografico de \texttt{s + s} es la menor rotacion; la posicion se obtiene con un solo pase de Duval.
	\item \textbf{Algoritmo de Booth}: alternativa para rotacion minima, tambien $O(n)$.
	\item \textbf{Validar}: si necesitas verificar que un string es de Lyndon, compararlo con sus rotaciones.
	\item \textbf{Encontrar todos los factores}: guardar la posicion y longitud de cada uno.
\end{itemize}

\paragraph{Trampas y errores frecuentes.}
\begin{itemize}\setlength\itemsep{0pt}
	\item La condicion \texttt{s[j] <= s[k]} (no \texttt{<}) en el while de Duval es crucial; cambiar a \texttt{<} da comportamiento distinto.
	\item El caso $s[k] < s[j]$ resetea $k$ a $i$ (inicio del factor actual).
	\item El bucle externo debe avanzar $i$ por la longitud del factor, no por 1.
\end{itemize}

\paragraph{Codigo.}
Ver \texttt{cpp.tex}, seccion \textit{Lyndon factorization (Duval)}.

\subsection{Palindromos}

\subsubsection{Manacher}

\paragraph{Idea intuitiva.}
Calcula en $O(n)$ los \textbf{radios} de todos los palindromos centrados en cada posicion (impares y pares). La idea: mantener una ventana $[l, r]$ que es el palindromo centrado en un punto anterior mas largo encontrado. Si el nuevo centro $i$ cae dentro de $[l, r]$, entonces el palindromo centrado en $i$ tiene al menos el radio del palindromo ``espejo'' $l + r - i$; sino, empezar desde 1 (impar) o 0 (par). Despues, expandir lo que falte. La demostracion de $O(n)$: cada caracter se compara a lo sumo 2 veces (extension + limite por ventana).

\paragraph{Propiedades clave (invariantes).}
\begin{itemize}\setlength\itemsep{0pt}
	\item Conveccion del snippet: \texttt{odd[i]} = $k$ tal que el palindromo impar mas largo centrado en $i$ tiene longitud $2k - 1$. \texttt{even[i]} = $k$ para par centrado entre $i-1$ e $i$, longitud $2k$.
	\item $k$ cuenta tambien cuantos palindromos hay centrados ahi.
	\item Las queries \texttt{isPalindrome(l, r)} son $O(1)$.
	\item La ventana $[l, r]$ siempre contiene el palindromo mas largo encontrado hasta ahora.
\end{itemize}

\paragraph{Codigo clave.}
Manacher para palindromos impares:
\begin{verbatim}
int l = 0, r = -1;
for (int i = 0; i < n; ++i) {
    int k = (i > r) ? 1 : min(odd[l + r - i], r - i + 1);
    while (i - k >= 0 && i + k < n && s[i-k] == s[i+k]) ++k;
    odd[i] = k;
    if (i + k - 1 > r) { l = i - k + 1; r = i + k - 1; }
}
\end{verbatim}

\paragraph{Complejidad.}
\begin{itemize}\setlength\itemsep{0pt}
	\item Construccion $O(n)$, queries $O(1)$.
	\item Cada extension se cuenta una sola vez.
\end{itemize}

\paragraph{Donde tocar para modificar.}
\begin{itemize}\setlength\itemsep{0pt}
	\item \textbf{Solo contar palindromos}: ya implementado en \texttt{countPalindromes}.
	\item \textbf{Longest palindromic substring}: ya en \texttt{longestPalindrome}.
	\item \textbf{Modificar la conveccion de radios}: si queres ``radio = mitad de la longitud'' en vez de ``k palindromos'', ajustar el offset.
	\item \textbf{Queries con k no precomputado}: aniadir helpers especificos.
	\item \textbf{Modificar el caracter de relleno}: para multiples separadores, aniadir un caracter unico.
\end{itemize}

\paragraph{Trampas y errores frecuentes.}
\begin{itemize}\setlength\itemsep{0pt}
	\item La conveccion del snippet es \textbf{distinta} de la ``clasica de radio'' (donde radio = $(L-1)/2$); tener cuidado al copiar a otro codigo.
	\item Si los palindromos se solapan, la condicion \texttt{i + k - 1 > r} actualiza la ventana.
	\item En palindromos pares, el centro esta ``entre'' dos indices; aniadir offset en queries.
\end{itemize}

\paragraph{Codigo.}
Ver \texttt{cpp.tex}, seccion \textit{Manacher}.

% ============================================================
\section{Grafos}
% ============================================================

\subsection{Caminos minimos}

\subsubsection{Dijkstra}

\paragraph{Idea intuitiva.}
Encuentra el shortest path desde un origen $s$ a todos los vertices en $O((V + E) \log V)$ asumiendo \textbf{pesos no negativos}. La idea: usar una cola de prioridad (min-heap) ordenada por distancia tentativa. En cada paso, extraer el nodo con menor distancia; si ya fue procesado, saltar. Si al actualizar un vecino la distancia mejora, insertar/actualizar en la cola. La demostracion: cuando un nodo se extrae por primera vez, no hay forma de llegar a el con menor distancia (cualquier camino alternativo deberia pasar por nodos no procesados, con distancia $\ge$ la actual).

\paragraph{Propiedades clave (invariantes).}
\begin{itemize}\setlength\itemsep{0pt}
	\item Asume pesos $\ge 0$; con pesos negativos usar Bellman-Ford.
	\item Greedy correctness: la primera vez que se extrae un nodo, su distancia es la definitiva.
	\item Si se quiere reconstruir el camino, guardar el \texttt{parent[]}.
	\item Cada nodo se procesa a lo sumo $V$ veces (mas relajarions, pero solo $V$ definitivas).
\end{itemize}

\paragraph{Codigo clave.}
El bucle principal de Dijkstra:
\begin{verbatim}
priority_queue<pair<long long, int>, vector<...>, greater<...>> pq;
pq.push({0, source});
while (!pq.empty()) {
    auto [d, v] = pq.top(); pq.pop();
    if (d > dist[v]) continue;  // entrada vieja
    for (auto [w, u] : adj[v])
        if (dist[v] + w < dist[u]) {
            dist[u] = dist[v] + w;
            pq.push({dist[u], u});
        }
}
\end{verbatim}

\paragraph{Complejidad.}
\begin{itemize}\setlength\itemsep{0pt}
	\item $O((V + E) \log V)$ con heap binario, $O(V + E)$ con heap de Fibonacci.
	\item En la practica, $\sim 1.5 \cdot 10^{-7}$ segundos por operacion para $V = 10^5$.
\end{itemize}

\paragraph{Donde tocar para modificar.}
\begin{itemize}\setlength\itemsep{0pt}
	\item \textbf{Camino mas corto entre todos los pares}: si se necesita desde todos los origenes, usar Floyd o ejecutar Dijkstra $V$ veces.
	\item \textbf{Shortest path con aristas negativas limitadas}: Bellman-Ford.
	\item \textbf{Reconstruir camino}: aniadir \texttt{parent[v]} actualizado junto a \texttt{dist[v]}.
	\item \textbf{Dijkstra con potentials (Johnson)}: para aristas negativas si se puede reescalar.
	\item \textbf{Multi-source}: aniadir varios nodos fuentes iniciales al heap.
\end{itemize}

\paragraph{Trampas y errores frecuentes.}
\begin{itemize}\setlength\itemsep{0pt}
	\item Distancias no inicializadas a \texttt{INF}; el snippet usa \texttt{1e18}.
	\item El \texttt{if(dist[adjN] != 1e18)} para borrar de la cola es importante si usas \texttt{set} (no asi con \texttt{priority\_queue}).
	\item Con pesos grandes, \texttt{dist[v] + w} puede overflow; usar \texttt{LLONG\_MAX/2} como INF.
	\item No usar con pesos negativos: el algoritmo puede dar respuestas incorrectas.
\end{itemize}

\paragraph{Codigo.}
Ver \texttt{cpp.tex}, seccion \textit{Dijkstra}.

\vspace{0.5em}

\subsubsection{Bellman-Ford}

\paragraph{Idea intuitiva.}
Encuentra shortest path desde $s$ incluso con \textbf{pesos negativos} en $O(VE)$. La idea: relajar todas las aristas $V - 1$ veces. Cada relajacion propaga la mejor distancia a ``un paso mas''. Despues de $V - 1$ pasadas, si alguna arista aun se puede relajar, hay un \textbf{ciclo negativo} alcanzable. La demostracion: el camino mas corto tiene a lo sumo $V - 1$ aristas (sin ciclos); tras $V - 1$ relajaciones todas las distancias estan en su valor optimo.

\paragraph{Propiedades clave (invariantes).}
\begin{itemize}\setlength\itemsep{0pt}
	\item Funciona con pesos negativos (sin ciclos negativos).
	\item Tras $V - 1$ relajaciones, las distancias son definitivas (si no hay ciclo negativo).
	\item Una relajacion adicional detecta ciclos negativos: si \texttt{dist[u] + w < dist[v]}, hay ciclo negativo alcanzable desde $s$.
	\item Cada relajacion solo mejora distancias, asi que termina (no oscila).
\end{itemize}

\paragraph{Codigo clave.}
El doble bucle de relajacion:
\begin{verbatim}
for (int i = 0; i < n - 1; ++i)
    for (auto [u, v, w] : edges)
        if (dist[u] + w < dist[v])
            dist[v] = dist[u] + w;
// Detectar ciclo negativo
for (auto [u, v, w] : edges)
    if (dist[u] + w < dist[v]) { /* ciclo negativo */ }
\end{verbatim}

\paragraph{Complejidad.}
\begin{itemize}\setlength\itemsep{0pt}
	\item $O(VE)$.
	\item SPFA (con cola) en promedio es mucho mas rapido, peor caso igual.
\end{itemize}

\paragraph{Donde tocar para modificar.}
\begin{itemize}\setlength\itemsep{0pt}
	\item \textbf{SPFA}: optimizacion que relaja solo nodos cuya distancia cambio; en el peor caso sigue siendo $O(VE)$.
	\item \textbf{Detectar ciclo negativo global}: ejecutar Bellman-Ford desde un super-origen conectado a todos.
	\item \textbf{Camino mas corto con K aristas exactas}: $K$ iteraciones de Bellman-Ford.
	\item \textbf{Johnson's}: reescalar pesos con potentials para usar Dijkstra.
\end{itemize}

\paragraph{Trampas y errores frecuentes.}
\begin{itemize}\setlength\itemsep{0pt}
	\item Overflow si los pesos son $\ge 2^{52}$; usar \texttt{long long} y \texttt{\_\_int128} si hace falta.
	\item Si el grafo es denso ($E \approx V^2$), $O(VE)$ es prohibitivo.
	\item El chequeo de ciclo negativo debe hacerse tras las $V - 1$ iteraciones, no durante.
	\item Distancias no inicializadas a \texttt{INF}; el primer nodo (source) tiene distancia 0.
\end{itemize}

\paragraph{Codigo.}
Ver \texttt{cpp.tex}, seccion \textit{Bellman-Ford}.

\vspace{0.5em}

\subsubsection{Floyd-Warshall}

\paragraph{Idea intuitiva.}
Shortest path entre \textbf{todos los pares} de vertices en $O(V^3)$. DP: $d[i][j]$ = shortest path usando solo nodos intermedios en $\{0, \dots, k\}$. Transicion: $d_{k+1}[i][j] = \min(d_k[i][j], d_k[i][k+1] + d_k[k+1][j])$. La demostracion: tras $k$ iteraciones, $d[i][j]$ es el shortest path con intermedios en $\{0, \dots, k-1\}$; aniadir el nodo $k$ da la recurrencia clasica.

\paragraph{Propiedades clave (invariantes).}
\begin{itemize}\setlength\itemsep{0pt}
	\item Detecta ciclos negativos: tras el algoritmo, $d[i][i] < 0$ indica ciclo negativo alcanzable desde $i$.
	\item Funciona con pesos negativos (sin ciclos negativos).
	\item Memoria $O(V^2)$; con $V \le 400$ es OK, con mas conviene Dijkstra $V$ veces.
	\item Es la unica opcion cuando se quieren shortest paths entre todos los pares y el grafo es denso.
\end{itemize}

\paragraph{Codigo clave.}
El triple bucle clasico:
\begin{verbatim}
for (int k = 0; k < n; ++k)
    for (int i = 0; i < n; ++i)
        for (int j = 0; j < n; ++j)
            d[i][j] = min(d[i][j], d[i][k] + d[k][j]);
// detectar ciclo negativo
for (int i = 0; i < n; ++i) if (d[i][i] < 0) tiene_ciclo_negativo = true;
\end{verbatim}

\paragraph{Complejidad.}
\begin{itemize}\setlength\itemsep{0pt}
	\item $O(V^3)$ tiempo, $O(V^2)$ memoria.
	\item Para $V \le 400$: $\sim 6.4 \cdot 10^7$ operaciones (factible).
\end{itemize}

\paragraph{Donde tocar para modificar.}
\begin{itemize}\setlength\itemsep{0pt}
	\item \textbf{Reconstruir camino}: aniadir \texttt{next[i][j]} y backtrack.
	\item \textbf{Transitividad}: el algoritmo tambien calcula transitividad (\texttt{d[i][j] != INF} significa alcanzable).
	\item \textbf{Pesos doubles}: aniadir tolerancia (\texttt{d[i][j] > d[i][k] + d[k][j] - eps}).
	\item \textbf{Grafo no denso}: usar Dijkstra $V$ veces en vez de Floyd.
\end{itemize}

\paragraph{Trampas y errores frecuentes.}
\begin{itemize}\setlength\itemsep{0pt}
	\item Para detectar ciclo negativo correctamente, aniadir un chequeo \texttt{if (ans[i][i] < 0)} tras el triple bucle.
	\item El orden de los bucles \texttt{k, i, j} es importante; invertir da resultados diferentes.
	\item Con \texttt{double}, usar tolerancia \texttt{< eps} en vez de \texttt{<} para evitar loops infinitos.
\end{itemize}

\paragraph{Codigo.}
Ver \texttt{cpp.tex}, seccion \textit{Floyd-Warshall}.

\subsection{Componentes y cortes}

\subsubsection{Kosaraju SCC}

\paragraph{Idea intuitiva.}
Encuentra las \textbf{componentes fuertemente conexas} (SCCs) en $O(V + E)$. Algoritmo en 3 pasos:
\begin{itemize}\setlength\itemsep{0pt}
	\item DFS en $G$ rellenando una pila en \textbf{post-orden} (los nodos que terminan tarde quedan arriba).
	\item Construir $G^T$ (grafo transpuesto).
	\item DFS en $G^T$ procesando nodos en orden de la pila (de mayor a menor tiempo); cada DFS da una SCC.
\end{itemize}
La demostracion: dos nodos $u, v$ estan en la misma SCC sii hay camino $u \leadsto v$ y $v \leadsto u$; Kosaraju captura exactamente esto al explorar $G^T$ desde los nodos con mayor tiempo de finalizacion en $G$.

\paragraph{Propiedades clave (invariantes).}
\begin{itemize}\setlength\itemsep{0pt}
	\item Cada SCC es maximal fuertemente conexa.
	\item El \textbf{grafo de SCCs} (metagraf) es un DAG.
	\item Kosaraju y Tarjan dan el mismo resultado; Tarjan es ``mas rapido'' en practica (un solo DFS).
	\item Cada nodo pertenece a exactamente una SCC.
\end{itemize}

\paragraph{Codigo clave.}
La estructura de Kosaraju:
\begin{verbatim}
vector<int> order;
vector<bool> used(n);
function<void(int)> dfs1 = [&](int v) {
    used[v] = true;
    for (int u : g[v]) if (!used[u]) dfs1(u);
    order.push_back(v);  // post-orden
};
function<void(int, int)> dfs2 = [&](int v, int root) {
    comp[v] = root;
    for (int u : gt[v]) if (comp[u] == -1) dfs2(u, root);
};
// 1: DFS en G para llenar order
// 2: DFS en G^T en orden inverso para asignar SCCs
\end{verbatim}

\paragraph{Complejidad.}
\begin{itemize}\setlength\itemsep{0pt}
	\item $O(V + E)$.
	\item Constante 2 (dos DFS + el grafo transpuesto).
\end{itemize}

\paragraph{Donde tocar para modificar.}
\begin{itemize}\setlength\itemsep{0pt}
	\item \textbf{2-SAT}: SCC sobre grafo de implicaciones.
	\item \textbf{Construir el DAG de SCCs}: aniadir aristas entre SCCs diferentes.
	\item \textbf{Tamano de cada SCC}: aniadir contador durante el segundo DFS.
	\item \textbf{Grafo grande}: usar Tarjan para evitar construir $G^T$.
\end{itemize}

\paragraph{Trampas y errores frecuentes.}
\begin{itemize}\setlength\itemsep{0pt}
	\item El snippet usa \texttt{vector<...> adj[n]} (VLA); reemplazable por \texttt{vector<vector<pair<int,int>>} para portabilidad.
	\item El orden del segundo DFS debe ser \textbf{inverso} al de finalizacion del primero.
	\item Si el grafo no es conexo, aniadir un loop sobre todos los nodos en el primer DFS.
	\item Confundir SCC con componentes conexas (las primeras son para dirigidos, las segundas para no dirigidos).
\end{itemize}

\paragraph{Codigo.}
Ver \texttt{cpp.tex}, seccion \textit{Kosaraju SCC}.

\vspace{0.5em}

\subsubsection{Tarjan (puentes)}

\paragraph{Idea intuitiva.}
Encuentra los \textbf{puentes} (bridges): aristas cuya remocion \textbf{desconecta} el grafo. Usa un solo DFS manteniendo \texttt{tin[v]} (tiempo de descubrimiento) y \texttt{low[v]} (el menor \texttt{tin} alcanzable desde $v$ via back-edges). Una arista $v - u$ es puente si $low[u] > tin[v]$. La demostracion: si $low[u] > tin[v]$, entonces $u$ (y su subarbol) no pueden ``escapar'' de $v$; quitar la arista $v - u$ parte el grafo.

\paragraph{Propiedades clave (invariantes).}
\begin{itemize}\setlength\itemsep{0pt}
	\item \texttt{low[v] = min(tin[v], tin[w])} para todo back-edge $(v, w)$, y \texttt{min(low[u])} para hijos en el DFS.
	\item Solo se consideran las aristas del \textbf{arbol DFS} (no back-edges) para \texttt{low}.
	\item Cada nodo tiene exactamente un \texttt{tin}; cada back-edge se procesa una vez.
\end{itemize}

\paragraph{Codigo clave.}
El DFS de Tarjan para bridges:
\begin{verbatim}
void dfs(int v, int pEdge) {
    used[v] = true;
    tin[v] = low[v] = timer++;
    for (auto [u, id] : adj[v]) {
        if (id == pEdge) continue;  // no contar el edge al padre
        if (used[u]) {
            low[v] = min(low[v], tin[u]);  // back-edge
        } else {
            dfs(u, id);
            low[v] = min(low[v], low[u]);
            if (low[u] > tin[v]) bridges.push_back(id);  // es puente
        }
    }
}
\end{verbatim}

\paragraph{Complejidad.}
\begin{itemize}\setlength\itemsep{0pt}
	\item $O(V + E)$.
	\item Una sola pasada de DFS.
\end{itemize}

\paragraph{Donde tocar para modificar.}
\begin{itemize}\setlength\itemsep{0pt}
	\item \textbf{Bridge tree}: contractar cada 2-edge-connected component a un nodo; el resultado es un arbol.
	\item \textbf{Articulation points}: variante con condicion diferente (ver el siguiente modulo).
	\item \textbf{Grafo no dirigido}: solo funciona en no dirigidos; en dirigidos hay conceptos analogos (strong bridges, dominator tree).
	\item \textbf{Reconstruir componentes 2-edge-connected}: BFS desde cada nodo excluyendo bridges.
\end{itemize}

\paragraph{Trampas y errores frecuentes.}
\begin{itemize}\setlength\itemsep{0pt}
	\item No confundir el parent con un back-edge al padre; chequear \texttt{if (adjN == parent) continue;} y contar solo las veces que se ``visita'' como hijo.
	\item Para multigrafos, los edges deben tener IDs unicos para distinguirlos.
	\item Si el grafo no es conexo, aniadir loop externo sobre todos los nodos.
\end{itemize}

\paragraph{Codigo.}
Ver \texttt{cpp.tex}, seccion \textit{Tarjan (puentes)}.

\vspace{0.5em}

\subsubsection{Tarjan (puntos de articulacion)}

\paragraph{Idea intuitiva.}
Encuentra los \textbf{puntos de articulacion} (cut vertices): vertices cuya remocion \textbf{desconecta} el grafo. Mismo setup que para bridges: \texttt{tin[v]}, \texttt{low[v]}. Un nodo $v$ (no raiz) es de articulacion si existe un hijo $u$ con $low[u] \ge tin[v]$. La \textbf{raiz} es de articulacion si tiene $\ge 2$ hijos en el DFS tree. La razon: para $v$ no raiz, la condicion $\ge$ indica que el subarbol de $u$ depende de $v$; quitar $v$ los desconecta. Para la raiz, perder $\ge 2$ hijos parte el DFS en $\ge 2$ componentes.

\paragraph{Propiedades clave (invariantes).}
\begin{itemize}\setlength\itemsep{0pt}
	\item $low[u] \ge tin[v]$ significa que el subarbol de $u$ no tiene back-edge ``por encima'' de $v$.
	\item La raiz es caso especial: tiene que perder mas de un hijo para que el grafo se parta.
	\item Cada nodo se evalua una sola vez.
\end{itemize}

\paragraph{Codigo clave.}
La condicion de articulacion:
\begin{verbatim}
void dfs(int v, int p) {
    used[v] = true;
    tin[v] = low[v] = timer++;
    int children = 0;
    for (int u : adj[v]) {
        if (u == p) continue;
        if (used[u]) low[v] = min(low[v], tin[u]);
        else {
            dfs(u, v);
            low[v] = min(low[v], low[u]);
            if (low[u] >= tin[v] && p != -1)  // no raiz
                isArt[v] = true;
            ++children;
        }
    }
    if (p == -1 && children > 1) isArt[v] = true;  // raiz especial
}
\end{verbatim}

\paragraph{Complejidad.}
\begin{itemize}\setlength\itemsep{0pt}
	\item $O(V + E)$.
	\item Una sola pasada de DFS.
\end{itemize}

\paragraph{Donde tocar para modificar.}
\begin{itemize}\setlength\itemsep{0pt}
	\item \textbf{Block-cut tree}: estructura bipartita con BCCs (2-vertex-connected components) y articulation points.
	\item \textbf{Verificar biconectividad}: chequear que no hay puntos de articulacion.
	\item \textbf{Componentes biconnected}: extender Tarjan para enumerar BCCs.
\end{itemize}

\paragraph{Trampas y errores frecuentes.}
\begin{itemize}\setlength\itemsep{0pt}
	\item No olvidar el caso especial de la raiz.
	\item La condicion usa \texttt{>=}, no \texttt{>}; un hijo con \texttt{low[u] == tin[v]} tambien cuenta.
	\item Multigrafos pueden tener ``falsos'' articulation points; verificar caso a caso.
\end{itemize}

\paragraph{Codigo.}
Ver \texttt{cpp.tex}, seccion \textit{Tarjan (puntos de articulacion)}.

\vspace{0.5em}

\subsubsection{Tarjan SCC}

\paragraph{Idea intuitiva.}
Variante del SCC que usa \textbf{un solo DFS} con una pila explicita. Mantiene \texttt{disc[v]}, \texttt{low[v]}, y una pila de nodos ``en la SCC actual''. Cuando \texttt{low[v] == disc[v]}, $v$ es la raiz de una SCC; popear hasta $v$ inclusive. La demostracion: $low[v] == disc[v]$ indica que $v$ no tiene back-edges a ancestros; entonces $v$ y los nodos en la pila hasta $v$ forman una SCC maximal.

\paragraph{Propiedades clave (invariantes).}
\begin{itemize}\setlength\itemsep{0pt}
	\item Mas eficiente que Kosaraju (un solo DFS).
	\item \texttt{low[v] = min(low[u], disc[w])} para back-edges o recursion.
	\item La pila asegura que solo los nodos en la SCC actual se popean.
	\item Cada nodo se pushea y popea a lo sumo una vez.
\end{itemize}

\paragraph{Codigo clave.}
El DFS de Tarjan SCC:
\begin{verbatim}
void dfs(int v) {
    disc[v] = low[v] = timer++;
    st.push(v); inStack[v] = true;
    for (int u : g[v]) {
        if (disc[u] == -1) { dfs(u); low[v] = min(low[v], low[u]); }
        else if (inStack[u]) low[v] = min(low[v], disc[u]);
    }
    if (low[v] == disc[v]) {
        // popear hasta v
        while (true) {
            int u = st.top(); st.pop(); inStack[u] = false;
            comp[u] = current_scc;
            if (u == v) break;
        }
        ++current_scc;
    }
}
\end{verbatim}

\paragraph{Complejidad.}
\begin{itemize}\setlength\itemsep{0pt}
	\item $O(V + E)$.
	\item Constante similar a Kosaraju pero sin construir $G^T$.
\end{itemize}

\paragraph{Donde tocar para modificar.}
\begin{itemize}\setlength\itemsep{0pt}
	\item \textbf{Tamano de cada SCC}: aniadir contador durante el pop.
	\item \textbf{Grafo condensation}: aniadir aristas entre SCCs en el DAG resultante.
	\item \textbf{2-SAT}: aniadir la construccion del grafo de implicaciones.
\end{itemize}

\paragraph{Trampas y errores frecuentes.}
\begin{itemize}\setlength\itemsep{0pt}
	\item \texttt{inStack[u]} debe chequearse antes de usar \texttt{low[v] = min(low[v], disc[u])} en back-edges.
	\item La pila global es importante; sin ella no se puede extraer la SCC.
	\item Si el grafo es muy grande, recursion puede overflow; aniadir version iterativa.
\end{itemize}

\paragraph{Codigo.}
Ver \texttt{cpp.tex}, seccion \textit{Tarjan SCC}.

\subsection{Matching}

\subsubsection{Kuhn (bipartite matching)}

\paragraph{Idea intuitiva.}
Encuentra un \textbf{matching maximo} en un grafo bipartito en $O(VE)$. Para cada nodo $u$ en $L$, intentar encontrar un vecino $v$ en $R$ que este libre o cuyo match se pueda ``reorganizar'' recursivamente. Si se encuentra, matchear $u - v$. La idea: si $u$ quiere un vecino $v$ ya ocupado, intentamos ``mover'' el match de $v$ a otro nodo; recursivamente, esto encuentra un augmenting path si existe.

\paragraph{Propiedades clave (invariantes).}
\begin{itemize}\setlength\itemsep{0pt}
	\item \texttt{matchR[v]} = el nodo de $L$ matcheado a $v$ (o $-1$ si libre).
	\item \texttt{vis[]} se resetea \textbf{cada intento} de $u$.
	\item Si el matching es perfecto, $|L| = |R| = $ tamano del matching.
	\item Kuhn no garantiza maximalidad greedy; puede haber augmenting paths no encontrados.
\end{itemize}

\paragraph{Codigo clave.}
La recursion del augmenting path:
\begin{verbatim}
bool tryKuhn(int v) {
    if (vis[v]) return false;
    vis[v] = true;
    for (int u : g[v]) {
        if (matchR[u] == -1 || tryKuhn(matchR[u])) {
            matchR[u] = v;
            return true;
        }
    }
    return false;
}
// Por cada v en L:
fill(vis, vis+n, false);
if (tryKuhn(v)) ++matches;
\end{verbatim}

\paragraph{Complejidad.}
\begin{itemize}\setlength\itemsep{0pt}
	\item $O(VE)$ en el peor caso.
	\item Con Hopcroft-Karp: $O(\sqrt{V} \cdot E)$, mucho mejor para grafos densos.
\end{itemize}

\paragraph{Donde tocar para modificar.}
\begin{itemize}\setlength\itemsep{0pt}
	\item \textbf{Matching bipartito minimo (min vertex cover)}: Konig's theorem.
	\item \textbf{Matching general}: blossom algorithm (no incluido).
	\item \textbf{Aniadir capacidades}: si los nodos de $R$ tienen capacidad $> 1$, partir cada uno en capacidad copias.
	\item \textbf{Output del matching}: el array \texttt{matchR} da los pares finales.
\end{itemize}

\paragraph{Trampas y errores frecuentes.}
\begin{itemize}\setlength\itemsep{0pt}
	\item Olvidar resetear \texttt{vis} por intento; eso es lo que evita ``ciclos infinitos'' en la recursion.
	\item Para grafos grandes ($V > 10^4$), Kuhn puede ser muy lento; usar Hopcroft-Karp.
	\item Si el matching es maximo pero no perfecto, $|L|$ o $|R|$ puede ser mayor; el matching es solo maximo.
\end{itemize}

\paragraph{Codigo.}
Ver \texttt{cpp.tex}, seccion \textit{Kuhn (bipartite matching)}.

\vspace{0.5em}

\subsubsection{Hopcroft-Karp}

\paragraph{Idea intuitiva.}
Matching bipartito maximo en $O(\sqrt{V} E)$. Algoritmo en \textbf{fases}: en cada fase, hacer BFS para encontrar los \textbf{caminos de aumento mas cortos} desde nodos libres de $L$; luego DFS solo por esos caminos. Cada fase aumenta el matching en $\ge 1$ y toma $O(E)$. La demostracion: tras una fase, todos los augmenting paths tienen longitud $\ge$ la nueva minima; eso sucede cada $\sqrt{V}$ fases (longitud minima crece geometricamente).

\paragraph{Propiedades clave (invariantes).}
\begin{itemize}\setlength\itemsep{0pt}
	\item \texttt{dist[u]} = distancia en niveles desde un nodo libre de $L$.
	\item BFS termina cuando no hay mas camino a un nodo libre de $R$.
	\item DFS respeta los niveles (solo sigue \texttt{dist[matchV[v]] == dist[u] + 1}).
	\item Cada fase aumenta el matching en $\ge 1$ arista.
\end{itemize}

\paragraph{Codigo clave.}
El BFS de Hopcroft-Karp:
\begin{verbatim}
bool bfs() {
    queue<int> q;
    for (int v : L) {
        if (matchU[v] == -1) { dist[v] = 0; q.push(v); }
        else dist[v] = -1;
    }
    bool found = false;
    while (!q.empty()) {
        int v = q.front(); q.pop();
        for (int u : g[v]) {
            if (matchV[u] != -1 && dist[matchV[u]] == -1) {
                dist[matchV[u]] = dist[v] + 1;
                q.push(matchV[u]);
            }
            if (matchV[u] == -1) found = true;  // llegamos a libre
        }
    }
    return found;
}
\end{verbatim}

\paragraph{Complejidad.}
\begin{itemize}\setlength\itemsep{0pt}
	\item $O(\sqrt{V} E)$.
	\item Para $V = 10^4$, $E = 10^5$: $\sim 3 \cdot 10^7$ operaciones.
\end{itemize}

\paragraph{Donde tocar para modificar.}
\begin{itemize}\setlength\itemsep{0pt}
	\item \textbf{Reconstruir el matching}: ya viene en \texttt{matchU}/\texttt{matchV}.
	\item \textbf{Aniadir pesos}: Hungarian algorithm.
	\item \textbf{Densidad alta}: si $E \approx V^2$, sigue siendo util pero verificar constantes.
	\item \textbf{Aniadir capacidades}: partir los nodos con capacidad $> 1$ en varias copias.
\end{itemize}

\paragraph{Trampas y errores frecuentes.}
\begin{itemize}\setlength\itemsep{0pt}
	\item \texttt{dist[matchV[v]] == dist[u] + 1} (no \texttt{==} a secas); sin esto, el DFS no respeta los niveles y se pierde la complejidad.
	\item El BFS debe marcar los niveles correctamente; resetear \texttt{dist} para nodos alcanzables solo via matching.
	\item El DFS debe respetar los niveles; cualquier ``atajo'' invalida la cota.
\end{itemize}

\paragraph{Codigo.}
Ver \texttt{cpp.tex}, seccion \textit{Hopcroft-Karp}.

\subsection{Basicos}

\subsubsection{DFS / BFS / componentes}

\paragraph{Idea intuitiva.}
\textbf{DFS} (Depth-First Search) y \textbf{BFS} (Breadth-First Search) son los dos recorridos basicos de grafos. DFS usa una pila (o recursion) y explora tan profundo como puede antes de retroceder; BFS usa una cola y explora por ``niveles'' de distancia. La cantidad de \textbf{componentes conexas} se cuenta iniciando un BFS/DFS desde cada nodo no visitado. La diferencia clave: BFS da el camino mas corto en aristas; DFS da un orden de finalizacion util para problemas topologicos.

\paragraph{Propiedades clave (invariantes).}
\begin{itemize}\setlength\itemsep{0pt}
	\item BFS da la \textbf{distancia en aristas} mas corta desde el origen.
	\item DFS da un \textbf{orden topologico} (si se aniaden al final).
	\item Ambos son $O(V + E)$.
	\item Un grafo con $C$ componentes conexas satisface $\sum \text{subtree\_size}_v = V$ y $\sum E_v = 2E$.
	\item DFS en arbol da un spanning tree; las ``back edges'' corresponden a ciclos.
\end{itemize}

\paragraph{Codigo clave.}
BFS clasico:
\begin{verbatim}
queue<int> q;
q.push(source);
dist[source] = 0;
while (!q.empty()) {
    int v = q.front(); q.pop();
    for (int u : adj[v])
        if (dist[u] == -1) {
            dist[u] = dist[v] + 1;
            q.push(u);
        }
}
\end{verbatim}

\paragraph{Complejidad.}
\begin{itemize}\setlength\itemsep{0pt}
	\item $O(V + E)$ tiempo, $O(V)$ memoria.
	\item Cada arista se procesa exactamente 2 veces (ida y vuelta) en no dirigidos.
\end{itemize}

\paragraph{Donde tocar para modificar.}
\begin{itemize}\setlength\itemsep{0pt}
	\item \textbf{Aniadir pesos}: si las aristas tienen peso y queres shortest path, BFS ya no sirve; usar Dijkstra.
	\item \textbf{Grafo dirigido}: en DFS, aniadir check \texttt{if (state[u] == 1)} para detectar ciclos.
	\item \textbf{Bipartite check}: BFS alternando colores (ver el siguiente modulo).
	\item \textbf{Componentes fuertemente conexas (dirigidos)}: usar Tarjan o Kosaraju.
	\item \textbf{BFS en matriz}: aniadir offsets por direccion para grid 2D.
\end{itemize}

\paragraph{Trampas y errores frecuentes.}
\begin{itemize}\setlength\itemsep{0pt}
	\item Stack overflow con DFS recursivo para $V > 10^5$; usar version iterativa con pila explicita.
	\item En BFS, olvidar marcar \texttt{dist[u] = dist[v] + 1} antes de pushear da visitas duplicadas.
	\item Para grafos dirigidos, ``componente conexa'' no esta definida; usar SCC.
\end{itemize}

\paragraph{Codigo.}
Ver \texttt{cpp.tex}, seccion \textit{DFS / BFS / componentes}.

\vspace{0.5em}

\subsubsection{Topological sort (Kahn + DFS)}

\paragraph{Idea intuitiva.}
Un \textbf{orden topologico} de un DAG es una permutacion de los vertices tal que toda arista $u \to v$ tiene $u$ antes que $v$. Dos algoritmos clasicos:
\begin{itemize}\setlength\itemsep{0pt}
	\item \textbf{Kahn}: BFS sobre nodos con \texttt{indeg == 0}; al ``consumir'' un nodo, decrementar el indeg de sus vecinos; los que lleguen a 0 se encolan.
	\item \textbf{DFS}: hacer DFS; al terminar un nodo (estado = 2), aniadirlo al final; el orden es el reverso.
\end{itemize}
La demostracion: en un DAG siempre hay al menos un nodo con indeg 0; quitarlo mantiene el DAG; por induccion, se obtiene el orden.

\paragraph{Propiedades clave (invariantes).}
\begin{itemize}\setlength\itemsep{0pt}
	\item Existe sii el grafo es un DAG.
	\item Si al final \texttt{order.size() < V}, hay un ciclo.
	\item Kahn da el orden ``natural'' (los que pueden ir primero); DFS da el orden ``al reves'' (los que pueden ir al final primero).
	\item La cantidad de ordenes topologicos puede ser exponencial.
\end{itemize}

\paragraph{Codigo clave.}
Kahn con priority queue (orden lexicografico):
\begin{verbatim}
priority_queue<int, vector<int>, greater<int>> pq;
for (int v = 0; v < n; ++v) if (indeg[v] == 0) pq.push(v);
while (!pq.empty()) {
    int v = pq.top(); pq.pop();
    order.push_back(v);
    for (int u : adj[v]) {
        if (--indeg[u] == 0) pq.push(u);
    }
}
\end{verbatim}

\paragraph{Complejidad.}
\begin{itemize}\setlength\itemsep{0pt}
	\item $O(V + E)$ ambos.
	\item Con priority queue, $O((V + E) \log V)$.
\end{itemize}

\paragraph{Donde tocar para modificar.}
\begin{itemize}\setlength\itemsep{0pt}
	\item \textbf{Detectar ciclos}: comparar \texttt{order.size()} con \texttt{V}.
	\item \textbf{Orden lexicografico minimo}: usar priority queue en Kahn.
	\item \textbf{Topological sort con pesos}: combinar con shortest path en DAG (un solo pase en orden topologico).
	\item \textbf{Encontrar todos los ordenes}: DP sobre subsets (solo para $V$ pequeno).
\end{itemize}

\paragraph{Trampas y errores frecuentes.}
\begin{itemize}\setlength\itemsep{0pt}
	\item En el snippet, Kahn retorna orden vacio si hay ciclo (util para detectar).
	\item En DFS, no olvidar \texttt{reverse(order)}.
	\item Confundir indeg con outdeg; el orden Kahn quita primero los nodos ``fuente''.
	\item Para grafos no DAG, el algoritmo termina con \texttt{order.size() < V}.
\end{itemize}

\paragraph{Codigo.}
Ver \texttt{cpp.tex}, seccion \textit{Topological sort (Kahn + DFS)}.

\vspace{0.5em}

\subsubsection{Bipartite check + 2-coloring}

\paragraph{Idea intuitiva.}
Un grafo es \textbf{bipartito} si sus vertices se pueden dividir en dos conjuntos $L$ y $R$ tales que toda arista va de $L$ a $R$. Esto equivale a decir que se puede \textbf{2-colorear} sin que dos vertices adyacentes tengan el mismo color. BFS alternando colores detecta esto: si en algun momento dos adyacentes tienen el mismo color, no es bipartito. La demostracion: por cada componente, fijar el color de un nodo determina todos los demas; si en algun paso hay conflicto, hay un ciclo impar (no bipartito).

\paragraph{Propiedades clave (invariantes).}
\begin{itemize}\setlength\itemsep{0pt}
	\item Grafo bipartito $\Leftrightarrow$ sin ciclos impares.
	\item Conexo: si un componente no es bipartito, todo el grafo no lo es.
	\item Si hay \textbf{multiple componentes}, cada uno se colorea independientemente.
	\item La 2-coloracion es unica salvo swap de colores.
\end{itemize}

\paragraph{Codigo clave.}
BFS con 2-coloreo:
\begin{verbatim}
vector<int> color(n, -1);
for (int v = 0; v < n; ++v) if (color[v] == -1) {
    color[v] = 0;
    queue<int> q; q.push(v);
    while (!q.empty()) {
        int u = q.front(); q.pop();
        for (int w : adj[u]) {
            if (color[w] == -1) { color[w] = 1 - color[u]; q.push(w); }
            else if (color[w] == color[u]) return false;
        }
    }
}
return true;
\end{verbatim}

\paragraph{Complejidad.}
\begin{itemize}\setlength\itemsep{0pt}
	\item $O(V + E)$.
	\item Memoria $O(V)$.
\end{itemize}

\paragraph{Donde tocar para modificar.}
\begin{itemize}\setlength\itemsep{0pt}
	\item \textbf{Colorear desde multiples fuentes}: BFS desde cada nodo no coloreado.
	\item \textbf{Bipartito + matching}: Hopcroft-Karp.
	\item \textbf{Grafo bipartito ponderado}: Hungarian algorithm.
	\item \textbf{Verificar bipartito en subgrafo}: solo aplicar BFS a los nodos relevantes.
\end{itemize}

\paragraph{Trampas y errores frecuentes.}
\begin{itemize}\setlength\itemsep{0pt}
	\item Verificar \textbf{todos} los componentes, no solo el primero.
	\item Para grafos con auto-loops, un auto-loop hace al grafo no bipartito.
	\item En multigrafos, las aristas multiples no afectan el chequeo de bipartito.
\end{itemize}

\paragraph{Codigo.}
Ver \texttt{cpp.tex}, seccion \textit{Bipartite check + 2-coloring}.

\vspace{0.5em}

\subsubsection{0-1 BFS}

\paragraph{Idea intuitiva.}
Variante de BFS para grafos con \textbf{pesos 0 o 1}. Usa un \texttt{deque}: aristas de peso 0 se encolan al frente (\texttt{push\_front}), aristas de peso 1 al final (\texttt{push\_back}). Asi, la primera vez que se extrae un nodo, su distancia es la definitiva, igual que en Dijkstra pero en $O(V + E)$. La intuicion: el deque mantiene los nodos ordenados por distancia; las aristas de peso 0 no cambian la distancia (entran al frente), las de peso 1 la incrementan (entran al final).

\paragraph{Propiedades clave (invariantes).}
\begin{itemize}\setlength\itemsep{0pt}
	\item Asume pesos \textbf{solo} 0 o 1.
	\item Mejor caso que Dijkstra (sin factor $\log V$).
	\item No funciona con pesos $\ge 2$.
	\item La cola mantiene los nodos ordenados por distancia no decreciente.
\end{itemize}

\paragraph{Codigo clave.}
El deque doble:
\begin{verbatim}
deque<int> q;
q.push_front(source);
dist[source] = 0;
while (!q.empty()) {
    int v = q.front(); q.pop_front();
    for (auto [w, u] : adj[v]) {
        if (dist[v] + w < dist[u]) {
            dist[u] = dist[v] + w;
            if (w == 0) q.push_front(u);
            else q.push_back(u);
        }
    }
}
\end{verbatim}

\paragraph{Complejidad.}
\begin{itemize}\setlength\itemsep{0pt}
	\item $O(V + E)$.
	\item Cada nodo se inserta a lo sumo $O(1)$ veces (con check de distancia).
\end{itemize}

\paragraph{Donde tocar para modificar.}
\begin{itemize}\setlength\itemsep{0pt}
	\item \textbf{Pesos en $[0, k]$ small}: usar small k-ary BFS o Dijkstra.
	\item \textbf{Reconstruir camino}: aniadir \texttt{parent[]}.
	\item \textbf{Pesos negativos}: no funciona, usar Bellman-Ford.
\end{itemize}

\paragraph{Trampas y errores frecuentes.}
\begin{itemize}\setlength\itemsep{0pt}
	\item Si los pesos son otro rango, el algoritmo da respuestas incorrectas. Verificar que las aristas son solo 0/1.
	\item El chequeo \texttt{dist[v] + w < dist[u]} es crucial; sino, el algoritmo procesa nodos multiples veces.
	\item Si una arista tiene peso $\ge 2$, aniadir verificaciones o cambiar a Dijkstra.
\end{itemize}

\paragraph{Codigo.}
Ver \texttt{cpp.tex}, seccion \textit{0-1 BFS}.

\subsection{MST}

\subsubsection{Kruskal}

\paragraph{Idea intuitiva.}
Encuentra el \textbf{arbol de expansion minima} (MST) en $O(E \log E)$. Ordena las aristas por peso ascendente y las aniade al MST si no forman ciclo (chequear con DSU). El arbol tiene exactamente $V - 1$ aristas. La demostracion (``Cut Property''): para cualquier corte, la arista minima cruzando el corte pertenece a \emph{algun} MST; por lo tanto, elegir la minima disponible que no forma ciclo es optimo.

\paragraph{Propiedades clave (invariantes).}
\begin{itemize}\setlength\itemsep{0pt}
	\item Algoritmo \textbf{greedy} correcto (Cut Property).
	\item Funciona en grafos \textbf{no conexos}; produce un \textbf{minimum spanning forest}.
	\item Si las aristas tienen pesos unicos, el MST es unico.
	\item El MST tiene exactamente $V - 1$ aristas (o menos si el grafo es disconexo).
\end{itemize}

\paragraph{Codigo clave.}
Kruskal con DSU:
\begin{verbatim}
sort(edges.begin(), edges.end());
DSU dsu(n);
long long mst = 0;
int edgesUsed = 0;
for (auto [w, u, v] : edges)
    if (dsu.unite(u, v)) {
        mst += w;
        if (++edgesUsed == n - 1) break;
    }
\end{verbatim}

\paragraph{Complejidad.}
\begin{itemize}\setlength\itemsep{0pt}
	\item $O(E \log E)$ (ordenamiento) o $O(E \log V)$ con DSU optimizations.
	\item DSU es casi $O(1)$ amortizado, asi que el sort domina.
\end{itemize}

\paragraph{Donde tocar para modificar.}
\begin{itemize}\setlength\itemsep{0pt}
	\item \textbf{Second best MST}: enumerar aristas del MST, sacar una, aniadir una no-MST; recalcular.
	\item \textbf{Kruskal randomizado}: si el grafo es casi-completo, agregar shuffle.
	\item \textbf{Aniadir info al DSU}: guardar el peso maximo en el camino, etc.
	\item \textbf{DSU con rollback}: si necesitas Kruskal reversible (para queries offline).
\end{itemize}

\paragraph{Trampas y errores frecuentes.}
\begin{itemize}\setlength\itemsep{0pt}
	\item El DSU debe estar limpio para cada invocacion.
	\item Si las aristas tienen pesos negativos, Kruskal sigue funcionando.
	\item Si el grafo es disconexo, el ``MST'' es un bosque con varias componentes.
	\item El check \texttt{if (++edgesUsed == n - 1) break;} optimiza la salida temprana.
\end{itemize}

\paragraph{Codigo.}
Ver \texttt{cpp.tex}, seccion \textit{Kruskal}.

\vspace{0.5em}

\subsubsection{Prim (priority queue)}

\paragraph{Idea intuitiva.}
Otra forma de calcular el MST: empezar desde un nodo y, en cada paso, aniadir la arista de menor peso que conecta el conjunto construido con el resto. Usando priority queue, es $O((V + E) \log V)$. La demostracion: cada vez que se aniade una arista, es la minima cruzando el ``corte'' entre el conjunto actual y el resto (Cut Property); eso preserva optimalidad.

\paragraph{Propiedades clave (invariantes).}
\begin{itemize}\setlength\itemsep{0pt}
	\item Similar a Dijkstra pero sin acumular distancias, solo eligiendo minimo por arista.
	\item Funciona mejor que Kruskal en \textbf{grafos densos} ($E \approx V^2$).
	\item Si el grafo es disconexo, genera el bosque (solo llega a los alcanzables).
	\item Cada nodo se inserta al heap una vez; cada arista se procesa una vez.
\end{itemize}

\paragraph{Codigo clave.}
Prim clasico con heap:
\begin{verbatim}
priority_queue<pair<int, int>, vector<...>, greater<...>> pq;
vector<bool> used(n, false);
vector<int> minEdge(n, INF);
minEdge[0] = 0;
pq.push({0, 0});
while (!pq.empty()) {
    auto [w, v] = pq.top(); pq.pop();
    if (used[v]) continue;
    used[v] = true;
    mst += w;
    for (auto [peso, u] : adj[v])
        if (!used[u] && peso < minEdge[u]) {
            minEdge[u] = peso;
            pq.push({peso, u});
        }
}
\end{verbatim}

\paragraph{Complejidad.}
\begin{itemize}\setlength\itemsep{0pt}
	\item $O((V + E) \log V)$.
	\item Con matriz de adyacencia: $O(V^2)$ sin heap.
\end{itemize}

\paragraph{Donde tocar para modificar.}
\begin{itemize}\setlength\itemsep{0pt}
	\item \textbf{Prim con matriz de adyacencia}: $O(V^2)$ sin heap, util en grafos densos pequenos.
	\item \textbf{Reconstruir el MST}: aniadir \texttt{parent[]}.
	\item \textbf{Prim en streaming}: variante para grafos que llegan incrementalmente.
\end{itemize}

\paragraph{Trampas y errores frecuentes.}
\begin{itemize}\setlength\itemsep{0pt}
	\item El \texttt{if (used[v]) continue;} es crucial: si el nodo ya fue procesado, saltar.
	\item En la version matriz, elije el minimo por linea (no heap).
	\item Si el grafo es disconexo, aniadir un loop externo sobre todos los nodos no usados.
\end{itemize}

\paragraph{Codigo.}
Ver \texttt{cpp.tex}, seccion \textit{Prim (priority queue)}.

\subsection{Flujos}

\subsubsection{Edmonds-Karp (max flow)}

\paragraph{Idea intuitiva.}
Implementacion BFS del algoritmo de \textbf{Ford-Fulkerson}. En cada iteracion, hacer BFS desde $s$ hasta $t$ siguiendo aristas con capacidad residual $> 0$; si se llega a $t$, se encontro un \textbf{camino de aumento}. Enviar el minimo de las capacidades en el camino; actualizar capacidades residuales. La demostracion de $O(VE^2)$: cada BFS encuentra el augmenting path mas corto, y la longitud minima crece monotona; como hay $\le E$ longitudes posibles, hay $\le E$ fases.

\paragraph{Propiedades clave (invariantes).}
\begin{itemize}\setlength\itemsep{0pt}
	\item BFS garantiza caminos de aumento mas cortos.
	\item Capacidad residual: forward edge baja, backward edge sube.
	\item Complejidad $O(VE^2)$.
	\item El flujo maximo respeta la conservacion en nodos (in = out salvo $s, t$).
\end{itemize}

\paragraph{Codigo clave.}
BFS + augmenting path:
\begin{verbatim}
while (bfs(s, t, parent)) {  // bfs encuentra camino o retorna false
    int pathFlow = INF;
    for (int v = t; v != s; v = parent[v])
        pathFlow = min(pathFlow, capacity[parent[v]][v]);
    for (int v = t; v != s; v = parent[v]) {
        capacity[parent[v]][v] -= pathFlow;
        capacity[v][parent[v]] += pathFlow;
    }
    maxFlow += pathFlow;
}
\end{verbatim}

\paragraph{Complejidad.}
\begin{itemize}\setlength\itemsep{0pt}
	\item $O(VE^2)$.
	\item En la practica, mas lento que Dinic para grafos grandes.
\end{itemize}

\paragraph{Donde tocar para modificar.}
\begin{itemize}\setlength\itemsep{0pt}
	\item \textbf{Grafo bipartito}: $O(\sqrt{V} E)$ (equivalente a Hopcroft-Karp).
	\item \textbf{Dinic} es preferible en la mayoria de los casos.
	\item \textbf{Capacidades doubles}: usar tolerancia para comparacion.
	\item \textbf{Min-cost flow}: aniadir costo por arista y modificar el algoritmo.
\end{itemize}

\paragraph{Trampas y errores frecuentes.}
\begin{itemize}\setlength\itemsep{0pt}
	\item Las capacidades deben ser $> 0$ para entrar a la cola (no $>= 0$); sino, ciclos infinitos.
	\item Para grafos bipartitos, Dinic o Hopcroft-Karp son mejores.
	\item Si el flujo maximo es muy grande, los \texttt{int} pueden overflow; usar \texttt{long long}.
\end{itemize}

\paragraph{Codigo.}
Ver \texttt{cpp.tex}, seccion \textit{Edmonds-Karp (max flow)}.

\vspace{0.5em}

\subsubsection{Dinic (max flow)}

\paragraph{Idea intuitiva.}
Algoritmo de flujo mas eficiente que Edmonds-Karp: combina \textbf{BFS para level graph} con \textbf{DFS bloqueante} (multiple augmenting paths en una sola pasada). Cada fase cuesta $O(E)$ y hay $O(V)$ fases en el peor caso; para bipartitos, $O(\sqrt{V} E)$. La razon: cada fase ``satura'' al menos una arista, asi que en $O(E)$ fases el flujo maximo se alcanza; con $O(V)$ niveles en el BFS, hay $O(VE)$ operaciones totales.

\paragraph{Propiedades clave (invariantes).}
\begin{itemize}\setlength\itemsep{0pt}
	\item \texttt{level[v]} = distancia desde $s$ en el level graph (solo aristas con cap $> 0$).
	\item \texttt{it[v]} = iterador actual para evitar reprocesar aristas.
	\item DFS solo sigue aristas con \texttt{level[e.to] == level[v] + 1}.
	\item Back-edges tienen \texttt{rev} apuntando a la forward edge original.
\end{itemize}

\paragraph{Codigo clave.}
El DFS bloqueante de Dinic:
\begin{verbatim}
long long dfs(int v, int t, long long f) {
    if (v == t) return f;
    for (int& i = it[v]; i < adj[v].size(); ++i) {
        Edge& e = adj[v][i];
        if (e.cap > 0 && level[e.to] == level[v] + 1) {
            long long ret = dfs(e.to, t, min(f, e.cap));
            if (ret > 0) {
                e.cap -= ret;
                adj[e.to][e.rev].cap += ret;
                return ret;
            }
        }
    }
    return 0;
}
\end{verbatim}

\paragraph{Complejidad.}
\begin{itemize}\setlength\itemsep{0pt}
	\item $O(V^2 E)$ general, $O(\sqrt{V} E)$ para bipartitos.
	\item En la practica, mucho mas rapido que Edmonds-Karp.
\end{itemize}

\paragraph{Donde tocar para modificar.}
\begin{itemize}\setlength\itemsep{0pt}
	\item \textbf{Dinic con scaling}: aniadir factor de escala para capacidades grandes.
	\item \textbf{Lower bounds}: aniadir un super-source/sink para forzar flujo minimo.
	\item \textbf{Matching bipartito}: el grafo se construye como source -> L -> R -> sink, todos con capacidad 1.
	\item \textbf{Push-relabel}: alternativa para grafos muy densos.
\end{itemize}

\paragraph{Trampas y errores frecuentes.}
\begin{itemize}\setlength\itemsep{0pt}
	\item \texttt{it[v]} debe \textbf{incrementarse} en cada iteracion; sino, el DFS se queda atascado en la misma arista.
	\item Si el BFS no encuentra $t$, el flujo es maximo; terminar.
	\item Las back-edges son cruciales para ``cancelar'' elecciones previas.
\end{itemize}

\paragraph{Codigo.}
Ver \texttt{cpp.tex}, seccion \textit{Dinic (max flow)}.

\vspace{0.5em}

\subsubsection{Min-Cost Max-Flow}

\paragraph{Idea intuitiva.}
Encuentra el flujo maximo con \textbf{costo minimo} (asumiendo costos no negativos). Usa \textbf{SPFA} (o Dijkstra con potentials) para encontrar el shortest path en el \textbf{costo} desde $s$ hasta $t$ considerando capacidades residuales. Cada augmenting path aporta su costo $\times$ flujo. La idea: los potentials reescalan las aristas para que todas tengan costo no negativo, permitiendo usar Dijkstra.

\paragraph{Propiedades clave (invariantes).}
\begin{itemize}\setlength\itemsep{0pt}
	\item \texttt{pot[]} = potentials para evitar aristas negativas (Johnson's trick).
	\item \texttt{dist[v]} = costo minimo desde $s$ con potentials.
	\item Solo se aumentan aristas con \texttt{cap > 0} (forward) y costo es \texttt{+cost} (forward) o \texttt{-cost} (backward).
	\item Tras cada augmenting, los potentials se actualizan para mantener la consistencia.
\end{itemize}

\paragraph{Codigo clave.}
El ciclo de min-cost flow:
\begin{verbatim}
while (spfa(s, t, dist, parent)) {
    int pathFlow = INF;
    for (int v = t; v != s; v = parent[v])
        pathFlow = min(pathFlow, capacity[parent[v]][v]);
    for (int v = t; v != s; v = parent[v]) {
        capacity[parent[v]][v] -= pathFlow;
        capacity[v][parent[v]] += pathFlow;
        cost += pathFlow * edge_cost[parent[v]][v];
    }
}
\end{verbatim}

\paragraph{Complejidad.}
\begin{itemize}\setlength\itemsep{0pt}
	\item $O(F \cdot V E)$ con SPFA, $O(F \cdot E \log V)$ con Dijkstra + potentials.
	\item $F$ = flujo maximo total.
\end{itemize}

\paragraph{Donde tocar para modificar.}
\begin{itemize}\setlength\itemsep{0pt}
	\item \textbf{Dijkstra + potentials}: si los costos son enteros no negativos, esto es mas estable.
	\item \textbf{Costos negativos en aristas}: requieren Bellman-Ford inicial o reorden.
	\item \textbf{Flujo minimo con costo}: encontrar el minimo flujo posible y su costo.
	\item \textbf{Costo por unidad}: si los costos son muy grandes, aniadir verificacion.
\end{itemize}

\paragraph{Trampas y errores frecuentes.}
\begin{itemize}\setlength\itemsep{0pt}
	\item El \texttt{pot[]} se actualiza \textbf{al final} de cada iteracion (no durante).
	\item Las back-edges tienen \textbf{costo negativo}, esencial para corregir elecciones previas.
	\item Si los costos son muy grandes, considerar modular arithmetic.
	\item Para problemas con flujo exacto, aniadir verificacion post-algoritmo.
\end{itemize}

\paragraph{Codigo.}
Ver \texttt{cpp.tex}, seccion \textit{Min-Cost Max-Flow}.

\subsection{Especales}

\subsubsection{2-SAT (implication graph + SCC)}

\paragraph{Idea intuitiva.}
Resuelve formulas booleanas en \textbf{CNF} con clausulas de \textbf{2 literales}: $(l_1 \lor l_2)$. Cada variable $x$ se representa con dos nodos: $x$ (verdadero) y $\neg x$ (falso). Cada clausula $l_1 \lor l_2$ se descompone en dos implicaciones: $\neg l_1 \to l_2$ y $\neg l_2 \to l_1$. Si en el grafo de implicaciones $\neg x$ y $x$ estan en la misma SCC, la formula es \textbf{UNSAT}. La demostracion: si $\neg x \to x$ en el grafo, entonces $\neg x$ implica $x$, asi que $\neg x$ debe ser falso y $x$ verdadero; pero entonces $\neg x$ es falso implica una contradiccion.

\paragraph{Propiedades clave (invariantes).}
\begin{itemize}\setlength\itemsep{0pt}
	\item Cada literal es un nodo; total $2n$ nodos.
	\item Kosaraju/Tarjan dan las SCCs.
	\item Si SAT: asignar $x = $ (topological order del condensation).
	\item Clausulas \texttt{setTrue(x)} se modelan como \texttt{implies(!x, x)}.
	\item La formula es SAT sii no hay variable $x$ con $x$ y $\neg x$ en la misma SCC.
\end{itemize}

\paragraph{Codigo clave.}
Conversion de clausulas a implicaciones:
\begin{verbatim}
// Clausula (a OR b) -> implica ~a -> b Y ~b -> a
addImplication(not(a), b);
addImplication(not(b), a);
// Forzar x: implica ~x -> x
if (force) addImplication(not(x), x);
\end{verbatim}

\paragraph{Complejidad.}
\begin{itemize}\setlength\itemsep{0pt}
	\item $O(V + E) = O(N + M)$ con Kosaraju o Tarjan.
	\item Memoria $O(V + E)$.
\end{itemize}

\paragraph{Donde tocar para modificar.}
\begin{itemize}\setlength\itemsep{0pt}
	\item \textbf{Forzar variables}: aniadir \texttt{setTrue} / \texttt{setFalse}.
	\item \textbf{3-SAT}: NP-completo; no usar este algoritmo.
	\item \textbf{Contar soluciones}: las SCCs condensation dan una formula para contar; multiplicar las formas de asignar variables en cada componente no-trivial.
	\item \textbf{Output de la asignacion}: aniadir el vector \texttt{assignment[]}.
\end{itemize}

\paragraph{Trampas y errores frecuentes.}
\begin{itemize}\setlength\itemsep{0pt}
	\item La asignacion usa el \textbf{orden topologico inverso} del condensation: \texttt{val[i] = comp[2i] < comp[2i+1]}.
	\item El ID de $\neg x$ es \texttt{2*x + 1} (o similar); verificar consistencia.
	\item Si hay clausulas con literales repetidos (e.g., $(x \lor x)$), tratar como $(x)$ simple.
\end{itemize}

\paragraph{Codigo.}
Ver \texttt{cpp.tex}, seccion \textit{2-SAT (implication graph + SCC)}.

\vspace{0.5em}

\subsubsection{Euler tour / Hierholzer (camino y circuito)}

\paragraph{Idea intuitiva.}
Encuentra un \textbf{camino o circuito euleriano} (que usa cada arista exactamente una vez). \textbf{Condiciones}:
\begin{itemize}\setlength\itemsep{0pt}
	\item \textbf{No dirigido}: todos los grados pares (circuito) o exactamente 2 impares (camino, empieza en uno, termina en el otro). Ademas, conexo sobre el subgrafo con aristas.
	\item \textbf{Dirigido}: $\text{indeg}(v) = \text{outdeg}(v)$ para todos (circuito) o exactamente un nodo con $\text{outdeg} = \text{indeg} + 1$ y otro con $\text{indeg} = \text{outdeg} + 1$ (camino).
\end{itemize}

\paragraph{Hierholzer}: elegir un ciclo desde el inicio, mientras haya ciclos ``laterales'' desde un nodo del camino, mergearlos. Implementacion con pila: avanzar por aristas, al no poder mas aniadir a la respuesta y backtrack. La demostracion de $O(V + E)$: cada arista se procesa una vez (entra y sale del stack una vez).

\paragraph{Propiedades clave (invariantes).}
\begin{itemize}\setlength\itemsep{0pt}
	\item Multigrafos permitidos; las aristas se marcan como usadas.
	\item No dirigido: cada arista aparece \textbf{dos veces} en \texttt{adj} (ida y vuelta).
	\item La construccion del path con ids consistentes es importante.
	\item El resultado es unico salvo orden de las aristas en multigrafos.
\end{itemize}

\paragraph{Codigo clave.}
Hierholzer con pila:
\begin{verbatim}
vector<int> path;
stack<int> st;
st.push(start);
while (!st.empty()) {
    int v = st.top();
    while (!adj[v].empty() && used[adj[v].back()]) adj[v].pop_back();
    if (adj[v].empty()) {
        path.push_back(v);
        st.pop();
    } else {
        int e = adj[v].back();
        used[e] = true;
        int u = edges[e].to(v);
        adj[v].pop_back();
        st.push(u);
    }
}
reverse(path.begin(), path.end());
\end{verbatim}

\paragraph{Complejidad.}
\begin{itemize}\setlength\itemsep{0pt}
	\item $O(V + E)$.
	\item Memoria $O(V + E)$.
\end{itemize}

\paragraph{Donde tocar para modificar.}
\begin{itemize}\setlength\itemsep{0pt}
	\item \textbf{Reconstruir con multigrafo}: aniadir ids consistentes antes del algoritmo.
	\item \textbf{Aplicar a ``de Bruijn'' o problemas de strings}: modelar como multigrafo y buscar Euler.
	\item \textbf{Dirigido vs no dirigido}: el snippet tiene ambas versiones.
	\item \textbf{Letter addition}: modelar palabras como aristas y buscar el superstring.
\end{itemize}

\paragraph{Trampas y errores frecuentes.}
\begin{itemize}\setlength\itemsep{0pt}
	\item Si el grafo no es conexo (sobre las aristas), no existe euleriano. Verificar primero.
	\item Para multigrafos, los IDs de aristas son cruciales para distinguirlas.
	\item Si el problema es ``encuentra todos los euler tours'', el algoritmo es exponencial.
\end{itemize}

\paragraph{Codigo.}
Ver \texttt{cpp.tex}, seccion \textit{Euler tour / Hierholzer (camino y circuito)}.

\vspace{0.5em}

\subsubsection{Hamiltonian path / cycle (bitmask DP + backtracking)}

\paragraph{Idea intuitiva.}
Un \textbf{camino hamiltoniano} visita cada vertice \textbf{exactamente una vez}. \textbf{Determinar} si existe es NP-completo en general, pero con $V \le 20$ se puede hacer DP con bitmask en $O(V^2 \cdot 2^V)$. \textbf{Teoremas suficientes}: \textbf{Dirac} (grado $\ge n/2$) y \textbf{Ore} ($\text{deg}(u) + \text{deg}(v) \ge n$ para todo par no adyacente). La idea de la DP: cada estado es (conjunto de visitados, ultimo vertice); la transicion es agregar un vecino no visitado.

\paragraph{Propiedades clave (invariantes).}
\begin{itemize}\setlength\itemsep{0pt}
	\item DP: $dp[mask][v]$ = true si hay un camino que visita exactamente \texttt{mask} y termina en $v$.
	\item TSP: variante con pesos minimos (Hamilton cycle minimo).
	\item Reconstruccion: guardar \texttt{par[mask][v]} (predecesor).
	\item $2^V$ estados con $V$ bits cada uno; factible para $V \le 20$.
\end{itemize}

\paragraph{Codigo clave.}
DP con bitmask:
\begin{verbatim}
dp[1][0] = 0;  // empezar en 0
for (int mask = 1; mask < (1<<n); ++mask)
    for (int v = 0; v < n; ++v) if (dp[mask][v] < INF)
        for (int u = 0; u < n; ++u)
            if (!(mask & (1<<u)))
                dp[mask | (1<<u)][u] = min(dp[mask | (1<<u)][u],
                                            dp[mask][v] + w[v][u]);
// TSP: respuesta = min(dp[(1<<n)-1][v] + w[v][0])
\end{verbatim}

\paragraph{Complejidad.}
\begin{itemize}\setlength\itemsep{0pt}
	\item $O(V^2 \cdot 2^V)$ tiempo, $O(V \cdot 2^V)$ memoria.
	\item Para $V = 20$: $\sim 4 \cdot 10^8$ operaciones (lento pero factible).
\end{itemize}

\paragraph{Donde tocar para modificar.}
\begin{itemize}\setlength\itemsep{0pt}
	\item \textbf{TSP asimetrico}: aniadir pesos $w[v][u]$ (no necesariamente simetricos).
	\item \textbf{Hamilton con inicio obligatorio}: fijar \texttt{src} en el DP inicial.
	\item \textbf{Bitmask DP mas compacto}: usar \texttt{long long} como mascara si $V \le 64$.
	\item \textbf{TSP optimo aproximado}: MST, 2-opt, simulated annealing.
	\item \textbf{Reconstruir el path}: usar \texttt{par[mask][v]} y backtrack.
\end{itemize}

\paragraph{Trampas y errores frecuentes.}
\begin{itemize}\setlength\itemsep{0pt}
	\item TSP asume grafo \textbf{completo} (con \texttt{LINF} para aristas faltantes); si no, aniadir chequeo.
	\item La condicion ``\texttt{mask == (1<<n) - 1}'' cierra el ciclo retornando al origen.
	\item Para $V > 20$, la memoria $2^V$ se vuelve prohibitiva; usar meet-in-the-middle o heuristicas.
\end{itemize}

\paragraph{Codigo.}
Ver \texttt{cpp.tex}, seccion \textit{Hamiltonian path / cycle (bitmask DP + backtracking)}.

% ============================================================
\section{DP}

\subsubsection{Knapsack 0/1}

\paragraph{Idea intuitiva.}
Dado un conjunto de $N$ objetos con peso $w_i$ y valor $v_i$, maximizar $\sum v_i$ con la restriccion $\sum w_i \le W$. La recurrencia clasica:
\[ dp[i][j] = \max(dp[i-1][j],\, dp[i-1][j - w_i] + v_i) \quad \text{si } j \ge w_i. \]
La optimizacion clave: la DP es 1-dimensional si iteramos $j$ \textbf{de mayor a menor} (asi no usamos el objeto $i$ dos veces). La razon de la iteracion inversa: en cada iteracion queremos el estado del ``item anterior''; iterar de menor a mayor contaminaria \texttt{dp[j - w_i]} con el item actual.

\paragraph{Propiedades clave (invariantes).}
\begin{itemize}\setlength\itemsep{0pt}
	\item $dp[j]$ al final = maximo valor con peso $\le j$.
	\item El bucle inverso \texttt{for (j = W; j >= w[i]; j--)} es lo que distingue 0/1 de unbounded.
	\item Solucion: maximo valor con cualquier peso $\le W$.
	\item La DP es exacta; no hay aproximacion.
\end{itemize}

\paragraph{Codigo clave.}
La version 1D optimizada:
\begin{verbatim}
vector<long long> dp(W + 1, 0);
for (int i = 0; i < N; ++i)
    for (int j = W; j >= w[i]; --j)  // invertido
        dp[j] = max(dp[j], dp[j - w[i]] + v[i]);
return dp[W];
\end{verbatim}

\paragraph{Complejidad.}
\begin{itemize}\setlength\itemsep{0pt}
	\item $O(N \cdot W)$ tiempo, $O(W)$ memoria.
	\item Constante muy pequena (solo sumas y max).
\end{itemize}

\paragraph{Donde tocar para modificar.}
\begin{itemize}\setlength\itemsep{0pt}
	\item \textbf{Reconstruir la seleccion}: aniadir \texttt{take[i][j]} (booleano o bool) en una version 2D.
	\item \textbf{Items con peso 0}: edge case; aniadir guard.
	\item \textbf{Knapsack con restricciones de grupos}: ``elegir a lo sumo 1 de cada grupo'' se modela como 0/1 con un peso extra.
	\item \textbf{Tamano de $W$}: si $W \le 10^9$ pero $N$ es chico, usar meet-in-the-middle.
	\item \textbf{Cambiar max por min}: aniadir \texttt{INF} en la inicializacion.
\end{itemize}

\paragraph{Trampas y errores frecuentes.}
\begin{itemize}\setlength\itemsep{0pt}
	\item Overflow en \texttt{dp[j - w[i]] + v[i]}; usar \texttt{long long} si los valores son grandes.
	\item El bucle \textbf{debe} ir de mayor a menor.
	\item Si los pesos son grandes y $N$ pequeno, conviene meet-in-the-middle.
	\item Para Knapsack exacto ($\sum w_i = W$), aniadir check al final.
\end{itemize}

\paragraph{Codigo.}
Ver \texttt{cpp.tex}, seccion \textit{Knapsack 0/1}.

\vspace{0.5em}

\subsubsection{Knapsack unbounded}

\paragraph{Idea intuitiva.}
Igual que Knapsack 0/1 pero \textbf{cada item se puede tomar multiples veces}. La recurrencia es similar:
\[ dp[j] = \max(dp[j],\, dp[j - w_i] + v_i), \]
pero ahora el bucle va de \textbf{menor a mayor} (porque tomar el item $i$ no afecta corridas futuras del mismo item). La idea: tras tomar el item $i$, queremos volver a tener la opcion de tomarlo; iterar de menor a mayor permite que el item ``se acumule''.

\paragraph{Propiedades clave (invariantes).}
\begin{itemize}\setlength\itemsep{0pt}
	\item El bucle \texttt{for (j = w[i]; j <= W; j++)} es lo que distingue unbounded.
	\item A diferencia de 0/1, no hay riesgo de ``usar dos veces el mismo item''.
	\item La DP es exacta; cada ``subconjunto'' se cuenta correctamente.
\end{itemize}

\paragraph{Codigo clave.}
La version 1D con bucle directo:
\begin{verbatim}
vector<long long> dp(W + 1, 0);
for (int i = 0; i < N; ++i)
    for (int j = w[i]; j <= W; ++j)  // directo (no invertido)
        dp[j] = max(dp[j], dp[j - w[i]] + v[i]);
return dp[W];
\end{verbatim}

\paragraph{Complejidad.}
\begin{itemize}\setlength\itemsep{0pt}
	\item $O(N \cdot W)$ tiempo, $O(W)$ memoria.
	\item Constante muy pequena.
\end{itemize}

\paragraph{Donde tocar para modificar.}
\begin{itemize}\setlength\itemsep{0pt}
	\item \textbf{Reconstruir cuantos de cada item}: aniadir \texttt{cnt[i]} y backtrack.
	\item \textbf{Cambiar moneda minima por maxima}: el snippet maximiza valor; para minimizar monedas, cambiar a \texttt{min} con inicializacion en \texttt{INF}.
	\item \textbf{Coin change}: si todos los pesos son positivos y queremos ``minimo numero de monedas'', es unbounded con min.
\end{itemize}

\paragraph{Trampas y errores frecuentes.}
\begin{itemize}\setlength\itemsep{0pt}
	\item Misma que Knapsack 0/1 con overflow.
	\item El orden de iteracion de items no afecta el resultado (convexo).
	\item Si se quiere ``maximo valor con exactamente $k$ items'', aniadir segunda dimension.
\end{itemize}

\paragraph{Codigo.}
Ver \texttt{cpp.tex}, seccion \textit{Knapsack unbounded}.

\vspace{0.5em}

\subsubsection{Knapsack meet-in-the-middle}

\paragraph{Idea intuitiva.}
Para $N \le 40$ y $W$ moderado, no se puede hacer $O(NW)$ directamente. Partir en dos mitades, enumerar todos los $2^{N/2}$ subconjuntos de cada mitad, y combinar. Para cada subconjunto de la segunda mitad con peso $sw_2$, encontrar el maximo valor de la primera mitad con peso $\le W - sw_2$. La idea: ``divide y venceras'' sobre el tamano de la entrada; cada mitad tiene mitad del espacio, asi que la combinacion es geometrica en lugar de multiplicativa.

\paragraph{Propiedades clave (invariantes).}
\begin{itemize}\setlength\itemsep{0pt}
	\item Hay $2^{N/2}$ subconjuntos por mitad; total $O(2^{N/2})$.
	\item \texttt{sums} guarda los valores de subconjuntos de la primera mitad ordenados (o preprocesados).
	\item Para cada subconjunto de la segunda mitad, busqueda lineal o binaria en \texttt{sums}.
	\item La mitad mas grande (con $N/2 + 1$ items para $N$ impar) tiene $2^{N/2+1}$ subconjuntos.
\end{itemize}

\paragraph{Codigo clave.}
Enumerar subconjuntos de la primera mitad:
\begin{verbatim}
int n1 = N / 2, n2 = N - n1;
vector<pair<long long, long long>> sums1;
for (int mask = 0; mask < (1 << n1); ++mask) {
    long long peso = 0, valor = 0;
    for (int i = 0; i < n1; ++i) if (mask & (1 << i)) {
        peso += w[i]; valor += v[i];
    }
    sums1.push_back({peso, valor});
}
sort(sums1.begin(), sums1.end());  // por peso
// eliminar duplicados dominados
\end{verbatim}

\paragraph{Complejidad.}
\begin{itemize}\setlength\itemsep{0pt}
	\item $O(2^{N/2})$ tiempo y memoria (mejor que $O(NW)$ cuando $W > 2^{N/2}$).
	\item Para $N = 40$: $\sim 2 \cdot 10^6$ subconjuntos, factible.
\end{itemize}

\paragraph{Donde tocar para modificar.}
\begin{itemize}\setlength\itemsep{0pt}
	\item \textbf{Tamano mayor}: si $N$ es mayor pero $W$ es chico, conviene iterar por peso, no por subconjuntos.
	\item \textbf{Varias mochilas}: adaptar para multiples $W$.
	\item \textbf{Optimizar el ``combinado''}: usar \texttt{upper\_bound} para encontrar el maximo valor con peso $\le \text{rem}.
	\item \textbf{Eliminar pares dominados}: si un par (peso, valor) tiene otro con menos peso y mas valor, eliminarlo.
\end{itemize}

\paragraph{Trampas y errores frecuentes.}
\begin{itemize}\setlength\itemsep{0pt}
	\item $N$ impar: una mitad tiene $N/2 + 1$ items.
	\item Overflow al sumar pesos/valores en \texttt{long long}.
	\item Sin deduplicacion, el algoritmo es correcto pero mas lento.
\end{itemize}

\paragraph{Codigo.}
Ver \texttt{cpp.tex}, seccion \textit{Knapsack meet-in-the-middle}.

\vspace{0.5em}

\subsubsection{LIS O(n log n)}

\paragraph{Idea intuitiva.}
La \textbf{Longest Increasing Subsequence} se calcula en $O(n \log n)$ con el truco del \textbf{patience sorting}: mantener un arreglo \texttt{dp} donde \texttt{dp[len]} es el minimo ultimo valor de una IS de largo \texttt{len}. Para cada $x$, encontrar el primer \texttt{dp[len] >= x} y reemplazarlo. El largo final es la LIS. La demostracion: si hay dos IS del mismo largo, la que termina en un valor menor es ``mejor'' (puede extenderse mas facilmente).

\paragraph{Propiedades clave (invariantes).}
\begin{itemize}\setlength\itemsep{0pt}
	\item \texttt{dp} es estrictamente creciente.
	\item \texttt{lower\_bound} da LIS estricta; \texttt{upper\_bound} da no-estricta ($\le$).
	\item El algoritmo es \textbf{online}: procesa elementos en orden.
	\item La longitud de \texttt{dp} es la LIS al final.
\end{itemize}

\paragraph{Codigo clave.}
La version online:
\begin{verbatim}
vector<int> dp;
for (int x : a) {
    auto it = lower_bound(dp.begin(), dp.end(), x);
    if (it == dp.end()) dp.push_back(x);
    else *it = x;
}
// dp.size() = longitud de LIS estricta
\end{verbatim}

\paragraph{Complejidad.}
\begin{itemize}\setlength\itemsep{0pt}
	\item $O(n \log n)$.
	\item Memoria $O(n)$.
\end{itemize}

\paragraph{Donde tocar para modificar.}
\begin{itemize}\setlength\itemsep{0pt}
	\item \textbf{Reconstruir la LIS}: aniadir \texttt{parent[]} y \texttt{posInDp[]} para backtrack.
	\item \textbf{Contar LIS}: DP clasica $O(n^2)$ o segment tree sobre $dp$.
	\item \textbf{LDS / Longest non-increasing}: invertir la condicion en \texttt{lower\_bound} o multiplicar por -1.
	\item \textbf{Longest bitonic subsequence}: combinar LIS desde izq y desde der.
	\item \textbf{LIS en 2D}: ordenar por una dimension y aplicar LIS en la otra.
\end{itemize}

\paragraph{Trampas y errores frecuentes.}
\begin{itemize}\setlength\itemsep{0pt}
	\item \texttt{lower\_bound} vs \texttt{upper\_bound}: si tu comparacion es $\le$ vs $<$, cambia.
	\item Si los elementos son \texttt{int} pero podrian ser negativos, aniadir offset para usar como indices.
	\item Para LIS reconstruida, aniadir el array \texttt{parent} que apunta al anterior en la IS.
\end{itemize}

\paragraph{Codigo.}
Ver \texttt{cpp.tex}, seccion \textit{LIS O(n log n)}.

\vspace{0.5em}

\subsubsection{LCS / Edit distance}

\paragraph{Idea intuitiva.}
\textbf{LCS} (Longest Common Subsequence): para dos strings $a$, $b$, $dp[i][j]$ = LCS de $a[0..i)$ y $b[0..j)$. Recurrencia:
\[ dp[i][j] = \begin{cases} dp[i-1][j-1] + 1 & \text{si } a[i-1] = b[j-1] \\ \max(dp[i-1][j], dp[i][j-1]) & \text{si no} \end{cases} \]

\textbf{Edit distance} (Levenshtein): numero minimo de operaciones (insert, delete, replace) para convertir $a$ en $b$. Recurrencia similar con un $+1$ adicional. La idea: ``alinear'' los strings eligiendo matches/mismatches/insert/delete.

\paragraph{Propiedades clave (invariantes).}
\begin{itemize}\setlength\itemsep{0pt}
	\item LCS y Edit Distance son DP $O(nm)$.
	\item Memoria se puede reducir a $O(\min(n, m))$ con rolling array.
	\item Edit distance con peso por operacion distinto: ajustar el \texttt{+1} por el peso.
	\item La subestructura optima: la respuesta para prefijos se construye de prefijos menores.
\end{itemize}

\paragraph{Codigo clave.}
La recurrencia clasica de LCS:
\begin{verbatim}
vector<vector<int>> dp(n + 1, vector<int>(m + 1, 0));
for (int i = 1; i <= n; ++i)
    for (int j = 1; j <= m; ++j)
        if (a[i-1] == b[j-1]) dp[i][j] = dp[i-1][j-1] + 1;
        else dp[i][j] = max(dp[i-1][j], dp[i][j-1]);
return dp[n][m];
\end{verbatim}

\paragraph{Complejidad.}
\begin{itemize}\setlength\itemsep{0pt}
	\item $O(NM)$ tiempo y memoria.
	\item Para $N, M \le 5000$: $\sim 2.5 \cdot 10^7$ operaciones.
\end{itemize}

\paragraph{Donde tocar para modificar.}
\begin{itemize}\setlength\itemsep{0pt}
	\item \textbf{Reconstruir la LCS / las operaciones}: aniadir \texttt{prev[i][j]} para backtrack.
	\item \textbf{LCS de varios strings}: NP-hard para $\ge 3$ strings; en la practica se hace con bitmask.
	\item \textbf{Edit distance con costos distintos}: cambiar el \texttt{+1} por el costo apropiado.
	\item \textbf{Rolling array}: si memoria es problema, mantener solo 2 filas.
	\item \textbf{Hirschberg}: algoritmo $O(NM)$ tiempo, $O(\min(N,M))$ memoria, con reconstruccion.
\end{itemize}

\paragraph{Trampas y errores frecuentes.}
\begin{itemize}\setlength\itemsep{0pt}
	\item Indices en la recurrencia: \texttt{a[i-1]} y \texttt{b[j-1]} (no \texttt{a[i]}).
	\item Para reconstruccion, el caso \texttt{a[i-1] == b[j-1]} tiene prioridad.
	\item Si los strings son muy largos, considerar Hirschberg o bitset optimization.
\end{itemize}

\paragraph{Codigo.}
Ver \texttt{cpp.tex}, seccion \textit{LCS / Edit distance}.

\vspace{0.5em}

\subsubsection{Bitmask DP over subsets}

\paragraph{Idea intuitiva.}
Cuando el estado es un \textbf{subconjunto} de $\{0, \dots, N-1\}$, se representa como una mascara de $N$ bits. Para iterar sobre todos los subconjuntos de \texttt{mask}:
\[ \text{for (sub = mask; sub; sub = (sub - 1) \& mask)} \{ \dots \} \]
Para iterar sobre todos los superconjuntos de \texttt{mask} en $[0, U)$:
\[ \text{for (sup = mask; sup < U; sup = (sup + 1) | mask)} \{ \dots \} \]
La idea: las mascaras forman un lattice de subconjuntos; los trucos aritmeticos (\texttt{(sub-1) \& mask} y \texttt{(sup+1) | mask}) enumeran todo el lattice eficientemente.

\paragraph{Propiedades clave (invariantes).}
\begin{itemize}\setlength\itemsep{0pt}
	\item Iterar subsets: $O(2^k)$ para una mascara con $k$ bits.
	\item \texttt{sub \& mask == sub} siempre; \texttt{mask - sub} no siempre da un subconjunto.
	\item \texttt{\_\_builtin\_popcount(mask)} cuenta los bits.
	\item \texttt{(sub - 1) \& mask} itera subsets de \texttt{mask} en orden decreciente de inclusion.
\end{itemize}

\paragraph{Codigo clave.}
Iteracion sobre subsets de una mascara:
\begin{verbatim}
// Todos los subsets de mask
for (int sub = mask; sub; sub = (sub - 1) & mask) {
    // procesar sub
}
// Incluir sub = 0 si es necesario
sub = 0;  // procesar primero
for (int sup = mask; sup < (1 << N); sup = (sup + 1) | mask) {
    // procesar sup (todos los superconjuntos)
}
\end{verbatim}

\paragraph{Complejidad.}
\begin{itemize}\setlength\itemsep{0pt}
	\item $O(N \cdot 2^N)$ para una DP estandar sobre todos los subsets.
	\item Iterar subsets de \texttt{mask}: $O(2^{|mask|})$.
\end{itemize}

\paragraph{Donde tocar para modificar.}
\begin{itemize}\setlength\itemsep{0pt}
	\item \textbf{TSP / Hamiltonian path}: bitmask DP (ver tambien \textit{Hamiltonian path} en \texttt{cpp.tex} seccion Grafos).
	\item \textbf{Set cover}: iterar subsets.
	\item \textbf{DP con supermask}: si el estado incluye ``todos los elementos que \textbf{no} hemos usado'', usar supermask iteration.
	\item \textbf{Bitmasks grandes}: con $N > 20$, aniadir compresion o simulacion con segmentos.
\end{itemize}

\paragraph{Trampas y errores frecuentes.}
\begin{itemize}\setlength\itemsep{0pt}
	\item El orden de iteracion de subsets de \texttt{mask} es decreciente (de \texttt{mask} a $0$).
	\item \texttt{mask = 0} no tiene subsets no triviales; el bucle lo saltea correctamente.
	\item El truco \texttt{(sub - 1) \& mask} requiere cuidado con \texttt{sub = 0} (termina el bucle).
	\item Superconjuntos: \texttt{sup < (1 << N)} debe estar bien definido.
\end{itemize}

\paragraph{Codigo.}
Ver \texttt{cpp.tex}, seccion \textit{Bitmask DP over subsets}.

\vspace{0.5em}

\subsubsection{SOS DP}

\paragraph{Idea intuitiva.}
Para calcular, para cada \texttt{mask}, la suma de $f(\text{submask})$ sobre todos los $\text{submask} \subseteq \text{mask}$:
\[ dp[mask] = \sum_{\text{sub} \subseteq mask} f(\text{sub}). \]
Se computa en $O(N \cdot 2^N)$ mediante un loop por bit: para cada $i$, sumar a $dp[mask]$ la contribucion de $dp[mask \oplus (1 << i)]$ cuando $mask$ tiene el bit $i$. La idea: cada subconjunto $s \subseteq m$ se obtiene quitando bits uno a uno; asi, sumar incrementalmente captura todos los subconjuntos en $O(N \cdot 2^N)$.

\paragraph{Propiedades clave (invariantes).}
\begin{itemize}\setlength\itemsep{0pt}
	\item Inicializar $dp = f$.
	\item Para cada bit $i$: \texttt{if (mask \& (1 << i)) dp[mask] += dp[mask \^{} (1 << i)]}.
	\item Variantes: supersets DP (sumar sobre superconjuntos).
	\item El orden de iteracion (bit externo, mascara interno) es crucial.
\end{itemize}

\paragraph{Codigo clave.}
SOS DP (Sum over Subsets):
\begin{verbatim}
vector<long long> dp(1 << N);
for (int i = 0; i < (1 << N); ++i) dp[i] = f[i];
for (int bit = 0; bit < N; ++bit)
    for (int mask = 0; mask < (1 << N); ++mask)
        if (mask & (1 << bit))
            dp[mask] += dp[mask ^ (1 << bit)];
// dp[mask] = suma de f(sub) para sub ⊆ mask
\end{verbatim}

\paragraph{Complejidad.}
\begin{itemize}\setlength\itemsep{0pt}
	\item $O(N \cdot 2^N)$.
	\item Memoria $O(2^N)$.
\end{itemize}

\paragraph{Donde tocar para modificar.}
\begin{itemize}\setlength\itemsep{0pt}
	\item \textbf{Max/min en lugar de suma}: cambiar \texttt{+=} por la operacion correspondiente.
	\item \textbf{Supersets DP}: iterar \texttt{mask} de $0$ a $2^N$ y aniadir \texttt{dp[mask] += dp[mask | (1<<i)]}.
	\item \textbf{XOR convolution} (subset convolution): version mas avanzada.
	\item \textbf{Aplicaciones}: contar subconjuntos con propiedades especificas.
\end{itemize}

\paragraph{Trampas y errores frecuentes.}
\begin{itemize}\setlength\itemsep{0pt}
	\item Si $N$ es 30 o mas, $2^N$ no entra en memoria; usar trucos o reducir $N$.
	\item El sentido del XOR (sumar a \texttt{mask \^{} bit}) es para ``agregar el bit $i$ al subconjunto''.
	\item En supersets DP, la condicion es \texttt{(mask \& (1<<bit)) == 0}.
\end{itemize}

\paragraph{Codigo.}
Ver \texttt{cpp.tex}, seccion \textit{SOS DP}.

\vspace{0.5em}

\subsubsection{Digit DP template}

\paragraph{Idea intuitiva.}
Contar numeros en $[0, N]$ que cumplen una propiedad sobre sus digitos. La idea: DP por posicion, manteniendo:
\begin{itemize}\setlength\itemsep{0pt}
	\item \texttt{pos}: posicion actual (de izquierda a derecha).
	\item \texttt{tight}: si los digitos ya colocados son exactamente los de $N$ (limite superior estricto).
	\item \texttt{state}: estado adicional (digitos usados, suma, etc.).
\end{itemize}
La clave: la condicion \texttt{tight} permite acotar correctamente el siguiente digito; sin ella, la DP contaria incorrectamente.

\paragraph{Propiedades clave (invariantes).}
\begin{itemize}\setlength\itemsep{0pt}
	\item Si \texttt{tight} es true, el siguiente digito esta acotado por el de $N$; si no, puede ser 0-9.
	\item Memoria: $O(D \cdot 2 \cdot \text{states})$.
	\item Caso base: si \texttt{pos == D}, devolver 1 si el estado cumple la propiedad, sino 0.
	\item El bit \texttt{started} distingue ``numero con menos digitos'' (con leading zeros).
\end{itemize}

\paragraph{Codigo clave.}
La estructura basica:
\begin{verbatim}
long long dp(int pos, bool tight, State s) {
    if (pos == D) return check(s) ? 1 : 0;
    if (memo[pos][tight][s] != -1) return memo;
    int limit = tight ? digits[pos] : 9;
    long long ans = 0;
    for (int d = 0; d <= limit; ++d)
        ans += dp(pos + 1, tight && (d == limit), update(s, d, pos));
    return memo[pos][tight][s] = ans;
}
\end{verbatim}

\paragraph{Complejidad.}
\begin{itemize}\setlength\itemsep{0pt}
	\item $O(D \cdot 10 \cdot \text{states})$ donde $D$ es el numero de digitos.
	\item Para $D \le 19$ (hasta $10^{18}$): $\sim 200 \cdot \text{states}$.
\end{itemize}

\paragraph{Donde tocar para modificar.}
\begin{itemize}\setlength\itemsep{0pt}
	\item \textbf{Cambiar la propiedad}: modificar \texttt{updateState} y la condicion base.
	\item \textbf{Multiples estados}: aniadir mas dimensiones al \texttt{memo}.
	\item \textbf{Leading zeros}: aniadir un flag \texttt{started} para distinguir numeros con menos digitos.
	\item \textbf{Rango $[L, R]$}: calcular $f(R) - f(L-1)$.
\end{itemize}

\paragraph{Trampas y errores frecuentes.}
\begin{itemize}\setlength\itemsep{0pt}
	\item Olvidar resetear \texttt{memo} entre queries.
	\item Si la cantidad de estados es grande, el \texttt{memo[pos][tight][state]} puede ser muy grande; usar \texttt{map}.
	\item La condicion \texttt{started} es crucial para numeros con menos digitos que $D$.
	\item Si $N$ es muy grande ($> 10^{18}$), aniadir soporte para long long.
\end{itemize}

\paragraph{Codigo.}
Ver \texttt{cpp.tex}, seccion \textit{Digit DP template}.

\vspace{0.5em}

\subsubsection{Tree DP (reroot)}

\paragraph{Idea intuitiva.}
Calcular alguna propiedad asumiendo que \textbf{cada nodo es la raiz} del arbol. La tecnica clasica de \textbf{rerooting} en dos pasadas:
\begin{itemize}\setlength\itemsep{0pt}
	\item \textbf{DFS 1}: desde una raiz arbitraria, calcular \texttt{dp[v]} para el subarbol de $v$.
	\item \textbf{DFS 2}: para cada hijo $u$ de $v$, calcular \texttt{up[u]} = respuesta si la raiz fuera $u$.
\end{itemize}
La demostracion: tras la primera DFS, cada nodo conoce ``lo que aporta su subarbol''; en la segunda, se transmite ``lo que aporta el resto del arbol''.

\paragraph{Propiedades clave (invariantes).}
\begin{itemize}\setlength\itemsep{0pt}
	\item \texttt{dp[v]} depende solo del subarbol de $v$ (con la raiz original).
	\item \texttt{up[u]} se calcula de manera que cuando la raiz es $u$, ``todo lo que estaba en $v$ se ajusta''.
	\item Formula para reroot depende del problema; el snippet usa el ejemplo ``suma de distancias''.
	\item La combinacion \texttt{dp[u] + up[u]} da la respuesta con raiz $u$.
\end{itemize}

\paragraph{Codigo clave.}
Reroot para suma de distancias:
\begin{verbatim}
// dfs1: dp[v] = suma de distancias desde v al subarbol
void dfs1(int v, int p) {
    dp[v] = 0;
    for (int u : adj[v]) if (u != p) {
        dfs1(u, v);
        dp[v] += dp[u] + sz[u];  // cada hijo aporta su subarbol
    }
}
// dfs2: up[u] = distancia desde el resto del arbol a u
void dfs2(int v, int p) {
    for (int u : adj[v]) if (u != p) {
        up[u] = up[v] + dp[v] - dp[u] - sz[u] + (n - sz[u]);
        dfs2(u, v);
    }
}
\end{verbatim}

\paragraph{Complejidad.}
\begin{itemize}\setlength\itemsep{0pt}
	\item $O(N)$ para ambas pasadas.
	\item Cada arista se procesa un numero constante de veces.
\end{itemize}

\paragraph{Donde tocar para modificar.}
\begin{itemize}\setlength\itemsep{0pt}
	\item \textbf{Cambiar la cantidad a calcular}: modificar la operacion en \texttt{dfs1} y la formula en \texttt{dfs2}.
	\item \textbf{Arbol con pesos}: aniadir \texttt{weights[v]} y ajustar.
	\item \textbf{Re-root en arboles con ciclos}: requiere otra tecnica (DSU on tree).
	\item \textbf{Rooted trees diferentes}: si la propiedad no es invariante por raiz, usar otra tecnica.
\end{itemize}

\paragraph{Trampas y errores frecuentes.}
\begin{itemize}\setlength\itemsep{0pt}
	\item La formula de \texttt{up[u]} debe sumar y restar correctamente las contribuciones de $u$ y del resto.
	\item Para arboles con pesos, aniadir el peso en las formulas.
	\item Verificar con ejemplos pequenos antes de generalizar.
\end{itemize}

\paragraph{Codigo.}
Ver \texttt{cpp.tex}, seccion \textit{Tree DP (reroot)}.

% ============================================================
\section{Geometria}
% ============================================================

\subsection{Puntos y segmentos}

\subsubsection{Point + cross + dot}

\paragraph{Idea intuitiva.}
La estructura \texttt{Point} representa un punto en 2D con coordenadas enteras. Las dos operaciones vectoriales clave son:
\begin{itemize}\setlength\itemsep{0pt}
	\item \textbf{Cross product} $P \times Q = P_x Q_y - P_y Q_x$. Da el ``area con signo'' del paralelogramo formado por $P$ y $Q$. Usado para orientacion.
	\item \textbf{Dot product} $P \cdot Q = P_x Q_x + P_y Q_y$. Da el producto de las longitudes por el coseno del angulo. Usado para proyeccion.
\end{itemize}
La interpretacion geometrica del cross product: es el area con signo del paralelogramo; el signo indica la orientacion relativa.

\paragraph{Propiedades clave (invariantes).}
\begin{itemize}\setlength\itemsep{0pt}
	\item \texttt{cross(a, b) > 0} si $b$ esta a la izquierda de $a$ (sistema de coordenadas usual).
	\item \texttt{cross(a, b) = 0} si son colineales.
	\item \texttt{dot(a, b) > 0} si el angulo es agudo.
	\item El snippet define \texttt{cross} respecto a \texttt{this}, util para orientacion de 3 puntos.
	\item $|cross|^2 + |dot|^2 = |P|^2 \cdot |Q|^2$ (identidad de Lagrange).
\end{itemize}

\paragraph{Codigo clave.}
Las dos operaciones vectoriales:
\begin{verbatim}
long long cross(const Point& a, const Point& b) {
    return (long long)a.x * b.y - (long long)a.y * b.x;
}
long long dot(const Point& a, const Point& b) {
    return (long long)a.x * b.x + (long long)a.y * b.y;
}
\end{verbatim}

\paragraph{Complejidad.}
\begin{itemize}\setlength\itemsep{0pt}
	\item $O(1)$ por operacion.
	\item Constante muy pequena (4 multiplicaciones + 2 sumas).
\end{itemize}

\paragraph{Donde tocar para modificar.}
\begin{itemize}\setlength\itemsep{0pt}
	\item \textbf{Coordenadas doubles}: cambiar \texttt{int} por \texttt{double} o \texttt{long double}. Cuidado con la precision.
	\item \textbf{3D}: aniadir \texttt{z} y adaptar cross/dot.
	\item \textbf{Operadores adicionales}: aniadir \texttt{operator+}, \texttt{operator-}, \texttt{operator*} (escalar).
	\item \textbf{Distancia al cuadrado}: aniadir \texttt{len2()} que retorna $x^2 + y^2$.
\end{itemize}

\paragraph{Trampas y errores frecuentes.}
\begin{itemize}\setlength\itemsep{0pt}
	\item Cross product \textbf{no es} conmutativo; $P \times Q = -Q \times P$.
	\item Overflow con coordenadas grandes ($|x|, |y| > 10^4$) si se hace \texttt{x * y}.
	\item El signo de cross depende del sistema de coordenadas (en computacion, usualmente CCW positivo).
	\item Para ``igualdad'' de doubles, aniadir tolerancia \texttt{eps}.
\end{itemize}

\paragraph{Codigo.}
Ver \texttt{cpp.tex}, seccion \textit{Point + cross + dot}.

\vspace{0.5em}

\subsubsection{Interseccion de segmentos}

\paragraph{Idea intuitiva.}
Dados dos segmentos $[a, b]$ y $[c, d]$, determinar si se intersectan. La idea clasica:
\begin{itemize}\setlength\itemsep{0pt}
	\item Verificar orientaciones: $a, b, c$ y $a, b, d$ deben estar en lados opuestos de la recta $ab$; igual para $cd$ y $ab$.
	\item Caso colineal: si las rectas son paralelas (\texttt{cross} $= 0$) y los segmentos son colineales, verificar que los bounding boxes se intersectan.
\end{itemize}
La demostracion: dos segmentos se intersectan sii cada segmento ``cruza'' la linea que contiene al otro (estricto) o ambos son colineales con solapamiento.

\paragraph{Propiedades clave (invariantes).}
\begin{itemize}\setlength\itemsep{0pt}
	\item Funciona con segmentos en cualquier orientacion.
	\item Caso colineal: requiere check adicional de bounding box.
	\item Para intersectar en un \textbf{punto} especifico: aniadir division con cuidado de precision.
	\item Funciona con segmentos verticales/horizontales sin casos especiales.
\end{itemize}

\paragraph{Codigo clave.}
Test clasico de interseccion:
\begin{verbatim}
int sign(long long x) { return (x > 0) - (x < 0); }
bool intersect(Point a, Point b, Point c, Point d) {
    long long o1 = cross(b - a, c - a);
    long long o2 = cross(b - a, d - a);
    long long o3 = cross(d - c, a - c);
    long long o4 = cross(d - c, b - c);
    if (sign(o1) != sign(o2) && sign(o3) != sign(o4)) return true;
    // casos colineales
    if (o1 == 0 && onSegment(a, b, c)) return true;
    if (o2 == 0 && onSegment(a, b, d)) return true;
    if (o3 == 0 && onSegment(c, d, a)) return true;
    if (o4 == 0 && onSegment(c, d, b)) return true;
    return false;
}
\end{verbatim}

\paragraph{Complejidad.}
\begin{itemize}\setlength\itemsep{0pt}
	\item $O(1)$.
	\item Constante pequena.
\end{itemize}

\paragraph{Donde tocar para modificar.}
\begin{itemize}\setlength\itemsep{0pt}
	\item \textbf{Interseccion con punto}: aniadir la division despues de verificar orientacion.
	\item \textbf{Interseccion rayo-segmento}: ajustar la primera condicion.
	\item \textbf{Poligonos convexos}: usar separacion de ejes (SAT).
	\item \textbf{Multiples queries}: precomputar bounding boxes o usar segment tree.
\end{itemize}

\paragraph{Trampas y errores frecuentes.}
\begin{itemize}\setlength\itemsep{0pt}
	\item Si los segmentos tocan solo en un extremo, el snippet lo considera interseccion (algunos problemas quieren estricto). Ajustar segun el problema.
	\item El caso colineal es comun; olvidar el check adicional de bounding box da falsos negativos.
	\item En problemas con muchos segmentos, considerar sweep line para $O(n \log n)$.
\end{itemize}

\paragraph{Codigo.}
Ver \texttt{cpp.tex}, seccion \textit{Interseccion de segmentos}.

\vspace{0.5em}

\subsubsection{Contains colineal}

\paragraph{Idea intuitiva.}
Verifica si un punto $c$ esta sobre el segmento $[a, b]$ (asumiendo que ya son colineales). Si \texttt{cross(a, b, c) = 0} y $c$ esta dentro del bounding box de $[a, b]$, entonces $c$ esta en el segmento. La razon: si los tres puntos son colineales, $c$ esta en el segmento sii sus coordenadas estan entre las de $a$ y $b$.

\paragraph{Propiedades clave (invariantes).}
\begin{itemize}\setlength\itemsep{0pt}
	\item Asume colinealidad (no la verifica mas alla de \texttt{cross == 0}).
	\item Usa comparaciones con \texttt{<=} para extremos \textbf{inclusivos}.
	\item Funciona con segmentos verticales (la condicion se reduce a una sola coordenada).
\end{itemize}

\paragraph{Codigo clave.}
Test de contencion en segmento (asume colinealidad):
\begin{verbatim}
bool onSegment(Point a, Point b, Point c) {
    return min(a.x, b.x) <= c.x && c.x <= max(a.x, b.x) &&
           min(a.y, b.y) <= c.y && c.y <= max(a.y, b.y);
}
\end{verbatim}

\paragraph{Complejidad.}
\begin{itemize}\setlength\itemsep{0pt}
	\item $O(1)$.
	\item Constante minima.
\end{itemize}

\paragraph{Donde tocar para modificar.}
\begin{itemize}\setlength\itemsep{0pt}
	\item \textbf{Contener estricto}: cambiar \texttt{<=} a \texttt{<} para excluir extremos.
	\item \textbf{Coordenadas doubles}: usar tolerancia (\texttt{eps}).
	\item \textbf{Solo x o solo y}: aniadir check especifico para segmentos verticales/horizontales.
\end{itemize}

\paragraph{Trampas y errores frecuentes.}
\begin{itemize}\setlength\itemsep{0pt}
	\item No usar este snippet para puntos \textbf{no colineales}; siempre verificar primero la colinealidad.
	\item Si los puntos tienen doubles, el chequeo \texttt{cross == 0} puede fallar; aniadir tolerancia.
	\item El bounding box incluye los extremos; verificar si el problema quiere exclusivo o inclusivo.
\end{itemize}

\paragraph{Codigo.}
Ver \texttt{cpp.tex}, seccion \textit{Contains colineal}.

\subsection{Poligonos}

\subsubsection{Convex Hull (monotono)}

\paragraph{Idea intuitiva.}
Algoritmo de \textbf{Andrew's monotone chain}: ordenar puntos por $(x, y)$, luego construir la hull inferior y superior por separado. Para cada punto, aniadirlo al hull y, mientras el ultimo giro no sea ``counter-clockwise'' (\texttt{cross > 0}), popear. La demostracion: si los puntos se procesan en orden, el hull inferior y superior cubren la frontera completa sin auto-intersecciones.

\paragraph{Propiedades clave (invariantes).}
\begin{itemize}\setlength\itemsep{0pt}
	\item Ordena los puntos en $O(n \log n)$.
	\item Construye la hull en $O(n)$ despues.
	\item Output: poligono convexo en counter-clockwise, sin punto final duplicado.
	\item El \texttt{<= 0} incluye colineales; usar \texttt{< 0} para excluirlos.
	\item Cada punto aparece en la hull inferior o superior (no ambos).
\end{itemize}

\paragraph{Codigo clave.}
Construccion del hull inferior y superior:
\begin{verbatim}
sort(p.begin(), p.end());
vector<Point> lower, upper;
for (auto& p : points) {
    while (lower.size() >= 2 && cross(lower.back() - lower[lower.size()-2],
                                       p - lower.back()) <= 0)
        lower.pop_back();
    lower.push_back(p);
}
for (int i = points.size() - 1; i >= 0; --i) {
    auto& p = points[i];
    while (upper.size() >= 2 && cross(upper.back() - upper[upper.size()-2],
                                       p - upper.back()) <= 0)
        upper.pop_back();
    upper.push_back(p);
}
vector<Point> hull = lower;
hull.insert(hull.end(), upper.begin() + 1, upper.end() - 1);
\end{verbatim}

\paragraph{Complejidad.}
\begin{itemize}\setlength\itemsep{0pt}
	\item $O(n \log n)$.
	\item Constante muy pequena (solo cross products).
\end{itemize}

\paragraph{Donde tocar para modificar.}
\begin{itemize}\setlength\itemsep{0pt}
	\item \textbf{Colineales}: cambiar \texttt{<= 0} por \texttt{< 0} para excluirlos del hull.
	\item \textbf{Puntos repetidos}: \texttt{unique} despues de sort.
	\item \textbf{Hull 3D}: gift wrapping es $O(n^2)$; convex hull 3D es mas complejo.
	\item \textbf{Output con los puntos}: aniadir el hull completo (no solo upper/lower).
\end{itemize}

\paragraph{Trampas y errores frecuentes.}
\begin{itemize}\setlength\itemsep{0pt}
	\item El \texttt{pop\_back} y el \texttt{reverse} son esenciales para no duplicar el ultimo punto.
	\item Para excluir colineales, ajustar el \texttt{<= 0} vs \texttt{< 0}.
	\item Si los puntos estan duplicados, aniadir \texttt{unique} antes del sort.
\end{itemize}

\paragraph{Codigo.}
Ver \texttt{cpp.tex}, seccion \textit{Convex Hull (monotono)}.

\vspace{0.5em}

\subsubsection{Polygon Area}

\paragraph{Idea intuitiva.}
El area (con signo) de un poligono se calcula con la \textbf{Shoelace formula}: $A = \frac{1}{2} \sum_{i=0}^{n-1} (x_i y_{i+1} - x_{i+1} y_i)$, donde el indice es modulo $n$. El snippet retorna el \textbf{doble} del area (en valor absoluto). La formula sale del ``sumar areas de trapezoides'': cada lado del poligono define un trapezoide con el eje $x$; la suma con signo da el area.

\paragraph{Propiedades clave (invariantes).}
\begin{itemize}\setlength\itemsep{0pt}
	\item El signo indica la orientacion: positivo si counter-clockwise.
	\item Para el doble: no dividir por 2 si no es necesario (evita fracciones).
	\item Funciona para poligonos simples (no auto-intersectantes).
	\item Si el poligono es no convexo, la formula da el area ``con signo''.
\end{itemize}

\paragraph{Codigo clave.}
La suma de trapezoides:
\begin{verbatim}
long long shoelace(const vector<Point>& p) {
    long long s = 0;
    int n = p.size();
    for (int i = 0; i < n; ++i)
        s += (long long)p[i].x * p[(i+1)%n].y - (long long)p[(i+1)%n].x * p[i].y;
    return abs(s);  // doble del area
}
\end{verbatim}

\paragraph{Complejidad.}
\begin{itemize}\setlength\itemsep{0pt}
	\item $O(n)$.
	\item Constante minima.
\end{itemize}

\paragraph{Donde tocar para modificar.}
\begin{itemize}\setlength\itemsep{0pt}
	\item \textbf{Area exacta}: dividir por 2 (con cuidado si da fraccion; usar \texttt{double}).
	\item \textbf{Area firmada}: quitar \texttt{abs} y dejar el signo.
	\item \textbf{Punto dentro de poligono convexo}: con la hull, chequear que esta a un mismo lado de todas las aristas.
	\item \textbf{Triangulacion}: descomponer en triangulos y sumar areas.
\end{itemize}

\paragraph{Trampas y errores frecuentes.}
\begin{itemize}\setlength\itemsep{0pt}
	\item Si los puntos no estan en orden (CCW o CW), el area firmada puede ser incorrecta. Para signed area consistente, usar siempre el mismo orden.
	\item Para poligonos con holes, sumar las areas con signo apropiado.
	\item Con doubles, aniadir tolerancia \texttt{eps}.
\end{itemize}

\paragraph{Codigo.}
Ver \texttt{cpp.tex}, seccion \textit{Polygon Area}.

\vspace{0.5em}

\subsubsection{Distance point-segment / point-line}

\paragraph{Idea intuitiva.}
Distancia minima de un punto $p$ a un segmento $[a, b]$ o a la recta que pasa por $a$ y $b$.
\begin{itemize}\setlength\itemsep{0pt}
	\item \textbf{Linea}: $d = \frac{|ab \times ap|}{|ab|}$.
	\item \textbf{Segmento}: si la proyeccion de $p$ sobre $ab$ cae fuera del segmento, la distancia es al extremo mas cercano.
\end{itemize}
La razon geometrica: la distancia a una linea es el area del paralelogramo dividido por la base; la proyeccion ortogonal indica si estamos ``dentro'' o ``fuera'' del segmento.

\paragraph{Propiedades clave (invariantes).}
\begin{itemize}\setlength\itemsep{0pt}
	\item \texttt{dist2PointLine} retorna distancia \textbf{al cuadrado}.
	\item \texttt{dist2PointSegment} distingue 3 casos: $p$ antes de $a$, despues de $b$, o proyectado sobre el segmento.
	\item Si \texttt{len2 = 0} (degenerate), la distancia es $|ap|$ o $|bp|$.
	\item La proyeccion se calcula con dot products y el parametro $t \in [0, 1]$.
\end{itemize}

\paragraph{Codigo clave.}
Distancia al cuadrado a segmento:
\begin{verbatim}
long long dist2Segment(Point p, Point a, Point b) {
    long long len2 = (a-b).len2();
    if (len2 == 0) return (p-a).len2();  // degenerate
    long long t = max(0LL, min(1LL, dot(p-a, b-a) / len2));
    Point proj = {a.x + t*(b.x - a.x), a.y + t*(b.y - a.y)};
    return (p - proj).len2();
}
\end{verbatim}

\paragraph{Complejidad.}
\begin{itemize}\setlength\itemsep{0pt}
	\item $O(1)$.
	\item Constante minima.
\end{itemize}

\paragraph{Donde tocar para modificar.}
\begin{itemize}\setlength\itemsep{0pt}
	\item \textbf{Distancia (no al cuadrado)}: aniadir \texttt{sqrt} al final.
	\item \textbf{Coordenadas doubles}: aniadir \texttt{long double} y tolerancia.
	\item \textbf{3D}: aniadir coordenada $z$ y adaptar.
	\item \textbf{Multiples queries}: usar KD-tree o segment tree para acelerar.
\end{itemize}

\paragraph{Trampas y errores frecuentes.}
\begin{itemize}\setlength\itemsep{0pt}
	\item Overflow en \texttt{area2 * area2} si el area es grande; \texttt{long long} alcanza para $|x|, |y| \le 10^9$ pero no mas.
	\item Si $a = b$ (segmento degenerado), aniadir check explicito.
	\item En doubles, aniadir tolerancia \texttt{eps} para comparar.
\end{itemize}

\paragraph{Codigo.}
Ver \texttt{cpp.tex}, seccion \textit{Distance point-segment / point-line}.

\vspace{0.5em}

\subsubsection{Projection of point onto line}

\paragraph{Idea intuitiva.}
Proyeccion ortogonal de $p$ sobre la recta que pasa por $a$ y $b$. El parametro $t$ tal que la proyeccion es $a + t \cdot (b - a)$ es:
\[ t = \frac{(p - a) \cdot (b - a)}{|b - a|^2}. \]
La idea: descomponemos $\vec{ap}$ en una componente paralela a $\vec{ab}$ y una perpendicular; $t$ mide la posicion en la primera.

\paragraph{Propiedades clave (invariantes).}
\begin{itemize}\setlength\itemsep{0pt}
	\item Si $t \in [0, 1]$, la proyeccion cae sobre el \textbf{segmento}; si no, cae sobre la extension.
	\item Retorna \texttt{double} (puede tener precision issues).
	\item $t < 0$: $p$ esta ``antes'' de $a$; $t > 1$: $p$ esta ``despues'' de $b$.
\end{itemize}

\paragraph{Codigo clave.}
Proyeccion ortogonal:
\begin{verbatim}
Point projection(Point p, Point a, Point b) {
    Point ab = b - a, ap = p - a;
    double t = (double)dot(ap, ab) / ab.len2();
    return {a.x + t * ab.x, a.y + t * ab.y};
}
\end{verbatim}

\paragraph{Complejidad.}
\begin{itemize}\setlength\itemsep{0pt}
	\item $O(1)$.
	\item Constante minima (dos dot products y una division).
\end{itemize}

\paragraph{Donde tocar para modificar.}
\begin{itemize}\setlength\itemsep{0pt}
	\item \textbf{Proyeccion sobre segmento}: clipar $t$ a $[0, 1]$ antes de calcular el punto.
	\item \textbf{Distancia al segmento via proyeccion}: combinar con el modulo anterior.
	\item \textbf{Solo el parametro $t$}: si solo necesitas saber si esta dentro del segmento, retorna $t$ sin construir el punto.
\end{itemize}

\paragraph{Trampas y errores frecuentes.}
\begin{itemize}\setlength\itemsep{0pt}
	\item Si $a = b$ ($|ab| = 0$), la proyeccion no esta definida; aniadir guard.
	\item En doubles, aniadir tolerancia \texttt{eps} para comparar $t$ con 0 o 1.
	\item La division por \texttt{len2} puede ser cara si las coordenadas son grandes; usar \texttt{double} o precomputar.
\end{itemize}

\paragraph{Codigo.}
Ver \texttt{cpp.tex}, seccion \textit{Projection of point onto line}.

\vspace{0.5em}

\subsubsection{Point in polygon (ray casting)}

\paragraph{Idea intuitiva.}
Algoritmo de \textbf{ray casting}: contar cuantas veces un rayo horizontal desde $p$ cruza los bordes del poligono. Si es impar, $p$ esta dentro; si par, fuera. El snippet implementa la version clasica recorriendo aristas. La razon: cada cruce cambia el lado del rayo respecto al poligono; un numero impar de cambios significa ``terminamos dentro''.

\paragraph{Propiedades clave (invariantes).}
\begin{itemize}\setlength\itemsep{0pt}
	\item Funciona con poligonos simples (no necesariamente convexos).
	\item Las intersecciones en vertices cuentan especialmente.
	\item Para poligonos con \textbf{huecos}, sigue funcionando.
	\item La respuesta es estrictamente 0 (fuera) o 1 (dentro) para poligonos simples.
\end{itemize}

\paragraph{Codigo clave.}
El test clasico de ray casting:
\begin{verbatim}
bool pointInPolygon(Point p, const vector<Point>& poly) {
    int n = poly.size();
    bool inside = false;
    for (int i = 0, j = n - 1; i < n; j = i++) {
        if ((poly[i].y > p.y) != (poly[j].y > p.y) &&
            p.x < (poly[j].x - poly[i].x) * (p.y - poly[i].y) /
                    (poly[j].y - poly[i].y) + poly[i].x)
            inside = !inside;
    }
    return inside;
}
\end{verbatim}

\paragraph{Complejidad.}
\begin{itemize}\setlength\itemsep{0pt}
	\item $O(n)$ por query.
	\item Memoria $O(1)$ (no se almacena el rayo).
\end{itemize}

\paragraph{Donde tocar para modificar.}
\begin{itemize}\setlength\itemsep{0pt}
	\item \textbf{Winding number}: version alternativa que cuenta el numero de vueltas.
	\item \textbf{Poligono convexo}: check lineal contra las aristas (mas rapido).
	\item \textbf{Multiples queries}: precomputar o usar segment tree sobre las aristas.
	\item \textbf{Punto en frontera}: aniadir check adicional si es colineal con alguna arista.
\end{itemize}

\paragraph{Trampas y errores frecuentes.}
\begin{itemize}\setlength\itemsep{0pt}
	\item El manejo del vertice (cuando el rayo pasa por un vertice) puede causar double-counting; el snippet lo evita con la condicion \texttt{(poly[i].y > p.y) != (poly[j].y > p.y)}.
	\item Si el poligono tiene holes, el algoritmo sigue funcionando.
	\item Para poligonos no simples (auto-intersectantes), el resultado puede ser ambiguo.
\end{itemize}

\paragraph{Codigo.}
Ver \texttt{cpp.tex}, seccion \textit{Point in polygon (ray casting)}.

\vspace{0.5em}

\subsubsection{Minkowski sum (convex polygons)}

\paragraph{Idea intuitiva.}
La \textbf{suma de Minkowski} $A \oplus B$ de dos poligonos convexos es otro poligono convexo que representa $\{a + b : a \in A, b \in B\}$. Algoritmo: ordenar las aristas por angulo polar (CCW) y ``mergearlas'' como en merge sort; el resultado es un poligono CCW sin repetir el ultimo punto. La idea: el resultado tiene una arista paralela a cada arista de los originales; la suma se obtiene barriendo los angulos.

\paragraph{Propiedades clave (invariantes).}
\begin{itemize}\setlength\itemsep{0pt}
	\item Si $A, B$ son convexos con $n, m$ vertices, $A \oplus B$ tiene $\le n + m$ vertices.
	\item Util para \textbf{detectar colision}: $A$ colisiona con $B$ si $0 \in A \oplus (-B)$.
	\item Asume poligonos en orden \textbf{counter-clockwise}.
	\item El resultado es convexo (suma de convexos es convexo).
\end{itemize}

\paragraph{Codigo clave.}
La suma de Minkowski:
\begin{verbatim}
vector<Point> minkowski(const vector<Point>& A, const vector<Point>& B) {
    vector<Point> result;
    int i = 0, j = 0;
    int n = A.size(), m = B.size();
    do {
        result.push_back(A[i] + B[j]);
        long long cross = (A[(i+1)%n] - A[i]) ^ (B[(j+1)%m] - B[j]);
        if (cross >= 0) i = (i + 1) % n;
        if (cross <= 0) j = (j + 1) % m;
    } while (i || j);
    return result;
}
\end{verbatim}

\paragraph{Complejidad.}
\begin{itemize}\setlength\itemsep{0pt}
	\item $O(n + m)$.
	\item Memoria $O(n + m)$.
\end{itemize}

\paragraph{Donde tocar para modificar.}
\begin{itemize}\setlength\itemsep{0pt}
	\item \textbf{Deteccion de colision}: aniadir $(-B)$ (reflejar) y verificar si $0$ esta en el resultado.
	\item \textbf{Poligonos no convexos}: la suma puede no ser poligono simple; algoritmo distinto.
	\item \textbf{Multiples queries}: precomputar las aristas una vez.
	\item \textbf{Robot motion planning}: usar para planear trayectorias.
\end{itemize}

\paragraph{Trampas y errores frecuentes.}
\begin{itemize}\setlength\itemsep{0pt}
	\item Si los poligonos no son CCW, la suma es incorrecta. Normalizar antes.
	\item El bucle termina cuando ambos indices vuelven a 0; aniadir check.
	\item Para la deteccion de colision, aniadir el test \texttt{0 in A \^{} (-B)}.
\end{itemize}

\paragraph{Codigo.}
Ver \texttt{cpp.tex}, seccion \textit{Minkowski sum (convex polygons)}.

\subsection{Herramientas}

\subsubsection{Li Chao Tree (long long)}

\paragraph{Idea intuitiva.}
Mantiene un \textbf{lower envelope} (o upper) de un conjunto de rectas $y = mx + b$ sobre un dominio $[X_{LO}, X_{HI}]$ y responde ``minimo $y$ en un $x$ dado'' en $O(\log X)$. La estructura es un segment tree donde cada nodo guarda la recta ``mejor'' para su segmento medio; al aniadir una recta, se compara con la actual y se decide cual es mejor en cada mitad, descendiendo recursivamente. La demostracion: en cada nivel, solo se guarda la mejor recta para el ``punto medio''; las otras se transmiten a los hijos.

\paragraph{Propiedades clave (invariantes).}
\begin{itemize}\setlength\itemsep{0pt}
	\item \textbf{Online}: cada \texttt{addLine} en $O(\log X)$.
	\item Solo insercion; no hay borrar (para borrar, usar Li Chao segmentado con reset).
	\item La recta guardada en un nodo es la ``mejor'' en su punto medio.
	\item Las rectas ``perdedoras'' se transmiten a los hijos para futuros inserts.
\end{itemize}

\paragraph{Codigo clave.}
La insercion en Li Chao:
\begin{verbatim}
void addLine(Line nw, int v = 1, int l = X_LO, int r = X_HI) {
    int m = (l + r) / 2;
    bool left = nw.get(l) < tree[v].get(l);
    bool mid = nw.get(m) < tree[v].get(m);
    if (mid) swap(tree[v], nw);
    if (r - l == 1) return;
    if (left != mid) addLine(nw, v*2, l, m);
    else             addLine(nw, v*2+1, m, r);
}
\end{verbatim}

\paragraph{Complejidad.}
\begin{itemize}\setlength\itemsep{0pt}
	\item $O(\log X)$ por insercion y query, donde $X$ es el tamano del dominio.
	\item Memoria $O(\log X)$ por insercion, $O(X)$ total.
\end{itemize}

\paragraph{Donde tocar para modificar.}
\begin{itemize}\setlength\itemsep{0pt}
	\item \textbf{Upper envelope en vez de lower}: invertir el signo (\texttt{>} en vez de \texttt{<}).
	\item \textbf{Queries de segmento} (no full line): si la recta solo esta activa en un rango, aniadir bandera \texttt{active}.
	\item \textbf{Eliminar rectas}: usar Li Chao segmentado (mas complejo).
	\item \textbf{Pendientes grandes}: si las pendientes tienen magnitud $> 10^9$, aniadir \texttt{\_\_int128} para evitar overflow.
\end{itemize}

\paragraph{Trampas y errores frecuentes.}
\begin{itemize}\setlength\itemsep{0pt}
	\item El parametro \texttt{lo == hi} debe retornar el valor del nodo (caso base).
	\item La comparacion \texttt{newLine.get(mid) < node->line.get(mid)} decide quien es mejor en el medio.
	\item Si las rectas tienen $m$ iguales pero $b$ distintos, se queda con la primera (tie-break).
\end{itemize}

\paragraph{Codigo.}
Ver \texttt{cpp.tex}, seccion \textit{Li Chao Tree (long long)}.

\vspace{0.5em}

\subsubsection{Sweep line template}

\paragraph{Idea intuitiva.}
Patron de diseno para problemas geometricos donde se barre una linea sobre el plano y se mantiene una estructura de datos sobre el otro eje. Eventos (add, query, remove) se ordenan por la coordenada de barrido y se procesan secuencialmente.

\paragraph{Propiedades clave (invariantes).}
\begin{itemize}\setlength\itemsep{0pt}
	\item Eventos: tuplas \texttt{(x, type, id)}.
	\item La estructura auxiliar (set, multiset, segtree) se mantiene sobre el otro eje.
	\item Procesar eventos en orden de $x$ ascendente.
\end{itemize}

\paragraph{Complejidad.}
\begin{itemize}\setlength\itemsep{0pt}
	\item Depende del problema; tipico $O((N + E) \log N)$.
\end{itemize}

\paragraph{Donde tocar para modificar.}
\begin{itemize}\setlength\itemsep{0pt}
	\item \textbf{Aniadir/quitar segmentos}: aniadir eventos add/remove al barrido.
	\item \textbf{Aniadir queries}: aniadir eventos query.
	\item \textbf{Cambiar el eje de barrido}: si barres en $y$ en vez de $x$, aniadir comparador.
\end{itemize}

\paragraph{Trampas y errores frecuentes.}
\begin{itemize}\setlength\itemsep{0pt}
	\item El snippet es solo el \textbf{esqueleto}; los handlers \texttt{add/query/remove} se llenan por problema.
	\item Cuidado con eventos duplicados.
\end{itemize}

\paragraph{Codigo.}
Ver \texttt{cpp.tex}, seccion \textit{Sweep line template}.

\vspace{0.5em}

\subsubsection{Closest pair of points}

\paragraph{Idea intuitiva.}
Encuentra el par de puntos mas cercano en $O(n \log n)$ con \textbf{divide y venceras}. Algoritmo:
\begin{itemize}\setlength\itemsep{0pt}
	\item Sort por $x$.
	\item Dividir en mitades; recursivo.
	\item Combinar: dado el minimo $d$ de las dos mitades, strip = puntos en $[midX - d, midX + d]$; ordenar por $y$; cada punto solo se compara con los siguientes hasta $y$ difiera en $< d$ (maximo 7).
\end{itemize}

\paragraph{Propiedades clave (invariantes).}
\begin{itemize}\setlength\itemsep{0pt}
	\item Requiere ordenar por $y$ despues (en el snippet se hace durante el merge).
	\item Strip tiene a lo sumo $O(n)$ puntos en total, pero la comparacion es $O(1)$ amortizada por punto.
	\item Resultado en \texttt{double}.
\end{itemize}

\paragraph{Complejidad.}
\begin{itemize}\setlength\itemsep{0pt}
	\item $O(n \log n)$.
\end{itemize}

\paragraph{Donde tocar para modificar.}
\begin{itemize}\setlength\itemsep{0pt}
	\item \textbf{Reconstruir el par}: aniadir indices en la recursion.
	\item \textbf{Distancia maxima}: cambiar \texttt{min} por \texttt{max} (mas facil; no requiere D\&C).
	\item \textbf{Puntos duplicados}: aniadir check especial (distancia 0).
\end{itemize}

\paragraph{Trampas y errores frecuentes.}
\begin{itemize}\setlength\itemsep{0pt}
	\item El merge requiere ordenar por $y$; el snippet lo hace in-place con un \texttt{tmp}.
	\item Si se hace en cada recursion, es $O(n \log n)$ extra.
\end{itemize}

\paragraph{Codigo.}
Ver \texttt{cpp.tex}, seccion \textit{Closest pair of points}.

\vspace{0.5em}

\subsubsection{Rotating calipers (polygon diameter)}

\paragraph{Idea intuitiva.}
Encuentra el \textbf{diametro} de un poligono convexo (par de vertices mas distantes) en $O(n)$. Asume poligono en orden counter-clockwise. Algoritmo: dos indices $i$ (sobre las aristas) y $k$ (sobre los vertices opuestos); en cada paso, rotar $k$ mientras el area del paralelogramo aumenta. La distancia $i$-$k$ puede ser maxima en cualquier momento del barrido.

\paragraph{Propiedades clave (invariantes).}
\begin{itemize}\setlength\itemsep{0pt}
	\item Cada indice avanza como maximo $n$ veces total.
	\item El ``diametro'' es la maxima distancia entre dos vertices.
	\item Retorna \textbf{distancia al cuadrado}.
\end{itemize}

\paragraph{Complejidad.}
\begin{itemize}\setlength\itemsep{0pt}
	\item $O(n)$.
\end{itemize}

\paragraph{Donde tocar para modificar.}
\begin{itemize}\setlength\itemsep{0pt}
	\item \textbf{Reconstruir el par de vertices}: aniadir indices.
	\item \textbf{Anchura minima / maxima}: variantes de calipers.
	\item \textbf{Distancia minima entre poligonos convexos}: otra variante.
\end{itemize}

\paragraph{Trampas y errores frecuentes.}
\begin{itemize}\setlength\itemsep{0pt}
	\item Asume poligono \textbf{estrictamente} convexo (sin vertices colineales). Si hay colineales, aniadir perturbacion o filtrar.
\end{itemize}

\paragraph{Codigo.}
Ver \texttt{cpp.tex}, seccion \textit{Rotating calipers (polygon diameter)}.

% ============================================================
\section{Miscelanea}
% ============================================================

\subsection{Utilidades del usuario}

\subsubsection{Hora}

\paragraph{Idea intuitiva.}
Una pequena clase para representar \textbf{tiempos en segundos} y hacer aritmetica basica sobre ellos. Internamente almacena todo en segundos (un solo \texttt{int}), lo cual simplifica comparaciones y diferencias. El constructor acepta $(h, m, s)$ y los convierte a segundos. Hay un metodo \texttt{cut()} que fuerza un minimo (util para ``horas trabajadas'').

\paragraph{Propiedades clave (invariantes).}
\begin{itemize}\setlength\itemsep{0pt}
	\item Representacion canonica: \textbf{solo segundos}. No hay zonas horarias, no hay DST.
	\item Comparadores \texttt{<} y \texttt{>} operan sobre los segundos.
	\item \texttt{operator-} retorna la \textbf{diferencia absoluta} en segundos (siempre no negativa).
	\item \texttt{cut()} aplica un piso (por defecto $8\text{h }30\text{m} = 30600$ segundos).
\end{itemize}

\paragraph{Complejidad.}
\begin{itemize}\setlength\itemsep{0pt}
	\item $O(1)$ por operacion.
\end{itemize}

\paragraph{Donde tocar para modificar.}
\begin{itemize}\setlength\itemsep{0pt}
	\item \textbf{Aniadir suma}: \texttt{operator+} que devuelve un \texttt{Hour} con la suma de segundos.
	\item \textbf{Aniadir output formateado}: aniadir \texttt{show\_hms()} que imprime $H$:$M$:$S$.
	\item \textbf{Cambiar el ``corte''}: modificar el valor magico \texttt{30'600} en \texttt{cut()} o pasarlo como parametro.
	\item \textbf{Tamano maximo}: si los tiempos exceden $\sim 68$ anos, el \texttt{int} desborda; usar \texttt{long long}.
	\item \textbf{Precision sub-segundo}: aniadir \texttt{ms} (milisegundos) y guardar en \texttt{long long} con factor 1000.
\end{itemize}

\paragraph{Trampas y errores frecuentes.}
\begin{itemize}\setlength\itemsep{0pt}
	\item El operador \texttt{-} retorna \textbf{valor absoluto}, no diferencia con signo. Si necesitas diferencia con signo, aniadir un nuevo metodo o sobrecargar.
	\item \texttt{cut()} modifica \texttt{s} \textbf{in-place}; si el llamador no espera esto, aniadir una version \texttt{cut(min)} que retorna un nuevo \texttt{Hour}.
	\item El umbral magico \texttt{30'600} (8h 30m) esta hardcoded; documentar bien que es ``jornada minima''.
\end{itemize}

\paragraph{Codigo.}
Ver \texttt{cpp.tex}, seccion \textit{Hora}.

\end{document}
', aniadir un puntero a ambas raices.
\end{itemize}

\paragraph{Codigo.}
Ver \texttt{cpp.tex}, seccion \textit{Segment Tree persistente}.

\vspace{0.5em}

\subsubsection{Segment Tree iterativo}

\paragraph{Idea intuitiva.}
Representacion \textbf{sin recursion}: los nodos se guardan en un \texttt{vector} de tamano $2n$. Los hijos del nodo $i$ son $2i$ y $2i+1$. Las hojas ocupan las posiciones $n, n+1, \dots, 2n-1$. Update: subir el cambio desde la hoja hasta la raiz. Query: subir dos indices $l$ y $r$ hasta que se ``crucen''. La ventaja principal: localidad de memoria (los nodos estan en posiciones contiguas del vector), que mejora el cache miss rate y la constante practica. La desventaja: implementar lazy propagation es mas complejo.

\paragraph{Propiedades clave (invariantes).}
\begin{itemize}\setlength\itemsep{0pt}
	\item $O(\log n)$ por operacion, \textbf{mas rapido} en practica que el recursivo (mejor uso de cache, no hay overhead de llamadas).
	\item Build simple: copiar las hojas y combinar de abajo hacia arriba.
	\item Las hojas tienen indices en $[n, 2n)$; los nodos internos en $[1, n)$.
\end{itemize}

\paragraph{Codigo clave.}
El bucle de query ascendente:
\begin{verbatim}
long long query(int l, int r) {  // [l, r] inclusivo
    long long res = 0;
    for (l += n, r += n + 1; l < r; l >>= 1, r >>= 1) {
        if (l & 1) res += st[l++];
        if (r & 1) res += st[--r];
    }
    return res;
}
\end{verbatim}

\paragraph{Complejidad.}
\begin{itemize}\setlength\itemsep{0pt}
	\item $O(\log n)$ por update/query, $O(n)$ build.
	\item En la practica, $\sim 1.5x$ mas rapido que el recursivo para $n \ge 10^5$.
\end{itemize}

\paragraph{Donde tocar para modificar.}
\begin{itemize}\setlength\itemsep{0pt}
	\item \textbf{Suma -> max/min/xor}: cambiar la operacion en \texttt{set} y \texttt{query}.
	\item \textbf{Aniadir update en rango con lazy}: implementar version lazy por separado (no trivial en iterativo).
	\item \textbf{Tamano no potencia de 2}: si $n$ no es potencia de 2, aniadir padding.
	\item \textbf{Persistencia}: usar \texttt{vector<array<int, 2>>} y duplicar nodos.
\end{itemize}

\paragraph{Trampas y errores frecuentes.}
\begin{itemize}\setlength\itemsep{0pt}
	\item En el query, \texttt{r += n+1} (no \texttt{r += n}) porque \texttt{query(l, r)} es \textbf{inclusivo} en ambos extremos.
	\item El indice \texttt{l & 1} detecta si \texttt{l} es hijo derecho; aniadir \texttt{res += st[l++]} en ese caso.
	\item El \texttt{--r} en la condicion derecha es para incluir el limite.
	\item Si las operaciones no son idempotentes, cuidado con el orden.
\end{itemize}

\paragraph{Codigo.}
Ver \texttt{cpp.tex}, seccion \textit{Segment Tree iterativo}.

\vspace{0.5em}

\subsection{Fenwick / BIT}

\subsubsection{BIT point update + prefix sum}

\paragraph{Idea intuitiva.}
El \textbf{Fenwick Tree} (BIT) almacena sums parciales de manera que \texttt{sum(i)} (suma de $[1, i]$) y \texttt{add(i, delta)} cuestan $O(\log n)$. La representacion es ingeniosa: el nodo $i$ cubre el rango $[i - \text{lowbit}(i) + 1, i]$. Asi, \texttt{sum(i)} itera \texttt{i -= lowbit(i)}, sumando los nodos que cubren un prefijo creciente, y \texttt{add(i, delta)} itera \texttt{i += lowbit(i)}, propagando el cambio a todos los nodos cuyo rango contiene a $i$. La idea: cada indice $i$ se descompone en una suma de bloques disjuntos de tamano potencias de 2 (su ``lowbit representation'').

\paragraph{Propiedades clave (invariantes).}
\begin{itemize}\setlength\itemsep{0pt}
	\item Indice \textbf{1-based} (no usar $i=0$).
	\item \texttt{lowbit(i) = i \& (-i)}.
	\item Cada nodo cubre exactamente un rango de tamano \texttt{lowbit(i)}.
	\item La operacion debe ser \textbf{invertible} (suma, xor, no min/max).
\end{itemize}

\paragraph{Codigo clave.}
Las dos operaciones basicas:
\begin{verbatim}
void add(int i, long long d) {
    for (; i <= n; i += i & -i) bit[i] += d;
}
long long sum(int i) {
    long long s = 0;
    for (; i > 0; i -= i & -i) s += bit[i];
    return s;
}
\end{verbatim}

\paragraph{Complejidad.}
\begin{itemize}\setlength\itemsep{0pt}
	\item $O(\log n)$ por operacion. Memoria $O(n)$.
	\item Constante muy pequena ($\sim 1$): en la practica, $\sim 3x$ mas rapido que segment tree.
\end{itemize}

\paragraph{Donde tocar para modificar.}
\begin{itemize}\setlength\itemsep{0pt}
	\item \textbf{Suma -> max/min/xor}: cambiar el \texttt{+=} por la operacion (con identidad apropriada). \textbf{No} funciona para \texttt{min} de forma natural (no es invertible).
	\item \textbf{Tamano dinamico}: aniadir un metodo \texttt{resize(newSize)} que preserve valores.
	\item \textbf{Tamano grande}: con $n = 10^7$, mejor comprimir coordenadas.
	\item \textbf{Operador no conmutativo}: usar segment tree en su lugar.
\end{itemize}

\paragraph{Trampas y errores frecuentes.}
\begin{itemize}\setlength\itemsep{0pt}
	\item Indices 0 vs 1: el BIT es \textbf{1-based}. Si el problema da indices 0-based, aniadir \texttt{+1} en cada llamada.
	\item El metodo \texttt{sum} itera \texttt{i -= i \& -i}; aniadir o quitar 1 rompe el invariante.
	\item Si el tipo es \texttt{int}, posibles overflows; usar \texttt{long long}.
	\item Para rango $[l, r]$, el query es \texttt{sum(r) - sum(l-1)} (cuidado con $l = 1$).
\end{itemize}

\paragraph{Codigo.}
Ver \texttt{cpp.tex}, seccion \textit{BIT point update + prefix sum}.

\vspace{0.5em}

\subsubsection{BIT range add + point query}

\paragraph{Idea intuitiva.}
Truco de \textbf{diferencias}: para sumar $d$ a $[l, r]$, aniadir $d$ en $l$ y $-d$ en $r+1$. Asi, el valor en $i$ es la suma de los prefijos hasta $i$, lo que se calcula con un BIT normal. La idea algebraica: si mantenemos un arreglo de diferencias $\Delta$ donde $a_i = \sum_{j \le i} \Delta_j$, entonces aniadir $d$ a $a[l..r]$ equivale a $\Delta_l \mathrel{+}= d$ y $\Delta_{r+1} \mathrel{-}= d$, y reconstruir $a$ es un prefix sum.

\paragraph{Propiedades clave (invariantes).}
\begin{itemize}\setlength\itemsep{0pt}
	\item Internamente es un BIT 5.1 con la convencion de diferencias.
	\item \texttt{rangeAdd(l, r, delta)} cuesta $O(\log n)$ (dos \texttt{add}).
	\item \texttt{pointQuery(i)} cuesta $O(\log n)$.
	\item Si $r+1 > n$, no se aniade el $-d$ (fuera del rango).
\end{itemize}

\paragraph{Codigo clave.}
El truco de diferencias:
\begin{verbatim}
void rangeAdd(int l, int r, long long d) {
    add(l, d);
    if (r + 1 <= n) add(r + 1, -d);
}
long long pointQuery(int i) { return sum(i); }
\end{verbatim}

\paragraph{Complejidad.}
\begin{itemize}\setlength\itemsep{0pt}
	\item $O(\log n)$ por operacion.
	\item Memoria $O(n)$.
\end{itemize}

\paragraph{Donde tocar para modificar.}
\begin{itemize}\setlength\itemsep{0pt}
	\item \textbf{Persistencia}: aniadir versionado (cada update crea una nueva ``copia'' del estado; en BIT no es trivial, mejor usar persistent segtree).
	\item \textbf{Coordenada compression}: si $n$ es grande pero los updates son pocos, comprimir.
	\item \textbf{2D range add + point query}: aniadir una dimension con la misma tecnica.
\end{itemize}

\paragraph{Trampas y errores frecuentes.}
\begin{itemize}\setlength\itemsep{0pt}
	\item Si $r+1 > n$, \textbf{no} aniadir \texttt{-delta} (no hay efecto fuera del rango); el snippet ya lo chequea.
	\item Para reconstruir el array original tras $k$ range adds, hacer $k$ point queries (no hay shortcut).
	\item Confundir ``suma de diferencias'' con ``suma acumulada''; son cosas distintas.
\end{itemize}

\paragraph{Codigo.}
Ver \texttt{cpp.tex}, seccion \textit{BIT range add + point query}.

\vspace{0.5em}

\subsubsection{BIT range add + range sum}

\paragraph{Idea intuitiva.}
Para queries $\sum_{i=1}^{x} a_i$ despues de varios range adds, el truco de diferencias no basta (necesitas la suma acumulada de las diferencias). Solucion: dos BITs, $t_1$ y $t_2$. Tras aniadir $d$ a $[l, r]$: $t_1$ recibe $+d$ en $l$, $-d$ en $r+1$; $t_2$ recibe $+d(l-1)$ en $l$, $-d \cdot r$ en $r+1$. Entonces $\text{prefixSum}(x) = t_1.\text{sum}(x) \cdot x - t_2.\text{sum}(x)$. La razon matematica: el valor en $i$ es $\sum_{j \le i} \Delta_j$, asi que la suma hasta $i$ es $\sum_{k \le i} (i - k + 1) \Delta_k$; expandir da la formula con dos BITs.

\paragraph{Propiedades clave (invariantes).}
\begin{itemize}\setlength\itemsep{0pt}
	\item Verificacion: para $x \ge r$, $\text{prefixSum}(x)$ acumula correctamente todas las contribuciones.
	\item La formula sale de expandir la suma de diferencias.
	\item La identidad clave: $\sum_{i=1}^{x} i = x(x+1)/2$, asi que la suma de diferencias ponderadas requiere BITs adicionales.
\end{itemize}

\paragraph{Codigo clave.}
Range add con range sum (dos BITs):
\begin{verbatim}
void rangeAdd(int l, int r, long long d) {
    t1.add(l, d); t1.add(r + 1, -d);
    t2.add(l, d * (l - 1)); t2.add(r + 1, -d * r);
}
long long prefixSum(int x) {
    return t1.sum(x) * x - t2.sum(x);
}
\end{verbatim}

\paragraph{Complejidad.}
\begin{itemize}\setlength\itemsep{0pt}
	\item $O(\log n)$ por operacion (2 \texttt{add} o 2 \texttt{sum} por cada lado).
	\item Memoria $O(n)$.
\end{itemize}

\paragraph{Donde tocar para modificar.}
\begin{itemize}\setlength\itemsep{0pt}
	\item \textbf{2D range add + range sum}: aniadir otra dimension con la misma tecnica (queda 4 BITs).
	\item \textbf{Cambiar mod}: aniadir modulo en cada operacion.
	\item \textbf{Output del array}: aniadir un metodo que reconstruya $a_i$ point query por cada $i$.
\end{itemize}

\paragraph{Trampas y errores frecuentes.}
\begin{itemize}\setlength\itemsep{0pt}
	\item Overflow en \texttt{delta * (l - 1)} si los valores son grandes; usar \texttt{long long} y \texttt{\% MOD} al final.
	\item Confundir la convencion $0$-based vs $1$-based del BIT; las dos afectan las formulas.
	\item En 2D, el codigo es $\sim 4x$ mas largo y propenso a errores; verificar con casos pequenos.
\end{itemize}

\paragraph{Codigo.}
Ver \texttt{cpp.tex}, seccion \textit{BIT range add + range sum}.

\subsection{RMQ y sparse table}

\subsubsection{Sparse Table}

\paragraph{Idea intuitiva.}
Precomputa respuestas a queries de tamano $2^k$ para $k = 0, 1, \dots, \log n$. Asi, \texttt{st[k][i]} guarda la respuesta sobre $[i, i + 2^k)$. Para query $[l, r]$, se usan dos bloques de tamano $2^{\lfloor \log_2(r-l+1) \rfloor}$: $[l, l + 2^k)$ y $[r - 2^k + 1, r)$, y se combinan. Esto da $O(1)$ por query \textbf{solo si la operacion es idempotente} (min, max, gcd; \textbf{no} suma). La razon: con idempotencia, solapar los dos bloques no afecta el resultado.

\paragraph{Propiedades clave (invariantes).}
\begin{itemize}\setlength\itemsep{0pt}
	\item \texttt{st[0][i] = a[i]}.
	\item \texttt{st[k][i] = op(st[k-1][i], st[k-1][i + 2\^{}(k-1)])}.
	\item \texttt{LOG = floor(log2(n)) + 1}.
	\item La operacion debe ser \textbf{idempotente} ($f(x, x) = x$) y asociativa.
\end{itemize}

\paragraph{Codigo clave.}
Query $O(1)$ en sparse table:
\begin{verbatim}
long long query(int l, int r) {
    int k = __builtin_clz(1) - __builtin_clz(r - l + 1);  // log2
    return min(st[k][l], st[k][r - (1 << k) + 1]);
}
\end{verbatim}

\paragraph{Complejidad.}
\begin{itemize}\setlength\itemsep{0pt}
	\item Build $O(n \log n)$, query $O(1)$.
	\item Memoria $O(n \log n)$.
\end{itemize}

\paragraph{Donde tocar para modificar.}
\begin{itemize}\setlength\itemsep{0pt}
	\item \textbf{Cambiar min por otra operacion idempotente}: cambiar las dos ocurrencias de \texttt{min(...)} por la operacion (max, gcd). Verificar identidad y asociatividad.
	\item \textbf{Suma en rango}: NO sirve con este metodo; usar prefix sums.
	\item \textbf{Sparse table 2D}: aniadir otra dimension (queda $O(n^2 \log^2 n)$ memoria).
	\item \textbf{LCA queries}: combinar con binary lifting en arboles.
\end{itemize}

\paragraph{Trampas y errores frecuentes.}
\begin{itemize}\setlength\itemsep{0pt}
	\item El query toma $[l, r]$ \textbf{inclusivo} (no $[l, r)$).
	\item El \texttt{LOG} debe estar bien calculado; con $n = 1$, \texttt{LOG = 1}.
	\item Usar sparse table para suma es un error; necesitas prefix sums o segment tree.
	\item Para queries $O(1)$ con operaciones no idempotentes, usar Disjoint Sparse Table.
\end{itemize}

\paragraph{Codigo.}
Ver \texttt{cpp.tex}, seccion \textit{Sparse Table}.

\vspace{0.5em}

\subsubsection{Disjoint Sparse Table}

\paragraph{Idea intuitiva.}
Variante que tambien da $O(1)$ por query pero \textbf{no requiere idempotencia} para algunas operaciones (sirve para suma, xor, etc. en linea con la mayoria de operaciones asociativas). La idea: para cada nivel $k$, los rangos ``se parten en mitades''. \texttt{st[k][i]} cubre un rango que \textbf{empieza} en $i$ y se extiende a izquierda o derecha hasta encontrar un ``punto de division''. El truco: para query $[l, r]$, el nivel optimo es el bit mas alto donde $l$ y $r$ difieren; los rangos $[l, ?]$ y $[?, r]$ cubren exactamente el query sin solape.

\paragraph{Propiedades clave (invariantes).}
\begin{itemize}\setlength\itemsep{0pt}
	\item Construccion diferente al Sparse Table clasico: en cada nivel, se extiende desde cada posicion hacia ambos lados hasta encontrar la mitad de nivel.
	\item \texttt{query(l, r)} usa \texttt{k = log2(l \^{} r)} y combina \texttt{st[k][l]} y \texttt{st[k][r]}.
	\item La operacion debe ser \textbf{asociativa} pero no necesariamente idempotente.
\end{itemize}

\paragraph{Codigo clave.}
Query con DST:
\begin{verbatim}
long long query(int l, int r) {
    if (l == r) return st[0][l];
    int k = 31 - __builtin_clz(l ^ r);  // bit mas alto que difiere
    return op(st[k][l], st[k][r]);
}
\end{verbatim}

\paragraph{Complejidad.}
\begin{itemize}\setlength\itemsep{0pt}
	\item Build $O(n \log n)$, query $O(1)$.
	\item Memoria $O(n \log n)$.
\end{itemize}

\paragraph{Donde tocar para modificar.}
\begin{itemize}\setlength\itemsep{0pt}
	\item \textbf{Operacion no asociativa}: no funciona; el Disjoint Sparse Table asume asociatividad.
	\item \textbf{Memory compression}: con $n \le 10^6$ alcanza; con mas, comprimir.
	\item \textbf{Operaciones no conmutativas}: aniadir el cuidado de que el orden es importante.
\end{itemize}

\paragraph{Trampas y errores frecuentes.}
\begin{itemize}\setlength\itemsep{0pt}
	\item La construccion es sutil; verificar con casos pequenos antes de usar.
	\item El bit ``mas alto que difiere'' requiere cuidado con $l = r$; aniadir caso base.
	\item Si la operacion no es asociativa, los dos bloques no se pueden combinar en cualquier orden.
\end{itemize}

\paragraph{Codigo.}
Ver \texttt{cpp.tex}, seccion \textit{Disjoint Sparse Table}.

\subsection{Trees y DP en arbol}

\subsubsection{Binary Lifting / LCA}

\paragraph{Idea intuitiva.}
Precomputa para cada nodo $v$ sus ancestros a distancias $2^k$ (los ``saltos''). Asi, para subir $d$ niveles, se descompone $d$ en binario y se aplican los saltos correspondientes. El \textbf{LCA} (Lowest Common Ancestor) se obtiene subiendo el nodo mas profundo hasta igualar profundidades y luego subiendo ambos hasta que sus padres coincidan. La intuicion: cualquier entero $\le n$ se representa en binario con $\le \log n$ bits, asi que subir $d$ niveles requiere $\le \log n$ saltos.

\paragraph{Propiedades clave (invariantes).}
\begin{itemize}\setlength\itemsep{0pt}
	\item \texttt{up[k][v]} = ancestro de $v$ a $2^k$ niveles.
	\item \texttt{LOG = floor(log2(n)) + 1}.
	\item La recursion de \texttt{dfs} llena \texttt{up[0][v]} (el padre directo) y por transitividad los demas.
	\item El LCA esta definido como el ancestro comun mas profundo de $u$ y $v$.
\end{itemize}

\paragraph{Codigo clave.}
LCA con binary lifting:
\begin{verbatim}
int lca(int u, int v) {
    if (depth[u] < depth[v]) swap(u, v);
    int diff = depth[u] - depth[v];
    for (int k = LOG - 1; k >= 0; --k)
        if (diff & (1 << k)) u = up[k][u];
    if (u == v) return u;
    for (int k = LOG - 1; k >= 0; --k)
        if (up[k][u] != up[k][v])
            u = up[k][u], v = up[k][v];
    return up[0][u];
}
\end{verbatim}

\paragraph{Complejidad.}
\begin{itemize}\setlength\itemsep{0pt}
	\item Build $O(n \log n)$, query LCA $O(\log n)$.
	\item Memoria $O(n \log n)$.
\end{itemize}

\paragraph{Donde tocar para modificar.}
\begin{itemize}\setlength\itemsep{0pt}
	\item \textbf{Distancia en el arbol}: combinar con \texttt{depth[]} para calcular \texttt{depth[u] + depth[v] - 2*depth[lca]}.
	\item \textbf{K-esimo ancestro}: subir $k$ niveles con el mismo truco.
	\item \textbf{Aniadir peso a las aristas}: guardar \texttt{dist[2\^{}k][v]} (distancia acumulada a $2^k$ saltos) ademas de \texttt{up}.
	\item \textbf{Sparse table sobre depth}: alternativa $O(1)$ a LCA (mas complejo).
\end{itemize}

\paragraph{Trampas y errores frecuentes.}
\begin{itemize}\setlength\itemsep{0pt}
	\item El arbol debe ser \textbf{rooted} (el snippet asume eso); si no, hacer un \texttt{dfs} desde cualquier nodo.
	\item Para grafos no conexos, correr \texttt{dfs} desde cada componente.
	\item En la segunda fase del LCA, subir ambos nodos mientras \texttt{up[k][u] != up[k][v]}; el LCA final es \texttt{up[0][u]}.
	\item El padre de la raiz es 0 (o -1); aniadir check si se sale del rango.
\end{itemize}

\paragraph{Codigo.}
Ver \texttt{cpp.tex}, seccion \textit{Binary Lifting / LCA}.

\vspace{0.5em}

\subsubsection{Heavy-Light Decomposition (HLD)}

\paragraph{Idea intuitiva.}
Descomponer un arbol en \textbf{cadenas pesadas} (paths donde cada nodo (salvo la raiz) tiene al hijo ``pesado'' como hijo de tamano maximo de su subarbol). Cada cadena se mapea a un rango contiguo en un arreglo base; asi, queries sobre \textbf{paths} se descomponen en $O(\log n)$ queries sobre rangos contiguos (cada cadena es un rango), que se resuelven con un BIT o segment tree. La intuicion: por cada nivel de recursion, el subarbol del hijo pesado es al menos la mitad, asi que un path cruza $\le 2 \log n$ cadenas.

\paragraph{Propiedades clave (invariantes).}
\begin{itemize}\setlength\itemsep{0pt}
	\item \texttt{heavy[v]} = hijo pesado de $v$ (el de mayor subarbol).
	\item \texttt{head[v]} = cabeza de la cadena que contiene a $v$.
	\item \texttt{pos[v]} = posicion en el arreglo base.
	\item Cada nodo pertenece a exactamente una cadena.
\end{itemize}

\paragraph{Codigo clave.}
El query sobre un path $u \to v$:
\begin{verbatim}
long long pathQuery(int u, int v) {
    long long res = 0;
    while (head[u] != head[v]) {
        if (depth[head[u]] < depth[head[v]]) swap(u, v);
        res += seg.query(pos[head[u]], pos[u]);
        u = parent[head[u]];
    }
    if (depth[u] > depth[v]) swap(u, v);
    res += seg.query(pos[u], pos[v]);
    return res;
}
\end{verbatim}

\paragraph{Complejidad.}
\begin{itemize}\setlength\itemsep{0pt}
	\item Build $O(n)$, query sobre path $O(\log^2 n)$ (con segtree) o $O(\log n)$ (con BIT si solo se hace suma).
	\item Memoria $O(n)$.
\end{itemize}

\paragraph{Donde tocar para modificar.}
\begin{itemize}\setlength\itemsep{0pt}
	\item \textbf{Usar con segment tree}: aniadir un segtree sobre el arreglo base; \texttt{pathSegments} da los rangos a actualizar/query.
	\item \textbf{Cambiar operacion}: ajustar el segtree (suma, max, etc.).
	\item \textbf{Subtree query}: \texttt{subtree[v]} es el rango \texttt{[pos[v], pos[v] + size[v] - 1]} en el arreglo base.
	\item \textbf{Updates en edges}: aniadir un offset para que el edge $(u, v)$ quede en el nodo mas profundo.
\end{itemize}

\paragraph{Trampas y errores frecuentes.}
\begin{itemize}\setlength\itemsep{0pt}
	\item \texttt{pathSegments} requiere un \texttt{segtree} o \texttt{BIT} externo para ser util; por si solo solo da los rangos.
	\item Para edges vs nodos: si la operacion es sobre edges, guardar el valor en el hijo mas profundo.
	\item Olvidar la recursion \texttt{while (head[u] != head[v])} da resultados incorrectos.
	\item En problemas con updates, asegurarse de que el segtree soporte lazy propagation.
\end{itemize}

\paragraph{Codigo.}
Ver \texttt{cpp.tex}, seccion \textit{Heavy-Light Decomposition (HLD)}.

\vspace{0.5em}

\subsubsection{Centroid Decomposition}

\paragraph{Idea intuitiva.}
Descomponer recursivamente el arbol en \textbf{centroides}. Un centroide de un arbol es un nodo tal que cada subtree (al removerlo) tiene tamano $\le n/2$. Existe siempre y se encuentra en $O(n)$. Esto da una estructura de \textbf{arbol de centroides} de profundidad $O(\log n)$, util para problemas ``para cada nodo, encontrar el mas cercano que cumple X''. La razon de la profundidad $O(\log n)$: cada recursion reduce el problema a subarboles de tamano $\le n/2$, asi que la recursion da una cadena geometrica decreciente.

\paragraph{Propiedades clave (invariantes).}
\begin{itemize}\setlength\itemsep{0pt}
	\item En cada nivel, los subarboles se procesan recursivamente.
	\item \texttt{removed[]} marca nodos ya usados como centroide.
	\item \texttt{centroidParent[]} da el padre en el arbol de centroides.
	\item La altura del arbol de centroides es $O(\log n)$.
\end{itemize}

\paragraph{Codigo clave.}
Encontrar el centroide de un subarbol:
\begin{verbatim}
int findCentroid(int v, int p, int sz) {
    for (int u : adj[v])
        if (u != p && !removed[u] && subsize[u] > sz / 2)
            return findCentroid(u, v, sz);
    return v;
}
\end{verbatim}

\paragraph{Complejidad.}
\begin{itemize}\setlength\itemsep{0pt}
	\item Build $O(n \log n)$, queries dependen del problema.
	\item Cada nivel de recursion es $O(n)$ para calcular subsize.
\end{itemize}

\paragraph{Donde tocar para modificar.}
\begin{itemize}\setlength\itemsep{0pt}
	\item \textbf{Query tipica}: en cada nivel, ``expandir'' desde el centroide contando caminos validos y usar una estructura auxiliar (set, multiset).
	\item \textbf{Distancia maxima a un nodo ``malo''}: mantener un multiset de distancias en cada centroide.
	\item \textbf{Cambiar a otro tipo de descomposicion}: si el problema se adapta mejor a small-to-large o DSU on tree, usar esos.
	\item \textbf{Dynamic updates}: combinacion con link-cut trees.
\end{itemize}

\paragraph{Trampas y errores frecuentes.}
\begin{itemize}\setlength\itemsep{0pt}
	\item El snippet actual \textbf{solo construye} la descomposicion; las queries se aniaden segun el problema.
	\item Olvidar resetear \texttt{removed[]} causa recursion infinita.
	\item El calculo de subsize debe ignorar nodos ya ``removed''.
	\item El centroide es \textbf{unico} salvo empates en tamano; cualquiera sirve.
\end{itemize}

\paragraph{Codigo.}
Ver \texttt{cpp.tex}, seccion \textit{Centroid Decomposition}.

\vspace{0.5em}

\subsubsection{Euler Tour + range queries}

\paragraph{Idea intuitiva.}
Un \textbf{Euler Tour} (first-time visit) numera los nodos en el orden en que se visitan. El subarbol de $v$ corresponde a un \textbf{rango contiguo} $[tin[v], tout[v]]$ en este orden. Asi, queries sobre subarboles se resuelven con un BIT o segment tree sobre el arreglo de tiempos. La razon del rango contiguo: cuando entramos a un subarbol, no salimos hasta haber recorrido todos sus descendientes; entonces las posiciones consecutivas corresponden exactamente al subarbol.

\paragraph{Propiedades clave (invariantes).}
\begin{itemize}\setlength\itemsep{0pt}
	\item \texttt{tin[v]} = tiempo de primera visita a $v$.
	\item \texttt{tout[v]} = tiempo al \textbf{terminar} de procesar el subarbol de $v$.
	\item El subarbol de $v$ ocupa las posiciones $[tin[v], tout[v]]$.
	\item Para verificar si $u$ es ancestro de $v$: \texttt{tin[u] <= tin[v] \&\& tout[v] <= tout[u]}.
\end{itemize}

\paragraph{Codigo clave.}
El DFS del Euler Tour:
\begin{verbatim}
void dfs(int v, int p) {
    tin[v] = ++timer;
    for (int u : adj[v])
        if (u != p) dfs(u, v);
    tout[v] = timer;
}
\end{verbatim}

\paragraph{Complejidad.}
\begin{itemize}\setlength\itemsep{0pt}
	\item Build $O(n)$, query sobre subarbol $O(\log n)$ con BIT/segtree.
	\item Memoria $O(n)$.
\end{itemize}

\paragraph{Donde tocar para modificar.}
\begin{itemize}\setlength\itemsep{0pt}
	\item \textbf{Aniadir valor por nodo}: guardar el valor en \texttt{tin[v]} y operar sobre el rango.
	\item \textbf{Path query}: el Euler Tour \textbf{no} sirve para paths; usar HLD o LCA.
	\item \textbf{Multiple queries offline}: procesar en orden de tiempo y mantener una estructura por profundidad.
	\item \textbf{Verificar ancestro}: $O(1)$ con la condicion \texttt{tin}.
\end{itemize}

\paragraph{Trampas y errores frecuentes.}
\begin{itemize}\setlength\itemsep{0pt}
	\item El snippet solo construye el tour; las queries concretas se aniaden.
	\item \texttt{tout[v]} esta en el momento de salir, asi que \texttt{tout[v] - tin[v] + 1} = tamano del subarbol.
	\item Si el arbol no es conexo, aniadir un loop externo sobre cada componente.
	\item Para queries en tiempo (no subtree), usar el orden del tour y offline processing.
\end{itemize}

\paragraph{Codigo.}
Ver \texttt{cpp.tex}, seccion \textit{Euler Tour + range queries}.

\subsection{BST balanceado}

\subsubsection{Treap implicito}

\paragraph{Idea intuitiva.}
Un \textbf{treap} es un arbol binario de busqueda donde cada nodo tiene una \textbf{clave} (el valor de la secuencia) y una \textbf{prioridad} aleatoria. La invariante es: BST por clave, heap por prioridad. Esto da un arbol \textbf{esperado balanceado} sin necesidad de rebalanceo explicito. Un \textbf{treap implicito} no almacena una clave; la posicion de un nodo en la secuencia esta dada por el \textbf{tamano del subarbol izquierdo}, permitiendo split/merge por posicion. La estructura es esencialmente un RBS (randomized BST): la altura esperada es $O(\log n)$ por la propiedad de heap + BST.

\paragraph{Propiedades clave (invariantes).}
\begin{itemize}\setlength\itemsep{0pt}
	\item \texttt{sz(node)} = tamano del subarbol (1 + izquierdo + derecho).
	\item \texttt{split(t, k, l, r)}: parte \texttt{t} en \texttt{l} (primeros $k$ nodos) y \texttt{r} (resto).
	\item \texttt{merge(l, r)}: une dos treaps (todos los de \texttt{l} antes que los de \texttt{r}).
	\item La prioridad aleatoria asegura equilibrio esperado.
\end{itemize}

\paragraph{Codigo clave.}
Split por posicion (la base del treap implicito):
\begin{verbatim}
void split(Node* t, int k, Node*& l, Node*& r) {
    if (!t) { l = r = nullptr; return; }
    push(t);
    if (sz(t->l) >= k) {
        split(t->l, k, l, t->l);
        r = update(t);
    } else {
        split(t->r, k - sz(t->l) - 1, t->r, r);
        l = update(t);
    }
}
\end{verbatim}

\paragraph{Complejidad.}
\begin{itemize}\setlength\itemsep{0pt}
	\item $O(\log n)$ \textbf{esperado} por operacion (split, merge, insert, erase).
	\item Peor caso exponencial (probabilidad astronomicamente pequena).
\end{itemize}

\paragraph{Donde tocar para modificar.}
\begin{itemize}\setlength\itemsep{0pt}
	\item \textbf{Treap con clave (no implicito)}: usar \texttt{key} en vez de posicion; cambiar \texttt{split} por \texttt{splitByKey}.
	\item \textbf{Lazy propagation}: aniadir campos \texttt{lazy} y aplicarlos en \texttt{push} (recursivo).
	\item \textbf{Persistencia}: aniadir versionado clonando nodos.
	\item \textbf{Mejor aleatoriedad}: en lugar de \texttt{rand()}, usar \texttt{rng()} (mejor estado).
	\item \textbf{Operacion de reversa}: aniadir \texttt{lazyReverse} y swap de hijos en push.
\end{itemize}

\paragraph{Trampas y errores frecuentes.}
\begin{itemize}\setlength\itemsep{0pt}
	\item En el snippet actual, \texttt{Node(int v)} usa \texttt{rand()} que en algunos compiladores da maxima prioridad 0; aniadir \texttt{rng()} para mejor distribucion.
	\item Olvidar \texttt{update(t)} tras modificar hijos da tamano incorrecto.
	\item Si lazy propagation no se aplica en push, las queries dan valores viejos.
	\item Memory leak con \texttt{new}; en contests OK, en produccion aniadir destructor.
\end{itemize}

\paragraph{Codigo.}
Ver \texttt{cpp.tex}, seccion \textit{Treap implicito}.

\vspace{0.5em}

\subsubsection{Mo con updates}

\paragraph{Idea intuitiva.}
Extension del algoritmo de Mo para queries que mezclan \textbf{rango} y \textbf{updates}. Las queries se ordenan por $(L/\text{blockSize}, R/\text{blockSize}, T)$ donde $T$ es el ``tiempo'' (numero de updates). Asi, moverse entre queries consecutivas cuesta $O(n^{2/3})$ amortizado. La intuicion: el sort por bloques reduce los movimientos totales, y la tercera coordenada $T$ se optimiza con alternancia similar a $R$.

\paragraph{Propiedades clave (invariantes).}
\begin{itemize}\setlength\itemsep{0pt}
	\item Tres punteros: $L$, $R$, $T$.
	\item \texttt{applyUpdate} y \texttt{revertUpdate} son \textbf{inversos} exactos.
	\item \texttt{addPos} y \texttt{removePos} mantienen el estado.
	\item El tamano de bloque optimo es $n^{2/3}$ (donde $n$ = tamano del array, $Q$ = numero de queries).
\end{itemize}

\paragraph{Codigo clave.}
El sort con tres coordenadas:
\begin{verbatim}
int blockSize = (int)pow(n, 2.0/3.0);
sort(queries.begin(), queries.end(), [&](const Q& a, const Q& b) {
    int blockA = a.l / blockSize, blockB = b.l / blockSize;
    if (blockA != blockB) return blockA < blockB;
    int blockAR = a.r / blockSize, blockBR = b.r / blockSize;
    if (blockAR != blockBR) return blockAR < blockBR;
    return a.t < b.t;
});
\end{verbatim}

\paragraph{Complejidad.}
\begin{itemize}\setlength\itemsep{0pt}
	\item $O((N + Q) \cdot n^{2/3})$ donde $n^{2/3}$ es el tamano de bloque optimo.
	\item En problemas reales, funciona bien hasta $N, Q \sim 10^5$.
\end{itemize}

\paragraph{Donde tocar para modificar.}
\begin{itemize}\setlength\itemsep{0pt}
	\item \textbf{El corazon es el ``add/remove''}: aniadir la logica del problema (ej. contar colores distintos, suma de frecuencias, etc.).
	\item \textbf{Cambiar blockSize}: si la mayoria de updates vs queries cambia, ajustar.
	\item \textbf{Rollback}: aniadir \texttt{save()} / \texttt{restore()} para volver a estados anteriores.
	\item \textbf{Hilbert order}: alternativa con mejor locality.
\end{itemize}

\paragraph{Trampas y errores frecuentes.}
\begin{itemize}\setlength\itemsep{0pt}
	\item El snippet actual es solo el \textbf{esqueleto}; \texttt{addPos}/\texttt{removePos}/\texttt{applyUpdate}/\texttt{revertUpdate} deben ser implementados por problema.
	\item Olvidar \texttt{revertUpdate} es el error mas comun.
	\item Si el numero de updates es muy grande, Mo con updates se vuelve lento; considerar otra estructura.
	\item Mantener consistencia entre \texttt{T} y la posicion de los updates.
\end{itemize}

\paragraph{Codigo.}
Ver \texttt{cpp.tex}, seccion \textit{Mo con updates}.

\vspace{0.5em}

\subsubsection{sqrt-decomposition generico}

\paragraph{Idea intuitiva.}
Dividir el arreglo en \textbf{bloques} de tamano $B \approx \sqrt{n}$. Mantener informacion precomputada por bloque (suma, max, etc.). Queries de rango se resuelven combinando bloques enteros ($O(\sqrt{n})$ en total) y elementos sueltos en los bloques parciales. Updates puntuales actualizan el valor y el resumen del bloque. La intuicion: las queries ``normales'' cruzan $O(n/B)$ bloques; tomar $B = \sqrt{n}$ balancea el costo de bloques enteros vs elementos sueltos.

\paragraph{Propiedades clave (invariantes).}
\begin{itemize}\setlength\itemsep{0pt}
	\item $B = \lfloor \sqrt{n} \rfloor + 1$.
	\item \texttt{bucket[i]} = resumen del bloque $i$-esimo (suma, max, etc.).
	\item \texttt{rangeSum(l, r)}: 3 casos segun $l$ y $r$ esten en el mismo bloque o no.
	\item Cada query cruza $\le 2 + (r - l)/B$ bloques.
\end{itemize}

\paragraph{Codigo clave.}
Estructura basica:
\begin{verbatim}
int B = sqrt(n) + 1;
vector<T> bucket(n / B + 1, IDENTITY);
T query(int l, int r) {
    T res = IDENTITY;
    int bl = l / B, br = r / B;
    if (bl == br) {
        for (int i = l; i <= r; ++i) res = combine(res, a[i]);
    } else {
        for (int i = l; i < (bl+1)*B; ++i) res = combine(res, a[i]);
        for (int i = bl+1; i < br; ++i) res = combine(res, bucket[i]);
        for (int i = br*B; i <= r; ++i) res = combine(res, a[i]);
    }
    return res;
}
\end{verbatim}

\paragraph{Complejidad.}
\begin{itemize}\setlength\itemsep{0pt}
	\item $O(\sqrt{n})$ por operacion.
	\item Memoria $O(n)$.
\end{itemize}

\paragraph{Donde tocar para modificar.}
\begin{itemize}\setlength\itemsep{0pt}
	\item \textbf{Cambiar suma por max/min}: ajustar \texttt{bucket} y la combinacion.
	\item \textbf{Aniadir update en rango}: si la operacion es asociativa, sumar al resumen del bloque + ajustar elementos sueltos.
	\item \textbf{Tamano de bloque optimo}: si las queries son muy largas, $B = n^{2/3}$ (sirve para Mo con updates).
	\item \textbf{Operacion no asociativa}: usar segment tree en su lugar.
\end{itemize}

\paragraph{Trampas y errores frecuentes.}
\begin{itemize}\setlength\itemsep{0pt}
	\item Si $n$ no es multiplo de $B$, el ultimo bloque es mas chico; aniadir check.
	\item La identidad de la operacion debe estar bien definida (0 para suma, $-\infty$ para max).
	\item Si la operacion no es asociativa, los bloques no se pueden combinar; cambiar a segment tree.
\end{itemize}

\paragraph{Codigo.}
Ver \texttt{cpp.tex}, seccion \textit{sqrt-decomposition generico}.

% ============================================================
\section{Strings}
% ============================================================

\subsection{Hashing}

\subsubsection{Rolling Hash (doble modulo)}

\paragraph{Idea intuitiva.}
Un \textbf{rolling hash} asigna a cada prefijo del string un valor numerico tal que el hash de un substring se puede obtener en $O(1)$ a partir de los prefijos. La forma clasica: $h[i] = h[i-1] \cdot A + s[i-1]$ (donde $A$ es una base y todo modulo $M$). Asi, $h[r] - h[l-1] \cdot A^{r-l+1}$ da el hash del substring $[l, r]$ (con $A^{r-l+1}$ precomputado). Usar \textbf{dos modulos} distintos reduce drasticamente la probabilidad de colision. La idea algebraica: el string se interpreta como un polinomio en $A$ con coeficientes los caracteres; el hash es su evaluacion modulo $M$.

\paragraph{Propiedades clave (invariantes).}
\begin{itemize}\setlength\itemsep{0pt}
	\item $h[0] = 0$, $h[i]$ = hash de los primeros $i$ caracteres.
	\item $p[i] = A^i \bmod M$ precomputado.
	\item Comparar dos substrings: comparar sus pares de hashes.
	\item \textbf{Colisiones}: posibles (probabilidad $1/M$ por par). Doble hash: $\approx 1/(M_1 M_2)$.
	\item Si $A$ y $M$ son coprimos, el hash es bijectivo modulo $M$ para strings de longitud fija.
\end{itemize}

\paragraph{Codigo clave.}
La formula del hash de substring:
\begin{verbatim}
pair<long long, long long> get(int l, int r) {
    // h[r] - h[l-1] * A^(r-l+1)
    long long x1 = (h1[r] - h1[l-1] * p1[r-l+1] % M1 + M1) % M1;
    long long x2 = (h2[r] - h2[l-1] * p2[r-l+1] % M2 + M2) % M2;
    return {x1, x2};
}
\end{verbatim}

\paragraph{Complejidad.}
\begin{itemize}\setlength\itemsep{0pt}
	\item Precompute $O(n)$, query hash $O(1)$.
	\item Memoria $O(n)$.
\end{itemize}

\paragraph{Donde tocar para modificar.}
\begin{itemize}\setlength\itemsep{0pt}
	\item \textbf{Caracteres con valor $> 26$}: si los strings incluyen simbolos, aniadir \texttt{if (c >= 'a' \&\& c <= 'z') v = c - 'a'; else v = 26 + ...}.
	\item \textbf{Hash de substrings con updates}: aniadir un Fenwick tree / segment tree sobre los hashes (no es trivial).
	\item \textbf{Cambiar mod}: usar dos primos grandes ($10^9+7$ y $10^9+9$) o un \texttt{u128} para un solo hash sin mod.
	\item \textbf{Base aleatoria}: aniadir random para evitar ataques (en contests no es necesario; en produccion si).
	\item \textbf{Hash polinomial}: $h = \sum s_i \cdot A^i$ en vez de $s_i \cdot A^{i-1}$; algunos problemas requieren este.
\end{itemize}

\paragraph{Trampas y errores frecuentes.}
\begin{itemize}\setlength\itemsep{0pt}
	\item Si los substrings a comparar estan vacios (\texttt{l > r}), su hash es 0; aniadir check.
	\item \textbf{Overflow} en \texttt{h[i-1] * A}: usar \texttt{long long} y aplicar \texttt{\% M} en cada operacion (no al final).
	\item La sustraccion \texttt{h[r] - h[l-1]*A^(r-l+1)} puede dar negativo; aniadir \texttt{+M} antes de \texttt{\% M}.
	\item Con doble hash, verificar \textbf{ambos} componentes; si solo uno coincide, es colision.
\end{itemize}

\paragraph{Codigo.}
Ver \texttt{cpp.tex}, seccion \textit{Rolling Hash (doble modulo)}.

\subsection{Patrones}

\subsubsection{KMP + LPS}

\paragraph{Idea intuitiva.}
El algoritmo \textbf{Knuth-Morris-Pratt} busca todas las ocurrencias de un patron $P$ en un texto $T$ en $O(|T| + |P|)$. La clave es la tabla \textbf{LPS} (Longest Proper Prefix which is Suffix) que, para cada posicion $i$ del patron, dice el largo del prefijo mas largo que tambien es sufijo \textbf{propio}. Durante el match, cuando hay mismatch, en vez de volver al inicio, saltamos al LPS anterior, evitando trabajo redundante. La demostracion de $O(|T| + |P|)$: cada iteracion del bucle externo avanza $i$ o decrece $j$ (con $j$ acotado por $|P|$), asi que en total hay $O(|T| + |P|)$ operaciones.

\paragraph{Propiedades clave (invariantes).}
\begin{itemize}\setlength\itemsep{0pt}
	\item \texttt{lps[i]} = largo del prefijo mas largo que es sufijo propio de \texttt{p[0..i]}.
	\item El match usa dos punteros \texttt{i} (texto) y \texttt{j} (patron); en mismatch, $j = lps[j-1]$.
	\item Cuando $j == |P|$, se encontro una ocurrencia; continuar con $j = lps[j-1]$ para buscar mas.
	\item El LPS se construye con una auto-aplicacion de KMP al patron.
\end{itemize}

\paragraph{Codigo clave.}
Construccion de la tabla LPS:
\begin{verbatim}
for (int i = 1; i < m; ++i) {
    int j = lps[i - 1];
    while (j > 0 && p[i] != p[j]) j = lps[j - 1];
    if (p[i] == p[j]) ++j;
    lps[i] = j;
}
\end{verbatim}

\paragraph{Complejidad.}
\begin{itemize}\setlength\itemsep{0pt}
	\item $O(|T| + |P|)$.
	\item En la practica, rapido y robusto.
\end{itemize}

\paragraph{Donde tocar para modificar.}
\begin{itemize}\setlength\itemsep{0pt}
	\item \textbf{Contar ocurrencias}: el snippet ya cuenta; para devolver posiciones, aniadir un \texttt{vector<int>}.
	\item \textbf{Varios patrones}: usar Aho-Corasick en su lugar.
	\item \textbf{Matcher exacto con wildcard}: aniadir logica para que \texttt{?} matchee cualquier caracter.
	\item \textbf{Tamano pequeno}: si el patron es $\le 5$ caracteres, Rabin-Karp es mas simple.
	\item \textbf{2D KMP}: aniadir una dimension; util en problemas de matrices.
\end{itemize}

\paragraph{Trampas y errores frecuentes.}
\begin{itemize}\setlength\itemsep{0pt}
	\item En el snippet, \texttt{i} y \texttt{j} son 0-based. La condicion \texttt{j == sz\_pattern} para contar funciona.
	\item La construccion de LPS usa el mismo patron como ``texto''.
	\item Si el patron tiene caracteres repetidos, el LPS puede ser no trivial; verificar con casos pequenos.
\end{itemize}

\paragraph{Codigo.}
Ver \texttt{cpp.tex}, seccion \textit{KMP + LPS}.

\vspace{0.5em}

\subsubsection{Z-algorithm}

\paragraph{Idea intuitiva.}
La \textbf{Z-function} $Z[i]$ de un string $S$ es el largo del prefijo mas largo de $S$ que \textbf{tambien} empieza en $S[i]$. Asi, $Z[0] = n$ por convencion. Se calcula en $O(n)$ manteniendo una ventana $[l, r]$ que es el match mas largo encontrado hasta ahora. La consulta $S[l..r] = S[0..r-l]$ da informacion sobre $Z[i]$ cuando $i \le r$. La intuicion: cualquier $i$ dentro de la ventana ``ve'' los caracteres desde una posicion conocida, asi que $Z[i]$ se puede copiar en lugar de recalcular.

\paragraph{Propiedades clave (invariantes).}
\begin{itemize}\setlength\itemsep{0pt}
	\item Construir $Z$ cuesta $O(n)$ con el truco de la ventana.
	\item Para matching de patron $P$ en texto $T$: concatenar $P + \# + T$ y aniadir todos los $i$ con $Z[i] = |P|$.
	\item La ventana $[l, r]$ es siempre la del match mas largo a la derecha.
\end{itemize}

\paragraph{Codigo clave.}
Construccion de Z-function en $O(n)$:
\begin{verbatim}
Z[0] = n;
int l = 0, r = 0;
for (int i = 1; i < n; ++i) {
    if (i <= r) Z[i] = min(r - i + 1, Z[i - l]);
    while (i + Z[i] < n && s[Z[i]] == s[i + Z[i]]) ++Z[i];
    if (i + Z[i] - 1 > r) { l = i; r = i + Z[i] - 1; }
}
\end{verbatim}

\paragraph{Complejidad.}
\begin{itemize}\setlength\itemsep{0pt}
	\item $O(n)$ para construir; matching $O(n + m)$.
	\item Cada caracter se compara a lo sumo 2 veces (extension + reduccion).
\end{itemize}

\paragraph{Donde tocar para modificar.}
\begin{itemize}\setlength\itemsep{0pt}
	\item \textbf{Detectar bordes (border)}: $Z[i] = n - i$ indica que $S[0..n-i-1]$ es borde.
	\item \textbf{Prefijo comun mas largo de dos strings}: aniadir $S_1 + \# + S_2$; $Z[|S_1|+1]$ es el prefijo comun.
	\item \textbf{Comparar prefijos y sufijos}: combinar $S + \# + \text{reverse}(S)$.
	\item \textbf{Matching aproximado}: aniadir tolerancia a mismatches.
\end{itemize}

\paragraph{Trampas y errores frecuentes.}
\begin{itemize}\setlength\itemsep{0pt}
	\item $Z[0]$ siempre vale $n$ (no se calcula).
	\item Si $i > r$, la ventana no aporta; empezar desde 0.
	\item El caracter separador en $P + \# + T$ debe ser uno que no aparezca en $P$ o $T$.
\end{itemize}

\paragraph{Codigo.}
Ver \texttt{cpp.tex}, seccion \textit{Z-algorithm}.

\vspace{0.5em}

\subsubsection{Rabin-Karp (single hash)}

\paragraph{Idea intuitiva.}
Rolling hash aplicado a pattern matching: calcular el hash del patron $P$ y de cada ventana de tamano $|P|$ en $T$. Si los hashes coinciden, verificar caracter por caracter (para descartar colisiones). Esto da $O(n + m)$ \textbf{esperado}. La razon de la eficiencia esperada: cada ventana tiene probabilidad $1/M$ de tener un hash coincidente por casualidad, asi que solo se verifica $O(n/M)$ ventanas en promedio.

\paragraph{Propiedades clave (invariantes).}
\begin{itemize}\setlength\itemsep{0pt}
	\item Si los hashes difieren, los substrings son distintos (sin colision).
	\item Si coinciden, \textbf{verificar} (la verificacion es $O(m)$ pero rara vez se ejecuta).
	\item Peor caso (muchas colisiones): $O(nm)$.
	\item Para $M \ge 10^9$, las colisiones son raras; en la practica nunca pasan.
\end{itemize}

\paragraph{Codigo clave.}
Rolling hash del texto (deslizar ventana):
\begin{verbatim}
long long hash = 0;
for (int i = 0; i < m; ++i) hash = (hash * BASE + text[i]) % MOD;
for (int i = m; i < n; ++i) {
    if (hash == target) check(i - m);
    hash = (hash - text[i - m] * h) * BASE + text[i];  // mod correctamente
    hash = (hash % MOD + MOD) % MOD;
}
\end{verbatim}

\paragraph{Complejidad.}
\begin{itemize}\setlength\itemsep{0pt}
	\item Esperado $O(n + m)$, peor $O(nm)$.
	\item Con doble hash y $M_1, M_2 \ge 10^9$, el peor caso es astronomicamente improbable.
\end{itemize}

\paragraph{Donde tocar para modificar.}
\begin{itemize}\setlength\itemsep{0pt}
	\item \textbf{Doble hash}: aniadir un segundo par \texttt{(BASE2, MOD2)} para reducir colisiones.
	\item \textbf{Varios patrones}: aniadir un set de hashes y buscar match con cualquiera.
	\item \textbf{Tamano de ventana variable}: si los patrones tienen largos distintos, usar el largo especifico por patron.
	\item \textbf{Multi-pattern search}: usar Aho-Corasick en su lugar.
\end{itemize}

\paragraph{Trampas y errores frecuentes.}
\begin{itemize}\setlength\itemsep{0pt}
	\item Overflow en \texttt{(tHash - text[i] * h) * BASE}: aplicar mod correctamente y aniadir \texttt{+MOD} antes de multiplicar.
	\item El factor $h = \text{BASE}^{m-1} \bmod \text{MOD}$ se precomputa.
	\item Si \texttt{MOD} no es primo, las colisiones son mas frecuentes; usar primo.
	\item En el peor caso, Rabin-Karp puede ser lento; preferir KMP si se requiere worst-case.
\end{itemize}

\paragraph{Codigo.}
Ver \texttt{cpp.tex}, seccion \textit{Rabin-Karp (single hash)}.

\vspace{0.5em}

\subsubsection{Aho-Corasick (basic)}

\paragraph{Idea intuitiva.}
Construye un \textbf{trie} con todos los patrones y aniade \textbf{failure links} analogos a los de KMP: el failure link de un nodo $v$ apunta al nodo mas largo que es sufijo del string representado por $v$ y sigue siendo prefijo de algun patron. Asi, al recorrer el texto, si hay mismatch, saltamos al failure link. Esto permite buscar \textbf{todos} los patrones en $O(|T| + \sum |P_i|)$. La idea: unificar varios KMP en una sola estructura.

\paragraph{Propiedades clave (invariantes).}
\begin{itemize}\setlength\itemsep{0pt}
	\item El trie almacena los patrones.
	\item \texttt{link[v]} (failure link) se calcula en BFS en orden de profundidad creciente.
	\item \texttt{out[v]} lista los IDs de patrones que terminan en $v$.
	\item Para queries eficientes, aniadir \textbf{output links}: propagar \texttt{out} de los failure links a los hijos.
\end{itemize}

\paragraph{Codigo clave.}
BFS para construir los failure links:
\begin{verbatim}
queue<int> q;
for (int c = 0; c < SIGMA; ++c)
    if (trie[0].next[c] != -1) {
        q.push(trie[0].next[c]);
        trie[trie[0].next[c]].link = 0;
    } else trie[0].next[c] = 0;  // goto a la raiz
while (!q.empty()) {
    int v = q.front(); q.pop();
    for (int pid : trie[trie[v].link].out)
        trie[v].out.push_back(pid);
    for (int c = 0; c < SIGMA; ++c) {
        if (trie[v].next[c] != -1) {
            trie[trie[v].next[c]].link = trie[trie[v].link].next[c];
            q.push(trie[v].next[c]);
        } else trie[v].next[c] = trie[trie[v].link].next[c];
    }
}
\end{verbatim}

\paragraph{Complejidad.}
\begin{itemize}\setlength\itemsep{0pt}
	\item Build $O(\sum |P_i| \cdot \Sigma)$, query $O(|T| + \text{ocurrencias})$.
\end{itemize}

\paragraph{Donde tocar para modificar.}
\begin{itemize}\setlength\itemsep{0pt}
	\item \textbf{Output links en build}: ya hecho en el snippet (\texttt{for(int pid : t[t[u].link].out) t[u].out.push\_back(pid);}).
	\item \textbf{Transiciones automaticas}: aniadir el goto \texttt{t[v].next[k] = t[t[v].link].next[k]} cuando el enlace directo es \texttt{-1}; asi el bucle de matching es un solo paso por caracter.
	\item \textbf{Aplicar a problemas tipo DP}: contar formas de construir strings evitando los patrones (ver el siguiente modulo).
	\item \textbf{2D Aho-Corasick}: aniadir una dimension para busqueda en matrices.
\end{itemize}

\paragraph{Trampas y errores frecuentes.}
\begin{itemize}\setlength\itemsep{0pt}
	\item La BFS debe empezar con la raiz en la cola y los hijos de la raiz con \texttt{link = 0}.
	\item Olvidar propagar \texttt{out} da ocurrencias perdidas.
	\item Si no se aniaden las transiciones automaticas (goto), el matching es mas lento.
	\item Con muchos patrones, el output puede ser enorme; aniadir contador en lugar de lista.
\end{itemize}

\paragraph{Codigo.}
Ver \texttt{cpp.tex}, seccion \textit{Aho-Corasick (basic)}.

\vspace{0.5em}

\subsubsection{Aho-Corasick DP for forbidden patterns}

\paragraph{Idea intuitiva.}
Contar strings de longitud $L$ que \textbf{evitan} un conjunto de patrones. Se modela como una DP sobre el \textbf{estado del automaton} de Aho-Corasick: $dp[v]$ = numero de formas de llegar al estado $v$ despues de cierto numero de caracteres. Transicion: por cada caracter $c$ en el alfabeto, ir al estado \texttt{next[v][c]}, pero \textbf{solo} si ese estado no tiene un patron (de lo contrario, el string contendria el patron prohibido). La idea: cada estado del automaton representa un ``prefijo conocido'' y las transiciones simulan ``agregar un caracter''.

\paragraph{Propiedades clave (invariantes).}
\begin{itemize}\setlength\itemsep{0pt}
	\item $dp$ se mantiene por nivel (longitud del string construido).
	\item Transicion: $ndp[u] = \sum_{c} dp[v]$ donde $u = next[v][c]$ y $u$ no es ``terminal'' (no tiene patron que termina ahi).
	\item Respuesta: suma de $dp[v]$ sobre todos los estados $v$ no terminales al final.
	\item Si todos los caracteres son permitidos, $|\Sigma|$ entra como factor.
\end{itemize}

\paragraph{Codigo clave.}
DP sobre el automaton:
\begin{verbatim}
vector<long long> dp(numStates, 0), ndp(numStates);
dp[0] = 1;  // empezar en la raiz
for (int step = 0; step < L; ++step) {
    fill(ndp.begin(), ndp.end(), 0);
    for (int v = 0; v < numStates; ++v) if (dp[v] && !isTerminal[v])
        for (int c = 0; c < SIGMA; ++c) {
            int u = next[v][c];
            if (!isTerminal[u]) ndp[u] = (ndp[u] + dp[v]) % MOD;
        }
    swap(dp, ndp);
}
long long ans = 0;
for (int v = 0; v < numStates; ++v) if (!isTerminal[v])
    ans = (ans + dp[v]) % MOD;
\end{verbatim}

\paragraph{Complejidad.}
\begin{itemize}\setlength\itemsep{0pt}
	\item $O(L \cdot |S| \cdot |\Sigma|)$ donde $|S|$ es el numero de estados y $|\Sigma|$ el alfabeto.
	\item Para $|\Sigma| = 2$ y $|S| \le 10^4$: $\sim 2 \cdot 10^{10}$ para $L = 10^6$ (puede ser lento).
\end{itemize}

\paragraph{Donde tocar para modificar.}
\begin{itemize}\setlength\itemsep{0pt}
	\item \textbf{Alfabeto mas pequeno}: si el problema es solo $\{0, 1\}$, $|\Sigma| = 2$.
	\item \textbf{Permitir ocurrencias parciales}: cambiar la condicion a ``no estado terminal'' o relajar segun el problema.
	\item \textbf{Recuperar el string}: aniadir backtracking si el problema lo pide.
	\item \textbf{Matriz de transicion}: para queries rapidas, precomputar $M^c$ con exponentiation.
\end{itemize}

\paragraph{Trampas y errores frecuentes.}
\begin{itemize}\setlength\itemsep{0pt}
	\item El snippet asume Aho-Corasick ya construido con transiciones automaticas. Sin eso, aniadir el goto en el bucle.
	\item Si $L$ es muy grande, la DP lineal es lenta; usar matrix exponentiation.
	\item Estados terminales deben excluirse de \textbf{ambos} lados (estado actual y siguiente).
\end{itemize}

\paragraph{Codigo.}
Ver \texttt{cpp.tex}, seccion \textit{Aho-Corasick DP for forbidden patterns}.

\subsection{Estructuras avanzadas}

\subsubsection{Suffix Array}

\paragraph{Idea intuitiva.}
Un \textbf{suffix array} $SA$ es un arreglo que contiene las posiciones de inicio de los sufijos del string $S$, ordenados lexicograficamente. Se construye en $O(n \log n)$ con doubling: ordenar por el primer caracter, luego por los primeros 2, luego 4, etc. Para hacer cada paso lineal, se usa \textbf{counting sort} sobre los ranks. La idea: cada nivel dobla la cantidad de caracteres considerados, asi que en $\log n$ niveles los sufijos estan completamente ordenados.

\paragraph{Propiedades clave (invariantes).}
\begin{itemize}\setlength\itemsep{0pt}
	\item \texttt{p[i]} = posicion de inicio del $i$-esimo sufijo (ordenado).
	\item \texttt{c[i]} = rank del sufijo que empieza en $i$.
	\item Cada iteracion duplica la longitud de la ``ventana'' considerada.
	\item El algoritmo termina cuando $2^k \ge n$.
	\item La permutacion \texttt{p} es estable entre iteraciones.
\end{itemize}

\paragraph{Codigo clave.}
El paso de doubling (clasificar por los primeros $2^k$ caracteres):
\begin{verbatim}
for (int k = 0; (1 << k) < n; ++k) {
    for (int i = 0; i < n; ++i) pn[i] = (p[i] - (1 << k) + n) % n;
    // counting sort por c[pn[i]] segundo, c[i] primero
    // ...
    cn[pn[0]] = 0;
    for (int i = 1; i < n; ++i)
        cn[pn[i]] = cn[pn[i-1]] + (pair != pair_prev);
    c = cn;
}
\end{verbatim}

\paragraph{Complejidad.}
\begin{itemize}\setlength\itemsep{0pt}
	\item $O(n \log n)$ con counting sort.
	\item Constante pequena; en la practica eficiente.
\end{itemize}

\paragraph{Donde tocar para modificar.}
\begin{itemize}\setlength\itemsep{0pt}
	\item \textbf{Suffix array en $O(n)$}: usar SA-IS (mas complejo, no incluido).
	\item \textbf{Buscar patron}: con SA + LCP es $O(|P| \log n)$ por busqueda binaria.
	\item \textbf{Suffix array de varios strings}: aniadir centinelas (\texttt{\$}) entre ellos.
	\item \textbf{Aniadir longitud}: aniadir \texttt{n - p[i]} como clave secundaria (sufijos mas cortos primero).
\end{itemize}

\paragraph{Trampas y errores frecuentes.}
\begin{itemize}\setlength\itemsep{0pt}
	\item La ``ventana'' $2^k$ se hace circular: \texttt{p[i] = (p[i] - (1<<k) + n) \% n}.
	\item Para contar bien los ranks duplicados, comparar pares \texttt{(c[i], c[(i+2\^{}k) \% n])}.
	\item El counting sort necesita un array de acumulacion de tamano $\le n$.
	\item Si los caracteres son negativos o fuera del alfabeto esperado, normalizar antes.
\end{itemize}

\paragraph{Codigo.}
Ver \texttt{cpp.tex}, seccion \textit{Suffix Array}.

\vspace{0.5em}

\subsubsection{LCP Array (Kasai)}

\paragraph{Idea intuitiva.}
El \textbf{LCP array} (Longest Common Prefix) guarda, para cada par de sufijos adyacentes en el SA, el largo de su prefijo comun. El algoritmo de \textbf{Kasai} lo calcula en $O(n)$ iterando los sufijos \textbf{en orden del string original} (no del SA), y reutiliza el LCP previo. La demostracion: el sufijo en $i+1$ se obtiene del sufijo en $i$ ``avanzando un caracter''; los LCPs decrecen como maximo 1 entre iteraciones.

\paragraph{Propiedades clave (invariantes).}
\begin{itemize}\setlength\itemsep{0pt}
	\item \texttt{lcp[i]} = LCP entre los sufijos en \texttt{sa[i]} y \texttt{sa[i-1]}.
	\item Se necesita \texttt{rank\_[i]} (la inversa del SA) para iterar.
	\item El truco de Kasai: si \texttt{lcp > 0} para el sufijo anterior, el actual parte de al menos \texttt{lcp - 1}.
	\item La suma de todos los LCPs da informacion sobre repeticiones.
\end{itemize}

\paragraph{Codigo clave.}
Algoritmo de Kasai:
\begin{verbatim}
vector<int> rank(n);
for (int i = 0; i < n; ++i) rank[sa[i]] = i;
int k = 0;
vector<int> lcp(n - 1);
for (int i = 0; i < n; ++i) {
    if (rank[i] == n - 1) { k = 0; continue; }
    int j = sa[rank[i] + 1];
    while (i + k < n && j + k < n && s[i+k] == s[j+k]) ++k;
    lcp[rank[i]] = k;
    if (k) --k;
}
\end{verbatim}

\paragraph{Complejidad.}
\begin{itemize}\setlength\itemsep{0pt}
	\item $O(n)$.
	\item Cada comparacion reduce $k$ o avanza linealmente; amortizado $O(1)$.
\end{itemize}

\paragraph{Donde tocar para modificar.}
\begin{itemize}\setlength\itemsep{0pt}
	\item \textbf{RMQ sobre LCP}: para queries ``LCP de dos sufijos'', encontrar sus ranks y hacer RMQ en \texttt{lcp}.
	\item \textbf{Contar substrings distintas}: $\sum (n - sa[i] - lcp[i])$.
	\item \textbf{Longest common substring}: maximo en \texttt{lcp}.
	\item \textbf{String matching}: busqueda binaria sobre SA + LCP verification.
\end{itemize}

\paragraph{Trampas y errores frecuentes.}
\begin{itemize}\setlength\itemsep{0pt}
	\item Si dos sufijos son iguales (raro, solo si el string tiene period $< n$), \texttt{lcp} puede ser \texttt{n - sa[i]}.
	\item En la iteracion, aniadir el chequeo \texttt{i + k < n} para evitar overflow.
	\item El \texttt{--k} en la salida es crucial para la complejidad amortizada.
\end{itemize}

\paragraph{Codigo.}
Ver \texttt{cpp.tex}, seccion \textit{LCP Array (Kasai)}.

\vspace{0.5em}

\subsubsection{Suffix Automaton}

\paragraph{Idea intuitiva.}
Un \textbf{suffix automaton} es un DFA minimo que reconoce todos los substrings del string original. Tiene la propiedad magica de tener $\le 2n - 1$ estados. Cada estado representa una clase de equivalencia de substrings (los que tienen el mismo conjunto de posiciones finales). El \textbf{link} de un estado apunta al estado del sufijo propio mas largo.

\paragraph{Propiedades clave (invariantes).}
\begin{itemize}\setlength\itemsep{0pt}
	\item \texttt{len[v]}: largo maximo de los substrings de la clase.
	\item \texttt{link[v]}: padre en el arbol de sufijos (suffix link).
	\item \texttt{next[v][c]}: transicion por caracter.
	\item \texttt{occ[v]}: numero de posiciones finales (cuantas veces aparece el patron).
\end{itemize}

\paragraph{Complejidad.}
\begin{itemize}\setlength\itemsep{0pt}
	\item Build $O(n)$ con la operacion de \textbf{clonar} estados.
\end{itemize}

\paragraph{Donde tocar para modificar.}
\begin{itemize}\setlength\itemsep{0pt}
	\item \textbf{Contar substrings distintas}: $\sum_{v \ne 0} (len[v] - len[link[v]])$.
	\item \textbf{Longest common substring}: navegar el SAM con el segundo string, manteniendo el match actual.
	\item \textbf{Substring count}: despues del build, ordenar estados por \texttt{len} descendente y propagar \texttt{occ}.
	\item \textbf{Endpos queries}: el SAM es el DFA canonico de los substrings; queries ``cuantas veces aparece $P$'' se responden en $O(|P|)$.
\end{itemize}

\paragraph{Trampas y errores frecuentes.}
\begin{itemize}\setlength\itemsep{0pt}
	\item El \textbf{clone} debe tener \texttt{occ = 0} (no es una nueva posicion final); el \texttt{link} del estado siendo creado y del estado clonado deben apuntar al clon.
	\item Saltarse esto da SAM incorrecto.
\end{itemize}

\paragraph{Codigo.}
Ver \texttt{cpp.tex}, seccion \textit{Suffix Automaton}.

\vspace{0.5em}

\subsubsection{Lyndon factorization (Duval)}

\paragraph{Idea intuitiva.}
Un \textbf{string de Lyndon} es un string que es estrictamente menor que todos sus rotaciones. Toda cadena tiene una \textbf{factorizacion unica} de Lyndon (algoritmo de \textbf{Duval}): una secuencia de strings de Lyndon $L_1, L_2, \dots, L_k$ tales que $L_1 \ge L_2 \ge \dots \ge L_k$ y la concatenacion es el string original. Esto da $O(n)$ para encontrar los cortes. La demostracion de correctitud: cada factor es minimal en su clase de rotacion, asi que la concatenacion respeta el orden lexicografico.

\paragraph{Propiedades clave (invariantes).}
\begin{itemize}\setlength\itemsep{0pt}
	\item Cada factor de Lyndon es de la forma $uv$ donde $u$ se repite y $v$ es prefijo de $u$.
	\item La longitud del factor esta dada por la ``longitud del periodo''.
	\item Aplicacion: el ``Lyndon word'' es el menor en orden lexicografico entre todas las rotaciones.
	\item La factorizacion es \textbf{unica}.
\end{itemize}

\paragraph{Codigo clave.}
El algoritmo de Duval (iterativo):
\begin{verbatim}
int n = s.size();
int i = 0;
while (i < n) {
    int j = i + 1, k = i;
    while (j < n && s[k] <= s[j]) {
        if (s[k] < s[j]) k = i;
        else ++k;
        ++j;
    }
    int len = j - k;
    while (i < j) { cout << "factor: " << s.substr(i, len) << "\n"; i += len; }
}
\end{verbatim}

\paragraph{Complejidad.}
\begin{itemize}\setlength\itemsep{0pt}
	\item $O(n)$.
	\item Cada comparacion es $O(1)$; el puntero $k$ retrocede a lo sumo $n$ veces en total.
\end{itemize}

\paragraph{Donde tocar para modificar.}
\begin{itemize}\setlength\itemsep{0pt}
	\item \textbf{Encontrar la menor rotacion}: el menor sufijo lexicografico de \texttt{s + s} es la menor rotacion; la posicion se obtiene con un solo pase de Duval.
	\item \textbf{Algoritmo de Booth}: alternativa para rotacion minima, tambien $O(n)$.
	\item \textbf{Validar}: si necesitas verificar que un string es de Lyndon, compararlo con sus rotaciones.
	\item \textbf{Encontrar todos los factores}: guardar la posicion y longitud de cada uno.
\end{itemize}

\paragraph{Trampas y errores frecuentes.}
\begin{itemize}\setlength\itemsep{0pt}
	\item La condicion \texttt{s[j] <= s[k]} (no \texttt{<}) en el while de Duval es crucial; cambiar a \texttt{<} da comportamiento distinto.
	\item El caso $s[k] < s[j]$ resetea $k$ a $i$ (inicio del factor actual).
	\item El bucle externo debe avanzar $i$ por la longitud del factor, no por 1.
\end{itemize}

\paragraph{Codigo.}
Ver \texttt{cpp.tex}, seccion \textit{Lyndon factorization (Duval)}.

\subsection{Palindromos}

\subsubsection{Manacher}

\paragraph{Idea intuitiva.}
Calcula en $O(n)$ los \textbf{radios} de todos los palindromos centrados en cada posicion (impares y pares). La idea: mantener una ventana $[l, r]$ que es el palindromo centrado en un punto anterior mas largo encontrado. Si el nuevo centro $i$ cae dentro de $[l, r]$, entonces el palindromo centrado en $i$ tiene al menos el radio del palindromo ``espejo'' $l + r - i$; sino, empezar desde 1 (impar) o 0 (par). Despues, expandir lo que falte. La demostracion de $O(n)$: cada caracter se compara a lo sumo 2 veces (extension + limite por ventana).

\paragraph{Propiedades clave (invariantes).}
\begin{itemize}\setlength\itemsep{0pt}
	\item Conveccion del snippet: \texttt{odd[i]} = $k$ tal que el palindromo impar mas largo centrado en $i$ tiene longitud $2k - 1$. \texttt{even[i]} = $k$ para par centrado entre $i-1$ e $i$, longitud $2k$.
	\item $k$ cuenta tambien cuantos palindromos hay centrados ahi.
	\item Las queries \texttt{isPalindrome(l, r)} son $O(1)$.
	\item La ventana $[l, r]$ siempre contiene el palindromo mas largo encontrado hasta ahora.
\end{itemize}

\paragraph{Codigo clave.}
Manacher para palindromos impares:
\begin{verbatim}
int l = 0, r = -1;
for (int i = 0; i < n; ++i) {
    int k = (i > r) ? 1 : min(odd[l + r - i], r - i + 1);
    while (i - k >= 0 && i + k < n && s[i-k] == s[i+k]) ++k;
    odd[i] = k;
    if (i + k - 1 > r) { l = i - k + 1; r = i + k - 1; }
}
\end{verbatim}

\paragraph{Complejidad.}
\begin{itemize}\setlength\itemsep{0pt}
	\item Construccion $O(n)$, queries $O(1)$.
	\item Cada extension se cuenta una sola vez.
\end{itemize}

\paragraph{Donde tocar para modificar.}
\begin{itemize}\setlength\itemsep{0pt}
	\item \textbf{Solo contar palindromos}: ya implementado en \texttt{countPalindromes}.
	\item \textbf{Longest palindromic substring}: ya en \texttt{longestPalindrome}.
	\item \textbf{Modificar la conveccion de radios}: si queres ``radio = mitad de la longitud'' en vez de ``k palindromos'', ajustar el offset.
	\item \textbf{Queries con k no precomputado}: aniadir helpers especificos.
	\item \textbf{Modificar el caracter de relleno}: para multiples separadores, aniadir un caracter unico.
\end{itemize}

\paragraph{Trampas y errores frecuentes.}
\begin{itemize}\setlength\itemsep{0pt}
	\item La conveccion del snippet es \textbf{distinta} de la ``clasica de radio'' (donde radio = $(L-1)/2$); tener cuidado al copiar a otro codigo.
	\item Si los palindromos se solapan, la condicion \texttt{i + k - 1 > r} actualiza la ventana.
	\item En palindromos pares, el centro esta ``entre'' dos indices; aniadir offset en queries.
\end{itemize}

\paragraph{Codigo.}
Ver \texttt{cpp.tex}, seccion \textit{Manacher}.

% ============================================================
\section{Grafos}
% ============================================================

\subsection{Caminos minimos}

\subsubsection{Dijkstra}

\paragraph{Idea intuitiva.}
Encuentra el shortest path desde un origen $s$ a todos los vertices en $O((V + E) \log V)$ asumiendo \textbf{pesos no negativos}. La idea: usar una cola de prioridad (min-heap) ordenada por distancia tentativa. En cada paso, extraer el nodo con menor distancia; si ya fue procesado, saltar. Si al actualizar un vecino la distancia mejora, insertar/actualizar en la cola. La demostracion: cuando un nodo se extrae por primera vez, no hay forma de llegar a el con menor distancia (cualquier camino alternativo deberia pasar por nodos no procesados, con distancia $\ge$ la actual).

\paragraph{Propiedades clave (invariantes).}
\begin{itemize}\setlength\itemsep{0pt}
	\item Asume pesos $\ge 0$; con pesos negativos usar Bellman-Ford.
	\item Greedy correctness: la primera vez que se extrae un nodo, su distancia es la definitiva.
	\item Si se quiere reconstruir el camino, guardar el \texttt{parent[]}.
	\item Cada nodo se procesa a lo sumo $V$ veces (mas relajarions, pero solo $V$ definitivas).
\end{itemize}

\paragraph{Codigo clave.}
El bucle principal de Dijkstra:
\begin{verbatim}
priority_queue<pair<long long, int>, vector<...>, greater<...>> pq;
pq.push({0, source});
while (!pq.empty()) {
    auto [d, v] = pq.top(); pq.pop();
    if (d > dist[v]) continue;  // entrada vieja
    for (auto [w, u] : adj[v])
        if (dist[v] + w < dist[u]) {
            dist[u] = dist[v] + w;
            pq.push({dist[u], u});
        }
}
\end{verbatim}

\paragraph{Complejidad.}
\begin{itemize}\setlength\itemsep{0pt}
	\item $O((V + E) \log V)$ con heap binario, $O(V + E)$ con heap de Fibonacci.
	\item En la practica, $\sim 1.5 \cdot 10^{-7}$ segundos por operacion para $V = 10^5$.
\end{itemize}

\paragraph{Donde tocar para modificar.}
\begin{itemize}\setlength\itemsep{0pt}
	\item \textbf{Camino mas corto entre todos los pares}: si se necesita desde todos los origenes, usar Floyd o ejecutar Dijkstra $V$ veces.
	\item \textbf{Shortest path con aristas negativas limitadas}: Bellman-Ford.
	\item \textbf{Reconstruir camino}: aniadir \texttt{parent[v]} actualizado junto a \texttt{dist[v]}.
	\item \textbf{Dijkstra con potentials (Johnson)}: para aristas negativas si se puede reescalar.
	\item \textbf{Multi-source}: aniadir varios nodos fuentes iniciales al heap.
\end{itemize}

\paragraph{Trampas y errores frecuentes.}
\begin{itemize}\setlength\itemsep{0pt}
	\item Distancias no inicializadas a \texttt{INF}; el snippet usa \texttt{1e18}.
	\item El \texttt{if(dist[adjN] != 1e18)} para borrar de la cola es importante si usas \texttt{set} (no asi con \texttt{priority\_queue}).
	\item Con pesos grandes, \texttt{dist[v] + w} puede overflow; usar \texttt{LLONG\_MAX/2} como INF.
	\item No usar con pesos negativos: el algoritmo puede dar respuestas incorrectas.
\end{itemize}

\paragraph{Codigo.}
Ver \texttt{cpp.tex}, seccion \textit{Dijkstra}.

\vspace{0.5em}

\subsubsection{Bellman-Ford}

\paragraph{Idea intuitiva.}
Encuentra shortest path desde $s$ incluso con \textbf{pesos negativos} en $O(VE)$. La idea: relajar todas las aristas $V - 1$ veces. Cada relajacion propaga la mejor distancia a ``un paso mas''. Despues de $V - 1$ pasadas, si alguna arista aun se puede relajar, hay un \textbf{ciclo negativo} alcanzable. La demostracion: el camino mas corto tiene a lo sumo $V - 1$ aristas (sin ciclos); tras $V - 1$ relajaciones todas las distancias estan en su valor optimo.

\paragraph{Propiedades clave (invariantes).}
\begin{itemize}\setlength\itemsep{0pt}
	\item Funciona con pesos negativos (sin ciclos negativos).
	\item Tras $V - 1$ relajaciones, las distancias son definitivas (si no hay ciclo negativo).
	\item Una relajacion adicional detecta ciclos negativos: si \texttt{dist[u] + w < dist[v]}, hay ciclo negativo alcanzable desde $s$.
	\item Cada relajacion solo mejora distancias, asi que termina (no oscila).
\end{itemize}

\paragraph{Codigo clave.}
El doble bucle de relajacion:
\begin{verbatim}
for (int i = 0; i < n - 1; ++i)
    for (auto [u, v, w] : edges)
        if (dist[u] + w < dist[v])
            dist[v] = dist[u] + w;
// Detectar ciclo negativo
for (auto [u, v, w] : edges)
    if (dist[u] + w < dist[v]) { /* ciclo negativo */ }
\end{verbatim}

\paragraph{Complejidad.}
\begin{itemize}\setlength\itemsep{0pt}
	\item $O(VE)$.
	\item SPFA (con cola) en promedio es mucho mas rapido, peor caso igual.
\end{itemize}

\paragraph{Donde tocar para modificar.}
\begin{itemize}\setlength\itemsep{0pt}
	\item \textbf{SPFA}: optimizacion que relaja solo nodos cuya distancia cambio; en el peor caso sigue siendo $O(VE)$.
	\item \textbf{Detectar ciclo negativo global}: ejecutar Bellman-Ford desde un super-origen conectado a todos.
	\item \textbf{Camino mas corto con K aristas exactas}: $K$ iteraciones de Bellman-Ford.
	\item \textbf{Johnson's}: reescalar pesos con potentials para usar Dijkstra.
\end{itemize}

\paragraph{Trampas y errores frecuentes.}
\begin{itemize}\setlength\itemsep{0pt}
	\item Overflow si los pesos son $\ge 2^{52}$; usar \texttt{long long} y \texttt{\_\_int128} si hace falta.
	\item Si el grafo es denso ($E \approx V^2$), $O(VE)$ es prohibitivo.
	\item El chequeo de ciclo negativo debe hacerse tras las $V - 1$ iteraciones, no durante.
	\item Distancias no inicializadas a \texttt{INF}; el primer nodo (source) tiene distancia 0.
\end{itemize}

\paragraph{Codigo.}
Ver \texttt{cpp.tex}, seccion \textit{Bellman-Ford}.

\vspace{0.5em}

\subsubsection{Floyd-Warshall}

\paragraph{Idea intuitiva.}
Shortest path entre \textbf{todos los pares} de vertices en $O(V^3)$. DP: $d[i][j]$ = shortest path usando solo nodos intermedios en $\{0, \dots, k\}$. Transicion: $d_{k+1}[i][j] = \min(d_k[i][j], d_k[i][k+1] + d_k[k+1][j])$. La demostracion: tras $k$ iteraciones, $d[i][j]$ es el shortest path con intermedios en $\{0, \dots, k-1\}$; aniadir el nodo $k$ da la recurrencia clasica.

\paragraph{Propiedades clave (invariantes).}
\begin{itemize}\setlength\itemsep{0pt}
	\item Detecta ciclos negativos: tras el algoritmo, $d[i][i] < 0$ indica ciclo negativo alcanzable desde $i$.
	\item Funciona con pesos negativos (sin ciclos negativos).
	\item Memoria $O(V^2)$; con $V \le 400$ es OK, con mas conviene Dijkstra $V$ veces.
	\item Es la unica opcion cuando se quieren shortest paths entre todos los pares y el grafo es denso.
\end{itemize}

\paragraph{Codigo clave.}
El triple bucle clasico:
\begin{verbatim}
for (int k = 0; k < n; ++k)
    for (int i = 0; i < n; ++i)
        for (int j = 0; j < n; ++j)
            d[i][j] = min(d[i][j], d[i][k] + d[k][j]);
// detectar ciclo negativo
for (int i = 0; i < n; ++i) if (d[i][i] < 0) tiene_ciclo_negativo = true;
\end{verbatim}

\paragraph{Complejidad.}
\begin{itemize}\setlength\itemsep{0pt}
	\item $O(V^3)$ tiempo, $O(V^2)$ memoria.
	\item Para $V \le 400$: $\sim 6.4 \cdot 10^7$ operaciones (factible).
\end{itemize}

\paragraph{Donde tocar para modificar.}
\begin{itemize}\setlength\itemsep{0pt}
	\item \textbf{Reconstruir camino}: aniadir \texttt{next[i][j]} y backtrack.
	\item \textbf{Transitividad}: el algoritmo tambien calcula transitividad (\texttt{d[i][j] != INF} significa alcanzable).
	\item \textbf{Pesos doubles}: aniadir tolerancia (\texttt{d[i][j] > d[i][k] + d[k][j] - eps}).
	\item \textbf{Grafo no denso}: usar Dijkstra $V$ veces en vez de Floyd.
\end{itemize}

\paragraph{Trampas y errores frecuentes.}
\begin{itemize}\setlength\itemsep{0pt}
	\item Para detectar ciclo negativo correctamente, aniadir un chequeo \texttt{if (ans[i][i] < 0)} tras el triple bucle.
	\item El orden de los bucles \texttt{k, i, j} es importante; invertir da resultados diferentes.
	\item Con \texttt{double}, usar tolerancia \texttt{< eps} en vez de \texttt{<} para evitar loops infinitos.
\end{itemize}

\paragraph{Codigo.}
Ver \texttt{cpp.tex}, seccion \textit{Floyd-Warshall}.

\subsection{Componentes y cortes}

\subsubsection{Kosaraju SCC}

\paragraph{Idea intuitiva.}
Encuentra las \textbf{componentes fuertemente conexas} (SCCs) en $O(V + E)$. Algoritmo en 3 pasos:
\begin{itemize}\setlength\itemsep{0pt}
	\item DFS en $G$ rellenando una pila en \textbf{post-orden} (los nodos que terminan tarde quedan arriba).
	\item Construir $G^T$ (grafo transpuesto).
	\item DFS en $G^T$ procesando nodos en orden de la pila (de mayor a menor tiempo); cada DFS da una SCC.
\end{itemize}
La demostracion: dos nodos $u, v$ estan en la misma SCC sii hay camino $u \leadsto v$ y $v \leadsto u$; Kosaraju captura exactamente esto al explorar $G^T$ desde los nodos con mayor tiempo de finalizacion en $G$.

\paragraph{Propiedades clave (invariantes).}
\begin{itemize}\setlength\itemsep{0pt}
	\item Cada SCC es maximal fuertemente conexa.
	\item El \textbf{grafo de SCCs} (metagraf) es un DAG.
	\item Kosaraju y Tarjan dan el mismo resultado; Tarjan es ``mas rapido'' en practica (un solo DFS).
	\item Cada nodo pertenece a exactamente una SCC.
\end{itemize}

\paragraph{Codigo clave.}
La estructura de Kosaraju:
\begin{verbatim}
vector<int> order;
vector<bool> used(n);
function<void(int)> dfs1 = [&](int v) {
    used[v] = true;
    for (int u : g[v]) if (!used[u]) dfs1(u);
    order.push_back(v);  // post-orden
};
function<void(int, int)> dfs2 = [&](int v, int root) {
    comp[v] = root;
    for (int u : gt[v]) if (comp[u] == -1) dfs2(u, root);
};
// 1: DFS en G para llenar order
// 2: DFS en G^T en orden inverso para asignar SCCs
\end{verbatim}

\paragraph{Complejidad.}
\begin{itemize}\setlength\itemsep{0pt}
	\item $O(V + E)$.
	\item Constante 2 (dos DFS + el grafo transpuesto).
\end{itemize}

\paragraph{Donde tocar para modificar.}
\begin{itemize}\setlength\itemsep{0pt}
	\item \textbf{2-SAT}: SCC sobre grafo de implicaciones.
	\item \textbf{Construir el DAG de SCCs}: aniadir aristas entre SCCs diferentes.
	\item \textbf{Tamano de cada SCC}: aniadir contador durante el segundo DFS.
	\item \textbf{Grafo grande}: usar Tarjan para evitar construir $G^T$.
\end{itemize}

\paragraph{Trampas y errores frecuentes.}
\begin{itemize}\setlength\itemsep{0pt}
	\item El snippet usa \texttt{vector<...> adj[n]} (VLA); reemplazable por \texttt{vector<vector<pair<int,int>>} para portabilidad.
	\item El orden del segundo DFS debe ser \textbf{inverso} al de finalizacion del primero.
	\item Si el grafo no es conexo, aniadir un loop sobre todos los nodos en el primer DFS.
	\item Confundir SCC con componentes conexas (las primeras son para dirigidos, las segundas para no dirigidos).
\end{itemize}

\paragraph{Codigo.}
Ver \texttt{cpp.tex}, seccion \textit{Kosaraju SCC}.

\vspace{0.5em}

\subsubsection{Tarjan (puentes)}

\paragraph{Idea intuitiva.}
Encuentra los \textbf{puentes} (bridges): aristas cuya remocion \textbf{desconecta} el grafo. Usa un solo DFS manteniendo \texttt{tin[v]} (tiempo de descubrimiento) y \texttt{low[v]} (el menor \texttt{tin} alcanzable desde $v$ via back-edges). Una arista $v - u$ es puente si $low[u] > tin[v]$. La demostracion: si $low[u] > tin[v]$, entonces $u$ (y su subarbol) no pueden ``escapar'' de $v$; quitar la arista $v - u$ parte el grafo.

\paragraph{Propiedades clave (invariantes).}
\begin{itemize}\setlength\itemsep{0pt}
	\item \texttt{low[v] = min(tin[v], tin[w])} para todo back-edge $(v, w)$, y \texttt{min(low[u])} para hijos en el DFS.
	\item Solo se consideran las aristas del \textbf{arbol DFS} (no back-edges) para \texttt{low}.
	\item Cada nodo tiene exactamente un \texttt{tin}; cada back-edge se procesa una vez.
\end{itemize}

\paragraph{Codigo clave.}
El DFS de Tarjan para bridges:
\begin{verbatim}
void dfs(int v, int pEdge) {
    used[v] = true;
    tin[v] = low[v] = timer++;
    for (auto [u, id] : adj[v]) {
        if (id == pEdge) continue;  // no contar el edge al padre
        if (used[u]) {
            low[v] = min(low[v], tin[u]);  // back-edge
        } else {
            dfs(u, id);
            low[v] = min(low[v], low[u]);
            if (low[u] > tin[v]) bridges.push_back(id);  // es puente
        }
    }
}
\end{verbatim}

\paragraph{Complejidad.}
\begin{itemize}\setlength\itemsep{0pt}
	\item $O(V + E)$.
	\item Una sola pasada de DFS.
\end{itemize}

\paragraph{Donde tocar para modificar.}
\begin{itemize}\setlength\itemsep{0pt}
	\item \textbf{Bridge tree}: contractar cada 2-edge-connected component a un nodo; el resultado es un arbol.
	\item \textbf{Articulation points}: variante con condicion diferente (ver el siguiente modulo).
	\item \textbf{Grafo no dirigido}: solo funciona en no dirigidos; en dirigidos hay conceptos analogos (strong bridges, dominator tree).
	\item \textbf{Reconstruir componentes 2-edge-connected}: BFS desde cada nodo excluyendo bridges.
\end{itemize}

\paragraph{Trampas y errores frecuentes.}
\begin{itemize}\setlength\itemsep{0pt}
	\item No confundir el parent con un back-edge al padre; chequear \texttt{if (adjN == parent) continue;} y contar solo las veces que se ``visita'' como hijo.
	\item Para multigrafos, los edges deben tener IDs unicos para distinguirlos.
	\item Si el grafo no es conexo, aniadir loop externo sobre todos los nodos.
\end{itemize}

\paragraph{Codigo.}
Ver \texttt{cpp.tex}, seccion \textit{Tarjan (puentes)}.

\vspace{0.5em}

\subsubsection{Tarjan (puntos de articulacion)}

\paragraph{Idea intuitiva.}
Encuentra los \textbf{puntos de articulacion} (cut vertices): vertices cuya remocion \textbf{desconecta} el grafo. Mismo setup que para bridges: \texttt{tin[v]}, \texttt{low[v]}. Un nodo $v$ (no raiz) es de articulacion si existe un hijo $u$ con $low[u] \ge tin[v]$. La \textbf{raiz} es de articulacion si tiene $\ge 2$ hijos en el DFS tree. La razon: para $v$ no raiz, la condicion $\ge$ indica que el subarbol de $u$ depende de $v$; quitar $v$ los desconecta. Para la raiz, perder $\ge 2$ hijos parte el DFS en $\ge 2$ componentes.

\paragraph{Propiedades clave (invariantes).}
\begin{itemize}\setlength\itemsep{0pt}
	\item $low[u] \ge tin[v]$ significa que el subarbol de $u$ no tiene back-edge ``por encima'' de $v$.
	\item La raiz es caso especial: tiene que perder mas de un hijo para que el grafo se parta.
	\item Cada nodo se evalua una sola vez.
\end{itemize}

\paragraph{Codigo clave.}
La condicion de articulacion:
\begin{verbatim}
void dfs(int v, int p) {
    used[v] = true;
    tin[v] = low[v] = timer++;
    int children = 0;
    for (int u : adj[v]) {
        if (u == p) continue;
        if (used[u]) low[v] = min(low[v], tin[u]);
        else {
            dfs(u, v);
            low[v] = min(low[v], low[u]);
            if (low[u] >= tin[v] && p != -1)  // no raiz
                isArt[v] = true;
            ++children;
        }
    }
    if (p == -1 && children > 1) isArt[v] = true;  // raiz especial
}
\end{verbatim}

\paragraph{Complejidad.}
\begin{itemize}\setlength\itemsep{0pt}
	\item $O(V + E)$.
	\item Una sola pasada de DFS.
\end{itemize}

\paragraph{Donde tocar para modificar.}
\begin{itemize}\setlength\itemsep{0pt}
	\item \textbf{Block-cut tree}: estructura bipartita con BCCs (2-vertex-connected components) y articulation points.
	\item \textbf{Verificar biconectividad}: chequear que no hay puntos de articulacion.
	\item \textbf{Componentes biconnected}: extender Tarjan para enumerar BCCs.
\end{itemize}

\paragraph{Trampas y errores frecuentes.}
\begin{itemize}\setlength\itemsep{0pt}
	\item No olvidar el caso especial de la raiz.
	\item La condicion usa \texttt{>=}, no \texttt{>}; un hijo con \texttt{low[u] == tin[v]} tambien cuenta.
	\item Multigrafos pueden tener ``falsos'' articulation points; verificar caso a caso.
\end{itemize}

\paragraph{Codigo.}
Ver \texttt{cpp.tex}, seccion \textit{Tarjan (puntos de articulacion)}.

\vspace{0.5em}

\subsubsection{Tarjan SCC}

\paragraph{Idea intuitiva.}
Variante del SCC que usa \textbf{un solo DFS} con una pila explicita. Mantiene \texttt{disc[v]}, \texttt{low[v]}, y una pila de nodos ``en la SCC actual''. Cuando \texttt{low[v] == disc[v]}, $v$ es la raiz de una SCC; popear hasta $v$ inclusive. La demostracion: $low[v] == disc[v]$ indica que $v$ no tiene back-edges a ancestros; entonces $v$ y los nodos en la pila hasta $v$ forman una SCC maximal.

\paragraph{Propiedades clave (invariantes).}
\begin{itemize}\setlength\itemsep{0pt}
	\item Mas eficiente que Kosaraju (un solo DFS).
	\item \texttt{low[v] = min(low[u], disc[w])} para back-edges o recursion.
	\item La pila asegura que solo los nodos en la SCC actual se popean.
	\item Cada nodo se pushea y popea a lo sumo una vez.
\end{itemize}

\paragraph{Codigo clave.}
El DFS de Tarjan SCC:
\begin{verbatim}
void dfs(int v) {
    disc[v] = low[v] = timer++;
    st.push(v); inStack[v] = true;
    for (int u : g[v]) {
        if (disc[u] == -1) { dfs(u); low[v] = min(low[v], low[u]); }
        else if (inStack[u]) low[v] = min(low[v], disc[u]);
    }
    if (low[v] == disc[v]) {
        // popear hasta v
        while (true) {
            int u = st.top(); st.pop(); inStack[u] = false;
            comp[u] = current_scc;
            if (u == v) break;
        }
        ++current_scc;
    }
}
\end{verbatim}

\paragraph{Complejidad.}
\begin{itemize}\setlength\itemsep{0pt}
	\item $O(V + E)$.
	\item Constante similar a Kosaraju pero sin construir $G^T$.
\end{itemize}

\paragraph{Donde tocar para modificar.}
\begin{itemize}\setlength\itemsep{0pt}
	\item \textbf{Tamano de cada SCC}: aniadir contador durante el pop.
	\item \textbf{Grafo condensation}: aniadir aristas entre SCCs en el DAG resultante.
	\item \textbf{2-SAT}: aniadir la construccion del grafo de implicaciones.
\end{itemize}

\paragraph{Trampas y errores frecuentes.}
\begin{itemize}\setlength\itemsep{0pt}
	\item \texttt{inStack[u]} debe chequearse antes de usar \texttt{low[v] = min(low[v], disc[u])} en back-edges.
	\item La pila global es importante; sin ella no se puede extraer la SCC.
	\item Si el grafo es muy grande, recursion puede overflow; aniadir version iterativa.
\end{itemize}

\paragraph{Codigo.}
Ver \texttt{cpp.tex}, seccion \textit{Tarjan SCC}.

\subsection{Matching}

\subsubsection{Kuhn (bipartite matching)}

\paragraph{Idea intuitiva.}
Encuentra un \textbf{matching maximo} en un grafo bipartito en $O(VE)$. Para cada nodo $u$ en $L$, intentar encontrar un vecino $v$ en $R$ que este libre o cuyo match se pueda ``reorganizar'' recursivamente. Si se encuentra, matchear $u - v$. La idea: si $u$ quiere un vecino $v$ ya ocupado, intentamos ``mover'' el match de $v$ a otro nodo; recursivamente, esto encuentra un augmenting path si existe.

\paragraph{Propiedades clave (invariantes).}
\begin{itemize}\setlength\itemsep{0pt}
	\item \texttt{matchR[v]} = el nodo de $L$ matcheado a $v$ (o $-1$ si libre).
	\item \texttt{vis[]} se resetea \textbf{cada intento} de $u$.
	\item Si el matching es perfecto, $|L| = |R| = $ tamano del matching.
	\item Kuhn no garantiza maximalidad greedy; puede haber augmenting paths no encontrados.
\end{itemize}

\paragraph{Codigo clave.}
La recursion del augmenting path:
\begin{verbatim}
bool tryKuhn(int v) {
    if (vis[v]) return false;
    vis[v] = true;
    for (int u : g[v]) {
        if (matchR[u] == -1 || tryKuhn(matchR[u])) {
            matchR[u] = v;
            return true;
        }
    }
    return false;
}
// Por cada v en L:
fill(vis, vis+n, false);
if (tryKuhn(v)) ++matches;
\end{verbatim}

\paragraph{Complejidad.}
\begin{itemize}\setlength\itemsep{0pt}
	\item $O(VE)$ en el peor caso.
	\item Con Hopcroft-Karp: $O(\sqrt{V} \cdot E)$, mucho mejor para grafos densos.
\end{itemize}

\paragraph{Donde tocar para modificar.}
\begin{itemize}\setlength\itemsep{0pt}
	\item \textbf{Matching bipartito minimo (min vertex cover)}: Konig's theorem.
	\item \textbf{Matching general}: blossom algorithm (no incluido).
	\item \textbf{Aniadir capacidades}: si los nodos de $R$ tienen capacidad $> 1$, partir cada uno en capacidad copias.
	\item \textbf{Output del matching}: el array \texttt{matchR} da los pares finales.
\end{itemize}

\paragraph{Trampas y errores frecuentes.}
\begin{itemize}\setlength\itemsep{0pt}
	\item Olvidar resetear \texttt{vis} por intento; eso es lo que evita ``ciclos infinitos'' en la recursion.
	\item Para grafos grandes ($V > 10^4$), Kuhn puede ser muy lento; usar Hopcroft-Karp.
	\item Si el matching es maximo pero no perfecto, $|L|$ o $|R|$ puede ser mayor; el matching es solo maximo.
\end{itemize}

\paragraph{Codigo.}
Ver \texttt{cpp.tex}, seccion \textit{Kuhn (bipartite matching)}.

\vspace{0.5em}

\subsubsection{Hopcroft-Karp}

\paragraph{Idea intuitiva.}
Matching bipartito maximo en $O(\sqrt{V} E)$. Algoritmo en \textbf{fases}: en cada fase, hacer BFS para encontrar los \textbf{caminos de aumento mas cortos} desde nodos libres de $L$; luego DFS solo por esos caminos. Cada fase aumenta el matching en $\ge 1$ y toma $O(E)$. La demostracion: tras una fase, todos los augmenting paths tienen longitud $\ge$ la nueva minima; eso sucede cada $\sqrt{V}$ fases (longitud minima crece geometricamente).

\paragraph{Propiedades clave (invariantes).}
\begin{itemize}\setlength\itemsep{0pt}
	\item \texttt{dist[u]} = distancia en niveles desde un nodo libre de $L$.
	\item BFS termina cuando no hay mas camino a un nodo libre de $R$.
	\item DFS respeta los niveles (solo sigue \texttt{dist[matchV[v]] == dist[u] + 1}).
	\item Cada fase aumenta el matching en $\ge 1$ arista.
\end{itemize}

\paragraph{Codigo clave.}
El BFS de Hopcroft-Karp:
\begin{verbatim}
bool bfs() {
    queue<int> q;
    for (int v : L) {
        if (matchU[v] == -1) { dist[v] = 0; q.push(v); }
        else dist[v] = -1;
    }
    bool found = false;
    while (!q.empty()) {
        int v = q.front(); q.pop();
        for (int u : g[v]) {
            if (matchV[u] != -1 && dist[matchV[u]] == -1) {
                dist[matchV[u]] = dist[v] + 1;
                q.push(matchV[u]);
            }
            if (matchV[u] == -1) found = true;  // llegamos a libre
        }
    }
    return found;
}
\end{verbatim}

\paragraph{Complejidad.}
\begin{itemize}\setlength\itemsep{0pt}
	\item $O(\sqrt{V} E)$.
	\item Para $V = 10^4$, $E = 10^5$: $\sim 3 \cdot 10^7$ operaciones.
\end{itemize}

\paragraph{Donde tocar para modificar.}
\begin{itemize}\setlength\itemsep{0pt}
	\item \textbf{Reconstruir el matching}: ya viene en \texttt{matchU}/\texttt{matchV}.
	\item \textbf{Aniadir pesos}: Hungarian algorithm.
	\item \textbf{Densidad alta}: si $E \approx V^2$, sigue siendo util pero verificar constantes.
	\item \textbf{Aniadir capacidades}: partir los nodos con capacidad $> 1$ en varias copias.
\end{itemize}

\paragraph{Trampas y errores frecuentes.}
\begin{itemize}\setlength\itemsep{0pt}
	\item \texttt{dist[matchV[v]] == dist[u] + 1} (no \texttt{==} a secas); sin esto, el DFS no respeta los niveles y se pierde la complejidad.
	\item El BFS debe marcar los niveles correctamente; resetear \texttt{dist} para nodos alcanzables solo via matching.
	\item El DFS debe respetar los niveles; cualquier ``atajo'' invalida la cota.
\end{itemize}

\paragraph{Codigo.}
Ver \texttt{cpp.tex}, seccion \textit{Hopcroft-Karp}.

\subsection{Basicos}

\subsubsection{DFS / BFS / componentes}

\paragraph{Idea intuitiva.}
\textbf{DFS} (Depth-First Search) y \textbf{BFS} (Breadth-First Search) son los dos recorridos basicos de grafos. DFS usa una pila (o recursion) y explora tan profundo como puede antes de retroceder; BFS usa una cola y explora por ``niveles'' de distancia. La cantidad de \textbf{componentes conexas} se cuenta iniciando un BFS/DFS desde cada nodo no visitado. La diferencia clave: BFS da el camino mas corto en aristas; DFS da un orden de finalizacion util para problemas topologicos.

\paragraph{Propiedades clave (invariantes).}
\begin{itemize}\setlength\itemsep{0pt}
	\item BFS da la \textbf{distancia en aristas} mas corta desde el origen.
	\item DFS da un \textbf{orden topologico} (si se aniaden al final).
	\item Ambos son $O(V + E)$.
	\item Un grafo con $C$ componentes conexas satisface $\sum \text{subtree\_size}_v = V$ y $\sum E_v = 2E$.
	\item DFS en arbol da un spanning tree; las ``back edges'' corresponden a ciclos.
\end{itemize}

\paragraph{Codigo clave.}
BFS clasico:
\begin{verbatim}
queue<int> q;
q.push(source);
dist[source] = 0;
while (!q.empty()) {
    int v = q.front(); q.pop();
    for (int u : adj[v])
        if (dist[u] == -1) {
            dist[u] = dist[v] + 1;
            q.push(u);
        }
}
\end{verbatim}

\paragraph{Complejidad.}
\begin{itemize}\setlength\itemsep{0pt}
	\item $O(V + E)$ tiempo, $O(V)$ memoria.
	\item Cada arista se procesa exactamente 2 veces (ida y vuelta) en no dirigidos.
\end{itemize}

\paragraph{Donde tocar para modificar.}
\begin{itemize}\setlength\itemsep{0pt}
	\item \textbf{Aniadir pesos}: si las aristas tienen peso y queres shortest path, BFS ya no sirve; usar Dijkstra.
	\item \textbf{Grafo dirigido}: en DFS, aniadir check \texttt{if (state[u] == 1)} para detectar ciclos.
	\item \textbf{Bipartite check}: BFS alternando colores (ver el siguiente modulo).
	\item \textbf{Componentes fuertemente conexas (dirigidos)}: usar Tarjan o Kosaraju.
	\item \textbf{BFS en matriz}: aniadir offsets por direccion para grid 2D.
\end{itemize}

\paragraph{Trampas y errores frecuentes.}
\begin{itemize}\setlength\itemsep{0pt}
	\item Stack overflow con DFS recursivo para $V > 10^5$; usar version iterativa con pila explicita.
	\item En BFS, olvidar marcar \texttt{dist[u] = dist[v] + 1} antes de pushear da visitas duplicadas.
	\item Para grafos dirigidos, ``componente conexa'' no esta definida; usar SCC.
\end{itemize}

\paragraph{Codigo.}
Ver \texttt{cpp.tex}, seccion \textit{DFS / BFS / componentes}.

\vspace{0.5em}

\subsubsection{Topological sort (Kahn + DFS)}

\paragraph{Idea intuitiva.}
Un \textbf{orden topologico} de un DAG es una permutacion de los vertices tal que toda arista $u \to v$ tiene $u$ antes que $v$. Dos algoritmos clasicos:
\begin{itemize}\setlength\itemsep{0pt}
	\item \textbf{Kahn}: BFS sobre nodos con \texttt{indeg == 0}; al ``consumir'' un nodo, decrementar el indeg de sus vecinos; los que lleguen a 0 se encolan.
	\item \textbf{DFS}: hacer DFS; al terminar un nodo (estado = 2), aniadirlo al final; el orden es el reverso.
\end{itemize}
La demostracion: en un DAG siempre hay al menos un nodo con indeg 0; quitarlo mantiene el DAG; por induccion, se obtiene el orden.

\paragraph{Propiedades clave (invariantes).}
\begin{itemize}\setlength\itemsep{0pt}
	\item Existe sii el grafo es un DAG.
	\item Si al final \texttt{order.size() < V}, hay un ciclo.
	\item Kahn da el orden ``natural'' (los que pueden ir primero); DFS da el orden ``al reves'' (los que pueden ir al final primero).
	\item La cantidad de ordenes topologicos puede ser exponencial.
\end{itemize}

\paragraph{Codigo clave.}
Kahn con priority queue (orden lexicografico):
\begin{verbatim}
priority_queue<int, vector<int>, greater<int>> pq;
for (int v = 0; v < n; ++v) if (indeg[v] == 0) pq.push(v);
while (!pq.empty()) {
    int v = pq.top(); pq.pop();
    order.push_back(v);
    for (int u : adj[v]) {
        if (--indeg[u] == 0) pq.push(u);
    }
}
\end{verbatim}

\paragraph{Complejidad.}
\begin{itemize}\setlength\itemsep{0pt}
	\item $O(V + E)$ ambos.
	\item Con priority queue, $O((V + E) \log V)$.
\end{itemize}

\paragraph{Donde tocar para modificar.}
\begin{itemize}\setlength\itemsep{0pt}
	\item \textbf{Detectar ciclos}: comparar \texttt{order.size()} con \texttt{V}.
	\item \textbf{Orden lexicografico minimo}: usar priority queue en Kahn.
	\item \textbf{Topological sort con pesos}: combinar con shortest path en DAG (un solo pase en orden topologico).
	\item \textbf{Encontrar todos los ordenes}: DP sobre subsets (solo para $V$ pequeno).
\end{itemize}

\paragraph{Trampas y errores frecuentes.}
\begin{itemize}\setlength\itemsep{0pt}
	\item En el snippet, Kahn retorna orden vacio si hay ciclo (util para detectar).
	\item En DFS, no olvidar \texttt{reverse(order)}.
	\item Confundir indeg con outdeg; el orden Kahn quita primero los nodos ``fuente''.
	\item Para grafos no DAG, el algoritmo termina con \texttt{order.size() < V}.
\end{itemize}

\paragraph{Codigo.}
Ver \texttt{cpp.tex}, seccion \textit{Topological sort (Kahn + DFS)}.

\vspace{0.5em}

\subsubsection{Bipartite check + 2-coloring}

\paragraph{Idea intuitiva.}
Un grafo es \textbf{bipartito} si sus vertices se pueden dividir en dos conjuntos $L$ y $R$ tales que toda arista va de $L$ a $R$. Esto equivale a decir que se puede \textbf{2-colorear} sin que dos vertices adyacentes tengan el mismo color. BFS alternando colores detecta esto: si en algun momento dos adyacentes tienen el mismo color, no es bipartito. La demostracion: por cada componente, fijar el color de un nodo determina todos los demas; si en algun paso hay conflicto, hay un ciclo impar (no bipartito).

\paragraph{Propiedades clave (invariantes).}
\begin{itemize}\setlength\itemsep{0pt}
	\item Grafo bipartito $\Leftrightarrow$ sin ciclos impares.
	\item Conexo: si un componente no es bipartito, todo el grafo no lo es.
	\item Si hay \textbf{multiple componentes}, cada uno se colorea independientemente.
	\item La 2-coloracion es unica salvo swap de colores.
\end{itemize}

\paragraph{Codigo clave.}
BFS con 2-coloreo:
\begin{verbatim}
vector<int> color(n, -1);
for (int v = 0; v < n; ++v) if (color[v] == -1) {
    color[v] = 0;
    queue<int> q; q.push(v);
    while (!q.empty()) {
        int u = q.front(); q.pop();
        for (int w : adj[u]) {
            if (color[w] == -1) { color[w] = 1 - color[u]; q.push(w); }
            else if (color[w] == color[u]) return false;
        }
    }
}
return true;
\end{verbatim}

\paragraph{Complejidad.}
\begin{itemize}\setlength\itemsep{0pt}
	\item $O(V + E)$.
	\item Memoria $O(V)$.
\end{itemize}

\paragraph{Donde tocar para modificar.}
\begin{itemize}\setlength\itemsep{0pt}
	\item \textbf{Colorear desde multiples fuentes}: BFS desde cada nodo no coloreado.
	\item \textbf{Bipartito + matching}: Hopcroft-Karp.
	\item \textbf{Grafo bipartito ponderado}: Hungarian algorithm.
	\item \textbf{Verificar bipartito en subgrafo}: solo aplicar BFS a los nodos relevantes.
\end{itemize}

\paragraph{Trampas y errores frecuentes.}
\begin{itemize}\setlength\itemsep{0pt}
	\item Verificar \textbf{todos} los componentes, no solo el primero.
	\item Para grafos con auto-loops, un auto-loop hace al grafo no bipartito.
	\item En multigrafos, las aristas multiples no afectan el chequeo de bipartito.
\end{itemize}

\paragraph{Codigo.}
Ver \texttt{cpp.tex}, seccion \textit{Bipartite check + 2-coloring}.

\vspace{0.5em}

\subsubsection{0-1 BFS}

\paragraph{Idea intuitiva.}
Variante de BFS para grafos con \textbf{pesos 0 o 1}. Usa un \texttt{deque}: aristas de peso 0 se encolan al frente (\texttt{push\_front}), aristas de peso 1 al final (\texttt{push\_back}). Asi, la primera vez que se extrae un nodo, su distancia es la definitiva, igual que en Dijkstra pero en $O(V + E)$. La intuicion: el deque mantiene los nodos ordenados por distancia; las aristas de peso 0 no cambian la distancia (entran al frente), las de peso 1 la incrementan (entran al final).

\paragraph{Propiedades clave (invariantes).}
\begin{itemize}\setlength\itemsep{0pt}
	\item Asume pesos \textbf{solo} 0 o 1.
	\item Mejor caso que Dijkstra (sin factor $\log V$).
	\item No funciona con pesos $\ge 2$.
	\item La cola mantiene los nodos ordenados por distancia no decreciente.
\end{itemize}

\paragraph{Codigo clave.}
El deque doble:
\begin{verbatim}
deque<int> q;
q.push_front(source);
dist[source] = 0;
while (!q.empty()) {
    int v = q.front(); q.pop_front();
    for (auto [w, u] : adj[v]) {
        if (dist[v] + w < dist[u]) {
            dist[u] = dist[v] + w;
            if (w == 0) q.push_front(u);
            else q.push_back(u);
        }
    }
}
\end{verbatim}

\paragraph{Complejidad.}
\begin{itemize}\setlength\itemsep{0pt}
	\item $O(V + E)$.
	\item Cada nodo se inserta a lo sumo $O(1)$ veces (con check de distancia).
\end{itemize}

\paragraph{Donde tocar para modificar.}
\begin{itemize}\setlength\itemsep{0pt}
	\item \textbf{Pesos en $[0, k]$ small}: usar small k-ary BFS o Dijkstra.
	\item \textbf{Reconstruir camino}: aniadir \texttt{parent[]}.
	\item \textbf{Pesos negativos}: no funciona, usar Bellman-Ford.
\end{itemize}

\paragraph{Trampas y errores frecuentes.}
\begin{itemize}\setlength\itemsep{0pt}
	\item Si los pesos son otro rango, el algoritmo da respuestas incorrectas. Verificar que las aristas son solo 0/1.
	\item El chequeo \texttt{dist[v] + w < dist[u]} es crucial; sino, el algoritmo procesa nodos multiples veces.
	\item Si una arista tiene peso $\ge 2$, aniadir verificaciones o cambiar a Dijkstra.
\end{itemize}

\paragraph{Codigo.}
Ver \texttt{cpp.tex}, seccion \textit{0-1 BFS}.

\subsection{MST}

\subsubsection{Kruskal}

\paragraph{Idea intuitiva.}
Encuentra el \textbf{arbol de expansion minima} (MST) en $O(E \log E)$. Ordena las aristas por peso ascendente y las aniade al MST si no forman ciclo (chequear con DSU). El arbol tiene exactamente $V - 1$ aristas. La demostracion (``Cut Property''): para cualquier corte, la arista minima cruzando el corte pertenece a \emph{algun} MST; por lo tanto, elegir la minima disponible que no forma ciclo es optimo.

\paragraph{Propiedades clave (invariantes).}
\begin{itemize}\setlength\itemsep{0pt}
	\item Algoritmo \textbf{greedy} correcto (Cut Property).
	\item Funciona en grafos \textbf{no conexos}; produce un \textbf{minimum spanning forest}.
	\item Si las aristas tienen pesos unicos, el MST es unico.
	\item El MST tiene exactamente $V - 1$ aristas (o menos si el grafo es disconexo).
\end{itemize}

\paragraph{Codigo clave.}
Kruskal con DSU:
\begin{verbatim}
sort(edges.begin(), edges.end());
DSU dsu(n);
long long mst = 0;
int edgesUsed = 0;
for (auto [w, u, v] : edges)
    if (dsu.unite(u, v)) {
        mst += w;
        if (++edgesUsed == n - 1) break;
    }
\end{verbatim}

\paragraph{Complejidad.}
\begin{itemize}\setlength\itemsep{0pt}
	\item $O(E \log E)$ (ordenamiento) o $O(E \log V)$ con DSU optimizations.
	\item DSU es casi $O(1)$ amortizado, asi que el sort domina.
\end{itemize}

\paragraph{Donde tocar para modificar.}
\begin{itemize}\setlength\itemsep{0pt}
	\item \textbf{Second best MST}: enumerar aristas del MST, sacar una, aniadir una no-MST; recalcular.
	\item \textbf{Kruskal randomizado}: si el grafo es casi-completo, agregar shuffle.
	\item \textbf{Aniadir info al DSU}: guardar el peso maximo en el camino, etc.
	\item \textbf{DSU con rollback}: si necesitas Kruskal reversible (para queries offline).
\end{itemize}

\paragraph{Trampas y errores frecuentes.}
\begin{itemize}\setlength\itemsep{0pt}
	\item El DSU debe estar limpio para cada invocacion.
	\item Si las aristas tienen pesos negativos, Kruskal sigue funcionando.
	\item Si el grafo es disconexo, el ``MST'' es un bosque con varias componentes.
	\item El check \texttt{if (++edgesUsed == n - 1) break;} optimiza la salida temprana.
\end{itemize}

\paragraph{Codigo.}
Ver \texttt{cpp.tex}, seccion \textit{Kruskal}.

\vspace{0.5em}

\subsubsection{Prim (priority queue)}

\paragraph{Idea intuitiva.}
Otra forma de calcular el MST: empezar desde un nodo y, en cada paso, aniadir la arista de menor peso que conecta el conjunto construido con el resto. Usando priority queue, es $O((V + E) \log V)$. La demostracion: cada vez que se aniade una arista, es la minima cruzando el ``corte'' entre el conjunto actual y el resto (Cut Property); eso preserva optimalidad.

\paragraph{Propiedades clave (invariantes).}
\begin{itemize}\setlength\itemsep{0pt}
	\item Similar a Dijkstra pero sin acumular distancias, solo eligiendo minimo por arista.
	\item Funciona mejor que Kruskal en \textbf{grafos densos} ($E \approx V^2$).
	\item Si el grafo es disconexo, genera el bosque (solo llega a los alcanzables).
	\item Cada nodo se inserta al heap una vez; cada arista se procesa una vez.
\end{itemize}

\paragraph{Codigo clave.}
Prim clasico con heap:
\begin{verbatim}
priority_queue<pair<int, int>, vector<...>, greater<...>> pq;
vector<bool> used(n, false);
vector<int> minEdge(n, INF);
minEdge[0] = 0;
pq.push({0, 0});
while (!pq.empty()) {
    auto [w, v] = pq.top(); pq.pop();
    if (used[v]) continue;
    used[v] = true;
    mst += w;
    for (auto [peso, u] : adj[v])
        if (!used[u] && peso < minEdge[u]) {
            minEdge[u] = peso;
            pq.push({peso, u});
        }
}
\end{verbatim}

\paragraph{Complejidad.}
\begin{itemize}\setlength\itemsep{0pt}
	\item $O((V + E) \log V)$.
	\item Con matriz de adyacencia: $O(V^2)$ sin heap.
\end{itemize}

\paragraph{Donde tocar para modificar.}
\begin{itemize}\setlength\itemsep{0pt}
	\item \textbf{Prim con matriz de adyacencia}: $O(V^2)$ sin heap, util en grafos densos pequenos.
	\item \textbf{Reconstruir el MST}: aniadir \texttt{parent[]}.
	\item \textbf{Prim en streaming}: variante para grafos que llegan incrementalmente.
\end{itemize}

\paragraph{Trampas y errores frecuentes.}
\begin{itemize}\setlength\itemsep{0pt}
	\item El \texttt{if (used[v]) continue;} es crucial: si el nodo ya fue procesado, saltar.
	\item En la version matriz, elije el minimo por linea (no heap).
	\item Si el grafo es disconexo, aniadir un loop externo sobre todos los nodos no usados.
\end{itemize}

\paragraph{Codigo.}
Ver \texttt{cpp.tex}, seccion \textit{Prim (priority queue)}.

\subsection{Flujos}

\subsubsection{Edmonds-Karp (max flow)}

\paragraph{Idea intuitiva.}
Implementacion BFS del algoritmo de \textbf{Ford-Fulkerson}. En cada iteracion, hacer BFS desde $s$ hasta $t$ siguiendo aristas con capacidad residual $> 0$; si se llega a $t$, se encontro un \textbf{camino de aumento}. Enviar el minimo de las capacidades en el camino; actualizar capacidades residuales. La demostracion de $O(VE^2)$: cada BFS encuentra el augmenting path mas corto, y la longitud minima crece monotona; como hay $\le E$ longitudes posibles, hay $\le E$ fases.

\paragraph{Propiedades clave (invariantes).}
\begin{itemize}\setlength\itemsep{0pt}
	\item BFS garantiza caminos de aumento mas cortos.
	\item Capacidad residual: forward edge baja, backward edge sube.
	\item Complejidad $O(VE^2)$.
	\item El flujo maximo respeta la conservacion en nodos (in = out salvo $s, t$).
\end{itemize}

\paragraph{Codigo clave.}
BFS + augmenting path:
\begin{verbatim}
while (bfs(s, t, parent)) {  // bfs encuentra camino o retorna false
    int pathFlow = INF;
    for (int v = t; v != s; v = parent[v])
        pathFlow = min(pathFlow, capacity[parent[v]][v]);
    for (int v = t; v != s; v = parent[v]) {
        capacity[parent[v]][v] -= pathFlow;
        capacity[v][parent[v]] += pathFlow;
    }
    maxFlow += pathFlow;
}
\end{verbatim}

\paragraph{Complejidad.}
\begin{itemize}\setlength\itemsep{0pt}
	\item $O(VE^2)$.
	\item En la practica, mas lento que Dinic para grafos grandes.
\end{itemize}

\paragraph{Donde tocar para modificar.}
\begin{itemize}\setlength\itemsep{0pt}
	\item \textbf{Grafo bipartito}: $O(\sqrt{V} E)$ (equivalente a Hopcroft-Karp).
	\item \textbf{Dinic} es preferible en la mayoria de los casos.
	\item \textbf{Capacidades doubles}: usar tolerancia para comparacion.
	\item \textbf{Min-cost flow}: aniadir costo por arista y modificar el algoritmo.
\end{itemize}

\paragraph{Trampas y errores frecuentes.}
\begin{itemize}\setlength\itemsep{0pt}
	\item Las capacidades deben ser $> 0$ para entrar a la cola (no $>= 0$); sino, ciclos infinitos.
	\item Para grafos bipartitos, Dinic o Hopcroft-Karp son mejores.
	\item Si el flujo maximo es muy grande, los \texttt{int} pueden overflow; usar \texttt{long long}.
\end{itemize}

\paragraph{Codigo.}
Ver \texttt{cpp.tex}, seccion \textit{Edmonds-Karp (max flow)}.

\vspace{0.5em}

\subsubsection{Dinic (max flow)}

\paragraph{Idea intuitiva.}
Algoritmo de flujo mas eficiente que Edmonds-Karp: combina \textbf{BFS para level graph} con \textbf{DFS bloqueante} (multiple augmenting paths en una sola pasada). Cada fase cuesta $O(E)$ y hay $O(V)$ fases en el peor caso; para bipartitos, $O(\sqrt{V} E)$. La razon: cada fase ``satura'' al menos una arista, asi que en $O(E)$ fases el flujo maximo se alcanza; con $O(V)$ niveles en el BFS, hay $O(VE)$ operaciones totales.

\paragraph{Propiedades clave (invariantes).}
\begin{itemize}\setlength\itemsep{0pt}
	\item \texttt{level[v]} = distancia desde $s$ en el level graph (solo aristas con cap $> 0$).
	\item \texttt{it[v]} = iterador actual para evitar reprocesar aristas.
	\item DFS solo sigue aristas con \texttt{level[e.to] == level[v] + 1}.
	\item Back-edges tienen \texttt{rev} apuntando a la forward edge original.
\end{itemize}

\paragraph{Codigo clave.}
El DFS bloqueante de Dinic:
\begin{verbatim}
long long dfs(int v, int t, long long f) {
    if (v == t) return f;
    for (int& i = it[v]; i < adj[v].size(); ++i) {
        Edge& e = adj[v][i];
        if (e.cap > 0 && level[e.to] == level[v] + 1) {
            long long ret = dfs(e.to, t, min(f, e.cap));
            if (ret > 0) {
                e.cap -= ret;
                adj[e.to][e.rev].cap += ret;
                return ret;
            }
        }
    }
    return 0;
}
\end{verbatim}

\paragraph{Complejidad.}
\begin{itemize}\setlength\itemsep{0pt}
	\item $O(V^2 E)$ general, $O(\sqrt{V} E)$ para bipartitos.
	\item En la practica, mucho mas rapido que Edmonds-Karp.
\end{itemize}

\paragraph{Donde tocar para modificar.}
\begin{itemize}\setlength\itemsep{0pt}
	\item \textbf{Dinic con scaling}: aniadir factor de escala para capacidades grandes.
	\item \textbf{Lower bounds}: aniadir un super-source/sink para forzar flujo minimo.
	\item \textbf{Matching bipartito}: el grafo se construye como source -> L -> R -> sink, todos con capacidad 1.
	\item \textbf{Push-relabel}: alternativa para grafos muy densos.
\end{itemize}

\paragraph{Trampas y errores frecuentes.}
\begin{itemize}\setlength\itemsep{0pt}
	\item \texttt{it[v]} debe \textbf{incrementarse} en cada iteracion; sino, el DFS se queda atascado en la misma arista.
	\item Si el BFS no encuentra $t$, el flujo es maximo; terminar.
	\item Las back-edges son cruciales para ``cancelar'' elecciones previas.
\end{itemize}

\paragraph{Codigo.}
Ver \texttt{cpp.tex}, seccion \textit{Dinic (max flow)}.

\vspace{0.5em}

\subsubsection{Min-Cost Max-Flow}

\paragraph{Idea intuitiva.}
Encuentra el flujo maximo con \textbf{costo minimo} (asumiendo costos no negativos). Usa \textbf{SPFA} (o Dijkstra con potentials) para encontrar el shortest path en el \textbf{costo} desde $s$ hasta $t$ considerando capacidades residuales. Cada augmenting path aporta su costo $\times$ flujo. La idea: los potentials reescalan las aristas para que todas tengan costo no negativo, permitiendo usar Dijkstra.

\paragraph{Propiedades clave (invariantes).}
\begin{itemize}\setlength\itemsep{0pt}
	\item \texttt{pot[]} = potentials para evitar aristas negativas (Johnson's trick).
	\item \texttt{dist[v]} = costo minimo desde $s$ con potentials.
	\item Solo se aumentan aristas con \texttt{cap > 0} (forward) y costo es \texttt{+cost} (forward) o \texttt{-cost} (backward).
	\item Tras cada augmenting, los potentials se actualizan para mantener la consistencia.
\end{itemize}

\paragraph{Codigo clave.}
El ciclo de min-cost flow:
\begin{verbatim}
while (spfa(s, t, dist, parent)) {
    int pathFlow = INF;
    for (int v = t; v != s; v = parent[v])
        pathFlow = min(pathFlow, capacity[parent[v]][v]);
    for (int v = t; v != s; v = parent[v]) {
        capacity[parent[v]][v] -= pathFlow;
        capacity[v][parent[v]] += pathFlow;
        cost += pathFlow * edge_cost[parent[v]][v];
    }
}
\end{verbatim}

\paragraph{Complejidad.}
\begin{itemize}\setlength\itemsep{0pt}
	\item $O(F \cdot V E)$ con SPFA, $O(F \cdot E \log V)$ con Dijkstra + potentials.
	\item $F$ = flujo maximo total.
\end{itemize}

\paragraph{Donde tocar para modificar.}
\begin{itemize}\setlength\itemsep{0pt}
	\item \textbf{Dijkstra + potentials}: si los costos son enteros no negativos, esto es mas estable.
	\item \textbf{Costos negativos en aristas}: requieren Bellman-Ford inicial o reorden.
	\item \textbf{Flujo minimo con costo}: encontrar el minimo flujo posible y su costo.
	\item \textbf{Costo por unidad}: si los costos son muy grandes, aniadir verificacion.
\end{itemize}

\paragraph{Trampas y errores frecuentes.}
\begin{itemize}\setlength\itemsep{0pt}
	\item El \texttt{pot[]} se actualiza \textbf{al final} de cada iteracion (no durante).
	\item Las back-edges tienen \textbf{costo negativo}, esencial para corregir elecciones previas.
	\item Si los costos son muy grandes, considerar modular arithmetic.
	\item Para problemas con flujo exacto, aniadir verificacion post-algoritmo.
\end{itemize}

\paragraph{Codigo.}
Ver \texttt{cpp.tex}, seccion \textit{Min-Cost Max-Flow}.

\subsection{Especales}

\subsubsection{2-SAT (implication graph + SCC)}

\paragraph{Idea intuitiva.}
Resuelve formulas booleanas en \textbf{CNF} con clausulas de \textbf{2 literales}: $(l_1 \lor l_2)$. Cada variable $x$ se representa con dos nodos: $x$ (verdadero) y $\neg x$ (falso). Cada clausula $l_1 \lor l_2$ se descompone en dos implicaciones: $\neg l_1 \to l_2$ y $\neg l_2 \to l_1$. Si en el grafo de implicaciones $\neg x$ y $x$ estan en la misma SCC, la formula es \textbf{UNSAT}. La demostracion: si $\neg x \to x$ en el grafo, entonces $\neg x$ implica $x$, asi que $\neg x$ debe ser falso y $x$ verdadero; pero entonces $\neg x$ es falso implica una contradiccion.

\paragraph{Propiedades clave (invariantes).}
\begin{itemize}\setlength\itemsep{0pt}
	\item Cada literal es un nodo; total $2n$ nodos.
	\item Kosaraju/Tarjan dan las SCCs.
	\item Si SAT: asignar $x = $ (topological order del condensation).
	\item Clausulas \texttt{setTrue(x)} se modelan como \texttt{implies(!x, x)}.
	\item La formula es SAT sii no hay variable $x$ con $x$ y $\neg x$ en la misma SCC.
\end{itemize}

\paragraph{Codigo clave.}
Conversion de clausulas a implicaciones:
\begin{verbatim}
// Clausula (a OR b) -> implica ~a -> b Y ~b -> a
addImplication(not(a), b);
addImplication(not(b), a);
// Forzar x: implica ~x -> x
if (force) addImplication(not(x), x);
\end{verbatim}

\paragraph{Complejidad.}
\begin{itemize}\setlength\itemsep{0pt}
	\item $O(V + E) = O(N + M)$ con Kosaraju o Tarjan.
	\item Memoria $O(V + E)$.
\end{itemize}

\paragraph{Donde tocar para modificar.}
\begin{itemize}\setlength\itemsep{0pt}
	\item \textbf{Forzar variables}: aniadir \texttt{setTrue} / \texttt{setFalse}.
	\item \textbf{3-SAT}: NP-completo; no usar este algoritmo.
	\item \textbf{Contar soluciones}: las SCCs condensation dan una formula para contar; multiplicar las formas de asignar variables en cada componente no-trivial.
	\item \textbf{Output de la asignacion}: aniadir el vector \texttt{assignment[]}.
\end{itemize}

\paragraph{Trampas y errores frecuentes.}
\begin{itemize}\setlength\itemsep{0pt}
	\item La asignacion usa el \textbf{orden topologico inverso} del condensation: \texttt{val[i] = comp[2i] < comp[2i+1]}.
	\item El ID de $\neg x$ es \texttt{2*x + 1} (o similar); verificar consistencia.
	\item Si hay clausulas con literales repetidos (e.g., $(x \lor x)$), tratar como $(x)$ simple.
\end{itemize}

\paragraph{Codigo.}
Ver \texttt{cpp.tex}, seccion \textit{2-SAT (implication graph + SCC)}.

\vspace{0.5em}

\subsubsection{Euler tour / Hierholzer (camino y circuito)}

\paragraph{Idea intuitiva.}
Encuentra un \textbf{camino o circuito euleriano} (que usa cada arista exactamente una vez). \textbf{Condiciones}:
\begin{itemize}\setlength\itemsep{0pt}
	\item \textbf{No dirigido}: todos los grados pares (circuito) o exactamente 2 impares (camino, empieza en uno, termina en el otro). Ademas, conexo sobre el subgrafo con aristas.
	\item \textbf{Dirigido}: $\text{indeg}(v) = \text{outdeg}(v)$ para todos (circuito) o exactamente un nodo con $\text{outdeg} = \text{indeg} + 1$ y otro con $\text{indeg} = \text{outdeg} + 1$ (camino).
\end{itemize}

\paragraph{Hierholzer}: elegir un ciclo desde el inicio, mientras haya ciclos ``laterales'' desde un nodo del camino, mergearlos. Implementacion con pila: avanzar por aristas, al no poder mas aniadir a la respuesta y backtrack. La demostracion de $O(V + E)$: cada arista se procesa una vez (entra y sale del stack una vez).

\paragraph{Propiedades clave (invariantes).}
\begin{itemize}\setlength\itemsep{0pt}
	\item Multigrafos permitidos; las aristas se marcan como usadas.
	\item No dirigido: cada arista aparece \textbf{dos veces} en \texttt{adj} (ida y vuelta).
	\item La construccion del path con ids consistentes es importante.
	\item El resultado es unico salvo orden de las aristas en multigrafos.
\end{itemize}

\paragraph{Codigo clave.}
Hierholzer con pila:
\begin{verbatim}
vector<int> path;
stack<int> st;
st.push(start);
while (!st.empty()) {
    int v = st.top();
    while (!adj[v].empty() && used[adj[v].back()]) adj[v].pop_back();
    if (adj[v].empty()) {
        path.push_back(v);
        st.pop();
    } else {
        int e = adj[v].back();
        used[e] = true;
        int u = edges[e].to(v);
        adj[v].pop_back();
        st.push(u);
    }
}
reverse(path.begin(), path.end());
\end{verbatim}

\paragraph{Complejidad.}
\begin{itemize}\setlength\itemsep{0pt}
	\item $O(V + E)$.
	\item Memoria $O(V + E)$.
\end{itemize}

\paragraph{Donde tocar para modificar.}
\begin{itemize}\setlength\itemsep{0pt}
	\item \textbf{Reconstruir con multigrafo}: aniadir ids consistentes antes del algoritmo.
	\item \textbf{Aplicar a ``de Bruijn'' o problemas de strings}: modelar como multigrafo y buscar Euler.
	\item \textbf{Dirigido vs no dirigido}: el snippet tiene ambas versiones.
	\item \textbf{Letter addition}: modelar palabras como aristas y buscar el superstring.
\end{itemize}

\paragraph{Trampas y errores frecuentes.}
\begin{itemize}\setlength\itemsep{0pt}
	\item Si el grafo no es conexo (sobre las aristas), no existe euleriano. Verificar primero.
	\item Para multigrafos, los IDs de aristas son cruciales para distinguirlas.
	\item Si el problema es ``encuentra todos los euler tours'', el algoritmo es exponencial.
\end{itemize}

\paragraph{Codigo.}
Ver \texttt{cpp.tex}, seccion \textit{Euler tour / Hierholzer (camino y circuito)}.

\vspace{0.5em}

\subsubsection{Hamiltonian path / cycle (bitmask DP + backtracking)}

\paragraph{Idea intuitiva.}
Un \textbf{camino hamiltoniano} visita cada vertice \textbf{exactamente una vez}. \textbf{Determinar} si existe es NP-completo en general, pero con $V \le 20$ se puede hacer DP con bitmask en $O(V^2 \cdot 2^V)$. \textbf{Teoremas suficientes}: \textbf{Dirac} (grado $\ge n/2$) y \textbf{Ore} ($\text{deg}(u) + \text{deg}(v) \ge n$ para todo par no adyacente). La idea de la DP: cada estado es (conjunto de visitados, ultimo vertice); la transicion es agregar un vecino no visitado.

\paragraph{Propiedades clave (invariantes).}
\begin{itemize}\setlength\itemsep{0pt}
	\item DP: $dp[mask][v]$ = true si hay un camino que visita exactamente \texttt{mask} y termina en $v$.
	\item TSP: variante con pesos minimos (Hamilton cycle minimo).
	\item Reconstruccion: guardar \texttt{par[mask][v]} (predecesor).
	\item $2^V$ estados con $V$ bits cada uno; factible para $V \le 20$.
\end{itemize}

\paragraph{Codigo clave.}
DP con bitmask:
\begin{verbatim}
dp[1][0] = 0;  // empezar en 0
for (int mask = 1; mask < (1<<n); ++mask)
    for (int v = 0; v < n; ++v) if (dp[mask][v] < INF)
        for (int u = 0; u < n; ++u)
            if (!(mask & (1<<u)))
                dp[mask | (1<<u)][u] = min(dp[mask | (1<<u)][u],
                                            dp[mask][v] + w[v][u]);
// TSP: respuesta = min(dp[(1<<n)-1][v] + w[v][0])
\end{verbatim}

\paragraph{Complejidad.}
\begin{itemize}\setlength\itemsep{0pt}
	\item $O(V^2 \cdot 2^V)$ tiempo, $O(V \cdot 2^V)$ memoria.
	\item Para $V = 20$: $\sim 4 \cdot 10^8$ operaciones (lento pero factible).
\end{itemize}

\paragraph{Donde tocar para modificar.}
\begin{itemize}\setlength\itemsep{0pt}
	\item \textbf{TSP asimetrico}: aniadir pesos $w[v][u]$ (no necesariamente simetricos).
	\item \textbf{Hamilton con inicio obligatorio}: fijar \texttt{src} en el DP inicial.
	\item \textbf{Bitmask DP mas compacto}: usar \texttt{long long} como mascara si $V \le 64$.
	\item \textbf{TSP optimo aproximado}: MST, 2-opt, simulated annealing.
	\item \textbf{Reconstruir el path}: usar \texttt{par[mask][v]} y backtrack.
\end{itemize}

\paragraph{Trampas y errores frecuentes.}
\begin{itemize}\setlength\itemsep{0pt}
	\item TSP asume grafo \textbf{completo} (con \texttt{LINF} para aristas faltantes); si no, aniadir chequeo.
	\item La condicion ``\texttt{mask == (1<<n) - 1}'' cierra el ciclo retornando al origen.
	\item Para $V > 20$, la memoria $2^V$ se vuelve prohibitiva; usar meet-in-the-middle o heuristicas.
\end{itemize}

\paragraph{Codigo.}
Ver \texttt{cpp.tex}, seccion \textit{Hamiltonian path / cycle (bitmask DP + backtracking)}.

% ============================================================
\section{DP}

\subsubsection{Knapsack 0/1}

\paragraph{Idea intuitiva.}
Dado un conjunto de $N$ objetos con peso $w_i$ y valor $v_i$, maximizar $\sum v_i$ con la restriccion $\sum w_i \le W$. La recurrencia clasica:
\[ dp[i][j] = \max(dp[i-1][j],\, dp[i-1][j - w_i] + v_i) \quad \text{si } j \ge w_i. \]
La optimizacion clave: la DP es 1-dimensional si iteramos $j$ \textbf{de mayor a menor} (asi no usamos el objeto $i$ dos veces). La razon de la iteracion inversa: en cada iteracion queremos el estado del ``item anterior''; iterar de menor a mayor contaminaria \texttt{dp[j - w_i]} con el item actual.

\paragraph{Propiedades clave (invariantes).}
\begin{itemize}\setlength\itemsep{0pt}
	\item $dp[j]$ al final = maximo valor con peso $\le j$.
	\item El bucle inverso \texttt{for (j = W; j >= w[i]; j--)} es lo que distingue 0/1 de unbounded.
	\item Solucion: maximo valor con cualquier peso $\le W$.
	\item La DP es exacta; no hay aproximacion.
\end{itemize}

\paragraph{Codigo clave.}
La version 1D optimizada:
\begin{verbatim}
vector<long long> dp(W + 1, 0);
for (int i = 0; i < N; ++i)
    for (int j = W; j >= w[i]; --j)  // invertido
        dp[j] = max(dp[j], dp[j - w[i]] + v[i]);
return dp[W];
\end{verbatim}

\paragraph{Complejidad.}
\begin{itemize}\setlength\itemsep{0pt}
	\item $O(N \cdot W)$ tiempo, $O(W)$ memoria.
	\item Constante muy pequena (solo sumas y max).
\end{itemize}

\paragraph{Donde tocar para modificar.}
\begin{itemize}\setlength\itemsep{0pt}
	\item \textbf{Reconstruir la seleccion}: aniadir \texttt{take[i][j]} (booleano o bool) en una version 2D.
	\item \textbf{Items con peso 0}: edge case; aniadir guard.
	\item \textbf{Knapsack con restricciones de grupos}: ``elegir a lo sumo 1 de cada grupo'' se modela como 0/1 con un peso extra.
	\item \textbf{Tamano de $W$}: si $W \le 10^9$ pero $N$ es chico, usar meet-in-the-middle.
	\item \textbf{Cambiar max por min}: aniadir \texttt{INF} en la inicializacion.
\end{itemize}

\paragraph{Trampas y errores frecuentes.}
\begin{itemize}\setlength\itemsep{0pt}
	\item Overflow en \texttt{dp[j - w[i]] + v[i]}; usar \texttt{long long} si los valores son grandes.
	\item El bucle \textbf{debe} ir de mayor a menor.
	\item Si los pesos son grandes y $N$ pequeno, conviene meet-in-the-middle.
	\item Para Knapsack exacto ($\sum w_i = W$), aniadir check al final.
\end{itemize}

\paragraph{Codigo.}
Ver \texttt{cpp.tex}, seccion \textit{Knapsack 0/1}.

\vspace{0.5em}

\subsubsection{Knapsack unbounded}

\paragraph{Idea intuitiva.}
Igual que Knapsack 0/1 pero \textbf{cada item se puede tomar multiples veces}. La recurrencia es similar:
\[ dp[j] = \max(dp[j],\, dp[j - w_i] + v_i), \]
pero ahora el bucle va de \textbf{menor a mayor} (porque tomar el item $i$ no afecta corridas futuras del mismo item). La idea: tras tomar el item $i$, queremos volver a tener la opcion de tomarlo; iterar de menor a mayor permite que el item ``se acumule''.

\paragraph{Propiedades clave (invariantes).}
\begin{itemize}\setlength\itemsep{0pt}
	\item El bucle \texttt{for (j = w[i]; j <= W; j++)} es lo que distingue unbounded.
	\item A diferencia de 0/1, no hay riesgo de ``usar dos veces el mismo item''.
	\item La DP es exacta; cada ``subconjunto'' se cuenta correctamente.
\end{itemize}

\paragraph{Codigo clave.}
La version 1D con bucle directo:
\begin{verbatim}
vector<long long> dp(W + 1, 0);
for (int i = 0; i < N; ++i)
    for (int j = w[i]; j <= W; ++j)  // directo (no invertido)
        dp[j] = max(dp[j], dp[j - w[i]] + v[i]);
return dp[W];
\end{verbatim}

\paragraph{Complejidad.}
\begin{itemize}\setlength\itemsep{0pt}
	\item $O(N \cdot W)$ tiempo, $O(W)$ memoria.
	\item Constante muy pequena.
\end{itemize}

\paragraph{Donde tocar para modificar.}
\begin{itemize}\setlength\itemsep{0pt}
	\item \textbf{Reconstruir cuantos de cada item}: aniadir \texttt{cnt[i]} y backtrack.
	\item \textbf{Cambiar moneda minima por maxima}: el snippet maximiza valor; para minimizar monedas, cambiar a \texttt{min} con inicializacion en \texttt{INF}.
	\item \textbf{Coin change}: si todos los pesos son positivos y queremos ``minimo numero de monedas'', es unbounded con min.
\end{itemize}

\paragraph{Trampas y errores frecuentes.}
\begin{itemize}\setlength\itemsep{0pt}
	\item Misma que Knapsack 0/1 con overflow.
	\item El orden de iteracion de items no afecta el resultado (convexo).
	\item Si se quiere ``maximo valor con exactamente $k$ items'', aniadir segunda dimension.
\end{itemize}

\paragraph{Codigo.}
Ver \texttt{cpp.tex}, seccion \textit{Knapsack unbounded}.

\vspace{0.5em}

\subsubsection{Knapsack meet-in-the-middle}

\paragraph{Idea intuitiva.}
Para $N \le 40$ y $W$ moderado, no se puede hacer $O(NW)$ directamente. Partir en dos mitades, enumerar todos los $2^{N/2}$ subconjuntos de cada mitad, y combinar. Para cada subconjunto de la segunda mitad con peso $sw_2$, encontrar el maximo valor de la primera mitad con peso $\le W - sw_2$. La idea: ``divide y venceras'' sobre el tamano de la entrada; cada mitad tiene mitad del espacio, asi que la combinacion es geometrica en lugar de multiplicativa.

\paragraph{Propiedades clave (invariantes).}
\begin{itemize}\setlength\itemsep{0pt}
	\item Hay $2^{N/2}$ subconjuntos por mitad; total $O(2^{N/2})$.
	\item \texttt{sums} guarda los valores de subconjuntos de la primera mitad ordenados (o preprocesados).
	\item Para cada subconjunto de la segunda mitad, busqueda lineal o binaria en \texttt{sums}.
	\item La mitad mas grande (con $N/2 + 1$ items para $N$ impar) tiene $2^{N/2+1}$ subconjuntos.
\end{itemize}

\paragraph{Codigo clave.}
Enumerar subconjuntos de la primera mitad:
\begin{verbatim}
int n1 = N / 2, n2 = N - n1;
vector<pair<long long, long long>> sums1;
for (int mask = 0; mask < (1 << n1); ++mask) {
    long long peso = 0, valor = 0;
    for (int i = 0; i < n1; ++i) if (mask & (1 << i)) {
        peso += w[i]; valor += v[i];
    }
    sums1.push_back({peso, valor});
}
sort(sums1.begin(), sums1.end());  // por peso
// eliminar duplicados dominados
\end{verbatim}

\paragraph{Complejidad.}
\begin{itemize}\setlength\itemsep{0pt}
	\item $O(2^{N/2})$ tiempo y memoria (mejor que $O(NW)$ cuando $W > 2^{N/2}$).
	\item Para $N = 40$: $\sim 2 \cdot 10^6$ subconjuntos, factible.
\end{itemize}

\paragraph{Donde tocar para modificar.}
\begin{itemize}\setlength\itemsep{0pt}
	\item \textbf{Tamano mayor}: si $N$ es mayor pero $W$ es chico, conviene iterar por peso, no por subconjuntos.
	\item \textbf{Varias mochilas}: adaptar para multiples $W$.
	\item \textbf{Optimizar el ``combinado''}: usar \texttt{upper\_bound} para encontrar el maximo valor con peso $\le \text{rem}.
	\item \textbf{Eliminar pares dominados}: si un par (peso, valor) tiene otro con menos peso y mas valor, eliminarlo.
\end{itemize}

\paragraph{Trampas y errores frecuentes.}
\begin{itemize}\setlength\itemsep{0pt}
	\item $N$ impar: una mitad tiene $N/2 + 1$ items.
	\item Overflow al sumar pesos/valores en \texttt{long long}.
	\item Sin deduplicacion, el algoritmo es correcto pero mas lento.
\end{itemize}

\paragraph{Codigo.}
Ver \texttt{cpp.tex}, seccion \textit{Knapsack meet-in-the-middle}.

\vspace{0.5em}

\subsubsection{LIS O(n log n)}

\paragraph{Idea intuitiva.}
La \textbf{Longest Increasing Subsequence} se calcula en $O(n \log n)$ con el truco del \textbf{patience sorting}: mantener un arreglo \texttt{dp} donde \texttt{dp[len]} es el minimo ultimo valor de una IS de largo \texttt{len}. Para cada $x$, encontrar el primer \texttt{dp[len] >= x} y reemplazarlo. El largo final es la LIS. La demostracion: si hay dos IS del mismo largo, la que termina en un valor menor es ``mejor'' (puede extenderse mas facilmente).

\paragraph{Propiedades clave (invariantes).}
\begin{itemize}\setlength\itemsep{0pt}
	\item \texttt{dp} es estrictamente creciente.
	\item \texttt{lower\_bound} da LIS estricta; \texttt{upper\_bound} da no-estricta ($\le$).
	\item El algoritmo es \textbf{online}: procesa elementos en orden.
	\item La longitud de \texttt{dp} es la LIS al final.
\end{itemize}

\paragraph{Codigo clave.}
La version online:
\begin{verbatim}
vector<int> dp;
for (int x : a) {
    auto it = lower_bound(dp.begin(), dp.end(), x);
    if (it == dp.end()) dp.push_back(x);
    else *it = x;
}
// dp.size() = longitud de LIS estricta
\end{verbatim}

\paragraph{Complejidad.}
\begin{itemize}\setlength\itemsep{0pt}
	\item $O(n \log n)$.
	\item Memoria $O(n)$.
\end{itemize}

\paragraph{Donde tocar para modificar.}
\begin{itemize}\setlength\itemsep{0pt}
	\item \textbf{Reconstruir la LIS}: aniadir \texttt{parent[]} y \texttt{posInDp[]} para backtrack.
	\item \textbf{Contar LIS}: DP clasica $O(n^2)$ o segment tree sobre $dp$.
	\item \textbf{LDS / Longest non-increasing}: invertir la condicion en \texttt{lower\_bound} o multiplicar por -1.
	\item \textbf{Longest bitonic subsequence}: combinar LIS desde izq y desde der.
	\item \textbf{LIS en 2D}: ordenar por una dimension y aplicar LIS en la otra.
\end{itemize}

\paragraph{Trampas y errores frecuentes.}
\begin{itemize}\setlength\itemsep{0pt}
	\item \texttt{lower\_bound} vs \texttt{upper\_bound}: si tu comparacion es $\le$ vs $<$, cambia.
	\item Si los elementos son \texttt{int} pero podrian ser negativos, aniadir offset para usar como indices.
	\item Para LIS reconstruida, aniadir el array \texttt{parent} que apunta al anterior en la IS.
\end{itemize}

\paragraph{Codigo.}
Ver \texttt{cpp.tex}, seccion \textit{LIS O(n log n)}.

\vspace{0.5em}

\subsubsection{LCS / Edit distance}

\paragraph{Idea intuitiva.}
\textbf{LCS} (Longest Common Subsequence): para dos strings $a$, $b$, $dp[i][j]$ = LCS de $a[0..i)$ y $b[0..j)$. Recurrencia:
\[ dp[i][j] = \begin{cases} dp[i-1][j-1] + 1 & \text{si } a[i-1] = b[j-1] \\ \max(dp[i-1][j], dp[i][j-1]) & \text{si no} \end{cases} \]

\textbf{Edit distance} (Levenshtein): numero minimo de operaciones (insert, delete, replace) para convertir $a$ en $b$. Recurrencia similar con un $+1$ adicional. La idea: ``alinear'' los strings eligiendo matches/mismatches/insert/delete.

\paragraph{Propiedades clave (invariantes).}
\begin{itemize}\setlength\itemsep{0pt}
	\item LCS y Edit Distance son DP $O(nm)$.
	\item Memoria se puede reducir a $O(\min(n, m))$ con rolling array.
	\item Edit distance con peso por operacion distinto: ajustar el \texttt{+1} por el peso.
	\item La subestructura optima: la respuesta para prefijos se construye de prefijos menores.
\end{itemize}

\paragraph{Codigo clave.}
La recurrencia clasica de LCS:
\begin{verbatim}
vector<vector<int>> dp(n + 1, vector<int>(m + 1, 0));
for (int i = 1; i <= n; ++i)
    for (int j = 1; j <= m; ++j)
        if (a[i-1] == b[j-1]) dp[i][j] = dp[i-1][j-1] + 1;
        else dp[i][j] = max(dp[i-1][j], dp[i][j-1]);
return dp[n][m];
\end{verbatim}

\paragraph{Complejidad.}
\begin{itemize}\setlength\itemsep{0pt}
	\item $O(NM)$ tiempo y memoria.
	\item Para $N, M \le 5000$: $\sim 2.5 \cdot 10^7$ operaciones.
\end{itemize}

\paragraph{Donde tocar para modificar.}
\begin{itemize}\setlength\itemsep{0pt}
	\item \textbf{Reconstruir la LCS / las operaciones}: aniadir \texttt{prev[i][j]} para backtrack.
	\item \textbf{LCS de varios strings}: NP-hard para $\ge 3$ strings; en la practica se hace con bitmask.
	\item \textbf{Edit distance con costos distintos}: cambiar el \texttt{+1} por el costo apropiado.
	\item \textbf{Rolling array}: si memoria es problema, mantener solo 2 filas.
	\item \textbf{Hirschberg}: algoritmo $O(NM)$ tiempo, $O(\min(N,M))$ memoria, con reconstruccion.
\end{itemize}

\paragraph{Trampas y errores frecuentes.}
\begin{itemize}\setlength\itemsep{0pt}
	\item Indices en la recurrencia: \texttt{a[i-1]} y \texttt{b[j-1]} (no \texttt{a[i]}).
	\item Para reconstruccion, el caso \texttt{a[i-1] == b[j-1]} tiene prioridad.
	\item Si los strings son muy largos, considerar Hirschberg o bitset optimization.
\end{itemize}

\paragraph{Codigo.}
Ver \texttt{cpp.tex}, seccion \textit{LCS / Edit distance}.

\vspace{0.5em}

\subsubsection{Bitmask DP over subsets}

\paragraph{Idea intuitiva.}
Cuando el estado es un \textbf{subconjunto} de $\{0, \dots, N-1\}$, se representa como una mascara de $N$ bits. Para iterar sobre todos los subconjuntos de \texttt{mask}:
\[ \text{for (sub = mask; sub; sub = (sub - 1) \& mask)} \{ \dots \} \]
Para iterar sobre todos los superconjuntos de \texttt{mask} en $[0, U)$:
\[ \text{for (sup = mask; sup < U; sup = (sup + 1) | mask)} \{ \dots \} \]
La idea: las mascaras forman un lattice de subconjuntos; los trucos aritmeticos (\texttt{(sub-1) \& mask} y \texttt{(sup+1) | mask}) enumeran todo el lattice eficientemente.

\paragraph{Propiedades clave (invariantes).}
\begin{itemize}\setlength\itemsep{0pt}
	\item Iterar subsets: $O(2^k)$ para una mascara con $k$ bits.
	\item \texttt{sub \& mask == sub} siempre; \texttt{mask - sub} no siempre da un subconjunto.
	\item \texttt{\_\_builtin\_popcount(mask)} cuenta los bits.
	\item \texttt{(sub - 1) \& mask} itera subsets de \texttt{mask} en orden decreciente de inclusion.
\end{itemize}

\paragraph{Codigo clave.}
Iteracion sobre subsets de una mascara:
\begin{verbatim}
// Todos los subsets de mask
for (int sub = mask; sub; sub = (sub - 1) & mask) {
    // procesar sub
}
// Incluir sub = 0 si es necesario
sub = 0;  // procesar primero
for (int sup = mask; sup < (1 << N); sup = (sup + 1) | mask) {
    // procesar sup (todos los superconjuntos)
}
\end{verbatim}

\paragraph{Complejidad.}
\begin{itemize}\setlength\itemsep{0pt}
	\item $O(N \cdot 2^N)$ para una DP estandar sobre todos los subsets.
	\item Iterar subsets de \texttt{mask}: $O(2^{|mask|})$.
\end{itemize}

\paragraph{Donde tocar para modificar.}
\begin{itemize}\setlength\itemsep{0pt}
	\item \textbf{TSP / Hamiltonian path}: bitmask DP (ver tambien \textit{Hamiltonian path} en \texttt{cpp.tex} seccion Grafos).
	\item \textbf{Set cover}: iterar subsets.
	\item \textbf{DP con supermask}: si el estado incluye ``todos los elementos que \textbf{no} hemos usado'', usar supermask iteration.
	\item \textbf{Bitmasks grandes}: con $N > 20$, aniadir compresion o simulacion con segmentos.
\end{itemize}

\paragraph{Trampas y errores frecuentes.}
\begin{itemize}\setlength\itemsep{0pt}
	\item El orden de iteracion de subsets de \texttt{mask} es decreciente (de \texttt{mask} a $0$).
	\item \texttt{mask = 0} no tiene subsets no triviales; el bucle lo saltea correctamente.
	\item El truco \texttt{(sub - 1) \& mask} requiere cuidado con \texttt{sub = 0} (termina el bucle).
	\item Superconjuntos: \texttt{sup < (1 << N)} debe estar bien definido.
\end{itemize}

\paragraph{Codigo.}
Ver \texttt{cpp.tex}, seccion \textit{Bitmask DP over subsets}.

\vspace{0.5em}

\subsubsection{SOS DP}

\paragraph{Idea intuitiva.}
Para calcular, para cada \texttt{mask}, la suma de $f(\text{submask})$ sobre todos los $\text{submask} \subseteq \text{mask}$:
\[ dp[mask] = \sum_{\text{sub} \subseteq mask} f(\text{sub}). \]
Se computa en $O(N \cdot 2^N)$ mediante un loop por bit: para cada $i$, sumar a $dp[mask]$ la contribucion de $dp[mask \oplus (1 << i)]$ cuando $mask$ tiene el bit $i$. La idea: cada subconjunto $s \subseteq m$ se obtiene quitando bits uno a uno; asi, sumar incrementalmente captura todos los subconjuntos en $O(N \cdot 2^N)$.

\paragraph{Propiedades clave (invariantes).}
\begin{itemize}\setlength\itemsep{0pt}
	\item Inicializar $dp = f$.
	\item Para cada bit $i$: \texttt{if (mask \& (1 << i)) dp[mask] += dp[mask \^{} (1 << i)]}.
	\item Variantes: supersets DP (sumar sobre superconjuntos).
	\item El orden de iteracion (bit externo, mascara interno) es crucial.
\end{itemize}

\paragraph{Codigo clave.}
SOS DP (Sum over Subsets):
\begin{verbatim}
vector<long long> dp(1 << N);
for (int i = 0; i < (1 << N); ++i) dp[i] = f[i];
for (int bit = 0; bit < N; ++bit)
    for (int mask = 0; mask < (1 << N); ++mask)
        if (mask & (1 << bit))
            dp[mask] += dp[mask ^ (1 << bit)];
// dp[mask] = suma de f(sub) para sub ⊆ mask
\end{verbatim}

\paragraph{Complejidad.}
\begin{itemize}\setlength\itemsep{0pt}
	\item $O(N \cdot 2^N)$.
	\item Memoria $O(2^N)$.
\end{itemize}

\paragraph{Donde tocar para modificar.}
\begin{itemize}\setlength\itemsep{0pt}
	\item \textbf{Max/min en lugar de suma}: cambiar \texttt{+=} por la operacion correspondiente.
	\item \textbf{Supersets DP}: iterar \texttt{mask} de $0$ a $2^N$ y aniadir \texttt{dp[mask] += dp[mask | (1<<i)]}.
	\item \textbf{XOR convolution} (subset convolution): version mas avanzada.
	\item \textbf{Aplicaciones}: contar subconjuntos con propiedades especificas.
\end{itemize}

\paragraph{Trampas y errores frecuentes.}
\begin{itemize}\setlength\itemsep{0pt}
	\item Si $N$ es 30 o mas, $2^N$ no entra en memoria; usar trucos o reducir $N$.
	\item El sentido del XOR (sumar a \texttt{mask \^{} bit}) es para ``agregar el bit $i$ al subconjunto''.
	\item En supersets DP, la condicion es \texttt{(mask \& (1<<bit)) == 0}.
\end{itemize}

\paragraph{Codigo.}
Ver \texttt{cpp.tex}, seccion \textit{SOS DP}.

\vspace{0.5em}

\subsubsection{Digit DP template}

\paragraph{Idea intuitiva.}
Contar numeros en $[0, N]$ que cumplen una propiedad sobre sus digitos. La idea: DP por posicion, manteniendo:
\begin{itemize}\setlength\itemsep{0pt}
	\item \texttt{pos}: posicion actual (de izquierda a derecha).
	\item \texttt{tight}: si los digitos ya colocados son exactamente los de $N$ (limite superior estricto).
	\item \texttt{state}: estado adicional (digitos usados, suma, etc.).
\end{itemize}
La clave: la condicion \texttt{tight} permite acotar correctamente el siguiente digito; sin ella, la DP contaria incorrectamente.

\paragraph{Propiedades clave (invariantes).}
\begin{itemize}\setlength\itemsep{0pt}
	\item Si \texttt{tight} es true, el siguiente digito esta acotado por el de $N$; si no, puede ser 0-9.
	\item Memoria: $O(D \cdot 2 \cdot \text{states})$.
	\item Caso base: si \texttt{pos == D}, devolver 1 si el estado cumple la propiedad, sino 0.
	\item El bit \texttt{started} distingue ``numero con menos digitos'' (con leading zeros).
\end{itemize}

\paragraph{Codigo clave.}
La estructura basica:
\begin{verbatim}
long long dp(int pos, bool tight, State s) {
    if (pos == D) return check(s) ? 1 : 0;
    if (memo[pos][tight][s] != -1) return memo;
    int limit = tight ? digits[pos] : 9;
    long long ans = 0;
    for (int d = 0; d <= limit; ++d)
        ans += dp(pos + 1, tight && (d == limit), update(s, d, pos));
    return memo[pos][tight][s] = ans;
}
\end{verbatim}

\paragraph{Complejidad.}
\begin{itemize}\setlength\itemsep{0pt}
	\item $O(D \cdot 10 \cdot \text{states})$ donde $D$ es el numero de digitos.
	\item Para $D \le 19$ (hasta $10^{18}$): $\sim 200 \cdot \text{states}$.
\end{itemize}

\paragraph{Donde tocar para modificar.}
\begin{itemize}\setlength\itemsep{0pt}
	\item \textbf{Cambiar la propiedad}: modificar \texttt{updateState} y la condicion base.
	\item \textbf{Multiples estados}: aniadir mas dimensiones al \texttt{memo}.
	\item \textbf{Leading zeros}: aniadir un flag \texttt{started} para distinguir numeros con menos digitos.
	\item \textbf{Rango $[L, R]$}: calcular $f(R) - f(L-1)$.
\end{itemize}

\paragraph{Trampas y errores frecuentes.}
\begin{itemize}\setlength\itemsep{0pt}
	\item Olvidar resetear \texttt{memo} entre queries.
	\item Si la cantidad de estados es grande, el \texttt{memo[pos][tight][state]} puede ser muy grande; usar \texttt{map}.
	\item La condicion \texttt{started} es crucial para numeros con menos digitos que $D$.
	\item Si $N$ es muy grande ($> 10^{18}$), aniadir soporte para long long.
\end{itemize}

\paragraph{Codigo.}
Ver \texttt{cpp.tex}, seccion \textit{Digit DP template}.

\vspace{0.5em}

\subsubsection{Tree DP (reroot)}

\paragraph{Idea intuitiva.}
Calcular alguna propiedad asumiendo que \textbf{cada nodo es la raiz} del arbol. La tecnica clasica de \textbf{rerooting} en dos pasadas:
\begin{itemize}\setlength\itemsep{0pt}
	\item \textbf{DFS 1}: desde una raiz arbitraria, calcular \texttt{dp[v]} para el subarbol de $v$.
	\item \textbf{DFS 2}: para cada hijo $u$ de $v$, calcular \texttt{up[u]} = respuesta si la raiz fuera $u$.
\end{itemize}
La demostracion: tras la primera DFS, cada nodo conoce ``lo que aporta su subarbol''; en la segunda, se transmite ``lo que aporta el resto del arbol''.

\paragraph{Propiedades clave (invariantes).}
\begin{itemize}\setlength\itemsep{0pt}
	\item \texttt{dp[v]} depende solo del subarbol de $v$ (con la raiz original).
	\item \texttt{up[u]} se calcula de manera que cuando la raiz es $u$, ``todo lo que estaba en $v$ se ajusta''.
	\item Formula para reroot depende del problema; el snippet usa el ejemplo ``suma de distancias''.
	\item La combinacion \texttt{dp[u] + up[u]} da la respuesta con raiz $u$.
\end{itemize}

\paragraph{Codigo clave.}
Reroot para suma de distancias:
\begin{verbatim}
// dfs1: dp[v] = suma de distancias desde v al subarbol
void dfs1(int v, int p) {
    dp[v] = 0;
    for (int u : adj[v]) if (u != p) {
        dfs1(u, v);
        dp[v] += dp[u] + sz[u];  // cada hijo aporta su subarbol
    }
}
// dfs2: up[u] = distancia desde el resto del arbol a u
void dfs2(int v, int p) {
    for (int u : adj[v]) if (u != p) {
        up[u] = up[v] + dp[v] - dp[u] - sz[u] + (n - sz[u]);
        dfs2(u, v);
    }
}
\end{verbatim}

\paragraph{Complejidad.}
\begin{itemize}\setlength\itemsep{0pt}
	\item $O(N)$ para ambas pasadas.
	\item Cada arista se procesa un numero constante de veces.
\end{itemize}

\paragraph{Donde tocar para modificar.}
\begin{itemize}\setlength\itemsep{0pt}
	\item \textbf{Cambiar la cantidad a calcular}: modificar la operacion en \texttt{dfs1} y la formula en \texttt{dfs2}.
	\item \textbf{Arbol con pesos}: aniadir \texttt{weights[v]} y ajustar.
	\item \textbf{Re-root en arboles con ciclos}: requiere otra tecnica (DSU on tree).
	\item \textbf{Rooted trees diferentes}: si la propiedad no es invariante por raiz, usar otra tecnica.
\end{itemize}

\paragraph{Trampas y errores frecuentes.}
\begin{itemize}\setlength\itemsep{0pt}
	\item La formula de \texttt{up[u]} debe sumar y restar correctamente las contribuciones de $u$ y del resto.
	\item Para arboles con pesos, aniadir el peso en las formulas.
	\item Verificar con ejemplos pequenos antes de generalizar.
\end{itemize}

\paragraph{Codigo.}
Ver \texttt{cpp.tex}, seccion \textit{Tree DP (reroot)}.

% ============================================================
\section{Geometria}
% ============================================================

\subsection{Puntos y segmentos}

\subsubsection{Point + cross + dot}

\paragraph{Idea intuitiva.}
La estructura \texttt{Point} representa un punto en 2D con coordenadas enteras. Las dos operaciones vectoriales clave son:
\begin{itemize}\setlength\itemsep{0pt}
	\item \textbf{Cross product} $P \times Q = P_x Q_y - P_y Q_x$. Da el ``area con signo'' del paralelogramo formado por $P$ y $Q$. Usado para orientacion.
	\item \textbf{Dot product} $P \cdot Q = P_x Q_x + P_y Q_y$. Da el producto de las longitudes por el coseno del angulo. Usado para proyeccion.
\end{itemize}
La interpretacion geometrica del cross product: es el area con signo del paralelogramo; el signo indica la orientacion relativa.

\paragraph{Propiedades clave (invariantes).}
\begin{itemize}\setlength\itemsep{0pt}
	\item \texttt{cross(a, b) > 0} si $b$ esta a la izquierda de $a$ (sistema de coordenadas usual).
	\item \texttt{cross(a, b) = 0} si son colineales.
	\item \texttt{dot(a, b) > 0} si el angulo es agudo.
	\item El snippet define \texttt{cross} respecto a \texttt{this}, util para orientacion de 3 puntos.
	\item $|cross|^2 + |dot|^2 = |P|^2 \cdot |Q|^2$ (identidad de Lagrange).
\end{itemize}

\paragraph{Codigo clave.}
Las dos operaciones vectoriales:
\begin{verbatim}
long long cross(const Point& a, const Point& b) {
    return (long long)a.x * b.y - (long long)a.y * b.x;
}
long long dot(const Point& a, const Point& b) {
    return (long long)a.x * b.x + (long long)a.y * b.y;
}
\end{verbatim}

\paragraph{Complejidad.}
\begin{itemize}\setlength\itemsep{0pt}
	\item $O(1)$ por operacion.
	\item Constante muy pequena (4 multiplicaciones + 2 sumas).
\end{itemize}

\paragraph{Donde tocar para modificar.}
\begin{itemize}\setlength\itemsep{0pt}
	\item \textbf{Coordenadas doubles}: cambiar \texttt{int} por \texttt{double} o \texttt{long double}. Cuidado con la precision.
	\item \textbf{3D}: aniadir \texttt{z} y adaptar cross/dot.
	\item \textbf{Operadores adicionales}: aniadir \texttt{operator+}, \texttt{operator-}, \texttt{operator*} (escalar).
	\item \textbf{Distancia al cuadrado}: aniadir \texttt{len2()} que retorna $x^2 + y^2$.
\end{itemize}

\paragraph{Trampas y errores frecuentes.}
\begin{itemize}\setlength\itemsep{0pt}
	\item Cross product \textbf{no es} conmutativo; $P \times Q = -Q \times P$.
	\item Overflow con coordenadas grandes ($|x|, |y| > 10^4$) si se hace \texttt{x * y}.
	\item El signo de cross depende del sistema de coordenadas (en computacion, usualmente CCW positivo).
	\item Para ``igualdad'' de doubles, aniadir tolerancia \texttt{eps}.
\end{itemize}

\paragraph{Codigo.}
Ver \texttt{cpp.tex}, seccion \textit{Point + cross + dot}.

\vspace{0.5em}

\subsubsection{Interseccion de segmentos}

\paragraph{Idea intuitiva.}
Dados dos segmentos $[a, b]$ y $[c, d]$, determinar si se intersectan. La idea clasica:
\begin{itemize}\setlength\itemsep{0pt}
	\item Verificar orientaciones: $a, b, c$ y $a, b, d$ deben estar en lados opuestos de la recta $ab$; igual para $cd$ y $ab$.
	\item Caso colineal: si las rectas son paralelas (\texttt{cross} $= 0$) y los segmentos son colineales, verificar que los bounding boxes se intersectan.
\end{itemize}
La demostracion: dos segmentos se intersectan sii cada segmento ``cruza'' la linea que contiene al otro (estricto) o ambos son colineales con solapamiento.

\paragraph{Propiedades clave (invariantes).}
\begin{itemize}\setlength\itemsep{0pt}
	\item Funciona con segmentos en cualquier orientacion.
	\item Caso colineal: requiere check adicional de bounding box.
	\item Para intersectar en un \textbf{punto} especifico: aniadir division con cuidado de precision.
	\item Funciona con segmentos verticales/horizontales sin casos especiales.
\end{itemize}

\paragraph{Codigo clave.}
Test clasico de interseccion:
\begin{verbatim}
int sign(long long x) { return (x > 0) - (x < 0); }
bool intersect(Point a, Point b, Point c, Point d) {
    long long o1 = cross(b - a, c - a);
    long long o2 = cross(b - a, d - a);
    long long o3 = cross(d - c, a - c);
    long long o4 = cross(d - c, b - c);
    if (sign(o1) != sign(o2) && sign(o3) != sign(o4)) return true;
    // casos colineales
    if (o1 == 0 && onSegment(a, b, c)) return true;
    if (o2 == 0 && onSegment(a, b, d)) return true;
    if (o3 == 0 && onSegment(c, d, a)) return true;
    if (o4 == 0 && onSegment(c, d, b)) return true;
    return false;
}
\end{verbatim}

\paragraph{Complejidad.}
\begin{itemize}\setlength\itemsep{0pt}
	\item $O(1)$.
	\item Constante pequena.
\end{itemize}

\paragraph{Donde tocar para modificar.}
\begin{itemize}\setlength\itemsep{0pt}
	\item \textbf{Interseccion con punto}: aniadir la division despues de verificar orientacion.
	\item \textbf{Interseccion rayo-segmento}: ajustar la primera condicion.
	\item \textbf{Poligonos convexos}: usar separacion de ejes (SAT).
	\item \textbf{Multiples queries}: precomputar bounding boxes o usar segment tree.
\end{itemize}

\paragraph{Trampas y errores frecuentes.}
\begin{itemize}\setlength\itemsep{0pt}
	\item Si los segmentos tocan solo en un extremo, el snippet lo considera interseccion (algunos problemas quieren estricto). Ajustar segun el problema.
	\item El caso colineal es comun; olvidar el check adicional de bounding box da falsos negativos.
	\item En problemas con muchos segmentos, considerar sweep line para $O(n \log n)$.
\end{itemize}

\paragraph{Codigo.}
Ver \texttt{cpp.tex}, seccion \textit{Interseccion de segmentos}.

\vspace{0.5em}

\subsubsection{Contains colineal}

\paragraph{Idea intuitiva.}
Verifica si un punto $c$ esta sobre el segmento $[a, b]$ (asumiendo que ya son colineales). Si \texttt{cross(a, b, c) = 0} y $c$ esta dentro del bounding box de $[a, b]$, entonces $c$ esta en el segmento. La razon: si los tres puntos son colineales, $c$ esta en el segmento sii sus coordenadas estan entre las de $a$ y $b$.

\paragraph{Propiedades clave (invariantes).}
\begin{itemize}\setlength\itemsep{0pt}
	\item Asume colinealidad (no la verifica mas alla de \texttt{cross == 0}).
	\item Usa comparaciones con \texttt{<=} para extremos \textbf{inclusivos}.
	\item Funciona con segmentos verticales (la condicion se reduce a una sola coordenada).
\end{itemize}

\paragraph{Codigo clave.}
Test de contencion en segmento (asume colinealidad):
\begin{verbatim}
bool onSegment(Point a, Point b, Point c) {
    return min(a.x, b.x) <= c.x && c.x <= max(a.x, b.x) &&
           min(a.y, b.y) <= c.y && c.y <= max(a.y, b.y);
}
\end{verbatim}

\paragraph{Complejidad.}
\begin{itemize}\setlength\itemsep{0pt}
	\item $O(1)$.
	\item Constante minima.
\end{itemize}

\paragraph{Donde tocar para modificar.}
\begin{itemize}\setlength\itemsep{0pt}
	\item \textbf{Contener estricto}: cambiar \texttt{<=} a \texttt{<} para excluir extremos.
	\item \textbf{Coordenadas doubles}: usar tolerancia (\texttt{eps}).
	\item \textbf{Solo x o solo y}: aniadir check especifico para segmentos verticales/horizontales.
\end{itemize}

\paragraph{Trampas y errores frecuentes.}
\begin{itemize}\setlength\itemsep{0pt}
	\item No usar este snippet para puntos \textbf{no colineales}; siempre verificar primero la colinealidad.
	\item Si los puntos tienen doubles, el chequeo \texttt{cross == 0} puede fallar; aniadir tolerancia.
	\item El bounding box incluye los extremos; verificar si el problema quiere exclusivo o inclusivo.
\end{itemize}

\paragraph{Codigo.}
Ver \texttt{cpp.tex}, seccion \textit{Contains colineal}.

\subsection{Poligonos}

\subsubsection{Convex Hull (monotono)}

\paragraph{Idea intuitiva.}
Algoritmo de \textbf{Andrew's monotone chain}: ordenar puntos por $(x, y)$, luego construir la hull inferior y superior por separado. Para cada punto, aniadirlo al hull y, mientras el ultimo giro no sea ``counter-clockwise'' (\texttt{cross > 0}), popear. La demostracion: si los puntos se procesan en orden, el hull inferior y superior cubren la frontera completa sin auto-intersecciones.

\paragraph{Propiedades clave (invariantes).}
\begin{itemize}\setlength\itemsep{0pt}
	\item Ordena los puntos en $O(n \log n)$.
	\item Construye la hull en $O(n)$ despues.
	\item Output: poligono convexo en counter-clockwise, sin punto final duplicado.
	\item El \texttt{<= 0} incluye colineales; usar \texttt{< 0} para excluirlos.
	\item Cada punto aparece en la hull inferior o superior (no ambos).
\end{itemize}

\paragraph{Codigo clave.}
Construccion del hull inferior y superior:
\begin{verbatim}
sort(p.begin(), p.end());
vector<Point> lower, upper;
for (auto& p : points) {
    while (lower.size() >= 2 && cross(lower.back() - lower[lower.size()-2],
                                       p - lower.back()) <= 0)
        lower.pop_back();
    lower.push_back(p);
}
for (int i = points.size() - 1; i >= 0; --i) {
    auto& p = points[i];
    while (upper.size() >= 2 && cross(upper.back() - upper[upper.size()-2],
                                       p - upper.back()) <= 0)
        upper.pop_back();
    upper.push_back(p);
}
vector<Point> hull = lower;
hull.insert(hull.end(), upper.begin() + 1, upper.end() - 1);
\end{verbatim}

\paragraph{Complejidad.}
\begin{itemize}\setlength\itemsep{0pt}
	\item $O(n \log n)$.
	\item Constante muy pequena (solo cross products).
\end{itemize}

\paragraph{Donde tocar para modificar.}
\begin{itemize}\setlength\itemsep{0pt}
	\item \textbf{Colineales}: cambiar \texttt{<= 0} por \texttt{< 0} para excluirlos del hull.
	\item \textbf{Puntos repetidos}: \texttt{unique} despues de sort.
	\item \textbf{Hull 3D}: gift wrapping es $O(n^2)$; convex hull 3D es mas complejo.
	\item \textbf{Output con los puntos}: aniadir el hull completo (no solo upper/lower).
\end{itemize}

\paragraph{Trampas y errores frecuentes.}
\begin{itemize}\setlength\itemsep{0pt}
	\item El \texttt{pop\_back} y el \texttt{reverse} son esenciales para no duplicar el ultimo punto.
	\item Para excluir colineales, ajustar el \texttt{<= 0} vs \texttt{< 0}.
	\item Si los puntos estan duplicados, aniadir \texttt{unique} antes del sort.
\end{itemize}

\paragraph{Codigo.}
Ver \texttt{cpp.tex}, seccion \textit{Convex Hull (monotono)}.

\vspace{0.5em}

\subsubsection{Polygon Area}

\paragraph{Idea intuitiva.}
El area (con signo) de un poligono se calcula con la \textbf{Shoelace formula}: $A = \frac{1}{2} \sum_{i=0}^{n-1} (x_i y_{i+1} - x_{i+1} y_i)$, donde el indice es modulo $n$. El snippet retorna el \textbf{doble} del area (en valor absoluto). La formula sale del ``sumar areas de trapezoides'': cada lado del poligono define un trapezoide con el eje $x$; la suma con signo da el area.

\paragraph{Propiedades clave (invariantes).}
\begin{itemize}\setlength\itemsep{0pt}
	\item El signo indica la orientacion: positivo si counter-clockwise.
	\item Para el doble: no dividir por 2 si no es necesario (evita fracciones).
	\item Funciona para poligonos simples (no auto-intersectantes).
	\item Si el poligono es no convexo, la formula da el area ``con signo''.
\end{itemize}

\paragraph{Codigo clave.}
La suma de trapezoides:
\begin{verbatim}
long long shoelace(const vector<Point>& p) {
    long long s = 0;
    int n = p.size();
    for (int i = 0; i < n; ++i)
        s += (long long)p[i].x * p[(i+1)%n].y - (long long)p[(i+1)%n].x * p[i].y;
    return abs(s);  // doble del area
}
\end{verbatim}

\paragraph{Complejidad.}
\begin{itemize}\setlength\itemsep{0pt}
	\item $O(n)$.
	\item Constante minima.
\end{itemize}

\paragraph{Donde tocar para modificar.}
\begin{itemize}\setlength\itemsep{0pt}
	\item \textbf{Area exacta}: dividir por 2 (con cuidado si da fraccion; usar \texttt{double}).
	\item \textbf{Area firmada}: quitar \texttt{abs} y dejar el signo.
	\item \textbf{Punto dentro de poligono convexo}: con la hull, chequear que esta a un mismo lado de todas las aristas.
	\item \textbf{Triangulacion}: descomponer en triangulos y sumar areas.
\end{itemize}

\paragraph{Trampas y errores frecuentes.}
\begin{itemize}\setlength\itemsep{0pt}
	\item Si los puntos no estan en orden (CCW o CW), el area firmada puede ser incorrecta. Para signed area consistente, usar siempre el mismo orden.
	\item Para poligonos con holes, sumar las areas con signo apropiado.
	\item Con doubles, aniadir tolerancia \texttt{eps}.
\end{itemize}

\paragraph{Codigo.}
Ver \texttt{cpp.tex}, seccion \textit{Polygon Area}.

\vspace{0.5em}

\subsubsection{Distance point-segment / point-line}

\paragraph{Idea intuitiva.}
Distancia minima de un punto $p$ a un segmento $[a, b]$ o a la recta que pasa por $a$ y $b$.
\begin{itemize}\setlength\itemsep{0pt}
	\item \textbf{Linea}: $d = \frac{|ab \times ap|}{|ab|}$.
	\item \textbf{Segmento}: si la proyeccion de $p$ sobre $ab$ cae fuera del segmento, la distancia es al extremo mas cercano.
\end{itemize}
La razon geometrica: la distancia a una linea es el area del paralelogramo dividido por la base; la proyeccion ortogonal indica si estamos ``dentro'' o ``fuera'' del segmento.

\paragraph{Propiedades clave (invariantes).}
\begin{itemize}\setlength\itemsep{0pt}
	\item \texttt{dist2PointLine} retorna distancia \textbf{al cuadrado}.
	\item \texttt{dist2PointSegment} distingue 3 casos: $p$ antes de $a$, despues de $b$, o proyectado sobre el segmento.
	\item Si \texttt{len2 = 0} (degenerate), la distancia es $|ap|$ o $|bp|$.
	\item La proyeccion se calcula con dot products y el parametro $t \in [0, 1]$.
\end{itemize}

\paragraph{Codigo clave.}
Distancia al cuadrado a segmento:
\begin{verbatim}
long long dist2Segment(Point p, Point a, Point b) {
    long long len2 = (a-b).len2();
    if (len2 == 0) return (p-a).len2();  // degenerate
    long long t = max(0LL, min(1LL, dot(p-a, b-a) / len2));
    Point proj = {a.x + t*(b.x - a.x), a.y + t*(b.y - a.y)};
    return (p - proj).len2();
}
\end{verbatim}

\paragraph{Complejidad.}
\begin{itemize}\setlength\itemsep{0pt}
	\item $O(1)$.
	\item Constante minima.
\end{itemize}

\paragraph{Donde tocar para modificar.}
\begin{itemize}\setlength\itemsep{0pt}
	\item \textbf{Distancia (no al cuadrado)}: aniadir \texttt{sqrt} al final.
	\item \textbf{Coordenadas doubles}: aniadir \texttt{long double} y tolerancia.
	\item \textbf{3D}: aniadir coordenada $z$ y adaptar.
	\item \textbf{Multiples queries}: usar KD-tree o segment tree para acelerar.
\end{itemize}

\paragraph{Trampas y errores frecuentes.}
\begin{itemize}\setlength\itemsep{0pt}
	\item Overflow en \texttt{area2 * area2} si el area es grande; \texttt{long long} alcanza para $|x|, |y| \le 10^9$ pero no mas.
	\item Si $a = b$ (segmento degenerado), aniadir check explicito.
	\item En doubles, aniadir tolerancia \texttt{eps} para comparar.
\end{itemize}

\paragraph{Codigo.}
Ver \texttt{cpp.tex}, seccion \textit{Distance point-segment / point-line}.

\vspace{0.5em}

\subsubsection{Projection of point onto line}

\paragraph{Idea intuitiva.}
Proyeccion ortogonal de $p$ sobre la recta que pasa por $a$ y $b$. El parametro $t$ tal que la proyeccion es $a + t \cdot (b - a)$ es:
\[ t = \frac{(p - a) \cdot (b - a)}{|b - a|^2}. \]
La idea: descomponemos $\vec{ap}$ en una componente paralela a $\vec{ab}$ y una perpendicular; $t$ mide la posicion en la primera.

\paragraph{Propiedades clave (invariantes).}
\begin{itemize}\setlength\itemsep{0pt}
	\item Si $t \in [0, 1]$, la proyeccion cae sobre el \textbf{segmento}; si no, cae sobre la extension.
	\item Retorna \texttt{double} (puede tener precision issues).
	\item $t < 0$: $p$ esta ``antes'' de $a$; $t > 1$: $p$ esta ``despues'' de $b$.
\end{itemize}

\paragraph{Codigo clave.}
Proyeccion ortogonal:
\begin{verbatim}
Point projection(Point p, Point a, Point b) {
    Point ab = b - a, ap = p - a;
    double t = (double)dot(ap, ab) / ab.len2();
    return {a.x + t * ab.x, a.y + t * ab.y};
}
\end{verbatim}

\paragraph{Complejidad.}
\begin{itemize}\setlength\itemsep{0pt}
	\item $O(1)$.
	\item Constante minima (dos dot products y una division).
\end{itemize}

\paragraph{Donde tocar para modificar.}
\begin{itemize}\setlength\itemsep{0pt}
	\item \textbf{Proyeccion sobre segmento}: clipar $t$ a $[0, 1]$ antes de calcular el punto.
	\item \textbf{Distancia al segmento via proyeccion}: combinar con el modulo anterior.
	\item \textbf{Solo el parametro $t$}: si solo necesitas saber si esta dentro del segmento, retorna $t$ sin construir el punto.
\end{itemize}

\paragraph{Trampas y errores frecuentes.}
\begin{itemize}\setlength\itemsep{0pt}
	\item Si $a = b$ ($|ab| = 0$), la proyeccion no esta definida; aniadir guard.
	\item En doubles, aniadir tolerancia \texttt{eps} para comparar $t$ con 0 o 1.
	\item La division por \texttt{len2} puede ser cara si las coordenadas son grandes; usar \texttt{double} o precomputar.
\end{itemize}

\paragraph{Codigo.}
Ver \texttt{cpp.tex}, seccion \textit{Projection of point onto line}.

\vspace{0.5em}

\subsubsection{Point in polygon (ray casting)}

\paragraph{Idea intuitiva.}
Algoritmo de \textbf{ray casting}: contar cuantas veces un rayo horizontal desde $p$ cruza los bordes del poligono. Si es impar, $p$ esta dentro; si par, fuera. El snippet implementa la version clasica recorriendo aristas. La razon: cada cruce cambia el lado del rayo respecto al poligono; un numero impar de cambios significa ``terminamos dentro''.

\paragraph{Propiedades clave (invariantes).}
\begin{itemize}\setlength\itemsep{0pt}
	\item Funciona con poligonos simples (no necesariamente convexos).
	\item Las intersecciones en vertices cuentan especialmente.
	\item Para poligonos con \textbf{huecos}, sigue funcionando.
	\item La respuesta es estrictamente 0 (fuera) o 1 (dentro) para poligonos simples.
\end{itemize}

\paragraph{Codigo clave.}
El test clasico de ray casting:
\begin{verbatim}
bool pointInPolygon(Point p, const vector<Point>& poly) {
    int n = poly.size();
    bool inside = false;
    for (int i = 0, j = n - 1; i < n; j = i++) {
        if ((poly[i].y > p.y) != (poly[j].y > p.y) &&
            p.x < (poly[j].x - poly[i].x) * (p.y - poly[i].y) /
                    (poly[j].y - poly[i].y) + poly[i].x)
            inside = !inside;
    }
    return inside;
}
\end{verbatim}

\paragraph{Complejidad.}
\begin{itemize}\setlength\itemsep{0pt}
	\item $O(n)$ por query.
	\item Memoria $O(1)$ (no se almacena el rayo).
\end{itemize}

\paragraph{Donde tocar para modificar.}
\begin{itemize}\setlength\itemsep{0pt}
	\item \textbf{Winding number}: version alternativa que cuenta el numero de vueltas.
	\item \textbf{Poligono convexo}: check lineal contra las aristas (mas rapido).
	\item \textbf{Multiples queries}: precomputar o usar segment tree sobre las aristas.
	\item \textbf{Punto en frontera}: aniadir check adicional si es colineal con alguna arista.
\end{itemize}

\paragraph{Trampas y errores frecuentes.}
\begin{itemize}\setlength\itemsep{0pt}
	\item El manejo del vertice (cuando el rayo pasa por un vertice) puede causar double-counting; el snippet lo evita con la condicion \texttt{(poly[i].y > p.y) != (poly[j].y > p.y)}.
	\item Si el poligono tiene holes, el algoritmo sigue funcionando.
	\item Para poligonos no simples (auto-intersectantes), el resultado puede ser ambiguo.
\end{itemize}

\paragraph{Codigo.}
Ver \texttt{cpp.tex}, seccion \textit{Point in polygon (ray casting)}.

\vspace{0.5em}

\subsubsection{Minkowski sum (convex polygons)}

\paragraph{Idea intuitiva.}
La \textbf{suma de Minkowski} $A \oplus B$ de dos poligonos convexos es otro poligono convexo que representa $\{a + b : a \in A, b \in B\}$. Algoritmo: ordenar las aristas por angulo polar (CCW) y ``mergearlas'' como en merge sort; el resultado es un poligono CCW sin repetir el ultimo punto. La idea: el resultado tiene una arista paralela a cada arista de los originales; la suma se obtiene barriendo los angulos.

\paragraph{Propiedades clave (invariantes).}
\begin{itemize}\setlength\itemsep{0pt}
	\item Si $A, B$ son convexos con $n, m$ vertices, $A \oplus B$ tiene $\le n + m$ vertices.
	\item Util para \textbf{detectar colision}: $A$ colisiona con $B$ si $0 \in A \oplus (-B)$.
	\item Asume poligonos en orden \textbf{counter-clockwise}.
	\item El resultado es convexo (suma de convexos es convexo).
\end{itemize}

\paragraph{Codigo clave.}
La suma de Minkowski:
\begin{verbatim}
vector<Point> minkowski(const vector<Point>& A, const vector<Point>& B) {
    vector<Point> result;
    int i = 0, j = 0;
    int n = A.size(), m = B.size();
    do {
        result.push_back(A[i] + B[j]);
        long long cross = (A[(i+1)%n] - A[i]) ^ (B[(j+1)%m] - B[j]);
        if (cross >= 0) i = (i + 1) % n;
        if (cross <= 0) j = (j + 1) % m;
    } while (i || j);
    return result;
}
\end{verbatim}

\paragraph{Complejidad.}
\begin{itemize}\setlength\itemsep{0pt}
	\item $O(n + m)$.
	\item Memoria $O(n + m)$.
\end{itemize}

\paragraph{Donde tocar para modificar.}
\begin{itemize}\setlength\itemsep{0pt}
	\item \textbf{Deteccion de colision}: aniadir $(-B)$ (reflejar) y verificar si $0$ esta en el resultado.
	\item \textbf{Poligonos no convexos}: la suma puede no ser poligono simple; algoritmo distinto.
	\item \textbf{Multiples queries}: precomputar las aristas una vez.
	\item \textbf{Robot motion planning}: usar para planear trayectorias.
\end{itemize}

\paragraph{Trampas y errores frecuentes.}
\begin{itemize}\setlength\itemsep{0pt}
	\item Si los poligonos no son CCW, la suma es incorrecta. Normalizar antes.
	\item El bucle termina cuando ambos indices vuelven a 0; aniadir check.
	\item Para la deteccion de colision, aniadir el test \texttt{0 in A \^{} (-B)}.
\end{itemize}

\paragraph{Codigo.}
Ver \texttt{cpp.tex}, seccion \textit{Minkowski sum (convex polygons)}.

\subsection{Herramientas}

\subsubsection{Li Chao Tree (long long)}

\paragraph{Idea intuitiva.}
Mantiene un \textbf{lower envelope} (o upper) de un conjunto de rectas $y = mx + b$ sobre un dominio $[X_{LO}, X_{HI}]$ y responde ``minimo $y$ en un $x$ dado'' en $O(\log X)$. La estructura es un segment tree donde cada nodo guarda la recta ``mejor'' para su segmento medio; al aniadir una recta, se compara con la actual y se decide cual es mejor en cada mitad, descendiendo recursivamente. La demostracion: en cada nivel, solo se guarda la mejor recta para el ``punto medio''; las otras se transmiten a los hijos.

\paragraph{Propiedades clave (invariantes).}
\begin{itemize}\setlength\itemsep{0pt}
	\item \textbf{Online}: cada \texttt{addLine} en $O(\log X)$.
	\item Solo insercion; no hay borrar (para borrar, usar Li Chao segmentado con reset).
	\item La recta guardada en un nodo es la ``mejor'' en su punto medio.
	\item Las rectas ``perdedoras'' se transmiten a los hijos para futuros inserts.
\end{itemize}

\paragraph{Codigo clave.}
La insercion en Li Chao:
\begin{verbatim}
void addLine(Line nw, int v = 1, int l = X_LO, int r = X_HI) {
    int m = (l + r) / 2;
    bool left = nw.get(l) < tree[v].get(l);
    bool mid = nw.get(m) < tree[v].get(m);
    if (mid) swap(tree[v], nw);
    if (r - l == 1) return;
    if (left != mid) addLine(nw, v*2, l, m);
    else             addLine(nw, v*2+1, m, r);
}
\end{verbatim}

\paragraph{Complejidad.}
\begin{itemize}\setlength\itemsep{0pt}
	\item $O(\log X)$ por insercion y query, donde $X$ es el tamano del dominio.
	\item Memoria $O(\log X)$ por insercion, $O(X)$ total.
\end{itemize}

\paragraph{Donde tocar para modificar.}
\begin{itemize}\setlength\itemsep{0pt}
	\item \textbf{Upper envelope en vez de lower}: invertir el signo (\texttt{>} en vez de \texttt{<}).
	\item \textbf{Queries de segmento} (no full line): si la recta solo esta activa en un rango, aniadir bandera \texttt{active}.
	\item \textbf{Eliminar rectas}: usar Li Chao segmentado (mas complejo).
	\item \textbf{Pendientes grandes}: si las pendientes tienen magnitud $> 10^9$, aniadir \texttt{\_\_int128} para evitar overflow.
\end{itemize}

\paragraph{Trampas y errores frecuentes.}
\begin{itemize}\setlength\itemsep{0pt}
	\item El parametro \texttt{lo == hi} debe retornar el valor del nodo (caso base).
	\item La comparacion \texttt{newLine.get(mid) < node->line.get(mid)} decide quien es mejor en el medio.
	\item Si las rectas tienen $m$ iguales pero $b$ distintos, se queda con la primera (tie-break).
\end{itemize}

\paragraph{Codigo.}
Ver \texttt{cpp.tex}, seccion \textit{Li Chao Tree (long long)}.

\vspace{0.5em}

\subsubsection{Sweep line template}

\paragraph{Idea intuitiva.}
Patron de diseno para problemas geometricos donde se barre una linea sobre el plano y se mantiene una estructura de datos sobre el otro eje. Eventos (add, query, remove) se ordenan por la coordenada de barrido y se procesan secuencialmente.

\paragraph{Propiedades clave (invariantes).}
\begin{itemize}\setlength\itemsep{0pt}
	\item Eventos: tuplas \texttt{(x, type, id)}.
	\item La estructura auxiliar (set, multiset, segtree) se mantiene sobre el otro eje.
	\item Procesar eventos en orden de $x$ ascendente.
\end{itemize}

\paragraph{Complejidad.}
\begin{itemize}\setlength\itemsep{0pt}
	\item Depende del problema; tipico $O((N + E) \log N)$.
\end{itemize}

\paragraph{Donde tocar para modificar.}
\begin{itemize}\setlength\itemsep{0pt}
	\item \textbf{Aniadir/quitar segmentos}: aniadir eventos add/remove al barrido.
	\item \textbf{Aniadir queries}: aniadir eventos query.
	\item \textbf{Cambiar el eje de barrido}: si barres en $y$ en vez de $x$, aniadir comparador.
\end{itemize}

\paragraph{Trampas y errores frecuentes.}
\begin{itemize}\setlength\itemsep{0pt}
	\item El snippet es solo el \textbf{esqueleto}; los handlers \texttt{add/query/remove} se llenan por problema.
	\item Cuidado con eventos duplicados.
\end{itemize}

\paragraph{Codigo.}
Ver \texttt{cpp.tex}, seccion \textit{Sweep line template}.

\vspace{0.5em}

\subsubsection{Closest pair of points}

\paragraph{Idea intuitiva.}
Encuentra el par de puntos mas cercano en $O(n \log n)$ con \textbf{divide y venceras}. Algoritmo:
\begin{itemize}\setlength\itemsep{0pt}
	\item Sort por $x$.
	\item Dividir en mitades; recursivo.
	\item Combinar: dado el minimo $d$ de las dos mitades, strip = puntos en $[midX - d, midX + d]$; ordenar por $y$; cada punto solo se compara con los siguientes hasta $y$ difiera en $< d$ (maximo 7).
\end{itemize}

\paragraph{Propiedades clave (invariantes).}
\begin{itemize}\setlength\itemsep{0pt}
	\item Requiere ordenar por $y$ despues (en el snippet se hace durante el merge).
	\item Strip tiene a lo sumo $O(n)$ puntos en total, pero la comparacion es $O(1)$ amortizada por punto.
	\item Resultado en \texttt{double}.
\end{itemize}

\paragraph{Complejidad.}
\begin{itemize}\setlength\itemsep{0pt}
	\item $O(n \log n)$.
\end{itemize}

\paragraph{Donde tocar para modificar.}
\begin{itemize}\setlength\itemsep{0pt}
	\item \textbf{Reconstruir el par}: aniadir indices en la recursion.
	\item \textbf{Distancia maxima}: cambiar \texttt{min} por \texttt{max} (mas facil; no requiere D\&C).
	\item \textbf{Puntos duplicados}: aniadir check especial (distancia 0).
\end{itemize}

\paragraph{Trampas y errores frecuentes.}
\begin{itemize}\setlength\itemsep{0pt}
	\item El merge requiere ordenar por $y$; el snippet lo hace in-place con un \texttt{tmp}.
	\item Si se hace en cada recursion, es $O(n \log n)$ extra.
\end{itemize}

\paragraph{Codigo.}
Ver \texttt{cpp.tex}, seccion \textit{Closest pair of points}.

\vspace{0.5em}

\subsubsection{Rotating calipers (polygon diameter)}

\paragraph{Idea intuitiva.}
Encuentra el \textbf{diametro} de un poligono convexo (par de vertices mas distantes) en $O(n)$. Asume poligono en orden counter-clockwise. Algoritmo: dos indices $i$ (sobre las aristas) y $k$ (sobre los vertices opuestos); en cada paso, rotar $k$ mientras el area del paralelogramo aumenta. La distancia $i$-$k$ puede ser maxima en cualquier momento del barrido.

\paragraph{Propiedades clave (invariantes).}
\begin{itemize}\setlength\itemsep{0pt}
	\item Cada indice avanza como maximo $n$ veces total.
	\item El ``diametro'' es la maxima distancia entre dos vertices.
	\item Retorna \textbf{distancia al cuadrado}.
\end{itemize}

\paragraph{Complejidad.}
\begin{itemize}\setlength\itemsep{0pt}
	\item $O(n)$.
\end{itemize}

\paragraph{Donde tocar para modificar.}
\begin{itemize}\setlength\itemsep{0pt}
	\item \textbf{Reconstruir el par de vertices}: aniadir indices.
	\item \textbf{Anchura minima / maxima}: variantes de calipers.
	\item \textbf{Distancia minima entre poligonos convexos}: otra variante.
\end{itemize}

\paragraph{Trampas y errores frecuentes.}
\begin{itemize}\setlength\itemsep{0pt}
	\item Asume poligono \textbf{estrictamente} convexo (sin vertices colineales). Si hay colineales, aniadir perturbacion o filtrar.
\end{itemize}

\paragraph{Codigo.}
Ver \texttt{cpp.tex}, seccion \textit{Rotating calipers (polygon diameter)}.

% ============================================================
\section{Miscelanea}
% ============================================================

\subsection{Utilidades del usuario}

\subsubsection{Hora}

\paragraph{Idea intuitiva.}
Una pequena clase para representar \textbf{tiempos en segundos} y hacer aritmetica basica sobre ellos. Internamente almacena todo en segundos (un solo \texttt{int}), lo cual simplifica comparaciones y diferencias. El constructor acepta $(h, m, s)$ y los convierte a segundos. Hay un metodo \texttt{cut()} que fuerza un minimo (util para ``horas trabajadas'').

\paragraph{Propiedades clave (invariantes).}
\begin{itemize}\setlength\itemsep{0pt}
	\item Representacion canonica: \textbf{solo segundos}. No hay zonas horarias, no hay DST.
	\item Comparadores \texttt{<} y \texttt{>} operan sobre los segundos.
	\item \texttt{operator-} retorna la \textbf{diferencia absoluta} en segundos (siempre no negativa).
	\item \texttt{cut()} aplica un piso (por defecto $8\text{h }30\text{m} = 30600$ segundos).
\end{itemize}

\paragraph{Complejidad.}
\begin{itemize}\setlength\itemsep{0pt}
	\item $O(1)$ por operacion.
\end{itemize}

\paragraph{Donde tocar para modificar.}
\begin{itemize}\setlength\itemsep{0pt}
	\item \textbf{Aniadir suma}: \texttt{operator+} que devuelve un \texttt{Hour} con la suma de segundos.
	\item \textbf{Aniadir output formateado}: aniadir \texttt{show\_hms()} que imprime $H$:$M$:$S$.
	\item \textbf{Cambiar el ``corte''}: modificar el valor magico \texttt{30'600} en \texttt{cut()} o pasarlo como parametro.
	\item \textbf{Tamano maximo}: si los tiempos exceden $\sim 68$ anos, el \texttt{int} desborda; usar \texttt{long long}.
	\item \textbf{Precision sub-segundo}: aniadir \texttt{ms} (milisegundos) y guardar en \texttt{long long} con factor 1000.
\end{itemize}

\paragraph{Trampas y errores frecuentes.}
\begin{itemize}\setlength\itemsep{0pt}
	\item El operador \texttt{-} retorna \textbf{valor absoluto}, no diferencia con signo. Si necesitas diferencia con signo, aniadir un nuevo metodo o sobrecargar.
	\item \texttt{cut()} modifica \texttt{s} \textbf{in-place}; si el llamador no espera esto, aniadir una version \texttt{cut(min)} que retorna un nuevo \texttt{Hour}.
	\item El umbral magico \texttt{30'600} (8h 30m) esta hardcoded; documentar bien que es ``jornada minima''.
\end{itemize}

\paragraph{Codigo.}
Ver \texttt{cpp.tex}, seccion \textit{Hora}.

\end{document}
' se responden en $O(|P|)$.
\end{itemize}

\paragraph{Trampas y errores frecuentes.}
\begin{itemize}\setlength\itemsep{0pt}
	\item El \textbf{clone} debe tener \texttt{occ = 0} (no es una nueva posicion final); el \texttt{link} del estado siendo creado y del estado clonado deben apuntar al clon.
	\item Saltarse esto da SAM incorrecto.
	\item La condicion \texttt{st[p].len + 1 == st[q].len} detecta cuando no se necesita clon.
	\item Olvidar \texttt{while (p != -1 && !st[p].next.count(c))} da transiciones rotas.
\end{itemize}

\paragraph{Codigo.}
Ver \texttt{cpp.tex}, seccion \textit{Suffix Automaton}.

\vspace{0.5em}

\subsubsection{Lyndon factorization (Duval)}

\paragraph{Idea intuitiva.}
Un \textbf{string de Lyndon} es un string que es estrictamente menor que todos sus rotaciones. Toda cadena tiene una \textbf{factorizacion unica} de Lyndon (algoritmo de \textbf{Duval}): una secuencia de strings de Lyndon $L_1, L_2, \dots, L_k$ tales que $L_1 \ge L_2 \ge \dots \ge L_k$ y la concatenacion es el string original. Esto da $O(n)$ para encontrar los cortes. La demostracion de correctitud: cada factor es minimal en su clase de rotacion, asi que la concatenacion respeta el orden lexicografico.

\paragraph{Propiedades clave (invariantes).}
\begin{itemize}\setlength\itemsep{0pt}
	\item Cada factor de Lyndon es de la forma $uv$ donde $u$ se repite y $v$ es prefijo de $u$.
	\item La longitud del factor esta dada por la ``longitud del periodo''.
	\item Aplicacion: el ``Lyndon word'' es el menor en orden lexicografico entre todas las rotaciones.
	\item La factorizacion es \textbf{unica}.
\end{itemize}

\paragraph{Codigo clave.}
El algoritmo de Duval (iterativo):
\begin{verbatim}
int n = s.size();
int i = 0;
while (i < n) {
    int j = i + 1, k = i;
    while (j < n && s[k] <= s[j]) {
        if (s[k] < s[j]) k = i;
        else ++k;
        ++j;
    }
    int len = j - k;
    while (i < j) { cout << "factor: " << s.substr(i, len) << "\n"; i += len; }
}
\end{verbatim}

\paragraph{Complejidad.}
\begin{itemize}\setlength\itemsep{0pt}
	\item $O(n)$.
	\item Cada comparacion es $O(1)$; el puntero $k$ retrocede a lo sumo $n$ veces en total.
\end{itemize}

\paragraph{Donde tocar para modificar.}
\begin{itemize}\setlength\itemsep{0pt}
	\item \textbf{Encontrar la menor rotacion}: el menor sufijo lexicografico de \texttt{s + s} es la menor rotacion; la posicion se obtiene con un solo pase de Duval.
	\item \textbf{Algoritmo de Booth}: alternativa para rotacion minima, tambien $O(n)$.
	\item \textbf{Validar}: si necesitas verificar que un string es de Lyndon, compararlo con sus rotaciones.
	\item \textbf{Encontrar todos los factores}: guardar la posicion y longitud de cada uno.
\end{itemize}

\paragraph{Trampas y errores frecuentes.}
\begin{itemize}\setlength\itemsep{0pt}
	\item La condicion \texttt{s[j] <= s[k]} (no \texttt{<}) en el while de Duval es crucial; cambiar a \texttt{<} da comportamiento distinto.
	\item El caso $s[k] < s[j]$ resetea $k$ a $i$ (inicio del factor actual).
	\item El bucle externo debe avanzar $i$ por la longitud del factor, no por 1.
\end{itemize}

\paragraph{Codigo.}
Ver \texttt{cpp.tex}, seccion \textit{Lyndon factorization (Duval)}.

\subsection{Palindromos}

\subsubsection{Manacher}

\paragraph{Idea intuitiva.}
Calcula en $O(n)$ los \textbf{radios} de todos los palindromos centrados en cada posicion (impares y pares). La idea: mantener una ventana $[l, r]$ que es el palindromo centrado en un punto anterior mas largo encontrado. Si el nuevo centro $i$ cae dentro de $[l, r]$, entonces el palindromo centrado en $i$ tiene al menos el radio del palindromo ``espejo'' $l + r - i$; sino, empezar desde 1 (impar) o 0 (par). Despues, expandir lo que falte. La demostracion de $O(n)$: cada caracter se compara a lo sumo 2 veces (extension + limite por ventana).

\paragraph{Propiedades clave (invariantes).}
\begin{itemize}\setlength\itemsep{0pt}
	\item Conveccion del snippet: \texttt{odd[i]} = $k$ tal que el palindromo impar mas largo centrado en $i$ tiene longitud $2k - 1$. \texttt{even[i]} = $k$ para par centrado entre $i-1$ e $i$, longitud $2k$.
	\item $k$ cuenta tambien cuantos palindromos hay centrados ahi.
	\item Las queries \texttt{isPalindrome(l, r)} son $O(1)$.
	\item La ventana $[l, r]$ siempre contiene el palindromo mas largo encontrado hasta ahora.
\end{itemize}

\paragraph{Codigo clave.}
Manacher para palindromos impares:
\begin{verbatim}
int l = 0, r = -1;
for (int i = 0; i < n; ++i) {
    int k = (i > r) ? 1 : min(odd[l + r - i], r - i + 1);
    while (i - k >= 0 && i + k < n && s[i-k] == s[i+k]) ++k;
    odd[i] = k;
    if (i + k - 1 > r) { l = i - k + 1; r = i + k - 1; }
}
\end{verbatim}

\paragraph{Complejidad.}
\begin{itemize}\setlength\itemsep{0pt}
	\item Construccion $O(n)$, queries $O(1)$.
	\item Cada extension se cuenta una sola vez.
\end{itemize}

\paragraph{Donde tocar para modificar.}
\begin{itemize}\setlength\itemsep{0pt}
	\item \textbf{Solo contar palindromos}: ya implementado en \texttt{countPalindromes}.
	\item \textbf{Longest palindromic substring}: ya en \texttt{longestPalindrome}.
	\item \textbf{Modificar la conveccion de radios}: si queres ``radio = mitad de la longitud'' en vez de ``k palindromos'', ajustar el offset.
	\item \textbf{Queries con k no precomputado}: aniadir helpers especificos.
	\item \textbf{Modificar el caracter de relleno}: para multiples separadores, aniadir un caracter unico.
\end{itemize}

\paragraph{Trampas y errores frecuentes.}
\begin{itemize}\setlength\itemsep{0pt}
	\item La conveccion del snippet es \textbf{distinta} de la ``clasica de radio'' (donde radio = $(L-1)/2$); tener cuidado al copiar a otro codigo.
	\item Si los palindromos se solapan, la condicion \texttt{i + k - 1 > r} actualiza la ventana.
	\item En palindromos pares, el centro esta ``entre'' dos indices; aniadir offset en queries.
\end{itemize}

\paragraph{Codigo.}
Ver \texttt{cpp.tex}, seccion \textit{Manacher}.

% ============================================================
\section{Grafos}
% ============================================================

\subsection{Caminos minimos}

\subsubsection{Dijkstra}

\paragraph{Idea intuitiva.}
Encuentra el shortest path desde un origen $s$ a todos los vertices en $O((V + E) \log V)$ asumiendo \textbf{pesos no negativos}. La idea: usar una cola de prioridad (min-heap) ordenada por distancia tentativa. En cada paso, extraer el nodo con menor distancia; si ya fue procesado, saltar. Si al actualizar un vecino la distancia mejora, insertar/actualizar en la cola. La demostracion: cuando un nodo se extrae por primera vez, no hay forma de llegar a el con menor distancia (cualquier camino alternativo deberia pasar por nodos no procesados, con distancia $\ge$ la actual).

\paragraph{Propiedades clave (invariantes).}
\begin{itemize}\setlength\itemsep{0pt}
	\item Asume pesos $\ge 0$; con pesos negativos usar Bellman-Ford.
	\item Greedy correctness: la primera vez que se extrae un nodo, su distancia es la definitiva.
	\item Si se quiere reconstruir el camino, guardar el \texttt{parent[]}.
	\item Cada nodo se procesa a lo sumo $V$ veces (mas relajarions, pero solo $V$ definitivas).
\end{itemize}

\paragraph{Codigo clave.}
El bucle principal de Dijkstra:
\begin{verbatim}
priority_queue<pair<long long, int>, vector<...>, greater<...>> pq;
pq.push({0, source});
while (!pq.empty()) {
    auto [d, v] = pq.top(); pq.pop();
    if (d > dist[v]) continue;  // entrada vieja
    for (auto [w, u] : adj[v])
        if (dist[v] + w < dist[u]) {
            dist[u] = dist[v] + w;
            pq.push({dist[u], u});
        }
}
\end{verbatim}

\paragraph{Complejidad.}
\begin{itemize}\setlength\itemsep{0pt}
	\item $O((V + E) \log V)$ con heap binario, $O(V + E)$ con heap de Fibonacci.
	\item En la practica, $\sim 1.5 \cdot 10^{-7}$ segundos por operacion para $V = 10^5$.
\end{itemize}

\paragraph{Donde tocar para modificar.}
\begin{itemize}\setlength\itemsep{0pt}
	\item \textbf{Camino mas corto entre todos los pares}: si se necesita desde todos los origenes, usar Floyd o ejecutar Dijkstra $V$ veces.
	\item \textbf{Shortest path con aristas negativas limitadas}: Bellman-Ford.
	\item \textbf{Reconstruir camino}: aniadir \texttt{parent[v]} actualizado junto a \texttt{dist[v]}.
	\item \textbf{Dijkstra con potentials (Johnson)}: para aristas negativas si se puede reescalar.
	\item \textbf{Multi-source}: aniadir varios nodos fuentes iniciales al heap.
\end{itemize}

\paragraph{Trampas y errores frecuentes.}
\begin{itemize}\setlength\itemsep{0pt}
	\item Distancias no inicializadas a \texttt{INF}; el snippet usa \texttt{1e18}.
	\item El \texttt{if(dist[adjN] != 1e18)} para borrar de la cola es importante si usas \texttt{set} (no asi con \texttt{priority\_queue}).
	\item Con pesos grandes, \texttt{dist[v] + w} puede overflow; usar \texttt{LLONG\_MAX/2} como INF.
	\item No usar con pesos negativos: el algoritmo puede dar respuestas incorrectas.
\end{itemize}

\paragraph{Codigo.}
Ver \texttt{cpp.tex}, seccion \textit{Dijkstra}.

\vspace{0.5em}

\subsubsection{Bellman-Ford}

\paragraph{Idea intuitiva.}
Encuentra shortest path desde $s$ incluso con \textbf{pesos negativos} en $O(VE)$. La idea: relajar todas las aristas $V - 1$ veces. Cada relajacion propaga la mejor distancia a ``un paso mas''. Despues de $V - 1$ pasadas, si alguna arista aun se puede relajar, hay un \textbf{ciclo negativo} alcanzable. La demostracion: el camino mas corto tiene a lo sumo $V - 1$ aristas (sin ciclos); tras $V - 1$ relajaciones todas las distancias estan en su valor optimo.

\paragraph{Propiedades clave (invariantes).}
\begin{itemize}\setlength\itemsep{0pt}
	\item Funciona con pesos negativos (sin ciclos negativos).
	\item Tras $V - 1$ relajaciones, las distancias son definitivas (si no hay ciclo negativo).
	\item Una relajacion adicional detecta ciclos negativos: si \texttt{dist[u] + w < dist[v]}, hay ciclo negativo alcanzable desde $s$.
	\item Cada relajacion solo mejora distancias, asi que termina (no oscila).
\end{itemize}

\paragraph{Codigo clave.}
El doble bucle de relajacion:
\begin{verbatim}
for (int i = 0; i < n - 1; ++i)
    for (auto [u, v, w] : edges)
        if (dist[u] + w < dist[v])
            dist[v] = dist[u] + w;
// Detectar ciclo negativo
for (auto [u, v, w] : edges)
    if (dist[u] + w < dist[v]) { /* ciclo negativo */ }
\end{verbatim}

\paragraph{Complejidad.}
\begin{itemize}\setlength\itemsep{0pt}
	\item $O(VE)$.
	\item SPFA (con cola) en promedio es mucho mas rapido, peor caso igual.
\end{itemize}

\paragraph{Donde tocar para modificar.}
\begin{itemize}\setlength\itemsep{0pt}
	\item \textbf{SPFA}: optimizacion que relaja solo nodos cuya distancia cambio; en el peor caso sigue siendo $O(VE)$.
	\item \textbf{Detectar ciclo negativo global}: ejecutar Bellman-Ford desde un super-origen conectado a todos.
	\item \textbf{Camino mas corto con K aristas exactas}: $K$ iteraciones de Bellman-Ford.
	\item \textbf{Johnson's}: reescalar pesos con potentials para usar Dijkstra.
\end{itemize}

\paragraph{Trampas y errores frecuentes.}
\begin{itemize}\setlength\itemsep{0pt}
	\item Overflow si los pesos son $\ge 2^{52}$; usar \texttt{long long} y \texttt{\_\_int128} si hace falta.
	\item Si el grafo es denso ($E \approx V^2$), $O(VE)$ es prohibitivo.
	\item El chequeo de ciclo negativo debe hacerse tras las $V - 1$ iteraciones, no durante.
	\item Distancias no inicializadas a \texttt{INF}; el primer nodo (source) tiene distancia 0.
\end{itemize}

\paragraph{Codigo.}
Ver \texttt{cpp.tex}, seccion \textit{Bellman-Ford}.

\vspace{0.5em}

\subsubsection{Floyd-Warshall}

\paragraph{Idea intuitiva.}
Shortest path entre \textbf{todos los pares} de vertices en $O(V^3)$. DP: $d[i][j]$ = shortest path usando solo nodos intermedios en $\{0, \dots, k\}$. Transicion: $d_{k+1}[i][j] = \min(d_k[i][j], d_k[i][k+1] + d_k[k+1][j])$. La demostracion: tras $k$ iteraciones, $d[i][j]$ es el shortest path con intermedios en $\{0, \dots, k-1\}$; aniadir el nodo $k$ da la recurrencia clasica.

\paragraph{Propiedades clave (invariantes).}
\begin{itemize}\setlength\itemsep{0pt}
	\item Detecta ciclos negativos: tras el algoritmo, $d[i][i] < 0$ indica ciclo negativo alcanzable desde $i$.
	\item Funciona con pesos negativos (sin ciclos negativos).
	\item Memoria $O(V^2)$; con $V \le 400$ es OK, con mas conviene Dijkstra $V$ veces.
	\item Es la unica opcion cuando se quieren shortest paths entre todos los pares y el grafo es denso.
\end{itemize}

\paragraph{Codigo clave.}
El triple bucle clasico:
\begin{verbatim}
for (int k = 0; k < n; ++k)
    for (int i = 0; i < n; ++i)
        for (int j = 0; j < n; ++j)
            d[i][j] = min(d[i][j], d[i][k] + d[k][j]);
// detectar ciclo negativo
for (int i = 0; i < n; ++i) if (d[i][i] < 0) tiene_ciclo_negativo = true;
\end{verbatim}

\paragraph{Complejidad.}
\begin{itemize}\setlength\itemsep{0pt}
	\item $O(V^3)$ tiempo, $O(V^2)$ memoria.
	\item Para $V \le 400$: $\sim 6.4 \cdot 10^7$ operaciones (factible).
\end{itemize}

\paragraph{Donde tocar para modificar.}
\begin{itemize}\setlength\itemsep{0pt}
	\item \textbf{Reconstruir camino}: aniadir \texttt{next[i][j]} y backtrack.
	\item \textbf{Transitividad}: el algoritmo tambien calcula transitividad (\texttt{d[i][j] != INF} significa alcanzable).
	\item \textbf{Pesos doubles}: aniadir tolerancia (\texttt{d[i][j] > d[i][k] + d[k][j] - eps}).
	\item \textbf{Grafo no denso}: usar Dijkstra $V$ veces en vez de Floyd.
\end{itemize}

\paragraph{Trampas y errores frecuentes.}
\begin{itemize}\setlength\itemsep{0pt}
	\item Para detectar ciclo negativo correctamente, aniadir un chequeo \texttt{if (ans[i][i] < 0)} tras el triple bucle.
	\item El orden de los bucles \texttt{k, i, j} es importante; invertir da resultados diferentes.
	\item Con \texttt{double}, usar tolerancia \texttt{< eps} en vez de \texttt{<} para evitar loops infinitos.
\end{itemize}

\paragraph{Codigo.}
Ver \texttt{cpp.tex}, seccion \textit{Floyd-Warshall}.

\subsection{Componentes y cortes}

\subsubsection{Kosaraju SCC}

\paragraph{Idea intuitiva.}
Encuentra las \textbf{componentes fuertemente conexas} (SCCs) en $O(V + E)$. Algoritmo en 3 pasos:
\begin{itemize}\setlength\itemsep{0pt}
	\item DFS en $G$ rellenando una pila en \textbf{post-orden} (los nodos que terminan tarde quedan arriba).
	\item Construir $G^T$ (grafo transpuesto).
	\item DFS en $G^T$ procesando nodos en orden de la pila (de mayor a menor tiempo); cada DFS da una SCC.
\end{itemize}
La demostracion: dos nodos $u, v$ estan en la misma SCC sii hay camino $u \leadsto v$ y $v \leadsto u$; Kosaraju captura exactamente esto al explorar $G^T$ desde los nodos con mayor tiempo de finalizacion en $G$.

\paragraph{Propiedades clave (invariantes).}
\begin{itemize}\setlength\itemsep{0pt}
	\item Cada SCC es maximal fuertemente conexa.
	\item El \textbf{grafo de SCCs} (metagraf) es un DAG.
	\item Kosaraju y Tarjan dan el mismo resultado; Tarjan es ``mas rapido'' en practica (un solo DFS).
	\item Cada nodo pertenece a exactamente una SCC.
\end{itemize}

\paragraph{Codigo clave.}
La estructura de Kosaraju:
\begin{verbatim}
vector<int> order;
vector<bool> used(n);
function<void(int)> dfs1 = [&](int v) {
    used[v] = true;
    for (int u : g[v]) if (!used[u]) dfs1(u);
    order.push_back(v);  // post-orden
};
function<void(int, int)> dfs2 = [&](int v, int root) {
    comp[v] = root;
    for (int u : gt[v]) if (comp[u] == -1) dfs2(u, root);
};
// 1: DFS en G para llenar order
// 2: DFS en G^T en orden inverso para asignar SCCs
\end{verbatim}

\paragraph{Complejidad.}
\begin{itemize}\setlength\itemsep{0pt}
	\item $O(V + E)$.
	\item Constante 2 (dos DFS + el grafo transpuesto).
\end{itemize}

\paragraph{Donde tocar para modificar.}
\begin{itemize}\setlength\itemsep{0pt}
	\item \textbf{2-SAT}: SCC sobre grafo de implicaciones.
	\item \textbf{Construir el DAG de SCCs}: aniadir aristas entre SCCs diferentes.
	\item \textbf{Tamano de cada SCC}: aniadir contador durante el segundo DFS.
	\item \textbf{Grafo grande}: usar Tarjan para evitar construir $G^T$.
\end{itemize}

\paragraph{Trampas y errores frecuentes.}
\begin{itemize}\setlength\itemsep{0pt}
	\item El snippet usa \texttt{vector<...> adj[n]} (VLA); reemplazable por \texttt{vector<vector<pair<int,int>>} para portabilidad.
	\item El orden del segundo DFS debe ser \textbf{inverso} al de finalizacion del primero.
	\item Si el grafo no es conexo, aniadir un loop sobre todos los nodos en el primer DFS.
	\item Confundir SCC con componentes conexas (las primeras son para dirigidos, las segundas para no dirigidos).
\end{itemize}

\paragraph{Codigo.}
Ver \texttt{cpp.tex}, seccion \textit{Kosaraju SCC}.

\vspace{0.5em}

\subsubsection{Tarjan (puentes)}

\paragraph{Idea intuitiva.}
Encuentra los \textbf{puentes} (bridges): aristas cuya remocion \textbf{desconecta} el grafo. Usa un solo DFS manteniendo \texttt{tin[v]} (tiempo de descubrimiento) y \texttt{low[v]} (el menor \texttt{tin} alcanzable desde $v$ via back-edges). Una arista $v - u$ es puente si $low[u] > tin[v]$. La demostracion: si $low[u] > tin[v]$, entonces $u$ (y su subarbol) no pueden ``escapar'' de $v$; quitar la arista $v - u$ parte el grafo.

\paragraph{Propiedades clave (invariantes).}
\begin{itemize}\setlength\itemsep{0pt}
	\item \texttt{low[v] = min(tin[v], tin[w])} para todo back-edge $(v, w)$, y \texttt{min(low[u])} para hijos en el DFS.
	\item Solo se consideran las aristas del \textbf{arbol DFS} (no back-edges) para \texttt{low}.
	\item Cada nodo tiene exactamente un \texttt{tin}; cada back-edge se procesa una vez.
\end{itemize}

\paragraph{Codigo clave.}
El DFS de Tarjan para bridges:
\begin{verbatim}
void dfs(int v, int pEdge) {
    used[v] = true;
    tin[v] = low[v] = timer++;
    for (auto [u, id] : adj[v]) {
        if (id == pEdge) continue;  // no contar el edge al padre
        if (used[u]) {
            low[v] = min(low[v], tin[u]);  // back-edge
        } else {
            dfs(u, id);
            low[v] = min(low[v], low[u]);
            if (low[u] > tin[v]) bridges.push_back(id);  // es puente
        }
    }
}
\end{verbatim}

\paragraph{Complejidad.}
\begin{itemize}\setlength\itemsep{0pt}
	\item $O(V + E)$.
	\item Una sola pasada de DFS.
\end{itemize}

\paragraph{Donde tocar para modificar.}
\begin{itemize}\setlength\itemsep{0pt}
	\item \textbf{Bridge tree}: contractar cada 2-edge-connected component a un nodo; el resultado es un arbol.
	\item \textbf{Articulation points}: variante con condicion diferente (ver el siguiente modulo).
	\item \textbf{Grafo no dirigido}: solo funciona en no dirigidos; en dirigidos hay conceptos analogos (strong bridges, dominator tree).
	\item \textbf{Reconstruir componentes 2-edge-connected}: BFS desde cada nodo excluyendo bridges.
\end{itemize}

\paragraph{Trampas y errores frecuentes.}
\begin{itemize}\setlength\itemsep{0pt}
	\item No confundir el parent con un back-edge al padre; chequear \texttt{if (adjN == parent) continue;} y contar solo las veces que se ``visita'' como hijo.
	\item Para multigrafos, los edges deben tener IDs unicos para distinguirlos.
	\item Si el grafo no es conexo, aniadir loop externo sobre todos los nodos.
\end{itemize}

\paragraph{Codigo.}
Ver \texttt{cpp.tex}, seccion \textit{Tarjan (puentes)}.

\vspace{0.5em}

\subsubsection{Tarjan (puntos de articulacion)}

\paragraph{Idea intuitiva.}
Encuentra los \textbf{puntos de articulacion} (cut vertices): vertices cuya remocion \textbf{desconecta} el grafo. Mismo setup que para bridges: \texttt{tin[v]}, \texttt{low[v]}. Un nodo $v$ (no raiz) es de articulacion si existe un hijo $u$ con $low[u] \ge tin[v]$. La \textbf{raiz} es de articulacion si tiene $\ge 2$ hijos en el DFS tree. La razon: para $v$ no raiz, la condicion $\ge$ indica que el subarbol de $u$ depende de $v$; quitar $v$ los desconecta. Para la raiz, perder $\ge 2$ hijos parte el DFS en $\ge 2$ componentes.

\paragraph{Propiedades clave (invariantes).}
\begin{itemize}\setlength\itemsep{0pt}
	\item $low[u] \ge tin[v]$ significa que el subarbol de $u$ no tiene back-edge ``por encima'' de $v$.
	\item La raiz es caso especial: tiene que perder mas de un hijo para que el grafo se parta.
	\item Cada nodo se evalua una sola vez.
\end{itemize}

\paragraph{Codigo clave.}
La condicion de articulacion:
\begin{verbatim}
void dfs(int v, int p) {
    used[v] = true;
    tin[v] = low[v] = timer++;
    int children = 0;
    for (int u : adj[v]) {
        if (u == p) continue;
        if (used[u]) low[v] = min(low[v], tin[u]);
        else {
            dfs(u, v);
            low[v] = min(low[v], low[u]);
            if (low[u] >= tin[v] && p != -1)  // no raiz
                isArt[v] = true;
            ++children;
        }
    }
    if (p == -1 && children > 1) isArt[v] = true;  // raiz especial
}
\end{verbatim}

\paragraph{Complejidad.}
\begin{itemize}\setlength\itemsep{0pt}
	\item $O(V + E)$.
	\item Una sola pasada de DFS.
\end{itemize}

\paragraph{Donde tocar para modificar.}
\begin{itemize}\setlength\itemsep{0pt}
	\item \textbf{Block-cut tree}: estructura bipartita con BCCs (2-vertex-connected components) y articulation points.
	\item \textbf{Verificar biconectividad}: chequear que no hay puntos de articulacion.
	\item \textbf{Componentes biconnected}: extender Tarjan para enumerar BCCs.
\end{itemize}

\paragraph{Trampas y errores frecuentes.}
\begin{itemize}\setlength\itemsep{0pt}
	\item No olvidar el caso especial de la raiz.
	\item La condicion usa \texttt{>=}, no \texttt{>}; un hijo con \texttt{low[u] == tin[v]} tambien cuenta.
	\item Multigrafos pueden tener ``falsos'' articulation points; verificar caso a caso.
\end{itemize}

\paragraph{Codigo.}
Ver \texttt{cpp.tex}, seccion \textit{Tarjan (puntos de articulacion)}.

\vspace{0.5em}

\subsubsection{Tarjan SCC}

\paragraph{Idea intuitiva.}
Variante del SCC que usa \textbf{un solo DFS} con una pila explicita. Mantiene \texttt{disc[v]}, \texttt{low[v]}, y una pila de nodos ``en la SCC actual''. Cuando \texttt{low[v] == disc[v]}, $v$ es la raiz de una SCC; popear hasta $v$ inclusive. La demostracion: $low[v] == disc[v]$ indica que $v$ no tiene back-edges a ancestros; entonces $v$ y los nodos en la pila hasta $v$ forman una SCC maximal.

\paragraph{Propiedades clave (invariantes).}
\begin{itemize}\setlength\itemsep{0pt}
	\item Mas eficiente que Kosaraju (un solo DFS).
	\item \texttt{low[v] = min(low[u], disc[w])} para back-edges o recursion.
	\item La pila asegura que solo los nodos en la SCC actual se popean.
	\item Cada nodo se pushea y popea a lo sumo una vez.
\end{itemize}

\paragraph{Codigo clave.}
El DFS de Tarjan SCC:
\begin{verbatim}
void dfs(int v) {
    disc[v] = low[v] = timer++;
    st.push(v); inStack[v] = true;
    for (int u : g[v]) {
        if (disc[u] == -1) { dfs(u); low[v] = min(low[v], low[u]); }
        else if (inStack[u]) low[v] = min(low[v], disc[u]);
    }
    if (low[v] == disc[v]) {
        // popear hasta v
        while (true) {
            int u = st.top(); st.pop(); inStack[u] = false;
            comp[u] = current_scc;
            if (u == v) break;
        }
        ++current_scc;
    }
}
\end{verbatim}

\paragraph{Complejidad.}
\begin{itemize}\setlength\itemsep{0pt}
	\item $O(V + E)$.
	\item Constante similar a Kosaraju pero sin construir $G^T$.
\end{itemize}

\paragraph{Donde tocar para modificar.}
\begin{itemize}\setlength\itemsep{0pt}
	\item \textbf{Tamano de cada SCC}: aniadir contador durante el pop.
	\item \textbf{Grafo condensation}: aniadir aristas entre SCCs en el DAG resultante.
	\item \textbf{2-SAT}: aniadir la construccion del grafo de implicaciones.
\end{itemize}

\paragraph{Trampas y errores frecuentes.}
\begin{itemize}\setlength\itemsep{0pt}
	\item \texttt{inStack[u]} debe chequearse antes de usar \texttt{low[v] = min(low[v], disc[u])} en back-edges.
	\item La pila global es importante; sin ella no se puede extraer la SCC.
	\item Si el grafo es muy grande, recursion puede overflow; aniadir version iterativa.
\end{itemize}

\paragraph{Codigo.}
Ver \texttt{cpp.tex}, seccion \textit{Tarjan SCC}.

\subsection{Matching}

\subsubsection{Kuhn (bipartite matching)}

\paragraph{Idea intuitiva.}
Encuentra un \textbf{matching maximo} en un grafo bipartito en $O(VE)$. Para cada nodo $u$ en $L$, intentar encontrar un vecino $v$ en $R$ que este libre o cuyo match se pueda ``reorganizar'' recursivamente. Si se encuentra, matchear $u - v$. La idea: si $u$ quiere un vecino $v$ ya ocupado, intentamos ``mover'' el match de $v$ a otro nodo; recursivamente, esto encuentra un augmenting path si existe.

\paragraph{Propiedades clave (invariantes).}
\begin{itemize}\setlength\itemsep{0pt}
	\item \texttt{matchR[v]} = el nodo de $L$ matcheado a $v$ (o $-1$ si libre).
	\item \texttt{vis[]} se resetea \textbf{cada intento} de $u$.
	\item Si el matching es perfecto, $|L| = |R| = $ tamano del matching.
	\item Kuhn no garantiza maximalidad greedy; puede haber augmenting paths no encontrados.
\end{itemize}

\paragraph{Codigo clave.}
La recursion del augmenting path:
\begin{verbatim}
bool tryKuhn(int v) {
    if (vis[v]) return false;
    vis[v] = true;
    for (int u : g[v]) {
        if (matchR[u] == -1 || tryKuhn(matchR[u])) {
            matchR[u] = v;
            return true;
        }
    }
    return false;
}
// Por cada v en L:
fill(vis, vis+n, false);
if (tryKuhn(v)) ++matches;
\end{verbatim}

\paragraph{Complejidad.}
\begin{itemize}\setlength\itemsep{0pt}
	\item $O(VE)$ en el peor caso.
	\item Con Hopcroft-Karp: $O(\sqrt{V} \cdot E)$, mucho mejor para grafos densos.
\end{itemize}

\paragraph{Donde tocar para modificar.}
\begin{itemize}\setlength\itemsep{0pt}
	\item \textbf{Matching bipartito minimo (min vertex cover)}: Konig's theorem.
	\item \textbf{Matching general}: blossom algorithm (no incluido).
	\item \textbf{Aniadir capacidades}: si los nodos de $R$ tienen capacidad $> 1$, partir cada uno en capacidad copias.
	\item \textbf{Output del matching}: el array \texttt{matchR} da los pares finales.
\end{itemize}

\paragraph{Trampas y errores frecuentes.}
\begin{itemize}\setlength\itemsep{0pt}
	\item Olvidar resetear \texttt{vis} por intento; eso es lo que evita ``ciclos infinitos'' en la recursion.
	\item Para grafos grandes ($V > 10^4$), Kuhn puede ser muy lento; usar Hopcroft-Karp.
	\item Si el matching es maximo pero no perfecto, $|L|$ o $|R|$ puede ser mayor; el matching es solo maximo.
\end{itemize}

\paragraph{Codigo.}
Ver \texttt{cpp.tex}, seccion \textit{Kuhn (bipartite matching)}.

\vspace{0.5em}

\subsubsection{Hopcroft-Karp}

\paragraph{Idea intuitiva.}
Matching bipartito maximo en $O(\sqrt{V} E)$. Algoritmo en \textbf{fases}: en cada fase, hacer BFS para encontrar los \textbf{caminos de aumento mas cortos} desde nodos libres de $L$; luego DFS solo por esos caminos. Cada fase aumenta el matching en $\ge 1$ y toma $O(E)$. La demostracion: tras una fase, todos los augmenting paths tienen longitud $\ge$ la nueva minima; eso sucede cada $\sqrt{V}$ fases (longitud minima crece geometricamente).

\paragraph{Propiedades clave (invariantes).}
\begin{itemize}\setlength\itemsep{0pt}
	\item \texttt{dist[u]} = distancia en niveles desde un nodo libre de $L$.
	\item BFS termina cuando no hay mas camino a un nodo libre de $R$.
	\item DFS respeta los niveles (solo sigue \texttt{dist[matchV[v]] == dist[u] + 1}).
	\item Cada fase aumenta el matching en $\ge 1$ arista.
\end{itemize}

\paragraph{Codigo clave.}
El BFS de Hopcroft-Karp:
\begin{verbatim}
bool bfs() {
    queue<int> q;
    for (int v : L) {
        if (matchU[v] == -1) { dist[v] = 0; q.push(v); }
        else dist[v] = -1;
    }
    bool found = false;
    while (!q.empty()) {
        int v = q.front(); q.pop();
        for (int u : g[v]) {
            if (matchV[u] != -1 && dist[matchV[u]] == -1) {
                dist[matchV[u]] = dist[v] + 1;
                q.push(matchV[u]);
            }
            if (matchV[u] == -1) found = true;  // llegamos a libre
        }
    }
    return found;
}
\end{verbatim}

\paragraph{Complejidad.}
\begin{itemize}\setlength\itemsep{0pt}
	\item $O(\sqrt{V} E)$.
	\item Para $V = 10^4$, $E = 10^5$: $\sim 3 \cdot 10^7$ operaciones.
\end{itemize}

\paragraph{Donde tocar para modificar.}
\begin{itemize}\setlength\itemsep{0pt}
	\item \textbf{Reconstruir el matching}: ya viene en \texttt{matchU}/\texttt{matchV}.
	\item \textbf{Aniadir pesos}: Hungarian algorithm.
	\item \textbf{Densidad alta}: si $E \approx V^2$, sigue siendo util pero verificar constantes.
	\item \textbf{Aniadir capacidades}: partir los nodos con capacidad $> 1$ en varias copias.
\end{itemize}

\paragraph{Trampas y errores frecuentes.}
\begin{itemize}\setlength\itemsep{0pt}
	\item \texttt{dist[matchV[v]] == dist[u] + 1} (no \texttt{==} a secas); sin esto, el DFS no respeta los niveles y se pierde la complejidad.
	\item El BFS debe marcar los niveles correctamente; resetear \texttt{dist} para nodos alcanzables solo via matching.
	\item El DFS debe respetar los niveles; cualquier ``atajo'' invalida la cota.
\end{itemize}

\paragraph{Codigo.}
Ver \texttt{cpp.tex}, seccion \textit{Hopcroft-Karp}.

\subsection{Basicos}

\subsubsection{DFS / BFS / componentes}

\paragraph{Idea intuitiva.}
\textbf{DFS} (Depth-First Search) y \textbf{BFS} (Breadth-First Search) son los dos recorridos basicos de grafos. DFS usa una pila (o recursion) y explora tan profundo como puede antes de retroceder; BFS usa una cola y explora por ``niveles'' de distancia. La cantidad de \textbf{componentes conexas} se cuenta iniciando un BFS/DFS desde cada nodo no visitado. La diferencia clave: BFS da el camino mas corto en aristas; DFS da un orden de finalizacion util para problemas topologicos.

\paragraph{Propiedades clave (invariantes).}
\begin{itemize}\setlength\itemsep{0pt}
	\item BFS da la \textbf{distancia en aristas} mas corta desde el origen.
	\item DFS da un \textbf{orden topologico} (si se aniaden al final).
	\item Ambos son $O(V + E)$.
	\item Un grafo con $C$ componentes conexas satisface $\sum \text{subtree\_size}_v = V$ y $\sum E_v = 2E$.
	\item DFS en arbol da un spanning tree; las ``back edges'' corresponden a ciclos.
\end{itemize}

\paragraph{Codigo clave.}
BFS clasico:
\begin{verbatim}
queue<int> q;
q.push(source);
dist[source] = 0;
while (!q.empty()) {
    int v = q.front(); q.pop();
    for (int u : adj[v])
        if (dist[u] == -1) {
            dist[u] = dist[v] + 1;
            q.push(u);
        }
}
\end{verbatim}

\paragraph{Complejidad.}
\begin{itemize}\setlength\itemsep{0pt}
	\item $O(V + E)$ tiempo, $O(V)$ memoria.
	\item Cada arista se procesa exactamente 2 veces (ida y vuelta) en no dirigidos.
\end{itemize}

\paragraph{Donde tocar para modificar.}
\begin{itemize}\setlength\itemsep{0pt}
	\item \textbf{Aniadir pesos}: si las aristas tienen peso y queres shortest path, BFS ya no sirve; usar Dijkstra.
	\item \textbf{Grafo dirigido}: en DFS, aniadir check \texttt{if (state[u] == 1)} para detectar ciclos.
	\item \textbf{Bipartite check}: BFS alternando colores (ver el siguiente modulo).
	\item \textbf{Componentes fuertemente conexas (dirigidos)}: usar Tarjan o Kosaraju.
	\item \textbf{BFS en matriz}: aniadir offsets por direccion para grid 2D.
\end{itemize}

\paragraph{Trampas y errores frecuentes.}
\begin{itemize}\setlength\itemsep{0pt}
	\item Stack overflow con DFS recursivo para $V > 10^5$; usar version iterativa con pila explicita.
	\item En BFS, olvidar marcar \texttt{dist[u] = dist[v] + 1} antes de pushear da visitas duplicadas.
	\item Para grafos dirigidos, ``componente conexa'' no esta definida; usar SCC.
\end{itemize}

\paragraph{Codigo.}
Ver \texttt{cpp.tex}, seccion \textit{DFS / BFS / componentes}.

\vspace{0.5em}

\subsubsection{Topological sort (Kahn + DFS)}

\paragraph{Idea intuitiva.}
Un \textbf{orden topologico} de un DAG es una permutacion de los vertices tal que toda arista $u \to v$ tiene $u$ antes que $v$. Dos algoritmos clasicos:
\begin{itemize}\setlength\itemsep{0pt}
	\item \textbf{Kahn}: BFS sobre nodos con \texttt{indeg == 0}; al ``consumir'' un nodo, decrementar el indeg de sus vecinos; los que lleguen a 0 se encolan.
	\item \textbf{DFS}: hacer DFS; al terminar un nodo (estado = 2), aniadirlo al final; el orden es el reverso.
\end{itemize}
La demostracion: en un DAG siempre hay al menos un nodo con indeg 0; quitarlo mantiene el DAG; por induccion, se obtiene el orden.

\paragraph{Propiedades clave (invariantes).}
\begin{itemize}\setlength\itemsep{0pt}
	\item Existe sii el grafo es un DAG.
	\item Si al final \texttt{order.size() < V}, hay un ciclo.
	\item Kahn da el orden ``natural'' (los que pueden ir primero); DFS da el orden ``al reves'' (los que pueden ir al final primero).
	\item La cantidad de ordenes topologicos puede ser exponencial.
\end{itemize}

\paragraph{Codigo clave.}
Kahn con priority queue (orden lexicografico):
\begin{verbatim}
priority_queue<int, vector<int>, greater<int>> pq;
for (int v = 0; v < n; ++v) if (indeg[v] == 0) pq.push(v);
while (!pq.empty()) {
    int v = pq.top(); pq.pop();
    order.push_back(v);
    for (int u : adj[v]) {
        if (--indeg[u] == 0) pq.push(u);
    }
}
\end{verbatim}

\paragraph{Complejidad.}
\begin{itemize}\setlength\itemsep{0pt}
	\item $O(V + E)$ ambos.
	\item Con priority queue, $O((V + E) \log V)$.
\end{itemize}

\paragraph{Donde tocar para modificar.}
\begin{itemize}\setlength\itemsep{0pt}
	\item \textbf{Detectar ciclos}: comparar \texttt{order.size()} con \texttt{V}.
	\item \textbf{Orden lexicografico minimo}: usar priority queue en Kahn.
	\item \textbf{Topological sort con pesos}: combinar con shortest path en DAG (un solo pase en orden topologico).
	\item \textbf{Encontrar todos los ordenes}: DP sobre subsets (solo para $V$ pequeno).
\end{itemize}

\paragraph{Trampas y errores frecuentes.}
\begin{itemize}\setlength\itemsep{0pt}
	\item En el snippet, Kahn retorna orden vacio si hay ciclo (util para detectar).
	\item En DFS, no olvidar \texttt{reverse(order)}.
	\item Confundir indeg con outdeg; el orden Kahn quita primero los nodos ``fuente''.
	\item Para grafos no DAG, el algoritmo termina con \texttt{order.size() < V}.
\end{itemize}

\paragraph{Codigo.}
Ver \texttt{cpp.tex}, seccion \textit{Topological sort (Kahn + DFS)}.

\vspace{0.5em}

\subsubsection{Bipartite check + 2-coloring}

\paragraph{Idea intuitiva.}
Un grafo es \textbf{bipartito} si sus vertices se pueden dividir en dos conjuntos $L$ y $R$ tales que toda arista va de $L$ a $R$. Esto equivale a decir que se puede \textbf{2-colorear} sin que dos vertices adyacentes tengan el mismo color. BFS alternando colores detecta esto: si en algun momento dos adyacentes tienen el mismo color, no es bipartito. La demostracion: por cada componente, fijar el color de un nodo determina todos los demas; si en algun paso hay conflicto, hay un ciclo impar (no bipartito).

\paragraph{Propiedades clave (invariantes).}
\begin{itemize}\setlength\itemsep{0pt}
	\item Grafo bipartito $\Leftrightarrow$ sin ciclos impares.
	\item Conexo: si un componente no es bipartito, todo el grafo no lo es.
	\item Si hay \textbf{multiple componentes}, cada uno se colorea independientemente.
	\item La 2-coloracion es unica salvo swap de colores.
\end{itemize}

\paragraph{Codigo clave.}
BFS con 2-coloreo:
\begin{verbatim}
vector<int> color(n, -1);
for (int v = 0; v < n; ++v) if (color[v] == -1) {
    color[v] = 0;
    queue<int> q; q.push(v);
    while (!q.empty()) {
        int u = q.front(); q.pop();
        for (int w : adj[u]) {
            if (color[w] == -1) { color[w] = 1 - color[u]; q.push(w); }
            else if (color[w] == color[u]) return false;
        }
    }
}
return true;
\end{verbatim}

\paragraph{Complejidad.}
\begin{itemize}\setlength\itemsep{0pt}
	\item $O(V + E)$.
	\item Memoria $O(V)$.
\end{itemize}

\paragraph{Donde tocar para modificar.}
\begin{itemize}\setlength\itemsep{0pt}
	\item \textbf{Colorear desde multiples fuentes}: BFS desde cada nodo no coloreado.
	\item \textbf{Bipartito + matching}: Hopcroft-Karp.
	\item \textbf{Grafo bipartito ponderado}: Hungarian algorithm.
	\item \textbf{Verificar bipartito en subgrafo}: solo aplicar BFS a los nodos relevantes.
\end{itemize}

\paragraph{Trampas y errores frecuentes.}
\begin{itemize}\setlength\itemsep{0pt}
	\item Verificar \textbf{todos} los componentes, no solo el primero.
	\item Para grafos con auto-loops, un auto-loop hace al grafo no bipartito.
	\item En multigrafos, las aristas multiples no afectan el chequeo de bipartito.
\end{itemize}

\paragraph{Codigo.}
Ver \texttt{cpp.tex}, seccion \textit{Bipartite check + 2-coloring}.

\vspace{0.5em}

\subsubsection{0-1 BFS}

\paragraph{Idea intuitiva.}
Variante de BFS para grafos con \textbf{pesos 0 o 1}. Usa un \texttt{deque}: aristas de peso 0 se encolan al frente (\texttt{push\_front}), aristas de peso 1 al final (\texttt{push\_back}). Asi, la primera vez que se extrae un nodo, su distancia es la definitiva, igual que en Dijkstra pero en $O(V + E)$. La intuicion: el deque mantiene los nodos ordenados por distancia; las aristas de peso 0 no cambian la distancia (entran al frente), las de peso 1 la incrementan (entran al final).

\paragraph{Propiedades clave (invariantes).}
\begin{itemize}\setlength\itemsep{0pt}
	\item Asume pesos \textbf{solo} 0 o 1.
	\item Mejor caso que Dijkstra (sin factor $\log V$).
	\item No funciona con pesos $\ge 2$.
	\item La cola mantiene los nodos ordenados por distancia no decreciente.
\end{itemize}

\paragraph{Codigo clave.}
El deque doble:
\begin{verbatim}
deque<int> q;
q.push_front(source);
dist[source] = 0;
while (!q.empty()) {
    int v = q.front(); q.pop_front();
    for (auto [w, u] : adj[v]) {
        if (dist[v] + w < dist[u]) {
            dist[u] = dist[v] + w;
            if (w == 0) q.push_front(u);
            else q.push_back(u);
        }
    }
}
\end{verbatim}

\paragraph{Complejidad.}
\begin{itemize}\setlength\itemsep{0pt}
	\item $O(V + E)$.
	\item Cada nodo se inserta a lo sumo $O(1)$ veces (con check de distancia).
\end{itemize}

\paragraph{Donde tocar para modificar.}
\begin{itemize}\setlength\itemsep{0pt}
	\item \textbf{Pesos en $[0, k]$ small}: usar small k-ary BFS o Dijkstra.
	\item \textbf{Reconstruir camino}: aniadir \texttt{parent[]}.
	\item \textbf{Pesos negativos}: no funciona, usar Bellman-Ford.
\end{itemize}

\paragraph{Trampas y errores frecuentes.}
\begin{itemize}\setlength\itemsep{0pt}
	\item Si los pesos son otro rango, el algoritmo da respuestas incorrectas. Verificar que las aristas son solo 0/1.
	\item El chequeo \texttt{dist[v] + w < dist[u]} es crucial; sino, el algoritmo procesa nodos multiples veces.
	\item Si una arista tiene peso $\ge 2$, aniadir verificaciones o cambiar a Dijkstra.
\end{itemize}

\paragraph{Codigo.}
Ver \texttt{cpp.tex}, seccion \textit{0-1 BFS}.

\subsection{MST}

\subsubsection{Kruskal}

\paragraph{Idea intuitiva.}
Encuentra el \textbf{arbol de expansion minima} (MST) en $O(E \log E)$. Ordena las aristas por peso ascendente y las aniade al MST si no forman ciclo (chequear con DSU). El arbol tiene exactamente $V - 1$ aristas. La demostracion (``Cut Property''): para cualquier corte, la arista minima cruzando el corte pertenece a \emph{algun} MST; por lo tanto, elegir la minima disponible que no forma ciclo es optimo.

\paragraph{Propiedades clave (invariantes).}
\begin{itemize}\setlength\itemsep{0pt}
	\item Algoritmo \textbf{greedy} correcto (Cut Property).
	\item Funciona en grafos \textbf{no conexos}; produce un \textbf{minimum spanning forest}.
	\item Si las aristas tienen pesos unicos, el MST es unico.
	\item El MST tiene exactamente $V - 1$ aristas (o menos si el grafo es disconexo).
\end{itemize}

\paragraph{Codigo clave.}
Kruskal con DSU:
\begin{verbatim}
sort(edges.begin(), edges.end());
DSU dsu(n);
long long mst = 0;
int edgesUsed = 0;
for (auto [w, u, v] : edges)
    if (dsu.unite(u, v)) {
        mst += w;
        if (++edgesUsed == n - 1) break;
    }
\end{verbatim}

\paragraph{Complejidad.}
\begin{itemize}\setlength\itemsep{0pt}
	\item $O(E \log E)$ (ordenamiento) o $O(E \log V)$ con DSU optimizations.
	\item DSU es casi $O(1)$ amortizado, asi que el sort domina.
\end{itemize}

\paragraph{Donde tocar para modificar.}
\begin{itemize}\setlength\itemsep{0pt}
	\item \textbf{Second best MST}: enumerar aristas del MST, sacar una, aniadir una no-MST; recalcular.
	\item \textbf{Kruskal randomizado}: si el grafo es casi-completo, agregar shuffle.
	\item \textbf{Aniadir info al DSU}: guardar el peso maximo en el camino, etc.
	\item \textbf{DSU con rollback}: si necesitas Kruskal reversible (para queries offline).
\end{itemize}

\paragraph{Trampas y errores frecuentes.}
\begin{itemize}\setlength\itemsep{0pt}
	\item El DSU debe estar limpio para cada invocacion.
	\item Si las aristas tienen pesos negativos, Kruskal sigue funcionando.
	\item Si el grafo es disconexo, el ``MST'' es un bosque con varias componentes.
	\item El check \texttt{if (++edgesUsed == n - 1) break;} optimiza la salida temprana.
\end{itemize}

\paragraph{Codigo.}
Ver \texttt{cpp.tex}, seccion \textit{Kruskal}.

\vspace{0.5em}

\subsubsection{Prim (priority queue)}

\paragraph{Idea intuitiva.}
Otra forma de calcular el MST: empezar desde un nodo y, en cada paso, aniadir la arista de menor peso que conecta el conjunto construido con el resto. Usando priority queue, es $O((V + E) \log V)$. La demostracion: cada vez que se aniade una arista, es la minima cruzando el ``corte'' entre el conjunto actual y el resto (Cut Property); eso preserva optimalidad.

\paragraph{Propiedades clave (invariantes).}
\begin{itemize}\setlength\itemsep{0pt}
	\item Similar a Dijkstra pero sin acumular distancias, solo eligiendo minimo por arista.
	\item Funciona mejor que Kruskal en \textbf{grafos densos} ($E \approx V^2$).
	\item Si el grafo es disconexo, genera el bosque (solo llega a los alcanzables).
	\item Cada nodo se inserta al heap una vez; cada arista se procesa una vez.
\end{itemize}

\paragraph{Codigo clave.}
Prim clasico con heap:
\begin{verbatim}
priority_queue<pair<int, int>, vector<...>, greater<...>> pq;
vector<bool> used(n, false);
vector<int> minEdge(n, INF);
minEdge[0] = 0;
pq.push({0, 0});
while (!pq.empty()) {
    auto [w, v] = pq.top(); pq.pop();
    if (used[v]) continue;
    used[v] = true;
    mst += w;
    for (auto [peso, u] : adj[v])
        if (!used[u] && peso < minEdge[u]) {
            minEdge[u] = peso;
            pq.push({peso, u});
        }
}
\end{verbatim}

\paragraph{Complejidad.}
\begin{itemize}\setlength\itemsep{0pt}
	\item $O((V + E) \log V)$.
	\item Con matriz de adyacencia: $O(V^2)$ sin heap.
\end{itemize}

\paragraph{Donde tocar para modificar.}
\begin{itemize}\setlength\itemsep{0pt}
	\item \textbf{Prim con matriz de adyacencia}: $O(V^2)$ sin heap, util en grafos densos pequenos.
	\item \textbf{Reconstruir el MST}: aniadir \texttt{parent[]}.
	\item \textbf{Prim en streaming}: variante para grafos que llegan incrementalmente.
\end{itemize}

\paragraph{Trampas y errores frecuentes.}
\begin{itemize}\setlength\itemsep{0pt}
	\item El \texttt{if (used[v]) continue;} es crucial: si el nodo ya fue procesado, saltar.
	\item En la version matriz, elije el minimo por linea (no heap).
	\item Si el grafo es disconexo, aniadir un loop externo sobre todos los nodos no usados.
\end{itemize}

\paragraph{Codigo.}
Ver \texttt{cpp.tex}, seccion \textit{Prim (priority queue)}.

\subsection{Flujos}

\subsubsection{Edmonds-Karp (max flow)}

\paragraph{Idea intuitiva.}
Implementacion BFS del algoritmo de \textbf{Ford-Fulkerson}. En cada iteracion, hacer BFS desde $s$ hasta $t$ siguiendo aristas con capacidad residual $> 0$; si se llega a $t$, se encontro un \textbf{camino de aumento}. Enviar el minimo de las capacidades en el camino; actualizar capacidades residuales. La demostracion de $O(VE^2)$: cada BFS encuentra el augmenting path mas corto, y la longitud minima crece monotona; como hay $\le E$ longitudes posibles, hay $\le E$ fases.

\paragraph{Propiedades clave (invariantes).}
\begin{itemize}\setlength\itemsep{0pt}
	\item BFS garantiza caminos de aumento mas cortos.
	\item Capacidad residual: forward edge baja, backward edge sube.
	\item Complejidad $O(VE^2)$.
	\item El flujo maximo respeta la conservacion en nodos (in = out salvo $s, t$).
\end{itemize}

\paragraph{Codigo clave.}
BFS + augmenting path:
\begin{verbatim}
while (bfs(s, t, parent)) {  // bfs encuentra camino o retorna false
    int pathFlow = INF;
    for (int v = t; v != s; v = parent[v])
        pathFlow = min(pathFlow, capacity[parent[v]][v]);
    for (int v = t; v != s; v = parent[v]) {
        capacity[parent[v]][v] -= pathFlow;
        capacity[v][parent[v]] += pathFlow;
    }
    maxFlow += pathFlow;
}
\end{verbatim}

\paragraph{Complejidad.}
\begin{itemize}\setlength\itemsep{0pt}
	\item $O(VE^2)$.
	\item En la practica, mas lento que Dinic para grafos grandes.
\end{itemize}

\paragraph{Donde tocar para modificar.}
\begin{itemize}\setlength\itemsep{0pt}
	\item \textbf{Grafo bipartito}: $O(\sqrt{V} E)$ (equivalente a Hopcroft-Karp).
	\item \textbf{Dinic} es preferible en la mayoria de los casos.
	\item \textbf{Capacidades doubles}: usar tolerancia para comparacion.
	\item \textbf{Min-cost flow}: aniadir costo por arista y modificar el algoritmo.
\end{itemize}

\paragraph{Trampas y errores frecuentes.}
\begin{itemize}\setlength\itemsep{0pt}
	\item Las capacidades deben ser $> 0$ para entrar a la cola (no $>= 0$); sino, ciclos infinitos.
	\item Para grafos bipartitos, Dinic o Hopcroft-Karp son mejores.
	\item Si el flujo maximo es muy grande, los \texttt{int} pueden overflow; usar \texttt{long long}.
\end{itemize}

\paragraph{Codigo.}
Ver \texttt{cpp.tex}, seccion \textit{Edmonds-Karp (max flow)}.

\vspace{0.5em}

\subsubsection{Dinic (max flow)}

\paragraph{Idea intuitiva.}
Algoritmo de flujo mas eficiente que Edmonds-Karp: combina \textbf{BFS para level graph} con \textbf{DFS bloqueante} (multiple augmenting paths en una sola pasada). Cada fase cuesta $O(E)$ y hay $O(V)$ fases en el peor caso; para bipartitos, $O(\sqrt{V} E)$. La razon: cada fase ``satura'' al menos una arista, asi que en $O(E)$ fases el flujo maximo se alcanza; con $O(V)$ niveles en el BFS, hay $O(VE)$ operaciones totales.

\paragraph{Propiedades clave (invariantes).}
\begin{itemize}\setlength\itemsep{0pt}
	\item \texttt{level[v]} = distancia desde $s$ en el level graph (solo aristas con cap $> 0$).
	\item \texttt{it[v]} = iterador actual para evitar reprocesar aristas.
	\item DFS solo sigue aristas con \texttt{level[e.to] == level[v] + 1}.
	\item Back-edges tienen \texttt{rev} apuntando a la forward edge original.
\end{itemize}

\paragraph{Codigo clave.}
El DFS bloqueante de Dinic:
\begin{verbatim}
long long dfs(int v, int t, long long f) {
    if (v == t) return f;
    for (int& i = it[v]; i < adj[v].size(); ++i) {
        Edge& e = adj[v][i];
        if (e.cap > 0 && level[e.to] == level[v] + 1) {
            long long ret = dfs(e.to, t, min(f, e.cap));
            if (ret > 0) {
                e.cap -= ret;
                adj[e.to][e.rev].cap += ret;
                return ret;
            }
        }
    }
    return 0;
}
\end{verbatim}

\paragraph{Complejidad.}
\begin{itemize}\setlength\itemsep{0pt}
	\item $O(V^2 E)$ general, $O(\sqrt{V} E)$ para bipartitos.
	\item En la practica, mucho mas rapido que Edmonds-Karp.
\end{itemize}

\paragraph{Donde tocar para modificar.}
\begin{itemize}\setlength\itemsep{0pt}
	\item \textbf{Dinic con scaling}: aniadir factor de escala para capacidades grandes.
	\item \textbf{Lower bounds}: aniadir un super-source/sink para forzar flujo minimo.
	\item \textbf{Matching bipartito}: el grafo se construye como source -> L -> R -> sink, todos con capacidad 1.
	\item \textbf{Push-relabel}: alternativa para grafos muy densos.
\end{itemize}

\paragraph{Trampas y errores frecuentes.}
\begin{itemize}\setlength\itemsep{0pt}
	\item \texttt{it[v]} debe \textbf{incrementarse} en cada iteracion; sino, el DFS se queda atascado en la misma arista.
	\item Si el BFS no encuentra $t$, el flujo es maximo; terminar.
	\item Las back-edges son cruciales para ``cancelar'' elecciones previas.
\end{itemize}

\paragraph{Codigo.}
Ver \texttt{cpp.tex}, seccion \textit{Dinic (max flow)}.

\vspace{0.5em}

\subsubsection{Min-Cost Max-Flow}

\paragraph{Idea intuitiva.}
Encuentra el flujo maximo con \textbf{costo minimo} (asumiendo costos no negativos). Usa \textbf{SPFA} (o Dijkstra con potentials) para encontrar el shortest path en el \textbf{costo} desde $s$ hasta $t$ considerando capacidades residuales. Cada augmenting path aporta su costo $\times$ flujo. La idea: los potentials reescalan las aristas para que todas tengan costo no negativo, permitiendo usar Dijkstra.

\paragraph{Propiedades clave (invariantes).}
\begin{itemize}\setlength\itemsep{0pt}
	\item \texttt{pot[]} = potentials para evitar aristas negativas (Johnson's trick).
	\item \texttt{dist[v]} = costo minimo desde $s$ con potentials.
	\item Solo se aumentan aristas con \texttt{cap > 0} (forward) y costo es \texttt{+cost} (forward) o \texttt{-cost} (backward).
	\item Tras cada augmenting, los potentials se actualizan para mantener la consistencia.
\end{itemize}

\paragraph{Codigo clave.}
El ciclo de min-cost flow:
\begin{verbatim}
while (spfa(s, t, dist, parent)) {
    int pathFlow = INF;
    for (int v = t; v != s; v = parent[v])
        pathFlow = min(pathFlow, capacity[parent[v]][v]);
    for (int v = t; v != s; v = parent[v]) {
        capacity[parent[v]][v] -= pathFlow;
        capacity[v][parent[v]] += pathFlow;
        cost += pathFlow * edge_cost[parent[v]][v];
    }
}
\end{verbatim}

\paragraph{Complejidad.}
\begin{itemize}\setlength\itemsep{0pt}
	\item $O(F \cdot V E)$ con SPFA, $O(F \cdot E \log V)$ con Dijkstra + potentials.
	\item $F$ = flujo maximo total.
\end{itemize}

\paragraph{Donde tocar para modificar.}
\begin{itemize}\setlength\itemsep{0pt}
	\item \textbf{Dijkstra + potentials}: si los costos son enteros no negativos, esto es mas estable.
	\item \textbf{Costos negativos en aristas}: requieren Bellman-Ford inicial o reorden.
	\item \textbf{Flujo minimo con costo}: encontrar el minimo flujo posible y su costo.
	\item \textbf{Costo por unidad}: si los costos son muy grandes, aniadir verificacion.
\end{itemize}

\paragraph{Trampas y errores frecuentes.}
\begin{itemize}\setlength\itemsep{0pt}
	\item El \texttt{pot[]} se actualiza \textbf{al final} de cada iteracion (no durante).
	\item Las back-edges tienen \textbf{costo negativo}, esencial para corregir elecciones previas.
	\item Si los costos son muy grandes, considerar modular arithmetic.
	\item Para problemas con flujo exacto, aniadir verificacion post-algoritmo.
\end{itemize}

\paragraph{Codigo.}
Ver \texttt{cpp.tex}, seccion \textit{Min-Cost Max-Flow}.

\subsection{Especales}

\subsubsection{2-SAT (implication graph + SCC)}

\paragraph{Idea intuitiva.}
Resuelve formulas booleanas en \textbf{CNF} con clausulas de \textbf{2 literales}: $(l_1 \lor l_2)$. Cada variable $x$ se representa con dos nodos: $x$ (verdadero) y $\neg x$ (falso). Cada clausula $l_1 \lor l_2$ se descompone en dos implicaciones: $\neg l_1 \to l_2$ y $\neg l_2 \to l_1$. Si en el grafo de implicaciones $\neg x$ y $x$ estan en la misma SCC, la formula es \textbf{UNSAT}. La demostracion: si $\neg x \to x$ en el grafo, entonces $\neg x$ implica $x$, asi que $\neg x$ debe ser falso y $x$ verdadero; pero entonces $\neg x$ es falso implica una contradiccion.

\paragraph{Propiedades clave (invariantes).}
\begin{itemize}\setlength\itemsep{0pt}
	\item Cada literal es un nodo; total $2n$ nodos.
	\item Kosaraju/Tarjan dan las SCCs.
	\item Si SAT: asignar $x = $ (topological order del condensation).
	\item Clausulas \texttt{setTrue(x)} se modelan como \texttt{implies(!x, x)}.
	\item La formula es SAT sii no hay variable $x$ con $x$ y $\neg x$ en la misma SCC.
\end{itemize}

\paragraph{Codigo clave.}
Conversion de clausulas a implicaciones:
\begin{verbatim}
// Clausula (a OR b) -> implica ~a -> b Y ~b -> a
addImplication(not(a), b);
addImplication(not(b), a);
// Forzar x: implica ~x -> x
if (force) addImplication(not(x), x);
\end{verbatim}

\paragraph{Complejidad.}
\begin{itemize}\setlength\itemsep{0pt}
	\item $O(V + E) = O(N + M)$ con Kosaraju o Tarjan.
	\item Memoria $O(V + E)$.
\end{itemize}

\paragraph{Donde tocar para modificar.}
\begin{itemize}\setlength\itemsep{0pt}
	\item \textbf{Forzar variables}: aniadir \texttt{setTrue} / \texttt{setFalse}.
	\item \textbf{3-SAT}: NP-completo; no usar este algoritmo.
	\item \textbf{Contar soluciones}: las SCCs condensation dan una formula para contar; multiplicar las formas de asignar variables en cada componente no-trivial.
	\item \textbf{Output de la asignacion}: aniadir el vector \texttt{assignment[]}.
\end{itemize}

\paragraph{Trampas y errores frecuentes.}
\begin{itemize}\setlength\itemsep{0pt}
	\item La asignacion usa el \textbf{orden topologico inverso} del condensation: \texttt{val[i] = comp[2i] < comp[2i+1]}.
	\item El ID de $\neg x$ es \texttt{2*x + 1} (o similar); verificar consistencia.
	\item Si hay clausulas con literales repetidos (e.g., $(x \lor x)$), tratar como $(x)$ simple.
\end{itemize}

\paragraph{Codigo.}
Ver \texttt{cpp.tex}, seccion \textit{2-SAT (implication graph + SCC)}.

\vspace{0.5em}

\subsubsection{Euler tour / Hierholzer (camino y circuito)}

\paragraph{Idea intuitiva.}
Encuentra un \textbf{camino o circuito euleriano} (que usa cada arista exactamente una vez). \textbf{Condiciones}:
\begin{itemize}\setlength\itemsep{0pt}
	\item \textbf{No dirigido}: todos los grados pares (circuito) o exactamente 2 impares (camino, empieza en uno, termina en el otro). Ademas, conexo sobre el subgrafo con aristas.
	\item \textbf{Dirigido}: $\text{indeg}(v) = \text{outdeg}(v)$ para todos (circuito) o exactamente un nodo con $\text{outdeg} = \text{indeg} + 1$ y otro con $\text{indeg} = \text{outdeg} + 1$ (camino).
\end{itemize}

\paragraph{Hierholzer}: elegir un ciclo desde el inicio, mientras haya ciclos ``laterales'' desde un nodo del camino, mergearlos. Implementacion con pila: avanzar por aristas, al no poder mas aniadir a la respuesta y backtrack. La demostracion de $O(V + E)$: cada arista se procesa una vez (entra y sale del stack una vez).

\paragraph{Propiedades clave (invariantes).}
\begin{itemize}\setlength\itemsep{0pt}
	\item Multigrafos permitidos; las aristas se marcan como usadas.
	\item No dirigido: cada arista aparece \textbf{dos veces} en \texttt{adj} (ida y vuelta).
	\item La construccion del path con ids consistentes es importante.
	\item El resultado es unico salvo orden de las aristas en multigrafos.
\end{itemize}

\paragraph{Codigo clave.}
Hierholzer con pila:
\begin{verbatim}
vector<int> path;
stack<int> st;
st.push(start);
while (!st.empty()) {
    int v = st.top();
    while (!adj[v].empty() && used[adj[v].back()]) adj[v].pop_back();
    if (adj[v].empty()) {
        path.push_back(v);
        st.pop();
    } else {
        int e = adj[v].back();
        used[e] = true;
        int u = edges[e].to(v);
        adj[v].pop_back();
        st.push(u);
    }
}
reverse(path.begin(), path.end());
\end{verbatim}

\paragraph{Complejidad.}
\begin{itemize}\setlength\itemsep{0pt}
	\item $O(V + E)$.
	\item Memoria $O(V + E)$.
\end{itemize}

\paragraph{Donde tocar para modificar.}
\begin{itemize}\setlength\itemsep{0pt}
	\item \textbf{Reconstruir con multigrafo}: aniadir ids consistentes antes del algoritmo.
	\item \textbf{Aplicar a ``de Bruijn'' o problemas de strings}: modelar como multigrafo y buscar Euler.
	\item \textbf{Dirigido vs no dirigido}: el snippet tiene ambas versiones.
	\item \textbf{Letter addition}: modelar palabras como aristas y buscar el superstring.
\end{itemize}

\paragraph{Trampas y errores frecuentes.}
\begin{itemize}\setlength\itemsep{0pt}
	\item Si el grafo no es conexo (sobre las aristas), no existe euleriano. Verificar primero.
	\item Para multigrafos, los IDs de aristas son cruciales para distinguirlas.
	\item Si el problema es ``encuentra todos los euler tours'', el algoritmo es exponencial.
\end{itemize}

\paragraph{Codigo.}
Ver \texttt{cpp.tex}, seccion \textit{Euler tour / Hierholzer (camino y circuito)}.

\vspace{0.5em}

\subsubsection{Hamiltonian path / cycle (bitmask DP + backtracking)}

\paragraph{Idea intuitiva.}
Un \textbf{camino hamiltoniano} visita cada vertice \textbf{exactamente una vez}. \textbf{Determinar} si existe es NP-completo en general, pero con $V \le 20$ se puede hacer DP con bitmask en $O(V^2 \cdot 2^V)$. \textbf{Teoremas suficientes}: \textbf{Dirac} (grado $\ge n/2$) y \textbf{Ore} ($\text{deg}(u) + \text{deg}(v) \ge n$ para todo par no adyacente). La idea de la DP: cada estado es (conjunto de visitados, ultimo vertice); la transicion es agregar un vecino no visitado.

\paragraph{Propiedades clave (invariantes).}
\begin{itemize}\setlength\itemsep{0pt}
	\item DP: $dp[mask][v]$ = true si hay un camino que visita exactamente \texttt{mask} y termina en $v$.
	\item TSP: variante con pesos minimos (Hamilton cycle minimo).
	\item Reconstruccion: guardar \texttt{par[mask][v]} (predecesor).
	\item $2^V$ estados con $V$ bits cada uno; factible para $V \le 20$.
\end{itemize}

\paragraph{Codigo clave.}
DP con bitmask:
\begin{verbatim}
dp[1][0] = 0;  // empezar en 0
for (int mask = 1; mask < (1<<n); ++mask)
    for (int v = 0; v < n; ++v) if (dp[mask][v] < INF)
        for (int u = 0; u < n; ++u)
            if (!(mask & (1<<u)))
                dp[mask | (1<<u)][u] = min(dp[mask | (1<<u)][u],
                                            dp[mask][v] + w[v][u]);
// TSP: respuesta = min(dp[(1<<n)-1][v] + w[v][0])
\end{verbatim}

\paragraph{Complejidad.}
\begin{itemize}\setlength\itemsep{0pt}
	\item $O(V^2 \cdot 2^V)$ tiempo, $O(V \cdot 2^V)$ memoria.
	\item Para $V = 20$: $\sim 4 \cdot 10^8$ operaciones (lento pero factible).
\end{itemize}

\paragraph{Donde tocar para modificar.}
\begin{itemize}\setlength\itemsep{0pt}
	\item \textbf{TSP asimetrico}: aniadir pesos $w[v][u]$ (no necesariamente simetricos).
	\item \textbf{Hamilton con inicio obligatorio}: fijar \texttt{src} en el DP inicial.
	\item \textbf{Bitmask DP mas compacto}: usar \texttt{long long} como mascara si $V \le 64$.
	\item \textbf{TSP optimo aproximado}: MST, 2-opt, simulated annealing.
	\item \textbf{Reconstruir el path}: usar \texttt{par[mask][v]} y backtrack.
\end{itemize}

\paragraph{Trampas y errores frecuentes.}
\begin{itemize}\setlength\itemsep{0pt}
	\item TSP asume grafo \textbf{completo} (con \texttt{LINF} para aristas faltantes); si no, aniadir chequeo.
	\item La condicion ``\texttt{mask == (1<<n) - 1}'' cierra el ciclo retornando al origen.
	\item Para $V > 20$, la memoria $2^V$ se vuelve prohibitiva; usar meet-in-the-middle o heuristicas.
\end{itemize}

\paragraph{Codigo.}
Ver \texttt{cpp.tex}, seccion \textit{Hamiltonian path / cycle (bitmask DP + backtracking)}.

% ============================================================
\section{DP}

\subsubsection{Knapsack 0/1}

\paragraph{Idea intuitiva.}
Dado un conjunto de $N$ objetos con peso $w_i$ y valor $v_i$, maximizar $\sum v_i$ con la restriccion $\sum w_i \le W$. La recurrencia clasica:
\[ dp[i][j] = \max(dp[i-1][j],\, dp[i-1][j - w_i] + v_i) \quad \text{si } j \ge w_i. \]
La optimizacion clave: la DP es 1-dimensional si iteramos $j$ \textbf{de mayor a menor} (asi no usamos el objeto $i$ dos veces). La razon de la iteracion inversa: en cada iteracion queremos el estado del ``item anterior''; iterar de menor a mayor contaminaria \texttt{dp[j - w_i]} con el item actual.

\paragraph{Propiedades clave (invariantes).}
\begin{itemize}\setlength\itemsep{0pt}
	\item $dp[j]$ al final = maximo valor con peso $\le j$.
	\item El bucle inverso \texttt{for (j = W; j >= w[i]; j--)} es lo que distingue 0/1 de unbounded.
	\item Solucion: maximo valor con cualquier peso $\le W$.
	\item La DP es exacta; no hay aproximacion.
\end{itemize}

\paragraph{Codigo clave.}
La version 1D optimizada:
\begin{verbatim}
vector<long long> dp(W + 1, 0);
for (int i = 0; i < N; ++i)
    for (int j = W; j >= w[i]; --j)  // invertido
        dp[j] = max(dp[j], dp[j - w[i]] + v[i]);
return dp[W];
\end{verbatim}

\paragraph{Complejidad.}
\begin{itemize}\setlength\itemsep{0pt}
	\item $O(N \cdot W)$ tiempo, $O(W)$ memoria.
	\item Constante muy pequena (solo sumas y max).
\end{itemize}

\paragraph{Donde tocar para modificar.}
\begin{itemize}\setlength\itemsep{0pt}
	\item \textbf{Reconstruir la seleccion}: aniadir \texttt{take[i][j]} (booleano o bool) en una version 2D.
	\item \textbf{Items con peso 0}: edge case; aniadir guard.
	\item \textbf{Knapsack con restricciones de grupos}: ``elegir a lo sumo 1 de cada grupo'' se modela como 0/1 con un peso extra.
	\item \textbf{Tamano de $W$}: si $W \le 10^9$ pero $N$ es chico, usar meet-in-the-middle.
	\item \textbf{Cambiar max por min}: aniadir \texttt{INF} en la inicializacion.
\end{itemize}

\paragraph{Trampas y errores frecuentes.}
\begin{itemize}\setlength\itemsep{0pt}
	\item Overflow en \texttt{dp[j - w[i]] + v[i]}; usar \texttt{long long} si los valores son grandes.
	\item El bucle \textbf{debe} ir de mayor a menor.
	\item Si los pesos son grandes y $N$ pequeno, conviene meet-in-the-middle.
	\item Para Knapsack exacto ($\sum w_i = W$), aniadir check al final.
\end{itemize}

\paragraph{Codigo.}
Ver \texttt{cpp.tex}, seccion \textit{Knapsack 0/1}.

\vspace{0.5em}

\subsubsection{Knapsack unbounded}

\paragraph{Idea intuitiva.}
Igual que Knapsack 0/1 pero \textbf{cada item se puede tomar multiples veces}. La recurrencia es similar:
\[ dp[j] = \max(dp[j],\, dp[j - w_i] + v_i), \]
pero ahora el bucle va de \textbf{menor a mayor} (porque tomar el item $i$ no afecta corridas futuras del mismo item). La idea: tras tomar el item $i$, queremos volver a tener la opcion de tomarlo; iterar de menor a mayor permite que el item ``se acumule''.

\paragraph{Propiedades clave (invariantes).}
\begin{itemize}\setlength\itemsep{0pt}
	\item El bucle \texttt{for (j = w[i]; j <= W; j++)} es lo que distingue unbounded.
	\item A diferencia de 0/1, no hay riesgo de ``usar dos veces el mismo item''.
	\item La DP es exacta; cada ``subconjunto'' se cuenta correctamente.
\end{itemize}

\paragraph{Codigo clave.}
La version 1D con bucle directo:
\begin{verbatim}
vector<long long> dp(W + 1, 0);
for (int i = 0; i < N; ++i)
    for (int j = w[i]; j <= W; ++j)  // directo (no invertido)
        dp[j] = max(dp[j], dp[j - w[i]] + v[i]);
return dp[W];
\end{verbatim}

\paragraph{Complejidad.}
\begin{itemize}\setlength\itemsep{0pt}
	\item $O(N \cdot W)$ tiempo, $O(W)$ memoria.
	\item Constante muy pequena.
\end{itemize}

\paragraph{Donde tocar para modificar.}
\begin{itemize}\setlength\itemsep{0pt}
	\item \textbf{Reconstruir cuantos de cada item}: aniadir \texttt{cnt[i]} y backtrack.
	\item \textbf{Cambiar moneda minima por maxima}: el snippet maximiza valor; para minimizar monedas, cambiar a \texttt{min} con inicializacion en \texttt{INF}.
	\item \textbf{Coin change}: si todos los pesos son positivos y queremos ``minimo numero de monedas'', es unbounded con min.
\end{itemize}

\paragraph{Trampas y errores frecuentes.}
\begin{itemize}\setlength\itemsep{0pt}
	\item Misma que Knapsack 0/1 con overflow.
	\item El orden de iteracion de items no afecta el resultado (convexo).
	\item Si se quiere ``maximo valor con exactamente $k$ items'', aniadir segunda dimension.
\end{itemize}

\paragraph{Codigo.}
Ver \texttt{cpp.tex}, seccion \textit{Knapsack unbounded}.

\vspace{0.5em}

\subsubsection{Knapsack meet-in-the-middle}

\paragraph{Idea intuitiva.}
Para $N \le 40$ y $W$ moderado, no se puede hacer $O(NW)$ directamente. Partir en dos mitades, enumerar todos los $2^{N/2}$ subconjuntos de cada mitad, y combinar. Para cada subconjunto de la segunda mitad con peso $sw_2$, encontrar el maximo valor de la primera mitad con peso $\le W - sw_2$. La idea: ``divide y venceras'' sobre el tamano de la entrada; cada mitad tiene mitad del espacio, asi que la combinacion es geometrica en lugar de multiplicativa.

\paragraph{Propiedades clave (invariantes).}
\begin{itemize}\setlength\itemsep{0pt}
	\item Hay $2^{N/2}$ subconjuntos por mitad; total $O(2^{N/2})$.
	\item \texttt{sums} guarda los valores de subconjuntos de la primera mitad ordenados (o preprocesados).
	\item Para cada subconjunto de la segunda mitad, busqueda lineal o binaria en \texttt{sums}.
	\item La mitad mas grande (con $N/2 + 1$ items para $N$ impar) tiene $2^{N/2+1}$ subconjuntos.
\end{itemize}

\paragraph{Codigo clave.}
Enumerar subconjuntos de la primera mitad:
\begin{verbatim}
int n1 = N / 2, n2 = N - n1;
vector<pair<long long, long long>> sums1;
for (int mask = 0; mask < (1 << n1); ++mask) {
    long long peso = 0, valor = 0;
    for (int i = 0; i < n1; ++i) if (mask & (1 << i)) {
        peso += w[i]; valor += v[i];
    }
    sums1.push_back({peso, valor});
}
sort(sums1.begin(), sums1.end());  // por peso
// eliminar duplicados dominados
\end{verbatim}

\paragraph{Complejidad.}
\begin{itemize}\setlength\itemsep{0pt}
	\item $O(2^{N/2})$ tiempo y memoria (mejor que $O(NW)$ cuando $W > 2^{N/2}$).
	\item Para $N = 40$: $\sim 2 \cdot 10^6$ subconjuntos, factible.
\end{itemize}

\paragraph{Donde tocar para modificar.}
\begin{itemize}\setlength\itemsep{0pt}
	\item \textbf{Tamano mayor}: si $N$ es mayor pero $W$ es chico, conviene iterar por peso, no por subconjuntos.
	\item \textbf{Varias mochilas}: adaptar para multiples $W$.
	\item \textbf{Optimizar el ``combinado''}: usar \texttt{upper\_bound} para encontrar el maximo valor con peso $\le \text{rem}.
	\item \textbf{Eliminar pares dominados}: si un par (peso, valor) tiene otro con menos peso y mas valor, eliminarlo.
\end{itemize}

\paragraph{Trampas y errores frecuentes.}
\begin{itemize}\setlength\itemsep{0pt}
	\item $N$ impar: una mitad tiene $N/2 + 1$ items.
	\item Overflow al sumar pesos/valores en \texttt{long long}.
	\item Sin deduplicacion, el algoritmo es correcto pero mas lento.
\end{itemize}

\paragraph{Codigo.}
Ver \texttt{cpp.tex}, seccion \textit{Knapsack meet-in-the-middle}.

\vspace{0.5em}

\subsubsection{LIS O(n log n)}

\paragraph{Idea intuitiva.}
La \textbf{Longest Increasing Subsequence} se calcula en $O(n \log n)$ con el truco del \textbf{patience sorting}: mantener un arreglo \texttt{dp} donde \texttt{dp[len]} es el minimo ultimo valor de una IS de largo \texttt{len}. Para cada $x$, encontrar el primer \texttt{dp[len] >= x} y reemplazarlo. El largo final es la LIS. La demostracion: si hay dos IS del mismo largo, la que termina en un valor menor es ``mejor'' (puede extenderse mas facilmente).

\paragraph{Propiedades clave (invariantes).}
\begin{itemize}\setlength\itemsep{0pt}
	\item \texttt{dp} es estrictamente creciente.
	\item \texttt{lower\_bound} da LIS estricta; \texttt{upper\_bound} da no-estricta ($\le$).
	\item El algoritmo es \textbf{online}: procesa elementos en orden.
	\item La longitud de \texttt{dp} es la LIS al final.
\end{itemize}

\paragraph{Codigo clave.}
La version online:
\begin{verbatim}
vector<int> dp;
for (int x : a) {
    auto it = lower_bound(dp.begin(), dp.end(), x);
    if (it == dp.end()) dp.push_back(x);
    else *it = x;
}
// dp.size() = longitud de LIS estricta
\end{verbatim}

\paragraph{Complejidad.}
\begin{itemize}\setlength\itemsep{0pt}
	\item $O(n \log n)$.
	\item Memoria $O(n)$.
\end{itemize}

\paragraph{Donde tocar para modificar.}
\begin{itemize}\setlength\itemsep{0pt}
	\item \textbf{Reconstruir la LIS}: aniadir \texttt{parent[]} y \texttt{posInDp[]} para backtrack.
	\item \textbf{Contar LIS}: DP clasica $O(n^2)$ o segment tree sobre $dp$.
	\item \textbf{LDS / Longest non-increasing}: invertir la condicion en \texttt{lower\_bound} o multiplicar por -1.
	\item \textbf{Longest bitonic subsequence}: combinar LIS desde izq y desde der.
	\item \textbf{LIS en 2D}: ordenar por una dimension y aplicar LIS en la otra.
\end{itemize}

\paragraph{Trampas y errores frecuentes.}
\begin{itemize}\setlength\itemsep{0pt}
	\item \texttt{lower\_bound} vs \texttt{upper\_bound}: si tu comparacion es $\le$ vs $<$, cambia.
	\item Si los elementos son \texttt{int} pero podrian ser negativos, aniadir offset para usar como indices.
	\item Para LIS reconstruida, aniadir el array \texttt{parent} que apunta al anterior en la IS.
\end{itemize}

\paragraph{Codigo.}
Ver \texttt{cpp.tex}, seccion \textit{LIS O(n log n)}.

\vspace{0.5em}

\subsubsection{LCS / Edit distance}

\paragraph{Idea intuitiva.}
\textbf{LCS} (Longest Common Subsequence): para dos strings $a$, $b$, $dp[i][j]$ = LCS de $a[0..i)$ y $b[0..j)$. Recurrencia:
\[ dp[i][j] = \begin{cases} dp[i-1][j-1] + 1 & \text{si } a[i-1] = b[j-1] \\ \max(dp[i-1][j], dp[i][j-1]) & \text{si no} \end{cases} \]

\textbf{Edit distance} (Levenshtein): numero minimo de operaciones (insert, delete, replace) para convertir $a$ en $b$. Recurrencia similar con un $+1$ adicional. La idea: ``alinear'' los strings eligiendo matches/mismatches/insert/delete.

\paragraph{Propiedades clave (invariantes).}
\begin{itemize}\setlength\itemsep{0pt}
	\item LCS y Edit Distance son DP $O(nm)$.
	\item Memoria se puede reducir a $O(\min(n, m))$ con rolling array.
	\item Edit distance con peso por operacion distinto: ajustar el \texttt{+1} por el peso.
	\item La subestructura optima: la respuesta para prefijos se construye de prefijos menores.
\end{itemize}

\paragraph{Codigo clave.}
La recurrencia clasica de LCS:
\begin{verbatim}
vector<vector<int>> dp(n + 1, vector<int>(m + 1, 0));
for (int i = 1; i <= n; ++i)
    for (int j = 1; j <= m; ++j)
        if (a[i-1] == b[j-1]) dp[i][j] = dp[i-1][j-1] + 1;
        else dp[i][j] = max(dp[i-1][j], dp[i][j-1]);
return dp[n][m];
\end{verbatim}

\paragraph{Complejidad.}
\begin{itemize}\setlength\itemsep{0pt}
	\item $O(NM)$ tiempo y memoria.
	\item Para $N, M \le 5000$: $\sim 2.5 \cdot 10^7$ operaciones.
\end{itemize}

\paragraph{Donde tocar para modificar.}
\begin{itemize}\setlength\itemsep{0pt}
	\item \textbf{Reconstruir la LCS / las operaciones}: aniadir \texttt{prev[i][j]} para backtrack.
	\item \textbf{LCS de varios strings}: NP-hard para $\ge 3$ strings; en la practica se hace con bitmask.
	\item \textbf{Edit distance con costos distintos}: cambiar el \texttt{+1} por el costo apropiado.
	\item \textbf{Rolling array}: si memoria es problema, mantener solo 2 filas.
	\item \textbf{Hirschberg}: algoritmo $O(NM)$ tiempo, $O(\min(N,M))$ memoria, con reconstruccion.
\end{itemize}

\paragraph{Trampas y errores frecuentes.}
\begin{itemize}\setlength\itemsep{0pt}
	\item Indices en la recurrencia: \texttt{a[i-1]} y \texttt{b[j-1]} (no \texttt{a[i]}).
	\item Para reconstruccion, el caso \texttt{a[i-1] == b[j-1]} tiene prioridad.
	\item Si los strings son muy largos, considerar Hirschberg o bitset optimization.
\end{itemize}

\paragraph{Codigo.}
Ver \texttt{cpp.tex}, seccion \textit{LCS / Edit distance}.

\vspace{0.5em}

\subsubsection{Bitmask DP over subsets}

\paragraph{Idea intuitiva.}
Cuando el estado es un \textbf{subconjunto} de $\{0, \dots, N-1\}$, se representa como una mascara de $N$ bits. Para iterar sobre todos los subconjuntos de \texttt{mask}:
\[ \text{for (sub = mask; sub; sub = (sub - 1) \& mask)} \{ \dots \} \]
Para iterar sobre todos los superconjuntos de \texttt{mask} en $[0, U)$:
\[ \text{for (sup = mask; sup < U; sup = (sup + 1) | mask)} \{ \dots \} \]
La idea: las mascaras forman un lattice de subconjuntos; los trucos aritmeticos (\texttt{(sub-1) \& mask} y \texttt{(sup+1) | mask}) enumeran todo el lattice eficientemente.

\paragraph{Propiedades clave (invariantes).}
\begin{itemize}\setlength\itemsep{0pt}
	\item Iterar subsets: $O(2^k)$ para una mascara con $k$ bits.
	\item \texttt{sub \& mask == sub} siempre; \texttt{mask - sub} no siempre da un subconjunto.
	\item \texttt{\_\_builtin\_popcount(mask)} cuenta los bits.
	\item \texttt{(sub - 1) \& mask} itera subsets de \texttt{mask} en orden decreciente de inclusion.
\end{itemize}

\paragraph{Codigo clave.}
Iteracion sobre subsets de una mascara:
\begin{verbatim}
// Todos los subsets de mask
for (int sub = mask; sub; sub = (sub - 1) & mask) {
    // procesar sub
}
// Incluir sub = 0 si es necesario
sub = 0;  // procesar primero
for (int sup = mask; sup < (1 << N); sup = (sup + 1) | mask) {
    // procesar sup (todos los superconjuntos)
}
\end{verbatim}

\paragraph{Complejidad.}
\begin{itemize}\setlength\itemsep{0pt}
	\item $O(N \cdot 2^N)$ para una DP estandar sobre todos los subsets.
	\item Iterar subsets de \texttt{mask}: $O(2^{|mask|})$.
\end{itemize}

\paragraph{Donde tocar para modificar.}
\begin{itemize}\setlength\itemsep{0pt}
	\item \textbf{TSP / Hamiltonian path}: bitmask DP (ver tambien \textit{Hamiltonian path} en \texttt{cpp.tex} seccion Grafos).
	\item \textbf{Set cover}: iterar subsets.
	\item \textbf{DP con supermask}: si el estado incluye ``todos los elementos que \textbf{no} hemos usado'', usar supermask iteration.
	\item \textbf{Bitmasks grandes}: con $N > 20$, aniadir compresion o simulacion con segmentos.
\end{itemize}

\paragraph{Trampas y errores frecuentes.}
\begin{itemize}\setlength\itemsep{0pt}
	\item El orden de iteracion de subsets de \texttt{mask} es decreciente (de \texttt{mask} a $0$).
	\item \texttt{mask = 0} no tiene subsets no triviales; el bucle lo saltea correctamente.
	\item El truco \texttt{(sub - 1) \& mask} requiere cuidado con \texttt{sub = 0} (termina el bucle).
	\item Superconjuntos: \texttt{sup < (1 << N)} debe estar bien definido.
\end{itemize}

\paragraph{Codigo.}
Ver \texttt{cpp.tex}, seccion \textit{Bitmask DP over subsets}.

\vspace{0.5em}

\subsubsection{SOS DP}

\paragraph{Idea intuitiva.}
Para calcular, para cada \texttt{mask}, la suma de $f(\text{submask})$ sobre todos los $\text{submask} \subseteq \text{mask}$:
\[ dp[mask] = \sum_{\text{sub} \subseteq mask} f(\text{sub}). \]
Se computa en $O(N \cdot 2^N)$ mediante un loop por bit: para cada $i$, sumar a $dp[mask]$ la contribucion de $dp[mask \oplus (1 << i)]$ cuando $mask$ tiene el bit $i$. La idea: cada subconjunto $s \subseteq m$ se obtiene quitando bits uno a uno; asi, sumar incrementalmente captura todos los subconjuntos en $O(N \cdot 2^N)$.

\paragraph{Propiedades clave (invariantes).}
\begin{itemize}\setlength\itemsep{0pt}
	\item Inicializar $dp = f$.
	\item Para cada bit $i$: \texttt{if (mask \& (1 << i)) dp[mask] += dp[mask \^{} (1 << i)]}.
	\item Variantes: supersets DP (sumar sobre superconjuntos).
	\item El orden de iteracion (bit externo, mascara interno) es crucial.
\end{itemize}

\paragraph{Codigo clave.}
SOS DP (Sum over Subsets):
\begin{verbatim}
vector<long long> dp(1 << N);
for (int i = 0; i < (1 << N); ++i) dp[i] = f[i];
for (int bit = 0; bit < N; ++bit)
    for (int mask = 0; mask < (1 << N); ++mask)
        if (mask & (1 << bit))
            dp[mask] += dp[mask ^ (1 << bit)];
// dp[mask] = suma de f(sub) para sub ⊆ mask
\end{verbatim}

\paragraph{Complejidad.}
\begin{itemize}\setlength\itemsep{0pt}
	\item $O(N \cdot 2^N)$.
	\item Memoria $O(2^N)$.
\end{itemize}

\paragraph{Donde tocar para modificar.}
\begin{itemize}\setlength\itemsep{0pt}
	\item \textbf{Max/min en lugar de suma}: cambiar \texttt{+=} por la operacion correspondiente.
	\item \textbf{Supersets DP}: iterar \texttt{mask} de $0$ a $2^N$ y aniadir \texttt{dp[mask] += dp[mask | (1<<i)]}.
	\item \textbf{XOR convolution} (subset convolution): version mas avanzada.
	\item \textbf{Aplicaciones}: contar subconjuntos con propiedades especificas.
\end{itemize}

\paragraph{Trampas y errores frecuentes.}
\begin{itemize}\setlength\itemsep{0pt}
	\item Si $N$ es 30 o mas, $2^N$ no entra en memoria; usar trucos o reducir $N$.
	\item El sentido del XOR (sumar a \texttt{mask \^{} bit}) es para ``agregar el bit $i$ al subconjunto''.
	\item En supersets DP, la condicion es \texttt{(mask \& (1<<bit)) == 0}.
\end{itemize}

\paragraph{Codigo.}
Ver \texttt{cpp.tex}, seccion \textit{SOS DP}.

\vspace{0.5em}

\subsubsection{Digit DP template}

\paragraph{Idea intuitiva.}
Contar numeros en $[0, N]$ que cumplen una propiedad sobre sus digitos. La idea: DP por posicion, manteniendo:
\begin{itemize}\setlength\itemsep{0pt}
	\item \texttt{pos}: posicion actual (de izquierda a derecha).
	\item \texttt{tight}: si los digitos ya colocados son exactamente los de $N$ (limite superior estricto).
	\item \texttt{state}: estado adicional (digitos usados, suma, etc.).
\end{itemize}
La clave: la condicion \texttt{tight} permite acotar correctamente el siguiente digito; sin ella, la DP contaria incorrectamente.

\paragraph{Propiedades clave (invariantes).}
\begin{itemize}\setlength\itemsep{0pt}
	\item Si \texttt{tight} es true, el siguiente digito esta acotado por el de $N$; si no, puede ser 0-9.
	\item Memoria: $O(D \cdot 2 \cdot \text{states})$.
	\item Caso base: si \texttt{pos == D}, devolver 1 si el estado cumple la propiedad, sino 0.
	\item El bit \texttt{started} distingue ``numero con menos digitos'' (con leading zeros).
\end{itemize}

\paragraph{Codigo clave.}
La estructura basica:
\begin{verbatim}
long long dp(int pos, bool tight, State s) {
    if (pos == D) return check(s) ? 1 : 0;
    if (memo[pos][tight][s] != -1) return memo;
    int limit = tight ? digits[pos] : 9;
    long long ans = 0;
    for (int d = 0; d <= limit; ++d)
        ans += dp(pos + 1, tight && (d == limit), update(s, d, pos));
    return memo[pos][tight][s] = ans;
}
\end{verbatim}

\paragraph{Complejidad.}
\begin{itemize}\setlength\itemsep{0pt}
	\item $O(D \cdot 10 \cdot \text{states})$ donde $D$ es el numero de digitos.
	\item Para $D \le 19$ (hasta $10^{18}$): $\sim 200 \cdot \text{states}$.
\end{itemize}

\paragraph{Donde tocar para modificar.}
\begin{itemize}\setlength\itemsep{0pt}
	\item \textbf{Cambiar la propiedad}: modificar \texttt{updateState} y la condicion base.
	\item \textbf{Multiples estados}: aniadir mas dimensiones al \texttt{memo}.
	\item \textbf{Leading zeros}: aniadir un flag \texttt{started} para distinguir numeros con menos digitos.
	\item \textbf{Rango $[L, R]$}: calcular $f(R) - f(L-1)$.
\end{itemize}

\paragraph{Trampas y errores frecuentes.}
\begin{itemize}\setlength\itemsep{0pt}
	\item Olvidar resetear \texttt{memo} entre queries.
	\item Si la cantidad de estados es grande, el \texttt{memo[pos][tight][state]} puede ser muy grande; usar \texttt{map}.
	\item La condicion \texttt{started} es crucial para numeros con menos digitos que $D$.
	\item Si $N$ es muy grande ($> 10^{18}$), aniadir soporte para long long.
\end{itemize}

\paragraph{Codigo.}
Ver \texttt{cpp.tex}, seccion \textit{Digit DP template}.

\vspace{0.5em}

\subsubsection{Tree DP (reroot)}

\paragraph{Idea intuitiva.}
Calcular alguna propiedad asumiendo que \textbf{cada nodo es la raiz} del arbol. La tecnica clasica de \textbf{rerooting} en dos pasadas:
\begin{itemize}\setlength\itemsep{0pt}
	\item \textbf{DFS 1}: desde una raiz arbitraria, calcular \texttt{dp[v]} para el subarbol de $v$.
	\item \textbf{DFS 2}: para cada hijo $u$ de $v$, calcular \texttt{up[u]} = respuesta si la raiz fuera $u$.
\end{itemize}
La demostracion: tras la primera DFS, cada nodo conoce ``lo que aporta su subarbol''; en la segunda, se transmite ``lo que aporta el resto del arbol''.

\paragraph{Propiedades clave (invariantes).}
\begin{itemize}\setlength\itemsep{0pt}
	\item \texttt{dp[v]} depende solo del subarbol de $v$ (con la raiz original).
	\item \texttt{up[u]} se calcula de manera que cuando la raiz es $u$, ``todo lo que estaba en $v$ se ajusta''.
	\item Formula para reroot depende del problema; el snippet usa el ejemplo ``suma de distancias''.
	\item La combinacion \texttt{dp[u] + up[u]} da la respuesta con raiz $u$.
\end{itemize}

\paragraph{Codigo clave.}
Reroot para suma de distancias:
\begin{verbatim}
// dfs1: dp[v] = suma de distancias desde v al subarbol
void dfs1(int v, int p) {
    dp[v] = 0;
    for (int u : adj[v]) if (u != p) {
        dfs1(u, v);
        dp[v] += dp[u] + sz[u];  // cada hijo aporta su subarbol
    }
}
// dfs2: up[u] = distancia desde el resto del arbol a u
void dfs2(int v, int p) {
    for (int u : adj[v]) if (u != p) {
        up[u] = up[v] + dp[v] - dp[u] - sz[u] + (n - sz[u]);
        dfs2(u, v);
    }
}
\end{verbatim}

\paragraph{Complejidad.}
\begin{itemize}\setlength\itemsep{0pt}
	\item $O(N)$ para ambas pasadas.
	\item Cada arista se procesa un numero constante de veces.
\end{itemize}

\paragraph{Donde tocar para modificar.}
\begin{itemize}\setlength\itemsep{0pt}
	\item \textbf{Cambiar la cantidad a calcular}: modificar la operacion en \texttt{dfs1} y la formula en \texttt{dfs2}.
	\item \textbf{Arbol con pesos}: aniadir \texttt{weights[v]} y ajustar.
	\item \textbf{Re-root en arboles con ciclos}: requiere otra tecnica (DSU on tree).
	\item \textbf{Rooted trees diferentes}: si la propiedad no es invariante por raiz, usar otra tecnica.
\end{itemize}

\paragraph{Trampas y errores frecuentes.}
\begin{itemize}\setlength\itemsep{0pt}
	\item La formula de \texttt{up[u]} debe sumar y restar correctamente las contribuciones de $u$ y del resto.
	\item Para arboles con pesos, aniadir el peso en las formulas.
	\item Verificar con ejemplos pequenos antes de generalizar.
\end{itemize}

\paragraph{Codigo.}
Ver \texttt{cpp.tex}, seccion \textit{Tree DP (reroot)}.

% ============================================================
\section{Geometria}
% ============================================================

\subsection{Puntos y segmentos}

\subsubsection{Point + cross + dot}

\paragraph{Idea intuitiva.}
La estructura \texttt{Point} representa un punto en 2D con coordenadas enteras. Las dos operaciones vectoriales clave son:
\begin{itemize}\setlength\itemsep{0pt}
	\item \textbf{Cross product} $P \times Q = P_x Q_y - P_y Q_x$. Da el ``area con signo'' del paralelogramo formado por $P$ y $Q$. Usado para orientacion.
	\item \textbf{Dot product} $P \cdot Q = P_x Q_x + P_y Q_y$. Da el producto de las longitudes por el coseno del angulo. Usado para proyeccion.
\end{itemize}
La interpretacion geometrica del cross product: es el area con signo del paralelogramo; el signo indica la orientacion relativa.

\paragraph{Propiedades clave (invariantes).}
\begin{itemize}\setlength\itemsep{0pt}
	\item \texttt{cross(a, b) > 0} si $b$ esta a la izquierda de $a$ (sistema de coordenadas usual).
	\item \texttt{cross(a, b) = 0} si son colineales.
	\item \texttt{dot(a, b) > 0} si el angulo es agudo.
	\item El snippet define \texttt{cross} respecto a \texttt{this}, util para orientacion de 3 puntos.
	\item $|cross|^2 + |dot|^2 = |P|^2 \cdot |Q|^2$ (identidad de Lagrange).
\end{itemize}

\paragraph{Codigo clave.}
Las dos operaciones vectoriales:
\begin{verbatim}
long long cross(const Point& a, const Point& b) {
    return (long long)a.x * b.y - (long long)a.y * b.x;
}
long long dot(const Point& a, const Point& b) {
    return (long long)a.x * b.x + (long long)a.y * b.y;
}
\end{verbatim}

\paragraph{Complejidad.}
\begin{itemize}\setlength\itemsep{0pt}
	\item $O(1)$ por operacion.
	\item Constante muy pequena (4 multiplicaciones + 2 sumas).
\end{itemize}

\paragraph{Donde tocar para modificar.}
\begin{itemize}\setlength\itemsep{0pt}
	\item \textbf{Coordenadas doubles}: cambiar \texttt{int} por \texttt{double} o \texttt{long double}. Cuidado con la precision.
	\item \textbf{3D}: aniadir \texttt{z} y adaptar cross/dot.
	\item \textbf{Operadores adicionales}: aniadir \texttt{operator+}, \texttt{operator-}, \texttt{operator*} (escalar).
	\item \textbf{Distancia al cuadrado}: aniadir \texttt{len2()} que retorna $x^2 + y^2$.
\end{itemize}

\paragraph{Trampas y errores frecuentes.}
\begin{itemize}\setlength\itemsep{0pt}
	\item Cross product \textbf{no es} conmutativo; $P \times Q = -Q \times P$.
	\item Overflow con coordenadas grandes ($|x|, |y| > 10^4$) si se hace \texttt{x * y}.
	\item El signo de cross depende del sistema de coordenadas (en computacion, usualmente CCW positivo).
	\item Para ``igualdad'' de doubles, aniadir tolerancia \texttt{eps}.
\end{itemize}

\paragraph{Codigo.}
Ver \texttt{cpp.tex}, seccion \textit{Point + cross + dot}.

\vspace{0.5em}

\subsubsection{Interseccion de segmentos}

\paragraph{Idea intuitiva.}
Dados dos segmentos $[a, b]$ y $[c, d]$, determinar si se intersectan. La idea clasica:
\begin{itemize}\setlength\itemsep{0pt}
	\item Verificar orientaciones: $a, b, c$ y $a, b, d$ deben estar en lados opuestos de la recta $ab$; igual para $cd$ y $ab$.
	\item Caso colineal: si las rectas son paralelas (\texttt{cross} $= 0$) y los segmentos son colineales, verificar que los bounding boxes se intersectan.
\end{itemize}
La demostracion: dos segmentos se intersectan sii cada segmento ``cruza'' la linea que contiene al otro (estricto) o ambos son colineales con solapamiento.

\paragraph{Propiedades clave (invariantes).}
\begin{itemize}\setlength\itemsep{0pt}
	\item Funciona con segmentos en cualquier orientacion.
	\item Caso colineal: requiere check adicional de bounding box.
	\item Para intersectar en un \textbf{punto} especifico: aniadir division con cuidado de precision.
	\item Funciona con segmentos verticales/horizontales sin casos especiales.
\end{itemize}

\paragraph{Codigo clave.}
Test clasico de interseccion:
\begin{verbatim}
int sign(long long x) { return (x > 0) - (x < 0); }
bool intersect(Point a, Point b, Point c, Point d) {
    long long o1 = cross(b - a, c - a);
    long long o2 = cross(b - a, d - a);
    long long o3 = cross(d - c, a - c);
    long long o4 = cross(d - c, b - c);
    if (sign(o1) != sign(o2) && sign(o3) != sign(o4)) return true;
    // casos colineales
    if (o1 == 0 && onSegment(a, b, c)) return true;
    if (o2 == 0 && onSegment(a, b, d)) return true;
    if (o3 == 0 && onSegment(c, d, a)) return true;
    if (o4 == 0 && onSegment(c, d, b)) return true;
    return false;
}
\end{verbatim}

\paragraph{Complejidad.}
\begin{itemize}\setlength\itemsep{0pt}
	\item $O(1)$.
	\item Constante pequena.
\end{itemize}

\paragraph{Donde tocar para modificar.}
\begin{itemize}\setlength\itemsep{0pt}
	\item \textbf{Interseccion con punto}: aniadir la division despues de verificar orientacion.
	\item \textbf{Interseccion rayo-segmento}: ajustar la primera condicion.
	\item \textbf{Poligonos convexos}: usar separacion de ejes (SAT).
	\item \textbf{Multiples queries}: precomputar bounding boxes o usar segment tree.
\end{itemize}

\paragraph{Trampas y errores frecuentes.}
\begin{itemize}\setlength\itemsep{0pt}
	\item Si los segmentos tocan solo en un extremo, el snippet lo considera interseccion (algunos problemas quieren estricto). Ajustar segun el problema.
	\item El caso colineal es comun; olvidar el check adicional de bounding box da falsos negativos.
	\item En problemas con muchos segmentos, considerar sweep line para $O(n \log n)$.
\end{itemize}

\paragraph{Codigo.}
Ver \texttt{cpp.tex}, seccion \textit{Interseccion de segmentos}.

\vspace{0.5em}

\subsubsection{Contains colineal}

\paragraph{Idea intuitiva.}
Verifica si un punto $c$ esta sobre el segmento $[a, b]$ (asumiendo que ya son colineales). Si \texttt{cross(a, b, c) = 0} y $c$ esta dentro del bounding box de $[a, b]$, entonces $c$ esta en el segmento. La razon: si los tres puntos son colineales, $c$ esta en el segmento sii sus coordenadas estan entre las de $a$ y $b$.

\paragraph{Propiedades clave (invariantes).}
\begin{itemize}\setlength\itemsep{0pt}
	\item Asume colinealidad (no la verifica mas alla de \texttt{cross == 0}).
	\item Usa comparaciones con \texttt{<=} para extremos \textbf{inclusivos}.
	\item Funciona con segmentos verticales (la condicion se reduce a una sola coordenada).
\end{itemize}

\paragraph{Codigo clave.}
Test de contencion en segmento (asume colinealidad):
\begin{verbatim}
bool onSegment(Point a, Point b, Point c) {
    return min(a.x, b.x) <= c.x && c.x <= max(a.x, b.x) &&
           min(a.y, b.y) <= c.y && c.y <= max(a.y, b.y);
}
\end{verbatim}

\paragraph{Complejidad.}
\begin{itemize}\setlength\itemsep{0pt}
	\item $O(1)$.
	\item Constante minima.
\end{itemize}

\paragraph{Donde tocar para modificar.}
\begin{itemize}\setlength\itemsep{0pt}
	\item \textbf{Contener estricto}: cambiar \texttt{<=} a \texttt{<} para excluir extremos.
	\item \textbf{Coordenadas doubles}: usar tolerancia (\texttt{eps}).
	\item \textbf{Solo x o solo y}: aniadir check especifico para segmentos verticales/horizontales.
\end{itemize}

\paragraph{Trampas y errores frecuentes.}
\begin{itemize}\setlength\itemsep{0pt}
	\item No usar este snippet para puntos \textbf{no colineales}; siempre verificar primero la colinealidad.
	\item Si los puntos tienen doubles, el chequeo \texttt{cross == 0} puede fallar; aniadir tolerancia.
	\item El bounding box incluye los extremos; verificar si el problema quiere exclusivo o inclusivo.
\end{itemize}

\paragraph{Codigo.}
Ver \texttt{cpp.tex}, seccion \textit{Contains colineal}.

\subsection{Poligonos}

\subsubsection{Convex Hull (monotono)}

\paragraph{Idea intuitiva.}
Algoritmo de \textbf{Andrew's monotone chain}: ordenar puntos por $(x, y)$, luego construir la hull inferior y superior por separado. Para cada punto, aniadirlo al hull y, mientras el ultimo giro no sea ``counter-clockwise'' (\texttt{cross > 0}), popear. La demostracion: si los puntos se procesan en orden, el hull inferior y superior cubren la frontera completa sin auto-intersecciones.

\paragraph{Propiedades clave (invariantes).}
\begin{itemize}\setlength\itemsep{0pt}
	\item Ordena los puntos en $O(n \log n)$.
	\item Construye la hull en $O(n)$ despues.
	\item Output: poligono convexo en counter-clockwise, sin punto final duplicado.
	\item El \texttt{<= 0} incluye colineales; usar \texttt{< 0} para excluirlos.
	\item Cada punto aparece en la hull inferior o superior (no ambos).
\end{itemize}

\paragraph{Codigo clave.}
Construccion del hull inferior y superior:
\begin{verbatim}
sort(p.begin(), p.end());
vector<Point> lower, upper;
for (auto& p : points) {
    while (lower.size() >= 2 && cross(lower.back() - lower[lower.size()-2],
                                       p - lower.back()) <= 0)
        lower.pop_back();
    lower.push_back(p);
}
for (int i = points.size() - 1; i >= 0; --i) {
    auto& p = points[i];
    while (upper.size() >= 2 && cross(upper.back() - upper[upper.size()-2],
                                       p - upper.back()) <= 0)
        upper.pop_back();
    upper.push_back(p);
}
vector<Point> hull = lower;
hull.insert(hull.end(), upper.begin() + 1, upper.end() - 1);
\end{verbatim}

\paragraph{Complejidad.}
\begin{itemize}\setlength\itemsep{0pt}
	\item $O(n \log n)$.
	\item Constante muy pequena (solo cross products).
\end{itemize}

\paragraph{Donde tocar para modificar.}
\begin{itemize}\setlength\itemsep{0pt}
	\item \textbf{Colineales}: cambiar \texttt{<= 0} por \texttt{< 0} para excluirlos del hull.
	\item \textbf{Puntos repetidos}: \texttt{unique} despues de sort.
	\item \textbf{Hull 3D}: gift wrapping es $O(n^2)$; convex hull 3D es mas complejo.
	\item \textbf{Output con los puntos}: aniadir el hull completo (no solo upper/lower).
\end{itemize}

\paragraph{Trampas y errores frecuentes.}
\begin{itemize}\setlength\itemsep{0pt}
	\item El \texttt{pop\_back} y el \texttt{reverse} son esenciales para no duplicar el ultimo punto.
	\item Para excluir colineales, ajustar el \texttt{<= 0} vs \texttt{< 0}.
	\item Si los puntos estan duplicados, aniadir \texttt{unique} antes del sort.
\end{itemize}

\paragraph{Codigo.}
Ver \texttt{cpp.tex}, seccion \textit{Convex Hull (monotono)}.

\vspace{0.5em}

\subsubsection{Polygon Area}

\paragraph{Idea intuitiva.}
El area (con signo) de un poligono se calcula con la \textbf{Shoelace formula}: $A = \frac{1}{2} \sum_{i=0}^{n-1} (x_i y_{i+1} - x_{i+1} y_i)$, donde el indice es modulo $n$. El snippet retorna el \textbf{doble} del area (en valor absoluto). La formula sale del ``sumar areas de trapezoides'': cada lado del poligono define un trapezoide con el eje $x$; la suma con signo da el area.

\paragraph{Propiedades clave (invariantes).}
\begin{itemize}\setlength\itemsep{0pt}
	\item El signo indica la orientacion: positivo si counter-clockwise.
	\item Para el doble: no dividir por 2 si no es necesario (evita fracciones).
	\item Funciona para poligonos simples (no auto-intersectantes).
	\item Si el poligono es no convexo, la formula da el area ``con signo''.
\end{itemize}

\paragraph{Codigo clave.}
La suma de trapezoides:
\begin{verbatim}
long long shoelace(const vector<Point>& p) {
    long long s = 0;
    int n = p.size();
    for (int i = 0; i < n; ++i)
        s += (long long)p[i].x * p[(i+1)%n].y - (long long)p[(i+1)%n].x * p[i].y;
    return abs(s);  // doble del area
}
\end{verbatim}

\paragraph{Complejidad.}
\begin{itemize}\setlength\itemsep{0pt}
	\item $O(n)$.
	\item Constante minima.
\end{itemize}

\paragraph{Donde tocar para modificar.}
\begin{itemize}\setlength\itemsep{0pt}
	\item \textbf{Area exacta}: dividir por 2 (con cuidado si da fraccion; usar \texttt{double}).
	\item \textbf{Area firmada}: quitar \texttt{abs} y dejar el signo.
	\item \textbf{Punto dentro de poligono convexo}: con la hull, chequear que esta a un mismo lado de todas las aristas.
	\item \textbf{Triangulacion}: descomponer en triangulos y sumar areas.
\end{itemize}

\paragraph{Trampas y errores frecuentes.}
\begin{itemize}\setlength\itemsep{0pt}
	\item Si los puntos no estan en orden (CCW o CW), el area firmada puede ser incorrecta. Para signed area consistente, usar siempre el mismo orden.
	\item Para poligonos con holes, sumar las areas con signo apropiado.
	\item Con doubles, aniadir tolerancia \texttt{eps}.
\end{itemize}

\paragraph{Codigo.}
Ver \texttt{cpp.tex}, seccion \textit{Polygon Area}.

\vspace{0.5em}

\subsubsection{Distance point-segment / point-line}

\paragraph{Idea intuitiva.}
Distancia minima de un punto $p$ a un segmento $[a, b]$ o a la recta que pasa por $a$ y $b$.
\begin{itemize}\setlength\itemsep{0pt}
	\item \textbf{Linea}: $d = \frac{|ab \times ap|}{|ab|}$.
	\item \textbf{Segmento}: si la proyeccion de $p$ sobre $ab$ cae fuera del segmento, la distancia es al extremo mas cercano.
\end{itemize}
La razon geometrica: la distancia a una linea es el area del paralelogramo dividido por la base; la proyeccion ortogonal indica si estamos ``dentro'' o ``fuera'' del segmento.

\paragraph{Propiedades clave (invariantes).}
\begin{itemize}\setlength\itemsep{0pt}
	\item \texttt{dist2PointLine} retorna distancia \textbf{al cuadrado}.
	\item \texttt{dist2PointSegment} distingue 3 casos: $p$ antes de $a$, despues de $b$, o proyectado sobre el segmento.
	\item Si \texttt{len2 = 0} (degenerate), la distancia es $|ap|$ o $|bp|$.
	\item La proyeccion se calcula con dot products y el parametro $t \in [0, 1]$.
\end{itemize}

\paragraph{Codigo clave.}
Distancia al cuadrado a segmento:
\begin{verbatim}
long long dist2Segment(Point p, Point a, Point b) {
    long long len2 = (a-b).len2();
    if (len2 == 0) return (p-a).len2();  // degenerate
    long long t = max(0LL, min(1LL, dot(p-a, b-a) / len2));
    Point proj = {a.x + t*(b.x - a.x), a.y + t*(b.y - a.y)};
    return (p - proj).len2();
}
\end{verbatim}

\paragraph{Complejidad.}
\begin{itemize}\setlength\itemsep{0pt}
	\item $O(1)$.
	\item Constante minima.
\end{itemize}

\paragraph{Donde tocar para modificar.}
\begin{itemize}\setlength\itemsep{0pt}
	\item \textbf{Distancia (no al cuadrado)}: aniadir \texttt{sqrt} al final.
	\item \textbf{Coordenadas doubles}: aniadir \texttt{long double} y tolerancia.
	\item \textbf{3D}: aniadir coordenada $z$ y adaptar.
	\item \textbf{Multiples queries}: usar KD-tree o segment tree para acelerar.
\end{itemize}

\paragraph{Trampas y errores frecuentes.}
\begin{itemize}\setlength\itemsep{0pt}
	\item Overflow en \texttt{area2 * area2} si el area es grande; \texttt{long long} alcanza para $|x|, |y| \le 10^9$ pero no mas.
	\item Si $a = b$ (segmento degenerado), aniadir check explicito.
	\item En doubles, aniadir tolerancia \texttt{eps} para comparar.
\end{itemize}

\paragraph{Codigo.}
Ver \texttt{cpp.tex}, seccion \textit{Distance point-segment / point-line}.

\vspace{0.5em}

\subsubsection{Projection of point onto line}

\paragraph{Idea intuitiva.}
Proyeccion ortogonal de $p$ sobre la recta que pasa por $a$ y $b$. El parametro $t$ tal que la proyeccion es $a + t \cdot (b - a)$ es:
\[ t = \frac{(p - a) \cdot (b - a)}{|b - a|^2}. \]
La idea: descomponemos $\vec{ap}$ en una componente paralela a $\vec{ab}$ y una perpendicular; $t$ mide la posicion en la primera.

\paragraph{Propiedades clave (invariantes).}
\begin{itemize}\setlength\itemsep{0pt}
	\item Si $t \in [0, 1]$, la proyeccion cae sobre el \textbf{segmento}; si no, cae sobre la extension.
	\item Retorna \texttt{double} (puede tener precision issues).
	\item $t < 0$: $p$ esta ``antes'' de $a$; $t > 1$: $p$ esta ``despues'' de $b$.
\end{itemize}

\paragraph{Codigo clave.}
Proyeccion ortogonal:
\begin{verbatim}
Point projection(Point p, Point a, Point b) {
    Point ab = b - a, ap = p - a;
    double t = (double)dot(ap, ab) / ab.len2();
    return {a.x + t * ab.x, a.y + t * ab.y};
}
\end{verbatim}

\paragraph{Complejidad.}
\begin{itemize}\setlength\itemsep{0pt}
	\item $O(1)$.
	\item Constante minima (dos dot products y una division).
\end{itemize}

\paragraph{Donde tocar para modificar.}
\begin{itemize}\setlength\itemsep{0pt}
	\item \textbf{Proyeccion sobre segmento}: clipar $t$ a $[0, 1]$ antes de calcular el punto.
	\item \textbf{Distancia al segmento via proyeccion}: combinar con el modulo anterior.
	\item \textbf{Solo el parametro $t$}: si solo necesitas saber si esta dentro del segmento, retorna $t$ sin construir el punto.
\end{itemize}

\paragraph{Trampas y errores frecuentes.}
\begin{itemize}\setlength\itemsep{0pt}
	\item Si $a = b$ ($|ab| = 0$), la proyeccion no esta definida; aniadir guard.
	\item En doubles, aniadir tolerancia \texttt{eps} para comparar $t$ con 0 o 1.
	\item La division por \texttt{len2} puede ser cara si las coordenadas son grandes; usar \texttt{double} o precomputar.
\end{itemize}

\paragraph{Codigo.}
Ver \texttt{cpp.tex}, seccion \textit{Projection of point onto line}.

\vspace{0.5em}

\subsubsection{Point in polygon (ray casting)}

\paragraph{Idea intuitiva.}
Algoritmo de \textbf{ray casting}: contar cuantas veces un rayo horizontal desde $p$ cruza los bordes del poligono. Si es impar, $p$ esta dentro; si par, fuera. El snippet implementa la version clasica recorriendo aristas. La razon: cada cruce cambia el lado del rayo respecto al poligono; un numero impar de cambios significa ``terminamos dentro''.

\paragraph{Propiedades clave (invariantes).}
\begin{itemize}\setlength\itemsep{0pt}
	\item Funciona con poligonos simples (no necesariamente convexos).
	\item Las intersecciones en vertices cuentan especialmente.
	\item Para poligonos con \textbf{huecos}, sigue funcionando.
	\item La respuesta es estrictamente 0 (fuera) o 1 (dentro) para poligonos simples.
\end{itemize}

\paragraph{Codigo clave.}
El test clasico de ray casting:
\begin{verbatim}
bool pointInPolygon(Point p, const vector<Point>& poly) {
    int n = poly.size();
    bool inside = false;
    for (int i = 0, j = n - 1; i < n; j = i++) {
        if ((poly[i].y > p.y) != (poly[j].y > p.y) &&
            p.x < (poly[j].x - poly[i].x) * (p.y - poly[i].y) /
                    (poly[j].y - poly[i].y) + poly[i].x)
            inside = !inside;
    }
    return inside;
}
\end{verbatim}

\paragraph{Complejidad.}
\begin{itemize}\setlength\itemsep{0pt}
	\item $O(n)$ por query.
	\item Memoria $O(1)$ (no se almacena el rayo).
\end{itemize}

\paragraph{Donde tocar para modificar.}
\begin{itemize}\setlength\itemsep{0pt}
	\item \textbf{Winding number}: version alternativa que cuenta el numero de vueltas.
	\item \textbf{Poligono convexo}: check lineal contra las aristas (mas rapido).
	\item \textbf{Multiples queries}: precomputar o usar segment tree sobre las aristas.
	\item \textbf{Punto en frontera}: aniadir check adicional si es colineal con alguna arista.
\end{itemize}

\paragraph{Trampas y errores frecuentes.}
\begin{itemize}\setlength\itemsep{0pt}
	\item El manejo del vertice (cuando el rayo pasa por un vertice) puede causar double-counting; el snippet lo evita con la condicion \texttt{(poly[i].y > p.y) != (poly[j].y > p.y)}.
	\item Si el poligono tiene holes, el algoritmo sigue funcionando.
	\item Para poligonos no simples (auto-intersectantes), el resultado puede ser ambiguo.
\end{itemize}

\paragraph{Codigo.}
Ver \texttt{cpp.tex}, seccion \textit{Point in polygon (ray casting)}.

\vspace{0.5em}

\subsubsection{Minkowski sum (convex polygons)}

\paragraph{Idea intuitiva.}
La \textbf{suma de Minkowski} $A \oplus B$ de dos poligonos convexos es otro poligono convexo que representa $\{a + b : a \in A, b \in B\}$. Algoritmo: ordenar las aristas por angulo polar (CCW) y ``mergearlas'' como en merge sort; el resultado es un poligono CCW sin repetir el ultimo punto. La idea: el resultado tiene una arista paralela a cada arista de los originales; la suma se obtiene barriendo los angulos.

\paragraph{Propiedades clave (invariantes).}
\begin{itemize}\setlength\itemsep{0pt}
	\item Si $A, B$ son convexos con $n, m$ vertices, $A \oplus B$ tiene $\le n + m$ vertices.
	\item Util para \textbf{detectar colision}: $A$ colisiona con $B$ si $0 \in A \oplus (-B)$.
	\item Asume poligonos en orden \textbf{counter-clockwise}.
	\item El resultado es convexo (suma de convexos es convexo).
\end{itemize}

\paragraph{Codigo clave.}
La suma de Minkowski:
\begin{verbatim}
vector<Point> minkowski(const vector<Point>& A, const vector<Point>& B) {
    vector<Point> result;
    int i = 0, j = 0;
    int n = A.size(), m = B.size();
    do {
        result.push_back(A[i] + B[j]);
        long long cross = (A[(i+1)%n] - A[i]) ^ (B[(j+1)%m] - B[j]);
        if (cross >= 0) i = (i + 1) % n;
        if (cross <= 0) j = (j + 1) % m;
    } while (i || j);
    return result;
}
\end{verbatim}

\paragraph{Complejidad.}
\begin{itemize}\setlength\itemsep{0pt}
	\item $O(n + m)$.
	\item Memoria $O(n + m)$.
\end{itemize}

\paragraph{Donde tocar para modificar.}
\begin{itemize}\setlength\itemsep{0pt}
	\item \textbf{Deteccion de colision}: aniadir $(-B)$ (reflejar) y verificar si $0$ esta en el resultado.
	\item \textbf{Poligonos no convexos}: la suma puede no ser poligono simple; algoritmo distinto.
	\item \textbf{Multiples queries}: precomputar las aristas una vez.
	\item \textbf{Robot motion planning}: usar para planear trayectorias.
\end{itemize}

\paragraph{Trampas y errores frecuentes.}
\begin{itemize}\setlength\itemsep{0pt}
	\item Si los poligonos no son CCW, la suma es incorrecta. Normalizar antes.
	\item El bucle termina cuando ambos indices vuelven a 0; aniadir check.
	\item Para la deteccion de colision, aniadir el test \texttt{0 in A \^{} (-B)}.
\end{itemize}

\paragraph{Codigo.}
Ver \texttt{cpp.tex}, seccion \textit{Minkowski sum (convex polygons)}.

\subsection{Herramientas}

\subsubsection{Li Chao Tree (long long)}

\paragraph{Idea intuitiva.}
Mantiene un \textbf{lower envelope} (o upper) de un conjunto de rectas $y = mx + b$ sobre un dominio $[X_{LO}, X_{HI}]$ y responde ``minimo $y$ en un $x$ dado'' en $O(\log X)$. La estructura es un segment tree donde cada nodo guarda la recta ``mejor'' para su segmento medio; al aniadir una recta, se compara con la actual y se decide cual es mejor en cada mitad, descendiendo recursivamente. La demostracion: en cada nivel, solo se guarda la mejor recta para el ``punto medio''; las otras se transmiten a los hijos.

\paragraph{Propiedades clave (invariantes).}
\begin{itemize}\setlength\itemsep{0pt}
	\item \textbf{Online}: cada \texttt{addLine} en $O(\log X)$.
	\item Solo insercion; no hay borrar (para borrar, usar Li Chao segmentado con reset).
	\item La recta guardada en un nodo es la ``mejor'' en su punto medio.
	\item Las rectas ``perdedoras'' se transmiten a los hijos para futuros inserts.
\end{itemize}

\paragraph{Codigo clave.}
La insercion en Li Chao:
\begin{verbatim}
void addLine(Line nw, int v = 1, int l = X_LO, int r = X_HI) {
    int m = (l + r) / 2;
    bool left = nw.get(l) < tree[v].get(l);
    bool mid = nw.get(m) < tree[v].get(m);
    if (mid) swap(tree[v], nw);
    if (r - l == 1) return;
    if (left != mid) addLine(nw, v*2, l, m);
    else             addLine(nw, v*2+1, m, r);
}
\end{verbatim}

\paragraph{Complejidad.}
\begin{itemize}\setlength\itemsep{0pt}
	\item $O(\log X)$ por insercion y query, donde $X$ es el tamano del dominio.
	\item Memoria $O(\log X)$ por insercion, $O(X)$ total.
\end{itemize}

\paragraph{Donde tocar para modificar.}
\begin{itemize}\setlength\itemsep{0pt}
	\item \textbf{Upper envelope en vez de lower}: invertir el signo (\texttt{>} en vez de \texttt{<}).
	\item \textbf{Queries de segmento} (no full line): si la recta solo esta activa en un rango, aniadir bandera \texttt{active}.
	\item \textbf{Eliminar rectas}: usar Li Chao segmentado (mas complejo).
	\item \textbf{Pendientes grandes}: si las pendientes tienen magnitud $> 10^9$, aniadir \texttt{\_\_int128} para evitar overflow.
\end{itemize}

\paragraph{Trampas y errores frecuentes.}
\begin{itemize}\setlength\itemsep{0pt}
	\item El parametro \texttt{lo == hi} debe retornar el valor del nodo (caso base).
	\item La comparacion \texttt{newLine.get(mid) < node->line.get(mid)} decide quien es mejor en el medio.
	\item Si las rectas tienen $m$ iguales pero $b$ distintos, se queda con la primera (tie-break).
\end{itemize}

\paragraph{Codigo.}
Ver \texttt{cpp.tex}, seccion \textit{Li Chao Tree (long long)}.

\vspace{0.5em}

\subsubsection{Sweep line template}

\paragraph{Idea intuitiva.}
Patron de diseno para problemas geometricos donde se barre una linea sobre el plano y se mantiene una estructura de datos sobre el otro eje. Eventos (add, query, remove) se ordenan por la coordenada de barrido y se procesan secuencialmente.

\paragraph{Propiedades clave (invariantes).}
\begin{itemize}\setlength\itemsep{0pt}
	\item Eventos: tuplas \texttt{(x, type, id)}.
	\item La estructura auxiliar (set, multiset, segtree) se mantiene sobre el otro eje.
	\item Procesar eventos en orden de $x$ ascendente.
\end{itemize}

\paragraph{Complejidad.}
\begin{itemize}\setlength\itemsep{0pt}
	\item Depende del problema; tipico $O((N + E) \log N)$.
\end{itemize}

\paragraph{Donde tocar para modificar.}
\begin{itemize}\setlength\itemsep{0pt}
	\item \textbf{Aniadir/quitar segmentos}: aniadir eventos add/remove al barrido.
	\item \textbf{Aniadir queries}: aniadir eventos query.
	\item \textbf{Cambiar el eje de barrido}: si barres en $y$ en vez de $x$, aniadir comparador.
\end{itemize}

\paragraph{Trampas y errores frecuentes.}
\begin{itemize}\setlength\itemsep{0pt}
	\item El snippet es solo el \textbf{esqueleto}; los handlers \texttt{add/query/remove} se llenan por problema.
	\item Cuidado con eventos duplicados.
\end{itemize}

\paragraph{Codigo.}
Ver \texttt{cpp.tex}, seccion \textit{Sweep line template}.

\vspace{0.5em}

\subsubsection{Closest pair of points}

\paragraph{Idea intuitiva.}
Encuentra el par de puntos mas cercano en $O(n \log n)$ con \textbf{divide y venceras}. Algoritmo:
\begin{itemize}\setlength\itemsep{0pt}
	\item Sort por $x$.
	\item Dividir en mitades; recursivo.
	\item Combinar: dado el minimo $d$ de las dos mitades, strip = puntos en $[midX - d, midX + d]$; ordenar por $y$; cada punto solo se compara con los siguientes hasta $y$ difiera en $< d$ (maximo 7).
\end{itemize}

\paragraph{Propiedades clave (invariantes).}
\begin{itemize}\setlength\itemsep{0pt}
	\item Requiere ordenar por $y$ despues (en el snippet se hace durante el merge).
	\item Strip tiene a lo sumo $O(n)$ puntos en total, pero la comparacion es $O(1)$ amortizada por punto.
	\item Resultado en \texttt{double}.
\end{itemize}

\paragraph{Complejidad.}
\begin{itemize}\setlength\itemsep{0pt}
	\item $O(n \log n)$.
\end{itemize}

\paragraph{Donde tocar para modificar.}
\begin{itemize}\setlength\itemsep{0pt}
	\item \textbf{Reconstruir el par}: aniadir indices en la recursion.
	\item \textbf{Distancia maxima}: cambiar \texttt{min} por \texttt{max} (mas facil; no requiere D\&C).
	\item \textbf{Puntos duplicados}: aniadir check especial (distancia 0).
\end{itemize}

\paragraph{Trampas y errores frecuentes.}
\begin{itemize}\setlength\itemsep{0pt}
	\item El merge requiere ordenar por $y$; el snippet lo hace in-place con un \texttt{tmp}.
	\item Si se hace en cada recursion, es $O(n \log n)$ extra.
\end{itemize}

\paragraph{Codigo.}
Ver \texttt{cpp.tex}, seccion \textit{Closest pair of points}.

\vspace{0.5em}

\subsubsection{Rotating calipers (polygon diameter)}

\paragraph{Idea intuitiva.}
Encuentra el \textbf{diametro} de un poligono convexo (par de vertices mas distantes) en $O(n)$. Asume poligono en orden counter-clockwise. Algoritmo: dos indices $i$ (sobre las aristas) y $k$ (sobre los vertices opuestos); en cada paso, rotar $k$ mientras el area del paralelogramo aumenta. La distancia $i$-$k$ puede ser maxima en cualquier momento del barrido.

\paragraph{Propiedades clave (invariantes).}
\begin{itemize}\setlength\itemsep{0pt}
	\item Cada indice avanza como maximo $n$ veces total.
	\item El ``diametro'' es la maxima distancia entre dos vertices.
	\item Retorna \textbf{distancia al cuadrado}.
\end{itemize}

\paragraph{Complejidad.}
\begin{itemize}\setlength\itemsep{0pt}
	\item $O(n)$.
\end{itemize}

\paragraph{Donde tocar para modificar.}
\begin{itemize}\setlength\itemsep{0pt}
	\item \textbf{Reconstruir el par de vertices}: aniadir indices.
	\item \textbf{Anchura minima / maxima}: variantes de calipers.
	\item \textbf{Distancia minima entre poligonos convexos}: otra variante.
\end{itemize}

\paragraph{Trampas y errores frecuentes.}
\begin{itemize}\setlength\itemsep{0pt}
	\item Asume poligono \textbf{estrictamente} convexo (sin vertices colineales). Si hay colineales, aniadir perturbacion o filtrar.
\end{itemize}

\paragraph{Codigo.}
Ver \texttt{cpp.tex}, seccion \textit{Rotating calipers (polygon diameter)}.

% ============================================================
\section{Miscelanea}
% ============================================================

\subsection{Utilidades del usuario}

\subsubsection{Hora}

\paragraph{Idea intuitiva.}
Una pequena clase para representar \textbf{tiempos en segundos} y hacer aritmetica basica sobre ellos. Internamente almacena todo en segundos (un solo \texttt{int}), lo cual simplifica comparaciones y diferencias. El constructor acepta $(h, m, s)$ y los convierte a segundos. Hay un metodo \texttt{cut()} que fuerza un minimo (util para ``horas trabajadas'').

\paragraph{Propiedades clave (invariantes).}
\begin{itemize}\setlength\itemsep{0pt}
	\item Representacion canonica: \textbf{solo segundos}. No hay zonas horarias, no hay DST.
	\item Comparadores \texttt{<} y \texttt{>} operan sobre los segundos.
	\item \texttt{operator-} retorna la \textbf{diferencia absoluta} en segundos (siempre no negativa).
	\item \texttt{cut()} aplica un piso (por defecto $8\text{h }30\text{m} = 30600$ segundos).
\end{itemize}

\paragraph{Complejidad.}
\begin{itemize}\setlength\itemsep{0pt}
	\item $O(1)$ por operacion.
\end{itemize}

\paragraph{Donde tocar para modificar.}
\begin{itemize}\setlength\itemsep{0pt}
	\item \textbf{Aniadir suma}: \texttt{operator+} que devuelve un \texttt{Hour} con la suma de segundos.
	\item \textbf{Aniadir output formateado}: aniadir \texttt{show\_hms()} que imprime $H$:$M$:$S$.
	\item \textbf{Cambiar el ``corte''}: modificar el valor magico \texttt{30'600} en \texttt{cut()} o pasarlo como parametro.
	\item \textbf{Tamano maximo}: si los tiempos exceden $\sim 68$ anos, el \texttt{int} desborda; usar \texttt{long long}.
	\item \textbf{Precision sub-segundo}: aniadir \texttt{ms} (milisegundos) y guardar en \texttt{long long} con factor 1000.
\end{itemize}

\paragraph{Trampas y errores frecuentes.}
\begin{itemize}\setlength\itemsep{0pt}
	\item El operador \texttt{-} retorna \textbf{valor absoluto}, no diferencia con signo. Si necesitas diferencia con signo, aniadir un nuevo metodo o sobrecargar.
	\item \texttt{cut()} modifica \texttt{s} \textbf{in-place}; si el llamador no espera esto, aniadir una version \texttt{cut(min)} que retorna un nuevo \texttt{Hour}.
	\item El umbral magico \texttt{30'600} (8h 30m) esta hardcoded; documentar bien que es ``jornada minima''.
\end{itemize}

\paragraph{Codigo.}
Ver \texttt{cpp.tex}, seccion \textit{Hora}.

\end{document}
', aniadir un puntero a ambas raices.
\end{itemize}

\paragraph{Codigo.}
Ver \texttt{cpp.tex}, seccion \textit{Segment Tree persistente}.

\vspace{0.5em}

\subsubsection{Segment Tree iterativo}

\paragraph{Idea intuitiva.}
Representacion \textbf{sin recursion}: los nodos se guardan en un \texttt{vector} de tamano $2n$. Los hijos del nodo $i$ son $2i$ y $2i+1$. Las hojas ocupan las posiciones $n, n+1, \dots, 2n-1$. Update: subir el cambio desde la hoja hasta la raiz. Query: subir dos indices $l$ y $r$ hasta que se ``crucen''. La ventaja principal: localidad de memoria (los nodos estan en posiciones contiguas del vector), que mejora el cache miss rate y la constante practica. La desventaja: implementar lazy propagation es mas complejo.

\paragraph{Propiedades clave (invariantes).}
\begin{itemize}\setlength\itemsep{0pt}
	\item $O(\log n)$ por operacion, \textbf{mas rapido} en practica que el recursivo (mejor uso de cache, no hay overhead de llamadas).
	\item Build simple: copiar las hojas y combinar de abajo hacia arriba.
	\item Las hojas tienen indices en $[n, 2n)$; los nodos internos en $[1, n)$.
\end{itemize}

\paragraph{Codigo clave.}
El bucle de query ascendente:
\begin{verbatim}
long long query(int l, int r) {  // [l, r] inclusivo
    long long res = 0;
    for (l += n, r += n + 1; l < r; l >>= 1, r >>= 1) {
        if (l & 1) res += st[l++];
        if (r & 1) res += st[--r];
    }
    return res;
}
\end{verbatim}

\paragraph{Complejidad.}
\begin{itemize}\setlength\itemsep{0pt}
	\item $O(\log n)$ por update/query, $O(n)$ build.
	\item En la practica, $\sim 1.5x$ mas rapido que el recursivo para $n \ge 10^5$.
\end{itemize}

\paragraph{Donde tocar para modificar.}
\begin{itemize}\setlength\itemsep{0pt}
	\item \textbf{Suma -> max/min/xor}: cambiar la operacion en \texttt{set} y \texttt{query}.
	\item \textbf{Aniadir update en rango con lazy}: implementar version lazy por separado (no trivial en iterativo).
	\item \textbf{Tamano no potencia de 2}: si $n$ no es potencia de 2, aniadir padding.
	\item \textbf{Persistencia}: usar \texttt{vector<array<int, 2>>} y duplicar nodos.
\end{itemize}

\paragraph{Trampas y errores frecuentes.}
\begin{itemize}\setlength\itemsep{0pt}
	\item En el query, \texttt{r += n+1} (no \texttt{r += n}) porque \texttt{query(l, r)} es \textbf{inclusivo} en ambos extremos.
	\item El indice \texttt{l & 1} detecta si \texttt{l} es hijo derecho; aniadir \texttt{res += st[l++]} en ese caso.
	\item El \texttt{--r} en la condicion derecha es para incluir el limite.
	\item Si las operaciones no son idempotentes, cuidado con el orden.
\end{itemize}

\paragraph{Codigo.}
Ver \texttt{cpp.tex}, seccion \textit{Segment Tree iterativo}.

\subsection{Otras}

\subsubsection{MO}

\paragraph{Idea intuitiva.}
Algoritmo de Mo ``puro'' (sin updates): queries offline ordenadas por $(L/\sqrt{n}, R)$ (con $R$ alternando entre ascendente y descendente segun el bloque). Mover $L$ o $R$ entre queries consecutivas es $O(1)$; el total es $O((N + Q) \sqrt{n})$. La intuicion: al reordenar las queries, los punteros $L$ y $R$ se mueven ``lo menos posible''; cada elemento del array se visita $O(\sqrt{n} \cdot Q/N)$ veces.

\paragraph{Propiedades clave (invariantes).}
\begin{itemize}\setlength\itemsep{0pt}
	\item \texttt{Q} guarda \texttt{(l, r, idx)}.
	\item El sort prioriza bloque de $L$; dentro del bloque, alterna $R$ ascendente/descendente para mejorar la locality.
	\item Tres punteros \texttt{l}, \texttt{r} que se mueven paso a paso.
	\item Tamano de bloque optimo: $\sim N / \sqrt{Q}$.
\end{itemize}

\paragraph{Codigo clave.}
El sort clasico de Mo con alternancia:
\begin{verbatim}
int block = (int)sqrt(n);
sort(queries.begin(), queries.end(), [&](const Q& a, const Q& b) {
    int blockA = a.l / block, blockB = b.l / block;
    if (blockA != blockB) return blockA < blockB;
    return (blockA & 1) ? (a.r > b.r) : (a.r < b.r);
});
\end{verbatim}

\paragraph{Complejidad.}
\begin{itemize}\setlength\itemsep{0pt}
	\item $O((N + Q) \sqrt{n})$.
	\item Con Hilbert order, $\sim 30\%$ mas rapido en practica.
\end{itemize}

\paragraph{Donde tocar para modificar.}
\begin{itemize}\setlength\itemsep{0pt}
	\item \textbf{Hilbert order}: alternativa al sort clasico; da mejor locality y es mas rapido en practica.
	\item \textbf{Aniadir la operacion concreta}: el snippet tiene \texttt{// add} y \texttt{// delete} como comentarios; reemplazarlos con la logica del problema (contar distintos, suma de frecuencias, etc.).
	\item \textbf{Bloque optimo}: calcular \texttt{int block = max(1, (int)(n / sqrt(q)));}.
\end{itemize}

\paragraph{Trampas y errores frecuentes.}
\begin{itemize}\setlength\itemsep{0pt}
	\item Los indices en el sort son \textbf{0-based} (\texttt{q[i].l = l - 1}), consistente con el resto del codigo.
	\item \texttt{add} y \texttt{remove} deben ser \textbf{inversos exactos}; sino, las queries dan respuestas incorrectas.
	\item Para queries cerradas $[l, r]$, Mo suele usar $[l, r+1)$ (half-open); aniadir o restar 1 consistentemente.
\end{itemize}

\paragraph{Codigo.}
Ver \texttt{cpp.tex}, seccion \textit{MO}.

\vspace{0.5em}

\subsubsection{DSU con rollback}

\paragraph{Idea intuitiva.}
Variante del DSU donde cada \texttt{unite} \textbf{se puede deshacer} en orden LIFO. La idea: en vez de path compression (que es irreversible), solo se usa union by size. Cada \texttt{unite} guarda un par \texttt{(parent\_changed, old\_size)} en una pila. \texttt{rollback()} desapila hasta el ultimo ``snapshot'' (marcador) restaurando los valores. La razon: path compression modifica la estructura de multiples nodos, no solo uno, asi que es dificil de deshacer; union by size solo modifica dos nodos (el padre del menor y su tamano), asi que es trivial deshacer.

\paragraph{Propiedades clave (invariantes).}
\begin{itemize}\setlength\itemsep{0pt}
	\item \textbf{Sin path compression}: cada \texttt{find} es $O(\log n)$ (peor caso del union by size) en el peor caso, pero amortizado $O(\log n)$.
	\item \textbf{Snapshots}: \texttt{snapshot()} marca un punto; \texttt{rollback()} deshace hasta ahi.
	\item Cada \texttt{unite} exitoso apila \textbf{dos} pares (uno por cada cambio).
	\item La pila de historial crece monotona; \texttt{rollback} la reduce.
\end{itemize}

\paragraph{Codigo clave.}
El unite con guardado de historial:
\begin{verbatim}
void unite(int a, int b) {
    a = find(a); b = find(b);
    if (a == b) { history.push({-1, -1}); return; }
    if (size[a] < size[b]) swap(a, b);
    history.push({b, size[a]});
    parent[b] = a;
    size[a] += size[b];
    components--;
}
void rollback() {
    auto [b, sz] = history.top(); history.pop();
    if (b == -1) return;
    size[parent[b]] = sz;
    parent[b] = b;
    components++;
}
\end{verbatim}

\paragraph{Complejidad.}
\begin{itemize}\setlength\itemsep{0pt}
	\item $O(\log n)$ por \texttt{unite}/\texttt{find}, $O(1)$ por rollback de una operacion.
	\item Sin path compression: peor caso $O(\log n)$ por find.
\end{itemize}

\paragraph{Donde tocar para modificar.}
\begin{itemize}\setlength\itemsep{0pt}
	\item \textbf{Rollback parcial} (no todo): implementar snapshots parciales guardando el tamano del historial.
	\item \textbf{Mas informacion por conjunto}: aniadir mas campos al struct y guardarlos en \texttt{history}.
	\item \textbf{Persistente}: aniadir versionado completo (clonar el struct en cada cambio).
	\item \textbf{DSU de Arpa}: variante con $O(\log n)$ amortizado y rollback en $O(1)$.
\end{itemize}

\paragraph{Trampas y errores frecuentes.}
\begin{itemize}\setlength\itemsep{0pt}
	\item El snippet guarda dos pares por \texttt{unite}; al rollback \textbf{desapilar en orden correcto} (no invertir).
	\item Olvidar el \texttt{components--} en \texttt{unite} o el \texttt{components++} en rollback deja el contador mal.
	\item El sentinela \texttt{\{-1, -1\}} indica un unite trivial; importante para mantener consistencia.
	\item \texttt{find} puede ser $O(\log n)$ peor caso sin path compression; verificar que es OK para el problema.
\end{itemize}

\paragraph{Codigo.}
Ver \texttt{cpp.tex}, seccion \textit{DSU con rollback}.

\vspace{0.5em}

\subsubsection{Segment Tree 2D}

\paragraph{Idea intuitiva.}
Extender el segment tree a 2 dimensiones: la raiz cubre toda la matriz $[0, n) \times [0, m)$. Cada nodo es un \textbf{segment tree 1D} sobre la segunda dimension. Para update en $(x, y)$, se actualizan $O(\log n)$ nodos externos (uno por nivel vertical) y dentro de cada uno, $O(\log m)$ nodos internos. Query sobre un rectangulo: $O(\log n \cdot \log m)$. La estructura matematica: el arbol de primer nivel divide el espacio vertical; en cada nivel, los nodos internos gestionan el espacio horizontal. Asi, se reduce la consulta rectangular a $O(\log n \cdot \log m)$ consultas 1D parciales.

\paragraph{Propiedades clave (invariantes).}
\begin{itemize}\setlength\itemsep{0pt}
	\item Memoria: $O(n \cdot m)$ en el peor caso (matriz densa).
	\item Para matrices grandes con valores \textbf{cero} y unos pocos updates, se puede comprimir.
	\item Cada nodo externo tiene $4m$ nodos internos.
\end{itemize}

\paragraph{Codigo clave.}
Update en $(x, y)$ con recursion en cascada:
\begin{verbatim}
void updateX(int nodeX, int lx, int rx, int x, int y, int v) {
    if (lx != rx) {
        int mx = (lx + rx) / 2;
        if (x <= mx) updateX(nodeX*2, lx, mx, x, y, v);
        else         updateX(nodeX*2+1, mx+1, rx, x, y, v);
    }
    updateY(nodeX, lx, rx, y, v);
}
\end{verbatim}

\paragraph{Complejidad.}
\begin{itemize}\setlength\itemsep{0pt}
	\item Update y query $O(\log n \cdot \log m)$, memoria $O(4n \cdot 4m)$.
	\item Para $n = m = 1000$: $16 \cdot 10^6$ nodos ($\sim 64$MB con \texttt{int}).
\end{itemize}

\paragraph{Donde tocar para modificar.}
\begin{itemize}\setlength\itemsep{0pt}
	\item \textbf{Build desde una matriz}: aniadir constructor que tome \texttt{vector<vector<int>>}.
	\item \textbf{Lazy propagation 2D}: aniadir un campo \texttt{lazy} por nodo.
	\item \textbf{Suma + maximo simultaneos}: aniadir mas campos.
	\item \textbf{Sparse 2D}: si la mayoria son ceros, usar map en vez de vector.
\end{itemize}

\paragraph{Trampas y errores frecuentes.}
\begin{itemize}\setlength\itemsep{0pt}
	\item La recursion puede ser profunda; aniadir iterativo si $n, m > 10^5$.
	\item Memoria $O(4n \cdot 4m) = 16nm$; si $n = m = 10^3$, son $16M$ nodos; con $n = m = 10^4$, son $1.6 \cdot 10^9$, no factible.
	\item Olvidar propagar en cascada da updates incorrectos.
	\item Si la mayoria de queries son rectangulares grandes, considerar KD-tree u otra estructura.
\end{itemize}

\paragraph{Codigo.}
Ver \texttt{cpp.tex}, seccion \textit{Segment Tree 2D}.

\vspace{0.5em}

\subsubsection{Segment Tree persistente}

\paragraph{Idea intuitiva.}
Cada \texttt{update} crea una \textbf{nueva raiz} sin modificar los nodos no afectados. Asi puedes tener \textbf{multiples versiones} del arreglo y consultar cada una en $O(\log n)$. Memoria: $O(N \log N)$ nodos en total para $N$ updates. La tecnica se llama ``path copying'': solo el camino de la raiz a la hoja modificada se copia; el resto del arbol se comparte entre versiones. Asi, $N$ updates producen $\le N \log n$ nodos en total.

\paragraph{Propiedades clave (invariantes).}
\begin{itemize}\setlength\itemsep{0pt}
	\item Cada nodo es \textbf{inmutable} despues de creado; las versiones comparten estructura.
	\item \texttt{pUpdate} retorna la nueva raiz; la anterior sigue accesible.
	\item \texttt{pQuery} toma una raiz especifica (no la global).
	\item Las versiones viven independientemente; borrar una no afecta a otras.
\end{itemize}

\paragraph{Codigo clave.}
El corazon del persistent update:
\begin{verbatim}
PNode* pUpdate(PNode* v, int l, int r, int pos, int val) {
    PNode* u = newNode(v);  // copia con misma info
    if (l == r) { u->sum = val; return u; }
    int m = (l + r) / 2;
    if (pos <= m) u->left = pUpdate(v->left, l, m, pos, val);
    else          u->right = pUpdate(v->right, m+1, r, pos, val);
    u->sum = u->left->sum + u->right->sum;
    return u;
}
// versions.push_back(pUpdate(versions.back(), 0, n-1, pos, val));
\end{verbatim}

\paragraph{Complejidad.}
\begin{itemize}\setlength\itemsep{0pt}
	\item Update $O(\log n)$ nodos nuevos, query $O(\log n)$.
	\item Memoria $O(N \log n)$ para $N$ updates.
\end{itemize}

\paragraph{Donde tocar para modificar.}
\begin{itemize}\setlength\itemsep{0pt}
	\item \textbf{Persistencia completa (cada update crea un nodo)}: aniadir \texttt{version[v]} para acceder a cada version.
	\item \textbf{Cambiar suma por otro operacion}: ajustar \texttt{pUpdate} y \texttt{pQuery}.
	\item \textbf{Memoria grande}: en contests, el allocator por defecto puede ser lento; usar \texttt{vector<PNode>} con \texttt{reserve(2*N*logN)} y un \texttt{int nextNode = 0} para asignar en $O(1)$.
	\item \textbf{Lazy persistence}: aniadir campos \texttt{lazy} en cada nodo; los hijos ``ven'' el lazy a traves de la copia.
\end{itemize}

\paragraph{Trampas y errores frecuentes.}
\begin{itemize}\setlength\itemsep{0pt}
	\item Olvidar pasar la \textbf{vieja raiz} al query, usando la nueva por error.
	\item Si guardas multiples versiones, asegurate de pasar la raiz correcta a cada query.
	\item Memory leak en produccion (no en contests) si no liberas los nodos.
	\item Para queries ``diferencia entre version $a$ y version $b

\vspace{0.5em}

\subsubsection{Segment Tree iterativo}

\paragraph{Idea intuitiva.}
Representacion \textbf{sin recursion}: los nodos se guardan en un \texttt{vector} de tamano $2n$. Los hijos del nodo $i$ son $2i$ y $2i+1$. Las hojas ocupan las posiciones $n, n+1, \dots, 2n-1$. Update: subir el cambio desde la hoja hasta la raiz. Query: subir dos indices $l$ y $r$ hasta que se ``crucen''. La ventaja principal: localidad de memoria (los nodos estan en posiciones contiguas del vector), que mejora el cache miss rate y la constante practica. La desventaja: implementar lazy propagation es mas complejo.

\paragraph{Propiedades clave (invariantes).}
\begin{itemize}\setlength\itemsep{0pt}
	\item $O(\log n)$ por operacion, \textbf{mas rapido} en practica que el recursivo (mejor uso de cache, no hay overhead de llamadas).
	\item Build simple: copiar las hojas y combinar de abajo hacia arriba.
	\item Las hojas tienen indices en $[n, 2n)$; los nodos internos en $[1, n)$.
\end{itemize}

\paragraph{Codigo clave.}
El bucle de query ascendente:
\begin{verbatim}
long long query(int l, int r) {  // [l, r] inclusivo
    long long res = 0;
    for (l += n, r += n + 1; l < r; l >>= 1, r >>= 1) {
        if (l & 1) res += st[l++];
        if (r & 1) res += st[--r];
    }
    return res;
}
\end{verbatim}

\paragraph{Complejidad.}
\begin{itemize}\setlength\itemsep{0pt}
	\item $O(\log n)$ por update/query, $O(n)$ build.
	\item En la practica, $\sim 1.5x$ mas rapido que el recursivo para $n \ge 10^5$.
\end{itemize}

\paragraph{Donde tocar para modificar.}
\begin{itemize}\setlength\itemsep{0pt}
	\item \textbf{Suma -> max/min/xor}: cambiar la operacion en \texttt{set} y \texttt{query}.
	\item \textbf{Aniadir update en rango con lazy}: implementar version lazy por separado (no trivial en iterativo).
	\item \textbf{Tamano no potencia de 2}: si $n$ no es potencia de 2, aniadir padding.
	\item \textbf{Persistencia}: usar \texttt{vector<array<int, 2>>} y duplicar nodos.
\end{itemize}

\paragraph{Trampas y errores frecuentes.}
\begin{itemize}\setlength\itemsep{0pt}
	\item En el query, \texttt{r += n+1} (no \texttt{r += n}) porque \texttt{query(l, r)} es \textbf{inclusivo} en ambos extremos.
	\item El indice \texttt{l & 1} detecta si \texttt{l} es hijo derecho; aniadir \texttt{res += st[l++]} en ese caso.
	\item El \texttt{--r} en la condicion derecha es para incluir el limite.
	\item Si las operaciones no son idempotentes, cuidado con el orden.
\end{itemize}

\paragraph{Codigo.}
Ver \texttt{cpp.tex}, seccion \textit{Segment Tree iterativo}.

\vspace{0.5em}

\subsection{Fenwick / BIT}

\subsubsection{BIT point update + prefix sum}

\paragraph{Idea intuitiva.}
El \textbf{Fenwick Tree} (BIT) almacena sums parciales de manera que \texttt{sum(i)} (suma de $[1, i]$) y \texttt{add(i, delta)} cuestan $O(\log n)$. La representacion es ingeniosa: el nodo $i$ cubre el rango $[i - \text{lowbit}(i) + 1, i]$. Asi, \texttt{sum(i)} itera \texttt{i -= lowbit(i)}, sumando los nodos que cubren un prefijo creciente, y \texttt{add(i, delta)} itera \texttt{i += lowbit(i)}, propagando el cambio a todos los nodos cuyo rango contiene a $i$. La idea: cada indice $i$ se descompone en una suma de bloques disjuntos de tamano potencias de 2 (su ``lowbit representation'').

\paragraph{Propiedades clave (invariantes).}
\begin{itemize}\setlength\itemsep{0pt}
	\item Indice \textbf{1-based} (no usar $i=0$).
	\item \texttt{lowbit(i) = i \& (-i)}.
	\item Cada nodo cubre exactamente un rango de tamano \texttt{lowbit(i)}.
	\item La operacion debe ser \textbf{invertible} (suma, xor, no min/max).
\end{itemize}

\paragraph{Codigo clave.}
Las dos operaciones basicas:
\begin{verbatim}
void add(int i, long long d) {
    for (; i <= n; i += i & -i) bit[i] += d;
}
long long sum(int i) {
    long long s = 0;
    for (; i > 0; i -= i & -i) s += bit[i];
    return s;
}
\end{verbatim}

\paragraph{Complejidad.}
\begin{itemize}\setlength\itemsep{0pt}
	\item $O(\log n)$ por operacion. Memoria $O(n)$.
	\item Constante muy pequena ($\sim 1$): en la practica, $\sim 3x$ mas rapido que segment tree.
\end{itemize}

\paragraph{Donde tocar para modificar.}
\begin{itemize}\setlength\itemsep{0pt}
	\item \textbf{Suma -> max/min/xor}: cambiar el \texttt{+=} por la operacion (con identidad apropriada). \textbf{No} funciona para \texttt{min} de forma natural (no es invertible).
	\item \textbf{Tamano dinamico}: aniadir un metodo \texttt{resize(newSize)} que preserve valores.
	\item \textbf{Tamano grande}: con $n = 10^7$, mejor comprimir coordenadas.
	\item \textbf{Operador no conmutativo}: usar segment tree en su lugar.
\end{itemize}

\paragraph{Trampas y errores frecuentes.}
\begin{itemize}\setlength\itemsep{0pt}
	\item Indices 0 vs 1: el BIT es \textbf{1-based}. Si el problema da indices 0-based, aniadir \texttt{+1} en cada llamada.
	\item El metodo \texttt{sum} itera \texttt{i -= i \& -i}; aniadir o quitar 1 rompe el invariante.
	\item Si el tipo es \texttt{int}, posibles overflows; usar \texttt{long long}.
	\item Para rango $[l, r]$, el query es \texttt{sum(r) - sum(l-1)} (cuidado con $l = 1$).
\end{itemize}

\paragraph{Codigo.}
Ver \texttt{cpp.tex}, seccion \textit{BIT point update + prefix sum}.

\vspace{0.5em}

\subsubsection{BIT range add + point query}

\paragraph{Idea intuitiva.}
Truco de \textbf{diferencias}: para sumar $d$ a $[l, r]$, aniadir $d$ en $l$ y $-d$ en $r+1$. Asi, el valor en $i$ es la suma de los prefijos hasta $i$, lo que se calcula con un BIT normal. La idea algebraica: si mantenemos un arreglo de diferencias $\Delta$ donde $a_i = \sum_{j \le i} \Delta_j$, entonces aniadir $d$ a $a[l..r]$ equivale a $\Delta_l \mathrel{+}= d$ y $\Delta_{r+1} \mathrel{-}= d$, y reconstruir $a$ es un prefix sum.

\paragraph{Propiedades clave (invariantes).}
\begin{itemize}\setlength\itemsep{0pt}
	\item Internamente es un BIT 5.1 con la convencion de diferencias.
	\item \texttt{rangeAdd(l, r, delta)} cuesta $O(\log n)$ (dos \texttt{add}).
	\item \texttt{pointQuery(i)} cuesta $O(\log n)$.
	\item Si $r+1 > n$, no se aniade el $-d$ (fuera del rango).
\end{itemize}

\paragraph{Codigo clave.}
El truco de diferencias:
\begin{verbatim}
void rangeAdd(int l, int r, long long d) {
    add(l, d);
    if (r + 1 <= n) add(r + 1, -d);
}
long long pointQuery(int i) { return sum(i); }
\end{verbatim}

\paragraph{Complejidad.}
\begin{itemize}\setlength\itemsep{0pt}
	\item $O(\log n)$ por operacion.
	\item Memoria $O(n)$.
\end{itemize}

\paragraph{Donde tocar para modificar.}
\begin{itemize}\setlength\itemsep{0pt}
	\item \textbf{Persistencia}: aniadir versionado (cada update crea una nueva ``copia'' del estado; en BIT no es trivial, mejor usar persistent segtree).
	\item \textbf{Coordenada compression}: si $n$ es grande pero los updates son pocos, comprimir.
	\item \textbf{2D range add + point query}: aniadir una dimension con la misma tecnica.
\end{itemize}

\paragraph{Trampas y errores frecuentes.}
\begin{itemize}\setlength\itemsep{0pt}
	\item Si $r+1 > n$, \textbf{no} aniadir \texttt{-delta} (no hay efecto fuera del rango); el snippet ya lo chequea.
	\item Para reconstruir el array original tras $k$ range adds, hacer $k$ point queries (no hay shortcut).
	\item Confundir ``suma de diferencias'' con ``suma acumulada''; son cosas distintas.
\end{itemize}

\paragraph{Codigo.}
Ver \texttt{cpp.tex}, seccion \textit{BIT range add + point query}.

\vspace{0.5em}

\subsubsection{BIT range add + range sum}

\paragraph{Idea intuitiva.}
Para queries $\sum_{i=1}^{x} a_i$ despues de varios range adds, el truco de diferencias no basta (necesitas la suma acumulada de las diferencias). Solucion: dos BITs, $t_1$ y $t_2$. Tras aniadir $d$ a $[l, r]$: $t_1$ recibe $+d$ en $l$, $-d$ en $r+1$; $t_2$ recibe $+d(l-1)$ en $l$, $-d \cdot r$ en $r+1$. Entonces $\text{prefixSum}(x) = t_1.\text{sum}(x) \cdot x - t_2.\text{sum}(x)$. La razon matematica: el valor en $i$ es $\sum_{j \le i} \Delta_j$, asi que la suma hasta $i$ es $\sum_{k \le i} (i - k + 1) \Delta_k$; expandir da la formula con dos BITs.

\paragraph{Propiedades clave (invariantes).}
\begin{itemize}\setlength\itemsep{0pt}
	\item Verificacion: para $x \ge r$, $\text{prefixSum}(x)$ acumula correctamente todas las contribuciones.
	\item La formula sale de expandir la suma de diferencias.
	\item La identidad clave: $\sum_{i=1}^{x} i = x(x+1)/2$, asi que la suma de diferencias ponderadas requiere BITs adicionales.
\end{itemize}

\paragraph{Codigo clave.}
Range add con range sum (dos BITs):
\begin{verbatim}
void rangeAdd(int l, int r, long long d) {
    t1.add(l, d); t1.add(r + 1, -d);
    t2.add(l, d * (l - 1)); t2.add(r + 1, -d * r);
}
long long prefixSum(int x) {
    return t1.sum(x) * x - t2.sum(x);
}
\end{verbatim}

\paragraph{Complejidad.}
\begin{itemize}\setlength\itemsep{0pt}
	\item $O(\log n)$ por operacion (2 \texttt{add} o 2 \texttt{sum} por cada lado).
	\item Memoria $O(n)$.
\end{itemize}

\paragraph{Donde tocar para modificar.}
\begin{itemize}\setlength\itemsep{0pt}
	\item \textbf{2D range add + range sum}: aniadir otra dimension con la misma tecnica (queda 4 BITs).
	\item \textbf{Cambiar mod}: aniadir modulo en cada operacion.
	\item \textbf{Output del array}: aniadir un metodo que reconstruya $a_i$ point query por cada $i$.
\end{itemize}

\paragraph{Trampas y errores frecuentes.}
\begin{itemize}\setlength\itemsep{0pt}
	\item Overflow en \texttt{delta * (l - 1)} si los valores son grandes; usar \texttt{long long} y \texttt{\% MOD} al final.
	\item Confundir la convencion $0$-based vs $1$-based del BIT; las dos afectan las formulas.
	\item En 2D, el codigo es $\sim 4x$ mas largo y propenso a errores; verificar con casos pequenos.
\end{itemize}

\paragraph{Codigo.}
Ver \texttt{cpp.tex}, seccion \textit{BIT range add + range sum}.

\subsection{RMQ y sparse table}

\subsubsection{Sparse Table}

\paragraph{Idea intuitiva.}
Precomputa respuestas a queries de tamano $2^k$ para $k = 0, 1, \dots, \log n$. Asi, \texttt{st[k][i]} guarda la respuesta sobre $[i, i + 2^k)$. Para query $[l, r]$, se usan dos bloques de tamano $2^{\lfloor \log_2(r-l+1) \rfloor}$: $[l, l + 2^k)$ y $[r - 2^k + 1, r)$, y se combinan. Esto da $O(1)$ por query \textbf{solo si la operacion es idempotente} (min, max, gcd; \textbf{no} suma). La razon: con idempotencia, solapar los dos bloques no afecta el resultado.

\paragraph{Propiedades clave (invariantes).}
\begin{itemize}\setlength\itemsep{0pt}
	\item \texttt{st[0][i] = a[i]}.
	\item \texttt{st[k][i] = op(st[k-1][i], st[k-1][i + 2\^{}(k-1)])}.
	\item \texttt{LOG = floor(log2(n)) + 1}.
	\item La operacion debe ser \textbf{idempotente} ($f(x, x) = x$) y asociativa.
\end{itemize}

\paragraph{Codigo clave.}
Query $O(1)$ en sparse table:
\begin{verbatim}
long long query(int l, int r) {
    int k = __builtin_clz(1) - __builtin_clz(r - l + 1);  // log2
    return min(st[k][l], st[k][r - (1 << k) + 1]);
}
\end{verbatim}

\paragraph{Complejidad.}
\begin{itemize}\setlength\itemsep{0pt}
	\item Build $O(n \log n)$, query $O(1)$.
	\item Memoria $O(n \log n)$.
\end{itemize}

\paragraph{Donde tocar para modificar.}
\begin{itemize}\setlength\itemsep{0pt}
	\item \textbf{Cambiar min por otra operacion idempotente}: cambiar las dos ocurrencias de \texttt{min(...)} por la operacion (max, gcd). Verificar identidad y asociatividad.
	\item \textbf{Suma en rango}: NO sirve con este metodo; usar prefix sums.
	\item \textbf{Sparse table 2D}: aniadir otra dimension (queda $O(n^2 \log^2 n)$ memoria).
	\item \textbf{LCA queries}: combinar con binary lifting en arboles.
\end{itemize}

\paragraph{Trampas y errores frecuentes.}
\begin{itemize}\setlength\itemsep{0pt}
	\item El query toma $[l, r]$ \textbf{inclusivo} (no $[l, r)$).
	\item El \texttt{LOG} debe estar bien calculado; con $n = 1$, \texttt{LOG = 1}.
	\item Usar sparse table para suma es un error; necesitas prefix sums o segment tree.
	\item Para queries $O(1)$ con operaciones no idempotentes, usar Disjoint Sparse Table.
\end{itemize}

\paragraph{Codigo.}
Ver \texttt{cpp.tex}, seccion \textit{Sparse Table}.

\vspace{0.5em}

\subsubsection{Disjoint Sparse Table}

\paragraph{Idea intuitiva.}
Variante que tambien da $O(1)$ por query pero \textbf{no requiere idempotencia} para algunas operaciones (sirve para suma, xor, etc. en linea con la mayoria de operaciones asociativas). La idea: para cada nivel $k$, los rangos ``se parten en mitades''. \texttt{st[k][i]} cubre un rango que \textbf{empieza} en $i$ y se extiende a izquierda o derecha hasta encontrar un ``punto de division''. El truco: para query $[l, r]$, el nivel optimo es el bit mas alto donde $l$ y $r$ difieren; los rangos $[l, ?]$ y $[?, r]$ cubren exactamente el query sin solape.

\paragraph{Propiedades clave (invariantes).}
\begin{itemize}\setlength\itemsep{0pt}
	\item Construccion diferente al Sparse Table clasico: en cada nivel, se extiende desde cada posicion hacia ambos lados hasta encontrar la mitad de nivel.
	\item \texttt{query(l, r)} usa \texttt{k = log2(l \^{} r)} y combina \texttt{st[k][l]} y \texttt{st[k][r]}.
	\item La operacion debe ser \textbf{asociativa} pero no necesariamente idempotente.
\end{itemize}

\paragraph{Codigo clave.}
Query con DST:
\begin{verbatim}
long long query(int l, int r) {
    if (l == r) return st[0][l];
    int k = 31 - __builtin_clz(l ^ r);  // bit mas alto que difiere
    return op(st[k][l], st[k][r]);
}
\end{verbatim}

\paragraph{Complejidad.}
\begin{itemize}\setlength\itemsep{0pt}
	\item Build $O(n \log n)$, query $O(1)$.
	\item Memoria $O(n \log n)$.
\end{itemize}

\paragraph{Donde tocar para modificar.}
\begin{itemize}\setlength\itemsep{0pt}
	\item \textbf{Operacion no asociativa}: no funciona; el Disjoint Sparse Table asume asociatividad.
	\item \textbf{Memory compression}: con $n \le 10^6$ alcanza; con mas, comprimir.
	\item \textbf{Operaciones no conmutativas}: aniadir el cuidado de que el orden es importante.
\end{itemize}

\paragraph{Trampas y errores frecuentes.}
\begin{itemize}\setlength\itemsep{0pt}
	\item La construccion es sutil; verificar con casos pequenos antes de usar.
	\item El bit ``mas alto que difiere'' requiere cuidado con $l = r$; aniadir caso base.
	\item Si la operacion no es asociativa, los dos bloques no se pueden combinar en cualquier orden.
\end{itemize}

\paragraph{Codigo.}
Ver \texttt{cpp.tex}, seccion \textit{Disjoint Sparse Table}.

\subsection{Trees y DP en arbol}

\subsubsection{Binary Lifting / LCA}

\paragraph{Idea intuitiva.}
Precomputa para cada nodo $v$ sus ancestros a distancias $2^k$ (los ``saltos''). Asi, para subir $d$ niveles, se descompone $d$ en binario y se aplican los saltos correspondientes. El \textbf{LCA} (Lowest Common Ancestor) se obtiene subiendo el nodo mas profundo hasta igualar profundidades y luego subiendo ambos hasta que sus padres coincidan. La intuicion: cualquier entero $\le n$ se representa en binario con $\le \log n$ bits, asi que subir $d$ niveles requiere $\le \log n$ saltos.

\paragraph{Propiedades clave (invariantes).}
\begin{itemize}\setlength\itemsep{0pt}
	\item \texttt{up[k][v]} = ancestro de $v$ a $2^k$ niveles.
	\item \texttt{LOG = floor(log2(n)) + 1}.
	\item La recursion de \texttt{dfs} llena \texttt{up[0][v]} (el padre directo) y por transitividad los demas.
	\item El LCA esta definido como el ancestro comun mas profundo de $u$ y $v$.
\end{itemize}

\paragraph{Codigo clave.}
LCA con binary lifting:
\begin{verbatim}
int lca(int u, int v) {
    if (depth[u] < depth[v]) swap(u, v);
    int diff = depth[u] - depth[v];
    for (int k = LOG - 1; k >= 0; --k)
        if (diff & (1 << k)) u = up[k][u];
    if (u == v) return u;
    for (int k = LOG - 1; k >= 0; --k)
        if (up[k][u] != up[k][v])
            u = up[k][u], v = up[k][v];
    return up[0][u];
}
\end{verbatim}

\paragraph{Complejidad.}
\begin{itemize}\setlength\itemsep{0pt}
	\item Build $O(n \log n)$, query LCA $O(\log n)$.
	\item Memoria $O(n \log n)$.
\end{itemize}

\paragraph{Donde tocar para modificar.}
\begin{itemize}\setlength\itemsep{0pt}
	\item \textbf{Distancia en el arbol}: combinar con \texttt{depth[]} para calcular \texttt{depth[u] + depth[v] - 2*depth[lca]}.
	\item \textbf{K-esimo ancestro}: subir $k$ niveles con el mismo truco.
	\item \textbf{Aniadir peso a las aristas}: guardar \texttt{dist[2\^{}k][v]} (distancia acumulada a $2^k$ saltos) ademas de \texttt{up}.
	\item \textbf{Sparse table sobre depth}: alternativa $O(1)$ a LCA (mas complejo).
\end{itemize}

\paragraph{Trampas y errores frecuentes.}
\begin{itemize}\setlength\itemsep{0pt}
	\item El arbol debe ser \textbf{rooted} (el snippet asume eso); si no, hacer un \texttt{dfs} desde cualquier nodo.
	\item Para grafos no conexos, correr \texttt{dfs} desde cada componente.
	\item En la segunda fase del LCA, subir ambos nodos mientras \texttt{up[k][u] != up[k][v]}; el LCA final es \texttt{up[0][u]}.
	\item El padre de la raiz es 0 (o -1); aniadir check si se sale del rango.
\end{itemize}

\paragraph{Codigo.}
Ver \texttt{cpp.tex}, seccion \textit{Binary Lifting / LCA}.

\vspace{0.5em}

\subsubsection{Heavy-Light Decomposition (HLD)}

\paragraph{Idea intuitiva.}
Descomponer un arbol en \textbf{cadenas pesadas} (paths donde cada nodo (salvo la raiz) tiene al hijo ``pesado'' como hijo de tamano maximo de su subarbol). Cada cadena se mapea a un rango contiguo en un arreglo base; asi, queries sobre \textbf{paths} se descomponen en $O(\log n)$ queries sobre rangos contiguos (cada cadena es un rango), que se resuelven con un BIT o segment tree. La intuicion: por cada nivel de recursion, el subarbol del hijo pesado es al menos la mitad, asi que un path cruza $\le 2 \log n$ cadenas.

\paragraph{Propiedades clave (invariantes).}
\begin{itemize}\setlength\itemsep{0pt}
	\item \texttt{heavy[v]} = hijo pesado de $v$ (el de mayor subarbol).
	\item \texttt{head[v]} = cabeza de la cadena que contiene a $v$.
	\item \texttt{pos[v]} = posicion en el arreglo base.
	\item Cada nodo pertenece a exactamente una cadena.
\end{itemize}

\paragraph{Codigo clave.}
El query sobre un path $u \to v$:
\begin{verbatim}
long long pathQuery(int u, int v) {
    long long res = 0;
    while (head[u] != head[v]) {
        if (depth[head[u]] < depth[head[v]]) swap(u, v);
        res += seg.query(pos[head[u]], pos[u]);
        u = parent[head[u]];
    }
    if (depth[u] > depth[v]) swap(u, v);
    res += seg.query(pos[u], pos[v]);
    return res;
}
\end{verbatim}

\paragraph{Complejidad.}
\begin{itemize}\setlength\itemsep{0pt}
	\item Build $O(n)$, query sobre path $O(\log^2 n)$ (con segtree) o $O(\log n)$ (con BIT si solo se hace suma).
	\item Memoria $O(n)$.
\end{itemize}

\paragraph{Donde tocar para modificar.}
\begin{itemize}\setlength\itemsep{0pt}
	\item \textbf{Usar con segment tree}: aniadir un segtree sobre el arreglo base; \texttt{pathSegments} da los rangos a actualizar/query.
	\item \textbf{Cambiar operacion}: ajustar el segtree (suma, max, etc.).
	\item \textbf{Subtree query}: \texttt{subtree[v]} es el rango \texttt{[pos[v], pos[v] + size[v] - 1]} en el arreglo base.
	\item \textbf{Updates en edges}: aniadir un offset para que el edge $(u, v)$ quede en el nodo mas profundo.
\end{itemize}

\paragraph{Trampas y errores frecuentes.}
\begin{itemize}\setlength\itemsep{0pt}
	\item \texttt{pathSegments} requiere un \texttt{segtree} o \texttt{BIT} externo para ser util; por si solo solo da los rangos.
	\item Para edges vs nodos: si la operacion es sobre edges, guardar el valor en el hijo mas profundo.
	\item Olvidar la recursion \texttt{while (head[u] != head[v])} da resultados incorrectos.
	\item En problemas con updates, asegurarse de que el segtree soporte lazy propagation.
\end{itemize}

\paragraph{Codigo.}
Ver \texttt{cpp.tex}, seccion \textit{Heavy-Light Decomposition (HLD)}.

\vspace{0.5em}

\subsubsection{Centroid Decomposition}

\paragraph{Idea intuitiva.}
Descomponer recursivamente el arbol en \textbf{centroides}. Un centroide de un arbol es un nodo tal que cada subtree (al removerlo) tiene tamano $\le n/2$. Existe siempre y se encuentra en $O(n)$. Esto da una estructura de \textbf{arbol de centroides} de profundidad $O(\log n)$, util para problemas ``para cada nodo, encontrar el mas cercano que cumple X''. La razon de la profundidad $O(\log n)$: cada recursion reduce el problema a subarboles de tamano $\le n/2$, asi que la recursion da una cadena geometrica decreciente.

\paragraph{Propiedades clave (invariantes).}
\begin{itemize}\setlength\itemsep{0pt}
	\item En cada nivel, los subarboles se procesan recursivamente.
	\item \texttt{removed[]} marca nodos ya usados como centroide.
	\item \texttt{centroidParent[]} da el padre en el arbol de centroides.
	\item La altura del arbol de centroides es $O(\log n)$.
\end{itemize}

\paragraph{Codigo clave.}
Encontrar el centroide de un subarbol:
\begin{verbatim}
int findCentroid(int v, int p, int sz) {
    for (int u : adj[v])
        if (u != p && !removed[u] && subsize[u] > sz / 2)
            return findCentroid(u, v, sz);
    return v;
}
\end{verbatim}

\paragraph{Complejidad.}
\begin{itemize}\setlength\itemsep{0pt}
	\item Build $O(n \log n)$, queries dependen del problema.
	\item Cada nivel de recursion es $O(n)$ para calcular subsize.
\end{itemize}

\paragraph{Donde tocar para modificar.}
\begin{itemize}\setlength\itemsep{0pt}
	\item \textbf{Query tipica}: en cada nivel, ``expandir'' desde el centroide contando caminos validos y usar una estructura auxiliar (set, multiset).
	\item \textbf{Distancia maxima a un nodo ``malo''}: mantener un multiset de distancias en cada centroide.
	\item \textbf{Cambiar a otro tipo de descomposicion}: si el problema se adapta mejor a small-to-large o DSU on tree, usar esos.
	\item \textbf{Dynamic updates}: combinacion con link-cut trees.
\end{itemize}

\paragraph{Trampas y errores frecuentes.}
\begin{itemize}\setlength\itemsep{0pt}
	\item El snippet actual \textbf{solo construye} la descomposicion; las queries se aniaden segun el problema.
	\item Olvidar resetear \texttt{removed[]} causa recursion infinita.
	\item El calculo de subsize debe ignorar nodos ya ``removed''.
	\item El centroide es \textbf{unico} salvo empates en tamano; cualquiera sirve.
\end{itemize}

\paragraph{Codigo.}
Ver \texttt{cpp.tex}, seccion \textit{Centroid Decomposition}.

\vspace{0.5em}

\subsubsection{Euler Tour + range queries}

\paragraph{Idea intuitiva.}
Un \textbf{Euler Tour} (first-time visit) numera los nodos en el orden en que se visitan. El subarbol de $v$ corresponde a un \textbf{rango contiguo} $[tin[v], tout[v]]$ en este orden. Asi, queries sobre subarboles se resuelven con un BIT o segment tree sobre el arreglo de tiempos. La razon del rango contiguo: cuando entramos a un subarbol, no salimos hasta haber recorrido todos sus descendientes; entonces las posiciones consecutivas corresponden exactamente al subarbol.

\paragraph{Propiedades clave (invariantes).}
\begin{itemize}\setlength\itemsep{0pt}
	\item \texttt{tin[v]} = tiempo de primera visita a $v$.
	\item \texttt{tout[v]} = tiempo al \textbf{terminar} de procesar el subarbol de $v$.
	\item El subarbol de $v$ ocupa las posiciones $[tin[v], tout[v]]$.
	\item Para verificar si $u$ es ancestro de $v$: \texttt{tin[u] <= tin[v] \&\& tout[v] <= tout[u]}.
\end{itemize}

\paragraph{Codigo clave.}
El DFS del Euler Tour:
\begin{verbatim}
void dfs(int v, int p) {
    tin[v] = ++timer;
    for (int u : adj[v])
        if (u != p) dfs(u, v);
    tout[v] = timer;
}
\end{verbatim}

\paragraph{Complejidad.}
\begin{itemize}\setlength\itemsep{0pt}
	\item Build $O(n)$, query sobre subarbol $O(\log n)$ con BIT/segtree.
	\item Memoria $O(n)$.
\end{itemize}

\paragraph{Donde tocar para modificar.}
\begin{itemize}\setlength\itemsep{0pt}
	\item \textbf{Aniadir valor por nodo}: guardar el valor en \texttt{tin[v]} y operar sobre el rango.
	\item \textbf{Path query}: el Euler Tour \textbf{no} sirve para paths; usar HLD o LCA.
	\item \textbf{Multiple queries offline}: procesar en orden de tiempo y mantener una estructura por profundidad.
	\item \textbf{Verificar ancestro}: $O(1)$ con la condicion \texttt{tin}.
\end{itemize}

\paragraph{Trampas y errores frecuentes.}
\begin{itemize}\setlength\itemsep{0pt}
	\item El snippet solo construye el tour; las queries concretas se aniaden.
	\item \texttt{tout[v]} esta en el momento de salir, asi que \texttt{tout[v] - tin[v] + 1} = tamano del subarbol.
	\item Si el arbol no es conexo, aniadir un loop externo sobre cada componente.
	\item Para queries en tiempo (no subtree), usar el orden del tour y offline processing.
\end{itemize}

\paragraph{Codigo.}
Ver \texttt{cpp.tex}, seccion \textit{Euler Tour + range queries}.

\subsection{BST balanceado}

\subsubsection{Treap implicito}

\paragraph{Idea intuitiva.}
Un \textbf{treap} es un arbol binario de busqueda donde cada nodo tiene una \textbf{clave} (el valor de la secuencia) y una \textbf{prioridad} aleatoria. La invariante es: BST por clave, heap por prioridad. Esto da un arbol \textbf{esperado balanceado} sin necesidad de rebalanceo explicito. Un \textbf{treap implicito} no almacena una clave; la posicion de un nodo en la secuencia esta dada por el \textbf{tamano del subarbol izquierdo}, permitiendo split/merge por posicion. La estructura es esencialmente un RBS (randomized BST): la altura esperada es $O(\log n)$ por la propiedad de heap + BST.

\paragraph{Propiedades clave (invariantes).}
\begin{itemize}\setlength\itemsep{0pt}
	\item \texttt{sz(node)} = tamano del subarbol (1 + izquierdo + derecho).
	\item \texttt{split(t, k, l, r)}: parte \texttt{t} en \texttt{l} (primeros $k$ nodos) y \texttt{r} (resto).
	\item \texttt{merge(l, r)}: une dos treaps (todos los de \texttt{l} antes que los de \texttt{r}).
	\item La prioridad aleatoria asegura equilibrio esperado.
\end{itemize}

\paragraph{Codigo clave.}
Split por posicion (la base del treap implicito):
\begin{verbatim}
void split(Node* t, int k, Node*& l, Node*& r) {
    if (!t) { l = r = nullptr; return; }
    push(t);
    if (sz(t->l) >= k) {
        split(t->l, k, l, t->l);
        r = update(t);
    } else {
        split(t->r, k - sz(t->l) - 1, t->r, r);
        l = update(t);
    }
}
\end{verbatim}

\paragraph{Complejidad.}
\begin{itemize}\setlength\itemsep{0pt}
	\item $O(\log n)$ \textbf{esperado} por operacion (split, merge, insert, erase).
	\item Peor caso exponencial (probabilidad astronomicamente pequena).
\end{itemize}

\paragraph{Donde tocar para modificar.}
\begin{itemize}\setlength\itemsep{0pt}
	\item \textbf{Treap con clave (no implicito)}: usar \texttt{key} en vez de posicion; cambiar \texttt{split} por \texttt{splitByKey}.
	\item \textbf{Lazy propagation}: aniadir campos \texttt{lazy} y aplicarlos en \texttt{push} (recursivo).
	\item \textbf{Persistencia}: aniadir versionado clonando nodos.
	\item \textbf{Mejor aleatoriedad}: en lugar de \texttt{rand()}, usar \texttt{rng()} (mejor estado).
	\item \textbf{Operacion de reversa}: aniadir \texttt{lazyReverse} y swap de hijos en push.
\end{itemize}

\paragraph{Trampas y errores frecuentes.}
\begin{itemize}\setlength\itemsep{0pt}
	\item En el snippet actual, \texttt{Node(int v)} usa \texttt{rand()} que en algunos compiladores da maxima prioridad 0; aniadir \texttt{rng()} para mejor distribucion.
	\item Olvidar \texttt{update(t)} tras modificar hijos da tamano incorrecto.
	\item Si lazy propagation no se aplica en push, las queries dan valores viejos.
	\item Memory leak con \texttt{new}; en contests OK, en produccion aniadir destructor.
\end{itemize}

\paragraph{Codigo.}
Ver \texttt{cpp.tex}, seccion \textit{Treap implicito}.

\vspace{0.5em}

\subsubsection{Mo con updates}

\paragraph{Idea intuitiva.}
Extension del algoritmo de Mo para queries que mezclan \textbf{rango} y \textbf{updates}. Las queries se ordenan por $(L/\text{blockSize}, R/\text{blockSize}, T)$ donde $T$ es el ``tiempo'' (numero de updates). Asi, moverse entre queries consecutivas cuesta $O(n^{2/3})$ amortizado. La intuicion: el sort por bloques reduce los movimientos totales, y la tercera coordenada $T$ se optimiza con alternancia similar a $R$.

\paragraph{Propiedades clave (invariantes).}
\begin{itemize}\setlength\itemsep{0pt}
	\item Tres punteros: $L$, $R$, $T$.
	\item \texttt{applyUpdate} y \texttt{revertUpdate} son \textbf{inversos} exactos.
	\item \texttt{addPos} y \texttt{removePos} mantienen el estado.
	\item El tamano de bloque optimo es $n^{2/3}$ (donde $n$ = tamano del array, $Q$ = numero de queries).
\end{itemize}

\paragraph{Codigo clave.}
El sort con tres coordenadas:
\begin{verbatim}
int blockSize = (int)pow(n, 2.0/3.0);
sort(queries.begin(), queries.end(), [&](const Q& a, const Q& b) {
    int blockA = a.l / blockSize, blockB = b.l / blockSize;
    if (blockA != blockB) return blockA < blockB;
    int blockAR = a.r / blockSize, blockBR = b.r / blockSize;
    if (blockAR != blockBR) return blockAR < blockBR;
    return a.t < b.t;
});
\end{verbatim}

\paragraph{Complejidad.}
\begin{itemize}\setlength\itemsep{0pt}
	\item $O((N + Q) \cdot n^{2/3})$ donde $n^{2/3}$ es el tamano de bloque optimo.
	\item En problemas reales, funciona bien hasta $N, Q \sim 10^5$.
\end{itemize}

\paragraph{Donde tocar para modificar.}
\begin{itemize}\setlength\itemsep{0pt}
	\item \textbf{El corazon es el ``add/remove''}: aniadir la logica del problema (ej. contar colores distintos, suma de frecuencias, etc.).
	\item \textbf{Cambiar blockSize}: si la mayoria de updates vs queries cambia, ajustar.
	\item \textbf{Rollback}: aniadir \texttt{save()} / \texttt{restore()} para volver a estados anteriores.
	\item \textbf{Hilbert order}: alternativa con mejor locality.
\end{itemize}

\paragraph{Trampas y errores frecuentes.}
\begin{itemize}\setlength\itemsep{0pt}
	\item El snippet actual es solo el \textbf{esqueleto}; \texttt{addPos}/\texttt{removePos}/\texttt{applyUpdate}/\texttt{revertUpdate} deben ser implementados por problema.
	\item Olvidar \texttt{revertUpdate} es el error mas comun.
	\item Si el numero de updates es muy grande, Mo con updates se vuelve lento; considerar otra estructura.
	\item Mantener consistencia entre \texttt{T} y la posicion de los updates.
\end{itemize}

\paragraph{Codigo.}
Ver \texttt{cpp.tex}, seccion \textit{Mo con updates}.

\vspace{0.5em}

\subsubsection{sqrt-decomposition generico}

\paragraph{Idea intuitiva.}
Dividir el arreglo en \textbf{bloques} de tamano $B \approx \sqrt{n}$. Mantener informacion precomputada por bloque (suma, max, etc.). Queries de rango se resuelven combinando bloques enteros ($O(\sqrt{n})$ en total) y elementos sueltos en los bloques parciales. Updates puntuales actualizan el valor y el resumen del bloque. La intuicion: las queries ``normales'' cruzan $O(n/B)$ bloques; tomar $B = \sqrt{n}$ balancea el costo de bloques enteros vs elementos sueltos.

\paragraph{Propiedades clave (invariantes).}
\begin{itemize}\setlength\itemsep{0pt}
	\item $B = \lfloor \sqrt{n} \rfloor + 1$.
	\item \texttt{bucket[i]} = resumen del bloque $i$-esimo (suma, max, etc.).
	\item \texttt{rangeSum(l, r)}: 3 casos segun $l$ y $r$ esten en el mismo bloque o no.
	\item Cada query cruza $\le 2 + (r - l)/B$ bloques.
\end{itemize}

\paragraph{Codigo clave.}
Estructura basica:
\begin{verbatim}
int B = sqrt(n) + 1;
vector<T> bucket(n / B + 1, IDENTITY);
T query(int l, int r) {
    T res = IDENTITY;
    int bl = l / B, br = r / B;
    if (bl == br) {
        for (int i = l; i <= r; ++i) res = combine(res, a[i]);
    } else {
        for (int i = l; i < (bl+1)*B; ++i) res = combine(res, a[i]);
        for (int i = bl+1; i < br; ++i) res = combine(res, bucket[i]);
        for (int i = br*B; i <= r; ++i) res = combine(res, a[i]);
    }
    return res;
}
\end{verbatim}

\paragraph{Complejidad.}
\begin{itemize}\setlength\itemsep{0pt}
	\item $O(\sqrt{n})$ por operacion.
	\item Memoria $O(n)$.
\end{itemize}

\paragraph{Donde tocar para modificar.}
\begin{itemize}\setlength\itemsep{0pt}
	\item \textbf{Cambiar suma por max/min}: ajustar \texttt{bucket} y la combinacion.
	\item \textbf{Aniadir update en rango}: si la operacion es asociativa, sumar al resumen del bloque + ajustar elementos sueltos.
	\item \textbf{Tamano de bloque optimo}: si las queries son muy largas, $B = n^{2/3}$ (sirve para Mo con updates).
	\item \textbf{Operacion no asociativa}: usar segment tree en su lugar.
\end{itemize}

\paragraph{Trampas y errores frecuentes.}
\begin{itemize}\setlength\itemsep{0pt}
	\item Si $n$ no es multiplo de $B$, el ultimo bloque es mas chico; aniadir check.
	\item La identidad de la operacion debe estar bien definida (0 para suma, $-\infty$ para max).
	\item Si la operacion no es asociativa, los bloques no se pueden combinar; cambiar a segment tree.
\end{itemize}

\paragraph{Codigo.}
Ver \texttt{cpp.tex}, seccion \textit{sqrt-decomposition generico}.

% ============================================================
\section{Strings}
% ============================================================

\subsection{Hashing}

\subsubsection{Rolling Hash (doble modulo)}

\paragraph{Idea intuitiva.}
Un \textbf{rolling hash} asigna a cada prefijo del string un valor numerico tal que el hash de un substring se puede obtener en $O(1)$ a partir de los prefijos. La forma clasica: $h[i] = h[i-1] \cdot A + s[i-1]$ (donde $A$ es una base y todo modulo $M$). Asi, $h[r] - h[l-1] \cdot A^{r-l+1}$ da el hash del substring $[l, r]$ (con $A^{r-l+1}$ precomputado). Usar \textbf{dos modulos} distintos reduce drasticamente la probabilidad de colision. La idea algebraica: el string se interpreta como un polinomio en $A$ con coeficientes los caracteres; el hash es su evaluacion modulo $M$.

\paragraph{Propiedades clave (invariantes).}
\begin{itemize}\setlength\itemsep{0pt}
	\item $h[0] = 0$, $h[i]$ = hash de los primeros $i$ caracteres.
	\item $p[i] = A^i \bmod M$ precomputado.
	\item Comparar dos substrings: comparar sus pares de hashes.
	\item \textbf{Colisiones}: posibles (probabilidad $1/M$ por par). Doble hash: $\approx 1/(M_1 M_2)$.
	\item Si $A$ y $M$ son coprimos, el hash es bijectivo modulo $M$ para strings de longitud fija.
\end{itemize}

\paragraph{Codigo clave.}
La formula del hash de substring:
\begin{verbatim}
pair<long long, long long> get(int l, int r) {
    // h[r] - h[l-1] * A^(r-l+1)
    long long x1 = (h1[r] - h1[l-1] * p1[r-l+1] % M1 + M1) % M1;
    long long x2 = (h2[r] - h2[l-1] * p2[r-l+1] % M2 + M2) % M2;
    return {x1, x2};
}
\end{verbatim}

\paragraph{Complejidad.}
\begin{itemize}\setlength\itemsep{0pt}
	\item Precompute $O(n)$, query hash $O(1)$.
	\item Memoria $O(n)$.
\end{itemize}

\paragraph{Donde tocar para modificar.}
\begin{itemize}\setlength\itemsep{0pt}
	\item \textbf{Caracteres con valor $> 26$}: si los strings incluyen simbolos, aniadir \texttt{if (c >= 'a' \&\& c <= 'z') v = c - 'a'; else v = 26 + ...}.
	\item \textbf{Hash de substrings con updates}: aniadir un Fenwick tree / segment tree sobre los hashes (no es trivial).
	\item \textbf{Cambiar mod}: usar dos primos grandes ($10^9+7$ y $10^9+9$) o un \texttt{u128} para un solo hash sin mod.
	\item \textbf{Base aleatoria}: aniadir random para evitar ataques (en contests no es necesario; en produccion si).
	\item \textbf{Hash polinomial}: $h = \sum s_i \cdot A^i$ en vez de $s_i \cdot A^{i-1}$; algunos problemas requieren este.
\end{itemize}

\paragraph{Trampas y errores frecuentes.}
\begin{itemize}\setlength\itemsep{0pt}
	\item Si los substrings a comparar estan vacios (\texttt{l > r}), su hash es 0; aniadir check.
	\item \textbf{Overflow} en \texttt{h[i-1] * A}: usar \texttt{long long} y aplicar \texttt{\% M} en cada operacion (no al final).
	\item La sustraccion \texttt{h[r] - h[l-1]*A^(r-l+1)} puede dar negativo; aniadir \texttt{+M} antes de \texttt{\% M}.
	\item Con doble hash, verificar \textbf{ambos} componentes; si solo uno coincide, es colision.
\end{itemize}

\paragraph{Codigo.}
Ver \texttt{cpp.tex}, seccion \textit{Rolling Hash (doble modulo)}.

\subsection{Patrones}

\subsubsection{KMP + LPS}

\paragraph{Idea intuitiva.}
El algoritmo \textbf{Knuth-Morris-Pratt} busca todas las ocurrencias de un patron $P$ en un texto $T$ en $O(|T| + |P|)$. La clave es la tabla \textbf{LPS} (Longest Proper Prefix which is Suffix) que, para cada posicion $i$ del patron, dice el largo del prefijo mas largo que tambien es sufijo \textbf{propio}. Durante el match, cuando hay mismatch, en vez de volver al inicio, saltamos al LPS anterior, evitando trabajo redundante. La demostracion de $O(|T| + |P|)$: cada iteracion del bucle externo avanza $i$ o decrece $j$ (con $j$ acotado por $|P|$), asi que en total hay $O(|T| + |P|)$ operaciones.

\paragraph{Propiedades clave (invariantes).}
\begin{itemize}\setlength\itemsep{0pt}
	\item \texttt{lps[i]} = largo del prefijo mas largo que es sufijo propio de \texttt{p[0..i]}.
	\item El match usa dos punteros \texttt{i} (texto) y \texttt{j} (patron); en mismatch, $j = lps[j-1]$.
	\item Cuando $j == |P|$, se encontro una ocurrencia; continuar con $j = lps[j-1]$ para buscar mas.
	\item El LPS se construye con una auto-aplicacion de KMP al patron.
\end{itemize}

\paragraph{Codigo clave.}
Construccion de la tabla LPS:
\begin{verbatim}
for (int i = 1; i < m; ++i) {
    int j = lps[i - 1];
    while (j > 0 && p[i] != p[j]) j = lps[j - 1];
    if (p[i] == p[j]) ++j;
    lps[i] = j;
}
\end{verbatim}

\paragraph{Complejidad.}
\begin{itemize}\setlength\itemsep{0pt}
	\item $O(|T| + |P|)$.
	\item En la practica, rapido y robusto.
\end{itemize}

\paragraph{Donde tocar para modificar.}
\begin{itemize}\setlength\itemsep{0pt}
	\item \textbf{Contar ocurrencias}: el snippet ya cuenta; para devolver posiciones, aniadir un \texttt{vector<int>}.
	\item \textbf{Varios patrones}: usar Aho-Corasick en su lugar.
	\item \textbf{Matcher exacto con wildcard}: aniadir logica para que \texttt{?} matchee cualquier caracter.
	\item \textbf{Tamano pequeno}: si el patron es $\le 5$ caracteres, Rabin-Karp es mas simple.
	\item \textbf{2D KMP}: aniadir una dimension; util en problemas de matrices.
\end{itemize}

\paragraph{Trampas y errores frecuentes.}
\begin{itemize}\setlength\itemsep{0pt}
	\item En el snippet, \texttt{i} y \texttt{j} son 0-based. La condicion \texttt{j == sz\_pattern} para contar funciona.
	\item La construccion de LPS usa el mismo patron como ``texto''.
	\item Si el patron tiene caracteres repetidos, el LPS puede ser no trivial; verificar con casos pequenos.
\end{itemize}

\paragraph{Codigo.}
Ver \texttt{cpp.tex}, seccion \textit{KMP + LPS}.

\vspace{0.5em}

\subsubsection{Z-algorithm}

\paragraph{Idea intuitiva.}
La \textbf{Z-function} $Z[i]$ de un string $S$ es el largo del prefijo mas largo de $S$ que \textbf{tambien} empieza en $S[i]$. Asi, $Z[0] = n$ por convencion. Se calcula en $O(n)$ manteniendo una ventana $[l, r]$ que es el match mas largo encontrado hasta ahora. La consulta $S[l..r] = S[0..r-l]$ da informacion sobre $Z[i]$ cuando $i \le r$. La intuicion: cualquier $i$ dentro de la ventana ``ve'' los caracteres desde una posicion conocida, asi que $Z[i]$ se puede copiar en lugar de recalcular.

\paragraph{Propiedades clave (invariantes).}
\begin{itemize}\setlength\itemsep{0pt}
	\item Construir $Z$ cuesta $O(n)$ con el truco de la ventana.
	\item Para matching de patron $P$ en texto $T$: concatenar $P + \# + T$ y aniadir todos los $i$ con $Z[i] = |P|$.
	\item La ventana $[l, r]$ es siempre la del match mas largo a la derecha.
\end{itemize}

\paragraph{Codigo clave.}
Construccion de Z-function en $O(n)$:
\begin{verbatim}
Z[0] = n;
int l = 0, r = 0;
for (int i = 1; i < n; ++i) {
    if (i <= r) Z[i] = min(r - i + 1, Z[i - l]);
    while (i + Z[i] < n && s[Z[i]] == s[i + Z[i]]) ++Z[i];
    if (i + Z[i] - 1 > r) { l = i; r = i + Z[i] - 1; }
}
\end{verbatim}

\paragraph{Complejidad.}
\begin{itemize}\setlength\itemsep{0pt}
	\item $O(n)$ para construir; matching $O(n + m)$.
	\item Cada caracter se compara a lo sumo 2 veces (extension + reduccion).
\end{itemize}

\paragraph{Donde tocar para modificar.}
\begin{itemize}\setlength\itemsep{0pt}
	\item \textbf{Detectar bordes (border)}: $Z[i] = n - i$ indica que $S[0..n-i-1]$ es borde.
	\item \textbf{Prefijo comun mas largo de dos strings}: aniadir $S_1 + \# + S_2$; $Z[|S_1|+1]$ es el prefijo comun.
	\item \textbf{Comparar prefijos y sufijos}: combinar $S + \# + \text{reverse}(S)$.
	\item \textbf{Matching aproximado}: aniadir tolerancia a mismatches.
\end{itemize}

\paragraph{Trampas y errores frecuentes.}
\begin{itemize}\setlength\itemsep{0pt}
	\item $Z[0]$ siempre vale $n$ (no se calcula).
	\item Si $i > r$, la ventana no aporta; empezar desde 0.
	\item El caracter separador en $P + \# + T$ debe ser uno que no aparezca en $P$ o $T$.
\end{itemize}

\paragraph{Codigo.}
Ver \texttt{cpp.tex}, seccion \textit{Z-algorithm}.

\vspace{0.5em}

\subsubsection{Rabin-Karp (single hash)}

\paragraph{Idea intuitiva.}
Rolling hash aplicado a pattern matching: calcular el hash del patron $P$ y de cada ventana de tamano $|P|$ en $T$. Si los hashes coinciden, verificar caracter por caracter (para descartar colisiones). Esto da $O(n + m)$ \textbf{esperado}. La razon de la eficiencia esperada: cada ventana tiene probabilidad $1/M$ de tener un hash coincidente por casualidad, asi que solo se verifica $O(n/M)$ ventanas en promedio.

\paragraph{Propiedades clave (invariantes).}
\begin{itemize}\setlength\itemsep{0pt}
	\item Si los hashes difieren, los substrings son distintos (sin colision).
	\item Si coinciden, \textbf{verificar} (la verificacion es $O(m)$ pero rara vez se ejecuta).
	\item Peor caso (muchas colisiones): $O(nm)$.
	\item Para $M \ge 10^9$, las colisiones son raras; en la practica nunca pasan.
\end{itemize}

\paragraph{Codigo clave.}
Rolling hash del texto (deslizar ventana):
\begin{verbatim}
long long hash = 0;
for (int i = 0; i < m; ++i) hash = (hash * BASE + text[i]) % MOD;
for (int i = m; i < n; ++i) {
    if (hash == target) check(i - m);
    hash = (hash - text[i - m] * h) * BASE + text[i];  // mod correctamente
    hash = (hash % MOD + MOD) % MOD;
}
\end{verbatim}

\paragraph{Complejidad.}
\begin{itemize}\setlength\itemsep{0pt}
	\item Esperado $O(n + m)$, peor $O(nm)$.
	\item Con doble hash y $M_1, M_2 \ge 10^9$, el peor caso es astronomicamente improbable.
\end{itemize}

\paragraph{Donde tocar para modificar.}
\begin{itemize}\setlength\itemsep{0pt}
	\item \textbf{Doble hash}: aniadir un segundo par \texttt{(BASE2, MOD2)} para reducir colisiones.
	\item \textbf{Varios patrones}: aniadir un set de hashes y buscar match con cualquiera.
	\item \textbf{Tamano de ventana variable}: si los patrones tienen largos distintos, usar el largo especifico por patron.
	\item \textbf{Multi-pattern search}: usar Aho-Corasick en su lugar.
\end{itemize}

\paragraph{Trampas y errores frecuentes.}
\begin{itemize}\setlength\itemsep{0pt}
	\item Overflow en \texttt{(tHash - text[i] * h) * BASE}: aplicar mod correctamente y aniadir \texttt{+MOD} antes de multiplicar.
	\item El factor $h = \text{BASE}^{m-1} \bmod \text{MOD}$ se precomputa.
	\item Si \texttt{MOD} no es primo, las colisiones son mas frecuentes; usar primo.
	\item En el peor caso, Rabin-Karp puede ser lento; preferir KMP si se requiere worst-case.
\end{itemize}

\paragraph{Codigo.}
Ver \texttt{cpp.tex}, seccion \textit{Rabin-Karp (single hash)}.

\vspace{0.5em}

\subsubsection{Aho-Corasick (basic)}

\paragraph{Idea intuitiva.}
Construye un \textbf{trie} con todos los patrones y aniade \textbf{failure links} analogos a los de KMP: el failure link de un nodo $v$ apunta al nodo mas largo que es sufijo del string representado por $v$ y sigue siendo prefijo de algun patron. Asi, al recorrer el texto, si hay mismatch, saltamos al failure link. Esto permite buscar \textbf{todos} los patrones en $O(|T| + \sum |P_i|)$. La idea: unificar varios KMP en una sola estructura.

\paragraph{Propiedades clave (invariantes).}
\begin{itemize}\setlength\itemsep{0pt}
	\item El trie almacena los patrones.
	\item \texttt{link[v]} (failure link) se calcula en BFS en orden de profundidad creciente.
	\item \texttt{out[v]} lista los IDs de patrones que terminan en $v$.
	\item Para queries eficientes, aniadir \textbf{output links}: propagar \texttt{out} de los failure links a los hijos.
\end{itemize}

\paragraph{Codigo clave.}
BFS para construir los failure links:
\begin{verbatim}
queue<int> q;
for (int c = 0; c < SIGMA; ++c)
    if (trie[0].next[c] != -1) {
        q.push(trie[0].next[c]);
        trie[trie[0].next[c]].link = 0;
    } else trie[0].next[c] = 0;  // goto a la raiz
while (!q.empty()) {
    int v = q.front(); q.pop();
    for (int pid : trie[trie[v].link].out)
        trie[v].out.push_back(pid);
    for (int c = 0; c < SIGMA; ++c) {
        if (trie[v].next[c] != -1) {
            trie[trie[v].next[c]].link = trie[trie[v].link].next[c];
            q.push(trie[v].next[c]);
        } else trie[v].next[c] = trie[trie[v].link].next[c];
    }
}
\end{verbatim}

\paragraph{Complejidad.}
\begin{itemize}\setlength\itemsep{0pt}
	\item Build $O(\sum |P_i| \cdot \Sigma)$, query $O(|T| + \text{ocurrencias})$.
\end{itemize}

\paragraph{Donde tocar para modificar.}
\begin{itemize}\setlength\itemsep{0pt}
	\item \textbf{Output links en build}: ya hecho en el snippet (\texttt{for(int pid : t[t[u].link].out) t[u].out.push\_back(pid);}).
	\item \textbf{Transiciones automaticas}: aniadir el goto \texttt{t[v].next[k] = t[t[v].link].next[k]} cuando el enlace directo es \texttt{-1}; asi el bucle de matching es un solo paso por caracter.
	\item \textbf{Aplicar a problemas tipo DP}: contar formas de construir strings evitando los patrones (ver el siguiente modulo).
	\item \textbf{2D Aho-Corasick}: aniadir una dimension para busqueda en matrices.
\end{itemize}

\paragraph{Trampas y errores frecuentes.}
\begin{itemize}\setlength\itemsep{0pt}
	\item La BFS debe empezar con la raiz en la cola y los hijos de la raiz con \texttt{link = 0}.
	\item Olvidar propagar \texttt{out} da ocurrencias perdidas.
	\item Si no se aniaden las transiciones automaticas (goto), el matching es mas lento.
	\item Con muchos patrones, el output puede ser enorme; aniadir contador en lugar de lista.
\end{itemize}

\paragraph{Codigo.}
Ver \texttt{cpp.tex}, seccion \textit{Aho-Corasick (basic)}.

\vspace{0.5em}

\subsubsection{Aho-Corasick DP for forbidden patterns}

\paragraph{Idea intuitiva.}
Contar strings de longitud $L$ que \textbf{evitan} un conjunto de patrones. Se modela como una DP sobre el \textbf{estado del automaton} de Aho-Corasick: $dp[v]$ = numero de formas de llegar al estado $v$ despues de cierto numero de caracteres. Transicion: por cada caracter $c$ en el alfabeto, ir al estado \texttt{next[v][c]}, pero \textbf{solo} si ese estado no tiene un patron (de lo contrario, el string contendria el patron prohibido). La idea: cada estado del automaton representa un ``prefijo conocido'' y las transiciones simulan ``agregar un caracter''.

\paragraph{Propiedades clave (invariantes).}
\begin{itemize}\setlength\itemsep{0pt}
	\item $dp$ se mantiene por nivel (longitud del string construido).
	\item Transicion: $ndp[u] = \sum_{c} dp[v]$ donde $u = next[v][c]$ y $u$ no es ``terminal'' (no tiene patron que termina ahi).
	\item Respuesta: suma de $dp[v]$ sobre todos los estados $v$ no terminales al final.
	\item Si todos los caracteres son permitidos, $|\Sigma|$ entra como factor.
\end{itemize}

\paragraph{Codigo clave.}
DP sobre el automaton:
\begin{verbatim}
vector<long long> dp(numStates, 0), ndp(numStates);
dp[0] = 1;  // empezar en la raiz
for (int step = 0; step < L; ++step) {
    fill(ndp.begin(), ndp.end(), 0);
    for (int v = 0; v < numStates; ++v) if (dp[v] && !isTerminal[v])
        for (int c = 0; c < SIGMA; ++c) {
            int u = next[v][c];
            if (!isTerminal[u]) ndp[u] = (ndp[u] + dp[v]) % MOD;
        }
    swap(dp, ndp);
}
long long ans = 0;
for (int v = 0; v < numStates; ++v) if (!isTerminal[v])
    ans = (ans + dp[v]) % MOD;
\end{verbatim}

\paragraph{Complejidad.}
\begin{itemize}\setlength\itemsep{0pt}
	\item $O(L \cdot |S| \cdot |\Sigma|)$ donde $|S|$ es el numero de estados y $|\Sigma|$ el alfabeto.
	\item Para $|\Sigma| = 2$ y $|S| \le 10^4$: $\sim 2 \cdot 10^{10}$ para $L = 10^6$ (puede ser lento).
\end{itemize}

\paragraph{Donde tocar para modificar.}
\begin{itemize}\setlength\itemsep{0pt}
	\item \textbf{Alfabeto mas pequeno}: si el problema es solo $\{0, 1\}$, $|\Sigma| = 2$.
	\item \textbf{Permitir ocurrencias parciales}: cambiar la condicion a ``no estado terminal'' o relajar segun el problema.
	\item \textbf{Recuperar el string}: aniadir backtracking si el problema lo pide.
	\item \textbf{Matriz de transicion}: para queries rapidas, precomputar $M^c$ con exponentiation.
\end{itemize}

\paragraph{Trampas y errores frecuentes.}
\begin{itemize}\setlength\itemsep{0pt}
	\item El snippet asume Aho-Corasick ya construido con transiciones automaticas. Sin eso, aniadir el goto en el bucle.
	\item Si $L$ es muy grande, la DP lineal es lenta; usar matrix exponentiation.
	\item Estados terminales deben excluirse de \textbf{ambos} lados (estado actual y siguiente).
\end{itemize}

\paragraph{Codigo.}
Ver \texttt{cpp.tex}, seccion \textit{Aho-Corasick DP for forbidden patterns}.

\subsection{Estructuras avanzadas}

\subsubsection{Suffix Array}

\paragraph{Idea intuitiva.}
Un \textbf{suffix array} $SA$ es un arreglo que contiene las posiciones de inicio de los sufijos del string $S$, ordenados lexicograficamente. Se construye en $O(n \log n)$ con doubling: ordenar por el primer caracter, luego por los primeros 2, luego 4, etc. Para hacer cada paso lineal, se usa \textbf{counting sort} sobre los ranks. La idea: cada nivel dobla la cantidad de caracteres considerados, asi que en $\log n$ niveles los sufijos estan completamente ordenados.

\paragraph{Propiedades clave (invariantes).}
\begin{itemize}\setlength\itemsep{0pt}
	\item \texttt{p[i]} = posicion de inicio del $i$-esimo sufijo (ordenado).
	\item \texttt{c[i]} = rank del sufijo que empieza en $i$.
	\item Cada iteracion duplica la longitud de la ``ventana'' considerada.
	\item El algoritmo termina cuando $2^k \ge n$.
	\item La permutacion \texttt{p} es estable entre iteraciones.
\end{itemize}

\paragraph{Codigo clave.}
El paso de doubling (clasificar por los primeros $2^k$ caracteres):
\begin{verbatim}
for (int k = 0; (1 << k) < n; ++k) {
    for (int i = 0; i < n; ++i) pn[i] = (p[i] - (1 << k) + n) % n;
    // counting sort por c[pn[i]] segundo, c[i] primero
    // ...
    cn[pn[0]] = 0;
    for (int i = 1; i < n; ++i)
        cn[pn[i]] = cn[pn[i-1]] + (pair != pair_prev);
    c = cn;
}
\end{verbatim}

\paragraph{Complejidad.}
\begin{itemize}\setlength\itemsep{0pt}
	\item $O(n \log n)$ con counting sort.
	\item Constante pequena; en la practica eficiente.
\end{itemize}

\paragraph{Donde tocar para modificar.}
\begin{itemize}\setlength\itemsep{0pt}
	\item \textbf{Suffix array en $O(n)$}: usar SA-IS (mas complejo, no incluido).
	\item \textbf{Buscar patron}: con SA + LCP es $O(|P| \log n)$ por busqueda binaria.
	\item \textbf{Suffix array de varios strings}: aniadir centinelas (\texttt{\$}) entre ellos.
	\item \textbf{Aniadir longitud}: aniadir \texttt{n - p[i]} como clave secundaria (sufijos mas cortos primero).
\end{itemize}

\paragraph{Trampas y errores frecuentes.}
\begin{itemize}\setlength\itemsep{0pt}
	\item La ``ventana'' $2^k$ se hace circular: \texttt{p[i] = (p[i] - (1<<k) + n) \% n}.
	\item Para contar bien los ranks duplicados, comparar pares \texttt{(c[i], c[(i+2\^{}k) \% n])}.
	\item El counting sort necesita un array de acumulacion de tamano $\le n$.
	\item Si los caracteres son negativos o fuera del alfabeto esperado, normalizar antes.
\end{itemize}

\paragraph{Codigo.}
Ver \texttt{cpp.tex}, seccion \textit{Suffix Array}.

\vspace{0.5em}

\subsubsection{LCP Array (Kasai)}

\paragraph{Idea intuitiva.}
El \textbf{LCP array} (Longest Common Prefix) guarda, para cada par de sufijos adyacentes en el SA, el largo de su prefijo comun. El algoritmo de \textbf{Kasai} lo calcula en $O(n)$ iterando los sufijos \textbf{en orden del string original} (no del SA), y reutiliza el LCP previo. La demostracion: el sufijo en $i+1$ se obtiene del sufijo en $i$ ``avanzando un caracter''; los LCPs decrecen como maximo 1 entre iteraciones.

\paragraph{Propiedades clave (invariantes).}
\begin{itemize}\setlength\itemsep{0pt}
	\item \texttt{lcp[i]} = LCP entre los sufijos en \texttt{sa[i]} y \texttt{sa[i-1]}.
	\item Se necesita \texttt{rank\_[i]} (la inversa del SA) para iterar.
	\item El truco de Kasai: si \texttt{lcp > 0} para el sufijo anterior, el actual parte de al menos \texttt{lcp - 1}.
	\item La suma de todos los LCPs da informacion sobre repeticiones.
\end{itemize}

\paragraph{Codigo clave.}
Algoritmo de Kasai:
\begin{verbatim}
vector<int> rank(n);
for (int i = 0; i < n; ++i) rank[sa[i]] = i;
int k = 0;
vector<int> lcp(n - 1);
for (int i = 0; i < n; ++i) {
    if (rank[i] == n - 1) { k = 0; continue; }
    int j = sa[rank[i] + 1];
    while (i + k < n && j + k < n && s[i+k] == s[j+k]) ++k;
    lcp[rank[i]] = k;
    if (k) --k;
}
\end{verbatim}

\paragraph{Complejidad.}
\begin{itemize}\setlength\itemsep{0pt}
	\item $O(n)$.
	\item Cada comparacion reduce $k$ o avanza linealmente; amortizado $O(1)$.
\end{itemize}

\paragraph{Donde tocar para modificar.}
\begin{itemize}\setlength\itemsep{0pt}
	\item \textbf{RMQ sobre LCP}: para queries ``LCP de dos sufijos'', encontrar sus ranks y hacer RMQ en \texttt{lcp}.
	\item \textbf{Contar substrings distintas}: $\sum (n - sa[i] - lcp[i])$.
	\item \textbf{Longest common substring}: maximo en \texttt{lcp}.
	\item \textbf{String matching}: busqueda binaria sobre SA + LCP verification.
\end{itemize}

\paragraph{Trampas y errores frecuentes.}
\begin{itemize}\setlength\itemsep{0pt}
	\item Si dos sufijos son iguales (raro, solo si el string tiene period $< n$), \texttt{lcp} puede ser \texttt{n - sa[i]}.
	\item En la iteracion, aniadir el chequeo \texttt{i + k < n} para evitar overflow.
	\item El \texttt{--k} en la salida es crucial para la complejidad amortizada.
\end{itemize}

\paragraph{Codigo.}
Ver \texttt{cpp.tex}, seccion \textit{LCP Array (Kasai)}.

\vspace{0.5em}

\subsubsection{Suffix Automaton}

\paragraph{Idea intuitiva.}
Un \textbf{suffix automaton} es un DFA minimo que reconoce todos los substrings del string original. Tiene la propiedad magica de tener $\le 2n - 1$ estados. Cada estado representa una clase de equivalencia de substrings (los que tienen el mismo conjunto de posiciones finales). El \textbf{link} de un estado apunta al estado del sufijo propio mas largo. La demostracion de $\le 2n - 1$: cada nuevo caracter anade a lo sumo 2 estados (uno nuevo y un clon), asi que en $n$ pasos hay $\le 2n - 1$ estados.

\paragraph{Propiedades clave (invariantes).}
\begin{itemize}\setlength\itemsep{0pt}
	\item \texttt{len[v]}: largo maximo de los substrings de la clase.
	\item \texttt{link[v]}: padre en el arbol de sufijos (suffix link).
	\item \texttt{next[v][c]}: transicion por caracter.
	\item \texttt{occ[v]}: numero de posiciones finales (cuantas veces aparece el patron).
	\item Cada estado tiene $\texttt{len[v]} - \texttt{len[link[v]]}$ substrings ``propios''.
\end{itemize}

\paragraph{Codigo clave.}
El corazon de la extension SAM:
\begin{verbatim}
void extend(int c) {
    int cur = size++;
    st[cur].len = st[last].len + 1;
    st[cur].occ = 1;
    int p = last;
    while (p != -1 && !st[p].next.count(c)) {
        st[p].next[c] = cur;
        p = st[p].link;
    }
    if (p == -1) st[cur].link = 0;
    else {
        int q = st[p].next[c];
        if (st[p].len + 1 == st[q].len) st[cur].link = q;
        else {
            int clone = size++;
            st[clone] = st[q];
            st[clone].len = st[p].len + 1;
            st[clone].occ = 0;  // clon NO es nueva posicion final
            while (p != -1 && st[p].next[c] == q) {
                st[p].next[c] = clone;
                p = st[p].link;
            }
            st[q].link = st[cur].link = clone;
        }
    }
    last = cur;
}
\end{verbatim}

\paragraph{Complejidad.}
\begin{itemize}\setlength\itemsep{0pt}
	\item Build $O(n)$ con la operacion de \textbf{clonar} estados.
	\item Memoria $O(n)$ (con map) o $O(n \cdot \Sigma)$ (con array).
\end{itemize}

\paragraph{Donde tocar para modificar.}
\begin{itemize}\setlength\itemsep{0pt}
	\item \textbf{Contar substrings distintas}: $\sum_{v \ne 0} (len[v] - len[link[v]])$.
	\item \textbf{Longest common substring}: navegar el SAM con el segundo string, manteniendo el match actual.
	\item \textbf{Substring count}: despues del build, ordenar estados por \texttt{len} descendente y propagar \texttt{occ}.
	\item \textbf{Endpos queries}: el SAM es el DFA canonico de los substrings; queries ``cuantas veces aparece $P

\vspace{0.5em}

\subsubsection{Lyndon factorization (Duval)}

\paragraph{Idea intuitiva.}
Un \textbf{string de Lyndon} es un string que es estrictamente menor que todos sus rotaciones. Toda cadena tiene una \textbf{factorizacion unica} de Lyndon (algoritmo de \textbf{Duval}): una secuencia de strings de Lyndon $L_1, L_2, \dots, L_k$ tales que $L_1 \ge L_2 \ge \dots \ge L_k$ y la concatenacion es el string original. Esto da $O(n)$ para encontrar los cortes. La demostracion de correctitud: cada factor es minimal en su clase de rotacion, asi que la concatenacion respeta el orden lexicografico.

\paragraph{Propiedades clave (invariantes).}
\begin{itemize}\setlength\itemsep{0pt}
	\item Cada factor de Lyndon es de la forma $uv$ donde $u$ se repite y $v$ es prefijo de $u$.
	\item La longitud del factor esta dada por la ``longitud del periodo''.
	\item Aplicacion: el ``Lyndon word'' es el menor en orden lexicografico entre todas las rotaciones.
	\item La factorizacion es \textbf{unica}.
\end{itemize}

\paragraph{Codigo clave.}
El algoritmo de Duval (iterativo):
\begin{verbatim}
int n = s.size();
int i = 0;
while (i < n) {
    int j = i + 1, k = i;
    while (j < n && s[k] <= s[j]) {
        if (s[k] < s[j]) k = i;
        else ++k;
        ++j;
    }
    int len = j - k;
    while (i < j) { cout << "factor: " << s.substr(i, len) << "\n"; i += len; }
}
\end{verbatim}

\paragraph{Complejidad.}
\begin{itemize}\setlength\itemsep{0pt}
	\item $O(n)$.
	\item Cada comparacion es $O(1)$; el puntero $k$ retrocede a lo sumo $n$ veces en total.
\end{itemize}

\paragraph{Donde tocar para modificar.}
\begin{itemize}\setlength\itemsep{0pt}
	\item \textbf{Encontrar la menor rotacion}: el menor sufijo lexicografico de \texttt{s + s} es la menor rotacion; la posicion se obtiene con un solo pase de Duval.
	\item \textbf{Algoritmo de Booth}: alternativa para rotacion minima, tambien $O(n)$.
	\item \textbf{Validar}: si necesitas verificar que un string es de Lyndon, compararlo con sus rotaciones.
	\item \textbf{Encontrar todos los factores}: guardar la posicion y longitud de cada uno.
\end{itemize}

\paragraph{Trampas y errores frecuentes.}
\begin{itemize}\setlength\itemsep{0pt}
	\item La condicion \texttt{s[j] <= s[k]} (no \texttt{<}) en el while de Duval es crucial; cambiar a \texttt{<} da comportamiento distinto.
	\item El caso $s[k] < s[j]$ resetea $k$ a $i$ (inicio del factor actual).
	\item El bucle externo debe avanzar $i$ por la longitud del factor, no por 1.
\end{itemize}

\paragraph{Codigo.}
Ver \texttt{cpp.tex}, seccion \textit{Lyndon factorization (Duval)}.

\subsection{Palindromos}

\subsubsection{Manacher}

\paragraph{Idea intuitiva.}
Calcula en $O(n)$ los \textbf{radios} de todos los palindromos centrados en cada posicion (impares y pares). La idea: mantener una ventana $[l, r]$ que es el palindromo centrado en un punto anterior mas largo encontrado. Si el nuevo centro $i$ cae dentro de $[l, r]$, entonces el palindromo centrado en $i$ tiene al menos el radio del palindromo ``espejo'' $l + r - i$; sino, empezar desde 1 (impar) o 0 (par). Despues, expandir lo que falte. La demostracion de $O(n)$: cada caracter se compara a lo sumo 2 veces (extension + limite por ventana).

\paragraph{Propiedades clave (invariantes).}
\begin{itemize}\setlength\itemsep{0pt}
	\item Conveccion del snippet: \texttt{odd[i]} = $k$ tal que el palindromo impar mas largo centrado en $i$ tiene longitud $2k - 1$. \texttt{even[i]} = $k$ para par centrado entre $i-1$ e $i$, longitud $2k$.
	\item $k$ cuenta tambien cuantos palindromos hay centrados ahi.
	\item Las queries \texttt{isPalindrome(l, r)} son $O(1)$.
	\item La ventana $[l, r]$ siempre contiene el palindromo mas largo encontrado hasta ahora.
\end{itemize}

\paragraph{Codigo clave.}
Manacher para palindromos impares:
\begin{verbatim}
int l = 0, r = -1;
for (int i = 0; i < n; ++i) {
    int k = (i > r) ? 1 : min(odd[l + r - i], r - i + 1);
    while (i - k >= 0 && i + k < n && s[i-k] == s[i+k]) ++k;
    odd[i] = k;
    if (i + k - 1 > r) { l = i - k + 1; r = i + k - 1; }
}
\end{verbatim}

\paragraph{Complejidad.}
\begin{itemize}\setlength\itemsep{0pt}
	\item Construccion $O(n)$, queries $O(1)$.
	\item Cada extension se cuenta una sola vez.
\end{itemize}

\paragraph{Donde tocar para modificar.}
\begin{itemize}\setlength\itemsep{0pt}
	\item \textbf{Solo contar palindromos}: ya implementado en \texttt{countPalindromes}.
	\item \textbf{Longest palindromic substring}: ya en \texttt{longestPalindrome}.
	\item \textbf{Modificar la conveccion de radios}: si queres ``radio = mitad de la longitud'' en vez de ``k palindromos'', ajustar el offset.
	\item \textbf{Queries con k no precomputado}: aniadir helpers especificos.
	\item \textbf{Modificar el caracter de relleno}: para multiples separadores, aniadir un caracter unico.
\end{itemize}

\paragraph{Trampas y errores frecuentes.}
\begin{itemize}\setlength\itemsep{0pt}
	\item La conveccion del snippet es \textbf{distinta} de la ``clasica de radio'' (donde radio = $(L-1)/2$); tener cuidado al copiar a otro codigo.
	\item Si los palindromos se solapan, la condicion \texttt{i + k - 1 > r} actualiza la ventana.
	\item En palindromos pares, el centro esta ``entre'' dos indices; aniadir offset en queries.
\end{itemize}

\paragraph{Codigo.}
Ver \texttt{cpp.tex}, seccion \textit{Manacher}.

% ============================================================
\section{Grafos}
% ============================================================

\subsection{Caminos minimos}

\subsubsection{Dijkstra}

\paragraph{Idea intuitiva.}
Encuentra el shortest path desde un origen $s$ a todos los vertices en $O((V + E) \log V)$ asumiendo \textbf{pesos no negativos}. La idea: usar una cola de prioridad (min-heap) ordenada por distancia tentativa. En cada paso, extraer el nodo con menor distancia; si ya fue procesado, saltar. Si al actualizar un vecino la distancia mejora, insertar/actualizar en la cola. La demostracion: cuando un nodo se extrae por primera vez, no hay forma de llegar a el con menor distancia (cualquier camino alternativo deberia pasar por nodos no procesados, con distancia $\ge$ la actual).

\paragraph{Propiedades clave (invariantes).}
\begin{itemize}\setlength\itemsep{0pt}
	\item Asume pesos $\ge 0$; con pesos negativos usar Bellman-Ford.
	\item Greedy correctness: la primera vez que se extrae un nodo, su distancia es la definitiva.
	\item Si se quiere reconstruir el camino, guardar el \texttt{parent[]}.
	\item Cada nodo se procesa a lo sumo $V$ veces (mas relajarions, pero solo $V$ definitivas).
\end{itemize}

\paragraph{Codigo clave.}
El bucle principal de Dijkstra:
\begin{verbatim}
priority_queue<pair<long long, int>, vector<...>, greater<...>> pq;
pq.push({0, source});
while (!pq.empty()) {
    auto [d, v] = pq.top(); pq.pop();
    if (d > dist[v]) continue;  // entrada vieja
    for (auto [w, u] : adj[v])
        if (dist[v] + w < dist[u]) {
            dist[u] = dist[v] + w;
            pq.push({dist[u], u});
        }
}
\end{verbatim}

\paragraph{Complejidad.}
\begin{itemize}\setlength\itemsep{0pt}
	\item $O((V + E) \log V)$ con heap binario, $O(V + E)$ con heap de Fibonacci.
	\item En la practica, $\sim 1.5 \cdot 10^{-7}$ segundos por operacion para $V = 10^5$.
\end{itemize}

\paragraph{Donde tocar para modificar.}
\begin{itemize}\setlength\itemsep{0pt}
	\item \textbf{Camino mas corto entre todos los pares}: si se necesita desde todos los origenes, usar Floyd o ejecutar Dijkstra $V$ veces.
	\item \textbf{Shortest path con aristas negativas limitadas}: Bellman-Ford.
	\item \textbf{Reconstruir camino}: aniadir \texttt{parent[v]} actualizado junto a \texttt{dist[v]}.
	\item \textbf{Dijkstra con potentials (Johnson)}: para aristas negativas si se puede reescalar.
	\item \textbf{Multi-source}: aniadir varios nodos fuentes iniciales al heap.
\end{itemize}

\paragraph{Trampas y errores frecuentes.}
\begin{itemize}\setlength\itemsep{0pt}
	\item Distancias no inicializadas a \texttt{INF}; el snippet usa \texttt{1e18}.
	\item El \texttt{if(dist[adjN] != 1e18)} para borrar de la cola es importante si usas \texttt{set} (no asi con \texttt{priority\_queue}).
	\item Con pesos grandes, \texttt{dist[v] + w} puede overflow; usar \texttt{LLONG\_MAX/2} como INF.
	\item No usar con pesos negativos: el algoritmo puede dar respuestas incorrectas.
\end{itemize}

\paragraph{Codigo.}
Ver \texttt{cpp.tex}, seccion \textit{Dijkstra}.

\vspace{0.5em}

\subsubsection{Bellman-Ford}

\paragraph{Idea intuitiva.}
Encuentra shortest path desde $s$ incluso con \textbf{pesos negativos} en $O(VE)$. La idea: relajar todas las aristas $V - 1$ veces. Cada relajacion propaga la mejor distancia a ``un paso mas''. Despues de $V - 1$ pasadas, si alguna arista aun se puede relajar, hay un \textbf{ciclo negativo} alcanzable. La demostracion: el camino mas corto tiene a lo sumo $V - 1$ aristas (sin ciclos); tras $V - 1$ relajaciones todas las distancias estan en su valor optimo.

\paragraph{Propiedades clave (invariantes).}
\begin{itemize}\setlength\itemsep{0pt}
	\item Funciona con pesos negativos (sin ciclos negativos).
	\item Tras $V - 1$ relajaciones, las distancias son definitivas (si no hay ciclo negativo).
	\item Una relajacion adicional detecta ciclos negativos: si \texttt{dist[u] + w < dist[v]}, hay ciclo negativo alcanzable desde $s$.
	\item Cada relajacion solo mejora distancias, asi que termina (no oscila).
\end{itemize}

\paragraph{Codigo clave.}
El doble bucle de relajacion:
\begin{verbatim}
for (int i = 0; i < n - 1; ++i)
    for (auto [u, v, w] : edges)
        if (dist[u] + w < dist[v])
            dist[v] = dist[u] + w;
// Detectar ciclo negativo
for (auto [u, v, w] : edges)
    if (dist[u] + w < dist[v]) { /* ciclo negativo */ }
\end{verbatim}

\paragraph{Complejidad.}
\begin{itemize}\setlength\itemsep{0pt}
	\item $O(VE)$.
	\item SPFA (con cola) en promedio es mucho mas rapido, peor caso igual.
\end{itemize}

\paragraph{Donde tocar para modificar.}
\begin{itemize}\setlength\itemsep{0pt}
	\item \textbf{SPFA}: optimizacion que relaja solo nodos cuya distancia cambio; en el peor caso sigue siendo $O(VE)$.
	\item \textbf{Detectar ciclo negativo global}: ejecutar Bellman-Ford desde un super-origen conectado a todos.
	\item \textbf{Camino mas corto con K aristas exactas}: $K$ iteraciones de Bellman-Ford.
	\item \textbf{Johnson's}: reescalar pesos con potentials para usar Dijkstra.
\end{itemize}

\paragraph{Trampas y errores frecuentes.}
\begin{itemize}\setlength\itemsep{0pt}
	\item Overflow si los pesos son $\ge 2^{52}$; usar \texttt{long long} y \texttt{\_\_int128} si hace falta.
	\item Si el grafo es denso ($E \approx V^2$), $O(VE)$ es prohibitivo.
	\item El chequeo de ciclo negativo debe hacerse tras las $V - 1$ iteraciones, no durante.
	\item Distancias no inicializadas a \texttt{INF}; el primer nodo (source) tiene distancia 0.
\end{itemize}

\paragraph{Codigo.}
Ver \texttt{cpp.tex}, seccion \textit{Bellman-Ford}.

\vspace{0.5em}

\subsubsection{Floyd-Warshall}

\paragraph{Idea intuitiva.}
Shortest path entre \textbf{todos los pares} de vertices en $O(V^3)$. DP: $d[i][j]$ = shortest path usando solo nodos intermedios en $\{0, \dots, k\}$. Transicion: $d_{k+1}[i][j] = \min(d_k[i][j], d_k[i][k+1] + d_k[k+1][j])$. La demostracion: tras $k$ iteraciones, $d[i][j]$ es el shortest path con intermedios en $\{0, \dots, k-1\}$; aniadir el nodo $k$ da la recurrencia clasica.

\paragraph{Propiedades clave (invariantes).}
\begin{itemize}\setlength\itemsep{0pt}
	\item Detecta ciclos negativos: tras el algoritmo, $d[i][i] < 0$ indica ciclo negativo alcanzable desde $i$.
	\item Funciona con pesos negativos (sin ciclos negativos).
	\item Memoria $O(V^2)$; con $V \le 400$ es OK, con mas conviene Dijkstra $V$ veces.
	\item Es la unica opcion cuando se quieren shortest paths entre todos los pares y el grafo es denso.
\end{itemize}

\paragraph{Codigo clave.}
El triple bucle clasico:
\begin{verbatim}
for (int k = 0; k < n; ++k)
    for (int i = 0; i < n; ++i)
        for (int j = 0; j < n; ++j)
            d[i][j] = min(d[i][j], d[i][k] + d[k][j]);
// detectar ciclo negativo
for (int i = 0; i < n; ++i) if (d[i][i] < 0) tiene_ciclo_negativo = true;
\end{verbatim}

\paragraph{Complejidad.}
\begin{itemize}\setlength\itemsep{0pt}
	\item $O(V^3)$ tiempo, $O(V^2)$ memoria.
	\item Para $V \le 400$: $\sim 6.4 \cdot 10^7$ operaciones (factible).
\end{itemize}

\paragraph{Donde tocar para modificar.}
\begin{itemize}\setlength\itemsep{0pt}
	\item \textbf{Reconstruir camino}: aniadir \texttt{next[i][j]} y backtrack.
	\item \textbf{Transitividad}: el algoritmo tambien calcula transitividad (\texttt{d[i][j] != INF} significa alcanzable).
	\item \textbf{Pesos doubles}: aniadir tolerancia (\texttt{d[i][j] > d[i][k] + d[k][j] - eps}).
	\item \textbf{Grafo no denso}: usar Dijkstra $V$ veces en vez de Floyd.
\end{itemize}

\paragraph{Trampas y errores frecuentes.}
\begin{itemize}\setlength\itemsep{0pt}
	\item Para detectar ciclo negativo correctamente, aniadir un chequeo \texttt{if (ans[i][i] < 0)} tras el triple bucle.
	\item El orden de los bucles \texttt{k, i, j} es importante; invertir da resultados diferentes.
	\item Con \texttt{double}, usar tolerancia \texttt{< eps} en vez de \texttt{<} para evitar loops infinitos.
\end{itemize}

\paragraph{Codigo.}
Ver \texttt{cpp.tex}, seccion \textit{Floyd-Warshall}.

\subsection{Componentes y cortes}

\subsubsection{Kosaraju SCC}

\paragraph{Idea intuitiva.}
Encuentra las \textbf{componentes fuertemente conexas} (SCCs) en $O(V + E)$. Algoritmo en 3 pasos:
\begin{itemize}\setlength\itemsep{0pt}
	\item DFS en $G$ rellenando una pila en \textbf{post-orden} (los nodos que terminan tarde quedan arriba).
	\item Construir $G^T$ (grafo transpuesto).
	\item DFS en $G^T$ procesando nodos en orden de la pila (de mayor a menor tiempo); cada DFS da una SCC.
\end{itemize}
La demostracion: dos nodos $u, v$ estan en la misma SCC sii hay camino $u \leadsto v$ y $v \leadsto u$; Kosaraju captura exactamente esto al explorar $G^T$ desde los nodos con mayor tiempo de finalizacion en $G$.

\paragraph{Propiedades clave (invariantes).}
\begin{itemize}\setlength\itemsep{0pt}
	\item Cada SCC es maximal fuertemente conexa.
	\item El \textbf{grafo de SCCs} (metagraf) es un DAG.
	\item Kosaraju y Tarjan dan el mismo resultado; Tarjan es ``mas rapido'' en practica (un solo DFS).
	\item Cada nodo pertenece a exactamente una SCC.
\end{itemize}

\paragraph{Codigo clave.}
La estructura de Kosaraju:
\begin{verbatim}
vector<int> order;
vector<bool> used(n);
function<void(int)> dfs1 = [&](int v) {
    used[v] = true;
    for (int u : g[v]) if (!used[u]) dfs1(u);
    order.push_back(v);  // post-orden
};
function<void(int, int)> dfs2 = [&](int v, int root) {
    comp[v] = root;
    for (int u : gt[v]) if (comp[u] == -1) dfs2(u, root);
};
// 1: DFS en G para llenar order
// 2: DFS en G^T en orden inverso para asignar SCCs
\end{verbatim}

\paragraph{Complejidad.}
\begin{itemize}\setlength\itemsep{0pt}
	\item $O(V + E)$.
	\item Constante 2 (dos DFS + el grafo transpuesto).
\end{itemize}

\paragraph{Donde tocar para modificar.}
\begin{itemize}\setlength\itemsep{0pt}
	\item \textbf{2-SAT}: SCC sobre grafo de implicaciones.
	\item \textbf{Construir el DAG de SCCs}: aniadir aristas entre SCCs diferentes.
	\item \textbf{Tamano de cada SCC}: aniadir contador durante el segundo DFS.
	\item \textbf{Grafo grande}: usar Tarjan para evitar construir $G^T$.
\end{itemize}

\paragraph{Trampas y errores frecuentes.}
\begin{itemize}\setlength\itemsep{0pt}
	\item El snippet usa \texttt{vector<...> adj[n]} (VLA); reemplazable por \texttt{vector<vector<pair<int,int>>} para portabilidad.
	\item El orden del segundo DFS debe ser \textbf{inverso} al de finalizacion del primero.
	\item Si el grafo no es conexo, aniadir un loop sobre todos los nodos en el primer DFS.
	\item Confundir SCC con componentes conexas (las primeras son para dirigidos, las segundas para no dirigidos).
\end{itemize}

\paragraph{Codigo.}
Ver \texttt{cpp.tex}, seccion \textit{Kosaraju SCC}.

\vspace{0.5em}

\subsubsection{Tarjan (puentes)}

\paragraph{Idea intuitiva.}
Encuentra los \textbf{puentes} (bridges): aristas cuya remocion \textbf{desconecta} el grafo. Usa un solo DFS manteniendo \texttt{tin[v]} (tiempo de descubrimiento) y \texttt{low[v]} (el menor \texttt{tin} alcanzable desde $v$ via back-edges). Una arista $v - u$ es puente si $low[u] > tin[v]$. La demostracion: si $low[u] > tin[v]$, entonces $u$ (y su subarbol) no pueden ``escapar'' de $v$; quitar la arista $v - u$ parte el grafo.

\paragraph{Propiedades clave (invariantes).}
\begin{itemize}\setlength\itemsep{0pt}
	\item \texttt{low[v] = min(tin[v], tin[w])} para todo back-edge $(v, w)$, y \texttt{min(low[u])} para hijos en el DFS.
	\item Solo se consideran las aristas del \textbf{arbol DFS} (no back-edges) para \texttt{low}.
	\item Cada nodo tiene exactamente un \texttt{tin}; cada back-edge se procesa una vez.
\end{itemize}

\paragraph{Codigo clave.}
El DFS de Tarjan para bridges:
\begin{verbatim}
void dfs(int v, int pEdge) {
    used[v] = true;
    tin[v] = low[v] = timer++;
    for (auto [u, id] : adj[v]) {
        if (id == pEdge) continue;  // no contar el edge al padre
        if (used[u]) {
            low[v] = min(low[v], tin[u]);  // back-edge
        } else {
            dfs(u, id);
            low[v] = min(low[v], low[u]);
            if (low[u] > tin[v]) bridges.push_back(id);  // es puente
        }
    }
}
\end{verbatim}

\paragraph{Complejidad.}
\begin{itemize}\setlength\itemsep{0pt}
	\item $O(V + E)$.
	\item Una sola pasada de DFS.
\end{itemize}

\paragraph{Donde tocar para modificar.}
\begin{itemize}\setlength\itemsep{0pt}
	\item \textbf{Bridge tree}: contractar cada 2-edge-connected component a un nodo; el resultado es un arbol.
	\item \textbf{Articulation points}: variante con condicion diferente (ver el siguiente modulo).
	\item \textbf{Grafo no dirigido}: solo funciona en no dirigidos; en dirigidos hay conceptos analogos (strong bridges, dominator tree).
	\item \textbf{Reconstruir componentes 2-edge-connected}: BFS desde cada nodo excluyendo bridges.
\end{itemize}

\paragraph{Trampas y errores frecuentes.}
\begin{itemize}\setlength\itemsep{0pt}
	\item No confundir el parent con un back-edge al padre; chequear \texttt{if (adjN == parent) continue;} y contar solo las veces que se ``visita'' como hijo.
	\item Para multigrafos, los edges deben tener IDs unicos para distinguirlos.
	\item Si el grafo no es conexo, aniadir loop externo sobre todos los nodos.
\end{itemize}

\paragraph{Codigo.}
Ver \texttt{cpp.tex}, seccion \textit{Tarjan (puentes)}.

\vspace{0.5em}

\subsubsection{Tarjan (puntos de articulacion)}

\paragraph{Idea intuitiva.}
Encuentra los \textbf{puntos de articulacion} (cut vertices): vertices cuya remocion \textbf{desconecta} el grafo. Mismo setup que para bridges: \texttt{tin[v]}, \texttt{low[v]}. Un nodo $v$ (no raiz) es de articulacion si existe un hijo $u$ con $low[u] \ge tin[v]$. La \textbf{raiz} es de articulacion si tiene $\ge 2$ hijos en el DFS tree. La razon: para $v$ no raiz, la condicion $\ge$ indica que el subarbol de $u$ depende de $v$; quitar $v$ los desconecta. Para la raiz, perder $\ge 2$ hijos parte el DFS en $\ge 2$ componentes.

\paragraph{Propiedades clave (invariantes).}
\begin{itemize}\setlength\itemsep{0pt}
	\item $low[u] \ge tin[v]$ significa que el subarbol de $u$ no tiene back-edge ``por encima'' de $v$.
	\item La raiz es caso especial: tiene que perder mas de un hijo para que el grafo se parta.
	\item Cada nodo se evalua una sola vez.
\end{itemize}

\paragraph{Codigo clave.}
La condicion de articulacion:
\begin{verbatim}
void dfs(int v, int p) {
    used[v] = true;
    tin[v] = low[v] = timer++;
    int children = 0;
    for (int u : adj[v]) {
        if (u == p) continue;
        if (used[u]) low[v] = min(low[v], tin[u]);
        else {
            dfs(u, v);
            low[v] = min(low[v], low[u]);
            if (low[u] >= tin[v] && p != -1)  // no raiz
                isArt[v] = true;
            ++children;
        }
    }
    if (p == -1 && children > 1) isArt[v] = true;  // raiz especial
}
\end{verbatim}

\paragraph{Complejidad.}
\begin{itemize}\setlength\itemsep{0pt}
	\item $O(V + E)$.
	\item Una sola pasada de DFS.
\end{itemize}

\paragraph{Donde tocar para modificar.}
\begin{itemize}\setlength\itemsep{0pt}
	\item \textbf{Block-cut tree}: estructura bipartita con BCCs (2-vertex-connected components) y articulation points.
	\item \textbf{Verificar biconectividad}: chequear que no hay puntos de articulacion.
	\item \textbf{Componentes biconnected}: extender Tarjan para enumerar BCCs.
\end{itemize}

\paragraph{Trampas y errores frecuentes.}
\begin{itemize}\setlength\itemsep{0pt}
	\item No olvidar el caso especial de la raiz.
	\item La condicion usa \texttt{>=}, no \texttt{>}; un hijo con \texttt{low[u] == tin[v]} tambien cuenta.
	\item Multigrafos pueden tener ``falsos'' articulation points; verificar caso a caso.
\end{itemize}

\paragraph{Codigo.}
Ver \texttt{cpp.tex}, seccion \textit{Tarjan (puntos de articulacion)}.

\vspace{0.5em}

\subsubsection{Tarjan SCC}

\paragraph{Idea intuitiva.}
Variante del SCC que usa \textbf{un solo DFS} con una pila explicita. Mantiene \texttt{disc[v]}, \texttt{low[v]}, y una pila de nodos ``en la SCC actual''. Cuando \texttt{low[v] == disc[v]}, $v$ es la raiz de una SCC; popear hasta $v$ inclusive. La demostracion: $low[v] == disc[v]$ indica que $v$ no tiene back-edges a ancestros; entonces $v$ y los nodos en la pila hasta $v$ forman una SCC maximal.

\paragraph{Propiedades clave (invariantes).}
\begin{itemize}\setlength\itemsep{0pt}
	\item Mas eficiente que Kosaraju (un solo DFS).
	\item \texttt{low[v] = min(low[u], disc[w])} para back-edges o recursion.
	\item La pila asegura que solo los nodos en la SCC actual se popean.
	\item Cada nodo se pushea y popea a lo sumo una vez.
\end{itemize}

\paragraph{Codigo clave.}
El DFS de Tarjan SCC:
\begin{verbatim}
void dfs(int v) {
    disc[v] = low[v] = timer++;
    st.push(v); inStack[v] = true;
    for (int u : g[v]) {
        if (disc[u] == -1) { dfs(u); low[v] = min(low[v], low[u]); }
        else if (inStack[u]) low[v] = min(low[v], disc[u]);
    }
    if (low[v] == disc[v]) {
        // popear hasta v
        while (true) {
            int u = st.top(); st.pop(); inStack[u] = false;
            comp[u] = current_scc;
            if (u == v) break;
        }
        ++current_scc;
    }
}
\end{verbatim}

\paragraph{Complejidad.}
\begin{itemize}\setlength\itemsep{0pt}
	\item $O(V + E)$.
	\item Constante similar a Kosaraju pero sin construir $G^T$.
\end{itemize}

\paragraph{Donde tocar para modificar.}
\begin{itemize}\setlength\itemsep{0pt}
	\item \textbf{Tamano de cada SCC}: aniadir contador durante el pop.
	\item \textbf{Grafo condensation}: aniadir aristas entre SCCs en el DAG resultante.
	\item \textbf{2-SAT}: aniadir la construccion del grafo de implicaciones.
\end{itemize}

\paragraph{Trampas y errores frecuentes.}
\begin{itemize}\setlength\itemsep{0pt}
	\item \texttt{inStack[u]} debe chequearse antes de usar \texttt{low[v] = min(low[v], disc[u])} en back-edges.
	\item La pila global es importante; sin ella no se puede extraer la SCC.
	\item Si el grafo es muy grande, recursion puede overflow; aniadir version iterativa.
\end{itemize}

\paragraph{Codigo.}
Ver \texttt{cpp.tex}, seccion \textit{Tarjan SCC}.

\subsection{Matching}

\subsubsection{Kuhn (bipartite matching)}

\paragraph{Idea intuitiva.}
Encuentra un \textbf{matching maximo} en un grafo bipartito en $O(VE)$. Para cada nodo $u$ en $L$, intentar encontrar un vecino $v$ en $R$ que este libre o cuyo match se pueda ``reorganizar'' recursivamente. Si se encuentra, matchear $u - v$. La idea: si $u$ quiere un vecino $v$ ya ocupado, intentamos ``mover'' el match de $v$ a otro nodo; recursivamente, esto encuentra un augmenting path si existe.

\paragraph{Propiedades clave (invariantes).}
\begin{itemize}\setlength\itemsep{0pt}
	\item \texttt{matchR[v]} = el nodo de $L$ matcheado a $v$ (o $-1$ si libre).
	\item \texttt{vis[]} se resetea \textbf{cada intento} de $u$.
	\item Si el matching es perfecto, $|L| = |R| = $ tamano del matching.
	\item Kuhn no garantiza maximalidad greedy; puede haber augmenting paths no encontrados.
\end{itemize}

\paragraph{Codigo clave.}
La recursion del augmenting path:
\begin{verbatim}
bool tryKuhn(int v) {
    if (vis[v]) return false;
    vis[v] = true;
    for (int u : g[v]) {
        if (matchR[u] == -1 || tryKuhn(matchR[u])) {
            matchR[u] = v;
            return true;
        }
    }
    return false;
}
// Por cada v en L:
fill(vis, vis+n, false);
if (tryKuhn(v)) ++matches;
\end{verbatim}

\paragraph{Complejidad.}
\begin{itemize}\setlength\itemsep{0pt}
	\item $O(VE)$ en el peor caso.
	\item Con Hopcroft-Karp: $O(\sqrt{V} \cdot E)$, mucho mejor para grafos densos.
\end{itemize}

\paragraph{Donde tocar para modificar.}
\begin{itemize}\setlength\itemsep{0pt}
	\item \textbf{Matching bipartito minimo (min vertex cover)}: Konig's theorem.
	\item \textbf{Matching general}: blossom algorithm (no incluido).
	\item \textbf{Aniadir capacidades}: si los nodos de $R$ tienen capacidad $> 1$, partir cada uno en capacidad copias.
	\item \textbf{Output del matching}: el array \texttt{matchR} da los pares finales.
\end{itemize}

\paragraph{Trampas y errores frecuentes.}
\begin{itemize}\setlength\itemsep{0pt}
	\item Olvidar resetear \texttt{vis} por intento; eso es lo que evita ``ciclos infinitos'' en la recursion.
	\item Para grafos grandes ($V > 10^4$), Kuhn puede ser muy lento; usar Hopcroft-Karp.
	\item Si el matching es maximo pero no perfecto, $|L|$ o $|R|$ puede ser mayor; el matching es solo maximo.
\end{itemize}

\paragraph{Codigo.}
Ver \texttt{cpp.tex}, seccion \textit{Kuhn (bipartite matching)}.

\vspace{0.5em}

\subsubsection{Hopcroft-Karp}

\paragraph{Idea intuitiva.}
Matching bipartito maximo en $O(\sqrt{V} E)$. Algoritmo en \textbf{fases}: en cada fase, hacer BFS para encontrar los \textbf{caminos de aumento mas cortos} desde nodos libres de $L$; luego DFS solo por esos caminos. Cada fase aumenta el matching en $\ge 1$ y toma $O(E)$. La demostracion: tras una fase, todos los augmenting paths tienen longitud $\ge$ la nueva minima; eso sucede cada $\sqrt{V}$ fases (longitud minima crece geometricamente).

\paragraph{Propiedades clave (invariantes).}
\begin{itemize}\setlength\itemsep{0pt}
	\item \texttt{dist[u]} = distancia en niveles desde un nodo libre de $L$.
	\item BFS termina cuando no hay mas camino a un nodo libre de $R$.
	\item DFS respeta los niveles (solo sigue \texttt{dist[matchV[v]] == dist[u] + 1}).
	\item Cada fase aumenta el matching en $\ge 1$ arista.
\end{itemize}

\paragraph{Codigo clave.}
El BFS de Hopcroft-Karp:
\begin{verbatim}
bool bfs() {
    queue<int> q;
    for (int v : L) {
        if (matchU[v] == -1) { dist[v] = 0; q.push(v); }
        else dist[v] = -1;
    }
    bool found = false;
    while (!q.empty()) {
        int v = q.front(); q.pop();
        for (int u : g[v]) {
            if (matchV[u] != -1 && dist[matchV[u]] == -1) {
                dist[matchV[u]] = dist[v] + 1;
                q.push(matchV[u]);
            }
            if (matchV[u] == -1) found = true;  // llegamos a libre
        }
    }
    return found;
}
\end{verbatim}

\paragraph{Complejidad.}
\begin{itemize}\setlength\itemsep{0pt}
	\item $O(\sqrt{V} E)$.
	\item Para $V = 10^4$, $E = 10^5$: $\sim 3 \cdot 10^7$ operaciones.
\end{itemize}

\paragraph{Donde tocar para modificar.}
\begin{itemize}\setlength\itemsep{0pt}
	\item \textbf{Reconstruir el matching}: ya viene en \texttt{matchU}/\texttt{matchV}.
	\item \textbf{Aniadir pesos}: Hungarian algorithm.
	\item \textbf{Densidad alta}: si $E \approx V^2$, sigue siendo util pero verificar constantes.
	\item \textbf{Aniadir capacidades}: partir los nodos con capacidad $> 1$ en varias copias.
\end{itemize}

\paragraph{Trampas y errores frecuentes.}
\begin{itemize}\setlength\itemsep{0pt}
	\item \texttt{dist[matchV[v]] == dist[u] + 1} (no \texttt{==} a secas); sin esto, el DFS no respeta los niveles y se pierde la complejidad.
	\item El BFS debe marcar los niveles correctamente; resetear \texttt{dist} para nodos alcanzables solo via matching.
	\item El DFS debe respetar los niveles; cualquier ``atajo'' invalida la cota.
\end{itemize}

\paragraph{Codigo.}
Ver \texttt{cpp.tex}, seccion \textit{Hopcroft-Karp}.

\subsection{Basicos}

\subsubsection{DFS / BFS / componentes}

\paragraph{Idea intuitiva.}
\textbf{DFS} (Depth-First Search) y \textbf{BFS} (Breadth-First Search) son los dos recorridos basicos de grafos. DFS usa una pila (o recursion) y explora tan profundo como puede antes de retroceder; BFS usa una cola y explora por ``niveles'' de distancia. La cantidad de \textbf{componentes conexas} se cuenta iniciando un BFS/DFS desde cada nodo no visitado. La diferencia clave: BFS da el camino mas corto en aristas; DFS da un orden de finalizacion util para problemas topologicos.

\paragraph{Propiedades clave (invariantes).}
\begin{itemize}\setlength\itemsep{0pt}
	\item BFS da la \textbf{distancia en aristas} mas corta desde el origen.
	\item DFS da un \textbf{orden topologico} (si se aniaden al final).
	\item Ambos son $O(V + E)$.
	\item Un grafo con $C$ componentes conexas satisface $\sum \text{subtree\_size}_v = V$ y $\sum E_v = 2E$.
	\item DFS en arbol da un spanning tree; las ``back edges'' corresponden a ciclos.
\end{itemize}

\paragraph{Codigo clave.}
BFS clasico:
\begin{verbatim}
queue<int> q;
q.push(source);
dist[source] = 0;
while (!q.empty()) {
    int v = q.front(); q.pop();
    for (int u : adj[v])
        if (dist[u] == -1) {
            dist[u] = dist[v] + 1;
            q.push(u);
        }
}
\end{verbatim}

\paragraph{Complejidad.}
\begin{itemize}\setlength\itemsep{0pt}
	\item $O(V + E)$ tiempo, $O(V)$ memoria.
	\item Cada arista se procesa exactamente 2 veces (ida y vuelta) en no dirigidos.
\end{itemize}

\paragraph{Donde tocar para modificar.}
\begin{itemize}\setlength\itemsep{0pt}
	\item \textbf{Aniadir pesos}: si las aristas tienen peso y queres shortest path, BFS ya no sirve; usar Dijkstra.
	\item \textbf{Grafo dirigido}: en DFS, aniadir check \texttt{if (state[u] == 1)} para detectar ciclos.
	\item \textbf{Bipartite check}: BFS alternando colores (ver el siguiente modulo).
	\item \textbf{Componentes fuertemente conexas (dirigidos)}: usar Tarjan o Kosaraju.
	\item \textbf{BFS en matriz}: aniadir offsets por direccion para grid 2D.
\end{itemize}

\paragraph{Trampas y errores frecuentes.}
\begin{itemize}\setlength\itemsep{0pt}
	\item Stack overflow con DFS recursivo para $V > 10^5$; usar version iterativa con pila explicita.
	\item En BFS, olvidar marcar \texttt{dist[u] = dist[v] + 1} antes de pushear da visitas duplicadas.
	\item Para grafos dirigidos, ``componente conexa'' no esta definida; usar SCC.
\end{itemize}

\paragraph{Codigo.}
Ver \texttt{cpp.tex}, seccion \textit{DFS / BFS / componentes}.

\vspace{0.5em}

\subsubsection{Topological sort (Kahn + DFS)}

\paragraph{Idea intuitiva.}
Un \textbf{orden topologico} de un DAG es una permutacion de los vertices tal que toda arista $u \to v$ tiene $u$ antes que $v$. Dos algoritmos clasicos:
\begin{itemize}\setlength\itemsep{0pt}
	\item \textbf{Kahn}: BFS sobre nodos con \texttt{indeg == 0}; al ``consumir'' un nodo, decrementar el indeg de sus vecinos; los que lleguen a 0 se encolan.
	\item \textbf{DFS}: hacer DFS; al terminar un nodo (estado = 2), aniadirlo al final; el orden es el reverso.
\end{itemize}
La demostracion: en un DAG siempre hay al menos un nodo con indeg 0; quitarlo mantiene el DAG; por induccion, se obtiene el orden.

\paragraph{Propiedades clave (invariantes).}
\begin{itemize}\setlength\itemsep{0pt}
	\item Existe sii el grafo es un DAG.
	\item Si al final \texttt{order.size() < V}, hay un ciclo.
	\item Kahn da el orden ``natural'' (los que pueden ir primero); DFS da el orden ``al reves'' (los que pueden ir al final primero).
	\item La cantidad de ordenes topologicos puede ser exponencial.
\end{itemize}

\paragraph{Codigo clave.}
Kahn con priority queue (orden lexicografico):
\begin{verbatim}
priority_queue<int, vector<int>, greater<int>> pq;
for (int v = 0; v < n; ++v) if (indeg[v] == 0) pq.push(v);
while (!pq.empty()) {
    int v = pq.top(); pq.pop();
    order.push_back(v);
    for (int u : adj[v]) {
        if (--indeg[u] == 0) pq.push(u);
    }
}
\end{verbatim}

\paragraph{Complejidad.}
\begin{itemize}\setlength\itemsep{0pt}
	\item $O(V + E)$ ambos.
	\item Con priority queue, $O((V + E) \log V)$.
\end{itemize}

\paragraph{Donde tocar para modificar.}
\begin{itemize}\setlength\itemsep{0pt}
	\item \textbf{Detectar ciclos}: comparar \texttt{order.size()} con \texttt{V}.
	\item \textbf{Orden lexicografico minimo}: usar priority queue en Kahn.
	\item \textbf{Topological sort con pesos}: combinar con shortest path en DAG (un solo pase en orden topologico).
	\item \textbf{Encontrar todos los ordenes}: DP sobre subsets (solo para $V$ pequeno).
\end{itemize}

\paragraph{Trampas y errores frecuentes.}
\begin{itemize}\setlength\itemsep{0pt}
	\item En el snippet, Kahn retorna orden vacio si hay ciclo (util para detectar).
	\item En DFS, no olvidar \texttt{reverse(order)}.
	\item Confundir indeg con outdeg; el orden Kahn quita primero los nodos ``fuente''.
	\item Para grafos no DAG, el algoritmo termina con \texttt{order.size() < V}.
\end{itemize}

\paragraph{Codigo.}
Ver \texttt{cpp.tex}, seccion \textit{Topological sort (Kahn + DFS)}.

\vspace{0.5em}

\subsubsection{Bipartite check + 2-coloring}

\paragraph{Idea intuitiva.}
Un grafo es \textbf{bipartito} si sus vertices se pueden dividir en dos conjuntos $L$ y $R$ tales que toda arista va de $L$ a $R$. Esto equivale a decir que se puede \textbf{2-colorear} sin que dos vertices adyacentes tengan el mismo color. BFS alternando colores detecta esto: si en algun momento dos adyacentes tienen el mismo color, no es bipartito. La demostracion: por cada componente, fijar el color de un nodo determina todos los demas; si en algun paso hay conflicto, hay un ciclo impar (no bipartito).

\paragraph{Propiedades clave (invariantes).}
\begin{itemize}\setlength\itemsep{0pt}
	\item Grafo bipartito $\Leftrightarrow$ sin ciclos impares.
	\item Conexo: si un componente no es bipartito, todo el grafo no lo es.
	\item Si hay \textbf{multiple componentes}, cada uno se colorea independientemente.
	\item La 2-coloracion es unica salvo swap de colores.
\end{itemize}

\paragraph{Codigo clave.}
BFS con 2-coloreo:
\begin{verbatim}
vector<int> color(n, -1);
for (int v = 0; v < n; ++v) if (color[v] == -1) {
    color[v] = 0;
    queue<int> q; q.push(v);
    while (!q.empty()) {
        int u = q.front(); q.pop();
        for (int w : adj[u]) {
            if (color[w] == -1) { color[w] = 1 - color[u]; q.push(w); }
            else if (color[w] == color[u]) return false;
        }
    }
}
return true;
\end{verbatim}

\paragraph{Complejidad.}
\begin{itemize}\setlength\itemsep{0pt}
	\item $O(V + E)$.
	\item Memoria $O(V)$.
\end{itemize}

\paragraph{Donde tocar para modificar.}
\begin{itemize}\setlength\itemsep{0pt}
	\item \textbf{Colorear desde multiples fuentes}: BFS desde cada nodo no coloreado.
	\item \textbf{Bipartito + matching}: Hopcroft-Karp.
	\item \textbf{Grafo bipartito ponderado}: Hungarian algorithm.
	\item \textbf{Verificar bipartito en subgrafo}: solo aplicar BFS a los nodos relevantes.
\end{itemize}

\paragraph{Trampas y errores frecuentes.}
\begin{itemize}\setlength\itemsep{0pt}
	\item Verificar \textbf{todos} los componentes, no solo el primero.
	\item Para grafos con auto-loops, un auto-loop hace al grafo no bipartito.
	\item En multigrafos, las aristas multiples no afectan el chequeo de bipartito.
\end{itemize}

\paragraph{Codigo.}
Ver \texttt{cpp.tex}, seccion \textit{Bipartite check + 2-coloring}.

\vspace{0.5em}

\subsubsection{0-1 BFS}

\paragraph{Idea intuitiva.}
Variante de BFS para grafos con \textbf{pesos 0 o 1}. Usa un \texttt{deque}: aristas de peso 0 se encolan al frente (\texttt{push\_front}), aristas de peso 1 al final (\texttt{push\_back}). Asi, la primera vez que se extrae un nodo, su distancia es la definitiva, igual que en Dijkstra pero en $O(V + E)$. La intuicion: el deque mantiene los nodos ordenados por distancia; las aristas de peso 0 no cambian la distancia (entran al frente), las de peso 1 la incrementan (entran al final).

\paragraph{Propiedades clave (invariantes).}
\begin{itemize}\setlength\itemsep{0pt}
	\item Asume pesos \textbf{solo} 0 o 1.
	\item Mejor caso que Dijkstra (sin factor $\log V$).
	\item No funciona con pesos $\ge 2$.
	\item La cola mantiene los nodos ordenados por distancia no decreciente.
\end{itemize}

\paragraph{Codigo clave.}
El deque doble:
\begin{verbatim}
deque<int> q;
q.push_front(source);
dist[source] = 0;
while (!q.empty()) {
    int v = q.front(); q.pop_front();
    for (auto [w, u] : adj[v]) {
        if (dist[v] + w < dist[u]) {
            dist[u] = dist[v] + w;
            if (w == 0) q.push_front(u);
            else q.push_back(u);
        }
    }
}
\end{verbatim}

\paragraph{Complejidad.}
\begin{itemize}\setlength\itemsep{0pt}
	\item $O(V + E)$.
	\item Cada nodo se inserta a lo sumo $O(1)$ veces (con check de distancia).
\end{itemize}

\paragraph{Donde tocar para modificar.}
\begin{itemize}\setlength\itemsep{0pt}
	\item \textbf{Pesos en $[0, k]$ small}: usar small k-ary BFS o Dijkstra.
	\item \textbf{Reconstruir camino}: aniadir \texttt{parent[]}.
	\item \textbf{Pesos negativos}: no funciona, usar Bellman-Ford.
\end{itemize}

\paragraph{Trampas y errores frecuentes.}
\begin{itemize}\setlength\itemsep{0pt}
	\item Si los pesos son otro rango, el algoritmo da respuestas incorrectas. Verificar que las aristas son solo 0/1.
	\item El chequeo \texttt{dist[v] + w < dist[u]} es crucial; sino, el algoritmo procesa nodos multiples veces.
	\item Si una arista tiene peso $\ge 2$, aniadir verificaciones o cambiar a Dijkstra.
\end{itemize}

\paragraph{Codigo.}
Ver \texttt{cpp.tex}, seccion \textit{0-1 BFS}.

\subsection{MST}

\subsubsection{Kruskal}

\paragraph{Idea intuitiva.}
Encuentra el \textbf{arbol de expansion minima} (MST) en $O(E \log E)$. Ordena las aristas por peso ascendente y las aniade al MST si no forman ciclo (chequear con DSU). El arbol tiene exactamente $V - 1$ aristas. La demostracion (``Cut Property''): para cualquier corte, la arista minima cruzando el corte pertenece a \emph{algun} MST; por lo tanto, elegir la minima disponible que no forma ciclo es optimo.

\paragraph{Propiedades clave (invariantes).}
\begin{itemize}\setlength\itemsep{0pt}
	\item Algoritmo \textbf{greedy} correcto (Cut Property).
	\item Funciona en grafos \textbf{no conexos}; produce un \textbf{minimum spanning forest}.
	\item Si las aristas tienen pesos unicos, el MST es unico.
	\item El MST tiene exactamente $V - 1$ aristas (o menos si el grafo es disconexo).
\end{itemize}

\paragraph{Codigo clave.}
Kruskal con DSU:
\begin{verbatim}
sort(edges.begin(), edges.end());
DSU dsu(n);
long long mst = 0;
int edgesUsed = 0;
for (auto [w, u, v] : edges)
    if (dsu.unite(u, v)) {
        mst += w;
        if (++edgesUsed == n - 1) break;
    }
\end{verbatim}

\paragraph{Complejidad.}
\begin{itemize}\setlength\itemsep{0pt}
	\item $O(E \log E)$ (ordenamiento) o $O(E \log V)$ con DSU optimizations.
	\item DSU es casi $O(1)$ amortizado, asi que el sort domina.
\end{itemize}

\paragraph{Donde tocar para modificar.}
\begin{itemize}\setlength\itemsep{0pt}
	\item \textbf{Second best MST}: enumerar aristas del MST, sacar una, aniadir una no-MST; recalcular.
	\item \textbf{Kruskal randomizado}: si el grafo es casi-completo, agregar shuffle.
	\item \textbf{Aniadir info al DSU}: guardar el peso maximo en el camino, etc.
	\item \textbf{DSU con rollback}: si necesitas Kruskal reversible (para queries offline).
\end{itemize}

\paragraph{Trampas y errores frecuentes.}
\begin{itemize}\setlength\itemsep{0pt}
	\item El DSU debe estar limpio para cada invocacion.
	\item Si las aristas tienen pesos negativos, Kruskal sigue funcionando.
	\item Si el grafo es disconexo, el ``MST'' es un bosque con varias componentes.
	\item El check \texttt{if (++edgesUsed == n - 1) break;} optimiza la salida temprana.
\end{itemize}

\paragraph{Codigo.}
Ver \texttt{cpp.tex}, seccion \textit{Kruskal}.

\vspace{0.5em}

\subsubsection{Prim (priority queue)}

\paragraph{Idea intuitiva.}
Otra forma de calcular el MST: empezar desde un nodo y, en cada paso, aniadir la arista de menor peso que conecta el conjunto construido con el resto. Usando priority queue, es $O((V + E) \log V)$. La demostracion: cada vez que se aniade una arista, es la minima cruzando el ``corte'' entre el conjunto actual y el resto (Cut Property); eso preserva optimalidad.

\paragraph{Propiedades clave (invariantes).}
\begin{itemize}\setlength\itemsep{0pt}
	\item Similar a Dijkstra pero sin acumular distancias, solo eligiendo minimo por arista.
	\item Funciona mejor que Kruskal en \textbf{grafos densos} ($E \approx V^2$).
	\item Si el grafo es disconexo, genera el bosque (solo llega a los alcanzables).
	\item Cada nodo se inserta al heap una vez; cada arista se procesa una vez.
\end{itemize}

\paragraph{Codigo clave.}
Prim clasico con heap:
\begin{verbatim}
priority_queue<pair<int, int>, vector<...>, greater<...>> pq;
vector<bool> used(n, false);
vector<int> minEdge(n, INF);
minEdge[0] = 0;
pq.push({0, 0});
while (!pq.empty()) {
    auto [w, v] = pq.top(); pq.pop();
    if (used[v]) continue;
    used[v] = true;
    mst += w;
    for (auto [peso, u] : adj[v])
        if (!used[u] && peso < minEdge[u]) {
            minEdge[u] = peso;
            pq.push({peso, u});
        }
}
\end{verbatim}

\paragraph{Complejidad.}
\begin{itemize}\setlength\itemsep{0pt}
	\item $O((V + E) \log V)$.
	\item Con matriz de adyacencia: $O(V^2)$ sin heap.
\end{itemize}

\paragraph{Donde tocar para modificar.}
\begin{itemize}\setlength\itemsep{0pt}
	\item \textbf{Prim con matriz de adyacencia}: $O(V^2)$ sin heap, util en grafos densos pequenos.
	\item \textbf{Reconstruir el MST}: aniadir \texttt{parent[]}.
	\item \textbf{Prim en streaming}: variante para grafos que llegan incrementalmente.
\end{itemize}

\paragraph{Trampas y errores frecuentes.}
\begin{itemize}\setlength\itemsep{0pt}
	\item El \texttt{if (used[v]) continue;} es crucial: si el nodo ya fue procesado, saltar.
	\item En la version matriz, elije el minimo por linea (no heap).
	\item Si el grafo es disconexo, aniadir un loop externo sobre todos los nodos no usados.
\end{itemize}

\paragraph{Codigo.}
Ver \texttt{cpp.tex}, seccion \textit{Prim (priority queue)}.

\subsection{Flujos}

\subsubsection{Edmonds-Karp (max flow)}

\paragraph{Idea intuitiva.}
Implementacion BFS del algoritmo de \textbf{Ford-Fulkerson}. En cada iteracion, hacer BFS desde $s$ hasta $t$ siguiendo aristas con capacidad residual $> 0$; si se llega a $t$, se encontro un \textbf{camino de aumento}. Enviar el minimo de las capacidades en el camino; actualizar capacidades residuales. La demostracion de $O(VE^2)$: cada BFS encuentra el augmenting path mas corto, y la longitud minima crece monotona; como hay $\le E$ longitudes posibles, hay $\le E$ fases.

\paragraph{Propiedades clave (invariantes).}
\begin{itemize}\setlength\itemsep{0pt}
	\item BFS garantiza caminos de aumento mas cortos.
	\item Capacidad residual: forward edge baja, backward edge sube.
	\item Complejidad $O(VE^2)$.
	\item El flujo maximo respeta la conservacion en nodos (in = out salvo $s, t$).
\end{itemize}

\paragraph{Codigo clave.}
BFS + augmenting path:
\begin{verbatim}
while (bfs(s, t, parent)) {  // bfs encuentra camino o retorna false
    int pathFlow = INF;
    for (int v = t; v != s; v = parent[v])
        pathFlow = min(pathFlow, capacity[parent[v]][v]);
    for (int v = t; v != s; v = parent[v]) {
        capacity[parent[v]][v] -= pathFlow;
        capacity[v][parent[v]] += pathFlow;
    }
    maxFlow += pathFlow;
}
\end{verbatim}

\paragraph{Complejidad.}
\begin{itemize}\setlength\itemsep{0pt}
	\item $O(VE^2)$.
	\item En la practica, mas lento que Dinic para grafos grandes.
\end{itemize}

\paragraph{Donde tocar para modificar.}
\begin{itemize}\setlength\itemsep{0pt}
	\item \textbf{Grafo bipartito}: $O(\sqrt{V} E)$ (equivalente a Hopcroft-Karp).
	\item \textbf{Dinic} es preferible en la mayoria de los casos.
	\item \textbf{Capacidades doubles}: usar tolerancia para comparacion.
	\item \textbf{Min-cost flow}: aniadir costo por arista y modificar el algoritmo.
\end{itemize}

\paragraph{Trampas y errores frecuentes.}
\begin{itemize}\setlength\itemsep{0pt}
	\item Las capacidades deben ser $> 0$ para entrar a la cola (no $>= 0$); sino, ciclos infinitos.
	\item Para grafos bipartitos, Dinic o Hopcroft-Karp son mejores.
	\item Si el flujo maximo es muy grande, los \texttt{int} pueden overflow; usar \texttt{long long}.
\end{itemize}

\paragraph{Codigo.}
Ver \texttt{cpp.tex}, seccion \textit{Edmonds-Karp (max flow)}.

\vspace{0.5em}

\subsubsection{Dinic (max flow)}

\paragraph{Idea intuitiva.}
Algoritmo de flujo mas eficiente que Edmonds-Karp: combina \textbf{BFS para level graph} con \textbf{DFS bloqueante} (multiple augmenting paths en una sola pasada). Cada fase cuesta $O(E)$ y hay $O(V)$ fases en el peor caso; para bipartitos, $O(\sqrt{V} E)$. La razon: cada fase ``satura'' al menos una arista, asi que en $O(E)$ fases el flujo maximo se alcanza; con $O(V)$ niveles en el BFS, hay $O(VE)$ operaciones totales.

\paragraph{Propiedades clave (invariantes).}
\begin{itemize}\setlength\itemsep{0pt}
	\item \texttt{level[v]} = distancia desde $s$ en el level graph (solo aristas con cap $> 0$).
	\item \texttt{it[v]} = iterador actual para evitar reprocesar aristas.
	\item DFS solo sigue aristas con \texttt{level[e.to] == level[v] + 1}.
	\item Back-edges tienen \texttt{rev} apuntando a la forward edge original.
\end{itemize}

\paragraph{Codigo clave.}
El DFS bloqueante de Dinic:
\begin{verbatim}
long long dfs(int v, int t, long long f) {
    if (v == t) return f;
    for (int& i = it[v]; i < adj[v].size(); ++i) {
        Edge& e = adj[v][i];
        if (e.cap > 0 && level[e.to] == level[v] + 1) {
            long long ret = dfs(e.to, t, min(f, e.cap));
            if (ret > 0) {
                e.cap -= ret;
                adj[e.to][e.rev].cap += ret;
                return ret;
            }
        }
    }
    return 0;
}
\end{verbatim}

\paragraph{Complejidad.}
\begin{itemize}\setlength\itemsep{0pt}
	\item $O(V^2 E)$ general, $O(\sqrt{V} E)$ para bipartitos.
	\item En la practica, mucho mas rapido que Edmonds-Karp.
\end{itemize}

\paragraph{Donde tocar para modificar.}
\begin{itemize}\setlength\itemsep{0pt}
	\item \textbf{Dinic con scaling}: aniadir factor de escala para capacidades grandes.
	\item \textbf{Lower bounds}: aniadir un super-source/sink para forzar flujo minimo.
	\item \textbf{Matching bipartito}: el grafo se construye como source -> L -> R -> sink, todos con capacidad 1.
	\item \textbf{Push-relabel}: alternativa para grafos muy densos.
\end{itemize}

\paragraph{Trampas y errores frecuentes.}
\begin{itemize}\setlength\itemsep{0pt}
	\item \texttt{it[v]} debe \textbf{incrementarse} en cada iteracion; sino, el DFS se queda atascado en la misma arista.
	\item Si el BFS no encuentra $t$, el flujo es maximo; terminar.
	\item Las back-edges son cruciales para ``cancelar'' elecciones previas.
\end{itemize}

\paragraph{Codigo.}
Ver \texttt{cpp.tex}, seccion \textit{Dinic (max flow)}.

\vspace{0.5em}

\subsubsection{Min-Cost Max-Flow}

\paragraph{Idea intuitiva.}
Encuentra el flujo maximo con \textbf{costo minimo} (asumiendo costos no negativos). Usa \textbf{SPFA} (o Dijkstra con potentials) para encontrar el shortest path en el \textbf{costo} desde $s$ hasta $t$ considerando capacidades residuales. Cada augmenting path aporta su costo $\times$ flujo. La idea: los potentials reescalan las aristas para que todas tengan costo no negativo, permitiendo usar Dijkstra.

\paragraph{Propiedades clave (invariantes).}
\begin{itemize}\setlength\itemsep{0pt}
	\item \texttt{pot[]} = potentials para evitar aristas negativas (Johnson's trick).
	\item \texttt{dist[v]} = costo minimo desde $s$ con potentials.
	\item Solo se aumentan aristas con \texttt{cap > 0} (forward) y costo es \texttt{+cost} (forward) o \texttt{-cost} (backward).
	\item Tras cada augmenting, los potentials se actualizan para mantener la consistencia.
\end{itemize}

\paragraph{Codigo clave.}
El ciclo de min-cost flow:
\begin{verbatim}
while (spfa(s, t, dist, parent)) {
    int pathFlow = INF;
    for (int v = t; v != s; v = parent[v])
        pathFlow = min(pathFlow, capacity[parent[v]][v]);
    for (int v = t; v != s; v = parent[v]) {
        capacity[parent[v]][v] -= pathFlow;
        capacity[v][parent[v]] += pathFlow;
        cost += pathFlow * edge_cost[parent[v]][v];
    }
}
\end{verbatim}

\paragraph{Complejidad.}
\begin{itemize}\setlength\itemsep{0pt}
	\item $O(F \cdot V E)$ con SPFA, $O(F \cdot E \log V)$ con Dijkstra + potentials.
	\item $F$ = flujo maximo total.
\end{itemize}

\paragraph{Donde tocar para modificar.}
\begin{itemize}\setlength\itemsep{0pt}
	\item \textbf{Dijkstra + potentials}: si los costos son enteros no negativos, esto es mas estable.
	\item \textbf{Costos negativos en aristas}: requieren Bellman-Ford inicial o reorden.
	\item \textbf{Flujo minimo con costo}: encontrar el minimo flujo posible y su costo.
	\item \textbf{Costo por unidad}: si los costos son muy grandes, aniadir verificacion.
\end{itemize}

\paragraph{Trampas y errores frecuentes.}
\begin{itemize}\setlength\itemsep{0pt}
	\item El \texttt{pot[]} se actualiza \textbf{al final} de cada iteracion (no durante).
	\item Las back-edges tienen \textbf{costo negativo}, esencial para corregir elecciones previas.
	\item Si los costos son muy grandes, considerar modular arithmetic.
	\item Para problemas con flujo exacto, aniadir verificacion post-algoritmo.
\end{itemize}

\paragraph{Codigo.}
Ver \texttt{cpp.tex}, seccion \textit{Min-Cost Max-Flow}.

\subsection{Especales}

\subsubsection{2-SAT (implication graph + SCC)}

\paragraph{Idea intuitiva.}
Resuelve formulas booleanas en \textbf{CNF} con clausulas de \textbf{2 literales}: $(l_1 \lor l_2)$. Cada variable $x$ se representa con dos nodos: $x$ (verdadero) y $\neg x$ (falso). Cada clausula $l_1 \lor l_2$ se descompone en dos implicaciones: $\neg l_1 \to l_2$ y $\neg l_2 \to l_1$. Si en el grafo de implicaciones $\neg x$ y $x$ estan en la misma SCC, la formula es \textbf{UNSAT}. La demostracion: si $\neg x \to x$ en el grafo, entonces $\neg x$ implica $x$, asi que $\neg x$ debe ser falso y $x$ verdadero; pero entonces $\neg x$ es falso implica una contradiccion.

\paragraph{Propiedades clave (invariantes).}
\begin{itemize}\setlength\itemsep{0pt}
	\item Cada literal es un nodo; total $2n$ nodos.
	\item Kosaraju/Tarjan dan las SCCs.
	\item Si SAT: asignar $x = $ (topological order del condensation).
	\item Clausulas \texttt{setTrue(x)} se modelan como \texttt{implies(!x, x)}.
	\item La formula es SAT sii no hay variable $x$ con $x$ y $\neg x$ en la misma SCC.
\end{itemize}

\paragraph{Codigo clave.}
Conversion de clausulas a implicaciones:
\begin{verbatim}
// Clausula (a OR b) -> implica ~a -> b Y ~b -> a
addImplication(not(a), b);
addImplication(not(b), a);
// Forzar x: implica ~x -> x
if (force) addImplication(not(x), x);
\end{verbatim}

\paragraph{Complejidad.}
\begin{itemize}\setlength\itemsep{0pt}
	\item $O(V + E) = O(N + M)$ con Kosaraju o Tarjan.
	\item Memoria $O(V + E)$.
\end{itemize}

\paragraph{Donde tocar para modificar.}
\begin{itemize}\setlength\itemsep{0pt}
	\item \textbf{Forzar variables}: aniadir \texttt{setTrue} / \texttt{setFalse}.
	\item \textbf{3-SAT}: NP-completo; no usar este algoritmo.
	\item \textbf{Contar soluciones}: las SCCs condensation dan una formula para contar; multiplicar las formas de asignar variables en cada componente no-trivial.
	\item \textbf{Output de la asignacion}: aniadir el vector \texttt{assignment[]}.
\end{itemize}

\paragraph{Trampas y errores frecuentes.}
\begin{itemize}\setlength\itemsep{0pt}
	\item La asignacion usa el \textbf{orden topologico inverso} del condensation: \texttt{val[i] = comp[2i] < comp[2i+1]}.
	\item El ID de $\neg x$ es \texttt{2*x + 1} (o similar); verificar consistencia.
	\item Si hay clausulas con literales repetidos (e.g., $(x \lor x)$), tratar como $(x)$ simple.
\end{itemize}

\paragraph{Codigo.}
Ver \texttt{cpp.tex}, seccion \textit{2-SAT (implication graph + SCC)}.

\vspace{0.5em}

\subsubsection{Euler tour / Hierholzer (camino y circuito)}

\paragraph{Idea intuitiva.}
Encuentra un \textbf{camino o circuito euleriano} (que usa cada arista exactamente una vez). \textbf{Condiciones}:
\begin{itemize}\setlength\itemsep{0pt}
	\item \textbf{No dirigido}: todos los grados pares (circuito) o exactamente 2 impares (camino, empieza en uno, termina en el otro). Ademas, conexo sobre el subgrafo con aristas.
	\item \textbf{Dirigido}: $\text{indeg}(v) = \text{outdeg}(v)$ para todos (circuito) o exactamente un nodo con $\text{outdeg} = \text{indeg} + 1$ y otro con $\text{indeg} = \text{outdeg} + 1$ (camino).
\end{itemize}

\paragraph{Hierholzer}: elegir un ciclo desde el inicio, mientras haya ciclos ``laterales'' desde un nodo del camino, mergearlos. Implementacion con pila: avanzar por aristas, al no poder mas aniadir a la respuesta y backtrack. La demostracion de $O(V + E)$: cada arista se procesa una vez (entra y sale del stack una vez).

\paragraph{Propiedades clave (invariantes).}
\begin{itemize}\setlength\itemsep{0pt}
	\item Multigrafos permitidos; las aristas se marcan como usadas.
	\item No dirigido: cada arista aparece \textbf{dos veces} en \texttt{adj} (ida y vuelta).
	\item La construccion del path con ids consistentes es importante.
	\item El resultado es unico salvo orden de las aristas en multigrafos.
\end{itemize}

\paragraph{Codigo clave.}
Hierholzer con pila:
\begin{verbatim}
vector<int> path;
stack<int> st;
st.push(start);
while (!st.empty()) {
    int v = st.top();
    while (!adj[v].empty() && used[adj[v].back()]) adj[v].pop_back();
    if (adj[v].empty()) {
        path.push_back(v);
        st.pop();
    } else {
        int e = adj[v].back();
        used[e] = true;
        int u = edges[e].to(v);
        adj[v].pop_back();
        st.push(u);
    }
}
reverse(path.begin(), path.end());
\end{verbatim}

\paragraph{Complejidad.}
\begin{itemize}\setlength\itemsep{0pt}
	\item $O(V + E)$.
	\item Memoria $O(V + E)$.
\end{itemize}

\paragraph{Donde tocar para modificar.}
\begin{itemize}\setlength\itemsep{0pt}
	\item \textbf{Reconstruir con multigrafo}: aniadir ids consistentes antes del algoritmo.
	\item \textbf{Aplicar a ``de Bruijn'' o problemas de strings}: modelar como multigrafo y buscar Euler.
	\item \textbf{Dirigido vs no dirigido}: el snippet tiene ambas versiones.
	\item \textbf{Letter addition}: modelar palabras como aristas y buscar el superstring.
\end{itemize}

\paragraph{Trampas y errores frecuentes.}
\begin{itemize}\setlength\itemsep{0pt}
	\item Si el grafo no es conexo (sobre las aristas), no existe euleriano. Verificar primero.
	\item Para multigrafos, los IDs de aristas son cruciales para distinguirlas.
	\item Si el problema es ``encuentra todos los euler tours'', el algoritmo es exponencial.
\end{itemize}

\paragraph{Codigo.}
Ver \texttt{cpp.tex}, seccion \textit{Euler tour / Hierholzer (camino y circuito)}.

\vspace{0.5em}

\subsubsection{Hamiltonian path / cycle (bitmask DP + backtracking)}

\paragraph{Idea intuitiva.}
Un \textbf{camino hamiltoniano} visita cada vertice \textbf{exactamente una vez}. \textbf{Determinar} si existe es NP-completo en general, pero con $V \le 20$ se puede hacer DP con bitmask en $O(V^2 \cdot 2^V)$. \textbf{Teoremas suficientes}: \textbf{Dirac} (grado $\ge n/2$) y \textbf{Ore} ($\text{deg}(u) + \text{deg}(v) \ge n$ para todo par no adyacente). La idea de la DP: cada estado es (conjunto de visitados, ultimo vertice); la transicion es agregar un vecino no visitado.

\paragraph{Propiedades clave (invariantes).}
\begin{itemize}\setlength\itemsep{0pt}
	\item DP: $dp[mask][v]$ = true si hay un camino que visita exactamente \texttt{mask} y termina en $v$.
	\item TSP: variante con pesos minimos (Hamilton cycle minimo).
	\item Reconstruccion: guardar \texttt{par[mask][v]} (predecesor).
	\item $2^V$ estados con $V$ bits cada uno; factible para $V \le 20$.
\end{itemize}

\paragraph{Codigo clave.}
DP con bitmask:
\begin{verbatim}
dp[1][0] = 0;  // empezar en 0
for (int mask = 1; mask < (1<<n); ++mask)
    for (int v = 0; v < n; ++v) if (dp[mask][v] < INF)
        for (int u = 0; u < n; ++u)
            if (!(mask & (1<<u)))
                dp[mask | (1<<u)][u] = min(dp[mask | (1<<u)][u],
                                            dp[mask][v] + w[v][u]);
// TSP: respuesta = min(dp[(1<<n)-1][v] + w[v][0])
\end{verbatim}

\paragraph{Complejidad.}
\begin{itemize}\setlength\itemsep{0pt}
	\item $O(V^2 \cdot 2^V)$ tiempo, $O(V \cdot 2^V)$ memoria.
	\item Para $V = 20$: $\sim 4 \cdot 10^8$ operaciones (lento pero factible).
\end{itemize}

\paragraph{Donde tocar para modificar.}
\begin{itemize}\setlength\itemsep{0pt}
	\item \textbf{TSP asimetrico}: aniadir pesos $w[v][u]$ (no necesariamente simetricos).
	\item \textbf{Hamilton con inicio obligatorio}: fijar \texttt{src} en el DP inicial.
	\item \textbf{Bitmask DP mas compacto}: usar \texttt{long long} como mascara si $V \le 64$.
	\item \textbf{TSP optimo aproximado}: MST, 2-opt, simulated annealing.
	\item \textbf{Reconstruir el path}: usar \texttt{par[mask][v]} y backtrack.
\end{itemize}

\paragraph{Trampas y errores frecuentes.}
\begin{itemize}\setlength\itemsep{0pt}
	\item TSP asume grafo \textbf{completo} (con \texttt{LINF} para aristas faltantes); si no, aniadir chequeo.
	\item La condicion ``\texttt{mask == (1<<n) - 1}'' cierra el ciclo retornando al origen.
	\item Para $V > 20$, la memoria $2^V$ se vuelve prohibitiva; usar meet-in-the-middle o heuristicas.
\end{itemize}

\paragraph{Codigo.}
Ver \texttt{cpp.tex}, seccion \textit{Hamiltonian path / cycle (bitmask DP + backtracking)}.

% ============================================================
\section{DP}

\subsubsection{Knapsack 0/1}

\paragraph{Idea intuitiva.}
Dado un conjunto de $N$ objetos con peso $w_i$ y valor $v_i$, maximizar $\sum v_i$ con la restriccion $\sum w_i \le W$. La recurrencia clasica:
\[ dp[i][j] = \max(dp[i-1][j],\, dp[i-1][j - w_i] + v_i) \quad \text{si } j \ge w_i. \]
La optimizacion clave: la DP es 1-dimensional si iteramos $j$ \textbf{de mayor a menor} (asi no usamos el objeto $i$ dos veces). La razon de la iteracion inversa: en cada iteracion queremos el estado del ``item anterior''; iterar de menor a mayor contaminaria \texttt{dp[j - w_i]} con el item actual.

\paragraph{Propiedades clave (invariantes).}
\begin{itemize}\setlength\itemsep{0pt}
	\item $dp[j]$ al final = maximo valor con peso $\le j$.
	\item El bucle inverso \texttt{for (j = W; j >= w[i]; j--)} es lo que distingue 0/1 de unbounded.
	\item Solucion: maximo valor con cualquier peso $\le W$.
	\item La DP es exacta; no hay aproximacion.
\end{itemize}

\paragraph{Codigo clave.}
La version 1D optimizada:
\begin{verbatim}
vector<long long> dp(W + 1, 0);
for (int i = 0; i < N; ++i)
    for (int j = W; j >= w[i]; --j)  // invertido
        dp[j] = max(dp[j], dp[j - w[i]] + v[i]);
return dp[W];
\end{verbatim}

\paragraph{Complejidad.}
\begin{itemize}\setlength\itemsep{0pt}
	\item $O(N \cdot W)$ tiempo, $O(W)$ memoria.
	\item Constante muy pequena (solo sumas y max).
\end{itemize}

\paragraph{Donde tocar para modificar.}
\begin{itemize}\setlength\itemsep{0pt}
	\item \textbf{Reconstruir la seleccion}: aniadir \texttt{take[i][j]} (booleano o bool) en una version 2D.
	\item \textbf{Items con peso 0}: edge case; aniadir guard.
	\item \textbf{Knapsack con restricciones de grupos}: ``elegir a lo sumo 1 de cada grupo'' se modela como 0/1 con un peso extra.
	\item \textbf{Tamano de $W$}: si $W \le 10^9$ pero $N$ es chico, usar meet-in-the-middle.
	\item \textbf{Cambiar max por min}: aniadir \texttt{INF} en la inicializacion.
\end{itemize}

\paragraph{Trampas y errores frecuentes.}
\begin{itemize}\setlength\itemsep{0pt}
	\item Overflow en \texttt{dp[j - w[i]] + v[i]}; usar \texttt{long long} si los valores son grandes.
	\item El bucle \textbf{debe} ir de mayor a menor.
	\item Si los pesos son grandes y $N$ pequeno, conviene meet-in-the-middle.
	\item Para Knapsack exacto ($\sum w_i = W$), aniadir check al final.
\end{itemize}

\paragraph{Codigo.}
Ver \texttt{cpp.tex}, seccion \textit{Knapsack 0/1}.

\vspace{0.5em}

\subsubsection{Knapsack unbounded}

\paragraph{Idea intuitiva.}
Igual que Knapsack 0/1 pero \textbf{cada item se puede tomar multiples veces}. La recurrencia es similar:
\[ dp[j] = \max(dp[j],\, dp[j - w_i] + v_i), \]
pero ahora el bucle va de \textbf{menor a mayor} (porque tomar el item $i$ no afecta corridas futuras del mismo item). La idea: tras tomar el item $i$, queremos volver a tener la opcion de tomarlo; iterar de menor a mayor permite que el item ``se acumule''.

\paragraph{Propiedades clave (invariantes).}
\begin{itemize}\setlength\itemsep{0pt}
	\item El bucle \texttt{for (j = w[i]; j <= W; j++)} es lo que distingue unbounded.
	\item A diferencia de 0/1, no hay riesgo de ``usar dos veces el mismo item''.
	\item La DP es exacta; cada ``subconjunto'' se cuenta correctamente.
\end{itemize}

\paragraph{Codigo clave.}
La version 1D con bucle directo:
\begin{verbatim}
vector<long long> dp(W + 1, 0);
for (int i = 0; i < N; ++i)
    for (int j = w[i]; j <= W; ++j)  // directo (no invertido)
        dp[j] = max(dp[j], dp[j - w[i]] + v[i]);
return dp[W];
\end{verbatim}

\paragraph{Complejidad.}
\begin{itemize}\setlength\itemsep{0pt}
	\item $O(N \cdot W)$ tiempo, $O(W)$ memoria.
	\item Constante muy pequena.
\end{itemize}

\paragraph{Donde tocar para modificar.}
\begin{itemize}\setlength\itemsep{0pt}
	\item \textbf{Reconstruir cuantos de cada item}: aniadir \texttt{cnt[i]} y backtrack.
	\item \textbf{Cambiar moneda minima por maxima}: el snippet maximiza valor; para minimizar monedas, cambiar a \texttt{min} con inicializacion en \texttt{INF}.
	\item \textbf{Coin change}: si todos los pesos son positivos y queremos ``minimo numero de monedas'', es unbounded con min.
\end{itemize}

\paragraph{Trampas y errores frecuentes.}
\begin{itemize}\setlength\itemsep{0pt}
	\item Misma que Knapsack 0/1 con overflow.
	\item El orden de iteracion de items no afecta el resultado (convexo).
	\item Si se quiere ``maximo valor con exactamente $k$ items'', aniadir segunda dimension.
\end{itemize}

\paragraph{Codigo.}
Ver \texttt{cpp.tex}, seccion \textit{Knapsack unbounded}.

\vspace{0.5em}

\subsubsection{Knapsack meet-in-the-middle}

\paragraph{Idea intuitiva.}
Para $N \le 40$ y $W$ moderado, no se puede hacer $O(NW)$ directamente. Partir en dos mitades, enumerar todos los $2^{N/2}$ subconjuntos de cada mitad, y combinar. Para cada subconjunto de la segunda mitad con peso $sw_2$, encontrar el maximo valor de la primera mitad con peso $\le W - sw_2$. La idea: ``divide y venceras'' sobre el tamano de la entrada; cada mitad tiene mitad del espacio, asi que la combinacion es geometrica en lugar de multiplicativa.

\paragraph{Propiedades clave (invariantes).}
\begin{itemize}\setlength\itemsep{0pt}
	\item Hay $2^{N/2}$ subconjuntos por mitad; total $O(2^{N/2})$.
	\item \texttt{sums} guarda los valores de subconjuntos de la primera mitad ordenados (o preprocesados).
	\item Para cada subconjunto de la segunda mitad, busqueda lineal o binaria en \texttt{sums}.
	\item La mitad mas grande (con $N/2 + 1$ items para $N$ impar) tiene $2^{N/2+1}$ subconjuntos.
\end{itemize}

\paragraph{Codigo clave.}
Enumerar subconjuntos de la primera mitad:
\begin{verbatim}
int n1 = N / 2, n2 = N - n1;
vector<pair<long long, long long>> sums1;
for (int mask = 0; mask < (1 << n1); ++mask) {
    long long peso = 0, valor = 0;
    for (int i = 0; i < n1; ++i) if (mask & (1 << i)) {
        peso += w[i]; valor += v[i];
    }
    sums1.push_back({peso, valor});
}
sort(sums1.begin(), sums1.end());  // por peso
// eliminar duplicados dominados
\end{verbatim}

\paragraph{Complejidad.}
\begin{itemize}\setlength\itemsep{0pt}
	\item $O(2^{N/2})$ tiempo y memoria (mejor que $O(NW)$ cuando $W > 2^{N/2}$).
	\item Para $N = 40$: $\sim 2 \cdot 10^6$ subconjuntos, factible.
\end{itemize}

\paragraph{Donde tocar para modificar.}
\begin{itemize}\setlength\itemsep{0pt}
	\item \textbf{Tamano mayor}: si $N$ es mayor pero $W$ es chico, conviene iterar por peso, no por subconjuntos.
	\item \textbf{Varias mochilas}: adaptar para multiples $W$.
	\item \textbf{Optimizar el ``combinado''}: usar \texttt{upper\_bound} para encontrar el maximo valor con peso $\le \text{rem}.
	\item \textbf{Eliminar pares dominados}: si un par (peso, valor) tiene otro con menos peso y mas valor, eliminarlo.
\end{itemize}

\paragraph{Trampas y errores frecuentes.}
\begin{itemize}\setlength\itemsep{0pt}
	\item $N$ impar: una mitad tiene $N/2 + 1$ items.
	\item Overflow al sumar pesos/valores en \texttt{long long}.
	\item Sin deduplicacion, el algoritmo es correcto pero mas lento.
\end{itemize}

\paragraph{Codigo.}
Ver \texttt{cpp.tex}, seccion \textit{Knapsack meet-in-the-middle}.

\vspace{0.5em}

\subsubsection{LIS O(n log n)}

\paragraph{Idea intuitiva.}
La \textbf{Longest Increasing Subsequence} se calcula en $O(n \log n)$ con el truco del \textbf{patience sorting}: mantener un arreglo \texttt{dp} donde \texttt{dp[len]} es el minimo ultimo valor de una IS de largo \texttt{len}. Para cada $x$, encontrar el primer \texttt{dp[len] >= x} y reemplazarlo. El largo final es la LIS. La demostracion: si hay dos IS del mismo largo, la que termina en un valor menor es ``mejor'' (puede extenderse mas facilmente).

\paragraph{Propiedades clave (invariantes).}
\begin{itemize}\setlength\itemsep{0pt}
	\item \texttt{dp} es estrictamente creciente.
	\item \texttt{lower\_bound} da LIS estricta; \texttt{upper\_bound} da no-estricta ($\le$).
	\item El algoritmo es \textbf{online}: procesa elementos en orden.
	\item La longitud de \texttt{dp} es la LIS al final.
\end{itemize}

\paragraph{Codigo clave.}
La version online:
\begin{verbatim}
vector<int> dp;
for (int x : a) {
    auto it = lower_bound(dp.begin(), dp.end(), x);
    if (it == dp.end()) dp.push_back(x);
    else *it = x;
}
// dp.size() = longitud de LIS estricta
\end{verbatim}

\paragraph{Complejidad.}
\begin{itemize}\setlength\itemsep{0pt}
	\item $O(n \log n)$.
	\item Memoria $O(n)$.
\end{itemize}

\paragraph{Donde tocar para modificar.}
\begin{itemize}\setlength\itemsep{0pt}
	\item \textbf{Reconstruir la LIS}: aniadir \texttt{parent[]} y \texttt{posInDp[]} para backtrack.
	\item \textbf{Contar LIS}: DP clasica $O(n^2)$ o segment tree sobre $dp$.
	\item \textbf{LDS / Longest non-increasing}: invertir la condicion en \texttt{lower\_bound} o multiplicar por -1.
	\item \textbf{Longest bitonic subsequence}: combinar LIS desde izq y desde der.
	\item \textbf{LIS en 2D}: ordenar por una dimension y aplicar LIS en la otra.
\end{itemize}

\paragraph{Trampas y errores frecuentes.}
\begin{itemize}\setlength\itemsep{0pt}
	\item \texttt{lower\_bound} vs \texttt{upper\_bound}: si tu comparacion es $\le$ vs $<$, cambia.
	\item Si los elementos son \texttt{int} pero podrian ser negativos, aniadir offset para usar como indices.
	\item Para LIS reconstruida, aniadir el array \texttt{parent} que apunta al anterior en la IS.
\end{itemize}

\paragraph{Codigo.}
Ver \texttt{cpp.tex}, seccion \textit{LIS O(n log n)}.

\vspace{0.5em}

\subsubsection{LCS / Edit distance}

\paragraph{Idea intuitiva.}
\textbf{LCS} (Longest Common Subsequence): para dos strings $a$, $b$, $dp[i][j]$ = LCS de $a[0..i)$ y $b[0..j)$. Recurrencia:
\[ dp[i][j] = \begin{cases} dp[i-1][j-1] + 1 & \text{si } a[i-1] = b[j-1] \\ \max(dp[i-1][j], dp[i][j-1]) & \text{si no} \end{cases} \]

\textbf{Edit distance} (Levenshtein): numero minimo de operaciones (insert, delete, replace) para convertir $a$ en $b$. Recurrencia similar con un $+1$ adicional. La idea: ``alinear'' los strings eligiendo matches/mismatches/insert/delete.

\paragraph{Propiedades clave (invariantes).}
\begin{itemize}\setlength\itemsep{0pt}
	\item LCS y Edit Distance son DP $O(nm)$.
	\item Memoria se puede reducir a $O(\min(n, m))$ con rolling array.
	\item Edit distance con peso por operacion distinto: ajustar el \texttt{+1} por el peso.
	\item La subestructura optima: la respuesta para prefijos se construye de prefijos menores.
\end{itemize}

\paragraph{Codigo clave.}
La recurrencia clasica de LCS:
\begin{verbatim}
vector<vector<int>> dp(n + 1, vector<int>(m + 1, 0));
for (int i = 1; i <= n; ++i)
    for (int j = 1; j <= m; ++j)
        if (a[i-1] == b[j-1]) dp[i][j] = dp[i-1][j-1] + 1;
        else dp[i][j] = max(dp[i-1][j], dp[i][j-1]);
return dp[n][m];
\end{verbatim}

\paragraph{Complejidad.}
\begin{itemize}\setlength\itemsep{0pt}
	\item $O(NM)$ tiempo y memoria.
	\item Para $N, M \le 5000$: $\sim 2.5 \cdot 10^7$ operaciones.
\end{itemize}

\paragraph{Donde tocar para modificar.}
\begin{itemize}\setlength\itemsep{0pt}
	\item \textbf{Reconstruir la LCS / las operaciones}: aniadir \texttt{prev[i][j]} para backtrack.
	\item \textbf{LCS de varios strings}: NP-hard para $\ge 3$ strings; en la practica se hace con bitmask.
	\item \textbf{Edit distance con costos distintos}: cambiar el \texttt{+1} por el costo apropiado.
	\item \textbf{Rolling array}: si memoria es problema, mantener solo 2 filas.
	\item \textbf{Hirschberg}: algoritmo $O(NM)$ tiempo, $O(\min(N,M))$ memoria, con reconstruccion.
\end{itemize}

\paragraph{Trampas y errores frecuentes.}
\begin{itemize}\setlength\itemsep{0pt}
	\item Indices en la recurrencia: \texttt{a[i-1]} y \texttt{b[j-1]} (no \texttt{a[i]}).
	\item Para reconstruccion, el caso \texttt{a[i-1] == b[j-1]} tiene prioridad.
	\item Si los strings son muy largos, considerar Hirschberg o bitset optimization.
\end{itemize}

\paragraph{Codigo.}
Ver \texttt{cpp.tex}, seccion \textit{LCS / Edit distance}.

\vspace{0.5em}

\subsubsection{Bitmask DP over subsets}

\paragraph{Idea intuitiva.}
Cuando el estado es un \textbf{subconjunto} de $\{0, \dots, N-1\}$, se representa como una mascara de $N$ bits. Para iterar sobre todos los subconjuntos de \texttt{mask}:
\[ \text{for (sub = mask; sub; sub = (sub - 1) \& mask)} \{ \dots \} \]
Para iterar sobre todos los superconjuntos de \texttt{mask} en $[0, U)$:
\[ \text{for (sup = mask; sup < U; sup = (sup + 1) | mask)} \{ \dots \} \]
La idea: las mascaras forman un lattice de subconjuntos; los trucos aritmeticos (\texttt{(sub-1) \& mask} y \texttt{(sup+1) | mask}) enumeran todo el lattice eficientemente.

\paragraph{Propiedades clave (invariantes).}
\begin{itemize}\setlength\itemsep{0pt}
	\item Iterar subsets: $O(2^k)$ para una mascara con $k$ bits.
	\item \texttt{sub \& mask == sub} siempre; \texttt{mask - sub} no siempre da un subconjunto.
	\item \texttt{\_\_builtin\_popcount(mask)} cuenta los bits.
	\item \texttt{(sub - 1) \& mask} itera subsets de \texttt{mask} en orden decreciente de inclusion.
\end{itemize}

\paragraph{Codigo clave.}
Iteracion sobre subsets de una mascara:
\begin{verbatim}
// Todos los subsets de mask
for (int sub = mask; sub; sub = (sub - 1) & mask) {
    // procesar sub
}
// Incluir sub = 0 si es necesario
sub = 0;  // procesar primero
for (int sup = mask; sup < (1 << N); sup = (sup + 1) | mask) {
    // procesar sup (todos los superconjuntos)
}
\end{verbatim}

\paragraph{Complejidad.}
\begin{itemize}\setlength\itemsep{0pt}
	\item $O(N \cdot 2^N)$ para una DP estandar sobre todos los subsets.
	\item Iterar subsets de \texttt{mask}: $O(2^{|mask|})$.
\end{itemize}

\paragraph{Donde tocar para modificar.}
\begin{itemize}\setlength\itemsep{0pt}
	\item \textbf{TSP / Hamiltonian path}: bitmask DP (ver tambien \textit{Hamiltonian path} en \texttt{cpp.tex} seccion Grafos).
	\item \textbf{Set cover}: iterar subsets.
	\item \textbf{DP con supermask}: si el estado incluye ``todos los elementos que \textbf{no} hemos usado'', usar supermask iteration.
	\item \textbf{Bitmasks grandes}: con $N > 20$, aniadir compresion o simulacion con segmentos.
\end{itemize}

\paragraph{Trampas y errores frecuentes.}
\begin{itemize}\setlength\itemsep{0pt}
	\item El orden de iteracion de subsets de \texttt{mask} es decreciente (de \texttt{mask} a $0$).
	\item \texttt{mask = 0} no tiene subsets no triviales; el bucle lo saltea correctamente.
	\item El truco \texttt{(sub - 1) \& mask} requiere cuidado con \texttt{sub = 0} (termina el bucle).
	\item Superconjuntos: \texttt{sup < (1 << N)} debe estar bien definido.
\end{itemize}

\paragraph{Codigo.}
Ver \texttt{cpp.tex}, seccion \textit{Bitmask DP over subsets}.

\vspace{0.5em}

\subsubsection{SOS DP}

\paragraph{Idea intuitiva.}
Para calcular, para cada \texttt{mask}, la suma de $f(\text{submask})$ sobre todos los $\text{submask} \subseteq \text{mask}$:
\[ dp[mask] = \sum_{\text{sub} \subseteq mask} f(\text{sub}). \]
Se computa en $O(N \cdot 2^N)$ mediante un loop por bit: para cada $i$, sumar a $dp[mask]$ la contribucion de $dp[mask \oplus (1 << i)]$ cuando $mask$ tiene el bit $i$. La idea: cada subconjunto $s \subseteq m$ se obtiene quitando bits uno a uno; asi, sumar incrementalmente captura todos los subconjuntos en $O(N \cdot 2^N)$.

\paragraph{Propiedades clave (invariantes).}
\begin{itemize}\setlength\itemsep{0pt}
	\item Inicializar $dp = f$.
	\item Para cada bit $i$: \texttt{if (mask \& (1 << i)) dp[mask] += dp[mask \^{} (1 << i)]}.
	\item Variantes: supersets DP (sumar sobre superconjuntos).
	\item El orden de iteracion (bit externo, mascara interno) es crucial.
\end{itemize}

\paragraph{Codigo clave.}
SOS DP (Sum over Subsets):
\begin{verbatim}
vector<long long> dp(1 << N);
for (int i = 0; i < (1 << N); ++i) dp[i] = f[i];
for (int bit = 0; bit < N; ++bit)
    for (int mask = 0; mask < (1 << N); ++mask)
        if (mask & (1 << bit))
            dp[mask] += dp[mask ^ (1 << bit)];
// dp[mask] = suma de f(sub) para sub ⊆ mask
\end{verbatim}

\paragraph{Complejidad.}
\begin{itemize}\setlength\itemsep{0pt}
	\item $O(N \cdot 2^N)$.
	\item Memoria $O(2^N)$.
\end{itemize}

\paragraph{Donde tocar para modificar.}
\begin{itemize}\setlength\itemsep{0pt}
	\item \textbf{Max/min en lugar de suma}: cambiar \texttt{+=} por la operacion correspondiente.
	\item \textbf{Supersets DP}: iterar \texttt{mask} de $0$ a $2^N$ y aniadir \texttt{dp[mask] += dp[mask | (1<<i)]}.
	\item \textbf{XOR convolution} (subset convolution): version mas avanzada.
	\item \textbf{Aplicaciones}: contar subconjuntos con propiedades especificas.
\end{itemize}

\paragraph{Trampas y errores frecuentes.}
\begin{itemize}\setlength\itemsep{0pt}
	\item Si $N$ es 30 o mas, $2^N$ no entra en memoria; usar trucos o reducir $N$.
	\item El sentido del XOR (sumar a \texttt{mask \^{} bit}) es para ``agregar el bit $i$ al subconjunto''.
	\item En supersets DP, la condicion es \texttt{(mask \& (1<<bit)) == 0}.
\end{itemize}

\paragraph{Codigo.}
Ver \texttt{cpp.tex}, seccion \textit{SOS DP}.

\vspace{0.5em}

\subsubsection{Digit DP template}

\paragraph{Idea intuitiva.}
Contar numeros en $[0, N]$ que cumplen una propiedad sobre sus digitos. La idea: DP por posicion, manteniendo:
\begin{itemize}\setlength\itemsep{0pt}
	\item \texttt{pos}: posicion actual (de izquierda a derecha).
	\item \texttt{tight}: si los digitos ya colocados son exactamente los de $N$ (limite superior estricto).
	\item \texttt{state}: estado adicional (digitos usados, suma, etc.).
\end{itemize}
La clave: la condicion \texttt{tight} permite acotar correctamente el siguiente digito; sin ella, la DP contaria incorrectamente.

\paragraph{Propiedades clave (invariantes).}
\begin{itemize}\setlength\itemsep{0pt}
	\item Si \texttt{tight} es true, el siguiente digito esta acotado por el de $N$; si no, puede ser 0-9.
	\item Memoria: $O(D \cdot 2 \cdot \text{states})$.
	\item Caso base: si \texttt{pos == D}, devolver 1 si el estado cumple la propiedad, sino 0.
	\item El bit \texttt{started} distingue ``numero con menos digitos'' (con leading zeros).
\end{itemize}

\paragraph{Codigo clave.}
La estructura basica:
\begin{verbatim}
long long dp(int pos, bool tight, State s) {
    if (pos == D) return check(s) ? 1 : 0;
    if (memo[pos][tight][s] != -1) return memo;
    int limit = tight ? digits[pos] : 9;
    long long ans = 0;
    for (int d = 0; d <= limit; ++d)
        ans += dp(pos + 1, tight && (d == limit), update(s, d, pos));
    return memo[pos][tight][s] = ans;
}
\end{verbatim}

\paragraph{Complejidad.}
\begin{itemize}\setlength\itemsep{0pt}
	\item $O(D \cdot 10 \cdot \text{states})$ donde $D$ es el numero de digitos.
	\item Para $D \le 19$ (hasta $10^{18}$): $\sim 200 \cdot \text{states}$.
\end{itemize}

\paragraph{Donde tocar para modificar.}
\begin{itemize}\setlength\itemsep{0pt}
	\item \textbf{Cambiar la propiedad}: modificar \texttt{updateState} y la condicion base.
	\item \textbf{Multiples estados}: aniadir mas dimensiones al \texttt{memo}.
	\item \textbf{Leading zeros}: aniadir un flag \texttt{started} para distinguir numeros con menos digitos.
	\item \textbf{Rango $[L, R]$}: calcular $f(R) - f(L-1)$.
\end{itemize}

\paragraph{Trampas y errores frecuentes.}
\begin{itemize}\setlength\itemsep{0pt}
	\item Olvidar resetear \texttt{memo} entre queries.
	\item Si la cantidad de estados es grande, el \texttt{memo[pos][tight][state]} puede ser muy grande; usar \texttt{map}.
	\item La condicion \texttt{started} es crucial para numeros con menos digitos que $D$.
	\item Si $N$ es muy grande ($> 10^{18}$), aniadir soporte para long long.
\end{itemize}

\paragraph{Codigo.}
Ver \texttt{cpp.tex}, seccion \textit{Digit DP template}.

\vspace{0.5em}

\subsubsection{Tree DP (reroot)}

\paragraph{Idea intuitiva.}
Calcular alguna propiedad asumiendo que \textbf{cada nodo es la raiz} del arbol. La tecnica clasica de \textbf{rerooting} en dos pasadas:
\begin{itemize}\setlength\itemsep{0pt}
	\item \textbf{DFS 1}: desde una raiz arbitraria, calcular \texttt{dp[v]} para el subarbol de $v$.
	\item \textbf{DFS 2}: para cada hijo $u$ de $v$, calcular \texttt{up[u]} = respuesta si la raiz fuera $u$.
\end{itemize}
La demostracion: tras la primera DFS, cada nodo conoce ``lo que aporta su subarbol''; en la segunda, se transmite ``lo que aporta el resto del arbol''.

\paragraph{Propiedades clave (invariantes).}
\begin{itemize}\setlength\itemsep{0pt}
	\item \texttt{dp[v]} depende solo del subarbol de $v$ (con la raiz original).
	\item \texttt{up[u]} se calcula de manera que cuando la raiz es $u$, ``todo lo que estaba en $v$ se ajusta''.
	\item Formula para reroot depende del problema; el snippet usa el ejemplo ``suma de distancias''.
	\item La combinacion \texttt{dp[u] + up[u]} da la respuesta con raiz $u$.
\end{itemize}

\paragraph{Codigo clave.}
Reroot para suma de distancias:
\begin{verbatim}
// dfs1: dp[v] = suma de distancias desde v al subarbol
void dfs1(int v, int p) {
    dp[v] = 0;
    for (int u : adj[v]) if (u != p) {
        dfs1(u, v);
        dp[v] += dp[u] + sz[u];  // cada hijo aporta su subarbol
    }
}
// dfs2: up[u] = distancia desde el resto del arbol a u
void dfs2(int v, int p) {
    for (int u : adj[v]) if (u != p) {
        up[u] = up[v] + dp[v] - dp[u] - sz[u] + (n - sz[u]);
        dfs2(u, v);
    }
}
\end{verbatim}

\paragraph{Complejidad.}
\begin{itemize}\setlength\itemsep{0pt}
	\item $O(N)$ para ambas pasadas.
	\item Cada arista se procesa un numero constante de veces.
\end{itemize}

\paragraph{Donde tocar para modificar.}
\begin{itemize}\setlength\itemsep{0pt}
	\item \textbf{Cambiar la cantidad a calcular}: modificar la operacion en \texttt{dfs1} y la formula en \texttt{dfs2}.
	\item \textbf{Arbol con pesos}: aniadir \texttt{weights[v]} y ajustar.
	\item \textbf{Re-root en arboles con ciclos}: requiere otra tecnica (DSU on tree).
	\item \textbf{Rooted trees diferentes}: si la propiedad no es invariante por raiz, usar otra tecnica.
\end{itemize}

\paragraph{Trampas y errores frecuentes.}
\begin{itemize}\setlength\itemsep{0pt}
	\item La formula de \texttt{up[u]} debe sumar y restar correctamente las contribuciones de $u$ y del resto.
	\item Para arboles con pesos, aniadir el peso en las formulas.
	\item Verificar con ejemplos pequenos antes de generalizar.
\end{itemize}

\paragraph{Codigo.}
Ver \texttt{cpp.tex}, seccion \textit{Tree DP (reroot)}.

% ============================================================
\section{Geometria}
% ============================================================

\subsection{Puntos y segmentos}

\subsubsection{Point + cross + dot}

\paragraph{Idea intuitiva.}
La estructura \texttt{Point} representa un punto en 2D con coordenadas enteras. Las dos operaciones vectoriales clave son:
\begin{itemize}\setlength\itemsep{0pt}
	\item \textbf{Cross product} $P \times Q = P_x Q_y - P_y Q_x$. Da el ``area con signo'' del paralelogramo formado por $P$ y $Q$. Usado para orientacion.
	\item \textbf{Dot product} $P \cdot Q = P_x Q_x + P_y Q_y$. Da el producto de las longitudes por el coseno del angulo. Usado para proyeccion.
\end{itemize}
La interpretacion geometrica del cross product: es el area con signo del paralelogramo; el signo indica la orientacion relativa.

\paragraph{Propiedades clave (invariantes).}
\begin{itemize}\setlength\itemsep{0pt}
	\item \texttt{cross(a, b) > 0} si $b$ esta a la izquierda de $a$ (sistema de coordenadas usual).
	\item \texttt{cross(a, b) = 0} si son colineales.
	\item \texttt{dot(a, b) > 0} si el angulo es agudo.
	\item El snippet define \texttt{cross} respecto a \texttt{this}, util para orientacion de 3 puntos.
	\item $|cross|^2 + |dot|^2 = |P|^2 \cdot |Q|^2$ (identidad de Lagrange).
\end{itemize}

\paragraph{Codigo clave.}
Las dos operaciones vectoriales:
\begin{verbatim}
long long cross(const Point& a, const Point& b) {
    return (long long)a.x * b.y - (long long)a.y * b.x;
}
long long dot(const Point& a, const Point& b) {
    return (long long)a.x * b.x + (long long)a.y * b.y;
}
\end{verbatim}

\paragraph{Complejidad.}
\begin{itemize}\setlength\itemsep{0pt}
	\item $O(1)$ por operacion.
	\item Constante muy pequena (4 multiplicaciones + 2 sumas).
\end{itemize}

\paragraph{Donde tocar para modificar.}
\begin{itemize}\setlength\itemsep{0pt}
	\item \textbf{Coordenadas doubles}: cambiar \texttt{int} por \texttt{double} o \texttt{long double}. Cuidado con la precision.
	\item \textbf{3D}: aniadir \texttt{z} y adaptar cross/dot.
	\item \textbf{Operadores adicionales}: aniadir \texttt{operator+}, \texttt{operator-}, \texttt{operator*} (escalar).
	\item \textbf{Distancia al cuadrado}: aniadir \texttt{len2()} que retorna $x^2 + y^2$.
\end{itemize}

\paragraph{Trampas y errores frecuentes.}
\begin{itemize}\setlength\itemsep{0pt}
	\item Cross product \textbf{no es} conmutativo; $P \times Q = -Q \times P$.
	\item Overflow con coordenadas grandes ($|x|, |y| > 10^4$) si se hace \texttt{x * y}.
	\item El signo de cross depende del sistema de coordenadas (en computacion, usualmente CCW positivo).
	\item Para ``igualdad'' de doubles, aniadir tolerancia \texttt{eps}.
\end{itemize}

\paragraph{Codigo.}
Ver \texttt{cpp.tex}, seccion \textit{Point + cross + dot}.

\vspace{0.5em}

\subsubsection{Interseccion de segmentos}

\paragraph{Idea intuitiva.}
Dados dos segmentos $[a, b]$ y $[c, d]$, determinar si se intersectan. La idea clasica:
\begin{itemize}\setlength\itemsep{0pt}
	\item Verificar orientaciones: $a, b, c$ y $a, b, d$ deben estar en lados opuestos de la recta $ab$; igual para $cd$ y $ab$.
	\item Caso colineal: si las rectas son paralelas (\texttt{cross} $= 0$) y los segmentos son colineales, verificar que los bounding boxes se intersectan.
\end{itemize}
La demostracion: dos segmentos se intersectan sii cada segmento ``cruza'' la linea que contiene al otro (estricto) o ambos son colineales con solapamiento.

\paragraph{Propiedades clave (invariantes).}
\begin{itemize}\setlength\itemsep{0pt}
	\item Funciona con segmentos en cualquier orientacion.
	\item Caso colineal: requiere check adicional de bounding box.
	\item Para intersectar en un \textbf{punto} especifico: aniadir division con cuidado de precision.
	\item Funciona con segmentos verticales/horizontales sin casos especiales.
\end{itemize}

\paragraph{Codigo clave.}
Test clasico de interseccion:
\begin{verbatim}
int sign(long long x) { return (x > 0) - (x < 0); }
bool intersect(Point a, Point b, Point c, Point d) {
    long long o1 = cross(b - a, c - a);
    long long o2 = cross(b - a, d - a);
    long long o3 = cross(d - c, a - c);
    long long o4 = cross(d - c, b - c);
    if (sign(o1) != sign(o2) && sign(o3) != sign(o4)) return true;
    // casos colineales
    if (o1 == 0 && onSegment(a, b, c)) return true;
    if (o2 == 0 && onSegment(a, b, d)) return true;
    if (o3 == 0 && onSegment(c, d, a)) return true;
    if (o4 == 0 && onSegment(c, d, b)) return true;
    return false;
}
\end{verbatim}

\paragraph{Complejidad.}
\begin{itemize}\setlength\itemsep{0pt}
	\item $O(1)$.
	\item Constante pequena.
\end{itemize}

\paragraph{Donde tocar para modificar.}
\begin{itemize}\setlength\itemsep{0pt}
	\item \textbf{Interseccion con punto}: aniadir la division despues de verificar orientacion.
	\item \textbf{Interseccion rayo-segmento}: ajustar la primera condicion.
	\item \textbf{Poligonos convexos}: usar separacion de ejes (SAT).
	\item \textbf{Multiples queries}: precomputar bounding boxes o usar segment tree.
\end{itemize}

\paragraph{Trampas y errores frecuentes.}
\begin{itemize}\setlength\itemsep{0pt}
	\item Si los segmentos tocan solo en un extremo, el snippet lo considera interseccion (algunos problemas quieren estricto). Ajustar segun el problema.
	\item El caso colineal es comun; olvidar el check adicional de bounding box da falsos negativos.
	\item En problemas con muchos segmentos, considerar sweep line para $O(n \log n)$.
\end{itemize}

\paragraph{Codigo.}
Ver \texttt{cpp.tex}, seccion \textit{Interseccion de segmentos}.

\vspace{0.5em}

\subsubsection{Contains colineal}

\paragraph{Idea intuitiva.}
Verifica si un punto $c$ esta sobre el segmento $[a, b]$ (asumiendo que ya son colineales). Si \texttt{cross(a, b, c) = 0} y $c$ esta dentro del bounding box de $[a, b]$, entonces $c$ esta en el segmento. La razon: si los tres puntos son colineales, $c$ esta en el segmento sii sus coordenadas estan entre las de $a$ y $b$.

\paragraph{Propiedades clave (invariantes).}
\begin{itemize}\setlength\itemsep{0pt}
	\item Asume colinealidad (no la verifica mas alla de \texttt{cross == 0}).
	\item Usa comparaciones con \texttt{<=} para extremos \textbf{inclusivos}.
	\item Funciona con segmentos verticales (la condicion se reduce a una sola coordenada).
\end{itemize}

\paragraph{Codigo clave.}
Test de contencion en segmento (asume colinealidad):
\begin{verbatim}
bool onSegment(Point a, Point b, Point c) {
    return min(a.x, b.x) <= c.x && c.x <= max(a.x, b.x) &&
           min(a.y, b.y) <= c.y && c.y <= max(a.y, b.y);
}
\end{verbatim}

\paragraph{Complejidad.}
\begin{itemize}\setlength\itemsep{0pt}
	\item $O(1)$.
	\item Constante minima.
\end{itemize}

\paragraph{Donde tocar para modificar.}
\begin{itemize}\setlength\itemsep{0pt}
	\item \textbf{Contener estricto}: cambiar \texttt{<=} a \texttt{<} para excluir extremos.
	\item \textbf{Coordenadas doubles}: usar tolerancia (\texttt{eps}).
	\item \textbf{Solo x o solo y}: aniadir check especifico para segmentos verticales/horizontales.
\end{itemize}

\paragraph{Trampas y errores frecuentes.}
\begin{itemize}\setlength\itemsep{0pt}
	\item No usar este snippet para puntos \textbf{no colineales}; siempre verificar primero la colinealidad.
	\item Si los puntos tienen doubles, el chequeo \texttt{cross == 0} puede fallar; aniadir tolerancia.
	\item El bounding box incluye los extremos; verificar si el problema quiere exclusivo o inclusivo.
\end{itemize}

\paragraph{Codigo.}
Ver \texttt{cpp.tex}, seccion \textit{Contains colineal}.

\subsection{Poligonos}

\subsubsection{Convex Hull (monotono)}

\paragraph{Idea intuitiva.}
Algoritmo de \textbf{Andrew's monotone chain}: ordenar puntos por $(x, y)$, luego construir la hull inferior y superior por separado. Para cada punto, aniadirlo al hull y, mientras el ultimo giro no sea ``counter-clockwise'' (\texttt{cross > 0}), popear. La demostracion: si los puntos se procesan en orden, el hull inferior y superior cubren la frontera completa sin auto-intersecciones.

\paragraph{Propiedades clave (invariantes).}
\begin{itemize}\setlength\itemsep{0pt}
	\item Ordena los puntos en $O(n \log n)$.
	\item Construye la hull en $O(n)$ despues.
	\item Output: poligono convexo en counter-clockwise, sin punto final duplicado.
	\item El \texttt{<= 0} incluye colineales; usar \texttt{< 0} para excluirlos.
	\item Cada punto aparece en la hull inferior o superior (no ambos).
\end{itemize}

\paragraph{Codigo clave.}
Construccion del hull inferior y superior:
\begin{verbatim}
sort(p.begin(), p.end());
vector<Point> lower, upper;
for (auto& p : points) {
    while (lower.size() >= 2 && cross(lower.back() - lower[lower.size()-2],
                                       p - lower.back()) <= 0)
        lower.pop_back();
    lower.push_back(p);
}
for (int i = points.size() - 1; i >= 0; --i) {
    auto& p = points[i];
    while (upper.size() >= 2 && cross(upper.back() - upper[upper.size()-2],
                                       p - upper.back()) <= 0)
        upper.pop_back();
    upper.push_back(p);
}
vector<Point> hull = lower;
hull.insert(hull.end(), upper.begin() + 1, upper.end() - 1);
\end{verbatim}

\paragraph{Complejidad.}
\begin{itemize}\setlength\itemsep{0pt}
	\item $O(n \log n)$.
	\item Constante muy pequena (solo cross products).
\end{itemize}

\paragraph{Donde tocar para modificar.}
\begin{itemize}\setlength\itemsep{0pt}
	\item \textbf{Colineales}: cambiar \texttt{<= 0} por \texttt{< 0} para excluirlos del hull.
	\item \textbf{Puntos repetidos}: \texttt{unique} despues de sort.
	\item \textbf{Hull 3D}: gift wrapping es $O(n^2)$; convex hull 3D es mas complejo.
	\item \textbf{Output con los puntos}: aniadir el hull completo (no solo upper/lower).
\end{itemize}

\paragraph{Trampas y errores frecuentes.}
\begin{itemize}\setlength\itemsep{0pt}
	\item El \texttt{pop\_back} y el \texttt{reverse} son esenciales para no duplicar el ultimo punto.
	\item Para excluir colineales, ajustar el \texttt{<= 0} vs \texttt{< 0}.
	\item Si los puntos estan duplicados, aniadir \texttt{unique} antes del sort.
\end{itemize}

\paragraph{Codigo.}
Ver \texttt{cpp.tex}, seccion \textit{Convex Hull (monotono)}.

\vspace{0.5em}

\subsubsection{Polygon Area}

\paragraph{Idea intuitiva.}
El area (con signo) de un poligono se calcula con la \textbf{Shoelace formula}: $A = \frac{1}{2} \sum_{i=0}^{n-1} (x_i y_{i+1} - x_{i+1} y_i)$, donde el indice es modulo $n$. El snippet retorna el \textbf{doble} del area (en valor absoluto). La formula sale del ``sumar areas de trapezoides'': cada lado del poligono define un trapezoide con el eje $x$; la suma con signo da el area.

\paragraph{Propiedades clave (invariantes).}
\begin{itemize}\setlength\itemsep{0pt}
	\item El signo indica la orientacion: positivo si counter-clockwise.
	\item Para el doble: no dividir por 2 si no es necesario (evita fracciones).
	\item Funciona para poligonos simples (no auto-intersectantes).
	\item Si el poligono es no convexo, la formula da el area ``con signo''.
\end{itemize}

\paragraph{Codigo clave.}
La suma de trapezoides:
\begin{verbatim}
long long shoelace(const vector<Point>& p) {
    long long s = 0;
    int n = p.size();
    for (int i = 0; i < n; ++i)
        s += (long long)p[i].x * p[(i+1)%n].y - (long long)p[(i+1)%n].x * p[i].y;
    return abs(s);  // doble del area
}
\end{verbatim}

\paragraph{Complejidad.}
\begin{itemize}\setlength\itemsep{0pt}
	\item $O(n)$.
	\item Constante minima.
\end{itemize}

\paragraph{Donde tocar para modificar.}
\begin{itemize}\setlength\itemsep{0pt}
	\item \textbf{Area exacta}: dividir por 2 (con cuidado si da fraccion; usar \texttt{double}).
	\item \textbf{Area firmada}: quitar \texttt{abs} y dejar el signo.
	\item \textbf{Punto dentro de poligono convexo}: con la hull, chequear que esta a un mismo lado de todas las aristas.
	\item \textbf{Triangulacion}: descomponer en triangulos y sumar areas.
\end{itemize}

\paragraph{Trampas y errores frecuentes.}
\begin{itemize}\setlength\itemsep{0pt}
	\item Si los puntos no estan en orden (CCW o CW), el area firmada puede ser incorrecta. Para signed area consistente, usar siempre el mismo orden.
	\item Para poligonos con holes, sumar las areas con signo apropiado.
	\item Con doubles, aniadir tolerancia \texttt{eps}.
\end{itemize}

\paragraph{Codigo.}
Ver \texttt{cpp.tex}, seccion \textit{Polygon Area}.

\vspace{0.5em}

\subsubsection{Distance point-segment / point-line}

\paragraph{Idea intuitiva.}
Distancia minima de un punto $p$ a un segmento $[a, b]$ o a la recta que pasa por $a$ y $b$.
\begin{itemize}\setlength\itemsep{0pt}
	\item \textbf{Linea}: $d = \frac{|ab \times ap|}{|ab|}$.
	\item \textbf{Segmento}: si la proyeccion de $p$ sobre $ab$ cae fuera del segmento, la distancia es al extremo mas cercano.
\end{itemize}
La razon geometrica: la distancia a una linea es el area del paralelogramo dividido por la base; la proyeccion ortogonal indica si estamos ``dentro'' o ``fuera'' del segmento.

\paragraph{Propiedades clave (invariantes).}
\begin{itemize}\setlength\itemsep{0pt}
	\item \texttt{dist2PointLine} retorna distancia \textbf{al cuadrado}.
	\item \texttt{dist2PointSegment} distingue 3 casos: $p$ antes de $a$, despues de $b$, o proyectado sobre el segmento.
	\item Si \texttt{len2 = 0} (degenerate), la distancia es $|ap|$ o $|bp|$.
	\item La proyeccion se calcula con dot products y el parametro $t \in [0, 1]$.
\end{itemize}

\paragraph{Codigo clave.}
Distancia al cuadrado a segmento:
\begin{verbatim}
long long dist2Segment(Point p, Point a, Point b) {
    long long len2 = (a-b).len2();
    if (len2 == 0) return (p-a).len2();  // degenerate
    long long t = max(0LL, min(1LL, dot(p-a, b-a) / len2));
    Point proj = {a.x + t*(b.x - a.x), a.y + t*(b.y - a.y)};
    return (p - proj).len2();
}
\end{verbatim}

\paragraph{Complejidad.}
\begin{itemize}\setlength\itemsep{0pt}
	\item $O(1)$.
	\item Constante minima.
\end{itemize}

\paragraph{Donde tocar para modificar.}
\begin{itemize}\setlength\itemsep{0pt}
	\item \textbf{Distancia (no al cuadrado)}: aniadir \texttt{sqrt} al final.
	\item \textbf{Coordenadas doubles}: aniadir \texttt{long double} y tolerancia.
	\item \textbf{3D}: aniadir coordenada $z$ y adaptar.
	\item \textbf{Multiples queries}: usar KD-tree o segment tree para acelerar.
\end{itemize}

\paragraph{Trampas y errores frecuentes.}
\begin{itemize}\setlength\itemsep{0pt}
	\item Overflow en \texttt{area2 * area2} si el area es grande; \texttt{long long} alcanza para $|x|, |y| \le 10^9$ pero no mas.
	\item Si $a = b$ (segmento degenerado), aniadir check explicito.
	\item En doubles, aniadir tolerancia \texttt{eps} para comparar.
\end{itemize}

\paragraph{Codigo.}
Ver \texttt{cpp.tex}, seccion \textit{Distance point-segment / point-line}.

\vspace{0.5em}

\subsubsection{Projection of point onto line}

\paragraph{Idea intuitiva.}
Proyeccion ortogonal de $p$ sobre la recta que pasa por $a$ y $b$. El parametro $t$ tal que la proyeccion es $a + t \cdot (b - a)$ es:
\[ t = \frac{(p - a) \cdot (b - a)}{|b - a|^2}. \]
La idea: descomponemos $\vec{ap}$ en una componente paralela a $\vec{ab}$ y una perpendicular; $t$ mide la posicion en la primera.

\paragraph{Propiedades clave (invariantes).}
\begin{itemize}\setlength\itemsep{0pt}
	\item Si $t \in [0, 1]$, la proyeccion cae sobre el \textbf{segmento}; si no, cae sobre la extension.
	\item Retorna \texttt{double} (puede tener precision issues).
	\item $t < 0$: $p$ esta ``antes'' de $a$; $t > 1$: $p$ esta ``despues'' de $b$.
\end{itemize}

\paragraph{Codigo clave.}
Proyeccion ortogonal:
\begin{verbatim}
Point projection(Point p, Point a, Point b) {
    Point ab = b - a, ap = p - a;
    double t = (double)dot(ap, ab) / ab.len2();
    return {a.x + t * ab.x, a.y + t * ab.y};
}
\end{verbatim}

\paragraph{Complejidad.}
\begin{itemize}\setlength\itemsep{0pt}
	\item $O(1)$.
	\item Constante minima (dos dot products y una division).
\end{itemize}

\paragraph{Donde tocar para modificar.}
\begin{itemize}\setlength\itemsep{0pt}
	\item \textbf{Proyeccion sobre segmento}: clipar $t$ a $[0, 1]$ antes de calcular el punto.
	\item \textbf{Distancia al segmento via proyeccion}: combinar con el modulo anterior.
	\item \textbf{Solo el parametro $t$}: si solo necesitas saber si esta dentro del segmento, retorna $t$ sin construir el punto.
\end{itemize}

\paragraph{Trampas y errores frecuentes.}
\begin{itemize}\setlength\itemsep{0pt}
	\item Si $a = b$ ($|ab| = 0$), la proyeccion no esta definida; aniadir guard.
	\item En doubles, aniadir tolerancia \texttt{eps} para comparar $t$ con 0 o 1.
	\item La division por \texttt{len2} puede ser cara si las coordenadas son grandes; usar \texttt{double} o precomputar.
\end{itemize}

\paragraph{Codigo.}
Ver \texttt{cpp.tex}, seccion \textit{Projection of point onto line}.

\vspace{0.5em}

\subsubsection{Point in polygon (ray casting)}

\paragraph{Idea intuitiva.}
Algoritmo de \textbf{ray casting}: contar cuantas veces un rayo horizontal desde $p$ cruza los bordes del poligono. Si es impar, $p$ esta dentro; si par, fuera. El snippet implementa la version clasica recorriendo aristas. La razon: cada cruce cambia el lado del rayo respecto al poligono; un numero impar de cambios significa ``terminamos dentro''.

\paragraph{Propiedades clave (invariantes).}
\begin{itemize}\setlength\itemsep{0pt}
	\item Funciona con poligonos simples (no necesariamente convexos).
	\item Las intersecciones en vertices cuentan especialmente.
	\item Para poligonos con \textbf{huecos}, sigue funcionando.
	\item La respuesta es estrictamente 0 (fuera) o 1 (dentro) para poligonos simples.
\end{itemize}

\paragraph{Codigo clave.}
El test clasico de ray casting:
\begin{verbatim}
bool pointInPolygon(Point p, const vector<Point>& poly) {
    int n = poly.size();
    bool inside = false;
    for (int i = 0, j = n - 1; i < n; j = i++) {
        if ((poly[i].y > p.y) != (poly[j].y > p.y) &&
            p.x < (poly[j].x - poly[i].x) * (p.y - poly[i].y) /
                    (poly[j].y - poly[i].y) + poly[i].x)
            inside = !inside;
    }
    return inside;
}
\end{verbatim}

\paragraph{Complejidad.}
\begin{itemize}\setlength\itemsep{0pt}
	\item $O(n)$ por query.
	\item Memoria $O(1)$ (no se almacena el rayo).
\end{itemize}

\paragraph{Donde tocar para modificar.}
\begin{itemize}\setlength\itemsep{0pt}
	\item \textbf{Winding number}: version alternativa que cuenta el numero de vueltas.
	\item \textbf{Poligono convexo}: check lineal contra las aristas (mas rapido).
	\item \textbf{Multiples queries}: precomputar o usar segment tree sobre las aristas.
	\item \textbf{Punto en frontera}: aniadir check adicional si es colineal con alguna arista.
\end{itemize}

\paragraph{Trampas y errores frecuentes.}
\begin{itemize}\setlength\itemsep{0pt}
	\item El manejo del vertice (cuando el rayo pasa por un vertice) puede causar double-counting; el snippet lo evita con la condicion \texttt{(poly[i].y > p.y) != (poly[j].y > p.y)}.
	\item Si el poligono tiene holes, el algoritmo sigue funcionando.
	\item Para poligonos no simples (auto-intersectantes), el resultado puede ser ambiguo.
\end{itemize}

\paragraph{Codigo.}
Ver \texttt{cpp.tex}, seccion \textit{Point in polygon (ray casting)}.

\vspace{0.5em}

\subsubsection{Minkowski sum (convex polygons)}

\paragraph{Idea intuitiva.}
La \textbf{suma de Minkowski} $A \oplus B$ de dos poligonos convexos es otro poligono convexo que representa $\{a + b : a \in A, b \in B\}$. Algoritmo: ordenar las aristas por angulo polar (CCW) y ``mergearlas'' como en merge sort; el resultado es un poligono CCW sin repetir el ultimo punto. La idea: el resultado tiene una arista paralela a cada arista de los originales; la suma se obtiene barriendo los angulos.

\paragraph{Propiedades clave (invariantes).}
\begin{itemize}\setlength\itemsep{0pt}
	\item Si $A, B$ son convexos con $n, m$ vertices, $A \oplus B$ tiene $\le n + m$ vertices.
	\item Util para \textbf{detectar colision}: $A$ colisiona con $B$ si $0 \in A \oplus (-B)$.
	\item Asume poligonos en orden \textbf{counter-clockwise}.
	\item El resultado es convexo (suma de convexos es convexo).
\end{itemize}

\paragraph{Codigo clave.}
La suma de Minkowski:
\begin{verbatim}
vector<Point> minkowski(const vector<Point>& A, const vector<Point>& B) {
    vector<Point> result;
    int i = 0, j = 0;
    int n = A.size(), m = B.size();
    do {
        result.push_back(A[i] + B[j]);
        long long cross = (A[(i+1)%n] - A[i]) ^ (B[(j+1)%m] - B[j]);
        if (cross >= 0) i = (i + 1) % n;
        if (cross <= 0) j = (j + 1) % m;
    } while (i || j);
    return result;
}
\end{verbatim}

\paragraph{Complejidad.}
\begin{itemize}\setlength\itemsep{0pt}
	\item $O(n + m)$.
	\item Memoria $O(n + m)$.
\end{itemize}

\paragraph{Donde tocar para modificar.}
\begin{itemize}\setlength\itemsep{0pt}
	\item \textbf{Deteccion de colision}: aniadir $(-B)$ (reflejar) y verificar si $0$ esta en el resultado.
	\item \textbf{Poligonos no convexos}: la suma puede no ser poligono simple; algoritmo distinto.
	\item \textbf{Multiples queries}: precomputar las aristas una vez.
	\item \textbf{Robot motion planning}: usar para planear trayectorias.
\end{itemize}

\paragraph{Trampas y errores frecuentes.}
\begin{itemize}\setlength\itemsep{0pt}
	\item Si los poligonos no son CCW, la suma es incorrecta. Normalizar antes.
	\item El bucle termina cuando ambos indices vuelven a 0; aniadir check.
	\item Para la deteccion de colision, aniadir el test \texttt{0 in A \^{} (-B)}.
\end{itemize}

\paragraph{Codigo.}
Ver \texttt{cpp.tex}, seccion \textit{Minkowski sum (convex polygons)}.

\subsection{Herramientas}

\subsubsection{Li Chao Tree (long long)}

\paragraph{Idea intuitiva.}
Mantiene un \textbf{lower envelope} (o upper) de un conjunto de rectas $y = mx + b$ sobre un dominio $[X_{LO}, X_{HI}]$ y responde ``minimo $y$ en un $x$ dado'' en $O(\log X)$. La estructura es un segment tree donde cada nodo guarda la recta ``mejor'' para su segmento medio; al aniadir una recta, se compara con la actual y se decide cual es mejor en cada mitad, descendiendo recursivamente. La demostracion: en cada nivel, solo se guarda la mejor recta para el ``punto medio''; las otras se transmiten a los hijos.

\paragraph{Propiedades clave (invariantes).}
\begin{itemize}\setlength\itemsep{0pt}
	\item \textbf{Online}: cada \texttt{addLine} en $O(\log X)$.
	\item Solo insercion; no hay borrar (para borrar, usar Li Chao segmentado con reset).
	\item La recta guardada en un nodo es la ``mejor'' en su punto medio.
	\item Las rectas ``perdedoras'' se transmiten a los hijos para futuros inserts.
\end{itemize}

\paragraph{Codigo clave.}
La insercion en Li Chao:
\begin{verbatim}
void addLine(Line nw, int v = 1, int l = X_LO, int r = X_HI) {
    int m = (l + r) / 2;
    bool left = nw.get(l) < tree[v].get(l);
    bool mid = nw.get(m) < tree[v].get(m);
    if (mid) swap(tree[v], nw);
    if (r - l == 1) return;
    if (left != mid) addLine(nw, v*2, l, m);
    else             addLine(nw, v*2+1, m, r);
}
\end{verbatim}

\paragraph{Complejidad.}
\begin{itemize}\setlength\itemsep{0pt}
	\item $O(\log X)$ por insercion y query, donde $X$ es el tamano del dominio.
	\item Memoria $O(\log X)$ por insercion, $O(X)$ total.
\end{itemize}

\paragraph{Donde tocar para modificar.}
\begin{itemize}\setlength\itemsep{0pt}
	\item \textbf{Upper envelope en vez de lower}: invertir el signo (\texttt{>} en vez de \texttt{<}).
	\item \textbf{Queries de segmento} (no full line): si la recta solo esta activa en un rango, aniadir bandera \texttt{active}.
	\item \textbf{Eliminar rectas}: usar Li Chao segmentado (mas complejo).
	\item \textbf{Pendientes grandes}: si las pendientes tienen magnitud $> 10^9$, aniadir \texttt{\_\_int128} para evitar overflow.
\end{itemize}

\paragraph{Trampas y errores frecuentes.}
\begin{itemize}\setlength\itemsep{0pt}
	\item El parametro \texttt{lo == hi} debe retornar el valor del nodo (caso base).
	\item La comparacion \texttt{newLine.get(mid) < node->line.get(mid)} decide quien es mejor en el medio.
	\item Si las rectas tienen $m$ iguales pero $b$ distintos, se queda con la primera (tie-break).
\end{itemize}

\paragraph{Codigo.}
Ver \texttt{cpp.tex}, seccion \textit{Li Chao Tree (long long)}.

\vspace{0.5em}

\subsubsection{Sweep line template}

\paragraph{Idea intuitiva.}
Patron de diseno para problemas geometricos donde se barre una linea sobre el plano y se mantiene una estructura de datos sobre el otro eje. Eventos (add, query, remove) se ordenan por la coordenada de barrido y se procesan secuencialmente.

\paragraph{Propiedades clave (invariantes).}
\begin{itemize}\setlength\itemsep{0pt}
	\item Eventos: tuplas \texttt{(x, type, id)}.
	\item La estructura auxiliar (set, multiset, segtree) se mantiene sobre el otro eje.
	\item Procesar eventos en orden de $x$ ascendente.
\end{itemize}

\paragraph{Complejidad.}
\begin{itemize}\setlength\itemsep{0pt}
	\item Depende del problema; tipico $O((N + E) \log N)$.
\end{itemize}

\paragraph{Donde tocar para modificar.}
\begin{itemize}\setlength\itemsep{0pt}
	\item \textbf{Aniadir/quitar segmentos}: aniadir eventos add/remove al barrido.
	\item \textbf{Aniadir queries}: aniadir eventos query.
	\item \textbf{Cambiar el eje de barrido}: si barres en $y$ en vez de $x$, aniadir comparador.
\end{itemize}

\paragraph{Trampas y errores frecuentes.}
\begin{itemize}\setlength\itemsep{0pt}
	\item El snippet es solo el \textbf{esqueleto}; los handlers \texttt{add/query/remove} se llenan por problema.
	\item Cuidado con eventos duplicados.
\end{itemize}

\paragraph{Codigo.}
Ver \texttt{cpp.tex}, seccion \textit{Sweep line template}.

\vspace{0.5em}

\subsubsection{Closest pair of points}

\paragraph{Idea intuitiva.}
Encuentra el par de puntos mas cercano en $O(n \log n)$ con \textbf{divide y venceras}. Algoritmo:
\begin{itemize}\setlength\itemsep{0pt}
	\item Sort por $x$.
	\item Dividir en mitades; recursivo.
	\item Combinar: dado el minimo $d$ de las dos mitades, strip = puntos en $[midX - d, midX + d]$; ordenar por $y$; cada punto solo se compara con los siguientes hasta $y$ difiera en $< d$ (maximo 7).
\end{itemize}

\paragraph{Propiedades clave (invariantes).}
\begin{itemize}\setlength\itemsep{0pt}
	\item Requiere ordenar por $y$ despues (en el snippet se hace durante el merge).
	\item Strip tiene a lo sumo $O(n)$ puntos en total, pero la comparacion es $O(1)$ amortizada por punto.
	\item Resultado en \texttt{double}.
\end{itemize}

\paragraph{Complejidad.}
\begin{itemize}\setlength\itemsep{0pt}
	\item $O(n \log n)$.
\end{itemize}

\paragraph{Donde tocar para modificar.}
\begin{itemize}\setlength\itemsep{0pt}
	\item \textbf{Reconstruir el par}: aniadir indices en la recursion.
	\item \textbf{Distancia maxima}: cambiar \texttt{min} por \texttt{max} (mas facil; no requiere D\&C).
	\item \textbf{Puntos duplicados}: aniadir check especial (distancia 0).
\end{itemize}

\paragraph{Trampas y errores frecuentes.}
\begin{itemize}\setlength\itemsep{0pt}
	\item El merge requiere ordenar por $y$; el snippet lo hace in-place con un \texttt{tmp}.
	\item Si se hace en cada recursion, es $O(n \log n)$ extra.
\end{itemize}

\paragraph{Codigo.}
Ver \texttt{cpp.tex}, seccion \textit{Closest pair of points}.

\vspace{0.5em}

\subsubsection{Rotating calipers (polygon diameter)}

\paragraph{Idea intuitiva.}
Encuentra el \textbf{diametro} de un poligono convexo (par de vertices mas distantes) en $O(n)$. Asume poligono en orden counter-clockwise. Algoritmo: dos indices $i$ (sobre las aristas) y $k$ (sobre los vertices opuestos); en cada paso, rotar $k$ mientras el area del paralelogramo aumenta. La distancia $i$-$k$ puede ser maxima en cualquier momento del barrido.

\paragraph{Propiedades clave (invariantes).}
\begin{itemize}\setlength\itemsep{0pt}
	\item Cada indice avanza como maximo $n$ veces total.
	\item El ``diametro'' es la maxima distancia entre dos vertices.
	\item Retorna \textbf{distancia al cuadrado}.
\end{itemize}

\paragraph{Complejidad.}
\begin{itemize}\setlength\itemsep{0pt}
	\item $O(n)$.
\end{itemize}

\paragraph{Donde tocar para modificar.}
\begin{itemize}\setlength\itemsep{0pt}
	\item \textbf{Reconstruir el par de vertices}: aniadir indices.
	\item \textbf{Anchura minima / maxima}: variantes de calipers.
	\item \textbf{Distancia minima entre poligonos convexos}: otra variante.
\end{itemize}

\paragraph{Trampas y errores frecuentes.}
\begin{itemize}\setlength\itemsep{0pt}
	\item Asume poligono \textbf{estrictamente} convexo (sin vertices colineales). Si hay colineales, aniadir perturbacion o filtrar.
\end{itemize}

\paragraph{Codigo.}
Ver \texttt{cpp.tex}, seccion \textit{Rotating calipers (polygon diameter)}.

% ============================================================
\section{Miscelanea}
% ============================================================

\subsection{Utilidades del usuario}

\subsubsection{Hora}

\paragraph{Idea intuitiva.}
Una pequena clase para representar \textbf{tiempos en segundos} y hacer aritmetica basica sobre ellos. Internamente almacena todo en segundos (un solo \texttt{int}), lo cual simplifica comparaciones y diferencias. El constructor acepta $(h, m, s)$ y los convierte a segundos. Hay un metodo \texttt{cut()} que fuerza un minimo (util para ``horas trabajadas'').

\paragraph{Propiedades clave (invariantes).}
\begin{itemize}\setlength\itemsep{0pt}
	\item Representacion canonica: \textbf{solo segundos}. No hay zonas horarias, no hay DST.
	\item Comparadores \texttt{<} y \texttt{>} operan sobre los segundos.
	\item \texttt{operator-} retorna la \textbf{diferencia absoluta} en segundos (siempre no negativa).
	\item \texttt{cut()} aplica un piso (por defecto $8\text{h }30\text{m} = 30600$ segundos).
\end{itemize}

\paragraph{Complejidad.}
\begin{itemize}\setlength\itemsep{0pt}
	\item $O(1)$ por operacion.
\end{itemize}

\paragraph{Donde tocar para modificar.}
\begin{itemize}\setlength\itemsep{0pt}
	\item \textbf{Aniadir suma}: \texttt{operator+} que devuelve un \texttt{Hour} con la suma de segundos.
	\item \textbf{Aniadir output formateado}: aniadir \texttt{show\_hms()} que imprime $H$:$M$:$S$.
	\item \textbf{Cambiar el ``corte''}: modificar el valor magico \texttt{30'600} en \texttt{cut()} o pasarlo como parametro.
	\item \textbf{Tamano maximo}: si los tiempos exceden $\sim 68$ anos, el \texttt{int} desborda; usar \texttt{long long}.
	\item \textbf{Precision sub-segundo}: aniadir \texttt{ms} (milisegundos) y guardar en \texttt{long long} con factor 1000.
\end{itemize}

\paragraph{Trampas y errores frecuentes.}
\begin{itemize}\setlength\itemsep{0pt}
	\item El operador \texttt{-} retorna \textbf{valor absoluto}, no diferencia con signo. Si necesitas diferencia con signo, aniadir un nuevo metodo o sobrecargar.
	\item \texttt{cut()} modifica \texttt{s} \textbf{in-place}; si el llamador no espera esto, aniadir una version \texttt{cut(min)} que retorna un nuevo \texttt{Hour}.
	\item El umbral magico \texttt{30'600} (8h 30m) esta hardcoded; documentar bien que es ``jornada minima''.
\end{itemize}

\paragraph{Codigo.}
Ver \texttt{cpp.tex}, seccion \textit{Hora}.

\end{document}
', aniadir un puntero a ambas raices.
\end{itemize}

\paragraph{Codigo.}
Ver \texttt{cpp.tex}, seccion \textit{Segment Tree persistente}.

\vspace{0.5em}

\subsubsection{Segment Tree iterativo}

\paragraph{Idea intuitiva.}
Representacion \textbf{sin recursion}: los nodos se guardan en un \texttt{vector} de tamano $2n$. Los hijos del nodo $i$ son $2i$ y $2i+1$. Las hojas ocupan las posiciones $n, n+1, \dots, 2n-1$. Update: subir el cambio desde la hoja hasta la raiz. Query: subir dos indices $l$ y $r$ hasta que se ``crucen''. La ventaja principal: localidad de memoria (los nodos estan en posiciones contiguas del vector), que mejora el cache miss rate y la constante practica. La desventaja: implementar lazy propagation es mas complejo.

\paragraph{Propiedades clave (invariantes).}
\begin{itemize}\setlength\itemsep{0pt}
	\item $O(\log n)$ por operacion, \textbf{mas rapido} en practica que el recursivo (mejor uso de cache, no hay overhead de llamadas).
	\item Build simple: copiar las hojas y combinar de abajo hacia arriba.
	\item Las hojas tienen indices en $[n, 2n)$; los nodos internos en $[1, n)$.
\end{itemize}

\paragraph{Codigo clave.}
El bucle de query ascendente:
\begin{verbatim}
long long query(int l, int r) {  // [l, r] inclusivo
    long long res = 0;
    for (l += n, r += n + 1; l < r; l >>= 1, r >>= 1) {
        if (l & 1) res += st[l++];
        if (r & 1) res += st[--r];
    }
    return res;
}
\end{verbatim}

\paragraph{Complejidad.}
\begin{itemize}\setlength\itemsep{0pt}
	\item $O(\log n)$ por update/query, $O(n)$ build.
	\item En la practica, $\sim 1.5x$ mas rapido que el recursivo para $n \ge 10^5$.
\end{itemize}

\paragraph{Donde tocar para modificar.}
\begin{itemize}\setlength\itemsep{0pt}
	\item \textbf{Suma -> max/min/xor}: cambiar la operacion en \texttt{set} y \texttt{query}.
	\item \textbf{Aniadir update en rango con lazy}: implementar version lazy por separado (no trivial en iterativo).
	\item \textbf{Tamano no potencia de 2}: si $n$ no es potencia de 2, aniadir padding.
	\item \textbf{Persistencia}: usar \texttt{vector<array<int, 2>>} y duplicar nodos.
\end{itemize}

\paragraph{Trampas y errores frecuentes.}
\begin{itemize}\setlength\itemsep{0pt}
	\item En el query, \texttt{r += n+1} (no \texttt{r += n}) porque \texttt{query(l, r)} es \textbf{inclusivo} en ambos extremos.
	\item El indice \texttt{l & 1} detecta si \texttt{l} es hijo derecho; aniadir \texttt{res += st[l++]} en ese caso.
	\item El \texttt{--r} en la condicion derecha es para incluir el limite.
	\item Si las operaciones no son idempotentes, cuidado con el orden.
\end{itemize}

\paragraph{Codigo.}
Ver \texttt{cpp.tex}, seccion \textit{Segment Tree iterativo}.

\vspace{0.5em}

\subsection{Fenwick / BIT}

\subsubsection{BIT point update + prefix sum}

\paragraph{Idea intuitiva.}
El \textbf{Fenwick Tree} (BIT) almacena sums parciales de manera que \texttt{sum(i)} (suma de $[1, i]$) y \texttt{add(i, delta)} cuestan $O(\log n)$. La representacion es ingeniosa: el nodo $i$ cubre el rango $[i - \text{lowbit}(i) + 1, i]$. Asi, \texttt{sum(i)} itera \texttt{i -= lowbit(i)}, sumando los nodos que cubren un prefijo creciente, y \texttt{add(i, delta)} itera \texttt{i += lowbit(i)}, propagando el cambio a todos los nodos cuyo rango contiene a $i$. La idea: cada indice $i$ se descompone en una suma de bloques disjuntos de tamano potencias de 2 (su ``lowbit representation'').

\paragraph{Propiedades clave (invariantes).}
\begin{itemize}\setlength\itemsep{0pt}
	\item Indice \textbf{1-based} (no usar $i=0$).
	\item \texttt{lowbit(i) = i \& (-i)}.
	\item Cada nodo cubre exactamente un rango de tamano \texttt{lowbit(i)}.
	\item La operacion debe ser \textbf{invertible} (suma, xor, no min/max).
\end{itemize}

\paragraph{Codigo clave.}
Las dos operaciones basicas:
\begin{verbatim}
void add(int i, long long d) {
    for (; i <= n; i += i & -i) bit[i] += d;
}
long long sum(int i) {
    long long s = 0;
    for (; i > 0; i -= i & -i) s += bit[i];
    return s;
}
\end{verbatim}

\paragraph{Complejidad.}
\begin{itemize}\setlength\itemsep{0pt}
	\item $O(\log n)$ por operacion. Memoria $O(n)$.
	\item Constante muy pequena ($\sim 1$): en la practica, $\sim 3x$ mas rapido que segment tree.
\end{itemize}

\paragraph{Donde tocar para modificar.}
\begin{itemize}\setlength\itemsep{0pt}
	\item \textbf{Suma -> max/min/xor}: cambiar el \texttt{+=} por la operacion (con identidad apropriada). \textbf{No} funciona para \texttt{min} de forma natural (no es invertible).
	\item \textbf{Tamano dinamico}: aniadir un metodo \texttt{resize(newSize)} que preserve valores.
	\item \textbf{Tamano grande}: con $n = 10^7$, mejor comprimir coordenadas.
	\item \textbf{Operador no conmutativo}: usar segment tree en su lugar.
\end{itemize}

\paragraph{Trampas y errores frecuentes.}
\begin{itemize}\setlength\itemsep{0pt}
	\item Indices 0 vs 1: el BIT es \textbf{1-based}. Si el problema da indices 0-based, aniadir \texttt{+1} en cada llamada.
	\item El metodo \texttt{sum} itera \texttt{i -= i \& -i}; aniadir o quitar 1 rompe el invariante.
	\item Si el tipo es \texttt{int}, posibles overflows; usar \texttt{long long}.
	\item Para rango $[l, r]$, el query es \texttt{sum(r) - sum(l-1)} (cuidado con $l = 1$).
\end{itemize}

\paragraph{Codigo.}
Ver \texttt{cpp.tex}, seccion \textit{BIT point update + prefix sum}.

\vspace{0.5em}

\subsubsection{BIT range add + point query}

\paragraph{Idea intuitiva.}
Truco de \textbf{diferencias}: para sumar $d$ a $[l, r]$, aniadir $d$ en $l$ y $-d$ en $r+1$. Asi, el valor en $i$ es la suma de los prefijos hasta $i$, lo que se calcula con un BIT normal. La idea algebraica: si mantenemos un arreglo de diferencias $\Delta$ donde $a_i = \sum_{j \le i} \Delta_j$, entonces aniadir $d$ a $a[l..r]$ equivale a $\Delta_l \mathrel{+}= d$ y $\Delta_{r+1} \mathrel{-}= d$, y reconstruir $a$ es un prefix sum.

\paragraph{Propiedades clave (invariantes).}
\begin{itemize}\setlength\itemsep{0pt}
	\item Internamente es un BIT 5.1 con la convencion de diferencias.
	\item \texttt{rangeAdd(l, r, delta)} cuesta $O(\log n)$ (dos \texttt{add}).
	\item \texttt{pointQuery(i)} cuesta $O(\log n)$.
	\item Si $r+1 > n$, no se aniade el $-d$ (fuera del rango).
\end{itemize}

\paragraph{Codigo clave.}
El truco de diferencias:
\begin{verbatim}
void rangeAdd(int l, int r, long long d) {
    add(l, d);
    if (r + 1 <= n) add(r + 1, -d);
}
long long pointQuery(int i) { return sum(i); }
\end{verbatim}

\paragraph{Complejidad.}
\begin{itemize}\setlength\itemsep{0pt}
	\item $O(\log n)$ por operacion.
	\item Memoria $O(n)$.
\end{itemize}

\paragraph{Donde tocar para modificar.}
\begin{itemize}\setlength\itemsep{0pt}
	\item \textbf{Persistencia}: aniadir versionado (cada update crea una nueva ``copia'' del estado; en BIT no es trivial, mejor usar persistent segtree).
	\item \textbf{Coordenada compression}: si $n$ es grande pero los updates son pocos, comprimir.
	\item \textbf{2D range add + point query}: aniadir una dimension con la misma tecnica.
\end{itemize}

\paragraph{Trampas y errores frecuentes.}
\begin{itemize}\setlength\itemsep{0pt}
	\item Si $r+1 > n$, \textbf{no} aniadir \texttt{-delta} (no hay efecto fuera del rango); el snippet ya lo chequea.
	\item Para reconstruir el array original tras $k$ range adds, hacer $k$ point queries (no hay shortcut).
	\item Confundir ``suma de diferencias'' con ``suma acumulada''; son cosas distintas.
\end{itemize}

\paragraph{Codigo.}
Ver \texttt{cpp.tex}, seccion \textit{BIT range add + point query}.

\vspace{0.5em}

\subsubsection{BIT range add + range sum}

\paragraph{Idea intuitiva.}
Para queries $\sum_{i=1}^{x} a_i$ despues de varios range adds, el truco de diferencias no basta (necesitas la suma acumulada de las diferencias). Solucion: dos BITs, $t_1$ y $t_2$. Tras aniadir $d$ a $[l, r]$: $t_1$ recibe $+d$ en $l$, $-d$ en $r+1$; $t_2$ recibe $+d(l-1)$ en $l$, $-d \cdot r$ en $r+1$. Entonces $\text{prefixSum}(x) = t_1.\text{sum}(x) \cdot x - t_2.\text{sum}(x)$. La razon matematica: el valor en $i$ es $\sum_{j \le i} \Delta_j$, asi que la suma hasta $i$ es $\sum_{k \le i} (i - k + 1) \Delta_k$; expandir da la formula con dos BITs.

\paragraph{Propiedades clave (invariantes).}
\begin{itemize}\setlength\itemsep{0pt}
	\item Verificacion: para $x \ge r$, $\text{prefixSum}(x)$ acumula correctamente todas las contribuciones.
	\item La formula sale de expandir la suma de diferencias.
	\item La identidad clave: $\sum_{i=1}^{x} i = x(x+1)/2$, asi que la suma de diferencias ponderadas requiere BITs adicionales.
\end{itemize}

\paragraph{Codigo clave.}
Range add con range sum (dos BITs):
\begin{verbatim}
void rangeAdd(int l, int r, long long d) {
    t1.add(l, d); t1.add(r + 1, -d);
    t2.add(l, d * (l - 1)); t2.add(r + 1, -d * r);
}
long long prefixSum(int x) {
    return t1.sum(x) * x - t2.sum(x);
}
\end{verbatim}

\paragraph{Complejidad.}
\begin{itemize}\setlength\itemsep{0pt}
	\item $O(\log n)$ por operacion (2 \texttt{add} o 2 \texttt{sum} por cada lado).
	\item Memoria $O(n)$.
\end{itemize}

\paragraph{Donde tocar para modificar.}
\begin{itemize}\setlength\itemsep{0pt}
	\item \textbf{2D range add + range sum}: aniadir otra dimension con la misma tecnica (queda 4 BITs).
	\item \textbf{Cambiar mod}: aniadir modulo en cada operacion.
	\item \textbf{Output del array}: aniadir un metodo que reconstruya $a_i$ point query por cada $i$.
\end{itemize}

\paragraph{Trampas y errores frecuentes.}
\begin{itemize}\setlength\itemsep{0pt}
	\item Overflow en \texttt{delta * (l - 1)} si los valores son grandes; usar \texttt{long long} y \texttt{\% MOD} al final.
	\item Confundir la convencion $0$-based vs $1$-based del BIT; las dos afectan las formulas.
	\item En 2D, el codigo es $\sim 4x$ mas largo y propenso a errores; verificar con casos pequenos.
\end{itemize}

\paragraph{Codigo.}
Ver \texttt{cpp.tex}, seccion \textit{BIT range add + range sum}.

\subsection{RMQ y sparse table}

\subsubsection{Sparse Table}

\paragraph{Idea intuitiva.}
Precomputa respuestas a queries de tamano $2^k$ para $k = 0, 1, \dots, \log n$. Asi, \texttt{st[k][i]} guarda la respuesta sobre $[i, i + 2^k)$. Para query $[l, r]$, se usan dos bloques de tamano $2^{\lfloor \log_2(r-l+1) \rfloor}$: $[l, l + 2^k)$ y $[r - 2^k + 1, r)$, y se combinan. Esto da $O(1)$ por query \textbf{solo si la operacion es idempotente} (min, max, gcd; \textbf{no} suma). La razon: con idempotencia, solapar los dos bloques no afecta el resultado.

\paragraph{Propiedades clave (invariantes).}
\begin{itemize}\setlength\itemsep{0pt}
	\item \texttt{st[0][i] = a[i]}.
	\item \texttt{st[k][i] = op(st[k-1][i], st[k-1][i + 2\^{}(k-1)])}.
	\item \texttt{LOG = floor(log2(n)) + 1}.
	\item La operacion debe ser \textbf{idempotente} ($f(x, x) = x$) y asociativa.
\end{itemize}

\paragraph{Codigo clave.}
Query $O(1)$ en sparse table:
\begin{verbatim}
long long query(int l, int r) {
    int k = __builtin_clz(1) - __builtin_clz(r - l + 1);  // log2
    return min(st[k][l], st[k][r - (1 << k) + 1]);
}
\end{verbatim}

\paragraph{Complejidad.}
\begin{itemize}\setlength\itemsep{0pt}
	\item Build $O(n \log n)$, query $O(1)$.
	\item Memoria $O(n \log n)$.
\end{itemize}

\paragraph{Donde tocar para modificar.}
\begin{itemize}\setlength\itemsep{0pt}
	\item \textbf{Cambiar min por otra operacion idempotente}: cambiar las dos ocurrencias de \texttt{min(...)} por la operacion (max, gcd). Verificar identidad y asociatividad.
	\item \textbf{Suma en rango}: NO sirve con este metodo; usar prefix sums.
	\item \textbf{Sparse table 2D}: aniadir otra dimension (queda $O(n^2 \log^2 n)$ memoria).
	\item \textbf{LCA queries}: combinar con binary lifting en arboles.
\end{itemize}

\paragraph{Trampas y errores frecuentes.}
\begin{itemize}\setlength\itemsep{0pt}
	\item El query toma $[l, r]$ \textbf{inclusivo} (no $[l, r)$).
	\item El \texttt{LOG} debe estar bien calculado; con $n = 1$, \texttt{LOG = 1}.
	\item Usar sparse table para suma es un error; necesitas prefix sums o segment tree.
	\item Para queries $O(1)$ con operaciones no idempotentes, usar Disjoint Sparse Table.
\end{itemize}

\paragraph{Codigo.}
Ver \texttt{cpp.tex}, seccion \textit{Sparse Table}.

\vspace{0.5em}

\subsubsection{Disjoint Sparse Table}

\paragraph{Idea intuitiva.}
Variante que tambien da $O(1)$ por query pero \textbf{no requiere idempotencia} para algunas operaciones (sirve para suma, xor, etc. en linea con la mayoria de operaciones asociativas). La idea: para cada nivel $k$, los rangos ``se parten en mitades''. \texttt{st[k][i]} cubre un rango que \textbf{empieza} en $i$ y se extiende a izquierda o derecha hasta encontrar un ``punto de division''. El truco: para query $[l, r]$, el nivel optimo es el bit mas alto donde $l$ y $r$ difieren; los rangos $[l, ?]$ y $[?, r]$ cubren exactamente el query sin solape.

\paragraph{Propiedades clave (invariantes).}
\begin{itemize}\setlength\itemsep{0pt}
	\item Construccion diferente al Sparse Table clasico: en cada nivel, se extiende desde cada posicion hacia ambos lados hasta encontrar la mitad de nivel.
	\item \texttt{query(l, r)} usa \texttt{k = log2(l \^{} r)} y combina \texttt{st[k][l]} y \texttt{st[k][r]}.
	\item La operacion debe ser \textbf{asociativa} pero no necesariamente idempotente.
\end{itemize}

\paragraph{Codigo clave.}
Query con DST:
\begin{verbatim}
long long query(int l, int r) {
    if (l == r) return st[0][l];
    int k = 31 - __builtin_clz(l ^ r);  // bit mas alto que difiere
    return op(st[k][l], st[k][r]);
}
\end{verbatim}

\paragraph{Complejidad.}
\begin{itemize}\setlength\itemsep{0pt}
	\item Build $O(n \log n)$, query $O(1)$.
	\item Memoria $O(n \log n)$.
\end{itemize}

\paragraph{Donde tocar para modificar.}
\begin{itemize}\setlength\itemsep{0pt}
	\item \textbf{Operacion no asociativa}: no funciona; el Disjoint Sparse Table asume asociatividad.
	\item \textbf{Memory compression}: con $n \le 10^6$ alcanza; con mas, comprimir.
	\item \textbf{Operaciones no conmutativas}: aniadir el cuidado de que el orden es importante.
\end{itemize}

\paragraph{Trampas y errores frecuentes.}
\begin{itemize}\setlength\itemsep{0pt}
	\item La construccion es sutil; verificar con casos pequenos antes de usar.
	\item El bit ``mas alto que difiere'' requiere cuidado con $l = r$; aniadir caso base.
	\item Si la operacion no es asociativa, los dos bloques no se pueden combinar en cualquier orden.
\end{itemize}

\paragraph{Codigo.}
Ver \texttt{cpp.tex}, seccion \textit{Disjoint Sparse Table}.

\subsection{Trees y DP en arbol}

\subsubsection{Binary Lifting / LCA}

\paragraph{Idea intuitiva.}
Precomputa para cada nodo $v$ sus ancestros a distancias $2^k$ (los ``saltos''). Asi, para subir $d$ niveles, se descompone $d$ en binario y se aplican los saltos correspondientes. El \textbf{LCA} (Lowest Common Ancestor) se obtiene subiendo el nodo mas profundo hasta igualar profundidades y luego subiendo ambos hasta que sus padres coincidan. La intuicion: cualquier entero $\le n$ se representa en binario con $\le \log n$ bits, asi que subir $d$ niveles requiere $\le \log n$ saltos.

\paragraph{Propiedades clave (invariantes).}
\begin{itemize}\setlength\itemsep{0pt}
	\item \texttt{up[k][v]} = ancestro de $v$ a $2^k$ niveles.
	\item \texttt{LOG = floor(log2(n)) + 1}.
	\item La recursion de \texttt{dfs} llena \texttt{up[0][v]} (el padre directo) y por transitividad los demas.
	\item El LCA esta definido como el ancestro comun mas profundo de $u$ y $v$.
\end{itemize}

\paragraph{Codigo clave.}
LCA con binary lifting:
\begin{verbatim}
int lca(int u, int v) {
    if (depth[u] < depth[v]) swap(u, v);
    int diff = depth[u] - depth[v];
    for (int k = LOG - 1; k >= 0; --k)
        if (diff & (1 << k)) u = up[k][u];
    if (u == v) return u;
    for (int k = LOG - 1; k >= 0; --k)
        if (up[k][u] != up[k][v])
            u = up[k][u], v = up[k][v];
    return up[0][u];
}
\end{verbatim}

\paragraph{Complejidad.}
\begin{itemize}\setlength\itemsep{0pt}
	\item Build $O(n \log n)$, query LCA $O(\log n)$.
	\item Memoria $O(n \log n)$.
\end{itemize}

\paragraph{Donde tocar para modificar.}
\begin{itemize}\setlength\itemsep{0pt}
	\item \textbf{Distancia en el arbol}: combinar con \texttt{depth[]} para calcular \texttt{depth[u] + depth[v] - 2*depth[lca]}.
	\item \textbf{K-esimo ancestro}: subir $k$ niveles con el mismo truco.
	\item \textbf{Aniadir peso a las aristas}: guardar \texttt{dist[2\^{}k][v]} (distancia acumulada a $2^k$ saltos) ademas de \texttt{up}.
	\item \textbf{Sparse table sobre depth}: alternativa $O(1)$ a LCA (mas complejo).
\end{itemize}

\paragraph{Trampas y errores frecuentes.}
\begin{itemize}\setlength\itemsep{0pt}
	\item El arbol debe ser \textbf{rooted} (el snippet asume eso); si no, hacer un \texttt{dfs} desde cualquier nodo.
	\item Para grafos no conexos, correr \texttt{dfs} desde cada componente.
	\item En la segunda fase del LCA, subir ambos nodos mientras \texttt{up[k][u] != up[k][v]}; el LCA final es \texttt{up[0][u]}.
	\item El padre de la raiz es 0 (o -1); aniadir check si se sale del rango.
\end{itemize}

\paragraph{Codigo.}
Ver \texttt{cpp.tex}, seccion \textit{Binary Lifting / LCA}.

\vspace{0.5em}

\subsubsection{Heavy-Light Decomposition (HLD)}

\paragraph{Idea intuitiva.}
Descomponer un arbol en \textbf{cadenas pesadas} (paths donde cada nodo (salvo la raiz) tiene al hijo ``pesado'' como hijo de tamano maximo de su subarbol). Cada cadena se mapea a un rango contiguo en un arreglo base; asi, queries sobre \textbf{paths} se descomponen en $O(\log n)$ queries sobre rangos contiguos (cada cadena es un rango), que se resuelven con un BIT o segment tree. La intuicion: por cada nivel de recursion, el subarbol del hijo pesado es al menos la mitad, asi que un path cruza $\le 2 \log n$ cadenas.

\paragraph{Propiedades clave (invariantes).}
\begin{itemize}\setlength\itemsep{0pt}
	\item \texttt{heavy[v]} = hijo pesado de $v$ (el de mayor subarbol).
	\item \texttt{head[v]} = cabeza de la cadena que contiene a $v$.
	\item \texttt{pos[v]} = posicion en el arreglo base.
	\item Cada nodo pertenece a exactamente una cadena.
\end{itemize}

\paragraph{Codigo clave.}
El query sobre un path $u \to v$:
\begin{verbatim}
long long pathQuery(int u, int v) {
    long long res = 0;
    while (head[u] != head[v]) {
        if (depth[head[u]] < depth[head[v]]) swap(u, v);
        res += seg.query(pos[head[u]], pos[u]);
        u = parent[head[u]];
    }
    if (depth[u] > depth[v]) swap(u, v);
    res += seg.query(pos[u], pos[v]);
    return res;
}
\end{verbatim}

\paragraph{Complejidad.}
\begin{itemize}\setlength\itemsep{0pt}
	\item Build $O(n)$, query sobre path $O(\log^2 n)$ (con segtree) o $O(\log n)$ (con BIT si solo se hace suma).
	\item Memoria $O(n)$.
\end{itemize}

\paragraph{Donde tocar para modificar.}
\begin{itemize}\setlength\itemsep{0pt}
	\item \textbf{Usar con segment tree}: aniadir un segtree sobre el arreglo base; \texttt{pathSegments} da los rangos a actualizar/query.
	\item \textbf{Cambiar operacion}: ajustar el segtree (suma, max, etc.).
	\item \textbf{Subtree query}: \texttt{subtree[v]} es el rango \texttt{[pos[v], pos[v] + size[v] - 1]} en el arreglo base.
	\item \textbf{Updates en edges}: aniadir un offset para que el edge $(u, v)$ quede en el nodo mas profundo.
\end{itemize}

\paragraph{Trampas y errores frecuentes.}
\begin{itemize}\setlength\itemsep{0pt}
	\item \texttt{pathSegments} requiere un \texttt{segtree} o \texttt{BIT} externo para ser util; por si solo solo da los rangos.
	\item Para edges vs nodos: si la operacion es sobre edges, guardar el valor en el hijo mas profundo.
	\item Olvidar la recursion \texttt{while (head[u] != head[v])} da resultados incorrectos.
	\item En problemas con updates, asegurarse de que el segtree soporte lazy propagation.
\end{itemize}

\paragraph{Codigo.}
Ver \texttt{cpp.tex}, seccion \textit{Heavy-Light Decomposition (HLD)}.

\vspace{0.5em}

\subsubsection{Centroid Decomposition}

\paragraph{Idea intuitiva.}
Descomponer recursivamente el arbol en \textbf{centroides}. Un centroide de un arbol es un nodo tal que cada subtree (al removerlo) tiene tamano $\le n/2$. Existe siempre y se encuentra en $O(n)$. Esto da una estructura de \textbf{arbol de centroides} de profundidad $O(\log n)$, util para problemas ``para cada nodo, encontrar el mas cercano que cumple X''. La razon de la profundidad $O(\log n)$: cada recursion reduce el problema a subarboles de tamano $\le n/2$, asi que la recursion da una cadena geometrica decreciente.

\paragraph{Propiedades clave (invariantes).}
\begin{itemize}\setlength\itemsep{0pt}
	\item En cada nivel, los subarboles se procesan recursivamente.
	\item \texttt{removed[]} marca nodos ya usados como centroide.
	\item \texttt{centroidParent[]} da el padre en el arbol de centroides.
	\item La altura del arbol de centroides es $O(\log n)$.
\end{itemize}

\paragraph{Codigo clave.}
Encontrar el centroide de un subarbol:
\begin{verbatim}
int findCentroid(int v, int p, int sz) {
    for (int u : adj[v])
        if (u != p && !removed[u] && subsize[u] > sz / 2)
            return findCentroid(u, v, sz);
    return v;
}
\end{verbatim}

\paragraph{Complejidad.}
\begin{itemize}\setlength\itemsep{0pt}
	\item Build $O(n \log n)$, queries dependen del problema.
	\item Cada nivel de recursion es $O(n)$ para calcular subsize.
\end{itemize}

\paragraph{Donde tocar para modificar.}
\begin{itemize}\setlength\itemsep{0pt}
	\item \textbf{Query tipica}: en cada nivel, ``expandir'' desde el centroide contando caminos validos y usar una estructura auxiliar (set, multiset).
	\item \textbf{Distancia maxima a un nodo ``malo''}: mantener un multiset de distancias en cada centroide.
	\item \textbf{Cambiar a otro tipo de descomposicion}: si el problema se adapta mejor a small-to-large o DSU on tree, usar esos.
	\item \textbf{Dynamic updates}: combinacion con link-cut trees.
\end{itemize}

\paragraph{Trampas y errores frecuentes.}
\begin{itemize}\setlength\itemsep{0pt}
	\item El snippet actual \textbf{solo construye} la descomposicion; las queries se aniaden segun el problema.
	\item Olvidar resetear \texttt{removed[]} causa recursion infinita.
	\item El calculo de subsize debe ignorar nodos ya ``removed''.
	\item El centroide es \textbf{unico} salvo empates en tamano; cualquiera sirve.
\end{itemize}

\paragraph{Codigo.}
Ver \texttt{cpp.tex}, seccion \textit{Centroid Decomposition}.

\vspace{0.5em}

\subsubsection{Euler Tour + range queries}

\paragraph{Idea intuitiva.}
Un \textbf{Euler Tour} (first-time visit) numera los nodos en el orden en que se visitan. El subarbol de $v$ corresponde a un \textbf{rango contiguo} $[tin[v], tout[v]]$ en este orden. Asi, queries sobre subarboles se resuelven con un BIT o segment tree sobre el arreglo de tiempos. La razon del rango contiguo: cuando entramos a un subarbol, no salimos hasta haber recorrido todos sus descendientes; entonces las posiciones consecutivas corresponden exactamente al subarbol.

\paragraph{Propiedades clave (invariantes).}
\begin{itemize}\setlength\itemsep{0pt}
	\item \texttt{tin[v]} = tiempo de primera visita a $v$.
	\item \texttt{tout[v]} = tiempo al \textbf{terminar} de procesar el subarbol de $v$.
	\item El subarbol de $v$ ocupa las posiciones $[tin[v], tout[v]]$.
	\item Para verificar si $u$ es ancestro de $v$: \texttt{tin[u] <= tin[v] \&\& tout[v] <= tout[u]}.
\end{itemize}

\paragraph{Codigo clave.}
El DFS del Euler Tour:
\begin{verbatim}
void dfs(int v, int p) {
    tin[v] = ++timer;
    for (int u : adj[v])
        if (u != p) dfs(u, v);
    tout[v] = timer;
}
\end{verbatim}

\paragraph{Complejidad.}
\begin{itemize}\setlength\itemsep{0pt}
	\item Build $O(n)$, query sobre subarbol $O(\log n)$ con BIT/segtree.
	\item Memoria $O(n)$.
\end{itemize}

\paragraph{Donde tocar para modificar.}
\begin{itemize}\setlength\itemsep{0pt}
	\item \textbf{Aniadir valor por nodo}: guardar el valor en \texttt{tin[v]} y operar sobre el rango.
	\item \textbf{Path query}: el Euler Tour \textbf{no} sirve para paths; usar HLD o LCA.
	\item \textbf{Multiple queries offline}: procesar en orden de tiempo y mantener una estructura por profundidad.
	\item \textbf{Verificar ancestro}: $O(1)$ con la condicion \texttt{tin}.
\end{itemize}

\paragraph{Trampas y errores frecuentes.}
\begin{itemize}\setlength\itemsep{0pt}
	\item El snippet solo construye el tour; las queries concretas se aniaden.
	\item \texttt{tout[v]} esta en el momento de salir, asi que \texttt{tout[v] - tin[v] + 1} = tamano del subarbol.
	\item Si el arbol no es conexo, aniadir un loop externo sobre cada componente.
	\item Para queries en tiempo (no subtree), usar el orden del tour y offline processing.
\end{itemize}

\paragraph{Codigo.}
Ver \texttt{cpp.tex}, seccion \textit{Euler Tour + range queries}.

\subsection{BST balanceado}

\subsubsection{Treap implicito}

\paragraph{Idea intuitiva.}
Un \textbf{treap} es un arbol binario de busqueda donde cada nodo tiene una \textbf{clave} (el valor de la secuencia) y una \textbf{prioridad} aleatoria. La invariante es: BST por clave, heap por prioridad. Esto da un arbol \textbf{esperado balanceado} sin necesidad de rebalanceo explicito. Un \textbf{treap implicito} no almacena una clave; la posicion de un nodo en la secuencia esta dada por el \textbf{tamano del subarbol izquierdo}, permitiendo split/merge por posicion. La estructura es esencialmente un RBS (randomized BST): la altura esperada es $O(\log n)$ por la propiedad de heap + BST.

\paragraph{Propiedades clave (invariantes).}
\begin{itemize}\setlength\itemsep{0pt}
	\item \texttt{sz(node)} = tamano del subarbol (1 + izquierdo + derecho).
	\item \texttt{split(t, k, l, r)}: parte \texttt{t} en \texttt{l} (primeros $k$ nodos) y \texttt{r} (resto).
	\item \texttt{merge(l, r)}: une dos treaps (todos los de \texttt{l} antes que los de \texttt{r}).
	\item La prioridad aleatoria asegura equilibrio esperado.
\end{itemize}

\paragraph{Codigo clave.}
Split por posicion (la base del treap implicito):
\begin{verbatim}
void split(Node* t, int k, Node*& l, Node*& r) {
    if (!t) { l = r = nullptr; return; }
    push(t);
    if (sz(t->l) >= k) {
        split(t->l, k, l, t->l);
        r = update(t);
    } else {
        split(t->r, k - sz(t->l) - 1, t->r, r);
        l = update(t);
    }
}
\end{verbatim}

\paragraph{Complejidad.}
\begin{itemize}\setlength\itemsep{0pt}
	\item $O(\log n)$ \textbf{esperado} por operacion (split, merge, insert, erase).
	\item Peor caso exponencial (probabilidad astronomicamente pequena).
\end{itemize}

\paragraph{Donde tocar para modificar.}
\begin{itemize}\setlength\itemsep{0pt}
	\item \textbf{Treap con clave (no implicito)}: usar \texttt{key} en vez de posicion; cambiar \texttt{split} por \texttt{splitByKey}.
	\item \textbf{Lazy propagation}: aniadir campos \texttt{lazy} y aplicarlos en \texttt{push} (recursivo).
	\item \textbf{Persistencia}: aniadir versionado clonando nodos.
	\item \textbf{Mejor aleatoriedad}: en lugar de \texttt{rand()}, usar \texttt{rng()} (mejor estado).
	\item \textbf{Operacion de reversa}: aniadir \texttt{lazyReverse} y swap de hijos en push.
\end{itemize}

\paragraph{Trampas y errores frecuentes.}
\begin{itemize}\setlength\itemsep{0pt}
	\item En el snippet actual, \texttt{Node(int v)} usa \texttt{rand()} que en algunos compiladores da maxima prioridad 0; aniadir \texttt{rng()} para mejor distribucion.
	\item Olvidar \texttt{update(t)} tras modificar hijos da tamano incorrecto.
	\item Si lazy propagation no se aplica en push, las queries dan valores viejos.
	\item Memory leak con \texttt{new}; en contests OK, en produccion aniadir destructor.
\end{itemize}

\paragraph{Codigo.}
Ver \texttt{cpp.tex}, seccion \textit{Treap implicito}.

\vspace{0.5em}

\subsubsection{Mo con updates}

\paragraph{Idea intuitiva.}
Extension del algoritmo de Mo para queries que mezclan \textbf{rango} y \textbf{updates}. Las queries se ordenan por $(L/\text{blockSize}, R/\text{blockSize}, T)$ donde $T$ es el ``tiempo'' (numero de updates). Asi, moverse entre queries consecutivas cuesta $O(n^{2/3})$ amortizado. La intuicion: el sort por bloques reduce los movimientos totales, y la tercera coordenada $T$ se optimiza con alternancia similar a $R$.

\paragraph{Propiedades clave (invariantes).}
\begin{itemize}\setlength\itemsep{0pt}
	\item Tres punteros: $L$, $R$, $T$.
	\item \texttt{applyUpdate} y \texttt{revertUpdate} son \textbf{inversos} exactos.
	\item \texttt{addPos} y \texttt{removePos} mantienen el estado.
	\item El tamano de bloque optimo es $n^{2/3}$ (donde $n$ = tamano del array, $Q$ = numero de queries).
\end{itemize}

\paragraph{Codigo clave.}
El sort con tres coordenadas:
\begin{verbatim}
int blockSize = (int)pow(n, 2.0/3.0);
sort(queries.begin(), queries.end(), [&](const Q& a, const Q& b) {
    int blockA = a.l / blockSize, blockB = b.l / blockSize;
    if (blockA != blockB) return blockA < blockB;
    int blockAR = a.r / blockSize, blockBR = b.r / blockSize;
    if (blockAR != blockBR) return blockAR < blockBR;
    return a.t < b.t;
});
\end{verbatim}

\paragraph{Complejidad.}
\begin{itemize}\setlength\itemsep{0pt}
	\item $O((N + Q) \cdot n^{2/3})$ donde $n^{2/3}$ es el tamano de bloque optimo.
	\item En problemas reales, funciona bien hasta $N, Q \sim 10^5$.
\end{itemize}

\paragraph{Donde tocar para modificar.}
\begin{itemize}\setlength\itemsep{0pt}
	\item \textbf{El corazon es el ``add/remove''}: aniadir la logica del problema (ej. contar colores distintos, suma de frecuencias, etc.).
	\item \textbf{Cambiar blockSize}: si la mayoria de updates vs queries cambia, ajustar.
	\item \textbf{Rollback}: aniadir \texttt{save()} / \texttt{restore()} para volver a estados anteriores.
	\item \textbf{Hilbert order}: alternativa con mejor locality.
\end{itemize}

\paragraph{Trampas y errores frecuentes.}
\begin{itemize}\setlength\itemsep{0pt}
	\item El snippet actual es solo el \textbf{esqueleto}; \texttt{addPos}/\texttt{removePos}/\texttt{applyUpdate}/\texttt{revertUpdate} deben ser implementados por problema.
	\item Olvidar \texttt{revertUpdate} es el error mas comun.
	\item Si el numero de updates es muy grande, Mo con updates se vuelve lento; considerar otra estructura.
	\item Mantener consistencia entre \texttt{T} y la posicion de los updates.
\end{itemize}

\paragraph{Codigo.}
Ver \texttt{cpp.tex}, seccion \textit{Mo con updates}.

\vspace{0.5em}

\subsubsection{sqrt-decomposition generico}

\paragraph{Idea intuitiva.}
Dividir el arreglo en \textbf{bloques} de tamano $B \approx \sqrt{n}$. Mantener informacion precomputada por bloque (suma, max, etc.). Queries de rango se resuelven combinando bloques enteros ($O(\sqrt{n})$ en total) y elementos sueltos en los bloques parciales. Updates puntuales actualizan el valor y el resumen del bloque. La intuicion: las queries ``normales'' cruzan $O(n/B)$ bloques; tomar $B = \sqrt{n}$ balancea el costo de bloques enteros vs elementos sueltos.

\paragraph{Propiedades clave (invariantes).}
\begin{itemize}\setlength\itemsep{0pt}
	\item $B = \lfloor \sqrt{n} \rfloor + 1$.
	\item \texttt{bucket[i]} = resumen del bloque $i$-esimo (suma, max, etc.).
	\item \texttt{rangeSum(l, r)}: 3 casos segun $l$ y $r$ esten en el mismo bloque o no.
	\item Cada query cruza $\le 2 + (r - l)/B$ bloques.
\end{itemize}

\paragraph{Codigo clave.}
Estructura basica:
\begin{verbatim}
int B = sqrt(n) + 1;
vector<T> bucket(n / B + 1, IDENTITY);
T query(int l, int r) {
    T res = IDENTITY;
    int bl = l / B, br = r / B;
    if (bl == br) {
        for (int i = l; i <= r; ++i) res = combine(res, a[i]);
    } else {
        for (int i = l; i < (bl+1)*B; ++i) res = combine(res, a[i]);
        for (int i = bl+1; i < br; ++i) res = combine(res, bucket[i]);
        for (int i = br*B; i <= r; ++i) res = combine(res, a[i]);
    }
    return res;
}
\end{verbatim}

\paragraph{Complejidad.}
\begin{itemize}\setlength\itemsep{0pt}
	\item $O(\sqrt{n})$ por operacion.
	\item Memoria $O(n)$.
\end{itemize}

\paragraph{Donde tocar para modificar.}
\begin{itemize}\setlength\itemsep{0pt}
	\item \textbf{Cambiar suma por max/min}: ajustar \texttt{bucket} y la combinacion.
	\item \textbf{Aniadir update en rango}: si la operacion es asociativa, sumar al resumen del bloque + ajustar elementos sueltos.
	\item \textbf{Tamano de bloque optimo}: si las queries son muy largas, $B = n^{2/3}$ (sirve para Mo con updates).
	\item \textbf{Operacion no asociativa}: usar segment tree en su lugar.
\end{itemize}

\paragraph{Trampas y errores frecuentes.}
\begin{itemize}\setlength\itemsep{0pt}
	\item Si $n$ no es multiplo de $B$, el ultimo bloque es mas chico; aniadir check.
	\item La identidad de la operacion debe estar bien definida (0 para suma, $-\infty$ para max).
	\item Si la operacion no es asociativa, los bloques no se pueden combinar; cambiar a segment tree.
\end{itemize}

\paragraph{Codigo.}
Ver \texttt{cpp.tex}, seccion \textit{sqrt-decomposition generico}.

% ============================================================
\section{Strings}
% ============================================================

\subsection{Hashing}

\subsubsection{Rolling Hash (doble modulo)}

\paragraph{Idea intuitiva.}
Un \textbf{rolling hash} asigna a cada prefijo del string un valor numerico tal que el hash de un substring se puede obtener en $O(1)$ a partir de los prefijos. La forma clasica: $h[i] = h[i-1] \cdot A + s[i-1]$ (donde $A$ es una base y todo modulo $M$). Asi, $h[r] - h[l-1] \cdot A^{r-l+1}$ da el hash del substring $[l, r]$ (con $A^{r-l+1}$ precomputado). Usar \textbf{dos modulos} distintos reduce drasticamente la probabilidad de colision. La idea algebraica: el string se interpreta como un polinomio en $A$ con coeficientes los caracteres; el hash es su evaluacion modulo $M$.

\paragraph{Propiedades clave (invariantes).}
\begin{itemize}\setlength\itemsep{0pt}
	\item $h[0] = 0$, $h[i]$ = hash de los primeros $i$ caracteres.
	\item $p[i] = A^i \bmod M$ precomputado.
	\item Comparar dos substrings: comparar sus pares de hashes.
	\item \textbf{Colisiones}: posibles (probabilidad $1/M$ por par). Doble hash: $\approx 1/(M_1 M_2)$.
	\item Si $A$ y $M$ son coprimos, el hash es bijectivo modulo $M$ para strings de longitud fija.
\end{itemize}

\paragraph{Codigo clave.}
La formula del hash de substring:
\begin{verbatim}
pair<long long, long long> get(int l, int r) {
    // h[r] - h[l-1] * A^(r-l+1)
    long long x1 = (h1[r] - h1[l-1] * p1[r-l+1] % M1 + M1) % M1;
    long long x2 = (h2[r] - h2[l-1] * p2[r-l+1] % M2 + M2) % M2;
    return {x1, x2};
}
\end{verbatim}

\paragraph{Complejidad.}
\begin{itemize}\setlength\itemsep{0pt}
	\item Precompute $O(n)$, query hash $O(1)$.
	\item Memoria $O(n)$.
\end{itemize}

\paragraph{Donde tocar para modificar.}
\begin{itemize}\setlength\itemsep{0pt}
	\item \textbf{Caracteres con valor $> 26$}: si los strings incluyen simbolos, aniadir \texttt{if (c >= 'a' \&\& c <= 'z') v = c - 'a'; else v = 26 + ...}.
	\item \textbf{Hash de substrings con updates}: aniadir un Fenwick tree / segment tree sobre los hashes (no es trivial).
	\item \textbf{Cambiar mod}: usar dos primos grandes ($10^9+7$ y $10^9+9$) o un \texttt{u128} para un solo hash sin mod.
	\item \textbf{Base aleatoria}: aniadir random para evitar ataques (en contests no es necesario; en produccion si).
	\item \textbf{Hash polinomial}: $h = \sum s_i \cdot A^i$ en vez de $s_i \cdot A^{i-1}$; algunos problemas requieren este.
\end{itemize}

\paragraph{Trampas y errores frecuentes.}
\begin{itemize}\setlength\itemsep{0pt}
	\item Si los substrings a comparar estan vacios (\texttt{l > r}), su hash es 0; aniadir check.
	\item \textbf{Overflow} en \texttt{h[i-1] * A}: usar \texttt{long long} y aplicar \texttt{\% M} en cada operacion (no al final).
	\item La sustraccion \texttt{h[r] - h[l-1]*A^(r-l+1)} puede dar negativo; aniadir \texttt{+M} antes de \texttt{\% M}.
	\item Con doble hash, verificar \textbf{ambos} componentes; si solo uno coincide, es colision.
\end{itemize}

\paragraph{Codigo.}
Ver \texttt{cpp.tex}, seccion \textit{Rolling Hash (doble modulo)}.

\subsection{Patrones}

\subsubsection{KMP + LPS}

\paragraph{Idea intuitiva.}
El algoritmo \textbf{Knuth-Morris-Pratt} busca todas las ocurrencias de un patron $P$ en un texto $T$ en $O(|T| + |P|)$. La clave es la tabla \textbf{LPS} (Longest Proper Prefix which is Suffix) que, para cada posicion $i$ del patron, dice el largo del prefijo mas largo que tambien es sufijo \textbf{propio}. Durante el match, cuando hay mismatch, en vez de volver al inicio, saltamos al LPS anterior, evitando trabajo redundante. La demostracion de $O(|T| + |P|)$: cada iteracion del bucle externo avanza $i$ o decrece $j$ (con $j$ acotado por $|P|$), asi que en total hay $O(|T| + |P|)$ operaciones.

\paragraph{Propiedades clave (invariantes).}
\begin{itemize}\setlength\itemsep{0pt}
	\item \texttt{lps[i]} = largo del prefijo mas largo que es sufijo propio de \texttt{p[0..i]}.
	\item El match usa dos punteros \texttt{i} (texto) y \texttt{j} (patron); en mismatch, $j = lps[j-1]$.
	\item Cuando $j == |P|$, se encontro una ocurrencia; continuar con $j = lps[j-1]$ para buscar mas.
	\item El LPS se construye con una auto-aplicacion de KMP al patron.
\end{itemize}

\paragraph{Codigo clave.}
Construccion de la tabla LPS:
\begin{verbatim}
for (int i = 1; i < m; ++i) {
    int j = lps[i - 1];
    while (j > 0 && p[i] != p[j]) j = lps[j - 1];
    if (p[i] == p[j]) ++j;
    lps[i] = j;
}
\end{verbatim}

\paragraph{Complejidad.}
\begin{itemize}\setlength\itemsep{0pt}
	\item $O(|T| + |P|)$.
	\item En la practica, rapido y robusto.
\end{itemize}

\paragraph{Donde tocar para modificar.}
\begin{itemize}\setlength\itemsep{0pt}
	\item \textbf{Contar ocurrencias}: el snippet ya cuenta; para devolver posiciones, aniadir un \texttt{vector<int>}.
	\item \textbf{Varios patrones}: usar Aho-Corasick en su lugar.
	\item \textbf{Matcher exacto con wildcard}: aniadir logica para que \texttt{?} matchee cualquier caracter.
	\item \textbf{Tamano pequeno}: si el patron es $\le 5$ caracteres, Rabin-Karp es mas simple.
	\item \textbf{2D KMP}: aniadir una dimension; util en problemas de matrices.
\end{itemize}

\paragraph{Trampas y errores frecuentes.}
\begin{itemize}\setlength\itemsep{0pt}
	\item En el snippet, \texttt{i} y \texttt{j} son 0-based. La condicion \texttt{j == sz\_pattern} para contar funciona.
	\item La construccion de LPS usa el mismo patron como ``texto''.
	\item Si el patron tiene caracteres repetidos, el LPS puede ser no trivial; verificar con casos pequenos.
\end{itemize}

\paragraph{Codigo.}
Ver \texttt{cpp.tex}, seccion \textit{KMP + LPS}.

\vspace{0.5em}

\subsubsection{Z-algorithm}

\paragraph{Idea intuitiva.}
La \textbf{Z-function} $Z[i]$ de un string $S$ es el largo del prefijo mas largo de $S$ que \textbf{tambien} empieza en $S[i]$. Asi, $Z[0] = n$ por convencion. Se calcula en $O(n)$ manteniendo una ventana $[l, r]$ que es el match mas largo encontrado hasta ahora. La consulta $S[l..r] = S[0..r-l]$ da informacion sobre $Z[i]$ cuando $i \le r$. La intuicion: cualquier $i$ dentro de la ventana ``ve'' los caracteres desde una posicion conocida, asi que $Z[i]$ se puede copiar en lugar de recalcular.

\paragraph{Propiedades clave (invariantes).}
\begin{itemize}\setlength\itemsep{0pt}
	\item Construir $Z$ cuesta $O(n)$ con el truco de la ventana.
	\item Para matching de patron $P$ en texto $T$: concatenar $P + \# + T$ y aniadir todos los $i$ con $Z[i] = |P|$.
	\item La ventana $[l, r]$ es siempre la del match mas largo a la derecha.
\end{itemize}

\paragraph{Codigo clave.}
Construccion de Z-function en $O(n)$:
\begin{verbatim}
Z[0] = n;
int l = 0, r = 0;
for (int i = 1; i < n; ++i) {
    if (i <= r) Z[i] = min(r - i + 1, Z[i - l]);
    while (i + Z[i] < n && s[Z[i]] == s[i + Z[i]]) ++Z[i];
    if (i + Z[i] - 1 > r) { l = i; r = i + Z[i] - 1; }
}
\end{verbatim}

\paragraph{Complejidad.}
\begin{itemize}\setlength\itemsep{0pt}
	\item $O(n)$ para construir; matching $O(n + m)$.
	\item Cada caracter se compara a lo sumo 2 veces (extension + reduccion).
\end{itemize}

\paragraph{Donde tocar para modificar.}
\begin{itemize}\setlength\itemsep{0pt}
	\item \textbf{Detectar bordes (border)}: $Z[i] = n - i$ indica que $S[0..n-i-1]$ es borde.
	\item \textbf{Prefijo comun mas largo de dos strings}: aniadir $S_1 + \# + S_2$; $Z[|S_1|+1]$ es el prefijo comun.
	\item \textbf{Comparar prefijos y sufijos}: combinar $S + \# + \text{reverse}(S)$.
	\item \textbf{Matching aproximado}: aniadir tolerancia a mismatches.
\end{itemize}

\paragraph{Trampas y errores frecuentes.}
\begin{itemize}\setlength\itemsep{0pt}
	\item $Z[0]$ siempre vale $n$ (no se calcula).
	\item Si $i > r$, la ventana no aporta; empezar desde 0.
	\item El caracter separador en $P + \# + T$ debe ser uno que no aparezca en $P$ o $T$.
\end{itemize}

\paragraph{Codigo.}
Ver \texttt{cpp.tex}, seccion \textit{Z-algorithm}.

\vspace{0.5em}

\subsubsection{Rabin-Karp (single hash)}

\paragraph{Idea intuitiva.}
Rolling hash aplicado a pattern matching: calcular el hash del patron $P$ y de cada ventana de tamano $|P|$ en $T$. Si los hashes coinciden, verificar caracter por caracter (para descartar colisiones). Esto da $O(n + m)$ \textbf{esperado}. La razon de la eficiencia esperada: cada ventana tiene probabilidad $1/M$ de tener un hash coincidente por casualidad, asi que solo se verifica $O(n/M)$ ventanas en promedio.

\paragraph{Propiedades clave (invariantes).}
\begin{itemize}\setlength\itemsep{0pt}
	\item Si los hashes difieren, los substrings son distintos (sin colision).
	\item Si coinciden, \textbf{verificar} (la verificacion es $O(m)$ pero rara vez se ejecuta).
	\item Peor caso (muchas colisiones): $O(nm)$.
	\item Para $M \ge 10^9$, las colisiones son raras; en la practica nunca pasan.
\end{itemize}

\paragraph{Codigo clave.}
Rolling hash del texto (deslizar ventana):
\begin{verbatim}
long long hash = 0;
for (int i = 0; i < m; ++i) hash = (hash * BASE + text[i]) % MOD;
for (int i = m; i < n; ++i) {
    if (hash == target) check(i - m);
    hash = (hash - text[i - m] * h) * BASE + text[i];  // mod correctamente
    hash = (hash % MOD + MOD) % MOD;
}
\end{verbatim}

\paragraph{Complejidad.}
\begin{itemize}\setlength\itemsep{0pt}
	\item Esperado $O(n + m)$, peor $O(nm)$.
	\item Con doble hash y $M_1, M_2 \ge 10^9$, el peor caso es astronomicamente improbable.
\end{itemize}

\paragraph{Donde tocar para modificar.}
\begin{itemize}\setlength\itemsep{0pt}
	\item \textbf{Doble hash}: aniadir un segundo par \texttt{(BASE2, MOD2)} para reducir colisiones.
	\item \textbf{Varios patrones}: aniadir un set de hashes y buscar match con cualquiera.
	\item \textbf{Tamano de ventana variable}: si los patrones tienen largos distintos, usar el largo especifico por patron.
	\item \textbf{Multi-pattern search}: usar Aho-Corasick en su lugar.
\end{itemize}

\paragraph{Trampas y errores frecuentes.}
\begin{itemize}\setlength\itemsep{0pt}
	\item Overflow en \texttt{(tHash - text[i] * h) * BASE}: aplicar mod correctamente y aniadir \texttt{+MOD} antes de multiplicar.
	\item El factor $h = \text{BASE}^{m-1} \bmod \text{MOD}$ se precomputa.
	\item Si \texttt{MOD} no es primo, las colisiones son mas frecuentes; usar primo.
	\item En el peor caso, Rabin-Karp puede ser lento; preferir KMP si se requiere worst-case.
\end{itemize}

\paragraph{Codigo.}
Ver \texttt{cpp.tex}, seccion \textit{Rabin-Karp (single hash)}.

\vspace{0.5em}

\subsubsection{Aho-Corasick (basic)}

\paragraph{Idea intuitiva.}
Construye un \textbf{trie} con todos los patrones y aniade \textbf{failure links} analogos a los de KMP: el failure link de un nodo $v$ apunta al nodo mas largo que es sufijo del string representado por $v$ y sigue siendo prefijo de algun patron. Asi, al recorrer el texto, si hay mismatch, saltamos al failure link. Esto permite buscar \textbf{todos} los patrones en $O(|T| + \sum |P_i|)$. La idea: unificar varios KMP en una sola estructura.

\paragraph{Propiedades clave (invariantes).}
\begin{itemize}\setlength\itemsep{0pt}
	\item El trie almacena los patrones.
	\item \texttt{link[v]} (failure link) se calcula en BFS en orden de profundidad creciente.
	\item \texttt{out[v]} lista los IDs de patrones que terminan en $v$.
	\item Para queries eficientes, aniadir \textbf{output links}: propagar \texttt{out} de los failure links a los hijos.
\end{itemize}

\paragraph{Codigo clave.}
BFS para construir los failure links:
\begin{verbatim}
queue<int> q;
for (int c = 0; c < SIGMA; ++c)
    if (trie[0].next[c] != -1) {
        q.push(trie[0].next[c]);
        trie[trie[0].next[c]].link = 0;
    } else trie[0].next[c] = 0;  // goto a la raiz
while (!q.empty()) {
    int v = q.front(); q.pop();
    for (int pid : trie[trie[v].link].out)
        trie[v].out.push_back(pid);
    for (int c = 0; c < SIGMA; ++c) {
        if (trie[v].next[c] != -1) {
            trie[trie[v].next[c]].link = trie[trie[v].link].next[c];
            q.push(trie[v].next[c]);
        } else trie[v].next[c] = trie[trie[v].link].next[c];
    }
}
\end{verbatim}

\paragraph{Complejidad.}
\begin{itemize}\setlength\itemsep{0pt}
	\item Build $O(\sum |P_i| \cdot \Sigma)$, query $O(|T| + \text{ocurrencias})$.
\end{itemize}

\paragraph{Donde tocar para modificar.}
\begin{itemize}\setlength\itemsep{0pt}
	\item \textbf{Output links en build}: ya hecho en el snippet (\texttt{for(int pid : t[t[u].link].out) t[u].out.push\_back(pid);}).
	\item \textbf{Transiciones automaticas}: aniadir el goto \texttt{t[v].next[k] = t[t[v].link].next[k]} cuando el enlace directo es \texttt{-1}; asi el bucle de matching es un solo paso por caracter.
	\item \textbf{Aplicar a problemas tipo DP}: contar formas de construir strings evitando los patrones (ver el siguiente modulo).
	\item \textbf{2D Aho-Corasick}: aniadir una dimension para busqueda en matrices.
\end{itemize}

\paragraph{Trampas y errores frecuentes.}
\begin{itemize}\setlength\itemsep{0pt}
	\item La BFS debe empezar con la raiz en la cola y los hijos de la raiz con \texttt{link = 0}.
	\item Olvidar propagar \texttt{out} da ocurrencias perdidas.
	\item Si no se aniaden las transiciones automaticas (goto), el matching es mas lento.
	\item Con muchos patrones, el output puede ser enorme; aniadir contador en lugar de lista.
\end{itemize}

\paragraph{Codigo.}
Ver \texttt{cpp.tex}, seccion \textit{Aho-Corasick (basic)}.

\vspace{0.5em}

\subsubsection{Aho-Corasick DP for forbidden patterns}

\paragraph{Idea intuitiva.}
Contar strings de longitud $L$ que \textbf{evitan} un conjunto de patrones. Se modela como una DP sobre el \textbf{estado del automaton} de Aho-Corasick: $dp[v]$ = numero de formas de llegar al estado $v$ despues de cierto numero de caracteres. Transicion: por cada caracter $c$ en el alfabeto, ir al estado \texttt{next[v][c]}, pero \textbf{solo} si ese estado no tiene un patron (de lo contrario, el string contendria el patron prohibido). La idea: cada estado del automaton representa un ``prefijo conocido'' y las transiciones simulan ``agregar un caracter''.

\paragraph{Propiedades clave (invariantes).}
\begin{itemize}\setlength\itemsep{0pt}
	\item $dp$ se mantiene por nivel (longitud del string construido).
	\item Transicion: $ndp[u] = \sum_{c} dp[v]$ donde $u = next[v][c]$ y $u$ no es ``terminal'' (no tiene patron que termina ahi).
	\item Respuesta: suma de $dp[v]$ sobre todos los estados $v$ no terminales al final.
	\item Si todos los caracteres son permitidos, $|\Sigma|$ entra como factor.
\end{itemize}

\paragraph{Codigo clave.}
DP sobre el automaton:
\begin{verbatim}
vector<long long> dp(numStates, 0), ndp(numStates);
dp[0] = 1;  // empezar en la raiz
for (int step = 0; step < L; ++step) {
    fill(ndp.begin(), ndp.end(), 0);
    for (int v = 0; v < numStates; ++v) if (dp[v] && !isTerminal[v])
        for (int c = 0; c < SIGMA; ++c) {
            int u = next[v][c];
            if (!isTerminal[u]) ndp[u] = (ndp[u] + dp[v]) % MOD;
        }
    swap(dp, ndp);
}
long long ans = 0;
for (int v = 0; v < numStates; ++v) if (!isTerminal[v])
    ans = (ans + dp[v]) % MOD;
\end{verbatim}

\paragraph{Complejidad.}
\begin{itemize}\setlength\itemsep{0pt}
	\item $O(L \cdot |S| \cdot |\Sigma|)$ donde $|S|$ es el numero de estados y $|\Sigma|$ el alfabeto.
	\item Para $|\Sigma| = 2$ y $|S| \le 10^4$: $\sim 2 \cdot 10^{10}$ para $L = 10^6$ (puede ser lento).
\end{itemize}

\paragraph{Donde tocar para modificar.}
\begin{itemize}\setlength\itemsep{0pt}
	\item \textbf{Alfabeto mas pequeno}: si el problema es solo $\{0, 1\}$, $|\Sigma| = 2$.
	\item \textbf{Permitir ocurrencias parciales}: cambiar la condicion a ``no estado terminal'' o relajar segun el problema.
	\item \textbf{Recuperar el string}: aniadir backtracking si el problema lo pide.
	\item \textbf{Matriz de transicion}: para queries rapidas, precomputar $M^c$ con exponentiation.
\end{itemize}

\paragraph{Trampas y errores frecuentes.}
\begin{itemize}\setlength\itemsep{0pt}
	\item El snippet asume Aho-Corasick ya construido con transiciones automaticas. Sin eso, aniadir el goto en el bucle.
	\item Si $L$ es muy grande, la DP lineal es lenta; usar matrix exponentiation.
	\item Estados terminales deben excluirse de \textbf{ambos} lados (estado actual y siguiente).
\end{itemize}

\paragraph{Codigo.}
Ver \texttt{cpp.tex}, seccion \textit{Aho-Corasick DP for forbidden patterns}.

\subsection{Estructuras avanzadas}

\subsubsection{Suffix Array}

\paragraph{Idea intuitiva.}
Un \textbf{suffix array} $SA$ es un arreglo que contiene las posiciones de inicio de los sufijos del string $S$, ordenados lexicograficamente. Se construye en $O(n \log n)$ con doubling: ordenar por el primer caracter, luego por los primeros 2, luego 4, etc. Para hacer cada paso lineal, se usa \textbf{counting sort} sobre los ranks. La idea: cada nivel dobla la cantidad de caracteres considerados, asi que en $\log n$ niveles los sufijos estan completamente ordenados.

\paragraph{Propiedades clave (invariantes).}
\begin{itemize}\setlength\itemsep{0pt}
	\item \texttt{p[i]} = posicion de inicio del $i$-esimo sufijo (ordenado).
	\item \texttt{c[i]} = rank del sufijo que empieza en $i$.
	\item Cada iteracion duplica la longitud de la ``ventana'' considerada.
	\item El algoritmo termina cuando $2^k \ge n$.
	\item La permutacion \texttt{p} es estable entre iteraciones.
\end{itemize}

\paragraph{Codigo clave.}
El paso de doubling (clasificar por los primeros $2^k$ caracteres):
\begin{verbatim}
for (int k = 0; (1 << k) < n; ++k) {
    for (int i = 0; i < n; ++i) pn[i] = (p[i] - (1 << k) + n) % n;
    // counting sort por c[pn[i]] segundo, c[i] primero
    // ...
    cn[pn[0]] = 0;
    for (int i = 1; i < n; ++i)
        cn[pn[i]] = cn[pn[i-1]] + (pair != pair_prev);
    c = cn;
}
\end{verbatim}

\paragraph{Complejidad.}
\begin{itemize}\setlength\itemsep{0pt}
	\item $O(n \log n)$ con counting sort.
	\item Constante pequena; en la practica eficiente.
\end{itemize}

\paragraph{Donde tocar para modificar.}
\begin{itemize}\setlength\itemsep{0pt}
	\item \textbf{Suffix array en $O(n)$}: usar SA-IS (mas complejo, no incluido).
	\item \textbf{Buscar patron}: con SA + LCP es $O(|P| \log n)$ por busqueda binaria.
	\item \textbf{Suffix array de varios strings}: aniadir centinelas (\texttt{\$}) entre ellos.
	\item \textbf{Aniadir longitud}: aniadir \texttt{n - p[i]} como clave secundaria (sufijos mas cortos primero).
\end{itemize}

\paragraph{Trampas y errores frecuentes.}
\begin{itemize}\setlength\itemsep{0pt}
	\item La ``ventana'' $2^k$ se hace circular: \texttt{p[i] = (p[i] - (1<<k) + n) \% n}.
	\item Para contar bien los ranks duplicados, comparar pares \texttt{(c[i], c[(i+2\^{}k) \% n])}.
	\item El counting sort necesita un array de acumulacion de tamano $\le n$.
	\item Si los caracteres son negativos o fuera del alfabeto esperado, normalizar antes.
\end{itemize}

\paragraph{Codigo.}
Ver \texttt{cpp.tex}, seccion \textit{Suffix Array}.

\vspace{0.5em}

\subsubsection{LCP Array (Kasai)}

\paragraph{Idea intuitiva.}
El \textbf{LCP array} (Longest Common Prefix) guarda, para cada par de sufijos adyacentes en el SA, el largo de su prefijo comun. El algoritmo de \textbf{Kasai} lo calcula en $O(n)$ iterando los sufijos \textbf{en orden del string original} (no del SA), y reutiliza el LCP previo. La demostracion: el sufijo en $i+1$ se obtiene del sufijo en $i$ ``avanzando un caracter''; los LCPs decrecen como maximo 1 entre iteraciones.

\paragraph{Propiedades clave (invariantes).}
\begin{itemize}\setlength\itemsep{0pt}
	\item \texttt{lcp[i]} = LCP entre los sufijos en \texttt{sa[i]} y \texttt{sa[i-1]}.
	\item Se necesita \texttt{rank\_[i]} (la inversa del SA) para iterar.
	\item El truco de Kasai: si \texttt{lcp > 0} para el sufijo anterior, el actual parte de al menos \texttt{lcp - 1}.
	\item La suma de todos los LCPs da informacion sobre repeticiones.
\end{itemize}

\paragraph{Codigo clave.}
Algoritmo de Kasai:
\begin{verbatim}
vector<int> rank(n);
for (int i = 0; i < n; ++i) rank[sa[i]] = i;
int k = 0;
vector<int> lcp(n - 1);
for (int i = 0; i < n; ++i) {
    if (rank[i] == n - 1) { k = 0; continue; }
    int j = sa[rank[i] + 1];
    while (i + k < n && j + k < n && s[i+k] == s[j+k]) ++k;
    lcp[rank[i]] = k;
    if (k) --k;
}
\end{verbatim}

\paragraph{Complejidad.}
\begin{itemize}\setlength\itemsep{0pt}
	\item $O(n)$.
	\item Cada comparacion reduce $k$ o avanza linealmente; amortizado $O(1)$.
\end{itemize}

\paragraph{Donde tocar para modificar.}
\begin{itemize}\setlength\itemsep{0pt}
	\item \textbf{RMQ sobre LCP}: para queries ``LCP de dos sufijos'', encontrar sus ranks y hacer RMQ en \texttt{lcp}.
	\item \textbf{Contar substrings distintas}: $\sum (n - sa[i] - lcp[i])$.
	\item \textbf{Longest common substring}: maximo en \texttt{lcp}.
	\item \textbf{String matching}: busqueda binaria sobre SA + LCP verification.
\end{itemize}

\paragraph{Trampas y errores frecuentes.}
\begin{itemize}\setlength\itemsep{0pt}
	\item Si dos sufijos son iguales (raro, solo si el string tiene period $< n$), \texttt{lcp} puede ser \texttt{n - sa[i]}.
	\item En la iteracion, aniadir el chequeo \texttt{i + k < n} para evitar overflow.
	\item El \texttt{--k} en la salida es crucial para la complejidad amortizada.
\end{itemize}

\paragraph{Codigo.}
Ver \texttt{cpp.tex}, seccion \textit{LCP Array (Kasai)}.

\vspace{0.5em}

\subsubsection{Suffix Automaton}

\paragraph{Idea intuitiva.}
Un \textbf{suffix automaton} es un DFA minimo que reconoce todos los substrings del string original. Tiene la propiedad magica de tener $\le 2n - 1$ estados. Cada estado representa una clase de equivalencia de substrings (los que tienen el mismo conjunto de posiciones finales). El \textbf{link} de un estado apunta al estado del sufijo propio mas largo.

\paragraph{Propiedades clave (invariantes).}
\begin{itemize}\setlength\itemsep{0pt}
	\item \texttt{len[v]}: largo maximo de los substrings de la clase.
	\item \texttt{link[v]}: padre en el arbol de sufijos (suffix link).
	\item \texttt{next[v][c]}: transicion por caracter.
	\item \texttt{occ[v]}: numero de posiciones finales (cuantas veces aparece el patron).
\end{itemize}

\paragraph{Complejidad.}
\begin{itemize}\setlength\itemsep{0pt}
	\item Build $O(n)$ con la operacion de \textbf{clonar} estados.
\end{itemize}

\paragraph{Donde tocar para modificar.}
\begin{itemize}\setlength\itemsep{0pt}
	\item \textbf{Contar substrings distintas}: $\sum_{v \ne 0} (len[v] - len[link[v]])$.
	\item \textbf{Longest common substring}: navegar el SAM con el segundo string, manteniendo el match actual.
	\item \textbf{Substring count}: despues del build, ordenar estados por \texttt{len} descendente y propagar \texttt{occ}.
	\item \textbf{Endpos queries}: el SAM es el DFA canonico de los substrings; queries ``cuantas veces aparece $P$'' se responden en $O(|P|)$.
\end{itemize}

\paragraph{Trampas y errores frecuentes.}
\begin{itemize}\setlength\itemsep{0pt}
	\item El \textbf{clone} debe tener \texttt{occ = 0} (no es una nueva posicion final); el \texttt{link} del estado siendo creado y del estado clonado deben apuntar al clon.
	\item Saltarse esto da SAM incorrecto.
\end{itemize}

\paragraph{Codigo.}
Ver \texttt{cpp.tex}, seccion \textit{Suffix Automaton}.

\vspace{0.5em}

\subsubsection{Lyndon factorization (Duval)}

\paragraph{Idea intuitiva.}
Un \textbf{string de Lyndon} es un string que es estrictamente menor que todos sus rotaciones. Toda cadena tiene una \textbf{factorizacion unica} de Lyndon (algoritmo de \textbf{Duval}): una secuencia de strings de Lyndon $L_1, L_2, \dots, L_k$ tales que $L_1 \ge L_2 \ge \dots \ge L_k$ y la concatenacion es el string original. Esto da $O(n)$ para encontrar los cortes. La demostracion de correctitud: cada factor es minimal en su clase de rotacion, asi que la concatenacion respeta el orden lexicografico.

\paragraph{Propiedades clave (invariantes).}
\begin{itemize}\setlength\itemsep{0pt}
	\item Cada factor de Lyndon es de la forma $uv$ donde $u$ se repite y $v$ es prefijo de $u$.
	\item La longitud del factor esta dada por la ``longitud del periodo''.
	\item Aplicacion: el ``Lyndon word'' es el menor en orden lexicografico entre todas las rotaciones.
	\item La factorizacion es \textbf{unica}.
\end{itemize}

\paragraph{Codigo clave.}
El algoritmo de Duval (iterativo):
\begin{verbatim}
int n = s.size();
int i = 0;
while (i < n) {
    int j = i + 1, k = i;
    while (j < n && s[k] <= s[j]) {
        if (s[k] < s[j]) k = i;
        else ++k;
        ++j;
    }
    int len = j - k;
    while (i < j) { cout << "factor: " << s.substr(i, len) << "\n"; i += len; }
}
\end{verbatim}

\paragraph{Complejidad.}
\begin{itemize}\setlength\itemsep{0pt}
	\item $O(n)$.
	\item Cada comparacion es $O(1)$; el puntero $k$ retrocede a lo sumo $n$ veces en total.
\end{itemize}

\paragraph{Donde tocar para modificar.}
\begin{itemize}\setlength\itemsep{0pt}
	\item \textbf{Encontrar la menor rotacion}: el menor sufijo lexicografico de \texttt{s + s} es la menor rotacion; la posicion se obtiene con un solo pase de Duval.
	\item \textbf{Algoritmo de Booth}: alternativa para rotacion minima, tambien $O(n)$.
	\item \textbf{Validar}: si necesitas verificar que un string es de Lyndon, compararlo con sus rotaciones.
	\item \textbf{Encontrar todos los factores}: guardar la posicion y longitud de cada uno.
\end{itemize}

\paragraph{Trampas y errores frecuentes.}
\begin{itemize}\setlength\itemsep{0pt}
	\item La condicion \texttt{s[j] <= s[k]} (no \texttt{<}) en el while de Duval es crucial; cambiar a \texttt{<} da comportamiento distinto.
	\item El caso $s[k] < s[j]$ resetea $k$ a $i$ (inicio del factor actual).
	\item El bucle externo debe avanzar $i$ por la longitud del factor, no por 1.
\end{itemize}

\paragraph{Codigo.}
Ver \texttt{cpp.tex}, seccion \textit{Lyndon factorization (Duval)}.

\subsection{Palindromos}

\subsubsection{Manacher}

\paragraph{Idea intuitiva.}
Calcula en $O(n)$ los \textbf{radios} de todos los palindromos centrados en cada posicion (impares y pares). La idea: mantener una ventana $[l, r]$ que es el palindromo centrado en un punto anterior mas largo encontrado. Si el nuevo centro $i$ cae dentro de $[l, r]$, entonces el palindromo centrado en $i$ tiene al menos el radio del palindromo ``espejo'' $l + r - i$; sino, empezar desde 1 (impar) o 0 (par). Despues, expandir lo que falte. La demostracion de $O(n)$: cada caracter se compara a lo sumo 2 veces (extension + limite por ventana).

\paragraph{Propiedades clave (invariantes).}
\begin{itemize}\setlength\itemsep{0pt}
	\item Conveccion del snippet: \texttt{odd[i]} = $k$ tal que el palindromo impar mas largo centrado en $i$ tiene longitud $2k - 1$. \texttt{even[i]} = $k$ para par centrado entre $i-1$ e $i$, longitud $2k$.
	\item $k$ cuenta tambien cuantos palindromos hay centrados ahi.
	\item Las queries \texttt{isPalindrome(l, r)} son $O(1)$.
	\item La ventana $[l, r]$ siempre contiene el palindromo mas largo encontrado hasta ahora.
\end{itemize}

\paragraph{Codigo clave.}
Manacher para palindromos impares:
\begin{verbatim}
int l = 0, r = -1;
for (int i = 0; i < n; ++i) {
    int k = (i > r) ? 1 : min(odd[l + r - i], r - i + 1);
    while (i - k >= 0 && i + k < n && s[i-k] == s[i+k]) ++k;
    odd[i] = k;
    if (i + k - 1 > r) { l = i - k + 1; r = i + k - 1; }
}
\end{verbatim}

\paragraph{Complejidad.}
\begin{itemize}\setlength\itemsep{0pt}
	\item Construccion $O(n)$, queries $O(1)$.
	\item Cada extension se cuenta una sola vez.
\end{itemize}

\paragraph{Donde tocar para modificar.}
\begin{itemize}\setlength\itemsep{0pt}
	\item \textbf{Solo contar palindromos}: ya implementado en \texttt{countPalindromes}.
	\item \textbf{Longest palindromic substring}: ya en \texttt{longestPalindrome}.
	\item \textbf{Modificar la conveccion de radios}: si queres ``radio = mitad de la longitud'' en vez de ``k palindromos'', ajustar el offset.
	\item \textbf{Queries con k no precomputado}: aniadir helpers especificos.
	\item \textbf{Modificar el caracter de relleno}: para multiples separadores, aniadir un caracter unico.
\end{itemize}

\paragraph{Trampas y errores frecuentes.}
\begin{itemize}\setlength\itemsep{0pt}
	\item La conveccion del snippet es \textbf{distinta} de la ``clasica de radio'' (donde radio = $(L-1)/2$); tener cuidado al copiar a otro codigo.
	\item Si los palindromos se solapan, la condicion \texttt{i + k - 1 > r} actualiza la ventana.
	\item En palindromos pares, el centro esta ``entre'' dos indices; aniadir offset en queries.
\end{itemize}

\paragraph{Codigo.}
Ver \texttt{cpp.tex}, seccion \textit{Manacher}.

% ============================================================
\section{Grafos}
% ============================================================

\subsection{Caminos minimos}

\subsubsection{Dijkstra}

\paragraph{Idea intuitiva.}
Encuentra el shortest path desde un origen $s$ a todos los vertices en $O((V + E) \log V)$ asumiendo \textbf{pesos no negativos}. La idea: usar una cola de prioridad (min-heap) ordenada por distancia tentativa. En cada paso, extraer el nodo con menor distancia; si ya fue procesado, saltar. Si al actualizar un vecino la distancia mejora, insertar/actualizar en la cola. La demostracion: cuando un nodo se extrae por primera vez, no hay forma de llegar a el con menor distancia (cualquier camino alternativo deberia pasar por nodos no procesados, con distancia $\ge$ la actual).

\paragraph{Propiedades clave (invariantes).}
\begin{itemize}\setlength\itemsep{0pt}
	\item Asume pesos $\ge 0$; con pesos negativos usar Bellman-Ford.
	\item Greedy correctness: la primera vez que se extrae un nodo, su distancia es la definitiva.
	\item Si se quiere reconstruir el camino, guardar el \texttt{parent[]}.
	\item Cada nodo se procesa a lo sumo $V$ veces (mas relajarions, pero solo $V$ definitivas).
\end{itemize}

\paragraph{Codigo clave.}
El bucle principal de Dijkstra:
\begin{verbatim}
priority_queue<pair<long long, int>, vector<...>, greater<...>> pq;
pq.push({0, source});
while (!pq.empty()) {
    auto [d, v] = pq.top(); pq.pop();
    if (d > dist[v]) continue;  // entrada vieja
    for (auto [w, u] : adj[v])
        if (dist[v] + w < dist[u]) {
            dist[u] = dist[v] + w;
            pq.push({dist[u], u});
        }
}
\end{verbatim}

\paragraph{Complejidad.}
\begin{itemize}\setlength\itemsep{0pt}
	\item $O((V + E) \log V)$ con heap binario, $O(V + E)$ con heap de Fibonacci.
	\item En la practica, $\sim 1.5 \cdot 10^{-7}$ segundos por operacion para $V = 10^5$.
\end{itemize}

\paragraph{Donde tocar para modificar.}
\begin{itemize}\setlength\itemsep{0pt}
	\item \textbf{Camino mas corto entre todos los pares}: si se necesita desde todos los origenes, usar Floyd o ejecutar Dijkstra $V$ veces.
	\item \textbf{Shortest path con aristas negativas limitadas}: Bellman-Ford.
	\item \textbf{Reconstruir camino}: aniadir \texttt{parent[v]} actualizado junto a \texttt{dist[v]}.
	\item \textbf{Dijkstra con potentials (Johnson)}: para aristas negativas si se puede reescalar.
	\item \textbf{Multi-source}: aniadir varios nodos fuentes iniciales al heap.
\end{itemize}

\paragraph{Trampas y errores frecuentes.}
\begin{itemize}\setlength\itemsep{0pt}
	\item Distancias no inicializadas a \texttt{INF}; el snippet usa \texttt{1e18}.
	\item El \texttt{if(dist[adjN] != 1e18)} para borrar de la cola es importante si usas \texttt{set} (no asi con \texttt{priority\_queue}).
	\item Con pesos grandes, \texttt{dist[v] + w} puede overflow; usar \texttt{LLONG\_MAX/2} como INF.
	\item No usar con pesos negativos: el algoritmo puede dar respuestas incorrectas.
\end{itemize}

\paragraph{Codigo.}
Ver \texttt{cpp.tex}, seccion \textit{Dijkstra}.

\vspace{0.5em}

\subsubsection{Bellman-Ford}

\paragraph{Idea intuitiva.}
Encuentra shortest path desde $s$ incluso con \textbf{pesos negativos} en $O(VE)$. La idea: relajar todas las aristas $V - 1$ veces. Cada relajacion propaga la mejor distancia a ``un paso mas''. Despues de $V - 1$ pasadas, si alguna arista aun se puede relajar, hay un \textbf{ciclo negativo} alcanzable. La demostracion: el camino mas corto tiene a lo sumo $V - 1$ aristas (sin ciclos); tras $V - 1$ relajaciones todas las distancias estan en su valor optimo.

\paragraph{Propiedades clave (invariantes).}
\begin{itemize}\setlength\itemsep{0pt}
	\item Funciona con pesos negativos (sin ciclos negativos).
	\item Tras $V - 1$ relajaciones, las distancias son definitivas (si no hay ciclo negativo).
	\item Una relajacion adicional detecta ciclos negativos: si \texttt{dist[u] + w < dist[v]}, hay ciclo negativo alcanzable desde $s$.
	\item Cada relajacion solo mejora distancias, asi que termina (no oscila).
\end{itemize}

\paragraph{Codigo clave.}
El doble bucle de relajacion:
\begin{verbatim}
for (int i = 0; i < n - 1; ++i)
    for (auto [u, v, w] : edges)
        if (dist[u] + w < dist[v])
            dist[v] = dist[u] + w;
// Detectar ciclo negativo
for (auto [u, v, w] : edges)
    if (dist[u] + w < dist[v]) { /* ciclo negativo */ }
\end{verbatim}

\paragraph{Complejidad.}
\begin{itemize}\setlength\itemsep{0pt}
	\item $O(VE)$.
	\item SPFA (con cola) en promedio es mucho mas rapido, peor caso igual.
\end{itemize}

\paragraph{Donde tocar para modificar.}
\begin{itemize}\setlength\itemsep{0pt}
	\item \textbf{SPFA}: optimizacion que relaja solo nodos cuya distancia cambio; en el peor caso sigue siendo $O(VE)$.
	\item \textbf{Detectar ciclo negativo global}: ejecutar Bellman-Ford desde un super-origen conectado a todos.
	\item \textbf{Camino mas corto con K aristas exactas}: $K$ iteraciones de Bellman-Ford.
	\item \textbf{Johnson's}: reescalar pesos con potentials para usar Dijkstra.
\end{itemize}

\paragraph{Trampas y errores frecuentes.}
\begin{itemize}\setlength\itemsep{0pt}
	\item Overflow si los pesos son $\ge 2^{52}$; usar \texttt{long long} y \texttt{\_\_int128} si hace falta.
	\item Si el grafo es denso ($E \approx V^2$), $O(VE)$ es prohibitivo.
	\item El chequeo de ciclo negativo debe hacerse tras las $V - 1$ iteraciones, no durante.
	\item Distancias no inicializadas a \texttt{INF}; el primer nodo (source) tiene distancia 0.
\end{itemize}

\paragraph{Codigo.}
Ver \texttt{cpp.tex}, seccion \textit{Bellman-Ford}.

\vspace{0.5em}

\subsubsection{Floyd-Warshall}

\paragraph{Idea intuitiva.}
Shortest path entre \textbf{todos los pares} de vertices en $O(V^3)$. DP: $d[i][j]$ = shortest path usando solo nodos intermedios en $\{0, \dots, k\}$. Transicion: $d_{k+1}[i][j] = \min(d_k[i][j], d_k[i][k+1] + d_k[k+1][j])$. La demostracion: tras $k$ iteraciones, $d[i][j]$ es el shortest path con intermedios en $\{0, \dots, k-1\}$; aniadir el nodo $k$ da la recurrencia clasica.

\paragraph{Propiedades clave (invariantes).}
\begin{itemize}\setlength\itemsep{0pt}
	\item Detecta ciclos negativos: tras el algoritmo, $d[i][i] < 0$ indica ciclo negativo alcanzable desde $i$.
	\item Funciona con pesos negativos (sin ciclos negativos).
	\item Memoria $O(V^2)$; con $V \le 400$ es OK, con mas conviene Dijkstra $V$ veces.
	\item Es la unica opcion cuando se quieren shortest paths entre todos los pares y el grafo es denso.
\end{itemize}

\paragraph{Codigo clave.}
El triple bucle clasico:
\begin{verbatim}
for (int k = 0; k < n; ++k)
    for (int i = 0; i < n; ++i)
        for (int j = 0; j < n; ++j)
            d[i][j] = min(d[i][j], d[i][k] + d[k][j]);
// detectar ciclo negativo
for (int i = 0; i < n; ++i) if (d[i][i] < 0) tiene_ciclo_negativo = true;
\end{verbatim}

\paragraph{Complejidad.}
\begin{itemize}\setlength\itemsep{0pt}
	\item $O(V^3)$ tiempo, $O(V^2)$ memoria.
	\item Para $V \le 400$: $\sim 6.4 \cdot 10^7$ operaciones (factible).
\end{itemize}

\paragraph{Donde tocar para modificar.}
\begin{itemize}\setlength\itemsep{0pt}
	\item \textbf{Reconstruir camino}: aniadir \texttt{next[i][j]} y backtrack.
	\item \textbf{Transitividad}: el algoritmo tambien calcula transitividad (\texttt{d[i][j] != INF} significa alcanzable).
	\item \textbf{Pesos doubles}: aniadir tolerancia (\texttt{d[i][j] > d[i][k] + d[k][j] - eps}).
	\item \textbf{Grafo no denso}: usar Dijkstra $V$ veces en vez de Floyd.
\end{itemize}

\paragraph{Trampas y errores frecuentes.}
\begin{itemize}\setlength\itemsep{0pt}
	\item Para detectar ciclo negativo correctamente, aniadir un chequeo \texttt{if (ans[i][i] < 0)} tras el triple bucle.
	\item El orden de los bucles \texttt{k, i, j} es importante; invertir da resultados diferentes.
	\item Con \texttt{double}, usar tolerancia \texttt{< eps} en vez de \texttt{<} para evitar loops infinitos.
\end{itemize}

\paragraph{Codigo.}
Ver \texttt{cpp.tex}, seccion \textit{Floyd-Warshall}.

\subsection{Componentes y cortes}

\subsubsection{Kosaraju SCC}

\paragraph{Idea intuitiva.}
Encuentra las \textbf{componentes fuertemente conexas} (SCCs) en $O(V + E)$. Algoritmo en 3 pasos:
\begin{itemize}\setlength\itemsep{0pt}
	\item DFS en $G$ rellenando una pila en \textbf{post-orden} (los nodos que terminan tarde quedan arriba).
	\item Construir $G^T$ (grafo transpuesto).
	\item DFS en $G^T$ procesando nodos en orden de la pila (de mayor a menor tiempo); cada DFS da una SCC.
\end{itemize}
La demostracion: dos nodos $u, v$ estan en la misma SCC sii hay camino $u \leadsto v$ y $v \leadsto u$; Kosaraju captura exactamente esto al explorar $G^T$ desde los nodos con mayor tiempo de finalizacion en $G$.

\paragraph{Propiedades clave (invariantes).}
\begin{itemize}\setlength\itemsep{0pt}
	\item Cada SCC es maximal fuertemente conexa.
	\item El \textbf{grafo de SCCs} (metagraf) es un DAG.
	\item Kosaraju y Tarjan dan el mismo resultado; Tarjan es ``mas rapido'' en practica (un solo DFS).
	\item Cada nodo pertenece a exactamente una SCC.
\end{itemize}

\paragraph{Codigo clave.}
La estructura de Kosaraju:
\begin{verbatim}
vector<int> order;
vector<bool> used(n);
function<void(int)> dfs1 = [&](int v) {
    used[v] = true;
    for (int u : g[v]) if (!used[u]) dfs1(u);
    order.push_back(v);  // post-orden
};
function<void(int, int)> dfs2 = [&](int v, int root) {
    comp[v] = root;
    for (int u : gt[v]) if (comp[u] == -1) dfs2(u, root);
};
// 1: DFS en G para llenar order
// 2: DFS en G^T en orden inverso para asignar SCCs
\end{verbatim}

\paragraph{Complejidad.}
\begin{itemize}\setlength\itemsep{0pt}
	\item $O(V + E)$.
	\item Constante 2 (dos DFS + el grafo transpuesto).
\end{itemize}

\paragraph{Donde tocar para modificar.}
\begin{itemize}\setlength\itemsep{0pt}
	\item \textbf{2-SAT}: SCC sobre grafo de implicaciones.
	\item \textbf{Construir el DAG de SCCs}: aniadir aristas entre SCCs diferentes.
	\item \textbf{Tamano de cada SCC}: aniadir contador durante el segundo DFS.
	\item \textbf{Grafo grande}: usar Tarjan para evitar construir $G^T$.
\end{itemize}

\paragraph{Trampas y errores frecuentes.}
\begin{itemize}\setlength\itemsep{0pt}
	\item El snippet usa \texttt{vector<...> adj[n]} (VLA); reemplazable por \texttt{vector<vector<pair<int,int>>} para portabilidad.
	\item El orden del segundo DFS debe ser \textbf{inverso} al de finalizacion del primero.
	\item Si el grafo no es conexo, aniadir un loop sobre todos los nodos en el primer DFS.
	\item Confundir SCC con componentes conexas (las primeras son para dirigidos, las segundas para no dirigidos).
\end{itemize}

\paragraph{Codigo.}
Ver \texttt{cpp.tex}, seccion \textit{Kosaraju SCC}.

\vspace{0.5em}

\subsubsection{Tarjan (puentes)}

\paragraph{Idea intuitiva.}
Encuentra los \textbf{puentes} (bridges): aristas cuya remocion \textbf{desconecta} el grafo. Usa un solo DFS manteniendo \texttt{tin[v]} (tiempo de descubrimiento) y \texttt{low[v]} (el menor \texttt{tin} alcanzable desde $v$ via back-edges). Una arista $v - u$ es puente si $low[u] > tin[v]$. La demostracion: si $low[u] > tin[v]$, entonces $u$ (y su subarbol) no pueden ``escapar'' de $v$; quitar la arista $v - u$ parte el grafo.

\paragraph{Propiedades clave (invariantes).}
\begin{itemize}\setlength\itemsep{0pt}
	\item \texttt{low[v] = min(tin[v], tin[w])} para todo back-edge $(v, w)$, y \texttt{min(low[u])} para hijos en el DFS.
	\item Solo se consideran las aristas del \textbf{arbol DFS} (no back-edges) para \texttt{low}.
	\item Cada nodo tiene exactamente un \texttt{tin}; cada back-edge se procesa una vez.
\end{itemize}

\paragraph{Codigo clave.}
El DFS de Tarjan para bridges:
\begin{verbatim}
void dfs(int v, int pEdge) {
    used[v] = true;
    tin[v] = low[v] = timer++;
    for (auto [u, id] : adj[v]) {
        if (id == pEdge) continue;  // no contar el edge al padre
        if (used[u]) {
            low[v] = min(low[v], tin[u]);  // back-edge
        } else {
            dfs(u, id);
            low[v] = min(low[v], low[u]);
            if (low[u] > tin[v]) bridges.push_back(id);  // es puente
        }
    }
}
\end{verbatim}

\paragraph{Complejidad.}
\begin{itemize}\setlength\itemsep{0pt}
	\item $O(V + E)$.
	\item Una sola pasada de DFS.
\end{itemize}

\paragraph{Donde tocar para modificar.}
\begin{itemize}\setlength\itemsep{0pt}
	\item \textbf{Bridge tree}: contractar cada 2-edge-connected component a un nodo; el resultado es un arbol.
	\item \textbf{Articulation points}: variante con condicion diferente (ver el siguiente modulo).
	\item \textbf{Grafo no dirigido}: solo funciona en no dirigidos; en dirigidos hay conceptos analogos (strong bridges, dominator tree).
	\item \textbf{Reconstruir componentes 2-edge-connected}: BFS desde cada nodo excluyendo bridges.
\end{itemize}

\paragraph{Trampas y errores frecuentes.}
\begin{itemize}\setlength\itemsep{0pt}
	\item No confundir el parent con un back-edge al padre; chequear \texttt{if (adjN == parent) continue;} y contar solo las veces que se ``visita'' como hijo.
	\item Para multigrafos, los edges deben tener IDs unicos para distinguirlos.
	\item Si el grafo no es conexo, aniadir loop externo sobre todos los nodos.
\end{itemize}

\paragraph{Codigo.}
Ver \texttt{cpp.tex}, seccion \textit{Tarjan (puentes)}.

\vspace{0.5em}

\subsubsection{Tarjan (puntos de articulacion)}

\paragraph{Idea intuitiva.}
Encuentra los \textbf{puntos de articulacion} (cut vertices): vertices cuya remocion \textbf{desconecta} el grafo. Mismo setup que para bridges: \texttt{tin[v]}, \texttt{low[v]}. Un nodo $v$ (no raiz) es de articulacion si existe un hijo $u$ con $low[u] \ge tin[v]$. La \textbf{raiz} es de articulacion si tiene $\ge 2$ hijos en el DFS tree. La razon: para $v$ no raiz, la condicion $\ge$ indica que el subarbol de $u$ depende de $v$; quitar $v$ los desconecta. Para la raiz, perder $\ge 2$ hijos parte el DFS en $\ge 2$ componentes.

\paragraph{Propiedades clave (invariantes).}
\begin{itemize}\setlength\itemsep{0pt}
	\item $low[u] \ge tin[v]$ significa que el subarbol de $u$ no tiene back-edge ``por encima'' de $v$.
	\item La raiz es caso especial: tiene que perder mas de un hijo para que el grafo se parta.
	\item Cada nodo se evalua una sola vez.
\end{itemize}

\paragraph{Codigo clave.}
La condicion de articulacion:
\begin{verbatim}
void dfs(int v, int p) {
    used[v] = true;
    tin[v] = low[v] = timer++;
    int children = 0;
    for (int u : adj[v]) {
        if (u == p) continue;
        if (used[u]) low[v] = min(low[v], tin[u]);
        else {
            dfs(u, v);
            low[v] = min(low[v], low[u]);
            if (low[u] >= tin[v] && p != -1)  // no raiz
                isArt[v] = true;
            ++children;
        }
    }
    if (p == -1 && children > 1) isArt[v] = true;  // raiz especial
}
\end{verbatim}

\paragraph{Complejidad.}
\begin{itemize}\setlength\itemsep{0pt}
	\item $O(V + E)$.
	\item Una sola pasada de DFS.
\end{itemize}

\paragraph{Donde tocar para modificar.}
\begin{itemize}\setlength\itemsep{0pt}
	\item \textbf{Block-cut tree}: estructura bipartita con BCCs (2-vertex-connected components) y articulation points.
	\item \textbf{Verificar biconectividad}: chequear que no hay puntos de articulacion.
	\item \textbf{Componentes biconnected}: extender Tarjan para enumerar BCCs.
\end{itemize}

\paragraph{Trampas y errores frecuentes.}
\begin{itemize}\setlength\itemsep{0pt}
	\item No olvidar el caso especial de la raiz.
	\item La condicion usa \texttt{>=}, no \texttt{>}; un hijo con \texttt{low[u] == tin[v]} tambien cuenta.
	\item Multigrafos pueden tener ``falsos'' articulation points; verificar caso a caso.
\end{itemize}

\paragraph{Codigo.}
Ver \texttt{cpp.tex}, seccion \textit{Tarjan (puntos de articulacion)}.

\vspace{0.5em}

\subsubsection{Tarjan SCC}

\paragraph{Idea intuitiva.}
Variante del SCC que usa \textbf{un solo DFS} con una pila explicita. Mantiene \texttt{disc[v]}, \texttt{low[v]}, y una pila de nodos ``en la SCC actual''. Cuando \texttt{low[v] == disc[v]}, $v$ es la raiz de una SCC; popear hasta $v$ inclusive. La demostracion: $low[v] == disc[v]$ indica que $v$ no tiene back-edges a ancestros; entonces $v$ y los nodos en la pila hasta $v$ forman una SCC maximal.

\paragraph{Propiedades clave (invariantes).}
\begin{itemize}\setlength\itemsep{0pt}
	\item Mas eficiente que Kosaraju (un solo DFS).
	\item \texttt{low[v] = min(low[u], disc[w])} para back-edges o recursion.
	\item La pila asegura que solo los nodos en la SCC actual se popean.
	\item Cada nodo se pushea y popea a lo sumo una vez.
\end{itemize}

\paragraph{Codigo clave.}
El DFS de Tarjan SCC:
\begin{verbatim}
void dfs(int v) {
    disc[v] = low[v] = timer++;
    st.push(v); inStack[v] = true;
    for (int u : g[v]) {
        if (disc[u] == -1) { dfs(u); low[v] = min(low[v], low[u]); }
        else if (inStack[u]) low[v] = min(low[v], disc[u]);
    }
    if (low[v] == disc[v]) {
        // popear hasta v
        while (true) {
            int u = st.top(); st.pop(); inStack[u] = false;
            comp[u] = current_scc;
            if (u == v) break;
        }
        ++current_scc;
    }
}
\end{verbatim}

\paragraph{Complejidad.}
\begin{itemize}\setlength\itemsep{0pt}
	\item $O(V + E)$.
	\item Constante similar a Kosaraju pero sin construir $G^T$.
\end{itemize}

\paragraph{Donde tocar para modificar.}
\begin{itemize}\setlength\itemsep{0pt}
	\item \textbf{Tamano de cada SCC}: aniadir contador durante el pop.
	\item \textbf{Grafo condensation}: aniadir aristas entre SCCs en el DAG resultante.
	\item \textbf{2-SAT}: aniadir la construccion del grafo de implicaciones.
\end{itemize}

\paragraph{Trampas y errores frecuentes.}
\begin{itemize}\setlength\itemsep{0pt}
	\item \texttt{inStack[u]} debe chequearse antes de usar \texttt{low[v] = min(low[v], disc[u])} en back-edges.
	\item La pila global es importante; sin ella no se puede extraer la SCC.
	\item Si el grafo es muy grande, recursion puede overflow; aniadir version iterativa.
\end{itemize}

\paragraph{Codigo.}
Ver \texttt{cpp.tex}, seccion \textit{Tarjan SCC}.

\subsection{Matching}

\subsubsection{Kuhn (bipartite matching)}

\paragraph{Idea intuitiva.}
Encuentra un \textbf{matching maximo} en un grafo bipartito en $O(VE)$. Para cada nodo $u$ en $L$, intentar encontrar un vecino $v$ en $R$ que este libre o cuyo match se pueda ``reorganizar'' recursivamente. Si se encuentra, matchear $u - v$. La idea: si $u$ quiere un vecino $v$ ya ocupado, intentamos ``mover'' el match de $v$ a otro nodo; recursivamente, esto encuentra un augmenting path si existe.

\paragraph{Propiedades clave (invariantes).}
\begin{itemize}\setlength\itemsep{0pt}
	\item \texttt{matchR[v]} = el nodo de $L$ matcheado a $v$ (o $-1$ si libre).
	\item \texttt{vis[]} se resetea \textbf{cada intento} de $u$.
	\item Si el matching es perfecto, $|L| = |R| = $ tamano del matching.
	\item Kuhn no garantiza maximalidad greedy; puede haber augmenting paths no encontrados.
\end{itemize}

\paragraph{Codigo clave.}
La recursion del augmenting path:
\begin{verbatim}
bool tryKuhn(int v) {
    if (vis[v]) return false;
    vis[v] = true;
    for (int u : g[v]) {
        if (matchR[u] == -1 || tryKuhn(matchR[u])) {
            matchR[u] = v;
            return true;
        }
    }
    return false;
}
// Por cada v en L:
fill(vis, vis+n, false);
if (tryKuhn(v)) ++matches;
\end{verbatim}

\paragraph{Complejidad.}
\begin{itemize}\setlength\itemsep{0pt}
	\item $O(VE)$ en el peor caso.
	\item Con Hopcroft-Karp: $O(\sqrt{V} \cdot E)$, mucho mejor para grafos densos.
\end{itemize}

\paragraph{Donde tocar para modificar.}
\begin{itemize}\setlength\itemsep{0pt}
	\item \textbf{Matching bipartito minimo (min vertex cover)}: Konig's theorem.
	\item \textbf{Matching general}: blossom algorithm (no incluido).
	\item \textbf{Aniadir capacidades}: si los nodos de $R$ tienen capacidad $> 1$, partir cada uno en capacidad copias.
	\item \textbf{Output del matching}: el array \texttt{matchR} da los pares finales.
\end{itemize}

\paragraph{Trampas y errores frecuentes.}
\begin{itemize}\setlength\itemsep{0pt}
	\item Olvidar resetear \texttt{vis} por intento; eso es lo que evita ``ciclos infinitos'' en la recursion.
	\item Para grafos grandes ($V > 10^4$), Kuhn puede ser muy lento; usar Hopcroft-Karp.
	\item Si el matching es maximo pero no perfecto, $|L|$ o $|R|$ puede ser mayor; el matching es solo maximo.
\end{itemize}

\paragraph{Codigo.}
Ver \texttt{cpp.tex}, seccion \textit{Kuhn (bipartite matching)}.

\vspace{0.5em}

\subsubsection{Hopcroft-Karp}

\paragraph{Idea intuitiva.}
Matching bipartito maximo en $O(\sqrt{V} E)$. Algoritmo en \textbf{fases}: en cada fase, hacer BFS para encontrar los \textbf{caminos de aumento mas cortos} desde nodos libres de $L$; luego DFS solo por esos caminos. Cada fase aumenta el matching en $\ge 1$ y toma $O(E)$. La demostracion: tras una fase, todos los augmenting paths tienen longitud $\ge$ la nueva minima; eso sucede cada $\sqrt{V}$ fases (longitud minima crece geometricamente).

\paragraph{Propiedades clave (invariantes).}
\begin{itemize}\setlength\itemsep{0pt}
	\item \texttt{dist[u]} = distancia en niveles desde un nodo libre de $L$.
	\item BFS termina cuando no hay mas camino a un nodo libre de $R$.
	\item DFS respeta los niveles (solo sigue \texttt{dist[matchV[v]] == dist[u] + 1}).
	\item Cada fase aumenta el matching en $\ge 1$ arista.
\end{itemize}

\paragraph{Codigo clave.}
El BFS de Hopcroft-Karp:
\begin{verbatim}
bool bfs() {
    queue<int> q;
    for (int v : L) {
        if (matchU[v] == -1) { dist[v] = 0; q.push(v); }
        else dist[v] = -1;
    }
    bool found = false;
    while (!q.empty()) {
        int v = q.front(); q.pop();
        for (int u : g[v]) {
            if (matchV[u] != -1 && dist[matchV[u]] == -1) {
                dist[matchV[u]] = dist[v] + 1;
                q.push(matchV[u]);
            }
            if (matchV[u] == -1) found = true;  // llegamos a libre
        }
    }
    return found;
}
\end{verbatim}

\paragraph{Complejidad.}
\begin{itemize}\setlength\itemsep{0pt}
	\item $O(\sqrt{V} E)$.
	\item Para $V = 10^4$, $E = 10^5$: $\sim 3 \cdot 10^7$ operaciones.
\end{itemize}

\paragraph{Donde tocar para modificar.}
\begin{itemize}\setlength\itemsep{0pt}
	\item \textbf{Reconstruir el matching}: ya viene en \texttt{matchU}/\texttt{matchV}.
	\item \textbf{Aniadir pesos}: Hungarian algorithm.
	\item \textbf{Densidad alta}: si $E \approx V^2$, sigue siendo util pero verificar constantes.
	\item \textbf{Aniadir capacidades}: partir los nodos con capacidad $> 1$ en varias copias.
\end{itemize}

\paragraph{Trampas y errores frecuentes.}
\begin{itemize}\setlength\itemsep{0pt}
	\item \texttt{dist[matchV[v]] == dist[u] + 1} (no \texttt{==} a secas); sin esto, el DFS no respeta los niveles y se pierde la complejidad.
	\item El BFS debe marcar los niveles correctamente; resetear \texttt{dist} para nodos alcanzables solo via matching.
	\item El DFS debe respetar los niveles; cualquier ``atajo'' invalida la cota.
\end{itemize}

\paragraph{Codigo.}
Ver \texttt{cpp.tex}, seccion \textit{Hopcroft-Karp}.

\subsection{Basicos}

\subsubsection{DFS / BFS / componentes}

\paragraph{Idea intuitiva.}
\textbf{DFS} (Depth-First Search) y \textbf{BFS} (Breadth-First Search) son los dos recorridos basicos de grafos. DFS usa una pila (o recursion) y explora tan profundo como puede antes de retroceder; BFS usa una cola y explora por ``niveles'' de distancia. La cantidad de \textbf{componentes conexas} se cuenta iniciando un BFS/DFS desde cada nodo no visitado. La diferencia clave: BFS da el camino mas corto en aristas; DFS da un orden de finalizacion util para problemas topologicos.

\paragraph{Propiedades clave (invariantes).}
\begin{itemize}\setlength\itemsep{0pt}
	\item BFS da la \textbf{distancia en aristas} mas corta desde el origen.
	\item DFS da un \textbf{orden topologico} (si se aniaden al final).
	\item Ambos son $O(V + E)$.
	\item Un grafo con $C$ componentes conexas satisface $\sum \text{subtree\_size}_v = V$ y $\sum E_v = 2E$.
	\item DFS en arbol da un spanning tree; las ``back edges'' corresponden a ciclos.
\end{itemize}

\paragraph{Codigo clave.}
BFS clasico:
\begin{verbatim}
queue<int> q;
q.push(source);
dist[source] = 0;
while (!q.empty()) {
    int v = q.front(); q.pop();
    for (int u : adj[v])
        if (dist[u] == -1) {
            dist[u] = dist[v] + 1;
            q.push(u);
        }
}
\end{verbatim}

\paragraph{Complejidad.}
\begin{itemize}\setlength\itemsep{0pt}
	\item $O(V + E)$ tiempo, $O(V)$ memoria.
	\item Cada arista se procesa exactamente 2 veces (ida y vuelta) en no dirigidos.
\end{itemize}

\paragraph{Donde tocar para modificar.}
\begin{itemize}\setlength\itemsep{0pt}
	\item \textbf{Aniadir pesos}: si las aristas tienen peso y queres shortest path, BFS ya no sirve; usar Dijkstra.
	\item \textbf{Grafo dirigido}: en DFS, aniadir check \texttt{if (state[u] == 1)} para detectar ciclos.
	\item \textbf{Bipartite check}: BFS alternando colores (ver el siguiente modulo).
	\item \textbf{Componentes fuertemente conexas (dirigidos)}: usar Tarjan o Kosaraju.
	\item \textbf{BFS en matriz}: aniadir offsets por direccion para grid 2D.
\end{itemize}

\paragraph{Trampas y errores frecuentes.}
\begin{itemize}\setlength\itemsep{0pt}
	\item Stack overflow con DFS recursivo para $V > 10^5$; usar version iterativa con pila explicita.
	\item En BFS, olvidar marcar \texttt{dist[u] = dist[v] + 1} antes de pushear da visitas duplicadas.
	\item Para grafos dirigidos, ``componente conexa'' no esta definida; usar SCC.
\end{itemize}

\paragraph{Codigo.}
Ver \texttt{cpp.tex}, seccion \textit{DFS / BFS / componentes}.

\vspace{0.5em}

\subsubsection{Topological sort (Kahn + DFS)}

\paragraph{Idea intuitiva.}
Un \textbf{orden topologico} de un DAG es una permutacion de los vertices tal que toda arista $u \to v$ tiene $u$ antes que $v$. Dos algoritmos clasicos:
\begin{itemize}\setlength\itemsep{0pt}
	\item \textbf{Kahn}: BFS sobre nodos con \texttt{indeg == 0}; al ``consumir'' un nodo, decrementar el indeg de sus vecinos; los que lleguen a 0 se encolan.
	\item \textbf{DFS}: hacer DFS; al terminar un nodo (estado = 2), aniadirlo al final; el orden es el reverso.
\end{itemize}
La demostracion: en un DAG siempre hay al menos un nodo con indeg 0; quitarlo mantiene el DAG; por induccion, se obtiene el orden.

\paragraph{Propiedades clave (invariantes).}
\begin{itemize}\setlength\itemsep{0pt}
	\item Existe sii el grafo es un DAG.
	\item Si al final \texttt{order.size() < V}, hay un ciclo.
	\item Kahn da el orden ``natural'' (los que pueden ir primero); DFS da el orden ``al reves'' (los que pueden ir al final primero).
	\item La cantidad de ordenes topologicos puede ser exponencial.
\end{itemize}

\paragraph{Codigo clave.}
Kahn con priority queue (orden lexicografico):
\begin{verbatim}
priority_queue<int, vector<int>, greater<int>> pq;
for (int v = 0; v < n; ++v) if (indeg[v] == 0) pq.push(v);
while (!pq.empty()) {
    int v = pq.top(); pq.pop();
    order.push_back(v);
    for (int u : adj[v]) {
        if (--indeg[u] == 0) pq.push(u);
    }
}
\end{verbatim}

\paragraph{Complejidad.}
\begin{itemize}\setlength\itemsep{0pt}
	\item $O(V + E)$ ambos.
	\item Con priority queue, $O((V + E) \log V)$.
\end{itemize}

\paragraph{Donde tocar para modificar.}
\begin{itemize}\setlength\itemsep{0pt}
	\item \textbf{Detectar ciclos}: comparar \texttt{order.size()} con \texttt{V}.
	\item \textbf{Orden lexicografico minimo}: usar priority queue en Kahn.
	\item \textbf{Topological sort con pesos}: combinar con shortest path en DAG (un solo pase en orden topologico).
	\item \textbf{Encontrar todos los ordenes}: DP sobre subsets (solo para $V$ pequeno).
\end{itemize}

\paragraph{Trampas y errores frecuentes.}
\begin{itemize}\setlength\itemsep{0pt}
	\item En el snippet, Kahn retorna orden vacio si hay ciclo (util para detectar).
	\item En DFS, no olvidar \texttt{reverse(order)}.
	\item Confundir indeg con outdeg; el orden Kahn quita primero los nodos ``fuente''.
	\item Para grafos no DAG, el algoritmo termina con \texttt{order.size() < V}.
\end{itemize}

\paragraph{Codigo.}
Ver \texttt{cpp.tex}, seccion \textit{Topological sort (Kahn + DFS)}.

\vspace{0.5em}

\subsubsection{Bipartite check + 2-coloring}

\paragraph{Idea intuitiva.}
Un grafo es \textbf{bipartito} si sus vertices se pueden dividir en dos conjuntos $L$ y $R$ tales que toda arista va de $L$ a $R$. Esto equivale a decir que se puede \textbf{2-colorear} sin que dos vertices adyacentes tengan el mismo color. BFS alternando colores detecta esto: si en algun momento dos adyacentes tienen el mismo color, no es bipartito. La demostracion: por cada componente, fijar el color de un nodo determina todos los demas; si en algun paso hay conflicto, hay un ciclo impar (no bipartito).

\paragraph{Propiedades clave (invariantes).}
\begin{itemize}\setlength\itemsep{0pt}
	\item Grafo bipartito $\Leftrightarrow$ sin ciclos impares.
	\item Conexo: si un componente no es bipartito, todo el grafo no lo es.
	\item Si hay \textbf{multiple componentes}, cada uno se colorea independientemente.
	\item La 2-coloracion es unica salvo swap de colores.
\end{itemize}

\paragraph{Codigo clave.}
BFS con 2-coloreo:
\begin{verbatim}
vector<int> color(n, -1);
for (int v = 0; v < n; ++v) if (color[v] == -1) {
    color[v] = 0;
    queue<int> q; q.push(v);
    while (!q.empty()) {
        int u = q.front(); q.pop();
        for (int w : adj[u]) {
            if (color[w] == -1) { color[w] = 1 - color[u]; q.push(w); }
            else if (color[w] == color[u]) return false;
        }
    }
}
return true;
\end{verbatim}

\paragraph{Complejidad.}
\begin{itemize}\setlength\itemsep{0pt}
	\item $O(V + E)$.
	\item Memoria $O(V)$.
\end{itemize}

\paragraph{Donde tocar para modificar.}
\begin{itemize}\setlength\itemsep{0pt}
	\item \textbf{Colorear desde multiples fuentes}: BFS desde cada nodo no coloreado.
	\item \textbf{Bipartito + matching}: Hopcroft-Karp.
	\item \textbf{Grafo bipartito ponderado}: Hungarian algorithm.
	\item \textbf{Verificar bipartito en subgrafo}: solo aplicar BFS a los nodos relevantes.
\end{itemize}

\paragraph{Trampas y errores frecuentes.}
\begin{itemize}\setlength\itemsep{0pt}
	\item Verificar \textbf{todos} los componentes, no solo el primero.
	\item Para grafos con auto-loops, un auto-loop hace al grafo no bipartito.
	\item En multigrafos, las aristas multiples no afectan el chequeo de bipartito.
\end{itemize}

\paragraph{Codigo.}
Ver \texttt{cpp.tex}, seccion \textit{Bipartite check + 2-coloring}.

\vspace{0.5em}

\subsubsection{0-1 BFS}

\paragraph{Idea intuitiva.}
Variante de BFS para grafos con \textbf{pesos 0 o 1}. Usa un \texttt{deque}: aristas de peso 0 se encolan al frente (\texttt{push\_front}), aristas de peso 1 al final (\texttt{push\_back}). Asi, la primera vez que se extrae un nodo, su distancia es la definitiva, igual que en Dijkstra pero en $O(V + E)$. La intuicion: el deque mantiene los nodos ordenados por distancia; las aristas de peso 0 no cambian la distancia (entran al frente), las de peso 1 la incrementan (entran al final).

\paragraph{Propiedades clave (invariantes).}
\begin{itemize}\setlength\itemsep{0pt}
	\item Asume pesos \textbf{solo} 0 o 1.
	\item Mejor caso que Dijkstra (sin factor $\log V$).
	\item No funciona con pesos $\ge 2$.
	\item La cola mantiene los nodos ordenados por distancia no decreciente.
\end{itemize}

\paragraph{Codigo clave.}
El deque doble:
\begin{verbatim}
deque<int> q;
q.push_front(source);
dist[source] = 0;
while (!q.empty()) {
    int v = q.front(); q.pop_front();
    for (auto [w, u] : adj[v]) {
        if (dist[v] + w < dist[u]) {
            dist[u] = dist[v] + w;
            if (w == 0) q.push_front(u);
            else q.push_back(u);
        }
    }
}
\end{verbatim}

\paragraph{Complejidad.}
\begin{itemize}\setlength\itemsep{0pt}
	\item $O(V + E)$.
	\item Cada nodo se inserta a lo sumo $O(1)$ veces (con check de distancia).
\end{itemize}

\paragraph{Donde tocar para modificar.}
\begin{itemize}\setlength\itemsep{0pt}
	\item \textbf{Pesos en $[0, k]$ small}: usar small k-ary BFS o Dijkstra.
	\item \textbf{Reconstruir camino}: aniadir \texttt{parent[]}.
	\item \textbf{Pesos negativos}: no funciona, usar Bellman-Ford.
\end{itemize}

\paragraph{Trampas y errores frecuentes.}
\begin{itemize}\setlength\itemsep{0pt}
	\item Si los pesos son otro rango, el algoritmo da respuestas incorrectas. Verificar que las aristas son solo 0/1.
	\item El chequeo \texttt{dist[v] + w < dist[u]} es crucial; sino, el algoritmo procesa nodos multiples veces.
	\item Si una arista tiene peso $\ge 2$, aniadir verificaciones o cambiar a Dijkstra.
\end{itemize}

\paragraph{Codigo.}
Ver \texttt{cpp.tex}, seccion \textit{0-1 BFS}.

\subsection{MST}

\subsubsection{Kruskal}

\paragraph{Idea intuitiva.}
Encuentra el \textbf{arbol de expansion minima} (MST) en $O(E \log E)$. Ordena las aristas por peso ascendente y las aniade al MST si no forman ciclo (chequear con DSU). El arbol tiene exactamente $V - 1$ aristas. La demostracion (``Cut Property''): para cualquier corte, la arista minima cruzando el corte pertenece a \emph{algun} MST; por lo tanto, elegir la minima disponible que no forma ciclo es optimo.

\paragraph{Propiedades clave (invariantes).}
\begin{itemize}\setlength\itemsep{0pt}
	\item Algoritmo \textbf{greedy} correcto (Cut Property).
	\item Funciona en grafos \textbf{no conexos}; produce un \textbf{minimum spanning forest}.
	\item Si las aristas tienen pesos unicos, el MST es unico.
	\item El MST tiene exactamente $V - 1$ aristas (o menos si el grafo es disconexo).
\end{itemize}

\paragraph{Codigo clave.}
Kruskal con DSU:
\begin{verbatim}
sort(edges.begin(), edges.end());
DSU dsu(n);
long long mst = 0;
int edgesUsed = 0;
for (auto [w, u, v] : edges)
    if (dsu.unite(u, v)) {
        mst += w;
        if (++edgesUsed == n - 1) break;
    }
\end{verbatim}

\paragraph{Complejidad.}
\begin{itemize}\setlength\itemsep{0pt}
	\item $O(E \log E)$ (ordenamiento) o $O(E \log V)$ con DSU optimizations.
	\item DSU es casi $O(1)$ amortizado, asi que el sort domina.
\end{itemize}

\paragraph{Donde tocar para modificar.}
\begin{itemize}\setlength\itemsep{0pt}
	\item \textbf{Second best MST}: enumerar aristas del MST, sacar una, aniadir una no-MST; recalcular.
	\item \textbf{Kruskal randomizado}: si el grafo es casi-completo, agregar shuffle.
	\item \textbf{Aniadir info al DSU}: guardar el peso maximo en el camino, etc.
	\item \textbf{DSU con rollback}: si necesitas Kruskal reversible (para queries offline).
\end{itemize}

\paragraph{Trampas y errores frecuentes.}
\begin{itemize}\setlength\itemsep{0pt}
	\item El DSU debe estar limpio para cada invocacion.
	\item Si las aristas tienen pesos negativos, Kruskal sigue funcionando.
	\item Si el grafo es disconexo, el ``MST'' es un bosque con varias componentes.
	\item El check \texttt{if (++edgesUsed == n - 1) break;} optimiza la salida temprana.
\end{itemize}

\paragraph{Codigo.}
Ver \texttt{cpp.tex}, seccion \textit{Kruskal}.

\vspace{0.5em}

\subsubsection{Prim (priority queue)}

\paragraph{Idea intuitiva.}
Otra forma de calcular el MST: empezar desde un nodo y, en cada paso, aniadir la arista de menor peso que conecta el conjunto construido con el resto. Usando priority queue, es $O((V + E) \log V)$. La demostracion: cada vez que se aniade una arista, es la minima cruzando el ``corte'' entre el conjunto actual y el resto (Cut Property); eso preserva optimalidad.

\paragraph{Propiedades clave (invariantes).}
\begin{itemize}\setlength\itemsep{0pt}
	\item Similar a Dijkstra pero sin acumular distancias, solo eligiendo minimo por arista.
	\item Funciona mejor que Kruskal en \textbf{grafos densos} ($E \approx V^2$).
	\item Si el grafo es disconexo, genera el bosque (solo llega a los alcanzables).
	\item Cada nodo se inserta al heap una vez; cada arista se procesa una vez.
\end{itemize}

\paragraph{Codigo clave.}
Prim clasico con heap:
\begin{verbatim}
priority_queue<pair<int, int>, vector<...>, greater<...>> pq;
vector<bool> used(n, false);
vector<int> minEdge(n, INF);
minEdge[0] = 0;
pq.push({0, 0});
while (!pq.empty()) {
    auto [w, v] = pq.top(); pq.pop();
    if (used[v]) continue;
    used[v] = true;
    mst += w;
    for (auto [peso, u] : adj[v])
        if (!used[u] && peso < minEdge[u]) {
            minEdge[u] = peso;
            pq.push({peso, u});
        }
}
\end{verbatim}

\paragraph{Complejidad.}
\begin{itemize}\setlength\itemsep{0pt}
	\item $O((V + E) \log V)$.
	\item Con matriz de adyacencia: $O(V^2)$ sin heap.
\end{itemize}

\paragraph{Donde tocar para modificar.}
\begin{itemize}\setlength\itemsep{0pt}
	\item \textbf{Prim con matriz de adyacencia}: $O(V^2)$ sin heap, util en grafos densos pequenos.
	\item \textbf{Reconstruir el MST}: aniadir \texttt{parent[]}.
	\item \textbf{Prim en streaming}: variante para grafos que llegan incrementalmente.
\end{itemize}

\paragraph{Trampas y errores frecuentes.}
\begin{itemize}\setlength\itemsep{0pt}
	\item El \texttt{if (used[v]) continue;} es crucial: si el nodo ya fue procesado, saltar.
	\item En la version matriz, elije el minimo por linea (no heap).
	\item Si el grafo es disconexo, aniadir un loop externo sobre todos los nodos no usados.
\end{itemize}

\paragraph{Codigo.}
Ver \texttt{cpp.tex}, seccion \textit{Prim (priority queue)}.

\subsection{Flujos}

\subsubsection{Edmonds-Karp (max flow)}

\paragraph{Idea intuitiva.}
Implementacion BFS del algoritmo de \textbf{Ford-Fulkerson}. En cada iteracion, hacer BFS desde $s$ hasta $t$ siguiendo aristas con capacidad residual $> 0$; si se llega a $t$, se encontro un \textbf{camino de aumento}. Enviar el minimo de las capacidades en el camino; actualizar capacidades residuales. La demostracion de $O(VE^2)$: cada BFS encuentra el augmenting path mas corto, y la longitud minima crece monotona; como hay $\le E$ longitudes posibles, hay $\le E$ fases.

\paragraph{Propiedades clave (invariantes).}
\begin{itemize}\setlength\itemsep{0pt}
	\item BFS garantiza caminos de aumento mas cortos.
	\item Capacidad residual: forward edge baja, backward edge sube.
	\item Complejidad $O(VE^2)$.
	\item El flujo maximo respeta la conservacion en nodos (in = out salvo $s, t$).
\end{itemize}

\paragraph{Codigo clave.}
BFS + augmenting path:
\begin{verbatim}
while (bfs(s, t, parent)) {  // bfs encuentra camino o retorna false
    int pathFlow = INF;
    for (int v = t; v != s; v = parent[v])
        pathFlow = min(pathFlow, capacity[parent[v]][v]);
    for (int v = t; v != s; v = parent[v]) {
        capacity[parent[v]][v] -= pathFlow;
        capacity[v][parent[v]] += pathFlow;
    }
    maxFlow += pathFlow;
}
\end{verbatim}

\paragraph{Complejidad.}
\begin{itemize}\setlength\itemsep{0pt}
	\item $O(VE^2)$.
	\item En la practica, mas lento que Dinic para grafos grandes.
\end{itemize}

\paragraph{Donde tocar para modificar.}
\begin{itemize}\setlength\itemsep{0pt}
	\item \textbf{Grafo bipartito}: $O(\sqrt{V} E)$ (equivalente a Hopcroft-Karp).
	\item \textbf{Dinic} es preferible en la mayoria de los casos.
	\item \textbf{Capacidades doubles}: usar tolerancia para comparacion.
	\item \textbf{Min-cost flow}: aniadir costo por arista y modificar el algoritmo.
\end{itemize}

\paragraph{Trampas y errores frecuentes.}
\begin{itemize}\setlength\itemsep{0pt}
	\item Las capacidades deben ser $> 0$ para entrar a la cola (no $>= 0$); sino, ciclos infinitos.
	\item Para grafos bipartitos, Dinic o Hopcroft-Karp son mejores.
	\item Si el flujo maximo es muy grande, los \texttt{int} pueden overflow; usar \texttt{long long}.
\end{itemize}

\paragraph{Codigo.}
Ver \texttt{cpp.tex}, seccion \textit{Edmonds-Karp (max flow)}.

\vspace{0.5em}

\subsubsection{Dinic (max flow)}

\paragraph{Idea intuitiva.}
Algoritmo de flujo mas eficiente que Edmonds-Karp: combina \textbf{BFS para level graph} con \textbf{DFS bloqueante} (multiple augmenting paths en una sola pasada). Cada fase cuesta $O(E)$ y hay $O(V)$ fases en el peor caso; para bipartitos, $O(\sqrt{V} E)$. La razon: cada fase ``satura'' al menos una arista, asi que en $O(E)$ fases el flujo maximo se alcanza; con $O(V)$ niveles en el BFS, hay $O(VE)$ operaciones totales.

\paragraph{Propiedades clave (invariantes).}
\begin{itemize}\setlength\itemsep{0pt}
	\item \texttt{level[v]} = distancia desde $s$ en el level graph (solo aristas con cap $> 0$).
	\item \texttt{it[v]} = iterador actual para evitar reprocesar aristas.
	\item DFS solo sigue aristas con \texttt{level[e.to] == level[v] + 1}.
	\item Back-edges tienen \texttt{rev} apuntando a la forward edge original.
\end{itemize}

\paragraph{Codigo clave.}
El DFS bloqueante de Dinic:
\begin{verbatim}
long long dfs(int v, int t, long long f) {
    if (v == t) return f;
    for (int& i = it[v]; i < adj[v].size(); ++i) {
        Edge& e = adj[v][i];
        if (e.cap > 0 && level[e.to] == level[v] + 1) {
            long long ret = dfs(e.to, t, min(f, e.cap));
            if (ret > 0) {
                e.cap -= ret;
                adj[e.to][e.rev].cap += ret;
                return ret;
            }
        }
    }
    return 0;
}
\end{verbatim}

\paragraph{Complejidad.}
\begin{itemize}\setlength\itemsep{0pt}
	\item $O(V^2 E)$ general, $O(\sqrt{V} E)$ para bipartitos.
	\item En la practica, mucho mas rapido que Edmonds-Karp.
\end{itemize}

\paragraph{Donde tocar para modificar.}
\begin{itemize}\setlength\itemsep{0pt}
	\item \textbf{Dinic con scaling}: aniadir factor de escala para capacidades grandes.
	\item \textbf{Lower bounds}: aniadir un super-source/sink para forzar flujo minimo.
	\item \textbf{Matching bipartito}: el grafo se construye como source -> L -> R -> sink, todos con capacidad 1.
	\item \textbf{Push-relabel}: alternativa para grafos muy densos.
\end{itemize}

\paragraph{Trampas y errores frecuentes.}
\begin{itemize}\setlength\itemsep{0pt}
	\item \texttt{it[v]} debe \textbf{incrementarse} en cada iteracion; sino, el DFS se queda atascado en la misma arista.
	\item Si el BFS no encuentra $t$, el flujo es maximo; terminar.
	\item Las back-edges son cruciales para ``cancelar'' elecciones previas.
\end{itemize}

\paragraph{Codigo.}
Ver \texttt{cpp.tex}, seccion \textit{Dinic (max flow)}.

\vspace{0.5em}

\subsubsection{Min-Cost Max-Flow}

\paragraph{Idea intuitiva.}
Encuentra el flujo maximo con \textbf{costo minimo} (asumiendo costos no negativos). Usa \textbf{SPFA} (o Dijkstra con potentials) para encontrar el shortest path en el \textbf{costo} desde $s$ hasta $t$ considerando capacidades residuales. Cada augmenting path aporta su costo $\times$ flujo. La idea: los potentials reescalan las aristas para que todas tengan costo no negativo, permitiendo usar Dijkstra.

\paragraph{Propiedades clave (invariantes).}
\begin{itemize}\setlength\itemsep{0pt}
	\item \texttt{pot[]} = potentials para evitar aristas negativas (Johnson's trick).
	\item \texttt{dist[v]} = costo minimo desde $s$ con potentials.
	\item Solo se aumentan aristas con \texttt{cap > 0} (forward) y costo es \texttt{+cost} (forward) o \texttt{-cost} (backward).
	\item Tras cada augmenting, los potentials se actualizan para mantener la consistencia.
\end{itemize}

\paragraph{Codigo clave.}
El ciclo de min-cost flow:
\begin{verbatim}
while (spfa(s, t, dist, parent)) {
    int pathFlow = INF;
    for (int v = t; v != s; v = parent[v])
        pathFlow = min(pathFlow, capacity[parent[v]][v]);
    for (int v = t; v != s; v = parent[v]) {
        capacity[parent[v]][v] -= pathFlow;
        capacity[v][parent[v]] += pathFlow;
        cost += pathFlow * edge_cost[parent[v]][v];
    }
}
\end{verbatim}

\paragraph{Complejidad.}
\begin{itemize}\setlength\itemsep{0pt}
	\item $O(F \cdot V E)$ con SPFA, $O(F \cdot E \log V)$ con Dijkstra + potentials.
	\item $F$ = flujo maximo total.
\end{itemize}

\paragraph{Donde tocar para modificar.}
\begin{itemize}\setlength\itemsep{0pt}
	\item \textbf{Dijkstra + potentials}: si los costos son enteros no negativos, esto es mas estable.
	\item \textbf{Costos negativos en aristas}: requieren Bellman-Ford inicial o reorden.
	\item \textbf{Flujo minimo con costo}: encontrar el minimo flujo posible y su costo.
	\item \textbf{Costo por unidad}: si los costos son muy grandes, aniadir verificacion.
\end{itemize}

\paragraph{Trampas y errores frecuentes.}
\begin{itemize}\setlength\itemsep{0pt}
	\item El \texttt{pot[]} se actualiza \textbf{al final} de cada iteracion (no durante).
	\item Las back-edges tienen \textbf{costo negativo}, esencial para corregir elecciones previas.
	\item Si los costos son muy grandes, considerar modular arithmetic.
	\item Para problemas con flujo exacto, aniadir verificacion post-algoritmo.
\end{itemize}

\paragraph{Codigo.}
Ver \texttt{cpp.tex}, seccion \textit{Min-Cost Max-Flow}.

\subsection{Especales}

\subsubsection{2-SAT (implication graph + SCC)}

\paragraph{Idea intuitiva.}
Resuelve formulas booleanas en \textbf{CNF} con clausulas de \textbf{2 literales}: $(l_1 \lor l_2)$. Cada variable $x$ se representa con dos nodos: $x$ (verdadero) y $\neg x$ (falso). Cada clausula $l_1 \lor l_2$ se descompone en dos implicaciones: $\neg l_1 \to l_2$ y $\neg l_2 \to l_1$. Si en el grafo de implicaciones $\neg x$ y $x$ estan en la misma SCC, la formula es \textbf{UNSAT}. La demostracion: si $\neg x \to x$ en el grafo, entonces $\neg x$ implica $x$, asi que $\neg x$ debe ser falso y $x$ verdadero; pero entonces $\neg x$ es falso implica una contradiccion.

\paragraph{Propiedades clave (invariantes).}
\begin{itemize}\setlength\itemsep{0pt}
	\item Cada literal es un nodo; total $2n$ nodos.
	\item Kosaraju/Tarjan dan las SCCs.
	\item Si SAT: asignar $x = $ (topological order del condensation).
	\item Clausulas \texttt{setTrue(x)} se modelan como \texttt{implies(!x, x)}.
	\item La formula es SAT sii no hay variable $x$ con $x$ y $\neg x$ en la misma SCC.
\end{itemize}

\paragraph{Codigo clave.}
Conversion de clausulas a implicaciones:
\begin{verbatim}
// Clausula (a OR b) -> implica ~a -> b Y ~b -> a
addImplication(not(a), b);
addImplication(not(b), a);
// Forzar x: implica ~x -> x
if (force) addImplication(not(x), x);
\end{verbatim}

\paragraph{Complejidad.}
\begin{itemize}\setlength\itemsep{0pt}
	\item $O(V + E) = O(N + M)$ con Kosaraju o Tarjan.
	\item Memoria $O(V + E)$.
\end{itemize}

\paragraph{Donde tocar para modificar.}
\begin{itemize}\setlength\itemsep{0pt}
	\item \textbf{Forzar variables}: aniadir \texttt{setTrue} / \texttt{setFalse}.
	\item \textbf{3-SAT}: NP-completo; no usar este algoritmo.
	\item \textbf{Contar soluciones}: las SCCs condensation dan una formula para contar; multiplicar las formas de asignar variables en cada componente no-trivial.
	\item \textbf{Output de la asignacion}: aniadir el vector \texttt{assignment[]}.
\end{itemize}

\paragraph{Trampas y errores frecuentes.}
\begin{itemize}\setlength\itemsep{0pt}
	\item La asignacion usa el \textbf{orden topologico inverso} del condensation: \texttt{val[i] = comp[2i] < comp[2i+1]}.
	\item El ID de $\neg x$ es \texttt{2*x + 1} (o similar); verificar consistencia.
	\item Si hay clausulas con literales repetidos (e.g., $(x \lor x)$), tratar como $(x)$ simple.
\end{itemize}

\paragraph{Codigo.}
Ver \texttt{cpp.tex}, seccion \textit{2-SAT (implication graph + SCC)}.

\vspace{0.5em}

\subsubsection{Euler tour / Hierholzer (camino y circuito)}

\paragraph{Idea intuitiva.}
Encuentra un \textbf{camino o circuito euleriano} (que usa cada arista exactamente una vez). \textbf{Condiciones}:
\begin{itemize}\setlength\itemsep{0pt}
	\item \textbf{No dirigido}: todos los grados pares (circuito) o exactamente 2 impares (camino, empieza en uno, termina en el otro). Ademas, conexo sobre el subgrafo con aristas.
	\item \textbf{Dirigido}: $\text{indeg}(v) = \text{outdeg}(v)$ para todos (circuito) o exactamente un nodo con $\text{outdeg} = \text{indeg} + 1$ y otro con $\text{indeg} = \text{outdeg} + 1$ (camino).
\end{itemize}

\paragraph{Hierholzer}: elegir un ciclo desde el inicio, mientras haya ciclos ``laterales'' desde un nodo del camino, mergearlos. Implementacion con pila: avanzar por aristas, al no poder mas aniadir a la respuesta y backtrack. La demostracion de $O(V + E)$: cada arista se procesa una vez (entra y sale del stack una vez).

\paragraph{Propiedades clave (invariantes).}
\begin{itemize}\setlength\itemsep{0pt}
	\item Multigrafos permitidos; las aristas se marcan como usadas.
	\item No dirigido: cada arista aparece \textbf{dos veces} en \texttt{adj} (ida y vuelta).
	\item La construccion del path con ids consistentes es importante.
	\item El resultado es unico salvo orden de las aristas en multigrafos.
\end{itemize}

\paragraph{Codigo clave.}
Hierholzer con pila:
\begin{verbatim}
vector<int> path;
stack<int> st;
st.push(start);
while (!st.empty()) {
    int v = st.top();
    while (!adj[v].empty() && used[adj[v].back()]) adj[v].pop_back();
    if (adj[v].empty()) {
        path.push_back(v);
        st.pop();
    } else {
        int e = adj[v].back();
        used[e] = true;
        int u = edges[e].to(v);
        adj[v].pop_back();
        st.push(u);
    }
}
reverse(path.begin(), path.end());
\end{verbatim}

\paragraph{Complejidad.}
\begin{itemize}\setlength\itemsep{0pt}
	\item $O(V + E)$.
	\item Memoria $O(V + E)$.
\end{itemize}

\paragraph{Donde tocar para modificar.}
\begin{itemize}\setlength\itemsep{0pt}
	\item \textbf{Reconstruir con multigrafo}: aniadir ids consistentes antes del algoritmo.
	\item \textbf{Aplicar a ``de Bruijn'' o problemas de strings}: modelar como multigrafo y buscar Euler.
	\item \textbf{Dirigido vs no dirigido}: el snippet tiene ambas versiones.
	\item \textbf{Letter addition}: modelar palabras como aristas y buscar el superstring.
\end{itemize}

\paragraph{Trampas y errores frecuentes.}
\begin{itemize}\setlength\itemsep{0pt}
	\item Si el grafo no es conexo (sobre las aristas), no existe euleriano. Verificar primero.
	\item Para multigrafos, los IDs de aristas son cruciales para distinguirlas.
	\item Si el problema es ``encuentra todos los euler tours'', el algoritmo es exponencial.
\end{itemize}

\paragraph{Codigo.}
Ver \texttt{cpp.tex}, seccion \textit{Euler tour / Hierholzer (camino y circuito)}.

\vspace{0.5em}

\subsubsection{Hamiltonian path / cycle (bitmask DP + backtracking)}

\paragraph{Idea intuitiva.}
Un \textbf{camino hamiltoniano} visita cada vertice \textbf{exactamente una vez}. \textbf{Determinar} si existe es NP-completo en general, pero con $V \le 20$ se puede hacer DP con bitmask en $O(V^2 \cdot 2^V)$. \textbf{Teoremas suficientes}: \textbf{Dirac} (grado $\ge n/2$) y \textbf{Ore} ($\text{deg}(u) + \text{deg}(v) \ge n$ para todo par no adyacente). La idea de la DP: cada estado es (conjunto de visitados, ultimo vertice); la transicion es agregar un vecino no visitado.

\paragraph{Propiedades clave (invariantes).}
\begin{itemize}\setlength\itemsep{0pt}
	\item DP: $dp[mask][v]$ = true si hay un camino que visita exactamente \texttt{mask} y termina en $v$.
	\item TSP: variante con pesos minimos (Hamilton cycle minimo).
	\item Reconstruccion: guardar \texttt{par[mask][v]} (predecesor).
	\item $2^V$ estados con $V$ bits cada uno; factible para $V \le 20$.
\end{itemize}

\paragraph{Codigo clave.}
DP con bitmask:
\begin{verbatim}
dp[1][0] = 0;  // empezar en 0
for (int mask = 1; mask < (1<<n); ++mask)
    for (int v = 0; v < n; ++v) if (dp[mask][v] < INF)
        for (int u = 0; u < n; ++u)
            if (!(mask & (1<<u)))
                dp[mask | (1<<u)][u] = min(dp[mask | (1<<u)][u],
                                            dp[mask][v] + w[v][u]);
// TSP: respuesta = min(dp[(1<<n)-1][v] + w[v][0])
\end{verbatim}

\paragraph{Complejidad.}
\begin{itemize}\setlength\itemsep{0pt}
	\item $O(V^2 \cdot 2^V)$ tiempo, $O(V \cdot 2^V)$ memoria.
	\item Para $V = 20$: $\sim 4 \cdot 10^8$ operaciones (lento pero factible).
\end{itemize}

\paragraph{Donde tocar para modificar.}
\begin{itemize}\setlength\itemsep{0pt}
	\item \textbf{TSP asimetrico}: aniadir pesos $w[v][u]$ (no necesariamente simetricos).
	\item \textbf{Hamilton con inicio obligatorio}: fijar \texttt{src} en el DP inicial.
	\item \textbf{Bitmask DP mas compacto}: usar \texttt{long long} como mascara si $V \le 64$.
	\item \textbf{TSP optimo aproximado}: MST, 2-opt, simulated annealing.
	\item \textbf{Reconstruir el path}: usar \texttt{par[mask][v]} y backtrack.
\end{itemize}

\paragraph{Trampas y errores frecuentes.}
\begin{itemize}\setlength\itemsep{0pt}
	\item TSP asume grafo \textbf{completo} (con \texttt{LINF} para aristas faltantes); si no, aniadir chequeo.
	\item La condicion ``\texttt{mask == (1<<n) - 1}'' cierra el ciclo retornando al origen.
	\item Para $V > 20$, la memoria $2^V$ se vuelve prohibitiva; usar meet-in-the-middle o heuristicas.
\end{itemize}

\paragraph{Codigo.}
Ver \texttt{cpp.tex}, seccion \textit{Hamiltonian path / cycle (bitmask DP + backtracking)}.

% ============================================================
\section{DP}

\subsubsection{Knapsack 0/1}

\paragraph{Idea intuitiva.}
Dado un conjunto de $N$ objetos con peso $w_i$ y valor $v_i$, maximizar $\sum v_i$ con la restriccion $\sum w_i \le W$. La recurrencia clasica:
\[ dp[i][j] = \max(dp[i-1][j],\, dp[i-1][j - w_i] + v_i) \quad \text{si } j \ge w_i. \]
La optimizacion clave: la DP es 1-dimensional si iteramos $j$ \textbf{de mayor a menor} (asi no usamos el objeto $i$ dos veces). La razon de la iteracion inversa: en cada iteracion queremos el estado del ``item anterior''; iterar de menor a mayor contaminaria \texttt{dp[j - w_i]} con el item actual.

\paragraph{Propiedades clave (invariantes).}
\begin{itemize}\setlength\itemsep{0pt}
	\item $dp[j]$ al final = maximo valor con peso $\le j$.
	\item El bucle inverso \texttt{for (j = W; j >= w[i]; j--)} es lo que distingue 0/1 de unbounded.
	\item Solucion: maximo valor con cualquier peso $\le W$.
	\item La DP es exacta; no hay aproximacion.
\end{itemize}

\paragraph{Codigo clave.}
La version 1D optimizada:
\begin{verbatim}
vector<long long> dp(W + 1, 0);
for (int i = 0; i < N; ++i)
    for (int j = W; j >= w[i]; --j)  // invertido
        dp[j] = max(dp[j], dp[j - w[i]] + v[i]);
return dp[W];
\end{verbatim}

\paragraph{Complejidad.}
\begin{itemize}\setlength\itemsep{0pt}
	\item $O(N \cdot W)$ tiempo, $O(W)$ memoria.
	\item Constante muy pequena (solo sumas y max).
\end{itemize}

\paragraph{Donde tocar para modificar.}
\begin{itemize}\setlength\itemsep{0pt}
	\item \textbf{Reconstruir la seleccion}: aniadir \texttt{take[i][j]} (booleano o bool) en una version 2D.
	\item \textbf{Items con peso 0}: edge case; aniadir guard.
	\item \textbf{Knapsack con restricciones de grupos}: ``elegir a lo sumo 1 de cada grupo'' se modela como 0/1 con un peso extra.
	\item \textbf{Tamano de $W$}: si $W \le 10^9$ pero $N$ es chico, usar meet-in-the-middle.
	\item \textbf{Cambiar max por min}: aniadir \texttt{INF} en la inicializacion.
\end{itemize}

\paragraph{Trampas y errores frecuentes.}
\begin{itemize}\setlength\itemsep{0pt}
	\item Overflow en \texttt{dp[j - w[i]] + v[i]}; usar \texttt{long long} si los valores son grandes.
	\item El bucle \textbf{debe} ir de mayor a menor.
	\item Si los pesos son grandes y $N$ pequeno, conviene meet-in-the-middle.
	\item Para Knapsack exacto ($\sum w_i = W$), aniadir check al final.
\end{itemize}

\paragraph{Codigo.}
Ver \texttt{cpp.tex}, seccion \textit{Knapsack 0/1}.

\vspace{0.5em}

\subsubsection{Knapsack unbounded}

\paragraph{Idea intuitiva.}
Igual que Knapsack 0/1 pero \textbf{cada item se puede tomar multiples veces}. La recurrencia es similar:
\[ dp[j] = \max(dp[j],\, dp[j - w_i] + v_i), \]
pero ahora el bucle va de \textbf{menor a mayor} (porque tomar el item $i$ no afecta corridas futuras del mismo item). La idea: tras tomar el item $i$, queremos volver a tener la opcion de tomarlo; iterar de menor a mayor permite que el item ``se acumule''.

\paragraph{Propiedades clave (invariantes).}
\begin{itemize}\setlength\itemsep{0pt}
	\item El bucle \texttt{for (j = w[i]; j <= W; j++)} es lo que distingue unbounded.
	\item A diferencia de 0/1, no hay riesgo de ``usar dos veces el mismo item''.
	\item La DP es exacta; cada ``subconjunto'' se cuenta correctamente.
\end{itemize}

\paragraph{Codigo clave.}
La version 1D con bucle directo:
\begin{verbatim}
vector<long long> dp(W + 1, 0);
for (int i = 0; i < N; ++i)
    for (int j = w[i]; j <= W; ++j)  // directo (no invertido)
        dp[j] = max(dp[j], dp[j - w[i]] + v[i]);
return dp[W];
\end{verbatim}

\paragraph{Complejidad.}
\begin{itemize}\setlength\itemsep{0pt}
	\item $O(N \cdot W)$ tiempo, $O(W)$ memoria.
	\item Constante muy pequena.
\end{itemize}

\paragraph{Donde tocar para modificar.}
\begin{itemize}\setlength\itemsep{0pt}
	\item \textbf{Reconstruir cuantos de cada item}: aniadir \texttt{cnt[i]} y backtrack.
	\item \textbf{Cambiar moneda minima por maxima}: el snippet maximiza valor; para minimizar monedas, cambiar a \texttt{min} con inicializacion en \texttt{INF}.
	\item \textbf{Coin change}: si todos los pesos son positivos y queremos ``minimo numero de monedas'', es unbounded con min.
\end{itemize}

\paragraph{Trampas y errores frecuentes.}
\begin{itemize}\setlength\itemsep{0pt}
	\item Misma que Knapsack 0/1 con overflow.
	\item El orden de iteracion de items no afecta el resultado (convexo).
	\item Si se quiere ``maximo valor con exactamente $k$ items'', aniadir segunda dimension.
\end{itemize}

\paragraph{Codigo.}
Ver \texttt{cpp.tex}, seccion \textit{Knapsack unbounded}.

\vspace{0.5em}

\subsubsection{Knapsack meet-in-the-middle}

\paragraph{Idea intuitiva.}
Para $N \le 40$ y $W$ moderado, no se puede hacer $O(NW)$ directamente. Partir en dos mitades, enumerar todos los $2^{N/2}$ subconjuntos de cada mitad, y combinar. Para cada subconjunto de la segunda mitad con peso $sw_2$, encontrar el maximo valor de la primera mitad con peso $\le W - sw_2$. La idea: ``divide y venceras'' sobre el tamano de la entrada; cada mitad tiene mitad del espacio, asi que la combinacion es geometrica en lugar de multiplicativa.

\paragraph{Propiedades clave (invariantes).}
\begin{itemize}\setlength\itemsep{0pt}
	\item Hay $2^{N/2}$ subconjuntos por mitad; total $O(2^{N/2})$.
	\item \texttt{sums} guarda los valores de subconjuntos de la primera mitad ordenados (o preprocesados).
	\item Para cada subconjunto de la segunda mitad, busqueda lineal o binaria en \texttt{sums}.
	\item La mitad mas grande (con $N/2 + 1$ items para $N$ impar) tiene $2^{N/2+1}$ subconjuntos.
\end{itemize}

\paragraph{Codigo clave.}
Enumerar subconjuntos de la primera mitad:
\begin{verbatim}
int n1 = N / 2, n2 = N - n1;
vector<pair<long long, long long>> sums1;
for (int mask = 0; mask < (1 << n1); ++mask) {
    long long peso = 0, valor = 0;
    for (int i = 0; i < n1; ++i) if (mask & (1 << i)) {
        peso += w[i]; valor += v[i];
    }
    sums1.push_back({peso, valor});
}
sort(sums1.begin(), sums1.end());  // por peso
// eliminar duplicados dominados
\end{verbatim}

\paragraph{Complejidad.}
\begin{itemize}\setlength\itemsep{0pt}
	\item $O(2^{N/2})$ tiempo y memoria (mejor que $O(NW)$ cuando $W > 2^{N/2}$).
	\item Para $N = 40$: $\sim 2 \cdot 10^6$ subconjuntos, factible.
\end{itemize}

\paragraph{Donde tocar para modificar.}
\begin{itemize}\setlength\itemsep{0pt}
	\item \textbf{Tamano mayor}: si $N$ es mayor pero $W$ es chico, conviene iterar por peso, no por subconjuntos.
	\item \textbf{Varias mochilas}: adaptar para multiples $W$.
	\item \textbf{Optimizar el ``combinado''}: usar \texttt{upper\_bound} para encontrar el maximo valor con peso $\le \text{rem}.
	\item \textbf{Eliminar pares dominados}: si un par (peso, valor) tiene otro con menos peso y mas valor, eliminarlo.
\end{itemize}

\paragraph{Trampas y errores frecuentes.}
\begin{itemize}\setlength\itemsep{0pt}
	\item $N$ impar: una mitad tiene $N/2 + 1$ items.
	\item Overflow al sumar pesos/valores en \texttt{long long}.
	\item Sin deduplicacion, el algoritmo es correcto pero mas lento.
\end{itemize}

\paragraph{Codigo.}
Ver \texttt{cpp.tex}, seccion \textit{Knapsack meet-in-the-middle}.

\vspace{0.5em}

\subsubsection{LIS O(n log n)}

\paragraph{Idea intuitiva.}
La \textbf{Longest Increasing Subsequence} se calcula en $O(n \log n)$ con el truco del \textbf{patience sorting}: mantener un arreglo \texttt{dp} donde \texttt{dp[len]} es el minimo ultimo valor de una IS de largo \texttt{len}. Para cada $x$, encontrar el primer \texttt{dp[len] >= x} y reemplazarlo. El largo final es la LIS. La demostracion: si hay dos IS del mismo largo, la que termina en un valor menor es ``mejor'' (puede extenderse mas facilmente).

\paragraph{Propiedades clave (invariantes).}
\begin{itemize}\setlength\itemsep{0pt}
	\item \texttt{dp} es estrictamente creciente.
	\item \texttt{lower\_bound} da LIS estricta; \texttt{upper\_bound} da no-estricta ($\le$).
	\item El algoritmo es \textbf{online}: procesa elementos en orden.
	\item La longitud de \texttt{dp} es la LIS al final.
\end{itemize}

\paragraph{Codigo clave.}
La version online:
\begin{verbatim}
vector<int> dp;
for (int x : a) {
    auto it = lower_bound(dp.begin(), dp.end(), x);
    if (it == dp.end()) dp.push_back(x);
    else *it = x;
}
// dp.size() = longitud de LIS estricta
\end{verbatim}

\paragraph{Complejidad.}
\begin{itemize}\setlength\itemsep{0pt}
	\item $O(n \log n)$.
	\item Memoria $O(n)$.
\end{itemize}

\paragraph{Donde tocar para modificar.}
\begin{itemize}\setlength\itemsep{0pt}
	\item \textbf{Reconstruir la LIS}: aniadir \texttt{parent[]} y \texttt{posInDp[]} para backtrack.
	\item \textbf{Contar LIS}: DP clasica $O(n^2)$ o segment tree sobre $dp$.
	\item \textbf{LDS / Longest non-increasing}: invertir la condicion en \texttt{lower\_bound} o multiplicar por -1.
	\item \textbf{Longest bitonic subsequence}: combinar LIS desde izq y desde der.
	\item \textbf{LIS en 2D}: ordenar por una dimension y aplicar LIS en la otra.
\end{itemize}

\paragraph{Trampas y errores frecuentes.}
\begin{itemize}\setlength\itemsep{0pt}
	\item \texttt{lower\_bound} vs \texttt{upper\_bound}: si tu comparacion es $\le$ vs $<$, cambia.
	\item Si los elementos son \texttt{int} pero podrian ser negativos, aniadir offset para usar como indices.
	\item Para LIS reconstruida, aniadir el array \texttt{parent} que apunta al anterior en la IS.
\end{itemize}

\paragraph{Codigo.}
Ver \texttt{cpp.tex}, seccion \textit{LIS O(n log n)}.

\vspace{0.5em}

\subsubsection{LCS / Edit distance}

\paragraph{Idea intuitiva.}
\textbf{LCS} (Longest Common Subsequence): para dos strings $a$, $b$, $dp[i][j]$ = LCS de $a[0..i)$ y $b[0..j)$. Recurrencia:
\[ dp[i][j] = \begin{cases} dp[i-1][j-1] + 1 & \text{si } a[i-1] = b[j-1] \\ \max(dp[i-1][j], dp[i][j-1]) & \text{si no} \end{cases} \]

\textbf{Edit distance} (Levenshtein): numero minimo de operaciones (insert, delete, replace) para convertir $a$ en $b$. Recurrencia similar con un $+1$ adicional. La idea: ``alinear'' los strings eligiendo matches/mismatches/insert/delete.

\paragraph{Propiedades clave (invariantes).}
\begin{itemize}\setlength\itemsep{0pt}
	\item LCS y Edit Distance son DP $O(nm)$.
	\item Memoria se puede reducir a $O(\min(n, m))$ con rolling array.
	\item Edit distance con peso por operacion distinto: ajustar el \texttt{+1} por el peso.
	\item La subestructura optima: la respuesta para prefijos se construye de prefijos menores.
\end{itemize}

\paragraph{Codigo clave.}
La recurrencia clasica de LCS:
\begin{verbatim}
vector<vector<int>> dp(n + 1, vector<int>(m + 1, 0));
for (int i = 1; i <= n; ++i)
    for (int j = 1; j <= m; ++j)
        if (a[i-1] == b[j-1]) dp[i][j] = dp[i-1][j-1] + 1;
        else dp[i][j] = max(dp[i-1][j], dp[i][j-1]);
return dp[n][m];
\end{verbatim}

\paragraph{Complejidad.}
\begin{itemize}\setlength\itemsep{0pt}
	\item $O(NM)$ tiempo y memoria.
	\item Para $N, M \le 5000$: $\sim 2.5 \cdot 10^7$ operaciones.
\end{itemize}

\paragraph{Donde tocar para modificar.}
\begin{itemize}\setlength\itemsep{0pt}
	\item \textbf{Reconstruir la LCS / las operaciones}: aniadir \texttt{prev[i][j]} para backtrack.
	\item \textbf{LCS de varios strings}: NP-hard para $\ge 3$ strings; en la practica se hace con bitmask.
	\item \textbf{Edit distance con costos distintos}: cambiar el \texttt{+1} por el costo apropiado.
	\item \textbf{Rolling array}: si memoria es problema, mantener solo 2 filas.
	\item \textbf{Hirschberg}: algoritmo $O(NM)$ tiempo, $O(\min(N,M))$ memoria, con reconstruccion.
\end{itemize}

\paragraph{Trampas y errores frecuentes.}
\begin{itemize}\setlength\itemsep{0pt}
	\item Indices en la recurrencia: \texttt{a[i-1]} y \texttt{b[j-1]} (no \texttt{a[i]}).
	\item Para reconstruccion, el caso \texttt{a[i-1] == b[j-1]} tiene prioridad.
	\item Si los strings son muy largos, considerar Hirschberg o bitset optimization.
\end{itemize}

\paragraph{Codigo.}
Ver \texttt{cpp.tex}, seccion \textit{LCS / Edit distance}.

\vspace{0.5em}

\subsubsection{Bitmask DP over subsets}

\paragraph{Idea intuitiva.}
Cuando el estado es un \textbf{subconjunto} de $\{0, \dots, N-1\}$, se representa como una mascara de $N$ bits. Para iterar sobre todos los subconjuntos de \texttt{mask}:
\[ \text{for (sub = mask; sub; sub = (sub - 1) \& mask)} \{ \dots \} \]
Para iterar sobre todos los superconjuntos de \texttt{mask} en $[0, U)$:
\[ \text{for (sup = mask; sup < U; sup = (sup + 1) | mask)} \{ \dots \} \]
La idea: las mascaras forman un lattice de subconjuntos; los trucos aritmeticos (\texttt{(sub-1) \& mask} y \texttt{(sup+1) | mask}) enumeran todo el lattice eficientemente.

\paragraph{Propiedades clave (invariantes).}
\begin{itemize}\setlength\itemsep{0pt}
	\item Iterar subsets: $O(2^k)$ para una mascara con $k$ bits.
	\item \texttt{sub \& mask == sub} siempre; \texttt{mask - sub} no siempre da un subconjunto.
	\item \texttt{\_\_builtin\_popcount(mask)} cuenta los bits.
	\item \texttt{(sub - 1) \& mask} itera subsets de \texttt{mask} en orden decreciente de inclusion.
\end{itemize}

\paragraph{Codigo clave.}
Iteracion sobre subsets de una mascara:
\begin{verbatim}
// Todos los subsets de mask
for (int sub = mask; sub; sub = (sub - 1) & mask) {
    // procesar sub
}
// Incluir sub = 0 si es necesario
sub = 0;  // procesar primero
for (int sup = mask; sup < (1 << N); sup = (sup + 1) | mask) {
    // procesar sup (todos los superconjuntos)
}
\end{verbatim}

\paragraph{Complejidad.}
\begin{itemize}\setlength\itemsep{0pt}
	\item $O(N \cdot 2^N)$ para una DP estandar sobre todos los subsets.
	\item Iterar subsets de \texttt{mask}: $O(2^{|mask|})$.
\end{itemize}

\paragraph{Donde tocar para modificar.}
\begin{itemize}\setlength\itemsep{0pt}
	\item \textbf{TSP / Hamiltonian path}: bitmask DP (ver tambien \textit{Hamiltonian path} en \texttt{cpp.tex} seccion Grafos).
	\item \textbf{Set cover}: iterar subsets.
	\item \textbf{DP con supermask}: si el estado incluye ``todos los elementos que \textbf{no} hemos usado'', usar supermask iteration.
	\item \textbf{Bitmasks grandes}: con $N > 20$, aniadir compresion o simulacion con segmentos.
\end{itemize}

\paragraph{Trampas y errores frecuentes.}
\begin{itemize}\setlength\itemsep{0pt}
	\item El orden de iteracion de subsets de \texttt{mask} es decreciente (de \texttt{mask} a $0$).
	\item \texttt{mask = 0} no tiene subsets no triviales; el bucle lo saltea correctamente.
	\item El truco \texttt{(sub - 1) \& mask} requiere cuidado con \texttt{sub = 0} (termina el bucle).
	\item Superconjuntos: \texttt{sup < (1 << N)} debe estar bien definido.
\end{itemize}

\paragraph{Codigo.}
Ver \texttt{cpp.tex}, seccion \textit{Bitmask DP over subsets}.

\vspace{0.5em}

\subsubsection{SOS DP}

\paragraph{Idea intuitiva.}
Para calcular, para cada \texttt{mask}, la suma de $f(\text{submask})$ sobre todos los $\text{submask} \subseteq \text{mask}$:
\[ dp[mask] = \sum_{\text{sub} \subseteq mask} f(\text{sub}). \]
Se computa en $O(N \cdot 2^N)$ mediante un loop por bit: para cada $i$, sumar a $dp[mask]$ la contribucion de $dp[mask \oplus (1 << i)]$ cuando $mask$ tiene el bit $i$. La idea: cada subconjunto $s \subseteq m$ se obtiene quitando bits uno a uno; asi, sumar incrementalmente captura todos los subconjuntos en $O(N \cdot 2^N)$.

\paragraph{Propiedades clave (invariantes).}
\begin{itemize}\setlength\itemsep{0pt}
	\item Inicializar $dp = f$.
	\item Para cada bit $i$: \texttt{if (mask \& (1 << i)) dp[mask] += dp[mask \^{} (1 << i)]}.
	\item Variantes: supersets DP (sumar sobre superconjuntos).
	\item El orden de iteracion (bit externo, mascara interno) es crucial.
\end{itemize}

\paragraph{Codigo clave.}
SOS DP (Sum over Subsets):
\begin{verbatim}
vector<long long> dp(1 << N);
for (int i = 0; i < (1 << N); ++i) dp[i] = f[i];
for (int bit = 0; bit < N; ++bit)
    for (int mask = 0; mask < (1 << N); ++mask)
        if (mask & (1 << bit))
            dp[mask] += dp[mask ^ (1 << bit)];
// dp[mask] = suma de f(sub) para sub ⊆ mask
\end{verbatim}

\paragraph{Complejidad.}
\begin{itemize}\setlength\itemsep{0pt}
	\item $O(N \cdot 2^N)$.
	\item Memoria $O(2^N)$.
\end{itemize}

\paragraph{Donde tocar para modificar.}
\begin{itemize}\setlength\itemsep{0pt}
	\item \textbf{Max/min en lugar de suma}: cambiar \texttt{+=} por la operacion correspondiente.
	\item \textbf{Supersets DP}: iterar \texttt{mask} de $0$ a $2^N$ y aniadir \texttt{dp[mask] += dp[mask | (1<<i)]}.
	\item \textbf{XOR convolution} (subset convolution): version mas avanzada.
	\item \textbf{Aplicaciones}: contar subconjuntos con propiedades especificas.
\end{itemize}

\paragraph{Trampas y errores frecuentes.}
\begin{itemize}\setlength\itemsep{0pt}
	\item Si $N$ es 30 o mas, $2^N$ no entra en memoria; usar trucos o reducir $N$.
	\item El sentido del XOR (sumar a \texttt{mask \^{} bit}) es para ``agregar el bit $i$ al subconjunto''.
	\item En supersets DP, la condicion es \texttt{(mask \& (1<<bit)) == 0}.
\end{itemize}

\paragraph{Codigo.}
Ver \texttt{cpp.tex}, seccion \textit{SOS DP}.

\vspace{0.5em}

\subsubsection{Digit DP template}

\paragraph{Idea intuitiva.}
Contar numeros en $[0, N]$ que cumplen una propiedad sobre sus digitos. La idea: DP por posicion, manteniendo:
\begin{itemize}\setlength\itemsep{0pt}
	\item \texttt{pos}: posicion actual (de izquierda a derecha).
	\item \texttt{tight}: si los digitos ya colocados son exactamente los de $N$ (limite superior estricto).
	\item \texttt{state}: estado adicional (digitos usados, suma, etc.).
\end{itemize}
La clave: la condicion \texttt{tight} permite acotar correctamente el siguiente digito; sin ella, la DP contaria incorrectamente.

\paragraph{Propiedades clave (invariantes).}
\begin{itemize}\setlength\itemsep{0pt}
	\item Si \texttt{tight} es true, el siguiente digito esta acotado por el de $N$; si no, puede ser 0-9.
	\item Memoria: $O(D \cdot 2 \cdot \text{states})$.
	\item Caso base: si \texttt{pos == D}, devolver 1 si el estado cumple la propiedad, sino 0.
	\item El bit \texttt{started} distingue ``numero con menos digitos'' (con leading zeros).
\end{itemize}

\paragraph{Codigo clave.}
La estructura basica:
\begin{verbatim}
long long dp(int pos, bool tight, State s) {
    if (pos == D) return check(s) ? 1 : 0;
    if (memo[pos][tight][s] != -1) return memo;
    int limit = tight ? digits[pos] : 9;
    long long ans = 0;
    for (int d = 0; d <= limit; ++d)
        ans += dp(pos + 1, tight && (d == limit), update(s, d, pos));
    return memo[pos][tight][s] = ans;
}
\end{verbatim}

\paragraph{Complejidad.}
\begin{itemize}\setlength\itemsep{0pt}
	\item $O(D \cdot 10 \cdot \text{states})$ donde $D$ es el numero de digitos.
	\item Para $D \le 19$ (hasta $10^{18}$): $\sim 200 \cdot \text{states}$.
\end{itemize}

\paragraph{Donde tocar para modificar.}
\begin{itemize}\setlength\itemsep{0pt}
	\item \textbf{Cambiar la propiedad}: modificar \texttt{updateState} y la condicion base.
	\item \textbf{Multiples estados}: aniadir mas dimensiones al \texttt{memo}.
	\item \textbf{Leading zeros}: aniadir un flag \texttt{started} para distinguir numeros con menos digitos.
	\item \textbf{Rango $[L, R]$}: calcular $f(R) - f(L-1)$.
\end{itemize}

\paragraph{Trampas y errores frecuentes.}
\begin{itemize}\setlength\itemsep{0pt}
	\item Olvidar resetear \texttt{memo} entre queries.
	\item Si la cantidad de estados es grande, el \texttt{memo[pos][tight][state]} puede ser muy grande; usar \texttt{map}.
	\item La condicion \texttt{started} es crucial para numeros con menos digitos que $D$.
	\item Si $N$ es muy grande ($> 10^{18}$), aniadir soporte para long long.
\end{itemize}

\paragraph{Codigo.}
Ver \texttt{cpp.tex}, seccion \textit{Digit DP template}.

\vspace{0.5em}

\subsubsection{Tree DP (reroot)}

\paragraph{Idea intuitiva.}
Calcular alguna propiedad asumiendo que \textbf{cada nodo es la raiz} del arbol. La tecnica clasica de \textbf{rerooting} en dos pasadas:
\begin{itemize}\setlength\itemsep{0pt}
	\item \textbf{DFS 1}: desde una raiz arbitraria, calcular \texttt{dp[v]} para el subarbol de $v$.
	\item \textbf{DFS 2}: para cada hijo $u$ de $v$, calcular \texttt{up[u]} = respuesta si la raiz fuera $u$.
\end{itemize}
La demostracion: tras la primera DFS, cada nodo conoce ``lo que aporta su subarbol''; en la segunda, se transmite ``lo que aporta el resto del arbol''.

\paragraph{Propiedades clave (invariantes).}
\begin{itemize}\setlength\itemsep{0pt}
	\item \texttt{dp[v]} depende solo del subarbol de $v$ (con la raiz original).
	\item \texttt{up[u]} se calcula de manera que cuando la raiz es $u$, ``todo lo que estaba en $v$ se ajusta''.
	\item Formula para reroot depende del problema; el snippet usa el ejemplo ``suma de distancias''.
	\item La combinacion \texttt{dp[u] + up[u]} da la respuesta con raiz $u$.
\end{itemize}

\paragraph{Codigo clave.}
Reroot para suma de distancias:
\begin{verbatim}
// dfs1: dp[v] = suma de distancias desde v al subarbol
void dfs1(int v, int p) {
    dp[v] = 0;
    for (int u : adj[v]) if (u != p) {
        dfs1(u, v);
        dp[v] += dp[u] + sz[u];  // cada hijo aporta su subarbol
    }
}
// dfs2: up[u] = distancia desde el resto del arbol a u
void dfs2(int v, int p) {
    for (int u : adj[v]) if (u != p) {
        up[u] = up[v] + dp[v] - dp[u] - sz[u] + (n - sz[u]);
        dfs2(u, v);
    }
}
\end{verbatim}

\paragraph{Complejidad.}
\begin{itemize}\setlength\itemsep{0pt}
	\item $O(N)$ para ambas pasadas.
	\item Cada arista se procesa un numero constante de veces.
\end{itemize}

\paragraph{Donde tocar para modificar.}
\begin{itemize}\setlength\itemsep{0pt}
	\item \textbf{Cambiar la cantidad a calcular}: modificar la operacion en \texttt{dfs1} y la formula en \texttt{dfs2}.
	\item \textbf{Arbol con pesos}: aniadir \texttt{weights[v]} y ajustar.
	\item \textbf{Re-root en arboles con ciclos}: requiere otra tecnica (DSU on tree).
	\item \textbf{Rooted trees diferentes}: si la propiedad no es invariante por raiz, usar otra tecnica.
\end{itemize}

\paragraph{Trampas y errores frecuentes.}
\begin{itemize}\setlength\itemsep{0pt}
	\item La formula de \texttt{up[u]} debe sumar y restar correctamente las contribuciones de $u$ y del resto.
	\item Para arboles con pesos, aniadir el peso en las formulas.
	\item Verificar con ejemplos pequenos antes de generalizar.
\end{itemize}

\paragraph{Codigo.}
Ver \texttt{cpp.tex}, seccion \textit{Tree DP (reroot)}.

% ============================================================
\section{Geometria}
% ============================================================

\subsection{Puntos y segmentos}

\subsubsection{Point + cross + dot}

\paragraph{Idea intuitiva.}
La estructura \texttt{Point} representa un punto en 2D con coordenadas enteras. Las dos operaciones vectoriales clave son:
\begin{itemize}\setlength\itemsep{0pt}
	\item \textbf{Cross product} $P \times Q = P_x Q_y - P_y Q_x$. Da el ``area con signo'' del paralelogramo formado por $P$ y $Q$. Usado para orientacion.
	\item \textbf{Dot product} $P \cdot Q = P_x Q_x + P_y Q_y$. Da el producto de las longitudes por el coseno del angulo. Usado para proyeccion.
\end{itemize}
La interpretacion geometrica del cross product: es el area con signo del paralelogramo; el signo indica la orientacion relativa.

\paragraph{Propiedades clave (invariantes).}
\begin{itemize}\setlength\itemsep{0pt}
	\item \texttt{cross(a, b) > 0} si $b$ esta a la izquierda de $a$ (sistema de coordenadas usual).
	\item \texttt{cross(a, b) = 0} si son colineales.
	\item \texttt{dot(a, b) > 0} si el angulo es agudo.
	\item El snippet define \texttt{cross} respecto a \texttt{this}, util para orientacion de 3 puntos.
	\item $|cross|^2 + |dot|^2 = |P|^2 \cdot |Q|^2$ (identidad de Lagrange).
\end{itemize}

\paragraph{Codigo clave.}
Las dos operaciones vectoriales:
\begin{verbatim}
long long cross(const Point& a, const Point& b) {
    return (long long)a.x * b.y - (long long)a.y * b.x;
}
long long dot(const Point& a, const Point& b) {
    return (long long)a.x * b.x + (long long)a.y * b.y;
}
\end{verbatim}

\paragraph{Complejidad.}
\begin{itemize}\setlength\itemsep{0pt}
	\item $O(1)$ por operacion.
	\item Constante muy pequena (4 multiplicaciones + 2 sumas).
\end{itemize}

\paragraph{Donde tocar para modificar.}
\begin{itemize}\setlength\itemsep{0pt}
	\item \textbf{Coordenadas doubles}: cambiar \texttt{int} por \texttt{double} o \texttt{long double}. Cuidado con la precision.
	\item \textbf{3D}: aniadir \texttt{z} y adaptar cross/dot.
	\item \textbf{Operadores adicionales}: aniadir \texttt{operator+}, \texttt{operator-}, \texttt{operator*} (escalar).
	\item \textbf{Distancia al cuadrado}: aniadir \texttt{len2()} que retorna $x^2 + y^2$.
\end{itemize}

\paragraph{Trampas y errores frecuentes.}
\begin{itemize}\setlength\itemsep{0pt}
	\item Cross product \textbf{no es} conmutativo; $P \times Q = -Q \times P$.
	\item Overflow con coordenadas grandes ($|x|, |y| > 10^4$) si se hace \texttt{x * y}.
	\item El signo de cross depende del sistema de coordenadas (en computacion, usualmente CCW positivo).
	\item Para ``igualdad'' de doubles, aniadir tolerancia \texttt{eps}.
\end{itemize}

\paragraph{Codigo.}
Ver \texttt{cpp.tex}, seccion \textit{Point + cross + dot}.

\vspace{0.5em}

\subsubsection{Interseccion de segmentos}

\paragraph{Idea intuitiva.}
Dados dos segmentos $[a, b]$ y $[c, d]$, determinar si se intersectan. La idea clasica:
\begin{itemize}\setlength\itemsep{0pt}
	\item Verificar orientaciones: $a, b, c$ y $a, b, d$ deben estar en lados opuestos de la recta $ab$; igual para $cd$ y $ab$.
	\item Caso colineal: si las rectas son paralelas (\texttt{cross} $= 0$) y los segmentos son colineales, verificar que los bounding boxes se intersectan.
\end{itemize}
La demostracion: dos segmentos se intersectan sii cada segmento ``cruza'' la linea que contiene al otro (estricto) o ambos son colineales con solapamiento.

\paragraph{Propiedades clave (invariantes).}
\begin{itemize}\setlength\itemsep{0pt}
	\item Funciona con segmentos en cualquier orientacion.
	\item Caso colineal: requiere check adicional de bounding box.
	\item Para intersectar en un \textbf{punto} especifico: aniadir division con cuidado de precision.
	\item Funciona con segmentos verticales/horizontales sin casos especiales.
\end{itemize}

\paragraph{Codigo clave.}
Test clasico de interseccion:
\begin{verbatim}
int sign(long long x) { return (x > 0) - (x < 0); }
bool intersect(Point a, Point b, Point c, Point d) {
    long long o1 = cross(b - a, c - a);
    long long o2 = cross(b - a, d - a);
    long long o3 = cross(d - c, a - c);
    long long o4 = cross(d - c, b - c);
    if (sign(o1) != sign(o2) && sign(o3) != sign(o4)) return true;
    // casos colineales
    if (o1 == 0 && onSegment(a, b, c)) return true;
    if (o2 == 0 && onSegment(a, b, d)) return true;
    if (o3 == 0 && onSegment(c, d, a)) return true;
    if (o4 == 0 && onSegment(c, d, b)) return true;
    return false;
}
\end{verbatim}

\paragraph{Complejidad.}
\begin{itemize}\setlength\itemsep{0pt}
	\item $O(1)$.
	\item Constante pequena.
\end{itemize}

\paragraph{Donde tocar para modificar.}
\begin{itemize}\setlength\itemsep{0pt}
	\item \textbf{Interseccion con punto}: aniadir la division despues de verificar orientacion.
	\item \textbf{Interseccion rayo-segmento}: ajustar la primera condicion.
	\item \textbf{Poligonos convexos}: usar separacion de ejes (SAT).
	\item \textbf{Multiples queries}: precomputar bounding boxes o usar segment tree.
\end{itemize}

\paragraph{Trampas y errores frecuentes.}
\begin{itemize}\setlength\itemsep{0pt}
	\item Si los segmentos tocan solo en un extremo, el snippet lo considera interseccion (algunos problemas quieren estricto). Ajustar segun el problema.
	\item El caso colineal es comun; olvidar el check adicional de bounding box da falsos negativos.
	\item En problemas con muchos segmentos, considerar sweep line para $O(n \log n)$.
\end{itemize}

\paragraph{Codigo.}
Ver \texttt{cpp.tex}, seccion \textit{Interseccion de segmentos}.

\vspace{0.5em}

\subsubsection{Contains colineal}

\paragraph{Idea intuitiva.}
Verifica si un punto $c$ esta sobre el segmento $[a, b]$ (asumiendo que ya son colineales). Si \texttt{cross(a, b, c) = 0} y $c$ esta dentro del bounding box de $[a, b]$, entonces $c$ esta en el segmento. La razon: si los tres puntos son colineales, $c$ esta en el segmento sii sus coordenadas estan entre las de $a$ y $b$.

\paragraph{Propiedades clave (invariantes).}
\begin{itemize}\setlength\itemsep{0pt}
	\item Asume colinealidad (no la verifica mas alla de \texttt{cross == 0}).
	\item Usa comparaciones con \texttt{<=} para extremos \textbf{inclusivos}.
	\item Funciona con segmentos verticales (la condicion se reduce a una sola coordenada).
\end{itemize}

\paragraph{Codigo clave.}
Test de contencion en segmento (asume colinealidad):
\begin{verbatim}
bool onSegment(Point a, Point b, Point c) {
    return min(a.x, b.x) <= c.x && c.x <= max(a.x, b.x) &&
           min(a.y, b.y) <= c.y && c.y <= max(a.y, b.y);
}
\end{verbatim}

\paragraph{Complejidad.}
\begin{itemize}\setlength\itemsep{0pt}
	\item $O(1)$.
	\item Constante minima.
\end{itemize}

\paragraph{Donde tocar para modificar.}
\begin{itemize}\setlength\itemsep{0pt}
	\item \textbf{Contener estricto}: cambiar \texttt{<=} a \texttt{<} para excluir extremos.
	\item \textbf{Coordenadas doubles}: usar tolerancia (\texttt{eps}).
	\item \textbf{Solo x o solo y}: aniadir check especifico para segmentos verticales/horizontales.
\end{itemize}

\paragraph{Trampas y errores frecuentes.}
\begin{itemize}\setlength\itemsep{0pt}
	\item No usar este snippet para puntos \textbf{no colineales}; siempre verificar primero la colinealidad.
	\item Si los puntos tienen doubles, el chequeo \texttt{cross == 0} puede fallar; aniadir tolerancia.
	\item El bounding box incluye los extremos; verificar si el problema quiere exclusivo o inclusivo.
\end{itemize}

\paragraph{Codigo.}
Ver \texttt{cpp.tex}, seccion \textit{Contains colineal}.

\subsection{Poligonos}

\subsubsection{Convex Hull (monotono)}

\paragraph{Idea intuitiva.}
Algoritmo de \textbf{Andrew's monotone chain}: ordenar puntos por $(x, y)$, luego construir la hull inferior y superior por separado. Para cada punto, aniadirlo al hull y, mientras el ultimo giro no sea ``counter-clockwise'' (\texttt{cross > 0}), popear. La demostracion: si los puntos se procesan en orden, el hull inferior y superior cubren la frontera completa sin auto-intersecciones.

\paragraph{Propiedades clave (invariantes).}
\begin{itemize}\setlength\itemsep{0pt}
	\item Ordena los puntos en $O(n \log n)$.
	\item Construye la hull en $O(n)$ despues.
	\item Output: poligono convexo en counter-clockwise, sin punto final duplicado.
	\item El \texttt{<= 0} incluye colineales; usar \texttt{< 0} para excluirlos.
	\item Cada punto aparece en la hull inferior o superior (no ambos).
\end{itemize}

\paragraph{Codigo clave.}
Construccion del hull inferior y superior:
\begin{verbatim}
sort(p.begin(), p.end());
vector<Point> lower, upper;
for (auto& p : points) {
    while (lower.size() >= 2 && cross(lower.back() - lower[lower.size()-2],
                                       p - lower.back()) <= 0)
        lower.pop_back();
    lower.push_back(p);
}
for (int i = points.size() - 1; i >= 0; --i) {
    auto& p = points[i];
    while (upper.size() >= 2 && cross(upper.back() - upper[upper.size()-2],
                                       p - upper.back()) <= 0)
        upper.pop_back();
    upper.push_back(p);
}
vector<Point> hull = lower;
hull.insert(hull.end(), upper.begin() + 1, upper.end() - 1);
\end{verbatim}

\paragraph{Complejidad.}
\begin{itemize}\setlength\itemsep{0pt}
	\item $O(n \log n)$.
	\item Constante muy pequena (solo cross products).
\end{itemize}

\paragraph{Donde tocar para modificar.}
\begin{itemize}\setlength\itemsep{0pt}
	\item \textbf{Colineales}: cambiar \texttt{<= 0} por \texttt{< 0} para excluirlos del hull.
	\item \textbf{Puntos repetidos}: \texttt{unique} despues de sort.
	\item \textbf{Hull 3D}: gift wrapping es $O(n^2)$; convex hull 3D es mas complejo.
	\item \textbf{Output con los puntos}: aniadir el hull completo (no solo upper/lower).
\end{itemize}

\paragraph{Trampas y errores frecuentes.}
\begin{itemize}\setlength\itemsep{0pt}
	\item El \texttt{pop\_back} y el \texttt{reverse} son esenciales para no duplicar el ultimo punto.
	\item Para excluir colineales, ajustar el \texttt{<= 0} vs \texttt{< 0}.
	\item Si los puntos estan duplicados, aniadir \texttt{unique} antes del sort.
\end{itemize}

\paragraph{Codigo.}
Ver \texttt{cpp.tex}, seccion \textit{Convex Hull (monotono)}.

\vspace{0.5em}

\subsubsection{Polygon Area}

\paragraph{Idea intuitiva.}
El area (con signo) de un poligono se calcula con la \textbf{Shoelace formula}: $A = \frac{1}{2} \sum_{i=0}^{n-1} (x_i y_{i+1} - x_{i+1} y_i)$, donde el indice es modulo $n$. El snippet retorna el \textbf{doble} del area (en valor absoluto). La formula sale del ``sumar areas de trapezoides'': cada lado del poligono define un trapezoide con el eje $x$; la suma con signo da el area.

\paragraph{Propiedades clave (invariantes).}
\begin{itemize}\setlength\itemsep{0pt}
	\item El signo indica la orientacion: positivo si counter-clockwise.
	\item Para el doble: no dividir por 2 si no es necesario (evita fracciones).
	\item Funciona para poligonos simples (no auto-intersectantes).
	\item Si el poligono es no convexo, la formula da el area ``con signo''.
\end{itemize}

\paragraph{Codigo clave.}
La suma de trapezoides:
\begin{verbatim}
long long shoelace(const vector<Point>& p) {
    long long s = 0;
    int n = p.size();
    for (int i = 0; i < n; ++i)
        s += (long long)p[i].x * p[(i+1)%n].y - (long long)p[(i+1)%n].x * p[i].y;
    return abs(s);  // doble del area
}
\end{verbatim}

\paragraph{Complejidad.}
\begin{itemize}\setlength\itemsep{0pt}
	\item $O(n)$.
	\item Constante minima.
\end{itemize}

\paragraph{Donde tocar para modificar.}
\begin{itemize}\setlength\itemsep{0pt}
	\item \textbf{Area exacta}: dividir por 2 (con cuidado si da fraccion; usar \texttt{double}).
	\item \textbf{Area firmada}: quitar \texttt{abs} y dejar el signo.
	\item \textbf{Punto dentro de poligono convexo}: con la hull, chequear que esta a un mismo lado de todas las aristas.
	\item \textbf{Triangulacion}: descomponer en triangulos y sumar areas.
\end{itemize}

\paragraph{Trampas y errores frecuentes.}
\begin{itemize}\setlength\itemsep{0pt}
	\item Si los puntos no estan en orden (CCW o CW), el area firmada puede ser incorrecta. Para signed area consistente, usar siempre el mismo orden.
	\item Para poligonos con holes, sumar las areas con signo apropiado.
	\item Con doubles, aniadir tolerancia \texttt{eps}.
\end{itemize}

\paragraph{Codigo.}
Ver \texttt{cpp.tex}, seccion \textit{Polygon Area}.

\vspace{0.5em}

\subsubsection{Distance point-segment / point-line}

\paragraph{Idea intuitiva.}
Distancia minima de un punto $p$ a un segmento $[a, b]$ o a la recta que pasa por $a$ y $b$.
\begin{itemize}\setlength\itemsep{0pt}
	\item \textbf{Linea}: $d = \frac{|ab \times ap|}{|ab|}$.
	\item \textbf{Segmento}: si la proyeccion de $p$ sobre $ab$ cae fuera del segmento, la distancia es al extremo mas cercano.
\end{itemize}
La razon geometrica: la distancia a una linea es el area del paralelogramo dividido por la base; la proyeccion ortogonal indica si estamos ``dentro'' o ``fuera'' del segmento.

\paragraph{Propiedades clave (invariantes).}
\begin{itemize}\setlength\itemsep{0pt}
	\item \texttt{dist2PointLine} retorna distancia \textbf{al cuadrado}.
	\item \texttt{dist2PointSegment} distingue 3 casos: $p$ antes de $a$, despues de $b$, o proyectado sobre el segmento.
	\item Si \texttt{len2 = 0} (degenerate), la distancia es $|ap|$ o $|bp|$.
	\item La proyeccion se calcula con dot products y el parametro $t \in [0, 1]$.
\end{itemize}

\paragraph{Codigo clave.}
Distancia al cuadrado a segmento:
\begin{verbatim}
long long dist2Segment(Point p, Point a, Point b) {
    long long len2 = (a-b).len2();
    if (len2 == 0) return (p-a).len2();  // degenerate
    long long t = max(0LL, min(1LL, dot(p-a, b-a) / len2));
    Point proj = {a.x + t*(b.x - a.x), a.y + t*(b.y - a.y)};
    return (p - proj).len2();
}
\end{verbatim}

\paragraph{Complejidad.}
\begin{itemize}\setlength\itemsep{0pt}
	\item $O(1)$.
	\item Constante minima.
\end{itemize}

\paragraph{Donde tocar para modificar.}
\begin{itemize}\setlength\itemsep{0pt}
	\item \textbf{Distancia (no al cuadrado)}: aniadir \texttt{sqrt} al final.
	\item \textbf{Coordenadas doubles}: aniadir \texttt{long double} y tolerancia.
	\item \textbf{3D}: aniadir coordenada $z$ y adaptar.
	\item \textbf{Multiples queries}: usar KD-tree o segment tree para acelerar.
\end{itemize}

\paragraph{Trampas y errores frecuentes.}
\begin{itemize}\setlength\itemsep{0pt}
	\item Overflow en \texttt{area2 * area2} si el area es grande; \texttt{long long} alcanza para $|x|, |y| \le 10^9$ pero no mas.
	\item Si $a = b$ (segmento degenerado), aniadir check explicito.
	\item En doubles, aniadir tolerancia \texttt{eps} para comparar.
\end{itemize}

\paragraph{Codigo.}
Ver \texttt{cpp.tex}, seccion \textit{Distance point-segment / point-line}.

\vspace{0.5em}

\subsubsection{Projection of point onto line}

\paragraph{Idea intuitiva.}
Proyeccion ortogonal de $p$ sobre la recta que pasa por $a$ y $b$. El parametro $t$ tal que la proyeccion es $a + t \cdot (b - a)$ es:
\[ t = \frac{(p - a) \cdot (b - a)}{|b - a|^2}. \]
La idea: descomponemos $\vec{ap}$ en una componente paralela a $\vec{ab}$ y una perpendicular; $t$ mide la posicion en la primera.

\paragraph{Propiedades clave (invariantes).}
\begin{itemize}\setlength\itemsep{0pt}
	\item Si $t \in [0, 1]$, la proyeccion cae sobre el \textbf{segmento}; si no, cae sobre la extension.
	\item Retorna \texttt{double} (puede tener precision issues).
	\item $t < 0$: $p$ esta ``antes'' de $a$; $t > 1$: $p$ esta ``despues'' de $b$.
\end{itemize}

\paragraph{Codigo clave.}
Proyeccion ortogonal:
\begin{verbatim}
Point projection(Point p, Point a, Point b) {
    Point ab = b - a, ap = p - a;
    double t = (double)dot(ap, ab) / ab.len2();
    return {a.x + t * ab.x, a.y + t * ab.y};
}
\end{verbatim}

\paragraph{Complejidad.}
\begin{itemize}\setlength\itemsep{0pt}
	\item $O(1)$.
	\item Constante minima (dos dot products y una division).
\end{itemize}

\paragraph{Donde tocar para modificar.}
\begin{itemize}\setlength\itemsep{0pt}
	\item \textbf{Proyeccion sobre segmento}: clipar $t$ a $[0, 1]$ antes de calcular el punto.
	\item \textbf{Distancia al segmento via proyeccion}: combinar con el modulo anterior.
	\item \textbf{Solo el parametro $t$}: si solo necesitas saber si esta dentro del segmento, retorna $t$ sin construir el punto.
\end{itemize}

\paragraph{Trampas y errores frecuentes.}
\begin{itemize}\setlength\itemsep{0pt}
	\item Si $a = b$ ($|ab| = 0$), la proyeccion no esta definida; aniadir guard.
	\item En doubles, aniadir tolerancia \texttt{eps} para comparar $t$ con 0 o 1.
	\item La division por \texttt{len2} puede ser cara si las coordenadas son grandes; usar \texttt{double} o precomputar.
\end{itemize}

\paragraph{Codigo.}
Ver \texttt{cpp.tex}, seccion \textit{Projection of point onto line}.

\vspace{0.5em}

\subsubsection{Point in polygon (ray casting)}

\paragraph{Idea intuitiva.}
Algoritmo de \textbf{ray casting}: contar cuantas veces un rayo horizontal desde $p$ cruza los bordes del poligono. Si es impar, $p$ esta dentro; si par, fuera. El snippet implementa la version clasica recorriendo aristas. La razon: cada cruce cambia el lado del rayo respecto al poligono; un numero impar de cambios significa ``terminamos dentro''.

\paragraph{Propiedades clave (invariantes).}
\begin{itemize}\setlength\itemsep{0pt}
	\item Funciona con poligonos simples (no necesariamente convexos).
	\item Las intersecciones en vertices cuentan especialmente.
	\item Para poligonos con \textbf{huecos}, sigue funcionando.
	\item La respuesta es estrictamente 0 (fuera) o 1 (dentro) para poligonos simples.
\end{itemize}

\paragraph{Codigo clave.}
El test clasico de ray casting:
\begin{verbatim}
bool pointInPolygon(Point p, const vector<Point>& poly) {
    int n = poly.size();
    bool inside = false;
    for (int i = 0, j = n - 1; i < n; j = i++) {
        if ((poly[i].y > p.y) != (poly[j].y > p.y) &&
            p.x < (poly[j].x - poly[i].x) * (p.y - poly[i].y) /
                    (poly[j].y - poly[i].y) + poly[i].x)
            inside = !inside;
    }
    return inside;
}
\end{verbatim}

\paragraph{Complejidad.}
\begin{itemize}\setlength\itemsep{0pt}
	\item $O(n)$ por query.
	\item Memoria $O(1)$ (no se almacena el rayo).
\end{itemize}

\paragraph{Donde tocar para modificar.}
\begin{itemize}\setlength\itemsep{0pt}
	\item \textbf{Winding number}: version alternativa que cuenta el numero de vueltas.
	\item \textbf{Poligono convexo}: check lineal contra las aristas (mas rapido).
	\item \textbf{Multiples queries}: precomputar o usar segment tree sobre las aristas.
	\item \textbf{Punto en frontera}: aniadir check adicional si es colineal con alguna arista.
\end{itemize}

\paragraph{Trampas y errores frecuentes.}
\begin{itemize}\setlength\itemsep{0pt}
	\item El manejo del vertice (cuando el rayo pasa por un vertice) puede causar double-counting; el snippet lo evita con la condicion \texttt{(poly[i].y > p.y) != (poly[j].y > p.y)}.
	\item Si el poligono tiene holes, el algoritmo sigue funcionando.
	\item Para poligonos no simples (auto-intersectantes), el resultado puede ser ambiguo.
\end{itemize}

\paragraph{Codigo.}
Ver \texttt{cpp.tex}, seccion \textit{Point in polygon (ray casting)}.

\vspace{0.5em}

\subsubsection{Minkowski sum (convex polygons)}

\paragraph{Idea intuitiva.}
La \textbf{suma de Minkowski} $A \oplus B$ de dos poligonos convexos es otro poligono convexo que representa $\{a + b : a \in A, b \in B\}$. Algoritmo: ordenar las aristas por angulo polar (CCW) y ``mergearlas'' como en merge sort; el resultado es un poligono CCW sin repetir el ultimo punto. La idea: el resultado tiene una arista paralela a cada arista de los originales; la suma se obtiene barriendo los angulos.

\paragraph{Propiedades clave (invariantes).}
\begin{itemize}\setlength\itemsep{0pt}
	\item Si $A, B$ son convexos con $n, m$ vertices, $A \oplus B$ tiene $\le n + m$ vertices.
	\item Util para \textbf{detectar colision}: $A$ colisiona con $B$ si $0 \in A \oplus (-B)$.
	\item Asume poligonos en orden \textbf{counter-clockwise}.
	\item El resultado es convexo (suma de convexos es convexo).
\end{itemize}

\paragraph{Codigo clave.}
La suma de Minkowski:
\begin{verbatim}
vector<Point> minkowski(const vector<Point>& A, const vector<Point>& B) {
    vector<Point> result;
    int i = 0, j = 0;
    int n = A.size(), m = B.size();
    do {
        result.push_back(A[i] + B[j]);
        long long cross = (A[(i+1)%n] - A[i]) ^ (B[(j+1)%m] - B[j]);
        if (cross >= 0) i = (i + 1) % n;
        if (cross <= 0) j = (j + 1) % m;
    } while (i || j);
    return result;
}
\end{verbatim}

\paragraph{Complejidad.}
\begin{itemize}\setlength\itemsep{0pt}
	\item $O(n + m)$.
	\item Memoria $O(n + m)$.
\end{itemize}

\paragraph{Donde tocar para modificar.}
\begin{itemize}\setlength\itemsep{0pt}
	\item \textbf{Deteccion de colision}: aniadir $(-B)$ (reflejar) y verificar si $0$ esta en el resultado.
	\item \textbf{Poligonos no convexos}: la suma puede no ser poligono simple; algoritmo distinto.
	\item \textbf{Multiples queries}: precomputar las aristas una vez.
	\item \textbf{Robot motion planning}: usar para planear trayectorias.
\end{itemize}

\paragraph{Trampas y errores frecuentes.}
\begin{itemize}\setlength\itemsep{0pt}
	\item Si los poligonos no son CCW, la suma es incorrecta. Normalizar antes.
	\item El bucle termina cuando ambos indices vuelven a 0; aniadir check.
	\item Para la deteccion de colision, aniadir el test \texttt{0 in A \^{} (-B)}.
\end{itemize}

\paragraph{Codigo.}
Ver \texttt{cpp.tex}, seccion \textit{Minkowski sum (convex polygons)}.

\subsection{Herramientas}

\subsubsection{Li Chao Tree (long long)}

\paragraph{Idea intuitiva.}
Mantiene un \textbf{lower envelope} (o upper) de un conjunto de rectas $y = mx + b$ sobre un dominio $[X_{LO}, X_{HI}]$ y responde ``minimo $y$ en un $x$ dado'' en $O(\log X)$. La estructura es un segment tree donde cada nodo guarda la recta ``mejor'' para su segmento medio; al aniadir una recta, se compara con la actual y se decide cual es mejor en cada mitad, descendiendo recursivamente. La demostracion: en cada nivel, solo se guarda la mejor recta para el ``punto medio''; las otras se transmiten a los hijos.

\paragraph{Propiedades clave (invariantes).}
\begin{itemize}\setlength\itemsep{0pt}
	\item \textbf{Online}: cada \texttt{addLine} en $O(\log X)$.
	\item Solo insercion; no hay borrar (para borrar, usar Li Chao segmentado con reset).
	\item La recta guardada en un nodo es la ``mejor'' en su punto medio.
	\item Las rectas ``perdedoras'' se transmiten a los hijos para futuros inserts.
\end{itemize}

\paragraph{Codigo clave.}
La insercion en Li Chao:
\begin{verbatim}
void addLine(Line nw, int v = 1, int l = X_LO, int r = X_HI) {
    int m = (l + r) / 2;
    bool left = nw.get(l) < tree[v].get(l);
    bool mid = nw.get(m) < tree[v].get(m);
    if (mid) swap(tree[v], nw);
    if (r - l == 1) return;
    if (left != mid) addLine(nw, v*2, l, m);
    else             addLine(nw, v*2+1, m, r);
}
\end{verbatim}

\paragraph{Complejidad.}
\begin{itemize}\setlength\itemsep{0pt}
	\item $O(\log X)$ por insercion y query, donde $X$ es el tamano del dominio.
	\item Memoria $O(\log X)$ por insercion, $O(X)$ total.
\end{itemize}

\paragraph{Donde tocar para modificar.}
\begin{itemize}\setlength\itemsep{0pt}
	\item \textbf{Upper envelope en vez de lower}: invertir el signo (\texttt{>} en vez de \texttt{<}).
	\item \textbf{Queries de segmento} (no full line): si la recta solo esta activa en un rango, aniadir bandera \texttt{active}.
	\item \textbf{Eliminar rectas}: usar Li Chao segmentado (mas complejo).
	\item \textbf{Pendientes grandes}: si las pendientes tienen magnitud $> 10^9$, aniadir \texttt{\_\_int128} para evitar overflow.
\end{itemize}

\paragraph{Trampas y errores frecuentes.}
\begin{itemize}\setlength\itemsep{0pt}
	\item El parametro \texttt{lo == hi} debe retornar el valor del nodo (caso base).
	\item La comparacion \texttt{newLine.get(mid) < node->line.get(mid)} decide quien es mejor en el medio.
	\item Si las rectas tienen $m$ iguales pero $b$ distintos, se queda con la primera (tie-break).
\end{itemize}

\paragraph{Codigo.}
Ver \texttt{cpp.tex}, seccion \textit{Li Chao Tree (long long)}.

\vspace{0.5em}

\subsubsection{Sweep line template}

\paragraph{Idea intuitiva.}
Patron de diseno para problemas geometricos donde se barre una linea sobre el plano y se mantiene una estructura de datos sobre el otro eje. Eventos (add, query, remove) se ordenan por la coordenada de barrido y se procesan secuencialmente.

\paragraph{Propiedades clave (invariantes).}
\begin{itemize}\setlength\itemsep{0pt}
	\item Eventos: tuplas \texttt{(x, type, id)}.
	\item La estructura auxiliar (set, multiset, segtree) se mantiene sobre el otro eje.
	\item Procesar eventos en orden de $x$ ascendente.
\end{itemize}

\paragraph{Complejidad.}
\begin{itemize}\setlength\itemsep{0pt}
	\item Depende del problema; tipico $O((N + E) \log N)$.
\end{itemize}

\paragraph{Donde tocar para modificar.}
\begin{itemize}\setlength\itemsep{0pt}
	\item \textbf{Aniadir/quitar segmentos}: aniadir eventos add/remove al barrido.
	\item \textbf{Aniadir queries}: aniadir eventos query.
	\item \textbf{Cambiar el eje de barrido}: si barres en $y$ en vez de $x$, aniadir comparador.
\end{itemize}

\paragraph{Trampas y errores frecuentes.}
\begin{itemize}\setlength\itemsep{0pt}
	\item El snippet es solo el \textbf{esqueleto}; los handlers \texttt{add/query/remove} se llenan por problema.
	\item Cuidado con eventos duplicados.
\end{itemize}

\paragraph{Codigo.}
Ver \texttt{cpp.tex}, seccion \textit{Sweep line template}.

\vspace{0.5em}

\subsubsection{Closest pair of points}

\paragraph{Idea intuitiva.}
Encuentra el par de puntos mas cercano en $O(n \log n)$ con \textbf{divide y venceras}. Algoritmo:
\begin{itemize}\setlength\itemsep{0pt}
	\item Sort por $x$.
	\item Dividir en mitades; recursivo.
	\item Combinar: dado el minimo $d$ de las dos mitades, strip = puntos en $[midX - d, midX + d]$; ordenar por $y$; cada punto solo se compara con los siguientes hasta $y$ difiera en $< d$ (maximo 7).
\end{itemize}

\paragraph{Propiedades clave (invariantes).}
\begin{itemize}\setlength\itemsep{0pt}
	\item Requiere ordenar por $y$ despues (en el snippet se hace durante el merge).
	\item Strip tiene a lo sumo $O(n)$ puntos en total, pero la comparacion es $O(1)$ amortizada por punto.
	\item Resultado en \texttt{double}.
\end{itemize}

\paragraph{Complejidad.}
\begin{itemize}\setlength\itemsep{0pt}
	\item $O(n \log n)$.
\end{itemize}

\paragraph{Donde tocar para modificar.}
\begin{itemize}\setlength\itemsep{0pt}
	\item \textbf{Reconstruir el par}: aniadir indices en la recursion.
	\item \textbf{Distancia maxima}: cambiar \texttt{min} por \texttt{max} (mas facil; no requiere D\&C).
	\item \textbf{Puntos duplicados}: aniadir check especial (distancia 0).
\end{itemize}

\paragraph{Trampas y errores frecuentes.}
\begin{itemize}\setlength\itemsep{0pt}
	\item El merge requiere ordenar por $y$; el snippet lo hace in-place con un \texttt{tmp}.
	\item Si se hace en cada recursion, es $O(n \log n)$ extra.
\end{itemize}

\paragraph{Codigo.}
Ver \texttt{cpp.tex}, seccion \textit{Closest pair of points}.

\vspace{0.5em}

\subsubsection{Rotating calipers (polygon diameter)}

\paragraph{Idea intuitiva.}
Encuentra el \textbf{diametro} de un poligono convexo (par de vertices mas distantes) en $O(n)$. Asume poligono en orden counter-clockwise. Algoritmo: dos indices $i$ (sobre las aristas) y $k$ (sobre los vertices opuestos); en cada paso, rotar $k$ mientras el area del paralelogramo aumenta. La distancia $i$-$k$ puede ser maxima en cualquier momento del barrido.

\paragraph{Propiedades clave (invariantes).}
\begin{itemize}\setlength\itemsep{0pt}
	\item Cada indice avanza como maximo $n$ veces total.
	\item El ``diametro'' es la maxima distancia entre dos vertices.
	\item Retorna \textbf{distancia al cuadrado}.
\end{itemize}

\paragraph{Complejidad.}
\begin{itemize}\setlength\itemsep{0pt}
	\item $O(n)$.
\end{itemize}

\paragraph{Donde tocar para modificar.}
\begin{itemize}\setlength\itemsep{0pt}
	\item \textbf{Reconstruir el par de vertices}: aniadir indices.
	\item \textbf{Anchura minima / maxima}: variantes de calipers.
	\item \textbf{Distancia minima entre poligonos convexos}: otra variante.
\end{itemize}

\paragraph{Trampas y errores frecuentes.}
\begin{itemize}\setlength\itemsep{0pt}
	\item Asume poligono \textbf{estrictamente} convexo (sin vertices colineales). Si hay colineales, aniadir perturbacion o filtrar.
\end{itemize}

\paragraph{Codigo.}
Ver \texttt{cpp.tex}, seccion \textit{Rotating calipers (polygon diameter)}.

% ============================================================
\section{Miscelanea}
% ============================================================

\subsection{Utilidades del usuario}

\subsubsection{Hora}

\paragraph{Idea intuitiva.}
Una pequena clase para representar \textbf{tiempos en segundos} y hacer aritmetica basica sobre ellos. Internamente almacena todo en segundos (un solo \texttt{int}), lo cual simplifica comparaciones y diferencias. El constructor acepta $(h, m, s)$ y los convierte a segundos. Hay un metodo \texttt{cut()} que fuerza un minimo (util para ``horas trabajadas'').

\paragraph{Propiedades clave (invariantes).}
\begin{itemize}\setlength\itemsep{0pt}
	\item Representacion canonica: \textbf{solo segundos}. No hay zonas horarias, no hay DST.
	\item Comparadores \texttt{<} y \texttt{>} operan sobre los segundos.
	\item \texttt{operator-} retorna la \textbf{diferencia absoluta} en segundos (siempre no negativa).
	\item \texttt{cut()} aplica un piso (por defecto $8\text{h }30\text{m} = 30600$ segundos).
\end{itemize}

\paragraph{Complejidad.}
\begin{itemize}\setlength\itemsep{0pt}
	\item $O(1)$ por operacion.
\end{itemize}

\paragraph{Donde tocar para modificar.}
\begin{itemize}\setlength\itemsep{0pt}
	\item \textbf{Aniadir suma}: \texttt{operator+} que devuelve un \texttt{Hour} con la suma de segundos.
	\item \textbf{Aniadir output formateado}: aniadir \texttt{show\_hms()} que imprime $H$:$M$:$S$.
	\item \textbf{Cambiar el ``corte''}: modificar el valor magico \texttt{30'600} en \texttt{cut()} o pasarlo como parametro.
	\item \textbf{Tamano maximo}: si los tiempos exceden $\sim 68$ anos, el \texttt{int} desborda; usar \texttt{long long}.
	\item \textbf{Precision sub-segundo}: aniadir \texttt{ms} (milisegundos) y guardar en \texttt{long long} con factor 1000.
\end{itemize}

\paragraph{Trampas y errores frecuentes.}
\begin{itemize}\setlength\itemsep{0pt}
	\item El operador \texttt{-} retorna \textbf{valor absoluto}, no diferencia con signo. Si necesitas diferencia con signo, aniadir un nuevo metodo o sobrecargar.
	\item \texttt{cut()} modifica \texttt{s} \textbf{in-place}; si el llamador no espera esto, aniadir una version \texttt{cut(min)} que retorna un nuevo \texttt{Hour}.
	\item El umbral magico \texttt{30'600} (8h 30m) esta hardcoded; documentar bien que es ``jornada minima''.
\end{itemize}

\paragraph{Codigo.}
Ver \texttt{cpp.tex}, seccion \textit{Hora}.

\end{document}
' se responden en $O(|P|)$.
\end{itemize}

\paragraph{Trampas y errores frecuentes.}
\begin{itemize}\setlength\itemsep{0pt}
	\item El \textbf{clone} debe tener \texttt{occ = 0} (no es una nueva posicion final); el \texttt{link} del estado siendo creado y del estado clonado deben apuntar al clon.
	\item Saltarse esto da SAM incorrecto.
	\item La condicion \texttt{st[p].len + 1 == st[q].len} detecta cuando no se necesita clon.
	\item Olvidar \texttt{while (p != -1 && !st[p].next.count(c))} da transiciones rotas.
\end{itemize}

\paragraph{Codigo.}
Ver \texttt{cpp.tex}, seccion \textit{Suffix Automaton}.

\vspace{0.5em}

\subsubsection{Lyndon factorization (Duval)}

\paragraph{Idea intuitiva.}
Un \textbf{string de Lyndon} es un string que es estrictamente menor que todos sus rotaciones. Toda cadena tiene una \textbf{factorizacion unica} de Lyndon (algoritmo de \textbf{Duval}): una secuencia de strings de Lyndon $L_1, L_2, \dots, L_k$ tales que $L_1 \ge L_2 \ge \dots \ge L_k$ y la concatenacion es el string original. Esto da $O(n)$ para encontrar los cortes. La demostracion de correctitud: cada factor es minimal en su clase de rotacion, asi que la concatenacion respeta el orden lexicografico.

\paragraph{Propiedades clave (invariantes).}
\begin{itemize}\setlength\itemsep{0pt}
	\item Cada factor de Lyndon es de la forma $uv$ donde $u$ se repite y $v$ es prefijo de $u$.
	\item La longitud del factor esta dada por la ``longitud del periodo''.
	\item Aplicacion: el ``Lyndon word'' es el menor en orden lexicografico entre todas las rotaciones.
	\item La factorizacion es \textbf{unica}.
\end{itemize}

\paragraph{Codigo clave.}
El algoritmo de Duval (iterativo):
\begin{verbatim}
int n = s.size();
int i = 0;
while (i < n) {
    int j = i + 1, k = i;
    while (j < n && s[k] <= s[j]) {
        if (s[k] < s[j]) k = i;
        else ++k;
        ++j;
    }
    int len = j - k;
    while (i < j) { cout << "factor: " << s.substr(i, len) << "\n"; i += len; }
}
\end{verbatim}

\paragraph{Complejidad.}
\begin{itemize}\setlength\itemsep{0pt}
	\item $O(n)$.
	\item Cada comparacion es $O(1)$; el puntero $k$ retrocede a lo sumo $n$ veces en total.
\end{itemize}

\paragraph{Donde tocar para modificar.}
\begin{itemize}\setlength\itemsep{0pt}
	\item \textbf{Encontrar la menor rotacion}: el menor sufijo lexicografico de \texttt{s + s} es la menor rotacion; la posicion se obtiene con un solo pase de Duval.
	\item \textbf{Algoritmo de Booth}: alternativa para rotacion minima, tambien $O(n)$.
	\item \textbf{Validar}: si necesitas verificar que un string es de Lyndon, compararlo con sus rotaciones.
	\item \textbf{Encontrar todos los factores}: guardar la posicion y longitud de cada uno.
\end{itemize}

\paragraph{Trampas y errores frecuentes.}
\begin{itemize}\setlength\itemsep{0pt}
	\item La condicion \texttt{s[j] <= s[k]} (no \texttt{<}) en el while de Duval es crucial; cambiar a \texttt{<} da comportamiento distinto.
	\item El caso $s[k] < s[j]$ resetea $k$ a $i$ (inicio del factor actual).
	\item El bucle externo debe avanzar $i$ por la longitud del factor, no por 1.
\end{itemize}

\paragraph{Codigo.}
Ver \texttt{cpp.tex}, seccion \textit{Lyndon factorization (Duval)}.

\subsection{Palindromos}

\subsubsection{Manacher}

\paragraph{Idea intuitiva.}
Calcula en $O(n)$ los \textbf{radios} de todos los palindromos centrados en cada posicion (impares y pares). La idea: mantener una ventana $[l, r]$ que es el palindromo centrado en un punto anterior mas largo encontrado. Si el nuevo centro $i$ cae dentro de $[l, r]$, entonces el palindromo centrado en $i$ tiene al menos el radio del palindromo ``espejo'' $l + r - i$; sino, empezar desde 1 (impar) o 0 (par). Despues, expandir lo que falte. La demostracion de $O(n)$: cada caracter se compara a lo sumo 2 veces (extension + limite por ventana).

\paragraph{Propiedades clave (invariantes).}
\begin{itemize}\setlength\itemsep{0pt}
	\item Conveccion del snippet: \texttt{odd[i]} = $k$ tal que el palindromo impar mas largo centrado en $i$ tiene longitud $2k - 1$. \texttt{even[i]} = $k$ para par centrado entre $i-1$ e $i$, longitud $2k$.
	\item $k$ cuenta tambien cuantos palindromos hay centrados ahi.
	\item Las queries \texttt{isPalindrome(l, r)} son $O(1)$.
	\item La ventana $[l, r]$ siempre contiene el palindromo mas largo encontrado hasta ahora.
\end{itemize}

\paragraph{Codigo clave.}
Manacher para palindromos impares:
\begin{verbatim}
int l = 0, r = -1;
for (int i = 0; i < n; ++i) {
    int k = (i > r) ? 1 : min(odd[l + r - i], r - i + 1);
    while (i - k >= 0 && i + k < n && s[i-k] == s[i+k]) ++k;
    odd[i] = k;
    if (i + k - 1 > r) { l = i - k + 1; r = i + k - 1; }
}
\end{verbatim}

\paragraph{Complejidad.}
\begin{itemize}\setlength\itemsep{0pt}
	\item Construccion $O(n)$, queries $O(1)$.
	\item Cada extension se cuenta una sola vez.
\end{itemize}

\paragraph{Donde tocar para modificar.}
\begin{itemize}\setlength\itemsep{0pt}
	\item \textbf{Solo contar palindromos}: ya implementado en \texttt{countPalindromes}.
	\item \textbf{Longest palindromic substring}: ya en \texttt{longestPalindrome}.
	\item \textbf{Modificar la conveccion de radios}: si queres ``radio = mitad de la longitud'' en vez de ``k palindromos'', ajustar el offset.
	\item \textbf{Queries con k no precomputado}: aniadir helpers especificos.
	\item \textbf{Modificar el caracter de relleno}: para multiples separadores, aniadir un caracter unico.
\end{itemize}

\paragraph{Trampas y errores frecuentes.}
\begin{itemize}\setlength\itemsep{0pt}
	\item La conveccion del snippet es \textbf{distinta} de la ``clasica de radio'' (donde radio = $(L-1)/2$); tener cuidado al copiar a otro codigo.
	\item Si los palindromos se solapan, la condicion \texttt{i + k - 1 > r} actualiza la ventana.
	\item En palindromos pares, el centro esta ``entre'' dos indices; aniadir offset en queries.
\end{itemize}

\paragraph{Codigo.}
Ver \texttt{cpp.tex}, seccion \textit{Manacher}.

% ============================================================
\section{Grafos}
% ============================================================

\subsection{Caminos minimos}

\subsubsection{Dijkstra}

\paragraph{Idea intuitiva.}
Encuentra el shortest path desde un origen $s$ a todos los vertices en $O((V + E) \log V)$ asumiendo \textbf{pesos no negativos}. La idea: usar una cola de prioridad (min-heap) ordenada por distancia tentativa. En cada paso, extraer el nodo con menor distancia; si ya fue procesado, saltar. Si al actualizar un vecino la distancia mejora, insertar/actualizar en la cola. La demostracion: cuando un nodo se extrae por primera vez, no hay forma de llegar a el con menor distancia (cualquier camino alternativo deberia pasar por nodos no procesados, con distancia $\ge$ la actual).

\paragraph{Propiedades clave (invariantes).}
\begin{itemize}\setlength\itemsep{0pt}
	\item Asume pesos $\ge 0$; con pesos negativos usar Bellman-Ford.
	\item Greedy correctness: la primera vez que se extrae un nodo, su distancia es la definitiva.
	\item Si se quiere reconstruir el camino, guardar el \texttt{parent[]}.
	\item Cada nodo se procesa a lo sumo $V$ veces (mas relajarions, pero solo $V$ definitivas).
\end{itemize}

\paragraph{Codigo clave.}
El bucle principal de Dijkstra:
\begin{verbatim}
priority_queue<pair<long long, int>, vector<...>, greater<...>> pq;
pq.push({0, source});
while (!pq.empty()) {
    auto [d, v] = pq.top(); pq.pop();
    if (d > dist[v]) continue;  // entrada vieja
    for (auto [w, u] : adj[v])
        if (dist[v] + w < dist[u]) {
            dist[u] = dist[v] + w;
            pq.push({dist[u], u});
        }
}
\end{verbatim}

\paragraph{Complejidad.}
\begin{itemize}\setlength\itemsep{0pt}
	\item $O((V + E) \log V)$ con heap binario, $O(V + E)$ con heap de Fibonacci.
	\item En la practica, $\sim 1.5 \cdot 10^{-7}$ segundos por operacion para $V = 10^5$.
\end{itemize}

\paragraph{Donde tocar para modificar.}
\begin{itemize}\setlength\itemsep{0pt}
	\item \textbf{Camino mas corto entre todos los pares}: si se necesita desde todos los origenes, usar Floyd o ejecutar Dijkstra $V$ veces.
	\item \textbf{Shortest path con aristas negativas limitadas}: Bellman-Ford.
	\item \textbf{Reconstruir camino}: aniadir \texttt{parent[v]} actualizado junto a \texttt{dist[v]}.
	\item \textbf{Dijkstra con potentials (Johnson)}: para aristas negativas si se puede reescalar.
	\item \textbf{Multi-source}: aniadir varios nodos fuentes iniciales al heap.
\end{itemize}

\paragraph{Trampas y errores frecuentes.}
\begin{itemize}\setlength\itemsep{0pt}
	\item Distancias no inicializadas a \texttt{INF}; el snippet usa \texttt{1e18}.
	\item El \texttt{if(dist[adjN] != 1e18)} para borrar de la cola es importante si usas \texttt{set} (no asi con \texttt{priority\_queue}).
	\item Con pesos grandes, \texttt{dist[v] + w} puede overflow; usar \texttt{LLONG\_MAX/2} como INF.
	\item No usar con pesos negativos: el algoritmo puede dar respuestas incorrectas.
\end{itemize}

\paragraph{Codigo.}
Ver \texttt{cpp.tex}, seccion \textit{Dijkstra}.

\vspace{0.5em}

\subsubsection{Bellman-Ford}

\paragraph{Idea intuitiva.}
Encuentra shortest path desde $s$ incluso con \textbf{pesos negativos} en $O(VE)$. La idea: relajar todas las aristas $V - 1$ veces. Cada relajacion propaga la mejor distancia a ``un paso mas''. Despues de $V - 1$ pasadas, si alguna arista aun se puede relajar, hay un \textbf{ciclo negativo} alcanzable. La demostracion: el camino mas corto tiene a lo sumo $V - 1$ aristas (sin ciclos); tras $V - 1$ relajaciones todas las distancias estan en su valor optimo.

\paragraph{Propiedades clave (invariantes).}
\begin{itemize}\setlength\itemsep{0pt}
	\item Funciona con pesos negativos (sin ciclos negativos).
	\item Tras $V - 1$ relajaciones, las distancias son definitivas (si no hay ciclo negativo).
	\item Una relajacion adicional detecta ciclos negativos: si \texttt{dist[u] + w < dist[v]}, hay ciclo negativo alcanzable desde $s$.
	\item Cada relajacion solo mejora distancias, asi que termina (no oscila).
\end{itemize}

\paragraph{Codigo clave.}
El doble bucle de relajacion:
\begin{verbatim}
for (int i = 0; i < n - 1; ++i)
    for (auto [u, v, w] : edges)
        if (dist[u] + w < dist[v])
            dist[v] = dist[u] + w;
// Detectar ciclo negativo
for (auto [u, v, w] : edges)
    if (dist[u] + w < dist[v]) { /* ciclo negativo */ }
\end{verbatim}

\paragraph{Complejidad.}
\begin{itemize}\setlength\itemsep{0pt}
	\item $O(VE)$.
	\item SPFA (con cola) en promedio es mucho mas rapido, peor caso igual.
\end{itemize}

\paragraph{Donde tocar para modificar.}
\begin{itemize}\setlength\itemsep{0pt}
	\item \textbf{SPFA}: optimizacion que relaja solo nodos cuya distancia cambio; en el peor caso sigue siendo $O(VE)$.
	\item \textbf{Detectar ciclo negativo global}: ejecutar Bellman-Ford desde un super-origen conectado a todos.
	\item \textbf{Camino mas corto con K aristas exactas}: $K$ iteraciones de Bellman-Ford.
	\item \textbf{Johnson's}: reescalar pesos con potentials para usar Dijkstra.
\end{itemize}

\paragraph{Trampas y errores frecuentes.}
\begin{itemize}\setlength\itemsep{0pt}
	\item Overflow si los pesos son $\ge 2^{52}$; usar \texttt{long long} y \texttt{\_\_int128} si hace falta.
	\item Si el grafo es denso ($E \approx V^2$), $O(VE)$ es prohibitivo.
	\item El chequeo de ciclo negativo debe hacerse tras las $V - 1$ iteraciones, no durante.
	\item Distancias no inicializadas a \texttt{INF}; el primer nodo (source) tiene distancia 0.
\end{itemize}

\paragraph{Codigo.}
Ver \texttt{cpp.tex}, seccion \textit{Bellman-Ford}.

\vspace{0.5em}

\subsubsection{Floyd-Warshall}

\paragraph{Idea intuitiva.}
Shortest path entre \textbf{todos los pares} de vertices en $O(V^3)$. DP: $d[i][j]$ = shortest path usando solo nodos intermedios en $\{0, \dots, k\}$. Transicion: $d_{k+1}[i][j] = \min(d_k[i][j], d_k[i][k+1] + d_k[k+1][j])$. La demostracion: tras $k$ iteraciones, $d[i][j]$ es el shortest path con intermedios en $\{0, \dots, k-1\}$; aniadir el nodo $k$ da la recurrencia clasica.

\paragraph{Propiedades clave (invariantes).}
\begin{itemize}\setlength\itemsep{0pt}
	\item Detecta ciclos negativos: tras el algoritmo, $d[i][i] < 0$ indica ciclo negativo alcanzable desde $i$.
	\item Funciona con pesos negativos (sin ciclos negativos).
	\item Memoria $O(V^2)$; con $V \le 400$ es OK, con mas conviene Dijkstra $V$ veces.
	\item Es la unica opcion cuando se quieren shortest paths entre todos los pares y el grafo es denso.
\end{itemize}

\paragraph{Codigo clave.}
El triple bucle clasico:
\begin{verbatim}
for (int k = 0; k < n; ++k)
    for (int i = 0; i < n; ++i)
        for (int j = 0; j < n; ++j)
            d[i][j] = min(d[i][j], d[i][k] + d[k][j]);
// detectar ciclo negativo
for (int i = 0; i < n; ++i) if (d[i][i] < 0) tiene_ciclo_negativo = true;
\end{verbatim}

\paragraph{Complejidad.}
\begin{itemize}\setlength\itemsep{0pt}
	\item $O(V^3)$ tiempo, $O(V^2)$ memoria.
	\item Para $V \le 400$: $\sim 6.4 \cdot 10^7$ operaciones (factible).
\end{itemize}

\paragraph{Donde tocar para modificar.}
\begin{itemize}\setlength\itemsep{0pt}
	\item \textbf{Reconstruir camino}: aniadir \texttt{next[i][j]} y backtrack.
	\item \textbf{Transitividad}: el algoritmo tambien calcula transitividad (\texttt{d[i][j] != INF} significa alcanzable).
	\item \textbf{Pesos doubles}: aniadir tolerancia (\texttt{d[i][j] > d[i][k] + d[k][j] - eps}).
	\item \textbf{Grafo no denso}: usar Dijkstra $V$ veces en vez de Floyd.
\end{itemize}

\paragraph{Trampas y errores frecuentes.}
\begin{itemize}\setlength\itemsep{0pt}
	\item Para detectar ciclo negativo correctamente, aniadir un chequeo \texttt{if (ans[i][i] < 0)} tras el triple bucle.
	\item El orden de los bucles \texttt{k, i, j} es importante; invertir da resultados diferentes.
	\item Con \texttt{double}, usar tolerancia \texttt{< eps} en vez de \texttt{<} para evitar loops infinitos.
\end{itemize}

\paragraph{Codigo.}
Ver \texttt{cpp.tex}, seccion \textit{Floyd-Warshall}.

\subsection{Componentes y cortes}

\subsubsection{Kosaraju SCC}

\paragraph{Idea intuitiva.}
Encuentra las \textbf{componentes fuertemente conexas} (SCCs) en $O(V + E)$. Algoritmo en 3 pasos:
\begin{itemize}\setlength\itemsep{0pt}
	\item DFS en $G$ rellenando una pila en \textbf{post-orden} (los nodos que terminan tarde quedan arriba).
	\item Construir $G^T$ (grafo transpuesto).
	\item DFS en $G^T$ procesando nodos en orden de la pila (de mayor a menor tiempo); cada DFS da una SCC.
\end{itemize}
La demostracion: dos nodos $u, v$ estan en la misma SCC sii hay camino $u \leadsto v$ y $v \leadsto u$; Kosaraju captura exactamente esto al explorar $G^T$ desde los nodos con mayor tiempo de finalizacion en $G$.

\paragraph{Propiedades clave (invariantes).}
\begin{itemize}\setlength\itemsep{0pt}
	\item Cada SCC es maximal fuertemente conexa.
	\item El \textbf{grafo de SCCs} (metagraf) es un DAG.
	\item Kosaraju y Tarjan dan el mismo resultado; Tarjan es ``mas rapido'' en practica (un solo DFS).
	\item Cada nodo pertenece a exactamente una SCC.
\end{itemize}

\paragraph{Codigo clave.}
La estructura de Kosaraju:
\begin{verbatim}
vector<int> order;
vector<bool> used(n);
function<void(int)> dfs1 = [&](int v) {
    used[v] = true;
    for (int u : g[v]) if (!used[u]) dfs1(u);
    order.push_back(v);  // post-orden
};
function<void(int, int)> dfs2 = [&](int v, int root) {
    comp[v] = root;
    for (int u : gt[v]) if (comp[u] == -1) dfs2(u, root);
};
// 1: DFS en G para llenar order
// 2: DFS en G^T en orden inverso para asignar SCCs
\end{verbatim}

\paragraph{Complejidad.}
\begin{itemize}\setlength\itemsep{0pt}
	\item $O(V + E)$.
	\item Constante 2 (dos DFS + el grafo transpuesto).
\end{itemize}

\paragraph{Donde tocar para modificar.}
\begin{itemize}\setlength\itemsep{0pt}
	\item \textbf{2-SAT}: SCC sobre grafo de implicaciones.
	\item \textbf{Construir el DAG de SCCs}: aniadir aristas entre SCCs diferentes.
	\item \textbf{Tamano de cada SCC}: aniadir contador durante el segundo DFS.
	\item \textbf{Grafo grande}: usar Tarjan para evitar construir $G^T$.
\end{itemize}

\paragraph{Trampas y errores frecuentes.}
\begin{itemize}\setlength\itemsep{0pt}
	\item El snippet usa \texttt{vector<...> adj[n]} (VLA); reemplazable por \texttt{vector<vector<pair<int,int>>} para portabilidad.
	\item El orden del segundo DFS debe ser \textbf{inverso} al de finalizacion del primero.
	\item Si el grafo no es conexo, aniadir un loop sobre todos los nodos en el primer DFS.
	\item Confundir SCC con componentes conexas (las primeras son para dirigidos, las segundas para no dirigidos).
\end{itemize}

\paragraph{Codigo.}
Ver \texttt{cpp.tex}, seccion \textit{Kosaraju SCC}.

\vspace{0.5em}

\subsubsection{Tarjan (puentes)}

\paragraph{Idea intuitiva.}
Encuentra los \textbf{puentes} (bridges): aristas cuya remocion \textbf{desconecta} el grafo. Usa un solo DFS manteniendo \texttt{tin[v]} (tiempo de descubrimiento) y \texttt{low[v]} (el menor \texttt{tin} alcanzable desde $v$ via back-edges). Una arista $v - u$ es puente si $low[u] > tin[v]$. La demostracion: si $low[u] > tin[v]$, entonces $u$ (y su subarbol) no pueden ``escapar'' de $v$; quitar la arista $v - u$ parte el grafo.

\paragraph{Propiedades clave (invariantes).}
\begin{itemize}\setlength\itemsep{0pt}
	\item \texttt{low[v] = min(tin[v], tin[w])} para todo back-edge $(v, w)$, y \texttt{min(low[u])} para hijos en el DFS.
	\item Solo se consideran las aristas del \textbf{arbol DFS} (no back-edges) para \texttt{low}.
	\item Cada nodo tiene exactamente un \texttt{tin}; cada back-edge se procesa una vez.
\end{itemize}

\paragraph{Codigo clave.}
El DFS de Tarjan para bridges:
\begin{verbatim}
void dfs(int v, int pEdge) {
    used[v] = true;
    tin[v] = low[v] = timer++;
    for (auto [u, id] : adj[v]) {
        if (id == pEdge) continue;  // no contar el edge al padre
        if (used[u]) {
            low[v] = min(low[v], tin[u]);  // back-edge
        } else {
            dfs(u, id);
            low[v] = min(low[v], low[u]);
            if (low[u] > tin[v]) bridges.push_back(id);  // es puente
        }
    }
}
\end{verbatim}

\paragraph{Complejidad.}
\begin{itemize}\setlength\itemsep{0pt}
	\item $O(V + E)$.
	\item Una sola pasada de DFS.
\end{itemize}

\paragraph{Donde tocar para modificar.}
\begin{itemize}\setlength\itemsep{0pt}
	\item \textbf{Bridge tree}: contractar cada 2-edge-connected component a un nodo; el resultado es un arbol.
	\item \textbf{Articulation points}: variante con condicion diferente (ver el siguiente modulo).
	\item \textbf{Grafo no dirigido}: solo funciona en no dirigidos; en dirigidos hay conceptos analogos (strong bridges, dominator tree).
	\item \textbf{Reconstruir componentes 2-edge-connected}: BFS desde cada nodo excluyendo bridges.
\end{itemize}

\paragraph{Trampas y errores frecuentes.}
\begin{itemize}\setlength\itemsep{0pt}
	\item No confundir el parent con un back-edge al padre; chequear \texttt{if (adjN == parent) continue;} y contar solo las veces que se ``visita'' como hijo.
	\item Para multigrafos, los edges deben tener IDs unicos para distinguirlos.
	\item Si el grafo no es conexo, aniadir loop externo sobre todos los nodos.
\end{itemize}

\paragraph{Codigo.}
Ver \texttt{cpp.tex}, seccion \textit{Tarjan (puentes)}.

\vspace{0.5em}

\subsubsection{Tarjan (puntos de articulacion)}

\paragraph{Idea intuitiva.}
Encuentra los \textbf{puntos de articulacion} (cut vertices): vertices cuya remocion \textbf{desconecta} el grafo. Mismo setup que para bridges: \texttt{tin[v]}, \texttt{low[v]}. Un nodo $v$ (no raiz) es de articulacion si existe un hijo $u$ con $low[u] \ge tin[v]$. La \textbf{raiz} es de articulacion si tiene $\ge 2$ hijos en el DFS tree. La razon: para $v$ no raiz, la condicion $\ge$ indica que el subarbol de $u$ depende de $v$; quitar $v$ los desconecta. Para la raiz, perder $\ge 2$ hijos parte el DFS en $\ge 2$ componentes.

\paragraph{Propiedades clave (invariantes).}
\begin{itemize}\setlength\itemsep{0pt}
	\item $low[u] \ge tin[v]$ significa que el subarbol de $u$ no tiene back-edge ``por encima'' de $v$.
	\item La raiz es caso especial: tiene que perder mas de un hijo para que el grafo se parta.
	\item Cada nodo se evalua una sola vez.
\end{itemize}

\paragraph{Codigo clave.}
La condicion de articulacion:
\begin{verbatim}
void dfs(int v, int p) {
    used[v] = true;
    tin[v] = low[v] = timer++;
    int children = 0;
    for (int u : adj[v]) {
        if (u == p) continue;
        if (used[u]) low[v] = min(low[v], tin[u]);
        else {
            dfs(u, v);
            low[v] = min(low[v], low[u]);
            if (low[u] >= tin[v] && p != -1)  // no raiz
                isArt[v] = true;
            ++children;
        }
    }
    if (p == -1 && children > 1) isArt[v] = true;  // raiz especial
}
\end{verbatim}

\paragraph{Complejidad.}
\begin{itemize}\setlength\itemsep{0pt}
	\item $O(V + E)$.
	\item Una sola pasada de DFS.
\end{itemize}

\paragraph{Donde tocar para modificar.}
\begin{itemize}\setlength\itemsep{0pt}
	\item \textbf{Block-cut tree}: estructura bipartita con BCCs (2-vertex-connected components) y articulation points.
	\item \textbf{Verificar biconectividad}: chequear que no hay puntos de articulacion.
	\item \textbf{Componentes biconnected}: extender Tarjan para enumerar BCCs.
\end{itemize}

\paragraph{Trampas y errores frecuentes.}
\begin{itemize}\setlength\itemsep{0pt}
	\item No olvidar el caso especial de la raiz.
	\item La condicion usa \texttt{>=}, no \texttt{>}; un hijo con \texttt{low[u] == tin[v]} tambien cuenta.
	\item Multigrafos pueden tener ``falsos'' articulation points; verificar caso a caso.
\end{itemize}

\paragraph{Codigo.}
Ver \texttt{cpp.tex}, seccion \textit{Tarjan (puntos de articulacion)}.

\vspace{0.5em}

\subsubsection{Tarjan SCC}

\paragraph{Idea intuitiva.}
Variante del SCC que usa \textbf{un solo DFS} con una pila explicita. Mantiene \texttt{disc[v]}, \texttt{low[v]}, y una pila de nodos ``en la SCC actual''. Cuando \texttt{low[v] == disc[v]}, $v$ es la raiz de una SCC; popear hasta $v$ inclusive. La demostracion: $low[v] == disc[v]$ indica que $v$ no tiene back-edges a ancestros; entonces $v$ y los nodos en la pila hasta $v$ forman una SCC maximal.

\paragraph{Propiedades clave (invariantes).}
\begin{itemize}\setlength\itemsep{0pt}
	\item Mas eficiente que Kosaraju (un solo DFS).
	\item \texttt{low[v] = min(low[u], disc[w])} para back-edges o recursion.
	\item La pila asegura que solo los nodos en la SCC actual se popean.
	\item Cada nodo se pushea y popea a lo sumo una vez.
\end{itemize}

\paragraph{Codigo clave.}
El DFS de Tarjan SCC:
\begin{verbatim}
void dfs(int v) {
    disc[v] = low[v] = timer++;
    st.push(v); inStack[v] = true;
    for (int u : g[v]) {
        if (disc[u] == -1) { dfs(u); low[v] = min(low[v], low[u]); }
        else if (inStack[u]) low[v] = min(low[v], disc[u]);
    }
    if (low[v] == disc[v]) {
        // popear hasta v
        while (true) {
            int u = st.top(); st.pop(); inStack[u] = false;
            comp[u] = current_scc;
            if (u == v) break;
        }
        ++current_scc;
    }
}
\end{verbatim}

\paragraph{Complejidad.}
\begin{itemize}\setlength\itemsep{0pt}
	\item $O(V + E)$.
	\item Constante similar a Kosaraju pero sin construir $G^T$.
\end{itemize}

\paragraph{Donde tocar para modificar.}
\begin{itemize}\setlength\itemsep{0pt}
	\item \textbf{Tamano de cada SCC}: aniadir contador durante el pop.
	\item \textbf{Grafo condensation}: aniadir aristas entre SCCs en el DAG resultante.
	\item \textbf{2-SAT}: aniadir la construccion del grafo de implicaciones.
\end{itemize}

\paragraph{Trampas y errores frecuentes.}
\begin{itemize}\setlength\itemsep{0pt}
	\item \texttt{inStack[u]} debe chequearse antes de usar \texttt{low[v] = min(low[v], disc[u])} en back-edges.
	\item La pila global es importante; sin ella no se puede extraer la SCC.
	\item Si el grafo es muy grande, recursion puede overflow; aniadir version iterativa.
\end{itemize}

\paragraph{Codigo.}
Ver \texttt{cpp.tex}, seccion \textit{Tarjan SCC}.

\subsection{Matching}

\subsubsection{Kuhn (bipartite matching)}

\paragraph{Idea intuitiva.}
Encuentra un \textbf{matching maximo} en un grafo bipartito en $O(VE)$. Para cada nodo $u$ en $L$, intentar encontrar un vecino $v$ en $R$ que este libre o cuyo match se pueda ``reorganizar'' recursivamente. Si se encuentra, matchear $u - v$. La idea: si $u$ quiere un vecino $v$ ya ocupado, intentamos ``mover'' el match de $v$ a otro nodo; recursivamente, esto encuentra un augmenting path si existe.

\paragraph{Propiedades clave (invariantes).}
\begin{itemize}\setlength\itemsep{0pt}
	\item \texttt{matchR[v]} = el nodo de $L$ matcheado a $v$ (o $-1$ si libre).
	\item \texttt{vis[]} se resetea \textbf{cada intento} de $u$.
	\item Si el matching es perfecto, $|L| = |R| = $ tamano del matching.
	\item Kuhn no garantiza maximalidad greedy; puede haber augmenting paths no encontrados.
\end{itemize}

\paragraph{Codigo clave.}
La recursion del augmenting path:
\begin{verbatim}
bool tryKuhn(int v) {
    if (vis[v]) return false;
    vis[v] = true;
    for (int u : g[v]) {
        if (matchR[u] == -1 || tryKuhn(matchR[u])) {
            matchR[u] = v;
            return true;
        }
    }
    return false;
}
// Por cada v en L:
fill(vis, vis+n, false);
if (tryKuhn(v)) ++matches;
\end{verbatim}

\paragraph{Complejidad.}
\begin{itemize}\setlength\itemsep{0pt}
	\item $O(VE)$ en el peor caso.
	\item Con Hopcroft-Karp: $O(\sqrt{V} \cdot E)$, mucho mejor para grafos densos.
\end{itemize}

\paragraph{Donde tocar para modificar.}
\begin{itemize}\setlength\itemsep{0pt}
	\item \textbf{Matching bipartito minimo (min vertex cover)}: Konig's theorem.
	\item \textbf{Matching general}: blossom algorithm (no incluido).
	\item \textbf{Aniadir capacidades}: si los nodos de $R$ tienen capacidad $> 1$, partir cada uno en capacidad copias.
	\item \textbf{Output del matching}: el array \texttt{matchR} da los pares finales.
\end{itemize}

\paragraph{Trampas y errores frecuentes.}
\begin{itemize}\setlength\itemsep{0pt}
	\item Olvidar resetear \texttt{vis} por intento; eso es lo que evita ``ciclos infinitos'' en la recursion.
	\item Para grafos grandes ($V > 10^4$), Kuhn puede ser muy lento; usar Hopcroft-Karp.
	\item Si el matching es maximo pero no perfecto, $|L|$ o $|R|$ puede ser mayor; el matching es solo maximo.
\end{itemize}

\paragraph{Codigo.}
Ver \texttt{cpp.tex}, seccion \textit{Kuhn (bipartite matching)}.

\vspace{0.5em}

\subsubsection{Hopcroft-Karp}

\paragraph{Idea intuitiva.}
Matching bipartito maximo en $O(\sqrt{V} E)$. Algoritmo en \textbf{fases}: en cada fase, hacer BFS para encontrar los \textbf{caminos de aumento mas cortos} desde nodos libres de $L$; luego DFS solo por esos caminos. Cada fase aumenta el matching en $\ge 1$ y toma $O(E)$. La demostracion: tras una fase, todos los augmenting paths tienen longitud $\ge$ la nueva minima; eso sucede cada $\sqrt{V}$ fases (longitud minima crece geometricamente).

\paragraph{Propiedades clave (invariantes).}
\begin{itemize}\setlength\itemsep{0pt}
	\item \texttt{dist[u]} = distancia en niveles desde un nodo libre de $L$.
	\item BFS termina cuando no hay mas camino a un nodo libre de $R$.
	\item DFS respeta los niveles (solo sigue \texttt{dist[matchV[v]] == dist[u] + 1}).
	\item Cada fase aumenta el matching en $\ge 1$ arista.
\end{itemize}

\paragraph{Codigo clave.}
El BFS de Hopcroft-Karp:
\begin{verbatim}
bool bfs() {
    queue<int> q;
    for (int v : L) {
        if (matchU[v] == -1) { dist[v] = 0; q.push(v); }
        else dist[v] = -1;
    }
    bool found = false;
    while (!q.empty()) {
        int v = q.front(); q.pop();
        for (int u : g[v]) {
            if (matchV[u] != -1 && dist[matchV[u]] == -1) {
                dist[matchV[u]] = dist[v] + 1;
                q.push(matchV[u]);
            }
            if (matchV[u] == -1) found = true;  // llegamos a libre
        }
    }
    return found;
}
\end{verbatim}

\paragraph{Complejidad.}
\begin{itemize}\setlength\itemsep{0pt}
	\item $O(\sqrt{V} E)$.
	\item Para $V = 10^4$, $E = 10^5$: $\sim 3 \cdot 10^7$ operaciones.
\end{itemize}

\paragraph{Donde tocar para modificar.}
\begin{itemize}\setlength\itemsep{0pt}
	\item \textbf{Reconstruir el matching}: ya viene en \texttt{matchU}/\texttt{matchV}.
	\item \textbf{Aniadir pesos}: Hungarian algorithm.
	\item \textbf{Densidad alta}: si $E \approx V^2$, sigue siendo util pero verificar constantes.
	\item \textbf{Aniadir capacidades}: partir los nodos con capacidad $> 1$ en varias copias.
\end{itemize}

\paragraph{Trampas y errores frecuentes.}
\begin{itemize}\setlength\itemsep{0pt}
	\item \texttt{dist[matchV[v]] == dist[u] + 1} (no \texttt{==} a secas); sin esto, el DFS no respeta los niveles y se pierde la complejidad.
	\item El BFS debe marcar los niveles correctamente; resetear \texttt{dist} para nodos alcanzables solo via matching.
	\item El DFS debe respetar los niveles; cualquier ``atajo'' invalida la cota.
\end{itemize}

\paragraph{Codigo.}
Ver \texttt{cpp.tex}, seccion \textit{Hopcroft-Karp}.

\subsection{Basicos}

\subsubsection{DFS / BFS / componentes}

\paragraph{Idea intuitiva.}
\textbf{DFS} (Depth-First Search) y \textbf{BFS} (Breadth-First Search) son los dos recorridos basicos de grafos. DFS usa una pila (o recursion) y explora tan profundo como puede antes de retroceder; BFS usa una cola y explora por ``niveles'' de distancia. La cantidad de \textbf{componentes conexas} se cuenta iniciando un BFS/DFS desde cada nodo no visitado. La diferencia clave: BFS da el camino mas corto en aristas; DFS da un orden de finalizacion util para problemas topologicos.

\paragraph{Propiedades clave (invariantes).}
\begin{itemize}\setlength\itemsep{0pt}
	\item BFS da la \textbf{distancia en aristas} mas corta desde el origen.
	\item DFS da un \textbf{orden topologico} (si se aniaden al final).
	\item Ambos son $O(V + E)$.
	\item Un grafo con $C$ componentes conexas satisface $\sum \text{subtree\_size}_v = V$ y $\sum E_v = 2E$.
	\item DFS en arbol da un spanning tree; las ``back edges'' corresponden a ciclos.
\end{itemize}

\paragraph{Codigo clave.}
BFS clasico:
\begin{verbatim}
queue<int> q;
q.push(source);
dist[source] = 0;
while (!q.empty()) {
    int v = q.front(); q.pop();
    for (int u : adj[v])
        if (dist[u] == -1) {
            dist[u] = dist[v] + 1;
            q.push(u);
        }
}
\end{verbatim}

\paragraph{Complejidad.}
\begin{itemize}\setlength\itemsep{0pt}
	\item $O(V + E)$ tiempo, $O(V)$ memoria.
	\item Cada arista se procesa exactamente 2 veces (ida y vuelta) en no dirigidos.
\end{itemize}

\paragraph{Donde tocar para modificar.}
\begin{itemize}\setlength\itemsep{0pt}
	\item \textbf{Aniadir pesos}: si las aristas tienen peso y queres shortest path, BFS ya no sirve; usar Dijkstra.
	\item \textbf{Grafo dirigido}: en DFS, aniadir check \texttt{if (state[u] == 1)} para detectar ciclos.
	\item \textbf{Bipartite check}: BFS alternando colores (ver el siguiente modulo).
	\item \textbf{Componentes fuertemente conexas (dirigidos)}: usar Tarjan o Kosaraju.
	\item \textbf{BFS en matriz}: aniadir offsets por direccion para grid 2D.
\end{itemize}

\paragraph{Trampas y errores frecuentes.}
\begin{itemize}\setlength\itemsep{0pt}
	\item Stack overflow con DFS recursivo para $V > 10^5$; usar version iterativa con pila explicita.
	\item En BFS, olvidar marcar \texttt{dist[u] = dist[v] + 1} antes de pushear da visitas duplicadas.
	\item Para grafos dirigidos, ``componente conexa'' no esta definida; usar SCC.
\end{itemize}

\paragraph{Codigo.}
Ver \texttt{cpp.tex}, seccion \textit{DFS / BFS / componentes}.

\vspace{0.5em}

\subsubsection{Topological sort (Kahn + DFS)}

\paragraph{Idea intuitiva.}
Un \textbf{orden topologico} de un DAG es una permutacion de los vertices tal que toda arista $u \to v$ tiene $u$ antes que $v$. Dos algoritmos clasicos:
\begin{itemize}\setlength\itemsep{0pt}
	\item \textbf{Kahn}: BFS sobre nodos con \texttt{indeg == 0}; al ``consumir'' un nodo, decrementar el indeg de sus vecinos; los que lleguen a 0 se encolan.
	\item \textbf{DFS}: hacer DFS; al terminar un nodo (estado = 2), aniadirlo al final; el orden es el reverso.
\end{itemize}
La demostracion: en un DAG siempre hay al menos un nodo con indeg 0; quitarlo mantiene el DAG; por induccion, se obtiene el orden.

\paragraph{Propiedades clave (invariantes).}
\begin{itemize}\setlength\itemsep{0pt}
	\item Existe sii el grafo es un DAG.
	\item Si al final \texttt{order.size() < V}, hay un ciclo.
	\item Kahn da el orden ``natural'' (los que pueden ir primero); DFS da el orden ``al reves'' (los que pueden ir al final primero).
	\item La cantidad de ordenes topologicos puede ser exponencial.
\end{itemize}

\paragraph{Codigo clave.}
Kahn con priority queue (orden lexicografico):
\begin{verbatim}
priority_queue<int, vector<int>, greater<int>> pq;
for (int v = 0; v < n; ++v) if (indeg[v] == 0) pq.push(v);
while (!pq.empty()) {
    int v = pq.top(); pq.pop();
    order.push_back(v);
    for (int u : adj[v]) {
        if (--indeg[u] == 0) pq.push(u);
    }
}
\end{verbatim}

\paragraph{Complejidad.}
\begin{itemize}\setlength\itemsep{0pt}
	\item $O(V + E)$ ambos.
	\item Con priority queue, $O((V + E) \log V)$.
\end{itemize}

\paragraph{Donde tocar para modificar.}
\begin{itemize}\setlength\itemsep{0pt}
	\item \textbf{Detectar ciclos}: comparar \texttt{order.size()} con \texttt{V}.
	\item \textbf{Orden lexicografico minimo}: usar priority queue en Kahn.
	\item \textbf{Topological sort con pesos}: combinar con shortest path en DAG (un solo pase en orden topologico).
	\item \textbf{Encontrar todos los ordenes}: DP sobre subsets (solo para $V$ pequeno).
\end{itemize}

\paragraph{Trampas y errores frecuentes.}
\begin{itemize}\setlength\itemsep{0pt}
	\item En el snippet, Kahn retorna orden vacio si hay ciclo (util para detectar).
	\item En DFS, no olvidar \texttt{reverse(order)}.
	\item Confundir indeg con outdeg; el orden Kahn quita primero los nodos ``fuente''.
	\item Para grafos no DAG, el algoritmo termina con \texttt{order.size() < V}.
\end{itemize}

\paragraph{Codigo.}
Ver \texttt{cpp.tex}, seccion \textit{Topological sort (Kahn + DFS)}.

\vspace{0.5em}

\subsubsection{Bipartite check + 2-coloring}

\paragraph{Idea intuitiva.}
Un grafo es \textbf{bipartito} si sus vertices se pueden dividir en dos conjuntos $L$ y $R$ tales que toda arista va de $L$ a $R$. Esto equivale a decir que se puede \textbf{2-colorear} sin que dos vertices adyacentes tengan el mismo color. BFS alternando colores detecta esto: si en algun momento dos adyacentes tienen el mismo color, no es bipartito. La demostracion: por cada componente, fijar el color de un nodo determina todos los demas; si en algun paso hay conflicto, hay un ciclo impar (no bipartito).

\paragraph{Propiedades clave (invariantes).}
\begin{itemize}\setlength\itemsep{0pt}
	\item Grafo bipartito $\Leftrightarrow$ sin ciclos impares.
	\item Conexo: si un componente no es bipartito, todo el grafo no lo es.
	\item Si hay \textbf{multiple componentes}, cada uno se colorea independientemente.
	\item La 2-coloracion es unica salvo swap de colores.
\end{itemize}

\paragraph{Codigo clave.}
BFS con 2-coloreo:
\begin{verbatim}
vector<int> color(n, -1);
for (int v = 0; v < n; ++v) if (color[v] == -1) {
    color[v] = 0;
    queue<int> q; q.push(v);
    while (!q.empty()) {
        int u = q.front(); q.pop();
        for (int w : adj[u]) {
            if (color[w] == -1) { color[w] = 1 - color[u]; q.push(w); }
            else if (color[w] == color[u]) return false;
        }
    }
}
return true;
\end{verbatim}

\paragraph{Complejidad.}
\begin{itemize}\setlength\itemsep{0pt}
	\item $O(V + E)$.
	\item Memoria $O(V)$.
\end{itemize}

\paragraph{Donde tocar para modificar.}
\begin{itemize}\setlength\itemsep{0pt}
	\item \textbf{Colorear desde multiples fuentes}: BFS desde cada nodo no coloreado.
	\item \textbf{Bipartito + matching}: Hopcroft-Karp.
	\item \textbf{Grafo bipartito ponderado}: Hungarian algorithm.
	\item \textbf{Verificar bipartito en subgrafo}: solo aplicar BFS a los nodos relevantes.
\end{itemize}

\paragraph{Trampas y errores frecuentes.}
\begin{itemize}\setlength\itemsep{0pt}
	\item Verificar \textbf{todos} los componentes, no solo el primero.
	\item Para grafos con auto-loops, un auto-loop hace al grafo no bipartito.
	\item En multigrafos, las aristas multiples no afectan el chequeo de bipartito.
\end{itemize}

\paragraph{Codigo.}
Ver \texttt{cpp.tex}, seccion \textit{Bipartite check + 2-coloring}.

\vspace{0.5em}

\subsubsection{0-1 BFS}

\paragraph{Idea intuitiva.}
Variante de BFS para grafos con \textbf{pesos 0 o 1}. Usa un \texttt{deque}: aristas de peso 0 se encolan al frente (\texttt{push\_front}), aristas de peso 1 al final (\texttt{push\_back}). Asi, la primera vez que se extrae un nodo, su distancia es la definitiva, igual que en Dijkstra pero en $O(V + E)$. La intuicion: el deque mantiene los nodos ordenados por distancia; las aristas de peso 0 no cambian la distancia (entran al frente), las de peso 1 la incrementan (entran al final).

\paragraph{Propiedades clave (invariantes).}
\begin{itemize}\setlength\itemsep{0pt}
	\item Asume pesos \textbf{solo} 0 o 1.
	\item Mejor caso que Dijkstra (sin factor $\log V$).
	\item No funciona con pesos $\ge 2$.
	\item La cola mantiene los nodos ordenados por distancia no decreciente.
\end{itemize}

\paragraph{Codigo clave.}
El deque doble:
\begin{verbatim}
deque<int> q;
q.push_front(source);
dist[source] = 0;
while (!q.empty()) {
    int v = q.front(); q.pop_front();
    for (auto [w, u] : adj[v]) {
        if (dist[v] + w < dist[u]) {
            dist[u] = dist[v] + w;
            if (w == 0) q.push_front(u);
            else q.push_back(u);
        }
    }
}
\end{verbatim}

\paragraph{Complejidad.}
\begin{itemize}\setlength\itemsep{0pt}
	\item $O(V + E)$.
	\item Cada nodo se inserta a lo sumo $O(1)$ veces (con check de distancia).
\end{itemize}

\paragraph{Donde tocar para modificar.}
\begin{itemize}\setlength\itemsep{0pt}
	\item \textbf{Pesos en $[0, k]$ small}: usar small k-ary BFS o Dijkstra.
	\item \textbf{Reconstruir camino}: aniadir \texttt{parent[]}.
	\item \textbf{Pesos negativos}: no funciona, usar Bellman-Ford.
\end{itemize}

\paragraph{Trampas y errores frecuentes.}
\begin{itemize}\setlength\itemsep{0pt}
	\item Si los pesos son otro rango, el algoritmo da respuestas incorrectas. Verificar que las aristas son solo 0/1.
	\item El chequeo \texttt{dist[v] + w < dist[u]} es crucial; sino, el algoritmo procesa nodos multiples veces.
	\item Si una arista tiene peso $\ge 2$, aniadir verificaciones o cambiar a Dijkstra.
\end{itemize}

\paragraph{Codigo.}
Ver \texttt{cpp.tex}, seccion \textit{0-1 BFS}.

\subsection{MST}

\subsubsection{Kruskal}

\paragraph{Idea intuitiva.}
Encuentra el \textbf{arbol de expansion minima} (MST) en $O(E \log E)$. Ordena las aristas por peso ascendente y las aniade al MST si no forman ciclo (chequear con DSU). El arbol tiene exactamente $V - 1$ aristas. La demostracion (``Cut Property''): para cualquier corte, la arista minima cruzando el corte pertenece a \emph{algun} MST; por lo tanto, elegir la minima disponible que no forma ciclo es optimo.

\paragraph{Propiedades clave (invariantes).}
\begin{itemize}\setlength\itemsep{0pt}
	\item Algoritmo \textbf{greedy} correcto (Cut Property).
	\item Funciona en grafos \textbf{no conexos}; produce un \textbf{minimum spanning forest}.
	\item Si las aristas tienen pesos unicos, el MST es unico.
	\item El MST tiene exactamente $V - 1$ aristas (o menos si el grafo es disconexo).
\end{itemize}

\paragraph{Codigo clave.}
Kruskal con DSU:
\begin{verbatim}
sort(edges.begin(), edges.end());
DSU dsu(n);
long long mst = 0;
int edgesUsed = 0;
for (auto [w, u, v] : edges)
    if (dsu.unite(u, v)) {
        mst += w;
        if (++edgesUsed == n - 1) break;
    }
\end{verbatim}

\paragraph{Complejidad.}
\begin{itemize}\setlength\itemsep{0pt}
	\item $O(E \log E)$ (ordenamiento) o $O(E \log V)$ con DSU optimizations.
	\item DSU es casi $O(1)$ amortizado, asi que el sort domina.
\end{itemize}

\paragraph{Donde tocar para modificar.}
\begin{itemize}\setlength\itemsep{0pt}
	\item \textbf{Second best MST}: enumerar aristas del MST, sacar una, aniadir una no-MST; recalcular.
	\item \textbf{Kruskal randomizado}: si el grafo es casi-completo, agregar shuffle.
	\item \textbf{Aniadir info al DSU}: guardar el peso maximo en el camino, etc.
	\item \textbf{DSU con rollback}: si necesitas Kruskal reversible (para queries offline).
\end{itemize}

\paragraph{Trampas y errores frecuentes.}
\begin{itemize}\setlength\itemsep{0pt}
	\item El DSU debe estar limpio para cada invocacion.
	\item Si las aristas tienen pesos negativos, Kruskal sigue funcionando.
	\item Si el grafo es disconexo, el ``MST'' es un bosque con varias componentes.
	\item El check \texttt{if (++edgesUsed == n - 1) break;} optimiza la salida temprana.
\end{itemize}

\paragraph{Codigo.}
Ver \texttt{cpp.tex}, seccion \textit{Kruskal}.

\vspace{0.5em}

\subsubsection{Prim (priority queue)}

\paragraph{Idea intuitiva.}
Otra forma de calcular el MST: empezar desde un nodo y, en cada paso, aniadir la arista de menor peso que conecta el conjunto construido con el resto. Usando priority queue, es $O((V + E) \log V)$. La demostracion: cada vez que se aniade una arista, es la minima cruzando el ``corte'' entre el conjunto actual y el resto (Cut Property); eso preserva optimalidad.

\paragraph{Propiedades clave (invariantes).}
\begin{itemize}\setlength\itemsep{0pt}
	\item Similar a Dijkstra pero sin acumular distancias, solo eligiendo minimo por arista.
	\item Funciona mejor que Kruskal en \textbf{grafos densos} ($E \approx V^2$).
	\item Si el grafo es disconexo, genera el bosque (solo llega a los alcanzables).
	\item Cada nodo se inserta al heap una vez; cada arista se procesa una vez.
\end{itemize}

\paragraph{Codigo clave.}
Prim clasico con heap:
\begin{verbatim}
priority_queue<pair<int, int>, vector<...>, greater<...>> pq;
vector<bool> used(n, false);
vector<int> minEdge(n, INF);
minEdge[0] = 0;
pq.push({0, 0});
while (!pq.empty()) {
    auto [w, v] = pq.top(); pq.pop();
    if (used[v]) continue;
    used[v] = true;
    mst += w;
    for (auto [peso, u] : adj[v])
        if (!used[u] && peso < minEdge[u]) {
            minEdge[u] = peso;
            pq.push({peso, u});
        }
}
\end{verbatim}

\paragraph{Complejidad.}
\begin{itemize}\setlength\itemsep{0pt}
	\item $O((V + E) \log V)$.
	\item Con matriz de adyacencia: $O(V^2)$ sin heap.
\end{itemize}

\paragraph{Donde tocar para modificar.}
\begin{itemize}\setlength\itemsep{0pt}
	\item \textbf{Prim con matriz de adyacencia}: $O(V^2)$ sin heap, util en grafos densos pequenos.
	\item \textbf{Reconstruir el MST}: aniadir \texttt{parent[]}.
	\item \textbf{Prim en streaming}: variante para grafos que llegan incrementalmente.
\end{itemize}

\paragraph{Trampas y errores frecuentes.}
\begin{itemize}\setlength\itemsep{0pt}
	\item El \texttt{if (used[v]) continue;} es crucial: si el nodo ya fue procesado, saltar.
	\item En la version matriz, elije el minimo por linea (no heap).
	\item Si el grafo es disconexo, aniadir un loop externo sobre todos los nodos no usados.
\end{itemize}

\paragraph{Codigo.}
Ver \texttt{cpp.tex}, seccion \textit{Prim (priority queue)}.

\subsection{Flujos}

\subsubsection{Edmonds-Karp (max flow)}

\paragraph{Idea intuitiva.}
Implementacion BFS del algoritmo de \textbf{Ford-Fulkerson}. En cada iteracion, hacer BFS desde $s$ hasta $t$ siguiendo aristas con capacidad residual $> 0$; si se llega a $t$, se encontro un \textbf{camino de aumento}. Enviar el minimo de las capacidades en el camino; actualizar capacidades residuales. La demostracion de $O(VE^2)$: cada BFS encuentra el augmenting path mas corto, y la longitud minima crece monotona; como hay $\le E$ longitudes posibles, hay $\le E$ fases.

\paragraph{Propiedades clave (invariantes).}
\begin{itemize}\setlength\itemsep{0pt}
	\item BFS garantiza caminos de aumento mas cortos.
	\item Capacidad residual: forward edge baja, backward edge sube.
	\item Complejidad $O(VE^2)$.
	\item El flujo maximo respeta la conservacion en nodos (in = out salvo $s, t$).
\end{itemize}

\paragraph{Codigo clave.}
BFS + augmenting path:
\begin{verbatim}
while (bfs(s, t, parent)) {  // bfs encuentra camino o retorna false
    int pathFlow = INF;
    for (int v = t; v != s; v = parent[v])
        pathFlow = min(pathFlow, capacity[parent[v]][v]);
    for (int v = t; v != s; v = parent[v]) {
        capacity[parent[v]][v] -= pathFlow;
        capacity[v][parent[v]] += pathFlow;
    }
    maxFlow += pathFlow;
}
\end{verbatim}

\paragraph{Complejidad.}
\begin{itemize}\setlength\itemsep{0pt}
	\item $O(VE^2)$.
	\item En la practica, mas lento que Dinic para grafos grandes.
\end{itemize}

\paragraph{Donde tocar para modificar.}
\begin{itemize}\setlength\itemsep{0pt}
	\item \textbf{Grafo bipartito}: $O(\sqrt{V} E)$ (equivalente a Hopcroft-Karp).
	\item \textbf{Dinic} es preferible en la mayoria de los casos.
	\item \textbf{Capacidades doubles}: usar tolerancia para comparacion.
	\item \textbf{Min-cost flow}: aniadir costo por arista y modificar el algoritmo.
\end{itemize}

\paragraph{Trampas y errores frecuentes.}
\begin{itemize}\setlength\itemsep{0pt}
	\item Las capacidades deben ser $> 0$ para entrar a la cola (no $>= 0$); sino, ciclos infinitos.
	\item Para grafos bipartitos, Dinic o Hopcroft-Karp son mejores.
	\item Si el flujo maximo es muy grande, los \texttt{int} pueden overflow; usar \texttt{long long}.
\end{itemize}

\paragraph{Codigo.}
Ver \texttt{cpp.tex}, seccion \textit{Edmonds-Karp (max flow)}.

\vspace{0.5em}

\subsubsection{Dinic (max flow)}

\paragraph{Idea intuitiva.}
Algoritmo de flujo mas eficiente que Edmonds-Karp: combina \textbf{BFS para level graph} con \textbf{DFS bloqueante} (multiple augmenting paths en una sola pasada). Cada fase cuesta $O(E)$ y hay $O(V)$ fases en el peor caso; para bipartitos, $O(\sqrt{V} E)$. La razon: cada fase ``satura'' al menos una arista, asi que en $O(E)$ fases el flujo maximo se alcanza; con $O(V)$ niveles en el BFS, hay $O(VE)$ operaciones totales.

\paragraph{Propiedades clave (invariantes).}
\begin{itemize}\setlength\itemsep{0pt}
	\item \texttt{level[v]} = distancia desde $s$ en el level graph (solo aristas con cap $> 0$).
	\item \texttt{it[v]} = iterador actual para evitar reprocesar aristas.
	\item DFS solo sigue aristas con \texttt{level[e.to] == level[v] + 1}.
	\item Back-edges tienen \texttt{rev} apuntando a la forward edge original.
\end{itemize}

\paragraph{Codigo clave.}
El DFS bloqueante de Dinic:
\begin{verbatim}
long long dfs(int v, int t, long long f) {
    if (v == t) return f;
    for (int& i = it[v]; i < adj[v].size(); ++i) {
        Edge& e = adj[v][i];
        if (e.cap > 0 && level[e.to] == level[v] + 1) {
            long long ret = dfs(e.to, t, min(f, e.cap));
            if (ret > 0) {
                e.cap -= ret;
                adj[e.to][e.rev].cap += ret;
                return ret;
            }
        }
    }
    return 0;
}
\end{verbatim}

\paragraph{Complejidad.}
\begin{itemize}\setlength\itemsep{0pt}
	\item $O(V^2 E)$ general, $O(\sqrt{V} E)$ para bipartitos.
	\item En la practica, mucho mas rapido que Edmonds-Karp.
\end{itemize}

\paragraph{Donde tocar para modificar.}
\begin{itemize}\setlength\itemsep{0pt}
	\item \textbf{Dinic con scaling}: aniadir factor de escala para capacidades grandes.
	\item \textbf{Lower bounds}: aniadir un super-source/sink para forzar flujo minimo.
	\item \textbf{Matching bipartito}: el grafo se construye como source -> L -> R -> sink, todos con capacidad 1.
	\item \textbf{Push-relabel}: alternativa para grafos muy densos.
\end{itemize}

\paragraph{Trampas y errores frecuentes.}
\begin{itemize}\setlength\itemsep{0pt}
	\item \texttt{it[v]} debe \textbf{incrementarse} en cada iteracion; sino, el DFS se queda atascado en la misma arista.
	\item Si el BFS no encuentra $t$, el flujo es maximo; terminar.
	\item Las back-edges son cruciales para ``cancelar'' elecciones previas.
\end{itemize}

\paragraph{Codigo.}
Ver \texttt{cpp.tex}, seccion \textit{Dinic (max flow)}.

\vspace{0.5em}

\subsubsection{Min-Cost Max-Flow}

\paragraph{Idea intuitiva.}
Encuentra el flujo maximo con \textbf{costo minimo} (asumiendo costos no negativos). Usa \textbf{SPFA} (o Dijkstra con potentials) para encontrar el shortest path en el \textbf{costo} desde $s$ hasta $t$ considerando capacidades residuales. Cada augmenting path aporta su costo $\times$ flujo. La idea: los potentials reescalan las aristas para que todas tengan costo no negativo, permitiendo usar Dijkstra.

\paragraph{Propiedades clave (invariantes).}
\begin{itemize}\setlength\itemsep{0pt}
	\item \texttt{pot[]} = potentials para evitar aristas negativas (Johnson's trick).
	\item \texttt{dist[v]} = costo minimo desde $s$ con potentials.
	\item Solo se aumentan aristas con \texttt{cap > 0} (forward) y costo es \texttt{+cost} (forward) o \texttt{-cost} (backward).
	\item Tras cada augmenting, los potentials se actualizan para mantener la consistencia.
\end{itemize}

\paragraph{Codigo clave.}
El ciclo de min-cost flow:
\begin{verbatim}
while (spfa(s, t, dist, parent)) {
    int pathFlow = INF;
    for (int v = t; v != s; v = parent[v])
        pathFlow = min(pathFlow, capacity[parent[v]][v]);
    for (int v = t; v != s; v = parent[v]) {
        capacity[parent[v]][v] -= pathFlow;
        capacity[v][parent[v]] += pathFlow;
        cost += pathFlow * edge_cost[parent[v]][v];
    }
}
\end{verbatim}

\paragraph{Complejidad.}
\begin{itemize}\setlength\itemsep{0pt}
	\item $O(F \cdot V E)$ con SPFA, $O(F \cdot E \log V)$ con Dijkstra + potentials.
	\item $F$ = flujo maximo total.
\end{itemize}

\paragraph{Donde tocar para modificar.}
\begin{itemize}\setlength\itemsep{0pt}
	\item \textbf{Dijkstra + potentials}: si los costos son enteros no negativos, esto es mas estable.
	\item \textbf{Costos negativos en aristas}: requieren Bellman-Ford inicial o reorden.
	\item \textbf{Flujo minimo con costo}: encontrar el minimo flujo posible y su costo.
	\item \textbf{Costo por unidad}: si los costos son muy grandes, aniadir verificacion.
\end{itemize}

\paragraph{Trampas y errores frecuentes.}
\begin{itemize}\setlength\itemsep{0pt}
	\item El \texttt{pot[]} se actualiza \textbf{al final} de cada iteracion (no durante).
	\item Las back-edges tienen \textbf{costo negativo}, esencial para corregir elecciones previas.
	\item Si los costos son muy grandes, considerar modular arithmetic.
	\item Para problemas con flujo exacto, aniadir verificacion post-algoritmo.
\end{itemize}

\paragraph{Codigo.}
Ver \texttt{cpp.tex}, seccion \textit{Min-Cost Max-Flow}.

\subsection{Especales}

\subsubsection{2-SAT (implication graph + SCC)}

\paragraph{Idea intuitiva.}
Resuelve formulas booleanas en \textbf{CNF} con clausulas de \textbf{2 literales}: $(l_1 \lor l_2)$. Cada variable $x$ se representa con dos nodos: $x$ (verdadero) y $\neg x$ (falso). Cada clausula $l_1 \lor l_2$ se descompone en dos implicaciones: $\neg l_1 \to l_2$ y $\neg l_2 \to l_1$. Si en el grafo de implicaciones $\neg x$ y $x$ estan en la misma SCC, la formula es \textbf{UNSAT}. La demostracion: si $\neg x \to x$ en el grafo, entonces $\neg x$ implica $x$, asi que $\neg x$ debe ser falso y $x$ verdadero; pero entonces $\neg x$ es falso implica una contradiccion.

\paragraph{Propiedades clave (invariantes).}
\begin{itemize}\setlength\itemsep{0pt}
	\item Cada literal es un nodo; total $2n$ nodos.
	\item Kosaraju/Tarjan dan las SCCs.
	\item Si SAT: asignar $x = $ (topological order del condensation).
	\item Clausulas \texttt{setTrue(x)} se modelan como \texttt{implies(!x, x)}.
	\item La formula es SAT sii no hay variable $x$ con $x$ y $\neg x$ en la misma SCC.
\end{itemize}

\paragraph{Codigo clave.}
Conversion de clausulas a implicaciones:
\begin{verbatim}
// Clausula (a OR b) -> implica ~a -> b Y ~b -> a
addImplication(not(a), b);
addImplication(not(b), a);
// Forzar x: implica ~x -> x
if (force) addImplication(not(x), x);
\end{verbatim}

\paragraph{Complejidad.}
\begin{itemize}\setlength\itemsep{0pt}
	\item $O(V + E) = O(N + M)$ con Kosaraju o Tarjan.
	\item Memoria $O(V + E)$.
\end{itemize}

\paragraph{Donde tocar para modificar.}
\begin{itemize}\setlength\itemsep{0pt}
	\item \textbf{Forzar variables}: aniadir \texttt{setTrue} / \texttt{setFalse}.
	\item \textbf{3-SAT}: NP-completo; no usar este algoritmo.
	\item \textbf{Contar soluciones}: las SCCs condensation dan una formula para contar; multiplicar las formas de asignar variables en cada componente no-trivial.
	\item \textbf{Output de la asignacion}: aniadir el vector \texttt{assignment[]}.
\end{itemize}

\paragraph{Trampas y errores frecuentes.}
\begin{itemize}\setlength\itemsep{0pt}
	\item La asignacion usa el \textbf{orden topologico inverso} del condensation: \texttt{val[i] = comp[2i] < comp[2i+1]}.
	\item El ID de $\neg x$ es \texttt{2*x + 1} (o similar); verificar consistencia.
	\item Si hay clausulas con literales repetidos (e.g., $(x \lor x)$), tratar como $(x)$ simple.
\end{itemize}

\paragraph{Codigo.}
Ver \texttt{cpp.tex}, seccion \textit{2-SAT (implication graph + SCC)}.

\vspace{0.5em}

\subsubsection{Euler tour / Hierholzer (camino y circuito)}

\paragraph{Idea intuitiva.}
Encuentra un \textbf{camino o circuito euleriano} (que usa cada arista exactamente una vez). \textbf{Condiciones}:
\begin{itemize}\setlength\itemsep{0pt}
	\item \textbf{No dirigido}: todos los grados pares (circuito) o exactamente 2 impares (camino, empieza en uno, termina en el otro). Ademas, conexo sobre el subgrafo con aristas.
	\item \textbf{Dirigido}: $\text{indeg}(v) = \text{outdeg}(v)$ para todos (circuito) o exactamente un nodo con $\text{outdeg} = \text{indeg} + 1$ y otro con $\text{indeg} = \text{outdeg} + 1$ (camino).
\end{itemize}

\paragraph{Hierholzer}: elegir un ciclo desde el inicio, mientras haya ciclos ``laterales'' desde un nodo del camino, mergearlos. Implementacion con pila: avanzar por aristas, al no poder mas aniadir a la respuesta y backtrack. La demostracion de $O(V + E)$: cada arista se procesa una vez (entra y sale del stack una vez).

\paragraph{Propiedades clave (invariantes).}
\begin{itemize}\setlength\itemsep{0pt}
	\item Multigrafos permitidos; las aristas se marcan como usadas.
	\item No dirigido: cada arista aparece \textbf{dos veces} en \texttt{adj} (ida y vuelta).
	\item La construccion del path con ids consistentes es importante.
	\item El resultado es unico salvo orden de las aristas en multigrafos.
\end{itemize}

\paragraph{Codigo clave.}
Hierholzer con pila:
\begin{verbatim}
vector<int> path;
stack<int> st;
st.push(start);
while (!st.empty()) {
    int v = st.top();
    while (!adj[v].empty() && used[adj[v].back()]) adj[v].pop_back();
    if (adj[v].empty()) {
        path.push_back(v);
        st.pop();
    } else {
        int e = adj[v].back();
        used[e] = true;
        int u = edges[e].to(v);
        adj[v].pop_back();
        st.push(u);
    }
}
reverse(path.begin(), path.end());
\end{verbatim}

\paragraph{Complejidad.}
\begin{itemize}\setlength\itemsep{0pt}
	\item $O(V + E)$.
	\item Memoria $O(V + E)$.
\end{itemize}

\paragraph{Donde tocar para modificar.}
\begin{itemize}\setlength\itemsep{0pt}
	\item \textbf{Reconstruir con multigrafo}: aniadir ids consistentes antes del algoritmo.
	\item \textbf{Aplicar a ``de Bruijn'' o problemas de strings}: modelar como multigrafo y buscar Euler.
	\item \textbf{Dirigido vs no dirigido}: el snippet tiene ambas versiones.
	\item \textbf{Letter addition}: modelar palabras como aristas y buscar el superstring.
\end{itemize}

\paragraph{Trampas y errores frecuentes.}
\begin{itemize}\setlength\itemsep{0pt}
	\item Si el grafo no es conexo (sobre las aristas), no existe euleriano. Verificar primero.
	\item Para multigrafos, los IDs de aristas son cruciales para distinguirlas.
	\item Si el problema es ``encuentra todos los euler tours'', el algoritmo es exponencial.
\end{itemize}

\paragraph{Codigo.}
Ver \texttt{cpp.tex}, seccion \textit{Euler tour / Hierholzer (camino y circuito)}.

\vspace{0.5em}

\subsubsection{Hamiltonian path / cycle (bitmask DP + backtracking)}

\paragraph{Idea intuitiva.}
Un \textbf{camino hamiltoniano} visita cada vertice \textbf{exactamente una vez}. \textbf{Determinar} si existe es NP-completo en general, pero con $V \le 20$ se puede hacer DP con bitmask en $O(V^2 \cdot 2^V)$. \textbf{Teoremas suficientes}: \textbf{Dirac} (grado $\ge n/2$) y \textbf{Ore} ($\text{deg}(u) + \text{deg}(v) \ge n$ para todo par no adyacente). La idea de la DP: cada estado es (conjunto de visitados, ultimo vertice); la transicion es agregar un vecino no visitado.

\paragraph{Propiedades clave (invariantes).}
\begin{itemize}\setlength\itemsep{0pt}
	\item DP: $dp[mask][v]$ = true si hay un camino que visita exactamente \texttt{mask} y termina en $v$.
	\item TSP: variante con pesos minimos (Hamilton cycle minimo).
	\item Reconstruccion: guardar \texttt{par[mask][v]} (predecesor).
	\item $2^V$ estados con $V$ bits cada uno; factible para $V \le 20$.
\end{itemize}

\paragraph{Codigo clave.}
DP con bitmask:
\begin{verbatim}
dp[1][0] = 0;  // empezar en 0
for (int mask = 1; mask < (1<<n); ++mask)
    for (int v = 0; v < n; ++v) if (dp[mask][v] < INF)
        for (int u = 0; u < n; ++u)
            if (!(mask & (1<<u)))
                dp[mask | (1<<u)][u] = min(dp[mask | (1<<u)][u],
                                            dp[mask][v] + w[v][u]);
// TSP: respuesta = min(dp[(1<<n)-1][v] + w[v][0])
\end{verbatim}

\paragraph{Complejidad.}
\begin{itemize}\setlength\itemsep{0pt}
	\item $O(V^2 \cdot 2^V)$ tiempo, $O(V \cdot 2^V)$ memoria.
	\item Para $V = 20$: $\sim 4 \cdot 10^8$ operaciones (lento pero factible).
\end{itemize}

\paragraph{Donde tocar para modificar.}
\begin{itemize}\setlength\itemsep{0pt}
	\item \textbf{TSP asimetrico}: aniadir pesos $w[v][u]$ (no necesariamente simetricos).
	\item \textbf{Hamilton con inicio obligatorio}: fijar \texttt{src} en el DP inicial.
	\item \textbf{Bitmask DP mas compacto}: usar \texttt{long long} como mascara si $V \le 64$.
	\item \textbf{TSP optimo aproximado}: MST, 2-opt, simulated annealing.
	\item \textbf{Reconstruir el path}: usar \texttt{par[mask][v]} y backtrack.
\end{itemize}

\paragraph{Trampas y errores frecuentes.}
\begin{itemize}\setlength\itemsep{0pt}
	\item TSP asume grafo \textbf{completo} (con \texttt{LINF} para aristas faltantes); si no, aniadir chequeo.
	\item La condicion ``\texttt{mask == (1<<n) - 1}'' cierra el ciclo retornando al origen.
	\item Para $V > 20$, la memoria $2^V$ se vuelve prohibitiva; usar meet-in-the-middle o heuristicas.
\end{itemize}

\paragraph{Codigo.}
Ver \texttt{cpp.tex}, seccion \textit{Hamiltonian path / cycle (bitmask DP + backtracking)}.

% ============================================================
\section{DP}

\subsubsection{Knapsack 0/1}

\paragraph{Idea intuitiva.}
Dado un conjunto de $N$ objetos con peso $w_i$ y valor $v_i$, maximizar $\sum v_i$ con la restriccion $\sum w_i \le W$. La recurrencia clasica:
\[ dp[i][j] = \max(dp[i-1][j],\, dp[i-1][j - w_i] + v_i) \quad \text{si } j \ge w_i. \]
La optimizacion clave: la DP es 1-dimensional si iteramos $j$ \textbf{de mayor a menor} (asi no usamos el objeto $i$ dos veces). La razon de la iteracion inversa: en cada iteracion queremos el estado del ``item anterior''; iterar de menor a mayor contaminaria \texttt{dp[j - w_i]} con el item actual.

\paragraph{Propiedades clave (invariantes).}
\begin{itemize}\setlength\itemsep{0pt}
	\item $dp[j]$ al final = maximo valor con peso $\le j$.
	\item El bucle inverso \texttt{for (j = W; j >= w[i]; j--)} es lo que distingue 0/1 de unbounded.
	\item Solucion: maximo valor con cualquier peso $\le W$.
	\item La DP es exacta; no hay aproximacion.
\end{itemize}

\paragraph{Codigo clave.}
La version 1D optimizada:
\begin{verbatim}
vector<long long> dp(W + 1, 0);
for (int i = 0; i < N; ++i)
    for (int j = W; j >= w[i]; --j)  // invertido
        dp[j] = max(dp[j], dp[j - w[i]] + v[i]);
return dp[W];
\end{verbatim}

\paragraph{Complejidad.}
\begin{itemize}\setlength\itemsep{0pt}
	\item $O(N \cdot W)$ tiempo, $O(W)$ memoria.
	\item Constante muy pequena (solo sumas y max).
\end{itemize}

\paragraph{Donde tocar para modificar.}
\begin{itemize}\setlength\itemsep{0pt}
	\item \textbf{Reconstruir la seleccion}: aniadir \texttt{take[i][j]} (booleano o bool) en una version 2D.
	\item \textbf{Items con peso 0}: edge case; aniadir guard.
	\item \textbf{Knapsack con restricciones de grupos}: ``elegir a lo sumo 1 de cada grupo'' se modela como 0/1 con un peso extra.
	\item \textbf{Tamano de $W$}: si $W \le 10^9$ pero $N$ es chico, usar meet-in-the-middle.
	\item \textbf{Cambiar max por min}: aniadir \texttt{INF} en la inicializacion.
\end{itemize}

\paragraph{Trampas y errores frecuentes.}
\begin{itemize}\setlength\itemsep{0pt}
	\item Overflow en \texttt{dp[j - w[i]] + v[i]}; usar \texttt{long long} si los valores son grandes.
	\item El bucle \textbf{debe} ir de mayor a menor.
	\item Si los pesos son grandes y $N$ pequeno, conviene meet-in-the-middle.
	\item Para Knapsack exacto ($\sum w_i = W$), aniadir check al final.
\end{itemize}

\paragraph{Codigo.}
Ver \texttt{cpp.tex}, seccion \textit{Knapsack 0/1}.

\vspace{0.5em}

\subsubsection{Knapsack unbounded}

\paragraph{Idea intuitiva.}
Igual que Knapsack 0/1 pero \textbf{cada item se puede tomar multiples veces}. La recurrencia es similar:
\[ dp[j] = \max(dp[j],\, dp[j - w_i] + v_i), \]
pero ahora el bucle va de \textbf{menor a mayor} (porque tomar el item $i$ no afecta corridas futuras del mismo item). La idea: tras tomar el item $i$, queremos volver a tener la opcion de tomarlo; iterar de menor a mayor permite que el item ``se acumule''.

\paragraph{Propiedades clave (invariantes).}
\begin{itemize}\setlength\itemsep{0pt}
	\item El bucle \texttt{for (j = w[i]; j <= W; j++)} es lo que distingue unbounded.
	\item A diferencia de 0/1, no hay riesgo de ``usar dos veces el mismo item''.
	\item La DP es exacta; cada ``subconjunto'' se cuenta correctamente.
\end{itemize}

\paragraph{Codigo clave.}
La version 1D con bucle directo:
\begin{verbatim}
vector<long long> dp(W + 1, 0);
for (int i = 0; i < N; ++i)
    for (int j = w[i]; j <= W; ++j)  // directo (no invertido)
        dp[j] = max(dp[j], dp[j - w[i]] + v[i]);
return dp[W];
\end{verbatim}

\paragraph{Complejidad.}
\begin{itemize}\setlength\itemsep{0pt}
	\item $O(N \cdot W)$ tiempo, $O(W)$ memoria.
	\item Constante muy pequena.
\end{itemize}

\paragraph{Donde tocar para modificar.}
\begin{itemize}\setlength\itemsep{0pt}
	\item \textbf{Reconstruir cuantos de cada item}: aniadir \texttt{cnt[i]} y backtrack.
	\item \textbf{Cambiar moneda minima por maxima}: el snippet maximiza valor; para minimizar monedas, cambiar a \texttt{min} con inicializacion en \texttt{INF}.
	\item \textbf{Coin change}: si todos los pesos son positivos y queremos ``minimo numero de monedas'', es unbounded con min.
\end{itemize}

\paragraph{Trampas y errores frecuentes.}
\begin{itemize}\setlength\itemsep{0pt}
	\item Misma que Knapsack 0/1 con overflow.
	\item El orden de iteracion de items no afecta el resultado (convexo).
	\item Si se quiere ``maximo valor con exactamente $k$ items'', aniadir segunda dimension.
\end{itemize}

\paragraph{Codigo.}
Ver \texttt{cpp.tex}, seccion \textit{Knapsack unbounded}.

\vspace{0.5em}

\subsubsection{Knapsack meet-in-the-middle}

\paragraph{Idea intuitiva.}
Para $N \le 40$ y $W$ moderado, no se puede hacer $O(NW)$ directamente. Partir en dos mitades, enumerar todos los $2^{N/2}$ subconjuntos de cada mitad, y combinar. Para cada subconjunto de la segunda mitad con peso $sw_2$, encontrar el maximo valor de la primera mitad con peso $\le W - sw_2$. La idea: ``divide y venceras'' sobre el tamano de la entrada; cada mitad tiene mitad del espacio, asi que la combinacion es geometrica en lugar de multiplicativa.

\paragraph{Propiedades clave (invariantes).}
\begin{itemize}\setlength\itemsep{0pt}
	\item Hay $2^{N/2}$ subconjuntos por mitad; total $O(2^{N/2})$.
	\item \texttt{sums} guarda los valores de subconjuntos de la primera mitad ordenados (o preprocesados).
	\item Para cada subconjunto de la segunda mitad, busqueda lineal o binaria en \texttt{sums}.
	\item La mitad mas grande (con $N/2 + 1$ items para $N$ impar) tiene $2^{N/2+1}$ subconjuntos.
\end{itemize}

\paragraph{Codigo clave.}
Enumerar subconjuntos de la primera mitad:
\begin{verbatim}
int n1 = N / 2, n2 = N - n1;
vector<pair<long long, long long>> sums1;
for (int mask = 0; mask < (1 << n1); ++mask) {
    long long peso = 0, valor = 0;
    for (int i = 0; i < n1; ++i) if (mask & (1 << i)) {
        peso += w[i]; valor += v[i];
    }
    sums1.push_back({peso, valor});
}
sort(sums1.begin(), sums1.end());  // por peso
// eliminar duplicados dominados
\end{verbatim}

\paragraph{Complejidad.}
\begin{itemize}\setlength\itemsep{0pt}
	\item $O(2^{N/2})$ tiempo y memoria (mejor que $O(NW)$ cuando $W > 2^{N/2}$).
	\item Para $N = 40$: $\sim 2 \cdot 10^6$ subconjuntos, factible.
\end{itemize}

\paragraph{Donde tocar para modificar.}
\begin{itemize}\setlength\itemsep{0pt}
	\item \textbf{Tamano mayor}: si $N$ es mayor pero $W$ es chico, conviene iterar por peso, no por subconjuntos.
	\item \textbf{Varias mochilas}: adaptar para multiples $W$.
	\item \textbf{Optimizar el ``combinado''}: usar \texttt{upper\_bound} para encontrar el maximo valor con peso $\le \text{rem}.
	\item \textbf{Eliminar pares dominados}: si un par (peso, valor) tiene otro con menos peso y mas valor, eliminarlo.
\end{itemize}

\paragraph{Trampas y errores frecuentes.}
\begin{itemize}\setlength\itemsep{0pt}
	\item $N$ impar: una mitad tiene $N/2 + 1$ items.
	\item Overflow al sumar pesos/valores en \texttt{long long}.
	\item Sin deduplicacion, el algoritmo es correcto pero mas lento.
\end{itemize}

\paragraph{Codigo.}
Ver \texttt{cpp.tex}, seccion \textit{Knapsack meet-in-the-middle}.

\vspace{0.5em}

\subsubsection{LIS O(n log n)}

\paragraph{Idea intuitiva.}
La \textbf{Longest Increasing Subsequence} se calcula en $O(n \log n)$ con el truco del \textbf{patience sorting}: mantener un arreglo \texttt{dp} donde \texttt{dp[len]} es el minimo ultimo valor de una IS de largo \texttt{len}. Para cada $x$, encontrar el primer \texttt{dp[len] >= x} y reemplazarlo. El largo final es la LIS. La demostracion: si hay dos IS del mismo largo, la que termina en un valor menor es ``mejor'' (puede extenderse mas facilmente).

\paragraph{Propiedades clave (invariantes).}
\begin{itemize}\setlength\itemsep{0pt}
	\item \texttt{dp} es estrictamente creciente.
	\item \texttt{lower\_bound} da LIS estricta; \texttt{upper\_bound} da no-estricta ($\le$).
	\item El algoritmo es \textbf{online}: procesa elementos en orden.
	\item La longitud de \texttt{dp} es la LIS al final.
\end{itemize}

\paragraph{Codigo clave.}
La version online:
\begin{verbatim}
vector<int> dp;
for (int x : a) {
    auto it = lower_bound(dp.begin(), dp.end(), x);
    if (it == dp.end()) dp.push_back(x);
    else *it = x;
}
// dp.size() = longitud de LIS estricta
\end{verbatim}

\paragraph{Complejidad.}
\begin{itemize}\setlength\itemsep{0pt}
	\item $O(n \log n)$.
	\item Memoria $O(n)$.
\end{itemize}

\paragraph{Donde tocar para modificar.}
\begin{itemize}\setlength\itemsep{0pt}
	\item \textbf{Reconstruir la LIS}: aniadir \texttt{parent[]} y \texttt{posInDp[]} para backtrack.
	\item \textbf{Contar LIS}: DP clasica $O(n^2)$ o segment tree sobre $dp$.
	\item \textbf{LDS / Longest non-increasing}: invertir la condicion en \texttt{lower\_bound} o multiplicar por -1.
	\item \textbf{Longest bitonic subsequence}: combinar LIS desde izq y desde der.
	\item \textbf{LIS en 2D}: ordenar por una dimension y aplicar LIS en la otra.
\end{itemize}

\paragraph{Trampas y errores frecuentes.}
\begin{itemize}\setlength\itemsep{0pt}
	\item \texttt{lower\_bound} vs \texttt{upper\_bound}: si tu comparacion es $\le$ vs $<$, cambia.
	\item Si los elementos son \texttt{int} pero podrian ser negativos, aniadir offset para usar como indices.
	\item Para LIS reconstruida, aniadir el array \texttt{parent} que apunta al anterior en la IS.
\end{itemize}

\paragraph{Codigo.}
Ver \texttt{cpp.tex}, seccion \textit{LIS O(n log n)}.

\vspace{0.5em}

\subsubsection{LCS / Edit distance}

\paragraph{Idea intuitiva.}
\textbf{LCS} (Longest Common Subsequence): para dos strings $a$, $b$, $dp[i][j]$ = LCS de $a[0..i)$ y $b[0..j)$. Recurrencia:
\[ dp[i][j] = \begin{cases} dp[i-1][j-1] + 1 & \text{si } a[i-1] = b[j-1] \\ \max(dp[i-1][j], dp[i][j-1]) & \text{si no} \end{cases} \]

\textbf{Edit distance} (Levenshtein): numero minimo de operaciones (insert, delete, replace) para convertir $a$ en $b$. Recurrencia similar con un $+1$ adicional. La idea: ``alinear'' los strings eligiendo matches/mismatches/insert/delete.

\paragraph{Propiedades clave (invariantes).}
\begin{itemize}\setlength\itemsep{0pt}
	\item LCS y Edit Distance son DP $O(nm)$.
	\item Memoria se puede reducir a $O(\min(n, m))$ con rolling array.
	\item Edit distance con peso por operacion distinto: ajustar el \texttt{+1} por el peso.
	\item La subestructura optima: la respuesta para prefijos se construye de prefijos menores.
\end{itemize}

\paragraph{Codigo clave.}
La recurrencia clasica de LCS:
\begin{verbatim}
vector<vector<int>> dp(n + 1, vector<int>(m + 1, 0));
for (int i = 1; i <= n; ++i)
    for (int j = 1; j <= m; ++j)
        if (a[i-1] == b[j-1]) dp[i][j] = dp[i-1][j-1] + 1;
        else dp[i][j] = max(dp[i-1][j], dp[i][j-1]);
return dp[n][m];
\end{verbatim}

\paragraph{Complejidad.}
\begin{itemize}\setlength\itemsep{0pt}
	\item $O(NM)$ tiempo y memoria.
	\item Para $N, M \le 5000$: $\sim 2.5 \cdot 10^7$ operaciones.
\end{itemize}

\paragraph{Donde tocar para modificar.}
\begin{itemize}\setlength\itemsep{0pt}
	\item \textbf{Reconstruir la LCS / las operaciones}: aniadir \texttt{prev[i][j]} para backtrack.
	\item \textbf{LCS de varios strings}: NP-hard para $\ge 3$ strings; en la practica se hace con bitmask.
	\item \textbf{Edit distance con costos distintos}: cambiar el \texttt{+1} por el costo apropiado.
	\item \textbf{Rolling array}: si memoria es problema, mantener solo 2 filas.
	\item \textbf{Hirschberg}: algoritmo $O(NM)$ tiempo, $O(\min(N,M))$ memoria, con reconstruccion.
\end{itemize}

\paragraph{Trampas y errores frecuentes.}
\begin{itemize}\setlength\itemsep{0pt}
	\item Indices en la recurrencia: \texttt{a[i-1]} y \texttt{b[j-1]} (no \texttt{a[i]}).
	\item Para reconstruccion, el caso \texttt{a[i-1] == b[j-1]} tiene prioridad.
	\item Si los strings son muy largos, considerar Hirschberg o bitset optimization.
\end{itemize}

\paragraph{Codigo.}
Ver \texttt{cpp.tex}, seccion \textit{LCS / Edit distance}.

\vspace{0.5em}

\subsubsection{Bitmask DP over subsets}

\paragraph{Idea intuitiva.}
Cuando el estado es un \textbf{subconjunto} de $\{0, \dots, N-1\}$, se representa como una mascara de $N$ bits. Para iterar sobre todos los subconjuntos de \texttt{mask}:
\[ \text{for (sub = mask; sub; sub = (sub - 1) \& mask)} \{ \dots \} \]
Para iterar sobre todos los superconjuntos de \texttt{mask} en $[0, U)$:
\[ \text{for (sup = mask; sup < U; sup = (sup + 1) | mask)} \{ \dots \} \]
La idea: las mascaras forman un lattice de subconjuntos; los trucos aritmeticos (\texttt{(sub-1) \& mask} y \texttt{(sup+1) | mask}) enumeran todo el lattice eficientemente.

\paragraph{Propiedades clave (invariantes).}
\begin{itemize}\setlength\itemsep{0pt}
	\item Iterar subsets: $O(2^k)$ para una mascara con $k$ bits.
	\item \texttt{sub \& mask == sub} siempre; \texttt{mask - sub} no siempre da un subconjunto.
	\item \texttt{\_\_builtin\_popcount(mask)} cuenta los bits.
	\item \texttt{(sub - 1) \& mask} itera subsets de \texttt{mask} en orden decreciente de inclusion.
\end{itemize}

\paragraph{Codigo clave.}
Iteracion sobre subsets de una mascara:
\begin{verbatim}
// Todos los subsets de mask
for (int sub = mask; sub; sub = (sub - 1) & mask) {
    // procesar sub
}
// Incluir sub = 0 si es necesario
sub = 0;  // procesar primero
for (int sup = mask; sup < (1 << N); sup = (sup + 1) | mask) {
    // procesar sup (todos los superconjuntos)
}
\end{verbatim}

\paragraph{Complejidad.}
\begin{itemize}\setlength\itemsep{0pt}
	\item $O(N \cdot 2^N)$ para una DP estandar sobre todos los subsets.
	\item Iterar subsets de \texttt{mask}: $O(2^{|mask|})$.
\end{itemize}

\paragraph{Donde tocar para modificar.}
\begin{itemize}\setlength\itemsep{0pt}
	\item \textbf{TSP / Hamiltonian path}: bitmask DP (ver tambien \textit{Hamiltonian path} en \texttt{cpp.tex} seccion Grafos).
	\item \textbf{Set cover}: iterar subsets.
	\item \textbf{DP con supermask}: si el estado incluye ``todos los elementos que \textbf{no} hemos usado'', usar supermask iteration.
	\item \textbf{Bitmasks grandes}: con $N > 20$, aniadir compresion o simulacion con segmentos.
\end{itemize}

\paragraph{Trampas y errores frecuentes.}
\begin{itemize}\setlength\itemsep{0pt}
	\item El orden de iteracion de subsets de \texttt{mask} es decreciente (de \texttt{mask} a $0$).
	\item \texttt{mask = 0} no tiene subsets no triviales; el bucle lo saltea correctamente.
	\item El truco \texttt{(sub - 1) \& mask} requiere cuidado con \texttt{sub = 0} (termina el bucle).
	\item Superconjuntos: \texttt{sup < (1 << N)} debe estar bien definido.
\end{itemize}

\paragraph{Codigo.}
Ver \texttt{cpp.tex}, seccion \textit{Bitmask DP over subsets}.

\vspace{0.5em}

\subsubsection{SOS DP}

\paragraph{Idea intuitiva.}
Para calcular, para cada \texttt{mask}, la suma de $f(\text{submask})$ sobre todos los $\text{submask} \subseteq \text{mask}$:
\[ dp[mask] = \sum_{\text{sub} \subseteq mask} f(\text{sub}). \]
Se computa en $O(N \cdot 2^N)$ mediante un loop por bit: para cada $i$, sumar a $dp[mask]$ la contribucion de $dp[mask \oplus (1 << i)]$ cuando $mask$ tiene el bit $i$. La idea: cada subconjunto $s \subseteq m$ se obtiene quitando bits uno a uno; asi, sumar incrementalmente captura todos los subconjuntos en $O(N \cdot 2^N)$.

\paragraph{Propiedades clave (invariantes).}
\begin{itemize}\setlength\itemsep{0pt}
	\item Inicializar $dp = f$.
	\item Para cada bit $i$: \texttt{if (mask \& (1 << i)) dp[mask] += dp[mask \^{} (1 << i)]}.
	\item Variantes: supersets DP (sumar sobre superconjuntos).
	\item El orden de iteracion (bit externo, mascara interno) es crucial.
\end{itemize}

\paragraph{Codigo clave.}
SOS DP (Sum over Subsets):
\begin{verbatim}
vector<long long> dp(1 << N);
for (int i = 0; i < (1 << N); ++i) dp[i] = f[i];
for (int bit = 0; bit < N; ++bit)
    for (int mask = 0; mask < (1 << N); ++mask)
        if (mask & (1 << bit))
            dp[mask] += dp[mask ^ (1 << bit)];
// dp[mask] = suma de f(sub) para sub ⊆ mask
\end{verbatim}

\paragraph{Complejidad.}
\begin{itemize}\setlength\itemsep{0pt}
	\item $O(N \cdot 2^N)$.
	\item Memoria $O(2^N)$.
\end{itemize}

\paragraph{Donde tocar para modificar.}
\begin{itemize}\setlength\itemsep{0pt}
	\item \textbf{Max/min en lugar de suma}: cambiar \texttt{+=} por la operacion correspondiente.
	\item \textbf{Supersets DP}: iterar \texttt{mask} de $0$ a $2^N$ y aniadir \texttt{dp[mask] += dp[mask | (1<<i)]}.
	\item \textbf{XOR convolution} (subset convolution): version mas avanzada.
	\item \textbf{Aplicaciones}: contar subconjuntos con propiedades especificas.
\end{itemize}

\paragraph{Trampas y errores frecuentes.}
\begin{itemize}\setlength\itemsep{0pt}
	\item Si $N$ es 30 o mas, $2^N$ no entra en memoria; usar trucos o reducir $N$.
	\item El sentido del XOR (sumar a \texttt{mask \^{} bit}) es para ``agregar el bit $i$ al subconjunto''.
	\item En supersets DP, la condicion es \texttt{(mask \& (1<<bit)) == 0}.
\end{itemize}

\paragraph{Codigo.}
Ver \texttt{cpp.tex}, seccion \textit{SOS DP}.

\vspace{0.5em}

\subsubsection{Digit DP template}

\paragraph{Idea intuitiva.}
Contar numeros en $[0, N]$ que cumplen una propiedad sobre sus digitos. La idea: DP por posicion, manteniendo:
\begin{itemize}\setlength\itemsep{0pt}
	\item \texttt{pos}: posicion actual (de izquierda a derecha).
	\item \texttt{tight}: si los digitos ya colocados son exactamente los de $N$ (limite superior estricto).
	\item \texttt{state}: estado adicional (digitos usados, suma, etc.).
\end{itemize}
La clave: la condicion \texttt{tight} permite acotar correctamente el siguiente digito; sin ella, la DP contaria incorrectamente.

\paragraph{Propiedades clave (invariantes).}
\begin{itemize}\setlength\itemsep{0pt}
	\item Si \texttt{tight} es true, el siguiente digito esta acotado por el de $N$; si no, puede ser 0-9.
	\item Memoria: $O(D \cdot 2 \cdot \text{states})$.
	\item Caso base: si \texttt{pos == D}, devolver 1 si el estado cumple la propiedad, sino 0.
	\item El bit \texttt{started} distingue ``numero con menos digitos'' (con leading zeros).
\end{itemize}

\paragraph{Codigo clave.}
La estructura basica:
\begin{verbatim}
long long dp(int pos, bool tight, State s) {
    if (pos == D) return check(s) ? 1 : 0;
    if (memo[pos][tight][s] != -1) return memo;
    int limit = tight ? digits[pos] : 9;
    long long ans = 0;
    for (int d = 0; d <= limit; ++d)
        ans += dp(pos + 1, tight && (d == limit), update(s, d, pos));
    return memo[pos][tight][s] = ans;
}
\end{verbatim}

\paragraph{Complejidad.}
\begin{itemize}\setlength\itemsep{0pt}
	\item $O(D \cdot 10 \cdot \text{states})$ donde $D$ es el numero de digitos.
	\item Para $D \le 19$ (hasta $10^{18}$): $\sim 200 \cdot \text{states}$.
\end{itemize}

\paragraph{Donde tocar para modificar.}
\begin{itemize}\setlength\itemsep{0pt}
	\item \textbf{Cambiar la propiedad}: modificar \texttt{updateState} y la condicion base.
	\item \textbf{Multiples estados}: aniadir mas dimensiones al \texttt{memo}.
	\item \textbf{Leading zeros}: aniadir un flag \texttt{started} para distinguir numeros con menos digitos.
	\item \textbf{Rango $[L, R]$}: calcular $f(R) - f(L-1)$.
\end{itemize}

\paragraph{Trampas y errores frecuentes.}
\begin{itemize}\setlength\itemsep{0pt}
	\item Olvidar resetear \texttt{memo} entre queries.
	\item Si la cantidad de estados es grande, el \texttt{memo[pos][tight][state]} puede ser muy grande; usar \texttt{map}.
	\item La condicion \texttt{started} es crucial para numeros con menos digitos que $D$.
	\item Si $N$ es muy grande ($> 10^{18}$), aniadir soporte para long long.
\end{itemize}

\paragraph{Codigo.}
Ver \texttt{cpp.tex}, seccion \textit{Digit DP template}.

\vspace{0.5em}

\subsubsection{Tree DP (reroot)}

\paragraph{Idea intuitiva.}
Calcular alguna propiedad asumiendo que \textbf{cada nodo es la raiz} del arbol. La tecnica clasica de \textbf{rerooting} en dos pasadas:
\begin{itemize}\setlength\itemsep{0pt}
	\item \textbf{DFS 1}: desde una raiz arbitraria, calcular \texttt{dp[v]} para el subarbol de $v$.
	\item \textbf{DFS 2}: para cada hijo $u$ de $v$, calcular \texttt{up[u]} = respuesta si la raiz fuera $u$.
\end{itemize}
La demostracion: tras la primera DFS, cada nodo conoce ``lo que aporta su subarbol''; en la segunda, se transmite ``lo que aporta el resto del arbol''.

\paragraph{Propiedades clave (invariantes).}
\begin{itemize}\setlength\itemsep{0pt}
	\item \texttt{dp[v]} depende solo del subarbol de $v$ (con la raiz original).
	\item \texttt{up[u]} se calcula de manera que cuando la raiz es $u$, ``todo lo que estaba en $v$ se ajusta''.
	\item Formula para reroot depende del problema; el snippet usa el ejemplo ``suma de distancias''.
	\item La combinacion \texttt{dp[u] + up[u]} da la respuesta con raiz $u$.
\end{itemize}

\paragraph{Codigo clave.}
Reroot para suma de distancias:
\begin{verbatim}
// dfs1: dp[v] = suma de distancias desde v al subarbol
void dfs1(int v, int p) {
    dp[v] = 0;
    for (int u : adj[v]) if (u != p) {
        dfs1(u, v);
        dp[v] += dp[u] + sz[u];  // cada hijo aporta su subarbol
    }
}
// dfs2: up[u] = distancia desde el resto del arbol a u
void dfs2(int v, int p) {
    for (int u : adj[v]) if (u != p) {
        up[u] = up[v] + dp[v] - dp[u] - sz[u] + (n - sz[u]);
        dfs2(u, v);
    }
}
\end{verbatim}

\paragraph{Complejidad.}
\begin{itemize}\setlength\itemsep{0pt}
	\item $O(N)$ para ambas pasadas.
	\item Cada arista se procesa un numero constante de veces.
\end{itemize}

\paragraph{Donde tocar para modificar.}
\begin{itemize}\setlength\itemsep{0pt}
	\item \textbf{Cambiar la cantidad a calcular}: modificar la operacion en \texttt{dfs1} y la formula en \texttt{dfs2}.
	\item \textbf{Arbol con pesos}: aniadir \texttt{weights[v]} y ajustar.
	\item \textbf{Re-root en arboles con ciclos}: requiere otra tecnica (DSU on tree).
	\item \textbf{Rooted trees diferentes}: si la propiedad no es invariante por raiz, usar otra tecnica.
\end{itemize}

\paragraph{Trampas y errores frecuentes.}
\begin{itemize}\setlength\itemsep{0pt}
	\item La formula de \texttt{up[u]} debe sumar y restar correctamente las contribuciones de $u$ y del resto.
	\item Para arboles con pesos, aniadir el peso en las formulas.
	\item Verificar con ejemplos pequenos antes de generalizar.
\end{itemize}

\paragraph{Codigo.}
Ver \texttt{cpp.tex}, seccion \textit{Tree DP (reroot)}.

% ============================================================
\section{Geometria}
% ============================================================

\subsection{Puntos y segmentos}

\subsubsection{Point + cross + dot}

\paragraph{Idea intuitiva.}
La estructura \texttt{Point} representa un punto en 2D con coordenadas enteras. Las dos operaciones vectoriales clave son:
\begin{itemize}\setlength\itemsep{0pt}
	\item \textbf{Cross product} $P \times Q = P_x Q_y - P_y Q_x$. Da el ``area con signo'' del paralelogramo formado por $P$ y $Q$. Usado para orientacion.
	\item \textbf{Dot product} $P \cdot Q = P_x Q_x + P_y Q_y$. Da el producto de las longitudes por el coseno del angulo. Usado para proyeccion.
\end{itemize}
La interpretacion geometrica del cross product: es el area con signo del paralelogramo; el signo indica la orientacion relativa.

\paragraph{Propiedades clave (invariantes).}
\begin{itemize}\setlength\itemsep{0pt}
	\item \texttt{cross(a, b) > 0} si $b$ esta a la izquierda de $a$ (sistema de coordenadas usual).
	\item \texttt{cross(a, b) = 0} si son colineales.
	\item \texttt{dot(a, b) > 0} si el angulo es agudo.
	\item El snippet define \texttt{cross} respecto a \texttt{this}, util para orientacion de 3 puntos.
	\item $|cross|^2 + |dot|^2 = |P|^2 \cdot |Q|^2$ (identidad de Lagrange).
\end{itemize}

\paragraph{Codigo clave.}
Las dos operaciones vectoriales:
\begin{verbatim}
long long cross(const Point& a, const Point& b) {
    return (long long)a.x * b.y - (long long)a.y * b.x;
}
long long dot(const Point& a, const Point& b) {
    return (long long)a.x * b.x + (long long)a.y * b.y;
}
\end{verbatim}

\paragraph{Complejidad.}
\begin{itemize}\setlength\itemsep{0pt}
	\item $O(1)$ por operacion.
	\item Constante muy pequena (4 multiplicaciones + 2 sumas).
\end{itemize}

\paragraph{Donde tocar para modificar.}
\begin{itemize}\setlength\itemsep{0pt}
	\item \textbf{Coordenadas doubles}: cambiar \texttt{int} por \texttt{double} o \texttt{long double}. Cuidado con la precision.
	\item \textbf{3D}: aniadir \texttt{z} y adaptar cross/dot.
	\item \textbf{Operadores adicionales}: aniadir \texttt{operator+}, \texttt{operator-}, \texttt{operator*} (escalar).
	\item \textbf{Distancia al cuadrado}: aniadir \texttt{len2()} que retorna $x^2 + y^2$.
\end{itemize}

\paragraph{Trampas y errores frecuentes.}
\begin{itemize}\setlength\itemsep{0pt}
	\item Cross product \textbf{no es} conmutativo; $P \times Q = -Q \times P$.
	\item Overflow con coordenadas grandes ($|x|, |y| > 10^4$) si se hace \texttt{x * y}.
	\item El signo de cross depende del sistema de coordenadas (en computacion, usualmente CCW positivo).
	\item Para ``igualdad'' de doubles, aniadir tolerancia \texttt{eps}.
\end{itemize}

\paragraph{Codigo.}
Ver \texttt{cpp.tex}, seccion \textit{Point + cross + dot}.

\vspace{0.5em}

\subsubsection{Interseccion de segmentos}

\paragraph{Idea intuitiva.}
Dados dos segmentos $[a, b]$ y $[c, d]$, determinar si se intersectan. La idea clasica:
\begin{itemize}\setlength\itemsep{0pt}
	\item Verificar orientaciones: $a, b, c$ y $a, b, d$ deben estar en lados opuestos de la recta $ab$; igual para $cd$ y $ab$.
	\item Caso colineal: si las rectas son paralelas (\texttt{cross} $= 0$) y los segmentos son colineales, verificar que los bounding boxes se intersectan.
\end{itemize}
La demostracion: dos segmentos se intersectan sii cada segmento ``cruza'' la linea que contiene al otro (estricto) o ambos son colineales con solapamiento.

\paragraph{Propiedades clave (invariantes).}
\begin{itemize}\setlength\itemsep{0pt}
	\item Funciona con segmentos en cualquier orientacion.
	\item Caso colineal: requiere check adicional de bounding box.
	\item Para intersectar en un \textbf{punto} especifico: aniadir division con cuidado de precision.
	\item Funciona con segmentos verticales/horizontales sin casos especiales.
\end{itemize}

\paragraph{Codigo clave.}
Test clasico de interseccion:
\begin{verbatim}
int sign(long long x) { return (x > 0) - (x < 0); }
bool intersect(Point a, Point b, Point c, Point d) {
    long long o1 = cross(b - a, c - a);
    long long o2 = cross(b - a, d - a);
    long long o3 = cross(d - c, a - c);
    long long o4 = cross(d - c, b - c);
    if (sign(o1) != sign(o2) && sign(o3) != sign(o4)) return true;
    // casos colineales
    if (o1 == 0 && onSegment(a, b, c)) return true;
    if (o2 == 0 && onSegment(a, b, d)) return true;
    if (o3 == 0 && onSegment(c, d, a)) return true;
    if (o4 == 0 && onSegment(c, d, b)) return true;
    return false;
}
\end{verbatim}

\paragraph{Complejidad.}
\begin{itemize}\setlength\itemsep{0pt}
	\item $O(1)$.
	\item Constante pequena.
\end{itemize}

\paragraph{Donde tocar para modificar.}
\begin{itemize}\setlength\itemsep{0pt}
	\item \textbf{Interseccion con punto}: aniadir la division despues de verificar orientacion.
	\item \textbf{Interseccion rayo-segmento}: ajustar la primera condicion.
	\item \textbf{Poligonos convexos}: usar separacion de ejes (SAT).
	\item \textbf{Multiples queries}: precomputar bounding boxes o usar segment tree.
\end{itemize}

\paragraph{Trampas y errores frecuentes.}
\begin{itemize}\setlength\itemsep{0pt}
	\item Si los segmentos tocan solo en un extremo, el snippet lo considera interseccion (algunos problemas quieren estricto). Ajustar segun el problema.
	\item El caso colineal es comun; olvidar el check adicional de bounding box da falsos negativos.
	\item En problemas con muchos segmentos, considerar sweep line para $O(n \log n)$.
\end{itemize}

\paragraph{Codigo.}
Ver \texttt{cpp.tex}, seccion \textit{Interseccion de segmentos}.

\vspace{0.5em}

\subsubsection{Contains colineal}

\paragraph{Idea intuitiva.}
Verifica si un punto $c$ esta sobre el segmento $[a, b]$ (asumiendo que ya son colineales). Si \texttt{cross(a, b, c) = 0} y $c$ esta dentro del bounding box de $[a, b]$, entonces $c$ esta en el segmento. La razon: si los tres puntos son colineales, $c$ esta en el segmento sii sus coordenadas estan entre las de $a$ y $b$.

\paragraph{Propiedades clave (invariantes).}
\begin{itemize}\setlength\itemsep{0pt}
	\item Asume colinealidad (no la verifica mas alla de \texttt{cross == 0}).
	\item Usa comparaciones con \texttt{<=} para extremos \textbf{inclusivos}.
	\item Funciona con segmentos verticales (la condicion se reduce a una sola coordenada).
\end{itemize}

\paragraph{Codigo clave.}
Test de contencion en segmento (asume colinealidad):
\begin{verbatim}
bool onSegment(Point a, Point b, Point c) {
    return min(a.x, b.x) <= c.x && c.x <= max(a.x, b.x) &&
           min(a.y, b.y) <= c.y && c.y <= max(a.y, b.y);
}
\end{verbatim}

\paragraph{Complejidad.}
\begin{itemize}\setlength\itemsep{0pt}
	\item $O(1)$.
	\item Constante minima.
\end{itemize}

\paragraph{Donde tocar para modificar.}
\begin{itemize}\setlength\itemsep{0pt}
	\item \textbf{Contener estricto}: cambiar \texttt{<=} a \texttt{<} para excluir extremos.
	\item \textbf{Coordenadas doubles}: usar tolerancia (\texttt{eps}).
	\item \textbf{Solo x o solo y}: aniadir check especifico para segmentos verticales/horizontales.
\end{itemize}

\paragraph{Trampas y errores frecuentes.}
\begin{itemize}\setlength\itemsep{0pt}
	\item No usar este snippet para puntos \textbf{no colineales}; siempre verificar primero la colinealidad.
	\item Si los puntos tienen doubles, el chequeo \texttt{cross == 0} puede fallar; aniadir tolerancia.
	\item El bounding box incluye los extremos; verificar si el problema quiere exclusivo o inclusivo.
\end{itemize}

\paragraph{Codigo.}
Ver \texttt{cpp.tex}, seccion \textit{Contains colineal}.

\subsection{Poligonos}

\subsubsection{Convex Hull (monotono)}

\paragraph{Idea intuitiva.}
Algoritmo de \textbf{Andrew's monotone chain}: ordenar puntos por $(x, y)$, luego construir la hull inferior y superior por separado. Para cada punto, aniadirlo al hull y, mientras el ultimo giro no sea ``counter-clockwise'' (\texttt{cross > 0}), popear. La demostracion: si los puntos se procesan en orden, el hull inferior y superior cubren la frontera completa sin auto-intersecciones.

\paragraph{Propiedades clave (invariantes).}
\begin{itemize}\setlength\itemsep{0pt}
	\item Ordena los puntos en $O(n \log n)$.
	\item Construye la hull en $O(n)$ despues.
	\item Output: poligono convexo en counter-clockwise, sin punto final duplicado.
	\item El \texttt{<= 0} incluye colineales; usar \texttt{< 0} para excluirlos.
	\item Cada punto aparece en la hull inferior o superior (no ambos).
\end{itemize}

\paragraph{Codigo clave.}
Construccion del hull inferior y superior:
\begin{verbatim}
sort(p.begin(), p.end());
vector<Point> lower, upper;
for (auto& p : points) {
    while (lower.size() >= 2 && cross(lower.back() - lower[lower.size()-2],
                                       p - lower.back()) <= 0)
        lower.pop_back();
    lower.push_back(p);
}
for (int i = points.size() - 1; i >= 0; --i) {
    auto& p = points[i];
    while (upper.size() >= 2 && cross(upper.back() - upper[upper.size()-2],
                                       p - upper.back()) <= 0)
        upper.pop_back();
    upper.push_back(p);
}
vector<Point> hull = lower;
hull.insert(hull.end(), upper.begin() + 1, upper.end() - 1);
\end{verbatim}

\paragraph{Complejidad.}
\begin{itemize}\setlength\itemsep{0pt}
	\item $O(n \log n)$.
	\item Constante muy pequena (solo cross products).
\end{itemize}

\paragraph{Donde tocar para modificar.}
\begin{itemize}\setlength\itemsep{0pt}
	\item \textbf{Colineales}: cambiar \texttt{<= 0} por \texttt{< 0} para excluirlos del hull.
	\item \textbf{Puntos repetidos}: \texttt{unique} despues de sort.
	\item \textbf{Hull 3D}: gift wrapping es $O(n^2)$; convex hull 3D es mas complejo.
	\item \textbf{Output con los puntos}: aniadir el hull completo (no solo upper/lower).
\end{itemize}

\paragraph{Trampas y errores frecuentes.}
\begin{itemize}\setlength\itemsep{0pt}
	\item El \texttt{pop\_back} y el \texttt{reverse} son esenciales para no duplicar el ultimo punto.
	\item Para excluir colineales, ajustar el \texttt{<= 0} vs \texttt{< 0}.
	\item Si los puntos estan duplicados, aniadir \texttt{unique} antes del sort.
\end{itemize}

\paragraph{Codigo.}
Ver \texttt{cpp.tex}, seccion \textit{Convex Hull (monotono)}.

\vspace{0.5em}

\subsubsection{Polygon Area}

\paragraph{Idea intuitiva.}
El area (con signo) de un poligono se calcula con la \textbf{Shoelace formula}: $A = \frac{1}{2} \sum_{i=0}^{n-1} (x_i y_{i+1} - x_{i+1} y_i)$, donde el indice es modulo $n$. El snippet retorna el \textbf{doble} del area (en valor absoluto). La formula sale del ``sumar areas de trapezoides'': cada lado del poligono define un trapezoide con el eje $x$; la suma con signo da el area.

\paragraph{Propiedades clave (invariantes).}
\begin{itemize}\setlength\itemsep{0pt}
	\item El signo indica la orientacion: positivo si counter-clockwise.
	\item Para el doble: no dividir por 2 si no es necesario (evita fracciones).
	\item Funciona para poligonos simples (no auto-intersectantes).
	\item Si el poligono es no convexo, la formula da el area ``con signo''.
\end{itemize}

\paragraph{Codigo clave.}
La suma de trapezoides:
\begin{verbatim}
long long shoelace(const vector<Point>& p) {
    long long s = 0;
    int n = p.size();
    for (int i = 0; i < n; ++i)
        s += (long long)p[i].x * p[(i+1)%n].y - (long long)p[(i+1)%n].x * p[i].y;
    return abs(s);  // doble del area
}
\end{verbatim}

\paragraph{Complejidad.}
\begin{itemize}\setlength\itemsep{0pt}
	\item $O(n)$.
	\item Constante minima.
\end{itemize}

\paragraph{Donde tocar para modificar.}
\begin{itemize}\setlength\itemsep{0pt}
	\item \textbf{Area exacta}: dividir por 2 (con cuidado si da fraccion; usar \texttt{double}).
	\item \textbf{Area firmada}: quitar \texttt{abs} y dejar el signo.
	\item \textbf{Punto dentro de poligono convexo}: con la hull, chequear que esta a un mismo lado de todas las aristas.
	\item \textbf{Triangulacion}: descomponer en triangulos y sumar areas.
\end{itemize}

\paragraph{Trampas y errores frecuentes.}
\begin{itemize}\setlength\itemsep{0pt}
	\item Si los puntos no estan en orden (CCW o CW), el area firmada puede ser incorrecta. Para signed area consistente, usar siempre el mismo orden.
	\item Para poligonos con holes, sumar las areas con signo apropiado.
	\item Con doubles, aniadir tolerancia \texttt{eps}.
\end{itemize}

\paragraph{Codigo.}
Ver \texttt{cpp.tex}, seccion \textit{Polygon Area}.

\vspace{0.5em}

\subsubsection{Distance point-segment / point-line}

\paragraph{Idea intuitiva.}
Distancia minima de un punto $p$ a un segmento $[a, b]$ o a la recta que pasa por $a$ y $b$.
\begin{itemize}\setlength\itemsep{0pt}
	\item \textbf{Linea}: $d = \frac{|ab \times ap|}{|ab|}$.
	\item \textbf{Segmento}: si la proyeccion de $p$ sobre $ab$ cae fuera del segmento, la distancia es al extremo mas cercano.
\end{itemize}
La razon geometrica: la distancia a una linea es el area del paralelogramo dividido por la base; la proyeccion ortogonal indica si estamos ``dentro'' o ``fuera'' del segmento.

\paragraph{Propiedades clave (invariantes).}
\begin{itemize}\setlength\itemsep{0pt}
	\item \texttt{dist2PointLine} retorna distancia \textbf{al cuadrado}.
	\item \texttt{dist2PointSegment} distingue 3 casos: $p$ antes de $a$, despues de $b$, o proyectado sobre el segmento.
	\item Si \texttt{len2 = 0} (degenerate), la distancia es $|ap|$ o $|bp|$.
	\item La proyeccion se calcula con dot products y el parametro $t \in [0, 1]$.
\end{itemize}

\paragraph{Codigo clave.}
Distancia al cuadrado a segmento:
\begin{verbatim}
long long dist2Segment(Point p, Point a, Point b) {
    long long len2 = (a-b).len2();
    if (len2 == 0) return (p-a).len2();  // degenerate
    long long t = max(0LL, min(1LL, dot(p-a, b-a) / len2));
    Point proj = {a.x + t*(b.x - a.x), a.y + t*(b.y - a.y)};
    return (p - proj).len2();
}
\end{verbatim}

\paragraph{Complejidad.}
\begin{itemize}\setlength\itemsep{0pt}
	\item $O(1)$.
	\item Constante minima.
\end{itemize}

\paragraph{Donde tocar para modificar.}
\begin{itemize}\setlength\itemsep{0pt}
	\item \textbf{Distancia (no al cuadrado)}: aniadir \texttt{sqrt} al final.
	\item \textbf{Coordenadas doubles}: aniadir \texttt{long double} y tolerancia.
	\item \textbf{3D}: aniadir coordenada $z$ y adaptar.
	\item \textbf{Multiples queries}: usar KD-tree o segment tree para acelerar.
\end{itemize}

\paragraph{Trampas y errores frecuentes.}
\begin{itemize}\setlength\itemsep{0pt}
	\item Overflow en \texttt{area2 * area2} si el area es grande; \texttt{long long} alcanza para $|x|, |y| \le 10^9$ pero no mas.
	\item Si $a = b$ (segmento degenerado), aniadir check explicito.
	\item En doubles, aniadir tolerancia \texttt{eps} para comparar.
\end{itemize}

\paragraph{Codigo.}
Ver \texttt{cpp.tex}, seccion \textit{Distance point-segment / point-line}.

\vspace{0.5em}

\subsubsection{Projection of point onto line}

\paragraph{Idea intuitiva.}
Proyeccion ortogonal de $p$ sobre la recta que pasa por $a$ y $b$. El parametro $t$ tal que la proyeccion es $a + t \cdot (b - a)$ es:
\[ t = \frac{(p - a) \cdot (b - a)}{|b - a|^2}. \]
La idea: descomponemos $\vec{ap}$ en una componente paralela a $\vec{ab}$ y una perpendicular; $t$ mide la posicion en la primera.

\paragraph{Propiedades clave (invariantes).}
\begin{itemize}\setlength\itemsep{0pt}
	\item Si $t \in [0, 1]$, la proyeccion cae sobre el \textbf{segmento}; si no, cae sobre la extension.
	\item Retorna \texttt{double} (puede tener precision issues).
	\item $t < 0$: $p$ esta ``antes'' de $a$; $t > 1$: $p$ esta ``despues'' de $b$.
\end{itemize}

\paragraph{Codigo clave.}
Proyeccion ortogonal:
\begin{verbatim}
Point projection(Point p, Point a, Point b) {
    Point ab = b - a, ap = p - a;
    double t = (double)dot(ap, ab) / ab.len2();
    return {a.x + t * ab.x, a.y + t * ab.y};
}
\end{verbatim}

\paragraph{Complejidad.}
\begin{itemize}\setlength\itemsep{0pt}
	\item $O(1)$.
	\item Constante minima (dos dot products y una division).
\end{itemize}

\paragraph{Donde tocar para modificar.}
\begin{itemize}\setlength\itemsep{0pt}
	\item \textbf{Proyeccion sobre segmento}: clipar $t$ a $[0, 1]$ antes de calcular el punto.
	\item \textbf{Distancia al segmento via proyeccion}: combinar con el modulo anterior.
	\item \textbf{Solo el parametro $t$}: si solo necesitas saber si esta dentro del segmento, retorna $t$ sin construir el punto.
\end{itemize}

\paragraph{Trampas y errores frecuentes.}
\begin{itemize}\setlength\itemsep{0pt}
	\item Si $a = b$ ($|ab| = 0$), la proyeccion no esta definida; aniadir guard.
	\item En doubles, aniadir tolerancia \texttt{eps} para comparar $t$ con 0 o 1.
	\item La division por \texttt{len2} puede ser cara si las coordenadas son grandes; usar \texttt{double} o precomputar.
\end{itemize}

\paragraph{Codigo.}
Ver \texttt{cpp.tex}, seccion \textit{Projection of point onto line}.

\vspace{0.5em}

\subsubsection{Point in polygon (ray casting)}

\paragraph{Idea intuitiva.}
Algoritmo de \textbf{ray casting}: contar cuantas veces un rayo horizontal desde $p$ cruza los bordes del poligono. Si es impar, $p$ esta dentro; si par, fuera. El snippet implementa la version clasica recorriendo aristas. La razon: cada cruce cambia el lado del rayo respecto al poligono; un numero impar de cambios significa ``terminamos dentro''.

\paragraph{Propiedades clave (invariantes).}
\begin{itemize}\setlength\itemsep{0pt}
	\item Funciona con poligonos simples (no necesariamente convexos).
	\item Las intersecciones en vertices cuentan especialmente.
	\item Para poligonos con \textbf{huecos}, sigue funcionando.
	\item La respuesta es estrictamente 0 (fuera) o 1 (dentro) para poligonos simples.
\end{itemize}

\paragraph{Codigo clave.}
El test clasico de ray casting:
\begin{verbatim}
bool pointInPolygon(Point p, const vector<Point>& poly) {
    int n = poly.size();
    bool inside = false;
    for (int i = 0, j = n - 1; i < n; j = i++) {
        if ((poly[i].y > p.y) != (poly[j].y > p.y) &&
            p.x < (poly[j].x - poly[i].x) * (p.y - poly[i].y) /
                    (poly[j].y - poly[i].y) + poly[i].x)
            inside = !inside;
    }
    return inside;
}
\end{verbatim}

\paragraph{Complejidad.}
\begin{itemize}\setlength\itemsep{0pt}
	\item $O(n)$ por query.
	\item Memoria $O(1)$ (no se almacena el rayo).
\end{itemize}

\paragraph{Donde tocar para modificar.}
\begin{itemize}\setlength\itemsep{0pt}
	\item \textbf{Winding number}: version alternativa que cuenta el numero de vueltas.
	\item \textbf{Poligono convexo}: check lineal contra las aristas (mas rapido).
	\item \textbf{Multiples queries}: precomputar o usar segment tree sobre las aristas.
	\item \textbf{Punto en frontera}: aniadir check adicional si es colineal con alguna arista.
\end{itemize}

\paragraph{Trampas y errores frecuentes.}
\begin{itemize}\setlength\itemsep{0pt}
	\item El manejo del vertice (cuando el rayo pasa por un vertice) puede causar double-counting; el snippet lo evita con la condicion \texttt{(poly[i].y > p.y) != (poly[j].y > p.y)}.
	\item Si el poligono tiene holes, el algoritmo sigue funcionando.
	\item Para poligonos no simples (auto-intersectantes), el resultado puede ser ambiguo.
\end{itemize}

\paragraph{Codigo.}
Ver \texttt{cpp.tex}, seccion \textit{Point in polygon (ray casting)}.

\vspace{0.5em}

\subsubsection{Minkowski sum (convex polygons)}

\paragraph{Idea intuitiva.}
La \textbf{suma de Minkowski} $A \oplus B$ de dos poligonos convexos es otro poligono convexo que representa $\{a + b : a \in A, b \in B\}$. Algoritmo: ordenar las aristas por angulo polar (CCW) y ``mergearlas'' como en merge sort; el resultado es un poligono CCW sin repetir el ultimo punto. La idea: el resultado tiene una arista paralela a cada arista de los originales; la suma se obtiene barriendo los angulos.

\paragraph{Propiedades clave (invariantes).}
\begin{itemize}\setlength\itemsep{0pt}
	\item Si $A, B$ son convexos con $n, m$ vertices, $A \oplus B$ tiene $\le n + m$ vertices.
	\item Util para \textbf{detectar colision}: $A$ colisiona con $B$ si $0 \in A \oplus (-B)$.
	\item Asume poligonos en orden \textbf{counter-clockwise}.
	\item El resultado es convexo (suma de convexos es convexo).
\end{itemize}

\paragraph{Codigo clave.}
La suma de Minkowski:
\begin{verbatim}
vector<Point> minkowski(const vector<Point>& A, const vector<Point>& B) {
    vector<Point> result;
    int i = 0, j = 0;
    int n = A.size(), m = B.size();
    do {
        result.push_back(A[i] + B[j]);
        long long cross = (A[(i+1)%n] - A[i]) ^ (B[(j+1)%m] - B[j]);
        if (cross >= 0) i = (i + 1) % n;
        if (cross <= 0) j = (j + 1) % m;
    } while (i || j);
    return result;
}
\end{verbatim}

\paragraph{Complejidad.}
\begin{itemize}\setlength\itemsep{0pt}
	\item $O(n + m)$.
	\item Memoria $O(n + m)$.
\end{itemize}

\paragraph{Donde tocar para modificar.}
\begin{itemize}\setlength\itemsep{0pt}
	\item \textbf{Deteccion de colision}: aniadir $(-B)$ (reflejar) y verificar si $0$ esta en el resultado.
	\item \textbf{Poligonos no convexos}: la suma puede no ser poligono simple; algoritmo distinto.
	\item \textbf{Multiples queries}: precomputar las aristas una vez.
	\item \textbf{Robot motion planning}: usar para planear trayectorias.
\end{itemize}

\paragraph{Trampas y errores frecuentes.}
\begin{itemize}\setlength\itemsep{0pt}
	\item Si los poligonos no son CCW, la suma es incorrecta. Normalizar antes.
	\item El bucle termina cuando ambos indices vuelven a 0; aniadir check.
	\item Para la deteccion de colision, aniadir el test \texttt{0 in A \^{} (-B)}.
\end{itemize}

\paragraph{Codigo.}
Ver \texttt{cpp.tex}, seccion \textit{Minkowski sum (convex polygons)}.

\subsection{Herramientas}

\subsubsection{Li Chao Tree (long long)}

\paragraph{Idea intuitiva.}
Mantiene un \textbf{lower envelope} (o upper) de un conjunto de rectas $y = mx + b$ sobre un dominio $[X_{LO}, X_{HI}]$ y responde ``minimo $y$ en un $x$ dado'' en $O(\log X)$. La estructura es un segment tree donde cada nodo guarda la recta ``mejor'' para su segmento medio; al aniadir una recta, se compara con la actual y se decide cual es mejor en cada mitad, descendiendo recursivamente. La demostracion: en cada nivel, solo se guarda la mejor recta para el ``punto medio''; las otras se transmiten a los hijos.

\paragraph{Propiedades clave (invariantes).}
\begin{itemize}\setlength\itemsep{0pt}
	\item \textbf{Online}: cada \texttt{addLine} en $O(\log X)$.
	\item Solo insercion; no hay borrar (para borrar, usar Li Chao segmentado con reset).
	\item La recta guardada en un nodo es la ``mejor'' en su punto medio.
	\item Las rectas ``perdedoras'' se transmiten a los hijos para futuros inserts.
\end{itemize}

\paragraph{Codigo clave.}
La insercion en Li Chao:
\begin{verbatim}
void addLine(Line nw, int v = 1, int l = X_LO, int r = X_HI) {
    int m = (l + r) / 2;
    bool left = nw.get(l) < tree[v].get(l);
    bool mid = nw.get(m) < tree[v].get(m);
    if (mid) swap(tree[v], nw);
    if (r - l == 1) return;
    if (left != mid) addLine(nw, v*2, l, m);
    else             addLine(nw, v*2+1, m, r);
}
\end{verbatim}

\paragraph{Complejidad.}
\begin{itemize}\setlength\itemsep{0pt}
	\item $O(\log X)$ por insercion y query, donde $X$ es el tamano del dominio.
	\item Memoria $O(\log X)$ por insercion, $O(X)$ total.
\end{itemize}

\paragraph{Donde tocar para modificar.}
\begin{itemize}\setlength\itemsep{0pt}
	\item \textbf{Upper envelope en vez de lower}: invertir el signo (\texttt{>} en vez de \texttt{<}).
	\item \textbf{Queries de segmento} (no full line): si la recta solo esta activa en un rango, aniadir bandera \texttt{active}.
	\item \textbf{Eliminar rectas}: usar Li Chao segmentado (mas complejo).
	\item \textbf{Pendientes grandes}: si las pendientes tienen magnitud $> 10^9$, aniadir \texttt{\_\_int128} para evitar overflow.
\end{itemize}

\paragraph{Trampas y errores frecuentes.}
\begin{itemize}\setlength\itemsep{0pt}
	\item El parametro \texttt{lo == hi} debe retornar el valor del nodo (caso base).
	\item La comparacion \texttt{newLine.get(mid) < node->line.get(mid)} decide quien es mejor en el medio.
	\item Si las rectas tienen $m$ iguales pero $b$ distintos, se queda con la primera (tie-break).
\end{itemize}

\paragraph{Codigo.}
Ver \texttt{cpp.tex}, seccion \textit{Li Chao Tree (long long)}.

\vspace{0.5em}

\subsubsection{Sweep line template}

\paragraph{Idea intuitiva.}
Patron de diseno para problemas geometricos donde se barre una linea sobre el plano y se mantiene una estructura de datos sobre el otro eje. Eventos (add, query, remove) se ordenan por la coordenada de barrido y se procesan secuencialmente.

\paragraph{Propiedades clave (invariantes).}
\begin{itemize}\setlength\itemsep{0pt}
	\item Eventos: tuplas \texttt{(x, type, id)}.
	\item La estructura auxiliar (set, multiset, segtree) se mantiene sobre el otro eje.
	\item Procesar eventos en orden de $x$ ascendente.
\end{itemize}

\paragraph{Complejidad.}
\begin{itemize}\setlength\itemsep{0pt}
	\item Depende del problema; tipico $O((N + E) \log N)$.
\end{itemize}

\paragraph{Donde tocar para modificar.}
\begin{itemize}\setlength\itemsep{0pt}
	\item \textbf{Aniadir/quitar segmentos}: aniadir eventos add/remove al barrido.
	\item \textbf{Aniadir queries}: aniadir eventos query.
	\item \textbf{Cambiar el eje de barrido}: si barres en $y$ en vez de $x$, aniadir comparador.
\end{itemize}

\paragraph{Trampas y errores frecuentes.}
\begin{itemize}\setlength\itemsep{0pt}
	\item El snippet es solo el \textbf{esqueleto}; los handlers \texttt{add/query/remove} se llenan por problema.
	\item Cuidado con eventos duplicados.
\end{itemize}

\paragraph{Codigo.}
Ver \texttt{cpp.tex}, seccion \textit{Sweep line template}.

\vspace{0.5em}

\subsubsection{Closest pair of points}

\paragraph{Idea intuitiva.}
Encuentra el par de puntos mas cercano en $O(n \log n)$ con \textbf{divide y venceras}. Algoritmo:
\begin{itemize}\setlength\itemsep{0pt}
	\item Sort por $x$.
	\item Dividir en mitades; recursivo.
	\item Combinar: dado el minimo $d$ de las dos mitades, strip = puntos en $[midX - d, midX + d]$; ordenar por $y$; cada punto solo se compara con los siguientes hasta $y$ difiera en $< d$ (maximo 7).
\end{itemize}

\paragraph{Propiedades clave (invariantes).}
\begin{itemize}\setlength\itemsep{0pt}
	\item Requiere ordenar por $y$ despues (en el snippet se hace durante el merge).
	\item Strip tiene a lo sumo $O(n)$ puntos en total, pero la comparacion es $O(1)$ amortizada por punto.
	\item Resultado en \texttt{double}.
\end{itemize}

\paragraph{Complejidad.}
\begin{itemize}\setlength\itemsep{0pt}
	\item $O(n \log n)$.
\end{itemize}

\paragraph{Donde tocar para modificar.}
\begin{itemize}\setlength\itemsep{0pt}
	\item \textbf{Reconstruir el par}: aniadir indices en la recursion.
	\item \textbf{Distancia maxima}: cambiar \texttt{min} por \texttt{max} (mas facil; no requiere D\&C).
	\item \textbf{Puntos duplicados}: aniadir check especial (distancia 0).
\end{itemize}

\paragraph{Trampas y errores frecuentes.}
\begin{itemize}\setlength\itemsep{0pt}
	\item El merge requiere ordenar por $y$; el snippet lo hace in-place con un \texttt{tmp}.
	\item Si se hace en cada recursion, es $O(n \log n)$ extra.
\end{itemize}

\paragraph{Codigo.}
Ver \texttt{cpp.tex}, seccion \textit{Closest pair of points}.

\vspace{0.5em}

\subsubsection{Rotating calipers (polygon diameter)}

\paragraph{Idea intuitiva.}
Encuentra el \textbf{diametro} de un poligono convexo (par de vertices mas distantes) en $O(n)$. Asume poligono en orden counter-clockwise. Algoritmo: dos indices $i$ (sobre las aristas) y $k$ (sobre los vertices opuestos); en cada paso, rotar $k$ mientras el area del paralelogramo aumenta. La distancia $i$-$k$ puede ser maxima en cualquier momento del barrido.

\paragraph{Propiedades clave (invariantes).}
\begin{itemize}\setlength\itemsep{0pt}
	\item Cada indice avanza como maximo $n$ veces total.
	\item El ``diametro'' es la maxima distancia entre dos vertices.
	\item Retorna \textbf{distancia al cuadrado}.
\end{itemize}

\paragraph{Complejidad.}
\begin{itemize}\setlength\itemsep{0pt}
	\item $O(n)$.
\end{itemize}

\paragraph{Donde tocar para modificar.}
\begin{itemize}\setlength\itemsep{0pt}
	\item \textbf{Reconstruir el par de vertices}: aniadir indices.
	\item \textbf{Anchura minima / maxima}: variantes de calipers.
	\item \textbf{Distancia minima entre poligonos convexos}: otra variante.
\end{itemize}

\paragraph{Trampas y errores frecuentes.}
\begin{itemize}\setlength\itemsep{0pt}
	\item Asume poligono \textbf{estrictamente} convexo (sin vertices colineales). Si hay colineales, aniadir perturbacion o filtrar.
\end{itemize}

\paragraph{Codigo.}
Ver \texttt{cpp.tex}, seccion \textit{Rotating calipers (polygon diameter)}.

% ============================================================
\section{Miscelanea}
% ============================================================

\subsection{Utilidades del usuario}

\subsubsection{Hora}

\paragraph{Idea intuitiva.}
Una pequena clase para representar \textbf{tiempos en segundos} y hacer aritmetica basica sobre ellos. Internamente almacena todo en segundos (un solo \texttt{int}), lo cual simplifica comparaciones y diferencias. El constructor acepta $(h, m, s)$ y los convierte a segundos. Hay un metodo \texttt{cut()} que fuerza un minimo (util para ``horas trabajadas'').

\paragraph{Propiedades clave (invariantes).}
\begin{itemize}\setlength\itemsep{0pt}
	\item Representacion canonica: \textbf{solo segundos}. No hay zonas horarias, no hay DST.
	\item Comparadores \texttt{<} y \texttt{>} operan sobre los segundos.
	\item \texttt{operator-} retorna la \textbf{diferencia absoluta} en segundos (siempre no negativa).
	\item \texttt{cut()} aplica un piso (por defecto $8\text{h }30\text{m} = 30600$ segundos).
\end{itemize}

\paragraph{Complejidad.}
\begin{itemize}\setlength\itemsep{0pt}
	\item $O(1)$ por operacion.
\end{itemize}

\paragraph{Donde tocar para modificar.}
\begin{itemize}\setlength\itemsep{0pt}
	\item \textbf{Aniadir suma}: \texttt{operator+} que devuelve un \texttt{Hour} con la suma de segundos.
	\item \textbf{Aniadir output formateado}: aniadir \texttt{show\_hms()} que imprime $H$:$M$:$S$.
	\item \textbf{Cambiar el ``corte''}: modificar el valor magico \texttt{30'600} en \texttt{cut()} o pasarlo como parametro.
	\item \textbf{Tamano maximo}: si los tiempos exceden $\sim 68$ anos, el \texttt{int} desborda; usar \texttt{long long}.
	\item \textbf{Precision sub-segundo}: aniadir \texttt{ms} (milisegundos) y guardar en \texttt{long long} con factor 1000.
\end{itemize}

\paragraph{Trampas y errores frecuentes.}
\begin{itemize}\setlength\itemsep{0pt}
	\item El operador \texttt{-} retorna \textbf{valor absoluto}, no diferencia con signo. Si necesitas diferencia con signo, aniadir un nuevo metodo o sobrecargar.
	\item \texttt{cut()} modifica \texttt{s} \textbf{in-place}; si el llamador no espera esto, aniadir una version \texttt{cut(min)} que retorna un nuevo \texttt{Hour}.
	\item El umbral magico \texttt{30'600} (8h 30m) esta hardcoded; documentar bien que es ``jornada minima''.
\end{itemize}

\paragraph{Codigo.}
Ver \texttt{cpp.tex}, seccion \textit{Hora}.

\end{document}
', aniadir un puntero a ambas raices.
\end{itemize}

\paragraph{Codigo.}
Ver \texttt{cpp.tex}, seccion \textit{Segment Tree persistente}.

\vspace{0.5em}

\subsubsection{Segment Tree iterativo}

\paragraph{Idea intuitiva.}
Representacion \textbf{sin recursion}: los nodos se guardan en un \texttt{vector} de tamano $2n$. Los hijos del nodo $i$ son $2i$ y $2i+1$. Las hojas ocupan las posiciones $n, n+1, \dots, 2n-1$. Update: subir el cambio desde la hoja hasta la raiz. Query: subir dos indices $l$ y $r$ hasta que se ``crucen''. La ventaja principal: localidad de memoria (los nodos estan en posiciones contiguas del vector), que mejora el cache miss rate y la constante practica. La desventaja: implementar lazy propagation es mas complejo.

\paragraph{Propiedades clave (invariantes).}
\begin{itemize}\setlength\itemsep{0pt}
	\item $O(\log n)$ por operacion, \textbf{mas rapido} en practica que el recursivo (mejor uso de cache, no hay overhead de llamadas).
	\item Build simple: copiar las hojas y combinar de abajo hacia arriba.
	\item Las hojas tienen indices en $[n, 2n)$; los nodos internos en $[1, n)$.
\end{itemize}

\paragraph{Codigo clave.}
El bucle de query ascendente:
\begin{verbatim}
long long query(int l, int r) {  // [l, r] inclusivo
    long long res = 0;
    for (l += n, r += n + 1; l < r; l >>= 1, r >>= 1) {
        if (l & 1) res += st[l++];
        if (r & 1) res += st[--r];
    }
    return res;
}
\end{verbatim}

\paragraph{Complejidad.}
\begin{itemize}\setlength\itemsep{0pt}
	\item $O(\log n)$ por update/query, $O(n)$ build.
	\item En la practica, $\sim 1.5x$ mas rapido que el recursivo para $n \ge 10^5$.
\end{itemize}

\paragraph{Donde tocar para modificar.}
\begin{itemize}\setlength\itemsep{0pt}
	\item \textbf{Suma -> max/min/xor}: cambiar la operacion en \texttt{set} y \texttt{query}.
	\item \textbf{Aniadir update en rango con lazy}: implementar version lazy por separado (no trivial en iterativo).
	\item \textbf{Tamano no potencia de 2}: si $n$ no es potencia de 2, aniadir padding.
	\item \textbf{Persistencia}: usar \texttt{vector<array<int, 2>>} y duplicar nodos.
\end{itemize}

\paragraph{Trampas y errores frecuentes.}
\begin{itemize}\setlength\itemsep{0pt}
	\item En el query, \texttt{r += n+1} (no \texttt{r += n}) porque \texttt{query(l, r)} es \textbf{inclusivo} en ambos extremos.
	\item El indice \texttt{l & 1} detecta si \texttt{l} es hijo derecho; aniadir \texttt{res += st[l++]} en ese caso.
	\item El \texttt{--r} en la condicion derecha es para incluir el limite.
	\item Si las operaciones no son idempotentes, cuidado con el orden.
\end{itemize}

\paragraph{Codigo.}
Ver \texttt{cpp.tex}, seccion \textit{Segment Tree iterativo}.

\vspace{0.5em}

\subsection{Fenwick / BIT}

\subsubsection{BIT point update + prefix sum}

\paragraph{Idea intuitiva.}
El \textbf{Fenwick Tree} (BIT) almacena sums parciales de manera que \texttt{sum(i)} (suma de $[1, i]$) y \texttt{add(i, delta)} cuestan $O(\log n)$. La representacion es ingeniosa: el nodo $i$ cubre el rango $[i - \text{lowbit}(i) + 1, i]$. Asi, \texttt{sum(i)} itera \texttt{i -= lowbit(i)}, sumando los nodos que cubren un prefijo creciente, y \texttt{add(i, delta)} itera \texttt{i += lowbit(i)}, propagando el cambio a todos los nodos cuyo rango contiene a $i$. La idea: cada indice $i$ se descompone en una suma de bloques disjuntos de tamano potencias de 2 (su ``lowbit representation'').

\paragraph{Propiedades clave (invariantes).}
\begin{itemize}\setlength\itemsep{0pt}
	\item Indice \textbf{1-based} (no usar $i=0$).
	\item \texttt{lowbit(i) = i \& (-i)}.
	\item Cada nodo cubre exactamente un rango de tamano \texttt{lowbit(i)}.
	\item La operacion debe ser \textbf{invertible} (suma, xor, no min/max).
\end{itemize}

\paragraph{Codigo clave.}
Las dos operaciones basicas:
\begin{verbatim}
void add(int i, long long d) {
    for (; i <= n; i += i & -i) bit[i] += d;
}
long long sum(int i) {
    long long s = 0;
    for (; i > 0; i -= i & -i) s += bit[i];
    return s;
}
\end{verbatim}

\paragraph{Complejidad.}
\begin{itemize}\setlength\itemsep{0pt}
	\item $O(\log n)$ por operacion. Memoria $O(n)$.
	\item Constante muy pequena ($\sim 1$): en la practica, $\sim 3x$ mas rapido que segment tree.
\end{itemize}

\paragraph{Donde tocar para modificar.}
\begin{itemize}\setlength\itemsep{0pt}
	\item \textbf{Suma -> max/min/xor}: cambiar el \texttt{+=} por la operacion (con identidad apropriada). \textbf{No} funciona para \texttt{min} de forma natural (no es invertible).
	\item \textbf{Tamano dinamico}: aniadir un metodo \texttt{resize(newSize)} que preserve valores.
	\item \textbf{Tamano grande}: con $n = 10^7$, mejor comprimir coordenadas.
	\item \textbf{Operador no conmutativo}: usar segment tree en su lugar.
\end{itemize}

\paragraph{Trampas y errores frecuentes.}
\begin{itemize}\setlength\itemsep{0pt}
	\item Indices 0 vs 1: el BIT es \textbf{1-based}. Si el problema da indices 0-based, aniadir \texttt{+1} en cada llamada.
	\item El metodo \texttt{sum} itera \texttt{i -= i \& -i}; aniadir o quitar 1 rompe el invariante.
	\item Si el tipo es \texttt{int}, posibles overflows; usar \texttt{long long}.
	\item Para rango $[l, r]$, el query es \texttt{sum(r) - sum(l-1)} (cuidado con $l = 1$).
\end{itemize}

\paragraph{Codigo.}
Ver \texttt{cpp.tex}, seccion \textit{BIT point update + prefix sum}.

\vspace{0.5em}

\subsubsection{BIT range add + point query}

\paragraph{Idea intuitiva.}
Truco de \textbf{diferencias}: para sumar $d$ a $[l, r]$, aniadir $d$ en $l$ y $-d$ en $r+1$. Asi, el valor en $i$ es la suma de los prefijos hasta $i$, lo que se calcula con un BIT normal. La idea algebraica: si mantenemos un arreglo de diferencias $\Delta$ donde $a_i = \sum_{j \le i} \Delta_j$, entonces aniadir $d$ a $a[l..r]$ equivale a $\Delta_l \mathrel{+}= d$ y $\Delta_{r+1} \mathrel{-}= d$, y reconstruir $a$ es un prefix sum.

\paragraph{Propiedades clave (invariantes).}
\begin{itemize}\setlength\itemsep{0pt}
	\item Internamente es un BIT 5.1 con la convencion de diferencias.
	\item \texttt{rangeAdd(l, r, delta)} cuesta $O(\log n)$ (dos \texttt{add}).
	\item \texttt{pointQuery(i)} cuesta $O(\log n)$.
	\item Si $r+1 > n$, no se aniade el $-d$ (fuera del rango).
\end{itemize}

\paragraph{Codigo clave.}
El truco de diferencias:
\begin{verbatim}
void rangeAdd(int l, int r, long long d) {
    add(l, d);
    if (r + 1 <= n) add(r + 1, -d);
}
long long pointQuery(int i) { return sum(i); }
\end{verbatim}

\paragraph{Complejidad.}
\begin{itemize}\setlength\itemsep{0pt}
	\item $O(\log n)$ por operacion.
	\item Memoria $O(n)$.
\end{itemize}

\paragraph{Donde tocar para modificar.}
\begin{itemize}\setlength\itemsep{0pt}
	\item \textbf{Persistencia}: aniadir versionado (cada update crea una nueva ``copia'' del estado; en BIT no es trivial, mejor usar persistent segtree).
	\item \textbf{Coordenada compression}: si $n$ es grande pero los updates son pocos, comprimir.
	\item \textbf{2D range add + point query}: aniadir una dimension con la misma tecnica.
\end{itemize}

\paragraph{Trampas y errores frecuentes.}
\begin{itemize}\setlength\itemsep{0pt}
	\item Si $r+1 > n$, \textbf{no} aniadir \texttt{-delta} (no hay efecto fuera del rango); el snippet ya lo chequea.
	\item Para reconstruir el array original tras $k$ range adds, hacer $k$ point queries (no hay shortcut).
	\item Confundir ``suma de diferencias'' con ``suma acumulada''; son cosas distintas.
\end{itemize}

\paragraph{Codigo.}
Ver \texttt{cpp.tex}, seccion \textit{BIT range add + point query}.

\vspace{0.5em}

\subsubsection{BIT range add + range sum}

\paragraph{Idea intuitiva.}
Para queries $\sum_{i=1}^{x} a_i$ despues de varios range adds, el truco de diferencias no basta (necesitas la suma acumulada de las diferencias). Solucion: dos BITs, $t_1$ y $t_2$. Tras aniadir $d$ a $[l, r]$: $t_1$ recibe $+d$ en $l$, $-d$ en $r+1$; $t_2$ recibe $+d(l-1)$ en $l$, $-d \cdot r$ en $r+1$. Entonces $\text{prefixSum}(x) = t_1.\text{sum}(x) \cdot x - t_2.\text{sum}(x)$. La razon matematica: el valor en $i$ es $\sum_{j \le i} \Delta_j$, asi que la suma hasta $i$ es $\sum_{k \le i} (i - k + 1) \Delta_k$; expandir da la formula con dos BITs.

\paragraph{Propiedades clave (invariantes).}
\begin{itemize}\setlength\itemsep{0pt}
	\item Verificacion: para $x \ge r$, $\text{prefixSum}(x)$ acumula correctamente todas las contribuciones.
	\item La formula sale de expandir la suma de diferencias.
	\item La identidad clave: $\sum_{i=1}^{x} i = x(x+1)/2$, asi que la suma de diferencias ponderadas requiere BITs adicionales.
\end{itemize}

\paragraph{Codigo clave.}
Range add con range sum (dos BITs):
\begin{verbatim}
void rangeAdd(int l, int r, long long d) {
    t1.add(l, d); t1.add(r + 1, -d);
    t2.add(l, d * (l - 1)); t2.add(r + 1, -d * r);
}
long long prefixSum(int x) {
    return t1.sum(x) * x - t2.sum(x);
}
\end{verbatim}

\paragraph{Complejidad.}
\begin{itemize}\setlength\itemsep{0pt}
	\item $O(\log n)$ por operacion (2 \texttt{add} o 2 \texttt{sum} por cada lado).
	\item Memoria $O(n)$.
\end{itemize}

\paragraph{Donde tocar para modificar.}
\begin{itemize}\setlength\itemsep{0pt}
	\item \textbf{2D range add + range sum}: aniadir otra dimension con la misma tecnica (queda 4 BITs).
	\item \textbf{Cambiar mod}: aniadir modulo en cada operacion.
	\item \textbf{Output del array}: aniadir un metodo que reconstruya $a_i$ point query por cada $i$.
\end{itemize}

\paragraph{Trampas y errores frecuentes.}
\begin{itemize}\setlength\itemsep{0pt}
	\item Overflow en \texttt{delta * (l - 1)} si los valores son grandes; usar \texttt{long long} y \texttt{\% MOD} al final.
	\item Confundir la convencion $0$-based vs $1$-based del BIT; las dos afectan las formulas.
	\item En 2D, el codigo es $\sim 4x$ mas largo y propenso a errores; verificar con casos pequenos.
\end{itemize}

\paragraph{Codigo.}
Ver \texttt{cpp.tex}, seccion \textit{BIT range add + range sum}.

\subsection{RMQ y sparse table}

\subsubsection{Sparse Table}

\paragraph{Idea intuitiva.}
Precomputa respuestas a queries de tamano $2^k$ para $k = 0, 1, \dots, \log n$. Asi, \texttt{st[k][i]} guarda la respuesta sobre $[i, i + 2^k)$. Para query $[l, r]$, se usan dos bloques de tamano $2^{\lfloor \log_2(r-l+1) \rfloor}$: $[l, l + 2^k)$ y $[r - 2^k + 1, r)$, y se combinan. Esto da $O(1)$ por query \textbf{solo si la operacion es idempotente} (min, max, gcd; \textbf{no} suma). La razon: con idempotencia, solapar los dos bloques no afecta el resultado.

\paragraph{Propiedades clave (invariantes).}
\begin{itemize}\setlength\itemsep{0pt}
	\item \texttt{st[0][i] = a[i]}.
	\item \texttt{st[k][i] = op(st[k-1][i], st[k-1][i + 2\^{}(k-1)])}.
	\item \texttt{LOG = floor(log2(n)) + 1}.
	\item La operacion debe ser \textbf{idempotente} ($f(x, x) = x$) y asociativa.
\end{itemize}

\paragraph{Codigo clave.}
Query $O(1)$ en sparse table:
\begin{verbatim}
long long query(int l, int r) {
    int k = __builtin_clz(1) - __builtin_clz(r - l + 1);  // log2
    return min(st[k][l], st[k][r - (1 << k) + 1]);
}
\end{verbatim}

\paragraph{Complejidad.}
\begin{itemize}\setlength\itemsep{0pt}
	\item Build $O(n \log n)$, query $O(1)$.
	\item Memoria $O(n \log n)$.
\end{itemize}

\paragraph{Donde tocar para modificar.}
\begin{itemize}\setlength\itemsep{0pt}
	\item \textbf{Cambiar min por otra operacion idempotente}: cambiar las dos ocurrencias de \texttt{min(...)} por la operacion (max, gcd). Verificar identidad y asociatividad.
	\item \textbf{Suma en rango}: NO sirve con este metodo; usar prefix sums.
	\item \textbf{Sparse table 2D}: aniadir otra dimension (queda $O(n^2 \log^2 n)$ memoria).
	\item \textbf{LCA queries}: combinar con binary lifting en arboles.
\end{itemize}

\paragraph{Trampas y errores frecuentes.}
\begin{itemize}\setlength\itemsep{0pt}
	\item El query toma $[l, r]$ \textbf{inclusivo} (no $[l, r)$).
	\item El \texttt{LOG} debe estar bien calculado; con $n = 1$, \texttt{LOG = 1}.
	\item Usar sparse table para suma es un error; necesitas prefix sums o segment tree.
	\item Para queries $O(1)$ con operaciones no idempotentes, usar Disjoint Sparse Table.
\end{itemize}

\paragraph{Codigo.}
Ver \texttt{cpp.tex}, seccion \textit{Sparse Table}.

\vspace{0.5em}

\subsubsection{Disjoint Sparse Table}

\paragraph{Idea intuitiva.}
Variante que tambien da $O(1)$ por query pero \textbf{no requiere idempotencia} para algunas operaciones (sirve para suma, xor, etc. en linea con la mayoria de operaciones asociativas). La idea: para cada nivel $k$, los rangos ``se parten en mitades''. \texttt{st[k][i]} cubre un rango que \textbf{empieza} en $i$ y se extiende a izquierda o derecha hasta encontrar un ``punto de division''. El truco: para query $[l, r]$, el nivel optimo es el bit mas alto donde $l$ y $r$ difieren; los rangos $[l, ?]$ y $[?, r]$ cubren exactamente el query sin solape.

\paragraph{Propiedades clave (invariantes).}
\begin{itemize}\setlength\itemsep{0pt}
	\item Construccion diferente al Sparse Table clasico: en cada nivel, se extiende desde cada posicion hacia ambos lados hasta encontrar la mitad de nivel.
	\item \texttt{query(l, r)} usa \texttt{k = log2(l \^{} r)} y combina \texttt{st[k][l]} y \texttt{st[k][r]}.
	\item La operacion debe ser \textbf{asociativa} pero no necesariamente idempotente.
\end{itemize}

\paragraph{Codigo clave.}
Query con DST:
\begin{verbatim}
long long query(int l, int r) {
    if (l == r) return st[0][l];
    int k = 31 - __builtin_clz(l ^ r);  // bit mas alto que difiere
    return op(st[k][l], st[k][r]);
}
\end{verbatim}

\paragraph{Complejidad.}
\begin{itemize}\setlength\itemsep{0pt}
	\item Build $O(n \log n)$, query $O(1)$.
	\item Memoria $O(n \log n)$.
\end{itemize}

\paragraph{Donde tocar para modificar.}
\begin{itemize}\setlength\itemsep{0pt}
	\item \textbf{Operacion no asociativa}: no funciona; el Disjoint Sparse Table asume asociatividad.
	\item \textbf{Memory compression}: con $n \le 10^6$ alcanza; con mas, comprimir.
	\item \textbf{Operaciones no conmutativas}: aniadir el cuidado de que el orden es importante.
\end{itemize}

\paragraph{Trampas y errores frecuentes.}
\begin{itemize}\setlength\itemsep{0pt}
	\item La construccion es sutil; verificar con casos pequenos antes de usar.
	\item El bit ``mas alto que difiere'' requiere cuidado con $l = r$; aniadir caso base.
	\item Si la operacion no es asociativa, los dos bloques no se pueden combinar en cualquier orden.
\end{itemize}

\paragraph{Codigo.}
Ver \texttt{cpp.tex}, seccion \textit{Disjoint Sparse Table}.

\subsection{Trees y DP en arbol}

\subsubsection{Binary Lifting / LCA}

\paragraph{Idea intuitiva.}
Precomputa para cada nodo $v$ sus ancestros a distancias $2^k$ (los ``saltos''). Asi, para subir $d$ niveles, se descompone $d$ en binario y se aplican los saltos correspondientes. El \textbf{LCA} (Lowest Common Ancestor) se obtiene subiendo el nodo mas profundo hasta igualar profundidades y luego subiendo ambos hasta que sus padres coincidan. La intuicion: cualquier entero $\le n$ se representa en binario con $\le \log n$ bits, asi que subir $d$ niveles requiere $\le \log n$ saltos.

\paragraph{Propiedades clave (invariantes).}
\begin{itemize}\setlength\itemsep{0pt}
	\item \texttt{up[k][v]} = ancestro de $v$ a $2^k$ niveles.
	\item \texttt{LOG = floor(log2(n)) + 1}.
	\item La recursion de \texttt{dfs} llena \texttt{up[0][v]} (el padre directo) y por transitividad los demas.
	\item El LCA esta definido como el ancestro comun mas profundo de $u$ y $v$.
\end{itemize}

\paragraph{Codigo clave.}
LCA con binary lifting:
\begin{verbatim}
int lca(int u, int v) {
    if (depth[u] < depth[v]) swap(u, v);
    int diff = depth[u] - depth[v];
    for (int k = LOG - 1; k >= 0; --k)
        if (diff & (1 << k)) u = up[k][u];
    if (u == v) return u;
    for (int k = LOG - 1; k >= 0; --k)
        if (up[k][u] != up[k][v])
            u = up[k][u], v = up[k][v];
    return up[0][u];
}
\end{verbatim}

\paragraph{Complejidad.}
\begin{itemize}\setlength\itemsep{0pt}
	\item Build $O(n \log n)$, query LCA $O(\log n)$.
	\item Memoria $O(n \log n)$.
\end{itemize}

\paragraph{Donde tocar para modificar.}
\begin{itemize}\setlength\itemsep{0pt}
	\item \textbf{Distancia en el arbol}: combinar con \texttt{depth[]} para calcular \texttt{depth[u] + depth[v] - 2*depth[lca]}.
	\item \textbf{K-esimo ancestro}: subir $k$ niveles con el mismo truco.
	\item \textbf{Aniadir peso a las aristas}: guardar \texttt{dist[2\^{}k][v]} (distancia acumulada a $2^k$ saltos) ademas de \texttt{up}.
	\item \textbf{Sparse table sobre depth}: alternativa $O(1)$ a LCA (mas complejo).
\end{itemize}

\paragraph{Trampas y errores frecuentes.}
\begin{itemize}\setlength\itemsep{0pt}
	\item El arbol debe ser \textbf{rooted} (el snippet asume eso); si no, hacer un \texttt{dfs} desde cualquier nodo.
	\item Para grafos no conexos, correr \texttt{dfs} desde cada componente.
	\item En la segunda fase del LCA, subir ambos nodos mientras \texttt{up[k][u] != up[k][v]}; el LCA final es \texttt{up[0][u]}.
	\item El padre de la raiz es 0 (o -1); aniadir check si se sale del rango.
\end{itemize}

\paragraph{Codigo.}
Ver \texttt{cpp.tex}, seccion \textit{Binary Lifting / LCA}.

\vspace{0.5em}

\subsubsection{Heavy-Light Decomposition (HLD)}

\paragraph{Idea intuitiva.}
Descomponer un arbol en \textbf{cadenas pesadas} (paths donde cada nodo (salvo la raiz) tiene al hijo ``pesado'' como hijo de tamano maximo de su subarbol). Cada cadena se mapea a un rango contiguo en un arreglo base; asi, queries sobre \textbf{paths} se descomponen en $O(\log n)$ queries sobre rangos contiguos (cada cadena es un rango), que se resuelven con un BIT o segment tree. La intuicion: por cada nivel de recursion, el subarbol del hijo pesado es al menos la mitad, asi que un path cruza $\le 2 \log n$ cadenas.

\paragraph{Propiedades clave (invariantes).}
\begin{itemize}\setlength\itemsep{0pt}
	\item \texttt{heavy[v]} = hijo pesado de $v$ (el de mayor subarbol).
	\item \texttt{head[v]} = cabeza de la cadena que contiene a $v$.
	\item \texttt{pos[v]} = posicion en el arreglo base.
	\item Cada nodo pertenece a exactamente una cadena.
\end{itemize}

\paragraph{Codigo clave.}
El query sobre un path $u \to v$:
\begin{verbatim}
long long pathQuery(int u, int v) {
    long long res = 0;
    while (head[u] != head[v]) {
        if (depth[head[u]] < depth[head[v]]) swap(u, v);
        res += seg.query(pos[head[u]], pos[u]);
        u = parent[head[u]];
    }
    if (depth[u] > depth[v]) swap(u, v);
    res += seg.query(pos[u], pos[v]);
    return res;
}
\end{verbatim}

\paragraph{Complejidad.}
\begin{itemize}\setlength\itemsep{0pt}
	\item Build $O(n)$, query sobre path $O(\log^2 n)$ (con segtree) o $O(\log n)$ (con BIT si solo se hace suma).
	\item Memoria $O(n)$.
\end{itemize}

\paragraph{Donde tocar para modificar.}
\begin{itemize}\setlength\itemsep{0pt}
	\item \textbf{Usar con segment tree}: aniadir un segtree sobre el arreglo base; \texttt{pathSegments} da los rangos a actualizar/query.
	\item \textbf{Cambiar operacion}: ajustar el segtree (suma, max, etc.).
	\item \textbf{Subtree query}: \texttt{subtree[v]} es el rango \texttt{[pos[v], pos[v] + size[v] - 1]} en el arreglo base.
	\item \textbf{Updates en edges}: aniadir un offset para que el edge $(u, v)$ quede en el nodo mas profundo.
\end{itemize}

\paragraph{Trampas y errores frecuentes.}
\begin{itemize}\setlength\itemsep{0pt}
	\item \texttt{pathSegments} requiere un \texttt{segtree} o \texttt{BIT} externo para ser util; por si solo solo da los rangos.
	\item Para edges vs nodos: si la operacion es sobre edges, guardar el valor en el hijo mas profundo.
	\item Olvidar la recursion \texttt{while (head[u] != head[v])} da resultados incorrectos.
	\item En problemas con updates, asegurarse de que el segtree soporte lazy propagation.
\end{itemize}

\paragraph{Codigo.}
Ver \texttt{cpp.tex}, seccion \textit{Heavy-Light Decomposition (HLD)}.

\vspace{0.5em}

\subsubsection{Centroid Decomposition}

\paragraph{Idea intuitiva.}
Descomponer recursivamente el arbol en \textbf{centroides}. Un centroide de un arbol es un nodo tal que cada subtree (al removerlo) tiene tamano $\le n/2$. Existe siempre y se encuentra en $O(n)$. Esto da una estructura de \textbf{arbol de centroides} de profundidad $O(\log n)$, util para problemas ``para cada nodo, encontrar el mas cercano que cumple X''. La razon de la profundidad $O(\log n)$: cada recursion reduce el problema a subarboles de tamano $\le n/2$, asi que la recursion da una cadena geometrica decreciente.

\paragraph{Propiedades clave (invariantes).}
\begin{itemize}\setlength\itemsep{0pt}
	\item En cada nivel, los subarboles se procesan recursivamente.
	\item \texttt{removed[]} marca nodos ya usados como centroide.
	\item \texttt{centroidParent[]} da el padre en el arbol de centroides.
	\item La altura del arbol de centroides es $O(\log n)$.
\end{itemize}

\paragraph{Codigo clave.}
Encontrar el centroide de un subarbol:
\begin{verbatim}
int findCentroid(int v, int p, int sz) {
    for (int u : adj[v])
        if (u != p && !removed[u] && subsize[u] > sz / 2)
            return findCentroid(u, v, sz);
    return v;
}
\end{verbatim}

\paragraph{Complejidad.}
\begin{itemize}\setlength\itemsep{0pt}
	\item Build $O(n \log n)$, queries dependen del problema.
	\item Cada nivel de recursion es $O(n)$ para calcular subsize.
\end{itemize}

\paragraph{Donde tocar para modificar.}
\begin{itemize}\setlength\itemsep{0pt}
	\item \textbf{Query tipica}: en cada nivel, ``expandir'' desde el centroide contando caminos validos y usar una estructura auxiliar (set, multiset).
	\item \textbf{Distancia maxima a un nodo ``malo''}: mantener un multiset de distancias en cada centroide.
	\item \textbf{Cambiar a otro tipo de descomposicion}: si el problema se adapta mejor a small-to-large o DSU on tree, usar esos.
	\item \textbf{Dynamic updates}: combinacion con link-cut trees.
\end{itemize}

\paragraph{Trampas y errores frecuentes.}
\begin{itemize}\setlength\itemsep{0pt}
	\item El snippet actual \textbf{solo construye} la descomposicion; las queries se aniaden segun el problema.
	\item Olvidar resetear \texttt{removed[]} causa recursion infinita.
	\item El calculo de subsize debe ignorar nodos ya ``removed''.
	\item El centroide es \textbf{unico} salvo empates en tamano; cualquiera sirve.
\end{itemize}

\paragraph{Codigo.}
Ver \texttt{cpp.tex}, seccion \textit{Centroid Decomposition}.

\vspace{0.5em}

\subsubsection{Euler Tour + range queries}

\paragraph{Idea intuitiva.}
Un \textbf{Euler Tour} (first-time visit) numera los nodos en el orden en que se visitan. El subarbol de $v$ corresponde a un \textbf{rango contiguo} $[tin[v], tout[v]]$ en este orden. Asi, queries sobre subarboles se resuelven con un BIT o segment tree sobre el arreglo de tiempos. La razon del rango contiguo: cuando entramos a un subarbol, no salimos hasta haber recorrido todos sus descendientes; entonces las posiciones consecutivas corresponden exactamente al subarbol.

\paragraph{Propiedades clave (invariantes).}
\begin{itemize}\setlength\itemsep{0pt}
	\item \texttt{tin[v]} = tiempo de primera visita a $v$.
	\item \texttt{tout[v]} = tiempo al \textbf{terminar} de procesar el subarbol de $v$.
	\item El subarbol de $v$ ocupa las posiciones $[tin[v], tout[v]]$.
	\item Para verificar si $u$ es ancestro de $v$: \texttt{tin[u] <= tin[v] \&\& tout[v] <= tout[u]}.
\end{itemize}

\paragraph{Codigo clave.}
El DFS del Euler Tour:
\begin{verbatim}
void dfs(int v, int p) {
    tin[v] = ++timer;
    for (int u : adj[v])
        if (u != p) dfs(u, v);
    tout[v] = timer;
}
\end{verbatim}

\paragraph{Complejidad.}
\begin{itemize}\setlength\itemsep{0pt}
	\item Build $O(n)$, query sobre subarbol $O(\log n)$ con BIT/segtree.
	\item Memoria $O(n)$.
\end{itemize}

\paragraph{Donde tocar para modificar.}
\begin{itemize}\setlength\itemsep{0pt}
	\item \textbf{Aniadir valor por nodo}: guardar el valor en \texttt{tin[v]} y operar sobre el rango.
	\item \textbf{Path query}: el Euler Tour \textbf{no} sirve para paths; usar HLD o LCA.
	\item \textbf{Multiple queries offline}: procesar en orden de tiempo y mantener una estructura por profundidad.
	\item \textbf{Verificar ancestro}: $O(1)$ con la condicion \texttt{tin}.
\end{itemize}

\paragraph{Trampas y errores frecuentes.}
\begin{itemize}\setlength\itemsep{0pt}
	\item El snippet solo construye el tour; las queries concretas se aniaden.
	\item \texttt{tout[v]} esta en el momento de salir, asi que \texttt{tout[v] - tin[v] + 1} = tamano del subarbol.
	\item Si el arbol no es conexo, aniadir un loop externo sobre cada componente.
	\item Para queries en tiempo (no subtree), usar el orden del tour y offline processing.
\end{itemize}

\paragraph{Codigo.}
Ver \texttt{cpp.tex}, seccion \textit{Euler Tour + range queries}.

\subsection{BST balanceado}

\subsubsection{Treap implicito}

\paragraph{Idea intuitiva.}
Un \textbf{treap} es un arbol binario de busqueda donde cada nodo tiene una \textbf{clave} (el valor de la secuencia) y una \textbf{prioridad} aleatoria. La invariante es: BST por clave, heap por prioridad. Esto da un arbol \textbf{esperado balanceado} sin necesidad de rebalanceo explicito. Un \textbf{treap implicito} no almacena una clave; la posicion de un nodo en la secuencia esta dada por el \textbf{tamano del subarbol izquierdo}, permitiendo split/merge por posicion. La estructura es esencialmente un RBS (randomized BST): la altura esperada es $O(\log n)$ por la propiedad de heap + BST.

\paragraph{Propiedades clave (invariantes).}
\begin{itemize}\setlength\itemsep{0pt}
	\item \texttt{sz(node)} = tamano del subarbol (1 + izquierdo + derecho).
	\item \texttt{split(t, k, l, r)}: parte \texttt{t} en \texttt{l} (primeros $k$ nodos) y \texttt{r} (resto).
	\item \texttt{merge(l, r)}: une dos treaps (todos los de \texttt{l} antes que los de \texttt{r}).
	\item La prioridad aleatoria asegura equilibrio esperado.
\end{itemize}

\paragraph{Codigo clave.}
Split por posicion (la base del treap implicito):
\begin{verbatim}
void split(Node* t, int k, Node*& l, Node*& r) {
    if (!t) { l = r = nullptr; return; }
    push(t);
    if (sz(t->l) >= k) {
        split(t->l, k, l, t->l);
        r = update(t);
    } else {
        split(t->r, k - sz(t->l) - 1, t->r, r);
        l = update(t);
    }
}
\end{verbatim}

\paragraph{Complejidad.}
\begin{itemize}\setlength\itemsep{0pt}
	\item $O(\log n)$ \textbf{esperado} por operacion (split, merge, insert, erase).
	\item Peor caso exponencial (probabilidad astronomicamente pequena).
\end{itemize}

\paragraph{Donde tocar para modificar.}
\begin{itemize}\setlength\itemsep{0pt}
	\item \textbf{Treap con clave (no implicito)}: usar \texttt{key} en vez de posicion; cambiar \texttt{split} por \texttt{splitByKey}.
	\item \textbf{Lazy propagation}: aniadir campos \texttt{lazy} y aplicarlos en \texttt{push} (recursivo).
	\item \textbf{Persistencia}: aniadir versionado clonando nodos.
	\item \textbf{Mejor aleatoriedad}: en lugar de \texttt{rand()}, usar \texttt{rng()} (mejor estado).
	\item \textbf{Operacion de reversa}: aniadir \texttt{lazyReverse} y swap de hijos en push.
\end{itemize}

\paragraph{Trampas y errores frecuentes.}
\begin{itemize}\setlength\itemsep{0pt}
	\item En el snippet actual, \texttt{Node(int v)} usa \texttt{rand()} que en algunos compiladores da maxima prioridad 0; aniadir \texttt{rng()} para mejor distribucion.
	\item Olvidar \texttt{update(t)} tras modificar hijos da tamano incorrecto.
	\item Si lazy propagation no se aplica en push, las queries dan valores viejos.
	\item Memory leak con \texttt{new}; en contests OK, en produccion aniadir destructor.
\end{itemize}

\paragraph{Codigo.}
Ver \texttt{cpp.tex}, seccion \textit{Treap implicito}.

\vspace{0.5em}

\subsubsection{Mo con updates}

\paragraph{Idea intuitiva.}
Extension del algoritmo de Mo para queries que mezclan \textbf{rango} y \textbf{updates}. Las queries se ordenan por $(L/\text{blockSize}, R/\text{blockSize}, T)$ donde $T$ es el ``tiempo'' (numero de updates). Asi, moverse entre queries consecutivas cuesta $O(n^{2/3})$ amortizado. La intuicion: el sort por bloques reduce los movimientos totales, y la tercera coordenada $T$ se optimiza con alternancia similar a $R$.

\paragraph{Propiedades clave (invariantes).}
\begin{itemize}\setlength\itemsep{0pt}
	\item Tres punteros: $L$, $R$, $T$.
	\item \texttt{applyUpdate} y \texttt{revertUpdate} son \textbf{inversos} exactos.
	\item \texttt{addPos} y \texttt{removePos} mantienen el estado.
	\item El tamano de bloque optimo es $n^{2/3}$ (donde $n$ = tamano del array, $Q$ = numero de queries).
\end{itemize}

\paragraph{Codigo clave.}
El sort con tres coordenadas:
\begin{verbatim}
int blockSize = (int)pow(n, 2.0/3.0);
sort(queries.begin(), queries.end(), [&](const Q& a, const Q& b) {
    int blockA = a.l / blockSize, blockB = b.l / blockSize;
    if (blockA != blockB) return blockA < blockB;
    int blockAR = a.r / blockSize, blockBR = b.r / blockSize;
    if (blockAR != blockBR) return blockAR < blockBR;
    return a.t < b.t;
});
\end{verbatim}

\paragraph{Complejidad.}
\begin{itemize}\setlength\itemsep{0pt}
	\item $O((N + Q) \cdot n^{2/3})$ donde $n^{2/3}$ es el tamano de bloque optimo.
	\item En problemas reales, funciona bien hasta $N, Q \sim 10^5$.
\end{itemize}

\paragraph{Donde tocar para modificar.}
\begin{itemize}\setlength\itemsep{0pt}
	\item \textbf{El corazon es el ``add/remove''}: aniadir la logica del problema (ej. contar colores distintos, suma de frecuencias, etc.).
	\item \textbf{Cambiar blockSize}: si la mayoria de updates vs queries cambia, ajustar.
	\item \textbf{Rollback}: aniadir \texttt{save()} / \texttt{restore()} para volver a estados anteriores.
	\item \textbf{Hilbert order}: alternativa con mejor locality.
\end{itemize}

\paragraph{Trampas y errores frecuentes.}
\begin{itemize}\setlength\itemsep{0pt}
	\item El snippet actual es solo el \textbf{esqueleto}; \texttt{addPos}/\texttt{removePos}/\texttt{applyUpdate}/\texttt{revertUpdate} deben ser implementados por problema.
	\item Olvidar \texttt{revertUpdate} es el error mas comun.
	\item Si el numero de updates es muy grande, Mo con updates se vuelve lento; considerar otra estructura.
	\item Mantener consistencia entre \texttt{T} y la posicion de los updates.
\end{itemize}

\paragraph{Codigo.}
Ver \texttt{cpp.tex}, seccion \textit{Mo con updates}.

\vspace{0.5em}

\subsubsection{sqrt-decomposition generico}

\paragraph{Idea intuitiva.}
Dividir el arreglo en \textbf{bloques} de tamano $B \approx \sqrt{n}$. Mantener informacion precomputada por bloque (suma, max, etc.). Queries de rango se resuelven combinando bloques enteros ($O(\sqrt{n})$ en total) y elementos sueltos en los bloques parciales. Updates puntuales actualizan el valor y el resumen del bloque. La intuicion: las queries ``normales'' cruzan $O(n/B)$ bloques; tomar $B = \sqrt{n}$ balancea el costo de bloques enteros vs elementos sueltos.

\paragraph{Propiedades clave (invariantes).}
\begin{itemize}\setlength\itemsep{0pt}
	\item $B = \lfloor \sqrt{n} \rfloor + 1$.
	\item \texttt{bucket[i]} = resumen del bloque $i$-esimo (suma, max, etc.).
	\item \texttt{rangeSum(l, r)}: 3 casos segun $l$ y $r$ esten en el mismo bloque o no.
	\item Cada query cruza $\le 2 + (r - l)/B$ bloques.
\end{itemize}

\paragraph{Codigo clave.}
Estructura basica:
\begin{verbatim}
int B = sqrt(n) + 1;
vector<T> bucket(n / B + 1, IDENTITY);
T query(int l, int r) {
    T res = IDENTITY;
    int bl = l / B, br = r / B;
    if (bl == br) {
        for (int i = l; i <= r; ++i) res = combine(res, a[i]);
    } else {
        for (int i = l; i < (bl+1)*B; ++i) res = combine(res, a[i]);
        for (int i = bl+1; i < br; ++i) res = combine(res, bucket[i]);
        for (int i = br*B; i <= r; ++i) res = combine(res, a[i]);
    }
    return res;
}
\end{verbatim}

\paragraph{Complejidad.}
\begin{itemize}\setlength\itemsep{0pt}
	\item $O(\sqrt{n})$ por operacion.
	\item Memoria $O(n)$.
\end{itemize}

\paragraph{Donde tocar para modificar.}
\begin{itemize}\setlength\itemsep{0pt}
	\item \textbf{Cambiar suma por max/min}: ajustar \texttt{bucket} y la combinacion.
	\item \textbf{Aniadir update en rango}: si la operacion es asociativa, sumar al resumen del bloque + ajustar elementos sueltos.
	\item \textbf{Tamano de bloque optimo}: si las queries son muy largas, $B = n^{2/3}$ (sirve para Mo con updates).
	\item \textbf{Operacion no asociativa}: usar segment tree en su lugar.
\end{itemize}

\paragraph{Trampas y errores frecuentes.}
\begin{itemize}\setlength\itemsep{0pt}
	\item Si $n$ no es multiplo de $B$, el ultimo bloque es mas chico; aniadir check.
	\item La identidad de la operacion debe estar bien definida (0 para suma, $-\infty$ para max).
	\item Si la operacion no es asociativa, los bloques no se pueden combinar; cambiar a segment tree.
\end{itemize}

\paragraph{Codigo.}
Ver \texttt{cpp.tex}, seccion \textit{sqrt-decomposition generico}.

% ============================================================
\section{Strings}
% ============================================================

\subsection{Hashing}

\subsubsection{Rolling Hash (doble modulo)}

\paragraph{Idea intuitiva.}
Un \textbf{rolling hash} asigna a cada prefijo del string un valor numerico tal que el hash de un substring se puede obtener en $O(1)$ a partir de los prefijos. La forma clasica: $h[i] = h[i-1] \cdot A + s[i-1]$ (donde $A$ es una base y todo modulo $M$). Asi, $h[r] - h[l-1] \cdot A^{r-l+1}$ da el hash del substring $[l, r]$ (con $A^{r-l+1}$ precomputado). Usar \textbf{dos modulos} distintos reduce drasticamente la probabilidad de colision. La idea algebraica: el string se interpreta como un polinomio en $A$ con coeficientes los caracteres; el hash es su evaluacion modulo $M$.

\paragraph{Propiedades clave (invariantes).}
\begin{itemize}\setlength\itemsep{0pt}
	\item $h[0] = 0$, $h[i]$ = hash de los primeros $i$ caracteres.
	\item $p[i] = A^i \bmod M$ precomputado.
	\item Comparar dos substrings: comparar sus pares de hashes.
	\item \textbf{Colisiones}: posibles (probabilidad $1/M$ por par). Doble hash: $\approx 1/(M_1 M_2)$.
	\item Si $A$ y $M$ son coprimos, el hash es bijectivo modulo $M$ para strings de longitud fija.
\end{itemize}

\paragraph{Codigo clave.}
La formula del hash de substring:
\begin{verbatim}
pair<long long, long long> get(int l, int r) {
    // h[r] - h[l-1] * A^(r-l+1)
    long long x1 = (h1[r] - h1[l-1] * p1[r-l+1] % M1 + M1) % M1;
    long long x2 = (h2[r] - h2[l-1] * p2[r-l+1] % M2 + M2) % M2;
    return {x1, x2};
}
\end{verbatim}

\paragraph{Complejidad.}
\begin{itemize}\setlength\itemsep{0pt}
	\item Precompute $O(n)$, query hash $O(1)$.
	\item Memoria $O(n)$.
\end{itemize}

\paragraph{Donde tocar para modificar.}
\begin{itemize}\setlength\itemsep{0pt}
	\item \textbf{Caracteres con valor $> 26$}: si los strings incluyen simbolos, aniadir \texttt{if (c >= 'a' \&\& c <= 'z') v = c - 'a'; else v = 26 + ...}.
	\item \textbf{Hash de substrings con updates}: aniadir un Fenwick tree / segment tree sobre los hashes (no es trivial).
	\item \textbf{Cambiar mod}: usar dos primos grandes ($10^9+7$ y $10^9+9$) o un \texttt{u128} para un solo hash sin mod.
	\item \textbf{Base aleatoria}: aniadir random para evitar ataques (en contests no es necesario; en produccion si).
	\item \textbf{Hash polinomial}: $h = \sum s_i \cdot A^i$ en vez de $s_i \cdot A^{i-1}$; algunos problemas requieren este.
\end{itemize}

\paragraph{Trampas y errores frecuentes.}
\begin{itemize}\setlength\itemsep{0pt}
	\item Si los substrings a comparar estan vacios (\texttt{l > r}), su hash es 0; aniadir check.
	\item \textbf{Overflow} en \texttt{h[i-1] * A}: usar \texttt{long long} y aplicar \texttt{\% M} en cada operacion (no al final).
	\item La sustraccion \texttt{h[r] - h[l-1]*A^(r-l+1)} puede dar negativo; aniadir \texttt{+M} antes de \texttt{\% M}.
	\item Con doble hash, verificar \textbf{ambos} componentes; si solo uno coincide, es colision.
\end{itemize}

\paragraph{Codigo.}
Ver \texttt{cpp.tex}, seccion \textit{Rolling Hash (doble modulo)}.

\subsection{Patrones}

\subsubsection{KMP + LPS}

\paragraph{Idea intuitiva.}
El algoritmo \textbf{Knuth-Morris-Pratt} busca todas las ocurrencias de un patron $P$ en un texto $T$ en $O(|T| + |P|)$. La clave es la tabla \textbf{LPS} (Longest Proper Prefix which is Suffix) que, para cada posicion $i$ del patron, dice el largo del prefijo mas largo que tambien es sufijo \textbf{propio}. Durante el match, cuando hay mismatch, en vez de volver al inicio, saltamos al LPS anterior, evitando trabajo redundante. La demostracion de $O(|T| + |P|)$: cada iteracion del bucle externo avanza $i$ o decrece $j$ (con $j$ acotado por $|P|$), asi que en total hay $O(|T| + |P|)$ operaciones.

\paragraph{Propiedades clave (invariantes).}
\begin{itemize}\setlength\itemsep{0pt}
	\item \texttt{lps[i]} = largo del prefijo mas largo que es sufijo propio de \texttt{p[0..i]}.
	\item El match usa dos punteros \texttt{i} (texto) y \texttt{j} (patron); en mismatch, $j = lps[j-1]$.
	\item Cuando $j == |P|$, se encontro una ocurrencia; continuar con $j = lps[j-1]$ para buscar mas.
	\item El LPS se construye con una auto-aplicacion de KMP al patron.
\end{itemize}

\paragraph{Codigo clave.}
Construccion de la tabla LPS:
\begin{verbatim}
for (int i = 1; i < m; ++i) {
    int j = lps[i - 1];
    while (j > 0 && p[i] != p[j]) j = lps[j - 1];
    if (p[i] == p[j]) ++j;
    lps[i] = j;
}
\end{verbatim}

\paragraph{Complejidad.}
\begin{itemize}\setlength\itemsep{0pt}
	\item $O(|T| + |P|)$.
	\item En la practica, rapido y robusto.
\end{itemize}

\paragraph{Donde tocar para modificar.}
\begin{itemize}\setlength\itemsep{0pt}
	\item \textbf{Contar ocurrencias}: el snippet ya cuenta; para devolver posiciones, aniadir un \texttt{vector<int>}.
	\item \textbf{Varios patrones}: usar Aho-Corasick en su lugar.
	\item \textbf{Matcher exacto con wildcard}: aniadir logica para que \texttt{?} matchee cualquier caracter.
	\item \textbf{Tamano pequeno}: si el patron es $\le 5$ caracteres, Rabin-Karp es mas simple.
	\item \textbf{2D KMP}: aniadir una dimension; util en problemas de matrices.
\end{itemize}

\paragraph{Trampas y errores frecuentes.}
\begin{itemize}\setlength\itemsep{0pt}
	\item En el snippet, \texttt{i} y \texttt{j} son 0-based. La condicion \texttt{j == sz\_pattern} para contar funciona.
	\item La construccion de LPS usa el mismo patron como ``texto''.
	\item Si el patron tiene caracteres repetidos, el LPS puede ser no trivial; verificar con casos pequenos.
\end{itemize}

\paragraph{Codigo.}
Ver \texttt{cpp.tex}, seccion \textit{KMP + LPS}.

\vspace{0.5em}

\subsubsection{Z-algorithm}

\paragraph{Idea intuitiva.}
La \textbf{Z-function} $Z[i]$ de un string $S$ es el largo del prefijo mas largo de $S$ que \textbf{tambien} empieza en $S[i]$. Asi, $Z[0] = n$ por convencion. Se calcula en $O(n)$ manteniendo una ventana $[l, r]$ que es el match mas largo encontrado hasta ahora. La consulta $S[l..r] = S[0..r-l]$ da informacion sobre $Z[i]$ cuando $i \le r$. La intuicion: cualquier $i$ dentro de la ventana ``ve'' los caracteres desde una posicion conocida, asi que $Z[i]$ se puede copiar en lugar de recalcular.

\paragraph{Propiedades clave (invariantes).}
\begin{itemize}\setlength\itemsep{0pt}
	\item Construir $Z$ cuesta $O(n)$ con el truco de la ventana.
	\item Para matching de patron $P$ en texto $T$: concatenar $P + \# + T$ y aniadir todos los $i$ con $Z[i] = |P|$.
	\item La ventana $[l, r]$ es siempre la del match mas largo a la derecha.
\end{itemize}

\paragraph{Codigo clave.}
Construccion de Z-function en $O(n)$:
\begin{verbatim}
Z[0] = n;
int l = 0, r = 0;
for (int i = 1; i < n; ++i) {
    if (i <= r) Z[i] = min(r - i + 1, Z[i - l]);
    while (i + Z[i] < n && s[Z[i]] == s[i + Z[i]]) ++Z[i];
    if (i + Z[i] - 1 > r) { l = i; r = i + Z[i] - 1; }
}
\end{verbatim}

\paragraph{Complejidad.}
\begin{itemize}\setlength\itemsep{0pt}
	\item $O(n)$ para construir; matching $O(n + m)$.
	\item Cada caracter se compara a lo sumo 2 veces (extension + reduccion).
\end{itemize}

\paragraph{Donde tocar para modificar.}
\begin{itemize}\setlength\itemsep{0pt}
	\item \textbf{Detectar bordes (border)}: $Z[i] = n - i$ indica que $S[0..n-i-1]$ es borde.
	\item \textbf{Prefijo comun mas largo de dos strings}: aniadir $S_1 + \# + S_2$; $Z[|S_1|+1]$ es el prefijo comun.
	\item \textbf{Comparar prefijos y sufijos}: combinar $S + \# + \text{reverse}(S)$.
	\item \textbf{Matching aproximado}: aniadir tolerancia a mismatches.
\end{itemize}

\paragraph{Trampas y errores frecuentes.}
\begin{itemize}\setlength\itemsep{0pt}
	\item $Z[0]$ siempre vale $n$ (no se calcula).
	\item Si $i > r$, la ventana no aporta; empezar desde 0.
	\item El caracter separador en $P + \# + T$ debe ser uno que no aparezca en $P$ o $T$.
\end{itemize}

\paragraph{Codigo.}
Ver \texttt{cpp.tex}, seccion \textit{Z-algorithm}.

\vspace{0.5em}

\subsubsection{Rabin-Karp (single hash)}

\paragraph{Idea intuitiva.}
Rolling hash aplicado a pattern matching: calcular el hash del patron $P$ y de cada ventana de tamano $|P|$ en $T$. Si los hashes coinciden, verificar caracter por caracter (para descartar colisiones). Esto da $O(n + m)$ \textbf{esperado}. La razon de la eficiencia esperada: cada ventana tiene probabilidad $1/M$ de tener un hash coincidente por casualidad, asi que solo se verifica $O(n/M)$ ventanas en promedio.

\paragraph{Propiedades clave (invariantes).}
\begin{itemize}\setlength\itemsep{0pt}
	\item Si los hashes difieren, los substrings son distintos (sin colision).
	\item Si coinciden, \textbf{verificar} (la verificacion es $O(m)$ pero rara vez se ejecuta).
	\item Peor caso (muchas colisiones): $O(nm)$.
	\item Para $M \ge 10^9$, las colisiones son raras; en la practica nunca pasan.
\end{itemize}

\paragraph{Codigo clave.}
Rolling hash del texto (deslizar ventana):
\begin{verbatim}
long long hash = 0;
for (int i = 0; i < m; ++i) hash = (hash * BASE + text[i]) % MOD;
for (int i = m; i < n; ++i) {
    if (hash == target) check(i - m);
    hash = (hash - text[i - m] * h) * BASE + text[i];  // mod correctamente
    hash = (hash % MOD + MOD) % MOD;
}
\end{verbatim}

\paragraph{Complejidad.}
\begin{itemize}\setlength\itemsep{0pt}
	\item Esperado $O(n + m)$, peor $O(nm)$.
	\item Con doble hash y $M_1, M_2 \ge 10^9$, el peor caso es astronomicamente improbable.
\end{itemize}

\paragraph{Donde tocar para modificar.}
\begin{itemize}\setlength\itemsep{0pt}
	\item \textbf{Doble hash}: aniadir un segundo par \texttt{(BASE2, MOD2)} para reducir colisiones.
	\item \textbf{Varios patrones}: aniadir un set de hashes y buscar match con cualquiera.
	\item \textbf{Tamano de ventana variable}: si los patrones tienen largos distintos, usar el largo especifico por patron.
	\item \textbf{Multi-pattern search}: usar Aho-Corasick en su lugar.
\end{itemize}

\paragraph{Trampas y errores frecuentes.}
\begin{itemize}\setlength\itemsep{0pt}
	\item Overflow en \texttt{(tHash - text[i] * h) * BASE}: aplicar mod correctamente y aniadir \texttt{+MOD} antes de multiplicar.
	\item El factor $h = \text{BASE}^{m-1} \bmod \text{MOD}$ se precomputa.
	\item Si \texttt{MOD} no es primo, las colisiones son mas frecuentes; usar primo.
	\item En el peor caso, Rabin-Karp puede ser lento; preferir KMP si se requiere worst-case.
\end{itemize}

\paragraph{Codigo.}
Ver \texttt{cpp.tex}, seccion \textit{Rabin-Karp (single hash)}.

\vspace{0.5em}

\subsubsection{Aho-Corasick (basic)}

\paragraph{Idea intuitiva.}
Construye un \textbf{trie} con todos los patrones y aniade \textbf{failure links} analogos a los de KMP: el failure link de un nodo $v$ apunta al nodo mas largo que es sufijo del string representado por $v$ y sigue siendo prefijo de algun patron. Asi, al recorrer el texto, si hay mismatch, saltamos al failure link. Esto permite buscar \textbf{todos} los patrones en $O(|T| + \sum |P_i|)$. La idea: unificar varios KMP en una sola estructura.

\paragraph{Propiedades clave (invariantes).}
\begin{itemize}\setlength\itemsep{0pt}
	\item El trie almacena los patrones.
	\item \texttt{link[v]} (failure link) se calcula en BFS en orden de profundidad creciente.
	\item \texttt{out[v]} lista los IDs de patrones que terminan en $v$.
	\item Para queries eficientes, aniadir \textbf{output links}: propagar \texttt{out} de los failure links a los hijos.
\end{itemize}

\paragraph{Codigo clave.}
BFS para construir los failure links:
\begin{verbatim}
queue<int> q;
for (int c = 0; c < SIGMA; ++c)
    if (trie[0].next[c] != -1) {
        q.push(trie[0].next[c]);
        trie[trie[0].next[c]].link = 0;
    } else trie[0].next[c] = 0;  // goto a la raiz
while (!q.empty()) {
    int v = q.front(); q.pop();
    for (int pid : trie[trie[v].link].out)
        trie[v].out.push_back(pid);
    for (int c = 0; c < SIGMA; ++c) {
        if (trie[v].next[c] != -1) {
            trie[trie[v].next[c]].link = trie[trie[v].link].next[c];
            q.push(trie[v].next[c]);
        } else trie[v].next[c] = trie[trie[v].link].next[c];
    }
}
\end{verbatim}

\paragraph{Complejidad.}
\begin{itemize}\setlength\itemsep{0pt}
	\item Build $O(\sum |P_i| \cdot \Sigma)$, query $O(|T| + \text{ocurrencias})$.
\end{itemize}

\paragraph{Donde tocar para modificar.}
\begin{itemize}\setlength\itemsep{0pt}
	\item \textbf{Output links en build}: ya hecho en el snippet (\texttt{for(int pid : t[t[u].link].out) t[u].out.push\_back(pid);}).
	\item \textbf{Transiciones automaticas}: aniadir el goto \texttt{t[v].next[k] = t[t[v].link].next[k]} cuando el enlace directo es \texttt{-1}; asi el bucle de matching es un solo paso por caracter.
	\item \textbf{Aplicar a problemas tipo DP}: contar formas de construir strings evitando los patrones (ver el siguiente modulo).
	\item \textbf{2D Aho-Corasick}: aniadir una dimension para busqueda en matrices.
\end{itemize}

\paragraph{Trampas y errores frecuentes.}
\begin{itemize}\setlength\itemsep{0pt}
	\item La BFS debe empezar con la raiz en la cola y los hijos de la raiz con \texttt{link = 0}.
	\item Olvidar propagar \texttt{out} da ocurrencias perdidas.
	\item Si no se aniaden las transiciones automaticas (goto), el matching es mas lento.
	\item Con muchos patrones, el output puede ser enorme; aniadir contador en lugar de lista.
\end{itemize}

\paragraph{Codigo.}
Ver \texttt{cpp.tex}, seccion \textit{Aho-Corasick (basic)}.

\vspace{0.5em}

\subsubsection{Aho-Corasick DP for forbidden patterns}

\paragraph{Idea intuitiva.}
Contar strings de longitud $L$ que \textbf{evitan} un conjunto de patrones. Se modela como una DP sobre el \textbf{estado del automaton} de Aho-Corasick: $dp[v]$ = numero de formas de llegar al estado $v$ despues de cierto numero de caracteres. Transicion: por cada caracter $c$ en el alfabeto, ir al estado \texttt{next[v][c]}, pero \textbf{solo} si ese estado no tiene un patron (de lo contrario, el string contendria el patron prohibido). La idea: cada estado del automaton representa un ``prefijo conocido'' y las transiciones simulan ``agregar un caracter''.

\paragraph{Propiedades clave (invariantes).}
\begin{itemize}\setlength\itemsep{0pt}
	\item $dp$ se mantiene por nivel (longitud del string construido).
	\item Transicion: $ndp[u] = \sum_{c} dp[v]$ donde $u = next[v][c]$ y $u$ no es ``terminal'' (no tiene patron que termina ahi).
	\item Respuesta: suma de $dp[v]$ sobre todos los estados $v$ no terminales al final.
	\item Si todos los caracteres son permitidos, $|\Sigma|$ entra como factor.
\end{itemize}

\paragraph{Codigo clave.}
DP sobre el automaton:
\begin{verbatim}
vector<long long> dp(numStates, 0), ndp(numStates);
dp[0] = 1;  // empezar en la raiz
for (int step = 0; step < L; ++step) {
    fill(ndp.begin(), ndp.end(), 0);
    for (int v = 0; v < numStates; ++v) if (dp[v] && !isTerminal[v])
        for (int c = 0; c < SIGMA; ++c) {
            int u = next[v][c];
            if (!isTerminal[u]) ndp[u] = (ndp[u] + dp[v]) % MOD;
        }
    swap(dp, ndp);
}
long long ans = 0;
for (int v = 0; v < numStates; ++v) if (!isTerminal[v])
    ans = (ans + dp[v]) % MOD;
\end{verbatim}

\paragraph{Complejidad.}
\begin{itemize}\setlength\itemsep{0pt}
	\item $O(L \cdot |S| \cdot |\Sigma|)$ donde $|S|$ es el numero de estados y $|\Sigma|$ el alfabeto.
	\item Para $|\Sigma| = 2$ y $|S| \le 10^4$: $\sim 2 \cdot 10^{10}$ para $L = 10^6$ (puede ser lento).
\end{itemize}

\paragraph{Donde tocar para modificar.}
\begin{itemize}\setlength\itemsep{0pt}
	\item \textbf{Alfabeto mas pequeno}: si el problema es solo $\{0, 1\}$, $|\Sigma| = 2$.
	\item \textbf{Permitir ocurrencias parciales}: cambiar la condicion a ``no estado terminal'' o relajar segun el problema.
	\item \textbf{Recuperar el string}: aniadir backtracking si el problema lo pide.
	\item \textbf{Matriz de transicion}: para queries rapidas, precomputar $M^c$ con exponentiation.
\end{itemize}

\paragraph{Trampas y errores frecuentes.}
\begin{itemize}\setlength\itemsep{0pt}
	\item El snippet asume Aho-Corasick ya construido con transiciones automaticas. Sin eso, aniadir el goto en el bucle.
	\item Si $L$ es muy grande, la DP lineal es lenta; usar matrix exponentiation.
	\item Estados terminales deben excluirse de \textbf{ambos} lados (estado actual y siguiente).
\end{itemize}

\paragraph{Codigo.}
Ver \texttt{cpp.tex}, seccion \textit{Aho-Corasick DP for forbidden patterns}.

\subsection{Estructuras avanzadas}

\subsubsection{Suffix Array}

\paragraph{Idea intuitiva.}
Un \textbf{suffix array} $SA$ es un arreglo que contiene las posiciones de inicio de los sufijos del string $S$, ordenados lexicograficamente. Se construye en $O(n \log n)$ con doubling: ordenar por el primer caracter, luego por los primeros 2, luego 4, etc. Para hacer cada paso lineal, se usa \textbf{counting sort} sobre los ranks. La idea: cada nivel dobla la cantidad de caracteres considerados, asi que en $\log n$ niveles los sufijos estan completamente ordenados.

\paragraph{Propiedades clave (invariantes).}
\begin{itemize}\setlength\itemsep{0pt}
	\item \texttt{p[i]} = posicion de inicio del $i$-esimo sufijo (ordenado).
	\item \texttt{c[i]} = rank del sufijo que empieza en $i$.
	\item Cada iteracion duplica la longitud de la ``ventana'' considerada.
	\item El algoritmo termina cuando $2^k \ge n$.
	\item La permutacion \texttt{p} es estable entre iteraciones.
\end{itemize}

\paragraph{Codigo clave.}
El paso de doubling (clasificar por los primeros $2^k$ caracteres):
\begin{verbatim}
for (int k = 0; (1 << k) < n; ++k) {
    for (int i = 0; i < n; ++i) pn[i] = (p[i] - (1 << k) + n) % n;
    // counting sort por c[pn[i]] segundo, c[i] primero
    // ...
    cn[pn[0]] = 0;
    for (int i = 1; i < n; ++i)
        cn[pn[i]] = cn[pn[i-1]] + (pair != pair_prev);
    c = cn;
}
\end{verbatim}

\paragraph{Complejidad.}
\begin{itemize}\setlength\itemsep{0pt}
	\item $O(n \log n)$ con counting sort.
	\item Constante pequena; en la practica eficiente.
\end{itemize}

\paragraph{Donde tocar para modificar.}
\begin{itemize}\setlength\itemsep{0pt}
	\item \textbf{Suffix array en $O(n)$}: usar SA-IS (mas complejo, no incluido).
	\item \textbf{Buscar patron}: con SA + LCP es $O(|P| \log n)$ por busqueda binaria.
	\item \textbf{Suffix array de varios strings}: aniadir centinelas (\texttt{\$}) entre ellos.
	\item \textbf{Aniadir longitud}: aniadir \texttt{n - p[i]} como clave secundaria (sufijos mas cortos primero).
\end{itemize}

\paragraph{Trampas y errores frecuentes.}
\begin{itemize}\setlength\itemsep{0pt}
	\item La ``ventana'' $2^k$ se hace circular: \texttt{p[i] = (p[i] - (1<<k) + n) \% n}.
	\item Para contar bien los ranks duplicados, comparar pares \texttt{(c[i], c[(i+2\^{}k) \% n])}.
	\item El counting sort necesita un array de acumulacion de tamano $\le n$.
	\item Si los caracteres son negativos o fuera del alfabeto esperado, normalizar antes.
\end{itemize}

\paragraph{Codigo.}
Ver \texttt{cpp.tex}, seccion \textit{Suffix Array}.

\vspace{0.5em}

\subsubsection{LCP Array (Kasai)}

\paragraph{Idea intuitiva.}
El \textbf{LCP array} (Longest Common Prefix) guarda, para cada par de sufijos adyacentes en el SA, el largo de su prefijo comun. El algoritmo de \textbf{Kasai} lo calcula en $O(n)$ iterando los sufijos \textbf{en orden del string original} (no del SA), y reutiliza el LCP previo. La demostracion: el sufijo en $i+1$ se obtiene del sufijo en $i$ ``avanzando un caracter''; los LCPs decrecen como maximo 1 entre iteraciones.

\paragraph{Propiedades clave (invariantes).}
\begin{itemize}\setlength\itemsep{0pt}
	\item \texttt{lcp[i]} = LCP entre los sufijos en \texttt{sa[i]} y \texttt{sa[i-1]}.
	\item Se necesita \texttt{rank\_[i]} (la inversa del SA) para iterar.
	\item El truco de Kasai: si \texttt{lcp > 0} para el sufijo anterior, el actual parte de al menos \texttt{lcp - 1}.
	\item La suma de todos los LCPs da informacion sobre repeticiones.
\end{itemize}

\paragraph{Codigo clave.}
Algoritmo de Kasai:
\begin{verbatim}
vector<int> rank(n);
for (int i = 0; i < n; ++i) rank[sa[i]] = i;
int k = 0;
vector<int> lcp(n - 1);
for (int i = 0; i < n; ++i) {
    if (rank[i] == n - 1) { k = 0; continue; }
    int j = sa[rank[i] + 1];
    while (i + k < n && j + k < n && s[i+k] == s[j+k]) ++k;
    lcp[rank[i]] = k;
    if (k) --k;
}
\end{verbatim}

\paragraph{Complejidad.}
\begin{itemize}\setlength\itemsep{0pt}
	\item $O(n)$.
	\item Cada comparacion reduce $k$ o avanza linealmente; amortizado $O(1)$.
\end{itemize}

\paragraph{Donde tocar para modificar.}
\begin{itemize}\setlength\itemsep{0pt}
	\item \textbf{RMQ sobre LCP}: para queries ``LCP de dos sufijos'', encontrar sus ranks y hacer RMQ en \texttt{lcp}.
	\item \textbf{Contar substrings distintas}: $\sum (n - sa[i] - lcp[i])$.
	\item \textbf{Longest common substring}: maximo en \texttt{lcp}.
	\item \textbf{String matching}: busqueda binaria sobre SA + LCP verification.
\end{itemize}

\paragraph{Trampas y errores frecuentes.}
\begin{itemize}\setlength\itemsep{0pt}
	\item Si dos sufijos son iguales (raro, solo si el string tiene period $< n$), \texttt{lcp} puede ser \texttt{n - sa[i]}.
	\item En la iteracion, aniadir el chequeo \texttt{i + k < n} para evitar overflow.
	\item El \texttt{--k} en la salida es crucial para la complejidad amortizada.
\end{itemize}

\paragraph{Codigo.}
Ver \texttt{cpp.tex}, seccion \textit{LCP Array (Kasai)}.

\vspace{0.5em}

\subsubsection{Suffix Automaton}

\paragraph{Idea intuitiva.}
Un \textbf{suffix automaton} es un DFA minimo que reconoce todos los substrings del string original. Tiene la propiedad magica de tener $\le 2n - 1$ estados. Cada estado representa una clase de equivalencia de substrings (los que tienen el mismo conjunto de posiciones finales). El \textbf{link} de un estado apunta al estado del sufijo propio mas largo. La demostracion de $\le 2n - 1$: cada nuevo caracter anade a lo sumo 2 estados (uno nuevo y un clon), asi que en $n$ pasos hay $\le 2n - 1$ estados.

\paragraph{Propiedades clave (invariantes).}
\begin{itemize}\setlength\itemsep{0pt}
	\item \texttt{len[v]}: largo maximo de los substrings de la clase.
	\item \texttt{link[v]}: padre en el arbol de sufijos (suffix link).
	\item \texttt{next[v][c]}: transicion por caracter.
	\item \texttt{occ[v]}: numero de posiciones finales (cuantas veces aparece el patron).
	\item Cada estado tiene $\texttt{len[v]} - \texttt{len[link[v]]}$ substrings ``propios''.
\end{itemize}

\paragraph{Codigo clave.}
El corazon de la extension SAM:
\begin{verbatim}
void extend(int c) {
    int cur = size++;
    st[cur].len = st[last].len + 1;
    st[cur].occ = 1;
    int p = last;
    while (p != -1 && !st[p].next.count(c)) {
        st[p].next[c] = cur;
        p = st[p].link;
    }
    if (p == -1) st[cur].link = 0;
    else {
        int q = st[p].next[c];
        if (st[p].len + 1 == st[q].len) st[cur].link = q;
        else {
            int clone = size++;
            st[clone] = st[q];
            st[clone].len = st[p].len + 1;
            st[clone].occ = 0;  // clon NO es nueva posicion final
            while (p != -1 && st[p].next[c] == q) {
                st[p].next[c] = clone;
                p = st[p].link;
            }
            st[q].link = st[cur].link = clone;
        }
    }
    last = cur;
}
\end{verbatim}

\paragraph{Complejidad.}
\begin{itemize}\setlength\itemsep{0pt}
	\item Build $O(n)$ con la operacion de \textbf{clonar} estados.
	\item Memoria $O(n)$ (con map) o $O(n \cdot \Sigma)$ (con array).
\end{itemize}

\paragraph{Donde tocar para modificar.}
\begin{itemize}\setlength\itemsep{0pt}
	\item \textbf{Contar substrings distintas}: $\sum_{v \ne 0} (len[v] - len[link[v]])$.
	\item \textbf{Longest common substring}: navegar el SAM con el segundo string, manteniendo el match actual.
	\item \textbf{Substring count}: despues del build, ordenar estados por \texttt{len} descendente y propagar \texttt{occ}.
	\item \textbf{Endpos queries}: el SAM es el DFA canonico de los substrings; queries ``cuantas veces aparece $P

\vspace{0.5em}

\subsubsection{Lyndon factorization (Duval)}

\paragraph{Idea intuitiva.}
Un \textbf{string de Lyndon} es un string que es estrictamente menor que todos sus rotaciones. Toda cadena tiene una \textbf{factorizacion unica} de Lyndon (algoritmo de \textbf{Duval}): una secuencia de strings de Lyndon $L_1, L_2, \dots, L_k$ tales que $L_1 \ge L_2 \ge \dots \ge L_k$ y la concatenacion es el string original. Esto da $O(n)$ para encontrar los cortes. La demostracion de correctitud: cada factor es minimal en su clase de rotacion, asi que la concatenacion respeta el orden lexicografico.

\paragraph{Propiedades clave (invariantes).}
\begin{itemize}\setlength\itemsep{0pt}
	\item Cada factor de Lyndon es de la forma $uv$ donde $u$ se repite y $v$ es prefijo de $u$.
	\item La longitud del factor esta dada por la ``longitud del periodo''.
	\item Aplicacion: el ``Lyndon word'' es el menor en orden lexicografico entre todas las rotaciones.
	\item La factorizacion es \textbf{unica}.
\end{itemize}

\paragraph{Codigo clave.}
El algoritmo de Duval (iterativo):
\begin{verbatim}
int n = s.size();
int i = 0;
while (i < n) {
    int j = i + 1, k = i;
    while (j < n && s[k] <= s[j]) {
        if (s[k] < s[j]) k = i;
        else ++k;
        ++j;
    }
    int len = j - k;
    while (i < j) { cout << "factor: " << s.substr(i, len) << "\n"; i += len; }
}
\end{verbatim}

\paragraph{Complejidad.}
\begin{itemize}\setlength\itemsep{0pt}
	\item $O(n)$.
	\item Cada comparacion es $O(1)$; el puntero $k$ retrocede a lo sumo $n$ veces en total.
\end{itemize}

\paragraph{Donde tocar para modificar.}
\begin{itemize}\setlength\itemsep{0pt}
	\item \textbf{Encontrar la menor rotacion}: el menor sufijo lexicografico de \texttt{s + s} es la menor rotacion; la posicion se obtiene con un solo pase de Duval.
	\item \textbf{Algoritmo de Booth}: alternativa para rotacion minima, tambien $O(n)$.
	\item \textbf{Validar}: si necesitas verificar que un string es de Lyndon, compararlo con sus rotaciones.
	\item \textbf{Encontrar todos los factores}: guardar la posicion y longitud de cada uno.
\end{itemize}

\paragraph{Trampas y errores frecuentes.}
\begin{itemize}\setlength\itemsep{0pt}
	\item La condicion \texttt{s[j] <= s[k]} (no \texttt{<}) en el while de Duval es crucial; cambiar a \texttt{<} da comportamiento distinto.
	\item El caso $s[k] < s[j]$ resetea $k$ a $i$ (inicio del factor actual).
	\item El bucle externo debe avanzar $i$ por la longitud del factor, no por 1.
\end{itemize}

\paragraph{Codigo.}
Ver \texttt{cpp.tex}, seccion \textit{Lyndon factorization (Duval)}.

\subsection{Palindromos}

\subsubsection{Manacher}

\paragraph{Idea intuitiva.}
Calcula en $O(n)$ los \textbf{radios} de todos los palindromos centrados en cada posicion (impares y pares). La idea: mantener una ventana $[l, r]$ que es el palindromo centrado en un punto anterior mas largo encontrado. Si el nuevo centro $i$ cae dentro de $[l, r]$, entonces el palindromo centrado en $i$ tiene al menos el radio del palindromo ``espejo'' $l + r - i$; sino, empezar desde 1 (impar) o 0 (par). Despues, expandir lo que falte. La demostracion de $O(n)$: cada caracter se compara a lo sumo 2 veces (extension + limite por ventana).

\paragraph{Propiedades clave (invariantes).}
\begin{itemize}\setlength\itemsep{0pt}
	\item Conveccion del snippet: \texttt{odd[i]} = $k$ tal que el palindromo impar mas largo centrado en $i$ tiene longitud $2k - 1$. \texttt{even[i]} = $k$ para par centrado entre $i-1$ e $i$, longitud $2k$.
	\item $k$ cuenta tambien cuantos palindromos hay centrados ahi.
	\item Las queries \texttt{isPalindrome(l, r)} son $O(1)$.
	\item La ventana $[l, r]$ siempre contiene el palindromo mas largo encontrado hasta ahora.
\end{itemize}

\paragraph{Codigo clave.}
Manacher para palindromos impares:
\begin{verbatim}
int l = 0, r = -1;
for (int i = 0; i < n; ++i) {
    int k = (i > r) ? 1 : min(odd[l + r - i], r - i + 1);
    while (i - k >= 0 && i + k < n && s[i-k] == s[i+k]) ++k;
    odd[i] = k;
    if (i + k - 1 > r) { l = i - k + 1; r = i + k - 1; }
}
\end{verbatim}

\paragraph{Complejidad.}
\begin{itemize}\setlength\itemsep{0pt}
	\item Construccion $O(n)$, queries $O(1)$.
	\item Cada extension se cuenta una sola vez.
\end{itemize}

\paragraph{Donde tocar para modificar.}
\begin{itemize}\setlength\itemsep{0pt}
	\item \textbf{Solo contar palindromos}: ya implementado en \texttt{countPalindromes}.
	\item \textbf{Longest palindromic substring}: ya en \texttt{longestPalindrome}.
	\item \textbf{Modificar la conveccion de radios}: si queres ``radio = mitad de la longitud'' en vez de ``k palindromos'', ajustar el offset.
	\item \textbf{Queries con k no precomputado}: aniadir helpers especificos.
	\item \textbf{Modificar el caracter de relleno}: para multiples separadores, aniadir un caracter unico.
\end{itemize}

\paragraph{Trampas y errores frecuentes.}
\begin{itemize}\setlength\itemsep{0pt}
	\item La conveccion del snippet es \textbf{distinta} de la ``clasica de radio'' (donde radio = $(L-1)/2$); tener cuidado al copiar a otro codigo.
	\item Si los palindromos se solapan, la condicion \texttt{i + k - 1 > r} actualiza la ventana.
	\item En palindromos pares, el centro esta ``entre'' dos indices; aniadir offset en queries.
\end{itemize}

\paragraph{Codigo.}
Ver \texttt{cpp.tex}, seccion \textit{Manacher}.

% ============================================================
\section{Grafos}
% ============================================================

\subsection{Caminos minimos}

\subsubsection{Dijkstra}

\paragraph{Idea intuitiva.}
Encuentra el shortest path desde un origen $s$ a todos los vertices en $O((V + E) \log V)$ asumiendo \textbf{pesos no negativos}. La idea: usar una cola de prioridad (min-heap) ordenada por distancia tentativa. En cada paso, extraer el nodo con menor distancia; si ya fue procesado, saltar. Si al actualizar un vecino la distancia mejora, insertar/actualizar en la cola. La demostracion: cuando un nodo se extrae por primera vez, no hay forma de llegar a el con menor distancia (cualquier camino alternativo deberia pasar por nodos no procesados, con distancia $\ge$ la actual).

\paragraph{Propiedades clave (invariantes).}
\begin{itemize}\setlength\itemsep{0pt}
	\item Asume pesos $\ge 0$; con pesos negativos usar Bellman-Ford.
	\item Greedy correctness: la primera vez que se extrae un nodo, su distancia es la definitiva.
	\item Si se quiere reconstruir el camino, guardar el \texttt{parent[]}.
	\item Cada nodo se procesa a lo sumo $V$ veces (mas relajarions, pero solo $V$ definitivas).
\end{itemize}

\paragraph{Codigo clave.}
El bucle principal de Dijkstra:
\begin{verbatim}
priority_queue<pair<long long, int>, vector<...>, greater<...>> pq;
pq.push({0, source});
while (!pq.empty()) {
    auto [d, v] = pq.top(); pq.pop();
    if (d > dist[v]) continue;  // entrada vieja
    for (auto [w, u] : adj[v])
        if (dist[v] + w < dist[u]) {
            dist[u] = dist[v] + w;
            pq.push({dist[u], u});
        }
}
\end{verbatim}

\paragraph{Complejidad.}
\begin{itemize}\setlength\itemsep{0pt}
	\item $O((V + E) \log V)$ con heap binario, $O(V + E)$ con heap de Fibonacci.
	\item En la practica, $\sim 1.5 \cdot 10^{-7}$ segundos por operacion para $V = 10^5$.
\end{itemize}

\paragraph{Donde tocar para modificar.}
\begin{itemize}\setlength\itemsep{0pt}
	\item \textbf{Camino mas corto entre todos los pares}: si se necesita desde todos los origenes, usar Floyd o ejecutar Dijkstra $V$ veces.
	\item \textbf{Shortest path con aristas negativas limitadas}: Bellman-Ford.
	\item \textbf{Reconstruir camino}: aniadir \texttt{parent[v]} actualizado junto a \texttt{dist[v]}.
	\item \textbf{Dijkstra con potentials (Johnson)}: para aristas negativas si se puede reescalar.
	\item \textbf{Multi-source}: aniadir varios nodos fuentes iniciales al heap.
\end{itemize}

\paragraph{Trampas y errores frecuentes.}
\begin{itemize}\setlength\itemsep{0pt}
	\item Distancias no inicializadas a \texttt{INF}; el snippet usa \texttt{1e18}.
	\item El \texttt{if(dist[adjN] != 1e18)} para borrar de la cola es importante si usas \texttt{set} (no asi con \texttt{priority\_queue}).
	\item Con pesos grandes, \texttt{dist[v] + w} puede overflow; usar \texttt{LLONG\_MAX/2} como INF.
	\item No usar con pesos negativos: el algoritmo puede dar respuestas incorrectas.
\end{itemize}

\paragraph{Codigo.}
Ver \texttt{cpp.tex}, seccion \textit{Dijkstra}.

\vspace{0.5em}

\subsubsection{Bellman-Ford}

\paragraph{Idea intuitiva.}
Encuentra shortest path desde $s$ incluso con \textbf{pesos negativos} en $O(VE)$. La idea: relajar todas las aristas $V - 1$ veces. Cada relajacion propaga la mejor distancia a ``un paso mas''. Despues de $V - 1$ pasadas, si alguna arista aun se puede relajar, hay un \textbf{ciclo negativo} alcanzable. La demostracion: el camino mas corto tiene a lo sumo $V - 1$ aristas (sin ciclos); tras $V - 1$ relajaciones todas las distancias estan en su valor optimo.

\paragraph{Propiedades clave (invariantes).}
\begin{itemize}\setlength\itemsep{0pt}
	\item Funciona con pesos negativos (sin ciclos negativos).
	\item Tras $V - 1$ relajaciones, las distancias son definitivas (si no hay ciclo negativo).
	\item Una relajacion adicional detecta ciclos negativos: si \texttt{dist[u] + w < dist[v]}, hay ciclo negativo alcanzable desde $s$.
	\item Cada relajacion solo mejora distancias, asi que termina (no oscila).
\end{itemize}

\paragraph{Codigo clave.}
El doble bucle de relajacion:
\begin{verbatim}
for (int i = 0; i < n - 1; ++i)
    for (auto [u, v, w] : edges)
        if (dist[u] + w < dist[v])
            dist[v] = dist[u] + w;
// Detectar ciclo negativo
for (auto [u, v, w] : edges)
    if (dist[u] + w < dist[v]) { /* ciclo negativo */ }
\end{verbatim}

\paragraph{Complejidad.}
\begin{itemize}\setlength\itemsep{0pt}
	\item $O(VE)$.
	\item SPFA (con cola) en promedio es mucho mas rapido, peor caso igual.
\end{itemize}

\paragraph{Donde tocar para modificar.}
\begin{itemize}\setlength\itemsep{0pt}
	\item \textbf{SPFA}: optimizacion que relaja solo nodos cuya distancia cambio; en el peor caso sigue siendo $O(VE)$.
	\item \textbf{Detectar ciclo negativo global}: ejecutar Bellman-Ford desde un super-origen conectado a todos.
	\item \textbf{Camino mas corto con K aristas exactas}: $K$ iteraciones de Bellman-Ford.
	\item \textbf{Johnson's}: reescalar pesos con potentials para usar Dijkstra.
\end{itemize}

\paragraph{Trampas y errores frecuentes.}
\begin{itemize}\setlength\itemsep{0pt}
	\item Overflow si los pesos son $\ge 2^{52}$; usar \texttt{long long} y \texttt{\_\_int128} si hace falta.
	\item Si el grafo es denso ($E \approx V^2$), $O(VE)$ es prohibitivo.
	\item El chequeo de ciclo negativo debe hacerse tras las $V - 1$ iteraciones, no durante.
	\item Distancias no inicializadas a \texttt{INF}; el primer nodo (source) tiene distancia 0.
\end{itemize}

\paragraph{Codigo.}
Ver \texttt{cpp.tex}, seccion \textit{Bellman-Ford}.

\vspace{0.5em}

\subsubsection{Floyd-Warshall}

\paragraph{Idea intuitiva.}
Shortest path entre \textbf{todos los pares} de vertices en $O(V^3)$. DP: $d[i][j]$ = shortest path usando solo nodos intermedios en $\{0, \dots, k\}$. Transicion: $d_{k+1}[i][j] = \min(d_k[i][j], d_k[i][k+1] + d_k[k+1][j])$. La demostracion: tras $k$ iteraciones, $d[i][j]$ es el shortest path con intermedios en $\{0, \dots, k-1\}$; aniadir el nodo $k$ da la recurrencia clasica.

\paragraph{Propiedades clave (invariantes).}
\begin{itemize}\setlength\itemsep{0pt}
	\item Detecta ciclos negativos: tras el algoritmo, $d[i][i] < 0$ indica ciclo negativo alcanzable desde $i$.
	\item Funciona con pesos negativos (sin ciclos negativos).
	\item Memoria $O(V^2)$; con $V \le 400$ es OK, con mas conviene Dijkstra $V$ veces.
	\item Es la unica opcion cuando se quieren shortest paths entre todos los pares y el grafo es denso.
\end{itemize}

\paragraph{Codigo clave.}
El triple bucle clasico:
\begin{verbatim}
for (int k = 0; k < n; ++k)
    for (int i = 0; i < n; ++i)
        for (int j = 0; j < n; ++j)
            d[i][j] = min(d[i][j], d[i][k] + d[k][j]);
// detectar ciclo negativo
for (int i = 0; i < n; ++i) if (d[i][i] < 0) tiene_ciclo_negativo = true;
\end{verbatim}

\paragraph{Complejidad.}
\begin{itemize}\setlength\itemsep{0pt}
	\item $O(V^3)$ tiempo, $O(V^2)$ memoria.
	\item Para $V \le 400$: $\sim 6.4 \cdot 10^7$ operaciones (factible).
\end{itemize}

\paragraph{Donde tocar para modificar.}
\begin{itemize}\setlength\itemsep{0pt}
	\item \textbf{Reconstruir camino}: aniadir \texttt{next[i][j]} y backtrack.
	\item \textbf{Transitividad}: el algoritmo tambien calcula transitividad (\texttt{d[i][j] != INF} significa alcanzable).
	\item \textbf{Pesos doubles}: aniadir tolerancia (\texttt{d[i][j] > d[i][k] + d[k][j] - eps}).
	\item \textbf{Grafo no denso}: usar Dijkstra $V$ veces en vez de Floyd.
\end{itemize}

\paragraph{Trampas y errores frecuentes.}
\begin{itemize}\setlength\itemsep{0pt}
	\item Para detectar ciclo negativo correctamente, aniadir un chequeo \texttt{if (ans[i][i] < 0)} tras el triple bucle.
	\item El orden de los bucles \texttt{k, i, j} es importante; invertir da resultados diferentes.
	\item Con \texttt{double}, usar tolerancia \texttt{< eps} en vez de \texttt{<} para evitar loops infinitos.
\end{itemize}

\paragraph{Codigo.}
Ver \texttt{cpp.tex}, seccion \textit{Floyd-Warshall}.

\subsection{Componentes y cortes}

\subsubsection{Kosaraju SCC}

\paragraph{Idea intuitiva.}
Encuentra las \textbf{componentes fuertemente conexas} (SCCs) en $O(V + E)$. Algoritmo en 3 pasos:
\begin{itemize}\setlength\itemsep{0pt}
	\item DFS en $G$ rellenando una pila en \textbf{post-orden} (los nodos que terminan tarde quedan arriba).
	\item Construir $G^T$ (grafo transpuesto).
	\item DFS en $G^T$ procesando nodos en orden de la pila (de mayor a menor tiempo); cada DFS da una SCC.
\end{itemize}
La demostracion: dos nodos $u, v$ estan en la misma SCC sii hay camino $u \leadsto v$ y $v \leadsto u$; Kosaraju captura exactamente esto al explorar $G^T$ desde los nodos con mayor tiempo de finalizacion en $G$.

\paragraph{Propiedades clave (invariantes).}
\begin{itemize}\setlength\itemsep{0pt}
	\item Cada SCC es maximal fuertemente conexa.
	\item El \textbf{grafo de SCCs} (metagraf) es un DAG.
	\item Kosaraju y Tarjan dan el mismo resultado; Tarjan es ``mas rapido'' en practica (un solo DFS).
	\item Cada nodo pertenece a exactamente una SCC.
\end{itemize}

\paragraph{Codigo clave.}
La estructura de Kosaraju:
\begin{verbatim}
vector<int> order;
vector<bool> used(n);
function<void(int)> dfs1 = [&](int v) {
    used[v] = true;
    for (int u : g[v]) if (!used[u]) dfs1(u);
    order.push_back(v);  // post-orden
};
function<void(int, int)> dfs2 = [&](int v, int root) {
    comp[v] = root;
    for (int u : gt[v]) if (comp[u] == -1) dfs2(u, root);
};
// 1: DFS en G para llenar order
// 2: DFS en G^T en orden inverso para asignar SCCs
\end{verbatim}

\paragraph{Complejidad.}
\begin{itemize}\setlength\itemsep{0pt}
	\item $O(V + E)$.
	\item Constante 2 (dos DFS + el grafo transpuesto).
\end{itemize}

\paragraph{Donde tocar para modificar.}
\begin{itemize}\setlength\itemsep{0pt}
	\item \textbf{2-SAT}: SCC sobre grafo de implicaciones.
	\item \textbf{Construir el DAG de SCCs}: aniadir aristas entre SCCs diferentes.
	\item \textbf{Tamano de cada SCC}: aniadir contador durante el segundo DFS.
	\item \textbf{Grafo grande}: usar Tarjan para evitar construir $G^T$.
\end{itemize}

\paragraph{Trampas y errores frecuentes.}
\begin{itemize}\setlength\itemsep{0pt}
	\item El snippet usa \texttt{vector<...> adj[n]} (VLA); reemplazable por \texttt{vector<vector<pair<int,int>>} para portabilidad.
	\item El orden del segundo DFS debe ser \textbf{inverso} al de finalizacion del primero.
	\item Si el grafo no es conexo, aniadir un loop sobre todos los nodos en el primer DFS.
	\item Confundir SCC con componentes conexas (las primeras son para dirigidos, las segundas para no dirigidos).
\end{itemize}

\paragraph{Codigo.}
Ver \texttt{cpp.tex}, seccion \textit{Kosaraju SCC}.

\vspace{0.5em}

\subsubsection{Tarjan (puentes)}

\paragraph{Idea intuitiva.}
Encuentra los \textbf{puentes} (bridges): aristas cuya remocion \textbf{desconecta} el grafo. Usa un solo DFS manteniendo \texttt{tin[v]} (tiempo de descubrimiento) y \texttt{low[v]} (el menor \texttt{tin} alcanzable desde $v$ via back-edges). Una arista $v - u$ es puente si $low[u] > tin[v]$. La demostracion: si $low[u] > tin[v]$, entonces $u$ (y su subarbol) no pueden ``escapar'' de $v$; quitar la arista $v - u$ parte el grafo.

\paragraph{Propiedades clave (invariantes).}
\begin{itemize}\setlength\itemsep{0pt}
	\item \texttt{low[v] = min(tin[v], tin[w])} para todo back-edge $(v, w)$, y \texttt{min(low[u])} para hijos en el DFS.
	\item Solo se consideran las aristas del \textbf{arbol DFS} (no back-edges) para \texttt{low}.
	\item Cada nodo tiene exactamente un \texttt{tin}; cada back-edge se procesa una vez.
\end{itemize}

\paragraph{Codigo clave.}
El DFS de Tarjan para bridges:
\begin{verbatim}
void dfs(int v, int pEdge) {
    used[v] = true;
    tin[v] = low[v] = timer++;
    for (auto [u, id] : adj[v]) {
        if (id == pEdge) continue;  // no contar el edge al padre
        if (used[u]) {
            low[v] = min(low[v], tin[u]);  // back-edge
        } else {
            dfs(u, id);
            low[v] = min(low[v], low[u]);
            if (low[u] > tin[v]) bridges.push_back(id);  // es puente
        }
    }
}
\end{verbatim}

\paragraph{Complejidad.}
\begin{itemize}\setlength\itemsep{0pt}
	\item $O(V + E)$.
	\item Una sola pasada de DFS.
\end{itemize}

\paragraph{Donde tocar para modificar.}
\begin{itemize}\setlength\itemsep{0pt}
	\item \textbf{Bridge tree}: contractar cada 2-edge-connected component a un nodo; el resultado es un arbol.
	\item \textbf{Articulation points}: variante con condicion diferente (ver el siguiente modulo).
	\item \textbf{Grafo no dirigido}: solo funciona en no dirigidos; en dirigidos hay conceptos analogos (strong bridges, dominator tree).
	\item \textbf{Reconstruir componentes 2-edge-connected}: BFS desde cada nodo excluyendo bridges.
\end{itemize}

\paragraph{Trampas y errores frecuentes.}
\begin{itemize}\setlength\itemsep{0pt}
	\item No confundir el parent con un back-edge al padre; chequear \texttt{if (adjN == parent) continue;} y contar solo las veces que se ``visita'' como hijo.
	\item Para multigrafos, los edges deben tener IDs unicos para distinguirlos.
	\item Si el grafo no es conexo, aniadir loop externo sobre todos los nodos.
\end{itemize}

\paragraph{Codigo.}
Ver \texttt{cpp.tex}, seccion \textit{Tarjan (puentes)}.

\vspace{0.5em}

\subsubsection{Tarjan (puntos de articulacion)}

\paragraph{Idea intuitiva.}
Encuentra los \textbf{puntos de articulacion} (cut vertices): vertices cuya remocion \textbf{desconecta} el grafo. Mismo setup que para bridges: \texttt{tin[v]}, \texttt{low[v]}. Un nodo $v$ (no raiz) es de articulacion si existe un hijo $u$ con $low[u] \ge tin[v]$. La \textbf{raiz} es de articulacion si tiene $\ge 2$ hijos en el DFS tree. La razon: para $v$ no raiz, la condicion $\ge$ indica que el subarbol de $u$ depende de $v$; quitar $v$ los desconecta. Para la raiz, perder $\ge 2$ hijos parte el DFS en $\ge 2$ componentes.

\paragraph{Propiedades clave (invariantes).}
\begin{itemize}\setlength\itemsep{0pt}
	\item $low[u] \ge tin[v]$ significa que el subarbol de $u$ no tiene back-edge ``por encima'' de $v$.
	\item La raiz es caso especial: tiene que perder mas de un hijo para que el grafo se parta.
	\item Cada nodo se evalua una sola vez.
\end{itemize}

\paragraph{Codigo clave.}
La condicion de articulacion:
\begin{verbatim}
void dfs(int v, int p) {
    used[v] = true;
    tin[v] = low[v] = timer++;
    int children = 0;
    for (int u : adj[v]) {
        if (u == p) continue;
        if (used[u]) low[v] = min(low[v], tin[u]);
        else {
            dfs(u, v);
            low[v] = min(low[v], low[u]);
            if (low[u] >= tin[v] && p != -1)  // no raiz
                isArt[v] = true;
            ++children;
        }
    }
    if (p == -1 && children > 1) isArt[v] = true;  // raiz especial
}
\end{verbatim}

\paragraph{Complejidad.}
\begin{itemize}\setlength\itemsep{0pt}
	\item $O(V + E)$.
	\item Una sola pasada de DFS.
\end{itemize}

\paragraph{Donde tocar para modificar.}
\begin{itemize}\setlength\itemsep{0pt}
	\item \textbf{Block-cut tree}: estructura bipartita con BCCs (2-vertex-connected components) y articulation points.
	\item \textbf{Verificar biconectividad}: chequear que no hay puntos de articulacion.
	\item \textbf{Componentes biconnected}: extender Tarjan para enumerar BCCs.
\end{itemize}

\paragraph{Trampas y errores frecuentes.}
\begin{itemize}\setlength\itemsep{0pt}
	\item No olvidar el caso especial de la raiz.
	\item La condicion usa \texttt{>=}, no \texttt{>}; un hijo con \texttt{low[u] == tin[v]} tambien cuenta.
	\item Multigrafos pueden tener ``falsos'' articulation points; verificar caso a caso.
\end{itemize}

\paragraph{Codigo.}
Ver \texttt{cpp.tex}, seccion \textit{Tarjan (puntos de articulacion)}.

\vspace{0.5em}

\subsubsection{Tarjan SCC}

\paragraph{Idea intuitiva.}
Variante del SCC que usa \textbf{un solo DFS} con una pila explicita. Mantiene \texttt{disc[v]}, \texttt{low[v]}, y una pila de nodos ``en la SCC actual''. Cuando \texttt{low[v] == disc[v]}, $v$ es la raiz de una SCC; popear hasta $v$ inclusive. La demostracion: $low[v] == disc[v]$ indica que $v$ no tiene back-edges a ancestros; entonces $v$ y los nodos en la pila hasta $v$ forman una SCC maximal.

\paragraph{Propiedades clave (invariantes).}
\begin{itemize}\setlength\itemsep{0pt}
	\item Mas eficiente que Kosaraju (un solo DFS).
	\item \texttt{low[v] = min(low[u], disc[w])} para back-edges o recursion.
	\item La pila asegura que solo los nodos en la SCC actual se popean.
	\item Cada nodo se pushea y popea a lo sumo una vez.
\end{itemize}

\paragraph{Codigo clave.}
El DFS de Tarjan SCC:
\begin{verbatim}
void dfs(int v) {
    disc[v] = low[v] = timer++;
    st.push(v); inStack[v] = true;
    for (int u : g[v]) {
        if (disc[u] == -1) { dfs(u); low[v] = min(low[v], low[u]); }
        else if (inStack[u]) low[v] = min(low[v], disc[u]);
    }
    if (low[v] == disc[v]) {
        // popear hasta v
        while (true) {
            int u = st.top(); st.pop(); inStack[u] = false;
            comp[u] = current_scc;
            if (u == v) break;
        }
        ++current_scc;
    }
}
\end{verbatim}

\paragraph{Complejidad.}
\begin{itemize}\setlength\itemsep{0pt}
	\item $O(V + E)$.
	\item Constante similar a Kosaraju pero sin construir $G^T$.
\end{itemize}

\paragraph{Donde tocar para modificar.}
\begin{itemize}\setlength\itemsep{0pt}
	\item \textbf{Tamano de cada SCC}: aniadir contador durante el pop.
	\item \textbf{Grafo condensation}: aniadir aristas entre SCCs en el DAG resultante.
	\item \textbf{2-SAT}: aniadir la construccion del grafo de implicaciones.
\end{itemize}

\paragraph{Trampas y errores frecuentes.}
\begin{itemize}\setlength\itemsep{0pt}
	\item \texttt{inStack[u]} debe chequearse antes de usar \texttt{low[v] = min(low[v], disc[u])} en back-edges.
	\item La pila global es importante; sin ella no se puede extraer la SCC.
	\item Si el grafo es muy grande, recursion puede overflow; aniadir version iterativa.
\end{itemize}

\paragraph{Codigo.}
Ver \texttt{cpp.tex}, seccion \textit{Tarjan SCC}.

\subsection{Matching}

\subsubsection{Kuhn (bipartite matching)}

\paragraph{Idea intuitiva.}
Encuentra un \textbf{matching maximo} en un grafo bipartito en $O(VE)$. Para cada nodo $u$ en $L$, intentar encontrar un vecino $v$ en $R$ que este libre o cuyo match se pueda ``reorganizar'' recursivamente. Si se encuentra, matchear $u - v$. La idea: si $u$ quiere un vecino $v$ ya ocupado, intentamos ``mover'' el match de $v$ a otro nodo; recursivamente, esto encuentra un augmenting path si existe.

\paragraph{Propiedades clave (invariantes).}
\begin{itemize}\setlength\itemsep{0pt}
	\item \texttt{matchR[v]} = el nodo de $L$ matcheado a $v$ (o $-1$ si libre).
	\item \texttt{vis[]} se resetea \textbf{cada intento} de $u$.
	\item Si el matching es perfecto, $|L| = |R| = $ tamano del matching.
	\item Kuhn no garantiza maximalidad greedy; puede haber augmenting paths no encontrados.
\end{itemize}

\paragraph{Codigo clave.}
La recursion del augmenting path:
\begin{verbatim}
bool tryKuhn(int v) {
    if (vis[v]) return false;
    vis[v] = true;
    for (int u : g[v]) {
        if (matchR[u] == -1 || tryKuhn(matchR[u])) {
            matchR[u] = v;
            return true;
        }
    }
    return false;
}
// Por cada v en L:
fill(vis, vis+n, false);
if (tryKuhn(v)) ++matches;
\end{verbatim}

\paragraph{Complejidad.}
\begin{itemize}\setlength\itemsep{0pt}
	\item $O(VE)$ en el peor caso.
	\item Con Hopcroft-Karp: $O(\sqrt{V} \cdot E)$, mucho mejor para grafos densos.
\end{itemize}

\paragraph{Donde tocar para modificar.}
\begin{itemize}\setlength\itemsep{0pt}
	\item \textbf{Matching bipartito minimo (min vertex cover)}: Konig's theorem.
	\item \textbf{Matching general}: blossom algorithm (no incluido).
	\item \textbf{Aniadir capacidades}: si los nodos de $R$ tienen capacidad $> 1$, partir cada uno en capacidad copias.
	\item \textbf{Output del matching}: el array \texttt{matchR} da los pares finales.
\end{itemize}

\paragraph{Trampas y errores frecuentes.}
\begin{itemize}\setlength\itemsep{0pt}
	\item Olvidar resetear \texttt{vis} por intento; eso es lo que evita ``ciclos infinitos'' en la recursion.
	\item Para grafos grandes ($V > 10^4$), Kuhn puede ser muy lento; usar Hopcroft-Karp.
	\item Si el matching es maximo pero no perfecto, $|L|$ o $|R|$ puede ser mayor; el matching es solo maximo.
\end{itemize}

\paragraph{Codigo.}
Ver \texttt{cpp.tex}, seccion \textit{Kuhn (bipartite matching)}.

\vspace{0.5em}

\subsubsection{Hopcroft-Karp}

\paragraph{Idea intuitiva.}
Matching bipartito maximo en $O(\sqrt{V} E)$. Algoritmo en \textbf{fases}: en cada fase, hacer BFS para encontrar los \textbf{caminos de aumento mas cortos} desde nodos libres de $L$; luego DFS solo por esos caminos. Cada fase aumenta el matching en $\ge 1$ y toma $O(E)$. La demostracion: tras una fase, todos los augmenting paths tienen longitud $\ge$ la nueva minima; eso sucede cada $\sqrt{V}$ fases (longitud minima crece geometricamente).

\paragraph{Propiedades clave (invariantes).}
\begin{itemize}\setlength\itemsep{0pt}
	\item \texttt{dist[u]} = distancia en niveles desde un nodo libre de $L$.
	\item BFS termina cuando no hay mas camino a un nodo libre de $R$.
	\item DFS respeta los niveles (solo sigue \texttt{dist[matchV[v]] == dist[u] + 1}).
	\item Cada fase aumenta el matching en $\ge 1$ arista.
\end{itemize}

\paragraph{Codigo clave.}
El BFS de Hopcroft-Karp:
\begin{verbatim}
bool bfs() {
    queue<int> q;
    for (int v : L) {
        if (matchU[v] == -1) { dist[v] = 0; q.push(v); }
        else dist[v] = -1;
    }
    bool found = false;
    while (!q.empty()) {
        int v = q.front(); q.pop();
        for (int u : g[v]) {
            if (matchV[u] != -1 && dist[matchV[u]] == -1) {
                dist[matchV[u]] = dist[v] + 1;
                q.push(matchV[u]);
            }
            if (matchV[u] == -1) found = true;  // llegamos a libre
        }
    }
    return found;
}
\end{verbatim}

\paragraph{Complejidad.}
\begin{itemize}\setlength\itemsep{0pt}
	\item $O(\sqrt{V} E)$.
	\item Para $V = 10^4$, $E = 10^5$: $\sim 3 \cdot 10^7$ operaciones.
\end{itemize}

\paragraph{Donde tocar para modificar.}
\begin{itemize}\setlength\itemsep{0pt}
	\item \textbf{Reconstruir el matching}: ya viene en \texttt{matchU}/\texttt{matchV}.
	\item \textbf{Aniadir pesos}: Hungarian algorithm.
	\item \textbf{Densidad alta}: si $E \approx V^2$, sigue siendo util pero verificar constantes.
	\item \textbf{Aniadir capacidades}: partir los nodos con capacidad $> 1$ en varias copias.
\end{itemize}

\paragraph{Trampas y errores frecuentes.}
\begin{itemize}\setlength\itemsep{0pt}
	\item \texttt{dist[matchV[v]] == dist[u] + 1} (no \texttt{==} a secas); sin esto, el DFS no respeta los niveles y se pierde la complejidad.
	\item El BFS debe marcar los niveles correctamente; resetear \texttt{dist} para nodos alcanzables solo via matching.
	\item El DFS debe respetar los niveles; cualquier ``atajo'' invalida la cota.
\end{itemize}

\paragraph{Codigo.}
Ver \texttt{cpp.tex}, seccion \textit{Hopcroft-Karp}.

\subsection{Basicos}

\subsubsection{DFS / BFS / componentes}

\paragraph{Idea intuitiva.}
\textbf{DFS} (Depth-First Search) y \textbf{BFS} (Breadth-First Search) son los dos recorridos basicos de grafos. DFS usa una pila (o recursion) y explora tan profundo como puede antes de retroceder; BFS usa una cola y explora por ``niveles'' de distancia. La cantidad de \textbf{componentes conexas} se cuenta iniciando un BFS/DFS desde cada nodo no visitado. La diferencia clave: BFS da el camino mas corto en aristas; DFS da un orden de finalizacion util para problemas topologicos.

\paragraph{Propiedades clave (invariantes).}
\begin{itemize}\setlength\itemsep{0pt}
	\item BFS da la \textbf{distancia en aristas} mas corta desde el origen.
	\item DFS da un \textbf{orden topologico} (si se aniaden al final).
	\item Ambos son $O(V + E)$.
	\item Un grafo con $C$ componentes conexas satisface $\sum \text{subtree\_size}_v = V$ y $\sum E_v = 2E$.
	\item DFS en arbol da un spanning tree; las ``back edges'' corresponden a ciclos.
\end{itemize}

\paragraph{Codigo clave.}
BFS clasico:
\begin{verbatim}
queue<int> q;
q.push(source);
dist[source] = 0;
while (!q.empty()) {
    int v = q.front(); q.pop();
    for (int u : adj[v])
        if (dist[u] == -1) {
            dist[u] = dist[v] + 1;
            q.push(u);
        }
}
\end{verbatim}

\paragraph{Complejidad.}
\begin{itemize}\setlength\itemsep{0pt}
	\item $O(V + E)$ tiempo, $O(V)$ memoria.
	\item Cada arista se procesa exactamente 2 veces (ida y vuelta) en no dirigidos.
\end{itemize}

\paragraph{Donde tocar para modificar.}
\begin{itemize}\setlength\itemsep{0pt}
	\item \textbf{Aniadir pesos}: si las aristas tienen peso y queres shortest path, BFS ya no sirve; usar Dijkstra.
	\item \textbf{Grafo dirigido}: en DFS, aniadir check \texttt{if (state[u] == 1)} para detectar ciclos.
	\item \textbf{Bipartite check}: BFS alternando colores (ver el siguiente modulo).
	\item \textbf{Componentes fuertemente conexas (dirigidos)}: usar Tarjan o Kosaraju.
	\item \textbf{BFS en matriz}: aniadir offsets por direccion para grid 2D.
\end{itemize}

\paragraph{Trampas y errores frecuentes.}
\begin{itemize}\setlength\itemsep{0pt}
	\item Stack overflow con DFS recursivo para $V > 10^5$; usar version iterativa con pila explicita.
	\item En BFS, olvidar marcar \texttt{dist[u] = dist[v] + 1} antes de pushear da visitas duplicadas.
	\item Para grafos dirigidos, ``componente conexa'' no esta definida; usar SCC.
\end{itemize}

\paragraph{Codigo.}
Ver \texttt{cpp.tex}, seccion \textit{DFS / BFS / componentes}.

\vspace{0.5em}

\subsubsection{Topological sort (Kahn + DFS)}

\paragraph{Idea intuitiva.}
Un \textbf{orden topologico} de un DAG es una permutacion de los vertices tal que toda arista $u \to v$ tiene $u$ antes que $v$. Dos algoritmos clasicos:
\begin{itemize}\setlength\itemsep{0pt}
	\item \textbf{Kahn}: BFS sobre nodos con \texttt{indeg == 0}; al ``consumir'' un nodo, decrementar el indeg de sus vecinos; los que lleguen a 0 se encolan.
	\item \textbf{DFS}: hacer DFS; al terminar un nodo (estado = 2), aniadirlo al final; el orden es el reverso.
\end{itemize}
La demostracion: en un DAG siempre hay al menos un nodo con indeg 0; quitarlo mantiene el DAG; por induccion, se obtiene el orden.

\paragraph{Propiedades clave (invariantes).}
\begin{itemize}\setlength\itemsep{0pt}
	\item Existe sii el grafo es un DAG.
	\item Si al final \texttt{order.size() < V}, hay un ciclo.
	\item Kahn da el orden ``natural'' (los que pueden ir primero); DFS da el orden ``al reves'' (los que pueden ir al final primero).
	\item La cantidad de ordenes topologicos puede ser exponencial.
\end{itemize}

\paragraph{Codigo clave.}
Kahn con priority queue (orden lexicografico):
\begin{verbatim}
priority_queue<int, vector<int>, greater<int>> pq;
for (int v = 0; v < n; ++v) if (indeg[v] == 0) pq.push(v);
while (!pq.empty()) {
    int v = pq.top(); pq.pop();
    order.push_back(v);
    for (int u : adj[v]) {
        if (--indeg[u] == 0) pq.push(u);
    }
}
\end{verbatim}

\paragraph{Complejidad.}
\begin{itemize}\setlength\itemsep{0pt}
	\item $O(V + E)$ ambos.
	\item Con priority queue, $O((V + E) \log V)$.
\end{itemize}

\paragraph{Donde tocar para modificar.}
\begin{itemize}\setlength\itemsep{0pt}
	\item \textbf{Detectar ciclos}: comparar \texttt{order.size()} con \texttt{V}.
	\item \textbf{Orden lexicografico minimo}: usar priority queue en Kahn.
	\item \textbf{Topological sort con pesos}: combinar con shortest path en DAG (un solo pase en orden topologico).
	\item \textbf{Encontrar todos los ordenes}: DP sobre subsets (solo para $V$ pequeno).
\end{itemize}

\paragraph{Trampas y errores frecuentes.}
\begin{itemize}\setlength\itemsep{0pt}
	\item En el snippet, Kahn retorna orden vacio si hay ciclo (util para detectar).
	\item En DFS, no olvidar \texttt{reverse(order)}.
	\item Confundir indeg con outdeg; el orden Kahn quita primero los nodos ``fuente''.
	\item Para grafos no DAG, el algoritmo termina con \texttt{order.size() < V}.
\end{itemize}

\paragraph{Codigo.}
Ver \texttt{cpp.tex}, seccion \textit{Topological sort (Kahn + DFS)}.

\vspace{0.5em}

\subsubsection{Bipartite check + 2-coloring}

\paragraph{Idea intuitiva.}
Un grafo es \textbf{bipartito} si sus vertices se pueden dividir en dos conjuntos $L$ y $R$ tales que toda arista va de $L$ a $R$. Esto equivale a decir que se puede \textbf{2-colorear} sin que dos vertices adyacentes tengan el mismo color. BFS alternando colores detecta esto: si en algun momento dos adyacentes tienen el mismo color, no es bipartito. La demostracion: por cada componente, fijar el color de un nodo determina todos los demas; si en algun paso hay conflicto, hay un ciclo impar (no bipartito).

\paragraph{Propiedades clave (invariantes).}
\begin{itemize}\setlength\itemsep{0pt}
	\item Grafo bipartito $\Leftrightarrow$ sin ciclos impares.
	\item Conexo: si un componente no es bipartito, todo el grafo no lo es.
	\item Si hay \textbf{multiple componentes}, cada uno se colorea independientemente.
	\item La 2-coloracion es unica salvo swap de colores.
\end{itemize}

\paragraph{Codigo clave.}
BFS con 2-coloreo:
\begin{verbatim}
vector<int> color(n, -1);
for (int v = 0; v < n; ++v) if (color[v] == -1) {
    color[v] = 0;
    queue<int> q; q.push(v);
    while (!q.empty()) {
        int u = q.front(); q.pop();
        for (int w : adj[u]) {
            if (color[w] == -1) { color[w] = 1 - color[u]; q.push(w); }
            else if (color[w] == color[u]) return false;
        }
    }
}
return true;
\end{verbatim}

\paragraph{Complejidad.}
\begin{itemize}\setlength\itemsep{0pt}
	\item $O(V + E)$.
	\item Memoria $O(V)$.
\end{itemize}

\paragraph{Donde tocar para modificar.}
\begin{itemize}\setlength\itemsep{0pt}
	\item \textbf{Colorear desde multiples fuentes}: BFS desde cada nodo no coloreado.
	\item \textbf{Bipartito + matching}: Hopcroft-Karp.
	\item \textbf{Grafo bipartito ponderado}: Hungarian algorithm.
	\item \textbf{Verificar bipartito en subgrafo}: solo aplicar BFS a los nodos relevantes.
\end{itemize}

\paragraph{Trampas y errores frecuentes.}
\begin{itemize}\setlength\itemsep{0pt}
	\item Verificar \textbf{todos} los componentes, no solo el primero.
	\item Para grafos con auto-loops, un auto-loop hace al grafo no bipartito.
	\item En multigrafos, las aristas multiples no afectan el chequeo de bipartito.
\end{itemize}

\paragraph{Codigo.}
Ver \texttt{cpp.tex}, seccion \textit{Bipartite check + 2-coloring}.

\vspace{0.5em}

\subsubsection{0-1 BFS}

\paragraph{Idea intuitiva.}
Variante de BFS para grafos con \textbf{pesos 0 o 1}. Usa un \texttt{deque}: aristas de peso 0 se encolan al frente (\texttt{push\_front}), aristas de peso 1 al final (\texttt{push\_back}). Asi, la primera vez que se extrae un nodo, su distancia es la definitiva, igual que en Dijkstra pero en $O(V + E)$. La intuicion: el deque mantiene los nodos ordenados por distancia; las aristas de peso 0 no cambian la distancia (entran al frente), las de peso 1 la incrementan (entran al final).

\paragraph{Propiedades clave (invariantes).}
\begin{itemize}\setlength\itemsep{0pt}
	\item Asume pesos \textbf{solo} 0 o 1.
	\item Mejor caso que Dijkstra (sin factor $\log V$).
	\item No funciona con pesos $\ge 2$.
	\item La cola mantiene los nodos ordenados por distancia no decreciente.
\end{itemize}

\paragraph{Codigo clave.}
El deque doble:
\begin{verbatim}
deque<int> q;
q.push_front(source);
dist[source] = 0;
while (!q.empty()) {
    int v = q.front(); q.pop_front();
    for (auto [w, u] : adj[v]) {
        if (dist[v] + w < dist[u]) {
            dist[u] = dist[v] + w;
            if (w == 0) q.push_front(u);
            else q.push_back(u);
        }
    }
}
\end{verbatim}

\paragraph{Complejidad.}
\begin{itemize}\setlength\itemsep{0pt}
	\item $O(V + E)$.
	\item Cada nodo se inserta a lo sumo $O(1)$ veces (con check de distancia).
\end{itemize}

\paragraph{Donde tocar para modificar.}
\begin{itemize}\setlength\itemsep{0pt}
	\item \textbf{Pesos en $[0, k]$ small}: usar small k-ary BFS o Dijkstra.
	\item \textbf{Reconstruir camino}: aniadir \texttt{parent[]}.
	\item \textbf{Pesos negativos}: no funciona, usar Bellman-Ford.
\end{itemize}

\paragraph{Trampas y errores frecuentes.}
\begin{itemize}\setlength\itemsep{0pt}
	\item Si los pesos son otro rango, el algoritmo da respuestas incorrectas. Verificar que las aristas son solo 0/1.
	\item El chequeo \texttt{dist[v] + w < dist[u]} es crucial; sino, el algoritmo procesa nodos multiples veces.
	\item Si una arista tiene peso $\ge 2$, aniadir verificaciones o cambiar a Dijkstra.
\end{itemize}

\paragraph{Codigo.}
Ver \texttt{cpp.tex}, seccion \textit{0-1 BFS}.

\subsection{MST}

\subsubsection{Kruskal}

\paragraph{Idea intuitiva.}
Encuentra el \textbf{arbol de expansion minima} (MST) en $O(E \log E)$. Ordena las aristas por peso ascendente y las aniade al MST si no forman ciclo (chequear con DSU). El arbol tiene exactamente $V - 1$ aristas. La demostracion (``Cut Property''): para cualquier corte, la arista minima cruzando el corte pertenece a \emph{algun} MST; por lo tanto, elegir la minima disponible que no forma ciclo es optimo.

\paragraph{Propiedades clave (invariantes).}
\begin{itemize}\setlength\itemsep{0pt}
	\item Algoritmo \textbf{greedy} correcto (Cut Property).
	\item Funciona en grafos \textbf{no conexos}; produce un \textbf{minimum spanning forest}.
	\item Si las aristas tienen pesos unicos, el MST es unico.
	\item El MST tiene exactamente $V - 1$ aristas (o menos si el grafo es disconexo).
\end{itemize}

\paragraph{Codigo clave.}
Kruskal con DSU:
\begin{verbatim}
sort(edges.begin(), edges.end());
DSU dsu(n);
long long mst = 0;
int edgesUsed = 0;
for (auto [w, u, v] : edges)
    if (dsu.unite(u, v)) {
        mst += w;
        if (++edgesUsed == n - 1) break;
    }
\end{verbatim}

\paragraph{Complejidad.}
\begin{itemize}\setlength\itemsep{0pt}
	\item $O(E \log E)$ (ordenamiento) o $O(E \log V)$ con DSU optimizations.
	\item DSU es casi $O(1)$ amortizado, asi que el sort domina.
\end{itemize}

\paragraph{Donde tocar para modificar.}
\begin{itemize}\setlength\itemsep{0pt}
	\item \textbf{Second best MST}: enumerar aristas del MST, sacar una, aniadir una no-MST; recalcular.
	\item \textbf{Kruskal randomizado}: si el grafo es casi-completo, agregar shuffle.
	\item \textbf{Aniadir info al DSU}: guardar el peso maximo en el camino, etc.
	\item \textbf{DSU con rollback}: si necesitas Kruskal reversible (para queries offline).
\end{itemize}

\paragraph{Trampas y errores frecuentes.}
\begin{itemize}\setlength\itemsep{0pt}
	\item El DSU debe estar limpio para cada invocacion.
	\item Si las aristas tienen pesos negativos, Kruskal sigue funcionando.
	\item Si el grafo es disconexo, el ``MST'' es un bosque con varias componentes.
	\item El check \texttt{if (++edgesUsed == n - 1) break;} optimiza la salida temprana.
\end{itemize}

\paragraph{Codigo.}
Ver \texttt{cpp.tex}, seccion \textit{Kruskal}.

\vspace{0.5em}

\subsubsection{Prim (priority queue)}

\paragraph{Idea intuitiva.}
Otra forma de calcular el MST: empezar desde un nodo y, en cada paso, aniadir la arista de menor peso que conecta el conjunto construido con el resto. Usando priority queue, es $O((V + E) \log V)$. La demostracion: cada vez que se aniade una arista, es la minima cruzando el ``corte'' entre el conjunto actual y el resto (Cut Property); eso preserva optimalidad.

\paragraph{Propiedades clave (invariantes).}
\begin{itemize}\setlength\itemsep{0pt}
	\item Similar a Dijkstra pero sin acumular distancias, solo eligiendo minimo por arista.
	\item Funciona mejor que Kruskal en \textbf{grafos densos} ($E \approx V^2$).
	\item Si el grafo es disconexo, genera el bosque (solo llega a los alcanzables).
	\item Cada nodo se inserta al heap una vez; cada arista se procesa una vez.
\end{itemize}

\paragraph{Codigo clave.}
Prim clasico con heap:
\begin{verbatim}
priority_queue<pair<int, int>, vector<...>, greater<...>> pq;
vector<bool> used(n, false);
vector<int> minEdge(n, INF);
minEdge[0] = 0;
pq.push({0, 0});
while (!pq.empty()) {
    auto [w, v] = pq.top(); pq.pop();
    if (used[v]) continue;
    used[v] = true;
    mst += w;
    for (auto [peso, u] : adj[v])
        if (!used[u] && peso < minEdge[u]) {
            minEdge[u] = peso;
            pq.push({peso, u});
        }
}
\end{verbatim}

\paragraph{Complejidad.}
\begin{itemize}\setlength\itemsep{0pt}
	\item $O((V + E) \log V)$.
	\item Con matriz de adyacencia: $O(V^2)$ sin heap.
\end{itemize}

\paragraph{Donde tocar para modificar.}
\begin{itemize}\setlength\itemsep{0pt}
	\item \textbf{Prim con matriz de adyacencia}: $O(V^2)$ sin heap, util en grafos densos pequenos.
	\item \textbf{Reconstruir el MST}: aniadir \texttt{parent[]}.
	\item \textbf{Prim en streaming}: variante para grafos que llegan incrementalmente.
\end{itemize}

\paragraph{Trampas y errores frecuentes.}
\begin{itemize}\setlength\itemsep{0pt}
	\item El \texttt{if (used[v]) continue;} es crucial: si el nodo ya fue procesado, saltar.
	\item En la version matriz, elije el minimo por linea (no heap).
	\item Si el grafo es disconexo, aniadir un loop externo sobre todos los nodos no usados.
\end{itemize}

\paragraph{Codigo.}
Ver \texttt{cpp.tex}, seccion \textit{Prim (priority queue)}.

\subsection{Flujos}

\subsubsection{Edmonds-Karp (max flow)}

\paragraph{Idea intuitiva.}
Implementacion BFS del algoritmo de \textbf{Ford-Fulkerson}. En cada iteracion, hacer BFS desde $s$ hasta $t$ siguiendo aristas con capacidad residual $> 0$; si se llega a $t$, se encontro un \textbf{camino de aumento}. Enviar el minimo de las capacidades en el camino; actualizar capacidades residuales. La demostracion de $O(VE^2)$: cada BFS encuentra el augmenting path mas corto, y la longitud minima crece monotona; como hay $\le E$ longitudes posibles, hay $\le E$ fases.

\paragraph{Propiedades clave (invariantes).}
\begin{itemize}\setlength\itemsep{0pt}
	\item BFS garantiza caminos de aumento mas cortos.
	\item Capacidad residual: forward edge baja, backward edge sube.
	\item Complejidad $O(VE^2)$.
	\item El flujo maximo respeta la conservacion en nodos (in = out salvo $s, t$).
\end{itemize}

\paragraph{Codigo clave.}
BFS + augmenting path:
\begin{verbatim}
while (bfs(s, t, parent)) {  // bfs encuentra camino o retorna false
    int pathFlow = INF;
    for (int v = t; v != s; v = parent[v])
        pathFlow = min(pathFlow, capacity[parent[v]][v]);
    for (int v = t; v != s; v = parent[v]) {
        capacity[parent[v]][v] -= pathFlow;
        capacity[v][parent[v]] += pathFlow;
    }
    maxFlow += pathFlow;
}
\end{verbatim}

\paragraph{Complejidad.}
\begin{itemize}\setlength\itemsep{0pt}
	\item $O(VE^2)$.
	\item En la practica, mas lento que Dinic para grafos grandes.
\end{itemize}

\paragraph{Donde tocar para modificar.}
\begin{itemize}\setlength\itemsep{0pt}
	\item \textbf{Grafo bipartito}: $O(\sqrt{V} E)$ (equivalente a Hopcroft-Karp).
	\item \textbf{Dinic} es preferible en la mayoria de los casos.
	\item \textbf{Capacidades doubles}: usar tolerancia para comparacion.
	\item \textbf{Min-cost flow}: aniadir costo por arista y modificar el algoritmo.
\end{itemize}

\paragraph{Trampas y errores frecuentes.}
\begin{itemize}\setlength\itemsep{0pt}
	\item Las capacidades deben ser $> 0$ para entrar a la cola (no $>= 0$); sino, ciclos infinitos.
	\item Para grafos bipartitos, Dinic o Hopcroft-Karp son mejores.
	\item Si el flujo maximo es muy grande, los \texttt{int} pueden overflow; usar \texttt{long long}.
\end{itemize}

\paragraph{Codigo.}
Ver \texttt{cpp.tex}, seccion \textit{Edmonds-Karp (max flow)}.

\vspace{0.5em}

\subsubsection{Dinic (max flow)}

\paragraph{Idea intuitiva.}
Algoritmo de flujo mas eficiente que Edmonds-Karp: combina \textbf{BFS para level graph} con \textbf{DFS bloqueante} (multiple augmenting paths en una sola pasada). Cada fase cuesta $O(E)$ y hay $O(V)$ fases en el peor caso; para bipartitos, $O(\sqrt{V} E)$. La razon: cada fase ``satura'' al menos una arista, asi que en $O(E)$ fases el flujo maximo se alcanza; con $O(V)$ niveles en el BFS, hay $O(VE)$ operaciones totales.

\paragraph{Propiedades clave (invariantes).}
\begin{itemize}\setlength\itemsep{0pt}
	\item \texttt{level[v]} = distancia desde $s$ en el level graph (solo aristas con cap $> 0$).
	\item \texttt{it[v]} = iterador actual para evitar reprocesar aristas.
	\item DFS solo sigue aristas con \texttt{level[e.to] == level[v] + 1}.
	\item Back-edges tienen \texttt{rev} apuntando a la forward edge original.
\end{itemize}

\paragraph{Codigo clave.}
El DFS bloqueante de Dinic:
\begin{verbatim}
long long dfs(int v, int t, long long f) {
    if (v == t) return f;
    for (int& i = it[v]; i < adj[v].size(); ++i) {
        Edge& e = adj[v][i];
        if (e.cap > 0 && level[e.to] == level[v] + 1) {
            long long ret = dfs(e.to, t, min(f, e.cap));
            if (ret > 0) {
                e.cap -= ret;
                adj[e.to][e.rev].cap += ret;
                return ret;
            }
        }
    }
    return 0;
}
\end{verbatim}

\paragraph{Complejidad.}
\begin{itemize}\setlength\itemsep{0pt}
	\item $O(V^2 E)$ general, $O(\sqrt{V} E)$ para bipartitos.
	\item En la practica, mucho mas rapido que Edmonds-Karp.
\end{itemize}

\paragraph{Donde tocar para modificar.}
\begin{itemize}\setlength\itemsep{0pt}
	\item \textbf{Dinic con scaling}: aniadir factor de escala para capacidades grandes.
	\item \textbf{Lower bounds}: aniadir un super-source/sink para forzar flujo minimo.
	\item \textbf{Matching bipartito}: el grafo se construye como source -> L -> R -> sink, todos con capacidad 1.
	\item \textbf{Push-relabel}: alternativa para grafos muy densos.
\end{itemize}

\paragraph{Trampas y errores frecuentes.}
\begin{itemize}\setlength\itemsep{0pt}
	\item \texttt{it[v]} debe \textbf{incrementarse} en cada iteracion; sino, el DFS se queda atascado en la misma arista.
	\item Si el BFS no encuentra $t$, el flujo es maximo; terminar.
	\item Las back-edges son cruciales para ``cancelar'' elecciones previas.
\end{itemize}

\paragraph{Codigo.}
Ver \texttt{cpp.tex}, seccion \textit{Dinic (max flow)}.

\vspace{0.5em}

\subsubsection{Min-Cost Max-Flow}

\paragraph{Idea intuitiva.}
Encuentra el flujo maximo con \textbf{costo minimo} (asumiendo costos no negativos). Usa \textbf{SPFA} (o Dijkstra con potentials) para encontrar el shortest path en el \textbf{costo} desde $s$ hasta $t$ considerando capacidades residuales. Cada augmenting path aporta su costo $\times$ flujo. La idea: los potentials reescalan las aristas para que todas tengan costo no negativo, permitiendo usar Dijkstra.

\paragraph{Propiedades clave (invariantes).}
\begin{itemize}\setlength\itemsep{0pt}
	\item \texttt{pot[]} = potentials para evitar aristas negativas (Johnson's trick).
	\item \texttt{dist[v]} = costo minimo desde $s$ con potentials.
	\item Solo se aumentan aristas con \texttt{cap > 0} (forward) y costo es \texttt{+cost} (forward) o \texttt{-cost} (backward).
	\item Tras cada augmenting, los potentials se actualizan para mantener la consistencia.
\end{itemize}

\paragraph{Codigo clave.}
El ciclo de min-cost flow:
\begin{verbatim}
while (spfa(s, t, dist, parent)) {
    int pathFlow = INF;
    for (int v = t; v != s; v = parent[v])
        pathFlow = min(pathFlow, capacity[parent[v]][v]);
    for (int v = t; v != s; v = parent[v]) {
        capacity[parent[v]][v] -= pathFlow;
        capacity[v][parent[v]] += pathFlow;
        cost += pathFlow * edge_cost[parent[v]][v];
    }
}
\end{verbatim}

\paragraph{Complejidad.}
\begin{itemize}\setlength\itemsep{0pt}
	\item $O(F \cdot V E)$ con SPFA, $O(F \cdot E \log V)$ con Dijkstra + potentials.
	\item $F$ = flujo maximo total.
\end{itemize}

\paragraph{Donde tocar para modificar.}
\begin{itemize}\setlength\itemsep{0pt}
	\item \textbf{Dijkstra + potentials}: si los costos son enteros no negativos, esto es mas estable.
	\item \textbf{Costos negativos en aristas}: requieren Bellman-Ford inicial o reorden.
	\item \textbf{Flujo minimo con costo}: encontrar el minimo flujo posible y su costo.
	\item \textbf{Costo por unidad}: si los costos son muy grandes, aniadir verificacion.
\end{itemize}

\paragraph{Trampas y errores frecuentes.}
\begin{itemize}\setlength\itemsep{0pt}
	\item El \texttt{pot[]} se actualiza \textbf{al final} de cada iteracion (no durante).
	\item Las back-edges tienen \textbf{costo negativo}, esencial para corregir elecciones previas.
	\item Si los costos son muy grandes, considerar modular arithmetic.
	\item Para problemas con flujo exacto, aniadir verificacion post-algoritmo.
\end{itemize}

\paragraph{Codigo.}
Ver \texttt{cpp.tex}, seccion \textit{Min-Cost Max-Flow}.

\subsection{Especales}

\subsubsection{2-SAT (implication graph + SCC)}

\paragraph{Idea intuitiva.}
Resuelve formulas booleanas en \textbf{CNF} con clausulas de \textbf{2 literales}: $(l_1 \lor l_2)$. Cada variable $x$ se representa con dos nodos: $x$ (verdadero) y $\neg x$ (falso). Cada clausula $l_1 \lor l_2$ se descompone en dos implicaciones: $\neg l_1 \to l_2$ y $\neg l_2 \to l_1$. Si en el grafo de implicaciones $\neg x$ y $x$ estan en la misma SCC, la formula es \textbf{UNSAT}. La demostracion: si $\neg x \to x$ en el grafo, entonces $\neg x$ implica $x$, asi que $\neg x$ debe ser falso y $x$ verdadero; pero entonces $\neg x$ es falso implica una contradiccion.

\paragraph{Propiedades clave (invariantes).}
\begin{itemize}\setlength\itemsep{0pt}
	\item Cada literal es un nodo; total $2n$ nodos.
	\item Kosaraju/Tarjan dan las SCCs.
	\item Si SAT: asignar $x = $ (topological order del condensation).
	\item Clausulas \texttt{setTrue(x)} se modelan como \texttt{implies(!x, x)}.
	\item La formula es SAT sii no hay variable $x$ con $x$ y $\neg x$ en la misma SCC.
\end{itemize}

\paragraph{Codigo clave.}
Conversion de clausulas a implicaciones:
\begin{verbatim}
// Clausula (a OR b) -> implica ~a -> b Y ~b -> a
addImplication(not(a), b);
addImplication(not(b), a);
// Forzar x: implica ~x -> x
if (force) addImplication(not(x), x);
\end{verbatim}

\paragraph{Complejidad.}
\begin{itemize}\setlength\itemsep{0pt}
	\item $O(V + E) = O(N + M)$ con Kosaraju o Tarjan.
	\item Memoria $O(V + E)$.
\end{itemize}

\paragraph{Donde tocar para modificar.}
\begin{itemize}\setlength\itemsep{0pt}
	\item \textbf{Forzar variables}: aniadir \texttt{setTrue} / \texttt{setFalse}.
	\item \textbf{3-SAT}: NP-completo; no usar este algoritmo.
	\item \textbf{Contar soluciones}: las SCCs condensation dan una formula para contar; multiplicar las formas de asignar variables en cada componente no-trivial.
	\item \textbf{Output de la asignacion}: aniadir el vector \texttt{assignment[]}.
\end{itemize}

\paragraph{Trampas y errores frecuentes.}
\begin{itemize}\setlength\itemsep{0pt}
	\item La asignacion usa el \textbf{orden topologico inverso} del condensation: \texttt{val[i] = comp[2i] < comp[2i+1]}.
	\item El ID de $\neg x$ es \texttt{2*x + 1} (o similar); verificar consistencia.
	\item Si hay clausulas con literales repetidos (e.g., $(x \lor x)$), tratar como $(x)$ simple.
\end{itemize}

\paragraph{Codigo.}
Ver \texttt{cpp.tex}, seccion \textit{2-SAT (implication graph + SCC)}.

\vspace{0.5em}

\subsubsection{Euler tour / Hierholzer (camino y circuito)}

\paragraph{Idea intuitiva.}
Encuentra un \textbf{camino o circuito euleriano} (que usa cada arista exactamente una vez). \textbf{Condiciones}:
\begin{itemize}\setlength\itemsep{0pt}
	\item \textbf{No dirigido}: todos los grados pares (circuito) o exactamente 2 impares (camino, empieza en uno, termina en el otro). Ademas, conexo sobre el subgrafo con aristas.
	\item \textbf{Dirigido}: $\text{indeg}(v) = \text{outdeg}(v)$ para todos (circuito) o exactamente un nodo con $\text{outdeg} = \text{indeg} + 1$ y otro con $\text{indeg} = \text{outdeg} + 1$ (camino).
\end{itemize}

\paragraph{Hierholzer}: elegir un ciclo desde el inicio, mientras haya ciclos ``laterales'' desde un nodo del camino, mergearlos. Implementacion con pila: avanzar por aristas, al no poder mas aniadir a la respuesta y backtrack. La demostracion de $O(V + E)$: cada arista se procesa una vez (entra y sale del stack una vez).

\paragraph{Propiedades clave (invariantes).}
\begin{itemize}\setlength\itemsep{0pt}
	\item Multigrafos permitidos; las aristas se marcan como usadas.
	\item No dirigido: cada arista aparece \textbf{dos veces} en \texttt{adj} (ida y vuelta).
	\item La construccion del path con ids consistentes es importante.
	\item El resultado es unico salvo orden de las aristas en multigrafos.
\end{itemize}

\paragraph{Codigo clave.}
Hierholzer con pila:
\begin{verbatim}
vector<int> path;
stack<int> st;
st.push(start);
while (!st.empty()) {
    int v = st.top();
    while (!adj[v].empty() && used[adj[v].back()]) adj[v].pop_back();
    if (adj[v].empty()) {
        path.push_back(v);
        st.pop();
    } else {
        int e = adj[v].back();
        used[e] = true;
        int u = edges[e].to(v);
        adj[v].pop_back();
        st.push(u);
    }
}
reverse(path.begin(), path.end());
\end{verbatim}

\paragraph{Complejidad.}
\begin{itemize}\setlength\itemsep{0pt}
	\item $O(V + E)$.
	\item Memoria $O(V + E)$.
\end{itemize}

\paragraph{Donde tocar para modificar.}
\begin{itemize}\setlength\itemsep{0pt}
	\item \textbf{Reconstruir con multigrafo}: aniadir ids consistentes antes del algoritmo.
	\item \textbf{Aplicar a ``de Bruijn'' o problemas de strings}: modelar como multigrafo y buscar Euler.
	\item \textbf{Dirigido vs no dirigido}: el snippet tiene ambas versiones.
	\item \textbf{Letter addition}: modelar palabras como aristas y buscar el superstring.
\end{itemize}

\paragraph{Trampas y errores frecuentes.}
\begin{itemize}\setlength\itemsep{0pt}
	\item Si el grafo no es conexo (sobre las aristas), no existe euleriano. Verificar primero.
	\item Para multigrafos, los IDs de aristas son cruciales para distinguirlas.
	\item Si el problema es ``encuentra todos los euler tours'', el algoritmo es exponencial.
\end{itemize}

\paragraph{Codigo.}
Ver \texttt{cpp.tex}, seccion \textit{Euler tour / Hierholzer (camino y circuito)}.

\vspace{0.5em}

\subsubsection{Hamiltonian path / cycle (bitmask DP + backtracking)}

\paragraph{Idea intuitiva.}
Un \textbf{camino hamiltoniano} visita cada vertice \textbf{exactamente una vez}. \textbf{Determinar} si existe es NP-completo en general, pero con $V \le 20$ se puede hacer DP con bitmask en $O(V^2 \cdot 2^V)$. \textbf{Teoremas suficientes}: \textbf{Dirac} (grado $\ge n/2$) y \textbf{Ore} ($\text{deg}(u) + \text{deg}(v) \ge n$ para todo par no adyacente). La idea de la DP: cada estado es (conjunto de visitados, ultimo vertice); la transicion es agregar un vecino no visitado.

\paragraph{Propiedades clave (invariantes).}
\begin{itemize}\setlength\itemsep{0pt}
	\item DP: $dp[mask][v]$ = true si hay un camino que visita exactamente \texttt{mask} y termina en $v$.
	\item TSP: variante con pesos minimos (Hamilton cycle minimo).
	\item Reconstruccion: guardar \texttt{par[mask][v]} (predecesor).
	\item $2^V$ estados con $V$ bits cada uno; factible para $V \le 20$.
\end{itemize}

\paragraph{Codigo clave.}
DP con bitmask:
\begin{verbatim}
dp[1][0] = 0;  // empezar en 0
for (int mask = 1; mask < (1<<n); ++mask)
    for (int v = 0; v < n; ++v) if (dp[mask][v] < INF)
        for (int u = 0; u < n; ++u)
            if (!(mask & (1<<u)))
                dp[mask | (1<<u)][u] = min(dp[mask | (1<<u)][u],
                                            dp[mask][v] + w[v][u]);
// TSP: respuesta = min(dp[(1<<n)-1][v] + w[v][0])
\end{verbatim}

\paragraph{Complejidad.}
\begin{itemize}\setlength\itemsep{0pt}
	\item $O(V^2 \cdot 2^V)$ tiempo, $O(V \cdot 2^V)$ memoria.
	\item Para $V = 20$: $\sim 4 \cdot 10^8$ operaciones (lento pero factible).
\end{itemize}

\paragraph{Donde tocar para modificar.}
\begin{itemize}\setlength\itemsep{0pt}
	\item \textbf{TSP asimetrico}: aniadir pesos $w[v][u]$ (no necesariamente simetricos).
	\item \textbf{Hamilton con inicio obligatorio}: fijar \texttt{src} en el DP inicial.
	\item \textbf{Bitmask DP mas compacto}: usar \texttt{long long} como mascara si $V \le 64$.
	\item \textbf{TSP optimo aproximado}: MST, 2-opt, simulated annealing.
	\item \textbf{Reconstruir el path}: usar \texttt{par[mask][v]} y backtrack.
\end{itemize}

\paragraph{Trampas y errores frecuentes.}
\begin{itemize}\setlength\itemsep{0pt}
	\item TSP asume grafo \textbf{completo} (con \texttt{LINF} para aristas faltantes); si no, aniadir chequeo.
	\item La condicion ``\texttt{mask == (1<<n) - 1}'' cierra el ciclo retornando al origen.
	\item Para $V > 20$, la memoria $2^V$ se vuelve prohibitiva; usar meet-in-the-middle o heuristicas.
\end{itemize}

\paragraph{Codigo.}
Ver \texttt{cpp.tex}, seccion \textit{Hamiltonian path / cycle (bitmask DP + backtracking)}.

% ============================================================
\section{DP}

\subsubsection{Knapsack 0/1}

\paragraph{Idea intuitiva.}
Dado un conjunto de $N$ objetos con peso $w_i$ y valor $v_i$, maximizar $\sum v_i$ con la restriccion $\sum w_i \le W$. La recurrencia clasica:
\[ dp[i][j] = \max(dp[i-1][j],\, dp[i-1][j - w_i] + v_i) \quad \text{si } j \ge w_i. \]
La optimizacion clave: la DP es 1-dimensional si iteramos $j$ \textbf{de mayor a menor} (asi no usamos el objeto $i$ dos veces). La razon de la iteracion inversa: en cada iteracion queremos el estado del ``item anterior''; iterar de menor a mayor contaminaria \texttt{dp[j - w_i]} con el item actual.

\paragraph{Propiedades clave (invariantes).}
\begin{itemize}\setlength\itemsep{0pt}
	\item $dp[j]$ al final = maximo valor con peso $\le j$.
	\item El bucle inverso \texttt{for (j = W; j >= w[i]; j--)} es lo que distingue 0/1 de unbounded.
	\item Solucion: maximo valor con cualquier peso $\le W$.
	\item La DP es exacta; no hay aproximacion.
\end{itemize}

\paragraph{Codigo clave.}
La version 1D optimizada:
\begin{verbatim}
vector<long long> dp(W + 1, 0);
for (int i = 0; i < N; ++i)
    for (int j = W; j >= w[i]; --j)  // invertido
        dp[j] = max(dp[j], dp[j - w[i]] + v[i]);
return dp[W];
\end{verbatim}

\paragraph{Complejidad.}
\begin{itemize}\setlength\itemsep{0pt}
	\item $O(N \cdot W)$ tiempo, $O(W)$ memoria.
	\item Constante muy pequena (solo sumas y max).
\end{itemize}

\paragraph{Donde tocar para modificar.}
\begin{itemize}\setlength\itemsep{0pt}
	\item \textbf{Reconstruir la seleccion}: aniadir \texttt{take[i][j]} (booleano o bool) en una version 2D.
	\item \textbf{Items con peso 0}: edge case; aniadir guard.
	\item \textbf{Knapsack con restricciones de grupos}: ``elegir a lo sumo 1 de cada grupo'' se modela como 0/1 con un peso extra.
	\item \textbf{Tamano de $W$}: si $W \le 10^9$ pero $N$ es chico, usar meet-in-the-middle.
	\item \textbf{Cambiar max por min}: aniadir \texttt{INF} en la inicializacion.
\end{itemize}

\paragraph{Trampas y errores frecuentes.}
\begin{itemize}\setlength\itemsep{0pt}
	\item Overflow en \texttt{dp[j - w[i]] + v[i]}; usar \texttt{long long} si los valores son grandes.
	\item El bucle \textbf{debe} ir de mayor a menor.
	\item Si los pesos son grandes y $N$ pequeno, conviene meet-in-the-middle.
	\item Para Knapsack exacto ($\sum w_i = W$), aniadir check al final.
\end{itemize}

\paragraph{Codigo.}
Ver \texttt{cpp.tex}, seccion \textit{Knapsack 0/1}.

\vspace{0.5em}

\subsubsection{Knapsack unbounded}

\paragraph{Idea intuitiva.}
Igual que Knapsack 0/1 pero \textbf{cada item se puede tomar multiples veces}. La recurrencia es similar:
\[ dp[j] = \max(dp[j],\, dp[j - w_i] + v_i), \]
pero ahora el bucle va de \textbf{menor a mayor} (porque tomar el item $i$ no afecta corridas futuras del mismo item). La idea: tras tomar el item $i$, queremos volver a tener la opcion de tomarlo; iterar de menor a mayor permite que el item ``se acumule''.

\paragraph{Propiedades clave (invariantes).}
\begin{itemize}\setlength\itemsep{0pt}
	\item El bucle \texttt{for (j = w[i]; j <= W; j++)} es lo que distingue unbounded.
	\item A diferencia de 0/1, no hay riesgo de ``usar dos veces el mismo item''.
	\item La DP es exacta; cada ``subconjunto'' se cuenta correctamente.
\end{itemize}

\paragraph{Codigo clave.}
La version 1D con bucle directo:
\begin{verbatim}
vector<long long> dp(W + 1, 0);
for (int i = 0; i < N; ++i)
    for (int j = w[i]; j <= W; ++j)  // directo (no invertido)
        dp[j] = max(dp[j], dp[j - w[i]] + v[i]);
return dp[W];
\end{verbatim}

\paragraph{Complejidad.}
\begin{itemize}\setlength\itemsep{0pt}
	\item $O(N \cdot W)$ tiempo, $O(W)$ memoria.
	\item Constante muy pequena.
\end{itemize}

\paragraph{Donde tocar para modificar.}
\begin{itemize}\setlength\itemsep{0pt}
	\item \textbf{Reconstruir cuantos de cada item}: aniadir \texttt{cnt[i]} y backtrack.
	\item \textbf{Cambiar moneda minima por maxima}: el snippet maximiza valor; para minimizar monedas, cambiar a \texttt{min} con inicializacion en \texttt{INF}.
	\item \textbf{Coin change}: si todos los pesos son positivos y queremos ``minimo numero de monedas'', es unbounded con min.
\end{itemize}

\paragraph{Trampas y errores frecuentes.}
\begin{itemize}\setlength\itemsep{0pt}
	\item Misma que Knapsack 0/1 con overflow.
	\item El orden de iteracion de items no afecta el resultado (convexo).
	\item Si se quiere ``maximo valor con exactamente $k$ items'', aniadir segunda dimension.
\end{itemize}

\paragraph{Codigo.}
Ver \texttt{cpp.tex}, seccion \textit{Knapsack unbounded}.

\vspace{0.5em}

\subsubsection{Knapsack meet-in-the-middle}

\paragraph{Idea intuitiva.}
Para $N \le 40$ y $W$ moderado, no se puede hacer $O(NW)$ directamente. Partir en dos mitades, enumerar todos los $2^{N/2}$ subconjuntos de cada mitad, y combinar. Para cada subconjunto de la segunda mitad con peso $sw_2$, encontrar el maximo valor de la primera mitad con peso $\le W - sw_2$. La idea: ``divide y venceras'' sobre el tamano de la entrada; cada mitad tiene mitad del espacio, asi que la combinacion es geometrica en lugar de multiplicativa.

\paragraph{Propiedades clave (invariantes).}
\begin{itemize}\setlength\itemsep{0pt}
	\item Hay $2^{N/2}$ subconjuntos por mitad; total $O(2^{N/2})$.
	\item \texttt{sums} guarda los valores de subconjuntos de la primera mitad ordenados (o preprocesados).
	\item Para cada subconjunto de la segunda mitad, busqueda lineal o binaria en \texttt{sums}.
	\item La mitad mas grande (con $N/2 + 1$ items para $N$ impar) tiene $2^{N/2+1}$ subconjuntos.
\end{itemize}

\paragraph{Codigo clave.}
Enumerar subconjuntos de la primera mitad:
\begin{verbatim}
int n1 = N / 2, n2 = N - n1;
vector<pair<long long, long long>> sums1;
for (int mask = 0; mask < (1 << n1); ++mask) {
    long long peso = 0, valor = 0;
    for (int i = 0; i < n1; ++i) if (mask & (1 << i)) {
        peso += w[i]; valor += v[i];
    }
    sums1.push_back({peso, valor});
}
sort(sums1.begin(), sums1.end());  // por peso
// eliminar duplicados dominados
\end{verbatim}

\paragraph{Complejidad.}
\begin{itemize}\setlength\itemsep{0pt}
	\item $O(2^{N/2})$ tiempo y memoria (mejor que $O(NW)$ cuando $W > 2^{N/2}$).
	\item Para $N = 40$: $\sim 2 \cdot 10^6$ subconjuntos, factible.
\end{itemize}

\paragraph{Donde tocar para modificar.}
\begin{itemize}\setlength\itemsep{0pt}
	\item \textbf{Tamano mayor}: si $N$ es mayor pero $W$ es chico, conviene iterar por peso, no por subconjuntos.
	\item \textbf{Varias mochilas}: adaptar para multiples $W$.
	\item \textbf{Optimizar el ``combinado''}: usar \texttt{upper\_bound} para encontrar el maximo valor con peso $\le \text{rem}.
	\item \textbf{Eliminar pares dominados}: si un par (peso, valor) tiene otro con menos peso y mas valor, eliminarlo.
\end{itemize}

\paragraph{Trampas y errores frecuentes.}
\begin{itemize}\setlength\itemsep{0pt}
	\item $N$ impar: una mitad tiene $N/2 + 1$ items.
	\item Overflow al sumar pesos/valores en \texttt{long long}.
	\item Sin deduplicacion, el algoritmo es correcto pero mas lento.
\end{itemize}

\paragraph{Codigo.}
Ver \texttt{cpp.tex}, seccion \textit{Knapsack meet-in-the-middle}.

\vspace{0.5em}

\subsubsection{LIS O(n log n)}

\paragraph{Idea intuitiva.}
La \textbf{Longest Increasing Subsequence} se calcula en $O(n \log n)$ con el truco del \textbf{patience sorting}: mantener un arreglo \texttt{dp} donde \texttt{dp[len]} es el minimo ultimo valor de una IS de largo \texttt{len}. Para cada $x$, encontrar el primer \texttt{dp[len] >= x} y reemplazarlo. El largo final es la LIS. La demostracion: si hay dos IS del mismo largo, la que termina en un valor menor es ``mejor'' (puede extenderse mas facilmente).

\paragraph{Propiedades clave (invariantes).}
\begin{itemize}\setlength\itemsep{0pt}
	\item \texttt{dp} es estrictamente creciente.
	\item \texttt{lower\_bound} da LIS estricta; \texttt{upper\_bound} da no-estricta ($\le$).
	\item El algoritmo es \textbf{online}: procesa elementos en orden.
	\item La longitud de \texttt{dp} es la LIS al final.
\end{itemize}

\paragraph{Codigo clave.}
La version online:
\begin{verbatim}
vector<int> dp;
for (int x : a) {
    auto it = lower_bound(dp.begin(), dp.end(), x);
    if (it == dp.end()) dp.push_back(x);
    else *it = x;
}
// dp.size() = longitud de LIS estricta
\end{verbatim}

\paragraph{Complejidad.}
\begin{itemize}\setlength\itemsep{0pt}
	\item $O(n \log n)$.
	\item Memoria $O(n)$.
\end{itemize}

\paragraph{Donde tocar para modificar.}
\begin{itemize}\setlength\itemsep{0pt}
	\item \textbf{Reconstruir la LIS}: aniadir \texttt{parent[]} y \texttt{posInDp[]} para backtrack.
	\item \textbf{Contar LIS}: DP clasica $O(n^2)$ o segment tree sobre $dp$.
	\item \textbf{LDS / Longest non-increasing}: invertir la condicion en \texttt{lower\_bound} o multiplicar por -1.
	\item \textbf{Longest bitonic subsequence}: combinar LIS desde izq y desde der.
	\item \textbf{LIS en 2D}: ordenar por una dimension y aplicar LIS en la otra.
\end{itemize}

\paragraph{Trampas y errores frecuentes.}
\begin{itemize}\setlength\itemsep{0pt}
	\item \texttt{lower\_bound} vs \texttt{upper\_bound}: si tu comparacion es $\le$ vs $<$, cambia.
	\item Si los elementos son \texttt{int} pero podrian ser negativos, aniadir offset para usar como indices.
	\item Para LIS reconstruida, aniadir el array \texttt{parent} que apunta al anterior en la IS.
\end{itemize}

\paragraph{Codigo.}
Ver \texttt{cpp.tex}, seccion \textit{LIS O(n log n)}.

\vspace{0.5em}

\subsubsection{LCS / Edit distance}

\paragraph{Idea intuitiva.}
\textbf{LCS} (Longest Common Subsequence): para dos strings $a$, $b$, $dp[i][j]$ = LCS de $a[0..i)$ y $b[0..j)$. Recurrencia:
\[ dp[i][j] = \begin{cases} dp[i-1][j-1] + 1 & \text{si } a[i-1] = b[j-1] \\ \max(dp[i-1][j], dp[i][j-1]) & \text{si no} \end{cases} \]

\textbf{Edit distance} (Levenshtein): numero minimo de operaciones (insert, delete, replace) para convertir $a$ en $b$. Recurrencia similar con un $+1$ adicional. La idea: ``alinear'' los strings eligiendo matches/mismatches/insert/delete.

\paragraph{Propiedades clave (invariantes).}
\begin{itemize}\setlength\itemsep{0pt}
	\item LCS y Edit Distance son DP $O(nm)$.
	\item Memoria se puede reducir a $O(\min(n, m))$ con rolling array.
	\item Edit distance con peso por operacion distinto: ajustar el \texttt{+1} por el peso.
	\item La subestructura optima: la respuesta para prefijos se construye de prefijos menores.
\end{itemize}

\paragraph{Codigo clave.}
La recurrencia clasica de LCS:
\begin{verbatim}
vector<vector<int>> dp(n + 1, vector<int>(m + 1, 0));
for (int i = 1; i <= n; ++i)
    for (int j = 1; j <= m; ++j)
        if (a[i-1] == b[j-1]) dp[i][j] = dp[i-1][j-1] + 1;
        else dp[i][j] = max(dp[i-1][j], dp[i][j-1]);
return dp[n][m];
\end{verbatim}

\paragraph{Complejidad.}
\begin{itemize}\setlength\itemsep{0pt}
	\item $O(NM)$ tiempo y memoria.
	\item Para $N, M \le 5000$: $\sim 2.5 \cdot 10^7$ operaciones.
\end{itemize}

\paragraph{Donde tocar para modificar.}
\begin{itemize}\setlength\itemsep{0pt}
	\item \textbf{Reconstruir la LCS / las operaciones}: aniadir \texttt{prev[i][j]} para backtrack.
	\item \textbf{LCS de varios strings}: NP-hard para $\ge 3$ strings; en la practica se hace con bitmask.
	\item \textbf{Edit distance con costos distintos}: cambiar el \texttt{+1} por el costo apropiado.
	\item \textbf{Rolling array}: si memoria es problema, mantener solo 2 filas.
	\item \textbf{Hirschberg}: algoritmo $O(NM)$ tiempo, $O(\min(N,M))$ memoria, con reconstruccion.
\end{itemize}

\paragraph{Trampas y errores frecuentes.}
\begin{itemize}\setlength\itemsep{0pt}
	\item Indices en la recurrencia: \texttt{a[i-1]} y \texttt{b[j-1]} (no \texttt{a[i]}).
	\item Para reconstruccion, el caso \texttt{a[i-1] == b[j-1]} tiene prioridad.
	\item Si los strings son muy largos, considerar Hirschberg o bitset optimization.
\end{itemize}

\paragraph{Codigo.}
Ver \texttt{cpp.tex}, seccion \textit{LCS / Edit distance}.

\vspace{0.5em}

\subsubsection{Bitmask DP over subsets}

\paragraph{Idea intuitiva.}
Cuando el estado es un \textbf{subconjunto} de $\{0, \dots, N-1\}$, se representa como una mascara de $N$ bits. Para iterar sobre todos los subconjuntos de \texttt{mask}:
\[ \text{for (sub = mask; sub; sub = (sub - 1) \& mask)} \{ \dots \} \]
Para iterar sobre todos los superconjuntos de \texttt{mask} en $[0, U)$:
\[ \text{for (sup = mask; sup < U; sup = (sup + 1) | mask)} \{ \dots \} \]
La idea: las mascaras forman un lattice de subconjuntos; los trucos aritmeticos (\texttt{(sub-1) \& mask} y \texttt{(sup+1) | mask}) enumeran todo el lattice eficientemente.

\paragraph{Propiedades clave (invariantes).}
\begin{itemize}\setlength\itemsep{0pt}
	\item Iterar subsets: $O(2^k)$ para una mascara con $k$ bits.
	\item \texttt{sub \& mask == sub} siempre; \texttt{mask - sub} no siempre da un subconjunto.
	\item \texttt{\_\_builtin\_popcount(mask)} cuenta los bits.
	\item \texttt{(sub - 1) \& mask} itera subsets de \texttt{mask} en orden decreciente de inclusion.
\end{itemize}

\paragraph{Codigo clave.}
Iteracion sobre subsets de una mascara:
\begin{verbatim}
// Todos los subsets de mask
for (int sub = mask; sub; sub = (sub - 1) & mask) {
    // procesar sub
}
// Incluir sub = 0 si es necesario
sub = 0;  // procesar primero
for (int sup = mask; sup < (1 << N); sup = (sup + 1) | mask) {
    // procesar sup (todos los superconjuntos)
}
\end{verbatim}

\paragraph{Complejidad.}
\begin{itemize}\setlength\itemsep{0pt}
	\item $O(N \cdot 2^N)$ para una DP estandar sobre todos los subsets.
	\item Iterar subsets de \texttt{mask}: $O(2^{|mask|})$.
\end{itemize}

\paragraph{Donde tocar para modificar.}
\begin{itemize}\setlength\itemsep{0pt}
	\item \textbf{TSP / Hamiltonian path}: bitmask DP (ver tambien \textit{Hamiltonian path} en \texttt{cpp.tex} seccion Grafos).
	\item \textbf{Set cover}: iterar subsets.
	\item \textbf{DP con supermask}: si el estado incluye ``todos los elementos que \textbf{no} hemos usado'', usar supermask iteration.
	\item \textbf{Bitmasks grandes}: con $N > 20$, aniadir compresion o simulacion con segmentos.
\end{itemize}

\paragraph{Trampas y errores frecuentes.}
\begin{itemize}\setlength\itemsep{0pt}
	\item El orden de iteracion de subsets de \texttt{mask} es decreciente (de \texttt{mask} a $0$).
	\item \texttt{mask = 0} no tiene subsets no triviales; el bucle lo saltea correctamente.
	\item El truco \texttt{(sub - 1) \& mask} requiere cuidado con \texttt{sub = 0} (termina el bucle).
	\item Superconjuntos: \texttt{sup < (1 << N)} debe estar bien definido.
\end{itemize}

\paragraph{Codigo.}
Ver \texttt{cpp.tex}, seccion \textit{Bitmask DP over subsets}.

\vspace{0.5em}

\subsubsection{SOS DP}

\paragraph{Idea intuitiva.}
Para calcular, para cada \texttt{mask}, la suma de $f(\text{submask})$ sobre todos los $\text{submask} \subseteq \text{mask}$:
\[ dp[mask] = \sum_{\text{sub} \subseteq mask} f(\text{sub}). \]
Se computa en $O(N \cdot 2^N)$ mediante un loop por bit: para cada $i$, sumar a $dp[mask]$ la contribucion de $dp[mask \oplus (1 << i)]$ cuando $mask$ tiene el bit $i$. La idea: cada subconjunto $s \subseteq m$ se obtiene quitando bits uno a uno; asi, sumar incrementalmente captura todos los subconjuntos en $O(N \cdot 2^N)$.

\paragraph{Propiedades clave (invariantes).}
\begin{itemize}\setlength\itemsep{0pt}
	\item Inicializar $dp = f$.
	\item Para cada bit $i$: \texttt{if (mask \& (1 << i)) dp[mask] += dp[mask \^{} (1 << i)]}.
	\item Variantes: supersets DP (sumar sobre superconjuntos).
	\item El orden de iteracion (bit externo, mascara interno) es crucial.
\end{itemize}

\paragraph{Codigo clave.}
SOS DP (Sum over Subsets):
\begin{verbatim}
vector<long long> dp(1 << N);
for (int i = 0; i < (1 << N); ++i) dp[i] = f[i];
for (int bit = 0; bit < N; ++bit)
    for (int mask = 0; mask < (1 << N); ++mask)
        if (mask & (1 << bit))
            dp[mask] += dp[mask ^ (1 << bit)];
// dp[mask] = suma de f(sub) para sub ⊆ mask
\end{verbatim}

\paragraph{Complejidad.}
\begin{itemize}\setlength\itemsep{0pt}
	\item $O(N \cdot 2^N)$.
	\item Memoria $O(2^N)$.
\end{itemize}

\paragraph{Donde tocar para modificar.}
\begin{itemize}\setlength\itemsep{0pt}
	\item \textbf{Max/min en lugar de suma}: cambiar \texttt{+=} por la operacion correspondiente.
	\item \textbf{Supersets DP}: iterar \texttt{mask} de $0$ a $2^N$ y aniadir \texttt{dp[mask] += dp[mask | (1<<i)]}.
	\item \textbf{XOR convolution} (subset convolution): version mas avanzada.
	\item \textbf{Aplicaciones}: contar subconjuntos con propiedades especificas.
\end{itemize}

\paragraph{Trampas y errores frecuentes.}
\begin{itemize}\setlength\itemsep{0pt}
	\item Si $N$ es 30 o mas, $2^N$ no entra en memoria; usar trucos o reducir $N$.
	\item El sentido del XOR (sumar a \texttt{mask \^{} bit}) es para ``agregar el bit $i$ al subconjunto''.
	\item En supersets DP, la condicion es \texttt{(mask \& (1<<bit)) == 0}.
\end{itemize}

\paragraph{Codigo.}
Ver \texttt{cpp.tex}, seccion \textit{SOS DP}.

\vspace{0.5em}

\subsubsection{Digit DP template}

\paragraph{Idea intuitiva.}
Contar numeros en $[0, N]$ que cumplen una propiedad sobre sus digitos. La idea: DP por posicion, manteniendo:
\begin{itemize}\setlength\itemsep{0pt}
	\item \texttt{pos}: posicion actual (de izquierda a derecha).
	\item \texttt{tight}: si los digitos ya colocados son exactamente los de $N$ (limite superior estricto).
	\item \texttt{state}: estado adicional (digitos usados, suma, etc.).
\end{itemize}
La clave: la condicion \texttt{tight} permite acotar correctamente el siguiente digito; sin ella, la DP contaria incorrectamente.

\paragraph{Propiedades clave (invariantes).}
\begin{itemize}\setlength\itemsep{0pt}
	\item Si \texttt{tight} es true, el siguiente digito esta acotado por el de $N$; si no, puede ser 0-9.
	\item Memoria: $O(D \cdot 2 \cdot \text{states})$.
	\item Caso base: si \texttt{pos == D}, devolver 1 si el estado cumple la propiedad, sino 0.
	\item El bit \texttt{started} distingue ``numero con menos digitos'' (con leading zeros).
\end{itemize}

\paragraph{Codigo clave.}
La estructura basica:
\begin{verbatim}
long long dp(int pos, bool tight, State s) {
    if (pos == D) return check(s) ? 1 : 0;
    if (memo[pos][tight][s] != -1) return memo;
    int limit = tight ? digits[pos] : 9;
    long long ans = 0;
    for (int d = 0; d <= limit; ++d)
        ans += dp(pos + 1, tight && (d == limit), update(s, d, pos));
    return memo[pos][tight][s] = ans;
}
\end{verbatim}

\paragraph{Complejidad.}
\begin{itemize}\setlength\itemsep{0pt}
	\item $O(D \cdot 10 \cdot \text{states})$ donde $D$ es el numero de digitos.
	\item Para $D \le 19$ (hasta $10^{18}$): $\sim 200 \cdot \text{states}$.
\end{itemize}

\paragraph{Donde tocar para modificar.}
\begin{itemize}\setlength\itemsep{0pt}
	\item \textbf{Cambiar la propiedad}: modificar \texttt{updateState} y la condicion base.
	\item \textbf{Multiples estados}: aniadir mas dimensiones al \texttt{memo}.
	\item \textbf{Leading zeros}: aniadir un flag \texttt{started} para distinguir numeros con menos digitos.
	\item \textbf{Rango $[L, R]$}: calcular $f(R) - f(L-1)$.
\end{itemize}

\paragraph{Trampas y errores frecuentes.}
\begin{itemize}\setlength\itemsep{0pt}
	\item Olvidar resetear \texttt{memo} entre queries.
	\item Si la cantidad de estados es grande, el \texttt{memo[pos][tight][state]} puede ser muy grande; usar \texttt{map}.
	\item La condicion \texttt{started} es crucial para numeros con menos digitos que $D$.
	\item Si $N$ es muy grande ($> 10^{18}$), aniadir soporte para long long.
\end{itemize}

\paragraph{Codigo.}
Ver \texttt{cpp.tex}, seccion \textit{Digit DP template}.

\vspace{0.5em}

\subsubsection{Tree DP (reroot)}

\paragraph{Idea intuitiva.}
Calcular alguna propiedad asumiendo que \textbf{cada nodo es la raiz} del arbol. La tecnica clasica de \textbf{rerooting} en dos pasadas:
\begin{itemize}\setlength\itemsep{0pt}
	\item \textbf{DFS 1}: desde una raiz arbitraria, calcular \texttt{dp[v]} para el subarbol de $v$.
	\item \textbf{DFS 2}: para cada hijo $u$ de $v$, calcular \texttt{up[u]} = respuesta si la raiz fuera $u$.
\end{itemize}
La demostracion: tras la primera DFS, cada nodo conoce ``lo que aporta su subarbol''; en la segunda, se transmite ``lo que aporta el resto del arbol''.

\paragraph{Propiedades clave (invariantes).}
\begin{itemize}\setlength\itemsep{0pt}
	\item \texttt{dp[v]} depende solo del subarbol de $v$ (con la raiz original).
	\item \texttt{up[u]} se calcula de manera que cuando la raiz es $u$, ``todo lo que estaba en $v$ se ajusta''.
	\item Formula para reroot depende del problema; el snippet usa el ejemplo ``suma de distancias''.
	\item La combinacion \texttt{dp[u] + up[u]} da la respuesta con raiz $u$.
\end{itemize}

\paragraph{Codigo clave.}
Reroot para suma de distancias:
\begin{verbatim}
// dfs1: dp[v] = suma de distancias desde v al subarbol
void dfs1(int v, int p) {
    dp[v] = 0;
    for (int u : adj[v]) if (u != p) {
        dfs1(u, v);
        dp[v] += dp[u] + sz[u];  // cada hijo aporta su subarbol
    }
}
// dfs2: up[u] = distancia desde el resto del arbol a u
void dfs2(int v, int p) {
    for (int u : adj[v]) if (u != p) {
        up[u] = up[v] + dp[v] - dp[u] - sz[u] + (n - sz[u]);
        dfs2(u, v);
    }
}
\end{verbatim}

\paragraph{Complejidad.}
\begin{itemize}\setlength\itemsep{0pt}
	\item $O(N)$ para ambas pasadas.
	\item Cada arista se procesa un numero constante de veces.
\end{itemize}

\paragraph{Donde tocar para modificar.}
\begin{itemize}\setlength\itemsep{0pt}
	\item \textbf{Cambiar la cantidad a calcular}: modificar la operacion en \texttt{dfs1} y la formula en \texttt{dfs2}.
	\item \textbf{Arbol con pesos}: aniadir \texttt{weights[v]} y ajustar.
	\item \textbf{Re-root en arboles con ciclos}: requiere otra tecnica (DSU on tree).
	\item \textbf{Rooted trees diferentes}: si la propiedad no es invariante por raiz, usar otra tecnica.
\end{itemize}

\paragraph{Trampas y errores frecuentes.}
\begin{itemize}\setlength\itemsep{0pt}
	\item La formula de \texttt{up[u]} debe sumar y restar correctamente las contribuciones de $u$ y del resto.
	\item Para arboles con pesos, aniadir el peso en las formulas.
	\item Verificar con ejemplos pequenos antes de generalizar.
\end{itemize}

\paragraph{Codigo.}
Ver \texttt{cpp.tex}, seccion \textit{Tree DP (reroot)}.

% ============================================================
\section{Geometria}
% ============================================================

\subsection{Puntos y segmentos}

\subsubsection{Point + cross + dot}

\paragraph{Idea intuitiva.}
La estructura \texttt{Point} representa un punto en 2D con coordenadas enteras. Las dos operaciones vectoriales clave son:
\begin{itemize}\setlength\itemsep{0pt}
	\item \textbf{Cross product} $P \times Q = P_x Q_y - P_y Q_x$. Da el ``area con signo'' del paralelogramo formado por $P$ y $Q$. Usado para orientacion.
	\item \textbf{Dot product} $P \cdot Q = P_x Q_x + P_y Q_y$. Da el producto de las longitudes por el coseno del angulo. Usado para proyeccion.
\end{itemize}
La interpretacion geometrica del cross product: es el area con signo del paralelogramo; el signo indica la orientacion relativa.

\paragraph{Propiedades clave (invariantes).}
\begin{itemize}\setlength\itemsep{0pt}
	\item \texttt{cross(a, b) > 0} si $b$ esta a la izquierda de $a$ (sistema de coordenadas usual).
	\item \texttt{cross(a, b) = 0} si son colineales.
	\item \texttt{dot(a, b) > 0} si el angulo es agudo.
	\item El snippet define \texttt{cross} respecto a \texttt{this}, util para orientacion de 3 puntos.
	\item $|cross|^2 + |dot|^2 = |P|^2 \cdot |Q|^2$ (identidad de Lagrange).
\end{itemize}

\paragraph{Codigo clave.}
Las dos operaciones vectoriales:
\begin{verbatim}
long long cross(const Point& a, const Point& b) {
    return (long long)a.x * b.y - (long long)a.y * b.x;
}
long long dot(const Point& a, const Point& b) {
    return (long long)a.x * b.x + (long long)a.y * b.y;
}
\end{verbatim}

\paragraph{Complejidad.}
\begin{itemize}\setlength\itemsep{0pt}
	\item $O(1)$ por operacion.
	\item Constante muy pequena (4 multiplicaciones + 2 sumas).
\end{itemize}

\paragraph{Donde tocar para modificar.}
\begin{itemize}\setlength\itemsep{0pt}
	\item \textbf{Coordenadas doubles}: cambiar \texttt{int} por \texttt{double} o \texttt{long double}. Cuidado con la precision.
	\item \textbf{3D}: aniadir \texttt{z} y adaptar cross/dot.
	\item \textbf{Operadores adicionales}: aniadir \texttt{operator+}, \texttt{operator-}, \texttt{operator*} (escalar).
	\item \textbf{Distancia al cuadrado}: aniadir \texttt{len2()} que retorna $x^2 + y^2$.
\end{itemize}

\paragraph{Trampas y errores frecuentes.}
\begin{itemize}\setlength\itemsep{0pt}
	\item Cross product \textbf{no es} conmutativo; $P \times Q = -Q \times P$.
	\item Overflow con coordenadas grandes ($|x|, |y| > 10^4$) si se hace \texttt{x * y}.
	\item El signo de cross depende del sistema de coordenadas (en computacion, usualmente CCW positivo).
	\item Para ``igualdad'' de doubles, aniadir tolerancia \texttt{eps}.
\end{itemize}

\paragraph{Codigo.}
Ver \texttt{cpp.tex}, seccion \textit{Point + cross + dot}.

\vspace{0.5em}

\subsubsection{Interseccion de segmentos}

\paragraph{Idea intuitiva.}
Dados dos segmentos $[a, b]$ y $[c, d]$, determinar si se intersectan. La idea clasica:
\begin{itemize}\setlength\itemsep{0pt}
	\item Verificar orientaciones: $a, b, c$ y $a, b, d$ deben estar en lados opuestos de la recta $ab$; igual para $cd$ y $ab$.
	\item Caso colineal: si las rectas son paralelas (\texttt{cross} $= 0$) y los segmentos son colineales, verificar que los bounding boxes se intersectan.
\end{itemize}
La demostracion: dos segmentos se intersectan sii cada segmento ``cruza'' la linea que contiene al otro (estricto) o ambos son colineales con solapamiento.

\paragraph{Propiedades clave (invariantes).}
\begin{itemize}\setlength\itemsep{0pt}
	\item Funciona con segmentos en cualquier orientacion.
	\item Caso colineal: requiere check adicional de bounding box.
	\item Para intersectar en un \textbf{punto} especifico: aniadir division con cuidado de precision.
	\item Funciona con segmentos verticales/horizontales sin casos especiales.
\end{itemize}

\paragraph{Codigo clave.}
Test clasico de interseccion:
\begin{verbatim}
int sign(long long x) { return (x > 0) - (x < 0); }
bool intersect(Point a, Point b, Point c, Point d) {
    long long o1 = cross(b - a, c - a);
    long long o2 = cross(b - a, d - a);
    long long o3 = cross(d - c, a - c);
    long long o4 = cross(d - c, b - c);
    if (sign(o1) != sign(o2) && sign(o3) != sign(o4)) return true;
    // casos colineales
    if (o1 == 0 && onSegment(a, b, c)) return true;
    if (o2 == 0 && onSegment(a, b, d)) return true;
    if (o3 == 0 && onSegment(c, d, a)) return true;
    if (o4 == 0 && onSegment(c, d, b)) return true;
    return false;
}
\end{verbatim}

\paragraph{Complejidad.}
\begin{itemize}\setlength\itemsep{0pt}
	\item $O(1)$.
	\item Constante pequena.
\end{itemize}

\paragraph{Donde tocar para modificar.}
\begin{itemize}\setlength\itemsep{0pt}
	\item \textbf{Interseccion con punto}: aniadir la division despues de verificar orientacion.
	\item \textbf{Interseccion rayo-segmento}: ajustar la primera condicion.
	\item \textbf{Poligonos convexos}: usar separacion de ejes (SAT).
	\item \textbf{Multiples queries}: precomputar bounding boxes o usar segment tree.
\end{itemize}

\paragraph{Trampas y errores frecuentes.}
\begin{itemize}\setlength\itemsep{0pt}
	\item Si los segmentos tocan solo en un extremo, el snippet lo considera interseccion (algunos problemas quieren estricto). Ajustar segun el problema.
	\item El caso colineal es comun; olvidar el check adicional de bounding box da falsos negativos.
	\item En problemas con muchos segmentos, considerar sweep line para $O(n \log n)$.
\end{itemize}

\paragraph{Codigo.}
Ver \texttt{cpp.tex}, seccion \textit{Interseccion de segmentos}.

\vspace{0.5em}

\subsubsection{Contains colineal}

\paragraph{Idea intuitiva.}
Verifica si un punto $c$ esta sobre el segmento $[a, b]$ (asumiendo que ya son colineales). Si \texttt{cross(a, b, c) = 0} y $c$ esta dentro del bounding box de $[a, b]$, entonces $c$ esta en el segmento. La razon: si los tres puntos son colineales, $c$ esta en el segmento sii sus coordenadas estan entre las de $a$ y $b$.

\paragraph{Propiedades clave (invariantes).}
\begin{itemize}\setlength\itemsep{0pt}
	\item Asume colinealidad (no la verifica mas alla de \texttt{cross == 0}).
	\item Usa comparaciones con \texttt{<=} para extremos \textbf{inclusivos}.
	\item Funciona con segmentos verticales (la condicion se reduce a una sola coordenada).
\end{itemize}

\paragraph{Codigo clave.}
Test de contencion en segmento (asume colinealidad):
\begin{verbatim}
bool onSegment(Point a, Point b, Point c) {
    return min(a.x, b.x) <= c.x && c.x <= max(a.x, b.x) &&
           min(a.y, b.y) <= c.y && c.y <= max(a.y, b.y);
}
\end{verbatim}

\paragraph{Complejidad.}
\begin{itemize}\setlength\itemsep{0pt}
	\item $O(1)$.
	\item Constante minima.
\end{itemize}

\paragraph{Donde tocar para modificar.}
\begin{itemize}\setlength\itemsep{0pt}
	\item \textbf{Contener estricto}: cambiar \texttt{<=} a \texttt{<} para excluir extremos.
	\item \textbf{Coordenadas doubles}: usar tolerancia (\texttt{eps}).
	\item \textbf{Solo x o solo y}: aniadir check especifico para segmentos verticales/horizontales.
\end{itemize}

\paragraph{Trampas y errores frecuentes.}
\begin{itemize}\setlength\itemsep{0pt}
	\item No usar este snippet para puntos \textbf{no colineales}; siempre verificar primero la colinealidad.
	\item Si los puntos tienen doubles, el chequeo \texttt{cross == 0} puede fallar; aniadir tolerancia.
	\item El bounding box incluye los extremos; verificar si el problema quiere exclusivo o inclusivo.
\end{itemize}

\paragraph{Codigo.}
Ver \texttt{cpp.tex}, seccion \textit{Contains colineal}.

\subsection{Poligonos}

\subsubsection{Convex Hull (monotono)}

\paragraph{Idea intuitiva.}
Algoritmo de \textbf{Andrew's monotone chain}: ordenar puntos por $(x, y)$, luego construir la hull inferior y superior por separado. Para cada punto, aniadirlo al hull y, mientras el ultimo giro no sea ``counter-clockwise'' (\texttt{cross > 0}), popear. La demostracion: si los puntos se procesan en orden, el hull inferior y superior cubren la frontera completa sin auto-intersecciones.

\paragraph{Propiedades clave (invariantes).}
\begin{itemize}\setlength\itemsep{0pt}
	\item Ordena los puntos en $O(n \log n)$.
	\item Construye la hull en $O(n)$ despues.
	\item Output: poligono convexo en counter-clockwise, sin punto final duplicado.
	\item El \texttt{<= 0} incluye colineales; usar \texttt{< 0} para excluirlos.
	\item Cada punto aparece en la hull inferior o superior (no ambos).
\end{itemize}

\paragraph{Codigo clave.}
Construccion del hull inferior y superior:
\begin{verbatim}
sort(p.begin(), p.end());
vector<Point> lower, upper;
for (auto& p : points) {
    while (lower.size() >= 2 && cross(lower.back() - lower[lower.size()-2],
                                       p - lower.back()) <= 0)
        lower.pop_back();
    lower.push_back(p);
}
for (int i = points.size() - 1; i >= 0; --i) {
    auto& p = points[i];
    while (upper.size() >= 2 && cross(upper.back() - upper[upper.size()-2],
                                       p - upper.back()) <= 0)
        upper.pop_back();
    upper.push_back(p);
}
vector<Point> hull = lower;
hull.insert(hull.end(), upper.begin() + 1, upper.end() - 1);
\end{verbatim}

\paragraph{Complejidad.}
\begin{itemize}\setlength\itemsep{0pt}
	\item $O(n \log n)$.
	\item Constante muy pequena (solo cross products).
\end{itemize}

\paragraph{Donde tocar para modificar.}
\begin{itemize}\setlength\itemsep{0pt}
	\item \textbf{Colineales}: cambiar \texttt{<= 0} por \texttt{< 0} para excluirlos del hull.
	\item \textbf{Puntos repetidos}: \texttt{unique} despues de sort.
	\item \textbf{Hull 3D}: gift wrapping es $O(n^2)$; convex hull 3D es mas complejo.
	\item \textbf{Output con los puntos}: aniadir el hull completo (no solo upper/lower).
\end{itemize}

\paragraph{Trampas y errores frecuentes.}
\begin{itemize}\setlength\itemsep{0pt}
	\item El \texttt{pop\_back} y el \texttt{reverse} son esenciales para no duplicar el ultimo punto.
	\item Para excluir colineales, ajustar el \texttt{<= 0} vs \texttt{< 0}.
	\item Si los puntos estan duplicados, aniadir \texttt{unique} antes del sort.
\end{itemize}

\paragraph{Codigo.}
Ver \texttt{cpp.tex}, seccion \textit{Convex Hull (monotono)}.

\vspace{0.5em}

\subsubsection{Polygon Area}

\paragraph{Idea intuitiva.}
El area (con signo) de un poligono se calcula con la \textbf{Shoelace formula}: $A = \frac{1}{2} \sum_{i=0}^{n-1} (x_i y_{i+1} - x_{i+1} y_i)$, donde el indice es modulo $n$. El snippet retorna el \textbf{doble} del area (en valor absoluto). La formula sale del ``sumar areas de trapezoides'': cada lado del poligono define un trapezoide con el eje $x$; la suma con signo da el area.

\paragraph{Propiedades clave (invariantes).}
\begin{itemize}\setlength\itemsep{0pt}
	\item El signo indica la orientacion: positivo si counter-clockwise.
	\item Para el doble: no dividir por 2 si no es necesario (evita fracciones).
	\item Funciona para poligonos simples (no auto-intersectantes).
	\item Si el poligono es no convexo, la formula da el area ``con signo''.
\end{itemize}

\paragraph{Codigo clave.}
La suma de trapezoides:
\begin{verbatim}
long long shoelace(const vector<Point>& p) {
    long long s = 0;
    int n = p.size();
    for (int i = 0; i < n; ++i)
        s += (long long)p[i].x * p[(i+1)%n].y - (long long)p[(i+1)%n].x * p[i].y;
    return abs(s);  // doble del area
}
\end{verbatim}

\paragraph{Complejidad.}
\begin{itemize}\setlength\itemsep{0pt}
	\item $O(n)$.
	\item Constante minima.
\end{itemize}

\paragraph{Donde tocar para modificar.}
\begin{itemize}\setlength\itemsep{0pt}
	\item \textbf{Area exacta}: dividir por 2 (con cuidado si da fraccion; usar \texttt{double}).
	\item \textbf{Area firmada}: quitar \texttt{abs} y dejar el signo.
	\item \textbf{Punto dentro de poligono convexo}: con la hull, chequear que esta a un mismo lado de todas las aristas.
	\item \textbf{Triangulacion}: descomponer en triangulos y sumar areas.
\end{itemize}

\paragraph{Trampas y errores frecuentes.}
\begin{itemize}\setlength\itemsep{0pt}
	\item Si los puntos no estan en orden (CCW o CW), el area firmada puede ser incorrecta. Para signed area consistente, usar siempre el mismo orden.
	\item Para poligonos con holes, sumar las areas con signo apropiado.
	\item Con doubles, aniadir tolerancia \texttt{eps}.
\end{itemize}

\paragraph{Codigo.}
Ver \texttt{cpp.tex}, seccion \textit{Polygon Area}.

\vspace{0.5em}

\subsubsection{Distance point-segment / point-line}

\paragraph{Idea intuitiva.}
Distancia minima de un punto $p$ a un segmento $[a, b]$ o a la recta que pasa por $a$ y $b$.
\begin{itemize}\setlength\itemsep{0pt}
	\item \textbf{Linea}: $d = \frac{|ab \times ap|}{|ab|}$.
	\item \textbf{Segmento}: si la proyeccion de $p$ sobre $ab$ cae fuera del segmento, la distancia es al extremo mas cercano.
\end{itemize}
La razon geometrica: la distancia a una linea es el area del paralelogramo dividido por la base; la proyeccion ortogonal indica si estamos ``dentro'' o ``fuera'' del segmento.

\paragraph{Propiedades clave (invariantes).}
\begin{itemize}\setlength\itemsep{0pt}
	\item \texttt{dist2PointLine} retorna distancia \textbf{al cuadrado}.
	\item \texttt{dist2PointSegment} distingue 3 casos: $p$ antes de $a$, despues de $b$, o proyectado sobre el segmento.
	\item Si \texttt{len2 = 0} (degenerate), la distancia es $|ap|$ o $|bp|$.
	\item La proyeccion se calcula con dot products y el parametro $t \in [0, 1]$.
\end{itemize}

\paragraph{Codigo clave.}
Distancia al cuadrado a segmento:
\begin{verbatim}
long long dist2Segment(Point p, Point a, Point b) {
    long long len2 = (a-b).len2();
    if (len2 == 0) return (p-a).len2();  // degenerate
    long long t = max(0LL, min(1LL, dot(p-a, b-a) / len2));
    Point proj = {a.x + t*(b.x - a.x), a.y + t*(b.y - a.y)};
    return (p - proj).len2();
}
\end{verbatim}

\paragraph{Complejidad.}
\begin{itemize}\setlength\itemsep{0pt}
	\item $O(1)$.
	\item Constante minima.
\end{itemize}

\paragraph{Donde tocar para modificar.}
\begin{itemize}\setlength\itemsep{0pt}
	\item \textbf{Distancia (no al cuadrado)}: aniadir \texttt{sqrt} al final.
	\item \textbf{Coordenadas doubles}: aniadir \texttt{long double} y tolerancia.
	\item \textbf{3D}: aniadir coordenada $z$ y adaptar.
	\item \textbf{Multiples queries}: usar KD-tree o segment tree para acelerar.
\end{itemize}

\paragraph{Trampas y errores frecuentes.}
\begin{itemize}\setlength\itemsep{0pt}
	\item Overflow en \texttt{area2 * area2} si el area es grande; \texttt{long long} alcanza para $|x|, |y| \le 10^9$ pero no mas.
	\item Si $a = b$ (segmento degenerado), aniadir check explicito.
	\item En doubles, aniadir tolerancia \texttt{eps} para comparar.
\end{itemize}

\paragraph{Codigo.}
Ver \texttt{cpp.tex}, seccion \textit{Distance point-segment / point-line}.

\vspace{0.5em}

\subsubsection{Projection of point onto line}

\paragraph{Idea intuitiva.}
Proyeccion ortogonal de $p$ sobre la recta que pasa por $a$ y $b$. El parametro $t$ tal que la proyeccion es $a + t \cdot (b - a)$ es:
\[ t = \frac{(p - a) \cdot (b - a)}{|b - a|^2}. \]
La idea: descomponemos $\vec{ap}$ en una componente paralela a $\vec{ab}$ y una perpendicular; $t$ mide la posicion en la primera.

\paragraph{Propiedades clave (invariantes).}
\begin{itemize}\setlength\itemsep{0pt}
	\item Si $t \in [0, 1]$, la proyeccion cae sobre el \textbf{segmento}; si no, cae sobre la extension.
	\item Retorna \texttt{double} (puede tener precision issues).
	\item $t < 0$: $p$ esta ``antes'' de $a$; $t > 1$: $p$ esta ``despues'' de $b$.
\end{itemize}

\paragraph{Codigo clave.}
Proyeccion ortogonal:
\begin{verbatim}
Point projection(Point p, Point a, Point b) {
    Point ab = b - a, ap = p - a;
    double t = (double)dot(ap, ab) / ab.len2();
    return {a.x + t * ab.x, a.y + t * ab.y};
}
\end{verbatim}

\paragraph{Complejidad.}
\begin{itemize}\setlength\itemsep{0pt}
	\item $O(1)$.
	\item Constante minima (dos dot products y una division).
\end{itemize}

\paragraph{Donde tocar para modificar.}
\begin{itemize}\setlength\itemsep{0pt}
	\item \textbf{Proyeccion sobre segmento}: clipar $t$ a $[0, 1]$ antes de calcular el punto.
	\item \textbf{Distancia al segmento via proyeccion}: combinar con el modulo anterior.
	\item \textbf{Solo el parametro $t$}: si solo necesitas saber si esta dentro del segmento, retorna $t$ sin construir el punto.
\end{itemize}

\paragraph{Trampas y errores frecuentes.}
\begin{itemize}\setlength\itemsep{0pt}
	\item Si $a = b$ ($|ab| = 0$), la proyeccion no esta definida; aniadir guard.
	\item En doubles, aniadir tolerancia \texttt{eps} para comparar $t$ con 0 o 1.
	\item La division por \texttt{len2} puede ser cara si las coordenadas son grandes; usar \texttt{double} o precomputar.
\end{itemize}

\paragraph{Codigo.}
Ver \texttt{cpp.tex}, seccion \textit{Projection of point onto line}.

\vspace{0.5em}

\subsubsection{Point in polygon (ray casting)}

\paragraph{Idea intuitiva.}
Algoritmo de \textbf{ray casting}: contar cuantas veces un rayo horizontal desde $p$ cruza los bordes del poligono. Si es impar, $p$ esta dentro; si par, fuera. El snippet implementa la version clasica recorriendo aristas. La razon: cada cruce cambia el lado del rayo respecto al poligono; un numero impar de cambios significa ``terminamos dentro''.

\paragraph{Propiedades clave (invariantes).}
\begin{itemize}\setlength\itemsep{0pt}
	\item Funciona con poligonos simples (no necesariamente convexos).
	\item Las intersecciones en vertices cuentan especialmente.
	\item Para poligonos con \textbf{huecos}, sigue funcionando.
	\item La respuesta es estrictamente 0 (fuera) o 1 (dentro) para poligonos simples.
\end{itemize}

\paragraph{Codigo clave.}
El test clasico de ray casting:
\begin{verbatim}
bool pointInPolygon(Point p, const vector<Point>& poly) {
    int n = poly.size();
    bool inside = false;
    for (int i = 0, j = n - 1; i < n; j = i++) {
        if ((poly[i].y > p.y) != (poly[j].y > p.y) &&
            p.x < (poly[j].x - poly[i].x) * (p.y - poly[i].y) /
                    (poly[j].y - poly[i].y) + poly[i].x)
            inside = !inside;
    }
    return inside;
}
\end{verbatim}

\paragraph{Complejidad.}
\begin{itemize}\setlength\itemsep{0pt}
	\item $O(n)$ por query.
	\item Memoria $O(1)$ (no se almacena el rayo).
\end{itemize}

\paragraph{Donde tocar para modificar.}
\begin{itemize}\setlength\itemsep{0pt}
	\item \textbf{Winding number}: version alternativa que cuenta el numero de vueltas.
	\item \textbf{Poligono convexo}: check lineal contra las aristas (mas rapido).
	\item \textbf{Multiples queries}: precomputar o usar segment tree sobre las aristas.
	\item \textbf{Punto en frontera}: aniadir check adicional si es colineal con alguna arista.
\end{itemize}

\paragraph{Trampas y errores frecuentes.}
\begin{itemize}\setlength\itemsep{0pt}
	\item El manejo del vertice (cuando el rayo pasa por un vertice) puede causar double-counting; el snippet lo evita con la condicion \texttt{(poly[i].y > p.y) != (poly[j].y > p.y)}.
	\item Si el poligono tiene holes, el algoritmo sigue funcionando.
	\item Para poligonos no simples (auto-intersectantes), el resultado puede ser ambiguo.
\end{itemize}

\paragraph{Codigo.}
Ver \texttt{cpp.tex}, seccion \textit{Point in polygon (ray casting)}.

\vspace{0.5em}

\subsubsection{Minkowski sum (convex polygons)}

\paragraph{Idea intuitiva.}
La \textbf{suma de Minkowski} $A \oplus B$ de dos poligonos convexos es otro poligono convexo que representa $\{a + b : a \in A, b \in B\}$. Algoritmo: ordenar las aristas por angulo polar (CCW) y ``mergearlas'' como en merge sort; el resultado es un poligono CCW sin repetir el ultimo punto. La idea: el resultado tiene una arista paralela a cada arista de los originales; la suma se obtiene barriendo los angulos.

\paragraph{Propiedades clave (invariantes).}
\begin{itemize}\setlength\itemsep{0pt}
	\item Si $A, B$ son convexos con $n, m$ vertices, $A \oplus B$ tiene $\le n + m$ vertices.
	\item Util para \textbf{detectar colision}: $A$ colisiona con $B$ si $0 \in A \oplus (-B)$.
	\item Asume poligonos en orden \textbf{counter-clockwise}.
	\item El resultado es convexo (suma de convexos es convexo).
\end{itemize}

\paragraph{Codigo clave.}
La suma de Minkowski:
\begin{verbatim}
vector<Point> minkowski(const vector<Point>& A, const vector<Point>& B) {
    vector<Point> result;
    int i = 0, j = 0;
    int n = A.size(), m = B.size();
    do {
        result.push_back(A[i] + B[j]);
        long long cross = (A[(i+1)%n] - A[i]) ^ (B[(j+1)%m] - B[j]);
        if (cross >= 0) i = (i + 1) % n;
        if (cross <= 0) j = (j + 1) % m;
    } while (i || j);
    return result;
}
\end{verbatim}

\paragraph{Complejidad.}
\begin{itemize}\setlength\itemsep{0pt}
	\item $O(n + m)$.
	\item Memoria $O(n + m)$.
\end{itemize}

\paragraph{Donde tocar para modificar.}
\begin{itemize}\setlength\itemsep{0pt}
	\item \textbf{Deteccion de colision}: aniadir $(-B)$ (reflejar) y verificar si $0$ esta en el resultado.
	\item \textbf{Poligonos no convexos}: la suma puede no ser poligono simple; algoritmo distinto.
	\item \textbf{Multiples queries}: precomputar las aristas una vez.
	\item \textbf{Robot motion planning}: usar para planear trayectorias.
\end{itemize}

\paragraph{Trampas y errores frecuentes.}
\begin{itemize}\setlength\itemsep{0pt}
	\item Si los poligonos no son CCW, la suma es incorrecta. Normalizar antes.
	\item El bucle termina cuando ambos indices vuelven a 0; aniadir check.
	\item Para la deteccion de colision, aniadir el test \texttt{0 in A \^{} (-B)}.
\end{itemize}

\paragraph{Codigo.}
Ver \texttt{cpp.tex}, seccion \textit{Minkowski sum (convex polygons)}.

\subsection{Herramientas}

\subsubsection{Li Chao Tree (long long)}

\paragraph{Idea intuitiva.}
Mantiene un \textbf{lower envelope} (o upper) de un conjunto de rectas $y = mx + b$ sobre un dominio $[X_{LO}, X_{HI}]$ y responde ``minimo $y$ en un $x$ dado'' en $O(\log X)$. La estructura es un segment tree donde cada nodo guarda la recta ``mejor'' para su segmento medio; al aniadir una recta, se compara con la actual y se decide cual es mejor en cada mitad, descendiendo recursivamente. La demostracion: en cada nivel, solo se guarda la mejor recta para el ``punto medio''; las otras se transmiten a los hijos.

\paragraph{Propiedades clave (invariantes).}
\begin{itemize}\setlength\itemsep{0pt}
	\item \textbf{Online}: cada \texttt{addLine} en $O(\log X)$.
	\item Solo insercion; no hay borrar (para borrar, usar Li Chao segmentado con reset).
	\item La recta guardada en un nodo es la ``mejor'' en su punto medio.
	\item Las rectas ``perdedoras'' se transmiten a los hijos para futuros inserts.
\end{itemize}

\paragraph{Codigo clave.}
La insercion en Li Chao:
\begin{verbatim}
void addLine(Line nw, int v = 1, int l = X_LO, int r = X_HI) {
    int m = (l + r) / 2;
    bool left = nw.get(l) < tree[v].get(l);
    bool mid = nw.get(m) < tree[v].get(m);
    if (mid) swap(tree[v], nw);
    if (r - l == 1) return;
    if (left != mid) addLine(nw, v*2, l, m);
    else             addLine(nw, v*2+1, m, r);
}
\end{verbatim}

\paragraph{Complejidad.}
\begin{itemize}\setlength\itemsep{0pt}
	\item $O(\log X)$ por insercion y query, donde $X$ es el tamano del dominio.
	\item Memoria $O(\log X)$ por insercion, $O(X)$ total.
\end{itemize}

\paragraph{Donde tocar para modificar.}
\begin{itemize}\setlength\itemsep{0pt}
	\item \textbf{Upper envelope en vez de lower}: invertir el signo (\texttt{>} en vez de \texttt{<}).
	\item \textbf{Queries de segmento} (no full line): si la recta solo esta activa en un rango, aniadir bandera \texttt{active}.
	\item \textbf{Eliminar rectas}: usar Li Chao segmentado (mas complejo).
	\item \textbf{Pendientes grandes}: si las pendientes tienen magnitud $> 10^9$, aniadir \texttt{\_\_int128} para evitar overflow.
\end{itemize}

\paragraph{Trampas y errores frecuentes.}
\begin{itemize}\setlength\itemsep{0pt}
	\item El parametro \texttt{lo == hi} debe retornar el valor del nodo (caso base).
	\item La comparacion \texttt{newLine.get(mid) < node->line.get(mid)} decide quien es mejor en el medio.
	\item Si las rectas tienen $m$ iguales pero $b$ distintos, se queda con la primera (tie-break).
\end{itemize}

\paragraph{Codigo.}
Ver \texttt{cpp.tex}, seccion \textit{Li Chao Tree (long long)}.

\vspace{0.5em}

\subsubsection{Sweep line template}

\paragraph{Idea intuitiva.}
Patron de diseno para problemas geometricos donde se barre una linea sobre el plano y se mantiene una estructura de datos sobre el otro eje. Eventos (add, query, remove) se ordenan por la coordenada de barrido y se procesan secuencialmente.

\paragraph{Propiedades clave (invariantes).}
\begin{itemize}\setlength\itemsep{0pt}
	\item Eventos: tuplas \texttt{(x, type, id)}.
	\item La estructura auxiliar (set, multiset, segtree) se mantiene sobre el otro eje.
	\item Procesar eventos en orden de $x$ ascendente.
\end{itemize}

\paragraph{Complejidad.}
\begin{itemize}\setlength\itemsep{0pt}
	\item Depende del problema; tipico $O((N + E) \log N)$.
\end{itemize}

\paragraph{Donde tocar para modificar.}
\begin{itemize}\setlength\itemsep{0pt}
	\item \textbf{Aniadir/quitar segmentos}: aniadir eventos add/remove al barrido.
	\item \textbf{Aniadir queries}: aniadir eventos query.
	\item \textbf{Cambiar el eje de barrido}: si barres en $y$ en vez de $x$, aniadir comparador.
\end{itemize}

\paragraph{Trampas y errores frecuentes.}
\begin{itemize}\setlength\itemsep{0pt}
	\item El snippet es solo el \textbf{esqueleto}; los handlers \texttt{add/query/remove} se llenan por problema.
	\item Cuidado con eventos duplicados.
\end{itemize}

\paragraph{Codigo.}
Ver \texttt{cpp.tex}, seccion \textit{Sweep line template}.

\vspace{0.5em}

\subsubsection{Closest pair of points}

\paragraph{Idea intuitiva.}
Encuentra el par de puntos mas cercano en $O(n \log n)$ con \textbf{divide y venceras}. Algoritmo:
\begin{itemize}\setlength\itemsep{0pt}
	\item Sort por $x$.
	\item Dividir en mitades; recursivo.
	\item Combinar: dado el minimo $d$ de las dos mitades, strip = puntos en $[midX - d, midX + d]$; ordenar por $y$; cada punto solo se compara con los siguientes hasta $y$ difiera en $< d$ (maximo 7).
\end{itemize}

\paragraph{Propiedades clave (invariantes).}
\begin{itemize}\setlength\itemsep{0pt}
	\item Requiere ordenar por $y$ despues (en el snippet se hace durante el merge).
	\item Strip tiene a lo sumo $O(n)$ puntos en total, pero la comparacion es $O(1)$ amortizada por punto.
	\item Resultado en \texttt{double}.
\end{itemize}

\paragraph{Complejidad.}
\begin{itemize}\setlength\itemsep{0pt}
	\item $O(n \log n)$.
\end{itemize}

\paragraph{Donde tocar para modificar.}
\begin{itemize}\setlength\itemsep{0pt}
	\item \textbf{Reconstruir el par}: aniadir indices en la recursion.
	\item \textbf{Distancia maxima}: cambiar \texttt{min} por \texttt{max} (mas facil; no requiere D\&C).
	\item \textbf{Puntos duplicados}: aniadir check especial (distancia 0).
\end{itemize}

\paragraph{Trampas y errores frecuentes.}
\begin{itemize}\setlength\itemsep{0pt}
	\item El merge requiere ordenar por $y$; el snippet lo hace in-place con un \texttt{tmp}.
	\item Si se hace en cada recursion, es $O(n \log n)$ extra.
\end{itemize}

\paragraph{Codigo.}
Ver \texttt{cpp.tex}, seccion \textit{Closest pair of points}.

\vspace{0.5em}

\subsubsection{Rotating calipers (polygon diameter)}

\paragraph{Idea intuitiva.}
Encuentra el \textbf{diametro} de un poligono convexo (par de vertices mas distantes) en $O(n)$. Asume poligono en orden counter-clockwise. Algoritmo: dos indices $i$ (sobre las aristas) y $k$ (sobre los vertices opuestos); en cada paso, rotar $k$ mientras el area del paralelogramo aumenta. La distancia $i$-$k$ puede ser maxima en cualquier momento del barrido.

\paragraph{Propiedades clave (invariantes).}
\begin{itemize}\setlength\itemsep{0pt}
	\item Cada indice avanza como maximo $n$ veces total.
	\item El ``diametro'' es la maxima distancia entre dos vertices.
	\item Retorna \textbf{distancia al cuadrado}.
\end{itemize}

\paragraph{Complejidad.}
\begin{itemize}\setlength\itemsep{0pt}
	\item $O(n)$.
\end{itemize}

\paragraph{Donde tocar para modificar.}
\begin{itemize}\setlength\itemsep{0pt}
	\item \textbf{Reconstruir el par de vertices}: aniadir indices.
	\item \textbf{Anchura minima / maxima}: variantes de calipers.
	\item \textbf{Distancia minima entre poligonos convexos}: otra variante.
\end{itemize}

\paragraph{Trampas y errores frecuentes.}
\begin{itemize}\setlength\itemsep{0pt}
	\item Asume poligono \textbf{estrictamente} convexo (sin vertices colineales). Si hay colineales, aniadir perturbacion o filtrar.
\end{itemize}

\paragraph{Codigo.}
Ver \texttt{cpp.tex}, seccion \textit{Rotating calipers (polygon diameter)}.

% ============================================================
\section{Miscelanea}
% ============================================================

\subsection{Utilidades del usuario}

\subsubsection{Hora}

\paragraph{Idea intuitiva.}
Una pequena clase para representar \textbf{tiempos en segundos} y hacer aritmetica basica sobre ellos. Internamente almacena todo en segundos (un solo \texttt{int}), lo cual simplifica comparaciones y diferencias. El constructor acepta $(h, m, s)$ y los convierte a segundos. Hay un metodo \texttt{cut()} que fuerza un minimo (util para ``horas trabajadas'').

\paragraph{Propiedades clave (invariantes).}
\begin{itemize}\setlength\itemsep{0pt}
	\item Representacion canonica: \textbf{solo segundos}. No hay zonas horarias, no hay DST.
	\item Comparadores \texttt{<} y \texttt{>} operan sobre los segundos.
	\item \texttt{operator-} retorna la \textbf{diferencia absoluta} en segundos (siempre no negativa).
	\item \texttt{cut()} aplica un piso (por defecto $8\text{h }30\text{m} = 30600$ segundos).
\end{itemize}

\paragraph{Complejidad.}
\begin{itemize}\setlength\itemsep{0pt}
	\item $O(1)$ por operacion.
\end{itemize}

\paragraph{Donde tocar para modificar.}
\begin{itemize}\setlength\itemsep{0pt}
	\item \textbf{Aniadir suma}: \texttt{operator+} que devuelve un \texttt{Hour} con la suma de segundos.
	\item \textbf{Aniadir output formateado}: aniadir \texttt{show\_hms()} que imprime $H$:$M$:$S$.
	\item \textbf{Cambiar el ``corte''}: modificar el valor magico \texttt{30'600} en \texttt{cut()} o pasarlo como parametro.
	\item \textbf{Tamano maximo}: si los tiempos exceden $\sim 68$ anos, el \texttt{int} desborda; usar \texttt{long long}.
	\item \textbf{Precision sub-segundo}: aniadir \texttt{ms} (milisegundos) y guardar en \texttt{long long} con factor 1000.
\end{itemize}

\paragraph{Trampas y errores frecuentes.}
\begin{itemize}\setlength\itemsep{0pt}
	\item El operador \texttt{-} retorna \textbf{valor absoluto}, no diferencia con signo. Si necesitas diferencia con signo, aniadir un nuevo metodo o sobrecargar.
	\item \texttt{cut()} modifica \texttt{s} \textbf{in-place}; si el llamador no espera esto, aniadir una version \texttt{cut(min)} que retorna un nuevo \texttt{Hour}.
	\item El umbral magico \texttt{30'600} (8h 30m) esta hardcoded; documentar bien que es ``jornada minima''.
\end{itemize}

\paragraph{Codigo.}
Ver \texttt{cpp.tex}, seccion \textit{Hora}.

\end{document}
' se responden en $O(|P|)$.
\end{itemize}

\paragraph{Trampas y errores frecuentes.}
\begin{itemize}\setlength\itemsep{0pt}
	\item El \textbf{clone} debe tener \texttt{occ = 0} (no es una nueva posicion final); el \texttt{link} del estado siendo creado y del estado clonado deben apuntar al clon.
	\item Saltarse esto da SAM incorrecto.
	\item La condicion \texttt{st[p].len + 1 == st[q].len} detecta cuando no se necesita clon.
	\item Olvidar \texttt{while (p != -1 && !st[p].next.count(c))} da transiciones rotas.
\end{itemize}

\paragraph{Codigo.}
Ver \texttt{cpp.tex}, seccion \textit{Suffix Automaton}.

\vspace{0.5em}

\subsubsection{Lyndon factorization (Duval)}

\paragraph{Idea intuitiva.}
Un \textbf{string de Lyndon} es un string que es estrictamente menor que todos sus rotaciones. Toda cadena tiene una \textbf{factorizacion unica} de Lyndon (algoritmo de \textbf{Duval}): una secuencia de strings de Lyndon $L_1, L_2, \dots, L_k$ tales que $L_1 \ge L_2 \ge \dots \ge L_k$ y la concatenacion es el string original. Esto da $O(n)$ para encontrar los cortes. La demostracion de correctitud: cada factor es minimal en su clase de rotacion, asi que la concatenacion respeta el orden lexicografico.

\paragraph{Propiedades clave (invariantes).}
\begin{itemize}\setlength\itemsep{0pt}
	\item Cada factor de Lyndon es de la forma $uv$ donde $u$ se repite y $v$ es prefijo de $u$.
	\item La longitud del factor esta dada por la ``longitud del periodo''.
	\item Aplicacion: el ``Lyndon word'' es el menor en orden lexicografico entre todas las rotaciones.
	\item La factorizacion es \textbf{unica}.
\end{itemize}

\paragraph{Codigo clave.}
El algoritmo de Duval (iterativo):
\begin{verbatim}
int n = s.size();
int i = 0;
while (i < n) {
    int j = i + 1, k = i;
    while (j < n && s[k] <= s[j]) {
        if (s[k] < s[j]) k = i;
        else ++k;
        ++j;
    }
    int len = j - k;
    while (i < j) { cout << "factor: " << s.substr(i, len) << "\n"; i += len; }
}
\end{verbatim}

\paragraph{Complejidad.}
\begin{itemize}\setlength\itemsep{0pt}
	\item $O(n)$.
	\item Cada comparacion es $O(1)$; el puntero $k$ retrocede a lo sumo $n$ veces en total.
\end{itemize}

\paragraph{Donde tocar para modificar.}
\begin{itemize}\setlength\itemsep{0pt}
	\item \textbf{Encontrar la menor rotacion}: el menor sufijo lexicografico de \texttt{s + s} es la menor rotacion; la posicion se obtiene con un solo pase de Duval.
	\item \textbf{Algoritmo de Booth}: alternativa para rotacion minima, tambien $O(n)$.
	\item \textbf{Validar}: si necesitas verificar que un string es de Lyndon, compararlo con sus rotaciones.
	\item \textbf{Encontrar todos los factores}: guardar la posicion y longitud de cada uno.
\end{itemize}

\paragraph{Trampas y errores frecuentes.}
\begin{itemize}\setlength\itemsep{0pt}
	\item La condicion \texttt{s[j] <= s[k]} (no \texttt{<}) en el while de Duval es crucial; cambiar a \texttt{<} da comportamiento distinto.
	\item El caso $s[k] < s[j]$ resetea $k$ a $i$ (inicio del factor actual).
	\item El bucle externo debe avanzar $i$ por la longitud del factor, no por 1.
\end{itemize}

\paragraph{Codigo.}
Ver \texttt{cpp.tex}, seccion \textit{Lyndon factorization (Duval)}.

\subsection{Palindromos}

\subsubsection{Manacher}

\paragraph{Idea intuitiva.}
Calcula en $O(n)$ los \textbf{radios} de todos los palindromos centrados en cada posicion (impares y pares). La idea: mantener una ventana $[l, r]$ que es el palindromo centrado en un punto anterior mas largo encontrado. Si el nuevo centro $i$ cae dentro de $[l, r]$, entonces el palindromo centrado en $i$ tiene al menos el radio del palindromo ``espejo'' $l + r - i$; sino, empezar desde 1 (impar) o 0 (par). Despues, expandir lo que falte. La demostracion de $O(n)$: cada caracter se compara a lo sumo 2 veces (extension + limite por ventana).

\paragraph{Propiedades clave (invariantes).}
\begin{itemize}\setlength\itemsep{0pt}
	\item Conveccion del snippet: \texttt{odd[i]} = $k$ tal que el palindromo impar mas largo centrado en $i$ tiene longitud $2k - 1$. \texttt{even[i]} = $k$ para par centrado entre $i-1$ e $i$, longitud $2k$.
	\item $k$ cuenta tambien cuantos palindromos hay centrados ahi.
	\item Las queries \texttt{isPalindrome(l, r)} son $O(1)$.
	\item La ventana $[l, r]$ siempre contiene el palindromo mas largo encontrado hasta ahora.
\end{itemize}

\paragraph{Codigo clave.}
Manacher para palindromos impares:
\begin{verbatim}
int l = 0, r = -1;
for (int i = 0; i < n; ++i) {
    int k = (i > r) ? 1 : min(odd[l + r - i], r - i + 1);
    while (i - k >= 0 && i + k < n && s[i-k] == s[i+k]) ++k;
    odd[i] = k;
    if (i + k - 1 > r) { l = i - k + 1; r = i + k - 1; }
}
\end{verbatim}

\paragraph{Complejidad.}
\begin{itemize}\setlength\itemsep{0pt}
	\item Construccion $O(n)$, queries $O(1)$.
	\item Cada extension se cuenta una sola vez.
\end{itemize}

\paragraph{Donde tocar para modificar.}
\begin{itemize}\setlength\itemsep{0pt}
	\item \textbf{Solo contar palindromos}: ya implementado en \texttt{countPalindromes}.
	\item \textbf{Longest palindromic substring}: ya en \texttt{longestPalindrome}.
	\item \textbf{Modificar la conveccion de radios}: si queres ``radio = mitad de la longitud'' en vez de ``k palindromos'', ajustar el offset.
	\item \textbf{Queries con k no precomputado}: aniadir helpers especificos.
	\item \textbf{Modificar el caracter de relleno}: para multiples separadores, aniadir un caracter unico.
\end{itemize}

\paragraph{Trampas y errores frecuentes.}
\begin{itemize}\setlength\itemsep{0pt}
	\item La conveccion del snippet es \textbf{distinta} de la ``clasica de radio'' (donde radio = $(L-1)/2$); tener cuidado al copiar a otro codigo.
	\item Si los palindromos se solapan, la condicion \texttt{i + k - 1 > r} actualiza la ventana.
	\item En palindromos pares, el centro esta ``entre'' dos indices; aniadir offset en queries.
\end{itemize}

\paragraph{Codigo.}
Ver \texttt{cpp.tex}, seccion \textit{Manacher}.

% ============================================================
\section{Grafos}
% ============================================================

\subsection{Caminos minimos}

\subsubsection{Dijkstra}

\paragraph{Idea intuitiva.}
Encuentra el shortest path desde un origen $s$ a todos los vertices en $O((V + E) \log V)$ asumiendo \textbf{pesos no negativos}. La idea: usar una cola de prioridad (min-heap) ordenada por distancia tentativa. En cada paso, extraer el nodo con menor distancia; si ya fue procesado, saltar. Si al actualizar un vecino la distancia mejora, insertar/actualizar en la cola. La demostracion: cuando un nodo se extrae por primera vez, no hay forma de llegar a el con menor distancia (cualquier camino alternativo deberia pasar por nodos no procesados, con distancia $\ge$ la actual).

\paragraph{Propiedades clave (invariantes).}
\begin{itemize}\setlength\itemsep{0pt}
	\item Asume pesos $\ge 0$; con pesos negativos usar Bellman-Ford.
	\item Greedy correctness: la primera vez que se extrae un nodo, su distancia es la definitiva.
	\item Si se quiere reconstruir el camino, guardar el \texttt{parent[]}.
	\item Cada nodo se procesa a lo sumo $V$ veces (mas relajarions, pero solo $V$ definitivas).
\end{itemize}

\paragraph{Codigo clave.}
El bucle principal de Dijkstra:
\begin{verbatim}
priority_queue<pair<long long, int>, vector<...>, greater<...>> pq;
pq.push({0, source});
while (!pq.empty()) {
    auto [d, v] = pq.top(); pq.pop();
    if (d > dist[v]) continue;  // entrada vieja
    for (auto [w, u] : adj[v])
        if (dist[v] + w < dist[u]) {
            dist[u] = dist[v] + w;
            pq.push({dist[u], u});
        }
}
\end{verbatim}

\paragraph{Complejidad.}
\begin{itemize}\setlength\itemsep{0pt}
	\item $O((V + E) \log V)$ con heap binario, $O(V + E)$ con heap de Fibonacci.
	\item En la practica, $\sim 1.5 \cdot 10^{-7}$ segundos por operacion para $V = 10^5$.
\end{itemize}

\paragraph{Donde tocar para modificar.}
\begin{itemize}\setlength\itemsep{0pt}
	\item \textbf{Camino mas corto entre todos los pares}: si se necesita desde todos los origenes, usar Floyd o ejecutar Dijkstra $V$ veces.
	\item \textbf{Shortest path con aristas negativas limitadas}: Bellman-Ford.
	\item \textbf{Reconstruir camino}: aniadir \texttt{parent[v]} actualizado junto a \texttt{dist[v]}.
	\item \textbf{Dijkstra con potentials (Johnson)}: para aristas negativas si se puede reescalar.
	\item \textbf{Multi-source}: aniadir varios nodos fuentes iniciales al heap.
\end{itemize}

\paragraph{Trampas y errores frecuentes.}
\begin{itemize}\setlength\itemsep{0pt}
	\item Distancias no inicializadas a \texttt{INF}; el snippet usa \texttt{1e18}.
	\item El \texttt{if(dist[adjN] != 1e18)} para borrar de la cola es importante si usas \texttt{set} (no asi con \texttt{priority\_queue}).
	\item Con pesos grandes, \texttt{dist[v] + w} puede overflow; usar \texttt{LLONG\_MAX/2} como INF.
	\item No usar con pesos negativos: el algoritmo puede dar respuestas incorrectas.
\end{itemize}

\paragraph{Codigo.}
Ver \texttt{cpp.tex}, seccion \textit{Dijkstra}.

\vspace{0.5em}

\subsubsection{Bellman-Ford}

\paragraph{Idea intuitiva.}
Encuentra shortest path desde $s$ incluso con \textbf{pesos negativos} en $O(VE)$. La idea: relajar todas las aristas $V - 1$ veces. Cada relajacion propaga la mejor distancia a ``un paso mas''. Despues de $V - 1$ pasadas, si alguna arista aun se puede relajar, hay un \textbf{ciclo negativo} alcanzable. La demostracion: el camino mas corto tiene a lo sumo $V - 1$ aristas (sin ciclos); tras $V - 1$ relajaciones todas las distancias estan en su valor optimo.

\paragraph{Propiedades clave (invariantes).}
\begin{itemize}\setlength\itemsep{0pt}
	\item Funciona con pesos negativos (sin ciclos negativos).
	\item Tras $V - 1$ relajaciones, las distancias son definitivas (si no hay ciclo negativo).
	\item Una relajacion adicional detecta ciclos negativos: si \texttt{dist[u] + w < dist[v]}, hay ciclo negativo alcanzable desde $s$.
	\item Cada relajacion solo mejora distancias, asi que termina (no oscila).
\end{itemize}

\paragraph{Codigo clave.}
El doble bucle de relajacion:
\begin{verbatim}
for (int i = 0; i < n - 1; ++i)
    for (auto [u, v, w] : edges)
        if (dist[u] + w < dist[v])
            dist[v] = dist[u] + w;
// Detectar ciclo negativo
for (auto [u, v, w] : edges)
    if (dist[u] + w < dist[v]) { /* ciclo negativo */ }
\end{verbatim}

\paragraph{Complejidad.}
\begin{itemize}\setlength\itemsep{0pt}
	\item $O(VE)$.
	\item SPFA (con cola) en promedio es mucho mas rapido, peor caso igual.
\end{itemize}

\paragraph{Donde tocar para modificar.}
\begin{itemize}\setlength\itemsep{0pt}
	\item \textbf{SPFA}: optimizacion que relaja solo nodos cuya distancia cambio; en el peor caso sigue siendo $O(VE)$.
	\item \textbf{Detectar ciclo negativo global}: ejecutar Bellman-Ford desde un super-origen conectado a todos.
	\item \textbf{Camino mas corto con K aristas exactas}: $K$ iteraciones de Bellman-Ford.
	\item \textbf{Johnson's}: reescalar pesos con potentials para usar Dijkstra.
\end{itemize}

\paragraph{Trampas y errores frecuentes.}
\begin{itemize}\setlength\itemsep{0pt}
	\item Overflow si los pesos son $\ge 2^{52}$; usar \texttt{long long} y \texttt{\_\_int128} si hace falta.
	\item Si el grafo es denso ($E \approx V^2$), $O(VE)$ es prohibitivo.
	\item El chequeo de ciclo negativo debe hacerse tras las $V - 1$ iteraciones, no durante.
	\item Distancias no inicializadas a \texttt{INF}; el primer nodo (source) tiene distancia 0.
\end{itemize}

\paragraph{Codigo.}
Ver \texttt{cpp.tex}, seccion \textit{Bellman-Ford}.

\vspace{0.5em}

\subsubsection{Floyd-Warshall}

\paragraph{Idea intuitiva.}
Shortest path entre \textbf{todos los pares} de vertices en $O(V^3)$. DP: $d[i][j]$ = shortest path usando solo nodos intermedios en $\{0, \dots, k\}$. Transicion: $d_{k+1}[i][j] = \min(d_k[i][j], d_k[i][k+1] + d_k[k+1][j])$. La demostracion: tras $k$ iteraciones, $d[i][j]$ es el shortest path con intermedios en $\{0, \dots, k-1\}$; aniadir el nodo $k$ da la recurrencia clasica.

\paragraph{Propiedades clave (invariantes).}
\begin{itemize}\setlength\itemsep{0pt}
	\item Detecta ciclos negativos: tras el algoritmo, $d[i][i] < 0$ indica ciclo negativo alcanzable desde $i$.
	\item Funciona con pesos negativos (sin ciclos negativos).
	\item Memoria $O(V^2)$; con $V \le 400$ es OK, con mas conviene Dijkstra $V$ veces.
	\item Es la unica opcion cuando se quieren shortest paths entre todos los pares y el grafo es denso.
\end{itemize}

\paragraph{Codigo clave.}
El triple bucle clasico:
\begin{verbatim}
for (int k = 0; k < n; ++k)
    for (int i = 0; i < n; ++i)
        for (int j = 0; j < n; ++j)
            d[i][j] = min(d[i][j], d[i][k] + d[k][j]);
// detectar ciclo negativo
for (int i = 0; i < n; ++i) if (d[i][i] < 0) tiene_ciclo_negativo = true;
\end{verbatim}

\paragraph{Complejidad.}
\begin{itemize}\setlength\itemsep{0pt}
	\item $O(V^3)$ tiempo, $O(V^2)$ memoria.
	\item Para $V \le 400$: $\sim 6.4 \cdot 10^7$ operaciones (factible).
\end{itemize}

\paragraph{Donde tocar para modificar.}
\begin{itemize}\setlength\itemsep{0pt}
	\item \textbf{Reconstruir camino}: aniadir \texttt{next[i][j]} y backtrack.
	\item \textbf{Transitividad}: el algoritmo tambien calcula transitividad (\texttt{d[i][j] != INF} significa alcanzable).
	\item \textbf{Pesos doubles}: aniadir tolerancia (\texttt{d[i][j] > d[i][k] + d[k][j] - eps}).
	\item \textbf{Grafo no denso}: usar Dijkstra $V$ veces en vez de Floyd.
\end{itemize}

\paragraph{Trampas y errores frecuentes.}
\begin{itemize}\setlength\itemsep{0pt}
	\item Para detectar ciclo negativo correctamente, aniadir un chequeo \texttt{if (ans[i][i] < 0)} tras el triple bucle.
	\item El orden de los bucles \texttt{k, i, j} es importante; invertir da resultados diferentes.
	\item Con \texttt{double}, usar tolerancia \texttt{< eps} en vez de \texttt{<} para evitar loops infinitos.
\end{itemize}

\paragraph{Codigo.}
Ver \texttt{cpp.tex}, seccion \textit{Floyd-Warshall}.

\subsection{Componentes y cortes}

\subsubsection{Kosaraju SCC}

\paragraph{Idea intuitiva.}
Encuentra las \textbf{componentes fuertemente conexas} (SCCs) en $O(V + E)$. Algoritmo en 3 pasos:
\begin{itemize}\setlength\itemsep{0pt}
	\item DFS en $G$ rellenando una pila en \textbf{post-orden} (los nodos que terminan tarde quedan arriba).
	\item Construir $G^T$ (grafo transpuesto).
	\item DFS en $G^T$ procesando nodos en orden de la pila (de mayor a menor tiempo); cada DFS da una SCC.
\end{itemize}
La demostracion: dos nodos $u, v$ estan en la misma SCC sii hay camino $u \leadsto v$ y $v \leadsto u$; Kosaraju captura exactamente esto al explorar $G^T$ desde los nodos con mayor tiempo de finalizacion en $G$.

\paragraph{Propiedades clave (invariantes).}
\begin{itemize}\setlength\itemsep{0pt}
	\item Cada SCC es maximal fuertemente conexa.
	\item El \textbf{grafo de SCCs} (metagraf) es un DAG.
	\item Kosaraju y Tarjan dan el mismo resultado; Tarjan es ``mas rapido'' en practica (un solo DFS).
	\item Cada nodo pertenece a exactamente una SCC.
\end{itemize}

\paragraph{Codigo clave.}
La estructura de Kosaraju:
\begin{verbatim}
vector<int> order;
vector<bool> used(n);
function<void(int)> dfs1 = [&](int v) {
    used[v] = true;
    for (int u : g[v]) if (!used[u]) dfs1(u);
    order.push_back(v);  // post-orden
};
function<void(int, int)> dfs2 = [&](int v, int root) {
    comp[v] = root;
    for (int u : gt[v]) if (comp[u] == -1) dfs2(u, root);
};
// 1: DFS en G para llenar order
// 2: DFS en G^T en orden inverso para asignar SCCs
\end{verbatim}

\paragraph{Complejidad.}
\begin{itemize}\setlength\itemsep{0pt}
	\item $O(V + E)$.
	\item Constante 2 (dos DFS + el grafo transpuesto).
\end{itemize}

\paragraph{Donde tocar para modificar.}
\begin{itemize}\setlength\itemsep{0pt}
	\item \textbf{2-SAT}: SCC sobre grafo de implicaciones.
	\item \textbf{Construir el DAG de SCCs}: aniadir aristas entre SCCs diferentes.
	\item \textbf{Tamano de cada SCC}: aniadir contador durante el segundo DFS.
	\item \textbf{Grafo grande}: usar Tarjan para evitar construir $G^T$.
\end{itemize}

\paragraph{Trampas y errores frecuentes.}
\begin{itemize}\setlength\itemsep{0pt}
	\item El snippet usa \texttt{vector<...> adj[n]} (VLA); reemplazable por \texttt{vector<vector<pair<int,int>>} para portabilidad.
	\item El orden del segundo DFS debe ser \textbf{inverso} al de finalizacion del primero.
	\item Si el grafo no es conexo, aniadir un loop sobre todos los nodos en el primer DFS.
	\item Confundir SCC con componentes conexas (las primeras son para dirigidos, las segundas para no dirigidos).
\end{itemize}

\paragraph{Codigo.}
Ver \texttt{cpp.tex}, seccion \textit{Kosaraju SCC}.

\vspace{0.5em}

\subsubsection{Tarjan (puentes)}

\paragraph{Idea intuitiva.}
Encuentra los \textbf{puentes} (bridges): aristas cuya remocion \textbf{desconecta} el grafo. Usa un solo DFS manteniendo \texttt{tin[v]} (tiempo de descubrimiento) y \texttt{low[v]} (el menor \texttt{tin} alcanzable desde $v$ via back-edges). Una arista $v - u$ es puente si $low[u] > tin[v]$. La demostracion: si $low[u] > tin[v]$, entonces $u$ (y su subarbol) no pueden ``escapar'' de $v$; quitar la arista $v - u$ parte el grafo.

\paragraph{Propiedades clave (invariantes).}
\begin{itemize}\setlength\itemsep{0pt}
	\item \texttt{low[v] = min(tin[v], tin[w])} para todo back-edge $(v, w)$, y \texttt{min(low[u])} para hijos en el DFS.
	\item Solo se consideran las aristas del \textbf{arbol DFS} (no back-edges) para \texttt{low}.
	\item Cada nodo tiene exactamente un \texttt{tin}; cada back-edge se procesa una vez.
\end{itemize}

\paragraph{Codigo clave.}
El DFS de Tarjan para bridges:
\begin{verbatim}
void dfs(int v, int pEdge) {
    used[v] = true;
    tin[v] = low[v] = timer++;
    for (auto [u, id] : adj[v]) {
        if (id == pEdge) continue;  // no contar el edge al padre
        if (used[u]) {
            low[v] = min(low[v], tin[u]);  // back-edge
        } else {
            dfs(u, id);
            low[v] = min(low[v], low[u]);
            if (low[u] > tin[v]) bridges.push_back(id);  // es puente
        }
    }
}
\end{verbatim}

\paragraph{Complejidad.}
\begin{itemize}\setlength\itemsep{0pt}
	\item $O(V + E)$.
	\item Una sola pasada de DFS.
\end{itemize}

\paragraph{Donde tocar para modificar.}
\begin{itemize}\setlength\itemsep{0pt}
	\item \textbf{Bridge tree}: contractar cada 2-edge-connected component a un nodo; el resultado es un arbol.
	\item \textbf{Articulation points}: variante con condicion diferente (ver el siguiente modulo).
	\item \textbf{Grafo no dirigido}: solo funciona en no dirigidos; en dirigidos hay conceptos analogos (strong bridges, dominator tree).
	\item \textbf{Reconstruir componentes 2-edge-connected}: BFS desde cada nodo excluyendo bridges.
\end{itemize}

\paragraph{Trampas y errores frecuentes.}
\begin{itemize}\setlength\itemsep{0pt}
	\item No confundir el parent con un back-edge al padre; chequear \texttt{if (adjN == parent) continue;} y contar solo las veces que se ``visita'' como hijo.
	\item Para multigrafos, los edges deben tener IDs unicos para distinguirlos.
	\item Si el grafo no es conexo, aniadir loop externo sobre todos los nodos.
\end{itemize}

\paragraph{Codigo.}
Ver \texttt{cpp.tex}, seccion \textit{Tarjan (puentes)}.

\vspace{0.5em}

\subsubsection{Tarjan (puntos de articulacion)}

\paragraph{Idea intuitiva.}
Encuentra los \textbf{puntos de articulacion} (cut vertices): vertices cuya remocion \textbf{desconecta} el grafo. Mismo setup que para bridges: \texttt{tin[v]}, \texttt{low[v]}. Un nodo $v$ (no raiz) es de articulacion si existe un hijo $u$ con $low[u] \ge tin[v]$. La \textbf{raiz} es de articulacion si tiene $\ge 2$ hijos en el DFS tree. La razon: para $v$ no raiz, la condicion $\ge$ indica que el subarbol de $u$ depende de $v$; quitar $v$ los desconecta. Para la raiz, perder $\ge 2$ hijos parte el DFS en $\ge 2$ componentes.

\paragraph{Propiedades clave (invariantes).}
\begin{itemize}\setlength\itemsep{0pt}
	\item $low[u] \ge tin[v]$ significa que el subarbol de $u$ no tiene back-edge ``por encima'' de $v$.
	\item La raiz es caso especial: tiene que perder mas de un hijo para que el grafo se parta.
	\item Cada nodo se evalua una sola vez.
\end{itemize}

\paragraph{Codigo clave.}
La condicion de articulacion:
\begin{verbatim}
void dfs(int v, int p) {
    used[v] = true;
    tin[v] = low[v] = timer++;
    int children = 0;
    for (int u : adj[v]) {
        if (u == p) continue;
        if (used[u]) low[v] = min(low[v], tin[u]);
        else {
            dfs(u, v);
            low[v] = min(low[v], low[u]);
            if (low[u] >= tin[v] && p != -1)  // no raiz
                isArt[v] = true;
            ++children;
        }
    }
    if (p == -1 && children > 1) isArt[v] = true;  // raiz especial
}
\end{verbatim}

\paragraph{Complejidad.}
\begin{itemize}\setlength\itemsep{0pt}
	\item $O(V + E)$.
	\item Una sola pasada de DFS.
\end{itemize}

\paragraph{Donde tocar para modificar.}
\begin{itemize}\setlength\itemsep{0pt}
	\item \textbf{Block-cut tree}: estructura bipartita con BCCs (2-vertex-connected components) y articulation points.
	\item \textbf{Verificar biconectividad}: chequear que no hay puntos de articulacion.
	\item \textbf{Componentes biconnected}: extender Tarjan para enumerar BCCs.
\end{itemize}

\paragraph{Trampas y errores frecuentes.}
\begin{itemize}\setlength\itemsep{0pt}
	\item No olvidar el caso especial de la raiz.
	\item La condicion usa \texttt{>=}, no \texttt{>}; un hijo con \texttt{low[u] == tin[v]} tambien cuenta.
	\item Multigrafos pueden tener ``falsos'' articulation points; verificar caso a caso.
\end{itemize}

\paragraph{Codigo.}
Ver \texttt{cpp.tex}, seccion \textit{Tarjan (puntos de articulacion)}.

\vspace{0.5em}

\subsubsection{Tarjan SCC}

\paragraph{Idea intuitiva.}
Variante del SCC que usa \textbf{un solo DFS} con una pila explicita. Mantiene \texttt{disc[v]}, \texttt{low[v]}, y una pila de nodos ``en la SCC actual''. Cuando \texttt{low[v] == disc[v]}, $v$ es la raiz de una SCC; popear hasta $v$ inclusive. La demostracion: $low[v] == disc[v]$ indica que $v$ no tiene back-edges a ancestros; entonces $v$ y los nodos en la pila hasta $v$ forman una SCC maximal.

\paragraph{Propiedades clave (invariantes).}
\begin{itemize}\setlength\itemsep{0pt}
	\item Mas eficiente que Kosaraju (un solo DFS).
	\item \texttt{low[v] = min(low[u], disc[w])} para back-edges o recursion.
	\item La pila asegura que solo los nodos en la SCC actual se popean.
	\item Cada nodo se pushea y popea a lo sumo una vez.
\end{itemize}

\paragraph{Codigo clave.}
El DFS de Tarjan SCC:
\begin{verbatim}
void dfs(int v) {
    disc[v] = low[v] = timer++;
    st.push(v); inStack[v] = true;
    for (int u : g[v]) {
        if (disc[u] == -1) { dfs(u); low[v] = min(low[v], low[u]); }
        else if (inStack[u]) low[v] = min(low[v], disc[u]);
    }
    if (low[v] == disc[v]) {
        // popear hasta v
        while (true) {
            int u = st.top(); st.pop(); inStack[u] = false;
            comp[u] = current_scc;
            if (u == v) break;
        }
        ++current_scc;
    }
}
\end{verbatim}

\paragraph{Complejidad.}
\begin{itemize}\setlength\itemsep{0pt}
	\item $O(V + E)$.
	\item Constante similar a Kosaraju pero sin construir $G^T$.
\end{itemize}

\paragraph{Donde tocar para modificar.}
\begin{itemize}\setlength\itemsep{0pt}
	\item \textbf{Tamano de cada SCC}: aniadir contador durante el pop.
	\item \textbf{Grafo condensation}: aniadir aristas entre SCCs en el DAG resultante.
	\item \textbf{2-SAT}: aniadir la construccion del grafo de implicaciones.
\end{itemize}

\paragraph{Trampas y errores frecuentes.}
\begin{itemize}\setlength\itemsep{0pt}
	\item \texttt{inStack[u]} debe chequearse antes de usar \texttt{low[v] = min(low[v], disc[u])} en back-edges.
	\item La pila global es importante; sin ella no se puede extraer la SCC.
	\item Si el grafo es muy grande, recursion puede overflow; aniadir version iterativa.
\end{itemize}

\paragraph{Codigo.}
Ver \texttt{cpp.tex}, seccion \textit{Tarjan SCC}.

\subsection{Matching}

\subsubsection{Kuhn (bipartite matching)}

\paragraph{Idea intuitiva.}
Encuentra un \textbf{matching maximo} en un grafo bipartito en $O(VE)$. Para cada nodo $u$ en $L$, intentar encontrar un vecino $v$ en $R$ que este libre o cuyo match se pueda ``reorganizar'' recursivamente. Si se encuentra, matchear $u - v$. La idea: si $u$ quiere un vecino $v$ ya ocupado, intentamos ``mover'' el match de $v$ a otro nodo; recursivamente, esto encuentra un augmenting path si existe.

\paragraph{Propiedades clave (invariantes).}
\begin{itemize}\setlength\itemsep{0pt}
	\item \texttt{matchR[v]} = el nodo de $L$ matcheado a $v$ (o $-1$ si libre).
	\item \texttt{vis[]} se resetea \textbf{cada intento} de $u$.
	\item Si el matching es perfecto, $|L| = |R| = $ tamano del matching.
	\item Kuhn no garantiza maximalidad greedy; puede haber augmenting paths no encontrados.
\end{itemize}

\paragraph{Codigo clave.}
La recursion del augmenting path:
\begin{verbatim}
bool tryKuhn(int v) {
    if (vis[v]) return false;
    vis[v] = true;
    for (int u : g[v]) {
        if (matchR[u] == -1 || tryKuhn(matchR[u])) {
            matchR[u] = v;
            return true;
        }
    }
    return false;
}
// Por cada v en L:
fill(vis, vis+n, false);
if (tryKuhn(v)) ++matches;
\end{verbatim}

\paragraph{Complejidad.}
\begin{itemize}\setlength\itemsep{0pt}
	\item $O(VE)$ en el peor caso.
	\item Con Hopcroft-Karp: $O(\sqrt{V} \cdot E)$, mucho mejor para grafos densos.
\end{itemize}

\paragraph{Donde tocar para modificar.}
\begin{itemize}\setlength\itemsep{0pt}
	\item \textbf{Matching bipartito minimo (min vertex cover)}: Konig's theorem.
	\item \textbf{Matching general}: blossom algorithm (no incluido).
	\item \textbf{Aniadir capacidades}: si los nodos de $R$ tienen capacidad $> 1$, partir cada uno en capacidad copias.
	\item \textbf{Output del matching}: el array \texttt{matchR} da los pares finales.
\end{itemize}

\paragraph{Trampas y errores frecuentes.}
\begin{itemize}\setlength\itemsep{0pt}
	\item Olvidar resetear \texttt{vis} por intento; eso es lo que evita ``ciclos infinitos'' en la recursion.
	\item Para grafos grandes ($V > 10^4$), Kuhn puede ser muy lento; usar Hopcroft-Karp.
	\item Si el matching es maximo pero no perfecto, $|L|$ o $|R|$ puede ser mayor; el matching es solo maximo.
\end{itemize}

\paragraph{Codigo.}
Ver \texttt{cpp.tex}, seccion \textit{Kuhn (bipartite matching)}.

\vspace{0.5em}

\subsubsection{Hopcroft-Karp}

\paragraph{Idea intuitiva.}
Matching bipartito maximo en $O(\sqrt{V} E)$. Algoritmo en \textbf{fases}: en cada fase, hacer BFS para encontrar los \textbf{caminos de aumento mas cortos} desde nodos libres de $L$; luego DFS solo por esos caminos. Cada fase aumenta el matching en $\ge 1$ y toma $O(E)$. La demostracion: tras una fase, todos los augmenting paths tienen longitud $\ge$ la nueva minima; eso sucede cada $\sqrt{V}$ fases (longitud minima crece geometricamente).

\paragraph{Propiedades clave (invariantes).}
\begin{itemize}\setlength\itemsep{0pt}
	\item \texttt{dist[u]} = distancia en niveles desde un nodo libre de $L$.
	\item BFS termina cuando no hay mas camino a un nodo libre de $R$.
	\item DFS respeta los niveles (solo sigue \texttt{dist[matchV[v]] == dist[u] + 1}).
	\item Cada fase aumenta el matching en $\ge 1$ arista.
\end{itemize}

\paragraph{Codigo clave.}
El BFS de Hopcroft-Karp:
\begin{verbatim}
bool bfs() {
    queue<int> q;
    for (int v : L) {
        if (matchU[v] == -1) { dist[v] = 0; q.push(v); }
        else dist[v] = -1;
    }
    bool found = false;
    while (!q.empty()) {
        int v = q.front(); q.pop();
        for (int u : g[v]) {
            if (matchV[u] != -1 && dist[matchV[u]] == -1) {
                dist[matchV[u]] = dist[v] + 1;
                q.push(matchV[u]);
            }
            if (matchV[u] == -1) found = true;  // llegamos a libre
        }
    }
    return found;
}
\end{verbatim}

\paragraph{Complejidad.}
\begin{itemize}\setlength\itemsep{0pt}
	\item $O(\sqrt{V} E)$.
	\item Para $V = 10^4$, $E = 10^5$: $\sim 3 \cdot 10^7$ operaciones.
\end{itemize}

\paragraph{Donde tocar para modificar.}
\begin{itemize}\setlength\itemsep{0pt}
	\item \textbf{Reconstruir el matching}: ya viene en \texttt{matchU}/\texttt{matchV}.
	\item \textbf{Aniadir pesos}: Hungarian algorithm.
	\item \textbf{Densidad alta}: si $E \approx V^2$, sigue siendo util pero verificar constantes.
	\item \textbf{Aniadir capacidades}: partir los nodos con capacidad $> 1$ en varias copias.
\end{itemize}

\paragraph{Trampas y errores frecuentes.}
\begin{itemize}\setlength\itemsep{0pt}
	\item \texttt{dist[matchV[v]] == dist[u] + 1} (no \texttt{==} a secas); sin esto, el DFS no respeta los niveles y se pierde la complejidad.
	\item El BFS debe marcar los niveles correctamente; resetear \texttt{dist} para nodos alcanzables solo via matching.
	\item El DFS debe respetar los niveles; cualquier ``atajo'' invalida la cota.
\end{itemize}

\paragraph{Codigo.}
Ver \texttt{cpp.tex}, seccion \textit{Hopcroft-Karp}.

\subsection{Basicos}

\subsubsection{DFS / BFS / componentes}

\paragraph{Idea intuitiva.}
\textbf{DFS} (Depth-First Search) y \textbf{BFS} (Breadth-First Search) son los dos recorridos basicos de grafos. DFS usa una pila (o recursion) y explora tan profundo como puede antes de retroceder; BFS usa una cola y explora por ``niveles'' de distancia. La cantidad de \textbf{componentes conexas} se cuenta iniciando un BFS/DFS desde cada nodo no visitado. La diferencia clave: BFS da el camino mas corto en aristas; DFS da un orden de finalizacion util para problemas topologicos.

\paragraph{Propiedades clave (invariantes).}
\begin{itemize}\setlength\itemsep{0pt}
	\item BFS da la \textbf{distancia en aristas} mas corta desde el origen.
	\item DFS da un \textbf{orden topologico} (si se aniaden al final).
	\item Ambos son $O(V + E)$.
	\item Un grafo con $C$ componentes conexas satisface $\sum \text{subtree\_size}_v = V$ y $\sum E_v = 2E$.
	\item DFS en arbol da un spanning tree; las ``back edges'' corresponden a ciclos.
\end{itemize}

\paragraph{Codigo clave.}
BFS clasico:
\begin{verbatim}
queue<int> q;
q.push(source);
dist[source] = 0;
while (!q.empty()) {
    int v = q.front(); q.pop();
    for (int u : adj[v])
        if (dist[u] == -1) {
            dist[u] = dist[v] + 1;
            q.push(u);
        }
}
\end{verbatim}

\paragraph{Complejidad.}
\begin{itemize}\setlength\itemsep{0pt}
	\item $O(V + E)$ tiempo, $O(V)$ memoria.
	\item Cada arista se procesa exactamente 2 veces (ida y vuelta) en no dirigidos.
\end{itemize}

\paragraph{Donde tocar para modificar.}
\begin{itemize}\setlength\itemsep{0pt}
	\item \textbf{Aniadir pesos}: si las aristas tienen peso y queres shortest path, BFS ya no sirve; usar Dijkstra.
	\item \textbf{Grafo dirigido}: en DFS, aniadir check \texttt{if (state[u] == 1)} para detectar ciclos.
	\item \textbf{Bipartite check}: BFS alternando colores (ver el siguiente modulo).
	\item \textbf{Componentes fuertemente conexas (dirigidos)}: usar Tarjan o Kosaraju.
	\item \textbf{BFS en matriz}: aniadir offsets por direccion para grid 2D.
\end{itemize}

\paragraph{Trampas y errores frecuentes.}
\begin{itemize}\setlength\itemsep{0pt}
	\item Stack overflow con DFS recursivo para $V > 10^5$; usar version iterativa con pila explicita.
	\item En BFS, olvidar marcar \texttt{dist[u] = dist[v] + 1} antes de pushear da visitas duplicadas.
	\item Para grafos dirigidos, ``componente conexa'' no esta definida; usar SCC.
\end{itemize}

\paragraph{Codigo.}
Ver \texttt{cpp.tex}, seccion \textit{DFS / BFS / componentes}.

\vspace{0.5em}

\subsubsection{Topological sort (Kahn + DFS)}

\paragraph{Idea intuitiva.}
Un \textbf{orden topologico} de un DAG es una permutacion de los vertices tal que toda arista $u \to v$ tiene $u$ antes que $v$. Dos algoritmos clasicos:
\begin{itemize}\setlength\itemsep{0pt}
	\item \textbf{Kahn}: BFS sobre nodos con \texttt{indeg == 0}; al ``consumir'' un nodo, decrementar el indeg de sus vecinos; los que lleguen a 0 se encolan.
	\item \textbf{DFS}: hacer DFS; al terminar un nodo (estado = 2), aniadirlo al final; el orden es el reverso.
\end{itemize}
La demostracion: en un DAG siempre hay al menos un nodo con indeg 0; quitarlo mantiene el DAG; por induccion, se obtiene el orden.

\paragraph{Propiedades clave (invariantes).}
\begin{itemize}\setlength\itemsep{0pt}
	\item Existe sii el grafo es un DAG.
	\item Si al final \texttt{order.size() < V}, hay un ciclo.
	\item Kahn da el orden ``natural'' (los que pueden ir primero); DFS da el orden ``al reves'' (los que pueden ir al final primero).
	\item La cantidad de ordenes topologicos puede ser exponencial.
\end{itemize}

\paragraph{Codigo clave.}
Kahn con priority queue (orden lexicografico):
\begin{verbatim}
priority_queue<int, vector<int>, greater<int>> pq;
for (int v = 0; v < n; ++v) if (indeg[v] == 0) pq.push(v);
while (!pq.empty()) {
    int v = pq.top(); pq.pop();
    order.push_back(v);
    for (int u : adj[v]) {
        if (--indeg[u] == 0) pq.push(u);
    }
}
\end{verbatim}

\paragraph{Complejidad.}
\begin{itemize}\setlength\itemsep{0pt}
	\item $O(V + E)$ ambos.
	\item Con priority queue, $O((V + E) \log V)$.
\end{itemize}

\paragraph{Donde tocar para modificar.}
\begin{itemize}\setlength\itemsep{0pt}
	\item \textbf{Detectar ciclos}: comparar \texttt{order.size()} con \texttt{V}.
	\item \textbf{Orden lexicografico minimo}: usar priority queue en Kahn.
	\item \textbf{Topological sort con pesos}: combinar con shortest path en DAG (un solo pase en orden topologico).
	\item \textbf{Encontrar todos los ordenes}: DP sobre subsets (solo para $V$ pequeno).
\end{itemize}

\paragraph{Trampas y errores frecuentes.}
\begin{itemize}\setlength\itemsep{0pt}
	\item En el snippet, Kahn retorna orden vacio si hay ciclo (util para detectar).
	\item En DFS, no olvidar \texttt{reverse(order)}.
	\item Confundir indeg con outdeg; el orden Kahn quita primero los nodos ``fuente''.
	\item Para grafos no DAG, el algoritmo termina con \texttt{order.size() < V}.
\end{itemize}

\paragraph{Codigo.}
Ver \texttt{cpp.tex}, seccion \textit{Topological sort (Kahn + DFS)}.

\vspace{0.5em}

\subsubsection{Bipartite check + 2-coloring}

\paragraph{Idea intuitiva.}
Un grafo es \textbf{bipartito} si sus vertices se pueden dividir en dos conjuntos $L$ y $R$ tales que toda arista va de $L$ a $R$. Esto equivale a decir que se puede \textbf{2-colorear} sin que dos vertices adyacentes tengan el mismo color. BFS alternando colores detecta esto: si en algun momento dos adyacentes tienen el mismo color, no es bipartito. La demostracion: por cada componente, fijar el color de un nodo determina todos los demas; si en algun paso hay conflicto, hay un ciclo impar (no bipartito).

\paragraph{Propiedades clave (invariantes).}
\begin{itemize}\setlength\itemsep{0pt}
	\item Grafo bipartito $\Leftrightarrow$ sin ciclos impares.
	\item Conexo: si un componente no es bipartito, todo el grafo no lo es.
	\item Si hay \textbf{multiple componentes}, cada uno se colorea independientemente.
	\item La 2-coloracion es unica salvo swap de colores.
\end{itemize}

\paragraph{Codigo clave.}
BFS con 2-coloreo:
\begin{verbatim}
vector<int> color(n, -1);
for (int v = 0; v < n; ++v) if (color[v] == -1) {
    color[v] = 0;
    queue<int> q; q.push(v);
    while (!q.empty()) {
        int u = q.front(); q.pop();
        for (int w : adj[u]) {
            if (color[w] == -1) { color[w] = 1 - color[u]; q.push(w); }
            else if (color[w] == color[u]) return false;
        }
    }
}
return true;
\end{verbatim}

\paragraph{Complejidad.}
\begin{itemize}\setlength\itemsep{0pt}
	\item $O(V + E)$.
	\item Memoria $O(V)$.
\end{itemize}

\paragraph{Donde tocar para modificar.}
\begin{itemize}\setlength\itemsep{0pt}
	\item \textbf{Colorear desde multiples fuentes}: BFS desde cada nodo no coloreado.
	\item \textbf{Bipartito + matching}: Hopcroft-Karp.
	\item \textbf{Grafo bipartito ponderado}: Hungarian algorithm.
	\item \textbf{Verificar bipartito en subgrafo}: solo aplicar BFS a los nodos relevantes.
\end{itemize}

\paragraph{Trampas y errores frecuentes.}
\begin{itemize}\setlength\itemsep{0pt}
	\item Verificar \textbf{todos} los componentes, no solo el primero.
	\item Para grafos con auto-loops, un auto-loop hace al grafo no bipartito.
	\item En multigrafos, las aristas multiples no afectan el chequeo de bipartito.
\end{itemize}

\paragraph{Codigo.}
Ver \texttt{cpp.tex}, seccion \textit{Bipartite check + 2-coloring}.

\vspace{0.5em}

\subsubsection{0-1 BFS}

\paragraph{Idea intuitiva.}
Variante de BFS para grafos con \textbf{pesos 0 o 1}. Usa un \texttt{deque}: aristas de peso 0 se encolan al frente (\texttt{push\_front}), aristas de peso 1 al final (\texttt{push\_back}). Asi, la primera vez que se extrae un nodo, su distancia es la definitiva, igual que en Dijkstra pero en $O(V + E)$. La intuicion: el deque mantiene los nodos ordenados por distancia; las aristas de peso 0 no cambian la distancia (entran al frente), las de peso 1 la incrementan (entran al final).

\paragraph{Propiedades clave (invariantes).}
\begin{itemize}\setlength\itemsep{0pt}
	\item Asume pesos \textbf{solo} 0 o 1.
	\item Mejor caso que Dijkstra (sin factor $\log V$).
	\item No funciona con pesos $\ge 2$.
	\item La cola mantiene los nodos ordenados por distancia no decreciente.
\end{itemize}

\paragraph{Codigo clave.}
El deque doble:
\begin{verbatim}
deque<int> q;
q.push_front(source);
dist[source] = 0;
while (!q.empty()) {
    int v = q.front(); q.pop_front();
    for (auto [w, u] : adj[v]) {
        if (dist[v] + w < dist[u]) {
            dist[u] = dist[v] + w;
            if (w == 0) q.push_front(u);
            else q.push_back(u);
        }
    }
}
\end{verbatim}

\paragraph{Complejidad.}
\begin{itemize}\setlength\itemsep{0pt}
	\item $O(V + E)$.
	\item Cada nodo se inserta a lo sumo $O(1)$ veces (con check de distancia).
\end{itemize}

\paragraph{Donde tocar para modificar.}
\begin{itemize}\setlength\itemsep{0pt}
	\item \textbf{Pesos en $[0, k]$ small}: usar small k-ary BFS o Dijkstra.
	\item \textbf{Reconstruir camino}: aniadir \texttt{parent[]}.
	\item \textbf{Pesos negativos}: no funciona, usar Bellman-Ford.
\end{itemize}

\paragraph{Trampas y errores frecuentes.}
\begin{itemize}\setlength\itemsep{0pt}
	\item Si los pesos son otro rango, el algoritmo da respuestas incorrectas. Verificar que las aristas son solo 0/1.
	\item El chequeo \texttt{dist[v] + w < dist[u]} es crucial; sino, el algoritmo procesa nodos multiples veces.
	\item Si una arista tiene peso $\ge 2$, aniadir verificaciones o cambiar a Dijkstra.
\end{itemize}

\paragraph{Codigo.}
Ver \texttt{cpp.tex}, seccion \textit{0-1 BFS}.

\subsection{MST}

\subsubsection{Kruskal}

\paragraph{Idea intuitiva.}
Encuentra el \textbf{arbol de expansion minima} (MST) en $O(E \log E)$. Ordena las aristas por peso ascendente y las aniade al MST si no forman ciclo (chequear con DSU). El arbol tiene exactamente $V - 1$ aristas. La demostracion (``Cut Property''): para cualquier corte, la arista minima cruzando el corte pertenece a \emph{algun} MST; por lo tanto, elegir la minima disponible que no forma ciclo es optimo.

\paragraph{Propiedades clave (invariantes).}
\begin{itemize}\setlength\itemsep{0pt}
	\item Algoritmo \textbf{greedy} correcto (Cut Property).
	\item Funciona en grafos \textbf{no conexos}; produce un \textbf{minimum spanning forest}.
	\item Si las aristas tienen pesos unicos, el MST es unico.
	\item El MST tiene exactamente $V - 1$ aristas (o menos si el grafo es disconexo).
\end{itemize}

\paragraph{Codigo clave.}
Kruskal con DSU:
\begin{verbatim}
sort(edges.begin(), edges.end());
DSU dsu(n);
long long mst = 0;
int edgesUsed = 0;
for (auto [w, u, v] : edges)
    if (dsu.unite(u, v)) {
        mst += w;
        if (++edgesUsed == n - 1) break;
    }
\end{verbatim}

\paragraph{Complejidad.}
\begin{itemize}\setlength\itemsep{0pt}
	\item $O(E \log E)$ (ordenamiento) o $O(E \log V)$ con DSU optimizations.
	\item DSU es casi $O(1)$ amortizado, asi que el sort domina.
\end{itemize}

\paragraph{Donde tocar para modificar.}
\begin{itemize}\setlength\itemsep{0pt}
	\item \textbf{Second best MST}: enumerar aristas del MST, sacar una, aniadir una no-MST; recalcular.
	\item \textbf{Kruskal randomizado}: si el grafo es casi-completo, agregar shuffle.
	\item \textbf{Aniadir info al DSU}: guardar el peso maximo en el camino, etc.
	\item \textbf{DSU con rollback}: si necesitas Kruskal reversible (para queries offline).
\end{itemize}

\paragraph{Trampas y errores frecuentes.}
\begin{itemize}\setlength\itemsep{0pt}
	\item El DSU debe estar limpio para cada invocacion.
	\item Si las aristas tienen pesos negativos, Kruskal sigue funcionando.
	\item Si el grafo es disconexo, el ``MST'' es un bosque con varias componentes.
	\item El check \texttt{if (++edgesUsed == n - 1) break;} optimiza la salida temprana.
\end{itemize}

\paragraph{Codigo.}
Ver \texttt{cpp.tex}, seccion \textit{Kruskal}.

\vspace{0.5em}

\subsubsection{Prim (priority queue)}

\paragraph{Idea intuitiva.}
Otra forma de calcular el MST: empezar desde un nodo y, en cada paso, aniadir la arista de menor peso que conecta el conjunto construido con el resto. Usando priority queue, es $O((V + E) \log V)$. La demostracion: cada vez que se aniade una arista, es la minima cruzando el ``corte'' entre el conjunto actual y el resto (Cut Property); eso preserva optimalidad.

\paragraph{Propiedades clave (invariantes).}
\begin{itemize}\setlength\itemsep{0pt}
	\item Similar a Dijkstra pero sin acumular distancias, solo eligiendo minimo por arista.
	\item Funciona mejor que Kruskal en \textbf{grafos densos} ($E \approx V^2$).
	\item Si el grafo es disconexo, genera el bosque (solo llega a los alcanzables).
	\item Cada nodo se inserta al heap una vez; cada arista se procesa una vez.
\end{itemize}

\paragraph{Codigo clave.}
Prim clasico con heap:
\begin{verbatim}
priority_queue<pair<int, int>, vector<...>, greater<...>> pq;
vector<bool> used(n, false);
vector<int> minEdge(n, INF);
minEdge[0] = 0;
pq.push({0, 0});
while (!pq.empty()) {
    auto [w, v] = pq.top(); pq.pop();
    if (used[v]) continue;
    used[v] = true;
    mst += w;
    for (auto [peso, u] : adj[v])
        if (!used[u] && peso < minEdge[u]) {
            minEdge[u] = peso;
            pq.push({peso, u});
        }
}
\end{verbatim}

\paragraph{Complejidad.}
\begin{itemize}\setlength\itemsep{0pt}
	\item $O((V + E) \log V)$.
	\item Con matriz de adyacencia: $O(V^2)$ sin heap.
\end{itemize}

\paragraph{Donde tocar para modificar.}
\begin{itemize}\setlength\itemsep{0pt}
	\item \textbf{Prim con matriz de adyacencia}: $O(V^2)$ sin heap, util en grafos densos pequenos.
	\item \textbf{Reconstruir el MST}: aniadir \texttt{parent[]}.
	\item \textbf{Prim en streaming}: variante para grafos que llegan incrementalmente.
\end{itemize}

\paragraph{Trampas y errores frecuentes.}
\begin{itemize}\setlength\itemsep{0pt}
	\item El \texttt{if (used[v]) continue;} es crucial: si el nodo ya fue procesado, saltar.
	\item En la version matriz, elije el minimo por linea (no heap).
	\item Si el grafo es disconexo, aniadir un loop externo sobre todos los nodos no usados.
\end{itemize}

\paragraph{Codigo.}
Ver \texttt{cpp.tex}, seccion \textit{Prim (priority queue)}.

\subsection{Flujos}

\subsubsection{Edmonds-Karp (max flow)}

\paragraph{Idea intuitiva.}
Implementacion BFS del algoritmo de \textbf{Ford-Fulkerson}. En cada iteracion, hacer BFS desde $s$ hasta $t$ siguiendo aristas con capacidad residual $> 0$; si se llega a $t$, se encontro un \textbf{camino de aumento}. Enviar el minimo de las capacidades en el camino; actualizar capacidades residuales. La demostracion de $O(VE^2)$: cada BFS encuentra el augmenting path mas corto, y la longitud minima crece monotona; como hay $\le E$ longitudes posibles, hay $\le E$ fases.

\paragraph{Propiedades clave (invariantes).}
\begin{itemize}\setlength\itemsep{0pt}
	\item BFS garantiza caminos de aumento mas cortos.
	\item Capacidad residual: forward edge baja, backward edge sube.
	\item Complejidad $O(VE^2)$.
	\item El flujo maximo respeta la conservacion en nodos (in = out salvo $s, t$).
\end{itemize}

\paragraph{Codigo clave.}
BFS + augmenting path:
\begin{verbatim}
while (bfs(s, t, parent)) {  // bfs encuentra camino o retorna false
    int pathFlow = INF;
    for (int v = t; v != s; v = parent[v])
        pathFlow = min(pathFlow, capacity[parent[v]][v]);
    for (int v = t; v != s; v = parent[v]) {
        capacity[parent[v]][v] -= pathFlow;
        capacity[v][parent[v]] += pathFlow;
    }
    maxFlow += pathFlow;
}
\end{verbatim}

\paragraph{Complejidad.}
\begin{itemize}\setlength\itemsep{0pt}
	\item $O(VE^2)$.
	\item En la practica, mas lento que Dinic para grafos grandes.
\end{itemize}

\paragraph{Donde tocar para modificar.}
\begin{itemize}\setlength\itemsep{0pt}
	\item \textbf{Grafo bipartito}: $O(\sqrt{V} E)$ (equivalente a Hopcroft-Karp).
	\item \textbf{Dinic} es preferible en la mayoria de los casos.
	\item \textbf{Capacidades doubles}: usar tolerancia para comparacion.
	\item \textbf{Min-cost flow}: aniadir costo por arista y modificar el algoritmo.
\end{itemize}

\paragraph{Trampas y errores frecuentes.}
\begin{itemize}\setlength\itemsep{0pt}
	\item Las capacidades deben ser $> 0$ para entrar a la cola (no $>= 0$); sino, ciclos infinitos.
	\item Para grafos bipartitos, Dinic o Hopcroft-Karp son mejores.
	\item Si el flujo maximo es muy grande, los \texttt{int} pueden overflow; usar \texttt{long long}.
\end{itemize}

\paragraph{Codigo.}
Ver \texttt{cpp.tex}, seccion \textit{Edmonds-Karp (max flow)}.

\vspace{0.5em}

\subsubsection{Dinic (max flow)}

\paragraph{Idea intuitiva.}
Algoritmo de flujo mas eficiente que Edmonds-Karp: combina \textbf{BFS para level graph} con \textbf{DFS bloqueante} (multiple augmenting paths en una sola pasada). Cada fase cuesta $O(E)$ y hay $O(V)$ fases en el peor caso; para bipartitos, $O(\sqrt{V} E)$. La razon: cada fase ``satura'' al menos una arista, asi que en $O(E)$ fases el flujo maximo se alcanza; con $O(V)$ niveles en el BFS, hay $O(VE)$ operaciones totales.

\paragraph{Propiedades clave (invariantes).}
\begin{itemize}\setlength\itemsep{0pt}
	\item \texttt{level[v]} = distancia desde $s$ en el level graph (solo aristas con cap $> 0$).
	\item \texttt{it[v]} = iterador actual para evitar reprocesar aristas.
	\item DFS solo sigue aristas con \texttt{level[e.to] == level[v] + 1}.
	\item Back-edges tienen \texttt{rev} apuntando a la forward edge original.
\end{itemize}

\paragraph{Codigo clave.}
El DFS bloqueante de Dinic:
\begin{verbatim}
long long dfs(int v, int t, long long f) {
    if (v == t) return f;
    for (int& i = it[v]; i < adj[v].size(); ++i) {
        Edge& e = adj[v][i];
        if (e.cap > 0 && level[e.to] == level[v] + 1) {
            long long ret = dfs(e.to, t, min(f, e.cap));
            if (ret > 0) {
                e.cap -= ret;
                adj[e.to][e.rev].cap += ret;
                return ret;
            }
        }
    }
    return 0;
}
\end{verbatim}

\paragraph{Complejidad.}
\begin{itemize}\setlength\itemsep{0pt}
	\item $O(V^2 E)$ general, $O(\sqrt{V} E)$ para bipartitos.
	\item En la practica, mucho mas rapido que Edmonds-Karp.
\end{itemize}

\paragraph{Donde tocar para modificar.}
\begin{itemize}\setlength\itemsep{0pt}
	\item \textbf{Dinic con scaling}: aniadir factor de escala para capacidades grandes.
	\item \textbf{Lower bounds}: aniadir un super-source/sink para forzar flujo minimo.
	\item \textbf{Matching bipartito}: el grafo se construye como source -> L -> R -> sink, todos con capacidad 1.
	\item \textbf{Push-relabel}: alternativa para grafos muy densos.
\end{itemize}

\paragraph{Trampas y errores frecuentes.}
\begin{itemize}\setlength\itemsep{0pt}
	\item \texttt{it[v]} debe \textbf{incrementarse} en cada iteracion; sino, el DFS se queda atascado en la misma arista.
	\item Si el BFS no encuentra $t$, el flujo es maximo; terminar.
	\item Las back-edges son cruciales para ``cancelar'' elecciones previas.
\end{itemize}

\paragraph{Codigo.}
Ver \texttt{cpp.tex}, seccion \textit{Dinic (max flow)}.

\vspace{0.5em}

\subsubsection{Min-Cost Max-Flow}

\paragraph{Idea intuitiva.}
Encuentra el flujo maximo con \textbf{costo minimo} (asumiendo costos no negativos). Usa \textbf{SPFA} (o Dijkstra con potentials) para encontrar el shortest path en el \textbf{costo} desde $s$ hasta $t$ considerando capacidades residuales. Cada augmenting path aporta su costo $\times$ flujo. La idea: los potentials reescalan las aristas para que todas tengan costo no negativo, permitiendo usar Dijkstra.

\paragraph{Propiedades clave (invariantes).}
\begin{itemize}\setlength\itemsep{0pt}
	\item \texttt{pot[]} = potentials para evitar aristas negativas (Johnson's trick).
	\item \texttt{dist[v]} = costo minimo desde $s$ con potentials.
	\item Solo se aumentan aristas con \texttt{cap > 0} (forward) y costo es \texttt{+cost} (forward) o \texttt{-cost} (backward).
	\item Tras cada augmenting, los potentials se actualizan para mantener la consistencia.
\end{itemize}

\paragraph{Codigo clave.}
El ciclo de min-cost flow:
\begin{verbatim}
while (spfa(s, t, dist, parent)) {
    int pathFlow = INF;
    for (int v = t; v != s; v = parent[v])
        pathFlow = min(pathFlow, capacity[parent[v]][v]);
    for (int v = t; v != s; v = parent[v]) {
        capacity[parent[v]][v] -= pathFlow;
        capacity[v][parent[v]] += pathFlow;
        cost += pathFlow * edge_cost[parent[v]][v];
    }
}
\end{verbatim}

\paragraph{Complejidad.}
\begin{itemize}\setlength\itemsep{0pt}
	\item $O(F \cdot V E)$ con SPFA, $O(F \cdot E \log V)$ con Dijkstra + potentials.
	\item $F$ = flujo maximo total.
\end{itemize}

\paragraph{Donde tocar para modificar.}
\begin{itemize}\setlength\itemsep{0pt}
	\item \textbf{Dijkstra + potentials}: si los costos son enteros no negativos, esto es mas estable.
	\item \textbf{Costos negativos en aristas}: requieren Bellman-Ford inicial o reorden.
	\item \textbf{Flujo minimo con costo}: encontrar el minimo flujo posible y su costo.
	\item \textbf{Costo por unidad}: si los costos son muy grandes, aniadir verificacion.
\end{itemize}

\paragraph{Trampas y errores frecuentes.}
\begin{itemize}\setlength\itemsep{0pt}
	\item El \texttt{pot[]} se actualiza \textbf{al final} de cada iteracion (no durante).
	\item Las back-edges tienen \textbf{costo negativo}, esencial para corregir elecciones previas.
	\item Si los costos son muy grandes, considerar modular arithmetic.
	\item Para problemas con flujo exacto, aniadir verificacion post-algoritmo.
\end{itemize}

\paragraph{Codigo.}
Ver \texttt{cpp.tex}, seccion \textit{Min-Cost Max-Flow}.

\subsection{Especales}

\subsubsection{2-SAT (implication graph + SCC)}

\paragraph{Idea intuitiva.}
Resuelve formulas booleanas en \textbf{CNF} con clausulas de \textbf{2 literales}: $(l_1 \lor l_2)$. Cada variable $x$ se representa con dos nodos: $x$ (verdadero) y $\neg x$ (falso). Cada clausula $l_1 \lor l_2$ se descompone en dos implicaciones: $\neg l_1 \to l_2$ y $\neg l_2 \to l_1$. Si en el grafo de implicaciones $\neg x$ y $x$ estan en la misma SCC, la formula es \textbf{UNSAT}. La demostracion: si $\neg x \to x$ en el grafo, entonces $\neg x$ implica $x$, asi que $\neg x$ debe ser falso y $x$ verdadero; pero entonces $\neg x$ es falso implica una contradiccion.

\paragraph{Propiedades clave (invariantes).}
\begin{itemize}\setlength\itemsep{0pt}
	\item Cada literal es un nodo; total $2n$ nodos.
	\item Kosaraju/Tarjan dan las SCCs.
	\item Si SAT: asignar $x = $ (topological order del condensation).
	\item Clausulas \texttt{setTrue(x)} se modelan como \texttt{implies(!x, x)}.
	\item La formula es SAT sii no hay variable $x$ con $x$ y $\neg x$ en la misma SCC.
\end{itemize}

\paragraph{Codigo clave.}
Conversion de clausulas a implicaciones:
\begin{verbatim}
// Clausula (a OR b) -> implica ~a -> b Y ~b -> a
addImplication(not(a), b);
addImplication(not(b), a);
// Forzar x: implica ~x -> x
if (force) addImplication(not(x), x);
\end{verbatim}

\paragraph{Complejidad.}
\begin{itemize}\setlength\itemsep{0pt}
	\item $O(V + E) = O(N + M)$ con Kosaraju o Tarjan.
	\item Memoria $O(V + E)$.
\end{itemize}

\paragraph{Donde tocar para modificar.}
\begin{itemize}\setlength\itemsep{0pt}
	\item \textbf{Forzar variables}: aniadir \texttt{setTrue} / \texttt{setFalse}.
	\item \textbf{3-SAT}: NP-completo; no usar este algoritmo.
	\item \textbf{Contar soluciones}: las SCCs condensation dan una formula para contar; multiplicar las formas de asignar variables en cada componente no-trivial.
	\item \textbf{Output de la asignacion}: aniadir el vector \texttt{assignment[]}.
\end{itemize}

\paragraph{Trampas y errores frecuentes.}
\begin{itemize}\setlength\itemsep{0pt}
	\item La asignacion usa el \textbf{orden topologico inverso} del condensation: \texttt{val[i] = comp[2i] < comp[2i+1]}.
	\item El ID de $\neg x$ es \texttt{2*x + 1} (o similar); verificar consistencia.
	\item Si hay clausulas con literales repetidos (e.g., $(x \lor x)$), tratar como $(x)$ simple.
\end{itemize}

\paragraph{Codigo.}
Ver \texttt{cpp.tex}, seccion \textit{2-SAT (implication graph + SCC)}.

\vspace{0.5em}

\subsubsection{Euler tour / Hierholzer (camino y circuito)}

\paragraph{Idea intuitiva.}
Encuentra un \textbf{camino o circuito euleriano} (que usa cada arista exactamente una vez). \textbf{Condiciones}:
\begin{itemize}\setlength\itemsep{0pt}
	\item \textbf{No dirigido}: todos los grados pares (circuito) o exactamente 2 impares (camino, empieza en uno, termina en el otro). Ademas, conexo sobre el subgrafo con aristas.
	\item \textbf{Dirigido}: $\text{indeg}(v) = \text{outdeg}(v)$ para todos (circuito) o exactamente un nodo con $\text{outdeg} = \text{indeg} + 1$ y otro con $\text{indeg} = \text{outdeg} + 1$ (camino).
\end{itemize}

\paragraph{Hierholzer}: elegir un ciclo desde el inicio, mientras haya ciclos ``laterales'' desde un nodo del camino, mergearlos. Implementacion con pila: avanzar por aristas, al no poder mas aniadir a la respuesta y backtrack. La demostracion de $O(V + E)$: cada arista se procesa una vez (entra y sale del stack una vez).

\paragraph{Propiedades clave (invariantes).}
\begin{itemize}\setlength\itemsep{0pt}
	\item Multigrafos permitidos; las aristas se marcan como usadas.
	\item No dirigido: cada arista aparece \textbf{dos veces} en \texttt{adj} (ida y vuelta).
	\item La construccion del path con ids consistentes es importante.
	\item El resultado es unico salvo orden de las aristas en multigrafos.
\end{itemize}

\paragraph{Codigo clave.}
Hierholzer con pila:
\begin{verbatim}
vector<int> path;
stack<int> st;
st.push(start);
while (!st.empty()) {
    int v = st.top();
    while (!adj[v].empty() && used[adj[v].back()]) adj[v].pop_back();
    if (adj[v].empty()) {
        path.push_back(v);
        st.pop();
    } else {
        int e = adj[v].back();
        used[e] = true;
        int u = edges[e].to(v);
        adj[v].pop_back();
        st.push(u);
    }
}
reverse(path.begin(), path.end());
\end{verbatim}

\paragraph{Complejidad.}
\begin{itemize}\setlength\itemsep{0pt}
	\item $O(V + E)$.
	\item Memoria $O(V + E)$.
\end{itemize}

\paragraph{Donde tocar para modificar.}
\begin{itemize}\setlength\itemsep{0pt}
	\item \textbf{Reconstruir con multigrafo}: aniadir ids consistentes antes del algoritmo.
	\item \textbf{Aplicar a ``de Bruijn'' o problemas de strings}: modelar como multigrafo y buscar Euler.
	\item \textbf{Dirigido vs no dirigido}: el snippet tiene ambas versiones.
	\item \textbf{Letter addition}: modelar palabras como aristas y buscar el superstring.
\end{itemize}

\paragraph{Trampas y errores frecuentes.}
\begin{itemize}\setlength\itemsep{0pt}
	\item Si el grafo no es conexo (sobre las aristas), no existe euleriano. Verificar primero.
	\item Para multigrafos, los IDs de aristas son cruciales para distinguirlas.
	\item Si el problema es ``encuentra todos los euler tours'', el algoritmo es exponencial.
\end{itemize}

\paragraph{Codigo.}
Ver \texttt{cpp.tex}, seccion \textit{Euler tour / Hierholzer (camino y circuito)}.

\vspace{0.5em}

\subsubsection{Hamiltonian path / cycle (bitmask DP + backtracking)}

\paragraph{Idea intuitiva.}
Un \textbf{camino hamiltoniano} visita cada vertice \textbf{exactamente una vez}. \textbf{Determinar} si existe es NP-completo en general, pero con $V \le 20$ se puede hacer DP con bitmask en $O(V^2 \cdot 2^V)$. \textbf{Teoremas suficientes}: \textbf{Dirac} (grado $\ge n/2$) y \textbf{Ore} ($\text{deg}(u) + \text{deg}(v) \ge n$ para todo par no adyacente). La idea de la DP: cada estado es (conjunto de visitados, ultimo vertice); la transicion es agregar un vecino no visitado.

\paragraph{Propiedades clave (invariantes).}
\begin{itemize}\setlength\itemsep{0pt}
	\item DP: $dp[mask][v]$ = true si hay un camino que visita exactamente \texttt{mask} y termina en $v$.
	\item TSP: variante con pesos minimos (Hamilton cycle minimo).
	\item Reconstruccion: guardar \texttt{par[mask][v]} (predecesor).
	\item $2^V$ estados con $V$ bits cada uno; factible para $V \le 20$.
\end{itemize}

\paragraph{Codigo clave.}
DP con bitmask:
\begin{verbatim}
dp[1][0] = 0;  // empezar en 0
for (int mask = 1; mask < (1<<n); ++mask)
    for (int v = 0; v < n; ++v) if (dp[mask][v] < INF)
        for (int u = 0; u < n; ++u)
            if (!(mask & (1<<u)))
                dp[mask | (1<<u)][u] = min(dp[mask | (1<<u)][u],
                                            dp[mask][v] + w[v][u]);
// TSP: respuesta = min(dp[(1<<n)-1][v] + w[v][0])
\end{verbatim}

\paragraph{Complejidad.}
\begin{itemize}\setlength\itemsep{0pt}
	\item $O(V^2 \cdot 2^V)$ tiempo, $O(V \cdot 2^V)$ memoria.
	\item Para $V = 20$: $\sim 4 \cdot 10^8$ operaciones (lento pero factible).
\end{itemize}

\paragraph{Donde tocar para modificar.}
\begin{itemize}\setlength\itemsep{0pt}
	\item \textbf{TSP asimetrico}: aniadir pesos $w[v][u]$ (no necesariamente simetricos).
	\item \textbf{Hamilton con inicio obligatorio}: fijar \texttt{src} en el DP inicial.
	\item \textbf{Bitmask DP mas compacto}: usar \texttt{long long} como mascara si $V \le 64$.
	\item \textbf{TSP optimo aproximado}: MST, 2-opt, simulated annealing.
	\item \textbf{Reconstruir el path}: usar \texttt{par[mask][v]} y backtrack.
\end{itemize}

\paragraph{Trampas y errores frecuentes.}
\begin{itemize}\setlength\itemsep{0pt}
	\item TSP asume grafo \textbf{completo} (con \texttt{LINF} para aristas faltantes); si no, aniadir chequeo.
	\item La condicion ``\texttt{mask == (1<<n) - 1}'' cierra el ciclo retornando al origen.
	\item Para $V > 20$, la memoria $2^V$ se vuelve prohibitiva; usar meet-in-the-middle o heuristicas.
\end{itemize}

\paragraph{Codigo.}
Ver \texttt{cpp.tex}, seccion \textit{Hamiltonian path / cycle (bitmask DP + backtracking)}.

% ============================================================
\section{DP}

\subsubsection{Knapsack 0/1}

\paragraph{Idea intuitiva.}
Dado un conjunto de $N$ objetos con peso $w_i$ y valor $v_i$, maximizar $\sum v_i$ con la restriccion $\sum w_i \le W$. La recurrencia clasica:
\[ dp[i][j] = \max(dp[i-1][j],\, dp[i-1][j - w_i] + v_i) \quad \text{si } j \ge w_i. \]
La optimizacion clave: la DP es 1-dimensional si iteramos $j$ \textbf{de mayor a menor} (asi no usamos el objeto $i$ dos veces). La razon de la iteracion inversa: en cada iteracion queremos el estado del ``item anterior''; iterar de menor a mayor contaminaria \texttt{dp[j - w_i]} con el item actual.

\paragraph{Propiedades clave (invariantes).}
\begin{itemize}\setlength\itemsep{0pt}
	\item $dp[j]$ al final = maximo valor con peso $\le j$.
	\item El bucle inverso \texttt{for (j = W; j >= w[i]; j--)} es lo que distingue 0/1 de unbounded.
	\item Solucion: maximo valor con cualquier peso $\le W$.
	\item La DP es exacta; no hay aproximacion.
\end{itemize}

\paragraph{Codigo clave.}
La version 1D optimizada:
\begin{verbatim}
vector<long long> dp(W + 1, 0);
for (int i = 0; i < N; ++i)
    for (int j = W; j >= w[i]; --j)  // invertido
        dp[j] = max(dp[j], dp[j - w[i]] + v[i]);
return dp[W];
\end{verbatim}

\paragraph{Complejidad.}
\begin{itemize}\setlength\itemsep{0pt}
	\item $O(N \cdot W)$ tiempo, $O(W)$ memoria.
	\item Constante muy pequena (solo sumas y max).
\end{itemize}

\paragraph{Donde tocar para modificar.}
\begin{itemize}\setlength\itemsep{0pt}
	\item \textbf{Reconstruir la seleccion}: aniadir \texttt{take[i][j]} (booleano o bool) en una version 2D.
	\item \textbf{Items con peso 0}: edge case; aniadir guard.
	\item \textbf{Knapsack con restricciones de grupos}: ``elegir a lo sumo 1 de cada grupo'' se modela como 0/1 con un peso extra.
	\item \textbf{Tamano de $W$}: si $W \le 10^9$ pero $N$ es chico, usar meet-in-the-middle.
	\item \textbf{Cambiar max por min}: aniadir \texttt{INF} en la inicializacion.
\end{itemize}

\paragraph{Trampas y errores frecuentes.}
\begin{itemize}\setlength\itemsep{0pt}
	\item Overflow en \texttt{dp[j - w[i]] + v[i]}; usar \texttt{long long} si los valores son grandes.
	\item El bucle \textbf{debe} ir de mayor a menor.
	\item Si los pesos son grandes y $N$ pequeno, conviene meet-in-the-middle.
	\item Para Knapsack exacto ($\sum w_i = W$), aniadir check al final.
\end{itemize}

\paragraph{Codigo.}
Ver \texttt{cpp.tex}, seccion \textit{Knapsack 0/1}.

\vspace{0.5em}

\subsubsection{Knapsack unbounded}

\paragraph{Idea intuitiva.}
Igual que Knapsack 0/1 pero \textbf{cada item se puede tomar multiples veces}. La recurrencia es similar:
\[ dp[j] = \max(dp[j],\, dp[j - w_i] + v_i), \]
pero ahora el bucle va de \textbf{menor a mayor} (porque tomar el item $i$ no afecta corridas futuras del mismo item). La idea: tras tomar el item $i$, queremos volver a tener la opcion de tomarlo; iterar de menor a mayor permite que el item ``se acumule''.

\paragraph{Propiedades clave (invariantes).}
\begin{itemize}\setlength\itemsep{0pt}
	\item El bucle \texttt{for (j = w[i]; j <= W; j++)} es lo que distingue unbounded.
	\item A diferencia de 0/1, no hay riesgo de ``usar dos veces el mismo item''.
	\item La DP es exacta; cada ``subconjunto'' se cuenta correctamente.
\end{itemize}

\paragraph{Codigo clave.}
La version 1D con bucle directo:
\begin{verbatim}
vector<long long> dp(W + 1, 0);
for (int i = 0; i < N; ++i)
    for (int j = w[i]; j <= W; ++j)  // directo (no invertido)
        dp[j] = max(dp[j], dp[j - w[i]] + v[i]);
return dp[W];
\end{verbatim}

\paragraph{Complejidad.}
\begin{itemize}\setlength\itemsep{0pt}
	\item $O(N \cdot W)$ tiempo, $O(W)$ memoria.
	\item Constante muy pequena.
\end{itemize}

\paragraph{Donde tocar para modificar.}
\begin{itemize}\setlength\itemsep{0pt}
	\item \textbf{Reconstruir cuantos de cada item}: aniadir \texttt{cnt[i]} y backtrack.
	\item \textbf{Cambiar moneda minima por maxima}: el snippet maximiza valor; para minimizar monedas, cambiar a \texttt{min} con inicializacion en \texttt{INF}.
	\item \textbf{Coin change}: si todos los pesos son positivos y queremos ``minimo numero de monedas'', es unbounded con min.
\end{itemize}

\paragraph{Trampas y errores frecuentes.}
\begin{itemize}\setlength\itemsep{0pt}
	\item Misma que Knapsack 0/1 con overflow.
	\item El orden de iteracion de items no afecta el resultado (convexo).
	\item Si se quiere ``maximo valor con exactamente $k$ items'', aniadir segunda dimension.
\end{itemize}

\paragraph{Codigo.}
Ver \texttt{cpp.tex}, seccion \textit{Knapsack unbounded}.

\vspace{0.5em}

\subsubsection{Knapsack meet-in-the-middle}

\paragraph{Idea intuitiva.}
Para $N \le 40$ y $W$ moderado, no se puede hacer $O(NW)$ directamente. Partir en dos mitades, enumerar todos los $2^{N/2}$ subconjuntos de cada mitad, y combinar. Para cada subconjunto de la segunda mitad con peso $sw_2$, encontrar el maximo valor de la primera mitad con peso $\le W - sw_2$. La idea: ``divide y venceras'' sobre el tamano de la entrada; cada mitad tiene mitad del espacio, asi que la combinacion es geometrica en lugar de multiplicativa.

\paragraph{Propiedades clave (invariantes).}
\begin{itemize}\setlength\itemsep{0pt}
	\item Hay $2^{N/2}$ subconjuntos por mitad; total $O(2^{N/2})$.
	\item \texttt{sums} guarda los valores de subconjuntos de la primera mitad ordenados (o preprocesados).
	\item Para cada subconjunto de la segunda mitad, busqueda lineal o binaria en \texttt{sums}.
	\item La mitad mas grande (con $N/2 + 1$ items para $N$ impar) tiene $2^{N/2+1}$ subconjuntos.
\end{itemize}

\paragraph{Codigo clave.}
Enumerar subconjuntos de la primera mitad:
\begin{verbatim}
int n1 = N / 2, n2 = N - n1;
vector<pair<long long, long long>> sums1;
for (int mask = 0; mask < (1 << n1); ++mask) {
    long long peso = 0, valor = 0;
    for (int i = 0; i < n1; ++i) if (mask & (1 << i)) {
        peso += w[i]; valor += v[i];
    }
    sums1.push_back({peso, valor});
}
sort(sums1.begin(), sums1.end());  // por peso
// eliminar duplicados dominados
\end{verbatim}

\paragraph{Complejidad.}
\begin{itemize}\setlength\itemsep{0pt}
	\item $O(2^{N/2})$ tiempo y memoria (mejor que $O(NW)$ cuando $W > 2^{N/2}$).
	\item Para $N = 40$: $\sim 2 \cdot 10^6$ subconjuntos, factible.
\end{itemize}

\paragraph{Donde tocar para modificar.}
\begin{itemize}\setlength\itemsep{0pt}
	\item \textbf{Tamano mayor}: si $N$ es mayor pero $W$ es chico, conviene iterar por peso, no por subconjuntos.
	\item \textbf{Varias mochilas}: adaptar para multiples $W$.
	\item \textbf{Optimizar el ``combinado''}: usar \texttt{upper\_bound} para encontrar el maximo valor con peso $\le \text{rem}.
	\item \textbf{Eliminar pares dominados}: si un par (peso, valor) tiene otro con menos peso y mas valor, eliminarlo.
\end{itemize}

\paragraph{Trampas y errores frecuentes.}
\begin{itemize}\setlength\itemsep{0pt}
	\item $N$ impar: una mitad tiene $N/2 + 1$ items.
	\item Overflow al sumar pesos/valores en \texttt{long long}.
	\item Sin deduplicacion, el algoritmo es correcto pero mas lento.
\end{itemize}

\paragraph{Codigo.}
Ver \texttt{cpp.tex}, seccion \textit{Knapsack meet-in-the-middle}.

\vspace{0.5em}

\subsubsection{LIS O(n log n)}

\paragraph{Idea intuitiva.}
La \textbf{Longest Increasing Subsequence} se calcula en $O(n \log n)$ con el truco del \textbf{patience sorting}: mantener un arreglo \texttt{dp} donde \texttt{dp[len]} es el minimo ultimo valor de una IS de largo \texttt{len}. Para cada $x$, encontrar el primer \texttt{dp[len] >= x} y reemplazarlo. El largo final es la LIS. La demostracion: si hay dos IS del mismo largo, la que termina en un valor menor es ``mejor'' (puede extenderse mas facilmente).

\paragraph{Propiedades clave (invariantes).}
\begin{itemize}\setlength\itemsep{0pt}
	\item \texttt{dp} es estrictamente creciente.
	\item \texttt{lower\_bound} da LIS estricta; \texttt{upper\_bound} da no-estricta ($\le$).
	\item El algoritmo es \textbf{online}: procesa elementos en orden.
	\item La longitud de \texttt{dp} es la LIS al final.
\end{itemize}

\paragraph{Codigo clave.}
La version online:
\begin{verbatim}
vector<int> dp;
for (int x : a) {
    auto it = lower_bound(dp.begin(), dp.end(), x);
    if (it == dp.end()) dp.push_back(x);
    else *it = x;
}
// dp.size() = longitud de LIS estricta
\end{verbatim}

\paragraph{Complejidad.}
\begin{itemize}\setlength\itemsep{0pt}
	\item $O(n \log n)$.
	\item Memoria $O(n)$.
\end{itemize}

\paragraph{Donde tocar para modificar.}
\begin{itemize}\setlength\itemsep{0pt}
	\item \textbf{Reconstruir la LIS}: aniadir \texttt{parent[]} y \texttt{posInDp[]} para backtrack.
	\item \textbf{Contar LIS}: DP clasica $O(n^2)$ o segment tree sobre $dp$.
	\item \textbf{LDS / Longest non-increasing}: invertir la condicion en \texttt{lower\_bound} o multiplicar por -1.
	\item \textbf{Longest bitonic subsequence}: combinar LIS desde izq y desde der.
	\item \textbf{LIS en 2D}: ordenar por una dimension y aplicar LIS en la otra.
\end{itemize}

\paragraph{Trampas y errores frecuentes.}
\begin{itemize}\setlength\itemsep{0pt}
	\item \texttt{lower\_bound} vs \texttt{upper\_bound}: si tu comparacion es $\le$ vs $<$, cambia.
	\item Si los elementos son \texttt{int} pero podrian ser negativos, aniadir offset para usar como indices.
	\item Para LIS reconstruida, aniadir el array \texttt{parent} que apunta al anterior en la IS.
\end{itemize}

\paragraph{Codigo.}
Ver \texttt{cpp.tex}, seccion \textit{LIS O(n log n)}.

\vspace{0.5em}

\subsubsection{LCS / Edit distance}

\paragraph{Idea intuitiva.}
\textbf{LCS} (Longest Common Subsequence): para dos strings $a$, $b$, $dp[i][j]$ = LCS de $a[0..i)$ y $b[0..j)$. Recurrencia:
\[ dp[i][j] = \begin{cases} dp[i-1][j-1] + 1 & \text{si } a[i-1] = b[j-1] \\ \max(dp[i-1][j], dp[i][j-1]) & \text{si no} \end{cases} \]

\textbf{Edit distance} (Levenshtein): numero minimo de operaciones (insert, delete, replace) para convertir $a$ en $b$. Recurrencia similar con un $+1$ adicional. La idea: ``alinear'' los strings eligiendo matches/mismatches/insert/delete.

\paragraph{Propiedades clave (invariantes).}
\begin{itemize}\setlength\itemsep{0pt}
	\item LCS y Edit Distance son DP $O(nm)$.
	\item Memoria se puede reducir a $O(\min(n, m))$ con rolling array.
	\item Edit distance con peso por operacion distinto: ajustar el \texttt{+1} por el peso.
	\item La subestructura optima: la respuesta para prefijos se construye de prefijos menores.
\end{itemize}

\paragraph{Codigo clave.}
La recurrencia clasica de LCS:
\begin{verbatim}
vector<vector<int>> dp(n + 1, vector<int>(m + 1, 0));
for (int i = 1; i <= n; ++i)
    for (int j = 1; j <= m; ++j)
        if (a[i-1] == b[j-1]) dp[i][j] = dp[i-1][j-1] + 1;
        else dp[i][j] = max(dp[i-1][j], dp[i][j-1]);
return dp[n][m];
\end{verbatim}

\paragraph{Complejidad.}
\begin{itemize}\setlength\itemsep{0pt}
	\item $O(NM)$ tiempo y memoria.
	\item Para $N, M \le 5000$: $\sim 2.5 \cdot 10^7$ operaciones.
\end{itemize}

\paragraph{Donde tocar para modificar.}
\begin{itemize}\setlength\itemsep{0pt}
	\item \textbf{Reconstruir la LCS / las operaciones}: aniadir \texttt{prev[i][j]} para backtrack.
	\item \textbf{LCS de varios strings}: NP-hard para $\ge 3$ strings; en la practica se hace con bitmask.
	\item \textbf{Edit distance con costos distintos}: cambiar el \texttt{+1} por el costo apropiado.
	\item \textbf{Rolling array}: si memoria es problema, mantener solo 2 filas.
	\item \textbf{Hirschberg}: algoritmo $O(NM)$ tiempo, $O(\min(N,M))$ memoria, con reconstruccion.
\end{itemize}

\paragraph{Trampas y errores frecuentes.}
\begin{itemize}\setlength\itemsep{0pt}
	\item Indices en la recurrencia: \texttt{a[i-1]} y \texttt{b[j-1]} (no \texttt{a[i]}).
	\item Para reconstruccion, el caso \texttt{a[i-1] == b[j-1]} tiene prioridad.
	\item Si los strings son muy largos, considerar Hirschberg o bitset optimization.
\end{itemize}

\paragraph{Codigo.}
Ver \texttt{cpp.tex}, seccion \textit{LCS / Edit distance}.

\vspace{0.5em}

\subsubsection{Bitmask DP over subsets}

\paragraph{Idea intuitiva.}
Cuando el estado es un \textbf{subconjunto} de $\{0, \dots, N-1\}$, se representa como una mascara de $N$ bits. Para iterar sobre todos los subconjuntos de \texttt{mask}:
\[ \text{for (sub = mask; sub; sub = (sub - 1) \& mask)} \{ \dots \} \]
Para iterar sobre todos los superconjuntos de \texttt{mask} en $[0, U)$:
\[ \text{for (sup = mask; sup < U; sup = (sup + 1) | mask)} \{ \dots \} \]
La idea: las mascaras forman un lattice de subconjuntos; los trucos aritmeticos (\texttt{(sub-1) \& mask} y \texttt{(sup+1) | mask}) enumeran todo el lattice eficientemente.

\paragraph{Propiedades clave (invariantes).}
\begin{itemize}\setlength\itemsep{0pt}
	\item Iterar subsets: $O(2^k)$ para una mascara con $k$ bits.
	\item \texttt{sub \& mask == sub} siempre; \texttt{mask - sub} no siempre da un subconjunto.
	\item \texttt{\_\_builtin\_popcount(mask)} cuenta los bits.
	\item \texttt{(sub - 1) \& mask} itera subsets de \texttt{mask} en orden decreciente de inclusion.
\end{itemize}

\paragraph{Codigo clave.}
Iteracion sobre subsets de una mascara:
\begin{verbatim}
// Todos los subsets de mask
for (int sub = mask; sub; sub = (sub - 1) & mask) {
    // procesar sub
}
// Incluir sub = 0 si es necesario
sub = 0;  // procesar primero
for (int sup = mask; sup < (1 << N); sup = (sup + 1) | mask) {
    // procesar sup (todos los superconjuntos)
}
\end{verbatim}

\paragraph{Complejidad.}
\begin{itemize}\setlength\itemsep{0pt}
	\item $O(N \cdot 2^N)$ para una DP estandar sobre todos los subsets.
	\item Iterar subsets de \texttt{mask}: $O(2^{|mask|})$.
\end{itemize}

\paragraph{Donde tocar para modificar.}
\begin{itemize}\setlength\itemsep{0pt}
	\item \textbf{TSP / Hamiltonian path}: bitmask DP (ver tambien \textit{Hamiltonian path} en \texttt{cpp.tex} seccion Grafos).
	\item \textbf{Set cover}: iterar subsets.
	\item \textbf{DP con supermask}: si el estado incluye ``todos los elementos que \textbf{no} hemos usado'', usar supermask iteration.
	\item \textbf{Bitmasks grandes}: con $N > 20$, aniadir compresion o simulacion con segmentos.
\end{itemize}

\paragraph{Trampas y errores frecuentes.}
\begin{itemize}\setlength\itemsep{0pt}
	\item El orden de iteracion de subsets de \texttt{mask} es decreciente (de \texttt{mask} a $0$).
	\item \texttt{mask = 0} no tiene subsets no triviales; el bucle lo saltea correctamente.
	\item El truco \texttt{(sub - 1) \& mask} requiere cuidado con \texttt{sub = 0} (termina el bucle).
	\item Superconjuntos: \texttt{sup < (1 << N)} debe estar bien definido.
\end{itemize}

\paragraph{Codigo.}
Ver \texttt{cpp.tex}, seccion \textit{Bitmask DP over subsets}.

\vspace{0.5em}

\subsubsection{SOS DP}

\paragraph{Idea intuitiva.}
Para calcular, para cada \texttt{mask}, la suma de $f(\text{submask})$ sobre todos los $\text{submask} \subseteq \text{mask}$:
\[ dp[mask] = \sum_{\text{sub} \subseteq mask} f(\text{sub}). \]
Se computa en $O(N \cdot 2^N)$ mediante un loop por bit: para cada $i$, sumar a $dp[mask]$ la contribucion de $dp[mask \oplus (1 << i)]$ cuando $mask$ tiene el bit $i$. La idea: cada subconjunto $s \subseteq m$ se obtiene quitando bits uno a uno; asi, sumar incrementalmente captura todos los subconjuntos en $O(N \cdot 2^N)$.

\paragraph{Propiedades clave (invariantes).}
\begin{itemize}\setlength\itemsep{0pt}
	\item Inicializar $dp = f$.
	\item Para cada bit $i$: \texttt{if (mask \& (1 << i)) dp[mask] += dp[mask \^{} (1 << i)]}.
	\item Variantes: supersets DP (sumar sobre superconjuntos).
	\item El orden de iteracion (bit externo, mascara interno) es crucial.
\end{itemize}

\paragraph{Codigo clave.}
SOS DP (Sum over Subsets):
\begin{verbatim}
vector<long long> dp(1 << N);
for (int i = 0; i < (1 << N); ++i) dp[i] = f[i];
for (int bit = 0; bit < N; ++bit)
    for (int mask = 0; mask < (1 << N); ++mask)
        if (mask & (1 << bit))
            dp[mask] += dp[mask ^ (1 << bit)];
// dp[mask] = suma de f(sub) para sub ⊆ mask
\end{verbatim}

\paragraph{Complejidad.}
\begin{itemize}\setlength\itemsep{0pt}
	\item $O(N \cdot 2^N)$.
	\item Memoria $O(2^N)$.
\end{itemize}

\paragraph{Donde tocar para modificar.}
\begin{itemize}\setlength\itemsep{0pt}
	\item \textbf{Max/min en lugar de suma}: cambiar \texttt{+=} por la operacion correspondiente.
	\item \textbf{Supersets DP}: iterar \texttt{mask} de $0$ a $2^N$ y aniadir \texttt{dp[mask] += dp[mask | (1<<i)]}.
	\item \textbf{XOR convolution} (subset convolution): version mas avanzada.
	\item \textbf{Aplicaciones}: contar subconjuntos con propiedades especificas.
\end{itemize}

\paragraph{Trampas y errores frecuentes.}
\begin{itemize}\setlength\itemsep{0pt}
	\item Si $N$ es 30 o mas, $2^N$ no entra en memoria; usar trucos o reducir $N$.
	\item El sentido del XOR (sumar a \texttt{mask \^{} bit}) es para ``agregar el bit $i$ al subconjunto''.
	\item En supersets DP, la condicion es \texttt{(mask \& (1<<bit)) == 0}.
\end{itemize}

\paragraph{Codigo.}
Ver \texttt{cpp.tex}, seccion \textit{SOS DP}.

\vspace{0.5em}

\subsubsection{Digit DP template}

\paragraph{Idea intuitiva.}
Contar numeros en $[0, N]$ que cumplen una propiedad sobre sus digitos. La idea: DP por posicion, manteniendo:
\begin{itemize}\setlength\itemsep{0pt}
	\item \texttt{pos}: posicion actual (de izquierda a derecha).
	\item \texttt{tight}: si los digitos ya colocados son exactamente los de $N$ (limite superior estricto).
	\item \texttt{state}: estado adicional (digitos usados, suma, etc.).
\end{itemize}
La clave: la condicion \texttt{tight} permite acotar correctamente el siguiente digito; sin ella, la DP contaria incorrectamente.

\paragraph{Propiedades clave (invariantes).}
\begin{itemize}\setlength\itemsep{0pt}
	\item Si \texttt{tight} es true, el siguiente digito esta acotado por el de $N$; si no, puede ser 0-9.
	\item Memoria: $O(D \cdot 2 \cdot \text{states})$.
	\item Caso base: si \texttt{pos == D}, devolver 1 si el estado cumple la propiedad, sino 0.
	\item El bit \texttt{started} distingue ``numero con menos digitos'' (con leading zeros).
\end{itemize}

\paragraph{Codigo clave.}
La estructura basica:
\begin{verbatim}
long long dp(int pos, bool tight, State s) {
    if (pos == D) return check(s) ? 1 : 0;
    if (memo[pos][tight][s] != -1) return memo;
    int limit = tight ? digits[pos] : 9;
    long long ans = 0;
    for (int d = 0; d <= limit; ++d)
        ans += dp(pos + 1, tight && (d == limit), update(s, d, pos));
    return memo[pos][tight][s] = ans;
}
\end{verbatim}

\paragraph{Complejidad.}
\begin{itemize}\setlength\itemsep{0pt}
	\item $O(D \cdot 10 \cdot \text{states})$ donde $D$ es el numero de digitos.
	\item Para $D \le 19$ (hasta $10^{18}$): $\sim 200 \cdot \text{states}$.
\end{itemize}

\paragraph{Donde tocar para modificar.}
\begin{itemize}\setlength\itemsep{0pt}
	\item \textbf{Cambiar la propiedad}: modificar \texttt{updateState} y la condicion base.
	\item \textbf{Multiples estados}: aniadir mas dimensiones al \texttt{memo}.
	\item \textbf{Leading zeros}: aniadir un flag \texttt{started} para distinguir numeros con menos digitos.
	\item \textbf{Rango $[L, R]$}: calcular $f(R) - f(L-1)$.
\end{itemize}

\paragraph{Trampas y errores frecuentes.}
\begin{itemize}\setlength\itemsep{0pt}
	\item Olvidar resetear \texttt{memo} entre queries.
	\item Si la cantidad de estados es grande, el \texttt{memo[pos][tight][state]} puede ser muy grande; usar \texttt{map}.
	\item La condicion \texttt{started} es crucial para numeros con menos digitos que $D$.
	\item Si $N$ es muy grande ($> 10^{18}$), aniadir soporte para long long.
\end{itemize}

\paragraph{Codigo.}
Ver \texttt{cpp.tex}, seccion \textit{Digit DP template}.

\vspace{0.5em}

\subsubsection{Tree DP (reroot)}

\paragraph{Idea intuitiva.}
Calcular alguna propiedad asumiendo que \textbf{cada nodo es la raiz} del arbol. La tecnica clasica de \textbf{rerooting} en dos pasadas:
\begin{itemize}\setlength\itemsep{0pt}
	\item \textbf{DFS 1}: desde una raiz arbitraria, calcular \texttt{dp[v]} para el subarbol de $v$.
	\item \textbf{DFS 2}: para cada hijo $u$ de $v$, calcular \texttt{up[u]} = respuesta si la raiz fuera $u$.
\end{itemize}
La demostracion: tras la primera DFS, cada nodo conoce ``lo que aporta su subarbol''; en la segunda, se transmite ``lo que aporta el resto del arbol''.

\paragraph{Propiedades clave (invariantes).}
\begin{itemize}\setlength\itemsep{0pt}
	\item \texttt{dp[v]} depende solo del subarbol de $v$ (con la raiz original).
	\item \texttt{up[u]} se calcula de manera que cuando la raiz es $u$, ``todo lo que estaba en $v$ se ajusta''.
	\item Formula para reroot depende del problema; el snippet usa el ejemplo ``suma de distancias''.
	\item La combinacion \texttt{dp[u] + up[u]} da la respuesta con raiz $u$.
\end{itemize}

\paragraph{Codigo clave.}
Reroot para suma de distancias:
\begin{verbatim}
// dfs1: dp[v] = suma de distancias desde v al subarbol
void dfs1(int v, int p) {
    dp[v] = 0;
    for (int u : adj[v]) if (u != p) {
        dfs1(u, v);
        dp[v] += dp[u] + sz[u];  // cada hijo aporta su subarbol
    }
}
// dfs2: up[u] = distancia desde el resto del arbol a u
void dfs2(int v, int p) {
    for (int u : adj[v]) if (u != p) {
        up[u] = up[v] + dp[v] - dp[u] - sz[u] + (n - sz[u]);
        dfs2(u, v);
    }
}
\end{verbatim}

\paragraph{Complejidad.}
\begin{itemize}\setlength\itemsep{0pt}
	\item $O(N)$ para ambas pasadas.
	\item Cada arista se procesa un numero constante de veces.
\end{itemize}

\paragraph{Donde tocar para modificar.}
\begin{itemize}\setlength\itemsep{0pt}
	\item \textbf{Cambiar la cantidad a calcular}: modificar la operacion en \texttt{dfs1} y la formula en \texttt{dfs2}.
	\item \textbf{Arbol con pesos}: aniadir \texttt{weights[v]} y ajustar.
	\item \textbf{Re-root en arboles con ciclos}: requiere otra tecnica (DSU on tree).
	\item \textbf{Rooted trees diferentes}: si la propiedad no es invariante por raiz, usar otra tecnica.
\end{itemize}

\paragraph{Trampas y errores frecuentes.}
\begin{itemize}\setlength\itemsep{0pt}
	\item La formula de \texttt{up[u]} debe sumar y restar correctamente las contribuciones de $u$ y del resto.
	\item Para arboles con pesos, aniadir el peso en las formulas.
	\item Verificar con ejemplos pequenos antes de generalizar.
\end{itemize}

\paragraph{Codigo.}
Ver \texttt{cpp.tex}, seccion \textit{Tree DP (reroot)}.

% ============================================================
\section{Geometria}
% ============================================================

\subsection{Puntos y segmentos}

\subsubsection{Point + cross + dot}

\paragraph{Idea intuitiva.}
La estructura \texttt{Point} representa un punto en 2D con coordenadas enteras. Las dos operaciones vectoriales clave son:
\begin{itemize}\setlength\itemsep{0pt}
	\item \textbf{Cross product} $P \times Q = P_x Q_y - P_y Q_x$. Da el ``area con signo'' del paralelogramo formado por $P$ y $Q$. Usado para orientacion.
	\item \textbf{Dot product} $P \cdot Q = P_x Q_x + P_y Q_y$. Da el producto de las longitudes por el coseno del angulo. Usado para proyeccion.
\end{itemize}
La interpretacion geometrica del cross product: es el area con signo del paralelogramo; el signo indica la orientacion relativa.

\paragraph{Propiedades clave (invariantes).}
\begin{itemize}\setlength\itemsep{0pt}
	\item \texttt{cross(a, b) > 0} si $b$ esta a la izquierda de $a$ (sistema de coordenadas usual).
	\item \texttt{cross(a, b) = 0} si son colineales.
	\item \texttt{dot(a, b) > 0} si el angulo es agudo.
	\item El snippet define \texttt{cross} respecto a \texttt{this}, util para orientacion de 3 puntos.
	\item $|cross|^2 + |dot|^2 = |P|^2 \cdot |Q|^2$ (identidad de Lagrange).
\end{itemize}

\paragraph{Codigo clave.}
Las dos operaciones vectoriales:
\begin{verbatim}
long long cross(const Point& a, const Point& b) {
    return (long long)a.x * b.y - (long long)a.y * b.x;
}
long long dot(const Point& a, const Point& b) {
    return (long long)a.x * b.x + (long long)a.y * b.y;
}
\end{verbatim}

\paragraph{Complejidad.}
\begin{itemize}\setlength\itemsep{0pt}
	\item $O(1)$ por operacion.
	\item Constante muy pequena (4 multiplicaciones + 2 sumas).
\end{itemize}

\paragraph{Donde tocar para modificar.}
\begin{itemize}\setlength\itemsep{0pt}
	\item \textbf{Coordenadas doubles}: cambiar \texttt{int} por \texttt{double} o \texttt{long double}. Cuidado con la precision.
	\item \textbf{3D}: aniadir \texttt{z} y adaptar cross/dot.
	\item \textbf{Operadores adicionales}: aniadir \texttt{operator+}, \texttt{operator-}, \texttt{operator*} (escalar).
	\item \textbf{Distancia al cuadrado}: aniadir \texttt{len2()} que retorna $x^2 + y^2$.
\end{itemize}

\paragraph{Trampas y errores frecuentes.}
\begin{itemize}\setlength\itemsep{0pt}
	\item Cross product \textbf{no es} conmutativo; $P \times Q = -Q \times P$.
	\item Overflow con coordenadas grandes ($|x|, |y| > 10^4$) si se hace \texttt{x * y}.
	\item El signo de cross depende del sistema de coordenadas (en computacion, usualmente CCW positivo).
	\item Para ``igualdad'' de doubles, aniadir tolerancia \texttt{eps}.
\end{itemize}

\paragraph{Codigo.}
Ver \texttt{cpp.tex}, seccion \textit{Point + cross + dot}.

\vspace{0.5em}

\subsubsection{Interseccion de segmentos}

\paragraph{Idea intuitiva.}
Dados dos segmentos $[a, b]$ y $[c, d]$, determinar si se intersectan. La idea clasica:
\begin{itemize}\setlength\itemsep{0pt}
	\item Verificar orientaciones: $a, b, c$ y $a, b, d$ deben estar en lados opuestos de la recta $ab$; igual para $cd$ y $ab$.
	\item Caso colineal: si las rectas son paralelas (\texttt{cross} $= 0$) y los segmentos son colineales, verificar que los bounding boxes se intersectan.
\end{itemize}
La demostracion: dos segmentos se intersectan sii cada segmento ``cruza'' la linea que contiene al otro (estricto) o ambos son colineales con solapamiento.

\paragraph{Propiedades clave (invariantes).}
\begin{itemize}\setlength\itemsep{0pt}
	\item Funciona con segmentos en cualquier orientacion.
	\item Caso colineal: requiere check adicional de bounding box.
	\item Para intersectar en un \textbf{punto} especifico: aniadir division con cuidado de precision.
	\item Funciona con segmentos verticales/horizontales sin casos especiales.
\end{itemize}

\paragraph{Codigo clave.}
Test clasico de interseccion:
\begin{verbatim}
int sign(long long x) { return (x > 0) - (x < 0); }
bool intersect(Point a, Point b, Point c, Point d) {
    long long o1 = cross(b - a, c - a);
    long long o2 = cross(b - a, d - a);
    long long o3 = cross(d - c, a - c);
    long long o4 = cross(d - c, b - c);
    if (sign(o1) != sign(o2) && sign(o3) != sign(o4)) return true;
    // casos colineales
    if (o1 == 0 && onSegment(a, b, c)) return true;
    if (o2 == 0 && onSegment(a, b, d)) return true;
    if (o3 == 0 && onSegment(c, d, a)) return true;
    if (o4 == 0 && onSegment(c, d, b)) return true;
    return false;
}
\end{verbatim}

\paragraph{Complejidad.}
\begin{itemize}\setlength\itemsep{0pt}
	\item $O(1)$.
	\item Constante pequena.
\end{itemize}

\paragraph{Donde tocar para modificar.}
\begin{itemize}\setlength\itemsep{0pt}
	\item \textbf{Interseccion con punto}: aniadir la division despues de verificar orientacion.
	\item \textbf{Interseccion rayo-segmento}: ajustar la primera condicion.
	\item \textbf{Poligonos convexos}: usar separacion de ejes (SAT).
	\item \textbf{Multiples queries}: precomputar bounding boxes o usar segment tree.
\end{itemize}

\paragraph{Trampas y errores frecuentes.}
\begin{itemize}\setlength\itemsep{0pt}
	\item Si los segmentos tocan solo en un extremo, el snippet lo considera interseccion (algunos problemas quieren estricto). Ajustar segun el problema.
	\item El caso colineal es comun; olvidar el check adicional de bounding box da falsos negativos.
	\item En problemas con muchos segmentos, considerar sweep line para $O(n \log n)$.
\end{itemize}

\paragraph{Codigo.}
Ver \texttt{cpp.tex}, seccion \textit{Interseccion de segmentos}.

\vspace{0.5em}

\subsubsection{Contains colineal}

\paragraph{Idea intuitiva.}
Verifica si un punto $c$ esta sobre el segmento $[a, b]$ (asumiendo que ya son colineales). Si \texttt{cross(a, b, c) = 0} y $c$ esta dentro del bounding box de $[a, b]$, entonces $c$ esta en el segmento. La razon: si los tres puntos son colineales, $c$ esta en el segmento sii sus coordenadas estan entre las de $a$ y $b$.

\paragraph{Propiedades clave (invariantes).}
\begin{itemize}\setlength\itemsep{0pt}
	\item Asume colinealidad (no la verifica mas alla de \texttt{cross == 0}).
	\item Usa comparaciones con \texttt{<=} para extremos \textbf{inclusivos}.
	\item Funciona con segmentos verticales (la condicion se reduce a una sola coordenada).
\end{itemize}

\paragraph{Codigo clave.}
Test de contencion en segmento (asume colinealidad):
\begin{verbatim}
bool onSegment(Point a, Point b, Point c) {
    return min(a.x, b.x) <= c.x && c.x <= max(a.x, b.x) &&
           min(a.y, b.y) <= c.y && c.y <= max(a.y, b.y);
}
\end{verbatim}

\paragraph{Complejidad.}
\begin{itemize}\setlength\itemsep{0pt}
	\item $O(1)$.
	\item Constante minima.
\end{itemize}

\paragraph{Donde tocar para modificar.}
\begin{itemize}\setlength\itemsep{0pt}
	\item \textbf{Contener estricto}: cambiar \texttt{<=} a \texttt{<} para excluir extremos.
	\item \textbf{Coordenadas doubles}: usar tolerancia (\texttt{eps}).
	\item \textbf{Solo x o solo y}: aniadir check especifico para segmentos verticales/horizontales.
\end{itemize}

\paragraph{Trampas y errores frecuentes.}
\begin{itemize}\setlength\itemsep{0pt}
	\item No usar este snippet para puntos \textbf{no colineales}; siempre verificar primero la colinealidad.
	\item Si los puntos tienen doubles, el chequeo \texttt{cross == 0} puede fallar; aniadir tolerancia.
	\item El bounding box incluye los extremos; verificar si el problema quiere exclusivo o inclusivo.
\end{itemize}

\paragraph{Codigo.}
Ver \texttt{cpp.tex}, seccion \textit{Contains colineal}.

\subsection{Poligonos}

\subsubsection{Convex Hull (monotono)}

\paragraph{Idea intuitiva.}
Algoritmo de \textbf{Andrew's monotone chain}: ordenar puntos por $(x, y)$, luego construir la hull inferior y superior por separado. Para cada punto, aniadirlo al hull y, mientras el ultimo giro no sea ``counter-clockwise'' (\texttt{cross > 0}), popear. La demostracion: si los puntos se procesan en orden, el hull inferior y superior cubren la frontera completa sin auto-intersecciones.

\paragraph{Propiedades clave (invariantes).}
\begin{itemize}\setlength\itemsep{0pt}
	\item Ordena los puntos en $O(n \log n)$.
	\item Construye la hull en $O(n)$ despues.
	\item Output: poligono convexo en counter-clockwise, sin punto final duplicado.
	\item El \texttt{<= 0} incluye colineales; usar \texttt{< 0} para excluirlos.
	\item Cada punto aparece en la hull inferior o superior (no ambos).
\end{itemize}

\paragraph{Codigo clave.}
Construccion del hull inferior y superior:
\begin{verbatim}
sort(p.begin(), p.end());
vector<Point> lower, upper;
for (auto& p : points) {
    while (lower.size() >= 2 && cross(lower.back() - lower[lower.size()-2],
                                       p - lower.back()) <= 0)
        lower.pop_back();
    lower.push_back(p);
}
for (int i = points.size() - 1; i >= 0; --i) {
    auto& p = points[i];
    while (upper.size() >= 2 && cross(upper.back() - upper[upper.size()-2],
                                       p - upper.back()) <= 0)
        upper.pop_back();
    upper.push_back(p);
}
vector<Point> hull = lower;
hull.insert(hull.end(), upper.begin() + 1, upper.end() - 1);
\end{verbatim}

\paragraph{Complejidad.}
\begin{itemize}\setlength\itemsep{0pt}
	\item $O(n \log n)$.
	\item Constante muy pequena (solo cross products).
\end{itemize}

\paragraph{Donde tocar para modificar.}
\begin{itemize}\setlength\itemsep{0pt}
	\item \textbf{Colineales}: cambiar \texttt{<= 0} por \texttt{< 0} para excluirlos del hull.
	\item \textbf{Puntos repetidos}: \texttt{unique} despues de sort.
	\item \textbf{Hull 3D}: gift wrapping es $O(n^2)$; convex hull 3D es mas complejo.
	\item \textbf{Output con los puntos}: aniadir el hull completo (no solo upper/lower).
\end{itemize}

\paragraph{Trampas y errores frecuentes.}
\begin{itemize}\setlength\itemsep{0pt}
	\item El \texttt{pop\_back} y el \texttt{reverse} son esenciales para no duplicar el ultimo punto.
	\item Para excluir colineales, ajustar el \texttt{<= 0} vs \texttt{< 0}.
	\item Si los puntos estan duplicados, aniadir \texttt{unique} antes del sort.
\end{itemize}

\paragraph{Codigo.}
Ver \texttt{cpp.tex}, seccion \textit{Convex Hull (monotono)}.

\vspace{0.5em}

\subsubsection{Polygon Area}

\paragraph{Idea intuitiva.}
El area (con signo) de un poligono se calcula con la \textbf{Shoelace formula}: $A = \frac{1}{2} \sum_{i=0}^{n-1} (x_i y_{i+1} - x_{i+1} y_i)$, donde el indice es modulo $n$. El snippet retorna el \textbf{doble} del area (en valor absoluto). La formula sale del ``sumar areas de trapezoides'': cada lado del poligono define un trapezoide con el eje $x$; la suma con signo da el area.

\paragraph{Propiedades clave (invariantes).}
\begin{itemize}\setlength\itemsep{0pt}
	\item El signo indica la orientacion: positivo si counter-clockwise.
	\item Para el doble: no dividir por 2 si no es necesario (evita fracciones).
	\item Funciona para poligonos simples (no auto-intersectantes).
	\item Si el poligono es no convexo, la formula da el area ``con signo''.
\end{itemize}

\paragraph{Codigo clave.}
La suma de trapezoides:
\begin{verbatim}
long long shoelace(const vector<Point>& p) {
    long long s = 0;
    int n = p.size();
    for (int i = 0; i < n; ++i)
        s += (long long)p[i].x * p[(i+1)%n].y - (long long)p[(i+1)%n].x * p[i].y;
    return abs(s);  // doble del area
}
\end{verbatim}

\paragraph{Complejidad.}
\begin{itemize}\setlength\itemsep{0pt}
	\item $O(n)$.
	\item Constante minima.
\end{itemize}

\paragraph{Donde tocar para modificar.}
\begin{itemize}\setlength\itemsep{0pt}
	\item \textbf{Area exacta}: dividir por 2 (con cuidado si da fraccion; usar \texttt{double}).
	\item \textbf{Area firmada}: quitar \texttt{abs} y dejar el signo.
	\item \textbf{Punto dentro de poligono convexo}: con la hull, chequear que esta a un mismo lado de todas las aristas.
	\item \textbf{Triangulacion}: descomponer en triangulos y sumar areas.
\end{itemize}

\paragraph{Trampas y errores frecuentes.}
\begin{itemize}\setlength\itemsep{0pt}
	\item Si los puntos no estan en orden (CCW o CW), el area firmada puede ser incorrecta. Para signed area consistente, usar siempre el mismo orden.
	\item Para poligonos con holes, sumar las areas con signo apropiado.
	\item Con doubles, aniadir tolerancia \texttt{eps}.
\end{itemize}

\paragraph{Codigo.}
Ver \texttt{cpp.tex}, seccion \textit{Polygon Area}.

\vspace{0.5em}

\subsubsection{Distance point-segment / point-line}

\paragraph{Idea intuitiva.}
Distancia minima de un punto $p$ a un segmento $[a, b]$ o a la recta que pasa por $a$ y $b$.
\begin{itemize}\setlength\itemsep{0pt}
	\item \textbf{Linea}: $d = \frac{|ab \times ap|}{|ab|}$.
	\item \textbf{Segmento}: si la proyeccion de $p$ sobre $ab$ cae fuera del segmento, la distancia es al extremo mas cercano.
\end{itemize}
La razon geometrica: la distancia a una linea es el area del paralelogramo dividido por la base; la proyeccion ortogonal indica si estamos ``dentro'' o ``fuera'' del segmento.

\paragraph{Propiedades clave (invariantes).}
\begin{itemize}\setlength\itemsep{0pt}
	\item \texttt{dist2PointLine} retorna distancia \textbf{al cuadrado}.
	\item \texttt{dist2PointSegment} distingue 3 casos: $p$ antes de $a$, despues de $b$, o proyectado sobre el segmento.
	\item Si \texttt{len2 = 0} (degenerate), la distancia es $|ap|$ o $|bp|$.
	\item La proyeccion se calcula con dot products y el parametro $t \in [0, 1]$.
\end{itemize}

\paragraph{Codigo clave.}
Distancia al cuadrado a segmento:
\begin{verbatim}
long long dist2Segment(Point p, Point a, Point b) {
    long long len2 = (a-b).len2();
    if (len2 == 0) return (p-a).len2();  // degenerate
    long long t = max(0LL, min(1LL, dot(p-a, b-a) / len2));
    Point proj = {a.x + t*(b.x - a.x), a.y + t*(b.y - a.y)};
    return (p - proj).len2();
}
\end{verbatim}

\paragraph{Complejidad.}
\begin{itemize}\setlength\itemsep{0pt}
	\item $O(1)$.
	\item Constante minima.
\end{itemize}

\paragraph{Donde tocar para modificar.}
\begin{itemize}\setlength\itemsep{0pt}
	\item \textbf{Distancia (no al cuadrado)}: aniadir \texttt{sqrt} al final.
	\item \textbf{Coordenadas doubles}: aniadir \texttt{long double} y tolerancia.
	\item \textbf{3D}: aniadir coordenada $z$ y adaptar.
	\item \textbf{Multiples queries}: usar KD-tree o segment tree para acelerar.
\end{itemize}

\paragraph{Trampas y errores frecuentes.}
\begin{itemize}\setlength\itemsep{0pt}
	\item Overflow en \texttt{area2 * area2} si el area es grande; \texttt{long long} alcanza para $|x|, |y| \le 10^9$ pero no mas.
	\item Si $a = b$ (segmento degenerado), aniadir check explicito.
	\item En doubles, aniadir tolerancia \texttt{eps} para comparar.
\end{itemize}

\paragraph{Codigo.}
Ver \texttt{cpp.tex}, seccion \textit{Distance point-segment / point-line}.

\vspace{0.5em}

\subsubsection{Projection of point onto line}

\paragraph{Idea intuitiva.}
Proyeccion ortogonal de $p$ sobre la recta que pasa por $a$ y $b$. El parametro $t$ tal que la proyeccion es $a + t \cdot (b - a)$ es:
\[ t = \frac{(p - a) \cdot (b - a)}{|b - a|^2}. \]
La idea: descomponemos $\vec{ap}$ en una componente paralela a $\vec{ab}$ y una perpendicular; $t$ mide la posicion en la primera.

\paragraph{Propiedades clave (invariantes).}
\begin{itemize}\setlength\itemsep{0pt}
	\item Si $t \in [0, 1]$, la proyeccion cae sobre el \textbf{segmento}; si no, cae sobre la extension.
	\item Retorna \texttt{double} (puede tener precision issues).
	\item $t < 0$: $p$ esta ``antes'' de $a$; $t > 1$: $p$ esta ``despues'' de $b$.
\end{itemize}

\paragraph{Codigo clave.}
Proyeccion ortogonal:
\begin{verbatim}
Point projection(Point p, Point a, Point b) {
    Point ab = b - a, ap = p - a;
    double t = (double)dot(ap, ab) / ab.len2();
    return {a.x + t * ab.x, a.y + t * ab.y};
}
\end{verbatim}

\paragraph{Complejidad.}
\begin{itemize}\setlength\itemsep{0pt}
	\item $O(1)$.
	\item Constante minima (dos dot products y una division).
\end{itemize}

\paragraph{Donde tocar para modificar.}
\begin{itemize}\setlength\itemsep{0pt}
	\item \textbf{Proyeccion sobre segmento}: clipar $t$ a $[0, 1]$ antes de calcular el punto.
	\item \textbf{Distancia al segmento via proyeccion}: combinar con el modulo anterior.
	\item \textbf{Solo el parametro $t$}: si solo necesitas saber si esta dentro del segmento, retorna $t$ sin construir el punto.
\end{itemize}

\paragraph{Trampas y errores frecuentes.}
\begin{itemize}\setlength\itemsep{0pt}
	\item Si $a = b$ ($|ab| = 0$), la proyeccion no esta definida; aniadir guard.
	\item En doubles, aniadir tolerancia \texttt{eps} para comparar $t$ con 0 o 1.
	\item La division por \texttt{len2} puede ser cara si las coordenadas son grandes; usar \texttt{double} o precomputar.
\end{itemize}

\paragraph{Codigo.}
Ver \texttt{cpp.tex}, seccion \textit{Projection of point onto line}.

\vspace{0.5em}

\subsubsection{Point in polygon (ray casting)}

\paragraph{Idea intuitiva.}
Algoritmo de \textbf{ray casting}: contar cuantas veces un rayo horizontal desde $p$ cruza los bordes del poligono. Si es impar, $p$ esta dentro; si par, fuera. El snippet implementa la version clasica recorriendo aristas. La razon: cada cruce cambia el lado del rayo respecto al poligono; un numero impar de cambios significa ``terminamos dentro''.

\paragraph{Propiedades clave (invariantes).}
\begin{itemize}\setlength\itemsep{0pt}
	\item Funciona con poligonos simples (no necesariamente convexos).
	\item Las intersecciones en vertices cuentan especialmente.
	\item Para poligonos con \textbf{huecos}, sigue funcionando.
	\item La respuesta es estrictamente 0 (fuera) o 1 (dentro) para poligonos simples.
\end{itemize}

\paragraph{Codigo clave.}
El test clasico de ray casting:
\begin{verbatim}
bool pointInPolygon(Point p, const vector<Point>& poly) {
    int n = poly.size();
    bool inside = false;
    for (int i = 0, j = n - 1; i < n; j = i++) {
        if ((poly[i].y > p.y) != (poly[j].y > p.y) &&
            p.x < (poly[j].x - poly[i].x) * (p.y - poly[i].y) /
                    (poly[j].y - poly[i].y) + poly[i].x)
            inside = !inside;
    }
    return inside;
}
\end{verbatim}

\paragraph{Complejidad.}
\begin{itemize}\setlength\itemsep{0pt}
	\item $O(n)$ por query.
	\item Memoria $O(1)$ (no se almacena el rayo).
\end{itemize}

\paragraph{Donde tocar para modificar.}
\begin{itemize}\setlength\itemsep{0pt}
	\item \textbf{Winding number}: version alternativa que cuenta el numero de vueltas.
	\item \textbf{Poligono convexo}: check lineal contra las aristas (mas rapido).
	\item \textbf{Multiples queries}: precomputar o usar segment tree sobre las aristas.
	\item \textbf{Punto en frontera}: aniadir check adicional si es colineal con alguna arista.
\end{itemize}

\paragraph{Trampas y errores frecuentes.}
\begin{itemize}\setlength\itemsep{0pt}
	\item El manejo del vertice (cuando el rayo pasa por un vertice) puede causar double-counting; el snippet lo evita con la condicion \texttt{(poly[i].y > p.y) != (poly[j].y > p.y)}.
	\item Si el poligono tiene holes, el algoritmo sigue funcionando.
	\item Para poligonos no simples (auto-intersectantes), el resultado puede ser ambiguo.
\end{itemize}

\paragraph{Codigo.}
Ver \texttt{cpp.tex}, seccion \textit{Point in polygon (ray casting)}.

\vspace{0.5em}

\subsubsection{Minkowski sum (convex polygons)}

\paragraph{Idea intuitiva.}
La \textbf{suma de Minkowski} $A \oplus B$ de dos poligonos convexos es otro poligono convexo que representa $\{a + b : a \in A, b \in B\}$. Algoritmo: ordenar las aristas por angulo polar (CCW) y ``mergearlas'' como en merge sort; el resultado es un poligono CCW sin repetir el ultimo punto. La idea: el resultado tiene una arista paralela a cada arista de los originales; la suma se obtiene barriendo los angulos.

\paragraph{Propiedades clave (invariantes).}
\begin{itemize}\setlength\itemsep{0pt}
	\item Si $A, B$ son convexos con $n, m$ vertices, $A \oplus B$ tiene $\le n + m$ vertices.
	\item Util para \textbf{detectar colision}: $A$ colisiona con $B$ si $0 \in A \oplus (-B)$.
	\item Asume poligonos en orden \textbf{counter-clockwise}.
	\item El resultado es convexo (suma de convexos es convexo).
\end{itemize}

\paragraph{Codigo clave.}
La suma de Minkowski:
\begin{verbatim}
vector<Point> minkowski(const vector<Point>& A, const vector<Point>& B) {
    vector<Point> result;
    int i = 0, j = 0;
    int n = A.size(), m = B.size();
    do {
        result.push_back(A[i] + B[j]);
        long long cross = (A[(i+1)%n] - A[i]) ^ (B[(j+1)%m] - B[j]);
        if (cross >= 0) i = (i + 1) % n;
        if (cross <= 0) j = (j + 1) % m;
    } while (i || j);
    return result;
}
\end{verbatim}

\paragraph{Complejidad.}
\begin{itemize}\setlength\itemsep{0pt}
	\item $O(n + m)$.
	\item Memoria $O(n + m)$.
\end{itemize}

\paragraph{Donde tocar para modificar.}
\begin{itemize}\setlength\itemsep{0pt}
	\item \textbf{Deteccion de colision}: aniadir $(-B)$ (reflejar) y verificar si $0$ esta en el resultado.
	\item \textbf{Poligonos no convexos}: la suma puede no ser poligono simple; algoritmo distinto.
	\item \textbf{Multiples queries}: precomputar las aristas una vez.
	\item \textbf{Robot motion planning}: usar para planear trayectorias.
\end{itemize}

\paragraph{Trampas y errores frecuentes.}
\begin{itemize}\setlength\itemsep{0pt}
	\item Si los poligonos no son CCW, la suma es incorrecta. Normalizar antes.
	\item El bucle termina cuando ambos indices vuelven a 0; aniadir check.
	\item Para la deteccion de colision, aniadir el test \texttt{0 in A \^{} (-B)}.
\end{itemize}

\paragraph{Codigo.}
Ver \texttt{cpp.tex}, seccion \textit{Minkowski sum (convex polygons)}.

\subsection{Herramientas}

\subsubsection{Li Chao Tree (long long)}

\paragraph{Idea intuitiva.}
Mantiene un \textbf{lower envelope} (o upper) de un conjunto de rectas $y = mx + b$ sobre un dominio $[X_{LO}, X_{HI}]$ y responde ``minimo $y$ en un $x$ dado'' en $O(\log X)$. La estructura es un segment tree donde cada nodo guarda la recta ``mejor'' para su segmento medio; al aniadir una recta, se compara con la actual y se decide cual es mejor en cada mitad, descendiendo recursivamente. La demostracion: en cada nivel, solo se guarda la mejor recta para el ``punto medio''; las otras se transmiten a los hijos.

\paragraph{Propiedades clave (invariantes).}
\begin{itemize}\setlength\itemsep{0pt}
	\item \textbf{Online}: cada \texttt{addLine} en $O(\log X)$.
	\item Solo insercion; no hay borrar (para borrar, usar Li Chao segmentado con reset).
	\item La recta guardada en un nodo es la ``mejor'' en su punto medio.
	\item Las rectas ``perdedoras'' se transmiten a los hijos para futuros inserts.
\end{itemize}

\paragraph{Codigo clave.}
La insercion en Li Chao:
\begin{verbatim}
void addLine(Line nw, int v = 1, int l = X_LO, int r = X_HI) {
    int m = (l + r) / 2;
    bool left = nw.get(l) < tree[v].get(l);
    bool mid = nw.get(m) < tree[v].get(m);
    if (mid) swap(tree[v], nw);
    if (r - l == 1) return;
    if (left != mid) addLine(nw, v*2, l, m);
    else             addLine(nw, v*2+1, m, r);
}
\end{verbatim}

\paragraph{Complejidad.}
\begin{itemize}\setlength\itemsep{0pt}
	\item $O(\log X)$ por insercion y query, donde $X$ es el tamano del dominio.
	\item Memoria $O(\log X)$ por insercion, $O(X)$ total.
\end{itemize}

\paragraph{Donde tocar para modificar.}
\begin{itemize}\setlength\itemsep{0pt}
	\item \textbf{Upper envelope en vez de lower}: invertir el signo (\texttt{>} en vez de \texttt{<}).
	\item \textbf{Queries de segmento} (no full line): si la recta solo esta activa en un rango, aniadir bandera \texttt{active}.
	\item \textbf{Eliminar rectas}: usar Li Chao segmentado (mas complejo).
	\item \textbf{Pendientes grandes}: si las pendientes tienen magnitud $> 10^9$, aniadir \texttt{\_\_int128} para evitar overflow.
\end{itemize}

\paragraph{Trampas y errores frecuentes.}
\begin{itemize}\setlength\itemsep{0pt}
	\item El parametro \texttt{lo == hi} debe retornar el valor del nodo (caso base).
	\item La comparacion \texttt{newLine.get(mid) < node->line.get(mid)} decide quien es mejor en el medio.
	\item Si las rectas tienen $m$ iguales pero $b$ distintos, se queda con la primera (tie-break).
\end{itemize}

\paragraph{Codigo.}
Ver \texttt{cpp.tex}, seccion \textit{Li Chao Tree (long long)}.

\vspace{0.5em}

\subsubsection{Sweep line template}

\paragraph{Idea intuitiva.}
Patron de diseno para problemas geometricos donde se barre una linea sobre el plano y se mantiene una estructura de datos sobre el otro eje. Eventos (add, query, remove) se ordenan por la coordenada de barrido y se procesan secuencialmente.

\paragraph{Propiedades clave (invariantes).}
\begin{itemize}\setlength\itemsep{0pt}
	\item Eventos: tuplas \texttt{(x, type, id)}.
	\item La estructura auxiliar (set, multiset, segtree) se mantiene sobre el otro eje.
	\item Procesar eventos en orden de $x$ ascendente.
\end{itemize}

\paragraph{Complejidad.}
\begin{itemize}\setlength\itemsep{0pt}
	\item Depende del problema; tipico $O((N + E) \log N)$.
\end{itemize}

\paragraph{Donde tocar para modificar.}
\begin{itemize}\setlength\itemsep{0pt}
	\item \textbf{Aniadir/quitar segmentos}: aniadir eventos add/remove al barrido.
	\item \textbf{Aniadir queries}: aniadir eventos query.
	\item \textbf{Cambiar el eje de barrido}: si barres en $y$ en vez de $x$, aniadir comparador.
\end{itemize}

\paragraph{Trampas y errores frecuentes.}
\begin{itemize}\setlength\itemsep{0pt}
	\item El snippet es solo el \textbf{esqueleto}; los handlers \texttt{add/query/remove} se llenan por problema.
	\item Cuidado con eventos duplicados.
\end{itemize}

\paragraph{Codigo.}
Ver \texttt{cpp.tex}, seccion \textit{Sweep line template}.

\vspace{0.5em}

\subsubsection{Closest pair of points}

\paragraph{Idea intuitiva.}
Encuentra el par de puntos mas cercano en $O(n \log n)$ con \textbf{divide y venceras}. Algoritmo:
\begin{itemize}\setlength\itemsep{0pt}
	\item Sort por $x$.
	\item Dividir en mitades; recursivo.
	\item Combinar: dado el minimo $d$ de las dos mitades, strip = puntos en $[midX - d, midX + d]$; ordenar por $y$; cada punto solo se compara con los siguientes hasta $y$ difiera en $< d$ (maximo 7).
\end{itemize}

\paragraph{Propiedades clave (invariantes).}
\begin{itemize}\setlength\itemsep{0pt}
	\item Requiere ordenar por $y$ despues (en el snippet se hace durante el merge).
	\item Strip tiene a lo sumo $O(n)$ puntos en total, pero la comparacion es $O(1)$ amortizada por punto.
	\item Resultado en \texttt{double}.
\end{itemize}

\paragraph{Complejidad.}
\begin{itemize}\setlength\itemsep{0pt}
	\item $O(n \log n)$.
\end{itemize}

\paragraph{Donde tocar para modificar.}
\begin{itemize}\setlength\itemsep{0pt}
	\item \textbf{Reconstruir el par}: aniadir indices en la recursion.
	\item \textbf{Distancia maxima}: cambiar \texttt{min} por \texttt{max} (mas facil; no requiere D\&C).
	\item \textbf{Puntos duplicados}: aniadir check especial (distancia 0).
\end{itemize}

\paragraph{Trampas y errores frecuentes.}
\begin{itemize}\setlength\itemsep{0pt}
	\item El merge requiere ordenar por $y$; el snippet lo hace in-place con un \texttt{tmp}.
	\item Si se hace en cada recursion, es $O(n \log n)$ extra.
\end{itemize}

\paragraph{Codigo.}
Ver \texttt{cpp.tex}, seccion \textit{Closest pair of points}.

\vspace{0.5em}

\subsubsection{Rotating calipers (polygon diameter)}

\paragraph{Idea intuitiva.}
Encuentra el \textbf{diametro} de un poligono convexo (par de vertices mas distantes) en $O(n)$. Asume poligono en orden counter-clockwise. Algoritmo: dos indices $i$ (sobre las aristas) y $k$ (sobre los vertices opuestos); en cada paso, rotar $k$ mientras el area del paralelogramo aumenta. La distancia $i$-$k$ puede ser maxima en cualquier momento del barrido.

\paragraph{Propiedades clave (invariantes).}
\begin{itemize}\setlength\itemsep{0pt}
	\item Cada indice avanza como maximo $n$ veces total.
	\item El ``diametro'' es la maxima distancia entre dos vertices.
	\item Retorna \textbf{distancia al cuadrado}.
\end{itemize}

\paragraph{Complejidad.}
\begin{itemize}\setlength\itemsep{0pt}
	\item $O(n)$.
\end{itemize}

\paragraph{Donde tocar para modificar.}
\begin{itemize}\setlength\itemsep{0pt}
	\item \textbf{Reconstruir el par de vertices}: aniadir indices.
	\item \textbf{Anchura minima / maxima}: variantes de calipers.
	\item \textbf{Distancia minima entre poligonos convexos}: otra variante.
\end{itemize}

\paragraph{Trampas y errores frecuentes.}
\begin{itemize}\setlength\itemsep{0pt}
	\item Asume poligono \textbf{estrictamente} convexo (sin vertices colineales). Si hay colineales, aniadir perturbacion o filtrar.
\end{itemize}

\paragraph{Codigo.}
Ver \texttt{cpp.tex}, seccion \textit{Rotating calipers (polygon diameter)}.

% ============================================================
\section{Miscelanea}
% ============================================================

\subsection{Utilidades del usuario}

\subsubsection{Hora}

\paragraph{Idea intuitiva.}
Una pequena clase para representar \textbf{tiempos en segundos} y hacer aritmetica basica sobre ellos. Internamente almacena todo en segundos (un solo \texttt{int}), lo cual simplifica comparaciones y diferencias. El constructor acepta $(h, m, s)$ y los convierte a segundos. Hay un metodo \texttt{cut()} que fuerza un minimo (util para ``horas trabajadas'').

\paragraph{Propiedades clave (invariantes).}
\begin{itemize}\setlength\itemsep{0pt}
	\item Representacion canonica: \textbf{solo segundos}. No hay zonas horarias, no hay DST.
	\item Comparadores \texttt{<} y \texttt{>} operan sobre los segundos.
	\item \texttt{operator-} retorna la \textbf{diferencia absoluta} en segundos (siempre no negativa).
	\item \texttt{cut()} aplica un piso (por defecto $8\text{h }30\text{m} = 30600$ segundos).
\end{itemize}

\paragraph{Complejidad.}
\begin{itemize}\setlength\itemsep{0pt}
	\item $O(1)$ por operacion.
\end{itemize}

\paragraph{Donde tocar para modificar.}
\begin{itemize}\setlength\itemsep{0pt}
	\item \textbf{Aniadir suma}: \texttt{operator+} que devuelve un \texttt{Hour} con la suma de segundos.
	\item \textbf{Aniadir output formateado}: aniadir \texttt{show\_hms()} que imprime $H$:$M$:$S$.
	\item \textbf{Cambiar el ``corte''}: modificar el valor magico \texttt{30'600} en \texttt{cut()} o pasarlo como parametro.
	\item \textbf{Tamano maximo}: si los tiempos exceden $\sim 68$ anos, el \texttt{int} desborda; usar \texttt{long long}.
	\item \textbf{Precision sub-segundo}: aniadir \texttt{ms} (milisegundos) y guardar en \texttt{long long} con factor 1000.
\end{itemize}

\paragraph{Trampas y errores frecuentes.}
\begin{itemize}\setlength\itemsep{0pt}
	\item El operador \texttt{-} retorna \textbf{valor absoluto}, no diferencia con signo. Si necesitas diferencia con signo, aniadir un nuevo metodo o sobrecargar.
	\item \texttt{cut()} modifica \texttt{s} \textbf{in-place}; si el llamador no espera esto, aniadir una version \texttt{cut(min)} que retorna un nuevo \texttt{Hour}.
	\item El umbral magico \texttt{30'600} (8h 30m) esta hardcoded; documentar bien que es ``jornada minima''.
\end{itemize}

\paragraph{Codigo.}
Ver \texttt{cpp.tex}, seccion \textit{Hora}.

\end{document}
